\documentclass[11pt]{article}

\usepackage[margin=1in]{geometry}
\usepackage{amsmath,amssymb,amsthm,mathtools,bm}
\usepackage{enumitem,booktabs,array}
\usepackage{microtype}
\usepackage[colorlinks=true,linkcolor=blue,citecolor=blue,urlcolor=blue]{hyperref}

\newtheorem{theorem}{Theorem}[section]
\newtheorem{proposition}[theorem]{Proposition}
\newtheorem{lemma}[theorem]{Lemma}
\newtheorem{corollary}[theorem]{Corollary}
\newtheorem{definition}[theorem]{Definition}
\newtheorem{remark}[theorem]{Remark}
\newtheorem{assumption}[theorem]{Assumption}
\newtheorem{obstruction}[theorem]{Obstruction}
\newtheorem{firewall}[theorem]{Firewall}

\newcommand{\cH}{\mathcal H}
\newcommand{\cL}{\mathcal L}
\newcommand{\cA}{\mathcal A}
\newcommand{\cD}{\mathcal D}
\newcommand{\cR}{\mathcal R}
\newcommand{\cP}{\mathcal P}
\newcommand{\cK}{\mathcal K}
\newcommand{\cS}{\mathcal S}
\newcommand{\cY}{\mathcal Y}
\newcommand{\cU}{\mathcal U}
\newcommand{\eps}{\varepsilon}
\newcommand{\1}{\mathbf 1}
\newcommand{\Tr}{\operatorname{Tr}}
\newcommand{\tr}{\operatorname{tr}}
\newcommand{\Det}{\operatorname{Det}}
\newcommand{\Pf}{\operatorname{Pf}}
\newcommand{\sgn}{\operatorname{sgn}}
\newcommand{\Ran}{\operatorname{Ran}}
\newcommand{\dist}{\operatorname{dist}}
\newcommand{\ad}{\operatorname{ad}}
\newcommand{\supp}{\operatorname{supp}}
\newcommand{\Lip}{\operatorname{Lip}}
\newcommand{\Span}{\operatorname{span}}
\newcommand{\Id}{\mathrm I}
\newcommand{\Norm}[1]{\left\lVert #1\right\rVert}
\newcommand{\Abs}[1]{\left\lvert #1\right\rvert}
\newcommand{\ip}[2]{\left\langle #1,#2\right\rangle}
\newcommand{\dd}{\,\mathrm d}
\newcommand{\UV}{\mathrm{UV}}
\newcommand{\IR}{\mathrm{IR}}
\newcommand{\SM}{\mathrm{SM}}
\newcommand{\BV}{\mathrm{BV}}
\newcommand{\BRST}{\mathrm{BRST}}
\newcommand{\ov}{\mathrm{ov}}
\newcommand{\hard}{\mathrm h}
\newcommand{\soft}{\mathrm s}
\newcommand{\wt}{\mathrm{wt}}
\newcommand{\loc}{\mathrm{loc}}
\newcommand{\remn}{\mathrm{rem}}
\newcommand{\ord}{\mathrm{ord}}
\newcommand{\Op}{\mathrm{Op}}
\newcommand{\Bell}{\mathrm B}
\newcommand{\Part}{\Pi}

\title{Renewal Perturbative Einstein--Standard Model Continuum\\
Research Notes V1\\[0.4em]
\large Native chiral hard-shell reduction: overlap functional calculus, local line frames, and the exact Wilson-symbol bottleneck}
\author{}
\date{21 September 2026}

\begin{document}
\maketitle
\tableofcontents

\section{Context, target, and the V1 attack}

\subsection{Ultimate coefficientwise target}

The programme is to establish a \emph{source-complete coefficientwise perturbative continuum} for the Renewal Einstein--Standard Model regulator, with gravity treated as an effective field theory.  At every prescribed finite loop order, EFT/derivative order, source order, and finite external observable/source panel, one seeks a single renormalized action/source/reference family whose normalized connected coefficients possess cutoff-independent continuum limits and satisfy the coupled BV/BRST/master identities and the measured shell-composition law.

The quantifiers are essential.  The target is coefficientwise at every fixed finite perturbative order.  It does \emph{not} assert convergence of the perturbation series, a uniform limit as loop order tends to infinity, finite-dimensional renormalizability of Einstein gravity, asymptotic safety, a nonperturbative quantum-gravity measure, a nonperturbative Standard-Model completion, or removal of a massless infrared reference scale unless that is separately proved.

\subsection{Checkpoint inherited into this new lineage}

The supplied V80 notes close the bosonic Einstein--$SU(3)\times SU(2)\times U(1)_Y$--Higgs two-loop local/range/source self-map at the declared finite weight and source order.  The programme statement further declares the next chiral checkpoint as already containing, at its stated scope:
\begin{enumerate}[label=\textup{(I\arabic*)},leftmargin=3.2em]
\item the three-generation Standard Model chiral representation;
\item cancellation of the five ordinary four-dimensional local gauge-anomaly coefficients in the declared hypercharge convention;
\item even $SU(2)$ doublet parity;
\item exact determinant Schur factorization and scale cocycle under finite regrouping of nested shells;
\item the analogous optional neutral-Majorana Pfaffian factorization/cocycle;
\item differentiated determinant/Pfaffian line-source identities;
\item the bosonic two-loop self-map inherited from V80.
\end{enumerate}
These statements are treated as \emph{inherited}.  They are not reproved below.

The Grand Tensor manuscript also supplies a literal protected overlap/Yukawa field shell.  In its notation one has a gauge-covariant Hermitian Wilson kernel $H_W$, a protected spectral window
\begin{equation}
 g_*\Id\preceq |H_W|\preceq M_*\Id,
 \label{eq:v1-GT-Wilson-window}
\end{equation}
the overlap sign and moving chiral projector
\begin{equation}
 \eps=\sgn H_W,
 \qquad
 D_{\ov}=a^{-1}(\Id+\gamma_5\eps),
 \qquad
 \widehat P_-={\tfrac12}(\Id+\eps),
 \label{eq:v1-GT-overlap}
\end{equation}
and the chiral overlap--Yukawa block
\begin{equation}
 D_\chi=P_+D_{\ov}\widehat P_-+\cY(H;\mathbf Y),
 \qquad
 s_{\min}(D_\chi)\ge \tau_*>0
 \label{eq:v1-GT-chiral-block}
\end{equation}
on a protected finite chart.  The manuscript also explicitly requires the matter, ghost, mediator and chiral operators to be differentiated with respect to one common coframe family when common stress/source responses are formed.  These are native finite-regulator data.  They are not continuum replacements.

\subsection{Programme E1--E6}

For orientation, the continuation programme is:
\begin{description}[leftmargin=3.6em,style=nextline]
\item[E1.] Prove the native chiral hard-shell theorem: the actual spin--Dirac/Yukawa annular symbol, Wilson/regulator completion, gauge and spin transport, geometric/gauge/Higgs/antifield/line derivatives, a local chiral line frame, a hard-pivot margin, finite-to-ordinary comparison, and rooted/quasilocal estimates.
\item[E2.] Use E1 plus the inherited anomaly/line algebra to restore the complete chiral two-loop shell and obtain one common two-loop ESM action/source/line update.
\item[E3.] Prove the complete common BV/master/source equation with one finite counterterm action, including diffeomorphism, Lorentz, internal gauge, Higgs/Yukawa, chiral line, ghost/antifield, measure/reference, stress and contact rows.
\item[E4.] Formulate and prove the arbitrary fixed-loop induction, including the finite EFT bank, source/contact tower, forests, line data, counterterms, rooted remainders and reference data.
\item[E5.] Compose arbitrarily many shells at fixed loop/EFT/source order and prove a coefficientwise multi-shell continuum for normalized connected functionals.
\item[E6.] Prove scheme/reference compatibility and order compatibility, then state the final coefficientwise ESM continuum theorem with its retained IR/reference hypotheses.
\end{description}

\subsection{What V1 attacks}

V1 attacks E1, but does so in an anti-circular order.  The protected overlap formula \eqref{eq:v1-GT-overlap} is a nonlinear functional calculus of the native Wilson kernel.  Consequently the first task is to prove, without replacing $H_W$ by a continuum Dirac operator, that a local gapped Wilson kernel transfers locality and all fixed source derivatives to $\eps$, $\widehat P_-$ and $D_\chi$.  We then construct a local Kato frame for the moving chiral source space, prove a hard-pivot stability theorem, and derive the determinant/Pfaffian derivative \emph{bounds} that are missing from the already-inherited exact factorization identities.

The finite-to-ordinary annular estimate has one further input: an explicit native flat Wilson symbol and its small-$a|p|$ comparison.  The attached Grand Tensor statement gives the gap \eqref{eq:v1-GT-Wilson-window} and covariance of $H_W$, but not an explicit stencil/symbol with the derivative estimates needed for this comparison.  V1 therefore proves a reduction theorem showing exactly which Wilson-symbol estimate remains load-bearing.  It does not silently choose a different Wilson operator.

\section{Native operator class and norms}

\subsection{Finite cell carrier and source coordinates}

Let $X_a$ be a finite four-dimensional periodic cell complex of mesh $a$, with uniformly bounded local geometry and polynomial ball growth
\begin{equation}
 \#\{y:d_a(x,y)\le R\}\le C_{\rm vol}(1+R/a)^4.
 \label{eq:v1-volume-growth}
\end{equation}
The one-particle carrier $\cH_a$ is the finite direct sum over cells of the spin, generation, colour, weak and hypercharge representation fibres required by the declared chiral Standard Model packet.  The optional real/Majorana carrier is denoted $\cH_a^M$.

Let $\zeta=(\zeta^A)$ denote a finite list of normalized local coordinates drawn from the retained coframe/spin connection, internal gauge links, Higgs/Yukawa variables, BV/antifield sources and line-source coordinates required at the fixed two-loop/source order.  A multi-index $I=(A_1,\ldots,A_r)$ denotes the mixed derivative
\[
 \partial_I=\partial_{A_1}\cdots\partial_{A_r}.
\]
The source coordinates are finite in number in V1.  No analytic infinite-source radius is claimed.

\subsection{Weighted Schur algebra}

For a matrix kernel $K=(K(x,y))_{x,y\in X_a}$ and $\eta\ge0$, define
\begin{equation}
 \Norm{K}_{\eta,1}
 :=\max\left\{
 \sup_x\sum_y e^{\eta d_a(x,y)}\Norm{K(x,y)},
 \sup_y\sum_x e^{\eta d_a(x,y)}\Norm{K(x,y)}
 \right\}.
 \label{eq:v1-weighted-Schur}
\end{equation}
Finite fibre-component sums are included in the matrix norm.  The same definition is used for rectangular blocks between fixed finite fibre types.

\begin{lemma}[Weighted Schur algebra]
\label{lem:v1-Schur-algebra}
For compatible kernels $K,L$,
\begin{equation}
 \Norm{KL}_{\eta,1}\le \Norm K_{\eta,1}\Norm L_{\eta,1}.
 \label{eq:v1-Schur-product}
\end{equation}
The adjoint has the same norm.  In particular every finite product of quasilocal kernels remains quasilocal with no volume factor.
\end{lemma}
\begin{proof}
Since $d_a(x,z)\le d_a(x,y)+d_a(y,z)$,
\[
 e^{\eta d_a(x,z)}\Norm{(KL)(x,z)}
 \le\sum_y e^{\eta d_a(x,y)}\Norm{K(x,y)}
            e^{\eta d_a(y,z)}\Norm{L(y,z)}.
\]
Sum first over $z$ and then over $y$ to obtain the row estimate; reverse the order for the column estimate.  Taking the maximum proves \eqref{eq:v1-Schur-product}.
\end{proof}

\begin{definition}[Native local Wilson class]
\label{def:v1-native-Wilson-class}
A family $H_W(\zeta)$ belongs to the V1 native local Wilson class on a protected parameter chart $\cU$ if:
\begin{enumerate}[label=\textup{(W\arabic*)},leftmargin=3.2em]
\item $H_W(\zeta)=H_W(\zeta)^*$ and \eqref{eq:v1-GT-Wilson-window} holds uniformly for $\zeta\in\cU$;
\item there is a fixed range $R_W$ in cell units such that $H_W(x,y)=0$ for $d_a(x,y)>R_Wa$;
\item for every nonempty source multi-index $I$ required at the fixed source order, $\partial_IH_W$ exists, has a fixed finite range $R_Ia$, and
\begin{equation}
 \sup_{\zeta\in\cU}\Norm{\partial_IH_W(\zeta)}_{0,1}\le h_I<\infty;
 \label{eq:v1-H-derivative-stock}
\end{equation}
\item the exact gauge/spin covariance of the native stencil is retained: local internal-gauge and spin-frame changes act by the declared unitary similarities on $H_W$ and its source derivatives;
\item the Yukawa block $\cY(H;\mathbf Y)$ and all required source derivatives are onsite or fixed-range kernels with finite stocks $y_I$.
\end{enumerate}
The constants may depend on the fixed perturbative/source order and on the protected coefficient chart, but not on volume or shell number.
\end{definition}

The finite-range assumption is the natural form of ``Wilson/spin completion actually used by the native discretization''.  It is deliberately separated from the still-unproved annular \emph{symbol comparison}; locality of the operator does not by itself identify its continuum principal symbol.

\section{Gapped Wilson functional calculus is uniformly quasilocal}

\subsection{A finite-volume Combes--Thomas estimate}

\begin{lemma}[Finite-range resolvent decay]
\label{lem:v1-CT}
Let $H=H^*$ have range at most $R_0a$ and $\Norm H_{0,1}\le M_0$.  If
$\delta=\dist(z,\sigma(H))>0$, then there are constants
$\eta_0=\eta_0(\delta,M_0,R_0)>0$ and $C_0=C_0(\delta,M_0,R_0,C_{\rm vol})$
such that, uniformly in the finite period,
\begin{equation}
 \Norm{(z-H)^{-1}}_{\eta,1}\le C_0
 \qquad(0\le\eta\le\eta_0/(2a)).
 \label{eq:v1-CT-Schur}
\end{equation}
One may take $\eta_0>0$ so small that
$M_0(e^{\eta_0R_0}-1)\le\delta/2$.
\end{lemma}
\begin{proof}
Fix a base cell $x_0$ and a truncated distance function
$\rho_R(x)=\min\{d_a(x,x_0),R\}$.  Conjugation by
$W_\theta=e^{\theta\rho_R}$ changes a nonzero matrix entry of $H$ by a factor whose modulus differs from one by at most $e^{\theta R_0a}-1$.  Hence
\[
 \Norm{W_\theta HW_\theta^{-1}-H}_{\ell^2\to\ell^2}
 \le \Norm{W_\theta HW_\theta^{-1}-H}_{0,1}
 \le M_0(e^{\theta R_0a}-1).
\]
For $\theta=\eta_0/a$ the right side is at most $\delta/2$.  The resolvent identity then gives
\[
 \Norm{W_\theta(z-H)^{-1}W_\theta^{-1}}_{2\to2}\le 2/\delta.
\]
Applying this estimate to a unit vector at $y$ and evaluating its component at $x$ yields the pointwise Combes--Thomas bound
\[
 \Norm{(z-H)^{-1}(x,y)}\le {2\over\delta}e^{-\eta_0d_a(x,y)/a}.
\]
Choose $0\le\eta<\eta_0/a$.  By \eqref{eq:v1-volume-growth}, summing
$e^{-(\eta_0/a-\eta)d_a(x,y)}$ over $y$ gives a constant depending only on the displayed stocks and the gap between the two exponents, not on the period.  The same argument for columns proves \eqref{eq:v1-CT-Schur}.  The truncation $R\to\infty$ is harmless because the finite carrier is finite.
\end{proof}

\subsection{Sign function and all fixed source derivatives}

Choose two positively oriented contours $\Gamma_+$ and $\Gamma_-$ surrounding respectively $[g_*,M_*]$ and $[-M_*,-g_*]$, each at distance at least $\delta_*>0$ from these intervals and from the other component.  Put $f(z)=+1$ inside $\Gamma_+$ and $f(z)=-1$ inside $\Gamma_-$.  Then
\begin{equation}
 \eps=f(H_W)
 ={1\over 2\pi i}\int_{\Gamma_+\cup\Gamma_-}
 f(z)(z-H_W)^{-1}\dd z.
 \label{eq:v1-sign-contour}
\end{equation}

\begin{theorem}[Uniform quasilocality of the native overlap sign]
\label{thm:v1-sign-locality}
Suppose $H_W$ lies in the native local Wilson class.  For every fixed source order $r$ there exists $\eta_r>0$ such that for every multi-index $I$ with $|I|\le r$,
\begin{equation}
 \boxed{
 \sup_{\zeta\in\cU}\Norm{\partial_I\eps(\zeta)}_{\eta_r,1}
 \le C_I(g_*,M_*,R_W,\{h_B: \varnothing\ne B\subseteq I\}).}
 \label{eq:v1-sign-derivative-bound}
\end{equation}
The constants are independent of the volume and shell number.  The same bound holds for $\widehat P_-=(\Id+\eps)/2$.  Gauge and local-spin covariance are preserved by the functional calculus and by every source derivative.
\end{theorem}
\begin{proof}
For $I=\varnothing$, use \eqref{eq:v1-sign-contour} and Lemma~\ref{lem:v1-CT} on the fixed contours.  Their distance from the spectrum is uniformly positive by \eqref{eq:v1-GT-Wilson-window}, so the contour resolvents possess a common weighted Schur bound.

For a nonempty ordered source list $I$, repeated differentiation of
$R(z)=(z-H_W)^{-1}$ gives a finite sum over ordered set partitions
$I=B_1\sqcup\cdots\sqcup B_k$ of terms
\begin{equation}
 R(z)(\partial_{B_1}H_W)R(z)\cdots
       (\partial_{B_k}H_W)R(z),
 \label{eq:v1-resolvent-derivative-word}
\end{equation}
with the usual multiplicities generated when derivatives hit an already differentiated $H_W$.  This follows inductively from
$\partial_AR=R(\partial_AH_W)R$.  Every derivative vertex in
\eqref{eq:v1-resolvent-derivative-word} is fixed range with finite weighted Schur norm for sufficiently small common $\eta_r$, and every resolvent has the bound of Lemma~\ref{lem:v1-CT}.  Lemma~\ref{lem:v1-Schur-algebra} therefore bounds each word independently of volume.  There are finitely many words at fixed $r$.  Integrating them over the fixed contours gives \eqref{eq:v1-sign-derivative-bound}.

If $H_W(g\cdot\zeta)=U_gH_W(\zeta)U_g^{-1}$, holomorphic functional calculus gives
$f(H_W(g\cdot\zeta))=U_gf(H_W(\zeta))U_g^{-1}$.  Differentiation gives the corresponding covariant source identities.  The same argument applies to local spin-frame similarities.
\end{proof}

\begin{corollary}[Quasilocal overlap/Yukawa derivatives before annular scaling]
\label{cor:v1-overlap-raw-locality}
At every fixed source order,
\begin{equation}
 \Norm{\partial_ID_{\ov}}_{\eta,1}
 \le a^{-1}C_I,
 \qquad
 \Norm{\partial_ID_\chi}_{\eta,1}
 \le a^{-1}C'_I+y_I.
 \label{eq:v1-raw-overlap-bound}
\end{equation}
The estimate is volume-uniform but is \emph{not} yet an acceptable continuum annular estimate because of the explicit $a^{-1}$ scale.
\end{corollary}
\begin{proof}
Differentiate \eqref{eq:v1-GT-overlap}--\eqref{eq:v1-GT-chiral-block}, use Theorem~\ref{thm:v1-sign-locality}, the fixed finite matrices $P_+$ and $\gamma_5$, and Lemma~\ref{lem:v1-Schur-algebra}.  The Yukawa derivatives use (W5).
\end{proof}

\begin{firewall}[Why the gap/locality theorem does not prove E1]
\label{fire:v1-a-inverse}
The $a^{-1}$ in \eqref{eq:v1-raw-overlap-bound} is not a defect in the overlap construction; on the physical low-momentum branch it is cancelled by the first-order zero of the overlap numerator.  That cancellation is a statement about the \emph{native flat symbol and its Wilson completion}.  It cannot be inferred from a spectral gap and finite range alone.  Therefore Theorem~\ref{thm:v1-sign-locality} proves native quasilocality and source differentiability but does not by itself prove the finite-to-ordinary Dirac comparison required by E1.
\end{firewall}

\section{A local chiral frame on the protected perturbative component}

The moving source space $\Ran\widehat P_-(\zeta)$ must be framed before a scalar determinant representative and its differentiated line sources are estimated.  A fixed basis in the ambient carrier is not a legitimate frame of a moving subspace.

\subsection{Kato transport}

Fix a base point $\zeta_0\in\cU$.  For $\zeta$ in a star-shaped subchart, write
$P(t)=\widehat P_-(\zeta_0+t(\zeta-\zeta_0))$ and solve
\begin{equation}
 \dot U(t)=[\dot P(t),P(t)]U(t),
 \qquad U(0)=\Id.
 \label{eq:v1-Kato-ODE}
\end{equation}

\begin{lemma}[Kato intertwining]
\label{lem:v1-Kato}
The solution of \eqref{eq:v1-Kato-ODE} is unitary and satisfies
\begin{equation}
 P(t)U(t)=U(t)P(0).
 \label{eq:v1-Kato-intertwine}
\end{equation}
Hence $U(1)$ transports an orthonormal basis of $\Ran P(0)$ to one of $\Ran P(1)$.
\end{lemma}
\begin{proof}
Since $P=P^*=P^2$, the commutator $[\dot P,P]$ is anti-Hermitian, so
$\frac d{dt}(U^*U)=0$.  Next differentiate
$F(t)=P(t)U(t)-U(t)P(0)$.  From $P\dot PP=0$, obtained by differentiating $P^2=P$ and multiplying by $P$ on both sides, one verifies
\[
 \dot P+P[\dot P,P]-[\dot P,P]P=0.
\]
Substitution into $\dot F$ gives
$\dot F=[\dot P,P]F$ with $F(0)=0$, hence $F=0$.
\end{proof}

\begin{theorem}[Protected local chiral line frame with source bounds]
\label{thm:v1-local-line-frame}
At every fixed source order $r$, after shrinking the perturbative coefficient chart if necessary, Kato transport defines a smooth local frame of
$\Ran\widehat P_-(\zeta)$ and hence a smooth nonzero frame $\lambda_\chi(\zeta)$ of
\begin{equation}
 \det(P_+\cH_a)\otimes\det(\widehat P_-(\zeta)\cH_a)^*.
 \label{eq:v1-det-line}
\end{equation}
Its normalized local connection/source coefficients through order $r$ are finite sums of products of derivatives of $\widehat P_-$ and obey volume-uniform quasilocal bounds inherited from \eqref{eq:v1-sign-derivative-bound} after one source root is fixed.

This is a \emph{local perturbative frame theorem}.  It makes no assertion about holonomy around arbitrary large gauge loops or about crossing the divisor $\det D_\chi=0$.
\end{theorem}
\begin{proof}
Theorem~\ref{thm:v1-sign-locality} gives smoothness of $P(\zeta)$ and finite bounds for every required derivative.  The ODE \eqref{eq:v1-Kato-ODE} has a unique smooth solution on a sufficiently small finite parameter chart.  Lemma~\ref{lem:v1-Kato} constructs the moving orthonormal frame.

Fix an oriented orthonormal basis $e_1^0,\ldots,e_m^0$ of $\Ran P(0)$ and a fixed oriented basis $f_1,\ldots,f_m$ of $P_+\cH_a$.  Put
\[
 e_k(\zeta)=U(1;\zeta)e_k^0,
 \qquad
 \lambda_\chi(\zeta)=
 (f_1\wedge\cdots\wedge f_m)\otimes
 (e_1(\zeta)\wedge\cdots\wedge e_m(\zeta))^*.
\]
This is a nonzero local line frame.  Differentiating the ODE in the source variables gives an inhomogeneous linear ODE whose coefficients are finite sums of products of $\partial_BP$.  Duhamel iteration therefore expresses every fixed derivative of $U$ as a finite time-ordered sum of such products.  The rooted source norm fixes one differentiated local insertion; the remaining kernels are composed in the weighted Schur algebra, so Theorem~\ref{thm:v1-sign-locality} and Lemma~\ref{lem:v1-Schur-algebra} give the claimed volume-uniform bounds.  The construction uses one chosen radial path in a local chart and therefore does not identify global line holonomy.
\end{proof}

\section{Hard-shell pivot: what the global chiral floor does and does not imply}

The global lower singular-value bound in \eqref{eq:v1-GT-chiral-block} does not automatically imply that an arbitrary compression to a hard subspace is invertible.  This distinction is load-bearing for exact Schur elimination.

\subsection{Compression can destroy invertibility}

\begin{lemma}[A global gap does not imply a compressed hard pivot]
\label{lem:v1-compression-counterexample}
There exists an invertible matrix $D$ with $s_{\min}(D)=1$ and an orthogonal projection $Q$ such that the compression $QDQ:\Ran Q\to\Ran Q$ is zero.
\end{lemma}
\begin{proof}
Take
$D=\begin{psmallmatrix}0&1\\1&0\end{psmallmatrix}$ and
$Q=\begin{psmallmatrix}1&0\\0&0\end{psmallmatrix}$.
Then $D$ is unitary while $QDQ=0$ on $\Ran Q$.
\end{proof}

Thus the Grand Tensor chiral floor $s_{\min}(D_\chi)\ge\tau_*$ protects the full finite chiral operator and its line section, but an E1 shell integration additionally requires a \emph{hard-block} pivot theorem.

\subsection{Annular hard-pivot stability}

Let $Q_j$ be the declared smooth hard annulus projector/multiplier at scale $\Lambda_j$, with complement $\bar Q_j=\Id-Q_j$.  Write the native chiral operator in block form
\begin{equation}
 D_\chi=
 \begin{pmatrix}
 A_j & B_j\\ C_j&E_j
 \end{pmatrix},
 \qquad
 A_j=Q_jD_\chi Q_j\big|_{\Ran Q_j}.
 \label{eq:v1-hard-block}
\end{equation}

\begin{theorem}[Hard-pivot inheritance from an annular principal bound]
\label{thm:v1-hard-pivot}
Suppose a reference annular operator $A_j^{(0)}$ on $\Ran Q_j$ satisfies
\begin{equation}
 s_{\min}(A_j^{(0)})\ge c_D\Lambda_j
 \label{eq:v1-reference-pivot}
\end{equation}
and the complete native perturbation, including gauge/spin transport, coframe variation, Yukawa insertion and the difference between the exact native hard block and the reference annular block, obeys
\begin{equation}
 \Norm{A_j-A_j^{(0)}}_{2\to2}
 \le \theta c_D\Lambda_j,
 \qquad 0\le\theta<1.
 \label{eq:v1-pivot-perturbation}
\end{equation}
Then
\begin{equation}
 \boxed{
 s_{\min}(A_j)\ge(1-\theta)c_D\Lambda_j,
 \qquad
 \Norm{A_j^{-1}}_{2\to2}\le{1\over(1-\theta)c_D\Lambda_j}.}
 \label{eq:v1-hard-pivot-bound}
\end{equation}
If, in addition, $A_j$ and all required source derivatives lie in one weighted Schur algebra and the inverse has a uniform weighted Schur bound at zero source, then every fixed mixed source derivative of $A_j^{-1}$ is bounded by
\begin{equation}
 \boxed{
 \Norm{\partial_IA_j^{-1}}_{\eta,1}
 \le
 \sum_{k=1}^{|I|}
 \sum_{I=B_1\sqcup\cdots\sqcup B_k}^{\rm ord}
 \Norm{A_j^{-1}}_{\eta,1}^{k+1}
 \prod_{q=1}^k\Norm{\partial_{B_q}A_j}_{\eta,1}.}
 \label{eq:v1-inverse-derivative-bound}
\end{equation}
For $I=\varnothing$ the right side is interpreted as $\Norm{A_j^{-1}}_{\eta,1}$.
\end{theorem}
\begin{proof}
For $v\in\Ran Q_j$,
\[
 \Norm{A_jv}
 \ge \Norm{A_j^{(0)}v}-\Norm{(A_j-A_j^{(0)})v}
 \ge(1-\theta)c_D\Lambda_j\Norm v,
\]
which proves \eqref{eq:v1-hard-pivot-bound}.  Differentiating
$A_jA_j^{-1}=\Id$ gives
$\partial_AA_j^{-1}=-A_j^{-1}(\partial_AA_j)A_j^{-1}$.
Repeated differentiation generates exactly the ordered set-partition words in
\eqref{eq:v1-inverse-derivative-bound}.  Lemma~\ref{lem:v1-Schur-algebra} bounds every word.
\end{proof}

\begin{remark}[What must be shown rather than assumed]
The nontrivial E1 content in \eqref{eq:v1-reference-pivot}--\eqref{eq:v1-pivot-perturbation} is that the \emph{native} overlap/Wilson symbol on the hard annulus has the correct first-order physical branch and that the gauge/coframe/Yukawa perturbations are small relative to $\Lambda_j$ on the protected perturbative component.  The full-operator floor $\tau_*$ alone cannot replace this annular statement.
\end{remark}

\section{Finite-to-ordinary comparison: a transfer theorem}

We now isolate the precise symbol statement that is sufficient for E1.  This section is formulated so that the only unproved input is a comparison for the native \emph{Wilson/overlap symbol itself}; once that input is established, the inverse and kernel estimates follow without further regulator substitution.

\subsection{Normalized annular symbol seminorms}

Let
\begin{equation}
 \cA_j=\{p:c_-\Lambda_j\le |p|\le c_+\Lambda_j\},
 \qquad
 a\Lambda_j\le c_0<1,
 \label{eq:v1-annulus}
\end{equation}
with a smooth profile $\chi_j(p)=\chi(p/\Lambda_j)$ supported in a slightly larger annulus.  Let $d=4$.

For a matrix symbol $F(p)$ define
\begin{equation}
 \Norm{F}_{j,m}^{(q)}
 :=\max_{|\alpha|\le m}
 \sup_{p\in\supp\chi_j}
 \Lambda_j^{|\alpha|-q}\Norm{\partial_p^\alpha F(p)}.
 \label{eq:v1-symbol-seminorm}
\end{equation}
Thus a first-order Dirac symbol has $q=1$, its inverse has $q=-1$, and a lattice-continuum error of size $a\Lambda_j^2$ has normalized size $a\Lambda_j$ in the $q=1$ seminorm.

\begin{assumption}[Native Wilson-to-ordinary annular row]
\label{ass:v1-native-symbol-row}
There exists an ordinary reference chiral symbol $d_{\chi,0}(p)$ on the same physical branch and, for every fixed source multi-index $I$ required by the two-loop packet, a native flat/background symbol $d_{\chi,a}^{(I)}(p)$ obtained by differentiating the \emph{actual} finite Wilson/overlap/Yukawa operator, such that for some $m>s+5$,
\begin{align}
 \Norm{d_{\chi,a}}_{j,m}^{(1)}
 +\Norm{d_{\chi,0}}_{j,m}^{(1)}&\le C_m,
 \label{eq:v1-symbol-size}\\
 \Norm{d_{\chi,a}-d_{\chi,0}}_{j,m}^{(1)}&\le C_m a\Lambda_j,
 \label{eq:v1-symbol-error}\\
 s_{\min}(d_{\chi,a}(p))+s_{\min}(d_{\chi,0}(p))&\ge 2c_D\Lambda_j
 \quad(p\in\supp\chi_j).
 \label{eq:v1-symbol-pivot}
\end{align}
For differentiated source vertices, the same statements hold after division by their declared normalized source stocks.  Gauge parallel transport, spin/coframe transport and Higgs/Yukawa insertions are those of the native stencil; no continuum transporter is inserted on the left side.
\end{assumption}

The assumption is intentionally explicit.  It is precisely the point at which the actual Wilson stencil and the actual geometric link/spin completion enter.

\subsection{Inverse symbol comparison}

\begin{lemma}[Annular inverse symbol bounds]
\label{lem:v1-inverse-symbol}
Under Assumption~\ref{ass:v1-native-symbol-row}, for every $|\alpha|\le m$,
\begin{align}
 \Norm{\partial_p^\alpha d_{\chi,a}^{-1}(p)}
 +\Norm{\partial_p^\alpha d_{\chi,0}^{-1}(p)}
 &\le C_\alpha\Lambda_j^{-1-|\alpha|},
 \label{eq:v1-inverse-symbol-size}\\
 \Norm{\partial_p^\alpha
 (d_{\chi,a}^{-1}-d_{\chi,0}^{-1})(p)}
 &\le C_\alpha a\Lambda_j^{-|\alpha|}.
 \label{eq:v1-inverse-symbol-error}
\end{align}
The analogous fixed source derivatives are bounded by finite polynomials of the normalized native vertex stocks.
\end{lemma}
\begin{proof}
The zeroth inverse bound follows from \eqref{eq:v1-symbol-pivot}.  Differentiate
$d^{-1}d=\Id$.  Every $p$-derivative of $d^{-1}$ is a finite sum of words
\[
 d^{-1}(\partial^{\beta_1}d)d^{-1}\cdots
 (\partial^{\beta_k}d)d^{-1},
 \qquad
 \beta_q\ne0,
 \quad \sum_q\beta_q=\alpha.
\]
Using \eqref{eq:v1-symbol-size} and the zeroth inverse bound gives one factor
$\Lambda_j^{-1-|\alpha|}$.

For the difference use the resolvent identity
\begin{equation}
 d_{\chi,a}^{-1}-d_{\chi,0}^{-1}
 =-d_{\chi,a}^{-1}(d_{\chi,a}-d_{\chi,0})d_{\chi,0}^{-1}.
 \label{eq:v1-inverse-resolvent-id}
\end{equation}
At zeroth order the right side is
$O(\Lambda_j^{-1})O(a\Lambda_j^2)O(\Lambda_j^{-1})=O(a)$.
Differentiate \eqref{eq:v1-inverse-resolvent-id}.  If a total of $|\alpha|$ momentum derivatives are distributed among the three factors, the engineering exponents add to $-|\alpha|$ after the two inverse factors and the one $q=1$ error factor are combined.  This proves \eqref{eq:v1-inverse-symbol-error}.  Source derivatives are handled by the same ordered-partition inverse formula as in \eqref{eq:v1-inverse-derivative-bound}; fixed source order leaves finitely many words.
\end{proof}

\subsection{Kernel comparison}

Let $S_{a,j}$ and $S_{0,j}$ be the inverse Fourier kernels of
$\chi_jd_{\chi,a}^{-1}$ and $\chi_jd_{\chi,0}^{-1}$ in the physical lattice and ordinary measures respectively, periodized on the same physical torus when a finite-volume comparison is made.

\begin{theorem}[Annular finite-to-ordinary quasilocal comparison]
\label{thm:v1-kernel-comparison}
Under Assumption~\ref{ass:v1-native-symbol-row}, for every fixed $s\ge0$ and sufficiently many symbol derivatives $m>s+5$,
\begin{align}
 a^4\sup_x\sum_y
 (1+\Lambda_jd_a(x,y))^s\Norm{S_{a,j}(x,y)}
 &\le C_s\Lambda_j^{-1},
 \label{eq:v1-shell-kernel-row}\\
 a^4\sup_x\sum_y
 (1+\Lambda_jd_a(x,y))^s
 \Norm{S_{a,j}(x,y)-S_{0,j}(x,y)}
 &\le C_s a.
 \label{eq:v1-shell-kernel-diff}
\end{align}
The corresponding column estimates hold.  Every fixed normalized source derivative obeys the same statements with constants replaced by finite polynomials of its vertex stocks.  The estimates are uniform in volume and shell number as long as $a\Lambda_j\le c_0$ and the annular pivot margin stays fixed.
\end{theorem}
\begin{proof}
Multiply \eqref{eq:v1-inverse-symbol-size}--\eqref{eq:v1-inverse-symbol-error} by the smooth profile $\chi_j$.  Leibniz' rule and
$\partial^\beta\chi_j=O(\Lambda_j^{-|\beta|})$ preserve the same derivative scales.  After $p=\Lambda_jq$, the native inverse multiplier has the form
$\Lambda_j^{-1}m_a(q)$ with $m_a$ supported in a fixed annulus and uniformly bounded in $C^m$; the difference has the form
$a\,r_a(q)$ with the same property.

On the infinite lattice, discrete Fourier inversion and $m$ integrations by parts in the dimensionless variable $q$ give
\begin{align*}
 \Norm{S_{a,j}(x,y)}
 &\le C_N\Lambda_j^3(1+\Lambda_jd_a(x,y))^{-N},\\
 \Norm{S_{a,j}(x,y)-S_{0,j}(x,y)}
 &\le C_Na\Lambda_j^4(1+\Lambda_jd_a(x,y))^{-N},
\end{align*}
for every $N\le m$.  The support lies strictly inside the Brillouin zone because $a\Lambda_j\le c_0<1$, so the integrations by parts have no seam contribution.  Choose $N>s+4$.  The physical lattice sum satisfies
\[
 a^4\sum_y(1+\Lambda_jd_a(x,y))^{-N+s}
 \le C\Lambda_j^{-4}
\]
uniformly in $a$ by \eqref{eq:v1-volume-growth}.  Multiplication gives
\eqref{eq:v1-shell-kernel-row}--\eqref{eq:v1-shell-kernel-diff}.

The finite torus kernel is the periodization of the absolutely summable infinite-lattice kernel.  Shortest periodic distance is bounded by the corresponding unwrapped distance, so the same row and column estimates hold without a volume factor.  Fixed source derivatives are finite sums of inverse/vertex words whose normalized symbol derivatives obey the same engineering estimates by Lemma~\ref{lem:v1-inverse-symbol}; apply the identical Fourier argument term by term.
\end{proof}

\begin{corollary}[Rooted connected-graph compatibility]
\label{cor:v1-rooted-graph}
The shell kernels of Theorem~\ref{thm:v1-kernel-comparison} satisfy the row/column input required by the rooted connected-contraction theorem of the earlier perturbative notes.  Hence, once the native source vertices have fixed rooted stocks, every fixed connected fermionic shell graph and every fixed mixed boson--fermion graph containing these hard fermion lines has a volume-uniform rooted bound.  Replacing one native hard fermion line by its ordinary comparison changes the graph by $O(a\Lambda_j)$ relative to the engineering size of that line, before the usual local Taylor subtraction.
\end{corollary}
\begin{proof}
Equation \eqref{eq:v1-shell-kernel-row} is exactly the weighted row/column summability required to orient a spanning tree through the graph.  The connected-contraction proof is insensitive to fermionic signs after absolute values are taken.  Equation \eqref{eq:v1-shell-kernel-diff} gives $O(a)$ in place of the native $O(\Lambda_j^{-1})$ propagator stock, hence the relative factor $a\Lambda_j$.  All remaining factors have the inherited fixed rooted stocks.
\end{proof}

\section{Gauge, coframe, Higgs and antifield derivatives do not create a new locality problem}

E1 lists several source types separately because their covariance and contact terms must all be retained.  Analytically, once the native Wilson stencil itself has the required local derivatives, they enter one common derivative calculus.

\begin{theorem}[Common fixed-order source calculus]
\label{thm:v1-common-source-calculus}
Assume the native Wilson-to-ordinary annular row, and let the source list contain any finite collection of:
\begin{enumerate}[label=\textup{(\alph*)},leftmargin=2.8em]
\item internal gauge-link or gauge-potential sources;
\item coframe, inverse-coframe, spin-connection, lapse/shift or metric-chart sources entering the same finite spin operator;
\item Higgs and Yukawa sources entering $\cY$;
\item ghost, antifield and composite/contact sources which enter the finite chiral operator or its blocking maps locally;
\item local determinant-line frame sources induced by Theorem~\ref{thm:v1-local-line-frame}.
\end{enumerate}
At every fixed mixed source order, the hard inverse, overlap projector, local Kato frame, Schur blocks and line logarithmic derivatives are finite sums of words built from:
\begin{equation}
 A_j^{-1},\quad
 \partial_BH_W,\quad
 \partial_B\cY,\quad
 \partial_BQ_j,\quad
 \partial_B\widehat P_-,
 \label{eq:v1-source-word-bank}
\end{equation}
with $B$ a nonempty sub-multi-index.  Consequently the weighted/rooted bounds close under one finite polynomial of the primitive native stocks.  No source type requires a separate nonlocal inverse provided its entry in \eqref{eq:v1-source-word-bank} is local/quasilocal on the same chart.
\end{theorem}
\begin{proof}
The sign/projector derivatives are Theorem~\ref{thm:v1-sign-locality}; inverse derivatives are Theorem~\ref{thm:v1-hard-pivot}; the line frame is Theorem~\ref{thm:v1-local-line-frame}.  Differentiating a block compression adds derivatives of $Q_j$ and of $D_\chi$.  Differentiating the Schur complement adds derivatives of $A_j^{-1}$ and the off-diagonal blocks.  Each operation is product, sum, or fixed finite contraction.  Lemma~\ref{lem:v1-Schur-algebra} therefore closes the weighted Schur norm, while the inherited rooted connected-contraction theorem closes the external source kernels.  The list is finite at fixed source order, so the resulting bound is a finite polynomial.
\end{proof}

\begin{remark}[Covariance is not optional]
The theorem is analytic.  The gauge and local-spin transformation laws still have to be those of the native stencil.  A small operator derivative with the wrong covariance would not be an admissible E1 vertex.  Definition~\ref{def:v1-native-Wilson-class}(W4) therefore remains a structural hypothesis, not a consequence of the norm estimate.
\end{remark}

\section{Differentiated determinant/Pfaffian line estimates without dropping mixed-shell contacts}

Exact Schur factorization and its differentiated line identities are inherited from the prior checkpoint.  Here we prove the quantitative bounds that follow once the hard pivot and source calculus are available.

\subsection{One hard determinant block}

Let $A(\zeta)$ be an invertible finite hard chiral block in the local line frame.  Put
$\ell_A=\log\det A$, with the branch fixed on the nonzero-pivot component.

\begin{lemma}[All fixed determinant derivatives are inverse--vertex words]
\label{lem:v1-logdet-words}
For every nonempty finite source set $I$,
\begin{equation}
 \partial_I\ell_A
 =\sum_{\pi\in\Part(I)}\ 
   \sum_{\operatorname{cyc}(\pi)}
   c_{\pi,\operatorname{cyc}}
   \Tr\!
   \left(A^{-1}(\partial_{B_1}A)\cdots
         A^{-1}(\partial_{B_k}A)\right),
 \label{eq:v1-logdet-word}
\end{equation}
where $\pi=\{B_1,\ldots,B_k\}$ runs over set partitions of $I$, the second sum runs over finitely many cyclic orderings, and the coefficients are integers depending only on the partition/order.  In particular
\begin{equation}
 \Abs{\partial_I\ell_A}_{\rm root}
 \le C_I\sum_{\pi\in\Part(I)}
 \Norm{A^{-1}}_{\eta,1}^{|\pi|}
 \prod_{B\in\pi}\Norm{\partial_BA}_{\eta,1},
 \label{eq:v1-logdet-bound}
\end{equation}
after one local source insertion is used as the root and the trace density is taken in the same normalized convention as the shell action.
\end{lemma}
\begin{proof}
The first derivative is Jacobi's identity
$\partial_A\ell_A=\Tr(A^{-1}\partial_AA)$.  Differentiate repeatedly.  Every derivative either joins an existing derivative block $\partial_BA$ or hits an inverse, in which case
$\partial A^{-1}=-A^{-1}(\partial A)A^{-1}$ creates a new inverse--vertex pair.  Induction gives precisely the finite partition/cyclic word family \eqref{eq:v1-logdet-word}.  To obtain \eqref{eq:v1-logdet-bound}, anchor one differentiated local vertex.  Sum the remaining internal coordinates cyclically using the row/column weighted Schur bounds.  This is the one-loop trace version of the rooted spanning-tree estimate and has no volume factor in the normalized local trace density.
\end{proof}

\subsection{Schur return and mixed hard/soft contacts}

With the block notation \eqref{eq:v1-hard-block}, the exact soft Schur operator is
\begin{equation}
 S_j=E_j-C_jA_j^{-1}B_j.
 \label{eq:v1-Schur-soft}
\end{equation}

\begin{proposition}[No differentiated Schur contact may be discarded]
\label{prop:v1-Schur-contacts}
For every nonempty source multi-index $I$,
\begin{equation}
 \partial_IS_j
 =\partial_IE_j
 -\sum_{I=I_C\sqcup I_0\sqcup I_B}
 c(I_C,I_0,I_B)
 (\partial_{I_C}C_j)(\partial_{I_0}A_j^{-1})
 (\partial_{I_B}B_j),
 \label{eq:v1-Schur-contact-expansion}
\end{equation}
where empty subindices are allowed and $\partial_{I_0}A_j^{-1}$ is the full ordered-partition expansion of Theorem~\ref{thm:v1-hard-pivot}.  Thus derivatives in which distinct source labels hit $C_j$, the interior hard inverse, and $B_j$ are genuine mixed-shell contact terms.  Under the E1 source stocks they satisfy one common finite polynomial weighted/rooted bound.
\end{proposition}
\begin{proof}
Equation \eqref{eq:v1-Schur-contact-expansion} is the multivariate Leibniz rule applied to \eqref{eq:v1-Schur-soft}.  The inverse derivative formula is exact.  Every term contains only factors from the common source word bank \eqref{eq:v1-source-word-bank}; apply Lemma~\ref{lem:v1-Schur-algebra} and Theorem~\ref{thm:v1-common-source-calculus}.  Since the expansion is an identity, deleting any distribution of source labels would change the differentiated Schur operator.
\end{proof}

\subsection{Many shells: exact Bell-polynomial product estimate}

Let $\sigma_j$ denote the nonzero local determinant-line section contributed by shell $j$ in the inherited scale cocycle, expressed in the local frame of Theorem~\ref{thm:v1-local-line-frame}.  Write
$g_j=\log\sigma_j$ on the chosen perturbative branch.  For a finite interval of shells $[j,n)$,
\begin{equation}
 \sigma_{[j,n)}=\bigotimes_{m=j}^{n-1}\sigma_m,
 \qquad
 g_{[j,n)}=\sum_{m=j}^{n-1}g_m.
 \label{eq:v1-line-product}
\end{equation}

For the following theorem fix the finite external source panel required by the coefficientwise target.  Denote by $\Norm{\cdot}_{\rm pan}$ any fixed cross norm on its finite tensor algebra, chosen submultiplicative.  Connected logarithmic line kernels continue to use the rooted norm.  This distinction is essential: the rooted norm of an extensive connected action is not an algebra norm, whereas the finite-panel tensor algebra may legitimately carry products of disconnected source contacts.

\begin{theorem}[Differentiated line-product bound with all mixed-shell contacts]
\label{thm:v1-line-product}
Fix source order $r$ and the finite source panel just declared.  Suppose that after the common local Taylor/EFT assignment at each shell the renormalized logarithmic line coefficients obey
\begin{equation}
 \Norm{\partial_I g_m^{\remn}}_{\rm pan}
 \le c_I\rho_m,
 \qquad
 0<|I|\le r,
 \qquad
 \sum_{m\ge j}\rho_m<\infty.
 \label{eq:v1-line-log-summable}
\end{equation}
Then every normalized derivative of the composed section is uniformly bounded in the number of shells:
\begin{equation}
 \boxed{
 \Norm{\sigma_{[j,n)}^{-1}\partial_I\sigma_{[j,n)}}_{\rm pan}
 \le
 \sum_{\pi\in\Part(I)}
 \prod_{B\in\pi}
 \left(c_B\sum_{m=j}^{n-1}\rho_m\right).}
 \label{eq:v1-line-Bell-bound}
\end{equation}
In particular the right side is bounded uniformly as $n\to\infty$.  The products over blocks $B\in\pi$ are precisely the mixed-shell line contacts generated by differentiating the exact tensor product; no contact is discarded.

If E1 supplies $\rho_m=C\mu/\Lambda_m$ with $\Lambda_m=b^m\lambda$, then
\begin{equation}
 \sup_{n\ge j}
 \Norm{\sigma_{[j,n)}^{-1}\partial_I\sigma_{[j,n)}}_{\rm pan}
 \le
 \sum_{\pi\in\Part(I)}
 \prod_{B\in\pi}
 \left(c_B{Cb\over b-1}{\mu\over\Lambda_j}\right).
 \label{eq:v1-line-geometric-bound}
\end{equation}
\end{theorem}
\begin{proof}
For scalar representatives in a fixed local line frame,
$\sigma_{[j,n)}=\exp(g_{[j,n)})$.  The multivariate Fa\`a di Bruno formula for the exponential gives the exact identity
\begin{equation}
 \sigma_{[j,n)}^{-1}\partial_I\sigma_{[j,n)}
 =\sum_{\pi\in\Part(I)}
 \prod_{B\in\pi}\partial_B g_{[j,n)}.
 \label{eq:v1-Bell-identity}
\end{equation}
But
$\partial_Bg_{[j,n)}=\sum_{m=j}^{n-1}\partial_Bg_m$.
Taking the submultiplicative finite-panel norm and applying \eqref{eq:v1-line-log-summable} gives \eqref{eq:v1-line-Bell-bound}.  If one instead retains source positions as free kernel variables, \eqref{eq:v1-Bell-identity} is to be read in the disconnected tensor algebra; its connected logarithm remains in the rooted action norm.  No claim that the rooted extensive-action norm itself is an algebra norm is used.  Expanding the product of shell sums in \eqref{eq:v1-Bell-identity} assigns each block of the source partition independently to a shell; these are exactly the same-shell and mixed-shell contacts.  Finally
$\sum_{m\ge j}\mu/\Lambda_m\le \frac b{b-1}\mu/\Lambda_j$ gives \eqref{eq:v1-line-geometric-bound}.
\end{proof}

\subsection{Optional real/Majorana Pfaffian branch}

\begin{proposition}[Pfaffian derivative estimate on an oriented nonzero component]
\label{prop:v1-Pfaffian}
Let $A^M(\zeta)$ be a real skew-symmetric hard block of even rank on a chosen oriented component with
$s_{\min}(A^M)\ge\kappa_j>0$.  Then, in that local orientation frame,
\begin{equation}
 \partial_A\log\Pf A^M
 ={\tfrac12}\Tr\bigl((A^M)^{-1}\partial_AA^M\bigr),
 \label{eq:v1-Pf-first}
\end{equation}
and every higher fixed source derivative obeys the analogue of
\eqref{eq:v1-logdet-bound} with an overall factor $1/2$ at the logarithmic level.  The multi-shell product estimate of Theorem~\ref{thm:v1-line-product} therefore applies verbatim to the local Pfaffian section.

No statement is made across a Pfaffian zero or an orientation-changing loop.
\end{proposition}
\begin{proof}
On an oriented nonzero component, $(\Pf A^M)^2=\det A^M$ and the chosen sign is continuous.  Differentiate
$2\log\Pf A^M=\log\det A^M$ and apply Lemma~\ref{lem:v1-logdet-words}.  The line-product proof uses only the logarithmic derivative and therefore carries over unchanged.
\end{proof}

\section{The exact remaining E1 bottleneck}

We now separate what is proved from what would be circular to assert.

\subsection{Source audit of the attached native data}

The protected field-shell construction in the Grand Tensor supplies:
\begin{enumerate}[label=\textup{(A\arabic*)},leftmargin=3em]
\item a gauge-covariant Hermitian Wilson kernel $H_W$;
\item a uniform protected Wilson spectral window $g_*\Id\preceq |H_W|\preceq M_*\Id$;
\item the overlap sign, overlap operator and moving chiral projector;
\item the native chiral overlap--Yukawa block $D_\chi$;
\item a nonzero full chiral floor $s_{\min}(D_\chi)\ge\tau_*$;
\item finite determinant-line Jacobi identities and a local protected line section;
\item a common coframe differentiation requirement for the assembled source/stress packet.
\end{enumerate}
These data are sufficient for the functional-calculus, local-frame and line-derivative arguments above once the local stencil/source-derivative properties are read from the actual finite action.

What is not supplied in the attached statement is an explicit formula for the native $H_W$ stencil on the common coframe/spin/gauge carrier from which one can prove the annular derivative estimates \eqref{eq:v1-symbol-size}--\eqref{eq:v1-symbol-pivot}.  In particular, the protected spectral window does not determine:
\begin{enumerate}[label=\textup{(B\arabic*)},leftmargin=3em]
\item which Wilson mass/parameter branch carries the physical zero;
\item the coefficient of the first-order $i\gamma^\mu p_\mu$ term after the overlap normalization;
\item the size and tensor structure of the $O(a|p|^2)$ Wilson/spin remainder;
\item the exact coframe/spin-link dependence of that remainder;
\item the differentiated comparison when gauge, geometric, Higgs and antifield sources are inserted;
\item the hard-annular compression margin needed for exact shell Schur elimination.
\end{enumerate}
These are not details that can be supplied by replacing the finite operator with a continuum Dirac operator.

\begin{obstruction}[Gap/locality does not identify the physical branch]
\label{obs:v1-gap-not-symbol}
There exist distinct finite-range Hermitian lattice kernels with the same two-sided spectral gap and the same uniform operator norm but different low-momentum winding, different first-order symbols, or no physical first-order zero after the overlap map.  Therefore \eqref{eq:v1-GT-Wilson-window} plus finite range cannot imply Assumption~\ref{ass:v1-native-symbol-row}.
\end{obstruction}
\begin{proof}
Already in a translation-invariant two-band model, conjugating a gapped reference symbol by a momentum-dependent unitary preserves its eigenvalues and hence its spectral gap while changing the eigenprojector map.  Alternatively, add to a gapped Wilson symbol a finite-range Hermitian trigonometric polynomial of size smaller than $g_*/2$ but with a nonzero first derivative at the physical point.  The gap survives by the singular-value perturbation inequality, while the first-order overlap symbol changes.  Spectral bounds therefore do not determine the required principal-symbol matching.
\end{proof}

\subsection{Reduction theorem}

\begin{theorem}[E1 reduction to the native Wilson-symbol theorem]
\label{thm:v1-E1-reduction}
Fix the finite two-loop/source packet and the protected nonzero-pivot component.  Suppose:
\begin{enumerate}[label=\textup{(R\arabic*)},leftmargin=3em]
\item the actual native $H_W(e,\omega,U)$ is a fixed-range gauge/spin-covariant stencil with all required local source derivatives;
\item its protected spectral window \eqref{eq:v1-GT-Wilson-window} holds;
\item its exact flat/coframe-normalized Wilson/overlap symbol satisfies Assumption~\ref{ass:v1-native-symbol-row}, including the differentiated gauge/geometric/Higgs/antifield rows;
\item the annular perturbation satisfies the hard-pivot inequality \eqref{eq:v1-pivot-perturbation} on the declared perturbative component.
\end{enumerate}
Then E1 is closed at that fixed source/derivative order in the following precise sense:
\begin{enumerate}[label=\textup{(E1.\arabic*)},leftmargin=3.4em]
\item the overlap sign, moving chiral projector and all required source derivatives are cutoff-uniformly quasilocal;
\item a local Kato chiral line frame exists with cutoff-uniform rooted source bounds;
\item the hard fermion block has a nonzero pivot of order $\Lambda_j$ and all inverse source derivatives are bounded;
\item the native annular hard inverse differs from the ordinary reference inverse by $O(a)$ in the physical weighted row/column norm, with the analogous fixed source-derivative estimates;
\item every fixed connected hard-fermion/mixed graph has a volume-uniform rooted bound;
\item the inherited exact determinant/Pfaffian Schur identities acquire the quantitative derivative bounds of Lemma~\ref{lem:v1-logdet-words} and Proposition~\ref{prop:v1-Pfaffian}, with all differentiated Schur contacts retained;
\item once the usual local Taylor/EFT assignment gives a summable one-excess-scale remainder $\rho_j=O(\mu/\Lambda_j)$, the differentiated determinant/Pfaffian line products are uniformly bounded through arbitrary finite shell composition by Theorem~\ref{thm:v1-line-product}.
\end{enumerate}
No global determinant-line holonomy, no zero-mode crossing theorem, and no nonperturbative chiral measure is implied.
\end{theorem}
\begin{proof}
Items (E1.1)--(E1.2) are Theorems~\ref{thm:v1-sign-locality} and
\ref{thm:v1-local-line-frame}.  Item (E1.3) is Theorem~\ref{thm:v1-hard-pivot}.  Item (E1.4) is Theorem~\ref{thm:v1-kernel-comparison}.  Item (E1.5) is Corollary~\ref{cor:v1-rooted-graph}.  Item (E1.6) follows from Lemma~\ref{lem:v1-logdet-words} and Proposition~\ref{prop:v1-Schur-contacts}; the Pfaffian variant is Proposition~\ref{prop:v1-Pfaffian}.  Item (E1.7) is Theorem~\ref{thm:v1-line-product}.  Every input is native except the ordinary reference used only on the right-hand side of a proved comparison.  The theorem never substitutes the ordinary operator for the native hard block in the exact shell algebra.
\end{proof}

\section{How to attack the remaining native Wilson-symbol theorem}

The reduction theorem shows that repeatedly polishing determinant identities would now be circular: their exact algebra is already inherited, and their estimates are controlled once the native annular pivot/comparison is known.  The next work should therefore move one level upstream to the finite Wilson stencil.

\subsection{Required exact stencil data}

For each oriented cell edge $e=(x,x+a\hat\mu)$, let
\begin{equation}
 \cU_e(e,\omega,U)
 =\cU_e^{\rm spin}(e,\omega)\otimes\rho_{\SM}(U_e)
 \label{eq:v1-combined-link}
\end{equation}
be the actual spin--internal parallel transport used by the finite action.  The Wilson kernel must be written in the finite form
\begin{equation}
 (H_W\psi)_x
 =\sum_{|y-x|\le R_Wa} h_{xy}(e,\omega,U)\psi_y,
 \label{eq:v1-native-stencil-template}
\end{equation}
including its diagonal Wilson mass/completion.  The next proof must not infer the coefficients $h_{xy}$ from the desired continuum answer; it must read them from the native finite action.

On the constant coframe, trivial spin connection, identity gauge-link and fixed Higgs background, Fourier transformation of \eqref{eq:v1-native-stencil-template} gives the literal finite trigonometric symbol
\begin{equation}
 h_{W,a}(p)=\sum_{|z|\le R_Wa}h_{0z}\,e^{ip\cdot z}.
 \label{eq:v1-native-flat-symbol}
\end{equation}
The protected overlap symbol is then
\begin{equation}
 \eps_a(p)=h_{W,a}(p)[h_{W,a}(p)^2]^{-1/2},
 \qquad
 d_{\ov,a}(p)=a^{-1}(\Id+\gamma_5\eps_a(p)),
 \label{eq:v1-native-overlap-symbol}
\end{equation}
with the branch fixed by the actual native Wilson mass/completion.

\subsection{The exact next load-bearing theorem}

\begin{theorem}[Target theorem for V2: native Wilson annular symbol]
\label{thm:v1-next-target}
The next load-bearing theorem is to prove directly from
\eqref{eq:v1-native-stencil-template} that on the physical index-zero branch there exist $c_0,c_D>0$ such that, whenever $a\Lambda_j\le c_0$,
\begin{align}
 d_{\chi,a}(p)
 &=d_{\chi,0}(p)+r_a(p),
 \label{eq:v1-next-expansion}\\
 \Norm{\partial_p^\alpha r_a(p)}
 &\le C_\alpha a\Lambda_j^{2-|\alpha|},
 \qquad p\in\supp\chi_j,
 \label{eq:v1-next-remainder}\\
 s_{\min}(d_{\chi,a}(p))
 &\ge c_D\Lambda_j,
 \label{eq:v1-next-pivot}
\end{align}
through the momentum derivative order required by the rooted kernel moment, and with the corresponding differentiated formulas for every gauge, coframe/spin, Higgs/Yukawa, antifield and line-source insertion in the fixed two-loop packet.  The same proof must give a full-zone bound near the lattice ultraviolet cap, where $p$ itself is not a good small variable and $x=ap$ must be used instead.
\end{theorem}

Theorem~\ref{thm:v1-next-target} is \emph{not proved in V1} because the attached sources, as supplied here, do not display the exact coefficient stencil needed to evaluate \eqref{eq:v1-native-flat-symbol}.  Once that stencil is made explicit, the proof should proceed by the following finite calculations:
\begin{enumerate}[label=\textup{(N\arabic*)},leftmargin=3em]
\item compute $h_{W,a}(p)$ exactly on the constant-coframe identity-link branch;
\item locate the physical Wilson branch and prove its gap from the doubler sectors using the actual Wilson mass/completion;
\item Taylor-expand the physical branch through the derivative order consumed by Theorem~\ref{thm:v1-kernel-comparison};
\item differentiate the finite link matrices in the common coframe/spin/gauge chart rather than inserting a continuum covariant derivative by hand;
\item add the onsite/fixed-range Higgs--Yukawa block and verify that it is a lower-order perturbation on a sufficiently hard annulus;
\item establish the $x=ap$ full-zone inverse bounds on the compact Brillouin region away from the physical small-$p$ chart;
\item combine the two charts by a smooth partition, producing the hard-pivot estimate \eqref{eq:v1-pivot-perturbation} and Assumption~\ref{ass:v1-native-symbol-row} as theorems.
\end{enumerate}

\section{Anti-circularity audit}

\begin{firewall}[No continuum substitution]
At no point in V1 is the native shell integral evaluated with a continuum Dirac operator.  The ordinary symbol appears only in Assumption~\ref{ass:v1-native-symbol-row} as the comparator in a finite-to-ordinary estimate.  All exact factorization, inverses, source derivatives and line sections remain those of the finite native $D_\chi$.
\end{firewall}

\begin{firewall}[No anomaly shortcut]
The inherited Standard Model anomaly cancellations remove the nontrivial local anomaly class only after the regulator matching problem is solved.  They do not imply Theorem~\ref{thm:v1-next-target}, hard-shell quasilocality, or a bound on the determinant-line connection.
\end{firewall}

\begin{firewall}[No global line claim]
The Kato frame is defined on one protected perturbative subchart.  V1 does not identify line holonomy across arbitrary gauge loops, across disconnected Wilson sectors, or through zeros of $D_\chi$ or of an optional Majorana skew block.
\end{firewall}

\begin{firewall}[No fixed-order to all-order promotion]
All constants may depend arbitrarily on the fixed source order, local derivative order and perturbative order.  V1 does not provide a uniform large-loop estimate or convergence of the perturbative series.
\end{firewall}

\section{Status ledger and next theorem}

\subsection*{Proved in V1}
\begin{enumerate}[leftmargin=2.2em]
\item A finite-range gapped native Hermitian Wilson kernel transfers cutoff-uniform quasilocality to its overlap sign $\eps=\sgn H_W$, moving projector $\widehat P_-$, and every fixed mixed source derivative, with exact gauge/spin covariance preserved.
\item On the protected perturbative component, Kato transport constructs a legitimate local frame of the moving chiral source space and of the determinant line, with fixed-order rooted source bounds.  This is explicitly local, not a global holonomy theorem.
\item The full chiral floor does not by itself imply an invertible shell compression.  A separate annular hard-pivot theorem was proved: an $O(\Lambda_j)$ reference singular-value margin survives every native perturbation smaller than that margin, and all fixed inverse source derivatives satisfy explicit ordered-partition bounds.
\item Assuming the precise native Wilson-to-ordinary annular symbol row, the hard fermion inverse has the physical weighted row/column bound $O(\Lambda_j^{-1})$ and differs from its ordinary comparator by $O(a)$, uniformly in volume and shell number for $a\Lambda_j\le c_0$.  The same holds for every fixed normalized source derivative.
\item These shell kernels plug into the inherited rooted connected-contraction estimate, yielding volume-uniform fixed-graph fermionic and mixed boson--fermion bounds.
\item The inherited exact determinant/Pfaffian Schur identities were upgraded with quantitative fixed-source derivative bounds.  Differentiation of the Schur complement retains every mixed hard/soft contact.
\item A Bell-polynomial line-product theorem was proved: if the renormalized logarithmic line remainder at shell $j$ is summable (in particular $O(\mu/\Lambda_j)$), then all fixed differentiated determinant/Pfaffian line products are uniformly bounded through arbitrarily many ultraviolet shell steps on every fixed finite source panel, with all mixed-shell source contacts retained; the connected logarithmic coefficients retain the rooted norm.
\item E1 was reduced to one explicit native Wilson-symbol theorem plus its hard-annular perturbation margin; no continuum substitution is needed after that point.
\end{enumerate}

\subsection*{Inherited}
\begin{enumerate}[leftmargin=2.2em]
\item The V80 bosonic Einstein--$SU(3)\times SU(2)\times U(1)_Y$--Higgs two-loop local/range/source self-map at the declared finite weight/source order.
\item The three-generation chiral representation, cancellation of the five ordinary local gauge-anomaly coefficients, and even $SU(2)$ doublet parity at the scope declared by the programme checkpoint.
\item Exact determinant Schur factorization, its finite regrouping cocycle, the optional neutral-Majorana Pfaffian analogue, and their exact differentiated line-source identities.
\item The Grand Tensor protected overlap/Yukawa data: a gauge-covariant Hermitian Wilson kernel with protected gap, overlap sign/projector, chiral block, full nonzero chiral floor, same-operator determinant section, and common coframe differentiation requirement.
\item The earlier rooted connected-contraction calculus, common source normalization, forest ownership and bosonic two-loop annular estimates.
\end{enumerate}

\subsection*{Conditional}
\begin{enumerate}[leftmargin=2.2em]
\item The finite-to-ordinary kernel estimate and the quantitative hard-shell E1 closure are conditional on the native Wilson annular symbol row, because the supplied Grand Tensor excerpt gives $H_W$ abstractly rather than displaying the exact Wilson/spin stencil needed for the symbol calculation.
\item The multi-shell line-product bound requires the locally subtracted logarithmic line remainder to satisfy a summable shell estimate; the desired two-loop value is $O(\mu/\Lambda_j)$.
\item The optional Pfaffian estimate is local to a chosen nonzero oriented real-skew component.
\end{enumerate}

\subsection*{Open}
\begin{enumerate}[leftmargin=2.2em]
\item Write the exact native $H_W(e,\omega,U)$ stencil, including its Wilson mass/completion and combined spin/internal parallel transports, directly from the finite action.
\item Prove Theorem~\ref{thm:v1-next-target}: the physical annular symbol, $O(a|p|^2)$ finite-to-ordinary remainder with all required derivatives, hard-pivot margin, and the complementary full-zone Brillouin estimate.
\item Use that theorem to prove the locally subtracted chiral determinant/Pfaffian line remainder is $O(\mu/\Lambda_j)$ in the same rooted/source norm as the bosonic shell.
\item Only then pass to E2: identify the native chiral BV breaking with the standard anomaly class plus a local BV-exact regulator term and construct one common two-loop ESM restoring action/source/line update.
\end{enumerate}

\subsection*{Exact next load-bearing theorem}

\begin{center}
\fbox{\parbox{0.92\textwidth}{
\textbf{Native Wilson annular symbol theorem.}
Starting from the actual finite $H_W(e,\omega,U)$ stencil used by the Renewal regulator, prove the physical-branch expansion
$d_{\chi,a}=d_{\chi,0}+O(a|p|^2)$ on $|p|\asymp\Lambda_j$, with the required momentum and gauge/coframe/spin/Higgs/antifield/line-source derivatives, an $O(\Lambda_j)$ hard singular-value margin, and a uniform $x=ap$ full-zone bound.  This theorem must be native: no replacement of $H_W$ by a continuum Dirac operator is allowed.
}}
\end{center}

\section*{Append-only continuation marker: V1}
\addcontentsline{toc}{section}{Append-only continuation marker: V1}
\textbf{End of V1.}  Future versions in this lineage must preserve the mathematical body above and append new work.  The next file is to be named
\texttt{Renewal\_Perturbative\_ESM\_Continuum\_Research\_Notes\_V2.tex}.


\clearpage
\section{V2 continuation: an explicit native Wilson completion and coefficientwise closure of E1}
\label{sec:v2-continuation}

V1 isolated the remaining fermionic obstruction at the correct place: the protected
spectral window of an abstract Hermitian Wilson kernel does not determine its
physical low-momentum branch.  There are two possible responses.  One may keep
asking for a stronger theorem about an unspecified kernel, or one may go upstream
and specify the regulator slot itself.  The present continuation takes the second
route.

There is also a quantifier correction.  The endpoint of the programme is a
\emph{coefficientwise} perturbative continuum at fixed source order.  Such a theorem
does not require the hard fermion block to remain invertible on a cutoff-independent
finite-amplitude ball of all gauge, coframe, Higgs and antifield sources.  It requires
an invertible reference hard block and controlled source jets of every fixed order.
The stronger finite-amplitude pivot theorem of V1 remains useful when available,
but it is not load-bearing for the coefficientwise theorem.  This observation avoids
paying a spurious $a^{-1}$ operator-norm cost for a finite coframe displacement when
only its fixed perturbative derivatives are needed.

The second upstream move is to fill the abstract Wilson slot by an explicit
covariant nearest-neighbour completion.  This is a \emph{regulator definition in the
present proof-note lineage}, not an inference that the earlier manuscript had
already written these coefficients.  If a different Wilson completion is intended
elsewhere, it must be compared with the one below before the final theorem is
identified with that different regulator.

\subsection{Coefficientwise inversion needs only the reference pivot}

Let $A(\zeta)$ be a finite-dimensional or finite-cutoff operator family with source
coordinates $\zeta=(\zeta^A)$ and Taylor jets
\begin{equation}
 A_I:=\partial_I A(0).
 \label{eq:v2-source-jets}
\end{equation}
No common source radius is assumed.

\begin{theorem}[Reference-pivot coefficient theorem]
\label{thm:v2-reference-pivot}
Assume only that $A_0:=A(0)$ is invertible.  For every nonempty finite source
multi-index $I$,
\begin{equation}
 \boxed{
 \partial_I A^{-1}(0)
 =\sum_{k=1}^{|I|}(-1)^k
 \sum_{I=B_1\sqcup\cdots\sqcup B_k}^{\rm ord}
 A_0^{-1}A_{B_1}A_0^{-1}\cdots A_{B_k}A_0^{-1}.}
 \label{eq:v2-reference-inverse-jet}
\end{equation}
Likewise, after a local logarithm is chosen at the reference point,
$\partial_I\log\det A(0)$ is a finite sum of cyclic traces of the same
reference inverse and source jets.  Consequently a fixed source coefficient of a
hard-shell graph, Schur complement, determinant line, or Pfaffian line needs only:
\begin{enumerate}[label=\textup{(C\arabic*)},leftmargin=3em]
\item a reference hard inverse $A_0^{-1}$ with the required annular bound;
\item the finitely many jets $A_B$ with $0<|B|\le |I|$;
\item the fixed local line orientation/frame needed to interpret the coefficient.
\end{enumerate}
It does not require a cutoff-independent ball on which $A(\zeta)^{-1}$ is uniformly
bounded.
\end{theorem}
\begin{proof}
Differentiate $A(\zeta)A(\zeta)^{-1}=I$ at $\zeta=0$.  For one derivative the
formula is
$\partial_AA^{-1}(0)=-A_0^{-1}A_AA_0^{-1}$.  Inductively, every new derivative
either joins one existing jet or differentiates an inverse factor, in which case it
creates one further $-A_0^{-1}A_BA_0^{-1}$ insertion.  The ordered set partitions
record exactly these possibilities and prove \eqref{eq:v2-reference-inverse-jet}.
Jacobi's formula at the reference point then gives the determinant statement.
The Pfaffian statement follows from $2\log\Pf=\log\det$ on a chosen nonzero
oriented component.  Every fixed coefficient is therefore an algebraic expression
in the displayed finite jet bank.  A uniform nonzero radius in $\zeta$ would be a
stronger analytic statement and is not used.
\end{proof}

\begin{remark}[Relation to the V1 hard-pivot theorem]
Theorem~\ref{thm:v1-hard-pivot} remains the correct theorem for finite-amplitude
shell integration on a fixed perturbative ball.  Theorem~\ref{thm:v2-reference-pivot}
shows that this stronger quantifier is unnecessary for E1 at the stated
coefficientwise endpoint.  In particular, source derivatives are to be treated as
vertices around the reference hard propagator rather than absorbed into a finite
background before inversion.
\end{remark}

\subsection{A concrete covariant Wilson slot}
\label{subsec:v2-Wilson-definition}

Work first on a periodic cubical regulator branch with spacing $a$.  The local
argument is later inserted into the general finite cell carrier by the already
inherited chart/partition machinery.  Use Hermitian Euclidean Clifford matrices
$\gamma_a$ satisfying
\begin{equation}
 \{\gamma_a,\gamma_b\}=2\delta_{ab}I,
 \qquad
 \gamma_5^*=\gamma_5,
 \qquad
 \{\gamma_5,\gamma_a\}=0.
 \label{eq:v2-Clifford}
\end{equation}
Let $\Gamma_\mu(e)$ be Hermitian coframe Clifford coefficients with
\begin{equation}
 \{\Gamma_\mu,\Gamma_\nu\}=2g_{\mu\nu}I,
 \qquad
 m_g I\preceq g\preceq M_g I,
 \label{eq:v2-coframe-elliptic}
\end{equation}
for fixed $0<m_g\le M_g<\infty$ on the protected reference coframe chart.

Let $\mathcal U_{x,\mu}$ be the combined spin--internal parallel transport on the
edge $x\to x+a\hat\mu$.  After absorbing the positive site Hodge weight into the
one-particle carrier, define the unitary covariant shift
\begin{equation}
 (T_\mu^{\mathcal U}\psi)_x
 =\mathcal U_{x,\mu}\psi_{x+a\hat\mu},
 \qquad
 (T_\mu^{\mathcal U})^*=(T_\mu^{\mathcal U})^{-1}.
 \label{eq:v2-covariant-shift}
\end{equation}
For a nonconstant positive site weight the standard square-root weight ratio is
inserted in \eqref{eq:v2-covariant-shift}; all formulas below are unchanged after
that unitary rescaling.  Put
\begin{equation}
 C_\mu^{\mathcal U}:={1\over2}
 (T_\mu^{\mathcal U}-(T_\mu^{\mathcal U})^*),
 \qquad
 L_\mu^{\mathcal U}:=I-{1\over2}
 (T_\mu^{\mathcal U}+(T_\mu^{\mathcal U})^*).
 \label{eq:v2-centered-Wilson-pieces}
\end{equation}
Thus $C_\mu^{\mathcal U}$ is anti-Hermitian and $L_\mu^{\mathcal U}$ is
Hermitian.  The V2 Wilson completion is
\begin{equation}
 \boxed{
 K_{W,a}(e,\omega,U)
 :=\sum_{\mu=1}^4 {1\over2}
   \{\Gamma_\mu(e),C_\mu^{\mathcal U}\}
   +\sum_{\mu=1}^4L_\mu^{\mathcal U}-I,
 \qquad
 H_{W,a}:=\gamma_5K_{W,a}.}
 \label{eq:v2-native-Wilson}
\end{equation}
This is the Wilson choice $r=\rho=1$.  The choice is deliberate: it gives an
exact all-zone gap identity on the flat/reference branch and normalizes the physical
overlap speed to one without an additional prefactor.

\begin{proposition}[Locality, covariance and Hermiticity of the V2 completion]
\label{prop:v2-Wilson-structure}
Assume the spin transport commutes with $\gamma_5$ and the local spin-frame action
intertwines $\Gamma_\mu$ in the usual way.  Then:
\begin{enumerate}[label=\textup{(W\arabic*)},leftmargin=3em]
\item $H_{W,a}=H_{W,a}^*$;
\item $H_{W,a}$ has range one in lattice units;
\item internal gauge and local spin-frame transformations act by exact unitary
similarity;
\item every fixed derivative with respect to a smooth gauge-link, spin-connection,
coframe or metric-chart coordinate is again a fixed-range kernel.
\end{enumerate}
Hence, on every reference chart on which a spectral gap is present,
\eqref{eq:v2-native-Wilson} belongs to the native local Wilson class of
Definition~\ref{def:v1-native-Wilson-class}.
\end{proposition}
\begin{proof}
The anticommutator of a Hermitian multiplication operator with an anti-Hermitian
shift difference is anti-Hermitian.  The Wilson Laplacian and the diagonal mass are
Hermitian.  Since $\gamma_5$ anticommutes with every $\Gamma_\mu$ and commutes
with the covariant shifts,
\[
 K_{W,a}^*=\gamma_5K_{W,a}\gamma_5.
\]
Therefore $(\gamma_5K_{W,a})^*=\gamma_5K_{W,a}$.  Each term contains at most
one nearest-neighbour shift, proving range one.  Gauge/spin covariance follows by
conjugating the edge transport and the local Clifford multiplication together.
Finally, differentiating a fixed product of local coefficients and one-edge
transports cannot enlarge the range.  Smoothness of the finite coordinate maps
gives the fixed derivative stocks on every compact source panel.
\end{proof}

\begin{firewall}[What has and has not been identified]
Equation~\eqref{eq:v2-native-Wilson} fills the Wilson slot that was abstract in the
supplied Grand Tensor field-shell statement.  The statements below are therefore
native theorems for this explicit regulator completion.  They are not a proof that
an unspecified alternative $H_W$ has the same symbol.  Such an alternative would
need a finite comparison theorem.
\end{firewall}

\subsection{Exact constant-coframe Wilson spectrum}
\label{subsec:v2-flat-spectrum}

Freeze a constant coframe, set the spin/internal links to the identity, and write
\begin{equation}
 x_\mu:=ap_\mu,
 \qquad
 s_\mu(x):=\sin x_\mu,
 \qquad
 c_\mu(x):=1-\cos x_\mu,
 \qquad
 W(x):=\sum_\mu c_\mu(x),
 \qquad
 B(x):=W(x)-1.
 \label{eq:v2-trig-defs}
\end{equation}
The exact dimensionless Wilson symbol is
\begin{equation}
 k_e(x)=i\sum_\mu\Gamma_\mu s_\mu(x)+B(x)I,
 \qquad
 h_e(x)=\gamma_5k_e(x).
 \label{eq:v2-Wilson-symbol}
\end{equation}
Put
\begin{equation}
 q_e(x):=\sum_{\mu,\nu}g_{\mu\nu}s_\mu(x)s_\nu(x),
 \qquad
 R_e(x):=\sqrt{q_e(x)+B(x)^2}.
 \label{eq:v2-R-def}
\end{equation}

\begin{lemma}[Exact trigonometric Wilson identity]
\label{lem:v2-trig-gap}
For the identity coframe,
\begin{equation}
 \boxed{
 \sum_\mu\sin^2x_\mu+B(x)^2
 =1+2\sum_{\mu<\nu}c_\mu(x)c_\nu(x)\ge1.}
 \label{eq:v2-trig-gap-identity}
\end{equation}
Consequently, for every constant coframe satisfying
\eqref{eq:v2-coframe-elliptic},
\begin{equation}
 \boxed{
 \sqrt{\min\{m_g,1\}}
 \le R_e(x)
 \le \sqrt{4M_g+49}}
 \qquad (x\in[-\pi,\pi]^4).
 \label{eq:v2-Wilson-gap-exact}
\end{equation}
\end{lemma}
\begin{proof}
Since $\sin^2x_\mu=2c_\mu-c_\mu^2$,
\[
 \sum_\mu\sin^2x_\mu+(W-1)^2
 =2W-\sum_\mu c_\mu^2+W^2-2W+1
 =1+2\sum_{\mu<\nu}c_\mu c_\nu.
\]
This proves the first identity.  Ellipticity gives
$q_e\ge m_g\sum_\mu s_\mu^2$, hence
$q_e+B^2\ge\min\{m_g,1\}(\sum s_\mu^2+B^2)$.
For the upper bound use $\sum s_\mu^2\le4$ and
$-1\le B\le7$.
\end{proof}

\begin{theorem}[Exact Wilson gap and physical overlap branch]
\label{thm:v2-exact-overlap-branch}
For the constant-coframe reference branch,
\begin{equation}
 h_e(x)^2=R_e(x)^2I,
 \qquad
 \varepsilon_e(x)=\sgn h_e(x)={h_e(x)\over R_e(x)},
 \label{eq:v2-exact-sign}
\end{equation}
and the overlap symbol is
\begin{equation}
 \boxed{
 d_{{\rm ov},a}(p;e)
 ={1\over a}\left[
 I+{i\Gamma\!\cdot s(x)+B(x)I\over R_e(x)}
 \right].}
 \label{eq:v2-exact-overlap-symbol}
\end{equation}
It has exactly one zero in the Brillouin zone, at $x=0$.  In particular no
Wilson corner doubler survives the overlap map.

Moreover
\begin{equation}
 \boxed{
 d_{{\rm ov},a}(p;e)^*d_{{\rm ov},a}(p;e)
 ={2\over a^2}\left(1+{B(x)\over R_e(x)}\right)I.}
 \label{eq:v2-overlap-singular-exact}
\end{equation}
\end{theorem}
\begin{proof}
The Clifford relation gives
$k_e^*k_e=k_ek_e^*=(q_e+B^2)I$, because the mixed scalar--Clifford terms cancel.
The $\gamma_5$-Hermiticity of $k_e$ then gives $h_e^2=R_e^2I$, and
Lemma~\ref{lem:v2-trig-gap} permits the exact sign formula.  Since
$\gamma_5\varepsilon_e=k_e/R_e$, the overlap formula follows.

A zero of $d_{{\rm ov},a}$ would require an eigenvalue of $k_e/R_e$ to equal
$-1$.  The eigenvalues of $k_e/R_e$ have real part $B/R_e$ and imaginary
magnitude $\sqrt{q_e}/R_e$.  Hence a zero requires $q_e=0$ and $B<0$.
Ellipticity implies $q_e=0$ exactly when every $\sin x_\mu=0$, so $x$ is a
Brillouin corner with $n$ coordinates equal to $\pi$.  There
$B=2n-1$, which is negative only for $n=0$.  Thus only $x=0$ is a zero.
Finally $V:=k_e/R_e$ is unitary and
$V+V^*=2B/R_e$, whence
$(I+V)^*(I+V)=2(1+B/R_e)I$.
\end{proof}

\begin{corollary}[A nonempty cutoff-uniform protected Wilson chart]
\label{cor:v2-protected-Wilson-chart}
Let
\[
 g_0:=\sqrt{\min\{m_g,1\}}.
\]
On the reference constant-coframe branch the V2 kernel satisfies
$|H_{W,a}|\succeq g_0I$.  If a gauge/spin/coframe perturbation obeys
\begin{equation}
 \Norm{H_{W,a}(\zeta)-H_{W,a}(0)}_{2\to2}\le {g_0\over2},
 \label{eq:v2-protected-chart-radius}
\end{equation}
then
\begin{equation}
 |H_{W,a}(\zeta)|\succeq {g_0\over2}I.
 \label{eq:v2-protected-chart-gap}
\end{equation}
The chart may be chosen uniformly in volume because the perturbation is
finite range and its operator norm is controlled by the corresponding local
row/column stocks.  Thus the protected Wilson chart required by the overlap
functional calculus is nonempty for the explicit V2 completion.
\end{corollary}
\begin{proof}
The reference gap is Lemma~\ref{lem:v2-trig-gap}.  The singular-value
perturbation inequality gives
$s_{\min}(H(\zeta))\ge s_{\min}(H(0))-\|H(\zeta)-H(0)\|$, proving
\eqref{eq:v2-protected-chart-gap}.  For a fixed-range kernel the
$2\to2$ norm is bounded by the geometric mean of its row and column Schur
bounds, which are local and volume independent.  Smooth dependence of the
one-edge transports and coframe coefficients then supplies a nonempty common
neighbourhood satisfying \eqref{eq:v2-protected-chart-radius}.
\end{proof}

\begin{firewall}[The physical zero is not the hard-shell pivot]
The exact overlap branch has its intended physical zero at $p=0$.  This does
not obstruct E1, whose inverse is always compressed to a hard annulus bounded
away from zero momentum.  Whenever a nowhere-zero \emph{full} determinant-line
section is needed, V2 retains the positive infrared/reference prescription
allowed by the programme (for example a fixed infrared spectral reference,
a compatible boundary-condition gap, or a protected Yukawa/Majorana floor).
No removal of that infrared reference is claimed here.
\end{firewall}

\subsection{Physical annulus and full Brillouin zone}

\begin{theorem}[Physical-branch expansion with all momentum derivatives]
\label{thm:v2-physical-expansion}
Fix a compact ellipticity set \eqref{eq:v2-coframe-elliptic} and an integer
$m\ge0$.  There exists $0<\delta<1$ and constants $C_\alpha,c_{\rm ph}>0$,
uniform on that coframe set, such that for $|ap|\le\delta$,
\begin{align}
 \Norm{\partial_p^\alpha
 (d_{{\rm ov},a}(p;e)-i\Gamma(e)\!\cdot p)}
 &\le C_\alpha a|p|^{2-|\alpha|},
 \qquad |\alpha|\le m,
 \label{eq:v2-overlap-expansion}\\
 \Norm{\partial_p^\alpha
 (\widehat P_{-,a}(p;e)-P_-)}
 &\le C_\alpha a^{|\alpha|}
 (a|p|)^{\max\{1-|\alpha|,0\}},
 \label{eq:v2-projector-expansion}\\
 s_{\min}(d_{{\rm ov},a}(p;e))
 &\ge c_{\rm ph}|p|.
 \label{eq:v2-physical-pivot}
\end{align}
Here $P_-=(I-\gamma_5)/2$ and
$\widehat P_{-,a}=(I+\varepsilon_e)/2$.  The power on the right side of
\eqref{eq:v2-overlap-expansion} is interpreted in the annular sense when
$|\alpha|>2$; because $a|p|\le\delta<1$, it is then a weaker bound than the
uniform analytic derivative bound and is therefore valid.
\end{theorem}
\begin{proof}
Define the dimensionless analytic matrix
\[
 F_e(x):=I+{i\Gamma\!\cdot s(x)+B(x)I\over R_e(x)}.
\]
The exact gap \eqref{eq:v2-Wilson-gap-exact} makes $F_e$ smooth, with all
fixed derivatives uniformly bounded on a common neighbourhood of $x=0$ and on
the compact coframe set.  Direct evaluation gives
$F_e(0)=0$ and
$D F_e(0)[x]=i\Gamma\!\cdot x$.  Hence
$G_e(x):=F_e(x)-i\Gamma\!\cdot x$ vanishes to order two.  Taylor's theorem
therefore gives
\[
 \Norm{\partial_x^\alpha G_e(x)}
 \le C_\alpha |x|^{\max\{2-|\alpha|,0\}}.
\]
For $|\alpha|>2$, the desired factor $|x|^{2-|\alpha|}$ is at least one on
$|x|\le1$, so the same statement implies the annular form.  Since
$d_{{\rm ov},a}=a^{-1}F_e(ap)$, scaling proves
\eqref{eq:v2-overlap-expansion}.  The projector statement follows identically
from the analytic map
$\widehat P_{-,a}(x)=(I+\gamma_5k_e(x)/R_e(x))/2$, whose value at zero is
$P_-$ and whose first nonzero displacement is linear in $x$.

For the pivot, shrink $\delta$ so that $B(x)<0$.  Then
\[
 1+{B\over R_e}={R_e+B\over R_e}
 ={q_e\over R_e(R_e-B)}.
\]
On the compact small ball the denominator is bounded above, while
$q_e\ge m_g|s(x)|^2$ and
$|s(x)|\ge c_\delta|x|$.  Equation
\eqref{eq:v2-overlap-singular-exact} therefore gives
$s_{\min}(d_{{\rm ov},a})\ge c_{\rm ph}|x|/a$.
\end{proof}

\begin{theorem}[Full-zone no-doubler pivot]
\label{thm:v2-full-zone-pivot}
For every $0<\delta<\pi$, there exists $c_{\rm UV}(\delta)>0$, uniform on the
compact coframe ellipticity set, such that
\begin{equation}
 \boxed{
 |x|\ge\delta
 \quad\Longrightarrow\quad
 s_{\min}(d_{{\rm ov},a}(x/a;e))
 \ge {c_{\rm UV}(\delta)\over a}.}
 \label{eq:v2-full-zone-pivot}
\end{equation}
Every fixed $x$-derivative of the corresponding inverse is bounded by $C a$ on
that compact full-zone region.  Consequently a smooth full-zone shell multiplier
has a lattice-scale quasilocal inverse kernel with weighted row/column norm
$O(a)$, uniformly in volume.
\end{theorem}
\begin{proof}
By Theorem~\ref{thm:v2-exact-overlap-branch}, the continuous singular-value
function
$\sqrt{2(1+B/R_e)}$ has no zero on the compact set
$\{(e,x): e\text{ in the ellipticity set},\ |x|\ge\delta\}$.
Its minimum is therefore positive, giving \eqref{eq:v2-full-zone-pivot}.
The inverse has the form $aG_e(x)$ with $G_e$ smooth on the same compact set,
so every fixed $x$-derivative is $O(a)$.  Multiplication by a smooth periodic
full-zone cutoff and repeated Fourier-series summation by parts give faster-than-
any-fixed-power decay in lattice distance.  The lattice row and column sums are
therefore $O(a)$ with no volume factor.
\end{proof}

\subsection{The moving chiral frame and the Yukawa block}

The exact overlap relation is particularly useful with the projector convention of
the Grand Tensor shell.

\begin{lemma}[Exact overlap chiral intertwining]
\label{lem:v2-chiral-intertwining}
For
$D_a=d_{{\rm ov},a}$ and
$\widehat P_-=(I+\varepsilon_e)/2$,
\begin{equation}
 \boxed{D_a\widehat P_-=P_+D_a,\qquad P_+=(I+\gamma_5)/2.}
 \label{eq:v2-chiral-intertwining}
\end{equation}
Hence $D_a$ maps $\Ran\widehat P_-$ into $\Ran P_+$ exactly.
\end{lemma}
\begin{proof}
Expand both products using
$D_a=a^{-1}(I+\gamma_5\varepsilon_e)$ and
$\varepsilon_e^2=I$:
\[
 D_a\widehat P_-
 ={1\over2a}(I+\varepsilon_e+\gamma_5+\gamma_5\varepsilon_e)
 =P_+D_a.
\]
\end{proof}

Let $U_e(x)$ be the Kato unitary obtained by transporting $P_-$ to
$\widehat P_-(x)$ along the radial path $t\mapsto tx$, with $U_e(0)=I$.
The V1 local-frame theorem applies; here the exact symbol permits a sharper scaled
form.

\begin{lemma}[Scaled Kato-frame estimate]
\label{lem:v2-Kato-scaled}
After shrinking $\delta$ if needed,
\begin{equation}
 U_e(x)P_-U_e(x)^*=\widehat P_-(x),
 \qquad
 \Norm{U_e(x)-I}\le C|x|,
 \label{eq:v2-Kato-first}
\end{equation}
and every fixed derivative satisfies
\begin{equation}
 \Norm{\partial_x^\alpha(U_e(x)-I)}
 \le C_\alpha |x|^{\max\{1-|\alpha|,0\}}.
 \label{eq:v2-Kato-derivative}
\end{equation}
The constants are uniform on the compact coframe set.
\end{lemma}
\begin{proof}
Along the radial path the Kato equation is
$\dot U=[\dot P,P]U$.  By
\eqref{eq:v2-projector-expansion},
$\dot P(tx)=x\cdot\nabla P(tx)=O(|x|)$ uniformly for $0\le t\le1$.
Gronwall's inequality gives $U-I=O(|x|)$.  Differentiating the finite-dimensional
ODE a fixed number of times gives uniform derivative bounds; the vanishing at zero
for the undifferentiated difference yields the displayed scaled form.
\end{proof}

Let $Y(H)$ denote the fixed-range Yukawa map on the reference Higgs chart, with
\begin{equation}
 M_Y:=\sup\Norm{Y(H)}<\infty
 \label{eq:v2-Yukawa-stock}
\end{equation}
on the finite panel.  In the Kato source frame define
\begin{align}
 \widetilde d_{\chi,a}(p;e,H)
 &:=\bigl(D_a(p;e)+Y(H)\bigr)U_e(ap)P_-:
 \Ran P_-\longrightarrow\Ran P_+,
 \label{eq:v2-fixed-frame-chiral}\\
 d_{\chi,0}(p;e,H)
 &:=\bigl(i\Gamma(e)\!\cdot p+Y(H)\bigr)P_-.
 \label{eq:v2-ordinary-chiral}
\end{align}
The target projection $P_+$ is suppressed because the kinetic term lands there by
Lemma~\ref{lem:v2-chiral-intertwining} and the Yukawa map is defined with that
target.

\begin{theorem}[Native chiral annular symbol theorem on the reference branch]
\label{thm:v2-native-chiral-annular}
Fix $m\ge0$.  There are constants $c_0,c_D,C_\alpha>0$ such that on every hard
annulus $c_-\Lambda_j\le|p|\le c_+\Lambda_j$ satisfying
\begin{equation}
 a\Lambda_j\le c_0,
 \qquad
 \Lambda_j\ge {2M_Y\over c_{\rm ph}c_-},
 \label{eq:v2-hard-scale-conditions}
\end{equation}
one has, for $|\alpha|\le m$,
\begin{align}
 \Norm{\partial_p^\alpha
 (\widetilde d_{\chi,a}-d_{\chi,0})}
 &\le C_\alpha a\Lambda_j^{2-|\alpha|},
 \label{eq:v2-chiral-error}\\
 s_{\min}(\widetilde d_{\chi,a}(p))
 +s_{\min}(d_{\chi,0}(p))
 &\ge 2c_D\Lambda_j.
 \label{eq:v2-chiral-pivot}
\end{align}
Thus the undifferentiated part of Assumption~\ref{ass:v1-native-symbol-row} is a
theorem for the explicit V2 Wilson completion.  On the complementary full-zone
region, the kinetic chiral inverse has $O(a)$ lattice-scale quasilocal norm; adding
the bounded Yukawa block preserves that bound for sufficiently small $a$.
\end{theorem}
\begin{proof}
By Theorem~\ref{thm:v2-physical-expansion} and
Lemma~\ref{lem:v2-Kato-scaled},
\[
 D_aU_eP_--i\Gamma\!\cdot pP_-
 =(D_a-i\Gamma\!\cdot p)U_eP_-
 +i\Gamma\!\cdot p(U_e-I)P_-
 =O(a|p|^2).
\]
The Yukawa-frame correction is
$Y(H)(U_e-I)P_-=O(M_Ya|p|)$, which is
$O(a\Lambda_j^2)$ under \eqref{eq:v2-hard-scale-conditions}.  Differentiating and
using the scaled derivative estimates proves \eqref{eq:v2-chiral-error}.

For the pivot, \eqref{eq:v2-overlap-singular-exact} shows that $D_a^*D_a$ is a
scalar multiple of the identity.  Since $U_e$ is unitary and maps $P_-$ onto
$\widehat P_-$, the kinetic chiral restriction has exactly the same nonzero
singular value as $D_a$.  Hence
\[
 s_{\min}(D_aU_eP_-+Y U_eP_-)
 \ge c_{\rm ph}c_-\Lambda_j-M_Y
 \ge {1\over2}c_{\rm ph}c_-\Lambda_j.
\]
The same elementary singular-value perturbation bound applies to
$i\Gamma\cdot p+Y$, whose kinetic singular value is at least
$\sqrt{m_g}\,c_-\Lambda_j$.  Decreasing $c_D$ if needed gives
\eqref{eq:v2-chiral-pivot}.  The full-zone statement follows from
Theorem~\ref{thm:v2-full-zone-pivot} and the bounded perturbation $Y$.
\end{proof}

\subsection{Finite external source jets}
\label{subsec:v2-source-jets}

We now treat gauge, spin, coframe and Higgs insertions as the fixed perturbative
jets required by the coefficientwise theorem.  Let the external source panel contain
only finitely many Fourier momenta $q$ with
\begin{equation}
 |q|\le\mu,
 \qquad
 \mu\le\kappa\Lambda_j
 \label{eq:v2-soft-source-window}
\end{equation}
for one fixed $0<\kappa<1$ on the hard shell under consideration.

For an internal gauge plus spin connection $\mathcal A_\mu$, use the actual edge
parallel transport
\begin{equation}
 \mathcal U_{x,\mu}(\zeta)
 =\mathcal P\exp\left(
 \int_0^a\mathcal A_\mu(x+t\hat\mu;\zeta)\,dt\right).
 \label{eq:v2-link-parallel-transport}
\end{equation}
Coframe sources enter through the sampled $\Gamma_\mu(e)$ in
\eqref{eq:v2-native-Wilson}; Higgs/Yukawa sources enter the onsite map $Y(H)$.
The primitive Wilson kernel has no antifield dependence in this regulator choice;
antifields enter through the local BV interaction/source vertices and blocking maps.
Likewise determinant-line coordinates act through the Kato/line frame, not by
changing the primitive Wilson stencil.

\begin{lemma}[Parallel-transport source jets]
\label{lem:v2-link-jets}
At the zero-connection reference point, every fixed $r$th source derivative of
\eqref{eq:v2-link-parallel-transport} is a finite sum of ordered simplex integrals
of total size $O(a^r)$.  If the inserted source modes satisfy
\eqref{eq:v2-soft-source-window}, then every fixed derivative with respect to their
momenta gains an additional factor $O(a)$.  In particular the exact edge jet is its
continuum midpoint/local jet plus a remainder smaller by at least one factor
$a(\Lambda_j+\mu)$ after it is inserted into a hard Wilson vertex.
\end{lemma}
\begin{proof}
Differentiate the path-ordered exponential at zero connection.  The $r$th
coefficient is a sum over orderings of integrals on
$0<t_1<\cdots<t_r<a$.  Their total simplex volume is $a^r/r!$.  A plane-wave
source contributes factors $e^{iq\cdot(x+t_k\hat\mu)}$; differentiating in $q$
produces powers of $t_k$, each bounded by $a$.  Taylor expansion of these phase
factors and of the hard exponentials $e^{\pm iap_\mu}$ then gives one further
factor $a(|p|+\sum|q_k|)$ beyond the leading local jet.
\end{proof}

For a source multi-index $I$, let $v_{I,a}(p;\mathbf q)$ denote the exact momentum
vertex obtained by differentiating the fixed-frame native chiral operator at the
reference point, including the Kato-frame derivative when the source space moves.
Let $v_{I,0}(p;\mathbf q)$ be the corresponding continuum covariant
Dirac--Yukawa Taylor jet.  Its hard-momentum degree is denoted
$\nu_I\in\{0,1\}$: coframe-type kinetic jets may have degree one, whereas
connection/Yukawa and the resulting local contact jets have degree zero.  If the
ordinary vertex vanishes, $\nu_I=0$ and $v_{I,0}=0$ by convention.

\begin{theorem}[All fixed native source vertices match their ordinary jets]
\label{thm:v2-source-vertex-transfer}
Fix source order $r$, momentum derivative order $m$, and the finite panel
\eqref{eq:v2-soft-source-window}.  For every $0<|I|\le r$ there are constants
$C_{I,\alpha,\beta}$, independent of volume, shell number and cutoff, such that
on $a\Lambda_j\le c_0$,
\begin{equation}
 \boxed{
 \Norm{
 \partial_p^\alpha\partial_{\mathbf q}^\beta
 (v_{I,a}-v_{I,0})}
 \le
 C_{I,\alpha,\beta}
 a\Lambda_j^{\nu_I+1-|\alpha|-|\beta|}}
 \label{eq:v2-source-vertex-error}
\end{equation}
for all fixed $|\alpha|+|\beta|\le m$, after division by the declared source
amplitudes/stocks.  The native and ordinary vertices themselves obey the same
bound without the prefactor $a\Lambda_j$ relative to their engineering scale.

The first nonzero ordinary jets are the expected ones:
\begin{enumerate}[label=\textup{(J\arabic*)},leftmargin=3em]
\item the zero-source kinetic jet is $i\Gamma^\mu p_\mu$;
\item a connection insertion is the local Clifford contraction
$\Gamma^\mu\mathcal A_\mu$ in the chosen anti-/Hermitian convention;
\item a coframe insertion differentiates $i\Gamma^\mu(e)p_\mu$;
\item a Higgs/Yukawa insertion is the corresponding derivative of $Y(H)$;
\item primitive antifield derivatives of $H_W$ vanish for
\eqref{eq:v2-native-Wilson}; their BV vertices remain separate local source
insertions;
\item determinant-line derivatives are derivatives of the Kato/line frame and are
controlled by the same analytic source calculus, not by an invented Wilson
coupling.
\end{enumerate}
Thus the differentiated source part of Assumption~\ref{ass:v1-native-symbol-row}
is proved at every fixed source order for the V2 regulator completion.
\end{theorem}
\begin{proof}
At the reference point the Wilson symbol has the uniform all-zone gap of
Lemma~\ref{lem:v2-trig-gap}.  Write the sign as the difference of its positive and
negative Riesz projectors.  A fixed source derivative of either projector is a finite
sum of contour-resolvent words of the form
\begin{equation}
 {1\over2\pi i}\oint_\Gamma
 R_{p+Q_k}(z)V_{B_k,a}\,
 R_{p+Q_{k-1}}(z)\cdots
 V_{B_1,a}R_p(z)\,dz,
 \label{eq:v2-resolvent-source-word}
\end{equation}
where $R_p(z)=(z-h_e(ap))^{-1}$, the $B_i$ form an ordered partition of $I$,
and $Q_i$ are partial sums of external source momenta.  The contour may be chosen
uniformly because the negative and positive Wilson spectra remain separated by the
fixed gap.  Every resolvent and every fixed momentum derivative is therefore a
uniform analytic function of the dimensionless variables
$ap,a q_1,\ldots,a q_r$ on the physical polydisk.

The primitive source vertices $V_{B,a}$ are obtained by differentiating the
one-edge formula \eqref{eq:v2-native-Wilson}.  Lemma~\ref{lem:v2-link-jets}
gives the connection scaling and its local leading jet.  Coframe derivatives act on
$\Gamma_\mu$ and are always accompanied, in the kinetic part, by the centred hard
shift, whose physical branch is $iap_\mu+O((ap)^2)$.  Higgs/Yukawa derivatives are
onsingular onsite jets.  Hence every word \eqref{eq:v2-resolvent-source-word} is a
smooth finite matrix function of the dimensionless hard and soft momenta.

At the origin of these variables its first nonzero Taylor polynomial is obtained by
replacing the centred edge difference by the ordinary covariant first derivative,
the edge parallel transport by its local connection jet, and the onsite Yukawa map
by itself.  Direct differentiation of \eqref{eq:v2-native-Wilson} at the physical
branch gives the six listed ordinary jets.  The Kato frame is an analytic functional
of the same projector on the protected local chart, so its derivatives satisfy the
same scaled Taylor calculus.  Subtracting the first nonzero ordinary Taylor jet
therefore leaves one extra factor
$a(|p|+\sum|q_i|)$.  Under
\eqref{eq:v2-soft-source-window} this is $O(a\Lambda_j)$.
Rescaling the remaining hard degree $\nu_I$ and differentiating in $p$ or $q$
proves \eqref{eq:v2-source-vertex-error}.  Fixed order leaves only finitely many
resolvent words, so no source-radius or volume factor enters.
\end{proof}

\begin{remark}[Why the coefficientwise route is stronger than a crude background norm]
A finite coframe displacement changes the Wilson operator by $O(1)$ in its global
dimensionless operator norm and would therefore produce an unhelpful $O(a^{-1})$
bound if one first perturbed the full overlap operator in operator norm.  The source
jet in Theorem~\ref{thm:v2-source-vertex-transfer} instead carries the hard
centred-difference factor $ap$, and hence has the correct $O(\Lambda_j)$ stress
scaling.  This is exactly the quantifier gain supplied by the coefficientwise
programme.
\end{remark}

\subsection{Full-zone source derivatives}

\begin{proposition}[Uniform ultraviolet-cap source calculus]
\label{prop:v2-full-zone-source}
Fix $\delta>0$ and a finite source order.  On the full-zone region
$|ap|\ge\delta$, every fixed reference inverse/source word generated by
Theorem~\ref{thm:v2-reference-pivot} has a smooth periodic dimensionless symbol.
After the physical factor of $a$ from each reference fermion inverse is restored,
its cutoff-profiled Fourier kernel obeys a volume-uniform lattice-scale weighted
row/column bound.  A derivative with respect to an external source momentum gains
one factor $a$, so the complementary source remainder on a panel $|q|\le\mu$ is
$O(a\mu)$.
\end{proposition}
\begin{proof}
Theorem~\ref{thm:v2-full-zone-pivot} gives an inverse $aG(x)$ with $G$ smooth on
the compact high-zone region.  Fixed source derivatives are finite resolvent/inverse
words with smooth periodic dimensionless factors by
Theorem~\ref{thm:v2-source-vertex-transfer}; hence every fixed $x$-derivative is
bounded.  Fourier-series summation by parts gives uniform rapid decay in lattice
units.  External source momentum enters only through the dimensionless combination
$aq$, so the mean-value theorem gives the additional factor $a\mu$.
\end{proof}

\subsection{The missing line summability follows from the same shell calculus}
\label{subsec:v2-line-summability}

V1 proved the exact Bell-polynomial composition theorem for determinant/Pfaffian
line sections under the sole remaining analytic input
$\rho_j=O(\mu/\Lambda_j)$.  We now prove that input for the coefficientwise V2
branch.

Let $g_{j,I}(\mathbf q)$ be any fixed connected logarithmic determinant-line
coefficient at shell $j$, after the exact source differentiation has been performed.
By Lemma~\ref{lem:v1-logdet-words} it is a finite sum of cyclic hard inverse--vertex
words.  Let $\mathsf T_j$ denote the common local Taylor/EFT assignment already
used by the bosonic shell; in particular it contains at least the zero-external-
momentum local jet of every superficially local line coefficient.

\begin{theorem}[One-excess-scale chiral line remainder]
\label{thm:v2-line-excess}
Fix the source order, the finite external panel $|q|\le\mu$, and the local/EFT weight
kept in the two-loop packet.  For every physical hard annulus and every connected
logarithmic determinant-line coefficient,
\begin{equation}
 \boxed{
 \Norm{\partial_I(g_j-\mathsf T_jg_j)}_{\rm pan}
 \le C_I{\mu\over\Lambda_j}.}
 \label{eq:v2-line-excess}
\end{equation}
The same estimate holds for the optional locally oriented Pfaffian logarithm.
For the ultraviolet-cap shells, the right side is read as $C_Ia\mu$, which is the
same estimate because $\Lambda_j\asymp a^{-1}$ there.

All mixed source contacts produced by differentiation of the Schur complement or
of the shell product are retained before $\mathsf T_j$ is applied.
\end{theorem}
\begin{proof}
On a physical annulus rescale the hard loop momentum by $p=\Lambda_jk$.  The
reference inverse bounds of Theorem~\ref{thm:v2-native-chiral-annular}, the source
vertex bounds of Theorem~\ref{thm:v2-source-vertex-transfer}, and the rooted trace
estimate of Lemma~\ref{lem:v1-logdet-words} show that, after division by the declared
engineering/local source stock, every fixed cyclic integrand is a $C^1$ function of
$k$ and of the dimensionless external variables $q_i/\Lambda_j$, with a bound
independent of $j$, $a$ and volume.  The local assignment $\mathsf T_j$ removes at
least the value at zero external momentum.  Therefore the multivariate mean-value
theorem gives
\[
 \Norm{F_j(\mathbf q)-F_j(0)}
 \le C_I\sum_i{|q_i|\over\Lambda_j}
 \le C_I'{\mu\over\Lambda_j}.
\]
If power counting requires a higher local Taylor polynomial, the remainder only
improves.  Integration of the normalized hard variable and the rooted trace
contraction preserve the same factor.  The ultraviolet-cap version is
Proposition~\ref{prop:v2-full-zone-source}.  The Pfaffian logarithm differs by the
fixed factor $1/2$ on the chosen oriented component.

Crucially, differentiation is performed before the Taylor assignment.  Hence every
mixed hard/soft Schur contact of Proposition~\ref{prop:v1-Schur-contacts} and every
mixed-shell contact of Theorem~\ref{thm:v1-line-product} remains present in the
coefficient being subtracted; no contact is discarded by the estimate.
\end{proof}

\begin{corollary}[Differentiated line products are summable]
\label{cor:v2-line-product-closed}
For geometric shell scales $\Lambda_j=b^j\lambda$, the hypothesis
\eqref{eq:v1-line-log-summable} of Theorem~\ref{thm:v1-line-product} holds with
\begin{equation}
 \rho_j=C{\mu\over\Lambda_j}.
 \label{eq:v2-rho-line}
\end{equation}
Hence every fixed differentiated determinant/Pfaffian line product is uniformly
bounded through arbitrarily many ultraviolet shell eliminations on the fixed source
panel, with all mixed-shell source contacts retained.
\end{corollary}
\begin{proof}
Theorem~\ref{thm:v2-line-excess} gives the shell estimate and the geometric series
$\sum_j\mu/\Lambda_j$ converges.  Apply
Theorem~\ref{thm:v1-line-product} and Proposition~\ref{prop:v1-Pfaffian}.
\end{proof}

\subsection{Coefficientwise E1 closure}

\begin{theorem}[E1 closes for the explicit V2 Wilson completion]
\label{thm:v2-E1-closure}
Fix:
\begin{enumerate}[label=\textup{(F\arabic*)},leftmargin=3em]
\item the explicit Wilson regulator \eqref{eq:v2-native-Wilson};
\item a compact uniformly elliptic reference coframe chart;
\item the protected local determinant/Pfaffian line component inherited from V1;
\item a finite loop/source packet and finite external source momentum window;
\item a positive IR/reference scale, with hard shells taken above the bounded
Yukawa/reference mass scale.
\end{enumerate}
Then E1 is proved at every fixed source/derivative order in the coefficientwise
sense required by the programme:
\begin{enumerate}[label=\textup{(E1.\arabic*)},leftmargin=3.5em]
\item the native Wilson kernel is finite range, Hermitian and exactly gauge/spin
covariant, with all fixed local source jets;
\item its reference overlap branch has an exact all-zone Wilson gap, one and only one
physical zero, and no doubler zero;
\item on $a\Lambda_j\ll1$ the fixed-frame native chiral symbol differs from the
ordinary covariant Weyl/Yukawa symbol by $O(a\Lambda_j^2)$ with every required
fixed momentum/source derivative, and has an $O(\Lambda_j)$ reference hard pivot;
\item on the complementary ultraviolet-cap region the reference inverse and all
fixed source words have $O(a)$ lattice-scale quasilocal bounds;
\item the V1 overlap functional-calculus, Kato-frame, inverse-word, Schur-contact and
rooted graph theorems therefore apply coefficientwise without a
cutoff-independent finite-amplitude source ball;
\item the locally subtracted determinant/Pfaffian logarithmic line remainder is
$O(\mu/\Lambda_j)$ and its differentiated products are summable through arbitrary
finite shell composition.
\end{enumerate}
No analytic finite-coupling RG trajectory, global determinant-line holonomy,
zero-crossing theorem, or nonperturbative chiral measure is asserted.
\end{theorem}
\begin{proof}
Items (E1.1)--(E1.2) are Proposition~\ref{prop:v2-Wilson-structure} and
Theorem~\ref{thm:v2-exact-overlap-branch}.  Item (E1.3) is
Theorems~\ref{thm:v2-native-chiral-annular} and
\ref{thm:v2-source-vertex-transfer}.  Item (E1.4) is
Theorem~\ref{thm:v2-full-zone-pivot} and
Proposition~\ref{prop:v2-full-zone-source}.  The reference-pivot replacement of the
stronger finite-amplitude inversion hypothesis is
Theorem~\ref{thm:v2-reference-pivot}; the remaining quasilocal, Schur and rooted
conclusions are exactly the V1 theorems fed by these now-proved native symbol rows.
Item (E1.6) is Theorem~\ref{thm:v2-line-excess} and
Corollary~\ref{cor:v2-line-product-closed}.
\end{proof}

\begin{firewall}[Scope of the E1 closure]
Theorem~\ref{thm:v2-E1-closure} is an existence/closure theorem for the explicit
Wilson completion fixed in \eqref{eq:v2-native-Wilson}.  If the final Grand Tensor
manuscript chooses a different Wilson mass, hopping stencil, or graph regulator, one
must either adopt \eqref{eq:v2-native-Wilson} as the native definition or prove a
finite regulator-comparison theorem.  The coefficientwise conclusion itself does
not license silently changing regulators afterward.
\end{firewall}

\section{V2 handoff to E2: the native chiral BV breaking}
\label{sec:v2-E2-handoff}

With E1 closed on the explicit regulator branch, continuing to refine the hard
fermion symbol would now be downstream polishing.  The next load-bearing issue is
cohomological but must be tied to the \emph{same native shell}: the local chiral
breaking generated by the regulated determinant/Pfaffian line must be decomposed
into the standard anomaly class and a local BV-exact regulator term, with all
source/antifield descendants carried by one common finite action.

Let $S_{0,j}$ denote the common classical BV action at shell $j$ and let
$\mathfrak s_j=(S_{0,j},\,\cdot\,)$ be its linearized BV differential on the finite
local bank.  Let $\mathcal B_{j}^{(\ell,r)}$ denote the ghost-number-one local
breaking at loop order $\ell$ and source order $r$ obtained from the coefficient of
\begin{equation}
 {1\over2}(S,S)-\hbar\Delta S
 \label{eq:v2-master-breaking}
\end{equation}
after the complete native hard shell, including the chiral line/reference term, has
been assembled and its quasilocal remainder has been separated by the E1 Taylor
assignment.

\begin{theorem}[Exact next load-bearing theorem for V3]
\label{thm:v2-next-E2}
At the fixed two-loop/EFT/source order, prove for the \emph{native V2 shell} that
there is one local decomposition
\begin{equation}
 \boxed{
 \mathcal B_j^{(\ell,r)}
 =\sum_A c_{A,j}^{(\ell,r)}\,\mathcal A_A^{\rm std}
 +\mathfrak s_j\Xi_j^{(\ell,r)},}
 \label{eq:v2-E2-target-decomposition}
\end{equation}
where:
\begin{enumerate}[label=\textup{(B\arabic*)},leftmargin=3em]
\item $\{\mathcal A_A^{\rm std}\}$ is the finite ordinary four-dimensional local
internal-gauge anomaly basis in the already fixed hypercharge convention;
\item the coefficients $c_{A,j}^{(\ell,r)}$ are computed from the complete native
chiral shell and shown to equal the inherited Standard-Model representation traces,
so they vanish after the full three-generation packet is assembled;
\item $\Xi_j^{(\ell,r)}$ is one local restoring primitive in the same finite
EFT/source bank, not a graph-dependent family of counterterms;
\item differentiating \eqref{eq:v2-E2-target-decomposition} generates all gauge,
diffeomorphism, local-Lorentz, Higgs/Yukawa, ghost, antifield, stress/contact and
line descendants required by the common two-loop source packet;
\item the resulting restored shell enlarges the inherited V80 bosonic self-map to
one common two-loop Einstein--Standard Model action/source/line update.
\end{enumerate}
The proof must use the E1 native comparison to identify the cohomology coefficient;
anomaly-trace cancellation alone is not a regulator-matching theorem.
\end{theorem}

The correct V3 attack is therefore not another anomaly-trace calculation.  Those
traces are inherited.  It is the finite local comparison map from the native
regulated BV breaking to the ordinary anomaly representative, followed by a common
local primitive for the exact remainder.

\section{V2 status ledger}
\label{sec:v2-status}

\subsection*{Proved in V2}
\begin{enumerate}[leftmargin=2.2em]
\item At fixed source order, a coefficientwise fermionic shell requires only the
reference hard inverse and finitely many source jets.  A cutoff-uniform
finite-amplitude source ball is a stronger analytic statement and is not required
for the programme's coefficientwise endpoint.
\item The previously abstract Wilson slot was filled by the explicit covariant
nearest-neighbour kernel \eqref{eq:v2-native-Wilson}.  It is finite range,
Hermitian and exactly gauge/spin covariant, with fixed-range source derivatives.
\item On every constant uniformly elliptic coframe, the Wilson symbol obeys the exact
all-zone gap identity \eqref{eq:v2-Wilson-gap-exact}.  The associated overlap
operator has exactly one physical zero and no doubler zeros.
\item The physical overlap branch satisfies
$d_{{\rm ov},a}=i\Gamma\cdot p+O(a|p|^2)$ with every fixed momentum derivative,
and has an $O(|p|)$ physical singular-value floor.  The complementary full
Brillouin zone has an $O(a^{-1})$ pivot and an $O(a)$ inverse kernel.
\item In the Kato chiral frame, the overlap--Yukawa block satisfies the V1 annular
finite-to-ordinary symbol row and hard reference pivot on sufficiently hard shells.
\item Every fixed gauge, spin, coframe and Higgs/Yukawa source jet of the native
overlap operator matches the corresponding ordinary covariant Dirac/Yukawa jet with
one extra factor $a\Lambda_j$ after normalized source scaling.  Primitive antifield
and line sources are correctly kept outside the Wilson stencil and enter through
the local BV/Kato source calculus.
\item The locally subtracted logarithmic determinant/Pfaffian line remainder is
$O(\mu/\Lambda_j)$, including the ultraviolet-cap interpretation
$O(a\mu)$.  Therefore the V1 differentiated line-product theorem is summable with
all mixed-shell contacts retained.
\item E1 is closed at every fixed source/derivative order for the explicit V2 Wilson
completion, at the coefficientwise scope of the programme.
\end{enumerate}

\subsection*{Inherited}
\begin{enumerate}[leftmargin=2.2em]
\item The V80 bosonic Einstein--gauge--Higgs two-loop self-map and its common
local/source normalization.
\item The three-generation Standard-Model chiral representation, exact cancellation
of the five ordinary local gauge-anomaly coefficients in the declared convention,
and even $SU(2)$ doublet parity.
\item Exact determinant/Pfaffian Schur factorization, finite regrouping cocycles,
and differentiated line-source identities.
\item The V1 weighted functional calculus, local Kato line frame, Schur-contact
calculus, rooted hard-graph bounds, and Bell-polynomial line-product identity.
\end{enumerate}

\subsection*{Conditional}
\begin{enumerate}[leftmargin=2.2em]
\item The E1 closure is native to the explicit Wilson completion
\eqref{eq:v2-native-Wilson}.  Identification with any other Wilson kernel intended
by a separate manuscript requires an explicit finite regulator-comparison theorem.
\item The construction remains on the protected perturbative determinant/Pfaffian
component and retains the positive IR/reference scale declared by the programme.
No statement is made across zero modes or global line-holonomy transitions.
\item Constants may depend on the fixed perturbative, EFT, source and derivative
orders and on the compact reference coframe/source panel.  No uniform
large-loop-order estimate is claimed.
\end{enumerate}

\subsection*{Open}
\begin{enumerate}[leftmargin=2.2em]
\item Prove the native two-loop chiral BV-breaking decomposition
\eqref{eq:v2-E2-target-decomposition} using the E1 finite-to-ordinary comparison,
not merely the already known anomaly traces.
\item Show that the nontrivial anomaly coefficients of that native breaking equal the
inherited Standard-Model trace vector and hence vanish only after the complete
chiral representation is assembled.
\item Construct one common local restoring primitive and all of its source,
antifield, stress/contact and determinant/Pfaffian-line descendants.
\item Combine that restored chiral row with the inherited V80 bosonic map to prove
one common two-loop Einstein--Standard Model action/source/line self-map.
\end{enumerate}

\subsection*{Exact next load-bearing theorem}
\begin{center}
\fbox{\parbox{0.92\textwidth}{
\textbf{Native chiral BV-breaking comparison and restoration.}
For the explicit E1-closed Wilson/overlap regulator, compute the complete local
ghost-number-one breaking of the two-loop hard shell in the common BV/source
scheme, prove that its non-exact projection is exactly the ordinary Standard-Model
anomaly trace vector, use the inherited trace cancellations to remove that
projection, and construct one local BV primitive with all source/antifield/line
descendants.  The output must be one common chiral extension of the V80 bosonic
two-loop shell self-map, not separate graphwise counterterms.
}}
\end{center}

\section*{Append-only continuation marker: V2}
\addcontentsline{toc}{section}{Append-only continuation marker: V2}
\textbf{End of V2.}  The complete V1 body above is retained.  Future versions must
append new work after this marker without replacing the V1--V2 mathematical body.
The next file is to be named
\texttt{Renewal\_Perturbative\_ESM\_Continuum\_Research\_Notes\_V3.tex}.



\clearpage
\part*{V3: Native anomaly-class cancellation and the common chiral two-loop restorer}
\addcontentsline{toc}{part}{V3: Native anomaly-class cancellation and the common chiral two-loop restorer}

\section{V3 upstream correction: the load-bearing target is zero native anomaly projection}
\label{sec:v3-upstream}

V2 ended with a deliberately strong version of E2: compare the complete native
chiral breaking with a conventionally normalized ordinary anomaly representative,
identify every coefficient with the corresponding Standard-Model representation
trace, and only then construct a restoring primitive.  That route is legitimate,
but for the actual coefficientwise E2 theorem it contains one unnecessary
normalization problem.

What the common local BV range theorem consumes is not a separate numerical
normalization of every anomalous species.  It consumes the statement
\begin{equation}
 \boxed{\mathsf A_{\rm std}\,\mathcal B^{(\ell,r)}_j=0}
 \label{eq:v3-load-bearing-zero}
\end{equation}
for the \emph{complete native Standard-Model packet}, where
$\mathsf A_{\rm std}$ is the ordinary finite anomaly-quotient projection already
constructed in V80.  Once \eqref{eq:v3-load-bearing-zero} is known, the inherited
matter-inclusive local cohomology theorem gives one local BV primitive.  The
coefficientwise continuum target does not require a separate theorem that the
uncancelled one-species anomaly has a preassigned continuum normalization.

There is a native way to prove \eqref{eq:v3-load-bearing-zero}.  The supplied
Grand Tensor field-shell construction already gives, on the protected nonzero
component, a coordinate-free chiral determinant section for the complete
anomaly-free generation and proves that this section descends under the finite
gauge group.  V1--V2 add precisely the missing analytic information: a legitimate
local line frame, fixed-order source differentiability, hard-shell quasilocality,
and the exact determinant/Pfaffian shell cocycle.  We therefore first build a
\emph{gauge-invariant local line frame for the complete packet itself}.  This
kills the nontrivial native internal-gauge line cocycle before any continuum
normalization is invoked.

\begin{firewall}[What V3 does and does not replace]
The argument below does not claim that an anomalous single Weyl representation
has zero anomaly, and it does not identify the universal numerical coefficient
of such an anomalous representation.  It proves the stronger statement needed
for E2: after the actual Standard-Model representation is assembled, the
complete native finite shell has zero projection onto the standard local anomaly
quotient.  The inherited trace cancellations and determinant-line descent are
used as inputs; they are not reproved.  No Adler--Bardeen theorem is imported.
\end{firewall}

\section{The source-extended finite ESM BV complex}
\label{sec:v3-complex}

Fix once and for all the two-loop perturbative order, local EFT weight six,
source order $r$, field/source arity $N$, the finite external panel, and the
explicit V2 Wilson/overlap branch.  Let $\Phi$ denote the complete retained
field list: coframe/metric and spin connection variables, the three internal
gauge fields, Higgs field, chiral quark and lepton multiplets, ghosts,
nonminimal variables, and all admitted composite/source writers.  Let $\Phi^*$
denote the exact antifield duals used by the common BV splitting.

Let $\mathfrak J_r$ be the finite truncated source-jet algebra generated by all
source coordinates retained through total source order $r$.  Source variables
which are members of BRST doublets or antifield/source multiplets carry their
actual differential; spectator physical sources carry zero differential.  Let
\begin{equation}
 \widehat{\mathfrak L}^{\rm ESM}_{N,6,r}
 :=\mathfrak L^{\rm ESM}_{N,6}\otimes \mathfrak J_r
 \label{eq:v3-source-extended-bank}
\end{equation}
be the complete finite local bank.  The notation means the quotient by the
chosen integration-by-parts, algebraic-symmetry, equation-of-motion and other
redundant directions, exactly as in V80; it does not mean that only a displayed
list of Warsaw-type representatives is retained.

Let $\mathsf P^{\rm ESM}_{N,6,r}$ denote the common local Taylor/EFT/source
projection.  It is applied only after every contact generated by source
differentiation has been assembled.  Let
\begin{equation}
 d_j:=\mathsf P^{\rm ESM}_{N,6,r}
       (S_{0,j},\,\cdot\,)
       \mathsf P^{\rm ESM}_{N,6,r}
 \label{eq:v3-local-differential}
\end{equation}
be the projected classical BV differential at shell $j$.

\begin{theorem}[Finite source-extended ESM chain complex]
\label{thm:v3-chain-complex}
On the protected perturbative chart and at the fixed $(N,6,r)$,
\begin{equation}
 \boxed{d_j^2=0,\qquad
 \mathsf P^{\rm ESM}_{N,6,r}d_j
 =d_j\mathsf P^{\rm ESM}_{N,6,r}.}
 \label{eq:v3-chain-identities}
\end{equation}
The nonminimal pairs remain contractible.  Adding the chiral fermion and
Yukawa rows does not enlarge the nontrivial ghost-number-one quotient beyond
the ordinary standard gauge-anomaly space already isolated in V80.
\end{theorem}

\begin{proof}
Before projection the complete classical ESM BV action satisfies the classical
master equation on the declared relative/protected domain.  Hence its
Hamiltonian vector field squares to zero.  The chiral fermion transformations
are linear in the corresponding gauge generators and the Yukawa terms are
gauge covariant; therefore the classical differential preserves canonical
weight, ghost number, representation type and the finite source order.  The
same paired-jet argument used in the V80 Einstein--gauge--Higgs bank then shows
that the complete Taylor/EFT/source projector commutes with the differential
provided every descendant at the retained weight and source order is kept.
This gives \eqref{eq:v3-chain-identities}.

The second statement is the matter-inclusive content of the cohomological
reduction already used in V80.  Hypercharge is not a free abelian generator on
the charged-Higgs branch, so the exceptional free-abelian antifield families
are absent.  In four dimensions the remaining nontrivial ghost-one quotient is
the ordinary gauge-anomaly quotient.  Adding further gauge-covariant matter
variables and their antifields enlarges the contractible and invariant parts of
the finite bank but does not create a new four-dimensional ghost-one anomaly
family outside that inherited classification.  The nonminimal doublets are
contractible by the same explicit homotopy as in V80.
\end{proof}

Denote the resulting anomaly quotient by
\begin{equation}
 \mathfrak A_{\rm std}
 =\Span\{\mathcal A_{333},\mathcal A_{Y33},\mathcal A_{Y22},
              \mathcal A_{YYY},\mathcal A_{YRR}\},
 \label{eq:v3-anomaly-space}
\end{equation}
where the labels refer to the consistent descendants of the five invariant
polynomials already fixed in V80.  Let
$\mathsf A_{\rm std}:\ker(d_j)|_{\operatorname{gh}=1}\to
\mathfrak A_{\rm std}$ be the inherited coefficient projection.

\section{Quantum-master consistency is a native finite identity}
\label{sec:v3-consistency}

The next step is upstream of anomaly classification.  A local anomaly
projection is meaningful only after the complete shell breaking has been shown
to be a cocycle.

Let $\Delta_j$ be the finite BV Laplacian in the same field--antifield frame
used by the shell integration, including the declared measure/reference
normalization exactly once.  For an even BV functional $S$ define
\begin{equation}
 \mathcal M_j(S)
 :=\frac12(S,S)-\hbar\Delta_jS.
 \label{eq:v3-master-functional}
\end{equation}
All reference terms which depend on retained sources are part of the same
finite functional and are not discarded as ``field independent.''

\begin{lemma}[Finite quantum consistency identity]
\label{lem:v3-quantum-consistency}
Assume $\Delta_j^2=0$ and the standard BV second-order identity for the
antibracket.  Then, identically at finite cutoff,
\begin{equation}
 \boxed{(S,\mathcal M_j(S))-\hbar\Delta_j\mathcal M_j(S)=0.}
 \label{eq:v3-quantum-WZ}
\end{equation}
\end{lemma}

\begin{proof}
The graded Jacobi identity gives
$(S,\frac12(S,S))=0$.  Since $S$ is even, the BV second-order identity gives
\[
 \Delta_j\frac12(S,S)=-(S,\Delta_jS).
\]
Therefore the contribution of $(S,-\hbar\Delta_jS)$ cancels
$-\hbar\Delta_j\frac12(S,S)$, while
$\Delta_j^2S=0$.  No limit and no formal change of regulator is used.
\end{proof}

Write the recursively normalized action through two loops as
\begin{equation}
 S_j^{[\le2]}=S_{0,j}+\hbar S_{1,j}+\hbar^2S_{2,j}.
 \label{eq:v3-action-expansion}
\end{equation}
At each loop order the symbol $S_{\ell,j}$ means the \emph{single common local
coefficient action} already containing every lower-order restoring term and
all of its source/antifield descendants.

\begin{theorem}[Cocycle property of the first un-restored shell coefficient]
\label{thm:v3-breaking-closed}
Suppose the common master equation has been restored through loop order
$\ell-1$.  Let $\widetilde{\mathcal B}^{(\ell,\le r)}_j$ be the complete
coefficient of $\hbar^\ell$ in the native shell master defect before the new
order-$\ell$ local counterterm is chosen, with the bosonic, chiral,
ghost/nonminimal, antifield, line, measure/reference, stress and contact rows
assembled.  Then
\begin{equation}
 (S_{0,j},\widetilde{\mathcal B}^{(\ell,\le r)}_j)=0.
 \label{eq:v3-unprojected-closed}
\end{equation}
Consequently its local coefficient
\begin{equation}
 \mathcal B^{(\ell,\le r)}_j
 :=\mathsf P^{\rm ESM}_{N,6,r}
   \widetilde{\mathcal B}^{(\ell,\le r)}_j
 \label{eq:v3-local-breaking}
\end{equation}
satisfies
\begin{equation}
 \boxed{d_j\mathcal B^{(\ell,\le r)}_j=0.}
 \label{eq:v3-local-breaking-closed}
\end{equation}
\end{theorem}

\begin{proof}
Expand \eqref{eq:v3-quantum-WZ} in powers of $\hbar$.  At the first order at
which the master defect has not yet been cancelled, every term in which a
positive-loop part of $S$ acts on a lower-order master defect vanishes by the
induction hypothesis.  The surviving term is therefore precisely
$(S_{0,j},\widetilde{\mathcal B}_j^{(\ell,\le r)})$.  This proves
\eqref{eq:v3-unprojected-closed}.  Apply the local projector and use
Theorem~\ref{thm:v3-chain-complex}.  Because all source contacts were assembled
before projection, no differentiated contact is omitted in this step.
\end{proof}

\begin{remark}[Explicit first two coefficients]
With the convention \eqref{eq:v3-action-expansion}, the formal master
coefficients are
\begin{align}
 \mathcal M_j^{(0)}&=\frac12(S_{0,j},S_{0,j}),\nonumber\\
 \mathcal M_j^{(1)}&=(S_{0,j},S_{1,j})-\Delta_jS_{0,j},
 \label{eq:v3-master-coefficients}\\
 \mathcal M_j^{(2)}&=(S_{0,j},S_{2,j})
  +\frac12(S_{1,j},S_{1,j})-\Delta_jS_{1,j}.\nonumber
\end{align}
Thus the order-two row cannot be repaired independently of the already chosen
order-one action.  V3 always forms the complete coefficient first and only then
applies the common local restorer.
\end{remark}

\section{A gauge-invariant local frame for the complete chiral packet}
\label{sec:v3-invariant-frame}

We now use the strongest native finite statement already available instead of
recomputing anomaly traces.  Let $\Sigma_n^{\rm SM}$ denote the tensor product
of the coordinate-free chiral determinant sections of the complete Standard
Model packet at cutoff $n$, including the generation multiplicity selected by
the programme.  On the protected component it is nowhere zero.  The supplied
Grand Tensor gauge-descent theorem and the inherited anomaly packet give
\begin{equation}
 g\cdot\Sigma_n^{\rm SM}=\Sigma_n^{\rm SM}
 \qquad(g\in\mathcal G_n)
 \label{eq:v3-full-section-invariant}
\end{equation}
for the connected finite gauge action on the declared chart.  Its Hermitian
norm is gauge invariant as well.

\begin{definition}[Invariant section frame]
\label{def:v3-invariant-frame}
On the protected nonzero component define
\begin{equation}
 \boxed{
 \lambda^{\rm inv}_n
 :=\frac{\Sigma_n^{\rm SM}}{\|\Sigma_n^{\rm SM}\|}.}
 \label{eq:v3-invariant-frame}
\end{equation}
This is a unit-norm frame of the complete determinant line.  It is a local
reference choice on the protected component, not a statement about global line
holonomy across zero modes or disconnected topological sectors.
\end{definition}

\begin{theorem}[Exact gauge invariance and fixed-order locality of the invariant frame]
\label{thm:v3-invariant-frame}
The frame \eqref{eq:v3-invariant-frame} satisfies
\begin{equation}
 \boxed{g\cdot\lambda^{\rm inv}_n=\lambda^{\rm inv}_n,
 \qquad s_j\lambda^{\rm inv}_n=0.}
 \label{eq:v3-invariant-frame-BRST}
\end{equation}
At every fixed source order required by the two-loop packet, its local source
connection coefficients and the transition coefficients between
$\lambda_n^{\rm inv}$ and the V1 Kato frame obey the same rooted/quasilocal
bounds as the determinant-line derivatives proved in V1--V2.
\end{theorem}

\begin{proof}
Equation \eqref{eq:v3-full-section-invariant} and gauge invariance of the norm
give the first identity.  Differentiating the finite gauge action at the
identity gives the BRST statement.

For regularity, write locally
$\Sigma_n^{\rm SM}=z_n\lambda_n^{\rm K}$ in the V1 Kato frame.  Then
\[
 \lambda_n^{\rm inv}
 =u_n\lambda_n^{\rm K},\qquad
 u_n=\frac{z_n}{|z_n|}\in U(1).
\]
Every fixed derivative of $\log z_n$ is a finite sum of the inverse--vertex,
Kato-frame and line words controlled in V1.  Derivatives of
$\log|z_n|=\operatorname{Re}\log z_n$ are the corresponding real parts.
Therefore derivatives of $\log u_n=i\operatorname{Im}\log z_n$ obey the same
fixed-order rooted/quasilocal estimates.  Multiplying a Kato frame by the
unit-modulus scalar $u_n$ preserves the asserted source bounds.  The argument
uses the nonzero divisor margin and is intentionally local in the protected
component.
\end{proof}

\begin{firewall}[The invariant frame is reference data, not a global physical connection]
The construction uses the complete nonzero determinant section to choose a
local frame.  It does not assert that an independently measured determinant-line
connection equals this frame connection, nor does it identify holonomy around
a large gauge loop.  Such global/physical line data remain separate exactly as
in V1.  E2 needs a local perturbative frame in which the shell master identity
can be restored; that is what Theorem~\ref{thm:v3-invariant-frame} supplies.
\end{firewall}

\section{Gauge invariance survives the exact shell cocycle}
\label{sec:v3-shell-line}

The anomaly question is shellwise, so the full-cutoff frame must be transported
through the exact determinant Schur factorization rather than simply assumed to
factor.

Let $\Sigma_{<j}^{\rm SM}$ denote the determinant section of the retained
soft block before eliminating shell $j$, and let $\sigma_j^{\rm SM}$ be the
coordinate-free shell section in the inherited exact Schur cocycle, so that
schematically
\begin{equation}
 \Sigma_{<j+1}^{\rm SM}
 =\sigma_j^{\rm SM}\otimes\Sigma_{<j}^{\rm SM}.
 \label{eq:v3-shell-section-cocycle}
\end{equation}
All tensor products are the actual determinant-line products; the notation is
not a multiplication of unrelated scalar determinants.

\begin{lemma}[Equivariance of the Schur factorization]
\label{lem:v3-Schur-equivariance}
Assume the hard/soft splitting is gauge covariant.  If a finite chiral block is
written
\[
 D=\begin{pmatrix}A&B\\ C&E\end{pmatrix}
\]
with invertible hard pivot $A$, then under the induced block gauge maps its
Schur complement
$S=E-CA^{-1}B$ transforms covariantly on the soft chiral spaces.  The exact
determinant-line isomorphism
\begin{equation}
 \Det D\simeq \Det A\otimes\Det S
 \label{eq:v3-det-Schur-line}
\end{equation}
is equivariant.
\end{lemma}

\begin{proof}
Write the source and target gauge maps as
$U_-^{\hard}\oplus U_-^{\soft}$ and
$U_+^{\hard}\oplus U_+^{\soft}$.  Covariance gives
\begin{align*}
 A'&=U_-^{\hard}A(U_+^{\hard})^{-1},&
 B'&=U_-^{\hard}B(U_+^{\soft})^{-1},\\
 C'&=U_-^{\soft}C(U_+^{\hard})^{-1},&
 E'&=U_-^{\soft}E(U_+^{\soft})^{-1}.
\end{align*}
Substitution gives
$S'=U_-^{\soft}S(U_+^{\soft})^{-1}$.
Taking top exterior powers yields the equivariant line isomorphism.  This is an
identity before any scalar frame is chosen.
\end{proof}

\begin{theorem}[The complete Standard-Model shell section is gauge invariant]
\label{thm:v3-shell-section-invariant}
On the protected component, every coordinate-free shell factor in
\eqref{eq:v3-shell-section-cocycle} obeys
\begin{equation}
 \boxed{g\cdot\sigma_j^{\rm SM}=\sigma_j^{\rm SM}.}
 \label{eq:v3-shell-section-invariant}
\end{equation}
The same statement holds for a locally oriented neutral-Majorana Pfaffian
factor when that optional block is included.
\end{theorem}

\begin{proof}
The full-cutoff complete Standard-Model sections on the two sides of one shell
are gauge invariant by the inherited gauge-descent theorem.  The splitting is
gauge covariant, and Lemma~\ref{lem:v3-Schur-equivariance} makes the exact shell
factorization equivariant.  Since the retained soft section is nonzero on the
protected component, tensor cancellation in the one-dimensional line gives
\eqref{eq:v3-shell-section-invariant}.  Equivalently, the shell section is the
line ratio of two invariant nonzero sections.  The Pfaffian statement is the
same argument on the chosen oriented nonzero skew component.
\end{proof}

Define the invariant shell frame
\begin{equation}
 \lambda_{j}^{\rm inv}
 :=\frac{\sigma_j^{\rm SM}}{\|\sigma_j^{\rm SM}\|}.
 \label{eq:v3-shell-invariant-frame}
\end{equation}
Then the scalar shell coefficient in this frame is the positive function
$\|\sigma_j^{\rm SM}\|$.  In particular its internal-gauge BRST variation is
identically zero at finite cutoff.

\begin{corollary}[Frame-change primitive in the V1 Kato trivialization]
\label{cor:v3-frame-change}
Let $\lambda_j^{\rm K}$ be the V1 Kato shell frame and write
\begin{equation}
 \lambda_j^{\rm inv}=e^{\vartheta_j}\lambda_j^{\rm K},
 \qquad \operatorname{Re}\vartheta_j=0
 \label{eq:v3-frame-transition}
\end{equation}
with a smooth logarithm on a sufficiently small star-shaped source chart.
Then the scalar shell actions in the two frames differ by the single local-line
reference term $\vartheta_j$, and every fixed differentiated coefficient of
$\vartheta_j$ belongs to the V1--V2 line source calculus.  After the common
local Taylor assignment,
\begin{equation}
 \|\partial_I(\vartheta_j-\mathsf T_j\vartheta_j)\|_{\rm pan}
 \le C_I\frac{\mu}{\Lambda_j}.
 \label{eq:v3-frame-transition-remainder}
\end{equation}
\end{corollary}

\begin{proof}
The transition is a smooth $U(1)$-valued function because both frames are
unit-norm and nonzero.  The source derivative statement is
Theorem~\ref{thm:v3-invariant-frame}.  The one-excess-scale remainder is exactly
the V2 line-excess theorem applied to the difference of two admissible local
line frames; its derivatives are finite linear combinations of the same
line-logarithmic words.  The term is retained as reference/source data and is
not discarded because it is field independent in any restricted field sector.
\end{proof}

\section{Native zero of the complete standard anomaly projection}
\label{sec:v3-anomaly-zero}

We can now prove the E2 obstruction statement at the level actually consumed by
the local BV range theorem.

\begin{proposition}[One-loop native trace factorization before assembly]
\label{prop:v3-trace-factorization}
For a single left-handed Weyl multiplet $R$ using the common V2 Wilson/overlap
kinetic profile, the ghost-linear zero-Higgs local projection of the one-loop
chiral shell breaking onto \eqref{eq:v3-anomaly-space} has the form
\begin{equation}
 \mathsf A_{\rm std}\mathcal B^{(1,0)}_{j,R}
 =\sum_A\kappa_{A,j}\,t_A(R)\,\mathcal A_A,
 \label{eq:v3-trace-factorization}
\end{equation}
where $\kappa_{A,j}$ depends on the common native momentum/spin regulator and
normalization convention but not on the internal representation $R$, while
\begin{equation}
 \begin{aligned}
 t_{333}(R)&=\tr_R T_3^{(a}\!T_3^bT_3^{c)},&
 t_{Y33}(R)&=\tr_R(YT_3^aT_3^b),\\
 t_{Y22}(R)&=\tr_R(YT_2^iT_2^k),&
 t_{YYY}(R)&=\tr_RY^3,\\
 t_{YRR}(R)&=\tr_RY.
 \end{aligned}
 \label{eq:v3-trace-tensors}
\end{equation}
The proposition makes no claim about the numerical ordinary-continuum
normalization of the $\kappa_{A,j}$.
\end{proposition}

\begin{proof}
Expand the common gauge-covariant Wilson links and the overlap functional
calculus in external gauge/ghost jets before applying the anomaly quotient.
The lattice/spin/momentum part is the same for every species because the native
kinetic profile is common.  Representation dependence enters only through the
matrices multiplying the external gauge and ghost insertions.  The standard
anomaly projection selects the ghost-one local jet with precisely the invariant
index tensor of the corresponding basis element.  Taking the finite internal
trace therefore factors the coefficient into a common kinematic number and the
symmetrized representation traces in \eqref{eq:v3-trace-tensors}.  The mixed
gravity row contains one hypercharge generator and the spin curvature trace,
so its representation factor is $\tr_RY$.

Gauge-covariant Yukawa/Higgs insertions either carry positive Higgs degree and
are killed by the zero-Higgs anomaly-coordinate evaluation or belong, by the
matter-inclusive cohomological reduction, to the exact complement of
$\mathfrak A_{\rm std}$.  The same is true of antifield/source contacts not
belonging to the five quotient coordinates.  Thus no additional
representation tensor contributes to \eqref{eq:v3-trace-factorization}.
\end{proof}

\begin{remark}[Why V3 does not normalize $\kappa_{A,j}$]
The inherited Standard-Model trace vector is zero componentwise.  Hence the
complete one-loop anomaly projection vanishes for arbitrary finite values of
the common $\kappa_{A,j}$.  Proving that each $\kappa_{A,j}$ equals a particular
ordinary-continuum convention is a regulator-normalization theorem for an
\emph{uncancelled} representation; it is not needed to prove the master identity
of the anomaly-free packet.  V3 therefore does not pay that extra constant.
\end{remark}

\begin{theorem}[Zero standard anomaly projection for the complete native shell]
\label{thm:v3-native-anomaly-zero}
For the complete Standard-Model packet on the V2 protected branch, and for the
complete first- and second-loop shell master breakings after all lower-order
restorers and all source/line/reference contacts have been assembled,
\begin{equation}
 \boxed{
 \mathsf A_{\rm std}\mathcal B_j^{(1,\le r)}=0,
 \qquad
 \mathsf A_{\rm std}\mathcal B_j^{(2,\le r)}=0.}
 \label{eq:v3-native-anomaly-zero}
\end{equation}
More generally, every coefficient obtained by differentiating the same finite
shell identity through the prescribed finite source order has zero projection.
\end{theorem}

\begin{proof}
At one loop, Proposition~\ref{prop:v3-trace-factorization} reduces the five
quotient coordinates to the five inherited Standard-Model trace sums.  Those
sums vanish componentwise for the complete generation packet, and generation
replication multiplies zero by the generation number.  The inherited even
$SU(2)$ doublet parity handles the separate global $SU(2)$ sign; it is not a
local ghost-one coordinate in \eqref{eq:v3-anomaly-space}.

For the complete native shell there is also a stronger finite statement.
Choose the invariant shell frame \eqref{eq:v3-shell-invariant-frame}.  The
chiral Berezin line factor is then internally gauge invariant exactly.  Every
interaction vertex in the finite Standard-Model action is gauge covariant, the
bosonic Haar/Lebesgue and ghost/nonminimal measures have the V80 finite
invariance properties, and the hard/soft blocking is gauge covariant.  Therefore
the complete finite hard pushforward, including non-Gaussian connected terms,
has zero internal-gauge variation before a loop expansion is taken.  Extracting
the coefficient of $\hbar^2$, differentiating in retained sources, and finally
applying the local anomaly-coordinate projection preserve this exact zero.
This proves the second identity without an Adler--Bardeen input.

If instead one works in the Kato frame, Corollary~\ref{cor:v3-frame-change}
shows that the change from the invariant frame is one explicit local-line
reference coboundary plus a quasilocal $O(\mu/\Lambda_j)$ remainder.  A change
of local frame therefore cannot create a nontrivial quotient coordinate.  All
mixed-shell line contacts remain present because the frame change is made
before the exact differentiated shell product is expanded.
\end{proof}

\begin{firewall}[Why exact finite gauge invariance does not trivialize the whole BV problem]
Theorem~\ref{thm:v3-native-anomaly-zero} removes only the nontrivial
internal-gauge quotient.  The cutoff shell can still generate local BV-exact
breaking in the diffeomorphism, local-Lorentz, source, antifield, gauge-fixing,
stress/contact and frame/reference rows.  Those terms must be removed by one
common primitive.  Gauge invariance of the chiral line is therefore not a
substitute for the remaining E2 restoration theorem.
\end{firewall}

\section{One common local primitive by finite Hodge contraction}
\label{sec:v3-hodge-restorer}

Theorem~\ref{thm:v3-breaking-closed} and
Theorem~\ref{thm:v3-native-anomaly-zero} put the complete breaking in the exact
part of the finite source-extended complex.  We now construct the restorer
without choosing counterterms graph by graph.

Choose once and for all a positive coefficient inner product on
$\widehat{\mathfrak L}^{\rm ESM}_{N,6,r}$, compatible with the finite
representation decomposition but otherwise arbitrary.  Let $d_j^\dagger$ be
the adjoint and
\begin{equation}
 \Box_j=d_jd_j^\dagger+d_j^\dagger d_j.
 \label{eq:v3-Hodge-Laplacian}
\end{equation}
Let $\mathsf H_j$ be the orthogonal projection onto $\ker\Box_j$ and let
$G_j$ be the inverse of $\Box_j$ on $(\ker\Box_j)^\perp$, extended by zero on
the harmonic space.  Define
\begin{equation}
 h_j:=d_j^\dagger G_j.
 \label{eq:v3-Hodge-homotopy}
\end{equation}

\begin{lemma}[Finite contraction identity]
\label{lem:v3-contraction}
On the complete finite source bank,
\begin{equation}
 \boxed{d_jh_j+h_jd_j=I-\mathsf H_j.}
 \label{eq:v3-contraction-identity}
\end{equation}
If $B$ has ghost number one, $d_jB=0$, and
$\mathsf A_{\rm std}B=0$, then
\begin{equation}
 \boxed{B=d_jh_jB.}
 \label{eq:v3-exact-hodge}
\end{equation}
\end{lemma}

\begin{proof}
Finite-dimensional Hodge decomposition gives
$\widehat{\mathfrak L}=\operatorname{Ran}d_j\oplus
\ker\Box_j\oplus\operatorname{Ran}d_j^\dagger$ orthogonally and the standard
Green identity gives \eqref{eq:v3-contraction-identity}.  By
Theorem~\ref{thm:v3-chain-complex}, the only nontrivial ghost-one harmonic
coordinates are the standard anomaly quotient, up to the already removed
contractible pairs and redundant directions.  Thus
$\mathsf A_{\rm std}B=0$ and $d_jB=0$ imply $\mathsf H_jB=0$.  Apply
\eqref{eq:v3-contraction-identity} to $B$.
\end{proof}

\begin{proposition}[Uniform finite-order restorer on a protected constant-rank chart]
\label{prop:v3-uniform-homotopy}
Shrink the finite perturbative normalization chart, if necessary, to one compact
component on which the ranks of $d_j$ in the retained degrees are constant.
Then there is $\delta_*>0$, independent of shell number on that chart, such that
the nonzero singular values of $d_j$ are at least $\delta_*$.  Consequently
\begin{equation}
 \|h_j\|\le \delta_*^{-1},
 \qquad
 \|G_j\|\le \delta_*^{-2}.
 \label{eq:v3-homotopy-bound}
\end{equation}
\end{proposition}

\begin{proof}
At fixed $(N,6,r)$ the matrices of $d_j$ are finite and depend continuously on
the finitely many normalized couplings/reference coordinates.  Constant rank
means the smallest nonzero singular value is a positive continuous function on
the chosen compact component.  Its minimum is therefore positive.  The two
operator bounds follow from the singular-value formulas for the Moore--Penrose
inverse and for the Green operator.
\end{proof}

\begin{theorem}[One common source-complete restoring primitive]
\label{thm:v3-common-primitive}
For $\ell=1,2$, let
$\mathcal B_j^{(\ell,\le r)}$ be the recursively assembled local shell breaking.
Define
\begin{equation}
 \boxed{
 \Xi_j^{(\ell,\le r)}:=h_j\mathcal B_j^{(\ell,\le r)}.}
 \label{eq:v3-common-primitive}
\end{equation}
Then
\begin{equation}
 \boxed{
 d_j\Xi_j^{(\ell,\le r)}
 =\mathcal B_j^{(\ell,\le r)}.}
 \label{eq:v3-common-primitive-equation}
\end{equation}
The object $\Xi_j^{(\ell,\le r)}$ is one element of the complete
source-extended local bank.  Its gauge, diffeomorphism, local-Lorentz,
Higgs/Yukawa, ghost, antifield, stress/contact and determinant/Pfaffian-line
coefficients are its source/field components; no descendant is chosen
independently.
\end{theorem}

\begin{proof}
Closure is Theorem~\ref{thm:v3-breaking-closed}; zero harmonic projection is
Theorem~\ref{thm:v3-native-anomaly-zero} together with the V80 cohomological
reduction.  Lemma~\ref{lem:v3-contraction} gives
\eqref{eq:v3-common-primitive-equation}.  Because the source jet algebra was
included in \eqref{eq:v3-source-extended-bank} before the Hodge problem was
solved, the primitive is a single polynomial in all retained source variables.
Differentiating it merely reads off coefficients already present in that one
element.  In particular, contacts generated because the differential itself
acts on a source/antifield multiplet are part of the same finite matrix $d_j$;
they are not added afterward.
\end{proof}

\begin{corollary}[Restoring counterterms through two loops]
\label{cor:v3-two-loop-counterterms}
At the first un-restored order set
\begin{equation}
 C_{\ell,j}:=-\Xi_j^{(\ell,\le r)}.
 \label{eq:v3-counterterm}
\end{equation}
Insert $C_{1,j}$ into the common action before the complete order-two master
coefficient is formed; then construct $C_{2,j}$ from that recursively corrected
order-two breaking.  The resulting action satisfies the projected local quantum
master equation through $O(\hbar^2)$ on the declared source bank.
\end{corollary}

\begin{proof}
Adding $\hbar^\ell C_{\ell,j}$ changes the first un-restored coefficient of the
master functional by $d_jC_{\ell,j}$; all terms involving lower master defects
vanish by recursive restoration.  Equation
\eqref{eq:v3-common-primitive-equation} cancels the coefficient.  At order two,
forming the breaking only after $C_{1,j}$ has been inserted automatically
retains the induced terms from
$\frac12(S_{1,j},S_{1,j})-\Delta_jS_{1,j}$, including every source contact.
\end{proof}

\section{The chiral two-loop hard remainder has the same one-excess-scale gain}
\label{sec:v3-two-loop-remainder}

Local solvability alone is not a shell self-map.  We also need the complementary
rooted kernels.  V1--V2 supply the first-order fermion hard-line row/column
bounds, every fixed source derivative, the full-zone cap estimate, and the
$O(\mu/\Lambda_j)$ line remainder.  The remaining point is that the inherited
V80 two-loop forest ownership still gives one excess hard scale when fermion
lines are admitted.

Let $G$ be a fixed connected two-loop ESM hard graph, possibly with mixed
bosonic and fermionic hard lines, after every proper divergent one-loop subgraph
has been assigned to the same common order-one action.  Let $\omega_G$ be its
superficial local degree in the declared EFT/source counting and let
$\mathsf T_G$ be the complete Taylor/EFT assignment through that degree.

\begin{theorem}[Mixed chiral two-loop hard-annulus remainder]
\label{thm:v3-chiral-annulus}
At fixed $(N,6,r)$ and on the V2 protected branch, every such graph obeys,
after the common forest subtraction,
\begin{equation}
 \boxed{
 \|\partial_K^\alpha(1-\mathsf T_G)\mathcal I_{G,j}(K)\|_{\rm root}
 \le C_{G,\alpha}\frac{\mu}{\Lambda_j}\,
       \mathfrak e_{G,\alpha}(\mu)}
 \label{eq:v3-chiral-graph-remainder}
\end{equation}
for the finite external panel $|K|\le\mu$, where
$\mathfrak e_{G,\alpha}(\mu)$ is the fixed engineering source stock of the
corresponding retained coefficient.  At the ultraviolet-cap shells the same
bound is read as $C a\mu$.  Summing the finite two-loop graph/forest family gives
\begin{equation}
 \boxed{
 \|(1-\mathsf P^{\rm ESM}_{N,6,r})\mathcal K^{\rm ESM}_{2,j}\|_{*\to *}
 \le C_{N,r}\frac{\mu}{\Lambda_j}.}
 \label{eq:v3-K2-bound}
\end{equation}
\end{theorem}

\begin{proof}
For each physical annulus, rescale every independent hard momentum by
$p=\Lambda_jk$.  Bosonic hard propagators and vertices have the inherited V80
stocks.  Every fermionic hard propagator has the weighted row/column stock
$O(\Lambda_j^{-1})$ and every fixed source derivative is a finite inverse--vertex
word with the V2 normalized engineering bounds.  Hence every fixed two-loop
integrand, after the proper one-loop forests have been subtracted, is a smooth
function of the dimensionless loop variables and of $K/\Lambda_j$, multiplied
by its declared engineering power.

The local bank contains the complete Taylor jet through the superficial degree
$\omega_G$.  The integral Taylor formula for
$(1-\mathsf T_G)\mathcal I_{G,j}$ therefore leaves at least one additional
factor $|K|/\Lambda_j$ beyond the last local coefficient.  The V1 rooted
connected-contraction theorem converts the weighted row/column bounds into a
volume-uniform rooted kernel estimate.  The inherited V80 forest ownership
ensures that overlapping proper subgraphs are subtracted by restrictions of the
same $S_{1,j}$ rather than by independent graph constants, so the preceding
Taylor gain applies to the complete forest sum.

There are only finitely many graph and forest types at fixed loop order, EFT
weight, arity and source order; summing their constants proves
\eqref{eq:v3-K2-bound}.  On the ultraviolet cap, V2 replaces the physical
annular variable by $x=ap$ and gives one factor $a\mu$, which is
$O(\mu/\Lambda_j)$ because $\Lambda_j\asymp a^{-1}$ there.
\end{proof}

\begin{remark}[No hidden fermionic positivity]
The proof uses the V1 weighted Schur/rooted bounds and ordinary trace-ideal
estimates for Grassmann contractions.  It never replaces a fermionic connected
trace by a positive bosonic variance and never assigns two independent small
factors to one connected contraction.
\end{remark}

\begin{corollary}[Common chiral source and line step]
\label{cor:v3-source-step}
Let $E^{\rm ESM}_{1,j}$ and $E^{\rm ESM}_{2,j}$ be the first- and second-loop
adjacent source/line steps after the local restorers have been applied.  Then
\begin{equation}
 \|E^{\rm ESM}_{1,j}\|_{*\to *}
 +\|E^{\rm ESM}_{2,j}\|_{*\to *}
 \le C_{N,r}\frac{\mu}{\Lambda_j},
 \label{eq:v3-source-step}
\end{equation}
with the quadratic product term bounded by
$C(\mu/\Lambda_j)^2$.  All differentiated determinant/Pfaffian products retain
the mixed-shell Bell contacts of V1.
\end{corollary}

\begin{proof}
The graph part is Theorem~\ref{thm:v3-chiral-annulus}; the line part is the V2
line-excess theorem and its Bell-polynomial corollary.  The local Hodge restorer
is finite dimensional and bounded by Proposition~\ref{prop:v3-uniform-homotopy},
so applying it to the local coefficient does not alter the complementary
one-excess-scale estimate.  Products of two first-order source steps have the
stated quadratic bound.
\end{proof}

\begin{corollary}[Uniform finite-order shell composition]
\label{cor:v3-source-product}
For geometric shell scales $\Lambda_j=b^j\lambda$,
\begin{equation}
 \sup_{n\ge j}
 \|\mathsf U^{\rm ESM,[\le2]}_{n,j}\|_{*\to *}
 \le
 \exp\!\left(C\sum_{m\ge j}\frac{\mu}{\Lambda_m}\right)
 \le
 \exp\!\left(\frac{Cb}{b-1}\frac{\mu}{\Lambda_j}\right).
 \label{eq:v3-composed-source}
\end{equation}
The bound is coefficientwise at the fixed two-loop/EFT/source order; it is not
an analytic finite-coupling RG trajectory.
\end{corollary}

\section{One common two-loop Einstein--Standard Model shell update}
\label{sec:v3-selfmap}

We can now assemble the E2 acceptance criterion.  Let the complete finite-order
state be
\begin{equation}
 \begin{split}
 \mathfrak S_j^{\rm ESM}=
 \bigl(&\mathfrak S_j^{\rm EGH,bos};
 Z_{\psi,j},\mathbf Y_j,w^{F}_{6,j};
 \mathcal N^{F}_{<,j},\mathcal G^{F}_{<,j};\\
 &\lambda^{\rm inv}_{\chi,j},\mathcal L_{\chi,j};
 \mathcal J^{\rm all}_j,\mathcal T^{\rm all}_j;
 \mathcal C_{1,j},\mathcal C_{2,j}\bigr).
 \end{split}
 \label{eq:v3-ESM-state}
\end{equation}
Here $w^F_{6,j}$ denotes the complete fermionic ghost-zero weight-six quotient,
not a hand-picked operator list; it contains all admitted dipole, current,
four-fermion, curvature--fermion, Higgs--fermion and redundant partner
directions selected by the finite quotient.  $\mathcal L_{\chi,j}$ denotes the
finite determinant/Pfaffian line-source/contact packet.  The symbols
$\mathcal C_{\ell,j}$ record the common local restoring coefficients generated
by \eqref{eq:v3-common-primitive}, not separate graphwise choices.

\begin{theorem}[Complete chiral ESM two-loop shell self-map]
\label{thm:v3-ESM-selfmap}
Fix the V2 Wilson/overlap branch, a protected nonzero determinant/Pfaffian
component, the positive IR/reference prescription, two-loop order, EFT weight
six, source order $r$, arity $N$, and the compact constant-rank BV chart of
Proposition~\ref{prop:v3-uniform-homotopy}.  After imposing one common set of
physical lower-weight normalization conditions, the normalized native shell
defines
\begin{equation}
 \boxed{
 \mathcal R_j^{\rm ESM,[\le2]}:
 \mathfrak C^{\rm ESM}_{N,6,r}
 \longrightarrow
 \mathfrak C^{\rm ESM}_{N,6,r}.}
 \label{eq:v3-ESM-selfmap}
\end{equation}
It has the following properties.
\begin{enumerate}[label=\textup{(E2.\arabic*)},leftmargin=3.4em]
\item The complete native internal-gauge anomaly projection is zero at one and
two loops, including all retained source descendants.
\item The complete local ghost-one breaking is cancelled by the single
source-extended primitives $\Xi_j^{(1,\le r)}$ and
$\Xi_j^{(2,\le r)}$; no graph or sector receives an independently chosen
restorer.
\item Independent ghost-zero ESM EFT coordinates are retained in the state and
are not removed merely because they occur in the same Taylor bank as a
restoring direction.
\item The mixed boson--fermion quasilocal two-loop complement obeys
\eqref{eq:v3-K2-bound}, and the determinant/Pfaffian line remainder obeys the
V2 $O(\mu/\Lambda_j)$ estimate with all Schur and mixed-shell contacts retained.
\item The adjacent source/line map obeys \eqref{eq:v3-source-step} and composes
uniformly as in \eqref{eq:v3-composed-source} at this fixed perturbative order.
\item Gravity, ghosts, Higgs, all three internal gauge fields, chiral fermions,
Yukawas, antifields, measure/reference terms and stress/contact sources occur
in the same action/source/line update.
\end{enumerate}
Thus the E2 acceptance criterion --- one common two-loop Einstein--Standard
Model action/source/line update --- is satisfied on this explicit native branch.
\end{theorem}

\begin{proof}
The inherited V80 theorem supplies the bosonic EGH self-map and its common
normalization.  Theorem~\ref{thm:v3-native-anomaly-zero} removes the only
nontrivial ghost-one quotient of the complete matter-inclusive local complex.
Theorem~\ref{thm:v3-common-primitive} and
Corollary~\ref{cor:v3-two-loop-counterterms} restore the remaining local
breaking with one finite action.  Theorem~\ref{thm:v3-chiral-annulus} closes the
new mixed chiral hard remainder; the V2 line theorem closes the line factor;
and Corollaries~\ref{cor:v3-source-step}--\ref{cor:v3-source-product} close the
source transport.  All of these constructions use the same source bank and the
same lower-order action before any sector is projected, so the result is one
shell map rather than a product of independently repaired sector maps.
\end{proof}

\begin{firewall}[Exact scope of the V3 E2 closure]
Theorem~\ref{thm:v3-ESM-selfmap} is a fixed-order coefficientwise theorem.  It
does not prove convergence of the perturbation series, uniformity in loop
order, an analytic finite-coupling RG trajectory, global determinant-line
holonomy, passage through chiral zero modes, removal of the retained IR scale,
or a nonperturbative ESM measure.  It is native to the explicit V2 Wilson
completion unless a separate regulator-comparison theorem is supplied.
\end{firewall}

\section{What remains for E3: the common master coefficient must be written as one object}
\label{sec:v3-E3-handoff}

E2 has now been closed at the shell-self-map level.  The next danger is a more
subtle kind of sectorwise reasoning: one could cite the bosonic V80 master row,
the chiral line invariance, and the Hodge restorer separately and declare the
complete ESM master equation solved.  E3 forbids that shortcut.

Define the common renormalized coefficients after the V3 shell update by
\begin{equation}
 \bar S_{1,j}=S_{1,j}+C_{1,j},
 \qquad
 \bar S_{2,j}=S_{2,j}+C_{2,j}.
 \label{eq:v3-renormalized-S12}
\end{equation}
The next theorem must write, in the same source/reference convention,
\begin{align}
 \mathcal M_{j,1}^{\rm all}
 &= (S_{0,j},\bar S_{1,j})-\Delta_jS_{0,j},
 \label{eq:v3-E3-M1}\\
 \mathcal M_{j,2}^{\rm all}
 &= (S_{0,j},\bar S_{2,j})
  +\frac12(\bar S_{1,j},\bar S_{1,j})
  -\Delta_j\bar S_{1,j},
 \label{eq:v3-E3-M2}
\end{align}
as one complete coefficient and verify every component with the same
$S_0,\bar S_1,\bar S_2$.  In particular, the cross antibrackets between
geometric, gauge, Higgs/Yukawa, chiral, ghost and source sectors must not be
replaced by separate Ward equations.

\begin{theorem}[Exact next load-bearing theorem for V4]
\label{thm:v3-next-E3}
Using the V3 common shell action, prove the full two-loop common
BV/master/source equation by expanding
\eqref{eq:v3-E3-M1}--\eqref{eq:v3-E3-M2} in the single source-extended bank and
checking simultaneously:
\begin{enumerate}[label=\textup{(M\arabic*)},leftmargin=3em]
\item diffeomorphism and retained gravitational constraint rows;
\item local Lorentz/spin rows;
\item all three internal gauge rows;
\item Higgs and Yukawa covariance rows;
\item chiral determinant/Pfaffian line and frame/reference rows;
\item ghost, nonminimal and antifield descendants;
\item measure/Jacobian/reference terms;
\item composite, stress and contact sources;
\item every cross antibracket/contact generated by the common $S_1$ in the
quadratic term $\frac12(\bar S_1,\bar S_1)$.
\end{enumerate}
The theorem must end with one coefficient identity
$\mathcal M_{j,1}^{\rm all}=\mathcal M_{j,2}^{\rm all}=0$ in the retained local
bank plus the already controlled quasilocal remainder.  It is not enough to
list sectorwise zeroes.
\end{theorem}

\section{V3 anti-circularity audit}
\label{sec:v3-audit}

The V3 route avoids four possible circularities.
\begin{enumerate}[leftmargin=2.4em]
\item It does not infer a native anomaly statement from the already-known
continuum trace cancellation.  The one-loop trace factorization is derived from
the native common kinetic profile, while the complete two-loop zero follows
from the exact finite gauge invariance of the assembled native shell.
\item It does not invoke an all-order anomaly nonrenormalization theorem.  The
finite gauge-invariance identity is expanded only after the complete shell is
assembled.
\item It does not solve a different counterterm equation for every graph or
source derivative.  The Hodge map $h_j$ acts once on the complete
source-extended breaking.
\item It does not erase line/reference data by choosing a convenient scalar
phase.  The invariant frame and its transition from the Kato frame are retained
as explicit line/reference coefficients with the same source estimates.
\end{enumerate}

\section{V3 status ledger}
\label{sec:v3-status}

\subsection*{Proved in V3}
\begin{enumerate}[leftmargin=2.2em]
\item The complete fixed-order ESM local/source bank is a finite BV chain
complex; at ghost number one its only nontrivial quotient is the standard
five-dimensional internal-gauge anomaly space inherited from V80.
\item The complete native shell master breaking is a BV cocycle at the first
un-restored loop order, including the complete source/contact packet.
\item The anomaly-free nonzero Standard-Model determinant section supplies a
smooth gauge-invariant local unit frame on the protected component.  Its
transition from the V1 Kato frame has the same fixed-order quasilocal source
bounds and an $O(\mu/\Lambda_j)$ locally subtracted shell remainder.
\item Exact Schur equivariance transports that gauge invariance to every shell
factor in the inherited determinant/Pfaffian cocycle.
\item The one-loop native anomaly seed factors into universal shell coefficients
times the five Standard-Model representation traces.  Their inherited
componentwise cancellation gives zero one-loop anomaly projection without
normalizing the universal coefficients.
\item The complete finite native shell is internally gauge invariant after the
full anomaly-free packet is assembled.  Hence its two-loop standard anomaly
projection and every retained source derivative of that projection vanish
exactly, without importing Adler--Bardeen nonrenormalization.
\item A single finite Hodge contraction on the complete source-extended bank
constructs the restoring primitive.  All gauge, diffeomorphism, Lorentz,
Higgs/Yukawa, ghost, antifield, stress/contact and line descendants belong to
that one primitive.
\item Mixed boson--fermion two-loop hard forests have the same one-excess-scale
$O(\mu/\Lambda_j)$ complementary rooted bound at fixed EFT/source order.
Together with the V2 line estimate this yields a uniformly composable adjacent
source/line step.
\item The inherited V80 bosonic self-map therefore enlarges to one complete
chiral Einstein--Standard Model two-loop action/source/line shell self-map on
the explicit V2 regulator branch.  E2 is closed at that scope.
\end{enumerate}

\subsection*{Inherited}
\begin{enumerate}[leftmargin=2.2em]
\item The V80 bosonic Einstein--gauge--Higgs two-loop self-map, finite local
cohomology reduction, zero native bosonic anomaly projection and common source
normalization.
\item The three-generation Standard-Model chiral representation, cancellation of
the five ordinary local anomaly trace sums, even $SU(2)$ doublet parity, and the
Grand Tensor gauge descent of the complete nonzero chiral determinant section.
\item Exact determinant/Pfaffian Schur factorization, finite regrouping cocycles
and differentiated line identities.
\item The V1 weighted functional calculus, Kato frame, hard-pivot and rooted
connected-contraction machinery.
\item The V2 explicit Wilson completion, all-zone gap/no-doubler theorem,
annular/full-zone finite-to-ordinary bounds, fixed source-jet calculus and
$O(\mu/\Lambda_j)$ differentiated line remainder.
\end{enumerate}

\subsection*{Conditional}
\begin{enumerate}[leftmargin=2.2em]
\item The result remains native to the explicit V2 Wilson/overlap completion.
A different microscopic Wilson kernel requires a finite regulator-comparison
theorem.
\item The Hodge-restorer norm is uniform on a compact constant-rank component of
the finite BV bank.  A cohomology-rank crossing requires a new local chart or a
separate analysis.
\item The determinant/Pfaffian construction remains on the declared nonzero
protected component and retains the positive IR/reference prescription.  No
global line holonomy or zero-crossing statement is made.
\item Constants may depend on fixed loop order, EFT degree, source order, arity,
representation data and the protected compact chart.  No uniform
$L\to\infty$ estimate is asserted.
\end{enumerate}

\subsection*{Open}
\begin{enumerate}[leftmargin=2.2em]
\item Write the \emph{actual common} first- and second-loop coefficients
\eqref{eq:v3-E3-M1}--\eqref{eq:v3-E3-M2} in one source-extended bank and audit
every cross antibracket/contact.  This is E3.
\item After E3, formulate the arbitrary fixed-loop induction state including the
finite EFT bank, complete contact tower, BV fields/antifields, line/frame data,
forest ownership, rooted remainders and reference/scheme data.  This is E4.
\item The ordinary-normalization constants $\kappa_{A,j}$ of an uncancelled
single Weyl representation were intentionally not identified; that comparison
is not load-bearing for the anomaly-free E2 theorem.
\end{enumerate}

\subsection*{Exact next load-bearing theorem}
\begin{center}
\fbox{\parbox{0.92\textwidth}{
\textbf{E3: one complete two-loop BV/master/source coefficient.}
Using the V3 common action $S_0+\hbar\bar S_1+\hbar^2\bar S_2$, expand the same
finite quantum master functional through order $\hbar^2$ and prove one common
identity containing the diffeomorphism, local-Lorentz, internal-gauge,
Higgs/Yukawa, chiral line, ghost/nonminimal, antifield, measure/reference,
stress/composite and contact rows, including every cross term in
$\frac12(\bar S_1,\bar S_1)$.  Do not replace this by sectorwise Ward identities.
}}
\end{center}

\section*{Append-only continuation marker: V3}
\addcontentsline{toc}{section}{Append-only continuation marker: V3}
\textbf{End of V3.}  The complete V1--V2 body above is retained verbatim.  Future
versions must append new work after this marker without replacing the V1--V3
mathematical body.  The next file is to be named
\texttt{Renewal\_Perturbative\_ESM\_Continuum\_Research\_Notes\_V4.tex}.


\clearpage
\part*{V4: The complete common two-loop master coefficient and the canonical-pair audit}
\addcontentsline{toc}{part}{V4: Complete common two-loop master coefficient and canonical-pair audit}
\label{sec:v4-start}

\section{E3 is a common-coefficient theorem, not another sectorwise restoration}
\label{sec:v4-E3-start}

V3 closes the two-loop shell self-map on the explicit V2 Wilson/overlap branch,
but deliberately leaves one bookkeeping theorem open.  The local Hodge
construction proves that the recursively assembled shell breaking is in the
range of one source-extended BV differential.  E3 asks for a stronger audit:
the actual coefficients of the \emph{same} finite master functional must be
written with the same $S_0,\bar S_1,\bar S_2$, including every cross
antibracket, source contact, determinant-line/reference contribution and
source-dependent BV-Laplacian term.

The correct upstream organization is by canonical BV pair.  The antibracket is
diagonal in canonical pairs.  Therefore every apparently ``mixed-sector'' term
is produced by one of finitely many canonical contractions, and there is no
additional cross Ward identity to impose after the common coefficient has been
formed.  This observation is what makes an exhaustive audit possible without
re-fitting graphs or writing unrelated sector equations.

Throughout V4 fix exactly the V3 domain: loop order through two, local weight
six, arity $N$, source order $r$, the finite external panel, the compact
constant-rank BV chart, the protected nonzero chiral line component, the
positive IR/reference prescription, and the explicit V2 Wilson completion.
All statements are coefficientwise on this finite domain.

\begin{firewall}[No strengthening of the global claim]
V4 does not assert that the full finite-coupling effective action exists as a
convergent analytic object.  The symbols $S_0,\bar S_1,\bar S_2$ are the fixed
coefficients of the perturbative shell action.  Nor does V4 remove the retained
IR/reference scale, cross a determinant zero, or establish uniformity as loop
order tends to infinity.
\end{firewall}

\section{Canonical carrier and source algebra}
\label{sec:v4-carrier}

\subsection{Central sources versus canonical source multiplets}

Let $\mathcal J_r^{\rm cen}$ be the truncated supercommutative algebra generated
by the retained \emph{central} physical source writers through total source
order $r$.  These are the sources with zero BV bracket.  Source variables which
belong to BRST doublets, antifield/source multiplets, or any other canonical
pair are \emph{not} placed in $\mathcal J_r^{\rm cen}$; they are included in the
canonical carrier below.  This is only a reorganization of the V3
source-extended bank and prevents a contact term from being counted both as a
source derivative and as a BV contraction.

Choose an ordered divided-power basis $J^I/I!$ of
$\mathcal J_r^{\rm cen}$.  For every source-polynomial functional $F(J)$ write
\begin{equation}
 F(J)=\sum_{|I|\le r}\frac{J^I}{I!}F_I,
 \qquad [F]_I:=F_I.
 \label{eq:v4-source-expansion}
\end{equation}
For odd central sources the notation means the ordered super divided-power
basis; permutations carry the corresponding Koszul sign.

\subsection{Exact canonical-pair inventory}

Let $\mathcal C_j$ be the finite list of canonical pairs retained at shell $j$.
We group it only for bookkeeping:
\begin{equation}
 \begin{split}
 \mathcal C_j={}&
 \mathcal C_{\rm geom}\sqcup
 \mathcal C_{\rm Lor}\sqcup
 \mathcal C_{3}\sqcup\mathcal C_{2}\sqcup\mathcal C_Y
 \sqcup\mathcal C_H\sqcup\mathcal C_F\\
 &\sqcup\mathcal C_{\rm diff-gh}\sqcup
 \mathcal C_{\rm Lor-gh}\sqcup
 \mathcal C_{\rm int-gh}\sqcup
 \mathcal C_{\rm nm}\sqcup
 \mathcal C_{\rm src}.
 \end{split}
 \label{eq:v4-canonical-inventory}
\end{equation}
Here $\mathcal C_{\rm geom}$ contains the coframe/metric variables and exact
antifield duals; $\mathcal C_{\rm Lor}$ contains the independent spin/local
Lorentz variables when present in the chosen chart; $\mathcal C_{3,2,Y}$ are
the three internal gauge pairs; $\mathcal C_H$ is the Higgs pair;
$\mathcal C_F$ contains every retained chiral quark/lepton pair (and the
optional neutral real/Majorana pair in its oriented chart);
$\mathcal C_{\rm diff-gh}$, $\mathcal C_{\rm Lor-gh}$ and
$\mathcal C_{\rm int-gh}$ are the ghost--ghost-antifield pairs;
$\mathcal C_{\rm nm}$ contains the contractible nonminimal pairs; and
$\mathcal C_{\rm src}$ contains the genuinely canonical source/marker
multiplets.  Spectator physical writers remain in $\mathcal J_r^{\rm cen}$.

For $a\in\mathcal C_j$ write the canonical coordinates as $(z^a,z_a^*)$.
With the inherited graded left/right derivative convention define the
pair contribution by
\begin{equation}
 (F,G)_a
 :=F\frac{\overleftarrow\delta}{\delta z^a}
      \frac{\overrightarrow\delta}{\delta z_a^*}G
   -F\frac{\overleftarrow\delta}{\delta z_a^*}
      \frac{\overrightarrow\delta}{\delta z^a}G,
 \qquad
 (F,G)=\sum_{a\in\mathcal C_j}(F,G)_a.
 \label{eq:v4-pair-bracket}
\end{equation}
All Koszul signs are contained in the inherited graded arrow-derivative
convention.  Every later argument uses only this exact Darboux-pair
decomposition and is therefore independent of a separate sign shorthand.

The finite BV Laplacian is kept in the same density/reference convention as V3.
Allow its source-dependent representation to have the finite expansion
\begin{equation}
 \Delta_j(J)=\sum_{|K|\le r}\frac{J^K}{K!}\Delta_{j,K}.
 \label{eq:v4-Delta-expansion}
\end{equation}
Thus source derivatives of the finite measure/reference density are retained
rather than silently set to zero.

\section{Exact source-contact calculus for the master functional}
\label{sec:v4-contact-calculus}

\begin{lemma}[Coefficient convolution for the antibracket and BV Laplacian]
\label{lem:v4-source-convolution}
Let $F(J),G(J)$ be finite source polynomials as in
\eqref{eq:v4-source-expansion}.  For every ordered source multi-index $I$,
\begin{equation}
 \boxed{
 [(F,G)]_I
 =\sum_{I_1\sqcup I_2=I}
 \varepsilon(I_1,I_2)
 \binom{I}{I_1}
 (F_{I_1},G_{I_2}),}
 \label{eq:v4-bracket-convolution}
\end{equation}
where $\varepsilon(I_1,I_2)$ is the Koszul sign required to move the odd
sources into the fixed order.  Likewise
\begin{equation}
 \boxed{
 [\Delta_j(J)F(J)]_I
 =\sum_{K\sqcup L=I}
 \varepsilon(K,L)
 \binom{I}{K}\,
 \Delta_{j,K}F_L.}
 \label{eq:v4-Delta-convolution}
\end{equation}
For a purely even source subpanel all signs are $+1$ and the sums are the
ordinary multinomial convolutions.
\end{lemma}

\begin{proof}
Expand both finite polynomials in the ordered divided-power basis.  The BV
functional derivatives act only on the canonical carrier because the writers in
$\mathcal J_r^{\rm cen}$ are central.  Multiplication of two basis monomials
$J^{I_1}/I_1!$ and $J^{I_2}/I_2!$ gives
$\varepsilon(I_1,I_2)\binom{I_1+I_2}{I_1}J^{I_1+I_2}/(I_1+I_2)!$.
Collecting the coefficient of $J^I/I!$ proves
\eqref{eq:v4-bracket-convolution}.  The same multiplication rule applied to the
operator expansion \eqref{eq:v4-Delta-expansion} proves
\eqref{eq:v4-Delta-convolution}.
\end{proof}

\begin{corollary}[The first nontrivial contact is forced]
\label{cor:v4-two-source-contact}
For two distinct even central writers $J^A,J^B$ and any even $F$,
\begin{equation}
 \left[\frac12(F,F)\right]_{AB}
 =(F_0,F_{AB})+(F_A,F_B).
 \label{eq:v4-two-source-contact}
\end{equation}
Thus the product of the two singly marked insertions is part of the same second
source derivative as the doubly marked contact coefficient.  Omitting either
term destroys the differentiated master identity.
\end{corollary}

\begin{proof}
Apply Lemma~\ref{lem:v4-source-convolution} to the four ordered splittings
$\varnothing\sqcup AB$, $A\sqcup B$, $B\sqcup A$ and
$AB\sqcup\varnothing$.  The antibracket of two even functionals is symmetric,
so the prefactor $1/2$ combines the first and last terms and the two mixed
terms.
\end{proof}

\begin{remark}[Why field-independent does not mean source-irrelevant]
A term independent of the canonical fields can have zero antibracket while still
having nonzero coefficients $F_I$ with $I\ne0$.  Such a term therefore remains
part of the normalized source/reference functional.  In addition,
\eqref{eq:v4-Delta-convolution} shows that source dependence of the BV density
itself produces genuine measure/reference contacts.  V4 never discards either
kind of term merely because its unmarked canonical derivative vanishes.
\end{remark}

\section{The canonical-pair incidence theorem}
\label{sec:v4-incidence}

The one-loop coefficient contains contributions with different physical
origins.  Fix once and for all the inherited once-counted decomposition
\begin{equation}
 \bar S_{1,j}=\sum_{\tau\in\mathcal T}\bar S_{1,j}^{\tau},
 \qquad
 \mathcal T=
 \{\mathrm{grav},\mathrm{gauge},H,\mathrm{ferm},Y,
   \mathrm{gh/nm},\mathrm{line},\mathrm{ref},\mathrm{comp}\}.
 \label{eq:v4-origin-decomposition}
\end{equation}
This is an ownership label only.  It does not split the BV differential or
permit separate normalization conditions.  The tags mean, respectively,
gravity/geometric, internal gauge, Higgs, fermion kinetic, Yukawa,
ghost/nonminimal, determinant/Pfaffian line, measure/reference, and
composite/stress/contact contributions.

For a functional $F$ define its canonical support
\begin{equation}
 \operatorname{Supp}_{\BV}(F)
 :=\{a\in\mathcal C_j:
       \delta F/\delta z^a\ne0\ \text{or}\
       \delta F/\delta z_a^*\ne0\}.
 \label{eq:v4-BV-support}
\end{equation}

\begin{lemma}[Canonical-pair support criterion]
\label{lem:v4-support-criterion}
For any two functionals $F,G$,
\begin{equation}
 \boxed{
 (F,G)
 =\sum_{a\in\operatorname{Supp}_{\BV}(F)\cap\operatorname{Supp}_{\BV}(G)}(F,G)_a.}
 \label{eq:v4-support-bracket}
\end{equation}
In particular, disjoint canonical support implies $(F,G)=0$ identically, before
projection or a cutoff limit.
\end{lemma}

\begin{proof}
This is immediate from the pairwise Darboux decomposition
\eqref{eq:v4-pair-bracket}.  A summand vanishes if either functional is
independent of both entries of that canonical pair.
\end{proof}

The following table is the E3 cross-sector audit.  It lists the canonical pair,
not a collection of separate Ward identities.  Every mixed origin in one row
can contract only through the pair displayed in that row.

\begin{center}
\small
\begin{tabular}{p{0.22\textwidth}p{0.70\textwidth}}
\toprule
Canonical pair block & Origin tags which can meet through that pair \\
\midrule
$\mathcal C_{\rm geom}$ &
geometry; gauge and Higgs metric/coframe dependence; fermion kinetic and
Yukawa density/coframe dependence; chiral line; measure/reference;
stress/composite writers; diffeomorphism antifield vertices \\
$\mathcal C_{\rm Lor}$ &
coframe/spin connection; fermion spin transport; Lorentz antifield vertices;
chiral line; geometric reference/source terms \\
$\mathcal C_{3,2,Y}$ &
internal gauge action; charged Higgs kinetic term; chiral fermion kinetic
term; gauge-covariant line determinant; gauge composite insertions; internal
BRST antifield vertices \\
$\mathcal C_H$ &
Higgs kinetic/potential action; Yukawa action; Higgs-dependent chiral line;
Higgs antifield transformation; Higgs/composite/contact writers \\
$\mathcal C_F$ &
chiral kinetic action; Yukawa action; fermion antifield BRST vertices;
fermionic composites and their contacts \\
$\mathcal C_{\rm diff-gh}$ &
all Lie-derivative antifield vertices, gravitational ghost action,
constraint/nonminimal descendants and transformed source multiplets \\
$\mathcal C_{\rm Lor-gh}$ &
frame/spin and chiral-fermion Lorentz transformations, Lorentz ghost and
nonminimal descendants \\
$\mathcal C_{\rm int-gh}$ &
internal gauge, Higgs and chiral-fermion BRST transformations, internal ghost
self-interactions and nonminimal descendants \\
$\mathcal C_{\rm nm}$ &
gauge-fixing fermion, antighost--multiplier packet, retained constraint
quartets and their marked source descendants \\
$\mathcal C_{\rm src}$ &
canonical source/antifield or BRST-doublet writers; all contacts produced by
their exact source differential \\
\bottomrule
\end{tabular}
\end{center}

The determinant/Pfaffian line, reference density and stress/composite rows do
not create mysterious extra antibracket directions.  They are functionals of
the canonical backgrounds and source multiplets, so their cross terms occur in
the appropriate rows of this table.  Their separate bookkeeping labels are
retained only to prevent their coefficients from being dropped.

\begin{theorem}[Exhaustive quadratic cross-term formula]
\label{thm:v4-quadratic-audit}
For every source multi-index $I$, the complete quadratic one-loop contribution
to the two-loop master coefficient is
\begin{equation}
 \boxed{
 \begin{aligned}
 \mathcal Q_{2,j;I}
 :=\left[\frac12(\bar S_{1,j},\bar S_{1,j})\right]_I
 =\frac12
 \sum_{I_1\sqcup I_2=I}
 &\varepsilon(I_1,I_2)\binom{I}{I_1}\\[-0.2em]
 &\times
 \sum_{a\in\mathcal C_j}
 \sum_{\tau,\upsilon\in\mathcal T}
 (\bar S_{1,j;I_1}^{\tau},
  \bar S_{1,j;I_2}^{\upsilon})_a .
 \end{aligned}}
 \label{eq:v4-full-quadratic}
\end{equation}
Every nonzero cross antibracket/contact between the E3 rows occurs exactly in
this finite sum.  There is no additional mixed master term outside
\eqref{eq:v4-full-quadratic}.
\end{theorem}

\begin{proof}
Apply Lemma~\ref{lem:v4-source-convolution}, then the origin decomposition
\eqref{eq:v4-origin-decomposition}, then the canonical-pair decomposition
\eqref{eq:v4-pair-bracket}.  These are identities, so the resulting sum is
exhaustive.  Conversely, every term in the displayed sum is an actual
contraction of one canonical field with its exact dual antifield and therefore
belongs to the master antibracket.  The factor $1/2$ removes the double count
under exchange of the two even one-loop action factors, with the inherited
Koszul convention covering odd marked coefficients.  Lemma
\ref{lem:v4-support-criterion} proves that no contraction can connect two
origin tags through a canonical pair absent from the incidence table.
\end{proof}

\begin{remark}[Examples of terms which are now automatically owned]
The geometric pair produces the stress cross terms between gravitational and
gauge/Higgs/fermion/line contributions.  The internal gauge pairs produce the
gauge--Higgs and gauge--fermion contacts.  The Higgs and fermion pairs produce
the Yukawa covariance contacts.  The diffeomorphism, Lorentz and internal ghost
pairs generate the corresponding antifield descendants.  The source-pair row
generates BRST-doublet contacts.  None of these is assigned an independent
finite constant in V4.
\end{remark}

\section{The actual common first- and second-loop coefficients}
\label{sec:v4-common-coefficients}

Let $S_{1,j}^{\rm pre}$ be the complete one-loop coefficient before the V3
one-loop restoring counterterm is installed.  Let
$S_{2,j}^{\rm pre}[\bar S_{1,j}]$ be the complete forest-prepared two-loop
coefficient in which every proper one-loop subtraction is owned by the same
already corrected $\bar S_{1,j}$.  This notation includes the inherited
canonical second-order packets, measure/reference pieces and all finite
normalization coordinates; it does not mean a new graphwise choice.

Define the un-restored local coefficients
\begin{align}
 B_{1,j}
 &:=\mathsf P^{\rm ESM}_{N,6,r}
 \bigl((S_{0,j},S_{1,j}^{\rm pre})-\Delta_jS_{0,j}\bigr),
 \label{eq:v4-B1}\\
 C_{1,j}&:=-h_jB_{1,j},
 \qquad
 \bar S_{1,j}:=S_{1,j}^{\rm pre}+C_{1,j}.
 \label{eq:v4-S1bar}
\end{align}
By V3, $d_jh_jB_{1,j}=B_{1,j}$.

The order-two breaking must be formed \emph{after} this replacement.  Set
\begin{equation}
 \boxed{
 B_{2,j}:=\mathsf P^{\rm ESM}_{N,6,r}
 \left
 \{
 (S_{0,j},S_{2,j}^{\rm pre}[\bar S_{1,j}])
 +\frac12(\bar S_{1,j},\bar S_{1,j})
 -\Delta_j\bar S_{1,j}
 \right\}.}
 \label{eq:v4-B2-recursive}
\end{equation}
Then
\begin{equation}
 C_{2,j}:=-h_jB_{2,j},
 \qquad
 \bar S_{2,j}:=S_{2,j}^{\rm pre}[\bar S_{1,j}]+C_{2,j}.
 \label{eq:v4-S2bar}
\end{equation}

\begin{lemma}[All one-loop-restorer descendants occur in $B_2$]
\label{lem:v4-C1-descendants}
If for notational comparison one writes
$S_{2,j}^{\rm pre}[\bar S_{1,j}]=S_{2,j}^{\rm pre,0}
 +\delta_{C_1}S_{2,j}^{\rm forest}$, then the complete order-two local breaking
contains
\begin{equation}
 \begin{split}
 B_{2,j}=\mathsf P^{\rm ESM}_{N,6,r}\bigl\{&
 (S_0,S_{2}^{\rm pre,0})
 +(S_0,\delta_{C_1}S_{2}^{\rm forest})
 +\tfrac12(S_1^{\rm pre},S_1^{\rm pre})-\Delta S_1^{\rm pre}\\
 &+(S_1^{\rm pre},C_1)
 +\tfrac12(C_1,C_1)-\Delta C_1
 \bigr\},
 \end{split}
 \label{eq:v4-C1-descendants}
\end{equation}
with all quantities at shell $j$.  Hence the order-two equation cannot use the
uncorrected $S_1^{\rm pre}$ without changing the problem.
\end{lemma}

\begin{proof}
Substitute \eqref{eq:v4-S1bar} into
\eqref{eq:v4-B2-recursive} and use bilinearity plus symmetry of the BV bracket
on even functionals.  The forest term
$\delta_{C_1}S_{2,j}^{\rm forest}$ is present because proper one-loop
subtractions are restrictions of the corrected common one-loop action, by the
inherited forest-ownership theorem.  No further term is available at order two.
\end{proof}

Now define the common master coefficients using the same action coefficients:
\begin{align}
 \mathcal M_{j,1}^{\rm all}
 &:=(S_{0,j},\bar S_{1,j})-\Delta_jS_{0,j},
 \label{eq:v4-M1}\\
 \mathcal M_{j,2}^{\rm all}
 &:=(S_{0,j},\bar S_{2,j})
 +\frac12(\bar S_{1,j},\bar S_{1,j})
 -\Delta_j\bar S_{1,j}.
 \label{eq:v4-M2}
\end{align}
These are exactly the coefficients requested at the end of V3.

\begin{theorem}[Common local two-loop master identity]
\label{thm:v4-common-local-master}
On the V3 protected constant-rank chart,
\begin{equation}
 \boxed{
 \mathsf P^{\rm ESM}_{N,6,r}\mathcal M_{j,1}^{\rm all}=0,
 \qquad
 \mathsf P^{\rm ESM}_{N,6,r}\mathcal M_{j,2}^{\rm all}=0.}
 \label{eq:v4-local-master-zero}
\end{equation}
The two equalities use one $C_{1,j}$ and one $C_{2,j}$ in the complete source
bank.  They are not sums of independently normalized sector equations.
\end{theorem}

\begin{proof}
For order one, use \eqref{eq:v4-S1bar}:
\[
 \mathsf P\mathcal M_{j,1}^{\rm all}
 =B_{1,j}+d_jC_{1,j}
 =B_{1,j}-d_jh_jB_{1,j}=0.
\]
The last equality is the V3 contraction identity, applicable because the
complete native anomaly projection of $B_{1,j}$ vanishes.

At order two, $B_{2,j}$ in \eqref{eq:v4-B2-recursive} is formed with the
already corrected $\bar S_{1,j}$.  V3 proves that it is a $d_j$-closed
source-extended cocycle and that its complete Standard-Model anomaly quotient
vanishes.  Therefore $d_jh_jB_{2,j}=B_{2,j}$.  Substituting
\eqref{eq:v4-S2bar} into \eqref{eq:v4-M2} gives
\[
 \mathsf P\mathcal M_{j,2}^{\rm all}
 =B_{2,j}+d_jC_{2,j}
 =B_{2,j}-d_jh_jB_{2,j}=0.
\]
Theorem~\ref{thm:v4-quadratic-audit} shows that the $B_{2,j}$ used here contains
every canonical-pair cross antibracket and every central-source contact.  Lemma
\ref{lem:v4-source-convolution} and \eqref{eq:v4-Delta-convolution} include all
source-dependent measure/reference contacts.  Thus no sector equation is being
silently omitted from either cancellation.
\end{proof}

\begin{proposition}[No hidden graphwise or rowwise restoring freedom]
\label{prop:v4-no-hidden-refit}
Let $C_{\ell,j}'$ be another local ghost-zero functional which restores the
same complete order-$\ell$ breaking with the same lower-order action.  Then
\begin{equation}
 d_j(C_{\ell,j}'-C_{\ell,j})=0.
 \label{eq:v4-closed-difference}
\end{equation}
On the finite Hodge decomposition its difference is a sum of a $d_j$-exact
redundant direction and a ghost-zero harmonic/EFT Wilson direction.  After the
common physical normalization and Wilson-coordinate convention are fixed, no
additional graphwise or rowwise finite parameter remains.
\end{proposition}

\begin{proof}
Subtract the two restoring equations.  Finite Hodge decomposition at ghost
number zero gives the exact plus harmonic splitting.  Exact directions are
canonical/redundant reparametrizations; harmonic ghost-zero directions are the
physical local Wilson coordinates already retained in the ESM state.  Fixing
the declared normalization and those coordinates fixes the ambiguity.  In
particular, there is no independent freedom indexed by Feynman graph or by E3
row.
\end{proof}

\section{The row audit is a projection of one zero vector}
\label{sec:v4-row-audit}

Choose coefficient projectors $\Pi_\rho$ which resolve the retained ghost-one
bank as a direct sum of bookkeeping rows,
\begin{equation}
 I=\sum_{\rho\in\mathcal R}\Pi_\rho,
 \qquad
 \mathcal R=
 \{\mathrm{diff},\mathrm{Lor},3,2,Y,H/Y,
   \chi\mathrm{-line},\mathrm{gh/nm},\mathrm{meas/ref},
   \mathrm{comp/stress}\}.
 \label{eq:v4-row-projectors}
\end{equation}
The projectors need not define independent subcomplexes; they are only a finite
coordinate resolution of the common coefficient.

\begin{theorem}[Simultaneous E3 row identity]
\label{thm:v4-row-identity}
For $\ell=1,2$ define the row vector
\begin{equation}
 \mathbf M_{j,\ell}
 :=\bigl(\Pi_\rho\mathsf P^{\rm ESM}_{N,6,r}
          \mathcal M_{j,\ell}^{\rm all}\bigr)_{\rho\in\mathcal R}.
 \label{eq:v4-row-vector}
\end{equation}
Then
\begin{equation}
 \boxed{\mathbf M_{j,1}=0,\qquad \mathbf M_{j,2}=0.}
 \label{eq:v4-row-vector-zero}
\end{equation}
More explicitly, the entries have the following common origin:
\begin{enumerate}[label=\textup{(M\arabic*)},leftmargin=3.2em]
\item the diffeomorphism/retained-constraint entry contains every matter,
gauge, line and source stress cross term through $\mathcal C_{\rm geom}$ and
$\mathcal C_{\rm diff-gh}$;
\item the Lorentz/spin entry contains the frame--spinor and line-frame contacts
through $\mathcal C_{\rm Lor}$ and $\mathcal C_{\rm Lor-gh}$;
\item the $SU(3)$, $SU(2)$ and $U(1)_Y$ entries contain the gauge--Higgs,
gauge--fermion and gauge-line contacts through $\mathcal C_{3,2,Y}$ and
$\mathcal C_{\rm int-gh}$;
\item the Higgs/Yukawa entry contains Higgs--fermion/Yukawa and marked Higgs
contacts through $\mathcal C_H$ and $\mathcal C_F$;
\item the chiral-line entry contains the determinant/Pfaffian frame and
line-source coefficients but uses the same background canonical pairs as the
preceding rows;
\item the ghost/nonminimal entry contains the exact contractible descendants
and retained constraint quartets;
\item the measure/reference entry contains every $\Delta_{j,K}$ contact in
\eqref{eq:v4-Delta-convolution} and every retained source-dependent reference
coefficient;
\item the composite/stress entry contains the marked coefficients and all
products/contacts forced by \eqref{eq:v4-bracket-convolution}.
\end{enumerate}
All these entries vanish because they are coordinates of the single zero vector
in \eqref{eq:v4-local-master-zero}; no entry has its own counterterm choice.
\end{theorem}

\begin{proof}
The direct-sum identity \eqref{eq:v4-row-projectors} and
Theorem~\ref{thm:v4-common-local-master} give
\eqref{eq:v4-row-vector-zero}.  The detailed ownership statements follow from
the incidence table and from the source/BV-Laplacian convolution formulas.
They are exhaustive by Theorem~\ref{thm:v4-quadratic-audit}.  Because the row
projectors are applied only after the full coefficient is assembled, a mixed
term is never omitted merely because its two endpoints carry different origin
tags.
\end{proof}

\section{Stress, composite and higher source derivatives of the same identity}
\label{sec:v4-source-descendants}

The source-complete content can be stated without introducing a new family of
Ward identities.  Let $I$ be any central-source multi-index with $|I|\le r$.
Coefficient extraction commutes with the finite local projection because both
act coefficientwise on the truncated source algebra.

\begin{theorem}[All retained source/contact descendants vanish in the same local master vector]
\label{thm:v4-source-descendants}
For $\ell=1,2$ and every $|I|\le r$,
\begin{equation}
 \boxed{
 \left[
 \mathsf P^{\rm ESM}_{N,6,r}
 \mathcal M_{j,\ell}^{\rm all}
 \right]_I=0.}
 \label{eq:v4-source-descendant-zero}
\end{equation}
In particular, stress-tensor insertions, composite insertions, operator-mixing
contacts, determinant/Pfaffian line insertions and coincident source contacts
are coefficients of this one identity.
\end{theorem}

\begin{proof}
Theorem~\ref{thm:v4-common-local-master} is an identity in the finite source
algebra, not merely at $J=0$.  Apply the linear coefficient map $[\cdot]_I$.
Lemma~\ref{lem:v4-source-convolution} shows that the coefficient of the
quadratic antibracket contains every partition of the source labels, while
\eqref{eq:v4-Delta-convolution} contains every source derivative of the finite
BV density/reference operator.  Hence coefficient extraction produces the
complete marked/contact equation and no additional correction is required.
\end{proof}

\begin{corollary}[Second marked derivative]
\label{cor:v4-second-marked}
For two even writers $A,B$, the order-two marked identity contains, among the
terms with the common one-loop action,
\begin{equation}
 (\bar S_{1,0},\bar S_{1,AB})
 +(\bar S_{1,A},\bar S_{1,B})
 -\Delta_{0}\bar S_{1,AB}
 -\Delta_A\bar S_{1,B}
 -\Delta_B\bar S_{1,A}
 -\Delta_{AB}\bar S_{1,0},
 \label{eq:v4-second-marked-explicit}
\end{equation}
together with the correspondingly marked $(S_0,\bar S_2)$ terms.  Thus the
mixed product contact $(\bar S_{1,A},\bar S_{1,B})$ and the source-dependent
measure/reference contacts are mandatory pieces of the same coefficient.
\end{corollary}

\section{The quasilocal complement of the common master coefficient}
\label{sec:v4-master-remainder}

E3 asks for the local master identity together with the already controlled
quasilocal remainder.  Put
\begin{equation}
 \mathsf Q_j:=I-\mathsf P^{\rm ESM}_{N,6,r},
 \qquad
 \mathcal R_{j,\ell}^{\rm M}
 :=\mathsf Q_j\mathcal M_{j,\ell}^{\rm all}.
 \label{eq:v4-master-remainder}
\end{equation}
Let $\varepsilon_j:=\mu/\Lambda_j$ on a physical annulus and
$\varepsilon_j:=a\mu$ on the ultraviolet cap.

\begin{lemma}[Bounded canonical contraction on the fixed rooted source bank]
\label{lem:v4-rooted-bracket}
At fixed $(N,6,r)$ there is a finite constant $C_{N,r}$ such that every
renormalized shell kernel occurring in V4 satisfies the rooted contraction
estimate
\begin{equation}
 \|(F,G)\|_*
 \le C_{N,r}\|F\|_*\|G\|_*,
 \label{eq:v4-rooted-bracket-bound}
\end{equation}
whenever the contraction is one of the canonical contractions already admitted
by the finite bank.  The corresponding renormalized complementary
BV-Laplacian contractions obey the same fixed-bank boundedness after their
local Taylor coefficient is assigned to $\mathsf P$.
\end{lemma}

\begin{proof}
There are finitely many canonical pairs, arities and source markings at fixed
$(N,6,r)$.  In the V1 rooted norm, closing one canonical field--antifield pair
is a finite sum of contractions in which one external/source kernel is placed
in the distinguished $L^1$ root and the remaining marked kernels are in the
corresponding $L^\infty$ stocks.  The weighted Schur estimate bounds each such
contraction with no volume multiplier.  Summing the finite incidence family
produces $C_{N,r}$.  For the BV Laplacian, the potentially local coincidence
piece is by definition assigned to the common Taylor/local bank; the remaining
renormalized kernel is one of the inherited shell-connection/source kernels
and obeys the same rooted estimate.  This statement does not assert an
unsubtracted trace bound for an arbitrary ultraviolet operator.
\end{proof}

\begin{theorem}[One-excess-scale master remainder through two loops]
\label{thm:v4-master-remainder}
On the declared V4 domain,
\begin{equation}
 \boxed{
 \|\mathcal R_{j,1}^{\rm M}\|_*
 \le C_{N,r}\varepsilon_j,
 \qquad
 \|\mathcal R_{j,2}^{\rm M}\|_*
 \le C_{N,r}\varepsilon_j.}
 \label{eq:v4-master-remainder-bound}
\end{equation}
The part of $\mathcal R_{j,2}^{\rm M}$ containing two complementary one-loop
remainders is $O(\varepsilon_j^2)$.
\end{theorem}

\begin{proof}
The one-loop statement is the inherited V1--V3 shell/source remainder estimate
applied to the complete first master coefficient after its local Taylor part is
removed.  At two loops, every primitive hard graph and every determinant-line
remainder has the one-excess-scale bound from the V3 mixed chiral annulus
 theorem and the V2 differentiated line theorem.  Proper one-loop subgraphs are
subtracted by the same corrected $\bar S_1$, so there is no additional
unowned local piece.

Expand the quadratic term using Theorem~\ref{thm:v4-quadratic-audit}.  A term
with one local bounded factor and one complementary $O(\varepsilon_j)$ factor
is $O(\varepsilon_j)$ by Lemma~\ref{lem:v4-rooted-bracket}; a term with two
complementary factors is $O(\varepsilon_j^2)$.  The same argument applies to
the renormalized complementary BV-Laplacian terms.  The finite Hodge restorer
acts only on the local coefficient and is uniformly bounded on the protected
constant-rank chart, so it does not worsen the complementary scale gain.
Summing the finite incidence and forest families proves
\eqref{eq:v4-master-remainder-bound}.
\end{proof}

\begin{corollary}[Summable master defect in the ultraviolet shell composition]
\label{cor:v4-master-summable}
For geometric shell scales $\Lambda_j=b^j\lambda$,
\begin{equation}
 \sum_{j\ge j_0}
 \bigl(\|\mathcal R_{j,1}^{\rm M}\|_*+
       \|\mathcal R_{j,2}^{\rm M}\|_*\bigr)
 <\infty.
 \label{eq:v4-master-sum}
\end{equation}
Hence the E3 master defect left outside the retained local bank is a summable
quasilocal shell remainder at the fixed two-loop/EFT/source order.
\end{corollary}

\begin{proof}
Use Theorem~\ref{thm:v4-master-remainder} and
$\sum_j\mu/\Lambda_j<\infty$; the ultraviolet-cap form $a\mu$ is the terminal
version of the same estimate.
\end{proof}

\section{E3 closure theorem}
\label{sec:v4-E3-close}

\begin{theorem}[Complete common two-loop BV/master/source equation]
\label{thm:v4-E3}
Fix the explicit V2 native Wilson/overlap branch and all V4 finite-order data.
Let
\begin{equation}
 S_j^{[\le2]}
 =S_{0,j}+\hbar\bar S_{1,j}+\hbar^2\bar S_{2,j}
 \label{eq:v4-common-action}
\end{equation}
with $\bar S_{1,j},\bar S_{2,j}$ defined recursively by
\eqref{eq:v4-S1bar} and \eqref{eq:v4-S2bar}.  Then the coefficient of the same
finite master functional
\begin{equation}
 \mathcal M_j(S_j^{[\le2]})
 =\frac12(S_j^{[\le2]},S_j^{[\le2]})
  -\hbar\Delta_jS_j^{[\le2]}
 \label{eq:v4-full-master-functional}
\end{equation}
obeys
\begin{equation}
 \boxed{
 \mathcal M_j(S_j^{[\le2]})
 =\hbar\mathcal R_{j,1}^{\rm M}
  +\hbar^2\mathcal R_{j,2}^{\rm M}
  +O(\hbar^3),}
 \label{eq:v4-master-final}
\end{equation}
where
\begin{equation}
 \boxed{
 \mathsf P^{\rm ESM}_{N,6,r}\mathcal R_{j,\ell}^{\rm M}=0,
 \qquad
 \|\mathcal R_{j,\ell}^{\rm M}\|_*
 \le C_{N,r}\frac{\mu}{\Lambda_j}
 \quad(\ell=1,2),}
 \label{eq:v4-E3-final-bounds}
\end{equation}
with $a\mu$ at the ultraviolet cap.

The single identity includes simultaneously the diffeomorphism,
local-Lorentz/spin, $SU(3)\times SU(2)\times U(1)_Y$, Higgs/Yukawa,
chiral determinant/Pfaffian line, ghost/nonminimal, antifield,
measure/Jacobian/reference, composite/stress and contact rows, together with all
cross terms in $\tfrac12(\bar S_1,\bar S_1)$.  No graph, row or source
derivative uses an independently selected restoring counterterm.
\end{theorem}

\begin{proof}
The coefficient expansion of \eqref{eq:v4-full-master-functional} is exactly
\eqref{eq:v4-M1}--\eqref{eq:v4-M2}.  The local part vanishes by
Theorem~\ref{thm:v4-common-local-master}.  The canonical-pair and source
convolution theorems prove that these coefficients contain all required cross
and contact terms.  The simultaneous row statement is
Theorem~\ref{thm:v4-row-identity}, and the complete source descendants are
Theorem~\ref{thm:v4-source-descendants}.  The complementary estimates are
Theorem~\ref{thm:v4-master-remainder}.  Combining local and complementary
parts yields \eqref{eq:v4-master-final}--\eqref{eq:v4-E3-final-bounds}.
\end{proof}

\begin{center}
\fbox{\parbox{0.90\textwidth}{
\textbf{E3 is closed at the stated two-loop coefficientwise scope.}
The conclusion is one common master/source identity for one common action,
not a collection of sectorwise finite Ward identities.}}
\end{center}

\section{Why the next theorem is now an induction theorem}
\label{sec:v4-E4-handoff}

There is no load-bearing reason to compute an isolated three-loop graph next.
The common two-loop coefficient reveals the general algebraic pattern.  If
\begin{equation}
 S^{[\le L]}=\sum_{q=0}^{L}\hbar^q\bar S_q,
 \label{eq:v4-L-action}
\end{equation}
then the order-$L$ master coefficient is
\begin{equation}
 \boxed{
 \mathcal M^{(L)}
 =(S_0,\bar S_L)
 +\frac12\sum_{p=1}^{L-1}(\bar S_p,\bar S_{L-p})
 -\Delta\bar S_{L-1}.}
 \label{eq:v4-general-master-coefficient}
\end{equation}
The source-contact formula of Lemma~\ref{lem:v4-source-convolution} and the
canonical-pair incidence theorem are independent of $L$.  The new work for E4
is therefore not algebraic expansion but the native all-forest estimate and the
finite local bank/range statement at an arbitrary \emph{fixed} $L$.

\begin{theorem}[Exact next load-bearing theorem for V5]
\label{thm:v4-next-E4}
For every prescribed fixed loop order $L$, EFT/derivative cap $\Delta$, source
order $r$, arity $N$ and finite external panel, construct an induction state
containing:
\begin{enumerate}[label=\textup{(I\arabic*)},leftmargin=3.2em]
\item the complete finite ghost-zero EFT/Wilson bank through the order required
by $L$;
\item the complete source/contact algebra through order $r$ and all exact BV
field--antifield/source pairs;
\item determinant/Pfaffian line/frame data on the protected component;
\item the already fixed common actions $\bar S_0,\ldots,\bar S_{L-1}$ and their
finite normalization/reference data;
\item ownership of every proper divergent subgraph by the restriction of the
appropriate previously fixed common $\bar S_q$;
\item rooted/quasilocal complementary kernels and the retained soft
constraint/nonminimal variables.
\end{enumerate}
Then prove that the forest-prepared order-$L$ coefficient decomposes as
\begin{equation}
 \text{complete local Taylor/EFT/source jet}
 +\text{quasilocal remainder},
 \label{eq:v4-E4-decomposition}
\end{equation}
with a finite local bank at that fixed order, that the first un-restored local
master coefficient obtained from
\eqref{eq:v4-general-master-coefficient} is a $d$-cocycle, that its complete
Standard-Model anomaly projection vanishes, and that one common finite Hodge
primitive restores it.  Finally prove a one-excess-scale rooted bound for the
complete order-$L$ forest remainder, with constants allowed to depend on $L$.
This yields a native shell self-map at arbitrary fixed $L$ without claiming any
uniform $L\to\infty$ estimate.
\end{theorem}

\section{V4 anti-circularity audit}
\label{sec:v4-audit}

\begin{enumerate}[leftmargin=2.4em]
\item The quadratic two-loop term is not inferred from separate sector Ward
identities.  It is expanded exactly by source partition, origin tag and
canonical pair in \eqref{eq:v4-full-quadratic}.
\item Source-dependent measure/reference terms are not suppressed.  Their
coefficients are the $\Delta_{j,K}$ terms of
\eqref{eq:v4-Delta-convolution}.
\item The order-two breaking is formed only after the one-loop restorer is in
the action.  Equation \eqref{eq:v4-C1-descendants} displays the induced
$(S_1,C_1)$, $(C_1,C_1)$, $\Delta C_1$ and forest-insertion terms explicitly.
\item No row or graph receives an independent restorer.  Proposition
\ref{prop:v4-no-hidden-refit} identifies the only remaining ambiguity as the
already-declared exact/redundant plus physical ghost-zero Wilson freedom.
\item The line/frame sector is not replaced by a scalar determinant.  It enters
through the V1--V3 protected line packet and through the same background
canonical pairs as the remaining action.
\item The quasilocal master defect is not promoted to an exact finite-cutoff
zero outside the retained bank.  It is recorded with the proved summable
$O(\mu/\Lambda_j)$ rooted bound.
\end{enumerate}

\section{V4 status ledger}
\label{sec:v4-status}

\subsection*{Proved in V4}
\begin{enumerate}[leftmargin=2.2em]
\item Central source writers can be separated from genuinely canonical source
multiplets without changing the V3 source-extended BV bank.  In that
organization, the exact source coefficient of an antibracket is the finite
super-binomial convolution \eqref{eq:v4-bracket-convolution}, and a
source-dependent BV density obeys \eqref{eq:v4-Delta-convolution}.
\item The canonical-pair support theorem gives an exhaustive finite incidence
audit of every mixed term in $\tfrac12(\bar S_1,\bar S_1)$.  Geometry--matter
stress, gauge--Higgs/fermion, Yukawa, ghost/antifield, line and source contacts
are entries of one quadratic master coefficient rather than separate Ward
problems.
\item The corrected one-loop action necessarily feeds the order-two master row
through its forest insertions, $(S_1,C_1)$, $\tfrac12(C_1,C_1)$ and
$-\Delta C_1$ descendants.  These terms are all retained in the recursive
$B_{2,j}$.
\item With the V3 common Hodge restorer, the complete local coefficients obey
one common identity
$\mathsf P\mathcal M_{j,1}^{\rm all}=
 \mathsf P\mathcal M_{j,2}^{\rm all}=0$ using the same
$S_0,\bar S_1,\bar S_2$.
\item Every E3 row is a coordinate projection of that one zero vector.  The
same is true after every retained source derivative, including stress,
composite, operator-mixing, determinant/Pfaffian and coincident contact
insertions.
\item The complementary first- and second-loop master coefficients retain the
one-excess-scale rooted bound $O(\mu/\Lambda_j)$, with the term quadratic in two
one-loop remainders $O((\mu/\Lambda_j)^2)$.  The ultraviolet shell sum is
finite at this fixed order.
\item Therefore E3 is closed on the explicit V2 branch at the stated
coefficientwise two-loop/EFT/source scope.
\end{enumerate}

\subsection*{Inherited}
\begin{enumerate}[leftmargin=2.2em]
\item The V80 bosonic Einstein--gauge--Higgs two-loop self-map, weight-six
finite BV range theorem, common one-loop forest ownership, retained soft
constraint architecture and rooted two-loop remainder estimate.
\item The three-generation Standard-Model representation, cancellation of the
five ordinary local anomaly trace sums, even $SU(2)$ doublet parity, finite
gauge descent of the complete nonzero chiral section and the exact
Schur/Pfaffian scale cocycles.
\item V1's overlap functional calculus, local Kato frame, hard pivot and rooted
contraction bounds.
\item V2's explicit native Wilson completion, all-zone gap/no-doubler result,
physical annular/full-zone comparison, fixed source-jet calculus and
differentiated line remainder.
\item V3's source-extended ESM complex, native anomaly-quotient zero, common
finite Hodge restorer, mixed chiral two-loop annulus theorem and complete E2
shell self-map.
\end{enumerate}

\subsection*{Conditional}
\begin{enumerate}[leftmargin=2.2em]
\item The theorem remains native to the explicit V2 Wilson/overlap completion;
a different microscopic fermion kernel requires the regulator-comparison gate
already isolated in V2.
\item The Hodge bound uses the compact constant-rank component of the finite BV
bank.  A rank crossing requires a new chart or a separate analysis.
\item Determinant/Pfaffian statements remain on the protected nonzero component
with the retained positive IR/reference prescription.  No global line holonomy
or zero crossing is claimed.
\item The complementary master estimate uses the inherited fixed-order native
hard-symbol, forest and rooted-source bounds.  Constants may depend on loop
order, EFT degree, source order, arity and the protected chart.
\end{enumerate}

\subsection*{Open}
\begin{enumerate}[leftmargin=2.2em]
\item E4: prove the arbitrary-fixed-loop native forest induction rather than
compute isolated higher-loop examples.  The decisive new analytic task is an
all-forest one-excess-scale theorem at each prescribed fixed $L$, with one
finite local Taylor/EFT bank and common lower-order counterterm ownership.
\item E5: after E4, compose arbitrarily many ultraviolet shells at fixed
$(L,\Delta,r)$ and prove Cauchy/unique cutoff limits for normalized connected
source coefficients, including operator mixing, stress/contact and line
insertions.
\item E6: prove finite-scheme/reference compatibility and order compatibility
between the $L$ and $L+1$ constructions.
\end{enumerate}

\subsection*{Exact next load-bearing theorem}
\begin{center}
\fbox{\parbox{0.92\textwidth}{
\textbf{E4: arbitrary fixed-loop native shell induction.}
For a prescribed fixed $L$, assemble the complete forest using the same
previously fixed actions $\bar S_0,\ldots,\bar S_{L-1}$, prove a finite local
Taylor/EFT/source decomposition plus a one-excess-scale rooted remainder, form
the common order-$L$ coefficient
\(
(S_0,\bar S_L)+\frac12\sum_{p=1}^{L-1}
(\bar S_p,\bar S_{L-p})-\Delta\bar S_{L-1}
\), prove its anomaly-free local cocycle lies in the common BV range, and
install one source-complete restoring primitive.  Constants may depend on $L$;
no uniform perturbative convergence claim is allowed.}}
\end{center}

\section*{Append-only continuation marker: V4}
\addcontentsline{toc}{section}{Append-only continuation marker: V4}
\textbf{End of V4.}  The complete V1--V3 body above is retained verbatim.  V4
closes E3 by writing and auditing one common first- and second-loop master
coefficient, including every canonical-pair cross term, every retained source
contact and every source-dependent measure/reference term.  Future versions
must append after this marker.  The next file is to be named
\texttt{Renewal\_Perturbative\_ESM\_Continuum\_Research\_Notes\_V5.tex}.

% =============================================================================
% V5 -- arbitrary fixed-loop native forest induction
% =============================================================================
\part*{V5: Arbitrary fixed-loop native forest induction and contraction-complete BV jet packets}
\addcontentsline{toc}{part}{V5: Arbitrary fixed-loop native forest induction and contraction-complete BV jet packets}
\label{sec:v5-start}

\section{Upstream correction: the induction state must remember future BV contractions}
\label{sec:v5-upstream}

V4 closed the common local master equation through two loops and retained a
quasilocal complementary master kernel of size $O(\mu/\Lambda_j)$.  That is
sufficient for the two-loop statement, but it is not by itself an induction
state.  At loop order three and above a BV Laplacian or a canonical
field--antifield contraction can act on a lower-order quasilocal kernel and
produce a new lower-arity local contact.  Therefore the implication
\[
 \text{``ordinary Taylor remainder''}
 \Longrightarrow
 \text{``irrelevant for every later local master coefficient''}
\]
is false without additional bookkeeping.

This is the first load-bearing issue in E4.  The correct upstream variable is
not a bare remainder $R$ but a \emph{contraction-complete jet packet}: the
remainder together with the local Taylor jets of every canonical/source
contraction which can still occur before the prescribed final loop order.  At
fixed loop order, derivative/EFT order, source order and arity there are only
finitely many such descendants.  The packet is therefore finite, while no
claim of uniformity as $L\to\infty$ is required.

Throughout V5 fix:
\begin{equation}
 L\in\mathbb N,\qquad
 \Delta<\infty,\qquad r<\infty,\qquad N<\infty,
 \qquad |p_{\rm ext}|\le\mu<c_0\Lambda_j,
 \label{eq:v5-fixed-data}
\end{equation}
and write
\begin{equation}
 \epsilon_j:=\frac{\mu}{\Lambda_j}
 \label{eq:v5-epsilon}
\end{equation}
on a physical annulus, with the terminal ultraviolet-cap replacement
$\epsilon_j=a\mu$ exactly as in V2--V4.  Constants in V5 may depend on
$(L,\Delta,r,N)$, on the protected chart and on the finite representation
packet.  They are not required to remain bounded with $L$.

\begin{firewall}[No naive exact-cocycle induction from V4 alone]
\label{fw:v5-no-naive-cocycle}
V3's exact statement that the first un-restored master coefficient is a
classical BV cocycle assumes that every lower master coefficient vanishes in
the algebra in which the quantum consistency identity is expanded.  V4 proves
an exact vanishing of the retained local projection but keeps a nonzero
quasilocal complement.  Hence one may not simply cite the V3 cocycle theorem
at arbitrary loop order while discarding the V4 complement.  V5 repairs this
by enlarging the induction state so that all local descendants of that
complement are retained before the next projection is formed.
\end{firewall}

\section{The fixed-order induction state}
\label{sec:v5-state}

Let $\widehat{\mathfrak L}^{\rm ESM}_{N,\Delta,r}$ be the complete finite local
ESM BV/source bank at arity at most $N$, EFT/derivative filtration at most
$\Delta$ and source order at most $r$.  It contains the same canonical
field--antifield pairs, determinant/Pfaffian line coordinates,
measure/reference coordinates, stress/composite writers and nonminimal pairs
as V3--V4, enlarged only as required by the chosen finite EFT filtration.
The word ``complete'' means that every local descendant under the classical
BV differential and every retained source/contact derivative is included.

The gravity sector is treated as an EFT.  Thus $\Delta$ is a filtration label,
not a statement that a finite list of gravitational couplings renormalizes all
loop orders.  When an order-$L$ graph requires a new local operator at the
chosen accuracy, that coordinate is added to the finite order-$L$ bank.  No
operator needed only at unbounded loop or derivative order is introduced.

\begin{definition}[Fixed-$L$ induction datum]
\label{def:v5-induction-datum}
For $0\le q<L$, the order-$q$ datum consists of
\begin{enumerate}[label=\textup{(D\arabic*)},leftmargin=3.2em]
\item the common normalized local action coefficient $\bar S_{q,j}$, including
its physical EFT/Wilson coordinates and every source/antifield descendant;
\item the determinant/Pfaffian line/frame packet on the protected nonzero
component;
\item the normalization and reference coordinates which define the common
finite scheme;
\item every rooted/quasilocal action kernel generated at order $q$;
\item the contraction-complete jet packet of each such quasilocal kernel up to
remaining contraction depth $L-q$;
\item the ancestry/ownership tag identifying the lower-order common action
which owns each proper local subgraph insertion.
\end{enumerate}
The collection of these data for $0\le q<L$ is denoted
$\mathfrak I_{<L}(\Delta,r,N)$.
\end{definition}

The fifth entry is new.  It is the device which makes the induction exact at
the level of every local coefficient that can be generated before loop order
$L$.

\section{Reference absorption and finiteness of the graph family}
\label{sec:v5-graph-finite}

At fixed loop order one must first eliminate a trivial source of infinitude:
arbitrarily many quadratic mass/kinetic insertions on one propagator.  The
native shell formalism already has a common normalized Gaussian/reference
family, so these insertions should not be counted as independent interaction
vertices.

\begin{lemma}[Quadratic/reference absorption]
\label{lem:v5-quadratic-absorption}
At every completed order $q<L$, all source-free local hard vertices of hard
valence $0,1,2$ may be assigned uniquely to the normalized background and
Gaussian/reference coordinates of the common state.  After this assignment,
every source-free perturbative hard interaction vertex has hard valence at
least three.  Source-marked vertices of hard valence at most two remain
explicit, but their number in a coefficient of source order at most $r$ is at
most $r$.
\end{lemma}

\begin{proof}
The vacuum coordinate is part of the reference normalization, the one-point
coordinate is fixed by the chosen background/tadpole condition, and the full
local two-point coordinate is part of the common quadratic block whose inverse
defines the native hard covariance.  Moving any of these coordinates between
``propagator'' and ``interaction'' without changing the normalization would
only be a redundant re-expansion of the same finite Gaussian identity.  We
choose the nonredundant convention in which they are absorbed once.  A marked
source insertion is not absorbed because it is a requested observable; every
such insertion consumes at least one of the at most $r$ source labels.  The
claim follows.
\end{proof}

\begin{proposition}[Finite graph family at fixed $(L,\Delta,r,N)$]
\label{prop:v5-finite-graphs}
For a fixed external arity at most $N$, source order at most $r$, loop number at
most $L$, and finite EFT filtration, only finitely many native hard graph
topologies, field/representation assignments, source decorations and
counterterm ancestry assignments can occur in one shell step.
\end{proposition}

\begin{proof}
Ignore first the at most $r$ marked low-valence source vertices.  If $V$ is the
number of remaining hard interaction vertices, $I$ the number of internal hard
lines and $E\le N+O(r)$ the number of external hard half-edges, Lemma
\ref{lem:v5-quadratic-absorption} gives
\[
 3V\le 2I+E.
\]
For each connected component contributing to an effective vertex,
$L_\Gamma=I-V+1$.  Hence
\[
 V\le 2L_\Gamma+E-2.
\]
Thus $V$ and $I$ are bounded in terms of the fixed data.  The local bank has
only finitely many field species, representation tensors, derivative monomials
and source labels at fixed $(\Delta,r,N)$.  A graph with bounded numbers of
vertices and half-edges therefore admits only finitely many assignments from
that bank.  Finally every proper-subgraph ancestry label lies in the finite set
$\{0,\ldots,L-1\}$.  This proves finiteness.
\end{proof}

\begin{lemma}[A proper 1PI ultraviolet subgraph has smaller loop number]
\label{lem:v5-loop-drop}
Let $\Gamma$ be a connected 1PI hard graph of loop number $L_\Gamma$ and let
$\gamma\subsetneq\Gamma$ be a proper connected 1PI subgraph which is admissible
as a local ultraviolet subtraction block.  Then
\begin{equation}
 \boxed{L_\gamma<L_\Gamma.}
 \label{eq:v5-loop-drop}
\end{equation}
\end{lemma}

\begin{proof}
Contract $\gamma$ to a single vertex.  The loop-number identity for connected
graphs gives
\[
 L_\Gamma=L_\gamma+L_{\Gamma/\gamma}.
\]
If equality $L_\gamma=L_\Gamma$ held, the quotient would be a connected tree.
Because $\gamma$ is proper, the quotient would contain at least one edge.  Any
edge of a tree is a bridge, and its preimage would be an internal bridge of
$\Gamma$, contradicting 1PI.  Hence $L_{\Gamma/\gamma}\ge1$ and
$L_\gamma<L_\Gamma$.
\end{proof}

This elementary lemma is what makes common counterterm ownership inductive:
no proper 1PI subtraction at loop $L$ is allowed to invent an order-$L$
finite normalization of its own.

\section{Decorated native forests and common lower-order ownership}
\label{sec:v5-forests}

\begin{definition}[Admissible decorated forest]
\label{def:v5-forest}
For a fixed native graph $\Gamma$, an admissible forest $F$ is a set of proper
connected 1PI local subtraction blocks such that any two elements are either
disjoint or nested.  Each $\gamma\in F$ carries:
\begin{enumerate}[label=\textup{(F\arabic*)},leftmargin=3.2em]
\item its field/representation and orientation data;
\item its source/contact multi-index;
\item its determinant/Pfaffian line/frame tag if present;
\item its loop number $L_\gamma$;
\item the unique local insertion
$C_\gamma=C_\gamma[\bar S_{L_\gamma,j}]$ obtained by restricting the already
fixed common action $\bar S_{L_\gamma,j}$ to the same local tensor/source
coordinate.
\end{enumerate}
Overlapping but nonnested subgraphs never occur in the same forest; they occur
in distinct forest terms, as in the V61/V80 two-loop overlapping-channel
analysis.
\end{definition}

Let $\Gamma/F$ denote simultaneous contraction of the mutually compatible
blocks in $F$.  The scheme-specific forest preparation is
\begin{equation}
 \mathbf R^{<L}_j\mathcal I_\Gamma
 :=\sum_{F\in\mathfrak F(\Gamma)}
 \mathcal I_{\Gamma/F,j}\star
 \prod_{\gamma\in F} C_\gamma,
 \label{eq:v5-forest-preparation}
\end{equation}
where the empty forest contributes the unprepared native amplitude and
$\star$ means insertion with the original routing, symmetry factor, source
labels and line/frame orientation.  Equation \eqref{eq:v5-forest-preparation}
is a finite algebraic identity at fixed cutoff; no continuum BPHZ theorem is
being imported.

\begin{theorem}[Common ownership of every proper subtraction]
\label{thm:v5-common-ownership}
For every order-$L$ graph in Proposition~\ref{prop:v5-finite-graphs}, every
proper local ultraviolet subtraction in
\eqref{eq:v5-forest-preparation} is owned by one of the previously fixed common
actions $\bar S_{q,j}$ with $q<L$.  In particular:
\begin{enumerate}[label=\textup{(O\arabic*)},leftmargin=3.2em]
\item no proper subgraph is independently refit graph by graph;
\item nested counterterms use the restriction of the already corrected lower
order action, including its earlier restoring terms;
\item all source/contact and line descendants of a counterterm are the
descendants of that same common coordinate;
\item changing the finite part of one $C_\gamma$ is a change of the declared
common renormalization scheme, not a harmless graphwise convention.
\end{enumerate}
\end{theorem}

\begin{proof}
By Lemma~\ref{lem:v5-loop-drop}, every proper 1PI subtraction block has loop
number strictly below $L$.  The induction datum therefore contains its common
local coefficient.  The source and line descendants are included in that
datum before restriction.  Forest compatibility ensures that nested
insertions are evaluated recursively from the innermost block outward, so a
counterterm already containing a lower restoring primitive is never replaced
by its uncorrected predecessor.  The uniqueness assertion is exactly the
fixed normalization/reference convention of the induction datum.
\end{proof}

\section{Contraction-complete jet packets}
\label{sec:v5-jet-packets}

We now formalize the upstream state enlargement.

Let $\mathcal C$ be the finite alphabet of elementary operations which can
appear in a later master/source coefficient at the fixed data
\eqref{eq:v5-fixed-data}.  Its letters are:
\begin{enumerate}[label=\textup{(C\arabic*)},leftmargin=3.2em]
\item one canonical field--antifield contraction appearing in an antibracket;
\item one self-contraction appearing in the finite BV Laplacian;
\item one retained source derivative or contact identification;
\item one determinant/Pfaffian line derivative/contact contraction.
\end{enumerate}
Each operation carries the precise signs and source labels of the V4
canonical-pair/source convolution formulas.  An admissible word
$\alpha=c_m\cdots c_1$ is one whose successive arities and ghost numbers are
defined.  Write $D_\alpha$ for the corresponding operator.

\begin{definition}[Depth-$m$ contraction-complete packet]
\label{def:v5-packet}
For a rooted/quasilocal kernel $K$, define
\begin{equation}
 \boxed{
 \mathbf J_j^{[m]}K
 :=\left(
   K,\ \bigl\{
   \mathsf T_{\Delta_\alpha,r_\alpha}
   D_\alpha K:
   0\le |\alpha|\le m
   \bigr\}
 \right).}
 \label{eq:v5-packet}
\end{equation}
Here $\mathsf T_{\Delta_\alpha,r_\alpha}$ is the complete local
Taylor/EFT/source projector appropriate to the descendant after the operations
in $\alpha$.  The empty word records the current local jet of $K$.  The
unprojected kernel $K$ is retained as the first component; the packet is not a
replacement of the quasilocal action by its Taylor series.
\end{definition}

\begin{lemma}[Finiteness of the packet]
\label{lem:v5-packet-finite}
For fixed $(m,\Delta,r,N)$, the packet
$\mathbf J_j^{[m]}K$ has finitely many local coordinates.
\end{lemma}

\begin{proof}
The alphabet $\mathcal C$ is finite because the field/antifield set, source
panel and line packet are finite.  There are finitely many words of length at
most $m$.  Every contraction lowers canonical arity or consumes a source/contact
slot, so inadmissible words terminate even earlier.  Each descendant local
projector has finite arity, source order and derivative/EFT degree; hence its
range is finite dimensional.  The finite union of these ranges is finite.
\end{proof}

\begin{lemma}[Shift identity]
\label{lem:v5-packet-shift}
Let $c\in\mathcal C$ be one admissible elementary operation and $m\ge1$.
Then every local coordinate of
$\mathbf J_j^{[m-1]}(D_cK)$ occurs verbatim as a coordinate of
$\mathbf J_j^{[m]}K$:
\begin{equation}
 \boxed{
 \mathsf T D_\alpha D_cK
 =\mathsf T D_{\alpha c}K,
 \qquad |\alpha|\le m-1.}
 \label{eq:v5-packet-shift}
\end{equation}
The equality includes the V4 graded source-partition signs and the declared
measure/reference contact convention.
\end{lemma}

\begin{proof}
This is associativity of the corresponding finite canonical/source
contractions, with the same fixed ordering convention used to define the V4
source convolution.  No limit is involved.  The descendant projector on the
left is exactly the projector attached to the composite word $\alpha c$ on the
right.  Because source-dependent reference terms are part of the finite BV
operator before differentiation, their contact descendants are present on
both sides.
\end{proof}

\begin{theorem}[Future-local-feedback completeness]
\label{thm:v5-feedback-complete}
Suppose an order-$q$ quasilocal kernel is stored with packet depth $L-q$.
Then every local term which that kernel can generate by BV, source/contact or
line contractions in any coefficient of loop order at most $L$ is already a
coordinate of its stored packet.  No later local master coefficient can receive
an unrecorded contribution from that kernel.
\end{theorem}

\begin{proof}
A contribution at loop order at most $L$ can contain at most $L-q$ additional
loop-generating self-contractions along an ancestry descending from an
order-$q$ kernel; antibracket and source/contact operations do not increase the
required depth beyond the number of remaining coefficient-building steps.
Read those operations from earliest to latest as a word $\alpha$ of length at
most $L-q$.  Its local descendant is one of the coordinates in
\eqref{eq:v5-packet}.  Lemma~\ref{lem:v5-packet-shift} shows inductively that
a later operation merely moves to a shallower coordinate subpacket.  Thus no
local feedback is lost.
\end{proof}

\begin{remark}[Why this is stronger than an ordinary Taylor remainder]
A kernel may have zero current Taylor jet and nevertheless acquire a nonzero
local descendant after a BV self-contraction.  Definition
\ref{def:v5-packet} records that descendant in advance.  This is exactly the
higher-loop generalization of V4's rule that one must assemble all contact and
measure/reference terms before projecting the master coefficient.
\end{remark}

\section{Rooted one-excess-scale calculus at fixed loop order}
\label{sec:v5-rooted-calculus}

The packet solves the bookkeeping problem.  We now prove the analytic estimate
which replaces an imported all-order continuum forest theorem.

For a kernel $K$ with current local jet $\mathsf T K$, write
\begin{equation}
 K=K^{\loc}+K^{\remn},
 \qquad K^{\loc}:=\mathsf T K,
 \qquad K^{\remn}:=(I-\mathsf T)K.
 \label{eq:v5-local-rem}
\end{equation}
Use the V1--V4 rooted norm $\|\cdot\|_*$ for quasilocal kernels and any fixed
coefficient norm on the finite local bank.

\begin{lemma}[Primitive annular Taylor gain at arbitrary fixed graph size]
\label{lem:v5-primitive-gain}
Let $\Gamma$ be one of the finitely many native hard graphs in Proposition
\ref{prop:v5-finite-graphs}.  Suppose all proper local subgraphs have already
been replaced by their owned lower-order insertions.  Let
$\omega_\Gamma$ be the finite Taylor degree required by the chosen EFT
filtration for the overall local part.  Then
\begin{equation}
 \boxed{
 \bigl\|(I-\mathsf T_{\omega_\Gamma})
 \mathbf R_j^{<L}\mathcal I_\Gamma\bigr\|_*
 \le C_{\Gamma,\Delta,r,N}\,\epsilon_j.}
 \label{eq:v5-primitive-gain}
\end{equation}
The same estimate holds after every admissible fixed source, geometric, gauge,
Higgs/Yukawa, antifield or line derivative needed by the packet.
\end{lemma}

\begin{proof}
After the proper local blocks have been replaced, every integrated native
propagator is supported on the hard annulus at scale $\Lambda_j$ and obeys the
fixed differentiated symbol/rooted bounds proved for the bosonic packets in
V80 and for the chiral packet in V1--V3.  The graph contains finitely many
lines and vertices by Proposition~\ref{prop:v5-finite-graphs}.  Taylor's
formula with integral remainder in the finite external momentum/source window
gives at least one explicit factor
$|p_{\rm ext}|/\Lambda_j\le\epsilon_j$ after subtraction through the complete
required local degree.  Every further fixed external derivative differentiates
a finite product of native symbols and is controlled by the same annular
bounds.  Fourier integration by parts an arbitrary fixed number of times gives
the weighted rooted $L^1$ moment required by $\|\cdot\|_*$.  The number of
Leibniz terms is finite and absorbed into $C_{\Gamma,\Delta,r,N}$.  At the
terminal full-zone cap the V2 finite-to-ordinary estimate gives the replacement
$a\mu$.
\end{proof}

\begin{lemma}[One-excess rooted ideal]
\label{lem:v5-rooted-ideal}
At fixed $(L,\Delta,r,N)$ the class of kernels satisfying
$\|R\|_*\le C\epsilon_j$ is stable, with a changed fixed constant, under every
allowed rooted contraction with bounded local kernels.  More generally, if a
contracted expression contains $s$ factors $R_1,\ldots,R_s$ with
$\|R_i\|_*\le C_i\epsilon_j$, then its complementary rooted part is
$O(\epsilon_j^s)$ after all newly created local jets are extracted into the
contraction-complete packet.
\end{lemma}

\begin{proof}
V4's canonical contraction estimate bounds one allowed contraction by a fixed
constant times the product of the rooted norms.  The same weighted Schur/root
argument applies to the additional finite arities present at fixed $L$.
Whenever a contraction creates a local Taylor/contact piece, Definition
\ref{def:v5-packet} assigns that piece to the appropriate local descendant
coordinate before the complement is estimated.  The remaining kernel is
therefore bounded by the product of the complementary norms.  Iterating the
finite number of contractions gives the claim.
\end{proof}

\begin{theorem}[Native fixed-$L$ forest decomposition]
\label{thm:v5-forest-decomposition}
For every graph/source coefficient in the fixed data
\eqref{eq:v5-fixed-data}, the scheme-specific forest preparation
\eqref{eq:v5-forest-preparation} has a decomposition
\begin{equation}
 \boxed{
 \mathbf R_j^{<L}\mathcal I_\Gamma
 =\mathsf L_{\Gamma,j}^{(L)}
  +\mathsf Q_{\Gamma,j}^{(L)},}
 \label{eq:v5-forest-decomp}
\end{equation}
where:
\begin{enumerate}[label=\textup{(A\arabic*)},leftmargin=3.2em]
\item $\mathsf L_{\Gamma,j}^{(L)}$ lies in a finite local
Taylor/EFT/source/contact bank determined by $(L,\Delta,r,N)$;
\item every proper subtraction entering $\mathsf L_{\Gamma,j}^{(L)}$ is owned
by the previously fixed common action of Theorem
\ref{thm:v5-common-ownership};
\item the full contraction-complete packet of
$\mathsf Q_{\Gamma,j}^{(L)}$ is retained in the induction state;
\item its current complementary kernel obeys
\begin{equation}
 \boxed{
 \|\mathsf Q_{\Gamma,j}^{(L)}\|_*
 \le C_{L,\Delta,r,N}\epsilon_j;}
 \label{eq:v5-forest-Q-bound}
\end{equation}
\item every term containing $s$ already complementary lower-order ancestors is
$O(\epsilon_j^s)$ after its newly generated local descendants are extracted.
\end{enumerate}
The theorem includes nested and overlapping proper subgraphs and all retained
source/contact markings.
\end{theorem}

\begin{proof}
The forest family is finite by Proposition~\ref{prop:v5-finite-graphs}.
Proceed by induction on the number of proper subtraction blocks.  For no proper
block, Lemma~\ref{lem:v5-primitive-gain} gives the result.  Suppose it holds for
all graphs with fewer proper blocks.  In each forest term choose its maximal
proper elements.  By Lemma~\ref{lem:v5-loop-drop} their loop orders are lower,
so their local insertions and quasilocal packets belong to the induction datum.
Replace each maximal block by its local coordinate plus its retained
complement.  The term in which all maximal blocks are local reduces to a
primitive quotient graph and is controlled by Lemma
\ref{lem:v5-primitive-gain}.  Every other term contains at least one
complementary factor and is $O(\epsilon_j)$ by Lemma
\ref{lem:v5-rooted-ideal}; terms with $s$ such factors have the stated higher
power.

Overlapping blocks do not require simultaneous subtraction: incompatible
blocks occur in distinct forest terms, exactly as dictated by Definition
\ref{def:v5-forest}.  Nested blocks are handled from inside outward, so their
source/contact labels and finite parts are inherited rather than refit.
Differentiating the finite forest sum distributes the source labels over graph,
subgraphs and counterterm insertions.  The resulting super-partitions are
precisely the V4 contact convolution, and every local descendant is stored by
Definition~\ref{def:v5-packet}.  Summing the finite forest family proves the
uniform fixed-$L$ bound.
\end{proof}

\begin{corollary}[No imported all-order forest theorem is needed]
\label{cor:v5-no-imported-forest}
At every prescribed fixed $L$, the native regulator itself supplies the
complete local-plus-quasilocal forest decomposition needed for the shell
update.  The argument uses the already proved native annular propagator/source
bounds, finite graph combinatorics, common lower-order ownership and the rooted
ideal estimate.  It does not assert that the constants are uniform in $L$ and
does not invoke a continuum forest theorem whose hypotheses have not been
matched to the native regulator.
\end{corollary}

\section{Finite closure of the local EFT/source bank}
\label{sec:v5-local-bank}

\begin{theorem}[Finite order-$L$ local closure bank]
\label{thm:v5-finite-local-bank}
For fixed $(L,\Delta,r,N)$ there is a finite-dimensional local bank
\begin{equation}
 \widehat{\mathfrak L}^{\rm ESM}_{L;N,\Delta,r}
 \label{eq:v5-L-bank}
\end{equation}
which contains:
\begin{enumerate}[label=\textup{(B\arabic*)},leftmargin=3.2em]
\item every local coefficient in Theorem~\ref{thm:v5-forest-decomposition};
\item all BV field--antifield descendants needed by the order-$L$ master
coefficient;
\item all source/contact and stress/composite descendants through order $r$;
\item all determinant/Pfaffian line descendants on the protected component;
\item all local descendant jets in the depth-$L$ contraction-complete packets;
\item the independent ghost-zero EFT/Wilson quotient at the chosen order.
\end{enumerate}
No finite-dimensional bank independent of $L$ or $\Delta$ is asserted.
\end{theorem}

\begin{proof}
The graph family is finite by Proposition~\ref{prop:v5-finite-graphs}.  Every
graph has finitely many forest terms and finitely many local Taylor
coefficients through its prescribed degree.  Lemma~\ref{lem:v5-packet-finite}
adds only finitely many contraction descendants.  The field and representation
packets are finite, while arity and source order are capped.  Taking the span
of this finite union gives \eqref{eq:v5-L-bank}.  The independent ghost-zero
cohomology is retained rather than quotienting it away, exactly as in V78--V80.
\end{proof}

\section{The local BV complex at arbitrary finite filtration}
\label{sec:v5-complex}

Let
\begin{equation}
 d_{j,L}:=
 \mathsf P_{L;N,\Delta,r}
 (S_{0,j},\,\cdot\,)
 \mathsf P_{L;N,\Delta,r}
 \label{eq:v5-differential}
\end{equation}
be the classical BV differential on the bank of Theorem
\ref{thm:v5-finite-local-bank}.

\begin{theorem}[Finite filtered ESM complex and ghost-one quotient]
\label{thm:v5-complex}
For every fixed $(L,\Delta,r,N)$,
\begin{equation}
 \boxed{d_{j,L}^2=0.}
 \label{eq:v5-d2}
\end{equation}
The nonminimal pairs remain contractible.  On the charged-Higgs branch, after
removal of the exceptional free-abelian families as in V80, the nontrivial
four-dimensional ghost-number-one local quotient relevant to the complete ESM
packet is the same standard gauge-anomaly quotient
\begin{equation}
 \mathfrak A_{\rm std}
 =\Span\{\mathcal A_{333},\mathcal A_{Y33},\mathcal A_{Y22},
              \mathcal A_{YYY},\mathcal A_{YRR}\}.
 \label{eq:v5-anomaly-space}
\end{equation}
Restriction to the finite order-$L$ bank creates no additional harmonic class.
\end{theorem}

\begin{proof}
Before projection, the complete classical ESM BV action satisfies the
classical master equation on the declared protected branch, so its Hamiltonian
vector field squares to zero.  The bank in Theorem
\ref{thm:v5-finite-local-bank} was defined to contain every retained descendant
under that field.  Hence its projector commutes with the differential and
\eqref{eq:v5-d2} follows.

For ghost number one, the ordinary local cohomology input is the same
matter-inclusive Einstein--Yang--Mills classification already matched in V79,
V80 and V3.  Hypercharge is not a free abelian generator because it acts on the
charged Higgs, so the additional free-abelian current/antifield families are
absent.  The exceptional gravitational/topological anomaly dimensions do not
include four.  Adding finitely many higher-derivative local representatives,
source writers and contractible BV pairs changes the finite representative
space but not the underlying four-dimensional anomaly quotient.  The finite
bank is a differential-stable restriction containing all descendants at the
chosen filtration, so the finite-jet extension argument of V80 identifies its
harmonic ghost-one quotient with the restriction of
\eqref{eq:v5-anomaly-space}.  No statement about global gauge loops or line
zero crossings is used.
\end{proof}

\begin{theorem}[Exact vanishing of the native internal-gauge harmonic projection at every fixed loop order]
\label{thm:v5-anomaly-zero}
For the complete three-generation Standard-Model shell on the explicit V2
Wilson/overlap branch, every coefficient produced by the finite forest
preparation and its contraction-complete descendants has zero projection onto
\eqref{eq:v5-anomaly-space}:
\begin{equation}
 \boxed{
 \mathsf A_{\rm std}
 \mathcal B_{j,L}^{\rm loc}=0.}
 \label{eq:v5-anomaly-zero}
\end{equation}
\end{theorem}

\begin{proof}
The complete chiral determinant section is exactly gauge invariant on the
protected finite gauge quotient by the inherited gauge-descent theorem.  The
bosonic, ghost and Higgs hard blocks transform by finite conjugation, and the
reference/source derivatives are taken inside that same covariant family.
Therefore every finite shell amplitude and every source derivative is gauge
equivariant before Taylor localization.  Forest contraction preserves this
property because each lower-order insertion is a restriction of a common
gauge-restored action, not an independently chosen graphwise counterterm.
The local projector and the packet descendant operations are representation
intertwiners.  Hence the complete native local breaking has zero coefficient
along every nontrivial internal-gauge descent class.  Equivalently, one may
decompose the one-fermion-loop seed into the five representation traces and use
the inherited Standard-Model trace cancellations; higher forest insertions
multiply gauge-covariant invariant tensors and cannot recreate a nonzero
quotient coefficient once the full finite section is invariant.  This proves
\eqref{eq:v5-anomaly-zero} without an Adler--Bardeen normalization theorem for
an individual uncancelled species.
\end{proof}

\section{Exact projected consistency from the contraction-complete packet}
\label{sec:v5-consistency}

Write the order-$q$ common coefficient as
\begin{equation}
 \widehat S_{q,j}
 =\bar S_{q,j}^{\loc}+K_{q,j}^{\remn},
 \label{eq:v5-full-coefficient}
\end{equation}
where $K_{q,j}^{\remn}$ is not discarded: its packet of the depth prescribed in
Definition~\ref{def:v5-induction-datum} is part of the state.  Let
\begin{equation}
 \widehat S_j^{[<L]}
 :=\sum_{q=0}^{L-1}\hbar^q\widehat S_{q,j}.
 \label{eq:v5-action-less-L}
\end{equation}

For any formal action $S=\sum_q\hbar^qS_q$, the finite quantum consistency
identity gives, at coefficient $n\ge1$,
\begin{equation}
 (S_0,\mathcal M_n)
 =\Delta_j\mathcal M_{n-1}
 -\sum_{p=1}^{n-1}(S_p,\mathcal M_{n-p}),
 \label{eq:v5-consistency-coefficient}
\end{equation}
with $\mathcal M_0=\frac12(S_0,S_0)=0$.  The point of the packet is that every
local descendant on the right of \eqref{eq:v5-consistency-coefficient} is
already a stored coordinate rather than an omitted remainder contact.

\begin{proposition}[Packet-complete lower-order restoration]
\label{prop:v5-packet-restoration}
Assume inductively that for every $q<L$ all local master coordinates through
loop order $q$ and all of their future descendants of depth $L-q$ vanish after
the common lower-order restorers have been installed.  Then the local
projection of the right-hand side of
\eqref{eq:v5-consistency-coefficient} at $n=L$ vanishes exactly.
\end{proposition}

\begin{proof}
Each term on the right is obtained from a lower master coefficient by one
letter of the contraction alphabet: a BV self-contraction for $\Delta_j$ or a
canonical contraction with a coefficient of the action for the antibracket.
By Theorem~\ref{thm:v5-feedback-complete}, its local descendant is one of the
stored packet coordinates of that lower master coefficient.  The induction
hypothesis sets precisely those coordinates to zero.  Therefore the full local
projection of the right-hand side is zero.  No assertion is made that the
unprojected quasilocal master kernel itself vanishes.
\end{proof}

\begin{theorem}[Exact local cocycle at arbitrary prescribed fixed loop order]
\label{thm:v5-L-cocycle}
Let $\widetilde{\mathcal B}_{j,L}$ be the complete order-$L$ master coefficient
formed from the forest-prepared native shell and the already fixed coefficients
$\widehat S_{0,j},\ldots,\widehat S_{L-1,j}$ before choosing the new order-$L$
restorer.  Define
\begin{equation}
 B_{j,L}
 :=\mathsf P_{L;N,\Delta,r}\widetilde{\mathcal B}_{j,L}.
 \label{eq:v5-BL}
\end{equation}
Then
\begin{equation}
 \boxed{d_{j,L}B_{j,L}=0.}
 \label{eq:v5-BL-closed}
\end{equation}
\end{theorem}

\begin{proof}
Apply \eqref{eq:v5-consistency-coefficient} at $n=L$ to the full
packet-carrying lower-order state.  Proposition
\ref{prop:v5-packet-restoration} makes the local projection of its right-hand
side vanish exactly.  The local projector commutes with the classical
differential by Theorem~\ref{thm:v5-complex}.  This gives
\eqref{eq:v5-BL-closed}.
\end{proof}

\begin{remark}[What V5 repaired]
Without the contraction-complete packet, the proof above would silently drop
terms of the form $\mathsf P\Delta K^{\remn}$ and
$\mathsf P(\bar S_p,K^{\remn})$.  V5 does not claim those terms vanish.  It
records them as lower-order future-contact coordinates and restores them when
they first become local.  This is the all-loop version of V4's insistence that
source-dependent measure/reference contacts be assembled before projection.
\end{remark}

\section{One common order-$L$ restoring primitive}
\label{sec:v5-hodge}

Choose the same type of positive coefficient inner product as in V3 on the
finite bank \eqref{eq:v5-L-bank}.  Let $d_{j,L}^\dagger$ be the adjoint,
\begin{equation}
 \Box_{j,L}=d_{j,L}d_{j,L}^\dagger+d_{j,L}^\dagger d_{j,L},
 \qquad
 h_{j,L}=d_{j,L}^\dagger G_{j,L},
 \label{eq:v5-hodge}
\end{equation}
where $G_{j,L}$ is the inverse of $\Box_{j,L}$ on the orthogonal complement of
its harmonic space.

\begin{lemma}[Fixed-$L$ Hodge contraction]
\label{lem:v5-hodge}
On a compact constant-rank component of the protected chart,
\begin{equation}
 d_{j,L}h_{j,L}+h_{j,L}d_{j,L}=I-\mathsf H_{j,L},
 \qquad
 \|h_{j,L}\|\le C_{L,\Delta,r,N}<\infty.
 \label{eq:v5-hodge-id}
\end{equation}
\end{lemma}

\begin{proof}
The bank is finite dimensional by Theorem~\ref{thm:v5-finite-local-bank}, so
finite-dimensional Hodge decomposition applies.  Constant rank on the compact
chart gives a positive lower bound for the nonzero singular values of
$d_{j,L}$, exactly as in V3.  The constant may deteriorate with $L$ and
$\Delta$; no uniform all-orders estimate is asserted.
\end{proof}

\begin{theorem}[Common order-$L$ master restorer]
\label{thm:v5-L-restorer}
The complete local breaking $B_{j,L}$ of Theorem~\ref{thm:v5-L-cocycle} has
zero harmonic projection and admits the single restoring functional
\begin{equation}
 \boxed{
 C_{j,L}^{\rm rest}:=-h_{j,L}B_{j,L},
 \qquad
 d_{j,L}C_{j,L}^{\rm rest}=-B_{j,L}.}
 \label{eq:v5-L-restorer}
\end{equation}
The same $C_{j,L}^{\rm rest}$ simultaneously restores the diffeomorphism,
local-Lorentz/spin, internal-gauge, Higgs/Yukawa, chiral-line,
ghost/nonminimal, antifield, measure/reference, stress/composite and retained
contact rows.  Its source descendants are not independently selected.
\end{theorem}

\begin{proof}
Theorem~\ref{thm:v5-L-cocycle} gives $d_{j,L}B_{j,L}=0$.  By Theorem
\ref{thm:v5-complex}, the only possible ghost-one harmonic coordinates are the
standard anomaly quotient, and Theorem~\ref{thm:v5-anomaly-zero} annihilates
that projection for the complete native Standard-Model shell.  Hence
$\mathsf H_{j,L}B_{j,L}=0$.  Apply the Hodge identity
\eqref{eq:v5-hodge-id}.  Because $B_{j,L}$ was assembled in the complete
source/contact bank before the homotopy was applied, the one primitive carries
all descendants at once.
\end{proof}

\section{The arbitrary fixed-loop shell theorem}
\label{sec:v5-shell-theorem}

Let $S_{j,L}^{\rm pre}$ be the common local order-$L$ coefficient obtained by
summing the local forest parts of Theorem~\ref{thm:v5-forest-decomposition} and
imposing the declared common physical normalization on the independent
ghost-zero EFT quotient.  Put
\begin{equation}
 \bar S_{j,L}
 :=S_{j,L}^{\rm pre}+C_{j,L}^{\rm rest}.
 \label{eq:v5-SLbar}
\end{equation}
Every quasilocal forest remainder and its contraction-complete packet is kept
as part of $K_{j,L}^{\remn}$.

\begin{theorem}[Arbitrary fixed-loop native ESM shell self-map]
\label{thm:v5-E4}
Fix the explicit V2 Wilson/overlap branch, a protected constant-rank chart and
finite data $(L,\Delta,r,N)$.  Assume the V1--V4 base case and the induction
datum $\mathfrak I_{<L}(\Delta,r,N)$.  Then the native order-$L$ shell step
produces a unique common update, modulo the declared physical finite
normalization coordinates,
\begin{equation}
 \boxed{
 \mathcal R_{j}^{[\le L]}:
 \mathfrak C^{\rm ESM}_{L;N,\Delta,r}
 \longrightarrow
 \mathfrak C^{\rm ESM}_{L;N,\Delta,r},}
 \label{eq:v5-selfmap}
\end{equation}
with the following properties:
\begin{enumerate}[label=\textup{(S\arabic*)},leftmargin=3.2em]
\item every proper forest subtraction is owned by a previously fixed common
coefficient $\bar S_{q,j}$, $q<L$;
\item the order-$L$ local action is the single coefficient
\eqref{eq:v5-SLbar}, with a finite EFT/Wilson bank at the prescribed order;
\item every retained source/contact and line coefficient is a descendant of
that same common update;
\item the complete local order-$L$ master coefficient vanishes after installing
$C_{j,L}^{\rm rest}$;
\item every quasilocal forest complement obeys
\begin{equation}
 \|K_{j,L}^{\remn}\|_*
 \le C_{L,\Delta,r,N}\epsilon_j,
 \label{eq:v5-KL-bound}
\end{equation}
and is carried with contraction depth sufficient for every later coefficient
through loop order $L$;
\item all induced lower-arity local descendants of earlier quasilocal kernels
which first appear at order $L$ are included before the master restorer is
formed.
\end{enumerate}
No uniform estimate as $L\to\infty$, no convergence of the perturbative series,
no finite-dimensional renormalizability of Einstein gravity and no
nonperturbative ESM continuum are claimed.
\end{theorem}

\begin{proof}
The finite graph family and proper-subgraph loop drop are Proposition
\ref{prop:v5-finite-graphs} and Lemma~\ref{lem:v5-loop-drop}.  Therefore all
proper local insertions are already owned, and Theorem
\ref{thm:v5-forest-decomposition} gives the finite local-plus-quasilocal
forest decomposition with the rooted bound.  Theorem
\ref{thm:v5-finite-local-bank} gives a finite target state.  The
contraction-complete packet makes every local feedback from lower quasilocal
kernels explicit, so Theorem~\ref{thm:v5-L-cocycle} gives an exact local
cocycle rather than a cocycle obtained by dropping remainder contacts.
Theorem~\ref{thm:v5-anomaly-zero} removes its harmonic obstruction and Theorem
\ref{thm:v5-L-restorer} supplies one common source-complete primitive.  Adding
that primitive to the forest-prepared local coefficient gives
\eqref{eq:v5-SLbar}; retaining the quasilocal packet gives the state on the
right of \eqref{eq:v5-selfmap}.  Every step depends only on the fixed finite
data, so the induction closes.
\end{proof}

\begin{corollary}[E4 induction]
\label{cor:v5-E4-induction}
Starting from the V1--V4 two-loop base, Theorem~\ref{thm:v5-E4} applies
recursively to every prescribed finite loop order.  Hence for each fixed
$L$ there is a source-complete native ESM shell state containing the common
local action through $L$, all required EFT/Wilson coordinates, determinant/
Pfaffian line data, complete source/contact descendants, common lower-order
forest ownership, and contraction-complete rooted/quasilocal remainders.
\end{corollary}

\begin{proof}
The base orders are the V1--V4 construction.  If the datum exists through
$L-1$, Theorem~\ref{thm:v5-E4} constructs the order-$L$ coefficient and the
future-safe remainder packets required by Definition
\ref{def:v5-induction-datum}.  This is exactly the induction hypothesis for the
next fixed order.
\end{proof}

\begin{center}
\fbox{\parbox{0.92\textwidth}{
\textbf{E4 is closed at arbitrary prescribed fixed loop order.}
The closure is coefficientwise and native: the local bank is finite only after
$(L,\Delta,r,N)$ are fixed, lower counterterms are owned by the same previously
normalized actions, and every quasilocal remainder is retained with the finite
set of future BV/source contraction jets which can feed a later local
coefficient.  Nothing here is uniform in $L$.}}
\end{center}

\section{A stronger anti-circularity audit}
\label{sec:v5-audit}

\begin{enumerate}[leftmargin=2.4em]
\item \textbf{No hidden infinite graph family.}
Quadratic and tadpole coordinates are absorbed into the common reference before
graph counting; source-marked low-valence vertices are bounded by $r$.
Therefore fixed loop number and arity genuinely give a finite topology family.

\item \textbf{No graphwise finite refits.}
Every proper 1PI ultraviolet block has smaller loop number and is replaced by
the restriction of the already fixed common action at that order.

\item \textbf{No continuum forest theorem is smuggled in.}
The forest formula is a finite native combinatorial identity.  Its estimate is
proved from native annular symbol bounds plus the rooted contraction ideal.

\item \textbf{Overlaps are not multiplied as if disjoint.}
Nonnested overlapping blocks belong to distinct compatible forests.  Nested
blocks are evaluated recursively inside-out.

\item \textbf{No discarded source or line contacts.}
Differentiation of a forest is performed before local projection.  The
contraction-complete packet stores every later local descendant generated by
source contacts, the BV Laplacian or determinant/Pfaffian line operations.

\item \textbf{The V4 remainder is not silently set to zero.}
Its current complementary kernel remains $O(\epsilon_j)$ and is carried in the
state.  What is set to zero inductively is the finite family of local descendant
coordinates required through the prescribed final loop order.

\item \textbf{No all-loop cohomology uniformity is claimed.}
The Hodge inverse is taken on a finite bank at each fixed $L$ and may have a
constant that grows arbitrarily badly with $L$.

\item \textbf{No nonrenormalizable-gravity sleight of hand.}
New ghost-zero gravitational EFT coordinates are admitted whenever the fixed
loop/EFT filtration requires them.  E4 proves finite closure at fixed order,
not a finite universal gravitational coupling set.
\end{enumerate}

\section{What E5 must now prove}
\label{sec:v5-E5-handoff}

E4 supplies a shell map at every fixed perturbative order.  The next problem is
not another loop induction.  It is an infinite-shell composition theorem at
fixed $(L,\Delta,r,N)$.

The contraction-complete packet clarifies the correct E5 state.  Let
$\mathfrak C_j^{(L,\Delta,r,N)}$ denote the full coefficient vector after the
shell at scale $\Lambda_j$, including local normalized coordinates and the
finite future-contact packet of its quasilocal remainders.  For two ultraviolet
cutoffs $J_2>J_1$, the difference of their induced low-scale coefficients must
be expressed as the sum of the shell tails
\begin{equation}
 \mathfrak C^{(J_2)}-\mathfrak C^{(J_1)}
 =\sum_{j=J_1+1}^{J_2}\delta_j\mathfrak C,
 \label{eq:v5-tail-sum}
\end{equation}
with the same finite normalization conditions imposed at every cutoff.
The forest theorem gives one-excess estimates for primitive complementary
pieces, but E5 must still prove that the \emph{composed} linearized transport
of those pieces through all lower shells has a cutoff-uniform bound.  That is
the genuine Cauchy theorem; it is not automatic from summability of the raw
single-shell remainder.

\begin{theorem}[Exact next load-bearing theorem for V6]
\label{thm:v5-next-E5}
Fix $(L,\Delta,r,N)$ and a positive reference/IR scale.  Prove a
cutoff-uniform bound on the derivative of the composed shell transport on the
complete contraction-closed state,
\begin{equation}
 \sup_{J>j\ge j_0}
 \left\|
 D\bigl(\mathcal R_{j_0}\circ\cdots\circ\mathcal R_{J}\bigr)
 \right\|_{L,\Delta,r,N}
 \le C_{L,\Delta,r,N},
 \label{eq:v5-composed-transport-goal}
\end{equation}
with the appropriate triangular interpretation on physical normalization
coordinates.  Then combine this with the fixed-shell excess estimate to prove
that every normalized connected coefficient and every stored BV/source/line
contact coefficient is Cauchy as the UV cutoff is removed.  The proof must also
show that the common local master identities and their contraction-descendant
coordinates pass to the same limit.  Scheme/reference compatibility and order
compatibility are reserved for E6.
\end{theorem}

\section{V5 status ledger}
\label{sec:v5-status}

\subsection*{Proved in V5}
\begin{enumerate}[leftmargin=2.2em]
\item The correct arbitrary-loop state cannot consist only of a local Taylor
bank plus an undifferentiated remainder.  A finite contraction-complete jet
packet is sufficient to retain every local feedback that the remainder can
produce before the prescribed final loop order.
\item After absorbing source-free valence-$0,1,2$ coordinates into the common
background/reference block, the graph, decoration and ancestry family is
finite at fixed $(L,\Delta,r,N)$.
\item Every proper connected 1PI ultraviolet subgraph has strictly smaller loop
number, so every proper subtraction is owned by a previously fixed common
action.
\item The native compatible-forest preparation, including overlapping and
nested blocks and all retained markings, decomposes into a finite local
Taylor/EFT/source part plus a rooted quasilocal complement of size
$O(\mu/\Lambda_j)$ (or $O(a\mu)$ at the ultraviolet cap).  Terms with $s$
complementary ancestors carry $O(\epsilon_j^s)$ after their new local descendants
are extracted.
\item The finite local bank closes under the classical BV differential and has
the same four-dimensional ghost-one anomaly quotient as the V80/V3 ESM
complex; the complete native Standard-Model shell has zero projection on that
quotient.
\item Because future local descendants of lower remainders are stored rather
than discarded, the order-$L$ local master breaking is an exact classical BV
cocycle.  One finite Hodge primitive therefore restores the complete local
order-$L$ master/source row.
\item These facts give a native source-complete shell self-map for every
prescribed fixed loop order, with constants permitted to depend on $L$.
Therefore E4 is closed at the stated coefficientwise shell scope.
\end{enumerate}

\subsection*{Inherited}
\begin{enumerate}[leftmargin=2.2em]
\item V80's bosonic two-loop forest ownership, local cohomology/range theorem,
retained soft-constraint architecture and rooted hard-annulus estimates.
\item The complete Standard-Model representation/anomaly trace cancellations,
even $SU(2)$ doublet parity and finite gauge descent of the protected chiral
section.
\item V1's overlap functional calculus, Kato frame and hard-pivot/rooted bounds;
V2's explicit Wilson branch and annular/full-zone comparison; V3's common
chiral Hodge restorer and E2 self-map; V4's complete two-loop canonical-pair and
source-contact master audit.
\end{enumerate}

\subsection*{Conditional}
\begin{enumerate}[leftmargin=2.2em]
\item The construction remains native to the explicit V2 Wilson/overlap
completion.  A different microscopic fermion kernel still requires the V2
regulator-comparison gate.
\item Hodge bounds are on compact constant-rank components of the finite local
complex.  Rank crossings require a new chart.
\item The determinant/Pfaffian packet remains on the protected nonzero component
with the retained positive IR/reference prescription; no global line holonomy
or zero-mode crossing theorem is claimed.
\item The native annular/rooted bounds used in the forest theorem are consumed
only at fixed graph size.  Their constants may grow without control as
$L\to\infty$.
\end{enumerate}

\subsection*{Open}
\begin{enumerate}[leftmargin=2.2em]
\item E5: prove cutoff-uniform composed transport and the Cauchy/unique
multi-shell continuum for every fixed $(L,\Delta,r,N)$, including the complete
contraction-descendant packet, operator mixing, stress/contact and line
insertions.
\item Prove that the limiting normalized connected generating-functional
coefficients inherit the common BV/master/source identities.  This must be done
on the composed state, not inferred merely from the single-shell
$O(\epsilon_j)$ estimate.
\item E6: finite scheme/reference compatibility, identification of
scheme-invariant normalized relations, and order compatibility between the
$L$ and $L+1$ constructions.
\end{enumerate}

\subsection*{Exact next load-bearing theorem}
\begin{center}
\fbox{\parbox{0.92\textwidth}{
\textbf{E5: fixed-order multi-shell Cauchy theorem with contraction-complete transport.}
At fixed $(L,\Delta,r,N)$ and positive reference scale, prove a uniform bound
for the composed derivative of the common shell maps on the full local plus
future-contact state.  Use that bound and the V5 one-excess forest estimate to
show that the normalized connected/source/BV/line coefficients form a Cauchy
family as the ultraviolet cutoff is removed, and that the common master/source
identities pass to the same limit.}}
\end{center}

\section*{Append-only continuation marker: V5}
\addcontentsline{toc}{section}{Append-only continuation marker: V5}
\textbf{End of V5.}  The complete V1--V4 body above is retained verbatim.  V5
closes E4 by replacing the naive ``local jet plus remainder'' state with a
finite contraction-complete jet packet, proving the native arbitrary-fixed-loop
forest decomposition, establishing common lower-order counterterm ownership,
and restoring the complete local order-$L$ master/source coefficient with one
Hodge primitive.  Future versions must append after this marker.  The next file
is to be named
\texttt{Renewal\_Perturbative\_ESM\_Continuum\_Research\_Notes\_V6.tex}.




\clearpage
\part*{V6: Fixed-order multi-shell continuum in a renormalized transport frame}
\addcontentsline{toc}{part}{V6: Fixed-order multi-shell continuum in a renormalized transport frame}
\label{part:v6}

\section{Purpose, inherited input and the E5 quantifier correction}
\label{sec:v6-purpose}

V5 closed the native shell problem at every prescribed finite loop order.  The
remaining E5 question is different: at fixed
\[
 (L,\Delta,r,N),
\]
one must compose arbitrarily many ultraviolet shells and prove that the
\emph{renormalized} connected/source/BV/line coefficients have a unique
cutoff-independent limit.  We keep throughout the explicit V2 Wilson/overlap
branch, a protected constant-rank chart, a positive reference/infrared scale
$\mu>0$, a finite external/source panel, and one fixed normalization/reference
prescription.  Scheme comparison is postponed to E6.

The exact next theorem stated at the end of V5 asked for a uniform bound on the
raw derivative of a long product of shell maps.  That is stronger than E5
requires and, in marginal directions, generally false.  The first task of V6
is therefore to correct the variable rather than to force a spurious
contraction estimate.

Let
\[
 \Lambda_j=b^j\mu_0,\qquad b>1,
 \qquad
 \epsilon_j:=\frac{\mu}{\Lambda_j},
\]
on the physical annular regime.  On the terminal ultraviolet cap we use the
already inherited equivalent form $\epsilon_j=a_j\mu$.  At fixed
$(L,\Delta,r,N)$ all coefficient spaces below are finite products of the local
bank and the finite contraction-complete rooted packet of V5.

\begin{proposition}[The raw composed-derivative bound is not an invariant E5 requirement]
\label{prop:v6-raw-obstruction}
There exist finite triangular shell transports which preserve a perfectly
well-defined renormalized coefficient but for which
\[
 \sup_{J>j}\left\|D(\mathcal R_j\circ\cdots\circ\mathcal R_J)\right\|=\infty.
\]
Consequently the raw version of \eqref{eq:v5-composed-transport-goal} cannot be
used as an intrinsic acceptance criterion for E5.
\end{proposition}

\begin{proof}
On $\mathbb R^2$ consider
\[
 \mathcal R_k(x,y)=(x,y+\beta x),\qquad \beta\ne0.
\]
Then
\[
 D\mathcal R_k=
 \begin{pmatrix}1&0\\ \beta&1\end{pmatrix},
 \qquad
 D(\mathcal R_j\circ\cdots\circ\mathcal R_J)=
 \begin{pmatrix}1&0\\ (J-j+1)\beta&1\end{pmatrix}.
\]
The raw derivative therefore grows linearly with the number of shells.  Yet
with the cutoff-dependent triangular coordinate
\[
 \widehat y_k:=y-k\beta x
\]
the same transport is exactly the identity.  This is the finite-dimensional
model of logarithmic operator mixing or marginal finite renormalization.  E5
must control the quotient transport after the local normalization cocycle has
been factored out; it must not demand bounded bare renormalization constants.
\end{proof}

\begin{remark}[What is and is not being changed]
\label{rem:v6-no-scheme-change}
The coordinate correction below is not a comparison of renormalization
schemes.  We keep the single common physical/source/reference normalization
already fixed in V1--V5 and merely express every cutoff in that same normalized
frame.  The possibly large triangular map from bare/local coordinates to this
frame is allowed to depend on the cutoff.  E6 will compare two different such
choices.
\end{remark}

\section{Normalization-complete local coordinates}
\label{sec:v6-normalization-frame}

Let $\mathcal L_{j}^{(L,\Delta,r,N)}$ be the finite local coefficient bank of
V5.  Split its coordinates, using the fixed V5 cohomology/Hodge decomposition,
as
\begin{equation}
 \mathcal L_j
 =\mathcal P\oplus\mathcal X_j\oplus\mathcal E_j.
 \label{eq:v6-local-split}
\end{equation}
Here $\mathcal P$ contains all independent renormalized physical/EFT and
source-normalization coordinates retained at the prescribed order,
$\mathcal X_j$ is the chosen finite BV-exact/redundant complement, and
$\mathcal E_j$ denotes the finite collection of local descendants fixed by the
common master/source equations.  This notation does not identify
$\mathcal E_j$ with a new counterterm space: its coordinates are determined by
the same common action and Hodge section.

Choose linear normalization functionals
\begin{equation}
 \mathfrak n_j:\mathcal L_j\longrightarrow\mathcal P
 \label{eq:v6-normalization-map}
\end{equation}
which read the prescribed masses, couplings, finite EFT/Wilson coordinates,
renormalized composite/source basis, and declared line/reference
normalizations at the fixed reference scale.  At fixed finite filtration we
may, and do, choose the local basis dual to these functionals on the quotient.
Thus the quotient block of $D\mathfrak n_j$ is the identity.  The V1--V5
finite-to-ordinary comparison implies that the remaining finite blocks differ
between adjacent ultraviolet cutoffs by $O(\epsilon_j)$.

\begin{lemma}[Uniform local normalization section at fixed filtration]
\label{lem:v6-normalization-section}
Fix a compact normalized parameter set $K_{\rm ren}\Subset\mathcal P$ inside
the protected component.  For all sufficiently large shell indices there is a
unique smooth local section
\begin{equation}
 \sigma_j:K_{\rm ren}\longrightarrow\mathcal L_j,
 \qquad
 \mathfrak n_j\circ\sigma_j=\operatorname{id},
 \label{eq:v6-section}
\end{equation}
compatible with the V5 Hodge choice and common reference convention, and
\begin{equation}
 \sup_j\|D\sigma_j\|<\infty,
 \qquad
 \|\sigma_{j+1}-\sigma_j\|_{C^1(K_{\rm ren})}
 \le C_{L,\Delta,r,N}\epsilon_j.
 \label{eq:v6-section-bound}
\end{equation}
The finitely many lower shells may be absorbed into the constant.
\end{lemma}

\begin{proof}
On the independent quotient block the chosen normalization coordinates are
literally the coefficient coordinates, so the Jacobian is the identity.  On
the exact/descendant block the section is obtained by the fixed finite Hodge
splitting of V5 and by the declared reference conditions.  Constant rank on
the compact protected chart gives a uniform lower singular-value bound for the
finite matrices entering that splitting.  Hence the implicit-function/Hodge
solution has a uniformly bounded derivative.

Between adjacent ultraviolet cutoffs, every coefficient entering these finite
matrices is either an exactly fixed normalization coordinate or a native
symbol/source/line coefficient.  V1--V5 give the latter an adjacent-cutoff
finite-to-ordinary discrepancy bounded by the one-excess factor
$C\epsilon_j$, including all source and contraction descendants.  The resolvent
identity for the finite inverse matrices then gives the same $C^1$ bound for
the sections.  No uniformity in $L$ is used.
\end{proof}

Let $\mathcal Q_j$ denote the finite V5 contraction-complete quasilocal packet.
The full state is
\begin{equation}
 \mathfrak C_j=\mathcal L_j\oplus\mathcal Q_j.
 \label{eq:v6-full-state}
\end{equation}
We use the section to identify its normalized version with a fixed finite
Banach coordinate space
\begin{equation}
 \widehat{\mathfrak C}
 :=K_{\rm ren}\times\widehat{\mathcal Q},
 \label{eq:v6-normalized-state}
\end{equation}
where the packet basis is transported by the same source/operator/frame
normalization.  The resulting cutoff-dependent identification will be denoted
$\Theta_j$.

\begin{definition}[Reduced or renormalized shell transport]
\label{def:v6-reduced-shell}
The E5 shell map in normalized coordinates is
\begin{equation}
 \widehat{\mathcal R}_j
 :=\Theta_{j-1}\circ\mathcal R_j^{[\le L]}\circ\Theta_j^{-1}.
 \label{eq:v6-reduced-shell}
\end{equation}
The local normalization cocycle, including triangular composite/operator
mixing and line-frame normalization, is therefore contained in the two
$\Theta$ factors rather than in the dynamical remainder of
$\widehat{\mathcal R}_j$.
\end{definition}

\section{One-shell normal form after removing the local cocycle}
\label{sec:v6-one-shell-normal-form}

\begin{theorem}[Summable normal form of the fixed-order shell map]
\label{thm:v6-one-shell-normal-form}
For every fixed $(L,\Delta,r,N)$ there is a neighborhood
$U\supset K_{\rm ren}\times\{0\}$ and constants $C_k$ such that, for every
fixed derivative order $k$ required by the retained source/contact panel,
\begin{equation}
 \boxed{
 \widehat{\mathcal R}_j=I+\mathcal E_j,
 \qquad
 \|\mathcal E_j\|_{C^k(U)}
 \le C_k\epsilon_j.}
 \label{eq:v6-I-plus-E}
\end{equation}
The statement includes the complete contraction-descendant packet.  In
particular, it is not merely a bound on the undifferentiated effective action.
\end{theorem}

\begin{proof}
Apply the V5 forest decomposition to one shell.  Every local Taylor/EFT/source
coefficient splits into an independent quotient coordinate and a dependent
exact/descendant coordinate.  The independent coordinates are reset by the
same normalization functionals at the two cutoffs; after conjugation by
$\Theta_j$ their shell transport is exactly the identity.  The dependent local
coordinates are values of the uniformly bounded Hodge/reference section of
Lemma~\ref{lem:v6-normalization-section}.  Their change between adjacent
cutoffs is therefore $O(\epsilon_j)$.

Every nonlocal forest term belongs to the V5 rooted complement and is
$O(\epsilon_j)$.  More importantly, V5 stores every local coefficient that can
be created from that complement by the finite contraction alphabet.  Hence a
source derivative, BV contraction, determinant/Pfaffian line operation or
allowed canonical contraction cannot expose an unrecorded $O(1)$ local term:
that term has already been extracted into the local normalization/Hodge block.
The residual differentiated packet therefore remains $O(\epsilon_j)$ at every
fixed derivative order in the finite alphabet.  The finite number of graph,
forest, source and canonical-pair incidences at fixed
$(L,\Delta,r,N)$ may be absorbed into $C_k$.
\end{proof}

\begin{corollary}[Summable $C^1$ size]
\label{cor:v6-summable-c1}
For geometric shells,
\begin{equation}
 \sum_{j\ge j_0}
 \|\widehat{\mathcal R}_j-I\|_{C^1(U)}<\infty.
 \label{eq:v6-summable-c1}
\end{equation}
\end{corollary}

\begin{proof}
Use Theorem~\ref{thm:v6-one-shell-normal-form} and
$\sum_{j\ge j_0}\mu/\Lambda_j<\infty$.
\end{proof}

\section{Infinite products of the reduced shell maps}
\label{sec:v6-product}

We record the elementary nonlinear product estimate explicitly because it is
the actual E5 transport theorem.

\begin{lemma}[Summable near-identity composition]
\label{lem:v6-near-identity-product}
Let $U_0\Subset U$ be compact and suppose
$F_j=I+E_j$ satisfy
\[
 \|E_j\|_{C^1(U)}\le\delta_j,
 \qquad
 \sum_j\delta_j<\infty.
\]
After enlarging a finite low-shell constant if necessary, all sufficiently
high-tail compositions are defined on $U_0$, and
\begin{equation}
 \left\|D(F_{j_0}\circ\cdots\circ F_J)\right\|
 \le
 \exp\!\left(\sum_{j=j_0}^J\delta_j\right).
 \label{eq:v6-product-derivative}
\end{equation}
Moreover the compositions are uniformly Cauchy on $U_0$ and
\begin{equation}
 \|F_{j_0}\circ\cdots\circ F_{J_2}
   -F_{j_0}\circ\cdots\circ F_{J_1}\|_{C^0(U_0)}
 \le C\sum_{j>J_1}\delta_j
 \label{eq:v6-product-tail}
\end{equation}
for $J_2>J_1$.
\end{lemma}

\begin{proof}
The derivative estimate is the chain rule followed by
$\|DF_j\|\le1+\delta_j\le e^{\delta_j}$.  The total displacement of a high
shell tail is bounded by $\sum_j\|E_j\|_{C^0}$, so a compact $U_0$ with a
slightly larger neighborhood in $U$ remains in the domain; finitely many lower
shells only multiply the final constant.  For the tail estimate insert the
additional maps one at a time and transport each one-shell displacement
through the lower composition.  The derivative bound gives
\[
 \|\Delta_{j\to j_0}\|
 \le C\delta_j.
\]
Summing $j=J_1+1,\ldots,J_2$ proves
\eqref{eq:v6-product-tail}.
\end{proof}

\begin{theorem}[Cutoff-uniform composed transport in the renormalized frame]
\label{thm:v6-composed-transport}
At fixed $(L,\Delta,r,N)$,
\begin{equation}
 \boxed{
 \sup_{J>j\ge j_0}
 \left\|
 D\bigl(
 \widehat{\mathcal R}_{j_0+1}\circ\cdots\circ
 \widehat{\mathcal R}_{J}
 \bigr)
 \right\|
 \le C_{L,\Delta,r,N}<\infty.}
 \label{eq:v6-composed-transport}
\end{equation}
This is the correct normalized version of the V5 transport goal.  No bound is
claimed for the unrenormalized triangular maps $\Theta_j^{-1}\Theta_{j+1}$.
\end{theorem}

\begin{proof}
Apply Corollary~\ref{cor:v6-summable-c1} and
Lemma~\ref{lem:v6-near-identity-product}.  The finite low-shell prefix and the
fixed compact normalization chart are absorbed into the constant.
\end{proof}

\section{The ultraviolet Cauchy theorem for the complete state}
\label{sec:v6-state-Cauchy}

For a UV cutoff at shell $J$, let
$\widehat{\mathfrak C}^{(J)}_{j_0}$ denote the normalized state at the fixed
reference shell $j_0$ obtained by starting with the same normalized physical,
source and reference data at the ultraviolet end and composing the reduced
shell maps downward.  The top packet is the native finite-cutoff packet in the
same normalized frame.

\begin{lemma}[Adjacent ultraviolet boundary data differ by one excess scale]
\label{lem:v6-top-boundary}
The normalized ultraviolet boundary states satisfy
\begin{equation}
 \|\widehat{\mathfrak C}^{(J+1)}_{J+1}
   -\widehat{\mathfrak C}^{(J)}_{J}\|_{\rm matched}
 \le C_{L,\Delta,r,N}\epsilon_J,
 \label{eq:v6-top-boundary}
\end{equation}
where ``matched'' means that both are expressed in the common normalized
coordinate and source/line frame.
\end{lemma}

\begin{proof}
Independent local quotient coordinates agree exactly by construction.  The
exact/redundant descendants differ only through the adjacent-cutoff change of
the finite Hodge/reference section, bounded in
Lemma~\ref{lem:v6-normalization-section}.  The remaining native hard kernels,
line factors and all contraction descendants differ by the V1--V5
finite-to-ordinary/full-zone estimate, whose ultraviolet-cap form is
$O(a_J\mu)=O(\epsilon_J)$.  Combining the finite coordinates gives the stated
bound.
\end{proof}

\begin{theorem}[Fixed-order multi-shell Cauchy theorem]
\label{thm:v6-state-Cauchy}
For every prescribed finite $(L,\Delta,r,N)$ and fixed positive reference
scale,
\begin{equation}
 \boxed{
 \widehat{\mathfrak C}^{(J)}_{j_0}
 \longrightarrow
 \widehat{\mathfrak C}^{(\infty)}_{j_0}}
 \label{eq:v6-state-limit}
\end{equation}
in the complete local plus contraction-descendant rooted topology.  More
precisely, for $J_2>J_1$,
\begin{equation}
 \|\widehat{\mathfrak C}^{(J_2)}_{j_0}
   -\widehat{\mathfrak C}^{(J_1)}_{j_0}\|_*
 \le
 C_{L,\Delta,r,N}
 \sum_{j>J_1}\epsilon_j
 \le
 C'_{L,\Delta,r,N}\frac{\mu}{\Lambda_{J_1}}.
 \label{eq:v6-state-Cauchy-rate}
\end{equation}
The limit is independent of the chosen cofinal sequence of UV cutoffs within
the fixed regulator branch and fixed normalization/reference prescription.
\end{theorem}

\begin{proof}
Compare two cutoffs by inserting the missing ultraviolet shells one at a time.
Each insertion changes the matched ultraviolet boundary data by
$O(\epsilon_j)$ by Lemma~\ref{lem:v6-top-boundary}; equivalently, after the
boundary is fixed, it inserts one reduced shell map
$I+O(\epsilon_j)$.  Transport this perturbation to the reference scale using
Theorem~\ref{thm:v6-composed-transport}.  Its norm remains
$C\epsilon_j$.  Sum the inserted shells.  The geometric tail is
\[
 \sum_{j>J_1}\frac{\mu}{b^j\mu_0}
 =\frac{\mu}{\Lambda_{J_1}}
   \frac{b^{-1}}{1-b^{-1}},
\]
up to the harmless indexing convention.  Hence the family is Cauchy in the
complete finite product norm.  Completeness gives the limit and the same tail
bound proves independence of a cofinal subsequence.
\end{proof}

\begin{remark}[Bare running is not asserted to converge]
\label{rem:v6-bare-running}
The theorem concerns $\widehat{\mathfrak C}$, i.e. coefficients in the fixed
renormalized local/source/line frame.  Bare marginal couplings, wave-function
factors, composite mixing matrices or determinant-line trivializations may
contain powers of $\log\Lambda_J$ or stronger allowed local dependence.  Those
belong to $\Theta_J$ and are not continuum observables by themselves.
\end{remark}

\section{Direct highest-scale form of the same Cauchy estimate}
\label{sec:v6-highest-scale}

For later use in the connected functional it is convenient to unfold the shell
composition into scale-decorated forests.  This also gives an independent
anti-circularity check on Theorem~\ref{thm:v6-state-Cauchy}.

\begin{definition}[Highest-scale label]
\label{def:v6-highest-scale}
For a V5 forest-prepared graph $\Gamma$ with internal shell labels
$\mathbf h=(h_e)_{e\in E(\Gamma)}$, put
\[
 H(\mathbf h):=\max_{e\in E(\Gamma)}h_e.
\]
All local proper-subgraph assignments use the already fixed lower-order common
actions and the common normalization section.
\end{definition}

\begin{lemma}[Highest-scale excess or local ownership]
\label{lem:v6-highest-excess}
Fix a scale-decorated forest term contributing to a normalized coefficient.
After all proper subtractions, exactly one of the following holds:
\begin{enumerate}[label=\textup{(\roman*)},leftmargin=2.7em]
\item its maximal-scale connected component is entirely local through the
prescribed filtration, in which case that contribution is owned by the common
local normalization/Hodge section and does not enter the cutoff difference;
\item its normalized remainder carries at least one Taylor/root separation
factor and is bounded by
\begin{equation}
 C_{\Gamma,L,\Delta,r,N}
 \frac{\mu}{\Lambda_{H(\mathbf h)}}.
 \label{eq:v6-highest-excess}
\end{equation}
\end{enumerate}
The same dichotomy holds after every stored BV/source/line contraction.
\end{lemma}

\begin{proof}
Apply the V5 forest decomposition to the connected component containing a line
at the maximal scale.  Its complete Taylor/source jet is, by definition,
inserted into the common local bank with its inherited lower-order ownership.
If the term is not exhausted by this local jet, V5's native Taylor remainder
theorem gives one factor of external/source momentum divided by the maximal
hard scale.  Rooted contraction with lower-scale components does not remove
that factor: any local descendant thereby generated is extracted into the
contraction-complete packet, leaving the complementary term in the one-excess
ideal.  This proves the dichotomy and its stability under the finite
contraction alphabet.
\end{proof}

\begin{proposition}[Absolute convergence of scale histories]
\label{prop:v6-scale-history}
For each fixed decorated graph topology $\Gamma$ with $I_\Gamma$ internal
lines, the sum of normalized forest terms over all shell assignments is
absolutely convergent after the common local normalization.  The contribution
of assignments with $H(\mathbf h)>J$ obeys
\begin{equation}
 \sum_{H>J}(1+H-j_0)^{I_\Gamma-1}
 \frac{C_\Gamma\mu}{\Lambda_H}
 \xrightarrow[J\to\infty]{}0.
 \label{eq:v6-scale-history-tail}
\end{equation}
\end{proposition}

\begin{proof}
For fixed maximal label $H$, there are at most
$C(1+H-j_0)^{I_\Gamma-1}$ assignments of the remaining finitely many line
labels.  By Lemma~\ref{lem:v6-highest-excess}, nonlocal terms have the factor
$\mu/\Lambda_H$; purely local terms are already removed by the common
normalization and do not enter the normalized cutoff tail.  Since a polynomial
times $b^{-H}$ is summable, the series converges absolutely.  The finite V5
graph family at fixed $(L,\Delta,r,N)$ allows summation over topologies as well.
\end{proof}

\section{Normalized connected generating-functional coefficients}
\label{sec:v6-W}

Let $Z_J[J]$ be the finite-cutoff source functional in the fixed normalized
source frame and define
\begin{equation}
 W_J[J]
 :=\log\frac{Z_J[J]}{Z_J[0]}.
 \label{eq:v6-W}
\end{equation}
We use $J$ both as the conventional source letter inside $W$ and as a cutoff
index only when the meaning is unambiguous; cutoff indices will otherwise be
written $J_{\rm UV}$.

At a fixed loop/source order and positive reference scale, coefficient
extraction from the terminal low-scale state is a finite sum of connected
rooted contractions.  Denote this finite polynomial/analytic map by
\begin{equation}
 \mathcal W_{L,\Delta,r,N}:
 \widehat{\mathfrak C}_{j_0}\longrightarrow\mathcal W_{L,\Delta,r,N}.
 \label{eq:v6-W-map}
\end{equation}

\begin{lemma}[Continuity of the terminal connected-functional map]
\label{lem:v6-W-continuity}
On the protected compact reference domain,
\begin{equation}
 \|D\mathcal W_{L,\Delta,r,N}\|
 \le C_{L,\Delta,r,N}<\infty.
 \label{eq:v6-W-Lipschitz}
\end{equation}
The map includes composite and stress insertions, source contacts, operator
mixing, ghosts/antifields and determinant/Pfaffian line insertions.
\end{lemma}

\begin{proof}
At fixed perturbative and source order, $W$ is the finite connected part of the
terminal Wick/forest expansion.  Every propagator on the retained reference
domain has the positive IR/hard-pivot bound inherited from V1--V5, and every
vertex/source/line derivative belongs to the finite rooted bank.  Hence every
coefficient is a finite multilinear contraction of state coordinates with a
uniform bound on the compact domain.  Differentiating that finite expression
again gives a finite sum of the same type.  Determinant/Pfaffian insertions use
the V1--V2 line inverse bounds and their stored derivatives.  The normalization
$Z[J]/Z[0]$ removes the source-free vacuum coordinate exactly; a field-independent
but source-dependent contact is not removed and remains one of the finite
coordinates of the map.
\end{proof}

\begin{theorem}[Source-complete coefficientwise ultraviolet continuum]
\label{thm:v6-W-continuum}
For every prescribed fixed $(L,\Delta,r,N)$ and finite external/source panel,
every coefficient of the normalized connected functional has a unique UV
limit
\begin{equation}
 \boxed{
 [W_\infty]_I^{(\ell)}
 :=\lim_{J_{\rm UV}\to\infty}
 [W_{J_{\rm UV}}]_I^{(\ell)},
 \qquad
 \ell\le L,\quad |I|\le r.}
 \label{eq:v6-W-limit}
\end{equation}
The convergence holds jointly for the retained momentum derivatives through
$\Delta$ and for the complete finite panel of composite, stress, contact,
operator-mixing, ghost/antifield and determinant/Pfaffian line insertions.
Moreover
\begin{equation}
 \|W_{J_2}-W_{J_1}\|_{L,\Delta,r,N}
 \le
 C_{L,\Delta,r,N}\frac{\mu}{\Lambda_{J_1}}
 \label{eq:v6-W-rate}
\end{equation}
for $J_2>J_1$ sufficiently large.
\end{theorem}

\begin{proof}
Apply Lemma~\ref{lem:v6-W-continuity} to the Cauchy estimate
\eqref{eq:v6-state-Cauchy-rate}.  This already proves the theorem.  Equivalently,
unfold each connected coefficient into its finite family of decorated graph
topologies and use Proposition~\ref{prop:v6-scale-history}; the logarithm
selects connected graphs and the vacuum normalization deletes exactly the
source-free vacuum family.  Both routes preserve all marked/contact and line
coordinates because these were part of the state before either composition or
connected extraction.
\end{proof}

\section{Operator mixing and the meaning of the limiting source basis}
\label{sec:v6-operator-mixing}

Let $\mathbf O$ be any finite set of composite/source operators retained in the
panel.  In raw coordinates one may have a triangular cutoff-dependent mixing
matrix $Z_j^{\mathbf O}$.  The normalization map $\Theta_j$ uses the prescribed
reference conditions to define
\begin{equation}
 \mathbf O_R=(Z_j^{\mathbf O})^{-1}\mathbf O_{\rm bare}.
 \label{eq:v6-renormalized-operator}
\end{equation}

\begin{corollary}[Coefficientwise continuum with operator mixing]
\label{cor:v6-operator-mixing}
All mixed connected coefficients of the finite renormalized operator panel
$\mathbf O_R$ converge as in Theorem~\ref{thm:v6-W-continuum}.  No convergence
of $Z_j^{\mathbf O}$ itself is required.
\end{corollary}

\begin{proof}
The renormalized operator/source basis is one block of the normalized state
coordinate $\widehat{\mathfrak C}$.  The possibly nonconvergent triangular
matrix is part of $\Theta_j$, already factored from the reduced shell map.
Apply Theorem~\ref{thm:v6-W-continuum} in that fixed normalized basis.
\end{proof}

\section{Passage of the common BV/master/source identities}
\label{sec:v6-master-limit}

The identity which can be passed without adding a new hypothesis is exactly the
fixed-order identity proved in V5: the complete local master coordinates and
all contraction descendants required by the prescribed final filtration vanish
in one common action/source state.  We do not silently promote an unprojected
finite-cutoff quasilocal master remainder into zero.

Let $\mathbf M_{L,\Delta,r,N}$ denote the finite vector consisting of:
\begin{enumerate}[label=\textup{(M\arabic*)},leftmargin=3em]
\item every local coefficient through the retained EFT/derivative order of
$\frac12(S,S)-\hbar\Delta S$;
\item every source/contact derivative of those coefficients through order $r$;
\item every future contraction descendant stored by the V5 packet through
final loop order $L$;
\item the corresponding finite connected Ward/Zinn--Justin coefficients in the
normalized source panel.
\end{enumerate}
By V5 these coordinates are zero for every finite cutoff after the common
restorers and normalizations are installed.

\begin{theorem}[Master/source identities commute with the UV limit]
\label{thm:v6-master-limit}
For the limiting normalized state of Theorem~\ref{thm:v6-state-Cauchy},
\begin{equation}
 \boxed{
 \mathbf M_{L,\Delta,r,N}
 \bigl(\widehat{\mathfrak C}^{(\infty)}_{j_0}\bigr)=0.}
 \label{eq:v6-master-limit}
\end{equation}
Equivalently, every retained coefficient of the common BV/BRST/master/source
identity and every stored stress/composite/contact/line descendant passes to
the same coefficientwise continuum limit.
\end{theorem}

\begin{proof}
The vector $\mathbf M_{L,\Delta,r,N}$ is a finite polynomial combination of
state coordinates, finite canonical contractions, the fixed finite BV
Laplacian blocks and source coefficient maps.  On the protected compact domain
it is continuous in the V5 rooted/product topology.  For every finite cutoff,
V5's exact projected consistency and common Hodge restoration give
\[
 \mathbf M_{L,\Delta,r,N}
 (\widehat{\mathfrak C}^{(J)}_{j_0})=0.
\]
Take $J\to\infty$ using Theorem~\ref{thm:v6-state-Cauchy}.  Because the
contraction-complete packet was included in the convergent state, no term of
the form $\mathsf P\Delta K^{\rm rem}$ or
$\mathsf P(S_p,K^{\rm rem})$ is lost in this passage.
\end{proof}

\begin{remark}[Scope of the master theorem]
\label{rem:v6-master-scope}
The theorem is coefficientwise at the prescribed finite EFT/derivative/source
order.  It does not claim that a finite-cutoff unprojected quasilocal master
kernel vanishes identically, nor does it claim an all-derivative or
nonperturbative Slavnov--Taylor functional identity in one step.  For any
larger finite derivative/source order the same construction may be rerun, and
E6 will impose compatibility between the resulting finite-order systems.
\end{remark}

\section{Separate thermodynamic comparison}
\label{sec:v6-thermodynamic}

The ultraviolet theorem above is formulated at fixed positive reference scale
and on the same-torus/nonwinding domain used in the inherited native shell
estimates.  The volume limit is logically separate.

\begin{lemma}[Uniform periodization tail for a finite rooted panel]
\label{lem:v6-volume-tail}
Suppose a retained infinite-volume kernel $K$ obeys, for some $m>4+s$,
\begin{equation}
 \int_{\mathbb R^4}(1+\mu|x|)^m|K(x)|\,dx\le C
 \label{eq:v6-weighted-kernel}
\end{equation}
uniformly in the UV cutoff and in the fixed source panel.  Let $K_R$ be its
periodization on a four-torus with minimal side $R$.  Then on the nonwinding
rooted panel with $s$ fixed relative coordinates,
\begin{equation}
 \|K_R-K\|_{*,s}
 \le C_{m,s}(\mu R)^{-(m-4-s)}.
 \label{eq:v6-volume-bound}
\end{equation}
\end{lemma}

\begin{proof}
Write the periodization error as the sum over nonzero lattice translates of the
fundamental cell.  On the nonwinding panel every such translate has distance
at least a fixed multiple of $R|n|$.  Pull
$(1+\mu R|n|)^{-(m-4-s)}$ from the weighted $L^1$ moment and sum the remaining
$|n|^{-(4+\eta)}$ tail.  Rooting removes the overall volume factor; the $s$
relative-coordinate weights cost at most $s$ powers.  This proves the stated
bound.
\end{proof}

\begin{theorem}[Staged thermodynamic continuum]
\label{thm:v6-volume-limit}
The V1--V5 arbitrary fixed polynomial-moment bounds imply the hypothesis of
Lemma~\ref{lem:v6-volume-tail} for every prescribed finite rooted/source panel.
Hence, after taking the ultraviolet limit of
Theorem~\ref{thm:v6-W-continuum}, the finite-volume normalized connected
coefficients have a unique thermodynamic limit along the declared nonwinding
panels.  The UV and volume limits commute whenever both are restricted to a
common compact source/momentum panel.
\end{theorem}

\begin{proof}
Choose $m$ larger than the fixed number of rooted relative-coordinate weights.
The inherited annular symbol estimates permit that arbitrary fixed number of
momentum integrations by parts, uniformly in the cutoff, and the V5 forest
contractions preserve the corresponding weighted rooted moment.  Lemma
\ref{lem:v6-volume-tail} therefore gives a volume error uniform in the UV
cutoff.  For cutoffs $J_1,J_2$ and volumes $R_1,R_2$, insert and subtract the
common infinite-volume or common finite-volume coefficient.  The triangle
inequality combines the UV tail
$C\mu/\Lambda_{J_1}$ with the volume tail
$C(\mu R_1)^{-q}$.  Both tend to zero independently, proving existence and
commutation of the staged limits.
\end{proof}

\section{E5 closure theorem}
\label{sec:v6-E5-close}

\begin{theorem}[Coefficientwise multi-shell Einstein--Standard Model continuum at fixed order]
\label{thm:v6-E5}
Fix the explicit V2 native Wilson/overlap branch, a protected constant-rank
component, a positive reference/IR scale, one common normalization/reference
scheme, finite $(L,\Delta,r,N)$, and a finite external/source panel.  Then:
\begin{enumerate}[label=\textup{(E5.\arabic*)},leftmargin=3.5em]
\item after factoring the finite-dimensional local normalization cocycle, every
native shell map is $I+O(\mu/\Lambda_j)$ in the complete
contraction-descendant $C^1$ state;
\item the reduced shell maps have cutoff-uniform composed transport and an
infinite-shell limit;
\item the common normalized local/source/line state is Cauchy with tail
$O(\mu/\Lambda_J)$;
\item every normalized connected coefficient through loop order $L$, derivative
order $\Delta$ and source order $r$ has a unique cutoff-independent limit,
including composite operators, stress insertions, contact terms, operator
mixing, ghosts/antifields and determinant/Pfaffian line insertions;
\item every retained common BV/BRST/master/source coefficient and every stored
contraction descendant passes to the same limit and remains zero;
\item the thermodynamic comparison can be taken separately, with uniform rooted
periodization tails, on the declared nonwinding finite panels.
\end{enumerate}
No convergence of the perturbative series, no uniform $L\to\infty$ estimate,
no finite-dimensional renormalizability of Einstein gravity, no asymptotic
safety, no nonperturbative ESM measure and no removal of the retained positive
IR/reference scale are claimed.
\end{theorem}

\begin{proof}
Item (1) is Theorem~\ref{thm:v6-one-shell-normal-form}; item (2) is Theorem
\ref{thm:v6-composed-transport} and Lemma~\ref{lem:v6-near-identity-product};
item (3) is Theorem~\ref{thm:v6-state-Cauchy}; item (4) is Theorem
\ref{thm:v6-W-continuum} and Corollary~\ref{cor:v6-operator-mixing}; item (5)
is Theorem~\ref{thm:v6-master-limit}; item (6) is Theorem
\ref{thm:v6-volume-limit}.  All statements use only fixed finite
$(L,\Delta,r,N)$ and the same normalization/reference prescription.
\end{proof}

\begin{center}
\fbox{\parbox{0.92\textwidth}{
\textbf{E5 is closed at the stated fixed-order coefficientwise scope.}
The essential correction is that the continuum theorem is formulated in the
common renormalized local/source/line frame.  Raw bare transport is allowed to
carry logarithmic or powerlike local renormalization factors; after this local
cocycle is factored out, the genuinely quasilocal shell evolution is a
summable near-identity product.}}
\end{center}

\section{Why E6 is now the only remaining programme gate}
\label{sec:v6-E6-handoff}

E5 fixes one admissible normalization/reference prescription and constructs its
coefficientwise continuum.  The remaining questions are compatibility
questions, not new ultraviolet estimates.

Let $\mathcal S$ and $\mathcal S'$ be two admissible finite schemes/reference
choices on the same native regulator branch.  At every finite filtration they
are related by a finite triangular local/source/line reparametrization
$F^{\mathcal S'\leftarrow\mathcal S}_J$.  E6 must prove that the corresponding
renormalized continuum coordinate change has a cutoff-independent limit,
identify which coefficient functions transform covariantly and which
normalized relations are invariant, and show that the construction at loop
order $L+1$ restricts exactly to the already fixed order-$L$ continuum rather
than refitting its lower coefficients.

\begin{theorem}[Exact next load-bearing theorem for V7]
\label{thm:v6-next-E6}
For every fixed finite filtration, prove:
\begin{enumerate}[label=\textup{(G\arabic*)},leftmargin=3em]
\item any two admissible normalization/source/reference schemes are related by
a finite local triangular map whose renormalized cutoff sequence is Cauchy and
has an invertible continuum limit;
\item changes of determinant/Pfaffian local line frame and positive reference
normalization act by the already stored finite cocycle/contact rules and do not
alter normalized gauge-invariant connected relations;
\item the order-$(L+1)$ local bank, forest ownership, Hodge section and reduced
transport restrict to the previously fixed order-$L$ data, so the continuum
coefficients form a projective family in $L$;
\item after these compatibility statements, formulate the final theorem: for
every prescribed finite perturbative, EFT/derivative and source order, the
native Renewal Einstein--Standard Model regulator admits a source-complete
renormalized coefficientwise continuum satisfying the corresponding common
BV/BRST/master identities and Standard-Model anomaly constraints, with gravity
interpreted as an EFT and all retained IR/reference assumptions stated
explicitly.
\end{enumerate}
\end{theorem}

\section{V6 anti-circularity audit}
\label{sec:v6-audit}

\begin{enumerate}[leftmargin=2.4em]
\item \textbf{No false raw contraction hypothesis.}
Proposition~\ref{prop:v6-raw-obstruction} shows why marginal/source mixing can
make the bare composed derivative grow.  V6 factors the already-declared local
normalization cocycle before asking for summable transport.
\item \textbf{No hidden scheme comparison.}
Only one fixed normalization/reference prescription is used.  Different
schemes are not identified in V6; that is E6.
\item \textbf{No dropped remainder contacts.}
The $I+O(\epsilon_j)$ statement is made on the full V5
contraction-complete packet, not on a bare local jet plus an undifferentiated
remainder.
\item \textbf{No graphwise retuning.}
Purely local maximal-scale pieces are owned by the same common local
normalization/Hodge section.  Proper forest blocks still use the previously
fixed common lower-order actions.
\item \textbf{No inference from one-shell summability alone.}
The nonlinear product lemma and normalized composed-transport theorem are
proved explicitly; the highest-scale forest expansion supplies an independent
absolute-convergence check.
\item \textbf{No loss of source-dependent constants.}
$W=\log(Z[J]/Z[0])$ removes only source-free vacuum coordinates.
Field-independent but source-dependent contacts remain in the convergent state
and in the limiting connected functional.
\item \textbf{No unproved all-derivative master identity.}
The identity passed to the limit is exactly the finite coefficient/descendant
system closed in V5.  Larger finite derivative/source orders require the
corresponding larger finite construction.
\item \textbf{No conflation of UV and volume limits.}
The ultraviolet Cauchy theorem is proved first.  Thermodynamic comparison is a
separate uniform rooted-tail argument.
\end{enumerate}

\section{V6 status ledger}
\label{sec:v6-status}

\subsection*{Proved in V6}
\begin{enumerate}[leftmargin=2.2em]
\item The raw long-product derivative requested at the end of V5 is not an
intrinsic continuum criterion; finite triangular marginal/source mixing gives
an explicit counterexample.
\item A common normalization-complete local/source/line frame exists on the
fixed compact protected filtration, with uniformly bounded section and
adjacent-cutoff $O(\epsilon_j)$ variation.
\item In that frame the arbitrary-fixed-loop native shell map is
$I+O(\epsilon_j)$ in $C^1$ on the complete contraction-descendant state.
\item Summability of those $C^1$ increments gives a cutoff-uniform derivative
bound for the composed \emph{renormalized} transport and a nonlinear infinite
product.
\item The complete normalized local plus quasilocal packet is Cauchy, with
ultraviolet tail $O(\mu/\Lambda_J)$, and the limit is independent of the
cofinal cutoff sequence within the fixed native branch and fixed scheme.
\item Unfolding into scale-decorated forests gives the equivalent
highest-scale theorem: every maximal-scale contribution is either already
owned locally or carries one excess factor, and the remaining polynomial scale
multiplicity is dominated by the geometric shell decay.
\item Every coefficient of the normalized connected generating functional at
the prescribed finite loop/EFT/source order converges, including composite,
stress, contact, operator-mixing, ghost/antifield and determinant/Pfaffian line
insertions.
\item The complete finite vector of common master/source coefficients and V5
future-contact descendants passes continuously to the same limit and remains
zero.
\item Uniform weighted rooted moments give a separate thermodynamic
periodization bound, so the staged volume limit exists on the declared
nonwinding panels and commutes with the UV limit there.
\item Therefore E5 is closed at the stated fixed-order coefficientwise scope.
\end{enumerate}

\subsection*{Inherited}
\begin{enumerate}[leftmargin=2.2em]
\item V80's common bosonic normalization/forest architecture and finite local
cohomology range theorem.
\item The Standard-Model anomaly cancellations, protected finite chiral gauge
descent and exact determinant/Pfaffian shell cocycles.
\item V1--V2's native fermion functional calculus, hard pivot,
annular/full-zone comparison, source jets and line remainder estimates.
\item V3--V4's complete chiral two-loop shell self-map and common
canonical-pair/source-contact master audit.
\item V5's arbitrary fixed-loop forest theorem, finite local bank, common
lower-order subtraction ownership, contraction-complete packet and exact local
master restoration at every prescribed fixed loop order.
\end{enumerate}

\subsection*{Conditional}
\begin{enumerate}[leftmargin=2.2em]
\item The result remains native to the explicit V2 Wilson/overlap completion;
a different microscopic fermion kernel requires the V2 regulator-comparison
gate.
\item Normalization/Hodge statements are on compact constant-rank protected
charts.  A rank crossing requires a new chart and transition analysis.
\item Determinant/Pfaffian limits remain on the protected nonzero component
with the retained positive IR/reference prescription; no global line holonomy
or zero-mode crossing theorem is claimed.
\item The thermodynamic theorem is on the nonwinding finite rooted/source
panels controlled by the inherited weighted moments.  A massless global IR
limit is not included.
\item Constants may depend arbitrarily badly on $L,\Delta,r,N$; no all-order
uniformity or perturbative convergence is asserted.
\end{enumerate}

\subsection*{Open}
\begin{enumerate}[leftmargin=2.2em]
\item E6: compare admissible finite renormalization/source/reference schemes and
identify the invariant normalized relations.
\item Prove order compatibility: the order-$(L+1)$ construction must restrict
to the already fixed order-$L$ normalized continuum coefficients, including
source/contact and line packets.
\item Assemble these compatibility statements into the final coefficientwise
Einstein--Standard Model continuum theorem, with every retained IR/reference
assumption explicit.
\end{enumerate}

\subsection*{Exact next load-bearing theorem}
\begin{center}
\fbox{\parbox{0.92\textwidth}{
\textbf{E6: scheme/reference and order compatibility.}
Fix two admissible finite normalization/source/reference choices.  Prove that
their finite triangular relation converges to an invertible continuum
reparametrization, characterize the normalized connected relations invariant
under that change, include determinant/Pfaffian frame/reference cocycles, and
prove projective compatibility in perturbative order: the $(L+1)$ continuum
must restrict to the already fixed $L$ continuum without any lower-order
refit.  Then state the final source-complete coefficientwise Renewal ESM
continuum theorem with gravity as an EFT and a retained positive IR/reference
scale.}}
\end{center}

\section*{Append-only continuation marker: V6}
\addcontentsline{toc}{section}{Append-only continuation marker: V6}
\textbf{End of V6.}  The complete V1--V5 body above is retained verbatim.  V6
closes E5 by correcting the raw transport quantifier, factoring the common
finite-dimensional local/source/line renormalization cocycle, proving that the
reduced shell maps form a summable $C^1$ near-identity product, establishing a
Cauchy continuum for the complete contraction-descendant state and normalized
connected coefficients, passing the finite common master/source identities to
the same limit, and separating the thermodynamic comparison from the UV
removal.  Future versions must append after this marker.  The next file is to
be named
\texttt{Renewal\_Perturbative\_ESM\_Continuum\_Research\_Notes\_V7.tex}.



% ============================================================================
% V7 APPEND-ONLY CONTINUATION
% ============================================================================
\clearpage
\part{V7 append-only continuation:\\
scheme/reference covariance, projective order compatibility, and the final coefficientwise continuum theorem}
\label{part:v7}

\section{Purpose and the final compatibility problem}
\label{sec:v7-purpose}

V6 closed E5 for one fixed admissible normalization/source/reference
prescription.  Thus no new ultraviolet graph estimate is load bearing for the
stated programme.  E6 is a compatibility problem: different admissible finite
coordinate choices must describe the same coefficientwise continuum up to an
invertible finite reparametrization, determinant/Pfaffian frame changes must be
carried with their source contacts rather than discarded, and enlarging the
perturbative order from $L$ to $L+1$ must not refit any coefficient already
fixed at orders $0,\ldots,L$.

We retain throughout the V6 hypotheses: the explicit V2 Wilson/overlap branch,
a protected constant-rank perturbative component, a protected nonzero
 determinant/Pfaffian line component, a positive reference/IR scale $\mu>0$,
a fixed finite external/source panel, and prescribed finite
$(L,\Delta,r,N)$.  Thermodynamic comparison remains the separate rooted-tail
limit of V6.  No $\mu\downarrow0$ statement is made.

The first point is an anti-circularity correction.  V5 used a finite Hodge
inverse at each fixed order.  If one were to recompute an orthogonal Hodge
inverse from scratch whenever the local bank is enlarged, exact order
compatibility would not follow.  E6 therefore fixes the chronology of the
finite choices rather than attempting to prove a false functoriality property
of Moore--Penrose/Hodge inverses.

\section{Why independently recomputed Hodge sections are not projective}
\label{sec:v7-hodge-obstruction}

\begin{proposition}[Orthogonal Hodge inversion is not functorial under bank enlargement]
\label{prop:v7-hodge-not-functorial}
There are finite complexes $C\subset \widetilde C$ for which the orthogonal
contracting homotopy on $\widetilde C$, formed using the enlarged inner
product, does not restrict to the orthogonal contracting homotopy on $C$.
Consequently exact order compatibility cannot be deduced merely by choosing a
new Moore--Penrose/Hodge inverse independently at every loop order.
\end{proposition}

\begin{proof}
Let the small complex have
\[
 C^0=\operatorname{span}\{e_1\},\qquad
 C^1=\operatorname{span}\{f\},\qquad d e_1=f,
\]
with the standard inner product.  Its degree-one contracting homotopy is
$h_C f=e_1$.
Enlarge degree zero to
\[
 \widetilde C^0=\operatorname{span}\{e_1,e_2\},\qquad
 \widetilde C^1=\operatorname{span}\{f\},\qquad
 d e_1=d e_2=f,
\]
again with the standard inner product.  The Moore--Penrose inverse of the row
matrix $(1\;1)$ sends
\[
 f\longmapsto \frac12(e_1+e_2).
\]
Hence $h_{\widetilde C}f\notin C^0$ and in particular
$h_{\widetilde C}|_{C^1}\ne h_C$.  The differential on the original
subcomplex has not changed; only the ambient orthogonality problem has.
\end{proof}

\begin{firewall}[Upstream correction for E6]
From V7 onward, ``the Hodge section'' means a \emph{chronologically fixed
finite contracting section}: once the order-$q$ primitive and its
normalization have been selected, they are frozen.  Enlarging the target order
adds only the new order and the additional future-descendant coordinates.  An
orthogonal Hodge inverse may be used to choose the new complement, but it is
never allowed to recompute an already fixed lower-order primitive.
\end{firewall}

\section{A nested ambient bank and exact perturbative truncations}
\label{sec:v7-nested-bank}

Fix $(\Delta,r,N)$ and the finite external/source panel.  Instead of allowing
the operator basis itself to depend on the terminal loop order, choose once the
finite ambient local operator/source/line bank
\begin{equation}
 \mathcal A_{\Delta,r,N}
 \label{eq:v7-ambient-bank}
\end{equation}
containing every ghost-zero EFT/Wilson coordinate through derivative degree
$\Delta$, every required BV field--antifield and nonminimal coordinate, every
source/contact coordinate through order $r$, and the protected local
 determinant/Pfaffian line/reference coordinates.  The loop filtration is then
formal:
\begin{equation}
 \mathcal A^{[\le L]}_{\Delta,r,N}
 :=\bigoplus_{q=0}^{L}\hbar^q\mathcal A_{\Delta,r,N}.
 \label{eq:v7-loop-filtered-bank}
\end{equation}
Coordinates which are not generated at a particular loop order simply have
zero coefficient there; the basis is not rediscovered when $L$ changes.

For the V5 contraction-complete remainder packet, let
$\mathcal Q_q^{[d]}$ denote the order-$q$ rooted packet stored to future
contraction depth $d$.  Define
\begin{equation}
 \mathfrak C_L
 :=\mathcal A^{[\le L]}_{\Delta,r,N}
   \oplus\bigoplus_{q=0}^{L}\hbar^q\mathcal Q_q^{[L-q]}.
 \label{eq:v7-projective-state}
\end{equation}
There are canonical inclusion and truncation maps
\begin{equation}
 \iota_L^{L+1}:\mathfrak C_L\hookrightarrow\mathfrak C_{L+1},
 \qquad
 \pi_L^{L+1}:\mathfrak C_{L+1}\twoheadrightarrow\mathfrak C_L,
 \qquad
 \pi_L^{L+1}\iota_L^{L+1}=I,
 \label{eq:v7-projections}
\end{equation}
where $\pi_L^{L+1}$ discards the $\hbar^{L+1}$ layer and discards exactly one
extra future-descendant depth from each older packet.

\begin{definition}[Coherent normalization/reference family]
\label{def:v7-coherent-scheme}
An admissible scheme family
$\mathcal S=\{\mathcal S_L\}_{L\ge0}$ is \emph{coherent} if:
\begin{enumerate}[label=\textup{(C\arabic*)},leftmargin=3em]
\item the physical/EFT normalization functionals at order $q$ are fixed when
$q$ first appears and are unchanged at every terminal order $L\ge q$;
\item the source/composite basis and its triangular mixing at source order
$\le r$ restrict under \eqref{eq:v7-projections};
\item every proper forest subtraction uses the already fixed common lower-loop
action coefficient, including its finite part;
\item the order-$q$ BV restoring primitive is frozen after it is constructed;
new Hodge/contracting choices act only on the new order-$L+1$ complement;
\item determinant/Pfaffian frame and positive reference normalizations are
carried with their complete local source/contact coordinates and restrict under
$\pi_L^{L+1}$.
\end{enumerate}
\end{definition}

\begin{lemma}[Every fixed-order admissible choice extends to a coherent family]
\label{lem:v7-coherent-extension}
Let an admissible normalization/reference prescription be fixed through order
$L_0$.  On the protected constant-rank component it can be extended
inductively to a coherent family at every higher prescribed finite order.
\end{lemma}

\begin{proof}
Assume the family has been fixed through order $L$.  The new local coefficient
at order $L+1$ lies in the finite bank \eqref{eq:v7-ambient-bank}.  Keep all
normalization functionals and restoring primitives at orders $q\le L$
unchanged.  V5 gives an exact local cocycle for the new order-$L+1$ breaking,
and its Standard-Model anomaly projection vanishes.  Therefore the new
breaking lies in the finite range of the classical differential.  Choose any
bounded right inverse on a complement to the already frozen lower-order data;
for example, choose an orthogonal complement only inside this new finite
layer.  Fix the independent order-$L+1$ physical/EFT coordinates by the new
normalization conditions and append the additional future-descendant depth.
All operations are finite dimensional at the prescribed filtration and
constant rank gives the required bounded inverses.  This extends the coherent
family without modifying any old coefficient.
\end{proof}

\section{Exact finite-cutoff order compatibility}
\label{sec:v7-finite-projective}

\begin{theorem}[Forest preparation commutes with perturbative truncation]
\label{thm:v7-forest-projective}
For every native shell $j$ and every $L\ge0$, the coherent forest-prepared
shell functional satisfies
\begin{equation}
 \boxed{
 \pi_L^{L+1}\,
 \mathbf R_{j}^{[\le L+1]}\,
 \iota_L^{L+1}
 =\mathbf R_{j}^{[\le L]}.}
 \label{eq:v7-forest-projective}
\end{equation}
This identity includes all retained source/contact, BV, and
 determinant/Pfaffian line descendants.
\end{theorem}

\begin{proof}
Take a graph contributing to a coefficient of loop order $q\le L$.  Loop
number is additive under contraction of a proper 1PI subtraction block, and
V5 proved that every proper subtraction block has strictly smaller loop number
than its parent graph.  Therefore every subtraction required by the order-$q$
graph is owned by action coefficients of orders $<q\le L$ and is already
frozen in the coherent family.  A counterterm first appearing at order
$L+1$ cannot enter a coefficient of order $q\le L$ because it carries an
explicit factor $\hbar^{L+1}$.

The local Taylor projector depends only on the fixed ambient derivative/source
bank and hence is the same on both sides after truncation.  Source partitions,
canonical BV contractions and line derivatives preserve the loop grading.
The V5 future-descendant packet at terminal order $L+1$ stores one additional
contraction depth, but applying $\pi_L^{L+1}$ discards precisely that extra
depth and leaves the packet previously stored at terminal order $L$.  Thus
every forest term, local assignment and retained complement at order $q\le L$
agrees term by term, proving \eqref{eq:v7-forest-projective}.
\end{proof}

\begin{theorem}[The common restoring action is projective in loop order]
\label{thm:v7-restorer-projective}
Let
\[
 S_{j}^{[\le L]}=\sum_{q=0}^{L}\hbar^q\bar S_{j,q}
\]
be the coherent common action.  Then the order-$(L+1)$ construction satisfies
\begin{equation}
 \boxed{
 \pi_L^{L+1}S_{j}^{[\le L+1]}=S_{j}^{[\le L]}.}
 \label{eq:v7-action-projective}
\end{equation}
In particular no lower-loop physical coefficient, finite forest part,
BV-exact restorer, source/contact descendant, or line/reference coordinate is
refitted when the terminal perturbative order is increased.
\end{theorem}

\begin{proof}
By Definition~\ref{def:v7-coherent-scheme}, every coefficient
$\bar S_{j,q}$ with $q\le L$ was frozen when order $q$ was first closed.
Theorem~\ref{thm:v7-forest-projective} shows that recomputing the larger forest
produces exactly the same order-$q$ pre-action and lower-order descendant data.
The chronological contracting section therefore returns the already stored
order-$q$ primitive rather than an independently recomputed Hodge primitive.
The only new local action is multiplied by $\hbar^{L+1}$ and is killed by
$\pi_L^{L+1}$.  This gives \eqref{eq:v7-action-projective}.
\end{proof}

\begin{corollary}[Exact shell-map projectivity]
\label{cor:v7-shell-projective}
For every shell $j$,
\begin{equation}
 \boxed{
 \pi_L^{L+1}\,
 \mathcal R_{j}^{[\le L+1]}\,
 \iota_L^{L+1}
 =\mathcal R_{j}^{[\le L]}.}
 \label{eq:v7-shell-projective}
\end{equation}
\end{corollary}

\begin{proof}
The shell map consists of the forest-prepared local and quasilocal update, the
common restoring action, and the stored contraction-descendant transport.
Apply Theorems~\ref{thm:v7-forest-projective} and
\ref{thm:v7-restorer-projective} to these three pieces.
\end{proof}

\section{Projective normalized frames and reduced transports}
\label{sec:v7-normalized-projective}

For a coherent scheme $\mathcal S$, choose the V6 normalization map
$\Theta_{j,L}^{\mathcal S}$ recursively: on the lower-order image of
$\iota_{L-1}^{L}$ it is the already fixed map
$\Theta_{j,L-1}^{\mathcal S}$, while only the new order-$L$ and extra packet
coordinates are normalized afresh.

\begin{lemma}[Normalized-frame projectivity]
\label{lem:v7-theta-projective}
The coherent normalization maps may be chosen so that
\begin{equation}
 \boxed{
 \pi_L^{L+1}\Theta_{j,L+1}^{\mathcal S}
 =\Theta_{j,L}^{\mathcal S}\pi_L^{L+1}.}
 \label{eq:v7-theta-projective}
\end{equation}
The same identity holds for their local inverses on the common protected
compact chart.
\end{lemma}

\begin{proof}
This is the recursive construction just described.  On old coordinates the
map is definitionally unchanged.  The new block is finite triangular with an
invertible diagonal quotient block by admissibility.  Hence the full map and
its local inverse are block triangular, with the old map as their upper-left
block.  Projection onto that block gives \eqref{eq:v7-theta-projective} and its
inverse analogue.
\end{proof}

\begin{theorem}[Reduced shell maps form a projective system]
\label{thm:v7-reduced-projective}
Let
\[
 \widehat{\mathcal R}_{j,L}^{\mathcal S}
 =\Theta_{j-1,L}^{\mathcal S}
  \mathcal R_{j}^{[\le L]}
  (\Theta_{j,L}^{\mathcal S})^{-1}.
\]
Then
\begin{equation}
 \boxed{
 \pi_L^{L+1}\widehat{\mathcal R}_{j,L+1}^{\mathcal S}
 \iota_L^{L+1}
 =\widehat{\mathcal R}_{j,L}^{\mathcal S}.}
 \label{eq:v7-reduced-projective}
\end{equation}
\end{theorem}

\begin{proof}
Insert the definition, commute $\pi_L^{L+1}$ through the two normalization
maps using Lemma~\ref{lem:v7-theta-projective}, and use
Corollary~\ref{cor:v7-shell-projective}.
\end{proof}

\section{Finite scheme changes as exact intertwiners}
\label{sec:v7-scheme-intertwiner}

We now compare two coherent admissible schemes $\mathcal S$ and $\mathcal S'$
on the same native regulator branch and protected chart.  ``Admissible'' here
means that both use complete local/source/line banks at the prescribed
filtration, their normalization Jacobians are uniformly nonsingular on the
chosen compact set, their BV-sector coordinate changes are canonical with the
corresponding density/reference transformation, and their source/composite
changes are finite triangular maps retaining all contact coordinates.

Define the same-cutoff transition
\begin{equation}
 F_{j,L}^{\mathcal S'\leftarrow\mathcal S}
 :=\Theta_{j,L}^{\mathcal S'}
   (\Theta_{j,L}^{\mathcal S})^{-1}.
 \label{eq:v7-Fj}
\end{equation}
It is a finite local/source/line map.  On a sufficiently small common compact
chart it and its inverse have uniformly bounded $C^1$ norm because both
normalization maps have uniformly nonsingular finite Jacobians.

\begin{theorem}[Exact shell intertwining of two schemes]
\label{thm:v7-intertwiner}
For every shell,
\begin{equation}
 \boxed{
 \widehat{\mathcal R}_{j,L}^{\mathcal S'}
 F_{j,L}^{\mathcal S'\leftarrow\mathcal S}
 =F_{j-1,L}^{\mathcal S'\leftarrow\mathcal S}
  \widehat{\mathcal R}_{j,L}^{\mathcal S}.}
 \label{eq:v7-intertwiner}
\end{equation}
\end{theorem}

\begin{proof}
Using the definitions and cancelling adjacent normalization maps,
\begin{align*}
 \widehat{\mathcal R}_{j,L}^{\mathcal S'}F_{j,L}
 &=\Theta_{j-1,L}^{\mathcal S'}
   \mathcal R_j
   (\Theta_{j,L}^{\mathcal S'})^{-1}
   \Theta_{j,L}^{\mathcal S'}
   (\Theta_{j,L}^{\mathcal S})^{-1}\\
 &=\Theta_{j-1,L}^{\mathcal S'}
   \mathcal R_j(\Theta_{j,L}^{\mathcal S})^{-1}\\
 &=F_{j-1,L}
   \Theta_{j-1,L}^{\mathcal S}\mathcal R_j
   (\Theta_{j,L}^{\mathcal S})^{-1},
\end{align*}
which is \eqref{eq:v7-intertwiner}.
\end{proof}

\begin{lemma}[Summable adjacent variation of the finite scheme map]
\label{lem:v7-F-adjacent}
At fixed finite $(L,\Delta,r,N)$,
\begin{equation}
 \boxed{
 \|F_{j-1,L}-F_{j,L}\|_{C^1}
 \le C_{L,\Delta,r,N}\epsilon_j.}
 \label{eq:v7-F-adjacent}
\end{equation}
\end{lemma}

\begin{proof}
Write the two V6 reduced shell maps as
$\widehat{\mathcal R}^{\mathcal S}_j=I+E_j$ and
$\widehat{\mathcal R}^{\mathcal S'}_j=I+E'_j$, with
$\|E_j\|_{C^1}+\|E'_j\|_{C^1}\le C\epsilon_j$.
Theorem~\ref{thm:v7-intertwiner} gives
\[
 F_{j-1,L}=(I+E'_j)F_{j,L}(I+E_j)^{-1}.
\]
On the common compact chart the maps $F_{j,L}$ and their inverses are uniformly
$C^1$ bounded, while
$(I+E_j)^{-1}=I+O_{C^1}(\epsilon_j)$ for all sufficiently large $j$.
Subtract $F_{j,L}$ and expand the finite composition.  Every resulting term
contains at least one factor $E_j$ or $E'_j$, proving
\eqref{eq:v7-F-adjacent}.  Finitely many low shells are absorbed into the
constant.
\end{proof}

\begin{theorem}[Invertible continuum scheme transition]
\label{thm:v7-F-limit}
The finite scheme maps have a cutoff-independent $C^1$ limit
\begin{equation}
 \boxed{
 F_{\infty,L}^{\mathcal S'\leftarrow\mathcal S}
 :=\lim_{j\to\infty}F_{j,L}^{\mathcal S'\leftarrow\mathcal S},}
 \label{eq:v7-F-limit}
\end{equation}
with tail
\begin{equation}
 \|F_{\infty,L}-F_{j,L}\|_{C^1}
 \le C_{L,\Delta,r,N}
 \sum_{k>j}\epsilon_k
 \le C'_{L,\Delta,r,N}\frac{\mu}{\Lambda_j}.
 \label{eq:v7-F-tail}
\end{equation}
The limit is locally invertible, and
\begin{equation}
 (F_{\infty,L}^{\mathcal S'\leftarrow\mathcal S})^{-1}
 =F_{\infty,L}^{\mathcal S\leftarrow\mathcal S'}.
 \label{eq:v7-F-inverse}
\end{equation}
\end{theorem}

\begin{proof}
Lemma~\ref{lem:v7-F-adjacent} and
$\sum_j\epsilon_j<\infty$ make $\{F_{j,L}\}$ a $C^1$ Cauchy sequence.
Apply the same argument to the inverse transitions
$G_{j,L}:=F_{j,L}^{\mathcal S\leftarrow\mathcal S'}$.  At every finite cutoff,
$G_{j,L}F_{j,L}=I=F_{j,L}G_{j,L}$.  Passing to the uniform $C^1$ limits gives
$G_{\infty,L}F_{\infty,L}=I=F_{\infty,L}G_{\infty,L}$ on the common local
chart.  The geometric tail bound is the same sum used in V6.
\end{proof}

\begin{corollary}[Cocycle law for continuum scheme transitions]
\label{cor:v7-scheme-cocycle}
For three coherent admissible schemes $\mathcal S,\mathcal S',\mathcal S''$,
\begin{equation}
 F_{\infty,L}^{\mathcal S''\leftarrow\mathcal S'}
 F_{\infty,L}^{\mathcal S'\leftarrow\mathcal S}
 =F_{\infty,L}^{\mathcal S''\leftarrow\mathcal S}.
 \label{eq:v7-scheme-cocycle}
\end{equation}
\end{corollary}

\begin{proof}
The identity holds exactly at every finite cutoff because it is cancellation of
three coordinate identifications.  Pass to the $C^1$ limit.
\end{proof}

\section{What is scheme dependent and what is invariant}
\label{sec:v7-scheme-invariants}

The continuum state is therefore not a preferred numerical coordinate vector;
it is an equivalence class under the finite maps
$F_{\infty,L}$.  It is useful to state the invariant content precisely.

\begin{definition}[Covariant continuum observable]
\label{def:v7-covariant-observable}
A family of coefficient functions $\mathcal O^{\mathcal S}$ is covariant if
for every pair of admissible schemes
\begin{equation}
 \mathcal O^{\mathcal S'}\circ
 F_{\infty,L}^{\mathcal S'\leftarrow\mathcal S}
 =\mathcal O^{\mathcal S}.
 \label{eq:v7-covariant-observable}
\end{equation}
A relation is scheme invariant if its zero set is carried into itself by all
such transitions.
\end{definition}

\begin{theorem}[Finite-scheme covariance of the continuum coefficient system]
\label{thm:v7-scheme-covariance}
At fixed finite filtration the following hold.
\begin{enumerate}[label=\textup{(SC\arabic*)},leftmargin=3.4em]
\item Independent local EFT/Wilson coefficients, composite-operator basis
coordinates, finite stress improvements, redundant/BV-exact coordinates,
coincident contact coefficients and scalar line/reference coordinates are in
general scheme dependent.
\item The complete normalized connected functional is a scalar after the
source/composite coordinates and local contact polynomial are transformed by
$F_{\infty,L}$.  Hence connected coefficients of matched renormalized
observables transform covariantly.
\item The common BV/BRST/master vector transforms by the induced canonical
ghost-one coordinate map.  Therefore its vanishing, the Standard-Model anomaly
cancellation, and every source/contact descendant relation proved in V5--V6
are scheme invariant.
\item Any separated-support connected relation, after the operator/source
basis is matched covariantly, is insensitive to a change consisting only of a
finite local source/contact polynomial.  Raw coincident contact coefficients
need not agree.
\end{enumerate}
\end{theorem}

\begin{proof}
Item (SC1) is the definition of a finite change of local coordinates: there is
no reason for coordinate components in two charts to agree numerically.
For (SC2), at finite cutoff both schemes are obtained from the same native
partition function by invertible source/local coordinate changes; the full
contact packet is included among those coordinates.  Thus the scalar
coordinate-free connected functional is unchanged, while its component array
is transformed by substitution and the finite local contact term.  Passing to
the V6 coefficientwise limit and using Theorem~\ref{thm:v7-F-limit} proves the
continuum statement.

For (SC3), admissibility requires the BV-sector transition to be canonical and
the finite BV density/reference term to be transformed simultaneously.  The
master functional therefore transforms covariantly at finite cutoff.  Its
retained coefficient vector is zero in one scheme if and only if it is zero in
the other.  The same is true after any retained source derivative because the
source/contact map is part of $F_{j,L}$.  Pass to the limit.  The anomaly
quotient is a cohomological coordinate-free obstruction, so its zero class is
unchanged by finite basis changes.

For (SC4), the difference generated by a local source polynomial is supported
on coincident source points after differentiation.  If the test/source
supports are pairwise separated, those diagonal distributions vanish.  A
change of composite basis is still applied contragrediently to the inserted
observables, which is precisely the matching in
\eqref{eq:v7-covariant-observable}.
\end{proof}

\section{Determinant/Pfaffian frame and positive-reference compatibility}
\label{sec:v7-line-frame}

The line sector deserves a separate statement because a scalar representative
of a determinant section is not itself coordinate free.

Let $\lambda$ and $\lambda'$ be two admissible nonzero local determinant-line
frames on the same protected chart and write
\begin{equation}
 \lambda'=e^{\chi}\lambda,
 \label{eq:v7-line-frame-change}
\end{equation}
where $\chi$ is a smooth local source function; for two unit frames,
$\Re\chi=0$.  If
$\sigma=s\lambda=s'\lambda'$, then
\begin{equation}
 \boxed{\log s'=\log s-\chi.}
 \label{eq:v7-line-scalar-change}
\end{equation}
For an oriented nonzero Pfaffian line the same formula holds on the chosen
orientation chart, with the inherited factor-$1/2$ relation to the determinant
logarithm where applicable.

\begin{theorem}[Frame changes are exact local coboundaries of the shell line cocycle]
\label{thm:v7-line-coboundary}
Under an admissible local line-frame change, the coordinate-free exact Schur
shell cocycle is unchanged.  Its scalar logarithmic representatives differ by
an endpoint coboundary.  Consequently a finite sequence of shell frame changes
contributes only the difference of the endpoint frame logs, together with all
of their retained source derivatives; no mixed-shell contact is lost.
\end{theorem}

\begin{proof}
The Schur factorization is an isomorphism of determinant lines before a frame
is selected.  Therefore changing frames cannot alter the coordinate-free
section.  Scalarizing the line identity in the two frames and using
\eqref{eq:v7-line-scalar-change} shows that the logarithmic shell factor changes
by the difference of the frame logs at its two endpoints.  Summing over nested
shells telescopes these differences.  Differentiation uses the ordinary
Leibniz/Faa--di--Bruno rule already encoded in the V1 Bell-polynomial line
packet, so every derivative of the endpoint logs and every mixed contact is
retained rather than discarded.
\end{proof}

Let a positive local reference normalization be changed by
\begin{equation}
 \rho'=e^{u}\rho,
 \label{eq:v7-positive-reference-change}
\end{equation}
with real smooth local source function $u$.  Source-free $u(0)$ disappears from
$W[J]=\log Z[J]-\log Z[0]$, but the difference $u(J)-u(0)$ is a genuine local
source/contact coordinate.

\begin{corollary}[Normalized line/reference functional]
\label{cor:v7-line-reference-W}
A simultaneous admissible line-frame/reference change alters a scalar
representative of the normalized connected functional only by the finite local
source/contact expression prescribed by
$\chi(J)-\chi(0)$ and $u(J)-u(0)$, with signs determined by the declared scalar
line convention.  After those contact coordinates are transformed as part of
$F_{\infty,L}$, the coordinate-free normalized determinant/Pfaffian insertion
relations are unchanged.
\end{corollary}

\begin{proof}
Apply Theorem~\ref{thm:v7-line-coboundary} to the determinant/Pfaffian factor
and subtract the $J=0$ normalization.  A source-free constant cancels, whereas
a source-dependent local function does not.  V1--V3 explicitly store its
source derivatives in the line/reference packet.  The continuum transition of
Theorem~\ref{thm:v7-F-limit} transports precisely those finite coordinates.
\end{proof}

\begin{firewall}[No global line-holonomy conclusion]
Theorems~\ref{thm:v7-line-coboundary}--\ref{cor:v7-line-reference-W} compare
local nonzero frames on the protected component.  They do not trivialize the
line bundle across zero crossings or large loops and do not prove global
 determinant/Pfaffian holonomy cancellation beyond the finite gauge-descent
statement already inherited.
\end{firewall}

\section{Projective compatibility of continuum limits}
\label{sec:v7-continuum-projective}

For a coherent scheme $\mathcal S$, let
$\widehat{\mathfrak C}_{L,j_0}^{(J),\mathcal S}$ be the V6 normalized state at
reference shell $j_0$ obtained with UV cutoff $J$, and let
$\widehat{\mathfrak C}_{L,j_0}^{(\infty),\mathcal S}$ denote its coefficientwise
limit.

\begin{theorem}[Finite-cutoff states are exactly projective]
\label{thm:v7-finite-state-projective}
For every $J$ and $j_0\le J$,
\begin{equation}
 \boxed{
 \pi_L^{L+1}
 \widehat{\mathfrak C}_{L+1,j_0}^{(J),\mathcal S}
 =
 \widehat{\mathfrak C}_{L,j_0}^{(J),\mathcal S}.}
 \label{eq:v7-finite-state-projective}
\end{equation}
\end{theorem}

\begin{proof}
The normalized UV boundary data are coherent by Definition
\ref{def:v7-coherent-scheme}.  Compose the reduced shell maps down to $j_0$.
Theorem~\ref{thm:v7-reduced-projective} lets $\pi_L^{L+1}$ pass through each
factor and turns every order-$(L+1)$ factor into its order-$L$ factor.  The
projected UV boundary is the order-$L$ boundary.  This proves the identity.
\end{proof}

\begin{theorem}[Exact projective compatibility in perturbative order]
\label{thm:v7-continuum-projective}
For every coherent admissible scheme,
\begin{equation}
 \boxed{
 \pi_L^{L+1}
 \widehat{\mathfrak C}_{L+1,j_0}^{(\infty),\mathcal S}
 =
 \widehat{\mathfrak C}_{L,j_0}^{(\infty),\mathcal S}.}
 \label{eq:v7-continuum-projective}
\end{equation}
The same identity holds coefficientwise for the normalized connected
functional, every retained source/contact descendant, and the complete finite
BV/master vector.
\end{theorem}

\begin{proof}
The projection $\pi_L^{L+1}$ is a bounded finite-coordinate map.  Apply it to
the V6 Cauchy limit and use the exact finite-cutoff identity
\eqref{eq:v7-finite-state-projective}:
\[
 \pi_L^{L+1}
 \widehat{\mathfrak C}_{L+1}^{(\infty)}
 =\lim_{J\to\infty}
  \pi_L^{L+1}\widehat{\mathfrak C}_{L+1}^{(J)}
 =\lim_{J\to\infty}\widehat{\mathfrak C}_{L}^{(J)}
 =\widehat{\mathfrak C}_{L}^{(\infty)}.
\]
The connected/source/master coefficient maps are finite polynomial or analytic
maps of the retained state at every prescribed finite order, so they commute
with the same projection and limit.
\end{proof}

\begin{theorem}[Scheme transitions are themselves projective in loop order]
\label{thm:v7-F-projective}
For two coherent schemes,
\begin{equation}
 \boxed{
 \pi_L^{L+1}
 F_{\infty,L+1}^{\mathcal S'\leftarrow\mathcal S}
 \iota_L^{L+1}
 =F_{\infty,L}^{\mathcal S'\leftarrow\mathcal S}.}
 \label{eq:v7-F-projective}
\end{equation}
\end{theorem}

\begin{proof}
The same identity holds at every finite cutoff by
Lemma~\ref{lem:v7-theta-projective} applied to both schemes.  The $C^1$ limits
of Theorem~\ref{thm:v7-F-limit} commute with the bounded finite projection.
\end{proof}

\section{Reference-scale covariance at positive scale}
\label{sec:v7-reference-scale}

A change of subtraction/reference point inside a compact positive interval is
an admissible finite scheme change provided the corresponding normalization
Jacobians remain nonsingular and the protected line/pivot component is not
crossed.  Therefore the preceding theorems also compare two positive reference
scales $\mu,\mu'>0$ under those hypotheses.

\begin{corollary}[Positive reference-scale covariance]
\label{cor:v7-reference-scale}
Let $\mu$ and $\mu'$ belong to a compact interval
$0<\mu_-\le\mu,\mu'\le\mu_+<\infty$ on which the protected chart and finite
normalization ranks persist.  The fixed-order continua normalized at $\mu$ and
$\mu'$ are related by an invertible finite continuum map
$F_{\infty,L}^{\mu'\leftarrow\mu}$ satisfying the cocycle law in the reference
point.  This does not imply existence of the limit $\mu\downarrow0$.
\end{corollary}

\begin{proof}
Changing the positive reference point changes only the finite normalization
conditions and the resulting local/source/line coordinates.  On the compact
positive interval their Jacobians are uniformly nonsingular by hypothesis, so
they form an admissible scheme pair.  Apply Theorem~\ref{thm:v7-F-limit} and
Corollary~\ref{cor:v7-scheme-cocycle}.  No estimate is uniform as
$\mu_-\downarrow0$.
\end{proof}

\section{E6 closure theorem}
\label{sec:v7-E6-close}

\begin{theorem}[Scheme/reference and perturbative-order compatibility]
\label{thm:v7-E6}
Fix the native V2 Wilson/overlap branch, a protected constant-rank/nonzero-line
component, positive IR/reference scale, finite external/source panel, and
finite $(L,\Delta,r,N)$.  Then:
\begin{enumerate}[label=\textup{(E6.\arabic*)},leftmargin=3.5em]
\item any two coherent admissible finite normalization/source/reference
schemes are related by an invertible cutoff-independent continuum transition
$F_{\infty,L}$ with the finite cocycle law;
\item determinant/Pfaffian local frame changes and positive reference
normalizations act by the stored local endpoint/contact coboundaries and leave
the coordinate-free normalized line relations unchanged;
\item scheme-dependent local coefficient arrays transform covariantly, while
the vanishing common BV/BRST/master relations, Standard-Model anomaly-zero
class, and matched normalized connected relations are invariant;
\item the native shell maps, normalized reduced transports, continuum states,
connected/source/master coefficients, and continuum scheme transitions form a
projective family in perturbative order, with exact restriction from $L+1$ to
$L$ and no lower-order refit.
\end{enumerate}
\end{theorem}

\begin{proof}
Item (1) is Theorems~\ref{thm:v7-intertwiner}--\ref{thm:v7-F-limit} and
Corollary~\ref{cor:v7-scheme-cocycle}.  Item (2) is
Theorem~\ref{thm:v7-line-coboundary} and
Corollary~\ref{cor:v7-line-reference-W}.  Item (3) is
Theorem~\ref{thm:v7-scheme-covariance}.  Item (4) is
Corollary~\ref{cor:v7-shell-projective},
Theorem~\ref{thm:v7-reduced-projective},
Theorem~\ref{thm:v7-continuum-projective}, and
Theorem~\ref{thm:v7-F-projective}.
\end{proof}

\begin{center}
\fbox{\parbox{0.92\textwidth}{
\textbf{E6 is closed at the stated coefficientwise scope.}
The essential upstream correction is chronological/projective normalization:
finite Hodge inverses are not independently recomputed when the perturbative
bank is enlarged.  With lower-order data frozen, both the native shell theorem
and its continuum limit restrict exactly in loop order.}}
\end{center}

\section{Final source-complete coefficientwise Renewal ESM continuum theorem}
\label{sec:v7-final-theorem}

We can now state the endpoint promised at the beginning of this lineage.  The
quantifiers are important: $L$, $\Delta$, $r$, $N$ and the external/source
panel are fixed before the cutoff is removed.  Constants may depend
arbitrarily badly on these finite choices.

\begin{theorem}[Source-complete coefficientwise perturbative Einstein--Standard Model continuum]
\label{thm:v7-final-continuum}
Assume:
\begin{enumerate}[label=\textup{(H\arabic*)},leftmargin=3em]
\item the explicit V2 native Wilson/overlap fermion completion, or another
microscopic fermion regulator for which the V2 native regulator-comparison gate
has separately been proved;
\item the inherited finite Renewal gravitational/gauge/Higgs regulator and
exact BV-compatible high/low splitting on a protected constant-rank
perturbative component;
\item the three-generation Standard-Model chiral representation with the
inherited cancellation of the ordinary local gauge-anomaly trace classes and
even $SU(2)$ doublet parity;
\item a protected nonzero determinant/Pfaffian line component and a positive
IR/reference scale $\mu>0$;
\item one coherent admissible normalization/source/reference family.
\end{enumerate}
Then for every prescribed finite loop order $L$, EFT/derivative order
$\Delta$, source/contact order $r$, arity $N$, and finite external
momentum/source panel, there exists a common renormalized native action/source/
reference/line family whose normalized connected coefficients have a unique
ultraviolet-cutoff-independent limit.  More precisely:
\begin{enumerate}[label=\textup{(F\arabic*)},leftmargin=3em]
\item every loop-$\le L$ native graph and forest is renormalized by the same
previously fixed common lower-order actions and decomposes into a finite local
EFT/source jet plus a quasilocal summable remainder;
\item the limit includes composite operators, operator mixing, stress-tensor
insertions, coincident source contacts, ghosts/nonminimal variables,
antifields, and determinant/Pfaffian line insertions at the prescribed finite
orders;
\item all retained coefficients of the common coupled BV/BRST/master/source
system and all contraction descendants required through loop order $L$ vanish
in the continuum, with the nontrivial local Standard-Model anomaly class
absent;
\item the finite-volume/nonwinding coefficients admit the separate V6
thermodynamic limit, and the staged UV and volume limits commute on the
prescribed compact panels;
\item different coherent admissible finite schemes/reference choices are
related by the invertible continuum maps of Theorem~\ref{thm:v7-E6}; raw local
EFT, contact, composite-basis and line-frame coordinates may change, while the
matched coordinate-free normalized relations transform covariantly;
\item the family is exactly compatible in perturbative order:
\begin{equation}
 \pi_L^{L+1}
 \widehat{\mathfrak C}_{L+1}^{(\infty)}
 =\widehat{\mathfrak C}_{L}^{(\infty)}.
 \label{eq:v7-final-projective}
\end{equation}
\end{enumerate}
Gravity is interpreted throughout as an effective field theory: at every fixed
$(L,\Delta)$ the required local Wilson bank is finite, but no finite-dimensional
all-orders gravitational coupling space is asserted.
\end{theorem}

\begin{proof}
E1 is supplied by V1--V2: native chiral hard-shell functional calculus,
hard-pivot/annular/full-zone comparison, fixed source jets and the complete
 determinant/Pfaffian line remainder.  E2 is V3: the complete native
Standard-Model anomaly projection vanishes and one common source-extended
primitive enlarges the V80 bosonic shell to the chiral ESM shell.  E3 is V4:
one common first- and second-loop master/source equation contains every
canonical-pair cross term and every retained contact.  E4 is V5: arbitrary
fixed-loop forest preparation, contraction-complete packets, exact local
cocycles and one common order-$L$ restorer.  E5 is V6: after factoring the
finite local normalization cocycle, the reduced shell maps form a summable
near-identity product, giving the unique coefficientwise UV limit and the
separate thermodynamic comparison.  E6 is
Theorem~\ref{thm:v7-E6}: finite schemes/reference choices are related by
invertible continuum maps and the construction is projective in loop order.
Combining these statements yields (F1)--(F6).
\end{proof}

\begin{firewall}[What the final theorem does not say]
Theorem~\ref{thm:v7-final-continuum} is coefficientwise.  It does not assert:
\begin{enumerate}[leftmargin=2.4em]
\item convergence or Borel summability of the perturbative series as
$L\to\infty$;
\item any estimate uniform in loop order;
\item finite-dimensional renormalizability of Einstein gravity;
\item asymptotic safety or a nonperturbative UV fixed point;
\item existence of a nonperturbative quantum-gravity/ESM measure;
\item removal of the positive IR/reference scale or a global massless IR
limit;
\item determinant/Pfaffian triviality across zero crossings or arbitrary large
line-bundle loops;
\item universality with respect to a different microscopic fermion regulator
unless the V2 regulator-comparison gate is separately proved.
\end{enumerate}
\end{firewall}

\section{A projective formulation of the endpoint}
\label{sec:v7-projective-endpoint}

For each finite $L$ let
$\mathfrak C_L^{\infty,\mathcal S}$ be the continuum state supplied by
Theorem~\ref{thm:v7-final-continuum}.  Theorem
\ref{thm:v7-continuum-projective} makes
\begin{equation}
 \left\{
 \mathfrak C_L^{\infty,\mathcal S},
 \pi_L^{L+1}
 \right\}_{L\ge0}
 \label{eq:v7-projective-family}
\end{equation}
a projective family of finite-order coefficient systems.  This notation is
occasionally convenient, but it must not be misread as constructing the
inverse-limit point as a convergent numerical perturbative series.  The
projective family means only that every finite order is compatible with every
lower finite order.

\begin{corollary}[Order-independent meaning of a fixed coefficient]
\label{cor:v7-fixed-coefficient}
Fix a loop index $q$, derivative/source label in the chosen finite bank, and an
external/source panel.  For every terminal order $L\ge q$, the corresponding
continuum coefficient is exactly the same coefficient already fixed at order
$q$.
\end{corollary}

\begin{proof}
Iterate \eqref{eq:v7-final-projective} from $L$ down to $q$ and read the chosen
coordinate.
\end{proof}

\section{V7 anti-circularity audit}
\label{sec:v7-audit}

\begin{enumerate}[leftmargin=2.4em]
\item \textbf{No independently recomputed lower-order Hodge inverse.}
Proposition~\ref{prop:v7-hodge-not-functorial} exhibits the obstruction.
Lower-order primitives are frozen and only the new loop layer is solved.
\item \textbf{No $L$-dependent rediscovery of the EFT basis.}
The ambient operator/source/line bank is fixed at $(\Delta,r,N)$; loop order is
a formal grading.  This is what makes the truncation maps literal projections.
\item \textbf{No graphwise finite retuning at higher order.}
Theorem~\ref{thm:v7-forest-projective} uses V5's strict proper-subgraph loop
drop, so all proper finite subtractions are owned by already fixed lower-order
actions.
\item \textbf{No claim that different schemes have equal coordinate arrays.}
They are related by the finite invertible maps $F_{\infty,L}$.  Local Wilson
coefficients, contact terms and composite bases may differ.
\item \textbf{No discarded source-dependent frame/reference term.}
Only a source-free normalization constant cancels in $W[J]-W[0]$.
$\chi(J)-\chi(0)$ and $u(J)-u(0)$ remain explicit contact coordinates.
\item \textbf{No global determinant-line conclusion.}
Frame compatibility is local on the protected nonzero component.
\item \textbf{No hidden massless limit.}
Reference-scale covariance is proved only on compact positive intervals.
\item \textbf{No inverse-limit convergence claim.}
Projectivity in $L$ is exact compatibility of finite coefficients, not
convergence of the perturbative series.
\end{enumerate}

\section{V7 status ledger}
\label{sec:v7-status}

\subsection*{Proved in V7}
\begin{enumerate}[leftmargin=2.2em]
\item Independently recomputed orthogonal Hodge inverses are not projective in
general; a finite explicit counterexample identifies the missing hypothesis.
\item A chronological coherent normalization/restoration family can be
extended to every higher prescribed finite loop order without altering any
previous coefficient.
\item With a fixed ambient EFT/source/line bank and the V5
contraction-complete packets, forest preparation, common restoring actions and
native shell maps commute exactly with loop-order truncation.
\item The V6 normalized coordinate maps and reduced shell transports can be
chosen projectively in loop order.
\item Two coherent admissible schemes are exact finite-cutoff intertwiners.
Their adjacent-cutoff transition maps differ by $O(\mu/\Lambda_j)$, hence
converge in $C^1$ to an invertible continuum transition satisfying the cocycle
law.
\item Scheme-dependent local coefficient arrays and scheme-invariant/covariant
relations are separated explicitly.  The common master/anomaly-zero relations
are invariant; matched connected observables are covariant; raw coincident
contact coefficients need not agree.
\item Determinant/Pfaffian local frame changes are endpoint coboundaries of the
exact line cocycle.  Positive reference changes retain every source-dependent
contact; only source-free normalization cancels in $W[J]/W[0]$.
\item Finite-cutoff normalized states and their V6 continuum limits are exactly
projective in perturbative order, as are the continuum scheme transitions.
\item Positive reference points in one compact protected interval are related
by finite invertible continuum maps, without claiming a $\mu\downarrow0$
limit.
\item E6 and therefore the stated E1--E6 coefficientwise programme are closed
on the declared native branch and assumptions.
\end{enumerate}

\subsection*{Inherited}
\begin{enumerate}[leftmargin=2.2em]
\item V80's bosonic Einstein--gauge--Higgs two-loop local/self-map theorem and
finite local BV range architecture.
\item The Standard-Model representation, cancellation of the five ordinary
local gauge-anomaly coefficients, even $SU(2)$ doublet parity, protected gauge
descent, and exact determinant/Pfaffian scale cocycles.
\item V1--V2's native chiral functional calculus, local Kato frame, hard pivot,
explicit Wilson/overlap completion, annular/full-zone comparison, source jets
and line remainder bounds.
\item V3's complete chiral anomaly-quotient zero and common ESM two-loop shell
restorer; V4's full common two-loop master/source audit.
\item V5's arbitrary-fixed-loop native forest theorem,
contraction-complete packet and fixed-order common master restoration.
\item V6's normalized multi-shell $C^1$ product, coefficientwise UV Cauchy
theorem, connected/source/master continuum, and separate thermodynamic
comparison.
\end{enumerate}

\subsection*{Conditional}
\begin{enumerate}[leftmargin=2.2em]
\item A microscopic fermion kernel different from the explicit V2 completion
still requires its own native regulator-comparison theorem before this closure
transfers to it.
\item Constant-rank/Hodge and normalization arguments remain local to compact
protected charts.  Rank changes require compatible chart transitions.
\item Determinant/Pfaffian results remain on a nonzero locally orientable
component and do not cross zeros.
\item The positive reference-scale comparison requires a compact interval
bounded away from zero; a massless IR limit is not proved.
\item Constants can depend arbitrarily on $(L,\Delta,r,N)$ and on the protected
chart; no all-orders uniformity is available.
\end{enumerate}

\subsection*{Open beyond the stated programme}
\begin{enumerate}[leftmargin=2.2em]
\item A native universality theorem comparing the explicit V2 fermion
completion with a genuinely different microscopic chiral regulator.
\item Removal of the retained positive IR/reference scale, if desired, with
uniform control of massless long-distance sectors.
\item Global determinant/Pfaffian line holonomy and zero-crossing analysis.
\item Any analytic finite-coupling RG trajectory, all-orders summability,
asymptotic-safety or nonperturbative quantum-gravity/ESM completion theorem.
These are strictly stronger problems and are not required for the
coefficientwise endpoint proved here.
\end{enumerate}

\subsection*{Exact next load-bearing theorem}
\begin{center}
\fbox{\parbox{0.92\textwidth}{
\textbf{None inside the stated E1--E6 coefficientwise programme.}
The source-complete fixed-order continuum, scheme/reference covariance and
projective compatibility in loop order are now closed under the declared
native-regulator, protected-chart and positive-IR assumptions.  Any further
attack should be labelled as a stronger programme.  The smallest natural one
is a regulator-universality theorem comparing the explicit V2
Wilson/overlap completion with an alternative microscopic chiral regulator;
IR removal and global line holonomy are independent stronger directions.}}
\end{center}

\section*{Append-only continuation marker: V7}
\addcontentsline{toc}{section}{Append-only continuation marker: V7}
\textbf{End of V7.}  The complete V1--V6 body above is retained verbatim.  V7
closes E6 by replacing nonfunctorial order-by-order Hodge recomputation with a
chronological projective construction, proving exact finite-cutoff and
continuum restriction in loop order, constructing invertible continuum
scheme/reference transitions, retaining all determinant/Pfaffian frame and
source-reference contacts, and stating the final source-complete
coefficientwise Renewal Einstein--Standard Model continuum theorem.  No
load-bearing gate remains inside the stated E1--E6 programme.

\clearpage
\section{V8 continuation: a stronger post-E1--E6 programme --- microscopic regulator universality}
\label{sec:v8-continuation}

V7 closed the stated E1--E6 coefficientwise continuum programme on the explicit
V2 Wilson/overlap branch.  The present continuation therefore does not reopen
any gate in that programme.  It attacks the first stronger question isolated in
the V7 ledger: whether the continuum just constructed depends on the
microscopic fermion completion used to realize the protected chiral shell.

The claim proved below is deliberately narrower than arbitrary regulator
independence.  We first formulate a local gapped-homotopy criterion and then
verify it for a genuinely different finite-range stencil obtained by adding a
range-two covariant Wilson deformation to the V2 range-one kernel.  At every
fixed loop/EFT/source order the two regulators are shown to differ by one finite
local scheme/source/line transformation plus a quasilocal remainder which
vanishes in the ultraviolet limit.  No statement is made between disconnected
Wilson phases, across determinant-line zeros, or for a regulator for which the
required gapped homotopy has not been verified.

\subsection{The relative question and the correct notion of universality}

Fix throughout a finite tuple
\begin{equation}
 \mathfrak q=(L,\Delta,r,N,\mathcal P_{\rm ext}),
 \label{eq:v8-fixed-data}
\end{equation}
where $\mathcal P_{\rm ext}$ is the finite external momentum/source panel, and
retain the positive reference scale $\mu>0$.  Let
$\mathfrak C^{\rm ESM}_{\mathfrak q}$ be the V5--V7
contraction-complete state at these fixed data.

\begin{definition}[Matched regulator universality at fixed order]
\label{def:v8-universality}
Two microscopic regulator families $A$ and $B$ are \emph{matched-universal at
$\mathfrak q$} if there are finite local invertible maps
\begin{equation}
 F_{a,\mathfrak q}^{B\leftarrow A}:
 \mathfrak C^{A}_{a,\mathfrak q}\longrightarrow
 \mathfrak C^{B}_{a,\mathfrak q}
 \label{eq:v8-finite-reg-map}
\end{equation}
which preserve the chosen physical normalization coordinates and transport the
composite/source/line basis covariantly, such that for the normalized connected
states at a fixed physical scale $\lambda$,
\begin{equation}
 \left\|
 \widehat{\mathfrak C}^{B}_{a,\lambda}
 -F_{a,\mathfrak q}^{B\leftarrow A}
  \widehat{\mathfrak C}^{A}_{a,\lambda}
 \right\|_{\mathfrak q}
 \xrightarrow[a\downarrow0]{}0.
 \label{eq:v8-universality-def}
\end{equation}
The maps must be projective in perturbative order in the sense of V7.
\end{definition}

This is the appropriate statement for an EFT.  Equality of bare Wilson
coordinates is neither expected nor required.  What is required is equality of
the normalized connected/source theory after one finite local matching map.

\begin{firewall}[Universality is not equality of bare actions]
A change of microscopic regulator generally changes power-divergent and finite
local coefficients.  Requiring those bare coordinates themselves to converge to
the same values would contradict ordinary scheme dependence.  The comparison
must first factor the finite local normalization cocycle, exactly as V6 factored
ordinary marginal running before proving the ultraviolet Cauchy theorem.
\end{firewall}

\subsection{An explicit genuinely different native stencil}

Retain the V2 covariant shifts $T_\mu^{\mathcal U}$ and Hermitian Wilson pieces
\begin{equation}
 L_\mu^{\mathcal U}
 =I-\frac12\left(T_\mu^{\mathcal U}+
 (T_\mu^{\mathcal U})^*\right).
 \label{eq:v8-Lmu}
\end{equation}
The V2 kernel is denoted $K_{W,a}^{(0)}$ and
$H_{W,a}^{(0)}=\gamma_5K_{W,a}^{(0)}$.

\begin{definition}[Range-two Wilson deformation]
\label{def:v8-range-two}
Put
\begin{equation}
 Q_a^{\mathcal U}:=
 \sum_{\mu=1}^4\left(L_\mu^{\mathcal U}\right)^2,
 \qquad
 K_{W,a}^{(\eta)}:=K_{W,a}^{(0)}+\eta Q_a^{\mathcal U},
 \qquad
 H_{W,a}^{(\eta)}:=\gamma_5K_{W,a}^{(\eta)}.
 \label{eq:v8-range-two-kernel}
\end{equation}
For $\eta\ne0$ this is a range-two microscopic stencil and is therefore not the
V2 nearest-neighbour kernel.
\end{definition}

\begin{lemma}[Uniform operator bound for the deformation]
\label{lem:v8-Q-bound}
For unitary gauge--spin parallel transports,
\begin{equation}
 \|L_\mu^{\mathcal U}\|\le2,
 \qquad
 \|Q_a^{\mathcal U}\|\le16.
 \label{eq:v8-Q-bound}
\end{equation}
Every fixed gauge/spin source derivative of $Q_a^{\mathcal U}$ is finite range
and obeys the same coefficientwise edge-jet calculus as V2.
\end{lemma}
\begin{proof}
Because $T_\mu^{\mathcal U}$ is unitary,
$\|(T_\mu^{\mathcal U}+(T_\mu^{\mathcal U})^*)/2\|\le1$, whence
$\|L_\mu^{\mathcal U}\|\le2$.  Thus
$\|(L_\mu^{\mathcal U})^2\|\le4$, and summing four directions gives
\eqref{eq:v8-Q-bound}.  Differentiation produces a finite Leibniz sum of
products of the one-edge derivatives already controlled in
Lemma~\ref{lem:v2-link-jets}; products of two range-one kernels have range at
most two.
\end{proof}

\begin{proposition}[Hermiticity, covariance and finite range]
\label{prop:v8-range-two-structure}
Assume the V2 spin transport hypotheses.  Then for every real $\eta$,
$H_{W,a}^{(\eta)}$ is Hermitian, exactly internal-gauge/spin covariant and of
range at most two.  Coframe dependence remains entirely in the V2 Clifford
kinetic term, and Higgs/Yukawa, antifield and determinant-line coordinates
enter exactly as in V2.
\end{proposition}
\begin{proof}
Each $L_\mu^{\mathcal U}$ is Hermitian and commutes with $\gamma_5$ under the
V2 spin-transport hypotheses.  Hence $Q_a^{\mathcal U}$ is Hermitian,
$\gamma_5$-even and transforms by the same unitary similarity as the original
kernel.  The range statement follows from two compositions of range-one
shifts.  The remaining claims are immediate from
\eqref{eq:v8-range-two-kernel}.
\end{proof}

Let $g_{\rm W}>0$ be the protected Wilson gap of the V2 kernel on the compact
reference chart used by V1--V7.  Shrink the chart once, if necessary, so this
same lower bound is valid for all reference source jets used below.

\begin{theorem}[A uniform gapped homotopy]
\label{thm:v8-gapped-path}
Choose
\begin{equation}
 0<\eta_*<\frac{g_{\rm W}}{32}.
 \label{eq:v8-eta-small}
\end{equation}
For $t\in[0,1]$ set $\eta(t)=t\eta_*$.  Then
\begin{equation}
 \boxed{
 |H_{W,a}^{(\eta(t))}|
 \succeq \frac{g_{\rm W}}2 I}
 \label{eq:v8-homotopy-gap}
\end{equation}
uniformly in $t$, cutoff and volume on the protected chart.  In particular the
path does not cross a Wilson-phase boundary.
\end{theorem}
\begin{proof}
By Lemma~\ref{lem:v8-Q-bound},
\[
 \|H_{W,a}^{(\eta(t))}-H_{W,a}^{(0)}\|
 =\eta(t)\|Q_a^{\mathcal U}\|
 \le16\eta_*<\frac{g_{\rm W}}2.
\]
The min--max/Weyl perturbation inequality for singular values gives
\eqref{eq:v8-homotopy-gap}.
\end{proof}

At the flat reference point the additional symbol is
\begin{equation}
 q(x)=\sum_{\mu=1}^4(1-\cos x_\mu)^2=O(|x|^4).
 \label{eq:v8-q-symbol}
\end{equation}
Thus the deformation changes neither the physical first-order Dirac symbol nor
the Standard-Model representation carried by the chiral block.

\subsection{Relative overlap functional calculus along the regulator path}

Write
\begin{equation}
 \varepsilon_t=\operatorname{sgn}H_{W,a}^{(\eta(t))},
 \qquad
 D_{{\rm ov},t}=a^{-1}(I+\gamma_5\varepsilon_t),
 \qquad
 \widehat P_{-,t}=\frac12(I+\varepsilon_t).
 \label{eq:v8-overlap-path}
\end{equation}
Let $D_{\chi,t}$ be the corresponding fixed-frame chiral/Yukawa block, with the
same Yukawa map as in V2.

\begin{lemma}[Quasilocal regulator transgression]
\label{lem:v8-sign-transgression}
For every fixed source multi-index $I$ and every fixed number of momentum
moments, the derivatives
\begin{equation}
 \partial_I\partial_t\varepsilon_t,
 \qquad
 \partial_I\partial_t\widehat P_{-,t}
 \label{eq:v8-sign-derivs}
\end{equation}
are cutoff-uniform quasilocal kernels on the protected chart.  More precisely,
with $\Gamma_\pm$ fixed contours around the positive/negative spectrum,
\begin{equation}
 \partial_t\varepsilon_t
 =\frac{1}{2\pi i}
 \left(\oint_{\Gamma_+}-\oint_{\Gamma_-}\right)
 (z-H_t)^{-1}\dot H_t(z-H_t)^{-1}\,dz,
 \label{eq:v8-sign-contour}
\end{equation}
where $\dot H_t=\eta_*\gamma_5Q_a^{\mathcal U}$.
\end{lemma}
\begin{proof}
The gap in Theorem~\ref{thm:v8-gapped-path} permits $t$-independent contours.
Differentiate the Riesz projectors.  Fixed source derivatives give finite sums
of resolvent words containing source derivatives of $H_t$ and one marked
$\dot H_t$.  The V1 Combes--Thomas/finite-range functional calculus applies
uniformly because the range, gap and finite source order are fixed.  Weighted
kernel moments are therefore uniform in $t$.
\end{proof}

\begin{proposition}[Common Kato frame along the regulator homotopy]
\label{prop:v8-regulator-Kato}
On a star-shaped protected chart there is a unitary two-parameter Kato frame
$U(t,\zeta)$ which simultaneously transports the moving source space
$\Ran\widehat P_{-,t}(\zeta)$ from the V2 base frame.  Every fixed mixed
$(t,\zeta)$ derivative is quasilocal and uniformly bounded in the V1 rooted
norms.
\end{proposition}
\begin{proof}
Apply the Kato ODE first in the $t$ direction and then in the star-shaped source
coordinates, or equivalently solve the path-ordered transport for the smooth
projector family $\widehat P_{-,t}(\zeta)$.  The generator is a commutator of
projector derivatives.  Lemma~\ref{lem:v8-sign-transgression} and the V1 Kato
bounds close the finite derivative algebra.  The parameter rectangle is
contractible, so no global line-holonomy assertion is needed.
\end{proof}

\begin{theorem}[Relative physical-annulus symbol estimate]
\label{thm:v8-relative-symbol}
Fix a source order and momentum-derivative order.  On a physical hard annulus
$|p|\asymp\Lambda_j$ with $a\Lambda_j\le c_0$, the fixed-frame chiral blocks
satisfy
\begin{equation}
 \boxed{
 \left\|
 \partial_p^\alpha\partial_{\mathbf q}^\beta\partial_I
 \bigl(D_{\chi,t}-D_{\chi,0}\bigr)
 \right\|
 \le C_{I,\alpha,\beta}\,
 a\Lambda_j^{\nu_I+1-|\alpha|-|\beta|}.}
 \label{eq:v8-relative-D}
\end{equation}
The estimate is uniform in $t\in[0,1]$.  Consequently, for the hard inverses
$G_{\chi,t}=D_{\chi,t}^{-1}$,
\begin{equation}
 \boxed{
 \left\|
 \partial_p^\alpha\partial_{\mathbf q}^\beta\partial_I
 \bigl(G_{\chi,t}-G_{\chi,0}\bigr)
 \right\|
 \le C_{I,\alpha,\beta}\,
 a\Lambda_j^{-|\alpha|-|\beta|}}
 \label{eq:v8-relative-G}
\end{equation}
after the usual engineering/source stocks are divided out.  Equivalently, the
relative change is $O(a\Lambda_j)$ compared with one ordinary hard fermion
propagator.
\end{theorem}
\begin{proof}
At zero source, \eqref{eq:v8-q-symbol} is $O(|ap|^4)$, hence the claimed weaker
$O(a|p|^2)$ relative overlap estimate follows immediately from the analytic
Riesz functional calculus and the prefactor $a^{-1}$.  For source derivatives,
Lemma~\ref{lem:v2-link-jets} applied to the products
$(L_\mu^{\mathcal U})^2$ shows that the first local covariant jet is unchanged
and that the deformation contributes only to the V2 remainder class.  Mixed
Kato-frame derivatives are controlled by
Proposition~\ref{prop:v8-regulator-Kato}.  This proves
\eqref{eq:v8-relative-D}.

The V2 hard pivot, decreased if necessary by a fixed factor using
Theorem~\ref{thm:v8-gapped-path}, gives
$\|G_{\chi,t}\|\le C/\Lambda_j$.  The resolvent identity
\[
 G_{\chi,t}-G_{\chi,0}
 =-G_{\chi,t}(D_{\chi,t}-D_{\chi,0})G_{\chi,0}
\]
and its fixed differentiated versions then give
\eqref{eq:v8-relative-G}.
\end{proof}

\begin{proposition}[Relative ultraviolet-cap calculus]
\label{prop:v8-relative-full-zone}
On every compact full-zone region $|ap|\ge\delta>0$, both hard inverses are
$a$ times smooth periodic dimensionless symbols.  Their difference and every
fixed source derivative therefore have the same $O(a)$ lattice-scale weighted
row/column bound.  An external source-momentum subtraction gains $a\mu$.
\end{proposition}
\begin{proof}
Use the uniform gap and no-physical-zero compactness argument of V2 for the
whole $t$ interval.  The difference of two $a$-scaled smooth periodic symbols
is again of that form.  Fourier-series summation by parts and the mean-value
theorem in the dimensionless variables $aq$ give the result.
\end{proof}

\subsection{Determinant/Pfaffian transgression along the same path}

On a hard block with nonzero pivot define its scalar representative in the
common Kato frame by $s_{j,t}$.  For an optional neutral Majorana block use the
corresponding Pfaffian representative.

\begin{lemma}[Exact line transgression]
\label{lem:v8-line-transgression}
Along the protected regulator path,
\begin{align}
 \partial_t\log s_{j,t}
 &=\Tr\bigl(G_{j,t}\,\partial_tD_{j,t}\bigr),
 \label{eq:v8-det-transgression}\\
 \partial_t\log\operatorname{Pf}A_{j,t}
 &=\frac12\Tr\bigl(A_{j,t}^{-1}\partial_tA_{j,t}\bigr)
 \label{eq:v8-pf-transgression}
\end{align}
whenever the Pfaffian block is present.  Every fixed source derivative is the
finite differentiated trace word obtained from these identities.  Gauge
variation of the complete Standard-Model product is zero at every $t$.
\end{lemma}
\begin{proof}
The first two formulas are Jacobi's determinant identity and its Pfaffian
analogue on the nonzero component.  Fixed source derivatives are finite by the
reference-pivot theorem.  Exact gauge covariance of $D_{\chi,t}$ and the same
complete Standard-Model representation at every $t$ give the same finite gauge
descent cocycle as in V3; that cocycle is one for the complete generation.
\end{proof}

\begin{theorem}[Local plus quasilocal line transgression]
\label{thm:v8-line-local-plus-rem}
For each fixed source order, the shell line transgression decomposes as
\begin{equation}
 \partial_t\log s_{j,t}
 =\mathsf T_j\partial_t\log s_{j,t}
  +\mathsf Q_j\partial_t\log s_{j,t},
 \label{eq:v8-line-decomp}
\end{equation}
where the first term belongs to the finite local source/reference bank and the
second obeys the same two-zone bounds as the relative hard kernels:
\begin{equation}
 \|\partial_I\mathsf Q_j\partial_t\log s_{j,t}\|
 \le
 \begin{cases}
 C_I a\Lambda_j,& a\Lambda_j\le c_0,\\
 C_I\mu/\Lambda_j,& \Lambda_j\ge M\text{ after local Taylor subtraction}.
 \end{cases}
 \label{eq:v8-line-two-zone}
\end{equation}
The same statement holds for the Pfaffian with the factor $1/2$ in
\eqref{eq:v8-pf-transgression}.
\end{theorem}
\begin{proof}
Differentiate the trace words of Lemma~\ref{lem:v8-line-transgression}.  On the
physical branch at least one factor is the marked relative insertion from
Theorem~\ref{thm:v8-relative-symbol}, which gives the relative
$a\Lambda_j$ gain.  On the high branch use the V2 full-zone source calculus and
subtract the complete finite local Taylor/source jet; the first omitted soft
momentum gives $\mu/\Lambda_j$.  The V1 Bell-polynomial line composition retains
all mixed-shell contacts, so the same decomposition is stable under shell
composition.
\end{proof}

\subsection{The relative forest problem}

A regulator comparison cannot be proved graph by graph with independently
chosen finite parts.  We therefore differentiate the \emph{already normalized
common forest} along the regulator path.

For every $t$, let $\mathbf R_{j,t}^{<L}$ be the V5 forest preparation formed
with the chronological lower-order actions already fixed at that same $t$.
The normalization conditions are held fixed in physical coordinates as $t$
varies.

\begin{definition}[Regulator-marked forest]
\label{def:v8-marked-forest}
For a forest-prepared graph $\Gamma$, a \emph{regulator mark} is either
\begin{enumerate}[label=\textup{(R\arabic*)},leftmargin=3em]
\item one primitive propagator or vertex differentiated with respect to $t$;
\item one derivative of a previously fixed lower-loop local counterterm;
\item one derivative of a determinant/Pfaffian line or source/reference local
coordinate.
\end{enumerate}
The derivative $\partial_t\mathbf R_{j,t}^{<L}\mathcal I_{\Gamma,t}$ is the
finite signed sum of all such marked forests.
\end{definition}

\begin{lemma}[No independent finite part appears in the relative forest]
\label{lem:v8-relative-ownership}
Every marked proper subtraction block in
$\partial_t\mathbf R_{j,t}^{<L}\mathcal I_{\Gamma,t}$ is the derivative of a
previously fixed common coefficient $\bar S_{q,t}$ with $q<L$.  In particular,
regulator comparison introduces no graph-dependent finite normalization.
\end{lemma}
\begin{proof}
By the V5 proper-subgraph loop-drop theorem, every proper 1PI subtraction block
has loop number strictly below that of $\Gamma$.  The chronological V7
construction freezes its finite normalization when that lower order is first
solved.  Differentiating in $t$ therefore differentiates that already-owned
coefficient; it does not reopen its finite part.
\end{proof}

\begin{theorem}[Relative fixed-loop forest decomposition]
\label{thm:v8-relative-forest}
Fix $\mathfrak q$ and a shell $j$.  After one common local Taylor/EFT/source
projection, the regulator derivative of every order-$\le L$ forest coefficient
has a decomposition
\begin{equation}
 \boxed{
 \partial_t\mathbf R_{j,t}^{<L}\mathcal I_{\Gamma,t}
 =\mathsf B_{\Gamma,j,t}^{\rm loc}
  +\mathsf E_{\Gamma,j,t}^{\rm rel}.}
 \label{eq:v8-relative-forest-decomp}
\end{equation}
The local term belongs to the same finite V5 bank.  For any intermediate scale
$M$ with
\begin{equation}
 \mu\ll M\ll a^{-1},
 \label{eq:v8-M-window}
\end{equation}
the complementary rooted kernel obeys
\begin{equation}
 \boxed{
 \|\mathsf E_{\Gamma,j,t}^{\rm rel}\|_*
 \le
 \begin{cases}
 C_{\mathfrak q}\,a\Lambda_j,
     & \Lambda_j\le M,\\[0.25em]
 C_{\mathfrak q}\,\mu/\Lambda_j,
     & \Lambda_j>M.
 \end{cases}}
 \label{eq:v8-relative-forest-bound}
\end{equation}
The statement includes all stored future BV/source contraction descendants.
\end{theorem}
\begin{proof}
Differentiate the finite V5 forest formula.  By
Definition~\ref{def:v8-marked-forest}, every term contains one regulator mark.
On shells below $M$, the mark is a physical-annulus relative propagator/vertex
or a lower-loop counterterm generated by such a mark.  Theorem
\ref{thm:v8-relative-symbol}, the bounded rooted contraction ideal of V5 and
induction in the loop number give the factor $a\Lambda_j$.  Nested and
overlapping proper subgraphs remain in their original compatible forest terms;
Lemma~\ref{lem:v8-relative-ownership} supplies their common finite parts.

For shells above $M$, no small-$ap$ expansion is used.  Apply the native V5
Taylor/forest decomposition separately at each $t$ and differentiate it.  All
nondecaying dependence on the microscopic stencil is local and is assigned to
$\mathsf B_{\Gamma,j,t}^{\rm loc}$.  The first omitted external/source Taylor
order carries one excess soft-to-hard ratio $\mu/\Lambda_j$.  The V5
contraction-complete packet guarantees that a later BV Laplacian or source
contraction of this remainder has its newly local descendant already stored;
there is therefore no hidden loss of the scale factor.  Finiteness of the graph
and source families at fixed $\mathfrak q$ completes the estimate.
\end{proof}

\subsection{The two-scale universality estimate}

The preceding theorem isolates the ultraviolet mechanism cleanly.  Microscopic
regulators are close on physical shells, whereas on cutoff-scale shells they
need not be close but their difference is local as seen from a fixed soft panel.

\begin{lemma}[Geometric two-scale sums]
\label{lem:v8-two-scale-sum}
For dyadic/geometric shell scales $\Lambda_j=b^j\lambda$, $b>1$,
\begin{align}
 \sum_{\Lambda_j\le M}a\Lambda_j
 &\le C_b aM,
 \label{eq:v8-low-sum}\\
 \sum_{\Lambda_j>M}\frac{\mu}{\Lambda_j}
 &\le C_b\frac{\mu}{M}.
 \label{eq:v8-high-sum}
\end{align}
\end{lemma}
\begin{proof}
Both are geometric series, summed respectively from below and from above the
unique shell comparable with $M$.
\end{proof}

Let $\mathcal N_{\mathfrak q}$ denote the fixed complete finite normalization
map used in V6--V7: physical couplings, redundant/BV-exact coordinates,
composite/source basis, stress improvements, contact coordinates and local line
normalization are all included.

\begin{theorem}[One common local regulator-matching velocity]
\label{thm:v8-local-matching}
Along the path $t\mapsto H_t$ there exists a unique chronological local velocity
$Y_{a,t,\mathfrak q}$ in the finite V5 local bank such that the normalized
coordinates remain fixed:
\begin{equation}
 \frac d{dt}\,
 \mathcal N_{\mathfrak q}
 \bigl(\widehat{\mathfrak C}_{a,t}\bigr)=0.
 \label{eq:v8-normalization-fixed}
\end{equation}
It is chosen once for the entire state, not separately for graphs or Ward rows,
and is projective in $L$.
\end{theorem}
\begin{proof}
V6 constructs a normalization section with an invertible finite Jacobian on the
protected constant-rank chart.  Differentiate the normalization equations in
$t$.  The implicit-function theorem gives a unique finite local velocity which
cancels the local vector assembled from
$\mathsf B_{\Gamma,j,t}^{\rm loc}$, line/reference local terms and the Hodge-exact
coordinates.  Because V7 freezes every lower-order choice chronologically, the
order-$q$ component of this linear system is unchanged when a higher terminal
loop order is appended.  Hence the resulting family is projective.
\end{proof}

\begin{theorem}[Two-scale regulator comparison at fixed order]
\label{thm:v8-two-scale-universality}
Let $\widehat{\mathfrak C}_{a,t,\lambda}$ be the normalized state at the fixed
physical scale $\lambda$, transported along the regulator homotopy with the
local velocity of Theorem~\ref{thm:v8-local-matching}.  Then for every
intermediate $M$ satisfying \eqref{eq:v8-M-window},
\begin{equation}
 \boxed{
 \left\|
 \partial_t\widehat{\mathfrak C}_{a,t,\lambda}
 \right\|_{\mathfrak q}
 \le
 C_{\mathfrak q}\left(aM+\frac{\mu}{M}\right).}
 \label{eq:v8-two-scale-bound}
\end{equation}
The same estimate holds after every retained source derivative and for the
complete BV/master and determinant/Pfaffian line packets.
\end{theorem}
\begin{proof}
The local part of every marked forest is cancelled by the one common
normalization velocity.  Sum the complementary contributions from
Theorem~\ref{thm:v8-relative-forest}.  Lemma~\ref{lem:v8-two-scale-sum} gives
$aM$ below $M$ and $\mu/M$ above $M$.  The V5 rooted ideal controls products
containing several complementary ancestors; those terms carry at least the same
single marked factor and higher powers only improve the bound.  V4--V5 source
convolution and contraction completeness show that source/BV descendants obey
the same estimate.  The line contribution is
Theorem~\ref{thm:v8-line-local-plus-rem}.  All families are finite at fixed
$\mathfrak q$.
\end{proof}

\begin{corollary}[Optimized regulator-difference rate]
\label{cor:v8-sqrt-rate}
For $a\mu<1$, choose
\begin{equation}
 M_a=\sqrt{\frac{\mu}{a}}.
 \label{eq:v8-optimal-M}
\end{equation}
Then $\mu\ll M_a\ll a^{-1}$ and
\begin{equation}
 \boxed{
 \left\|
 \partial_t\widehat{\mathfrak C}_{a,t,\lambda}
 \right\|_{\mathfrak q}
 \le C_{\mathfrak q}\sqrt{a\mu}.}
 \label{eq:v8-sqrt-rate}
\end{equation}
Consequently
\begin{equation}
 \left\|
 \widehat{\mathfrak C}_{a,1,\lambda}
 -F_{a,\mathfrak q}^{1\leftarrow0}
  \widehat{\mathfrak C}_{a,0,\lambda}
 \right\|_{\mathfrak q}
 \le C_{\mathfrak q}\eta_*\sqrt{a\mu},
 \label{eq:v8-endpoint-rate}
\end{equation}
where $F_{a,\mathfrak q}^{1\leftarrow0}$ is the integrated finite local
matching map.
\end{corollary}
\begin{proof}
At \eqref{eq:v8-optimal-M}, both terms in
\eqref{eq:v8-two-scale-bound} equal $\sqrt{a\mu}$.  Integrate in
$t\in[0,1]$.  The local flow generated by $Y_{a,t,\mathfrak q}$ is finite
dimensional and invertible on the compact chart, giving
$F_{a,\mathfrak q}^{1\leftarrow0}$.
\end{proof}

\subsection{Continuum universality and exact order compatibility}

\begin{theorem}[Microscopic universality for the explicit range-two deformation]
\label{thm:v8-explicit-universality}
Let regulator $A$ be the V2 nearest-neighbour Wilson/overlap completion and let
regulator $B$ be \eqref{eq:v8-range-two-kernel} with
$\eta=\eta_*$ satisfying \eqref{eq:v8-eta-small}.  Then for every fixed
$\mathfrak q$ they are matched-universal in the sense of
Definition~\ref{def:v8-universality}.  More precisely, the finite local maps
$F_{a,\mathfrak q}^{B\leftarrow A}$ have a projective continuum limit
$F_{\infty,\mathfrak q}^{B\leftarrow A}$ and
\begin{equation}
 \boxed{
 \widehat{\mathfrak C}^{B}_{\infty,\lambda}
 =F_{\infty,\mathfrak q}^{B\leftarrow A}
  \widehat{\mathfrak C}^{A}_{\infty,\lambda}.}
 \label{eq:v8-continuum-universality}
\end{equation}
After covariant source/operator matching, the normalized connected coefficients,
all retained BV/BRST/master relations, stress/composite/contact coefficients and
determinant/Pfaffian line insertions agree.
\end{theorem}
\begin{proof}
The V2 regulator has the V7 continuum.  The range-two regulator satisfies the
same E1 hypotheses by Propositions~\ref{prop:v8-range-two-structure},
Theorem~\ref{thm:v8-gapped-path}, Theorem~\ref{thm:v8-relative-symbol} and
Proposition~\ref{prop:v8-relative-full-zone}.  The Standard-Model representation,
anomaly cancellation and finite gauge descent are unchanged, so the V3--V7
cohomological, forest, normalization and Cauchy constructions apply with the
same fixed-order quantifiers.

For the relative statement, Corollary~\ref{cor:v8-sqrt-rate} shows that the two
normalized finite-cutoff states differ, after the common local matching map, by
a quantity tending to zero.  The local maps are built from the same compact
normalization section as V7; their scale-to-scale differences are summable once
the physical normalization is fixed, so the V7 finite-scheme argument gives an
invertible continuum limit $F_{\infty,\mathfrak q}^{B\leftarrow A}$.  Passing
$a\downarrow0$ in \eqref{eq:v8-endpoint-rate} proves
\eqref{eq:v8-continuum-universality}.  Every listed source/BV/line component is
a coordinate of the contraction-complete state, so the same limit identity
holds coefficientwise for it.
\end{proof}

\begin{theorem}[Strict perturbative-order projectivity of the regulator map]
\label{thm:v8-reg-projective}
Let $\pi_L^{L+1}$ be the V7 perturbative truncation.  The chronological
regulator-matching maps may be chosen so that at finite cutoff
\begin{equation}
 \boxed{
 \pi_L^{L+1}
 F_{a,L+1}^{B\leftarrow A}
 \iota_L^{L+1}
 =F_{a,L}^{B\leftarrow A},}
 \label{eq:v8-reg-projective-finite}
\end{equation}
and therefore
\begin{equation}
 \boxed{
 \pi_L^{L+1}
 F_{\infty,L+1}^{B\leftarrow A}
 \iota_L^{L+1}
 =F_{\infty,L}^{B\leftarrow A}.}
 \label{eq:v8-reg-projective-limit}
\end{equation}
Thus regulator universality is an equivalence of the projective coefficientwise
tower, not a separate finite redefinition at every terminal loop order.
\end{theorem}
\begin{proof}
Theorem~\ref{thm:v8-local-matching} solves the regulator-matching velocity
chronologically.  Once its order-$q$ component is fixed, increasing the
terminal loop order only appends a new equation in the new
$\hbar^{q+1}$ layer and deeper future-descendant slots.  Lower components are
not recomputed.  This proves \eqref{eq:v8-reg-projective-finite}; continuity of
truncation and the finite-dimensional scheme limits gives
\eqref{eq:v8-reg-projective-limit}.
\end{proof}

\subsection{A general gapped-homotopy criterion}

The explicit range-two example reveals which hypotheses were actually used.

\begin{definition}[Admissible regulator homotopy at fixed coefficient order]
\label{def:v8-admissible-homotopy}
A path $t\mapsto H_{a,t}$ is \emph{$\mathfrak q$-admissible} if:
\begin{enumerate}[label=\textup{(A\arabic*)},leftmargin=3em]
\item it is finite range (or uniformly exponentially local), Hermitian and
exactly gauge/spin covariant, with fixed source jets through the orders demanded
by $\mathfrak q$;
\item it has a cutoff- and $t$-uniform Wilson gap on one common protected chart;
\item its overlap/chiral physical branch has the same continuum principal
Dirac--Yukawa symbol for all $t$, and the $t$ derivative of every required
physical source vertex is one irrelevant order smaller, in the sense of the
relative estimate \eqref{eq:v8-relative-D};
\item on the complementary full zone, the hard inverse/source words are smooth
periodic dimensionless symbols with the V2 weighted kernel bounds;
\item the determinant/Pfaffian hard pivots do not cross zero along the path;
\item the same finite V5 local EFT/BV/source bank and normalization-rank chart
remain valid along the path.
\end{enumerate}
\end{definition}

\begin{theorem}[Gapped-homotopy universality criterion]
\label{thm:v8-general-homotopy}
Any two endpoints of a $\mathfrak q$-admissible regulator homotopy are
matched-universal at $\mathfrak q$.  If one coherent path is admissible for
every finite $L$ and the local matching is chosen chronologically, the endpoint
maps form a projective family in $L$.
\end{theorem}
\begin{proof}
Items (A1)--(A5) reproduce precisely the inputs used in
Lemmas~\ref{lem:v8-sign-transgression}--\ref{lem:v8-relative-ownership} and
Theorem~\ref{thm:v8-relative-forest}.  Item (A6) supplies the common finite
normalization velocity.  The two-scale proof
Theorem~\ref{thm:v8-two-scale-universality} and its optimization then apply
verbatim.  Chronological choice gives the projective statement exactly as in
Theorem~\ref{thm:v8-reg-projective}.
\end{proof}

\begin{firewall}[What the homotopy criterion does not prove]
Theorem~\ref{thm:v8-general-homotopy} is a criterion, not a connectedness theorem
for the entire space of chiral regulators.  Two kernels with different overlap
index, a forced Wilson-gap closing, incompatible anomaly representation, or a
determinant-line zero need not lie in one admissible component.  No such phase
transition is crossed in V8.
\end{firewall}

\subsection{V8 anti-circularity audit}
\label{sec:v8-audit}

\begin{enumerate}[leftmargin=2.4em]
\item The second regulator is explicit and microscopically different: it has a
range-two stencil.  Universality is not inferred by renaming the V2 kernel.
\item The regulator path is proved gapped by the operator bound
\eqref{eq:v8-homotopy-gap}; a continuum Dirac operator is not substituted for a
native kernel.
\item High-zone regulator differences are not assumed small.  Their nondecaying
part is placed in the finite local bank and matched by one common normalization
flow; only the Taylor complement is estimated by $\mu/\Lambda_j$.
\item Low-zone closeness and high-zone locality are combined with a moving
intermediate scale.  The comparison does not sum a nonuniform small-$ap$
expansion up to the cutoff.
\item Proper-subgraph finite parts remain owned by the previously fixed common
actions.  No graph is independently refitted while changing regulator.
\item Determinant/Pfaffian frame and source-reference terms are differentiated
explicitly and retained.  A field-independent but source-dependent endpoint
term is not discarded.
\item The order-$L$ regulator map is frozen when constructed.  Higher terminal
orders append new layers and do not recompute lower-order matching.
\end{enumerate}

\section{V8 status ledger}
\label{sec:v8-status}

\subsection*{Proved in V8}
\begin{enumerate}[leftmargin=2.2em]
\item The explicit range-two deformation
$K_{W,a}^{(\eta)}=K_{W,a}^{(0)}+\eta\sum_\mu(L_\mu^{\mathcal U})^2$ is
Hermitian, gauge/spin covariant and finite range, and for sufficiently small
fixed $\eta_*$ it is joined to the V2 regulator by a cutoff-uniform gapped
native homotopy.
\item Along that homotopy, overlap projectors, Kato frames, chiral hard
propagators and every required fixed source jet have uniform quasilocal
$t$ derivatives.  On physical shells the relative propagator is smaller by
$O(a\Lambda_j)$; on the full zone the ordinary lattice-scale bounds hold.
\item Determinant/Pfaffian regulator transgression is given by exact trace
identities, with all source descendants.  Its complete Standard-Model gauge
variation vanishes and its nonlocal shell remainder obeys the same two-zone
bounds.
\item The derivative of the complete native forest is a finite sum of marked
forests whose proper local subtractions remain owned by lower-loop common
actions.  After one common local projection the relative remainder obeys
\eqref{eq:v8-relative-forest-bound}.
\item A single finite normalization/source/line velocity cancels the complete
local regulator derivative.  Summing shells below and above an intermediate
scale gives $C(aM+\mu/M)$; the choice $M=\sqrt{\mu/a}$ yields the explicit
$O(\!\sqrt{a\mu}\!)$ regulator-difference bound.
\item Therefore the V2 nearest-neighbour regulator and the explicit V8
range-two regulator have the same fixed-order continuum after one finite local
matching map.  The equivalence includes normalized connected coefficients,
composites, stress/contact terms, ghosts/antifields, BV/master relations and
determinant/Pfaffian line insertions.
\item The regulator equivalence maps are strictly projective in perturbative
order.
\item More generally, endpoints of any explicitly verified admissible gapped
homotopy satisfying Definition~\ref{def:v8-admissible-homotopy} are
matched-universal at the prescribed fixed order.
\end{enumerate}

\subsection*{Inherited}
\begin{enumerate}[leftmargin=2.2em]
\item The complete V1--V7 coefficientwise continuum construction on the V2
branch, including the native hard-shell estimates, anomaly cancellation,
common BV restoration, arbitrary-fixed-loop forest theorem, ultraviolet Cauchy
theorem, scheme/reference covariance and exact order projectivity.
\item The V2 edge-parallel-transport source calculus and full-zone weighted
kernel estimates.
\item The protected nonzero determinant/Pfaffian component and positive
IR/reference prescription.
\end{enumerate}

\subsection*{Conditional}
\begin{enumerate}[leftmargin=2.2em]
\item The explicit second-regulator theorem requires
$0<\eta_*<g_{\rm W}/32$ and one common compact protected chart.  Larger
range-two deformations are not covered without a new gap proof.
\item The general criterion requires an explicitly verified gapped homotopy and
a common constant-rank normalization/BV chart.  V8 does not prove that every
reasonable chiral regulator admits such a path.
\item All statements remain coefficientwise at prescribed finite
$(L,\Delta,r,N)$, with constants allowed to depend on these data.
\item The positive IR/reference scale and local nonzero line component remain in
force.  No zero crossing or global line holonomy is treated.
\end{enumerate}

\subsection*{Open beyond V8}
\begin{enumerate}[leftmargin=2.2em]
\item Classify the connected components of the admissible native chiral
regulator space and prove a connectedness theorem within a fixed overlap-index
and anomaly-representation sector.  This would remove the need to exhibit the
homotopy by hand for each new microscopic stencil.
\item Extend the universality comparison to a qualitatively different native
fermion realization (for example a non-nearest-neighbour graph kernel or a
non-ultralocal exponentially local completion) by verifying the V8 admissible
homotopy gates.
\item The independent stronger problems from V7 remain: massless IR removal,
global determinant/Pfaffian zero-crossing/holonomy, analytic finite-coupling RG
trajectories and nonperturbative completion.
\end{enumerate}

\subsection*{Exact next load-bearing theorem}
\begin{center}
\fbox{\parbox{0.92\textwidth}{
\textbf{Regulator-phase connectedness / classification.}
Fix the Standard-Model representation and the physical overlap index.  Define
the space of finite-range or uniformly exponentially local, gauge/spin
covariant Hermitian kernels satisfying the common physical-principal-symbol,
source-jet and protected-gap conditions.  Prove that any two kernels in the same
topological phase can be joined by a finite chain of V8-admissible gapped
homotopies, or identify a computable additional obstruction.  Combined with
Theorem~\ref{thm:v8-general-homotopy}, this would upgrade the present explicit
range-two universality theorem to a regulator-class universality theorem.}}
\end{center}

\section*{Append-only continuation marker: V8}
\addcontentsline{toc}{section}{Append-only continuation marker: V8}
\textbf{End of V8.}  The complete V1--V7 body above is retained verbatim.  V8
starts a strictly stronger post-E1--E6 programme and proves microscopic
regulator universality for an explicit range-two covariant deformation of the
V2 Wilson/overlap kernel.  The proof uses a uniform native gap, relative overlap
functional calculus, regulator-marked common forests, a single finite local
normalization/source/line flow, and the two-scale estimate
$O(aM+\mu/M)$ optimized at $O(\sqrt{a\mu})$.  It also proves a general
universality criterion for any explicitly verified admissible gapped regulator
homotopy.  Future versions must append after this marker.



\clearpage
\section{V9 continuation: regulator-phase classification, weak topological obstructions, and a connected Clifford--Wilson sector}
\label{sec:v9-continuation}

V8 proved that the endpoints of any explicitly verified admissible gapped
regulator homotopy have the same fixed-order continuum after one finite local
matching map.  The remaining question is therefore topological before it is
analytic: when does such a homotopy exist?  The strongest possible guess,
``equal four-dimensional overlap index implies connectedness'', is false for a
general multiband Hermitian lattice symbol.  This continuation first identifies
the missing invariant exactly.  It then proves a positive connectedness theorem
for the Wilson-like five-Clifford/no-doubler sector actually containing the V2
kernel.

The distinction is load-bearing.  A general gapped Hermitian symbol defines a
vector bundle over the Brillouin torus, and its complete stable class belongs to
$K^0(\mathbb T^4)$; the four-dimensional strong number sees only one component
of that class.  By contrast, a five-Clifford symbol factors through $S^4$, so
all weak first-Chern classes vanish automatically.  In the no-doubler component
with the physical Dirac germ fixed, the complement of the physical south-pole
preimage is contractible in the target, and a direct relative homotopy can be
constructed.  A quantitative trigonometric approximation then converts the
smooth homotopy into a finite chain of finite-range native homotopies with the
same physical first jet.

Throughout V9, $x=ap\in\mathbb T^4=(-\pi,\pi]^4$ denotes the dimensionless
Brillouin variable.  The positive IR/reference scale and the protected
source/line chart of V1--V8 remain in force.

\subsection{The flattened bundle of a general gapped Hermitian symbol}

\begin{definition}[Flat gapped symbol and flattened projector]
\label{def:v9-flat-symbol}
Let $H:\mathbb T^4\to\operatorname{Herm}(M)$ be a smooth periodic Hermitian
matrix symbol satisfying
\begin{equation}
 \inf_{x\in\mathbb T^4}s_{\min}(H(x))\ge g>0.
 \label{eq:v9-gap}
\end{equation}
Its spectral flattening and positive projector are
\begin{equation}
 \varepsilon_H(x):=\operatorname{sgn}H(x),
 \qquad
 P_H(x):=\frac12(I+\varepsilon_H(x)).
 \label{eq:v9-flattening}
\end{equation}
The associated positive-energy bundle is
\begin{equation}
 E_H:=\operatorname{Ran}P_H\longrightarrow\mathbb T^4.
 \label{eq:v9-bundle}
\end{equation}
A finite-range translation-invariant kernel is the special case in which the
entries of $H$ are trigonometric polynomials.
\end{definition}

\begin{lemma}[Rank is constant on a gapped component]
\label{lem:v9-rank}
For a gapped symbol in Definition~\ref{def:v9-flat-symbol},
$\operatorname{rank}P_H(x)$ is independent of $x$.  Along a continuous gapped
homotopy $H_t$, it is also independent of $t$.
\end{lemma}
\begin{proof}
The rank is an integer-valued continuous function of $x$ because the gap keeps
zero out of the spectrum.  The torus is connected, hence the rank is constant.
The same argument on $\mathbb T^4\times[0,1]$ gives the homotopy statement.
\end{proof}

\begin{theorem}[The $K$-class is a necessary invariant of every gapped regulator homotopy]
\label{thm:v9-K-obstruction}
Let $H_t$, $0\le t\le1$, be a continuous gapped homotopy of smooth Hermitian
symbols.  Then
\begin{equation}
 \boxed{[E_{H_0}]=[E_{H_1}]\quad\text{in }K^0(\mathbb T^4).}
 \label{eq:v9-K-invariant}
\end{equation}
In particular, every Chern number of $E_{H_t}$ is independent of $t$.
\end{theorem}
\begin{proof}
Functional calculus makes $(x,t)\mapsto P_{H_t}(x)$ continuous, and smooth when
the homotopy is smooth.  Its range therefore defines one vector bundle
$\mathcal E\to\mathbb T^4\times[0,1]$.  Restriction to the two boundary faces
gives $E_{H_0}$ and $E_{H_1}$.  Since the inclusions
$\mathbb T^4\times\{0\},\mathbb T^4\times\{1\}\hookrightarrow
\mathbb T^4\times[0,1]$ are homotopy equivalences, both restrictions have the
same $K^0$ class.  Naturality of Chern classes gives the last statement.
\end{proof}

For later use, choose the four coordinate one-cycles and orient the six
coordinate two-tori $\mathbb T^2_{\mu\nu}$ by $dx_\mu\wedge dx_\nu$.  For a
complex bundle $E$ define
\begin{align}
 r(E)&:=\operatorname{rank}E,
 \label{eq:v9-rank-invariant}\\
 \nu_{\mu\nu}(E)&:=
 \left\langle c_1(E),[\mathbb T^2_{\mu\nu}]\right\rangle,
 \qquad 1\le\mu<\nu\le4,
 \label{eq:v9-weak-Chern}\\
 \nu_4(E)&:=
 \left\langle \operatorname{ch}_2(E),[\mathbb T^4]\right\rangle.
 \label{eq:v9-strong-Chern}
\end{align}
On the four-torus the last pairing is integral.  One may equivalently use
$c_2$ together with the six first-Chern numbers.

\begin{theorem}[Stable classification on the four-torus]
\label{thm:v9-KT4}
There is a natural group isomorphism
\begin{equation}
 \boxed{
 K^0(\mathbb T^4)
 \cong
 \mathbb Z
 \oplus \bigwedge
olimits^2\mathbb Z^4
 \oplus \bigwedge
olimits^4\mathbb Z^4
 \cong \mathbb Z^8.}
 \label{eq:v9-KT4}
\end{equation}
Under this identification a bundle class is determined by
\begin{equation}
 \boxed{
 \left(r(E),\{\nu_{\mu\nu}(E)\}_{\mu<\nu},\nu_4(E)\right).}
 \label{eq:v9-K-invariants}
\end{equation}
Consequently, after fixing the spectral rank, a general gapped Hermitian symbol
has seven independent stable topological integers: six weak two-dimensional
numbers and one strong four-dimensional number.
\end{theorem}
\begin{proof}
The circle has
$K^0(S^1)\cong\mathbb Z$ and $K^1(S^1)\cong\mathbb Z$.  Since these groups are
free abelian, the Kunneth sequence for complex $K$-theory has no torsion term.
Iterating it four times gives
\[
 K^0((S^1)^4)
 \cong\bigoplus_{k\ \mathrm{even}}\bigwedge^k\mathbb Z^4
 =\bigwedge^0\mathbb Z^4\oplus
  \bigwedge^2\mathbb Z^4\oplus
  \bigwedge^4\mathbb Z^4.
\]
The three summands have ranks $1,6,1$.  The Chern character of the exterior
products of the four $K^1(S^1)$ generators is the corresponding exterior
product of the normalized integral one-forms.  Hence the degree-zero, degree-two
and degree-four components are read exactly by rank, the six pairings of
$c_1$, and the degree-four Chern-character pairing.  Because the displayed
$K$-group is torsion free, these integer coordinates determine the class.
\end{proof}

\begin{corollary}[Equality of the strong index alone is not a connectedness criterion]
\label{cor:v9-strong-not-enough}
For general gapped Hermitian symbols, equality of rank and $\nu_4$ does not
imply the existence of a gapped homotopy.  Any mismatch in one of the six
integers $\nu_{\mu\nu}$ is an obstruction.
\end{corollary}
\begin{proof}
Immediate from Theorems~\ref{thm:v9-K-obstruction} and
\ref{thm:v9-KT4}.
\end{proof}

\subsection{A concrete weak-index obstruction at fixed strong number}

The preceding obstruction is not merely formal.  Let
$\pi_{12}:\mathbb T^4\to\mathbb T^2$ be projection onto the first two
coordinates.  Choose a complex line bundle $L\to\mathbb T^2$ with
$c_1(L)$ equal to the positive generator and put
\begin{equation}
 E_1:=\pi_{12}^*L,
 \qquad
 E_0:=\underline{\mathbb C}.
 \label{eq:v9-line-example}
\end{equation}
Then
\begin{equation}
 r(E_0)=r(E_1)=1,
 \qquad
 \nu_4(E_0)=\nu_4(E_1)=0,
 \qquad
 \nu_{12}(E_0)=0,
 \quad
 \nu_{12}(E_1)=1.
 \label{eq:v9-line-invariants}
\end{equation}
The equality of the four-dimensional number follows because
$c_1(E_1)$ is pulled back from a two-dimensional factor, hence
$c_1(E_1)^2=0$ and $c_2(E_1)=0$.

\begin{proposition}[The weak obstruction is realized by finite-range gapped symbols]
\label{prop:v9-finite-range-obstruction}
There exist finite-dimensional, finite-range translation-invariant Hermitian
symbols $H_0,H_1$ whose positive bundles are $E_0,E_1$ after stabilization, so
that they have the same spectral rank and strong number but cannot be joined by
any gapped Hermitian homotopy.
\end{proposition}
\begin{proof}
Every complex bundle on the compact torus embeds in a sufficiently large
trivial bundle.  Let $P_i(x)$ be smooth orthogonal projector representatives of
$E_i$ in one common trivial $\mathbb C^M$ bundle and put
$Q_i=2P_i-I$.  Then $Q_i^2=I$, so each $Q_i$ is Hermitian and has unit gap.
Trigonometric matrix polynomials are uniformly dense in smooth periodic matrix
functions together with any prescribed finite number of derivatives.  Choose a
Hermitian trigonometric polynomial $\widetilde Q_i$ with
$\|\widetilde Q_i-Q_i\|<1/4$.  The straight line from $Q_i$ to
$\widetilde Q_i$ stays invertible by the singular-value perturbation inequality,
so spectral flattening shows that $\widetilde Q_i$ has the same positive-bundle
class as $Q_i$.  Trigonometric polynomial symbols are precisely finite-range
translation-invariant kernels.  If a gapped homotopy joined
$\widetilde Q_0$ to $\widetilde Q_1$, Theorem~\ref{thm:v9-K-obstruction}
would force their $K$-classes, and hence $\nu_{12}$, to agree, a contradiction.
\end{proof}

\begin{firewall}[Scope of the counterexample]
Proposition~\ref{prop:v9-finite-range-obstruction} disproves a connectedness
claim for the \emph{general} gapped Hermitian symbol space.  It does not assert
that the extra weak phase is compatible with every additional native condition
imposed on the physical Standard-Model Wilson block.  Those additional
conditions are analyzed next.  In particular, the V2 kernel belongs to a much
smaller five-Clifford/no-doubler class in which the weak invariants vanish.
\end{firewall}

\subsection{The five-Clifford Wilson class}

Fix Hermitian matrices $\Gamma_1,\ldots,\Gamma_5$ satisfying
\begin{equation}
 \Gamma_A\Gamma_B+\Gamma_B\Gamma_A=2\delta_{AB}I.
 \label{eq:v9-Clifford}
\end{equation}
The V2 flat Hermitian Wilson symbol can be written in such a basis after the
standard identification
$\Gamma_\mu=i\gamma_5\gamma_\mu$ for $1\le\mu\le4$ and
$\Gamma_5=\gamma_5$.

\begin{definition}[Five-Clifford gapped symbol]
\label{def:v9-Clifford-symbol}
A five-Clifford symbol is
\begin{equation}
 H_h(x):=\sum_{A=1}^5h_A(x)\Gamma_A,
 \qquad
 h:\mathbb T^4\to\mathbb R^5,
 \label{eq:v9-Clifford-H}
\end{equation}
with $|h(x)|\ge g>0$.  Its normalized map is
\begin{equation}
 n_h(x):=\frac{h(x)}{|h(x)|}\in S^4.
 \label{eq:v9-nmap}
\end{equation}
Let $s:=-e_5\in S^4$ denote the physical south pole.  We call the symbol
\emph{single-cone/no-doubler} if
\begin{equation}
 n_h^{-1}(s)=\{0\},
 \label{eq:v9-no-doubler}
\end{equation}
and the tangent derivative of the first four components at $0$ is invertible.
Two such symbols have the same \emph{physical Dirac germ} if they have the same
value $h(0)$ up to the fixed positive radial normalization and the same first
jet that determines the continuum Dirac speed/coframe.
\end{definition}

The Clifford relation gives
\begin{equation}
 H_h(x)^2=|h(x)|^2I,
 \qquad
 \operatorname{sgn}H_h(x)=n_h(x)\cdot\Gamma.
 \label{eq:v9-Clifford-flat}
\end{equation}
Thus the topology of the flattened symbol is the topology of a map to $S^4$.

\begin{lemma}[Weak Chern numbers vanish in the five-Clifford class]
\label{lem:v9-weak-vanish}
For a five-Clifford gapped symbol, the positive bundle satisfies
\begin{equation}
 \boxed{c_1(E_{H_h})=0.}
 \label{eq:v9-c1zero}
\end{equation}
Hence all six weak numbers $\nu_{\mu\nu}$ vanish.
\end{lemma}
\begin{proof}
The positive projector is the pullback by $n_h$ of the positive eigenspace
bundle $\mathcal S_+\to S^4$ of $n\cdot\Gamma$.  Since
$H^2(S^4;\mathbb Z)=0$, one has $c_1(\mathcal S_+)=0$.  Naturality gives
$c_1(E_{H_h})=n_h^*c_1(\mathcal S_+)=0$.
\end{proof}

\begin{theorem}[Degree is the complete flattened homotopy invariant in the five-Clifford class]
\label{thm:v9-degree-classification}
Based homotopy classes of normalized five-Clifford symbols satisfy
\begin{equation}
 \boxed{
 [\mathbb T^4,S^4]_*\cong H^4(\mathbb T^4;\mathbb Z)\cong\mathbb Z.}
 \label{eq:v9-degree-classification}
\end{equation}
The integer is the degree of $n_h$, equivalently the strong Chern number of the
pulled-back positive bundle up to the fixed sign convention of the chosen
Clifford representation.
\end{theorem}
\begin{proof}
The sphere $S^4$ is $3$-connected and has
$\pi_4(S^4)\cong\mathbb Z$.  Give $\mathbb T^4$ a CW structure with one
zero-cell at the chosen base point.  Obstruction theory extends a based map
uniquely up to homotopy over the $3$-skeleton because
$\pi_k(S^4)=0$ for $k<4$.  The only obstruction/difference class appears on the
four-cells and lies in
$H^4(\mathbb T^4;\pi_4(S^4))\cong H^4(\mathbb T^4;\mathbb Z)$.  Evaluation on
the fundamental class is precisely the degree.  The Chern-number statement
follows because $H^4(S^4;\mathbb Z)$ is generated by the second Chern/Chern
character class of the positive Clifford eigenspace bundle.
\end{proof}

For the no-doubler sector one can do better than abstract obstruction theory.
The physical south pole is attained at exactly one point.  Once the two local
Dirac germs have been aligned, the rest of the target is
$S^4\setminus\{s\}\cong\mathbb R^4$, which is contractible.

\subsection{Local germ straightening without creating a doubler}

Write $B_\rho\subset\mathbb T^4$ for a sufficiently small coordinate ball
around $0$.  In geodesic coordinates
$\exp_s:T_sS^4\supset U\to S^4$ near the south pole, a single-cone map has
\begin{equation}
 \exp_s^{-1}(n_h(x))=Ax+R(x),
 \qquad
 R(x)=O(|x|^2),
 \label{eq:v9-local-germ}
\end{equation}
where $A:\mathbb R^4\to T_sS^4\cong\mathbb R^4$ is invertible.

\begin{lemma}[Quantitative local straightening]
\label{lem:v9-local-straightening}
Let $n_0,n_1$ be two smooth single-cone five-Clifford normalized symbols with
the same physical first jet $A$ at $0$.  There is $\rho>0$ and homotopies,
supported in $B_\rho$, which:
\begin{enumerate}[label=\textup{(S\arabic*)},leftmargin=3em]
\item preserve $n_i(0)=s$ and the derivative $A$ at all times;
\item never create another preimage of $s$;
\item leave each map unchanged outside $B_\rho$;
\item make both final maps equal to one common canonical germ on a smaller
ball $B_{\rho/3}$.
\end{enumerate}
\end{lemma}
\begin{proof}
Because the two maps have the same invertible derivative, in the exponential
chart they are $Ax+R_i(x)$ with
$|R_i(x)|\le C|x|^2$ and
$\|DR_i(x)\|\le C|x|$ after shrinking $\rho$.  Choose a smooth cutoff
$\chi$ equal to one on $B_{\rho/3}$ and zero outside $B_{2\rho/3}$, and define
\begin{equation}
 f_{i,u}(x):=Ax+(1-u\chi(x))R_i(x),
 \qquad 0\le u\le1.
 \label{eq:v9-local-straightening}
\end{equation}
For $\rho$ small enough,
$|f_{i,u}(x)|\ge c|x|$ uniformly for $x\ne0$, because the quadratic remainder
is dominated by the smallest singular value of $A$.  Thus
$\exp_s(f_{i,u}(x))=s$ only at $x=0$.  The correction is quadratic, so the
first derivative at the origin remains $A$.  At $u=1$ both maps equal
$\exp_s(Ax)$ on $B_{\rho/3}$ and are unchanged outside $B_{2\rho/3}$.
\end{proof}

\begin{theorem}[Smooth no-doubler connectedness in the five-Clifford sector]
\label{thm:v9-smooth-connectedness}
Let $n_0,n_1:\mathbb T^4\to S^4$ be smooth single-cone/no-doubler maps with the
same physical Dirac germ.  Then there exists a smooth homotopy $n_t$ such that
\begin{equation}
 \boxed{
 n_t^{-1}(s)=\{0\}
 \quad\text{for every }t\in[0,1],}
 \label{eq:v9-no-doubler-path}
\end{equation}
while the physical first jet at $0$ is fixed throughout.
\end{theorem}
\begin{proof}
Apply Lemma~\ref{lem:v9-local-straightening} to make $n_0$ and $n_1$ agree on a
closed ball $B\subset B_{\rho/3}$, without changing either map outside the
larger local ball and without producing a second south-pole preimage.  Put
$X:=\mathbb T^4\setminus\operatorname{int}B$.  Both restrictions map $X$ into
$S^4\setminus\{s\}$ and agree on $\partial X$.

Choose a fixed stereographic diffeomorphism
\begin{equation}
 \psi:S^4\setminus\{s\}\longrightarrow\mathbb R^4.
 \label{eq:v9-stereographic}
\end{equation}
Because the boundary restrictions agree, the linear homotopy
\begin{equation}
 F_t(x):=(1-t)\psi(n_0(x))+t\psi(n_1(x)),
 \qquad x\in X,
 \label{eq:v9-stereo-homotopy}
\end{equation}
is fixed on $\partial X$.  Define
$n_t=\psi^{-1}(F_t)$ on $X$ and use the common map on $B$.  The two pieces glue
smoothly after the local straightening is chosen constant on a collar of
$\partial B$.  The image on $X$ never contains $s$ because
$\psi^{-1}$ takes values in $S^4\setminus\{s\}$, while the common map on $B$
attains $s$ only at $0$.  The homotopy is stationary near $0$, so the physical
first jet is fixed.
\end{proof}

\begin{corollary}[The strong degree is forced by a single physical cone]
\label{cor:v9-degree-single-cone}
In the orientation convention fixed by the common physical Dirac germ, every
single-cone map of Theorem~\ref{thm:v9-smooth-connectedness} has the same degree,
namely the local degree of the physical south-pole preimage.  In particular the
five-Clifford no-doubler component has no independent weak or strong topological
label once the physical germ is fixed.
\end{corollary}
\begin{proof}
Take $s$ as a regular value after an arbitrarily small germ-preserving
perturbation if necessary.  The degree is the signed sum of the local degrees
at the preimages of $s$.  There is exactly one preimage, and its local degree is
the sign of the determinant of the fixed tangent derivative.  Homotopy
invariance removes the auxiliary perturbation.
\end{proof}

\subsection{From a smooth homotopy to a finite chain of finite-range homotopies}

Theorem~\ref{thm:v9-smooth-connectedness} is a continuum statement on the
Brillouin torus.  V8, however, requires a native finite-range or uniformly
exponentially local path with quantitative source jets.  The next step is a
relative approximation theorem which preserves the physical jet and the
no-doubler margin.

\begin{lemma}[Trigonometric approximation with exact finite jet]
\label{lem:v9-Hermite-Fourier}
Let $f:\mathbb T^4\to\mathbb R^m$ be $C^k$, $k\ge1$, and prescribe the value
and first derivative of $f$ at $0$.  For every $\epsilon>0$ there is a real
trigonometric polynomial $q$ such that
\begin{equation}
 q(0)=f(0),
 \qquad
 Dq(0)=Df(0),
 \qquad
 \|q-f\|_{C^1}<\epsilon.
 \label{eq:v9-Hermite-Fourier}
\end{equation}
The same statement holds for Hermitian matrix-valued functions while preserving
Hermiticity.
\end{lemma}
\begin{proof}
Ordinary trigonometric polynomials are dense in $C^1(\mathbb T^4)$.  Start with
a polynomial $q_0$ within $\delta$ of $f$.  The finite jet error
\[
 \bigl(f(0)-q_0(0),\,Df(0)-Dq_0(0)\bigr)
\]
is corrected by a fixed finite collection of trigonometric polynomials whose
value/first-derivative evaluation matrix at $0$ is the identity.  For example,
constants correct values and the functions $\sin x_\mu$ correct first
derivatives; additional components are corrected independently.  The
correction norm is bounded by a fixed constant times the jet error, which is
$O(\delta)$.  Choose $\delta$ sufficiently small.  For a Hermitian matrix,
approximate the independent real diagonal and real/imaginary off-diagonal
components and pair Fourier coefficients by adjunction.
\end{proof}

\begin{lemma}[Quantitative openness of the single-cone conditions]
\label{lem:v9-no-doubler-open}
Let $n_t$ be the smooth homotopy from
Theorem~\ref{thm:v9-smooth-connectedness}.  There are a ball $B_0$ around zero
and constants $c_0,\delta_0>0$ such that uniformly in $t$:
\begin{enumerate}[label=\textup{(O\arabic*)},leftmargin=3em]
\item in $B_0$, the first four tangent components satisfy
$|n_t^{\top}(x)|\ge c_0|x|$;
\item on $\mathbb T^4\setminus B_0$,
$\operatorname{dist}(n_t(x),s)\ge\delta_0$.
\end{enumerate}
Consequently every sufficiently small $C^1$ perturbation with the same value
and first derivative at $0$ is again single-cone/no-doubler.
\end{lemma}
\begin{proof}
The homotopy is stationary with invertible tangent derivative near $0$, giving
(O1) after shrinking $B_0$.  The compact set
$(\mathbb T^4\setminus B_0)\times[0,1]$ has image disjoint from $s$, so its
distance from $s$ has a positive minimum, proving (O2).  A sufficiently small
$C^1$ perturbation preserves the lower derivative bound in $B_0$ and the
positive target-space separation on the complement.
\end{proof}

\begin{theorem}[Finite-chain finite-range connectedness]
\label{thm:v9-finite-chain}
Let $H^{(0)}$ and $H^{(1)}$ be finite-range five-Clifford symbols satisfying:
\begin{enumerate}[label=\textup{(C\arabic*)},leftmargin=3em]
\item a positive flat Wilson gap;
\item the same physical Dirac first jet at $x=0$;
\item one and only one physical overlap zero, equivalently
$n_{H^{(i)}}^{-1}(s)=\{0\}$.
\end{enumerate}
Then there exist finitely many finite-range five-Clifford symbols
\begin{equation}
 H^{(0)}=K_0,K_1,\ldots,K_m=H^{(1)}
 \label{eq:v9-chain}
\end{equation}
and finite-range linear homotopies
$K_{j,u}=(1-u)K_j+uK_{j+1}$ such that every $K_{j,u}$ is gapped,
single-cone/no-doubler, and has the same physical first jet.
\end{theorem}
\begin{proof}
Use Theorem~\ref{thm:v9-smooth-connectedness} to obtain a smooth no-doubler
homotopy $n_t$.  Choose positive radial functions at the endpoints matching the
original symbols and extend them smoothly and positively along the homotopy;
this gives a smooth vector homotopy $h_t$ with
$H_t=h_t\cdot\Gamma$, a uniform gap
$\inf_{x,t}|h_t(x)|=:g_*>0$, and the same physical first jet.

Uniform continuity in $C^1$ allows a partition
$0=t_0<\cdots<t_m=1$ so fine that adjacent $h_{t_j}$ differ by less than a
small tolerance $\epsilon$ in $C^1$.  Keep the endpoints exact.  At each
interior node apply Lemma~\ref{lem:v9-Hermite-Fourier} to choose a real
trigonometric polynomial vector $q_j$ with the exact common value/first jet at
zero and
$\|q_j-h_{t_j}\|_{C^1}<\epsilon$.  Put $q_0=h_0$ and $q_m=h_1$.

Choose $\epsilon$ smaller than one quarter of the uniform radial gap and smaller
than the openness margin in Lemma~\ref{lem:v9-no-doubler-open}.  Every convex
segment
$q_{j,u}=(1-u)q_j+uq_{j+1}$ then stays uniformly close to the corresponding
short portion of $h_t$.  Hence $|q_{j,u}|\ge g_*/2$ and the normalized map
$q_{j,u}/|q_{j,u}|$ remains single-cone/no-doubler.  Since all node
approximations have the exact same value and first derivative at zero, every
linear segment preserves that physical jet.  Finally
$K_{j,u}=q_{j,u}\cdot\Gamma$ is a trigonometric polynomial symbol, hence a
finite-range kernel.  The union of the finitely many segments gives the
claimed chain.
\end{proof}

\subsection{Native gauge/spin covariant lifting of the finite chain}

A trigonometric polynomial is not yet a native background kernel.  We now fix a
canonical lifting.  For a displacement $n\in\mathbb Z^4$, choose once and for
all a lattice path $\wp(n)$ from $x$ to $x+n$ and let
$T_{\wp(n)}^{\mathcal U}$ be the ordered gauge--spin parallel transport followed
by the shift.  Reverse paths are chosen so that
$(T_{\wp(n)}^{\mathcal U})^*=T_{\wp(-n)}^{\mathcal U}$.

\begin{definition}[Canonical covariantization]
\label{def:v9-covariantization}
For a Hermitian trigonometric polynomial symbol
\begin{equation}
 K(x)=\sum_{n\in F}A_ne^{in\cdot x},
 \qquad A_{-n}=A_n^*,
 \label{eq:v9-Fourier-symbol}
\end{equation}
with finite $F$, define $\mathcal C(K)$ by replacing each Fourier shift
$e^{in\cdot x}$ with the corresponding gauge--spin parallel shift
$T_{\wp(n)}^{\mathcal U}$ and pairing every term with its Hermitian reverse.
Clifford-vector coefficients are interpreted in the local coframe/spin frame
exactly as in the V2 anticommutator prescription.  The zero-displacement scalar
part acts onsite.
\end{definition}

\begin{lemma}[Properties of canonical covariantization]
\label{lem:v9-covariantization}
For every finite polynomial $K$:
\begin{enumerate}[label=\textup{(V\arabic*)},leftmargin=3em]
\item $\mathcal C(K)$ is finite range and Hermitian;
\item it transforms by exact local gauge/spin unitary similarity;
\item at the flat identity-link/reference-coframe background its symbol is
exactly $K$;
\item every fixed gauge, spin-connection and coframe source derivative is again
a finite-range kernel with a bound depending only on the finite Fourier support
and derivative order.
\end{enumerate}
\end{lemma}
\begin{proof}
A finite path product has finite range and transforms with the gauge/spin matrix
at its two endpoints.  Pairing a path with its reverse gives Hermiticity.
At identity transport the path operator is the ordinary displacement shift and
therefore has Fourier multiplier $e^{in\cdot x}$.  Differentiating a finite path
product inserts finitely many differentiated link/spin factors; at fixed source
order the Leibniz expansion is finite.  Coframe derivatives act on the same
local Clifford coefficients used by V2.  This proves all four statements.
\end{proof}

\begin{theorem}[Native lifting theorem for the canonical Clifford--Wilson class]
\label{thm:v9-native-lift}
Let the two endpoints of Theorem~\ref{thm:v9-finite-chain} be equipped with the
canonical covariantization of Definition~\ref{def:v9-covariantization}.  There
exists a common compact perturbative background/source chart, obtained by one
finite shrinking of the V1--V8 chart if necessary, on which every segment of
the covariantized finite chain is a V8-admissible regulator homotopy at every
prescribed fixed coefficient order $\mathfrak q$.
\end{theorem}
\begin{proof}
There are only finitely many segments, each with finite range and a positive
flat Wilson gap.  Their minimum flat gap is therefore a positive number
$g_{\rm chain}$.  By Lemma~\ref{lem:v9-covariantization}, the operator depends
continuously on the finitely many background/source coordinates.  Shrink one
common compact chart until the operator perturbation from the flat symbol is
less than $g_{\rm chain}/4$ for every segment and homotopy parameter.  The
singular-value perturbation inequality then gives a uniform Wilson gap.

The exact common value and first jet at $x=0$ imply that the derivative with
respect to the segment parameter has zero constant and linear physical
momentum jet.  Hence its physical insertion is $O(|ap|^2)$, which is precisely
the one-irrelevant-order gain required by item (A3) of
Definition~\ref{def:v8-admissible-homotopy}.  The quantitative no-doubler
margin of Theorem~\ref{thm:v9-finite-chain}, compactness of the full zone away
from the physical ball, and the same chart shrinking give a uniform hard
inverse on the complementary zone.  Finite range and smooth periodic symbols
supply item (A4).

For the retained positive IR/reference scale, hard shells exclude the physical
zero.  The hard chiral pivot and determinant/Pfaffian nonzero condition are
open conditions.  Because the chain and source parameter sets are compact and
finite, shrink the protected line/source chart once more so the common pivot is
positive on every segment.

For the normalization gate, use the chronological V7 local coordinates in
which the finite normalization functionals are dual to the retained local
counterterm bank.  Adding a local coordinate $c^a\mathcal O_a$ changes the
corresponding normalization row by $c^a$; hence the Jacobian on the local
counterterm directions has the regulator-independent identity block.  Regulator
variation changes the inhomogeneous loop coefficients but not this coefficientwise
rank statement.  Redundant/BV-exact and source/line coordinates are carried in
the same finite triangular ordering.  Thus the finite coefficientwise
normalization/range problem retains constant rank along every segment, giving
item (A6) without assuming a finite-amplitude contraction of the regulator
flow.  If one subsequently chooses a nonlinear physical reparameterization of
these local coordinates, one restricts to the common nonsingular chart exactly
as in V7.
All V8 admissibility requirements are then satisfied simultaneously.
\end{proof}

\subsection{Regulator-class universality inside the canonical Clifford--Wilson component}

\begin{definition}[Canonical single-cone Clifford--Wilson regulator class]
\label{def:v9-canonical-class}
Let $\mathfrak W_{\rm Cliff}^{\rm sc}$ be the class of canonically covariantized
finite-range five-Clifford kernels whose flat symbols have a positive Wilson
gap, the fixed V2 physical Dirac germ, and exactly one physical overlap zero at
$x=0$, and which lie on the common protected source/line component after the
finite chart shrinking of Theorem~\ref{thm:v9-native-lift}.
\end{definition}

\begin{theorem}[Connectedness of the canonical single-cone Clifford--Wilson class]
\label{thm:v9-class-connected}
Any two regulators $A,B\in\mathfrak W_{\rm Cliff}^{\rm sc}$ are connected by a
finite chain of V8-admissible native gapped homotopies at every prescribed
fixed coefficient order $\mathfrak q$.
\end{theorem}
\begin{proof}
Apply Theorem~\ref{thm:v9-finite-chain} to their flat symbols and then
Theorem~\ref{thm:v9-native-lift} to every segment.  The finite chain may be
traversed consecutively; endpoint charts can be replaced by their finite
intersection, which still contains the common reference point and remains a
compact protected perturbative chart.
\end{proof}

\begin{theorem}[Regulator-class universality in the canonical Clifford--Wilson sector]
\label{thm:v9-class-universality}
Fix finite
$\mathfrak q=(L,\Delta,r,N,\mathcal P_{\rm ext})$.  For any two regulators
$A,B\in\mathfrak W_{\rm Cliff}^{\rm sc}$ there is a finite local invertible
matching map
\begin{equation}
 F_{\infty,\mathfrak q}^{B\leftarrow A}
 \label{eq:v9-class-map}
\end{equation}
such that
\begin{equation}
 \boxed{
 \widehat{\mathfrak C}^{B}_{\infty,\lambda}
 =F_{\infty,\mathfrak q}^{B\leftarrow A}
  \widehat{\mathfrak C}^{A}_{\infty,\lambda}.}
 \label{eq:v9-class-universality}
\end{equation}
After covariant operator/source/line matching, the normalized connected
coefficients, stress/composite/contact coefficients, ghosts/antifields,
BV/BRST/master relations and determinant/Pfaffian line insertions agree.  The
maps can be chosen strictly projective in perturbative order.
\end{theorem}
\begin{proof}
By Theorem~\ref{thm:v9-class-connected}, choose a finite chain
$A=R_0,R_1,\ldots,R_m=B$ of V8-admissible homotopies.  Apply
Theorem~\ref{thm:v8-general-homotopy} to each segment, obtaining invertible
continuum matching maps
$F^{R_{j+1}\leftarrow R_j}_{\infty,\mathfrak q}$.  Their finite composition is
$F^{B\leftarrow A}_{\infty,\mathfrak q}$ and satisfies
\eqref{eq:v9-class-universality}.  The V8 map for each segment is chronological
and projective in $L$; a finite composition of projective maps is projective.
The listed coefficient systems are exactly the coordinates transported in V8.
\end{proof}

\begin{corollary}[Path independence up to the V7 finite scheme groupoid]
\label{cor:v9-path-independence}
Two different admissible finite chains from $A$ to $B$ produce continuum
matching maps which differ, if at all, by an admissible finite continuum scheme
transition of V7 fixing the chosen physical normalization coordinates.  Hence
the coordinate-free matched continuum relations are independent of the chosen
chain.
\end{corollary}
\begin{proof}
Compose one chain with the reverse of the other.  This is a closed sequence of
finite regulator/scheme transports starting and ending at the same native
continuum theory.  V7 identifies any residual local coordinate transformation
with an invertible finite scheme transition satisfying the cocycle law.  The
coordinate-free normalized connected functional and master relations are
covariant under that groupoid, so they are unchanged.
\end{proof}

\subsection{What remains outside the Clifford component}

Theorems~\ref{thm:v9-KT4} and \ref{thm:v9-class-connected} together give the
correct classification picture at the present level:
\begin{equation}
 \boxed{
 \begin{array}{c}
 \text{general gapped Hermitian symbols: full }K^0(\mathbb T^4)
 \text{ data are necessary;}\\[0.25em]
 \text{canonical single-cone five-Clifford symbols: the weak data vanish and}
 \text{ the fixed physical germ gives one connected component.}
 \end{array}}
 \label{eq:v9-classification-summary}
\end{equation}

What has \emph{not} been proved is that equality of the stable $K^0$ class is
sufficient for a V8-admissible no-doubler homotopy in a completely general
multiband native regulator class.  Stable $K$-theory allows the addition of
trivial bands and forgets possible unstable finite-rank information.  Moreover,
V8 requires a fixed physical cone, hard-pivot margins and source/line
compatibility, not merely a homotopy of abstract projectors.

\begin{firewall}[Stable classification is not yet general native connectedness]
Equality in $K^0(\mathbb T^4)$ is a complete stable bundle invariant, but V9
does not identify it with connectedness of every finite-rank native
Standard-Model regulator.  A proof of such a statement would have to control
unstable Grassmannian homotopy, preserve the physical no-doubler germ, lift the
projector homotopy to a Hermitian kernel with the required local source jets,
and maintain the determinant/Pfaffian hard pivots.  None of those extra gates is
silently inferred from stable $K$-theory.
\end{firewall}

\subsection{V9 anti-circularity audit}
\label{sec:v9-audit}

\begin{enumerate}[leftmargin=2.4em]
\item \textbf{The four-dimensional index was not assumed complete.}
The full flattened projector bundle was examined first.  The six weak Chern
numbers give explicit additional obstructions in the general Hermitian class.
\item \textbf{The positive Clifford result uses extra structure openly.}
Weak invariants vanish because the projector factors through $S^4$, not because
they were ignored.
\item \textbf{No-doubler connectedness is proved directly.}
After aligning the local germ, the complement maps into
$S^4\setminus\{s\}\cong\mathbb R^4$ and is homotoped there, so the physical
south pole is never crossed away from $x=0$.
\item \textbf{Finite range is recovered, not assumed from a smooth homotopy.}
The Hermite--Fourier approximation preserves the physical value/first jet, and
a finite partition plus openness margins converts the smooth path to a finite
chain of trigonometric-polynomial paths.
\item \textbf{Native gauge/spin covariance is restored before invoking V8.}
The canonical path-transport lift is finite range and exactly covariant; the
common protected chart is then shrunk once to preserve all segment gaps and
hard pivots.
\item \textbf{Stable $K$-theory is not overpromoted.}
Outside the canonical Clifford--Wilson class, V9 records equality of the full
$K$-class as necessary and stable-complete, but does not claim it is sufficient
for a native no-doubler path at fixed finite rank.
\end{enumerate}

\section{V9 status ledger}
\label{sec:v9-status}

\subsection*{Proved in V9}
\begin{enumerate}[leftmargin=2.2em]
\item Every gapped Hermitian regulator homotopy preserves the $K^0$ class of the
flattened positive-energy bundle.
\item $K^0(\mathbb T^4)\cong\mathbb Z^8$; after fixing rank, the stable class has
six weak first-Chern integers and one strong four-dimensional integer.
\item Equality of the strong four-dimensional number alone is not sufficient in
the general Hermitian class.  A concrete weak-index obstruction is realized by
finite-range gapped translation-invariant symbols.
\item In the five-Clifford class all weak first-Chern numbers vanish, and based
flattened homotopy classes are classified by the degree
$[\mathbb T^4,S^4]_*\cong\mathbb Z$.
\item For the single-cone/no-doubler five-Clifford sector with a fixed physical
Dirac germ, any two smooth symbols are connected through a smooth homotopy that
keeps the physical first jet fixed and never creates another overlap zero.
\item That smooth homotopy can be replaced by a finite chain of finite-range
five-Clifford homotopies preserving a positive gap, the exact physical first
jet and the no-doubler condition.
\item Canonical gauge/spin covariantization lifts the finite chain to native
finite-range kernels.  After one finite shrinking of the protected chart, each
segment satisfies the V8 admissible-homotopy gates at every prescribed fixed
coefficient order.
\item Therefore the canonical single-cone Clifford--Wilson regulator class is
connected by finite chains of V8-admissible homotopies, and all regulators in
that class have the same matched coefficientwise continuum after one finite
local scheme/source/line map.  The equivalence is projective in perturbative
order.
\end{enumerate}

\subsection*{Inherited}
\begin{enumerate}[leftmargin=2.2em]
\item V1--V7: the source-complete fixed-order Renewal ESM continuum on the V2
branch, with native chiral shell control, anomaly cancellation, common BV
restoration, arbitrary-fixed-loop forests, UV Cauchy convergence,
scheme/reference covariance and exact order projectivity.
\item V8: microscopic universality for an explicit range-two deformation and
the general admissible-gapped-homotopy universality criterion.
\item The positive IR/reference scale and protected local determinant/Pfaffian
component.
\end{enumerate}

\subsection*{Conditional}
\begin{enumerate}[leftmargin=2.2em]
\item The regulator-class universality theorem is proved for the
\emph{canonical covariantized finite-range five-Clifford single-cone class}.
An arbitrary native kernel with additional multiband, curvature-dependent or
noncanonical source couplings requires an extension of the lifting theorem.
\item The common background/source/line chart may need to be shrunk to the
finite intersection required by a chosen regulator chain.  The theorem is local
on that protected component and does not cross a rank or determinant-line zero.
\item All continuum statements remain coefficientwise at fixed finite
$(L,\Delta,r,N)$ and positive reference scale.  No uniform $L\to\infty$ or
massless IR claim is made.
\end{enumerate}

\subsection*{Open beyond V9}
\begin{enumerate}[leftmargin=2.2em]
\item Extend the connectedness theorem from the five-Clifford class to a
general finite-rank multiband native class.  Determine whether equality of the
full $K^0(\mathbb T^4)$ class together with a fixed physical no-doubler germ is
sufficient, or identify unstable finite-rank invariants.
\item Extend the native lifting theorem to uniformly exponentially local
non-ultralocal kernels and to noncanonical curvature/source couplings while
preserving V8's relative physical-shell gain.
\item The independent stronger problems remain: removal of the positive IR
scale, global determinant/Pfaffian zero-crossing and holonomy, analytic
finite-coupling RG trajectories, and nonperturbative completion.
\end{enumerate}

\subsection*{Exact next load-bearing theorem}
\begin{center}
\fbox{\parbox{0.92\textwidth}{
\textbf{General multiband native phase lifting.}
Fix the Standard-Model representation, the physical single-cone germ and a
full flattened projector class in $K^0(\mathbb T^4)$.  Determine the unstable
finite-rank homotopy classes of the relevant Grassmannian-valued projectors and
prove that equality of all necessary invariants permits a finite chain of
native V8-admissible homotopies preserving the no-doubler margin, source jets
and determinant/Pfaffian pivots; otherwise exhibit the next computable
obstruction.  This is the precise gate between the V9 Clifford-class
universality theorem and universality for a genuinely general multiband native
fermion regulator.}}
\end{center}

\section*{Append-only continuation marker: V9}
\addcontentsline{toc}{section}{Append-only continuation marker: V9}
\textbf{End of V9.}  The complete V1--V8 body above is retained verbatim.  V9
shows that the strong four-dimensional overlap number is not a complete
invariant for general gapped Hermitian symbols: the flattened bundle carries a
full $K^0(\mathbb T^4)$ class with six additional weak Chern numbers at fixed
rank.  It then proves a positive classification theorem for the physically
relevant canonical single-cone five-Clifford Wilson sector: weak invariants
vanish, no-doubler maps with the same physical germ are smoothly connected,
the path admits a finite-range jet-preserving approximation and a native
covariant lift, and V8 therefore gives regulator-class universality throughout
that component.  Future versions must append after this marker.


\clearpage
\section{V10 continuation: the overlap big cell and general balanced multiband phase lifting}
\label{sec:v10-continuation}

V9 deliberately stopped before identifying unstable finite-rank
Grassmannian invariants for a general multiband projector.  That was the
correct firewall for an arbitrary gapped Hermitian symbol, but it is still too
large a target for the actual overlap problem.  The overlap/no-doubler
condition supplies additional geometry which changes the classification
question completely.

The decisive observation is elementary and exact.  On the index-zero overlap
sector the positive spectral projector of the Hermitian Wilson kernel and the
fixed negative-chirality projector have the same rank.  Away from an overlap
zero, these two planes must be transverse.  The set of transverse planes is
one Grassmannian big cell, canonically identified with a vector space of graph
maps and therefore contractible.  Consequently, once the physical cone germ
is fixed near the unique zero, neither the stable $K^0$ class nor any unstable
finite-rank Grassmannian class remains free on the punctured Brillouin torus.
The apparent V9 unstable-classification gate disappears in the perturbative
balanced single-cone sector.

This continuation proves that statement first at the flat projector level,
then lifts it to gapped Hermitian symbols, to finite-range/uniformly
exponentially local native kernels with all required source jets, and finally
to the V8 regulator-universality theorem.  The V9 $K$-theory obstruction is not
retracted: it remains a genuine obstruction in the larger space of arbitrary
gapped Hermitian symbols.  V10 proves that the obstructed examples cannot lie
in the same balanced overlap-compatible single-cone component with the same
physical germ.

Throughout, let
\begin{equation}
 \mathcal H=\mathcal H_+\oplus\mathcal H_-,
 \qquad
 \gamma_5|_{\mathcal H_+}=+I,
 \qquad
 \gamma_5|_{\mathcal H_-}=-I,
 \label{eq:v10-chiral-splitting}
\end{equation}
with
\begin{equation}
 \dim_{\mathbb C}\mathcal H_+
 =\dim_{\mathbb C}\mathcal H_-=:m,
 \qquad
 Q:=P_-:=\frac12(I-\gamma_5).
 \label{eq:v10-Q}
\end{equation}
The internal Standard-Model representation and any spectator flavour factors
are included in the finite-dimensional fibre $\mathcal H$.  The equality of
the two ranks is the perturbative index-zero sector already used to define the
chiral block in the inherited construction.

\subsection{The exact overlap-zero discriminant}

Let $H(x)$ be a smooth periodic Hermitian Wilson symbol with a nonzero Wilson
gap and put
\begin{equation}
 \varepsilon(x):=\operatorname{sgn}H(x),
 \qquad
 P(x):=\frac12(I+\varepsilon(x)).
 \label{eq:v10-P}
\end{equation}
Assume throughout this subsection that
$\operatorname{rank}P(x)=m$, as is automatic on the connected index-zero
component.  The corresponding overlap symbol is
\begin{equation}
 D_{\rm ov}(x)=a^{-1}\bigl(I+\gamma_5\varepsilon(x)\bigr).
 \label{eq:v10-Dov}
\end{equation}

\begin{lemma}[Exact projector form of the overlap operator]
\label{lem:v10-D-projector}
For every $x$,
\begin{equation}
 \boxed{
 \gamma_5D_{\rm ov}(x)
 =\frac{2}{a}\bigl(P(x)-Q\bigr).}
 \label{eq:v10-D-projector}
\end{equation}
Hence $D_{\rm ov}(x)$ is invertible if and only if $P(x)-Q$ is invertible.
\end{lemma}
\begin{proof}
Use $\varepsilon=2P-I$ and
$Q=(I-\gamma_5)/2$:
\[
 \gamma_5D_{\rm ov}
 =a^{-1}(\gamma_5+\varepsilon)
 =a^{-1}(\gamma_5+2P-I)
 =2a^{-1}(P-Q).
\]
Multiplication by the unitary involution $\gamma_5$ does not change
invertibility.
\end{proof}

\begin{lemma}[Difference of balanced orthogonal projectors]
\label{lem:v10-balanced-projectors}
Let $P,Q$ be orthogonal projectors of rank $m$ on a $2m$-dimensional Hermitian
space.  Then
\begin{equation}
 \boxed{
 P-Q\text{ is invertible}
 \quad\Longleftrightarrow\quad
 \operatorname{Ran}P\cap\operatorname{Ran}Q=\{0\}.}
 \label{eq:v10-projection-transverse}
\end{equation}
When these conditions hold,
$\operatorname{Ran}P\oplus\operatorname{Ran}Q=\mathcal H$.
\end{lemma}
\begin{proof}
If $0\ne v\in\operatorname{Ran}P\cap\operatorname{Ran}Q$, then
$(P-Q)v=0$, so $P-Q$ is singular.

Conversely assume the intersection is zero.  The two ranges have dimension
$m$, hence their direct sum has dimension $2m$ and equals $\mathcal H$.  If
$(P-Q)v=0$, put $w:=Pv=Qv$.  Then
$w\in\operatorname{Ran}P\cap\operatorname{Ran}Q$, so $w=0$.  Therefore
$v\in\ker P\cap\ker Q=(\operatorname{Ran}P+\operatorname{Ran}Q)^\perp=0$.
Thus $P-Q$ is injective and hence invertible in finite dimension.
\end{proof}

\begin{theorem}[Exact overlap-zero discriminant]
\label{thm:v10-zero-discriminant}
On the balanced index-zero sector,
\begin{equation}
 \boxed{
 \ker D_{\rm ov}(x)\ne0
 \quad\Longleftrightarrow\quad
 \operatorname{Ran}P(x)\cap\mathcal H_-\ne\{0\}.}
 \label{eq:v10-zero-discriminant}
\end{equation}
Equivalently, the no-overlap-zero locus in
$\operatorname{Gr}_m(\mathcal H)$ is
\begin{equation}
 \mathcal U_Q
 :=\left\{W\in\operatorname{Gr}_m(\mathcal H):W\cap\mathcal H_-=0\right\}.
 \label{eq:v10-big-cell-def}
\end{equation}
\end{theorem}
\begin{proof}
Combine Lemmas~\ref{lem:v10-D-projector} and
\ref{lem:v10-balanced-projectors}, using
$\operatorname{Ran}Q=\mathcal H_-$.
\end{proof}

This theorem is the upstream point missed by a classification of arbitrary
gapped projectors.  The relevant target away from the physical overlap zero is
not the whole Grassmannian but the open complement of the Schubert
discriminant in \eqref{eq:v10-big-cell-def}.

\subsection{The no-zero locus is one contractible big cell}

Let $\pi_\pm$ be the projections associated with
\eqref{eq:v10-chiral-splitting}.

\begin{theorem}[Graph parametrization of the overlap big cell]
\label{thm:v10-big-cell}
The map
\begin{equation}
 \operatorname{Hom}(\mathcal H_+,\mathcal H_-)
 \longrightarrow \mathcal U_Q,
 \qquad
 A\longmapsto \operatorname{Graph}(A)
 :=\{u+Au:u\in\mathcal H_+\},
 \label{eq:v10-graph-map}
\end{equation}
is a real-analytic diffeomorphism.  The corresponding orthogonal projector is
\begin{equation}
 \boxed{
 P_A=
 \begin{pmatrix}I\\ A\end{pmatrix}
 (I+A^*A)^{-1}
 \begin{pmatrix}I&A^*\end{pmatrix}.}
 \label{eq:v10-graph-projector}
\end{equation}
Consequently
\begin{equation}
 \boxed{\mathcal U_Q\simeq\mathbb C^{m^2}}
 \label{eq:v10-big-cell-contractible}
\end{equation}
and in particular $\mathcal U_Q$ is contractible.
\end{theorem}
\begin{proof}
If $W\cap\mathcal H_-=0$ and $\dim W=m$, then
$\pi_+|_W:W\to\mathcal H_+$ is injective between spaces of the same dimension,
hence an isomorphism.  Define
$A=\pi_-\circ(\pi_+|_W)^{-1}$.  Then $W=\operatorname{Graph}(A)$, and uniqueness
is immediate.  Conversely every graph is transverse to $\mathcal H_-$.

Writing the graph inclusion as
$J_Au=(u,Au)$ gives $J_A^*J_A=I+A^*A$.  The orthogonal projector onto the range
of an injective map $J_A$ is
$J_A(J_A^*J_A)^{-1}J_A^*$, which is
\eqref{eq:v10-graph-projector}.  Both directions of the graph correspondence
are analytic on the open transversality locus.  Finally
$\operatorname{Hom}(\mathcal H_+,\mathcal H_-)$ is a complex vector space, so
it is contractible by $A\mapsto(1-t)A$.
\end{proof}

\begin{corollary}[Equivariant big cell]
\label{cor:v10-equivariant-big-cell}
Let a compact group $G$ act unitarily on $\mathcal H$, preserve the chiral
splitting, and suppose $Q$ is $G$-invariant.  Then the $G$-invariant planes in
$\mathcal U_Q$ are identified with
\begin{equation}
 \operatorname{Hom}_G(\mathcal H_+,\mathcal H_-),
 \label{eq:v10-equivariant-Hom}
\end{equation}
a complex vector space.  Hence the equivariant no-zero big cell is also
contractible.
\end{corollary}
\begin{proof}
The graph of $A$ is $G$-invariant exactly when
$A\rho_+(g)=\rho_-(g)A$ for every $g\in G$.  The intertwiners form a linear
subspace of the full Hom space, and the graph homotopy stays inside it.
\end{proof}

For the Standard-Model Wilson block, $G$ in this corollary is applied
representation block by representation block to the internal compact gauge
group.  The spin/coframe covariance is restored at the native-lifting stage
below; the flat graph interpolation itself is performed in one fixed reference
spin frame.

\subsection{The physical point and the local multiband Dirac germ}

The physical point does not lie in the big cell: at $x=0$ the V2 convention has
\begin{equation}
 P(0)=Q,
 \qquad
 \varepsilon(0)=2Q-I=-\gamma_5,
 \qquad
 D_{\rm ov}(0)=0.
 \label{eq:v10-physical-point}
\end{equation}
Near $Q$, use the opposite Grassmannian chart.  Every plane sufficiently close
to $\mathcal H_-$ is the graph of a unique map
\begin{equation}
 B(x):\mathcal H_-\longrightarrow\mathcal H_+,
 \qquad
 \operatorname{Ran}P(x)=\{B(x)v+v:v\in\mathcal H_-\}.
 \label{eq:v10-local-B}
\end{equation}
In this chart
\begin{equation}
 \operatorname{Ran}P(x)\cap\mathcal H_-
 \cong\ker B(x).
 \label{eq:v10-B-kernel}
\end{equation}

\begin{definition}[Balanced elliptic physical overlap germ]
\label{def:v10-physical-germ}
A balanced projector family has the fixed physical germ $\mathcal B$ if, in
\eqref{eq:v10-local-B},
\begin{equation}
 B(x)=\mathcal B(x)+R(x),
 \qquad
 \mathcal B(x):=\sum_{\mu=1}^4x_\mu\mathcal B_\mu,
 \qquad
 R(x)=O(|x|^2),
 \label{eq:v10-germ}
\end{equation}
and there is $c_{\mathcal B}>0$ such that
\begin{equation}
 \boxed{
 s_{\min}(\mathcal B(x))\ge c_{\mathcal B}|x|
 \quad(x\in\mathbb R^4).}
 \label{eq:v10-elliptic-germ}
\end{equation}
The matrices $\mathcal B_\mu$ include the fixed coframe/Dirac speed and all
spectator identity factors.  Equality of physical germs means equality of this
first jet, not merely equality of its determinant or topological degree.
\end{definition}

The ordinary Weyl/Dirac principal symbol has exactly
\eqref{eq:v10-elliptic-germ}; after division by the engineering scale its
smallest singular value is proportional to $|p|$.

\begin{lemma}[Quantitative multiband germ straightening]
\label{lem:v10-germ-straightening}
Let $P_0,P_1$ be two smooth balanced single-cone projector maps with the same
physical germ $\mathcal B$.  There is a ball $B_\rho$ and homotopies supported
in $B_\rho$ such that:
\begin{enumerate}[label=\textup{(G\arabic*)},leftmargin=3em]
\item $P_i(0)=Q$ and the first physical jet is fixed throughout;
\item no new overlap zero is created in $B_\rho\setminus\{0\}$;
\item the maps are unchanged outside $B_\rho$;
\item after the homotopy the two projectors agree on a neighborhood of a
smaller closed ball $B\Subset B_\rho$.
\end{enumerate}
\end{lemma}
\begin{proof}
In the local chart \eqref{eq:v10-local-B}, write
$B_i(x)=\mathcal B(x)+R_i(x)$ with
$\|R_i(x)\|\le C|x|^2$ after shrinking $\rho$.  Choose a cutoff $\chi$ equal
to one on a neighborhood of a smaller ball and zero near
$\partial B_\rho$, and put
\begin{equation}
 B_{i,u}(x)=\mathcal B(x)+(1-u\chi(x))R_i(x).
 \label{eq:v10-B-straightening}
\end{equation}
The value and first derivative at zero do not change.  Moreover
\[
 s_{\min}(B_{i,u}(x))
 \ge c_{\mathcal B}|x|-C|x|^2.
\]
Choose $\rho<c_{\mathcal B}/(2C)$.  Then the right-hand side is positive for
$0<|x|<\rho$, so \eqref{eq:v10-B-kernel} and
Theorem~\ref{thm:v10-zero-discriminant} exclude every additional overlap zero.
Where $\chi=1$, both final maps are the graph of the same canonical linear germ
$\mathcal B(x)$; where $\chi=0$, each original map is unchanged.
\end{proof}

\subsection{General smooth multiband no-doubler connectedness}

We can now replace the anticipated unstable Grassmannian obstruction theory by
a direct relative homotopy.

\begin{theorem}[Smooth balanced multiband single-cone connectedness]
\label{thm:v10-smooth-multiband}
Let
\begin{equation}
 P_0,P_1:\mathbb T^4\longrightarrow
 \operatorname{Gr}_m(\mathcal H)
 \label{eq:v10-two-projectors}
\end{equation}
be smooth projector maps satisfying:
\begin{enumerate}[label=\textup{(M\arabic*)},leftmargin=3em]
\item $P_i(0)=Q$ and both have the same balanced elliptic physical germ
$\mathcal B$;
\item for every $x\ne0$,
$\operatorname{Ran}P_i(x)\cap\mathcal H_-=0$;
\item the same internal representation grading is preserved.
\end{enumerate}
Then there is a smooth homotopy $P_t$ with
\begin{equation}
 \boxed{
 P_t(0)=Q,
 \qquad
 dP_t(0)=dP_0(0),
 \qquad
 \operatorname{Ran}P_t(x)\cap\mathcal H_-=0
 \quad(x\ne0),}
 \label{eq:v10-projector-homotopy}
\end{equation}
for all $t\in[0,1]$.  The homotopy may be chosen within the internal-gauge
equivariant projector subspace.
\end{theorem}
\begin{proof}
Apply Lemma~\ref{lem:v10-germ-straightening} to both endpoints.  After this
preliminary homotopy there is a closed ball $B$ about $0$ and a collar of its
boundary on which the two maps agree.  Keep that common local map fixed.

Put
$X:=\mathbb T^4\setminus\operatorname{int}B$.  By (M2) and
Theorem~\ref{thm:v10-big-cell}, each restriction has a unique graph coordinate
\begin{equation}
 A_i:X\longrightarrow
 \operatorname{Hom}(\mathcal H_+,\mathcal H_-).
 \label{eq:v10-Ai}
\end{equation}
The two graph maps agree on a collar of $\partial X$.  Define
\begin{equation}
 A_t=(1-t)A_0+tA_1.
 \label{eq:v10-A-linear}
\end{equation}
Every $A_t$ defines a plane in the big cell, so
$P_t:=P_{A_t}$ has no overlap zero on $X$.  Since the interpolation is constant
on the boundary collar, it glues smoothly to the fixed common projector on
$B$.  The physical value and first jet are therefore unchanged.  If the
endpoints preserve the internal representation grading, their graph maps lie
in the intertwiner space of Corollary~\ref{cor:v10-equivariant-big-cell}, which
is linear; hence the same interpolation is equivariant.
\end{proof}

\begin{corollary}[No independent stable or unstable bundle invariant remains]
\label{cor:v10-no-extra-topology}
Within the balanced single-cone sector of
Theorem~\ref{thm:v10-smooth-multiband}, the complete homotopy class of the
flattened projector is fixed by the physical germ.  In particular,
\begin{equation}
 [E_{P_0}]=[E_{P_1}]\in K^0(\mathbb T^4),
 \label{eq:v10-K-forced}
\end{equation}
and any finite-rank unstable Grassmannian invariant takes the same value at the
two endpoints.
\end{corollary}
\begin{proof}
Theorem~\ref{thm:v10-smooth-multiband} gives an actual homotopy in the same
finite Grassmannian, which is stronger than stable equivalence.  Pullback
bundles and all homotopy invariants are therefore equal.
\end{proof}

\begin{corollary}[Location of the V9 weak-index counterexamples]
\label{cor:v10-v9-counterexamples}
The two finite-range symbols used in the V9 weak-index obstruction cannot both
satisfy the same balanced physical germ and the single-cone/no-doubler
condition relative to one fixed $Q$.  Any attempted interpolation into that
sector must either change the physical germ or meet the overlap-zero
discriminant away from $x=0$.
\end{corollary}
\begin{proof}
If both symbols satisfied the hypotheses of
Theorem~\ref{thm:v10-smooth-multiband}, Corollary
\ref{cor:v10-no-extra-topology} would force their $K^0$ classes to agree,
contradicting the weak Chern-number difference proved in V9.
\end{proof}

\begin{firewall}[V9 remains correct]
V9 classified the larger space of arbitrary gapped Hermitian symbols before
imposing the overlap no-doubler incidence relative to $Q$.  In that larger
space the full $K^0(\mathbb T^4)$ class is genuinely necessary.  V10 proves
that the actual balanced single-cone overlap component is a much smaller
relative mapping problem: the punctured torus lands in the contractible big
cell $\mathcal U_Q$.  The two statements concern different spaces and are
therefore compatible.
\end{firewall}

\subsection{From projector homotopy to a gapped Hermitian-kernel homotopy}

The overlap depends on a Wilson kernel only through its sign.  This permits a
clean separation between spectral radial data and projector topology.

\begin{lemma}[Radial flattening preserves the overlap operator exactly]
\label{lem:v10-radial-flatten}
Let $H$ be a Hermitian symbol with
$s_{\min}(H)\ge g>0$ and put
$\varepsilon=\operatorname{sgn}H$.  For
$0\le s\le1$ define
\begin{equation}
 H_s:=(1-s)H+s\varepsilon.
 \label{eq:v10-radial-path}
\end{equation}
Then
\begin{equation}
 s_{\min}(H_s)\ge\min\{g,1\},
 \qquad
 \operatorname{sgn}H_s=\varepsilon.
 \label{eq:v10-radial-gap}
\end{equation}
Consequently the overlap projector, overlap Dirac operator, physical zeros and
all determinant-line data depending only on the chiral sign are unchanged
along the radial path.
\end{lemma}
\begin{proof}
The operators $H$ and $\varepsilon$ commute and have the same spectral
projections.  On an eigenvector of $H$ with eigenvalue $\lambda$, the eigenvalue
of $H_s$ is
$(1-s)\lambda+s\operatorname{sgn}\lambda$, which has the same sign as
$\lambda$ and magnitude at least $\min\{|\lambda|,1\}$.  Hence no eigenvalue
crosses zero and the sign is constant.
\end{proof}

\begin{theorem}[Smooth gapped multiband Wilson homotopy]
\label{thm:v10-smooth-H-homotopy}
Let $H_0,H_1$ be two smooth gapped balanced Hermitian symbols whose flattened
projectors satisfy Theorem~\ref{thm:v10-smooth-multiband}.  Then they are joined
by a smooth gapped Hermitian homotopy which preserves the physical overlap germ
and the single-cone/no-doubler condition throughout.
\end{theorem}
\begin{proof}
Use Lemma~\ref{lem:v10-radial-flatten} to join $H_0$ to
$\varepsilon_0=2P_0-I$.  Use Theorem~\ref{thm:v10-smooth-multiband} and set
\begin{equation}
 \varepsilon_t:=2P_t-I.
 \label{eq:v10-epsilon-path}
\end{equation}
This middle path has unit Wilson gap.  Its overlap zeros are exactly those of
$P_t-Q$ by Theorem~\ref{thm:v10-zero-discriminant}, so there is only the fixed
physical zero.  Finally reverse the radial path for $H_1$.  The physical first
jet is unchanged on the middle path and the radial pieces have exactly constant
sign, hence exactly constant overlap operator.
\end{proof}

For finite-range $H_i$, their signs are generally not finite range, but the V1
gapped functional calculus gives uniform exponential locality and the same
fixed source-jet bounds.  Thus the two short radial pieces already lie in the
``finite range or uniformly exponentially local'' class permitted by the V8
admissibility definition.

\subsection{Finite-range approximation of the middle multiband path}

The middle involution path \eqref{eq:v10-epsilon-path} is smooth in the
Brillouin variable but need not itself be finite range.  We now approximate it
without losing its exact physical first jet or its no-doubler margin.

\begin{lemma}[First-jet rigidity of the matrix sign at an involution]
\label{lem:v10-sign-first-jet}
Let $E=E^*=E^{-1}$ and let $X=X^*$ satisfy
\begin{equation}
 EX+XE=0.
 \label{eq:v10-tangent-involution}
\end{equation}
Then the Fr\'echet derivative of the matrix sign at $E$ obeys
\begin{equation}
 \boxed{D(\operatorname{sgn})_E[X]=X.}
 \label{eq:v10-sign-first-jet}
\end{equation}
\end{lemma}
\begin{proof}
Let $P_\pm=(I\pm E)/2$.  The divided-difference formula for the derivative of a
spectral function gives zero on the $P_+XP_+$ and $P_-XP_-$ blocks and
coefficient
$(1-(-1))/(1-(-1))=1$ on both off-diagonal blocks.  Condition
\eqref{eq:v10-tangent-involution} is equivalent to
$P_+XP_+=P_-XP_-=0$.  Hence the derivative is precisely
$P_+XP_-+P_-XP_+=X$.
\end{proof}

\begin{lemma}[Uniform no-doubler margins for the smooth path]
\label{lem:v10-no-doubler-margins}
For the path $P_t$ of Theorem~\ref{thm:v10-smooth-multiband}, one can choose
$0<\rho<\pi$ and constants $c_0,\delta_0>0$ such that
\begin{align}
 s_{\min}B_t(x)&\ge c_0|x|,
 &&0<|x|\le\rho,
 \label{eq:v10-local-margin}\\
 s_{\min}(P_t(x)-Q)&\ge\delta_0,
 &&|x|\ge\rho/2,
 \label{eq:v10-global-margin}
\end{align}
uniformly in $t$.
\end{lemma}
\begin{proof}
The homotopy is stationary in a neighborhood of the physical point after the
germ straightening, so \eqref{eq:v10-local-margin} follows from
\eqref{eq:v10-elliptic-germ} after shrinking $\rho$.  On the compact set
$[0,1]\times\{x:|x|\ge\rho/2\}$, the matrices $P_t(x)-Q$ are invertible by
Theorem~\ref{thm:v10-smooth-multiband}; continuity of the least singular value
gives the positive minimum $\delta_0$.
\end{proof}

\begin{theorem}[Finite-chain approximation for a general balanced multiband path]
\label{thm:v10-finite-chain}
Assume the endpoint symbols are finite range or uniformly exponentially local.
The smooth middle path of Theorem~\ref{thm:v10-smooth-H-homotopy} can be
replaced by a finite chain of Hermitian trigonometric-polynomial segments
\begin{equation}
 K_{j,u}(x),
 \qquad
 1\le j\le J,
 \quad 0\le u\le1,
 \label{eq:v10-finite-chain}
\end{equation}
with the following properties:
\begin{enumerate}[label=\textup{(F\arabic*)},leftmargin=3em]
\item every segment has a common positive Wilson gap;
\item after flattening, every segment has exactly the same physical overlap
value and first jet at $x=0$;
\item no segment creates an overlap zero away from $x=0$;
\item all symbols preserve the fixed internal representation grading;
\item the first and last middle-segment endpoints can be joined to the exact
$\varepsilon_0,\varepsilon_1$ by uniformly exponentially local short segments
with the same properties.
\end{enumerate}
\end{theorem}
\begin{proof}
Regard the Hermitian matrices as a finite-dimensional real vector space.  Apply
the V9 Hermite--Fourier approximation lemma componentwise to the smooth
involution path $\varepsilon_t$, using $C^2$ rather than merely $C^1$
approximation and imposing at every node the exact common value and first
$x$-derivative at $0$.  If an internal compact gauge representation is
present, perform the approximation inside its finite-dimensional commutant;
equivalently average the approximant over the compact group.  This preserves
the exact jet constraints.

Choose a partition $0=t_0<\cdots<t_J=1$ so fine, and the node approximants so
close in $C^2$, that every linear segment $K_{j,u}$ remains within a fixed small
operator-norm neighborhood of the corresponding short portion of
$\varepsilon_t$.  Since the latter has unit gap, the singular-value perturbation
inequality gives a common positive Wilson gap.

Let
$\widetilde\varepsilon_{j,u}:=\operatorname{sgn}K_{j,u}$ and
$\widetilde P_{j,u}=(I+\widetilde\varepsilon_{j,u})/2$.  At $x=0$ the common
value is an involution and the prescribed first derivative is tangent to the
involution manifold, because it is the derivative of the exact path
$\varepsilon_t^2=I$.  Lemma~\ref{lem:v10-sign-first-jet} therefore shows that
flattening preserves the exact first physical jet.

On $|x|\ge\rho/2$, the Riesz functional calculus is locally Lipschitz in the
gapped operator norm.  Make the approximation small enough that
$\|\widetilde P_{j,u}-P_t\|<\delta_0/2$ for a nearby path parameter $t$.
Equation \eqref{eq:v10-global-margin} then keeps
$\widetilde P_{j,u}-Q$ invertible.  On $|x|<\rho$, the exact first jet and the
uniform $C^2$ bound give, in the local graph chart,
\[
 \widetilde B_{j,u}(x)=\mathcal B(x)+O(|x|^2)
\]
uniformly.  Shrinking $\rho$ once more if necessary gives
$s_{\min}\widetilde B_{j,u}(x)\ge(c_0/2)|x|$ for $x\ne0$.  Thus no extra zero
appears near the physical point either.

The exact endpoints $\varepsilon_i$ are uniformly exponentially local by the
V1 functional calculus.  Connect each to a sufficiently close first/last
trigonometric node by a short linear segment.  The same gap and no-doubler
openness estimates apply, and a convex combination of an exponentially local
kernel with a finite-range one is uniformly exponentially local.  This proves
all items.
\end{proof}

\subsection{Native multiband covariant lifting and source jets}

The remaining issue is not topology but covariance.  A general multiband
Fourier coefficient is not necessarily one of the five Clifford matrices used
in V9.  On the protected local coframe/spin chart, however, the whole fibre
endomorphism algebra can be lifted covariantly.

Fix a smooth local spin-frame trivialization
$S_x(e)$ on the protected coframe chart, normalized by $S_x(e_0)=I$ at the
reference coframe and chosen equivariantly under the local spin action.  Such a
trivialization exists on the contractible chart already used for the V1 Kato
and V2 source constructions.  Let
$\operatorname{End}_{G_{\rm int}}(\mathcal H)$ denote the finite-dimensional
internal-gauge intertwiner algebra appropriate to the fixed representation
block.

\begin{definition}[Chart-covariant multiband lift]
\label{def:v10-multiband-lift}
Let
\begin{equation}
 K(x)=\sum_{n\in F}A_ne^{in\cdot x},
 \qquad
 A_{-n}=A_n^*,
 \qquad
 A_n\in\operatorname{End}_{G_{\rm int}}(\mathcal H),
 \label{eq:v10-matrix-Fourier}
\end{equation}
be a finite Hermitian Fourier symbol.  For each nonzero displacement choose a
fixed lattice path $\wp(n)$ and let $T_{\wp(n)}^{\mathcal U}$ be the combined
internal-gauge/spin parallel shift.  Set
\begin{equation}
 A_n(e;x):=S_x(e)A_nS_x(e)^{-1}.
 \label{eq:v10-frame-conjugate}
\end{equation}
The chart-covariant lift is the Hermitian sum obtained from the oriented terms
$A_n(e;x)T_{\wp(n)}^{\mathcal U}$ together with their exact adjoints and the
onsifted onsite term.
\end{definition}

\begin{lemma}[Exact covariance and finite source calculus]
\label{lem:v10-lift-properties}
The lift of Definition~\ref{def:v10-multiband-lift} has the following
properties.
\begin{enumerate}[label=\textup{(C\arabic*)},leftmargin=3em]
\item it is Hermitian and finite range;
\item it transforms by exact local internal-gauge and spin similarity;
\item at the identity-link/reference-coframe background its Fourier symbol is
exactly $K(x)$;
\item every prescribed finite gauge, spin-connection and coframe derivative is
a finite-range kernel with a cutoff-independent coefficientwise bound depending
only on the finite path set and derivative order.
\end{enumerate}
The same construction applies to an exponentially summable Fourier family and
preserves a uniform exponential localization rate after a fixed small loss of
that rate.
\end{lemma}
\begin{proof}
Internal gauge covariance follows because the $A_n$ lie in the intertwiner
algebra and the parallel shift transforms at its two endpoints.  Spin
covariance follows from the equivariant frame rule
$S_x(g\cdot e)=g_xS_x(e)$ and the corresponding endpoint transformation of the
spin parallel transport.  Adding each oriented term together with its exact
adjoint gives Hermiticity.  At the reference frame $S_x=I$ and all parallel
transports are identity before the displacement shift, so the prescribed
Fourier multiplier is recovered.

A fixed source derivative differentiates only finitely many factors along each
fixed path and finitely many factors in the local frame conjugation.  The
Leibniz sum is therefore finite.  For an exponentially summable family, the
number of differentiated factors grows at most polynomially with path length;
this is absorbed by an arbitrarily small reduction of the exponential weight.
\end{proof}

The V8 relative theorem requires more than covariance: the regulator derivative
must not alter the physical principal symbol.  This is enforced before the
lift.

\begin{definition}[Principal-jet-preserving native lift]
\label{def:v10-principal-lift}
Fix once and for all a Hermitian trigonometric polynomial
$G_{\rm phys}(x)$ having the common value and first momentum jet of the
physical involution at $x=0$, and lift $G_{\rm phys}$ by the same native V2
Dirac/coframe/parallel-transport prescription for every regulator parameter.
For a middle-chain symbol $K_t$, write
\begin{equation}
 K_t=G_{\rm phys}+R_t,
 \qquad
 R_t(0)=0,
 \qquad
 \partial_{x_\mu}R_t(0)=0,
 \label{eq:v10-R-moments}
\end{equation}
and apply Definition~\ref{def:v10-multiband-lift} only to the remainder $R_t$.
\end{definition}

\begin{lemma}[The regulator velocity is irrelevant with all fixed source jets]
\label{lem:v10-native-irrelevant}
For the principal-jet-preserving lift, the derivative with respect to the
regulator parameter has zero constant and linear physical momentum jet.  At
every fixed source order required by $\mathfrak q$ it therefore belongs to the
same one-irrelevant-order source class used in V8: on a physical hard shell,
after the engineering/source stocks are divided out, the marked insertion has
an extra factor $a\Lambda_j$ relative to the ordinary hard kinetic insertion.
\end{lemma}
\begin{proof}
The common $G_{\rm phys}$ has no regulator derivative.  For $R_t$, the two
moment conditions in \eqref{eq:v10-R-moments} hold identically in $t$.
Frame conjugation is common to all Fourier coefficients, so differentiating a
coframe source preserves the vanishing sums corresponding to the zero value and
first momentum moment.  Gauge/spin-connection derivatives of path transports
are the same edge-insertion calculus used in V2.  A covariant finite-difference
operator whose flat symbol vanishes to second order has, after any fixed number
of such source insertions, at least one explicit lattice spacing beyond the
principal first-order Dirac vertex.  Equivalently, the V2 Taylor/source
counting gives the relative $a\Lambda_j$ factor on a shell
$|p|\asymp\Lambda_j$.  Since the source order is fixed, the finite Leibniz
multiplicity is absorbed in the constant.
\end{proof}

\subsection{The native balanced multiband phase-lifting theorem}

We can now verify the V8 gates without appealing to stable or unstable
Grassmannian classification.

\begin{definition}[Balanced single-cone multiband native class]
\label{def:v10-multiband-class}
Let $\mathfrak W_{\rm mb}^{\rm sc,0}$ consist of native Hermitian Wilson kernels
such that:
\begin{enumerate}[label=\textup{(B\arabic*)},leftmargin=3em]
\item they are finite range or uniformly exponentially local, exactly
internal-gauge/spin covariant, and possess all fixed source jets required by the
coefficientwise programme;
\item their flat reference symbols have a positive Wilson gap and balanced
positive spectral rank $m=\dim\mathcal H_-$;
\item the overlap symbol has exactly one physical zero at $x=0$ and the fixed
balanced elliptic physical germ $\mathcal B$;
\item they preserve the same Standard-Model representation grading and lie on
the same protected positive-IR determinant/Pfaffian source component.
\end{enumerate}
\end{definition}

\begin{theorem}[General balanced multiband native phase lifting]
\label{thm:v10-native-phase-lift}
Let $A,B\in\mathfrak W_{\rm mb}^{\rm sc,0}$.  For every prescribed finite
coefficient packet
$\mathfrak q=(L,\Delta,r,N,\mathcal P_{\rm ext})$ there exists a finite chain of
native regulator homotopies joining $A$ to $B$ such that every segment is
V8-admissible after one common finite shrinking of the protected
background/source/line chart.
\end{theorem}
\begin{proof}
At the flat reference background, flatten both endpoint kernels.  Their
balanced positive projectors satisfy
Theorem~\ref{thm:v10-smooth-multiband}, which gives a smooth projector homotopy
with the fixed physical germ and no extra overlap zero.  Theorem
\ref{thm:v10-finite-chain} replaces its middle part by a finite chain of
gapped finite-range or uniformly exponentially local Hermitian symbols whose
flattened projectors retain the same properties.  Lemma
\ref{lem:v10-radial-flatten} supplies the two endpoint radial segments.

Lift every finite middle segment by Definitions
\ref{def:v10-multiband-lift} and \ref{def:v10-principal-lift}.  Lemma
\ref{lem:v10-lift-properties} gives exact covariance and the required fixed
source calculus.  Lemma~\ref{lem:v10-native-irrelevant} gives the V8 relative
physical-shell gain.  The full-zone flat symbols are gapped and no-doubler by
construction; compactness of the finite chain provides a positive minimum
margin away from the physical ball.

There are only finitely many segments and regulator parameters.  Continuity of
the native operator and its fixed source derivatives therefore permits one
common shrinking of the V1--V9 perturbative chart so that the Wilson gap, the
hard chiral pivots above the retained positive IR scale, and all
Determinant/Pfaffian nonzero conditions remain uniform on every segment.  The
normalization-rank gate is the same chronological finite local-coordinate gate
used in V9: the counterterm coordinates give the regulator-independent identity
block, while regulator variation changes only the inhomogeneous loop data.
Thus no rank is lost along the chain.

The endpoint radial pieces are uniformly exponentially local by the inherited
gapped functional calculus and have exactly constant sign, so they trivially
satisfy the physical-principal-symbol part of V8 admissibility.  Every V8 gate
is therefore satisfied.
\end{proof}

\begin{theorem}[General balanced multiband regulator-class universality]
\label{thm:v10-multiband-universality}
Fix finite $\mathfrak q$.  Any two regulators
$A,B\in\mathfrak W_{\rm mb}^{\rm sc,0}$ are matched-universal.  There exists an
invertible finite local continuum map
\begin{equation}
 F_{\infty,\mathfrak q}^{B\leftarrow A}
 \label{eq:v10-multiband-map}
\end{equation}
such that
\begin{equation}
 \boxed{
 \widehat{\mathfrak C}^{B}_{\infty,\lambda}
 =F_{\infty,\mathfrak q}^{B\leftarrow A}
  \widehat{\mathfrak C}^{A}_{\infty,\lambda}.}
 \label{eq:v10-multiband-universality}
\end{equation}
After covariant matching, the complete normalized connected/source system,
operator mixing, stress/contact insertions, ghosts/antifields, coupled
BV/BRST/master identities and determinant/Pfaffian line insertions coincide.
The maps may be chosen strictly projective in perturbative order.
\end{theorem}
\begin{proof}
Apply Theorem~\ref{thm:v10-native-phase-lift} and then the V8
gapped-homotopy universality theorem to every segment of the finite chain.
Compose the resulting finite invertible local matching maps.  V8 proves the
coefficientwise equality for every stored coordinate of the contraction-complete
state.  Each segment uses the chronological V7--V8 construction, so the finite
composition remains projective in $L$.
\end{proof}

\begin{corollary}[The V9 Clifford class is a subcase]
\label{cor:v10-Clifford-subcase}
The canonical five-Clifford class
$\mathfrak W_{\rm Cliff}^{\rm sc}$ of V9 is contained in
$\mathfrak W_{\rm mb}^{\rm sc,0}$.  Hence the V9 regulator-class universality
theorem follows as a special case of
Theorem~\ref{thm:v10-multiband-universality}; the five-Clifford factorization
through $S^4$ is sufficient but not necessary for balanced single-cone
connectedness.
\end{corollary}
\begin{proof}
The V9 class has balanced chiral rank, the same physical Dirac germ, one
physical overlap zero and the required native covariance/source properties.
Thus it satisfies Definition~\ref{def:v10-multiband-class}.
\end{proof}

\subsection{Why the anticipated unstable Grassmannian computation was the wrong next variable}

For completeness, V10 records the logical relation among the three spaces which
have appeared in V8--V10:
\begin{equation}
 \boxed{
 \begin{array}{c}
 \text{all gapped Hermitian symbols}
 \supset
 \text{balanced overlap symbols}
 \supset
 \text{balanced single-cone symbols with fixed germ},\\[0.3em]
 K^0(\mathbb T^4)\text{ can obstruct the first space,}\qquad
 \mathcal U_Q\simeq\mathbb C^{m^2}\text{ controls the punctured target of the last.}
 \end{array}}
 \label{eq:v10-space-hierarchy}
\end{equation}

The finite-rank Grassmannian itself can possess unstable homotopy information
in other dimensions or under other rank constraints.  None of that information
is load-bearing here because the no-doubler condition never lets the punctured
Brillouin torus explore the whole Grassmannian.  It is confined to one affine
chart.  This is exactly the kind of upstream variable change demanded by the
programme when a route begins to classify more structure than the theorem uses.

\begin{proposition}[The physical germ is the clutching datum for the balanced single-cone class]
\label{prop:v10-germ-clutching}
Let $P$ be a balanced single-cone projector.  After choosing a sufficiently
small physical ball $B$, the restriction to
$X=\mathbb T^4\setminus\operatorname{int}B$ is null-homotopic inside the big
cell relative to any fixed boundary graph.  Therefore every global bundle or
projector homotopy invariant is determined entirely by how the physical local
germ is glued to that big-cell exterior.
\end{proposition}
\begin{proof}
The exterior graph coordinate is a map
$A:X\to\operatorname{Hom}(\mathcal H_+,\mathcal H_-)$.  Once its boundary
value is fixed, the affine target permits the straight-line homotopy to any
other extension with the same boundary, in particular to a chosen reference
extension.  Thus there is no independent exterior homotopy datum.  All
remaining information is in the fixed boundary/local germ.
\end{proof}

This proposition also explains why a weak Chern class can appear in the V9
unconstrained example but not here: producing such a class requires leaving
the single affine cell somewhere on the punctured torus, which by
Theorem~\ref{thm:v10-zero-discriminant} is precisely an additional overlap-zero
incidence.

\subsection{Boundary of the theorem: nonzero index and forced zero modes}

The balanced-rank hypothesis is essential.  If
\begin{equation}
 \operatorname{rank}P-\operatorname{rank}Q\ne0,
 \label{eq:v10-index-nonzero}
\end{equation}
then two equal-dimensional complementary planes are no longer available and
$P-Q$ cannot be invertible everywhere.  Zero modes are forced by dimension
before any local estimate is attempted.  Likewise, if a topologically
nontrivial gauge background forces the overlap index to jump, the determinant
line must be treated over a zero-mode stratum rather than by the nonzero local
frame used in V1--V10.

\begin{firewall}[No nonzero-index theorem]
Theorems~\ref{thm:v10-native-phase-lift} and
\ref{thm:v10-multiband-universality} are perturbative index-zero theorems on the
protected nonzero hard-shell component.  They do not connect distinct overlap
index sectors, remove topological zero modes, or trivialize the determinant or
Pfaffian line across an index-changing locus.
\end{firewall}

\subsection{V10 anti-circularity audit}
\label{sec:v10-audit}

\begin{enumerate}[leftmargin=2.4em]
\item \textbf{The unstable-classification problem was not assumed away.}
The exact identity \eqref{eq:v10-D-projector} was used first.  It shows that
the actual no-doubler target is the transverse-plane big cell, which is proved
directly to be affine and contractible.
\item \textbf{V9's $K$-theory obstruction is retained at its true scope.}
It applies to arbitrary gapped Hermitian symbols.  V10 proves that the
additional balanced no-doubler hypotheses force all such global invariants once
the physical germ is fixed.
\item \textbf{The physical point is not incorrectly placed inside the big cell.}
It is treated in the opposite local graph chart around $Q$; only the punctured
complement is parametrized by $A:\mathcal H_+\to\mathcal H_-$.
\item \textbf{No-doubler preservation is quantitative.}
Near the physical point it follows from the elliptic singular-value lower bound
on the common first jet; away from that point it follows from a compact positive
margin for $P-Q$.
\item \textbf{A smooth projector homotopy is not silently declared native.}
The middle path is approximated by a finite trigonometric chain with exact
physical first jet, then lifted with explicit gauge/spin parallel transport and
a local spin-frame conjugation.
\item \textbf{Leading geometric/gauge source vertices are not changed.}
The common principal jet is lifted once and frozen; only a remainder with zero
constant/linear momentum moments carries regulator dependence, so the V8
one-irrelevant-order source estimate remains valid.
\item \textbf{Determinant/Pfaffian pivots are not inferred from flat topology.}
They are preserved by one finite shrinking of the compact protected source
chart, using openness and the positive IR hard-shell prescription.
\item \textbf{The theorem does not cross index sectors.}
Balanced rank is used in Lemma~\ref{lem:v10-balanced-projectors}; when it fails,
forced zero modes require a different line-bundle/zero-mode analysis.
\end{enumerate}

\section{V10 status ledger}
\label{sec:v10-status}

\subsection*{Proved in V10}
\begin{enumerate}[leftmargin=2.2em]
\item On the perturbative index-zero overlap sector,
$\gamma_5D_{\rm ov}=2a^{-1}(P-Q)$ and the overlap operator is invertible exactly
when the positive Wilson plane is transverse to the fixed negative-chirality
plane.
\item The no-overlap-zero locus is one Grassmannian big cell,
analytically isomorphic to
$\operatorname{Hom}(\mathcal H_+,\mathcal H_-)$ and therefore contractible;
the equivariant big cell is the corresponding intertwiner vector space.
\item A fixed balanced elliptic physical germ can be straightened locally
without creating another zero.  After that straightening, any two smooth
balanced multiband single-cone projectors with the same germ are connected by
linear interpolation of their exterior graph coordinates, while the physical
ball is held fixed.
\item Hence no independent stable $K^0$ invariant or unstable finite-rank
Grassmannian invariant remains inside the balanced single-cone component once
the physical germ is fixed.  The V9 weak-index counterexamples necessarily
leave that component.
\item Radial spectral flattening connects an arbitrary gapped endpoint kernel
to its sign without changing the overlap operator at all.  The middle projector
homotopy therefore gives a smooth gapped Hermitian-kernel homotopy.
\item The middle path admits a finite Hermitian trigonometric approximation with
exact physical first jet, a uniform Wilson gap and no extra overlap zero; the
endpoint flattening pieces are uniformly exponentially local.
\item A chart-covariant multiband lift gives exact internal-gauge/spin
covariance and finite source calculus for arbitrary allowed matrix Fourier
coefficients.  Freezing the common principal native jet and lifting only the
zero-value/zero-first-moment remainder preserves the V8 one-irrelevant-order
gain for all required fixed source derivatives.
\item Therefore any two regulators in the balanced single-cone multiband native
class $\mathfrak W_{\rm mb}^{\rm sc,0}$ are joined by a finite chain of
V8-admissible homotopies and have the same matched coefficientwise continuum,
with strict perturbative-order projectivity.
\end{enumerate}

\subsection*{Inherited}
\begin{enumerate}[leftmargin=2.2em]
\item V1--V7: the source-complete fixed-order Renewal ESM continuum,
anomaly-free common BV restoration, arbitrary-fixed-loop forests, ultraviolet
Cauchy theorem, finite-scheme/reference covariance and exact loop-order
projectivity.
\item V8: regulator universality along every verified admissible gapped
homotopy, including determinant/Pfaffian source transgression and the two-scale
relative-forest estimate.
\item V9: the full $K^0(\mathbb T^4)$ obstruction for arbitrary gapped Hermitian
symbols, the Clifford single-cone connectedness theorem, Hermite--Fourier jet
approximation and the first native finite-range lifting construction.
\end{enumerate}

\subsection*{Conditional}
\begin{enumerate}[leftmargin=2.2em]
\item The general multiband theorem is local to the balanced perturbative
index-zero sector and to the fixed Standard-Model representation grading.
\item Native lifting uses one protected local spin/coframe trivialization and
may require a finite shrinking of the common source/line chart.  A transition
between distinct coframe/spin charts is not part of V10.
\item All statements remain coefficientwise at prescribed finite
$(L,\Delta,r,N)$ and positive IR/reference scale, with no uniform
$L\to\infty$ estimate.
\item Global determinant/Pfaffian zero crossings, nonzero-index sectors and a
massless IR limit remain outside the theorem.
\end{enumerate}

\subsection*{Open beyond V10}
\begin{enumerate}[leftmargin=2.2em]
\item Extend the phase-lifting/universality theorem to nonzero overlap-index or
topologically nontrivial gauge sectors, where zero modes are forced and the
relevant object is a determinant/Pfaffian line over a stratified Fredholm
Grassmannian rather than the nonzero big cell.
\item Glue the local spin/coframe regulator lifts across overlapping protected
geometric charts and identify the resulting global transition cocycle, if a
global regulator-class statement is desired.
\item The independent stronger problems remain: removal of the positive IR
scale, global determinant/Pfaffian holonomy and zero crossing, analytic
finite-coupling RG trajectories, perturbative summability and nonperturbative
completion.
\end{enumerate}

\subsection*{Exact next load-bearing theorem}
\begin{center}
\fbox{\parbox{0.92\textwidth}{
\textbf{Nonzero-index / zero-mode phase lifting.}
Replace the balanced big-cell argument by a stratified finite-dimensional
Fredholm analysis for
$\operatorname{rank}P-\operatorname{rank}Q=k\ne0$ (and for gauge backgrounds
where the overlap index changes).  Construct the exact zero-mode bundle and its
determinant/Pfaffian line, prove shell factorization and source-jet estimates in
a local zero-mode chart, and classify which regulator homotopies preserve the
index and line orientation.  The target is a V8-style universality theorem
within a fixed nonzero-index sector, with the perturbative V10 theorem recovered
at $k=0$.}}
\end{center}

\section*{Append-only continuation marker: V10}
\addcontentsline{toc}{section}{Append-only continuation marker: V10}
\textbf{End of V10.}  The complete V1--V9 body above is retained verbatim.  V10
goes upstream from the unstable-Grassmannian problem isolated in V9 and uses
the exact overlap identity
$\gamma_5D_{\rm ov}=2a^{-1}(P-P_-)$ to identify the true no-doubler target.
In the balanced perturbative index-zero sector this target is the contractible
Grassmannian big cell of planes transverse to $\mathcal H_-$.  With the
physical Dirac germ fixed, this removes all independent stable and unstable
projector invariants.  A quantitative local germ straightening, big-cell graph
homotopy, finite-range jet-preserving approximation, multiband native
covariant lift and V8 admissibility check then prove regulator-class
universality for the full balanced single-cone multiband native class.  Future
versions must append after this marker.



\clearpage
\section{V11 continuation: fixed-index zero-mode splitting, reduced shell calculus, and local universality}
\label{sec:v11-continuation}

V10 closed the balanced index-zero phase-lifting problem by identifying the
punctured no-doubler target with one affine Grassmannian big cell.  At nonzero
index that exact argument cannot be repeated because zero modes are forced.
The correct upstream variable is nevertheless simpler than the full stratified
Fredholm Grassmannian suggested in the V10 handoff.  On the minimal fixed-index
stratum, all nonzero-mode data again form an affine graph fibre; the only
noncontractible datum is the forced kernel or cokernel plane.  Thus the
nonzero-index problem separates canonically into
\begin{equation}
 \boxed{
 \text{zero-mode bundle/line}
 \quad\oplus\quad
 \text{invertible reduced chiral complement}.}
 \label{eq:v11-upstream-split}
\end{equation}

The finite zero-mode determinant-line factorization and exact Grassmann
saturation identity are already present in the Grand Tensor construction and
are inherited here.  They are not reproved.  The new work is to identify the
fixed-index overlap geometry explicitly, prove a source-differentiable reduced
inverse and shell factorization on a constant-rank zero-mode chart, and show
that the V8 regulator-universality argument survives after the zero-mode line is
carried as an independent finite source packet.

Throughout this continuation let
\begin{equation}
 \mathcal H=\mathcal H_+\oplus\mathcal H_-,
 \qquad
 \dim\mathcal H_+=\dim\mathcal H_-=m,
 \qquad
 Q:=P_-
 \label{eq:v11-chiral-split}
\end{equation}
with $Q$ the orthogonal projection onto $\mathcal H_-$.  Let
$P=(I+\varepsilon)/2$ be the positive spectral projector of a gapped Hermitian
Wilson kernel, and put $W=\operatorname{Ran}P$ and
$p=\dim W$.  We write
\begin{equation}
 k:=p-m.
 \label{eq:v11-index-k}
\end{equation}
The sign convention is the index of the finite chiral map below; changing the
usual physics convention merely changes the sign of $k$.

\subsection{The exact fixed-index overlap map}

Let
\begin{equation}
 D_{\rm ov}=a^{-1}(I+\gamma_5\varepsilon),
 \qquad
 P_+=I-Q.
 \label{eq:v11-overlap}
\end{equation}
The kinetic Weyl block before Yukawa insertion is the rectangular map
\begin{equation}
 D_P:=P_+D_{\rm ov}\big|_W:W\longrightarrow\mathcal H_+.
 \label{eq:v11-chiral-map}
\end{equation}

\begin{lemma}[Exact projection form of the overlap Weyl block]
\label{lem:v11-projection-form}
For every $v\in W$,
\begin{equation}
 \boxed{
 D_Pv=\frac{2}{a}P_+v.}
 \label{eq:v11-D-projection}
\end{equation}
Consequently
\begin{align}
 \ker D_P&=W\cap\mathcal H_-,
 \label{eq:v11-kernel}\\
 \operatorname{coker}D_P
 &\cong \mathcal H_+\ominus P_+W,
 \label{eq:v11-cokernel}
\end{align}
and
\begin{equation}
 \boxed{
 \dim\ker D_P-\dim\operatorname{coker}D_P
 =\dim W-m=k.}
 \label{eq:v11-index-formula}
\end{equation}
\end{lemma}

\begin{proof}
For $v\in W$, $\varepsilon v=v$.  Hence
\[
 D_{\rm ov}v=a^{-1}(I+\gamma_5)v
 =2a^{-1}P_+v,
\]
which proves \eqref{eq:v11-D-projection}.  The kernel and cokernel statements
follow immediately.  If
$a_0:=\dim(W\cap\mathcal H_-)$, then
$\operatorname{rank}(P_+|_W)=p-a_0$ and therefore
\[
 \dim\operatorname{coker}D_P=m-(p-a_0).
\]
Subtracting this from $a_0$ gives $p-m=k$.
\end{proof}

\begin{definition}[Minimal fixed-index stratum]
\label{def:v11-minimal-stratum}
The projector $P$, or its plane $W$, lies in the \emph{minimal index-$k$
stratum} if
\begin{equation}
 \dim\ker D_P=k_+:=\max(k,0),
 \qquad
 \dim\operatorname{coker}D_P=k_-:=\max(-k,0).
 \label{eq:v11-minimal-nullities}
\end{equation}
Equivalently, $D_P$ has maximal possible rank $\min(p,m)$ at fixed $p$.
Zero modes beyond $k_++k_-=|k|$ are called accidental paired zero modes.
\end{definition}

The minimal stratum is the natural fixed-index analogue of the V10 nonzero
big cell.  It excludes an index-changing event and also excludes a pair of
opposite-chirality zero modes appearing without changing the index.

\subsection{The minimal stratum is an affine bundle over the zero-mode Grassmannian}

\begin{theorem}[Geometry for $k\ge0$: kernel plane plus an affine graph]
\label{thm:v11-positive-stratum}
Let $k\ge0$ and $p=m+k$.  A plane
$W\in\operatorname{Gr}(m+k,\mathcal H)$ lies in the minimal index-$k$ stratum
if and only if there are unique data
\begin{equation}
 K\in\operatorname{Gr}(k,\mathcal H_-),
 \qquad
 A\in\operatorname{Hom}(\mathcal H_+,K^\perp)
 \label{eq:v11-positive-data}
\end{equation}
such that
\begin{equation}
 \boxed{
 W=K\oplus
 \{u+Au:u\in\mathcal H_+\}.}
 \label{eq:v11-positive-graph}
\end{equation}
Here $K=\ker D_P$.  Thus the minimal index-$k$ stratum is the total space of
the complex vector bundle
\begin{equation}
 \operatorname{Hom}(\underline{\mathcal H_+},\mathcal K^\perp)
 \longrightarrow
 \operatorname{Gr}(k,\mathcal H_-),
 \label{eq:v11-positive-bundle}
\end{equation}
where $\mathcal K$ is the tautological $k$-plane bundle.  In particular it
deformation retracts onto the kernel Grassmannian.
\end{theorem}

\begin{proof}
Minimality for $k\ge0$ means that $P_+|_W:W\to\mathcal H_+$ is surjective and
has kernel $K=W\cap\mathcal H_-$ of dimension $k$.  Orthogonally decompose
$\mathcal H_-=K\oplus K^\perp$.  For each $u\in\mathcal H_+$ choose
$w\in W$ with $P_+w=u$ and subtract its unique $K$ component.  The resulting
vector has the form $u+Au$ with $Au\in K^\perp$.  This defines a linear map
$A:\mathcal H_+\to K^\perp$.  Its graph is orthogonal to $K$, has dimension
$m$, and together with $K$ exhausts $W$.  Uniqueness follows because a graph
vector with vanishing $\mathcal H_+$ component is zero.

Conversely every plane in \eqref{eq:v11-positive-graph} has
$W\cap\mathcal H_-=K$ and projects surjectively to $\mathcal H_+$, so its
kernel and cokernel dimensions are $k$ and $0$.  The construction depends
smoothly on $K$ and $A$ in standard Grassmann charts.  Fibrewise contraction
$A\mapsto tA$, $0\le t\le1$, gives the deformation retraction.
\end{proof}

\begin{theorem}[Geometry for $k\le0$: cokernel plane plus an affine graph]
\label{thm:v11-negative-stratum}
Let $k=-s\le0$ and $p=m-s$.  A plane
$W\in\operatorname{Gr}(m-s,\mathcal H)$ lies in the minimal index-$k$ stratum
if and only if there are unique data
\begin{equation}
 C\in\operatorname{Gr}(s,\mathcal H_+),
 \qquad
 A\in\operatorname{Hom}(C^\perp,\mathcal H_-)
 \label{eq:v11-negative-data}
\end{equation}
such that
\begin{equation}
 \boxed{
 W=\{u+Au:u\in C^\perp\}.}
 \label{eq:v11-negative-graph}
\end{equation}
Here $C$ identifies canonically with $\operatorname{coker}D_P$.  Hence the
minimal index-$k$ stratum is the total space of the vector bundle
\begin{equation}
 \operatorname{Hom}(\mathcal C^\perp,\underline{\mathcal H_-})
 \longrightarrow
 \operatorname{Gr}(s,\mathcal H_+)
 \label{eq:v11-negative-bundle}
\end{equation}
and deformation retracts onto the cokernel Grassmannian.
\end{theorem}

\begin{proof}
Minimality for $k=-s<0$ means $P_+|_W$ is injective and its image
$S:=P_+W$ has codimension $s$.  Put $C=S^\perp$.  Then every element of $W$
can be written uniquely as $u+Au$ with $u\in S=C^\perp$ and
$Au\in\mathcal H_-$.  Conversely every such graph intersects
$\mathcal H_-$ trivially and projects isomorphically onto $C^\perp$, leaving
cokernel $C$.  Smoothness and the deformation retraction again follow from
$A\mapsto tA$.
\end{proof}

\begin{corollary}[All minimal-stratum topology is carried by the forced zero-mode plane]
\label{cor:v11-topology-zero-mode}
At fixed nonzero index, the affine graph coordinate carries no homotopy
obstruction.  The minimal stratum is homotopy equivalent to
\begin{equation}
 \boxed{
 \operatorname{Gr}(k,\mathcal H_-),\quad k>0;
 \qquad
 \operatorname{Gr}(-k,\mathcal H_+),\quad k<0.}
 \label{eq:v11-zero-mode-Grassmannian}
\end{equation}
Thus for a family over a background parameter space, every global obstruction
which survives inside the minimal stratum is zero-mode-bundle data.  On a
contractible protected chart the zero-mode bundle is trivializable and there is
no additional local projector obstruction.
\end{corollary}

\begin{proof}
The first statement is the pair of deformation retractions just proved.  A
complex vector bundle over a contractible base is trivial, so its classifying
map to the corresponding finite Grassmannian is homotopic to a constant map.
The affine graph fibre is contractible as well.
\end{proof}

\begin{firewall}[Fixed index is not the same as fixed zero-mode bundle globally]
Corollary~\ref{cor:v11-topology-zero-mode} is local on a contractible protected
background chart.  On a noncontractible background component, two families can
have the same numerical index but different kernel/cokernel bundles, determinant
line holonomy, or real Pfaffian orientation.  V11 does not identify such global
data merely from the integer $k$.
\end{firewall}

\subsection{The reduced kinetic block is explicitly invertible}

The affine description gives a canonical reduced inverse before any
perturbative estimate is used.

\begin{proposition}[Canonical reduced block on the minimal stratum]
\label{prop:v11-reduced-block}
In the notation of Theorems~\ref{thm:v11-positive-stratum} and
\ref{thm:v11-negative-stratum}, let
\begin{equation}
 J_Au=(u+Au)(I+A^*A)^{-1/2}
 \label{eq:v11-graph-isometry}
\end{equation}
be the graph isometry from the active $\mathcal H_+$ subspace onto the graph
part of $W$.  Then after deleting the forced kernel and cokernel, the kinetic
chiral map is
\begin{equation}
 \boxed{
 D_P^{\perp}J_A
 =\frac{2}{a}(I+A^*A)^{-1/2}.}
 \label{eq:v11-reduced-explicit}
\end{equation}
For $k\ge0$ the active source and target are both $\mathcal H_+$; for $k<0$
they are both $C^\perp$.  Consequently
\begin{align}
 s_{\min}(D_P^\perp)
 &=\frac{2}{a\sqrt{1+\|A\|^2}},
 \label{eq:v11-reduced-gap}\\
 \det D_P^\perp
 &=\left(\frac2a\right)^r
   \det(I+A^*A)^{-1/2}>0,
 \qquad r=\min(p,m),
 \label{eq:v11-reduced-det}
\end{align}
in these canonical graph frames.
\end{proposition}

\begin{proof}
The vector $J_Au$ is normalized because
$\|u+Au\|^2=\langle u,(I+A^*A)u\rangle$.  By
Lemma~\ref{lem:v11-projection-form}, $D_P$ keeps only the
$\mathcal H_+$ component and multiplies it by $2/a$.  This gives
\eqref{eq:v11-reduced-explicit}.  Functional calculus of the positive matrix
$I+A^*A$ gives the singular-value and determinant formulas.
\end{proof}

The positivity in \eqref{eq:v11-reduced-det} is a graph-frame statement for the
pure kinetic overlap block.  It does not orient a separate real Pfaffian line
and does not remove the zero-mode exterior factor.

\subsection{Constant-rank zero-mode charts for the full chiral/Yukawa block}

The shell theorem uses the full finite chiral block, including Yukawa and
source insertions.  We therefore formulate the analytic part independently of
the special projection coordinates.

\begin{definition}[Protected constant-rank zero-mode chart]
\label{def:v11-zero-chart}
Let $D(\zeta):E\to F$ be a smooth finite-cutoff chiral/Yukawa family on a
compact source/background chart $\mathcal U$.  It is a
\emph{protected index-$k$ zero-mode chart} if:
\begin{enumerate}[label=\textup{(Z\arabic*)},leftmargin=3em]
\item $\operatorname{ind}D(\zeta)=k$ and
$\dim\ker D(\zeta)=k_+$,
$\dim\operatorname{coker}D(\zeta)=k_-$ throughout $\mathcal U$;
\item every nonzero singular value obeys
\begin{equation}
 s_j(D(\zeta))\ge\tau_*>0;
 \label{eq:v11-reduced-floor}
\end{equation}
\item the fixed source derivatives required by the chosen coefficient packet
exist and obey the native finite-range/quasilocal bounds inherited from
V1--V10.
\end{enumerate}
The finite-dimensional zero modes are retained as soft variables and are never
integrated into a hard shell.
\end{definition}

\begin{lemma}[Smooth kernel/cokernel projectors and coefficientwise source jets]
\label{lem:v11-zero-projectors}
On a protected zero-mode chart, the orthogonal projectors
$\Pi_K$ onto $K=\ker D$ and $\Pi_C$ onto
$C=\ker D^*=\operatorname{coker}D$ are smooth.  For any contour
$\Gamma_0$ enclosing $0$ but no nonzero singular square,
\begin{align}
 \Pi_K
 &=\frac{1}{2\pi i}\oint_{\Gamma_0}
 (z-D^*D)^{-1}\,dz,
 \label{eq:v11-kernel-Riesz}\\
 \Pi_C
 &=\frac{1}{2\pi i}\oint_{\Gamma_0}
 (z-DD^*)^{-1}\,dz.
 \label{eq:v11-coker-Riesz}
\end{align}
Every prescribed fixed source derivative is a finite sum of resolvent words and
is bounded by a constant depending only on the fixed derivative order,
$\tau_*^{-1}$, and the declared source stocks.
\end{lemma}

\begin{proof}
By \eqref{eq:v11-reduced-floor}, the spectra of $D^*D$ and $DD^*$ have a
uniform gap between $0$ and $\tau_*^2$.  Choose a common contour inside this
gap.  Riesz functional calculus gives the projectors.  Differentiating a
resolvent uses
\[
 \partial(z-A)^{-1}=(z-A)^{-1}(\partial A)(z-A)^{-1};
\]
iterating produces a finite sum at every fixed derivative order.  The contour
resolvents are uniformly bounded by the inverse spectral separation.
\end{proof}

\begin{proposition}[Moore--Penrose inverse in a zero-mode chart]
\label{prop:v11-pseudoinverse}
Define
\begin{equation}
 \boxed{
 D^\dagger
 :=(D^*D+\Pi_K)^{-1}D^*
 =D^*(DD^*+\Pi_C)^{-1}.}
 \label{eq:v11-pseudoinverse}
\end{equation}
Then $D^\dagger$ is the inverse of
$D^\perp:K^\perp\to C^\perp$ and vanishes on $C$.  Moreover
\begin{equation}
 \|D^\dagger\|\le\tau_*^{-1},
 \label{eq:v11-pseudoinverse-bound}
\end{equation}
and every fixed source derivative is a finite expression in
$D,D^*,\Pi_K,\Pi_C$ and their source jets with uniformly bounded inverse factors.
\end{proposition}

\begin{proof}
On $K^\perp$, $D^*D$ is positive with lower edge $\tau_*^2$, while on $K$ the
added projector is the identity.  Hence $D^*D+\Pi_K$ is invertible.  Its inverse
times $D^*$ is exactly the singular-value definition of the Moore--Penrose
inverse.  The second expression is obtained similarly on the target.  The norm
is the reciprocal of the smallest nonzero singular value.  Differentiation uses
the ordinary inverse rule together with Lemma~\ref{lem:v11-zero-projectors}.
\end{proof}

\begin{corollary}[Local Kato frames for the zero-mode bundle]
\label{cor:v11-zero-Kato}
On every star-shaped subchart of $\mathcal U$, the kernel and cokernel bundles
admit unitary Kato frames with every fixed source derivative bounded.  Their top
exterior powers furnish local frames of the zero-mode determinant line.  A real
Pfaffian zero-mode line additionally requires its declared orientation.
\end{corollary}

\begin{proof}
Apply the V1 Kato ODE to the smooth finite-rank projectors
$\Pi_K$ and $\Pi_C$.  Lemma~\ref{lem:v11-zero-projectors} supplies the derivative
bounds.  Take top exterior powers.  Complex exterior powers are canonically
oriented as complex lines; a real Pfaffian line carries the separate orientation
data already distinguished in the inherited construction.
\end{proof}

\subsection{Exact zero-mode-preserving shell factorization}

The inherited finite zero-mode theorem factors the full determinant line into a
zero-mode factor and a nonzero reduced factor.  V11 now proves that the latter
factor admits the same exact shell Schur algebra as V1--V8 while the zero-mode
factor is transported unchanged.

Choose orthogonal decompositions
\begin{equation}
 E=K\oplus E_\ell\oplus E_h,
 \qquad
 F=C\oplus F_\ell\oplus F_h,
 \label{eq:v11-zero-shell-split}
\end{equation}
with the zero modes entirely in the retained soft sector.  Relative to the
reduced complements, write
\begin{equation}
 D^\perp=
 \begin{pmatrix}
 A&B\\ C_1&H
 \end{pmatrix}:
 E_\ell\oplus E_h\longrightarrow F_\ell\oplus F_h.
 \label{eq:v11-reduced-block-matrix}
\end{equation}

\begin{theorem}[Zero-mode-preserving Schur factorization]
\label{thm:v11-zero-Schur}
Assume the hard block $H$ is invertible.  Then the reduced Schur complement
\begin{equation}
 S=A-BH^{-1}C_1
 \label{eq:v11-zero-Schur-complement}
\end{equation}
is invertible whenever $D^\perp$ is invertible, and the full determinant-line
section factors canonically as
\begin{equation}
 \boxed{
 \mathcal Z_D\otimes\det(D^\perp)
 =\mathcal Z_D\otimes\det(H)\otimes\det(S),}
 \label{eq:v11-zero-Schur-line}
\end{equation}
where $\mathcal Z_D$ denotes the inherited kernel/cokernel exterior factor in
the chosen determinant-line convention.  In particular, integrating one hard
shell does not delete, divide by, or re-fit the zero-mode line.
\end{theorem}

\begin{proof}
The inherited zero-mode factorization first separates $K$ and $C$, leaving the
square invertible map $D^\perp:K^\perp\to C^\perp$.  The ordinary finite Schur
identity applied to \eqref{eq:v11-reduced-block-matrix} gives
\[
 \det D^\perp=\det H\,\det(A-BH^{-1}C_1).
\]
If $D^\perp$ and $H$ are invertible, the right-hand determinant forces $S$ to
be invertible.  Tensor the scalar reduced identity with the unchanged
zero-mode exterior factor.  No operation acts on $K$ or $C$ because those
coordinates were excluded from $E_h,F_h$ by construction.
\end{proof}

\begin{corollary}[Differentiated shell line with zero-mode contacts]
\label{cor:v11-zero-Schur-derivatives}
In local Kato zero-mode frames, every fixed source derivative of
\eqref{eq:v11-zero-Schur-line} is the finite Leibniz/Bell-polynomial expansion
of:
\begin{enumerate}[label=\textup{(D\arabic*)},leftmargin=3em]
\item the zero-mode frame/saturation factor;
\item the hard reduced determinant or Pfaffian factor;
\item the reduced Schur factor;
\item every mixed contact generated by derivatives hitting more than one of
these three factors.
\end{enumerate}
No source-dependent zero-mode endpoint term may be discarded as field
independent.
\end{corollary}

\begin{proof}
Differentiate the exact tensor-product identity.  The reduced inverse
insertions are controlled by Proposition~\ref{prop:v11-pseudoinverse}; the hard
block uses the V1--V8 determinant/Pfaffian derivative calculus.  Kato-frame
derivatives are finite by Corollary~\ref{cor:v11-zero-Kato}.  Repeated
differentiation is exactly the finite product-partition expansion already used
for the V1 line cocycle.
\end{proof}

\subsection{Saturation writers and scalar normalized amplitudes}

A nonzero-index vacuum determinant is not a scalar nonzero partition function.
The inherited finite zero-mode theorem states that a nonvanishing amplitude
requires a section of the dual zero-mode line.  We make that datum explicit in
the shell state.

\begin{definition}[Zero-mode saturation writer]
\label{def:v11-saturation-writer}
On a protected zero-mode chart let
\begin{equation}
 \omega_Z\in\Gamma(\mathcal Z_D^*)
 \label{eq:v11-saturation-writer}
\end{equation}
be a chosen smooth nonvanishing local saturation writer, with all required
source derivatives retained.  The corresponding scalarized reduced fermion
factor is
\begin{equation}
 \mathcal A_D^{\rm sat}
 :=\langle\omega_Z,z_D\rangle\,\det'(D),
 \label{eq:v11-saturated-amplitude}
\end{equation}
where $z_D$ is the zero-mode exterior section in the inherited
factorization and $\det'(D)=\det D^\perp$.
\end{definition}

The writer may represent explicit external fermion insertions, a boundary
state, or another physical saturation mechanism.  It is not supplied by the
pseudodeterminant itself.

\begin{theorem}[Exact shell factorization of the saturated amplitude]
\label{thm:v11-saturated-shell}
Under Theorem~\ref{thm:v11-zero-Schur},
\begin{equation}
 \boxed{
 \mathcal A_D^{\rm sat}
 =\langle\omega_Z,z_D\rangle
   \det(H)\det(S).}
 \label{eq:v11-saturated-shell}
\end{equation}
Every retained source derivative is obtained by differentiating this one
identity.  A change of local zero-mode frame or saturation writer acts by an
endpoint line coboundary/contact transformation of the same type as V7.
\end{theorem}

\begin{proof}
Pair the zero-mode factor in \eqref{eq:v11-zero-Schur-line} with
$\omega_Z$ and use the inherited zero-mode saturation theorem.  The frame-change
statement follows because changing a nonzero local frame multiplies
$z_D$ and the dual writer by reciprocal nonzero functions.  If only
one is reparametrized, the resulting local function and all of its source
derivatives are precisely finite source/contact coordinates and must be
transported as in V7.
\end{proof}

\begin{firewall}[The pseudodeterminant is not the amplitude]
The factor $\det'(D)$ controls the invertible complement but does not saturate
the Grassmann zero modes and does not orient a real Pfaffian kernel.  Every V11
scalar continuum statement below is therefore either line-valued or explicitly
conditioned on the saturation writer $\omega_Z$ (and, when relevant, a real
orientation).  This is inherited physical data, not a regulator counterterm.
\end{firewall}

\subsection{The fixed-index extension of the V8 hard-shell calculus}

The hard ultraviolet estimates never need an inverse on the zero-mode space.
They need only the inverse of the reduced hard block.  This observation is the
analytic reason the coefficientwise programme can be extended to a fixed-index
sector without lifting the zero modes.

\begin{definition}[Fixed-index admissible regulator homotopy]
\label{def:v11-index-admissible}
A native regulator path $t\mapsto H_{W,t}$ on a compact background/source chart
is \emph{$(\mathfrak q,k)$-admissible} if:
\begin{enumerate}[label=\textup{(I\arabic*)},leftmargin=3em]
\item it satisfies the V8 locality, covariance, physical-principal-symbol and
full-zone source-jet conditions at the prescribed fixed coefficient data
$\mathfrak q$;
\item the Wilson gap remains uniformly positive, so $\operatorname{rank}P_t$
and hence the index $k$ are constant;
\item the full chiral/Yukawa block stays in the minimal index-$k$ stratum and
its reduced nonzero singular values have a uniform floor $\tau_*>0$;
\item kernel/cokernel modes are retained outside every hard shell and their
Kato source jets are included in the finite source packet;
\item the zero-mode determinant/Pfaffian line and the chosen saturation writer
(or real orientation) are transported continuously along the path;
\item the same finite local BV/EFT/source normalization chart remains of
constant rank.
\end{enumerate}
\end{definition}

\begin{theorem}[Fixed-index hard-shell comparison]
\label{thm:v11-index-hard-shell}
Along a $(\mathfrak q,k)$-admissible homotopy, every V8 relative hard-shell
estimate remains valid after replacing the full chiral inverse by the reduced
inverse $D^\dagger$ and adjoining the finite zero-mode source packet.  In
particular, after the common local Taylor/source projection, a regulator-marked
order-$\le L$ forest obeys
\begin{equation}
 \boxed{
 \|\mathsf E_{\Gamma,j,t}^{\rm rel}\|_*
 \le
 \begin{cases}
 C_{\mathfrak q,k}\,a\Lambda_j,&\Lambda_j\le M,\\[0.2em]
 C_{\mathfrak q,k}\,\mu/\Lambda_j,&\Lambda_j>M,
 \end{cases}}
 \label{eq:v11-index-relative-forest}
\end{equation}
for every intermediate scale $\mu\ll M\ll a^{-1}$.  The same estimate holds
after every stored BV/source contraction descendant.
\end{theorem}

\begin{proof}
By Definition~\ref{def:v11-index-admissible}, the zero modes are never assigned
to the high shell.  Every internal hard fermion contraction therefore uses the
inverse of the reduced hard block, which is the restriction of $D^\dagger$ to
the relevant shell.  Proposition~\ref{prop:v11-pseudoinverse} supplies the
uniform fixed-source derivative calculus on the protected chart.  The relative
physical-shell insertion is the same irrelevant regulator mark as in V8,
because the microscopic principal Dirac--Yukawa symbol is unchanged; the finite
zero-mode projectors contribute only bounded finite-rank source factors.

For $\Lambda_j\le M$, the V8 rooted ideal therefore retains one factor
$a\Lambda_j$.  For $\Lambda_j>M$, subtract the complete local Taylor/source jet
of the reduced forest; the first omitted soft momentum again supplies
$\mu/\Lambda_j$.  Proper subgraphs remain owned by the same lower-loop common
actions.  Differentiating a zero-mode frame or saturation writer creates a
finite local/source insertion but no hard inverse, so it is assigned to the
same contraction-complete packet.  The V5 future-contact construction then
prevents a later BV contraction from converting an unrecorded zero-mode term
into a new local divergence.  Summing the finite forest family proves
\eqref{eq:v11-index-relative-forest}.
\end{proof}

\begin{corollary}[The two-scale estimate is unchanged by protected zero modes]
\label{cor:v11-index-two-scale}
With one common local normalization/source/zero-line velocity chosen to keep
the physical normalization and saturation coordinates fixed,
\begin{equation}
 \left\|
 \partial_t\widehat{\mathfrak C}_{a,t,\lambda}^{(k),\rm sat}
 \right\|_{\mathfrak q}
 \le C_{\mathfrak q,k}
 \left(aM+\frac{\mu}{M}\right).
 \label{eq:v11-index-two-scale}
\end{equation}
Choosing $M=\sqrt{\mu/a}$ gives
\begin{equation}
 \boxed{
 \left\|
 \partial_t\widehat{\mathfrak C}_{a,t,\lambda}^{(k),\rm sat}
 \right\|_{\mathfrak q}
 \le C_{\mathfrak q,k}\sqrt{a\mu}.}
 \label{eq:v11-index-sqrt}
\end{equation}
\end{corollary}

\begin{proof}
The proof of the V8 two-scale theorem applies to the reduced forests by
Theorem~\ref{thm:v11-index-hard-shell}.  The local velocity now also includes
the finite Kato zero-mode-frame and saturation coordinates.  They form a finite
triangular extension of the V7 normalization system and do not change the
geometric shell sums.  Optimize exactly as in V8.
\end{proof}

\subsection{Fixed-index regulator universality with zero-mode saturation}

\begin{theorem}[Local fixed-index regulator universality]
\label{thm:v11-index-universality}
Fix finite coefficient data $\mathfrak q$, an integer $k$, a compact protected
constant-rank zero-mode chart, and a nonvanishing local zero-mode saturation
writer (plus a real orientation if a Pfaffian sector is present).  Let regulators
$A$ and $B$ be the endpoints of a $(\mathfrak q,k)$-admissible native homotopy.
Then there is an invertible finite local matching map
\begin{equation}
 F_{\infty,\mathfrak q,k}^{B\leftarrow A}
 \label{eq:v11-index-matching}
\end{equation}
such that the saturated normalized continuum states obey
\begin{equation}
 \boxed{
 \widehat{\mathfrak C}_{\infty,\lambda}^{B,(k),\rm sat}
 =F_{\infty,\mathfrak q,k}^{B\leftarrow A}
  \widehat{\mathfrak C}_{\infty,\lambda}^{A,(k),\rm sat}.}
 \label{eq:v11-index-universality}
\end{equation}
The equivalence includes normalized connected coefficients with the declared
zero-mode insertions, composite/operator-mixing coefficients, stress/contact
terms, ghosts/antifields, the common BV/BRST/master relations on the saturated
source packet, and determinant/Pfaffian reduced-line insertions.  The matching
maps may be chosen projectively in perturbative order.
\end{theorem}

\begin{proof}
Corollary~\ref{cor:v11-index-two-scale} is the sole new analytic input required
in the V8 universality proof.  Integrating the regulator parameter gives a
finite-cutoff matching error $O(\sqrt{a\mu})$.  The local matching flow contains
the same finite EFT/BV/source coordinates as V8 together with the finite
zero-mode frame/saturation coordinates.  Its cutoff-to-cutoff variation is
summable by the V6--V8 argument, so it has an invertible continuum limit.
Equation \eqref{eq:v11-index-universality} follows as $a\downarrow0$.

The connected functional is taken only after the zero-mode line has been paired
with the declared saturation writer; Theorem~\ref{thm:v11-saturated-shell}
shows that this scalarization is compatible with every shell elimination.  The
BV/source descendants are included because the zero-mode Kato packet is part
of the same contraction-complete source state.  Chronological V7 normalization
freezes each lower-loop zero-mode/source matching when first chosen, proving
projectivity in terminal loop order.
\end{proof}

\begin{corollary}[The index-zero theorem is recovered without a separate case]
\label{cor:v11-kzero-recovery}
For $k=0$, $K=C=0$, the saturation line is trivial and the minimal stratum is
the V10 balanced big cell.  Theorem~\ref{thm:v11-index-universality} reduces to
the V8--V10 regulator-universality theorem.
\end{corollary}

\begin{proof}
All zero-mode projectors and line factors vanish from the construction, while
$D^\dagger=D^{-1}$.  The positive- and negative-index graph descriptions both
reduce to the balanced graph cell of V10.
\end{proof}

\subsection{A local phase-lifting theorem on a contractible zero-mode chart}

The preceding theorem is conditional on a native admissible path.  The
finite-dimensional topological part of constructing such a path can now be
closed completely on a local chart.

\begin{theorem}[Local minimal-stratum connectedness after zero-mode matching]
\label{thm:v11-local-stratum-connected}
Let $X$ be a contractible compact parameter chart.  Let
$P_0,P_1:X\to\operatorname{Gr}(m+k,\mathcal H)$ be two smooth families in the
minimal index-$k$ stratum.  Choose a smooth unitary identification of their
kernel bundles when $k>0$, or of their cokernel bundles when $k<0$.  Then there
is a smooth homotopy $P_t$ through the same minimal stratum which realizes this
zero-mode identification and linearly interpolates the remaining affine graph
coordinate.  In particular, after the zero-mode bundle is matched, there is no
further local projector-level phase obstruction.
\end{theorem}

\begin{proof}
Assume first $k>0$.  By Theorem~\ref{thm:v11-positive-stratum}, each family is
specified by a kernel subbundle $K_i\subset X\times\mathcal H_-$ and a section
$A_i$ of
$\operatorname{Hom}(\underline{\mathcal H_+},K_i^\perp)$.  The chosen unitary
bundle identification and contractibility of $X$ allow a smooth block-unitary
homotopy in $\mathcal H_-$ carrying $K_0$ to $K_1$.  Apply it to the first
projector family; this preserves the minimal stratum.  We may then assume the
kernel bundle is the same $K$.  Both graph coordinates are sections of the
same vector bundle, so
$A_t=(1-t)A_0+tA_1$ remains in the affine fibre and therefore in the minimal
stratum.

For $k<0$, use Theorem~\ref{thm:v11-negative-stratum} and the same argument on
the cokernel bundle $C\subset X\times\mathcal H_+$.  The case $k=0$ is the V10
big-cell interpolation.
\end{proof}

\begin{firewall}[Projector connectedness is not yet a global native lift]
Theorem~\ref{thm:v11-local-stratum-connected} closes the finite-dimensional
local topology.  It does not by itself prove that an arbitrary projector
homotopy is generated by a cutoff-uniform finite-range or exponentially local,
gauge/spin-covariant microscopic Wilson kernel on an entire nontrivial gauge
background component.  V11 regulator universality therefore uses the native
admissibility hypothesis of Definition~\ref{def:v11-index-admissible}.  A
global locality-preserving lift is the next separate gate.
\end{firewall}

\subsection{What becomes of the determinant line at fixed nonzero index}

The geometry above also identifies precisely where a global obstruction can
remain.  On the minimal stratum the reduced graph fibre is affine, while the
zero-mode plane is not.  The determinant/Pfaffian phase therefore lives on the
zero-mode bundle after the reduced phase has been normalized.

\begin{proposition}[Residual line is the zero-mode line]
\label{prop:v11-residual-line}
On a protected minimal index-$k$ chart, choose the canonical positive graph
frame of Proposition~\ref{prop:v11-reduced-block} for the pure overlap kinetic
complement and the inherited finite normalization for the interacting reduced
block.  Then every residual topological line factor is carried by the kernel
and cokernel exterior powers.  In particular, on a contractible chart the
complex determinant line is locally trivial, while on a noncontractible
background component its possible holonomy is exactly a zero-mode/Fredholm-line
problem and is not removed by the reduced pseudodeterminant.
\end{proposition}

\begin{proof}
The pure kinetic reduced determinant is positive in the canonical graph frame
by \eqref{eq:v11-reduced-det}.  For the interacting block, the inherited
zero-mode factorization writes the full line as the zero-mode exterior factor
tensored with the reduced invertible determinant line.  The latter is locally
scalarized by the V7/V11 normalization and Schur cocycle.  Therefore any
remaining line topology is in the zero-mode exterior factor and its transition
functions.  Local triviality on a contractible chart is standard for complex
line bundles.  No statement about loop holonomy follows without examining the
global background component.
\end{proof}

\subsection{V11 anti-circularity audit}
\label{sec:v11-audit}

\begin{enumerate}[leftmargin=2.4em]
\item \textbf{Zero modes were not deleted.}
They are retained as the kernel/cokernel bundle and its exterior line.  Only
the reduced complement is inverted.
\item \textbf{The pseudodeterminant was not promoted to a vacuum amplitude.}
A scalar amplitude requires the explicit dual zero-mode saturation writer of
Definition~\ref{def:v11-saturation-writer}.
\item \textbf{Fixed numerical index was not assumed to classify a global
family.}
The minimal stratum deformation retracts to the zero-mode Grassmannian; global
kernel/cokernel bundle data and line holonomy remain visible.
\item \textbf{Accidental paired zero modes are excluded by a stated reduced
singular-value floor.}
A crossing of that floor leaves the protected minimal stratum and is not hidden
inside a divergent inverse.
\item \textbf{Hard shells never integrate the zero modes.}
Every hard fermion contraction uses the reduced inverse $D^\dagger$ on
$K^\perp\to C^\perp$.
\item \textbf{Source derivatives of the zero-mode projectors are retained.}
They are Riesz-projector/Kato derivatives and enter the same finite contact
packet as the determinant-line frame.
\item \textbf{The V8 two-scale estimate is applied only after the exact local
zero-mode factorization.}
Finite-rank zero-mode factors do not supply the ultraviolet scale gain; the
gain still comes from one irrelevant microscopic regulator mark on the reduced
hard forest.
\item \textbf{Projector-level connectedness was not silently called a native
locality theorem.}
The native lift remains an explicit admissibility gate outside the local
finite-dimensional topology.
\end{enumerate}

\section{V11 status ledger}
\label{sec:v11-status}

\subsection*{Proved in V11}
\begin{enumerate}[leftmargin=2.2em]
\item The finite overlap Weyl kinetic map is exactly
$D_P=2a^{-1}P_+|_{\operatorname{Ran}P}$, and its Fredholm index is
$\operatorname{rank}P-m$.
\item The minimal fixed-index projector stratum is an affine vector bundle over
the forced kernel Grassmannian for $k>0$ or the forced cokernel Grassmannian for
$k<0$.  Hence all nonzero-mode projector data are contractible once the
zero-mode plane is fixed.
\item In canonical graph frames the reduced kinetic block is explicitly
$2a^{-1}(I+A^*A)^{-1/2}$, with a positive reduced determinant and an exact
singular-value floor.
\item On a protected constant-rank full chiral/Yukawa chart, the kernel and
cokernel projectors have Riesz formulas, Kato frames and fixed-order source jets;
the Moore--Penrose inverse is uniformly bounded on the reduced complement.
\item The inherited finite zero-mode determinant-line factorization is stable
under exact shell Schur elimination: the zero-mode exterior factor is carried
unchanged while the reduced determinant obeys the ordinary hard-block/Schur
product.  All mixed source contacts are retained.
\item A nonzero local saturation writer in the dual zero-mode line converts the
line-valued factor into a scalar saturated amplitude, and this scalarization
commutes exactly with shell elimination.
\item Along every fixed-index V11-admissible native regulator homotopy, the V8
relative forest bounds survive with $D^{-1}$ replaced by the reduced
pseudoinverse.  The same $aM+\mu/M$ two-scale bound and
$O(\!\sqrt{a\mu}\!)$ optimized regulator difference follow.
\item Therefore endpoints of such a homotopy have the same matched
coefficientwise saturated continuum, including the complete source/BV/contact
and reduced determinant/Pfaffian packets, with strict perturbative-order
projectivity.
\item On a contractible protected parameter chart, after matching the
kernel/cokernel bundle there is no additional local projector-level phase
obstruction: the remaining graph coordinate interpolates linearly in an affine
fibre.
\end{enumerate}

\subsection*{Inherited}
\begin{enumerate}[leftmargin=2.2em]
\item The Grand Tensor finite zero-mode factorization, saturation identity,
Fredholm determinant line and the statement that zero modes carry the residual
massless line obstruction.
\item V1--V7: source-complete fixed-order continuum, common BV restoration,
arbitrary-fixed-loop forest theorem, ultraviolet Cauchy convergence,
scheme/reference covariance and exact loop-order projectivity.
\item V8: regulator-marked relative forests, line transgression and the
two-scale universality theorem along admissible gapped homotopies.
\item V9--V10: regulator-class phase lifting in the index-zero
Clifford and general balanced multiband single-cone sectors.
\end{enumerate}

\subsection*{Conditional}
\begin{enumerate}[leftmargin=2.2em]
\item V11 works on a minimal constant-index stratum with a positive reduced
singular-value floor.  An accidental zero-mode pair or index-changing crossing
requires a different stratified analysis.
\item Scalar normalized amplitudes require a declared nonzero zero-mode
saturation writer; a real Pfaffian sector additionally requires its physical
orientation.
\item The local projector connectedness theorem does not itself manufacture a
uniformly local native microscopic path on an arbitrary global topological
gauge component.  The regulator-universality theorem assumes the path satisfies
Definition~\ref{def:v11-index-admissible}.
\item All continuum claims remain coefficientwise at fixed finite
$(L,\Delta,r,N)$ and positive IR/reference scale.  No uniform large-loop or
massless-IR limit is asserted.
\end{enumerate}

\subsection*{Open beyond V11}
\begin{enumerate}[leftmargin=2.2em]
\item Prove a \emph{global native fixed-index lifting theorem}: starting from a
projector homotopy in the minimal stratum on a nontrivial gauge-background
component, construct a cutoff-uniform finite-range or exponentially local,
gauge/spin-covariant Wilson-kernel homotopy with the required source jets and
reduced chiral gap.
\item Compute and trivialize, when physically appropriate, the global
zero-mode determinant/Pfaffian line holonomy.  Equal index and local anomaly
cancellation alone do not supply this.
\item Analyze crossings between minimal strata: creation/annihilation of an
opposite-chirality zero-mode pair and genuine index change.  This requires a
spectral-flow/divisor theorem rather than a bounded pseudoinverse.
\item The independent stronger problems remain: positive-IR removal, analytic
finite-coupling RG trajectories, perturbative summability and nonperturbative
completion.
\end{enumerate}

\subsection*{Exact next load-bearing theorem}
\begin{center}
\fbox{\parbox{0.92\textwidth}{
\textbf{Global zero-mode line transport and native fixed-index lifting.}
On a connected nonzero-index background component, promote the local V11
kernel/cokernel Kato bundles to a globally patched zero-mode determinant or
Pfaffian line, construct a locality-preserving native lift of minimal-stratum
projector homotopies, and compute the resulting loop holonomy/spectral-flow
cocycle.  Prove regulator-class universality precisely on the subcomponent
where this global line datum and the reduced chiral gap agree.  This is the
remaining gate between the local fixed-index theorem of V11 and a global
nonzero-index regulator-class theorem.}}
\end{center}

\section*{Append-only continuation marker: V11}
\addcontentsline{toc}{section}{Append-only continuation marker: V11}
\textbf{End of V11.}  The complete V1--V10 body above is retained verbatim.
V11 replaces the full stratified-Fredholm picture by the minimal fixed-index
geometry actually needed by the shell theorem.  The overlap Weyl map is the
projection $2a^{-1}P_+|_{\operatorname{Ran}P}$; its minimal index-$k$ stratum is
an affine bundle over the forced kernel or cokernel Grassmannian.  The reduced
complement has an explicit positive inverse, while the zero-mode bundle and its
exterior line are carried as finite source data.  Riesz/Kato source calculus,
Moore--Penrose hard propagation, exact zero-mode-preserving Schur factorization
and an explicit saturation writer extend the V8 two-scale universality theorem
to every native homotopy which preserves the fixed-index reduced gap.  The
remaining obstruction is global: locality-preserving lifting and holonomy of
the zero-mode determinant/Pfaffian line on noncontractible background
components.


\clearpage
\section{V12 continuation: global zero-mode bundles, faithful-spin line transport, and the native locality obstruction}
\label{sec:v12-continuation}

V11 closed the fixed-index shell calculus on a contractible protected chart.
The remaining gate has three logically distinct layers which should not be
collapsed into one ``global phase'' condition:
\begin{equation}
 \boxed{
 \text{determinant/Pfaffian line datum}
 \quad\subsetneq\quad
 \text{full kernel/cokernel bundle datum}
 \quad\subsetneq\quad
 \text{native local-operator path datum}.}
 \label{eq:v12-three-layers}
\end{equation}
The first layer controls scalarization and phase/orientation.  The second
controls the homotopy class of a family inside the V11 minimal fixed-index
stratum.  The third asks whether such a family homotopy is realized by a
cutoff-uniform local microscopic Wilson kernel with the required source jets.
The inclusions in \eqref{eq:v12-three-layers} are strict in general.

The Grand Tensor manuscript already contains two load-bearing inputs for the
first layer.  First, its finite zero-mode theorem identifies the Fredholm line
with the kernel/cokernel exterior factor tensored with the reduced nonzero
line.  Second, on the ordinary-spin faithful Standard-Model branch with
\(G_6=(\mathrm{SU}(3)\times\mathrm{SU}(2)\times\mathrm{U}(1))/\mathbb Z_6\),
it proves vanishing of the complete local anomaly polynomial together with the
five-dimensional faithful-spin bordism/eta character, and hence a unit
parallel section of the complete complex determinant line on every smooth
connected family to which its geometric hypotheses apply.  These statements
are inherited and are not reproved below.

The new work is instead:
\begin{enumerate}[label=\textup{(V12.\arabic*)},leftmargin=3.8em]
\item to patch the V11 local Kato zero-mode frames globally and identify their
exact Cech/connection cocycle;
\item to combine the inherited faithful-spin parallel determinant section with
V11's reduced nonzero determinant to construct a global zero-mode saturation
writer for the complex Standard-Model determinant packet;
\item to identify the full global projector obstruction with the homotopy class
of the kernel or cokernel classifying map, not merely with its determinant
line;
\item to prove a finite-dimensional stable-range theorem under which an
isomorphism of zero-mode bundles already gives the required projector homotopy;
\item to show by an explicit volume-growing example that an abstract zero-mode
projector homotopy need not be a native quasilocal regulator homotopy;
\item to reduce global fixed-index regulator universality to path connectedness
of an open protected subset of the native exponentially local operator space,
and to prove that every such path may be replaced by a finite polygonal native
chain to which the V11--V8 estimates apply.
\end{enumerate}

Throughout, let \(B\) be a connected compact smooth parameter space (or a
finite CW complex with smooth charts where derivatives are used) carrying a
finite-cutoff chiral/Yukawa family
\begin{equation}
 D_b:E_b\longrightarrow F_b,
 \qquad b\in B,
 \label{eq:v12-family}
\end{equation}
which stays in the V11 minimal index-\(k\) stratum and has a uniform reduced
singular-value floor
\begin{equation}
 s_{\rm nonzero}(D_b)\ge \tau_*>0.
 \label{eq:v12-reduced-floor}
\end{equation}
Write
\begin{equation}
 K_b:=\ker D_b,
 \qquad
 C_b:=\ker D_b^*,
 \qquad
 k_+=\dim K_b,
 \qquad
 k_-=\dim C_b.
 \label{eq:v12-KC}
\end{equation}
On the minimal stratum \(k_+=\max(k,0)\) and
\(k_-=\max(-k,0)\).

\subsection{Global Cech--Kato patching of the zero-mode bundles}
\label{sec:v12-Cech-Kato}

The V11 Riesz formulas are global even when a single Kato frame is not.
Consequently the kernel and cokernel themselves patch without additional
choices.

\begin{theorem}[Global kernel and cokernel bundles]
\label{thm:v12-global-KC}
Under \eqref{eq:v12-reduced-floor}, the V11 Riesz projectors
\(\Pi_K(b)\) and \(\Pi_C(b)\) define smooth Hermitian vector bundles
\begin{equation}
 K\longrightarrow B,
 \qquad
 C\longrightarrow B
 \label{eq:v12-global-KC}
\end{equation}
of ranks \(k_+\) and \(k_-\).  On every finite good cover
\((U_\alpha)_\alpha\) of \(B\), choose V11 unitary Kato frames
\(e^K_\alpha,e^C_\alpha\).  On an overlap,
\begin{equation}
 e^K_\beta=e^K_\alpha u^K_{\alpha\beta},
 \qquad
 e^C_\beta=e^C_\alpha u^C_{\alpha\beta},
 \label{eq:v12-frame-transition}
\end{equation}
with
\(u^K_{\alpha\beta}:U_{\alpha\beta}\to U(k_+)\) and
\(u^C_{\alpha\beta}:U_{\alpha\beta}\to U(k_-)\), satisfying the exact
Cech cocycle identities on triple overlaps.  Every fixed source derivative of
the transition maps required by the coefficient packet is bounded on the
compact cover.
\end{theorem}

\begin{proof}
The Riesz-projector formulas of V11 use a common contour separated from the
nonzero singular spectrum by \(\tau_*^2\).  They are therefore globally smooth
functions of \(b\), and their constant ranks give smooth subbundles of the
ambient finite-cutoff bundles.  On each contractible \(U_\alpha\), V11 Kato
transport gives orthonormal frames.  Two orthonormal frames of the same fibre
differ by a unique unitary matrix, giving
\eqref{eq:v12-frame-transition}.  Uniqueness immediately gives
\[
 u^K_{\alpha\beta}u^K_{\beta\gamma}
 =u^K_{\alpha\gamma},
 \qquad
 u^C_{\alpha\beta}u^C_{\beta\gamma}
 =u^C_{\alpha\gamma}.
\]
The fixed source derivatives of the frames are bounded by the V11 Riesz/Kato
calculus; differentiating
\(u^K_{\alpha\beta}=(e^K_\alpha)^*e^K_\beta\) and similarly for \(C\)
gives the claimed transition bounds.
\end{proof}

We use the V11 determinant convention for the zero-mode factor,
\begin{equation}
 \boxed{
 \mathcal Z_D:=\det C\otimes(\det K)^*.}
 \label{eq:v12-zero-line}
\end{equation}
If one uses the opposite Fredholm-line convention, every formula below is
dualized and nothing substantive changes.

\begin{corollary}[Exact zero-mode line cocycle]
\label{cor:v12-zero-line-cocycle}
Let
\begin{equation}
 z_\alpha:=
 \det(e^C_\alpha)\otimes\det(e^K_\alpha)^*
 \label{eq:v12-local-z}
\end{equation}
be the local unit frame of \(\mathcal Z_D\).  Then
\begin{equation}
 z_\beta=g_{\alpha\beta}z_\alpha,
 \qquad
 \boxed{
 g_{\alpha\beta}
 =\frac{\det u^C_{\alpha\beta}}
        {\det u^K_{\alpha\beta}}
 \in U(1),}
 \label{eq:v12-zero-line-transition}
\end{equation}
and
\begin{equation}
 g_{\alpha\beta}g_{\beta\gamma}=g_{\alpha\gamma}.
 \label{eq:v12-zero-line-cocycle}
\end{equation}
Thus the V11 local saturation writers patch as sections of
\(\mathcal Z_D^*\), not as unrelated scalar functions.
\end{corollary}

\begin{proof}
Take top exterior powers in
\eqref{eq:v12-frame-transition}.  The determinant of the cokernel frame
contributes \(\det u^C_{\alpha\beta}\), while dualizing the kernel exterior
power contributes \((\det u^K_{\alpha\beta})^{-1}\).  The triple-overlap
identity follows by taking determinants of the vector-bundle cocycles.
\end{proof}

\subsection{The zero-mode Berry connection and its exact obstruction}
\label{sec:v12-zero-connection}

The local Kato frames also produce a canonical unitary connection on the
zero-mode line.  Define the imaginary-valued one-forms
\begin{equation}
 \mathcal A_\alpha^K
 :=\operatorname{Tr}\bigl((e_\alpha^K)^*\,d e_\alpha^K\bigr),
 \qquad
 \mathcal A_\alpha^C
 :=\operatorname{Tr}\bigl((e_\alpha^C)^*\,d e_\alpha^C\bigr),
 \label{eq:v12-Berry-forms}
\end{equation}
and
\begin{equation}
 \mathcal A_\alpha^Z:=\mathcal A_\alpha^C-\mathcal A_\alpha^K.
 \label{eq:v12-zero-connection}
\end{equation}

\begin{proposition}[Connection patching and curvature]
\label{prop:v12-zero-connection}
On overlaps,
\begin{equation}
 \boxed{
 \mathcal A_\beta^Z
 =\mathcal A_\alpha^Z+d\log g_{\alpha\beta}.}
 \label{eq:v12-connection-transform}
\end{equation}
Hence the local forms define a unitary connection \(\nabla^Z\) on
\(\mathcal Z_D\).  Its curvature is the globally defined two-form
\begin{equation}
 \boxed{
 \mathcal F_Z
 =\operatorname{Tr}(\Pi_C\,d\Pi_C\wedge d\Pi_C)
  -\operatorname{Tr}(\Pi_K\,d\Pi_K\wedge d\Pi_K),}
 \label{eq:v12-zero-curvature}
\end{equation}
with the usual harmless convention-dependent factor of \(i\) if one uses
real-valued connection forms.  In particular
\begin{equation}
 c_1(\mathcal Z_D)
 =\left[\frac{i}{2\pi}\mathcal F_Z\right]
 \in H^2(B;\mathbb Z).
 \label{eq:v12-c1-zero-line}
\end{equation}
\end{proposition}

\begin{proof}
For a unitary frame change \(e_\beta=e_\alpha u\),
\[
 \operatorname{Tr}(e_\beta^*de_\beta)
 =\operatorname{Tr}(e_\alpha^*de_\alpha)
  +\operatorname{Tr}(u^{-1}du)
 =\operatorname{Tr}(e_\alpha^*de_\alpha)+d\log\det u.
\]
Subtract the kernel formula from the cokernel formula and use
\eqref{eq:v12-zero-line-transition}.  The curvature formula is the standard
finite-rank projector identity, obtained by differentiating the orthogonal
projector \(ee^*\), using \(e^*e=I\), and taking the trace.  Chern--Weil theory
then gives \eqref{eq:v12-c1-zero-line}.
\end{proof}

\begin{theorem}[Global saturation criterion]
\label{thm:v12-global-saturation-criterion}
A nowhere-zero global complex saturation writer
\begin{equation}
 \omega_Z\in\Gamma(\mathcal Z_D^*)
 \label{eq:v12-global-writer}
\end{equation}
exists if and only if \(\mathcal Z_D\) is topologically trivial.  For a
paracompact background this is equivalent to
\begin{equation}
 c_1(\mathcal Z_D)=0.
 \label{eq:v12-line-trivial-c1}
\end{equation}
If, in addition, \(\nabla^Z\) is flat, a global nonzero \emph{parallel}
writer exists if and only if its holonomy character
\begin{equation}
 \operatorname{Hol}_Z:\pi_1(B)\longrightarrow U(1)
 \label{eq:v12-hol-character}
\end{equation}
is trivial.
\end{theorem}

\begin{proof}
A nonzero section of a complex line bundle trivializes it, and a trivial line
has an everywhere nonzero section.  Complex line bundles over a paracompact
base are classified by their first Chern class.  If the connection is flat,
parallel transport of one nonzero fibre vector is path independent exactly
when every loop has unit holonomy; this then produces a global parallel
section.  The converse is immediate.
\end{proof}

\begin{firewall}[Local anomaly cancellation need not flatten the zero-mode Berry connection by itself]
The complete chiral determinant line factors into the zero-mode line and the
reduced invertible line.  Vanishing of the curvature of the \emph{complete}
Standard-Model determinant connection therefore implies cancellation between
the two factor curvatures; it does not require the raw Kato connection
\(\nabla^Z\) on the zero-mode line to have zero curvature separately.  Only
the combined line-valued amplitude is physical before a factorized
normalization is chosen.
\end{firewall}

\subsection{The faithful-spin Standard Model supplies a global complex zero-mode writer}
\label{sec:v12-faithful-global-writer}

The inherited Grand Tensor result now closes the complex line-level problem on
the ordinary-spin faithful Standard-Model branch.  We formulate only the new
consequence for the V11 factorization.

Let \(\mathcal L_{\rm full}\to B\) denote the complete complex Standard-Model
chiral determinant line, and let \(\mathcal L_{\rm red}\to B\) be the reduced
nonzero-mode determinant line.  V11 gives a canonical line isomorphism
\begin{equation}
 \Phi:\mathcal L_{\rm full}
 \xrightarrow{\ \simeq\ }
 \mathcal Z_D\otimes\mathcal L_{\rm red}.
 \label{eq:v12-line-factorization}
\end{equation}
The reduced invertible map \(D^\perp\) has the nowhere-zero determinant section
\begin{equation}
 s_{\rm red}:=\det(D^\perp)
 \in\Gamma(\mathcal L_{\rm red}).
 \label{eq:v12-red-section}
\end{equation}

\begin{theorem}[Global complex zero-mode saturation on the faithful ordinary-spin branch]
\label{thm:v12-global-SM-writer}
Assume the inherited ordinary-spin \(G_6\) Standard-Model hypotheses under
which the complete determinant line has a unit parallel section
\begin{equation}
 s_{\rm full}\in\Gamma(\mathcal L_{\rm full}),
 \qquad
 \nabla^{\rm full}s_{\rm full}=0,
 \label{eq:v12-full-parallel}
\end{equation}
unique up to one constant phase on the connected component.  On any V11
minimal fixed-index branch with reduced floor \eqref{eq:v12-reduced-floor},
define
\begin{equation}
 \boxed{
 z_Z:=\Phi(s_{\rm full})\otimes s_{\rm red}^{-1}
 \in\Gamma(\mathcal Z_D).}
 \label{eq:v12-global-z}
\end{equation}
Then \(z_Z\) is smooth and nowhere zero.  Consequently
\(\mathcal Z_D\) is globally trivial and there is a canonical dual saturation
writer \(\omega_Z\), unique after fixing the constant phase of
\(s_{\rm full}\), such that
\begin{equation}
 \boxed{\langle\omega_Z,z_Z\rangle=1.}
 \label{eq:v12-global-pairing}
\end{equation}
All retained source derivatives patch globally.  Under exact shell elimination
and cutoff transport, \(z_Z\) and \(\omega_Z\) obey the same exact line cocycle
as the inherited complete determinant section; no additional complex
zero-mode holonomy remains.
\end{theorem}

\begin{proof}
The inherited faithful-spin theorem supplies the nowhere-zero parallel section
\(s_{\rm full}\).  The V11 reduced floor makes \(D^\perp\) invertible at every
point of \(B\), so its top exterior map is a nowhere-zero smooth section
\(s_{\rm red}\) of the reduced line.  The line factorization
\eqref{eq:v12-line-factorization} therefore makes
\eqref{eq:v12-global-z} a nowhere-zero smooth section of \(\mathcal Z_D\).
Its fibrewise dual normalized by \eqref{eq:v12-global-pairing} is smooth and
nowhere zero.

On a Kato cover, the local representatives of \(s_{\rm full}\),
\(s_{\rm red}\), and \(z_Z\) multiply by the corresponding transition
functions.  Differentiation therefore produces exactly the Cech/contact terms
already stored in V7 and V11; the transition factors cancel in the pairing
\eqref{eq:v12-global-pairing}.  The inherited complete line comparisons compose
exactly, while V11's reduced Schur determinants compose exactly.  Taking their
quotient shows that the zero-mode factors compose exactly as well.  Thus no
new loop character is introduced by the V11 splitting.
\end{proof}

\begin{corollary}[Factorized holonomy cancellation]
\label{cor:v12-factor-holonomy}
For any loop \(\gamma\subset B\), the factorized connections satisfy
\begin{equation}
 \operatorname{Hol}_{\rm full}(\gamma)
 =\operatorname{Hol}_{Z}(\gamma)
  \operatorname{Hol}_{\rm red}(\gamma).
 \label{eq:v12-hol-factor}
\end{equation}
On the inherited faithful ordinary-spin Standard-Model branch,
\(\operatorname{Hol}_{\rm full}(\gamma)=1\), and hence
\begin{equation}
 \boxed{
 \operatorname{Hol}_{Z}(\gamma)
 =\operatorname{Hol}_{\rm red}(\gamma)^{-1}.}
 \label{eq:v12-hol-cancel}
\end{equation}
The scalar saturated product is therefore single valued even when the two
factor connections separately have nontrivial holonomy.
\end{corollary}

\begin{proof}
Holonomy is multiplicative under tensor products of line bundles and inverse
under dualization.  Apply this to \eqref{eq:v12-line-factorization} and the
inherited unit holonomy of the complete determinant line.
\end{proof}

\subsection{Relative line transport between regulators is path independent}
\label{sec:v12-relative-line}

Now let \(t\in[0,1]\) parametrize a smooth native fixed-index regulator family
on \(B\).  Assume the V11 reduced floor and minimal nullities hold uniformly,
so the kernel/cokernel bundles extend over the cylinder
\begin{equation}
 \widetilde B:=B\times[0,1].
 \label{eq:v12-cylinder}
\end{equation}

\begin{theorem}[Cylinder transport of the zero-mode bundle and line]
\label{thm:v12-cylinder-transport}
The endpoint zero-mode bundles are restrictions of one bundle over
\(\widetilde B\).  Kato transport in the regulator parameter gives unitary
bundle isomorphisms
\begin{equation}
 J_K:K_0\xrightarrow{\simeq}K_1,
 \qquad
 J_C:C_0\xrightarrow{\simeq}C_1,
 \label{eq:v12-KC-transport}
\end{equation}
and therefore an induced unitary zero-line map
\begin{equation}
 J_Z:\mathcal Z_{D_0}\xrightarrow{\simeq}\mathcal Z_{D_1}.
 \label{eq:v12-Z-transport}
\end{equation}
On the faithful ordinary-spin Standard-Model branch, the corresponding
\emph{complete} complex determinant-line identification is independent of the
chosen fixed-index native regulator path with the same endpoints.  The induced
saturated scalar comparison is therefore path independent as well.
\end{theorem}

\begin{proof}
The Riesz projectors are smooth on \(\widetilde B\), hence define bundles there.
The Kato ODE in \(t\) gives \eqref{eq:v12-KC-transport}, and taking determinant
exterior powers gives \eqref{eq:v12-Z-transport}.  For two regulator paths with
the same endpoints, traverse one path and the reverse of the other.  This gives
a closed loop in the enlarged faithful-spin determinant family.  The inherited
eta/bordism theorem gives unit holonomy for the complete Standard-Model line on
every such paired loop.  Hence the two complete line identifications agree.
The reduced determinant section and the factorization theorem then force the
same equality for the factorized saturated amplitude.
\end{proof}

\subsection{The determinant line is not the complete fixed-index projector invariant}
\label{sec:v12-full-zero-bundle}

The previous theorem closes the complex amplitude line on the faithful branch,
but it does not classify the family inside the V11 minimal stratum.  The full
kernel or cokernel bundle is the relevant projector datum.

Let \(\mathcal M_k\) denote the finite-dimensional minimal index-\(k\) stratum
of V11.  Its bundle projection is
\begin{equation}
 \rho_k:\mathcal M_k\longrightarrow
 \begin{cases}
 \operatorname{Gr}(k,\mathcal H_-),&k>0,\\
 \operatorname{Gr}(-k,\mathcal H_+),&k<0,
 \end{cases}
 \label{eq:v12-rho-k}
\end{equation}
with affine vector-space fibres.

\begin{theorem}[Exact global projector classification inside the minimal stratum]
\label{thm:v12-minimal-classification}
For every compact CW parameter space \(B\), the projection \(\rho_k\) induces a
bijection of homotopy classes
\begin{equation}
 \boxed{
 [B,\mathcal M_k]
 \cong
 \begin{cases}
 [B,\operatorname{Gr}(k,\mathcal H_-)],&k>0,\\
 [B,\operatorname{Gr}(-k,\mathcal H_+)],&k<0.
 \end{cases}}
 \label{eq:v12-homotopy-classification}
\end{equation}
Thus two minimal fixed-index projector families are homotopic through the
minimal stratum if and only if their kernel-plane classifying maps
(respectively cokernel-plane maps) are homotopic.  Once such a zero-mode
homotopy is fixed, the affine nonzero-mode graph coordinates create no further
obstruction.
\end{theorem}

\begin{proof}
V11 proves that \(\mathcal M_k\) is the total space of a vector bundle over the
Grassmannian in \eqref{eq:v12-rho-k}.  Fibrewise multiplication by \(s\in[0,1]\)
is a strong deformation retraction onto the zero section.  Hence \(\rho_k\) is
a homotopy equivalence, which induces the claimed bijection on homotopy classes
from every CW source.  The final statement is the explicit fibre contraction.
\end{proof}

\begin{corollary}[The zero-mode line is only a shadow of the full bundle datum]
\label{cor:v12-line-shadow}
The determinant-line class
\begin{equation}
 [\mathcal Z_D]
 =\bigl[\det C\otimes(\det K)^*\bigr]
 \label{eq:v12-line-shadow}
\end{equation}
is a functor of the zero-mode bundle class, but it need not determine that
class.  Consequently equality or triviality of the global complex determinant
line is not, by itself, a projector-level connectedness theorem when
\(|k|>1\).
\end{corollary}

\begin{proof}
Taking the top exterior power forgets all characteristic data not seen by the
determinant.  For example, over \(S^4\) a rank-two complex bundle obtained from
a clutching map \(S^3\to SU(2)\) of degree one has
\(c_1=0\) and hence trivial determinant line, but has nonzero \(c_2\) and is
not isomorphic, hence not homotopic as a classifying map, to the trivial
rank-two bundle.  Theorem~\ref{thm:v12-minimal-classification} therefore
distinguishes the corresponding minimal-stratum families although their
determinant lines agree.
\end{proof}

This observation is load-bearing: the inherited faithful-spin eta theorem
trivializes the \emph{complex amplitude line}; it does not erase the full
zero-mode bundle from the microscopic phase-lifting problem.

\subsection{A finite-cutoff stable range for zero-mode classifying maps}
\label{sec:v12-stable-range}

The full Grassmannian obstruction simplifies when the ambient finite-cutoff
space is large compared with the dimension of the parameter panel.  This is
the regime relevant to a fixed finite background/source panel as the cutoff is
refined.

\begin{lemma}[Connectivity of the finite complex Grassmannian]
\label{lem:v12-Gr-connectivity}
Let
\begin{equation}
 i:\operatorname{Gr}(r,\mathbb C^m)
 \longrightarrow
 \operatorname{Gr}(r,\mathbb C^\infty)=BU(r)
 \label{eq:v12-Gr-inclusion}
\end{equation}
be the standard inclusion.  Then \(i\) is
\((2(m-r)+1)\)-connected.  Consequently, if \(B\) is a finite CW complex of
dimension
\begin{equation}
 \dim B\le 2(m-r),
 \label{eq:v12-stable-range-dim}
\end{equation}
then
\begin{equation}
 [B,\operatorname{Gr}(r,\mathbb C^m)]
 \longrightarrow [B,BU(r)]
 \label{eq:v12-Gr-map-bijection}
\end{equation}
is a bijection.
\end{lemma}

\begin{proof}
Use the Schubert CW decomposition.  The finite Grassmannian contains exactly
the cells indexed by partitions \(\lambda\) with at most \(r\) rows and
\(\lambda_1\le m-r\).  The first cell present in \(BU(r)\) but absent from the
finite Grassmannian corresponds to \(\lambda=(m-r+1)\), of complex dimension
\(m-r+1\) and hence real dimension \(2(m-r+1)\).  Therefore the relative CW
pair has no cells below that real dimension, so the inclusion is
\((2(m-r)+1)\)-connected.  Standard cellular obstruction theory then gives a
bijection on homotopy classes from any CW complex of dimension strictly below
the connectivity, which is exactly
\eqref{eq:v12-stable-range-dim}.
\end{proof}

\begin{theorem}[Bundle isomorphism implies projector homotopy in the stable ambient range]
\label{thm:v12-bundle-to-projector}
Suppose \(B\) has finite CW dimension \(d\).  For \(k>0\), let the negative
chiral ambient space have dimension \(m\) and assume
\begin{equation}
 d\le 2(m-k).
 \label{eq:v12-positive-stable-range}
\end{equation}
If two V11 minimal index-\(k\) projector families have isomorphic kernel bundles
\(K_0\simeq K_1\), then they are homotopic through the minimal stratum.  For
\(k<0\), the analogous statement holds for isomorphic cokernel bundles under
\(d\le2(m-|k|)\).
\end{theorem}

\begin{proof}
An isomorphism of rank-\(|k|\) complex bundles means that their classifying
maps become homotopic in \(BU(|k|)\).  Under the stable-range hypothesis,
Lemma~\ref{lem:v12-Gr-connectivity} says that the finite Grassmannian
classifying maps are already homotopic before stabilization.  Apply
Theorem~\ref{thm:v12-minimal-classification} to lift that zero-mode homotopy
through the affine minimal stratum.
\end{proof}

\begin{corollary}[Why the finite-cutoff ambient rank is not the main topological bottleneck]
\label{cor:v12-large-cutoff-stable}
For a fixed finite-dimensional parameter panel \(B\) and fixed index \(k\), the
stable-range inequality of Theorem~\ref{thm:v12-bundle-to-projector} holds at
all sufficiently large cutoffs whenever the ambient chiral dimension grows
with the regulator.  In that regime the projector-level obstruction is the
ordinary isomorphism class of the finite-rank zero-mode bundle.
\end{corollary}

\begin{proof}
The left-hand side \(d\) and \(|k|\) are fixed while \(m\to\infty\).  Hence
\(2(m-|k|)\ge d\) eventually.
\end{proof}

\subsection{Why an abstract zero-mode homotopy is not a native locality theorem}
\label{sec:v12-locality-obstruction}

The preceding topology still does not solve the microscopic problem.  A
zero-mode projector is typically a global operator in spacetime.  Rotating it
directly can destroy every cutoff-uniform locality moment even though its rank
is finite.

\begin{proposition}[Countertheorem: finite rank does not imply cutoff-uniform quasilocality]
\label{cth:v12-finite-rank-nonlocal}
Let
\(\Lambda_L=(\mathbb Z/L\mathbb Z)^d\) and
\(V=L^d\).  Let
\begin{equation}
 u_0(x)=V^{-1/2},
 \qquad
 (P_0f)(x)=u_0(x)\sum_y\overline{u_0(y)}f(y).
 \label{eq:v12-constant-projector}
\end{equation}
Then \(P_0\) has rank one and operator norm one, but for every \(s>0\) its
weighted row locality norm obeys
\begin{equation}
 \boxed{
 \sup_x\sum_{y\in\Lambda_L}
 (1+d_L(x,y))^s|P_0(x,y)|
 \ge c_{d,s}L^s.}
 \label{eq:v12-projector-nonlocal}
\end{equation}
Hence the family is not uniformly quasilocal as \(L\to\infty\).
Moreover, if
\(u_1(x)=V^{-1/2}e^{2\pi ix_1/L}\) and
\(u_t=\cos t\,u_0+\sin t\,u_1\), then the Kato generator of the rank-one
projector \(P_t=|u_t\rangle\langle u_t|\) has the same volume-growing weighted
moments.  Thus a smooth finite-rank zero-mode rotation does not by itself
produce a native local regulator homotopy.
\end{proposition}

\begin{proof}
The kernel is \(P_0(x,y)=V^{-1}\).  A fixed positive fraction of the lattice
sites lie at periodic distance at least \(cL\) from any chosen \(x\).  Therefore
\[
 V^{-1}\sum_y(1+d_L(x,y))^s
 \ge c' (1+cL)^s,
\]
which proves \eqref{eq:v12-projector-nonlocal}.  For \(P_t\),
\(\dot P_t=|\dot u_t\rangle\langle u_t|+|u_t\rangle\langle\dot u_t|\), and the
Kato generator \([\dot P_t,P_t]\) is a finite sum of rank-one kernels built
from \(u_0,u_1\).  Their pointwise magnitudes are \(O(V^{-1})\) on a positive
fraction of pairs, so the same weighted-moment lower bound follows (up to a
fixed constant for generic \(t\)).
\end{proof}

\begin{firewall}[Do not build the microscopic path by adding zero-mode projectors]
The V11 Riesz/Kato projectors are the correct variables for differentiating the
\emph{finite zero-mode source packet}.  Countertheorem
\ref{cth:v12-finite-rank-nonlocal} shows that they are generally the wrong
variables for constructing the microscopic ultraviolet regulator itself.
A native lift must be built inside the local Wilson-operator algebra; it cannot
be obtained by inserting a finite-rank projector correction and appealing only
to its operator norm.
\end{firewall}

\subsection{The correct native object: the protected local-operator component}
\label{sec:v12-native-component}

We therefore formulate the microscopic problem directly in a locality norm.
Fix an exponential weight \(\alpha>0\), the finite source-derivative order of
\(\mathfrak q\), and a compact background/source panel \(B\).  Let
\(\mathfrak A_{\alpha,\mathfrak q}\) be the real Banach affine space of
Hermitian, exactly gauge/spin-covariant Wilson-kernel families on \(B\) whose
weighted row and column norms and all required source derivatives are finite,
with the common physical principal Dirac--Yukawa jet frozen.  Finite-range
kernels are contained in this space.

Let \(\mathfrak P_{k,\mathfrak q}\subset\mathfrak A_{\alpha,\mathfrak q}\)
be the subset for which:
\begin{enumerate}[label=\textup{(P\arabic*)},leftmargin=3em]
\item the Wilson gap is positive uniformly on \(B\);
\item the full chiral/Yukawa family has minimal index \(k\) and a positive
reduced singular-value floor;
\item the declared complex zero-mode determinant data and, when present, real
Pfaffian orientation stay on the chosen protected component;
\item the V7 finite normalization Jacobian has the required constant rank.
\end{enumerate}

\begin{lemma}[The protected native set is open]
\label{lem:v12-native-open}
Inside the affine space with the physical principal jet frozen,
\(\mathfrak P_{k,\mathfrak q}\) is open in the
\(\mathfrak A_{\alpha,\mathfrak q}\) norm, after fixing one connected Wilson
rank/index component and one real Pfaffian orientation component when relevant.
\end{lemma}

\begin{proof}
A positive Wilson spectral gap is stable under an operator-norm perturbation
smaller than the gap.  On a fixed Wilson-rank component the chiral Fredholm
index is locally constant.  A positive smallest nonzero singular value is
stable under perturbations smaller than that value, so no accidental
kernel--cokernel pair is created.  The finite normalization Jacobian remains of
constant rank under a sufficiently small perturbation because its relevant
minor is nonzero.  Complex line data vary continuously; a real orientation is
locally constant away from a zero.  The weighted locality norm dominates the
operator and the finitely many source norms used in these tests, proving
openness.
\end{proof}

\begin{theorem}[Finite polygonal native lifting inside one protected operator component]
\label{thm:v12-polygonal-native}
Let \(A,B\in\mathfrak P_{k,\mathfrak q}\) lie in the same path-connected
component.  Then there exist finitely many kernels
\begin{equation}
 A=H_0,H_1,\ldots,H_M=B
 \label{eq:v12-polygonal-nodes}
\end{equation}
such that every straight segment
\begin{equation}
 H_{j,t}=(1-t)H_j+tH_{j+1},
 \qquad 0\le t\le1,
 \label{eq:v12-polygonal-segment}
\end{equation}
remains in \(\mathfrak P_{k,\mathfrak q}\).  Every segment is therefore a
locality-preserving V11 fixed-index native homotopy.  If the original kernels
are finite range and all nodes are chosen from one fixed finite-range affine
subspace, the polygonal chain is finite range; in the general case it is
uniformly exponentially local with the same weight \(\alpha\).
\end{theorem}

\begin{proof}
Choose a continuous path \(\gamma:[0,1]\to\mathfrak P_{k,\mathfrak q}\)
from \(A\) to \(B\).  By Lemma~\ref{lem:v12-native-open}, for every \(t\) there
is a norm ball \(B(\gamma(t),r_t)\) contained in the protected set.  Compactness
of \([0,1]\) and uniform continuity of \(\gamma\) give a finite subdivision
\(0=t_0<\cdots<t_M=1\) such that consecutive path points lie in one common
protected ball whose convex hull is still inside that ball.  Set
\(H_j=\gamma(t_j)\).  Norm balls in an affine Banach space are convex, so every
segment \eqref{eq:v12-polygonal-segment} remains protected.

Hermiticity, covariance and the frozen physical principal jet are affine linear
constraints and hence are preserved by the segment.  The exponential locality
norm is convex.  If the path lives in a fixed finite-range coefficient space,
linear interpolation does not enlarge that finite support.
\end{proof}

This theorem is the correct form of a global native lifting statement that can
be proved without classifying the path components of the local operator
algebra.  It is stronger than merely assuming a smooth projector path, because
its hypotheses are posed in the same locality topology used by the shell
estimates.

\subsection{Global fixed-index regulator universality on a native component}
\label{sec:v12-global-universality}

\begin{theorem}[Global complex fixed-index universality on the faithful-spin native component]
\label{thm:v12-global-index-universality}
Fix finite coefficient data \(\mathfrak q\), a compact connected ordinary-spin
faithful \(G_6\) Standard-Model background component \(B\), a fixed index
\(k\), and the positive IR/reference scale.  Let
\(A,B\in\mathfrak P_{k,\mathfrak q}\) lie in the same native path component.
Assume the inherited faithful-family hypotheses for the complete complex
Standard-Model determinant line.  Then:
\begin{enumerate}[label=\textup{(U\arabic*)},leftmargin=3em]
\item the zero-mode complex line admits the global saturation writer of
Theorem~\ref{thm:v12-global-SM-writer};
\item a finite polygonal native chain joining \(A\) to \(B\) exists by
Theorem~\ref{thm:v12-polygonal-native};
\item the complete complex determinant-line comparison along the chain is
path independent and has trivial paired loop holonomy;
\item the V11 relative hard-shell theorem applies on every segment and the
segment matching maps compose to an invertible continuum map
\begin{equation}
 F_{\infty,\mathfrak q,k}^{B\leftarrow A};
 \label{eq:v12-global-map}
\end{equation}
\item the globally saturated coefficientwise continua satisfy
\begin{equation}
 \boxed{
 \widehat{\mathfrak C}_{\infty,\lambda}^{B,(k),\rm sat}
 =F_{\infty,\mathfrak q,k}^{B\leftarrow A}
  \widehat{\mathfrak C}_{\infty,\lambda}^{A,(k),\rm sat}.}
 \label{eq:v12-global-universality}
\end{equation}
\end{enumerate}
The equivalence contains all retained connected, composite, stress/contact,
ghost/antifield, BV/BRST/master and reduced determinant-line coefficients and
is strictly projective in perturbative order.
\end{theorem}

\begin{proof}
Item (U1) is Theorem~\ref{thm:v12-global-SM-writer}; item (U2) is
Theorem~\ref{thm:v12-polygonal-native}.  Every polygonal segment lies in the
V11 admissible fixed-index region and hence carries the reduced
\(aM+\mu/M\) relative-forest estimate and the optimized
\(O(\sqrt{a\mu})\) regulator difference.  The inherited faithful-spin
parallel determinant section and Theorem~\ref{thm:v12-cylinder-transport}
make the complex line identifications independent of the chain and compatible
with the global saturation writer.  Apply V11 local universality to each
segment and compose the finitely many invertible matching maps.  The V7
chronological construction makes each segment map projective in terminal loop
order, hence so is their finite composition.
\end{proof}

\begin{corollary}[The complex line is no longer a load-bearing obstruction on this branch]
\label{cor:v12-complex-line-closed}
Within the hypotheses of Theorem~\ref{thm:v12-global-index-universality}, the
remaining classification problem for general fixed-index regulator
universality is not the complex determinant-line holonomy.  It is whether two
microscopic kernels lie in the same protected path component of the native
local-operator space.  Projector-level zero-mode bundle topology supplies a
necessary obstruction to such a path, but not by itself a locality-preserving
construction.
\end{corollary}

\subsection{The real Pfaffian branch retains a mod-two obstruction}
\label{sec:v12-real-Pfaffian}

The preceding faithful-spin theorem concerns the complete complex chiral
determinant line.  A genuinely real Majorana/Pfaffian sector has a different
global obstruction which must not be erased by complexification.

\begin{proposition}[Real Pfaffian orientation local system]
\label{prop:v12-Pfaffian-w1}
Let \(\mathcal P\to B\) be a real one-dimensional Pfaffian line on a
constant-rank branch.  Its local orientation frames have transition signs
\(\epsilon_{\alpha\beta}\in\{\pm1\}\), defining
\begin{equation}
 w_1(\mathcal P)\in H^1(B;\mathbb Z_2).
 \label{eq:v12-Pf-w1}
\end{equation}
A global real orientation exists if and only if
\(w_1(\mathcal P)=0\).  For a generic loop of real skew operators with only
simple crossings, the orientation holonomy is
\begin{equation}
 \boxed{
 \operatorname{Hol}_{\mathcal P}(\gamma)
 =(-1)^{\operatorname{SF}_2(\gamma)},}
 \label{eq:v12-Pf-sf2}
\end{equation}
where \(\operatorname{SF}_2\) is the parity of the real spectral flow.  The
formula extends by homotopy to arbitrary loops avoiding an endpoint rank
jump.
\end{proposition}

\begin{proof}
The transition signs of a real line form a Cech \(\mathbb Z_2\)-cocycle whose
class is the first Stiefel--Whitney class.  It is a coboundary exactly when the
local orientations can be rephased to agree globally.  At a simple real
crossing one Pfaffian eigenvalue changes sign, reversing the local orientation.
Along a generic loop the total sign is therefore the parity of the number of
crossings.  Pair creation/annihilation under a homotopy changes that number by
an even integer, so the parity extends to the homotopy-invariant mod-two
spectral flow.
\end{proof}

\begin{firewall}[The faithful complex eta theorem does not orient an independent real Pfaffian]
The inherited vanishing of the faithful-spin complex determinant holonomy
cannot be used to set \(w_1(\mathcal P)=0\) for an independently physical real
Majorana line.  Global scalar universality in such a sector requires either a
separate proof that the mod-two character vanishes or a declared physical
orientation local system which is transported unchanged between regulators.
\end{firewall}

\subsection{What is and is not closed by V12}
\label{sec:v12-boundary}

The global picture can now be stated sharply.

\begin{theorem}[Three-level global obstruction theorem]
\label{thm:v12-three-level}
On a V11 minimal fixed-index family the following implications hold.
\begin{enumerate}[label=\textup{(G\arabic*)},leftmargin=3em]
\item A native local-operator homotopy implies a minimal-stratum projector
homotopy and hence equality of the kernel/cokernel classifying-map class.
\item A minimal-stratum projector homotopy implies isomorphism of the
zero-mode determinant lines and equality of all their characteristic and
holonomy data transported by that homotopy.
\item Neither converse holds as a formal consequence of the preceding level:
the determinant line forgets higher zero-mode-bundle data, and an abstract
projector homotopy need not have cutoff-uniform locality.
\end{enumerate}
On the faithful ordinary-spin complex Standard-Model branch the line-level
obstruction is trivialized by the inherited eta/bordism theorem, but the two
stronger levels remain distinct.
\end{theorem}

\begin{proof}
Item (G1) follows by applying spectral functional calculus continuously along
a gapped native Wilson path.  Item (G2) follows by taking top exterior powers
of the kernel/cokernel bundle over the homotopy cylinder.  Failure of the first
converse is Corollary~\ref{cor:v12-line-shadow}; failure of the second as a
proof principle is Countertheorem~\ref{cth:v12-finite-rank-nonlocal}.  The last
sentence is Theorem~\ref{thm:v12-global-SM-writer}.
\end{proof}

\begin{center}
\fbox{\parbox{0.92\textwidth}{
\textbf{V12 closes global complex line transport, but it does not identify
projector equivalence with microscopic locality.}
For the faithful ordinary-spin Standard-Model determinant packet, the
complete complex line has trivial paired holonomy and V11's reduced nonzero
section therefore produces a global zero-mode saturation writer.  Global
fixed-index universality follows for every pair of regulators in the same
protected path component of the native exponentially local operator space.
The remaining load-bearing problem is to characterize those native path
components from intrinsic local data rather than from an already-given native
path.}}
\end{center}

\section{V12 anti-circularity audit}
\label{sec:v12-audit}

\begin{enumerate}[leftmargin=2.4em]
\item \textbf{The Grand Tensor eta theorem is inherited, not rederived.}
V12 uses its unit parallel complex determinant section and exact cutoff cocycle
only to obtain the new zero-mode saturation consequence after V11
factorization.
\item \textbf{The raw Kato zero-mode connection is not silently declared
flat.}
Its curvature is \eqref{eq:v12-zero-curvature}; only the complete determinant
connection has the inherited anomaly/eta trivialization.
\item \textbf{Determinant-line triviality is not promoted to full bundle
triviality.}
The rank-two \(S^4\) example in Corollary~\ref{cor:v12-line-shadow} keeps the
higher zero-mode-bundle obstruction explicit.
\item \textbf{Stable bundle classification is used only in its stated finite
ambient range.}
Outside \eqref{eq:v12-stable-range-dim}, V12 retains the actual finite
Grassmannian classifying map as the projector invariant.
\item \textbf{A finite-rank zero-mode projector is not called local because its
operator norm is bounded.}
Countertheorem~\ref{cth:v12-finite-rank-nonlocal} exhibits the missing
weighted-locality control explicitly.
\item \textbf{The native lift is constructed only inside an already identified
protected local-operator component.}
The polygonal theorem is a Banach-space statement in the native locality norm,
not a reconstruction from an abstract Grassmannian homotopy.
\item \textbf{The real Pfaffian is kept separate.}
Its \(\mathbb Z_2\) orientation character is not inferred from the complex
faithful-spin determinant theorem.
\end{enumerate}

\section{V12 status ledger}
\label{sec:v12-status}

\subsection*{Proved in V12}
\begin{enumerate}[leftmargin=2.2em]
\item The local V11 kernel/cokernel Kato frames patch to global smooth bundles
on every protected constant-rank component.  Their transition matrices have
all required fixed source derivatives.
\item The zero-mode determinant factor
\(\mathcal Z_D=\det C\otimes(\det K)^*\) has the exact unitary Cech cocycle
\(g_{\alpha\beta}=\det u^C_{\alpha\beta}/\det u^K_{\alpha\beta}\), a global
Kato/Berry connection and curvature
\(\operatorname{Tr}(\Pi_Cd\Pi_C\wedge d\Pi_C)-
  \operatorname{Tr}(\Pi_Kd\Pi_K\wedge d\Pi_K)\).
\item On the inherited faithful ordinary-spin \(G_6\) Standard-Model branch,
the complete parallel determinant section divided by the global nonzero
reduced determinant section gives a nowhere-zero global zero-mode section and
hence a global saturation writer.  The factorized holonomies cancel exactly.
\item Along a smooth fixed-index native regulator cylinder, Kato transport
gives global kernel/cokernel and zero-line endpoint isomorphisms.  On the
faithful complex branch the complete line identification and saturated scalar
comparison are path independent.
\item The homotopy class of a family inside the V11 minimal stratum is exactly
the homotopy class of its kernel-plane map for \(k>0\) or cokernel-plane map
for \(k<0\).  The determinant line alone is not complete when higher
zero-mode-bundle data are present.
\item The inclusion
\(\operatorname{Gr}(r,\mathbb C^m)\to BU(r)\) is
\((2(m-r)+1)\)-connected.  Hence on a finite parameter panel of dimension
\(d\le2(m-r)\), isomorphic zero-mode bundles already give homotopic finite
Grassmannian classifying maps and therefore a minimal-stratum projector
homotopy.
\item Finite rank does not imply native locality: the constant-mode projector
has weighted row moments growing like \(L^s\).  Thus an abstract zero-mode
Kato rotation is not a valid microscopic regulator lift without an independent
locality theorem.
\item The protected fixed-index native operator set is open in the weighted
exponentially local/source norm.  Any two regulators in one path component are
connected by a finite polygonal chain of native protected homotopies.
\item Combining that polygonal chain with V11 and the inherited faithful-spin
line theorem yields global complex fixed-index regulator universality on every
protected native path component, with exact loop-order projectivity.
\item A real Pfaffian line carries the independent obstruction
\(w_1\in H^1(B;\mathbb Z_2)\), with loop sign given by mod-two spectral flow.
\end{enumerate}

\subsection*{Inherited}
\begin{enumerate}[leftmargin=2.2em]
\item The Grand Tensor finite zero-mode factorization/saturation theorem,
Fredholm determinant line, faithful-spin \(G_6\) eta/bordism trivialization,
and exact smooth cutoff line cocycle.
\item V1--V7: the source-complete fixed-order continuum, common BV restoration,
forest theorem, UV Cauchy convergence, scheme/reference covariance and exact
loop-order projectivity.
\item V8--V10: regulator-marked relative forests, the two-scale universality
estimate, and native regulator-class lifting on the protected index-zero
single-cone component.
\item V11: minimal fixed-index geometry, reduced pseudoinverse/source calculus,
zero-mode-preserving shell factorization, saturation writers and local
fixed-index universality.
\end{enumerate}

\subsection*{Conditional}
\begin{enumerate}[leftmargin=2.2em]
\item The global complex saturation theorem uses the inherited ordinary-spin
faithful \(G_6\) family hypotheses.  Other spin--gauge structures require their
own eta/holonomy calculation.
\item The finite-Grassmannian stable-range theorem assumes a finite-dimensional
CW parameter panel.  It is not a classification theorem for an unrestricted
infinite-dimensional gauge configuration space.
\item Global regulator-class universality is proved for regulators already in
the same path component of the protected native local-operator space.  V12 does
not prove that every projector-level zero-mode-bundle equivalence lifts to that
native component.
\item A real Pfaffian sector still requires a vanishing mod-two orientation
character or an explicitly supplied physical orientation local system.
\item All continuum claims remain coefficientwise at prescribed finite
\((L,\Delta,r,N)\) and positive IR/reference scale.
\end{enumerate}

\subsection*{Open beyond V12}
\begin{enumerate}[leftmargin=2.2em]
\item Classify the path components of the protected exponentially local Wilson
operator space relative to the fixed physical Dirac--Yukawa principal jet.
Equivalently, determine when the zero-mode-bundle/projector homotopy produced by
V11--V12 admits a uniformly local native lift without inserting nonlocal
finite-rank projectors.
\item For a genuinely real Majorana/Pfaffian sector, compute the mod-two
spectral-flow character on the relevant global background component and prove
(or falsify) orientability.
\item Analyze crossings out of the minimal fixed-index stratum, including
creation of accidental zero-mode pairs and true index change, with divisor and
spectral-flow source packets.
\item The independent stronger problems remain: removal of the positive IR
scale, analytic finite-coupling RG trajectories, perturbative summability and
nonperturbative completion.
\end{enumerate}

\subsection*{Exact next load-bearing theorem}
\begin{center}
\fbox{\parbox{0.92\textwidth}{
\textbf{Local-operator phase classification / native lifting.}
Fix the physical Standard-Model principal Dirac--Yukawa jet, index \(k\), the
minimal reduced-gap condition and the global zero-mode bundle class.  Determine
the path components of the corresponding gapped finite-range or uniformly
exponentially local, gauge/spin-covariant Wilson-operator space.  Prove that
projector-equivalent regulators lie in one protected native component, or
exhibit an additional invariant of the local operator algebra.  The proof must
avoid finite-rank zero-mode projector corrections, whose weighted locality
fails by Countertheorem~\ref{cth:v12-finite-rank-nonlocal}.  Combined with
Theorem~\ref{thm:v12-global-index-universality}, such a result would upgrade
V12 from universality within a native component to a complete global
fixed-index regulator-class theorem.}}
\end{center}

\section*{Append-only continuation marker: V12}
\addcontentsline{toc}{section}{Append-only continuation marker: V12}
\textbf{End of V12.}  The complete V1--V11 body above is retained verbatim.
V12 patches the V11 zero-mode bundles globally, identifies the exact Cech and
Berry data of their determinant line, and uses the inherited faithful-spin
Standard-Model eta/bordism trivialization together with the global reduced
determinant section to construct a global complex zero-mode saturation writer.
It proves that the full kernel/cokernel classifying map, not merely its
determinant line, is the projector-level fixed-index invariant; in a finite
stable ambient range, bundle isomorphism already gives projector homotopy.  It
then isolates the remaining microscopic obstruction by showing that finite-rank
zero-mode rotations need not be uniformly quasilocal.  Finally, it proves a
finite polygonal native lifting theorem inside every protected path component
of the exponentially local operator space and obtains global complex
fixed-index regulator universality on that component.  The next gate is the
intrinsic classification of those native local-operator components.



\clearpage
\section{V13 continuation: local Riesz projectors, controlled zero-mode modules, and stable native lifting}
\label{sec:v13-continuation}

V12 isolated the remaining fixed-index problem as a locality problem rather
than a determinant-line problem.  Its counterexample showed correctly that an
\emph{arbitrary} finite-rank projector can be completely nonlocal even when its
operator norm is one.  For the actual V11 protected zero modes, however, there
is one additional hypothesis which was not used in that counterexample: the
zero eigenspace is the isolated spectral subspace of a uniformly local operator
whose remaining singular spectrum has a cutoff-uniform floor.  Going upstream
to that spectral gap changes the local classification problem.

The first result below is therefore a correction of variable, not a
retraction of V12: Riesz projectors of the protected chiral operator are
uniformly exponentially local (with a possibly smaller exponential rate).
Hence Kato transport along an already native fixed-index path is itself local.
The genuine obstruction is whether the endpoint zero-mode bundles admit a
\emph{controlled} local identification.  Ordinary bundle isomorphism and
complex determinant-line triviality forget this metric/locality datum.

V13 proves four statements.
\begin{enumerate}[label=\textup{(V13.\arabic*)},leftmargin=3.8em]
\item A gap plus exponential locality implies exponential locality of the
resolvent, sign, zero-mode Riesz projectors, Moore--Penrose inverse, and every
fixed source derivative, after a harmless loss of exponent.
\item The V11 minimal fixed-index geometry can be reconstructed entirely from
its local zero-mode projector and a local affine graph coordinate.  Thus the
full native phase problem reduces to a finite-rank \emph{controlled zero-mode
module} problem.
\item Ordinary zero-mode bundle data are not complete for this controlled
problem: a pair of macroscopically separated site zero modes has the same rank,
trivial determinant line and uniform reduced gap, but any intertwining partial
isometry has exponentially growing locality norm.
\item If the two zero-mode packets are controlled-isomorphic, then after
adjoining one inert massive spectator copy there is an explicit protected
local homotopy between the two regulators.  The spectator contributes only a
source-independent factor and disappears exactly from the normalized connected
functional.  Therefore controlled-isomorphic regulators are \emph{stably}
regulator-universal at every fixed coefficient order.
\end{enumerate}

The remaining issue after V13 is an unstable cancellation problem: whether the
spectator can always be removed in the exact Standard-Model carrier, or whether
there is a genuine controlled/coarse obstruction to such cancellation.

\subsection{The weighted local kernel algebra}
\label{sec:v13-local-algebra}

Let $\Lambda$ be one of the finite periodic lattices in the regulator sequence,
with periodic distance $d(\cdot,\cdot)$ and uniformly polynomial volume growth.
For a matrix kernel $A=(A_{xy})_{x,y\in\Lambda}$ put
\begin{equation}
 \|A\|_{\alpha}
 :=\max\left\{
 \sup_x\sum_y e^{\alpha d(x,y)}\|A_{xy}\|,
 \sup_y\sum_x e^{\alpha d(x,y)}\|A_{xy}\|
 \right\}.
 \label{eq:v13-local-norm}
\end{equation}
For a compact background/source panel $B$, the same notation means the
supremum over $b\in B$ together with the finitely many source derivatives
required by the fixed coefficient packet $\mathfrak q$.  Exact gauge/spin
covariance and the fixed representation grading are imposed as closed linear
constraints.

\begin{lemma}[Weighted Schur algebra]
\label{lem:v13-Schur-algebra}
For every $\alpha>0$, kernels with finite norm
\eqref{eq:v13-local-norm} form a Banach $*$-algebra and
\begin{equation}
 \boxed{
 \|AB\|_{\alpha}\le \|A\|_{\alpha}\|B\|_{\alpha},
 \qquad
 \|A^*\|_{\alpha}=\|A\|_{\alpha}.}
 \label{eq:v13-Schur-product}
\end{equation}
Moreover $\|A\|_{\ell^2\to\ell^2}\le \|A\|_{\alpha}$.
The same statements hold for the finite source-jet norm.
\end{lemma}

\begin{proof}
The triangle inequality gives
$d(x,z)\le d(x,y)+d(y,z)$ and hence
$e^{\alpha d(x,z)}\le e^{\alpha d(x,y)}e^{\alpha d(y,z)}$.
Insert the kernel product
$(AB)_{xz}=\sum_yA_{xy}B_{yz}$ and sum first in $z$ and then in $y$.
The column estimate is identical.  Taking adjoints exchanges the two Schur
rows.  The ordinary Schur test gives the $\ell^2$ operator bound.  A fixed
source derivative is a finite Leibniz sum, so the same argument applies after
summing the finitely many derivative seminorms.
\end{proof}

The fixed exponent is not preserved by every inverse with the same numerical
value.  What is preserved under a uniform spectral gap is \emph{some} positive
exponent.  This is precisely what the shell proof needs.

\begin{theorem}[Gap-to-resolvent locality with exponent loss]
\label{thm:v13-resolvent-locality}
Let $A=A^*$ obey
\begin{equation}
 \|A\|_{\alpha}\le M,
 \qquad
 \operatorname{dist}(z,\operatorname{spec}A)\ge\delta>0
 \label{eq:v13-resolvent-hyp}
\end{equation}
with constants uniform in cutoff, volume and the chosen compact source panel.
Then there exists
\begin{equation}
 0<\beta=\beta(\alpha,M,\delta)<\alpha
 \label{eq:v13-beta}
\end{equation}
for which
\begin{equation}
 \boxed{
 \|(z-A)^{-1}\|_{\beta}
 \le C_{\alpha,M,\delta}.}
 \label{eq:v13-resolvent-local}
\end{equation}
The same conclusion holds for every prescribed fixed source derivative of the
resolvent if the corresponding source jets of $A$ are bounded in
$\|\cdot\|_\alpha$.
\end{theorem}

\begin{proof}
Fix a lattice point $x_0$ and let
$\phi_{x_0}(x)=d(x,x_0)$.  For $\eta>0$ define the diagonal multiplication
$W_{\eta,x_0}f(x)=e^{\eta\phi_{x_0}(x)}f(x)$.  The conjugated kernel is
\[
 (W_{\eta,x_0}AW_{\eta,x_0}^{-1})_{xy}
 =e^{\eta(\phi_{x_0}(x)-\phi_{x_0}(y))}A_{xy}.
\]
Since $|\phi_{x_0}(x)-\phi_{x_0}(y)|\le d(x,y)$, the Schur norm of the
difference from $A$ is bounded by
\[
 \rho_A(\eta):=
 \max\left\{
 \sup_x\sum_y (e^{\eta d(x,y)}-1)\|A_{xy}\|,
 \sup_y\sum_x (e^{\eta d(x,y)}-1)\|A_{xy}\|
 \right\}.
\]
Uniform exponential summability at rate $\alpha$ implies
$\rho_A(\eta)\to0$ as $\eta\downarrow0$, uniformly on the declared family:
split the sum at a fixed distance $R$, use ordinary dominated convergence on
the finite head, and use the $e^{\alpha d}$ tail beyond $R$.  Choose
$0<\eta<\alpha$ so that $\rho_A(\eta)<\delta/2$.
Then
\[
 \operatorname{dist}\bigl(z,
 \operatorname{spec}(W_{\eta,x_0}AW_{\eta,x_0}^{-1})\bigr)
 \ge\delta/2
\]
in the resolvent sense supplied by the Neumann identity, and therefore
\begin{equation}
 \|W_{\eta,x_0}(z-A)^{-1}W_{\eta,x_0}^{-1}\|_{2\to2}
 \le 2/\delta.
 \label{eq:v13-CT-l2}
\end{equation}
Applying \eqref{eq:v13-CT-l2} to a coordinate vector at $x_0$ gives a uniform
exponentially weighted $\ell^2$ column bound for the resolvent.  Apply the same
argument to the adjoint for rows.  Choose $0<\beta<\eta$.  Polynomial lattice
volume growth gives
\[
 \sup_x\sum_y e^{-2(\eta-\beta)d(x,y)}<\infty
\]
uniformly in the finite periodic volumes.  Cauchy--Schwarz therefore converts
the weighted $\ell^2$ row and column bounds into
\eqref{eq:v13-resolvent-local}.

For source derivatives use
\[
 \partial(z-A)^{-1}
 =(z-A)^{-1}(\partial A)(z-A)^{-1}
\]
and iterate.  Every fixed derivative is a finite sum of resolvent words.
Apply Lemma~\ref{lem:v13-Schur-algebra}, decreasing the exponent once if
necessary to a common positive value for the finite family of words.
\end{proof}

\begin{corollary}[Local holomorphic/Riesz calculus on a protected spectral split]
\label{cor:v13-local-functional-calculus}
Let $A=A^*$ be exponentially local and suppose its spectrum is separated into
two closed sets by a uniform gap.  Every Riesz projector defined by a contour
inside that gap is exponentially local at some uniform positive rate, together
with every prescribed fixed source derivative.  The same is true of
$\operatorname{sgn}A$ when zero is uniformly excluded from the spectrum.
\end{corollary}

\begin{proof}
Use the contour representation
\[
 P_\Gamma={1\over2\pi i}\oint_\Gamma(z-A)^{-1}\,dz
\]
and Theorem~\ref{thm:v13-resolvent-locality}.  A fixed compact contour permits
one common exponent and constant.  The sign is a difference of the positive
and negative Riesz projectors.
\end{proof}

\subsection{Protected zero-mode projectors are local}
\label{sec:v13-zero-local}

We now apply the preceding theorem to the actual V11 zero modes.  This is the
load-bearing distinction from an arbitrary finite-rank projector.

\begin{theorem}[Uniform locality of V11 kernel/cokernel projectors]
\label{thm:v13-zero-projector-local}
Let $D_b$ be a V11 protected constant-rank chiral/Yukawa family.  Assume its
kernel is generated by a native finite-range or exponentially local operator
family and its nonzero singular values obey
\begin{equation}
 s_{\rm nonzero}(D_b)\ge\tau_*>0.
 \label{eq:v13-tau}
\end{equation}
Then there is a cutoff-uniform exponent $\beta>0$ such that
\begin{equation}
 \boxed{
 \|\Pi_K\|_\beta+\|\Pi_C\|_\beta
 \le C,}
 \label{eq:v13-zero-projector-local}
\end{equation}
and the same estimate holds for every fixed source derivative required by
$\mathfrak q$.  The Moore--Penrose inverse $D^\dagger$ is exponentially local
at a possibly smaller positive exponent.
\end{theorem}

\begin{proof}
The products $D^*D$ and $DD^*$ are exponentially local by
Lemma~\ref{lem:v13-Schur-algebra}.  Their spectra are contained in
$\{0\}\cup[\tau_*^2,\infty)$.  Choose a circle around zero of radius
$\tau_*^2/2$.  The V11 Riesz formulas for $\Pi_K$ and $\Pi_C$ and
Corollary~\ref{cor:v13-local-functional-calculus} prove
\eqref{eq:v13-zero-projector-local}, including source derivatives.

For the pseudoinverse use the V11 formula
\[
 D^\dagger=(D^*D+\Pi_K)^{-1}D^*.
\]
The positive operator $D^*D+\Pi_K$ has spectral floor
$\min\{1,\tau_*^2\}$ and is exponentially local.  Apply
Theorem~\ref{thm:v13-resolvent-locality} at $z=0$ and multiply by $D^*$.
\end{proof}

\begin{firewall}[The exact scope of the V12 nonlocal projector counterexample]
\label{assess:v13-v12-counterexample}
Countertheorem~\ref{cth:v12-finite-rank-nonlocal} remains correct: finite rank
alone does not imply locality.  The constant-mode projector used there is not
the isolated kernel projector of a uniformly local family with a
cutoff-uniform reduced spectral floor.  Theorem
\ref{thm:v13-zero-projector-local} shows that the V11 protected zero-mode
packet has additional analytic structure which rules out that particular
failure mode.  It does \emph{not} imply that two endpoint zero-mode projectors
are connected by a uniformly local path; that is the separate controlled
module question treated below.
\end{firewall}

\subsection{Kato transport along a native path is itself native-local}
\label{sec:v13-local-Kato}

\begin{theorem}[Local Kato transport]
\label{thm:v13-local-Kato}
Let $t\mapsto D_t$ be a $C^1$ V11 fixed-index native homotopy on a compact
parameter interval with one common locality exponent at the microscopic level
and one common reduced singular-value floor $\tau_*>0$.  Then for some
$\beta>0$ the zero-mode projectors and their $t$ derivatives satisfy
\begin{equation}
 \sup_t\bigl(
 \|\Pi_t\|_\beta+\|\dot\Pi_t\|_\beta
 \bigr)<\infty.
 \label{eq:v13-Pi-dot-local}
\end{equation}
The Kato equation
\begin{equation}
 \dot U_t=[\dot\Pi_t,\Pi_t]U_t,
 \qquad U_0=I,
 \label{eq:v13-local-Kato-ODE}
\end{equation}
has a unitary solution with
\begin{equation}
 \boxed{
 \sup_t\bigl(
 \|U_t\|_\beta+\|U_t^{-1}\|_\beta
 \bigr)<\infty,
 \qquad
 \Pi_t=U_t\Pi_0U_t^*.}
 \label{eq:v13-local-Kato-result}
\end{equation}
All prescribed fixed source derivatives obey the same statement after one
finite exponent loss.
\end{theorem}

\begin{proof}
Theorem~\ref{thm:v13-zero-projector-local} applied to the two-parameter family
$(t,b)$ gives \eqref{eq:v13-Pi-dot-local}; equivalently one differentiates the
Riesz formula in $t$.  The generator
$G_t=[\dot\Pi_t,\Pi_t]$ is anti-Hermitian and bounded in the Banach algebra by
Lemma~\ref{lem:v13-Schur-algebra}.  The Dyson series for
\eqref{eq:v13-local-Kato-ODE} therefore converges in the same Banach algebra,
with
\[
 \|U_t\|_\beta\le
 \exp\left(\int_0^t\|G_s\|_\beta\,ds\right).
\]
Anti-Hermiticity gives unitarity on $\ell^2$.  The standard Kato computation
shows $U_t\Pi_0U_t^*=\Pi_t$.  Source derivatives satisfy finite inhomogeneous
Dyson equations with coefficients already controlled in the same algebra.
\end{proof}

Thus a native homotopy produces more than an abstract bundle isomorphism: it
produces an exponentially local one.  This motivates the correct invariant.

\subsection{Controlled zero-mode modules}
\label{sec:v13-controlled-module}

Let $\mathcal A_{\exp,\mathfrak q}$ denote the union over positive exponents of
the preceding gauge/spin-covariant weighted algebras, with the fixed source
jet package understood.  The union is used only so that the harmless exponent
losses of Theorem~\ref{thm:v13-resolvent-locality} do not create artificial
phase boundaries.

\begin{definition}[Controlled local equivalence of zero-mode packets]
\label{def:v13-controlled-equivalence}
Two finite-rank zero-mode projector families $\Pi_0,\Pi_1$ on the same compact
background panel are \emph{controlled-equivalent}, written
\begin{equation}
 \Pi_0\sim_{\rm ctl}\Pi_1,
 \label{eq:v13-ctl-equivalence}
\end{equation}
if there is a partial isometry
$V\in\mathcal A_{\exp,\mathfrak q}$ such that
\begin{equation}
 \boxed{
 V^*V=\Pi_0,
 \qquad
 VV^*=\Pi_1,}
 \label{eq:v13-controlled-partial-isometry}
\end{equation}
with all retained background/source derivatives in one common positive
locality class.  The isomorphism class of
$\Pi\mathcal A_{\exp,\mathfrak q}$ is called the
\emph{controlled zero-mode module class}.
\end{definition}

\begin{proposition}[Native homotopy preserves the controlled module class]
\label{prop:v13-native-implies-controlled}
Endpoints of every V11 fixed-index native homotopy have
controlled-equivalent kernel bundles for $k>0$ and controlled-equivalent
cokernel bundles for $k<0$.
\end{proposition}

\begin{proof}
Use the local Kato unitary $U_1$ of
Theorem~\ref{thm:v13-local-Kato}.  For the kernel case set
$V=U_1\Pi_{K,0}$.  Then
$V^*V=\Pi_{K,0}$ and
$VV^*=\Pi_{K,1}$.  The cokernel case is identical.
\end{proof}

There are natural forgetful maps
\begin{equation}
 \boxed{
 \text{controlled module class}
 \longrightarrow
 \text{ordinary zero-mode bundle class}
 \longrightarrow
 \text{determinant/Pfaffian line class}.}
 \label{eq:v13-forgetful-hierarchy}
\end{equation}
V12 proved that the second arrow loses information.  The next proposition shows
that the first arrow can also lose information even when the reduced gap is
uniform.

\begin{proposition}[Macroscopic transport obstruction invisible to the ordinary bundle]
\label{prop:v13-macro-obstruction}
Let $\Lambda_L$ be an even one-dimensional periodic lattice and choose sites
$x_L=0$ and $y_L=L/2$.  Put
\begin{equation}
 \Pi_{0,L}=|\delta_{x_L}\rangle\langle\delta_{x_L}|,
 \qquad
 \Pi_{1,L}=|\delta_{y_L}\rangle\langle\delta_{y_L}|.
 \label{eq:v13-site-projectors}
\end{equation}
Both are range-zero local rank-one projections with trivial complex determinant
line.  They are also exact kernel projectors of the range-zero local operators
\begin{equation}
 D_{i,L}:=I-\Pi_{i,L},
 \label{eq:v13-site-D}
\end{equation}
whose nonzero singular spectrum is identically one.  Nevertheless, for every
fixed $\beta>0$ there is no family of partial isometries $V_L$ satisfying
\eqref{eq:v13-controlled-partial-isometry} with
$\sup_L\|V_L\|_\beta<\infty$.
Indeed any such partial isometry is
\begin{equation}
 V_L=e^{i\theta_L}
 |\delta_{y_L}\rangle\langle\delta_{x_L}|
 \label{eq:v13-site-partial-isometry}
\end{equation}
and hence
\begin{equation}
 \boxed{
 \|V_L\|_\beta=e^{\beta L/2}.}
 \label{eq:v13-site-cost}
\end{equation}
Thus ordinary rank/bundle data and a uniform reduced spectral gap do not by
themselves determine the controlled native phase.
\end{proposition}

\begin{proof}
A partial isometry with initial projection onto the one-dimensional line
$\mathbb C\delta_{x_L}$ and final projection onto
$\mathbb C\delta_{y_L}$ is uniquely of the form
\eqref{eq:v13-site-partial-isometry} up to phase.  Its only nonzero matrix entry
connects sites at periodic distance $L/2$, so its weighted row and column norms
are exactly $e^{\beta L/2}$.  The remaining statements are immediate from the
diagonal formula \eqref{eq:v13-site-D}.
\end{proof}

\begin{firewall}[What the macroscopic example does and does not show]
Proposition~\ref{prop:v13-macro-obstruction} is a native-locality obstruction,
not an additional perturbative gauge anomaly.  It may be embedded as a
spectator defect block next to a common physical Dirac block, so it shows that
``same principal physical jet + same numerical index + same ordinary zero-mode
bundle'' is not a theorem of controlled path connectedness in the unrestricted
local-operator category.  Whether the exact Standard-Model representation and
background grammar exclude such coarse sectors is a further physical
restriction and is not assumed here.
\end{firewall}

\subsection{The affine reduced graph coordinate is local}
\label{sec:v13-local-graph}

The controlled zero-mode class is the only new finite-rank datum needed by the
V11 minimal-stratum geometry.  The remaining graph coordinate can be recovered
by local algebraic operations.

\begin{proposition}[Local graph coordinate for positive index]
\label{prop:v13-positive-local-graph}
Assume $k>0$.  Let $P$ be the positive Wilson projector, let
$\Pi_K$ be the V11 kernel projector, and put
\begin{equation}
 R:=P-\Pi_K.
 \label{eq:v13-R-positive}
\end{equation}
On $\mathcal H_+$ define
\begin{equation}
 S_+:=P_+RP_+\big|_{\mathcal H_+}.
 \label{eq:v13-Splus}
\end{equation}
On the minimal stratum $S_+$ is positive and invertible, and the V11 graph
coordinate is exactly
\begin{equation}
 \boxed{
 A=QRP_+S_+^{-1}:\mathcal H_+\longrightarrow K^\perp.}
 \label{eq:v13-A-positive}
\end{equation}
If $P$ is exponentially local and the reduced chiral singular value has a
uniform positive floor, then $A$ and every fixed source derivative are
exponentially local at some uniform positive rate.
\end{proposition}

\begin{proof}
In the V11 decomposition
$W=K\oplus\operatorname{Graph}(A)$, the projection $R$ onto the graph has the
block matrix
\[
 R=
 \begin{pmatrix}I\\A\end{pmatrix}
 (I+A^*A)^{-1}
 \begin{pmatrix}I&A^*\end{pmatrix}
\]

after ordering $\mathcal H_+\oplus K^\perp$.  Hence
$S_+=(I+A^*A)^{-1}$ and
$QRP_+=A(I+A^*A)^{-1}$, proving
\eqref{eq:v13-A-positive}.  The V11 reduced singular floor is equivalent to a
positive lower bound on $S_+$.  Theorem~\ref{thm:v13-zero-projector-local}
makes $\Pi_K$ local, so $R$ and $S_+$ are local.  Apply
Theorem~\ref{thm:v13-resolvent-locality} to $S_+^{-1}$ and use the weighted
algebra product bound.  Source derivatives follow by the inverse rule.
\end{proof}

\begin{proposition}[Local graph coordinate for negative index]
\label{prop:v13-negative-local-graph}
Assume $k<0$.  Let $\Pi_C$ be the cokernel projection in $\mathcal H_+$ and set
\begin{equation}
 E:=P_+-\Pi_C,
 \qquad
 \widehat S:=\Pi_C+EPE
 \quad\text{on }\mathcal H_+.
 \label{eq:v13-Snegative}
\end{equation}
Then $\widehat S$ is positive and invertible on the protected minimal stratum,
and the V11 graph coordinate is
\begin{equation}
 \boxed{
 A=QPE\widehat S^{-1}E:C^\perp\longrightarrow\mathcal H_-.}
 \label{eq:v13-A-negative}
\end{equation}
It is exponentially local with all prescribed fixed source derivatives.
\end{proposition}

\begin{proof}
On $C^\perp$, the plus--plus block of the graph projector is
$(I+A^*A)^{-1}$; on $C$ the added term $\Pi_C$ is the identity.  Thus
$\widehat S$ has a uniform positive floor.  The block identity
$QPE=A(I+A^*A)^{-1}$ on $C^\perp$ gives
\eqref{eq:v13-A-negative}.  Locality follows exactly as in
Proposition~\ref{prop:v13-positive-local-graph}.
\end{proof}

\subsection{Reduction of the native phase problem to the controlled zero-mode module}
\label{sec:v13-reduction}

\begin{theorem}[Controlled zero-mode reduction of the fixed-index phase problem]
\label{thm:v13-controlled-reduction}
Fix the physical principal Dirac--Yukawa jet, the Standard-Model representation
grading, one positive IR/reference scale, a fixed index $k$, and one protected
minimal reduced-gap component.  Then:
\begin{enumerate}[label=\textup{(R\arabic*)},leftmargin=3em]
\item a native local-operator homotopy induces a controlled equivalence of the
kernel packets for $k>0$ or cokernel packets for $k<0$;
\item after the controlled zero-mode packet has been transported, the remaining
V11 affine graph coordinate can be connected by a straight-line homotopy in
the local algebra, with a uniform reduced singular-value floor;
\item hence no additional nonzero-mode local invariant survives inside the
minimal fixed-index overlap geometry.
\end{enumerate}
The only remaining microscopic datum beyond the already-fixed principal/line
information is the controlled zero-mode module class.
\end{theorem}

\begin{proof}
Item (R1) is Proposition~\ref{prop:v13-native-implies-controlled}.  For (R2),
assume first $k>0$ and let $U_t$ be a local Kato transport of the zero-mode
projection.  Extend it by the identity on $\mathcal H_+$.  Transport the first
graph coordinate by
$A_0(t)=U_tA_0$; this remains a local map into $K_t^\perp$ and therefore defines
a minimal-stratum local projector path by V11.  At the endpoint the kernel is
$K_1$.  Keeping $K_1$ fixed, linearly interpolate
\[
 A_s=(1-s)U_1A_0+sA_1.
\]
The graph fibre is a vector space, so the path stays in the minimal stratum.
Its norm is uniformly bounded, and the explicit V11 formula
$2a^{-1}(I+A_s^*A_s)^{-1/2}$ gives a positive reduced kinetic floor.  On the
compact interacting source chart one shrinks once, as in V9--V11, to preserve
the full chiral/Yukawa reduced floor.  Propositions
\ref{prop:v13-positive-local-graph} and
\ref{prop:v13-negative-local-graph} give the required locality.  The negative
index case is identical with the Kato transport performed in $\mathcal H_+$ on
the cokernel bundle.

All zero-mode transports lie in the local compact/soft ideal and hence do not
alter the frozen physical principal quotient.  The endpoints have the same
principal jet, so the linear graph interpolation may be performed in the affine
subspace with that quotient fixed.  This proves (R2), and (R3) follows from the
affineness of the V11 graph fibre.
\end{proof}

The theorem is deliberately phrased in terms of a controlled zero-mode
transport.  Definition~\ref{def:v13-controlled-equivalence} supplies an endpoint
partial isometry but not automatically an unstabilized path of such projectors.
The next subsection gives an explicit stable construction which is sufficient
for the coefficientwise continuum.

\subsection{Stable controlled lifting with an inert massive spectator}
\label{sec:v13-stable-lift}

Adjoin one auxiliary copy of the same finite chiral carrier, with no physical
sources, no gauge/Yukawa couplings and fixed Wilson sign
\begin{equation}
 \varepsilon_{\rm sp}=\gamma_5.
 \label{eq:v13-spectator-sign}
\end{equation}
Then
\begin{equation}
 D_{{\rm ov},{\rm sp}}
 =a^{-1}(I+\gamma_5^2)=2a^{-1}I,
 \label{eq:v13-spectator-overlap}
\end{equation}
so the spectator is massive, has index zero and has no zero mode.

\begin{lemma}[Exact decoupling of the inert spectator]
\label{lem:v13-spectator-decouple}
At either endpoint, the spectator of
\eqref{eq:v13-spectator-sign} contributes only a source-independent positive
finite factor to the fermion determinant.  It has no source vertices,
BV/BRST descendants, anomaly coefficients, stress/composite writers or line
holonomy.  Consequently it cancels exactly from
\begin{equation}
 W[J]=\log Z[J]-\log Z[0]
 \label{eq:v13-W-normalized}
\end{equation}
and from every retained normalized connected/source/master coefficient.
Adding or deleting finitely many such decoupled spectators at the endpoints
therefore does not change the coefficientwise continuum proved in V7.
\end{lemma}

\begin{proof}
The spectator action is a fixed nonzero quadratic form independent of every
physical source and background writer.  Its Grassmann integral is its constant
nonzero determinant.  The same constant multiplies $Z[J]$ and $Z[0]$ and hence
cancels in \eqref{eq:v13-W-normalized}.  With no source or gauge couplings it
has no differentiated vertices, no BV bracket entries and no anomaly trace.
Its determinant section is the fixed positive scalar section, so its line
transport is trivial.
\end{proof}

Let $J:\mathcal H_-\to\mathcal H_-^{\rm sp}$ be the fixed onsite copy
identification.  Suppose now that
$V$ is a controlled partial isometry from the initial zero-mode projector
$\Pi_0$ to the final one $\Pi_1$.  Put
\begin{equation}
 \widetilde V:=JV:
 \mathcal H_-\longrightarrow\mathcal H_-^{\rm sp}.
 \label{eq:v13-Vtilde}
\end{equation}

\begin{lemma}[Controlled rotation into the spectator copy]
\label{lem:v13-spectator-rotation}
For $0\le\theta\le\pi/2$ the block operator on
$\mathcal H_-\oplus\mathcal H_-^{\rm sp}$
\begin{equation}
 \Pi_\theta=
 \begin{pmatrix}
 \cos^2\theta\,\Pi_0 &
 \cos\theta\sin\theta\,\widetilde V^*\\
 \cos\theta\sin\theta\,\widetilde V &
 \sin^2\theta\,J\Pi_1J^*
 \end{pmatrix}
 \label{eq:v13-spectator-rotation}
\end{equation}
is an orthogonal rank-$|k|$ projection, exponentially local uniformly in
$\theta$, with
\begin{equation}
 \Pi_{0}=\Pi_0\oplus0,
 \qquad
 \Pi_{\pi/2}=0\oplus J\Pi_1J^*.
 \label{eq:v13-rotation-endpoints}
\end{equation}
A second identical onsite rotation using $J^*$ moves
$0\oplus J\Pi_1J^*$ to $\Pi_1\oplus0$.  Hence controlled-equivalent zero-mode
packets are connected after one inert-copy stabilization by an explicit local
projection path of the same rank.
\end{lemma}

\begin{proof}
Use
$\widetilde V^*\widetilde V=\Pi_0$ and
$\widetilde V\widetilde V^*=J\Pi_1J^*$.  Direct block multiplication shows
$\Pi_\theta^2=\Pi_\theta$ and self-adjointness is manifest.  Locality follows
from Definition~\ref{def:v13-controlled-equivalence}, the onsite nature of $J$,
and Lemma~\ref{lem:v13-Schur-algebra}.  The endpoint formulas are immediate.
The second rotation is the same construction with the roles of the two copies
reversed and partial isometry $J^*J\Pi_1=\Pi_1$ on the relevant range.
\end{proof}

\begin{theorem}[Stable native lifting from controlled zero-mode equivalence]
\label{thm:v13-stable-native-lift}
Let two V11 protected minimal fixed-index regulators $A,B$ have the same
physical principal jet, representation grading and global complex line datum
(and the same real orientation datum when a Pfaffian sector is present).
Assume their kernel packets for $k>0$, or cokernel packets for $k<0$, are
controlled-equivalent.  After adjoining one inert massive spectator copy to
both endpoints, there is a finite concatenation of uniformly exponentially
local native homotopies joining the stabilized regulators such that:
\begin{enumerate}[label=\textup{(S\arabic*)},leftmargin=3em]
\item the index and minimal zero-mode rank remain fixed;
\item the physical principal Dirac--Yukawa jet remains fixed;
\item the reduced kinetic singular value has a positive uniform floor, and the
full chiral/Yukawa reduced floor remains positive after one common finite
shrinking of the protected source chart;
\item all required fixed source jets remain in one positive locality class;
\item the zero-mode line/saturation packet is transported continuously.
\end{enumerate}
Thus every segment is V11-admissible.
\end{theorem}

\begin{proof}
Treat $k>0$; the cokernel proof is the chiral dual.  Add the inert balanced copy
of \eqref{eq:v13-spectator-sign}.  The endpoint kernel projectors become
$\Pi_0\oplus0$ and $\Pi_1\oplus0$ in the enlarged negative-chirality carrier.
Lemma~\ref{lem:v13-spectator-rotation} supplies a controlled local path between
them, of constant rank $k$.

Use Theorem~\ref{thm:v13-local-Kato} on this explicit local projector path to
obtain a local unitary transport of the negative-chirality complement.  Extend
it by the identity on the enlarged positive-chirality carrier and transport the
initial local graph coordinate supplied by
Proposition~\ref{prop:v13-positive-local-graph}.  This gives the first stage of
a minimal-stratum projector homotopy.  Once the target kernel projection is
reached, linearly interpolate the transported graph coordinate to the target
graph coordinate (with zero graph on the spectator block).  The V11 affine
geometry keeps the index minimal throughout and gives the explicit reduced
kinetic floor
\[
 {2\over a\sqrt{1+\|A_t\|^2}}.
\]
All operators used are finite products, sums, inverses across a positive gap,
or Dyson evolutions inside the weighted local algebra; hence they retain some
uniform positive exponential rate by the preceding theorems.

The zero-mode rotation is a compact/soft local deformation and the common
principal quotient is frozen.  The endpoint graph coordinates have the same
physical first jet, so the affine interpolation is taken inside that fixed
principal-jet fibre.  The full Yukawa/source pivot is an open condition on the
compact path, so the usual one-time chart shrinking of V9--V12 supplies a
common positive reduced floor.  Kato exterior powers transport the zero-mode
line and saturation writer.
\end{proof}

\begin{theorem}[Stable controlled regulator universality]
\label{thm:v13-stable-universality}
Under the hypotheses of Theorem~\ref{thm:v13-stable-native-lift}, the original
unstabilized regulators have the same matched coefficientwise continuum.  More
precisely, for every prescribed finite coefficient packet $\mathfrak q$ there
is an invertible finite local map
\begin{equation}
 F_{\infty,\mathfrak q,k}^{B\leftarrow A}
 \label{eq:v13-stable-map}
\end{equation}
such that
\begin{equation}
 \boxed{
 \widehat{\mathfrak C}^{B,(k),\rm sat}_{\infty,\lambda}
 =F_{\infty,\mathfrak q,k}^{B\leftarrow A}
  \widehat{\mathfrak C}^{A,(k),\rm sat}_{\infty,\lambda}.}
 \label{eq:v13-stable-universality}
\end{equation}
The equality includes the complete normalized connected/source packet,
operator mixing, stress/contact terms, ghosts/antifields, BV/BRST/master
relations and determinant/Pfaffian line insertions, and the maps can be chosen
projectively in perturbative order.
\end{theorem}

\begin{proof}
Apply V11 fixed-index universality to the stabilized local homotopy of
Theorem~\ref{thm:v13-stable-native-lift}.  At both endpoints remove the inert
spectator using Lemma~\ref{lem:v13-spectator-decouple}.  Since this removal is
an exact equality of normalized coefficient systems rather than a limiting
argument, the matching map descends to the original regulators.  The V7
chronological construction and the spectator's trivial source packet preserve
strict perturbative-order projectivity.
\end{proof}

\subsection{A positive criterion: uniformly tight zero modes}
\label{sec:v13-tight}

The controlled module condition is stronger than ordinary bundle isomorphism,
but it is automatic for many localized defect sectors.

\begin{definition}[Uniformly tight zero-mode packet]
\label{def:v13-tight-packet}
A finite-rank zero-mode projector family $\Pi$ is \emph{uniformly tight} if
there is a family of multiplication cutoffs $\chi_R$ supported in sets of
uniformly bounded diameter (independent of cutoff and volume) such that
\begin{equation}
 \|(I-\chi_R)\Pi\|+\|\Pi(I-\chi_R)\|
 \xrightarrow[R\to\infty]{}0
 \label{eq:v13-tightness}
\end{equation}
uniformly on the compact background panel, and the same holds for the retained
source derivatives.
\end{definition}

\begin{theorem}[Bundle isomorphism plus uniform tightness gives a controlled isomorphism]
\label{thm:v13-tight-controlled}
Let $\Pi_0,\Pi_1$ be uniformly tight zero-mode packets of the same rank over a
compact background panel.  Assume their ordinary complex zero-mode bundles are
isomorphic by a smooth unitary bundle map.  Then
\begin{equation}
 \boxed{\Pi_0\sim_{\rm ctl}\Pi_1.}
 \label{eq:v13-tight-ctl}
\end{equation}
Consequently Theorem~\ref{thm:v13-stable-universality} applies.
\end{theorem}

\begin{proof}
Let $W$ be a smooth unitary bundle isomorphism
$\operatorname{Ran}\Pi_0\to\operatorname{Ran}\Pi_1$, extended by zero on the
orthogonal complement.  Because the panel is compact and the rank is finite,
$W$ is uniformly bounded with all retained source derivatives.  Choose a common
localizing cutoff $\chi_R$ from Definition~\ref{def:v13-tight-packet} so that
\begin{equation}
 T_R:=\Pi_1\chi_RW\chi_R\Pi_0
 \label{eq:v13-TR}
\end{equation}
differs from $W$ by less than $1/4$ on the initial zero-mode range, uniformly on
the panel.  The support of $\chi_R$ has uniformly bounded diameter, so $T_R$ is
a controlled local finite-rank kernel.  Its restriction to
$\operatorname{Ran}\Pi_0$ is invertible onto
$\operatorname{Ran}\Pi_1$.  Put
\begin{equation}
 V=T_R\bigl(T_R^*T_R+(I-\Pi_0)\bigr)^{-1/2}\Pi_0.
 \label{eq:v13-polar-V}
\end{equation}
The operator in parentheses has a uniform positive floor.  Theorem
\ref{thm:v13-resolvent-locality} gives a local inverse square root, so $V$ is
controlled.  Polar decomposition gives
$V^*V=\Pi_0$ and $VV^*=\Pi_1$.  Source derivatives follow from the same
functional calculus.
\end{proof}

\begin{corollary}[Localized fixed-index sectors have no extra controlled obstruction]
\label{cor:v13-localized-sector}
On a uniformly tight zero-mode sector, the V12 ordinary kernel/cokernel bundle
class together with the already-declared complex/Pfaffian line data is enough
to imply stable native regulator universality.
\end{corollary}

\subsection{Controlled phase hierarchy and exact boundary of V13}
\label{sec:v13-boundary}

The preceding results sharpen the V12 three-level hierarchy to
\begin{equation}
 \boxed{
 \text{line class}
 \ \leftarrow\ 
 \text{ordinary zero-mode bundle class}
 \ \leftarrow\ 
 \text{controlled zero-mode module class}
 \ \leftarrow\ 
 \text{unstabilized native component}.}
 \label{eq:v13-four-level}
\end{equation}
The first two arrows can lose information by V12.  Proposition
\ref{prop:v13-macro-obstruction} shows that the second-to-third arrow can also
lose information.  Theorem~\ref{thm:v13-stable-native-lift} shows that the
third level is already sufficient after one physically inert stabilization.
What remains is cancellation of that stabilization in the microscopic path,
not in the normalized continuum where it has already disappeared exactly.

\begin{theorem}[V13 phase-reduction theorem]
\label{thm:v13-phase-reduction}
For the protected minimal fixed-index overlap problem with a positive reduced
singular-value floor:
\begin{enumerate}[label=\textup{(C\arabic*)},leftmargin=3em]
\item zero-mode Riesz projectors and Kato transports are automatically
quasilocal; arbitrary nonlocal finite-rank projectors are not the correct
spectral variables;
\item the nonzero-mode minimal-stratum fibre contributes no independent local
phase invariant;
\item a necessary invariant of an unstabilized native component is the
controlled zero-mode module class;
\item equality of that class is sufficient for native lifting after one inert
massive stabilization and hence is sufficient for equality of the original
normalized coefficientwise continua;
\item in uniformly tight zero-mode sectors, ordinary bundle isomorphism already
implies controlled equivalence.
\end{enumerate}
\end{theorem}

\begin{proof}
Items (C1)--(C5) are respectively Theorems
\ref{thm:v13-zero-projector-local},
\ref{thm:v13-controlled-reduction},
Proposition~\ref{prop:v13-native-implies-controlled},
Theorems~\ref{thm:v13-stable-native-lift}--\ref{thm:v13-stable-universality},
and Theorem~\ref{thm:v13-tight-controlled}.
\end{proof}

\subsection{V13 anti-circularity audit}
\label{sec:v13-audit}

\begin{enumerate}[leftmargin=2.4em]
\item \textbf{V12's counterexample was not discarded.}
V13 identifies the extra hypothesis it lacked: a local parent operator with a
uniform spectral separation of the zero eigenspace.  Riesz functional calculus
then forces locality of the actual V11 zero-mode projector.
\item \textbf{Locality of each endpoint projector was not promoted to locality
of an endpoint identification.}
Proposition~\ref{prop:v13-macro-obstruction} exhibits two perfectly local,
uniformly gapped rank-one zero modes whose unique intertwiner has locality cost
$e^{\beta L/2}$.
\item \textbf{The reduced graph fibre is not treated as a new topological
sector.}
Its coordinate is reconstructed by the local formulas
\eqref{eq:v13-A-positive} and \eqref{eq:v13-A-negative} and interpolated in an
affine space.
\item \textbf{The stable lift does not add physical massless degrees of
freedom.}
The auxiliary copy has overlap operator $2a^{-1}I$, no zero modes and no
sources.  Its endpoint contribution cancels exactly from normalized
coefficients.
\item \textbf{Stable continuum equivalence is distinguished from an
unstabilized microscopic path.}
V13 proves the former.  Removing the spectator from the path itself is the next
cancellation problem.
\item \textbf{The complex and real line data remain separate.}
All V13 controlled constructions carry the V12 complex saturation writer and,
when present, require the same real Pfaffian orientation component.
\end{enumerate}

\section{V13 status ledger}
\label{sec:v13-status}

\subsection*{Proved in V13}
\begin{enumerate}[leftmargin=2.2em]
\item Weighted exponentially local kernels form a Banach $*$-algebra, and a
uniform spectral gap gives a cutoff-uniform exponentially local resolvent after
a finite loss of exponent.  Riesz projectors, signs and fixed source derivatives
inherit the same property.
\item The V11 kernel/cokernel Riesz projectors and reduced Moore--Penrose inverse
are therefore uniformly exponentially local on every protected minimal
fixed-index chart with a positive reduced singular-value floor.
\item Kato transport along an already-native fixed-index path is itself
exponentially local, with all retained source derivatives.  Hence native paths
preserve a controlled zero-mode module class, stronger than the ordinary bundle
class.
\item Two macroscopically separated onsite zero modes give an explicit
counterexample showing that ordinary rank/bundle/determinant data plus a uniform
reduced gap do not imply controlled equivalence: the unique partial isometry
has locality norm $e^{\beta L/2}$.
\item The V11 affine graph coordinate is reconstructed from local projectors
and a gapped inverse by the exact formulas
\eqref{eq:v13-A-positive}--\eqref{eq:v13-A-negative}; it carries no independent
nonzero-mode local phase invariant.
\item Controlled-equivalent zero-mode packets admit an explicit stable native
lift after adjoining one inert massive spectator copy.  The path preserves
index, minimal nullity, the physical principal jet, a positive reduced gap,
source locality and the zero-mode line packet.
\item The inert spectator cancels exactly from normalized connected/source/BV
coefficients.  Therefore controlled equivalence implies equality of the
original unstabilized coefficientwise continua, despite the temporary
microscopic stabilization used to construct the path.
\item In uniformly tight localized zero-mode sectors, an ordinary smooth bundle
isomorphism can be localized and polar-corrected to a controlled partial
isometry.  Such sectors therefore have no extra controlled obstruction beyond
the V12 bundle/line data.
\end{enumerate}

\subsection*{Inherited}
\begin{enumerate}[leftmargin=2.2em]
\item V1--V7: source-complete fixed-order continuum, anomaly-free common BV
restoration, arbitrary-fixed-loop forests, ultraviolet Cauchy convergence,
scheme/reference covariance and exact loop-order projectivity.
\item V8--V10: regulator-marked two-scale comparison and regulator-class
universality on the protected index-zero single-cone sectors.
\item V11: minimal fixed-index affine geometry, reduced pseudoinverse calculus,
zero-mode Schur factorization, saturation writers and local fixed-index
universality.
\item V12: global zero-mode bundle/line patching, faithful-spin complex line
trivialization, global saturation, the stable finite-Grassmannian bundle
classification, and global universality within a protected native component.
\end{enumerate}

\subsection*{Conditional}
\begin{enumerate}[leftmargin=2.2em]
\item The local resolvent/Riesz theorem uses a cutoff-uniform exponential
locality bound and a positive reduced singular-value floor.  If the floor
collapses with volume, the V12 nonlocal projector phenomenon can reappear.
\item Stable native lifting requires a controlled partial isometry of the
kernel/cokernel packets.  Ordinary bundle isomorphism supplies one automatically
only under additional locality hypotheses such as the uniform-tightness theorem.
\item The spectator argument proves equality of normalized continuum
coefficients, not an unstabilized microscopic homotopy in the original carrier.
\item The V12 real Pfaffian mod-two orientation requirement remains in force.
\item All statements remain coefficientwise at fixed finite
$(L,\Delta,r,N)$ and positive IR/reference scale.
\end{enumerate}

\subsection*{Open beyond V13}
\begin{enumerate}[leftmargin=2.2em]
\item Prove \emph{controlled cancellation}: determine whether a controlled
zero-mode module isomorphism always yields an unstabilized native homotopy in
the exact Standard-Model carrier, or exhibit a genuine obstruction to removing
the inert spectator from the microscopic path.
\item Classify controlled zero-mode modules on physically relevant
noncontractible gauge-background components.  The macroscopic-transport example
shows that the answer can contain coarse/locality information absent from
ordinary $K$-theory.
\item Compute the V12 real Pfaffian mod-two spectral-flow character on any
independently physical Majorana sector.
\item Analyze exits from the minimal fixed-index stratum, including accidental
zero-mode-pair creation and true index change.
\item The independent stronger problems remain: positive-IR removal, analytic
finite-coupling RG trajectories, perturbative summability and nonperturbative
completion.
\end{enumerate}

\subsection*{Exact next load-bearing theorem}
\begin{center}
\fbox{\parbox{0.92\textwidth}{
\textbf{Controlled cancellation / unstabilized native lifting.}
Fix the Standard-Model representation, physical principal Dirac--Yukawa jet,
index $k$, reduced-gap component and global zero-mode line data.  Given a
controlled partial isometry between the two local kernel/cokernel packets,
prove that it extends inside the same native gauge/spin-covariant exponentially
local carrier to a path of minimal fixed-index Wilson projectors without the
auxiliary massive copy, or identify the exact obstruction (for example a
controlled complement-unitary/$K_1$ or coarse-transport class).  A positive
cancellation theorem, combined with V13 and
Theorem~\ref{thm:v12-global-index-universality}, would upgrade stable controlled
universality to a complete unstabilized global fixed-index regulator-class
theorem.}}
\end{center}

\section*{Append-only continuation marker: V13}
\addcontentsline{toc}{section}{Append-only continuation marker: V13}
\textbf{End of V13.}  The complete V1--V12 body above is retained verbatim.
V13 goes upstream from the abstract-projector locality obstruction and proves
that the actual V11 zero-mode Riesz projectors are uniformly exponentially
local whenever their native parent operator has the declared reduced spectral
floor.  It identifies the remaining microscopic datum as a controlled local
zero-mode module class, shows by a macroscopically separated onsite example
that this datum is finer than ordinary bundle/determinant information, and
proves that the V11 reduced graph coordinate carries no further local phase
invariant.  Controlled-equivalent zero-mode packets admit an explicit stable
native lift after one inert massive spectator copy; the spectator cancels
exactly from all normalized coefficientwise observables, yielding stable
global regulator universality.  The next gate is controlled cancellation: can
the spectator be removed from the microscopic path itself?


\clearpage
\section{V14 continuation: controlled cancellation, complement defects, and unstabilized native lifting}
\label{sec:v14-continuation}

V13 reduced the protected fixed-index phase problem to a controlled local
zero-mode module.  It also proved that controlled-equivalent packets become
connected after adjoining one inert massive copy.  The remaining question is
therefore a genuine \emph{cancellation} problem: can the auxiliary copy be
removed from the microscopic path itself?

The strongest possible answer is false in an arbitrary local coefficient
algebra.  Murray--von Neumann equivalence of two zero-mode projections gives a
controlled partial isometry between their ranges, but an unstabilized unitary
extension additionally requires the two complementary modules to cancel.  Even
when that complement cancellation holds, the extending unitary must lie in the
identity component of the controlled unitary group in order to generate a
projection homotopy.  These are logically distinct requirements.

The present continuation makes that distinction exact.  It proves:
\begin{enumerate}[label=\textup{(V14.\arabic*)},leftmargin=3.8em]
\item a necessary-and-sufficient complement criterion for extending a
controlled partial isometry to a controlled unitary on the original carrier;
\item a necessary-and-sufficient identity-component criterion for an
unstabilized controlled projection homotopy;
\item a propagation-zero finite-dimensional counterexample showing that
controlled equivalence alone does not imply cancellation before stabilization;
\item a quantitative small-angle cancellation theorem and an explicit
``internal buffer'' theorem which remove the spectator whenever the original
carrier already contains one suitable local massive reserve;
\item an unstabilized fixed-index native lifting/universality theorem whenever
the zero-mode projections lie in the same controlled $U_0$ orbit;
\item a stable-range corollary showing that on fixed finite-dimensional
parameter panels the purely finite-dimensional complement-bundle obstruction
disappears at sufficiently large ambient rank, leaving a genuinely controlled
locality/unitary-component problem as the remaining gate.
\end{enumerate}

Throughout V14 let $\mathfrak A_\beta$ denote one of the V13 unital Banach
$*$-algebras of exponentially local, gauge/spin-covariant kernel families with
the prescribed finite source derivatives, at some fixed exponent
$0<\beta<\alpha$.  All projections, partial isometries and unitaries below are
required to obey the same representation/source constraints.  Write
\begin{equation}
 \mathcal U(\mathfrak A_\beta)
 :=\{u\in\mathfrak A_\beta:u^*u=uu^*=I\},
 \qquad
 \mathcal U_0(\mathfrak A_\beta)
 \label{eq:v14-U0}
\end{equation}
for its unitary group and identity component.

\subsection{Exact complement cancellation for a controlled partial isometry}
\label{sec:v14-complement-cancellation}

Let $p,q\in\mathfrak A_\beta$ be projections and let
\begin{equation}
 v^*v=p,
 \qquad
 vv^*=q
 \label{eq:v14-v}
\end{equation}
be a controlled partial isometry, as in V13.

\begin{theorem}[Controlled unitary extension if and only if the complements cancel]
\label{thm:v14-unitary-extension}
The following are equivalent.
\begin{enumerate}[label=\textup{(E\arabic*)},leftmargin=3em]
\item There exists $u\in\mathcal U(\mathfrak A_\beta)$ such that
\begin{equation}
 up=v.
 \label{eq:v14-up-v}
\end{equation}
\item There exists a controlled partial isometry $w\in\mathfrak A_\beta$ with
\begin{equation}
 \boxed{
 w^*w=I-p,
 \qquad
 ww^*=I-q.}
 \label{eq:v14-complement-MvN}
\end{equation}
\end{enumerate}
When \textup{(E2)} holds,
\begin{equation}
 \boxed{u=v+w}
 \label{eq:v14-u-vw}
\end{equation}
is a unitary extension.  Conversely every extension has
$w=u(I-p)$.
\end{theorem}

\begin{proof}
Assume \textup{(E2)}.  Since $v=qvp$ and
$w=(I-q)w(I-p)$, the initial and final projections of $v$ and $w$ are
orthogonal.  Hence
\[
 v^*w=w^*v=vw^*=wv^*=0.
\]
Therefore
\[
 (v+w)^*(v+w)=p+(I-p)=I,
 \qquad
 (v+w)(v+w)^*=q+(I-q)=I.
\]
Also $(v+w)p=v$.  Thus \eqref{eq:v14-u-vw} is the desired unitary.

Conversely suppose $u$ is unitary and $up=v$.  Then
$upu^*=vv^*=q$.  Put $w=u(I-p)$.  One has
\[
 w^*w=(I-p)u^*u(I-p)=I-p
\]
and
\[
 ww^*=u(I-p)u^*=I-upu^*=I-q.
\]
Thus \textup{(E2)} holds.
\end{proof}

This theorem isolates the first unstable obstruction exactly.  The endpoint
partial isometry of V13 proves $p\sim q$, but unstabilized extension asks in
addition for
\begin{equation}
 \boxed{I-p\sim I-q}
 \label{eq:v14-complement-cancel}
\end{equation}
in the \emph{same controlled algebra}.  Equality in stabilized $K_0$ does not,
by itself, imply this unstabilized Murray--von Neumann cancellation.

\begin{definition}[Controlled complement defect]
\label{def:v14-complement-defect}
For controlled-equivalent $p,q$ define the \emph{unstable complement defect}
to vanish when the complements are Murray--von Neumann equivalent
\emph{inside the original controlled carrier}:
\begin{equation}
 \boxed{
 \mathfrak c_\beta(p,q)=0
 \quad\Longleftrightarrow\quad
 \exists\,w\in\mathfrak A_\beta:\ 
 w^*w=I-p,\ ww^*=I-q.}
 \label{eq:v14-complement-defect}
\end{equation}
No matrix amplification or added free summand is allowed in this definition.
Thus $\mathfrak c_\beta$ is deliberately an unstable relation, stronger than
equality of the corresponding group-completed $K_0$ classes.
\end{definition}

By Theorem~\ref{thm:v14-unitary-extension},
$\mathfrak c_\beta(p,q)=0$ is exactly the existence of a controlled unitary
extension of a chosen controlled equivalence after its complement has been
filled in inside the same carrier.

\subsection{The second obstruction: the controlled unitary component}
\label{sec:v14-unitary-component}

A unitary extension does not yet give a path from the identity.  For a
\emph{fixed} controlled partial isometry $v$, the residual component is also
exact.

Suppose \eqref{eq:v14-complement-MvN} holds and fix one complement partial
isometry $w_0$.  Put
\begin{equation}
 u_0:=v+w_0\in\mathcal U(\mathfrak A_\beta).
 \label{eq:v14-u0}
\end{equation}
If $w_1$ is another complement equivalence, then
\begin{equation}
 z:=w_0^*w_1
 \label{eq:v14-zcorner}
\end{equation}
is a unitary in the corner
$(I-p)\mathfrak A_\beta(I-p)$ and
\begin{equation}
 v+w_1=u_0(p+z).
 \label{eq:v14-extension-difference}
\end{equation}

\begin{definition}[Relative controlled unitary-component class]
\label{def:v14-kappa}
When $\mathfrak c_\beta(p,q)=0$, define
\begin{equation}
 \kappa_\beta(v)
 :=[u_0]\,
 \operatorname{im}\!\left(
 \pi_0\mathcal U((I-p)\mathfrak A_\beta(I-p))
 \longrightarrow
 \pi_0\mathcal U(\mathfrak A_\beta)
 \right)
 \label{eq:v14-kappa}
\end{equation}
as a right coset in the component group.  Equation
\eqref{eq:v14-extension-difference} makes this coset independent of the chosen
$w_0$.
\end{definition}

\begin{proposition}[Exact component criterion for an extension of a fixed $v$]
\label{prop:v14-kappa-criterion}
Under $\mathfrak c_\beta(p,q)=0$, there exists a controlled unitary extension
$u$ of the fixed $v$ with
\begin{equation}
 u\in\mathcal U_0(\mathfrak A_\beta)
 \label{eq:v14-extension-U0}
\end{equation}
if and only if the coset $\kappa_\beta(v)$ contains the identity component.
\end{proposition}

\begin{proof}
Every extension has the form $u_0(p+z)$ with $z$ a unitary in the complement
corner, by \eqref{eq:v14-extension-difference}.  Hence one can move $u_0$ into
the identity component while keeping its restriction to $p$ equal to $v$
exactly when the component class of $u_0$ can be cancelled by the image of a
complement-corner unitary.  This is precisely the stated coset condition.
\end{proof}

If the endpoint zero-mode identification is not fixed a priori, one may also
change $v$ by a unitary in the $p$ corner.  The corresponding obstruction is a
double-coset version of \eqref{eq:v14-kappa}; V14 will not need separate
notation for it because the projection-homotopy criterion below is cleaner.

\subsection{Projection paths are exactly controlled $U_0$ orbits}
\label{sec:v14-U0-orbits}

We first record the local direct rotation between two nearby projections.

\begin{lemma}[Controlled direct rotation for a small projection angle]
\label{lem:v14-direct-rotation}
Let $p,q\in\mathfrak A_\beta$ be projections with
\begin{equation}
 \delta:=\|p-q\|_{\rm op}<1.
 \label{eq:v14-small-angle}
\end{equation}
Put
\begin{equation}
 X:=qp+(I-q)(I-p).
 \label{eq:v14-X}
\end{equation}
Then
\begin{align}
 X-I&=(q-p)(2p-I),
 \label{eq:v14-X-near-I}\\
 X^*X&=I-(p-q)^2,
 \label{eq:v14-XstarX}
\end{align}
so $X$ is invertible and
\begin{equation}
 \boxed{
 u_{q,p}:=X(X^*X)^{-1/2}}
 \label{eq:v14-direct-unitary}
\end{equation}
is a unitary in $\mathfrak A_{\beta'}$ for every sufficiently small fixed
loss $0<\beta'<\beta$.  It satisfies
\begin{equation}
 u_{q,p}pu_{q,p}^*=q,
 \qquad
 u_{q,p}\in\mathcal U_0(\mathfrak A_{\beta'}).
 \label{eq:v14-direct-properties}
\end{equation}
All prescribed fixed source derivatives obey uniform bounds when
$\delta\le1-\eta$ for a fixed $\eta>0$.
\end{lemma}

\begin{proof}
The first identity is direct expansion.  For the second,
$Xp=qX$, so the cross terms in $X^*X$ vanish between $p$ and $I-p$; expanding
gives
$X^*X=I-p-q+pq+qp=I-(p-q)^2$.  Thus
$X^*X\succeq(1-\delta^2)I$.

The V13 weighted Banach-algebra/resolvent functional calculus gives
$(X^*X)^{-1/2}\in\mathfrak A_{\beta'}$ with all fixed source derivatives.
The polar factor \eqref{eq:v14-direct-unitary} is unitary.  Since
$Xp=qX$ and $X^*X$ commutes with $p$, one has $u_{q,p}p=qu_{q,p}$, proving the
projection identity.

Finally \eqref{eq:v14-X-near-I} gives $\|X-I\|<1$.  Hence
$X_s=I+s(X-I)$ is invertible for every $s\in[0,1]$.  Its polar factors
$u_s=X_s(X_s^*X_s)^{-1/2}$ form a continuous controlled unitary path from
$I$ to $u_{q,p}$.  The uniform source bounds follow from the fixed lower
spectral margin $1-\delta^2\ge\eta(2-\eta)$.
\end{proof}

\begin{theorem}[Exact controlled projection-homotopy criterion]
\label{thm:v14-projection-U0}
For two projections $p,q\in\mathfrak A_\beta$, the following are equivalent
(after one harmless fixed loss of exponential weight if needed for functional
calculus).
\begin{enumerate}[label=\textup{(H\arabic*)},leftmargin=3em]
\item There exists a norm-continuous path of controlled projections
$p_t\in\mathfrak A_{\beta'}$, $0\le t\le1$, with $p_0=p$ and $p_1=q$.
\item There exists a controlled unitary
\begin{equation}
 u\in\mathcal U_0(\mathfrak A_{\beta'})
 \label{eq:v14-U0-orbit}
\end{equation}
such that
\begin{equation}
 \boxed{q=upu^*.}
 \label{eq:v14-U0-conjugacy}
\end{equation}
\end{enumerate}
The same equivalence holds in the source-differentiable gauge/spin-equivariant
subalgebra.
\end{theorem}

\begin{proof}
If \textup{(H2)} holds, choose a continuous controlled unitary path
$u_t$ from $I$ to $u$ and set $p_t=u_tpu_t^*$.  This gives \textup{(H1)}.

Conversely suppose \textup{(H1)} holds.  Uniform continuity of $t\mapsto p_t$
on the compact interval gives a finite partition
$0=t_0<\cdots<t_M=1$ such that
\[
 \|p_{t_{j+1}}-p_{t_j}\|_{\rm op}<1/2.
\]
Apply Lemma~\ref{lem:v14-direct-rotation} to each adjacent pair.  It gives
$u_j\in\mathcal U_0(\mathfrak A_{\beta'})$ with
$u_jp_{t_j}u_j^*=p_{t_{j+1}}$.  The finite product
$u=u_{M-1}\cdots u_0$ lies in the identity component and satisfies
$q=upu^*$.  Fixed source and covariance constraints are preserved at each
step because the direct-rotation formula uses only algebra operations and
functional calculus inside the same constrained subalgebra.
\end{proof}

\begin{definition}[Unstable controlled cancellation obstruction]
\label{def:v14-unstable-obstruction}
For controlled-equivalent zero-mode projections $p,q$, define
\begin{equation}
 \boxed{
 \operatorname{Can}_\beta(p,q)=0
 \quad\Longleftrightarrow\quad
 q=upu^*\text{ for some }
 u\in\mathcal U_0(\mathfrak A_\beta).}
 \label{eq:v14-Can}
\end{equation}
Thus $\operatorname{Can}_\beta$ is the exact obstruction to an unstabilized
controlled zero-mode projection path.  It contains both complement cancellation
and unitary-component information.
\end{definition}

\subsection{Controlled equivalence alone does not cancel in general}
\label{sec:v14-counterexample}

The next counterexample is deliberately propagation zero.  It shows that no
proof of V14 cancellation can follow from V13's endpoint partial isometry alone;
one must use additional structure of the actual Standard-Model carrier or a
stable-range/local-buffer hypothesis.

\begin{theorem}[Propagation-zero unstable cancellation failure]
\label{cth:v14-unstable-cancellation}
There exist a compact parameter space $B$, a finite-dimensional onsite carrier,
and smooth propagation-zero projections $p,q$ such that:
\begin{enumerate}[label=\textup{(C\arabic*)},leftmargin=3em]
\item $p$ and $q$ have the same finite rank and are connected by a smooth
propagation-zero partial isometry $v$;
\item their range bundles are isomorphic and their determinant lines agree;
\item nevertheless no continuous onsite unitary family $u(b)$ satisfies
$upu^*=q$.
\end{enumerate}
Consequently controlled Murray--von Neumann equivalence does not imply
$\operatorname{Can}_\beta(p,q)=0$ in an arbitrary native coefficient algebra.
The same example embeds into the V11 minimal positive-index projector geometry
by taking zero graph coordinate.
\end{theorem}

\begin{proof}
Take $B=S^6$.  Complex rank-two bundles over $S^6$ are classified by the
clutching group
\[
 \pi_5(U(2))\cong\pi_5(S^3)\cong\mathbb Z_2.
\]
Let $F\to S^6$ be the nontrivial bundle represented by the nonzero element.
Under stabilization, the clutching class maps to
$\pi_5(U)\cong\mathbb Z$ by Bott periodicity.  Every homomorphism
$\mathbb Z_2\to\mathbb Z$ is zero, so $F$ is stably trivial.  Hence for some
$r\ge1$ there is a unitary bundle isomorphism
\begin{equation}
 F\oplus\underline{\mathbb C}^{\,r}
 \simeq
 \underline{\mathbb C}^{\,r+2}.
 \label{eq:v14-stably-trivial-F}
\end{equation}
Use such an isomorphism to regard the trivial rank-$r$ summand as a subbundle
$E_0$ of the trivial rank-$(r+2)$ bundle whose orthogonal complement is $F$.
Let $E_1$ be the standard trivial rank-$r$ coordinate subbundle, whose
orthogonal complement is the trivial rank-two bundle.  Let $p,q$ be the
orthogonal projections onto $E_0,E_1$.  Since both range bundles are trivial,
choose a smooth unitary bundle isomorphism
$v:E_0\to E_1$ and extend it by zero.  Then
$v^*v=p$ and $vv^*=q$.

If a continuous unitary family $u$ of the ambient trivial bundle satisfied
$upu^*=q$, it would restrict to a unitary isomorphism
\[
 E_0^\perp\simeq E_1^\perp,
\]
that is, $F\simeq\underline{\mathbb C}^2$, contradicting the nontrivial
clutching class.  All matrices act onsite in the physical lattice variable, so
the locality cost is exactly zero propagation.

For the V11 positive-index embedding, take
$\mathcal H_-\cong\mathbb C^{r+2}$, a fixed
$\mathcal H_+$ of the same ambient dimension, kernel plane $K_i=E_i$, and
zero affine graph coordinate $A_i=0$.  Then
$W_i=K_i\oplus\mathcal H_+$ lies in the minimal index-$r$ stratum, while a
minimal-stratum homotopy would in particular homotope the kernel projections,
which has just been excluded.
\end{proof}

\begin{firewall}[The counterexample is unstable by construction]
The bundle $F$ in Countertheorem~\ref{cth:v14-unstable-cancellation} is
stably trivial.  Adding enough trivial ambient reserve removes this particular
obstruction.  This is exactly why V13's one-copy stabilization can succeed even
when an unstabilized finite carrier fails.  On a fixed finite-dimensional
parameter panel, the V12 stable-range theorem similarly removes such purely
finite-Grassmannian cancellation defects once the complement rank is large
enough.  The counterexample therefore identifies the logical obstruction; it
does not claim that this specific $S^6$ defect survives in every large-cutoff
Standard-Model carrier.
\end{firewall}

\subsection{Two constructive regimes where cancellation is automatic}
\label{sec:v14-positive-regimes}

The exact obstruction theorem is useful because it reveals simple sufficient
conditions which are native and quantitative.

\begin{corollary}[Small-angle unstabilized cancellation]
\label{cor:v14-small-angle}
If two protected zero-mode projectors satisfy
$\|p-q\|_{\rm op}<1$, then
\begin{equation}
 \operatorname{Can}_{\beta'}(p,q)=0
 \label{eq:v14-small-angle-cancel}
\end{equation}
for every sufficiently small fixed loss $\beta'<\beta$.  The path and its fixed
source derivatives are given explicitly by the direct rotation of
Lemma~\ref{lem:v14-direct-rotation}.
\end{corollary}

\begin{proof}
The canonical unitary $u_{q,p}$ lies in
$\mathcal U_0(\mathfrak A_{\beta'})$ and conjugates $p$ to $q$.
Apply Definition~\ref{def:v14-unstable-obstruction}.
\end{proof}

The second criterion is the unstabilized analogue of the V13 spectator
rotation, but with the reserve already present in the original carrier.

\begin{definition}[Controlled internal buffer]
\label{def:v14-buffer}
A projection $r\in\mathfrak A_\beta$ is a \emph{controlled internal buffer}
for $p,q$ if
\begin{equation}
 rp=rq=0,
 \label{eq:v14-buffer-orthogonal}
\end{equation}
and there are controlled partial isometries $a,b$ such that
\begin{equation}
 a^*a=p,
 \quad aa^*=r,
 \qquad
 b^*b=q,
 \quad bb^*=r.
 \label{eq:v14-buffer-isometries}
\end{equation}
The buffer is required to carry the same representation/source grading as the
zero-mode packet but may be a gapped massive subpacket of the original
microscopic carrier.
\end{definition}

\begin{lemma}[Exact controlled rotation through an internal buffer]
\label{lem:v14-buffer-rotation}
Let $p,r$ be orthogonal projections and let
$a^*a=p$, $aa^*=r$.  Then
\begin{equation}
 U_a(\theta)
 :=I-p-r+\cos\theta\,(p+r)+\sin\theta\,(a-a^*),
 \qquad 0\le\theta\le\frac\pi2,
 \label{eq:v14-buffer-U}
\end{equation}
is a controlled unitary path with
\begin{equation}
 U_a(0)=I,
 \qquad
 U_a(\pi/2)pU_a(\pi/2)^*=r.
 \label{eq:v14-buffer-endpoint}
\end{equation}
\end{lemma}

\begin{proof}
On the orthogonal sum $p\mathcal H\oplus r\mathcal H$, set
$J=a-a^*$.  Since $a^2=(a^*)^2=0$ and
$a^*a=p$, $aa^*=r$,
\[
 J^*=-J,
 \qquad
 J^2=-(p+r).
\]
Hence the restriction of \eqref{eq:v14-buffer-U} to this sum is
$e^{\theta J}$, while the operator is the identity on the orthogonal
complement.  It is therefore unitary and continuous in the same controlled
algebra.  Moreover $Jp=a$, so at $\theta=\pi/2$ the $p$-range is mapped onto
the $r$-range.
\end{proof}

\begin{theorem}[Internal-buffer cancellation theorem]
\label{thm:v14-buffer-cancellation}
If $p,q$ admit a controlled internal buffer $r$, then
\begin{equation}
 \boxed{\operatorname{Can}_\beta(p,q)=0.}
 \label{eq:v14-buffer-cancel}
\end{equation}
More precisely, concatenate the path moving $p$ to $r$ by
$U_a(\theta)$ with the reverse of the path moving $q$ to $r$ by
$U_b(\theta)$.  This gives an explicit controlled projection path from $p$ to
$q$ entirely inside the original carrier.
\end{theorem}

\begin{proof}
Lemma~\ref{lem:v14-buffer-rotation} gives a controlled path
$p\rightsquigarrow r$.  Applying the same lemma to $b$ gives a path
$q\rightsquigarrow r$; reversing it gives $r\rightsquigarrow q$.  Concatenation
proves a controlled projection path, so
Theorem~\ref{thm:v14-projection-U0} gives
$\operatorname{Can}_\beta(p,q)=0$.
\end{proof}

\begin{corollary}[When the V13 spectator is already internal]
\label{cor:v14-internal-spectator}
Suppose the original Standard-Model carrier contains, on the entire protected
background/source panel, a source-smooth exponentially local massive subpacket
$r$ of rank $|k|$ which is orthogonal to both endpoint zero-mode packets and is
controlled-equivalent to each of them in the appropriate chiral sector.  Then
the V13 stable rotation can be carried out inside the original carrier; no
external spectator copy is required.
\end{corollary}

\subsection{Unstabilized native lifting from the controlled $U_0$ criterion}
\label{sec:v14-native-lift}

The V13 controlled-reduction theorem already proved that the nonzero-mode graph
coordinate carries no additional phase invariant.  We can therefore promote
the exact projection criterion to the full native regulator.

\begin{theorem}[Unstabilized fixed-index native lifting]
\label{thm:v14-unstable-native-lift}
Let $A,B$ be two V11 protected minimal fixed-index native regulators with the
same physical principal Dirac--Yukawa jet, Standard-Model representation,
positive IR/reference prescription and global complex/Pfaffian line datum.
Let $p_A,p_B$ denote their kernel projectors for $k>0$ or cokernel projectors
for $k<0$.  Assume that, in the corresponding equivariant source-differentiable
local algebra,
\begin{equation}
 \boxed{
 p_B=u p_Au^*
 \quad\text{for some}\quad
 u\in\mathcal U_0(\mathfrak A_\beta).}
 \label{eq:v14-native-U0-hyp}
\end{equation}
Then $A$ and $B$ are joined \emph{without stabilization} by a finite
concatenation of V11-admissible native fixed-index homotopies after the usual
one-time shrinking of the compact protected source chart.
\end{theorem}

\begin{proof}
Choose a controlled unitary path $u_t$ from $I$ to $u$ and put
$p_t=u_tp_Au_t^*$.  This is a controlled zero-mode projector path with all
prescribed source derivatives.  Apply the constructive part of
Theorem~\ref{thm:v13-controlled-reduction}: for $k>0$ transport the initial
local graph coordinate by $u_t$ in $\mathcal H_-$, thereby obtaining a
minimal-stratum projector path whose kernel is $p_t$; when the target kernel is
reached, linearly interpolate the transported affine graph coordinate to the
target graph coordinate.  For $k<0$ use the chiral-dual construction in
$\mathcal H_+$.  Propositions~\ref{prop:v13-positive-local-graph} and
\ref{prop:v13-negative-local-graph} keep every graph coordinate in the local
algebra and give the explicit positive reduced kinetic floor.

As in V13, the zero-mode unitary transport is a controlled compact/soft
deformation and the physical principal quotient is frozen.  The endpoint graph
coordinates have the same prescribed principal jet, so the final affine
interpolation is taken inside that fixed-jet fibre.  The compact path and the
open reduced/Yukawa pivot conditions permit one finite shrinking of the source
chart to keep the full reduced singular floor positive.  Kato exterior powers
carry the global line/saturation packet.  Every segment therefore satisfies
Definition~\ref{def:v11-index-admissible} without adding a new carrier copy.
\end{proof}

\begin{theorem}[Unstabilized global fixed-index regulator universality]
\label{thm:v14-unstable-universality}
Under the hypotheses of Theorem~\ref{thm:v14-unstable-native-lift}, for every
fixed coefficient packet $\mathfrak q$ there is an invertible finite local
continuum matching map
\begin{equation}
 F_{\infty,\mathfrak q,k}^{B\leftarrow A}
 \label{eq:v14-unstable-map}
\end{equation}
such that
\begin{equation}
 \boxed{
 \widehat{\mathfrak C}^{B,(k),\rm sat}_{\infty,\lambda}
 =F_{\infty,\mathfrak q,k}^{B\leftarrow A}
  \widehat{\mathfrak C}^{A,(k),\rm sat}_{\infty,\lambda}.}
 \label{eq:v14-unstable-universality}
\end{equation}
No spectator is introduced at any stage.  The equality contains the complete
normalized connected/source packet, operator mixing, stress/contact terms,
ghosts/antifields, BV/BRST/master relations and determinant/Pfaffian insertions,
and remains strictly projective in perturbative order.
\end{theorem}

\begin{proof}
Apply V11 fixed-index universality to the unstabilized native path of
Theorem~\ref{thm:v14-unstable-native-lift}.  The V12 global line transport
supplies the saturated complex comparison on the faithful branch, with the
same separate real-orientation hypothesis where applicable.  The V7
chronological construction gives exact projectivity in the terminal loop order.
\end{proof}

\begin{corollary}[Small-angle and internal-buffer universality]
\label{cor:v14-positive-universality}
The conclusion of Theorem~\ref{thm:v14-unstable-universality} holds, without
stabilization, whenever either:
\begin{enumerate}[label=\textup{(P\arabic*)},leftmargin=3em]
\item the two endpoint zero-mode projectors have operator-norm distance less
than one; or
\item they admit a controlled internal buffer in the sense of
Definition~\ref{def:v14-buffer}.
\end{enumerate}
\end{corollary}

\begin{proof}
Use Corollary~\ref{cor:v14-small-angle} or
Theorem~\ref{thm:v14-buffer-cancellation} to verify
\eqref{eq:v14-native-U0-hyp}, then apply
Theorem~\ref{thm:v14-unstable-universality}.
\end{proof}

\subsection{Stable range removes the purely finite-dimensional complement defect}
\label{sec:v14-stable-range}

V12 already proved that on a finite parameter panel and sufficiently large
ambient rank, ordinary bundle isomorphism gives a finite-Grassmannian projector
homotopy.  Combined with Theorem~\ref{thm:v14-projection-U0}, this has a useful
interpretation for V14.

\begin{proposition}[Onsite stable-range cancellation]
\label{prop:v14-onsite-stable-range}
Let $B$ be a finite CW complex of dimension $d$, and consider the propagation-zero
algebra $C(B,M_m)$ (or its smooth analogue on a compact manifold).  Let $p,q$
be rank-$r$ projections whose range bundles are isomorphic.  If
\begin{equation}
 d\le2(m-r),
 \label{eq:v14-onsite-stable-range}
\end{equation}
then
\begin{equation}
 q=upu^*
 \end{equation}
for some $u$ in the identity component of the onsite unitary group.
Consequently the purely finite-dimensional complement defect and unitary
component obstruction vanish in this stable range.
\end{proposition}

\begin{proof}
The V12 finite-Grassmannian stable-range theorem says that the two classifying
maps $B\to\operatorname{Gr}(r,\mathbb C^m)$ are homotopic when their vector
bundles are isomorphic under \eqref{eq:v14-onsite-stable-range}.  This is a
path of onsite projections.  Apply
Theorem~\ref{thm:v14-projection-U0} in the onsite algebra to obtain the desired
identity-component unitary.
\end{proof}

\begin{corollary}[What remains at large cutoff on a fixed finite panel]
\label{cor:v14-large-cutoff-boundary}
Fix a finite-dimensional background/source panel and a fixed zero-mode rank.
When the cutoff makes the ambient chiral complement sufficiently large for
\eqref{eq:v14-onsite-stable-range}, an obstruction to unstabilized native
lifting cannot be attributed merely to an unstable finite-dimensional bundle
complement.  Any remaining obstruction must use the genuinely controlled
spacetime-local algebra (or the independent real Pfaffian orientation data).
\end{corollary}

\subsection{The exact relationship between V13 stability and V14 cancellation}
\label{sec:v14-stable-vs-unstable}

The previous results identify precisely what the inert V13 copy accomplishes.

\begin{theorem}[Stable versus unstabilized controlled equivalence]
\label{thm:v14-stable-unstable}
For protected zero-mode projections $p,q$:
\begin{enumerate}[label=\textup{(S\arabic*)},leftmargin=3em]
\item an unstabilized native projection path implies controlled equivalence
$p\sim_{\rm ctl}q$;
\item controlled equivalence always gives a stabilized controlled projection
path by the V13 two-copy rotation;
\item removing the stabilization is equivalent to
$\operatorname{Can}_\beta(p,q)=0$;
\item the obstruction $\operatorname{Can}_\beta$ can be analyzed in two stages:
first the complement defect $\mathfrak c_\beta$, then, when it vanishes, the
relative controlled unitary-component class of
Definition~\ref{def:v14-kappa}.
\end{enumerate}
Thus V13 already proves continuum universality at the stable controlled level,
while V14 identifies the exact extra datum required for a microscopic path in
the original carrier.
\end{theorem}

\begin{proof}
Item (S1) is V13 Kato transport.  Item (S2) is
Lemma~\ref{lem:v13-spectator-rotation}.  Item (S3) is
Theorem~\ref{thm:v14-projection-U0}.  The two-stage statement is
Theorem~\ref{thm:v14-unitary-extension} and
Proposition~\ref{prop:v14-kappa-criterion}.
\end{proof}

\begin{center}
\fbox{\parbox{0.92\textwidth}{
\textbf{The universal controlled-cancellation conjecture is false without
additional carrier hypotheses.}
V13's controlled partial isometry is a stable datum.  An unstabilized native
path exists exactly when the two zero-mode projections belong to the same
identity-component orbit of the controlled local unitary group.  Small-angle
pairs and pairs sharing a controlled internal buffer satisfy this criterion
constructively.  A propagation-zero unstable bundle example shows that no
proof can follow from controlled Murray--von Neumann equivalence alone.}}
\end{center}

\section{V14 anti-circularity audit}
\label{sec:v14-audit}

\begin{enumerate}[leftmargin=2.4em]
\item \textbf{Stable equivalence was not silently called cancellation.}
Theorem~\ref{thm:v14-unitary-extension} exhibits the missing complement
partial isometry explicitly.
\item \textbf{A unitary extension was not silently called a path.}
Theorem~\ref{thm:v14-projection-U0} isolates the identity-component condition.
\item \textbf{The obstruction is not inferred only from abstract $K_1$.}
The primary object is the actual controlled $U_0$ orbit.  The complement
semigroup defect and the relative unitary-component coset are exact diagnostics
of that orbit problem.
\item \textbf{The counterexample is local.}
All matrices in Countertheorem~\ref{cth:v14-unstable-cancellation} have zero
spatial propagation.  Its failure is unstable complement topology, not a hidden
long-range operator.
\item \textbf{The counterexample is not overpromoted to the high-cutoff
physical branch.}
It is stably trivial, and Proposition~\ref{prop:v14-onsite-stable-range} shows
that the purely finite-dimensional obstruction disappears in the appropriate
large ambient stable range.
\item \textbf{Positive cancellation criteria are constructive.}
The small-angle unitary is an explicit polar direct rotation; the internal
buffer theorem gives an explicit exponential rotation in the original carrier.
\item \textbf{The V11 hard-shell theorem is invoked only after an actual native
path has been constructed.}
Theorem~\ref{thm:v14-unstable-native-lift} uses the $U_0$ path first and only
then applies the inherited fixed-index shell/universality estimates.
\end{enumerate}

\section{V14 status ledger}
\label{sec:v14-status}

\subsection*{Proved in V14}
\begin{enumerate}[leftmargin=2.2em]
\item A controlled partial isometry $v:p\to q$ extends to a controlled unitary
on the original carrier if and only if the complementary projections
$I-p$ and $I-q$ are themselves controlled Murray--von Neumann equivalent.  An
extension is exactly $u=v+w$.
\item For a fixed $v$, once complement cancellation holds, the remaining
extension obstruction is the relative component coset
$\kappa_\beta(v)$; it vanishes exactly when $v$ has an extension in the
identity component of the controlled unitary group.
\item Two controlled projections are joined by a controlled projection path if
and only if they are conjugate by a unitary in
$\mathcal U_0(\mathfrak A_\beta)$.  The proof uses a finite subdivision and an
explicit local direct-rotation formula for pairs at norm distance less than
one.
\item Controlled zero-mode equivalence alone does not imply unstabilized
cancellation in an arbitrary local coefficient algebra.  A propagation-zero
$S^6$ clutching example gives equivalent zero-mode ranges with nonisomorphic
complements and therefore no ambient unitary extension.
\item The unstable obstruction vanishes constructively for small-angle
zero-mode packets and for packets sharing a controlled internal buffer already
inside the original carrier.
\item The controlled $U_0$ criterion lifts through the V13 affine reduced graph
coordinate to a complete unstabilized V11-admissible native fixed-index path.
Hence regulators satisfying this criterion have the same matched global
coefficientwise continuum without any spectator stabilization.
\item On a finite-dimensional parameter panel in the V12 stable ambient range,
purely onsite bundle isomorphism already implies identity-component unitary
conjugacy.  Any remaining high-cutoff obstruction must therefore be genuinely
controlled/spatial (or the separate real Pfaffian orientation class).
\item V13 stable controlled equivalence and V14 unstabilized equivalence differ
precisely by the cancellation obstruction
$\operatorname{Can}_\beta$.
\end{enumerate}

\subsection*{Inherited}
\begin{enumerate}[leftmargin=2.2em]
\item V1--V7: source-complete coefficientwise continuum, common BV restoration,
arbitrary-fixed-loop forest theorem, ultraviolet Cauchy convergence,
scheme/reference covariance and exact perturbative-order projectivity.
\item V8--V10: relative regulator forests, the two-scale universality estimate,
and regulator-class lifting on the protected index-zero branches.
\item V11--V12: minimal fixed-index geometry, reduced pseudoinverse/source
calculus, global zero-mode bundle/line transport and global fixed-index
universality inside one protected native component.
\item V13: exponential locality of protected Riesz zero modes, the controlled
zero-mode module invariant, local reconstruction of the affine reduced graph,
and stable native lifting/universality after one inert massive copy.
\end{enumerate}

\subsection*{Conditional}
\begin{enumerate}[leftmargin=2.2em]
\item The exact cancellation criteria are statements in the chosen weighted
local algebra.  A loss of the V13 exponential-resolvent bound or reduced
singular-value floor can change the controlled component problem.
\item Countertheorem~\ref{cth:v14-unstable-cancellation} proves that universal
cancellation is false in general, but its unstable finite-dimensional defect
may disappear in the large-cutoff stable range of a fixed physical panel.
\item The internal-buffer theorem is load-bearing only when such a buffer is
actually present in the original Standard-Model representation/source carrier;
V14 does not assume one exists universally.
\item The V12 real Pfaffian mod-two orientation condition remains independent.
\item All continuum statements remain coefficientwise at prescribed finite
$(L,\Delta,r,N)$ and positive IR/reference scale.
\end{enumerate}

\subsection*{Open beyond V14}
\begin{enumerate}[leftmargin=2.2em]
\item Prove or falsify \emph{automatic controlled cancellation in the actual
Renewal Standard-Model carrier}: show that the protected reduced hard
complement supplies a uniformly local internal buffer (or an equivalent
stable-rank/cancellation property) for every fixed finite source/background
panel at sufficiently large cutoff.
\item If automatic cancellation fails, construct a genuine physical witness of
$\operatorname{Can}_\beta\ne0$ in the native spacetime-local algebra, not only
an unstable finite-dimensional coefficient example.
\item Classify the controlled unitary component/corner coset on noncontractible
gauge-background components and relate it, where possible, to a computable
local/coarse $K_1$ or transport index.
\item Compute the V12 real Pfaffian mod-two spectral-flow character on any
independently physical Majorana sector.
\item Analyze exits from the minimal fixed-index stratum and the independent
positive-IR/nonperturbative problems retained by V13.
\end{enumerate}

\subsection*{Exact next load-bearing theorem}
\begin{center}
\fbox{\parbox{0.92\textwidth}{
\textbf{Automatic internal cancellation in the Renewal ESM carrier.}
For each prescribed finite coefficient/source/background panel and fixed index
$k$, analyze the actual protected reduced Standard-Model chiral carrier at large
cutoff.  Prove that it contains a source-smooth exponentially local internal
buffer of rank $|k|$ (or prove an equivalent cancellation/stable-rank theorem
for the relevant local operator corner) so that every V13 controlled zero-mode
equivalence has $\operatorname{Can}_\beta=0$ without adding a spectator.
Otherwise produce an explicit native controlled complement/unitary-component
witness.  A positive theorem, combined with
Theorem~\ref{thm:v14-unstable-universality}, would close the distinction between
stable and unstabilized global fixed-index regulator universality.}}
\end{center}

\section*{Append-only continuation marker: V14}
\addcontentsline{toc}{section}{Append-only continuation marker: V14}
\textbf{End of V14.}  The complete V1--V13 body above is retained verbatim.
V14 proves that controlled zero-mode equivalence is a stable
Murray--von Neumann statement and isolates the exact extra data needed for an
unstabilized microscopic path: controlled complement cancellation followed by
the identity-component condition in the local unitary group.  It proves the
exact $U_0$-orbit criterion for controlled projection homotopy, gives a
propagation-zero unstable cancellation counterexample, and supplies explicit
small-angle and internal-buffer cancellation constructions.  Whenever the
$U_0$ criterion holds, the V13 affine graph reduction produces an unstabilized
V11-admissible native path and hence global fixed-index regulator universality
without a spectator.  The next gate is to decide whether the actual protected
Renewal ESM carrier automatically supplies the required internal cancellation
at large cutoff.



\clearpage
\section{V15 continuation: finite-rank localization, coefficientwise internal reserve, and automatic unstabilized lifting}
\label{sec:v15-continuation}

V14 identified the exact distinction between stable and unstabilized native
phase lifting.  A controlled equivalence of two zero-mode projections need not
cancel in an arbitrary local coefficient algebra; an unstabilized path requires
membership in the same controlled $\mathcal U_0$ orbit.  The V14 handoff asked
whether the actual Renewal ESM carrier supplies the missing reserve
automatically at sufficiently large cutoff.

There are two reasons not to attack this by a blanket stable-rank assertion.
First, the microscopic spacetime algebra is four-dimensional and contains far
more projections than the fixed-rank protected zero-mode packet.  Cancellation
of the entire local operator algebra is therefore a much stronger question than
the one needed here.  Second, the continuum theorem is coefficientwise: for a
prescribed finite source/contact order one needs the corresponding finite source
jet of the regulator path, not a single cutoff-independent finite-amplitude
source ball.  The same quantifier correction already used in V2 is load-bearing
again.

The present continuation proves a positive theorem at exactly that scope.  A
rank-$\kappa$ protected Riesz projector with a cutoff-uniform exponential
kernel bound has a uniformly localized orthonormal frame.  A finite bank of
bounded gauge/spin-covariant lattice translations of that frame supplies, at
large enough cutoff, a finite local candidate space of arbitrarily large fixed
multiplicity.  Because the retained source algebra is finite, one can choose a
rank-$\kappa$ subpacket which is orthogonal to \emph{every} source jet of both
endpoint zero-mode packets through the requested order.  In the truncated
source algebra this is an exact V14 internal buffer.  Its additional translation distance from the original localized zero-mode
centres is bounded in lattice units independently of the cutoff and therefore
has physical transport scale $O(a)$.

Thus the V13 spectator is unnecessary for every fixed coefficient/source
packet on a protected local chart.  This is deliberately not advertised as
cancellation of the full four-dimensional Roe-type algebra, nor as a global
finite-amplitude path over an arbitrary noncontractible gauge component.

Throughout V15 let $X_a$ be the periodic regulator lattice at spacing $a$,
with its graph distance $d_a$ measured in lattice units and a fixed finite
onsite fermion fibre $\mathbb C^{d_f}$ in the relevant chiral/representation
block.  Bounded geometry means that for every fixed integer $R$,
\begin{equation}
 N_R:=\sup_{a,x}|B_R(x)|<\infty.
 \label{eq:v15-bounded-geometry}
\end{equation}
The V13 weighted kernel algebra at exponent $\alpha>0$ is denoted
$\mathfrak A_{\alpha,a}$.

\subsection{A finite-rank exponentially local projector has a localized frame}
\label{sec:v15-localized-frame}

We begin with a quantitative statement which is independent of the overlap
specialization.

\begin{lemma}[Finite-propagation truncation]
\label{lem:v15-truncation}
Let $A\in\mathfrak A_{\alpha,a}$ and define
\begin{equation}
 A^{(R)}_{xy}:=\mathbf 1_{\{d_a(x,y)\le R\}}A_{xy}.
 \label{eq:v15-truncate}
\end{equation}
Then
\begin{equation}
 \boxed{
 \|A-A^{(R)}\|_{\rm op}
 \le e^{-\alpha R}\|A\|_\alpha .}
 \label{eq:v15-truncation-bound}
\end{equation}
The same estimate holds coefficientwise for every retained source derivative.
\end{lemma}

\begin{proof}
For every row,
\[
 \sum_{d_a(x,y)>R}\|A_{xy}\|
 \le e^{-\alpha R}
      \sum_y e^{\alpha d_a(x,y)}\|A_{xy}\|
 \le e^{-\alpha R}\|A\|_\alpha,
\]
and the same estimate holds for columns.  The Schur bound gives
\eqref{eq:v15-truncation-bound}.  Source derivatives are separate elements of
the same weighted algebra and obey the identical argument.
\end{proof}

\begin{lemma}[Uniform diagonal floor for a finite-rank local projection]
\label{lem:v15-diagonal-floor}
Let $p=p^*=p^2\in\mathfrak A_{\alpha,a}$ be nonzero and assume
\begin{equation}
 \|p\|_\alpha\le C_p
 \label{eq:v15-p-local-bound}
\end{equation}
uniformly in the cutoff.  Choose $R$ so that
$C_pe^{-\alpha R}\le1/4$.  Then
\begin{equation}
 \boxed{
 \max_x\|p_{xx}\|_{\rm op}
 \ge c_0:=\frac{3}{4N_R}.}
 \label{eq:v15-diagonal-floor}
\end{equation}
Consequently there are a site $x_0$ and a unit vector
$\xi\in\mathbb C^{d_f}$ such that
\begin{equation}
 \|p\iota_{x_0}\xi\|^2
 =\langle\xi,p_{x_0x_0}\xi\rangle
 \ge c_0,
 \label{eq:v15-column-floor}
\end{equation}
where $\iota_x$ is the canonical inclusion of the onsite fibre.
\end{lemma}

\begin{proof}
Lemma~\ref{lem:v15-truncation} gives
$\|p-p^{(R)}\|\le1/4$.  Since a nonzero projection has operator norm one,
\begin{equation}
 \|p^{(R)}\|\ge3/4.
 \label{eq:v15-pr-trunc-norm}
\end{equation}
Put $M:=\max_x\|p_{xx}\|$.  Positivity of the $2\times2$ block compression of
$p$ to the onsite fibres at $x,y$ gives the operator Cauchy--Schwarz estimate
\begin{equation}
 \|p_{xy}\|^2\le\|p_{xx}\|\,\|p_{yy}\|\le M^2.
 \label{eq:v15-block-CS}
\end{equation}
Every row and column of $p^{(R)}$ has at most $N_R$ nonzero blocks, hence its
Schur norm is at most $N_RM$.  Therefore
$3/4\le\|p^{(R)}\|\le N_RM$, proving
\eqref{eq:v15-diagonal-floor}.  Taking a top eigenvector of the positive block
$p_{x_0x_0}$ gives \eqref{eq:v15-column-floor}.
\end{proof}

\begin{theorem}[Uniformly localized orthonormal frame]
\label{thm:v15-localized-frame}
Fix a rank $\kappa<\infty$.  Let
$p_a\in\mathfrak A_{\alpha,a}$ be rank-$\kappa$ projections with
$\|p_a\|_\alpha\le C_p$ uniformly in the cutoff.  Then there exist an exponent
$0<\beta<\alpha$, constants $C_{\rm fr},c_{\rm fr}>0$ depending only on
$(\kappa,\alpha,C_p,d_f)$ and the bounded-geometry constants, sites
$x_{1,a},\ldots,x_{\kappa,a}$, and an orthonormal basis
$\psi_{1,a},\ldots,\psi_{\kappa,a}$ of $\Ran p_a$ such that
\begin{equation}
 \boxed{
 \sum_x e^{\beta d_a(x,x_{j,a})}\|\psi_{j,a}(x)\|
 \le C_{\rm fr}}
 \label{eq:v15-frame-l1}
\end{equation}
and therefore
\begin{equation}
 \|\psi_{j,a}(x)\|
 \le C_{\rm fr}e^{-\beta d_a(x,x_{j,a})}.
 \label{eq:v15-frame-pointwise}
\end{equation}
All constants are independent of the volume and cutoff.

If $p_a(\zeta)$ lies on one star-shaped V11 protected source/background chart
and its prescribed source derivatives are uniformly local, the reference frame
may be transported by the V13 local Kato equation.  For every fixed source
multi-index $I$ in the packet,
\begin{equation}
 \sum_x e^{\beta' d_a(x,x_{j,a})}
 \|\partial_I\psi_{j,a}(0;x)\|
 \le C_I
 \label{eq:v15-frame-source-jets}
\end{equation}
for one fixed $0<\beta'<\beta$.
\end{theorem}

\begin{proof}
Apply Lemma~\ref{lem:v15-diagonal-floor} to $p_a$.  Choose $x_{1,a},\xi_1$
as there and put
\begin{equation}
 \psi_{1,a}:=
 \frac{p_a\iota_{x_{1,a}}\xi_1}
      {\|p_a\iota_{x_{1,a}}\xi_1\|}.
 \label{eq:v15-first-frame-vector}
\end{equation}
The weighted column bound for $p_a$ and the denominator floor give
\eqref{eq:v15-frame-l1} for $j=1$ at exponent $\alpha$.
The rank-one projection
$e_{1,a}=|\psi_{1,a}\rangle\langle\psi_{1,a}|$ therefore belongs to
$\mathfrak A_{\beta_1,a}$, for any fixed $\beta_1<\alpha$, with a uniform
bound: use
$d_a(x,y)\le d_a(x,x_{1,a})+d_a(y,x_{1,a})$ and
\eqref{eq:v15-frame-l1}.

Because $\psi_{1,a}\in\Ran p_a$, one has $e_{1,a}\le p_a$ and
$p_a-e_{1,a}$ is a projection of rank $\kappa-1$.  Its weighted norm at
$\beta_1$ is bounded by a constant depending only on the preceding uniform
bounds.  Apply Lemma~\ref{lem:v15-diagonal-floor} again to this residual
projection.  Iterating $\kappa$ times and choosing in advance a decreasing
finite exponent chain
\[
 \alpha>\beta_1>\cdots>\beta_\kappa=:\beta>0
\]
produces the claimed orthonormal localized frame with uniform constants.  The
rank is fixed, so no constant depends on the cutoff.

For the source statement, use this frame at the reference point and solve the
V13 Kato transport along the radial path in the star-shaped source chart.  V13
proved that the Kato generator and all prescribed fixed derivatives belong to a
weighted algebra after an arbitrarily small exponent loss.  Acting with a
weighted local operator on a vector satisfying \eqref{eq:v15-frame-l1}
preserves an exponential moment about the same centre.  Finite differentiation
of the Kato ODE therefore proves \eqref{eq:v15-frame-source-jets}.
\end{proof}

\subsection{Bounded covariant translations provide arbitrarily many local copies}
\label{sec:v15-translation-bank}

The V2 native lattice already contains gauge/spin parallel shifts.  For a fixed
integer displacement $s\in\mathbb Z^4$, choose one bounded path from $x$ to
$x+s$ and denote the corresponding unitary parallel shift by
$\mathsf T_s^{\mathcal U}$.  Its propagation is at most a constant multiple of
$|s|_1$, it is exactly gauge/spin covariant, and every prescribed finite source
derivative is finite range.

\begin{lemma}[Overlap decay of translated localized vectors]
\label{lem:v15-shift-overlap}
Let $\psi,\phi$ satisfy exponential bounds of the form
\eqref{eq:v15-frame-pointwise} about centres $x_\psi,x_\phi$.  For every
$0<\beta_0<\beta$ there is $C_{\beta_0}$ such that
\begin{equation}
 \boxed{
 |\langle\mathsf T_s^{\mathcal U}\psi,
          \mathsf T_t^{\mathcal U}\phi\rangle|
 \le C_{\beta_0}
 e^{-\beta_0 d_a(x_\psi+s,x_\phi+t)}.}
 \label{eq:v15-shift-overlap}
\end{equation}
The bound is uniform in the cutoff and in the unitary parallel transports.
\end{lemma}

\begin{proof}
Parallel transport does not change pointwise norms.  Hence the absolute value
of the inner product is bounded by a convolution of two exponentials centred
at $x_\psi+s$ and $x_\phi+t$.  The triangle inequality gives an overall factor
$e^{-\beta_0d_a(x_\psi+s,x_\phi+t)}$ after reducing the exponent from $\beta$
to $\beta_0$; the remaining lattice convolution is uniformly summable by
bounded geometry.
\end{proof}

\begin{lemma}[Finite bounded-displacement separation]
\label{lem:v15-shift-selection}
Fix integers $\kappa,M\ge1$ and a separation radius $D$.  For every sufficiently
large periodic lattice and every set of centres
$X=\{x_1,\ldots,x_\kappa\}$ there exist displacements
\begin{equation}
 s_1,\ldots,s_M\in\mathbb Z^4,
 \qquad
 |s_j|_1\le S_*(\kappa,M,D),
 \label{eq:v15-shifts}
\end{equation}
with $S_*$ independent of the cutoff and of $X$, such that
\begin{equation}
 d_a(x_b+s_i,x_c+s_j)\ge D
 \quad
 (i\\ne j,\\ 1\\le b,c\\le\\kappa).
 \label{eq:v15-shift-separated}
\end{equation}
\end{lemma}

\begin{proof}
Choose an integer $L>2D$ and consider the candidate displacements
$nLe_1$, $1\le n\le N$.  For a fixed candidate $n$ and a fixed ordered pair
$(b,c)$, there is at most one other candidate $m$ for which
$x_b+nLe_1$ and $x_c+mLe_1$ are within distance $D$, provided the torus side
length exceeds $2NL+2D$.  Hence the conflict graph on the $N$ candidates has
maximum degree at most $2\kappa^2$.  Taking
$N\ge M(2\kappa^2+1)$, the greedy independent-set bound gives $M$ candidates
with no conflicts.  Their maximum displacement is
$S_*:=NL$.  All these numbers depend only on $(\kappa,M,D)$.
\end{proof}

\begin{corollary}[A uniformly conditioned translated-copy bank]
\label{cor:v15-copy-bank}
Fix $M$.  For all sufficiently large cutoffs, the displacements in
Lemma~\ref{lem:v15-shift-selection} may be chosen so that the $M\kappa$ vectors
\begin{equation}
 \mathsf T_{s_j}^{\mathcal U}\psi_b,
 \qquad
 1\le j\le M,\ 1\le b\le\kappa,
 \label{eq:v15-copy-vectors}
\end{equation}
have Gram matrix $G_0$ satisfying
\begin{equation}
 \boxed{\frac12I\preceq G_0\preceq\frac32I.}
 \label{eq:v15-copy-Gram}
\end{equation}
The raw synthesis operator of these vectors and its orthonormalization
$G_0^{-1/2}$ are uniformly exponentially local with propagation/locality
constants independent of the cutoff.
\end{corollary}

\begin{proof}
Choose $D$ in Lemma~\ref{lem:v15-shift-selection} so large that the right-hand
side of \eqref{eq:v15-shift-overlap} is less than
$[2M\kappa]^{-1}$ whenever the two copy labels differ.  Within one translated
copy the vectors remain exactly orthonormal because
$\mathsf T_s^{\mathcal U}$ is unitary.  The row-sum estimate for the off-diagonal
part of the Gram matrix is therefore less than $1/2$, proving
\eqref{eq:v15-copy-Gram}.  The synthesis operator is a finite sum of bounded
propagation shifts acting on the localized frame.  Functional calculus of the
finite Gram matrix preserves all uniform fixed-source bounds.
\end{proof}

\subsection{The finite source-jet algebra and the exact annihilator count}
\label{sec:v15-source-jet}

Let $\mathbf R_\varrho$ denote the complete finite graded
source/BV algebra retained by the prescribed packet through total source degree
$\varrho$; it includes the even central writers and the odd ghost/antifield
writers already used in V4--V7.  Only its finite vector-space dimension is
load bearing:
\begin{equation}
 \boxed{\nu_\varrho:=\dim_{\mathbb C}\mathbf R_\varrho<\infty.}
 \label{eq:v15-source-monomials}
\end{equation}
For $s$ commuting even coordinates alone this specializes to
$\mathbf R_\varrho=\mathbb C[J^1,\ldots,J^s]/(J)^{\varrho+1}$ and
$\nu_\varrho=\binom{s+\varrho}{\varrho}$.  With odd writers present, the
ordinary monomial basis is replaced by the corresponding finite graded basis;
the argument below uses only the number $\nu_\varrho$ and the fact that
multiplication is truncated at total degree $\varrho$.

Fix two rank-$\kappa$ protected zero-mode packets $p(J),q(J)$ in this truncated
algebra and a V13 controlled equivalence
\begin{equation}
 v(J)^*v(J)=p(J),
 \qquad
 v(J)v(J)^*=q(J).
 \label{eq:v15-controlled-v}
\end{equation}
Choose the V15 localized Kato frame synthesis map
\begin{equation}
 \Psi(J):\mathbb C^\kappa\longrightarrow\mathcal H_a,
 \qquad
 \Psi(J)^*\Psi(J)=I,
 \qquad
 \Psi(J)\Psi(J)^*=p(J),
 \label{eq:v15-Psi}
\end{equation}
and put
\begin{equation}
 \mathsf X(J):=v(J)\Psi(J),
 \qquad
 \mathsf X(J)^*\mathsf X(J)=I,
 \qquad
 \mathsf X(J)\mathsf X(J)^*=q(J).
 \label{eq:v15-Chi}
\end{equation}
All identities are in $\mathbf R_\varrho$.

Choose
\begin{equation}
 M\ge2\nu_\varrho+1
 \label{eq:v15-M}
\end{equation}
and the bounded shifts of
Corollary~\ref{cor:v15-copy-bank}.  Let
$\widetilde W(J):\mathbb C^{M\kappa}\to\mathcal H_a$ be the raw synthesis
operator whose $j$th block is
$\mathsf T_{s_j}^{\mathcal U}(J)\Psi(J)$.  Since the constant coefficient of
its Gram matrix has the floor \eqref{eq:v15-copy-Gram}, its inverse square root
exists uniquely in the finite nilpotent extension of the positive constant
Gram matrix by analytic matrix functional calculus (equivalently, by the
finite recursive Sylvester equations for its source coefficients).
Define
\begin{equation}
 W(J):=\widetilde W(J)
 \bigl(\widetilde W(J)^*\widetilde W(J)\bigr)^{-1/2}.
 \label{eq:v15-W}
\end{equation}
Then
\begin{equation}
 \boxed{W(J)^*W(J)=I_{M\kappa}}
 \label{eq:v15-W-isometry}
\end{equation}
exactly in $\mathbf R_\varrho$.

\begin{theorem}[Finite-jet common-complement reserve]
\label{thm:v15-jet-reserve}
For every fixed finite source order $\varrho$, and every sufficiently large
cutoff, there is an isometry
\begin{equation}
 Z:\mathbb C^\kappa\longrightarrow\mathbb C^{M\kappa}
 \label{eq:v15-Z}
\end{equation}
with constant scalar coefficients such that, with
\begin{equation}
 \mathsf E(J):=W(J)Z,
 \qquad
 r(J):=\mathsf E(J)\mathsf E(J)^*,
 \label{eq:v15-Eta-r}
\end{equation}
one has in the truncated source algebra
\begin{align}
 \mathsf E(J)^*\mathsf E(J)&=I_\kappa,
 \label{eq:v15-Eta-isometry}\\
 p(J)r(J)&=r(J)p(J)=0,
 \label{eq:v15-pr-zero}\\
 q(J)r(J)&=r(J)q(J)=0.
 \label{eq:v15-qr-zero}
\end{align}
The projection $r(J)$ and every coefficient of its retained source jet are
uniformly exponentially local.  Its propagation/locality radius depends on
$(\kappa,\nu_\varrho)$ and the fixed locality constants, but not on the
cutoff.  Its added displacement from the original zero-mode localization
centres is bounded in lattice units, hence has physical transport scale $O(a)$.
\end{theorem}

\begin{proof}
Orthogonality to $p$ and $q$ is equivalent, respectively, to annihilation by
$\Psi^*$ and $\mathsf X^*$.  Define
\begin{equation}
 F(J):=
 \begin{pmatrix}
  \Psi(J)^*W(J)\\
  \mathsf X(J)^*W(J)
 \end{pmatrix}
 :\mathbb C^{M\kappa}\longrightarrow\mathbb C^{2\kappa}.
 \label{eq:v15-F}
\end{equation}
Choose once and for all a homogeneous vector-space basis
$(m_A)_{A=1}^{\nu_\varrho}$ of $\mathbf R_\varrho$ and write the finite graded
source expansion
\begin{equation}
 F=\sum_{A=1}^{\nu_\varrho}m_A F_A.
 \label{eq:v15-F-expansion}
\end{equation}
(The Koszul signs are already contained in the multiplication table of
$\mathbf R_\varrho$ and do not affect the following linear dimension count.)
Stack all coefficient maps:
\begin{equation}
 \mathbf F:
 \mathbb C^{M\kappa}
 \longrightarrow
 \bigoplus_{A=1}^{\nu_\varrho}\mathbb C^{2\kappa},
 \qquad
 \mathbf Fz=(F_Az)_A.
 \label{eq:v15-stacked-F}
\end{equation}
Pure dimension counting gives
\begin{equation}
 \dim\ker\mathbf F
 \ge M\kappa-2\kappa\nu_\varrho
 \ge\kappa
 \label{eq:v15-kernel-dim}
\end{equation}
by \eqref{eq:v15-M}.  Choose any orthonormal $\kappa$-frame in this kernel and
let $Z$ be its synthesis isometry.  Then every coefficient of $FZ$ in the retained graded basis vanishes; equivalently,
\begin{equation}
 FZ=0\quad\text{in }\mathbf R_\varrho.
 \label{eq:v15-FZ-zero}
\end{equation}
Since $W^*W=I$ and $Z^*Z=I$,
\eqref{eq:v15-Eta-isometry} follows.  Equation
\eqref{eq:v15-FZ-zero} says
$\Psi^*\mathsf E=0=\mathsf X^*\mathsf E$.  Multiplying by the frame synthesis maps gives
\eqref{eq:v15-pr-zero}--\eqref{eq:v15-qr-zero}.

Locality follows from Corollary~\ref{cor:v15-copy-bank}: $W$ is obtained from a
finite bank of fixed bounded-displacement native shifts and a finite Gram
inverse whose constant coefficient has a uniform floor.  Multiplication by the
fixed finite matrix $Z$ changes no spatial support estimate.  Every source
coefficient is a finite algebraic combination of the same bounded local jets.
Finally the largest displacement $S_*$ is independent of $a$, so its physical
length is $aS_*=O(a)$.
\end{proof}

\begin{remark}[Frame and gauge/spin covariance]
\label{rem:v15-frame-covariance}
The proof used one V11 Kato frame only to write the finite-dimensional
annihilator problem.  Under a change of zero-mode frame, the synthesis maps
$\Psi,\mathsf X$ and the translated-copy bank are right multiplied by the
corresponding finite unitary matrices.  Transporting the coefficient isometry
$Z$ by the same finite change leaves the physical projection $r=\mathsf E\mathsf E^*$
and the partial isometries below conjugated by the exact gauge/spin action.
Thus the construction defines a covariant buffer on the framed protected
chart; different local frame choices give V7/V12-equivalent auxiliary paths,
not different endpoint regulators.
\end{remark}

\subsection{Automatic coefficientwise internal buffer}
\label{sec:v15-buffer}

Define
\begin{equation}
 a(J):=\mathsf E(J)\Psi(J)^*,
 \qquad
 b(J):=\mathsf E(J)\mathsf X(J)^*.
 \label{eq:v15-buffer-partials}
\end{equation}

\begin{theorem}[Automatic internal reserve at fixed source order]
\label{thm:v15-automatic-buffer}
Under the hypotheses of Theorem~\ref{thm:v15-jet-reserve},
\begin{align}
 a^*a&=p,& aa^*&=r,
 \label{eq:v15-a-buffer}\\
 b^*b&=q,& bb^*&=r,
 \label{eq:v15-b-buffer}
\end{align}
exactly in $\mathbf R_\varrho$, while $r$ is orthogonal to both endpoint
zero-mode packets.  Hence $r$ is a V14 internal buffer of rank $\kappa$ for the
complete retained source jet.

Moreover, because $r$ lies in the orthogonal complement of the protected zero
modes, the V11 reduced singular-value floor makes it a nonzero/massive packet at
both endpoints.  No new fermion species or carrier copy has been added.
\end{theorem}

\begin{proof}
Use $\mathsf E^*\mathsf E=I$ and
$\Psi^*\Psi=\mathsf X^*\mathsf X=I$:
\[
 a^*a=\Psi\mathsf E^*\mathsf E\Psi^*=\Psi\Psi^*=p,
 \qquad
 aa^*=\mathsf E\Psi^*\Psi\mathsf E^*=\mathsf E\mathsf E^*=r,
\]
and similarly for $b$.  Orthogonality is
Theorem~\ref{thm:v15-jet-reserve}.  On a V11 protected zero-mode chart,
$D^*D\succeq\tau_*^2$ on the orthogonal complement of the kernel (and likewise
on the cokernel side).  Thus the buffer is contained in the reduced nonzero
packet at each endpoint.  It is an internal direction of the existing carrier.
\end{proof}

The V14 exponential rotation can now be performed coefficientwise.  Put
\begin{equation}
 U_a(\theta)
 =I-p-r+\cos\theta\,(p+r)+\sin\theta\,(a-a^*),
 \label{eq:v15-Ua}
\end{equation}
with the analogous $U_b$.  Because
\eqref{eq:v15-a-buffer} and the orthogonality relations are exact identities in
$\mathbf R_\varrho$, the ordinary V14 proof gives
\begin{equation}
 U_a(\theta)^*U_a(\theta)=I,
 \qquad
 U_a(\pi/2)pU_a(\pi/2)^*=r
 \label{eq:v15-Ua-exact}
\end{equation}
inside that same truncated algebra.  Concatenating with the reversed $q$ to
$r$ rotation gives an exact source-jet projection path from $p$ to $q$.

\begin{corollary}[Coefficientwise cancellation obstruction vanishes at large cutoff]
\label{cor:v15-Can-zero}
For a prescribed finite source order $\varrho$, define
$\operatorname{Can}^{[\varrho]}_\beta(p,q)$ by the V14 $\mathcal U_0$ criterion
inside the truncated source algebra.  Then for every V13 controlled-equivalent
rank-$\kappa$ protected pair there is $a_0>0$ such that
\begin{equation}
 \boxed{
 \operatorname{Can}^{[\varrho]}_\beta(p_a,q_a)=0
 \qquad(0<a<a_0).}
 \label{eq:v15-Can-zero}
\end{equation}
The required path is uniformly local and uses only the original Renewal ESM
fermion carrier.
\end{corollary}

\begin{proof}
The only cutoff threshold is the geometric requirement in
Lemma~\ref{lem:v15-shift-selection}: the torus must be large enough to contain
the finite bounded-displacement copy bank without wraparound collisions.
Theorem~\ref{thm:v15-automatic-buffer} then verifies the V14 internal-buffer
hypothesis in $\mathbf R_\varrho$.  Apply
Theorem~\ref{thm:v14-buffer-cancellation} coefficientwise.
\end{proof}

\subsection{Why the coefficientwise quantifier is sufficient}
\label{sec:v15-quantifier}

The preceding theorem does not construct one analytic buffer on a
cutoff-independent finite-amplitude ball of all sources.  Such a theorem is not
needed for the continuum target.

\begin{theorem}[Finite-jet native lifting is sufficient for a fixed coefficient packet]
\label{thm:v15-jet-suffices}
Fix a perturbative/source packet
\begin{equation}
 \mathfrak q=(L,\Delta,\varrho,N,\mathcal P_{\rm ext})
 \label{eq:v15-fixed-q}
\end{equation}
and the positive IR/reference prescription.  Suppose two V11 protected native
regulators have V13 controlled-equivalent zero-mode packets.  At sufficiently
large cutoff, the coefficientwise internal-buffer path constructed above,
together with the V13 affine reduced-graph interpolation, supplies every
regulator-path source jet required by the V8--V11 relative forest theorem
through order $\varrho$.  Therefore no analytic finite-amplitude source radius
is required to compare the order-$\mathfrak q$ continuum coefficients.
\end{theorem}

\begin{proof}
Every operation entering the fixed-order shell proof---inverse differentiation,
Riesz/Kato differentiation, determinant/Pfaffian transgression, forest
subtraction, BV brackets and source/contact extraction---is a finite polynomial
or rational jet operation whose denominators are evaluated on the protected
reference pivots.  This is exactly the V2 reference-pivot coefficient theorem
and the V5 contraction-complete induction.  The V15 rotation identities are
exact in $\mathbf R_\varrho$, hence all of their coefficients through source
order $\varrho$ agree with those of an exact path.  Terms of degree
$\varrho+1$ or higher cannot feed a retained coefficient.  The reduced graph
coordinate is treated by V13 in the same finite source algebra.  Thus the V8
marked-forest derivative has all required regulator/source coefficients without
postulating a common nonzero source radius.
\end{proof}

\begin{firewall}[No analytic-source promotion]
Theorem~\ref{thm:v15-jet-suffices} is coefficientwise.  It does not assert that
one V15 buffer remains exactly orthogonal to both zero-mode packets on an open
finite-amplitude source ball independent of the cutoff.  Proving such an
analytic-family cancellation theorem would be a stronger problem, analogous
to the stronger finite-coupling RG trajectory explicitly excluded from E1--E6.
\end{firewall}

\subsection{Unstabilized coefficientwise universality without a spectator}
\label{sec:v15-universality}

\begin{theorem}[Automatic unstabilized universality for the fixed-rank protected packet]
\label{thm:v15-unstabilized-universality}
Fix finite $\mathfrak q$ as in \eqref{eq:v15-fixed-q}.  Let $A$ and $B$ be two
V11 minimal fixed-index native regulators on one protected local source chart,
with the same physical Standard-Model principal Dirac--Yukawa jet,
representation, positive IR/reference prescription and global complex line
data.  Assume their rank-$\kappa=|k|$ kernel packets ($k>0$) or cokernel packets
($k<0$) are V13 controlled-equivalent with the required source jets.

Then, for all sufficiently large cutoffs, the complete regulator comparison can
be performed in the original carrier with no spectator at every source order
retained by $\mathfrak q$.  Consequently the matched continuum states obey
\begin{equation}
 \boxed{
 \widehat{\mathfrak C}^{B,(k),\rm sat}_{\infty,\lambda}
 =F_{\infty,\mathfrak q,k}^{B\leftarrow A}
  \widehat{\mathfrak C}^{A,(k),\rm sat}_{\infty,\lambda}}
 \label{eq:v15-unstabilized-universality}
\end{equation}
without adjoining the V13 inert massive copy.
The equality contains the complete retained connected/source system,
operator mixing, stress/contact coefficients, ghosts/antifields,
BV/BRST/master descendants and determinant/Pfaffian insertions, and remains
strictly projective in perturbative order.
\end{theorem}

\begin{proof}
Corollary~\ref{cor:v15-Can-zero} supplies the V14 unstabilized zero-mode path in
the finite source algebra of the chosen packet.  Theorem
\ref{thm:v15-jet-suffices} promotes it to exactly the regulator/source jets
consumed by the V8--V11 relative forest theorem.  V13 supplies the local affine
reduced-graph interpolation after the zero-mode packet has been transported.
The V11 reduced singular floor and the openness argument permit the usual
one-time shrinking of the protected chart.  Thus the V11 two-scale estimate
applies to the complete coefficient packet without stabilization.  Passing to
the V6 ultraviolet limit gives \eqref{eq:v15-unstabilized-universality}.  V12
supplies the global complex saturation on the faithful branch; the independent
real Pfaffian orientation condition is unchanged.  Chronological V7
normalization gives strict loop-order projectivity.
\end{proof}

\begin{corollary}[Stable and unstabilized coefficientwise universality coincide]
\label{cor:v15-stable-equals-unstable}
At fixed finite coefficient/source order and fixed finite zero-mode rank, the
V13 spectator does not enlarge the class of matched coefficientwise continua on
a protected local Renewal ESM chart.  Every V13 controlled-equivalent pair is
already unstabilized-equivalent for all sufficiently large cutoffs.
\end{corollary}

\begin{proof}
V13 proved the stable statement.  Theorem
\ref{thm:v15-unstabilized-universality} proves the same comparison without the
spectator under exactly the fixed-rank/source-jet hypotheses of the present
packet.
\end{proof}

\subsection{Why this is not cancellation of the full four-dimensional local algebra}
\label{sec:v15-Roe-firewall}

The fixed-rank/source-jet hypotheses are essential.  They must not be replaced
by the statement that every pair of projections in a four-dimensional local
operator algebra cancels.

\begin{remark}[External comparison: cancellation is a low-dimensional property]
\label{rem:v15-Roe-comparison}
For a bounded-geometry metric space $X$, Li and Willett proved that cancellation
of projections in the uniform Roe $C^*$-algebra $C_u^*(X)$ is equivalent to
asymptotic dimension zero; see K.~Li and R.~Willett,
\emph{Low-dimensional properties of uniform Roe algebras}, J. London Math.
Soc. 97 (2018), 98--124, Theorem~2.2.  A four-dimensional lattice has positive
asymptotic dimension, so no unrestricted Roe-algebra cancellation principle is
available.  This external theorem is not used in the proof above; it explains
why V15 keeps the fixed-rank and finite-source-jet quantifiers explicit.
\end{remark}

The V15 proof escapes that obstruction for a concrete reason.  Its candidate
multiplicity $M$ is finite and depends only on the fixed number of source
monomials and the fixed zero-mode rank.  The extra copies are not new fields:
they are bounded covariant translations by at most $S_*$ lattice steps of the
same finite zero-mode packet.  The large cutoff is used only to make room for
this finite local construction.

\subsection{What remains globally}
\label{sec:v15-global-boundary}

A noncontractible global background component can still carry information not
visible in one truncated source chart.  V12 already separated the full
zero-mode bundle and determinant/Pfaffian line data.  V14 further identified
the controlled unitary component.  V15 removes the latter locally at every
fixed source order by constructing an internal reserve, but it does not assert
that the independently chosen local reserves glue to one finite-amplitude
microscopic path around every large background loop.

\begin{proposition}[Overlap of two V15 buffers]
\label{prop:v15-buffer-overlap}
On an overlap of two protected source charts, two V15 buffer choices produce
unstabilized regulator comparison paths with the same endpoints.  Their
concatenation is a closed controlled native loop.  The induced normalized
continuum maps therefore differ, if at all, by the V7 finite scheme/source
transition together with the V12 global line/orientation transport.  On the
faithful complex Standard-Model branch the determinant-line part is trivial;
a real Pfaffian mod-two character remains separate.
\end{proposition}

\begin{proof}
Each local construction is a V11-admissible coefficientwise path with the same
endpoint regulators.  Concatenating one with the reverse of the other gives a
closed sequence of native regulator/scheme transports.  V7 identifies the
residual finite local-coordinate action with its scheme groupoid, and V12
identifies the independent global line transport.  The inherited faithful-spin
complex determinant section has trivial paired holonomy, so only the declared
finite local scheme change and any independent real orientation character can
remain.
\end{proof}

\begin{firewall}[Local coefficientwise cancellation is not a global $\mathcal U_0$ theorem]
V15 closes the spectator distinction for every fixed finite coefficient packet
on a protected local chart.  It does not prove that the full source-dependent
family of zero-mode projections on a noncontractible background component lies
in one identity-component orbit of a global exponentially local unitary group.
That stronger statement requires descent of the local buffer choices, or an
independent computation of the resulting controlled loop class.
\end{firewall}

\section{V15 anti-circularity audit}
\label{sec:v15-audit}

\begin{enumerate}[leftmargin=2.4em]
\item \textbf{The buffer is not assumed to exist.}
It is built from a localized frame of the actual protected Riesz projector and
a finite bank of native bounded parallel shifts.
\item \textbf{Uniform localization is proved from the native spectral/locality
hypotheses.}
The diagonal-floor argument uses only exponential kernel control, finite rank
and bounded geometry; it does not assume prelocalized zero modes.
\item \textbf{The large cutoff is used only geometrically.}
All required displacements are bounded in lattice units by
$S_*(\kappa,M,D)$, independent of the cutoff.  Their physical radius shrinks
like $a$.
\item \textbf{Only finitely many source constraints are paid.}
The rank count in \eqref{eq:v15-kernel-dim} is finite because source order and
source panel are fixed before the cutoff is removed.  No estimate uniform as
$\varrho\to\infty$ is asserted.
\item \textbf{Finite-amplitude source invertibility is not smuggled in.}
The internal-buffer identities are exact in the truncated source algebra; V2's
reference-pivot theorem is what makes that the correct quantifier for the
coefficientwise continuum.
\item \textbf{No global Roe-algebra cancellation is claimed.}
The theorem is a fixed-rank finite-jet result, consistent with the known failure
of unrestricted cancellation on positive-dimensional coarse spaces.
\item \textbf{The independent real Pfaffian orientation is untouched.}
A complex internal buffer does not calculate the V12 mod-two orientation
character.
\end{enumerate}

\section{V15 status ledger}
\label{sec:v15-status}

\subsection*{Proved in V15}
\begin{enumerate}[leftmargin=2.2em]
\item A uniformly exponentially local projection of fixed finite rank admits a
uniformly exponentially localized orthonormal frame.  The proof is constructive:
a finite-propagation truncation gives a cutoff-independent diagonal floor, one
localized projection column is extracted, and the argument iterates on the
orthogonal residual projection.
\item V13 local Kato transport preserves the localized-frame estimates for every
prescribed fixed source derivative after one harmless exponent loss.
\item Finitely many bounded gauge/spin-covariant lattice translations of that
frame can be chosen with a uniformly conditioned Gram matrix; the number and
maximum displacement depend only on the fixed zero-mode rank and desired copy
multiplicity, not on the cutoff.
\item At source order $\varrho$, taking
$M\ge2\nu_\varrho+1$, where $\nu_\varrho$ is the dimension of the complete
finite graded source/BV algebra (equal to $\binom{s+\varrho}{\varrho}$ in the
purely even $s$-writer case), leaves a
$\kappa$-dimensional common annihilator of all $p$- and $q$-overlap coefficients
through that order.  Orthogonalizing the copy bank and selecting this
annihilator yields a rank-$\kappa$ exponentially local projection $r$ which is
orthogonal to both endpoint zero-mode packets in the exact truncated source
algebra.
\item The explicit partial isometries
$a=\mathsf E\Psi^*$ and $b=\mathsf E\mathsf X^*$ make $r$ a V14 internal buffer.  Hence the
V14 cancellation obstruction vanishes coefficientwise through every prescribed
finite source order once the cutoff is sufficiently large.
\item The V14 buffer rotation and V13 reduced-graph interpolation therefore give
an unstabilized fixed-index native regulator path at the level of every retained
source jet.  V8--V11 require no stronger source quantifier.
\item Consequently stable and unstabilized regulator universality coincide for
every prescribed finite coefficient/source packet on a protected local Renewal
ESM chart; the V13 massive spectator is not needed for the continuum
coefficients.
\item The required reserve has bounded lattice radius and hence shrinking
physical radius $O(a)$; it is a genuine ultraviolet internal reserve rather
than an added macroscopic degree of freedom.
\end{enumerate}

\subsection*{Inherited}
\begin{enumerate}[leftmargin=2.2em]
\item V1--V7: source-complete coefficientwise continuum, common BV restoration,
fixed-loop forests, ultraviolet Cauchy convergence, scheme/reference covariance
and exact loop-order projectivity.
\item V8--V12: relative regulator forests, fixed-index reduced shell calculus,
global zero-mode bundle/line transport and complex faithful-spin saturation.
\item V13: exponential locality of protected Riesz zero modes, controlled
zero-mode equivalence, local affine graph reconstruction and stable spectator
universality.
\item V14: the exact complement/$\mathcal U_0$ cancellation criterion and the
explicit internal-buffer rotation giving unstabilized native lifting whenever a
buffer is present.
\end{enumerate}

\subsection*{Conditional}
\begin{enumerate}[leftmargin=2.2em]
\item The automatic reserve theorem is coefficientwise on a protected framed
source/background chart.  It does not construct one analytic finite-amplitude
buffer over an arbitrary large source ball.
\item The zero-mode rank $\kappa$, source order $\varrho$, source-panel dimension
$s$ and representation fibre are fixed before the cutoff is enlarged.  The
buffer radius may depend badly on these finite data.
\item The global noncontractible-background problem is reduced to descent of the
local buffer paths and the controlled loop/unitary component; V15 does not claim
a single global microscopic $\mathcal U_0$ path on every component.
\item The V12 real Pfaffian mod-two orientation condition remains independent.
\item All continuum conclusions retain the positive IR/reference scale and no
uniform $L\to\infty$ estimate.
\end{enumerate}

\subsection*{Open beyond V15}
\begin{enumerate}[leftmargin=2.2em]
\item Globalize the local V15 reserve over a noncontractible gauge/background
component.  The overlap of local reserves defines a controlled unitary-loop
Cech cocycle; compute whether this class vanishes on the faithful Standard-Model
native component or identify a genuine coarse/$K_1$ obstruction.
\item Compute the V12 real Pfaffian mod-two spectral-flow character on any
independently physical Majorana sector.
\item Analyze crossings out of the V11 minimal fixed-index stratum, including
accidental zero-mode pair creation and genuine index change, with the required
divisor/spectral-flow source packet.
\item The independent stronger problems remain: positive-IR removal, analytic
finite-coupling RG trajectories, perturbative summability and nonperturbative
completion.
\end{enumerate}

\subsection*{Exact next load-bearing theorem}
\begin{center}
\fbox{\parbox{0.92\textwidth}{
\textbf{Global descent of the coefficientwise internal reserve.}
Cover a connected protected fixed-index gauge/background component by V15
source charts and construct the overlap controlled-unitary cocycle relating
their internal buffers.  Prove that this cocycle is a coboundary in the native
exponentially local unitary group (hence the local unstabilized paths glue to one
global path), or identify its computable controlled $K_1$/coarse-transport
class.  On the faithful complex Standard-Model branch the determinant-line
holonomy is already trivial by V12, so any surviving obstruction is genuinely
a microscopic local-unitary descent obstruction rather than an anomaly phase.}}
\end{center}

\section*{Append-only continuation marker: V15}
\addcontentsline{toc}{section}{Append-only continuation marker: V15}
\textbf{End of V15.}  The complete V1--V14 body above is retained verbatim.
V15 proves automatic coefficientwise internal cancellation for the actual
fixed-rank protected zero-mode packet at sufficiently large cutoff.  It first
extracts a uniformly localized orthonormal frame from the exponentially local
Riesz projector, then uses finitely many bounded native parallel shifts to
build a uniformly conditioned copy bank.  Since only finitely many source jets
are retained, a dimension count produces a rank-$|k|$ common annihilator of all
endpoint overlap jets, yielding an exact V14 internal buffer in the truncated
source algebra.  The buffer lives in the existing carrier, uses only a bounded lattice
displacement from the original localized zero-mode centres, and therefore has
shrinking added physical transport scale; it removes the V13 spectator from all
fixed-order continuum coefficients.  The remaining stronger gate is global
descent of these local buffer choices around noncontractible background loops.


\clearpage
\section{V16 continuation: Cech descent of local buffer paths and global coefficientwise universality}
\label{sec:v16-continuation}

V15 closes the stable-versus-unstabilized distinction on every protected local
source/background chart at every prescribed finite coefficient order.  The
remaining handoff asked whether the resulting local internal-buffer paths glue
to one microscopic path over a noncontractible background component.  Going
upstream once more separates two logically different questions:
\begin{equation}
 \boxed{
 \text{global coefficientwise continuum equivalence}
 \quad\text{versus}\quad
 \text{one global microscopic regulator path}.}
 \label{eq:v16-two-global-questions}
\end{equation}
The first is the physical endpoint of the perturbative programme and can be
proved by descent of the already established V7 finite scheme/source/line
groupoid.  The second is a stronger locality/topology statement about the
microscopic path space and may retain a controlled loop obstruction even when
all normalized continuum observables already agree globally.

V16 therefore does not assume that the V15 buffers themselves can be chosen as
one global projection.  Instead, on a finite good cover it compares the local
continuum matching maps.  Their overlap differences form an exact nonabelian
Cech cocycle of finite local scheme automorphisms.  Coordinate-free normalized
observables and the common master/source relations are insensitive to that
cocycle, hence descend globally.  A single global coordinate matching exists
precisely when the cocycle is a coboundary.  Moreover, if the cocycle lies in
the finite unipotent subgroup acting trivially on the normalized quotient and
on the associated graded source/contact bank, the cocycle is automatically a
coboundary by an elementary partition-of-unity/central-series argument.

The microscopic V15 paths define a stronger loop cocycle whose image is the
scheme cocycle.  On the faithful complex Standard-Model branch its determinant
line character is already trivial by V12.  Any residual obstruction to one
global microscopic path is therefore a genuinely controlled native-loop class,
not an obstruction to the global coefficientwise continuum.

Throughout V16 let $B$ be a compact connected protected background/source
component at fixed finite coefficient data
\begin{equation}
 \mathfrak q=(L,\Delta,\varrho,N,\mathcal P_{\rm ext})
 \label{eq:v16-fixed-q}
\end{equation}
and positive reference scale.  Choose a finite good cover
$\mathfrak U=(U_\alpha)_{\alpha\in A}$ such that every nonempty finite
intersection is contractible and contained in a V15 protected framed chart.
Let $A$ and $B'$ denote two native regulators to be compared; the prime on
$B'$ avoids confusing the second regulator with the parameter space $B$.
Assume their endpoint zero-mode packets are V13 controlled-equivalent locally,
with the V12 global complex line data and any independent real orientation data
fixed as declared.

\subsection{Local V15 matchings and the finite scheme automorphism sheaf}
\label{sec:v16-scheme-sheaf}

For each chart $U_\alpha$, V15 gives an unstabilized coefficientwise native
comparison path and therefore a continuum matching map
\begin{equation}
 F_\alpha:
 \widehat{\mathfrak C}^{A}_{\infty,\lambda}|_{U_\alpha}
 \longrightarrow
 \widehat{\mathfrak C}^{B'}_{\infty,\lambda}|_{U_\alpha}.
 \label{eq:v16-local-F}
\end{equation}
Every $F_\alpha$ is invertible, finite/local in the V7 sense, preserves the
chosen independent physical normalization coordinates, carries the complete
finite source/contact packet, and is projective in perturbative order.

\begin{definition}[Finite normalized scheme automorphism sheaf]
\label{def:v16-G-sheaf}
For an open set $U\subset B$, let $\mathcal G_A(U)$ be the group of V7-admissible
finite local automorphisms of the normalized $A$-continuum coefficient system
over $U$ which:
\begin{enumerate}[label=\textup{(G\arabic*)},leftmargin=3em]
\item fix the declared independent physical normalization coordinates;
\item transform composite/source coordinates and local contacts by the allowed
finite triangular rules;
\item preserve the common zero vector of the retained BV/BRST/master/source
system;
\item preserve the declared complex determinant-line normalization and, when
present, the chosen real Pfaffian orientation component.
\end{enumerate}
Restriction of coefficients makes $U\mapsto\mathcal G_A(U)$ a sheaf of groups.
Define $\mathcal G_{B'}$ analogously.
\end{definition}

The definition is deliberately intrinsic to the finite coefficient bank.  No
claim is made that $\mathcal G_A$ is the full unitary group of the microscopic
local operator algebra.

\begin{theorem}[Exact overlap cocycle of local regulator matchings]
\label{thm:v16-overlap-cocycle}
On every nonempty overlap $U_{\alpha\beta}:=U_\alpha\cap U_\beta$ define
\begin{equation}
 h_{\alpha\beta}
 :=F_\alpha^{-1}F_\beta
 \in\mathcal G_A(U_{\alpha\beta}).
 \label{eq:v16-hab}
\end{equation}
Then
\begin{align}
 h_{\alpha\alpha}&=I,
 \label{eq:v16-h-aa}\\
 h_{\beta\alpha}&=h_{\alpha\beta}^{-1},
 \label{eq:v16-h-inv}\\
 h_{\alpha\beta}h_{\beta\gamma}&=h_{\alpha\gamma}
 \qquad\text{on }U_{\alpha\beta\gamma}.
 \label{eq:v16-h-cocycle}
\end{align}
Thus $(h_{\alpha\beta})$ is a nonabelian Cech $1$-cocycle of the finite scheme
automorphism sheaf.  The analogous $B'$-side cocycle is conjugate to it by the
local maps $F_\alpha$.
\end{theorem}

\begin{proof}
Each $F_\alpha$ and $F_\beta$ is an isomorphism between the same two normalized
coefficient systems on the overlap.  Their relative map is therefore a finite
local automorphism of the $A$ system.  It fixes the independent normalization
coordinates because both matchings use the same declared normalization, and it
preserves the master/source zero vector and line/orientation data because each
$F_\alpha$ does.  Hence it belongs to $\mathcal G_A$.
The first two identities are immediate, while
\[
 h_{\alpha\beta}h_{\beta\gamma}
 =F_\alpha^{-1}F_\beta F_\beta^{-1}F_\gamma
 =F_\alpha^{-1}F_\gamma
 =h_{\alpha\gamma}.
\]
The $B'$-side formula follows by conjugating with $F_\alpha$.
\end{proof}

\begin{definition}[Continuum descent class]
\label{def:v16-descent-class}
The nonabelian Cech class
\begin{equation}
 \boxed{
 \mathfrak D_{\mathfrak q}(B',A)
 :=[(h_{\alpha\beta})]
 \in\check H^1(B,\mathcal G_A)}
 \label{eq:v16-descent-class}
\end{equation}
is the \emph{continuum coordinate-descent class} of the local V15 comparison.
Here $\check H^1$ is understood as the pointed set of cocycles modulo
$0$-cochain conjugation.
\end{definition}

Changing a local path or a local finite matching replaces
$F_\alpha$ by $F_\alpha s_\alpha^{-1}$ with
$s_\alpha\in\mathcal G_A(U_\alpha)$, and therefore
\begin{equation}
 h_{\alpha\beta}\longmapsto
 s_\alpha h_{\alpha\beta}s_\beta^{-1}.
 \label{eq:v16-coboundary-change}
\end{equation}
Hence \eqref{eq:v16-descent-class} is independent of such local choices.

\subsection{A single global coordinate matching exists exactly when the Cech class vanishes}
\label{sec:v16-global-coordinate}

\begin{theorem}[Global coordinate-matching criterion]
\label{thm:v16-global-coordinate}
There exists one finite local continuum matching map
\begin{equation}
 F:\widehat{\mathfrak C}^{A}_{\infty,\lambda}
 \longrightarrow
 \widehat{\mathfrak C}^{B'}_{\infty,\lambda}
 \label{eq:v16-global-F}
\end{equation}
whose restriction to every $U_\alpha$ differs from the V15 map $F_\alpha$
only by an admissible $A$-scheme automorphism if and only if
\begin{equation}
 \boxed{\mathfrak D_{\mathfrak q}(B',A)=0.}
 \label{eq:v16-D-zero}
\end{equation}
More explicitly, if
\begin{equation}
 h_{\alpha\beta}=s_\alpha^{-1}s_\beta
 \label{eq:v16-h-coboundary}
\end{equation}
for a $0$-cochain $s_\alpha\in\mathcal G_A(U_\alpha)$, then
\begin{equation}
 \boxed{F|_{U_\alpha}=F_\alpha s_\alpha^{-1}}
 \label{eq:v16-glued-F}
\end{equation}
glues exactly on overlaps.  Conversely every global $F$ gives such a
trivialization.
\end{theorem}

\begin{proof}
Assume \eqref{eq:v16-h-coboundary}.  On an overlap,
\[
 (F_\alpha s_\alpha^{-1})^{-1}(F_\beta s_\beta^{-1})
 =s_\alpha F_\alpha^{-1}F_\beta s_\beta^{-1}
 =s_\alpha h_{\alpha\beta}s_\beta^{-1}=I.
\]
Therefore the local maps agree and glue to \eqref{eq:v16-global-F}.

Conversely, suppose $F$ is global and write on each chart
$F=F_\alpha s_\alpha^{-1}$ with
$s_\alpha:=F^{-1}F_\alpha\in\mathcal G_A(U_\alpha)$.  Equality of the two local
expressions on an overlap gives
$F_\alpha^{-1}F_\beta=s_\alpha^{-1}s_\beta$, proving
\eqref{eq:v16-h-coboundary}.
\end{proof}

\subsection{Global physical descent does not require a global coordinate trivialization}
\label{sec:v16-physical-descent}

The stronger-looking condition \eqref{eq:v16-D-zero} is not load bearing for
coordinate-free normalized physics.  The V7 finite scheme groupoid was
introduced precisely to separate coordinates from invariant relations.

\begin{definition}[Scheme-invariant coefficient functional]
\label{def:v16-invariant-functional}
A normalized coefficient functional $\mathcal O_A$ on the $A$-continuum is
\emph{scheme invariant} if
\begin{equation}
 \mathcal O_A\circ g=\mathcal O_A
 \qquad(g\in\mathcal G_A),
 \label{eq:v16-O-invariant}
\end{equation}
and analogously for $B'$.  A covariant finite operator/source tensor is treated
as a section of its V7 finite basis bundle; in that language
\eqref{eq:v16-O-invariant} means equality after the prescribed contragredient
basis transformation.
\end{definition}

This includes the coordinate-free normalized connected functional, separated
support correlators after matched operator/source transport, the statement that
the common BV/BRST/master vector vanishes, and the Standard-Model anomaly-zero
class.

\begin{theorem}[Global coefficientwise universality by scheme-groupoid descent]
\label{thm:v16-global-physical-descent}
Assume on every $U_\alpha$ the V15 local comparison gives
\begin{equation}
 \mathcal O_{B'}\circ F_\alpha=\mathcal O_A
 \label{eq:v16-local-O-match}
\end{equation}
for a V7 scheme-invariant or covariant normalized coefficient functional.
Then the local equalities patch to one global equality on $B$, independently of
the Cech class \eqref{eq:v16-descent-class}.  In particular:
\begin{enumerate}[label=\textup{(P\arabic*)},leftmargin=3em]
\item the complete normalized connected/source functional of $A$ and $B'$ is
globally matched on the whole protected component;
\item every retained composite, stress/contact, ghost/antifield and
operator-mixing relation descends with its V7 covariant finite basis change;
\item the coupled BV/BRST/master/source vector remains globally zero;
\item the Standard-Model local anomaly-zero relation remains globally zero.
\end{enumerate}
Thus a nontrivial coordinate-descent class is a chart/trivialization issue, not
a failure of global coefficientwise regulator universality.
\end{theorem}

\begin{proof}
On an overlap,
\[
 F_\beta=F_\alpha h_{\alpha\beta}.
\]
Hence, using scheme invariance,
\[
 \mathcal O_{B'}\circ F_\beta
 =\mathcal O_{B'}\circ F_\alpha\circ h_{\alpha\beta}
 =\mathcal O_A\circ h_{\alpha\beta}
 =\mathcal O_A.
\]
For covariant operator/source tensors the same computation is made in the
associated finite basis bundle, where the transition $h_{\alpha\beta}$ acts on
components and the tensor itself transforms contragrediently.  The result is a
single global section rather than a single global component array.  The master
vector and anomaly class are zero sections, so every admissible finite
transition fixes them.  Therefore the local V15 equalities satisfy descent on
all overlaps and define the stated global relations.
\end{proof}

\begin{corollary}[Global regulator universality as a scheme-groupoid isomorphism]
\label{cor:v16-groupoid-universality}
For every fixed finite coefficient packet $\mathfrak q$, the two regulator
continua define isomorphic global objects in the V7 groupoid of normalized
finite coefficient systems, even when
$\mathfrak D_{\mathfrak q}(B',A)\ne0$.  Choosing one global coordinate
representative is an additional trivialization problem, not part of the
coordinate-free continuum theorem.
\end{corollary}

\subsection{The faithful complex determinant-line component of the cocycle is trivial}
\label{sec:v16-line-component}

The V12 inherited faithful-spin theorem removes one possible component of the
descent class completely.

\begin{theorem}[No complex determinant-line contribution to the faithful V16 descent class]
\label{thm:v16-line-trivial}
On the ordinary-spin faithful $G_6$ Standard-Model branch of V12, choose the
global parallel complete determinant section and the induced global zero-mode
saturation writer.  Then the determinant-line/frame component of every overlap
automorphism $h_{\alpha\beta}$ is a Cech coboundary and can be set identically
to one by that global line choice.  Hence any nontrivial
$\mathfrak D_{\mathfrak q}(B',A)$ is carried by finite local
scheme/source/operator-basis data, not by complex chiral determinant holonomy.
A genuinely real Pfaffian sector retains its separate V12
$\mathbb Z_2$ orientation character.
\end{theorem}

\begin{proof}
V12 constructs one global nonzero complex determinant/zero-mode saturation
section from the inherited faithful-spin parallel section and the reduced
nonzero determinant.  Restrict it to each $U_\alpha$.  In this common global
frame every local line transition is the identity.  Any line factor which was
present in the independently chosen local V15 frames is therefore exactly the
ratio of two local frame functions and hence a Cech coboundary.  The real
Pfaffian orientation is a real line-local-system question and was explicitly
excluded from that complex argument in V12.
\end{proof}

\subsection{Automatic coordinate descent for the finite unipotent contact subgroup}
\label{sec:v16-unipotent}

Although a general finite scheme cocycle need not be trivial, the purely
redundant/contact part of the V7 scheme group has no Cech obstruction on a
paracompact smooth base.  This yields a useful positive global theorem.

\begin{definition}[Finite unipotent redundant/contact subgroup]
\label{def:v16-unipotent}
At fixed $\mathfrak q$, let $\mathcal N_A\subset\mathcal G_A$ be the subgroup
of finite automorphisms which:
\begin{enumerate}[label=\textup{(N\arabic*)},leftmargin=3em]
\item act identically on all independent normalized physical/EFT quotient
coordinates and on the chosen global operator/source basis;
\item preserve the global complex line frame;
\item are strictly upper triangular with respect to one finite filtration by
BV-exact/redundant degree and source/contact degree.
\end{enumerate}
Because the coefficient bank is finite, there is a finite central series
\begin{equation}
 \mathcal N_A=\mathcal N^1_A\supset\mathcal N^2_A\supset\cdots
 \supset\mathcal N^{s+1}_A=\{I\}
 \label{eq:v16-central-series}
\end{equation}
whose successive quotients are sheaves of smooth sections of finite-dimensional
complex vector bundles under addition.
\end{definition}

\begin{theorem}[Unipotent Cech descent]
\label{thm:v16-unipotent-H1}
Let $B$ be paracompact.  Every Cech $1$-cocycle with values in
$\mathcal N_A$ is a coboundary.  Equivalently,
\begin{equation}
 \boxed{\check H^1(B,\mathcal N_A)=\{0\}.}
 \label{eq:v16-unipotent-H1}
\end{equation}
Therefore, if the V15 overlap maps $h_{\alpha\beta}$ lie in the purely
unipotent redundant/contact subgroup after one coherent global choice of
physical quotient coordinates, operator/source basis and complex line frame,
then the local matching maps glue to one global finite coordinate map
\eqref{eq:v16-global-F}.
\end{theorem}

\begin{proof}
Project the cocycle to the abelian quotient
$\mathcal N_A^1/\mathcal N_A^2$.  This quotient is the sheaf of smooth sections
of a finite-dimensional vector bundle, hence a fine sheaf.  Its Cech
$H^1$ vanishes: explicitly, for an additive cocycle $c_{\alpha\beta}$ and a
partition of unity $(\rho_\gamma)$ subordinate to the cover, set
\begin{equation}
 s_\alpha:=\sum_\gamma\rho_\gamma c_{\gamma\alpha}.
 \label{eq:v16-partition-cochain}
\end{equation}
The cocycle identity gives
$c_{\alpha\beta}=s_\beta-s_\alpha$ on overlaps.  Lift this $0$-cochain to
$\mathcal N_A^1$ and modify the original nonabelian cocycle by the corresponding
coboundary.  The modified cocycle now takes values in $\mathcal N_A^2$.

Repeat the argument on
$\mathcal N_A^2/\mathcal N_A^3$, and continue through the finite central
series.  After finitely many steps the cocycle is the identity.  Composing the
finite lifted $0$-cochains gives a trivialization in $\mathcal N_A$.
\end{proof}

\begin{corollary}[A practical global-coordinate criterion]
\label{cor:v16-practical-global}
If one can choose the finite physical quotient coordinates, the composite/source
basis, and the faithful complex line frame globally so that differences between
local V15 buffer paths act only by BV-exact/redundant/contact redefinitions,
then
\begin{equation}
 \boxed{\mathfrak D_{\mathfrak q}(B',A)=0}
 \label{eq:v16-D-practical-zero}
\end{equation}
and a single global finite continuum matching map exists.
\end{corollary}

This corollary is stronger than what is required for physical descent, but it
pinpoints the only finite-coordinate ingredient that must be checked in any
specific global implementation.

\subsection{Projective compatibility in perturbative order survives descent}
\label{sec:v16-projective}

For each terminal loop order $L$, write
$F_{\alpha,L}$ and $h_{\alpha\beta,L}$ for the preceding maps.  V15 constructed
all local buffer paths chronologically, so
\begin{equation}
 \pi_L^{L+1}F_{\alpha,L+1}\iota_L^{L+1}=F_{\alpha,L}.
 \label{eq:v16-local-projective}
\end{equation}

\begin{theorem}[Projective Cech descent]
\label{thm:v16-projective-descent}
The overlap cocycles satisfy
\begin{equation}
 \boxed{
 \pi_L^{L+1}h_{\alpha\beta,L+1}\iota_L^{L+1}
 =h_{\alpha\beta,L}.}
 \label{eq:v16-h-projective}
\end{equation}
Thus the descent classes form a projective family in perturbative order.  The
global scheme-groupoid isomorphism of
Corollary~\ref{cor:v16-groupoid-universality} restricts exactly from $L+1$ to
$L$.  If the cocycle is trivialized by coherent $0$-cochains (in particular by
the unipotent construction above), the global coordinate maps may also be
chosen projectively.
\end{theorem}

\begin{proof}
Insert
$h_{\alpha\beta,L}=F_{\alpha,L}^{-1}F_{\beta,L}$ and commute the V7 bounded
truncation maps through the projective local matchings.  This gives
\eqref{eq:v16-h-projective}.  All coordinate-free observables already commute
with perturbative truncation by V7, so the descended groupoid isomorphism does
as well.  In the unipotent case choose the central-series trivializations
chronologically: at order $L+1$ leave the already fixed order-$\le L$
$0$-cochains untouched and solve only the new finite layer.  This is the same
V7 projective normalization argument used for the Hodge/restoring sections.
\end{proof}

\subsection{The stronger microscopic loop cocycle}
\label{sec:v16-microscopic-loop}

The scheme cocycle is only the continuum image of a stronger microscopic
object.  Choose the actual local V15 native path
\begin{equation}
 \Gamma_\alpha:
 [0,1]\times U_\alpha\longrightarrow\mathfrak P_{k,\mathfrak q}
 \label{eq:v16-Gamma-alpha}
\end{equation}
from regulator $A$ to regulator $B'$, where
$\mathfrak P_{k,\mathfrak q}$ is the protected native local-operator space of
V12--V15.  On an overlap, concatenate
\begin{equation}
 \ell_{\alpha\beta}
 :=\Gamma_\alpha*\overline{\Gamma_\beta}
 \label{eq:v16-loop}
\end{equation}
into a based closed native loop at $A$.

\begin{proposition}[Microscopic loop cocycle and its continuum image]
\label{prop:v16-loop-cocycle}
Follow $\Gamma_\alpha$ from $A$ to $B'$ and then
$\overline{\Gamma_\beta}$ back from $B'$ to $A$ as in
\eqref{eq:v16-loop}.  Modulo endpoint-fixed homotopy through protected native
loops, chronological concatenation satisfies
\begin{equation}
 [\ell_{\alpha\beta}]\star[\ell_{\beta\gamma}]
 =[\ell_{\alpha\gamma}]
 \label{eq:v16-loop-cocycle}
\end{equation}
on triple overlaps.  The V8--V15 finite regulator-matching construction defines
a homomorphism from chronological loop concatenation to composition of finite
scheme maps,
\begin{equation}
 \operatorname{Hol}_{\rm cont}:
 [\text{protected native loops}]
 \longrightarrow\mathcal G_A,
 \label{eq:v16-Hol-cont}
\end{equation}
and the exact index convention is
\begin{equation}
 \boxed{
 \operatorname{Hol}_{\rm cont}([\ell_{\alpha\beta}])
 =F_\beta^{-1}F_\alpha
 =h_{\beta\alpha}
 =h_{\alpha\beta}^{-1}.}
 \label{eq:v16-loop-image}
\end{equation}
On the faithful complex Standard-Model branch, the determinant-line character
of every $[\ell_{\alpha\beta}]$ is trivial by V12.  A real Pfaffian loop may
retain its independent mod-two character.
\end{proposition}

\begin{proof}
Chronological concatenation gives, with the middle path cancelling its reverse,
\[
 (\Gamma_\alpha\overline{\Gamma_\beta})\star
 (\Gamma_\beta\overline{\Gamma_\gamma})
 \simeq
 \Gamma_\alpha\overline{\Gamma_\gamma},
\]
which proves \eqref{eq:v16-loop-cocycle}.  The local regulator-matching map is
multiplicative under chronological concatenation because its finite
normalization/source/line flow is composed in the same order; reversing a path
gives the inverse map.  Hence it descends to endpoint-fixed homotopy classes.
The first leg $\Gamma_\alpha$ acts by $F_\alpha$ and the return leg
$\overline{\Gamma_\beta}$ acts by $F_\beta^{-1}$, so the closed loop acts by
$F_\beta^{-1}F_\alpha=h_{\beta\alpha}$, proving
\eqref{eq:v16-loop-image}.  The triple-overlap holonomy identity is then
\[
 h_{\gamma\beta}h_{\beta\alpha}=h_{\gamma\alpha},
\]
which is the inverse-index form of \eqref{eq:v16-h-cocycle}.  The final line
statement is the V12 faithful complex holonomy theorem and its separate real
Pfaffian firewall.
\end{proof}

\begin{definition}[Continuum-invisible native loop kernel]
\label{def:v16-loop-kernel}
The kernel
\begin{equation}
 \mathcal K_{\rm nat}
 :=\ker\operatorname{Hol}_{\rm cont}
 \label{eq:v16-loop-kernel}
\end{equation}
consists of protected microscopic regulator loops which can be nontrivial in
the local operator path space while acting trivially on the entire retained
coefficientwise continuum system.
\end{definition}

This kernel is exactly why a global microscopic path is stronger than the
global continuum theorem.  Triviality of
$\mathfrak D_{\mathfrak q}$ does not imply that the microscopic loop cocycle is
null-homotopic; it only says that its image in the finite continuum scheme
groupoid can be removed by local coordinate choices.

\begin{firewall}[Do not promote continuum descent into microscopic path connectedness]
V16 proves global equality of coordinate-free coefficientwise continua and,
under an explicit Cech-triviality condition, one global finite coordinate
matching.  It does not prove that the local V15 buffer paths themselves glue to
one microscopic path in the exponentially local operator space.  A closed
native loop can lie in $\mathcal K_{\rm nat}$ and therefore be invisible to all
retained continuum coefficients.  Classifying or eliminating that kernel is a
strictly stronger local-operator topology problem.
\end{firewall}

\subsection{Global universality theorem at the coefficientwise endpoint}
\label{sec:v16-global-theorem}

\begin{theorem}[Global coefficientwise regulator universality on a protected component]
\label{thm:v16-global-universality}
Fix the faithful complex Standard-Model branch, positive IR/reference scale,
fixed index/minimal reduced-gap component and finite coefficient packet
$\mathfrak q$.  Let two native regulators $A,B'$ be such that on a finite good
cover of a compact connected protected background/source component their
zero-mode packets are V13 controlled-equivalent and hence admit the V15 local
unstabilized comparison.

Then the two ultraviolet continua are globally isomorphic as objects of the V7
finite scheme/source/line groupoid.  Equivalently, every globally defined
scheme-invariant or covariant normalized connected/source/master relation
agrees between the two regulators on the whole component.  The result is exact
and projective in perturbative order.  No global microscopic regulator path is
required.

If, in addition, the Cech class
$\mathfrak D_{\mathfrak q}(B',A)$ vanishes---in particular under the unipotent
hypothesis of Theorem~\ref{thm:v16-unipotent-H1}---then the equivalence is
represented by one global finite local continuum matching map
\eqref{eq:v16-global-F}.
\end{theorem}

\begin{proof}
Apply V15 on each chart to obtain the local continuum maps $F_\alpha$.
Theorem~\ref{thm:v16-overlap-cocycle} produces the exact finite scheme cocycle.
Theorem~\ref{thm:v16-global-physical-descent} then patches every coordinate-free
or covariant normalized continuum relation globally, while
Theorem~\ref{thm:v16-line-trivial} removes the complex determinant-line part of
the cocycle.  Projective compatibility is
Theorem~\ref{thm:v16-projective-descent}.  If the Cech class vanishes,
Theorem~\ref{thm:v16-global-coordinate} gives one global coordinate map; the
unipotent sufficient condition is Theorem~\ref{thm:v16-unipotent-H1}.
\end{proof}

\begin{center}
\fbox{\parbox{0.92\textwidth}{
\textbf{The global coefficientwise continuum no longer depends on global
microscopic buffer descent.}
Local V15 unstabilized regulator paths suffice.  Their overlap mismatch is a
finite V7 scheme/source/line cocycle; coordinate-free normalized observables and
the common master/anomaly relations satisfy descent automatically.  A single
global component array is a separate Cech trivialization problem, and a single
global microscopic path is stronger still.}}
\end{center}

\section{V16 anti-circularity audit}
\label{sec:v16-audit}

\begin{enumerate}[leftmargin=2.4em]
\item \textbf{A global microscopic path was not assumed merely because local
paths exist.}  Their overlap difference is recorded explicitly as the native
loop cocycle \eqref{eq:v16-loop-cocycle}.
\item \textbf{A global coordinate map was not assumed merely because physical
observables agree.}  Its exact obstruction is the nonabelian Cech class
$\mathfrak D_{\mathfrak q}$.
\item \textbf{The continuum theorem is not weakened by a nontrivial coordinate
cocycle.}  V7 scheme invariance/covariance is precisely the descent datum for
the normalized connected and master/source systems.
\item \textbf{The faithful complex determinant line is not counted twice.}
V12 provides one global frame; V16 removes that factor from the overlap cocycle
before discussing residual scheme/native-loop data.
\item \textbf{Unipotent descent is proved, not asserted.}  The finite central
series reduces the nonabelian problem to additive fine sheaves killed by a
partition of unity.
\item \textbf{The native loop kernel is kept separate from finite scheme
holonomy.}  Loops invisible to all retained continuum coefficients can still be
nontrivial microscopic local-operator loops.
\item \textbf{No real Pfaffian orientation is inferred from the complex
faithful-spin theorem.}  Its V12 mod-two character remains an independent global
condition.
\end{enumerate}

\section{V16 status ledger}
\label{sec:v16-status}

\subsection*{Proved in V16}
\begin{enumerate}[leftmargin=2.2em]
\item The local V15 regulator matching maps on a good cover differ on overlaps
by exact finite V7 scheme/source/line automorphisms.  These automorphisms form a
nonabelian Cech $1$-cocycle whose class is independent of the local path choices.
\item A single global finite continuum coordinate matching exists if and only
if this Cech class is trivial.  An explicit $0$-cochain then glues the local
maps.
\item Coordinate-free or covariant normalized connected/source observables,
operator-mixing relations, the common BV/BRST/master zero vector and the
Standard-Model anomaly-zero class descend globally regardless of that
coordinate Cech class.  Hence global coefficientwise regulator universality is
an isomorphism in the V7 finite scheme groupoid, not a requirement of one
global component array.
\item On the inherited faithful complex Standard-Model branch, the determinant
line/frame component of the descent cocycle is trivial; any residual class is a
finite scheme/source/operator-basis or microscopic-locality datum.  A real
Pfaffian mod-two character remains separate.
\item Every cocycle valued in the finite nilpotent unipotent
redundant/contact subgroup is a coboundary.  Thus, after a coherent global
choice of physical quotient coordinates, source/operator basis and complex
line frame, purely exact/contact overlap differences always admit one global
finite coordinate matching.
\item The overlap cocycles and all resulting descended continuum relations are
strictly projective in perturbative order.
\item The actual local V15 native paths define a stronger Cech cocycle of
closed microscopic loops.  Its image under the coefficientwise matching
holonomy is exactly the finite scheme cocycle.  The kernel consists of native
loops invisible to the entire retained coefficientwise continuum.
\item Therefore the global physical coefficientwise continuum is already
closed on any component admitting the local V15 hypotheses; classification of
one global microscopic path is a strictly stronger post-continuum problem.
\end{enumerate}

\subsection*{Inherited}
\begin{enumerate}[leftmargin=2.2em]
\item V1--V7: source-complete coefficientwise ESM continuum, common BV
restoration, fixed-loop forests, UV Cauchy theorem, scheme/reference covariance
and exact perturbative-order projectivity.
\item V8--V12: regulator-marked relative forests, fixed-index reduced shell
calculus, global zero-mode bundle/line transport and faithful complex
saturation.
\item V13--V14: exponential locality of protected zero modes, controlled module
and exact controlled cancellation/$\mathcal U_0$ criteria.
\item V15: automatic coefficientwise internal buffers and local unstabilized
regulator universality in the original Renewal carrier at sufficiently large
cutoff.
\end{enumerate}

\subsection*{Conditional}
\begin{enumerate}[leftmargin=2.2em]
\item V16 requires a finite good cover by protected V15 charts for the fixed
coefficient packet.  It does not produce such a cover across an index-changing
or reduced-gap-closing locus.
\item Existence of one global finite component map requires triviality of the
nonabelian Cech class.  This is automatic for the unipotent exact/contact
subgroup but not proved for an arbitrary composite/operator-basis automorphism
cocycle.
\item Existence of one global microscopic regulator path is stronger still and
is not implied by trivial continuum descent; native loops in
$\mathcal K_{\rm nat}$ can be continuum invisible.
\item The V12 real Pfaffian mod-two orientation condition remains independent.
\item All conclusions remain coefficientwise at fixed finite
$(L,\Delta,\varrho,N)$ and positive IR/reference scale.
\end{enumerate}

\subsection*{Open beyond V16}
\begin{enumerate}[leftmargin=2.2em]
\item Compute the microscopic native-loop descent class of the local V15 paths
on noncontractible background components.  Determine whether the controlled
loop cocycle is null-homotopic, or identify a nonzero class in an appropriate
controlled/coarse $K_1$ or transport group.
\item Determine whether all physically relevant overlap automorphisms can be
reduced globally to the unipotent exact/contact subgroup after fixing the
natural Standard-Model composite/source basis.  A positive result would produce
one global continuum coordinate matching, not merely a groupoid isomorphism.
\item Compute the V12 real Pfaffian mod-two spectral-flow character on any
independently physical Majorana sector.
\item Analyze index-changing/reduced-gap-closing divisors and the independent
positive-IR/nonperturbative directions retained by V15.
\end{enumerate}

\subsection*{Exact next load-bearing theorem}
\begin{center}
\fbox{\parbox{0.92\textwidth}{
\textbf{Microscopic native-loop descent / controlled $K_1$.}
For the good-cover V15 construction on a connected noncontractible protected
background component, classify the closed overlap loops
$\ell_{\alpha\beta}$ in the gauge/spin-covariant exponentially local operator
space.  Prove that their Cech class is null in the protected native path
groupoid (hence the local buffer paths glue to one global microscopic path), or
construct a computable controlled/coarse $K_1$ or transport invariant which
survives while mapping to the identity under
$\operatorname{Hol}_{\rm cont}$.  On the faithful complex Standard-Model
branch any such surviving class is necessarily invisible to the determinant
eta character and to the entire fixed-order coefficientwise continuum proved
above.}}
\end{center}

\section*{Append-only continuation marker: V16}
\addcontentsline{toc}{section}{Append-only continuation marker: V16}
\textbf{End of V16.}  The complete V1--V15 body above is retained verbatim.
V16 goes upstream from the demand for one global microscopic V15 buffer and
separates global physical continuum descent from global path trivialization.
Local V15 matching maps define an exact nonabelian Cech cocycle of V7 finite
scheme/source/line automorphisms.  A single global component map exists exactly
when this cocycle is a coboundary, and cocycles in the finite unipotent
exact/contact subgroup are proved to be automatically trivial.  Regardless of
that coordinate class, all scheme-invariant/covariant normalized connected and
master/source relations satisfy descent and therefore give one global
coefficientwise regulator equivalence in the V7 scheme groupoid.  The faithful
complex determinant-line contribution is already trivial by V12.  The local
microscopic paths define a stronger closed-loop cocycle whose image is the
finite scheme cocycle; its continuum-invisible kernel is the remaining
strictly stronger local-operator topology problem.



\clearpage
\section{V17 continuation: controlled zero-mode descent comes before controlled $K_1$}
\label{sec:v17-continuation}

V16 isolated a microscopic loop cocycle left over after the global
coefficientwise continuum had already descended.  Its proposed next variable
was a controlled $K_1$/transport class.  That variable is not wrong, but it is
one layer too downstream.  Before a complement-unitary loop can be classified,
one must first know that the local zero-mode identifications themselves glue
to one global controlled module isomorphism.  The obstruction to that gluing
is a controlled Murray--von Neumann/$K_0$ datum.

This is the load-bearing upstream correction of V17.  Let $B$ be a compact
connected protected background/source component covered by a finite good cover
$(U_\alpha)_\alpha$.  Let $A$ and $B'$ denote the two endpoint native
regulators, and write $p$ and $q$ for their global V11 kernel projectors when
$k>0$ (or their cokernel projectors when $k<0$).  The local V15 paths on
$U_\alpha$ give controlled Kato identifications
\begin{equation}
 v_\alpha:p|_{U_\alpha}\longrightarrow q|_{U_\alpha},
 \qquad
 v_\alpha^*v_\alpha=p,
 \qquad
 v_\alpha v_\alpha^*=q.
 \label{eq:v17-local-v}
\end{equation}
All source derivatives through the prescribed finite coefficient packet are
understood.  We use the fixed-index kernel notation below; every statement has
the chiral-dual cokernel form for $k<0$.

V17 proves:
\begin{enumerate}[label=\textup{(V17.\arabic*)},leftmargin=3.8em]
\item the local identifications \eqref{eq:v17-local-v} define an exact
nonabelian Cech cocycle in the controlled unitary automorphism sheaf of the
zero-mode module;
\item its Cech class is precisely the obstruction to one global controlled
partial isometry $v:p\to q$;
\item the group-completed shadow of this class is the relative controlled
$K_0$ element $[q]-[p]$, while its ordinary bundle and determinant-line shadows
are obtained by forgetting locality and then taking the top exterior power;
\item the V16 microscopic overlap loop restricts on the zero-mode packet to the
inverse of this same Cech cocycle, so the controlled module class is the
primary microscopic descent invariant;
\item once this primary class vanishes, V13 gives one global stabilized native
path, and V14 gives a single global \emph{unstabilized} native path exactly when
the remaining global complement/$\mathcal U_0$ cancellation obstruction
vanishes;
\item the faithful complex determinant theorem of V12 removes only the
$U(1)$ determinant of the controlled zero-mode cocycle; for rank at least two,
higher bundle data can remain.  Thus a controlled $K_1$ or complement
transport invariant is meaningful only after this controlled $K_0$ layer has
been neutralized.
\end{enumerate}

\subsection{The global controlled coefficient algebra}
\label{sec:v17-global-algebra}

Fix one exponent $\beta>0$ small enough for all V13 Riesz/Kato/source estimates
on the finite cover.  Let
\begin{equation}
 \mathfrak A_\beta(B)
 \label{eq:v17-global-algebra}
\end{equation}
be the unital Banach $*$-algebra of continuous (smooth on each chart)
background/source families of gauge/spin-covariant kernels whose weighted row
and column norms, together with the finitely many retained source derivatives,
are uniformly bounded at exponent $\beta$.  Multiplication and involution are
the pointwise kernel operations.  The V13 weighted Schur theorem makes this a
Banach $*$-algebra after replacing the norm by the finite sum of its source-jet
seminorms.

Let $\mathcal V_{\rm ctl}(B)$ be the commutative semigroup of projections in
finite matrix algebras over $\mathfrak A_\beta(B)$ modulo controlled
Murray--von Neumann equivalence, with addition by block direct sum.  Its
Grothendieck group will be denoted
\begin{equation}
 K_{0,\rm ctl}(B):=G(\mathcal V_{\rm ctl}(B)).
 \label{eq:v17-controlled-K0}
\end{equation}
This definition is deliberately internal to the weighted native algebra; no
spectral-invariance theorem identifying it with the $K$-theory of a larger
$C^*$ completion is needed below.

\begin{definition}[Relative controlled zero-mode class]
\label{def:v17-relative-K0}
For two global protected zero-mode projectors $p,q\in\mathfrak A_\beta(B)$ of
the same finite rank, define
\begin{equation}
 \boxed{
 \Delta_{0,\rm ctl}(q,p):=[q]-[p]\in K_{0,\rm ctl}(B).}
 \label{eq:v17-relative-K0}
\end{equation}
The stronger semigroup statement $[p]_{\mathcal V}=[q]_{\mathcal V}$ means
that there is already a global controlled partial isometry without any matrix
stabilization.
\end{definition}

\begin{proposition}[Native homotopy annihilates the relative controlled class]
\label{prop:v17-native-K0-zero}
If $A$ and $B'$ are joined by one global V11-admissible native fixed-index
homotopy on $B$, then
\begin{equation}
 \boxed{
 [p]_{\mathcal V_{\rm ctl}(B)}=[q]_{\mathcal V_{\rm ctl}(B)},
 \qquad
 \Delta_{0,\rm ctl}(q,p)=0.}
 \label{eq:v17-native-K0-zero}
\end{equation}
\end{proposition}

\begin{proof}
V13 proves that Kato transport along a native fixed-index path is exponentially
local with all retained source derivatives.  Its endpoint partial isometry
$v=U_1p$ satisfies $v^*v=p$ and $vv^*=q$, proving controlled
Murray--von Neumann equivalence.  Equality in the Grothendieck group follows.
\end{proof}

Thus $\Delta_{0,\rm ctl}$ is already a computable necessary obstruction to a
global microscopic path.  It precedes every complement-unitary or $K_1$ issue.

\subsection{The exact Cech cocycle of local controlled zero-mode identifications}
\label{sec:v17-zero-Cech}

On an overlap $U_{\alpha\beta}$ define
\begin{equation}
 \boxed{
 g_{\alpha\beta}:=v_\alpha^*v_\beta
 \in\mathcal U\bigl(p\mathfrak A_\beta(U_{\alpha\beta})p\bigr).}
 \label{eq:v17-gab}
\end{equation}
Here the unit of the corner is $p$.  The inverse is
$g_{\beta\alpha}=g_{\alpha\beta}^*$.

\begin{theorem}[Controlled zero-mode descent cocycle]
\label{thm:v17-zero-Cech}
The maps \eqref{eq:v17-gab} obey the exact cocycle identity
\begin{equation}
 \boxed{
 g_{\alpha\beta}g_{\beta\gamma}=g_{\alpha\gamma}}
 \label{eq:v17-g-cocycle}
\end{equation}
on every triple overlap.  Replacing the local identifications by
\begin{equation}
 v_\alpha' = v_\alpha c_\alpha,
 \qquad
 c_\alpha\in
 \mathcal U\bigl(p\mathfrak A_\beta(U_\alpha)p\bigr),
 \label{eq:v17-v-gauge}
\end{equation}
changes the cocycle by
\begin{equation}
 g'_{\alpha\beta}
 =c_\alpha^*g_{\alpha\beta}c_\beta.
 \label{eq:v17-g-gauge}
\end{equation}
Hence the local V15 identifications define a well-defined nonabelian Cech
class
\begin{equation}
 \boxed{
 \mathfrak Z_{\rm ctl}(q,p)
 \in
 \check H^1\!\left(
 B,\mathcal U(p\mathfrak A_\beta p)
 \right).}
 \label{eq:v17-Zctl}
\end{equation}
This class is trivial if and only if the local $v_\alpha$ can be modified by
controlled zero-mode automorphisms to glue to one global controlled partial
isometry
\begin{equation}
 v\in\mathfrak A_\beta(B),
 \qquad v^*v=p,
 \qquad vv^*=q.
 \label{eq:v17-global-v}
\end{equation}
\end{theorem}

\begin{proof}
Since $v_\beta v_\beta^*=q$,
\[
 g_{\alpha\beta}g_{\beta\gamma}
 =v_\alpha^*v_\beta v_\beta^*v_\gamma
 =v_\alpha^*qv_\gamma
 =v_\alpha^*v_\gamma
 =g_{\alpha\gamma}.
\]
The frame-change formula \eqref{eq:v17-g-gauge} is immediate.  If the Cech
class is trivial, choose $c_\alpha$ so that
$g'_{\alpha\beta}=p$.  Then on overlaps
\[
 (v_\alpha c_\alpha)^*(v_\beta c_\beta)=p.
\]
Both maps are isometries from the same $p$-range onto the same $q$-range; the
preceding identity implies equality.  Thus they glue to a global $v$.  The
converse follows by taking $v_\alpha$ to be restrictions of a global $v$.
\end{proof}

\begin{corollary}[Exact global module obstruction]
\label{cor:v17-module-obstruction}
The class $\mathfrak Z_{\rm ctl}(q,p)$ is the exact obstruction to global
controlled equivalence of the endpoint zero-mode packets.  In particular,
\begin{equation}
 \mathfrak Z_{\rm ctl}(q,p)=0
 \quad\Longrightarrow\quad
 \Delta_{0,\rm ctl}(q,p)=0,
 \label{eq:v17-Z-implies-K0}
\end{equation}
but the converse is only a stable statement and need not give an
unstabilized global partial isometry.
\end{corollary}

\subsection{Forgetting control recovers the ordinary zero-mode bundle and line}
\label{sec:v17-forgetful}

The V17 class has the exact hierarchy anticipated abstractly in V12--V14.

\begin{theorem}[Three forgetful shadows of the controlled Cech class]
\label{thm:v17-forgetful}
There are natural forgetful maps
\begin{equation}
 \boxed{
 \mathfrak Z_{\rm ctl}(q,p)
 \longmapsto
 \mathfrak Z_{\rm bun}(K_{B'},K_A)
 \longmapsto
 \mathfrak Z_{\det}
 }
 \label{eq:v17-forgetful}
\end{equation}
with the following meanings.
\begin{enumerate}[label=\textup{(F\arabic*)},leftmargin=3em]
\item $\mathfrak Z_{\rm bun}$ is the ordinary unitary bundle-isomorphism
Cech class of the V12 zero-mode bundles.  It vanishes exactly when the two
kernel bundles (or cokernel bundles) are globally isomorphic.
\item $\mathfrak Z_{\det}$ is obtained by taking determinants of the local
finite-rank transition matrices and is precisely the relative zero-mode
complex line cocycle.
\item On the faithful complex Standard-Model branch, the inherited V12 global
saturation/eta theorem trivializes $\mathfrak Z_{\det}$, but it need not
trivialize $\mathfrak Z_{\rm bun}$ for rank at least two and says still less
about $\mathfrak Z_{\rm ctl}$.
\end{enumerate}
\end{theorem}

\begin{proof}
Choose local V11 Kato frames for $p$ and $q$.  The controlled partial isometry
$v_\alpha$ becomes a finite unitary matrix in these frames.  Forgetting its
spatial kernel leaves the ordinary bundle isomorphism matrix, and the cocycle
law \eqref{eq:v17-g-cocycle} becomes the usual unitary Cech descent condition.
Taking top exterior powers sends the matrix to its determinant, which is the
transition of the relative zero-mode determinant line.  The final assertion is
exactly the V12 distinction between the full zero-mode bundle and its top
exterior power, together with the inherited faithful complex line
trivialization.
\end{proof}

\begin{proposition}[A determinant-invisible primary obstruction]
\label{ex:v17-c2}
Let the background component contain a four-sphere $S^4$ and let $k=2$.  Take
a rank-two zero-mode bundle with $c_1=0$ and $c_2=1$, obtained from the
unit-degree clutching map $S^3\to SU(2)$, and compare it with the trivial
rank-two bundle.  Their determinant lines are both trivial, so
$\mathfrak Z_{\det}=0$, but the bundles are not isomorphic and hence
$\mathfrak Z_{\rm bun}\ne0$.  Therefore
$\mathfrak Z_{\rm ctl}\ne0$ as well.  This is the same finite-rank topology
already isolated in V12, now identified as the primary obstruction sitting
inside the V16 microscopic loop descent problem.
\end{proposition}

The example is topological and onsite in the background coefficient; it does
not assert that this particular class occurs on the faithful Renewal ESM
background component.

\subsection{The V16 microscopic overlap loop contains the same cocycle}
\label{sec:v17-loop-relation}

Let $\Gamma_\alpha$ be the local V15 native path from $A$ to $B'$ and let
$U_\alpha(t)$ denote the V13 local Kato transport of the relevant zero-mode
projection along that path, normalized by $U_\alpha(0)=I$.  Then
\begin{equation}
 v_\alpha=U_\alpha(1)p.
 \label{eq:v17-v-Kato}
\end{equation}
For the V16 closed loop
$\ell_{\alpha\beta}=\Gamma_\alpha*\overline{\Gamma_\beta}$, concatenate the
corresponding Kato transports.

\begin{theorem}[Zero-mode block of the V16 microscopic loop]
\label{thm:v17-loop-zero-block}
The endpoint Kato unitary of the closed loop $\ell_{\alpha\beta}$ commutes with
$p$, and its restriction to the zero-mode range is
\begin{equation}
 \boxed{
 pU_{\ell_{\alpha\beta}}(1)p
 =v_\beta^*v_\alpha
 =g_{\beta\alpha}
 =g_{\alpha\beta}^{-1}.}
 \label{eq:v17-loop-zero-block}
\end{equation}
Thus the controlled zero-mode Cech class is the primary finite-rank block of
the microscopic loop cocycle.  Contracting the V13 affine reduced graph fibre
and the V10 radial spectral coordinate does not change this block.
\end{theorem}

\begin{proof}
Kato transport along the reversed path is the inverse of the forward Kato
transport.  Therefore the loop transport is
$U_\beta(1)^*U_\alpha(1)$.  Since the loop returns to the same projection,
this endpoint commutes with $p$.  Restricting to $p\mathcal H$ gives
\[
 pU_\beta(1)^*U_\alpha(1)p
 =(U_\beta(1)p)^*(U_\alpha(1)p)
 =v_\beta^*v_\alpha.
\]
The remaining identities are \eqref{eq:v17-gab}.  V13 proves that the reduced
minimal-stratum graph coordinate is affine and V10 proves that radial
flattening is contractible with fixed sign, so these sectors cannot alter the
zero-mode endpoint block.
\end{proof}

\begin{corollary}[Why controlled $K_1$ is secondary]
\label{cor:v17-K1-secondary}
A complement-unitary/$K_1$ invariant of the V16 loop is meaningful only after
$\mathfrak Z_{\rm ctl}(q,p)$ has been trivialized.  If the zero-mode Cech class
is nonzero, there is no global controlled zero-mode identification relative to
which a single complement unitary can even be defined on all of $B$.
\end{corollary}

This is the precise upstream correction to the V16 handoff.

\subsection{What vanishing of the primary class buys}
\label{sec:v17-vanishing-primary}

Assume from now on
\begin{equation}
 \mathfrak Z_{\rm ctl}(q,p)=0
 \label{eq:v17-primary-zero}
\end{equation}
and choose the resulting global controlled partial isometry $v$ of
\eqref{eq:v17-global-v}.

\begin{theorem}[Global stabilized native path after controlled module descent]
\label{thm:v17-global-stable-path}
Under \eqref{eq:v17-primary-zero}, the V13 spectator construction may be carried
out globally on $B$.  Hence $A$ and $B'$ are joined by one global stabilized
V11-admissible native path, and after deleting the inert endpoint spectator
they have the same original normalized coefficientwise continuum.
\end{theorem}

\begin{proof}
V13's stable lift depends only on a controlled partial isometry between the
endpoint kernel/cokernel packets and on the V13 local affine graph calculus.
Equation \eqref{eq:v17-global-v} supplies that partial isometry globally.
Adjoin one global inert massive copy and apply the explicit V13 two-copy
rotation fibrewise; every coefficient lies in the global weighted algebra
because $v$ does.  The reduced graph interpolation is global by the V13
formulas and the protected reduced gap.  The endpoint spectator cancels
exactly from the normalized connected/source/BV system by V13.
\end{proof}

The unstabilized global path is now exactly the V14 cancellation problem in the
\emph{global} controlled algebra.

\begin{theorem}[Exact global unstabilized criterion after primary descent]
\label{thm:v17-global-unstable-criterion}
Assume \eqref{eq:v17-primary-zero} and fix one global controlled equivalence
$v:p\to q$.  A global unstabilized V11-admissible native path exists whenever
$v$ extends to a unitary in the identity component
\begin{equation}
 u\in\mathcal U_0(\mathfrak A_\beta(B)).
 \label{eq:v17-global-U0}
\end{equation}
Equivalently, in the terminology of V14, the global complement defect and the
relative component coset vanish.  Conversely every global native path forces
these V14 conditions for its Kato endpoint.  Therefore, after
$\mathfrak Z_{\rm ctl}$ is killed, the remaining microscopic obstruction is
precisely the global complement/$\mathcal U_0$ obstruction isolated in V14.
\end{theorem}

\begin{proof}
The positive statement is V14's unstabilized native lifting theorem applied in
the global algebra.  Conversely a global native path has a global local Kato
endpoint in $\mathcal U_0$ by V13, and its restriction to $p$ is the induced
controlled equivalence.  V14's unitary-extension theorem then gives complement
cancellation, while the path itself gives the identity-component condition.
\end{proof}

Thus the global microscopic problem has a strict order:
\begin{equation}
 \boxed{
 \text{controlled zero-mode descent }(K_0/\text{module})
 \quad\longrightarrow\quad
 \text{controlled complement cancellation/component }(K_1\text{-type}).}
 \label{eq:v17-obstruction-order}
\end{equation}

\subsection{The stable $K_0$ shadow and what it can certify}
\label{sec:v17-stable-shadow}

The exact Cech class is nonabelian and unstable.  Its group-completed shadow is
still useful because it is computable and automatically preserved by native
homotopy.

\begin{proposition}[Stable shadow of the zero-mode descent class]
\label{prop:v17-stable-shadow}
If $\mathfrak Z_{\rm ctl}(q,p)=0$, then
$\Delta_{0,\rm ctl}(q,p)=0$.  Conversely, if
\begin{equation}
 \Delta_{0,\rm ctl}(q,p)\ne0,
 \label{eq:v17-K0-nonzero}
\end{equation}
then no global native path exists, even after adding and subsequently removing
finitely many controlled trivial zero-mode complements within the same
representation/source category.
\end{proposition}

\begin{proof}
The first statement is Corollary~\ref{cor:v17-module-obstruction}.  For the
second, any stabilized native homotopy gives controlled Murray--von Neumann
equivalence after adding the same auxiliary controlled projection to both
endpoints.  This equality becomes $[p]+[r]=[q]+[r]$ in
$\mathcal V_{\rm ctl}$ and hence $[p]=[q]$ in its Grothendieck group, contrary
to \eqref{eq:v17-K0-nonzero}.
\end{proof}

Forgetting control gives the ordinary relative bundle class
\begin{equation}
 [K_{B'}]-[K_A]\in K^0(B)
 \label{eq:v17-ordinary-K0}
\end{equation}
(or the cokernel analogue), and its Chern classes give immediately computable
necessary obstructions.  On the faithful complex Standard-Model branch the
first Chern/determinant-line component is already removed by V12, so the first
possible ordinary complex obstruction at rank at least two is the higher
zero-mode-bundle data such as $c_2$.

\subsection{A finite-dimensional stable range}
\label{sec:v17-stable-range}

V12 proved that finite Grassmannian and stable bundle homotopy agree once the
ambient complement rank is large relative to the dimension of the parameter
panel.  Combining that theorem with V17 gives a clean separation between
ordinary finite-dimensional topology and genuinely controlled spacetime
locality.

\begin{corollary}[Large-cutoff interpretation on a finite panel]
\label{cor:v17-large-cutoff}
Let $B$ be a finite CW complex of dimension $d$ and let the zero-mode rank be
$\kappa$.  Once the finite-cutoff chiral ambient dimension $m$ obeys
\begin{equation}
 d\le2(m-\kappa),
 \label{eq:v17-stable-range}
\end{equation}
ordinary isomorphism of the endpoint zero-mode bundles is equivalent to
homotopy of their finite Grassmannian classifying maps.  Hence, in this range,
any surviving nontriviality of $\mathfrak Z_{\rm ctl}$ after the ordinary
bundle class has been matched is a genuinely locality-controlled obstruction,
not an unstable finite-dimensional Grassmannian defect.
\end{corollary}

\begin{proof}
Apply the V12 stable-range theorem to the kernel/cokernel classifying maps.
The forgetful implication of Theorem~\ref{thm:v17-forgetful} then isolates any
remaining failure of controlled descent in the first arrow of
\eqref{eq:v17-forgetful}, namely the spatial/source locality of the module
identification.
\end{proof}

\subsection{Relation to the V16 continuum scheme cocycle}
\label{sec:v17-scheme-relation}

There are now two separate Cech classes:
\begin{align}
 \mathfrak Z_{\rm ctl}(q,p)
 &\in\check H^1(B,\mathcal U(p\mathfrak A_\beta p)),
 \label{eq:v17-two-cocycles-a}\\
 \mathfrak D_{\mathfrak q}(B',A)
 &\in\check H^1(B,\mathcal G_A).
 \label{eq:v17-two-cocycles-b}
\end{align}
The first lives in the microscopic zero-mode module.  The second lives in the
finite continuum scheme/source/line groupoid.  Neither is obtained from the
other by simply taking a determinant.

\begin{theorem}[Primary microscopic descent versus continuum descent]
\label{thm:v17-micro-vs-cont}
The V8--V16 matching construction defines a natural map from endpoint-fixed
native path data to the finite continuum scheme groupoid, but global physical
coefficientwise universality does not require
$\mathfrak Z_{\rm ctl}=0$.  In particular:
\begin{enumerate}[label=\textup{(D\arabic*)},leftmargin=3em]
\item $\mathfrak Z_{\rm ctl}\ne0$ forbids one global microscopic fixed-index
path and one global controlled zero-mode identification;
\item V16 nevertheless gives a global groupoid equivalence of every
scheme-invariant/covariant normalized continuum relation from the local V15
paths;
\item on the faithful complex branch the determinant-line shadow of both local
descent systems is trivial, so a surviving $\mathfrak Z_{\rm ctl}$ is a
higher zero-mode/locality datum invisible to the anomaly/eta character.
\end{enumerate}
\end{theorem}

\begin{proof}
Item (D1) is Theorem~\ref{thm:v17-zero-Cech} and
Proposition~\ref{prop:v17-native-K0-zero}.  Item (D2) is the V16 global
coefficientwise descent theorem, whose hypotheses are only the local V15
comparisons.  Item (D3) follows from Theorem~\ref{thm:v17-forgetful} together
with the V12 faithful complex line trivialization used again in V16.
\end{proof}

\begin{firewall}[The phrase ``controlled $K_1$'' was one gate too early]
V16 correctly anticipated a possible controlled unitary/transport obstruction,
but V17 shows that it becomes a globally defined invariant only after the
endpoint zero-mode modules have first been descended.  Before
$\mathfrak Z_{\rm ctl}$ is trivialized, the complement unitaries live in
locally identified corners and do not define one global controlled $K_1$
problem.  The correct order is \eqref{eq:v17-obstruction-order}.
\end{firewall}

\subsection{V17 anti-circularity audit}
\label{sec:v17-audit}

\begin{enumerate}[leftmargin=2.4em]
\item \textbf{A loop invariant was not defined before its module was globally
identified.}  The zero-mode Cech class is solved first; only then does the V14
complement-unitary problem become global.
\item \textbf{The complex determinant line is not mistaken for the full
zero-mode bundle.}  The V12 rank-two $c_2$ phenomenon appears explicitly as a
possible determinant-invisible primary obstruction.
\item \textbf{Stable $K_0$ is not promoted to exact Murray--von Neumann
cancellation.}  The exact obstruction is the nonabelian controlled Cech class;
$\Delta_{0,\rm ctl}$ is only its group-completed shadow.
\item \textbf{Local V15 paths are not assumed to glue because their endpoint
regulators agree.}  Their Kato endpoint maps produce the explicit overlap
unitaries \eqref{eq:v17-gab}.
\item \textbf{The V16 global continuum theorem is not weakened.}  It remains
valid even when $\mathfrak Z_{\rm ctl}$ is nonzero because physical descent is
performed in the finite scheme/source/line groupoid.
\item \textbf{Controlled $K_1$ is retained, but at the correct stage.}  After
one global controlled partial isometry has been chosen, V14 gives the exact
complement/$\mathcal U_0$ obstruction which is the next unitary-level problem.
\end{enumerate}

\section{V17 status ledger}
\label{sec:v17-status}

\subsection*{Proved in V17}
\begin{enumerate}[leftmargin=2.2em]
\item Global protected zero-mode projectors define classes in the controlled
Murray--von Neumann semigroup and its Grothendieck group
$K_{0,\rm ctl}(B)$.  Every global native fixed-index homotopy annihilates the
relative controlled class $[q]-[p]$.
\item The local V15 Kato identifications
$v_\alpha:p\to q$ define the exact controlled zero-mode Cech cocycle
$g_{\alpha\beta}=v_\alpha^*v_\beta$.  Its nonabelian Cech class is independent
of local zero-mode frame/isomorphism choices and vanishes exactly when the
local controlled identifications glue to one global controlled partial
isometry.
\item Forgetting spatial control sends this class to the ordinary V12
kernel/cokernel bundle-isomorphism obstruction; taking determinants sends it to
the relative zero-mode line.  The faithful complex eta theorem kills only the
last of these automatically.
\item The zero-mode block of the V16 microscopic overlap loop is exactly
$g_{\beta\alpha}$.  The V13 affine reduced graph and V10 radial sign sectors do
not alter this primary finite-rank block.
\item Consequently the controlled zero-mode module descent is logically prior
to every global complement-unitary/$K_1$ invariant.  Before the module class is
trivialized there is no single global complement corner on which to formulate
that problem.
\item Vanishing of the controlled zero-mode Cech class gives one global
controlled partial isometry and therefore one global V13 stabilized native
path.  A global unstabilized path then exists exactly under the V14 global
complement-cancellation and identity-component conditions.
\item Nonvanishing of the relative controlled $K_0$ shadow forbids every global
native path even after finite controlled stabilization.  In the large-cutoff
finite-panel stable range, any residual obstruction after matching the ordinary
zero-mode bundle is genuinely controlled/spatial rather than a finite
Grassmannian defect.
\item The V16 global coefficientwise continuum equivalence remains valid even
when the primary microscopic controlled module class is nonzero; the latter is
a strictly stronger post-continuum regulator-path datum.
\end{enumerate}

\subsection*{Inherited}
\begin{enumerate}[leftmargin=2.2em]
\item V1--V7: source-complete coefficientwise ESM continuum, common BV
restoration, arbitrary-fixed-loop forests, ultraviolet Cauchy convergence,
scheme/reference covariance and exact perturbative-order projectivity.
\item V8--V12: regulator-marked two-scale comparison, fixed-index reduced shell
calculus, global zero-mode bundle/line transport and faithful complex
saturation.
\item V13--V14: exponential locality of protected zero modes, controlled module
and exact complement/$\mathcal U_0$ cancellation criteria.
\item V15: automatic coefficientwise local internal buffers and unstabilized
local regulator universality in the original carrier.
\item V16: global physical continuum descent through a finite scheme groupoid
and the stronger microscopic overlap-loop cocycle.
\end{enumerate}

\subsection*{Conditional}
\begin{enumerate}[leftmargin=2.2em]
\item The controlled Cech class is defined on a finite good cover of one
protected constant-index/reduced-gap component and at the fixed finite source
packet.  Index-changing or reduced-gap-closing loci require the separate
spectral-flow/divisor analysis retained since V11.
\item Equality of the group-completed controlled $K_0$ classes is only a stable
necessary condition; exact global controlled equivalence is the stronger Cech
triviality/Murray--von Neumann statement.
\item On the faithful complex branch the determinant shadow is trivial, but
higher zero-mode bundle or spatially controlled module data may remain.
\item The real Pfaffian mod-two orientation character remains independent.
\item All continuum conclusions retain the fixed finite
$(L,\Delta,\varrho,N)$ and positive IR/reference quantifiers.
\end{enumerate}

\subsection*{Open beyond V17}
\begin{enumerate}[leftmargin=2.2em]
\item After $\mathfrak Z_{\rm ctl}=0$, compute the global V14 complement
cancellation and identity-component obstruction in the controlled local unitary
algebra.  This is the correct stage for a controlled $K_1$/transport invariant.
\item Determine whether the actual faithful Renewal ESM zero-mode bundles on
physically relevant noncontractible background components are always globally
controlled-isomorphic, or exhibit a higher Chern/coarse-locality obstruction.
\item Compute the independent V12 real Pfaffian mod-two spectral-flow
character on any physical Majorana sector.
\item Analyze index-changing/reduced-gap-closing divisors and the independent
positive-IR/nonperturbative directions retained by V16.
\end{enumerate}

\subsection*{Exact next load-bearing theorem}
\begin{center}
\fbox{\parbox{0.92\textwidth}{
\textbf{Global complement descent / controlled $K_1$ after zero-mode module
trivialization.}
Assume the V17 primary class
$\mathfrak Z_{\rm ctl}(q,p)$ vanishes and fix one global controlled partial
isometry $v:p\to q$.  Analyze the complementary projections $I-p,I-q$ in the
global gauge/spin-covariant exponentially local algebra.  Prove global
controlled cancellation and construct a unitary extension in the identity
component, or define and compute the residual relative unitary-component/
controlled $K_1$ class.  The result must separate its ordinary finite-bundle
shadow from genuinely spatial/coarse transport data and must remain projective
in the fixed perturbative/source order.  A positive theorem, combined with
V17 and V14, would yield one global unstabilized microscopic regulator path;
a nonzero class would give the first genuine continuum-invisible microscopic
phase invariant beyond the zero-mode module itself.}}
\end{center}

\section*{Append-only continuation marker: V17}
\addcontentsline{toc}{section}{Append-only continuation marker: V17}
\textbf{End of V17.}  The complete V1--V16 body above is retained verbatim.
V17 goes upstream from the V16 controlled-$K_1$ handoff and proves that a more
basic controlled zero-mode module descent must be solved first.  Local V15
Kato identifications define an exact nonabelian Cech cocycle in the controlled
unitary automorphism sheaf of the zero-mode packet; its class vanishes exactly
when the local identifications glue to one global controlled partial isometry.
The group-completed shadow is a relative controlled $K_0$ class, while
forgetting locality and taking determinants recover respectively the ordinary
zero-mode bundle and line shadows of V12.  The V16 microscopic overlap loop has
this cocycle as its zero-mode endpoint block.  Only after the primary class is
trivialized does the V14 complement/unitary-component problem become one global
controlled $K_1$ question.  The V16 coefficientwise continuum remains globally
well defined regardless, so these classes belong strictly to the stronger
microscopic regulator-path topology.


\clearpage
\section{V18 continuation: global complement descent, unstable unitary components, and the exact microscopic obstruction tower}
\label{sec:v18-continuation}

V17 proved that the zero-mode module must be descended globally before a
unitary-level obstruction can even be formulated.  Assume from now on that
this primary class has been trivialized.  Thus the endpoint protected zero-mode
projections $p,q$ admit at least one global controlled partial isometry
\begin{equation}
 v^*v=p,
 \qquad
 vv^*=q
 \label{eq:v18-global-v}
\end{equation}
in the global gauge/spin-covariant exponentially local algebra appropriate to
the fixed coefficient packet.  The remaining problem is to extend $v$ over the
complement and then to place that extension in the identity component.

The V17 handoff called this a controlled $K_1$ problem.  That description is
correct only after one further refinement.  There are actually two secondary
layers:
\begin{equation}
 \boxed{
 \text{complement descent}
 \quad\Longrightarrow\quad
 \text{global unitary extension}
 \quad\Longrightarrow\quad
 \text{identity-component cancellation}.}
 \label{eq:v18-secondary-layers}
\end{equation}
The first is a nonabelian Cech problem in the complement corner.  The second is
an unstable component problem for a global extension.  Moreover, the
\emph{stable} projection-pair obstruction vanishes identically: one explicit
$2\times2$ controlled rotation connects the stabilized projections through the
identity component.  Hence any obstruction which survives V18 is genuinely
unstable and microscopic; it is not a stable $K_1$ phase of the coefficientwise
continuum.

Throughout V18 let
\begin{equation}
 \mathfrak A
 :=\mathfrak A_{\exp,\mathfrak q}(B)
 \label{eq:v18-global-algebra}
\end{equation}
denote the global V13 union of positive-exponent local algebras on the compact
protected background/source component $B$, with the fixed representation and
source-jet conditions understood.  When a single exponent is needed we work in
some $\mathfrak A_\beta$ after the finite exponent losses already allowed by
V13.  Put
\begin{equation}
 p^\perp:=I-p,
 \qquad
 q^\perp:=I-q,
 \qquad
 \mathfrak B_p:=p^\perp\mathfrak A p^\perp.
 \label{eq:v18-complement-corner}
\end{equation}

\subsection{Local unitary extensions with a fixed global zero-mode map}
\label{sec:v18-local-extensions}

The local V15 theorem gives unstabilized native paths on every protected chart.
The induced local endpoint unitaries need not initially restrict to the chosen
global partial isometry $v$, but they can be corrected without changing their
projection conjugacy.

\begin{lemma}[Local extensions may be normalized to a prescribed global $v$]
\label{lem:v18-normalize-v}
Let $U\subset B$ be a protected V15 chart and let
$u_U\in\mathcal U(\mathfrak A(U))$ satisfy
\begin{equation}
 u_Upu_U^*=q.
 \label{eq:v18-local-u}
\end{equation}
Define
\begin{equation}
 a_U:=v^*u_Up\in p\mathfrak A(U)p.
 \label{eq:v18-aU}
\end{equation}
Then $a_U$ is unitary in the $p$ corner and
\begin{equation}
 \boxed{
 \widetilde u_U
 :=u_U\bigl(a_U^*+p^\perp\bigr)}
 \label{eq:v18-utilde}
\end{equation}
is a controlled unitary with
\begin{equation}
 \widetilde u_Up=v,
 \qquad
 \widetilde u_Up\widetilde u_U^*=q.
 \label{eq:v18-utilde-properties}
\end{equation}
Consequently
\begin{equation}
 w_U:=\widetilde u_Up^\perp
 \label{eq:v18-wU}
\end{equation}
satisfies
\begin{equation}
 w_U^*w_U=p^\perp,
 \qquad
 w_Uw_U^*=q^\perp.
 \label{eq:v18-local-complement-iso}
\end{equation}
\end{lemma}

\begin{proof}
Using $u_Upu_U^*=q$ and \eqref{eq:v18-global-v},
\[
 a_U^*a_U
 =pu_U^*vv^*u_Up
 =pu_U^*qu_Up=p,
\]
and similarly $a_Ua_U^*=p$.  Thus
$a_U^*+p^\perp$ is unitary.  Moreover
\[
 \widetilde u_Up
 =u_Upa_U^*
 =u_Up\,pu_U^*v
 =qv=v.
\]
The second relation in \eqref{eq:v18-utilde-properties} follows immediately.
Finally the complement identities are
\[
 w_U^*w_U
 =p^\perp\widetilde u_U^*\widetilde u_Up^\perp=p^\perp,
\]
\[
 w_Uw_U^*
 =\widetilde u_Up^\perp\widetilde u_U^*
 =I-\widetilde u_Up\widetilde u_U^*=q^\perp.
\]
All operations take place in the same controlled algebra.
\end{proof}

Thus, once the primary V17 class has been killed, every local V15 comparison
may be put in a common gauge where the zero-mode block is literally the same
fixed $v$.  Only the complement maps $w_U$ remain to be descended.

\subsection{The complement Cech cocycle}
\label{sec:v18-complement-Cech}

Choose a finite good cover $(U_\alpha)_\alpha$ and local normalized extensions
$\widetilde u_\alpha=v+w_\alpha$ from
Lemma~\ref{lem:v18-normalize-v}.  On an overlap define
\begin{equation}
 \boxed{
 z_{\alpha\beta}
 :=w_\alpha^*w_\beta
 \in\mathcal U\bigl(\mathfrak B_p(U_{\alpha\beta})\bigr).}
 \label{eq:v18-zab}
\end{equation}

\begin{lemma}[Exact complement cocycle]
\label{lem:v18-complement-cocycle}
The overlap maps satisfy
\begin{equation}
 z_{\alpha\beta}z_{\beta\gamma}=z_{\alpha\gamma},
 \qquad
 z_{\beta\alpha}=z_{\alpha\beta}^{-1}.
 \label{eq:v18-z-cocycle}
\end{equation}
If another local complement family is chosen as
\begin{equation}
 w'_\alpha=w_\alpha c_\alpha,
 \qquad
 c_\alpha\in\mathcal U(\mathfrak B_p(U_\alpha)),
 \label{eq:v18-w-change}
\end{equation}
then
\begin{equation}
 z'_{\alpha\beta}
 =c_\alpha^*z_{\alpha\beta}c_\beta.
 \label{eq:v18-z-coboundary}
\end{equation}
Hence the local complement extensions define a well-defined nonabelian class
\begin{equation}
 \boxed{
 \mathfrak C_{\rm comp}(v)
 \in
 \check H^1\!\left(
 B,\mathcal U(\mathfrak B_p)
 \right).}
 \label{eq:v18-complement-class}
\end{equation}
\end{lemma}

\begin{proof}
Since every $w_\alpha$ has final projection $q^\perp$,
\[
 z_{\alpha\beta}z_{\beta\gamma}
 =w_\alpha^*w_\beta w_\beta^*w_\gamma
 =w_\alpha^*q^\perp w_\gamma
 =w_\alpha^*w_\gamma.
\]
The inverse formula is the same computation.  Equation
\eqref{eq:v18-z-coboundary} is direct substitution.
\end{proof}

\begin{theorem}[Complement descent criterion]
\label{thm:v18-complement-descent}
For the fixed global zero-mode isomorphism $v$, the following are equivalent:
\begin{enumerate}[label=\textup{(C\arabic*)},leftmargin=3em]
\item $\mathfrak C_{\rm comp}(v)=0$;
\item there exists a global controlled partial isometry
\begin{equation}
 w\in\mathfrak A,
 \qquad
 w^*w=p^\perp,
 \qquad
 ww^*=q^\perp;
 \label{eq:v18-global-w}
\end{equation}
\item there exists a global controlled unitary
\begin{equation}
 \boxed{u=v+w\in\mathcal U(\mathfrak A)}
 \label{eq:v18-global-u}
\end{equation}
whose restriction to the zero-mode packet is exactly $v$.
\end{enumerate}
\end{theorem}

\begin{proof}
If the Cech class vanishes, choose
$c_\alpha\in\mathcal U(\mathfrak B_p(U_\alpha))$ so that
$z_{\alpha\beta}=c_\alpha c_\beta^*$.  Then
$\widehat w_\alpha:=w_\alpha c_\alpha$ agree on overlaps and patch to a global
$w$ satisfying \eqref{eq:v18-global-w}.  Conversely, the restrictions of a
global $w$ give a trivial cocycle.  The equivalence of (C2) and (C3) is exactly
Theorem~\ref{thm:v14-unitary-extension}.
\end{proof}

\begin{definition}[Secondary complement-descent obstruction]
\label{def:v18-secondary-Cech}
For a primary trivialization $v$ of the V17 zero-mode class, the class
$\mathfrak C_{\rm comp}(v)$ is the \emph{secondary complement-descent
obstruction}.  The set of all primary trivializations is a torsor under the
global controlled zero-mode automorphism group
$\mathcal U(p\mathfrak A p)$; changing $v$ can change
$\mathfrak C_{\rm comp}(v)$.
\end{definition}

This dependence is essential.  A global microscopic path need not realize a
preselected zero-mode isomorphism; it need realize \emph{some} primary
trivialization for which the secondary obstruction vanishes.

\subsection{Ordinary and stable shadows of the complement class}
\label{sec:v18-shadows}

Forgetting spatial/source control sends the complement cocycle to the ordinary
unitary transition cocycle between the orthogonal complement bundles.  Hence
\begin{equation}
 \mathfrak C_{\rm comp}(v)=0
 \quad\Longrightarrow\quad
 p^\perp\mathcal H\simeq q^\perp\mathcal H
 \label{eq:v18-ordinary-shadow}
\end{equation}
as ordinary Hermitian bundles.  The converse can fail because the bundle
isomorphism need not be controlled.

\begin{proposition}[The V14 $S^6$ example is a pure secondary obstruction]
\label{prop:v18-S6-secondary}
In Countertheorem~\ref{cth:v14-unstable-cancellation}, the zero-mode range
bundles are globally isomorphic, so the V17 primary class may be chosen zero.
For every such global zero-mode isomorphism $v$, the complement class
$\mathfrak C_{\rm comp}(v)$ is nonzero because one complement is the nontrivial
rank-two bundle $F\to S^6$ and the other is trivial.  Nevertheless the
obstruction disappears after finite stabilization because $F$ is stably
trivial.
\end{proposition}

\begin{proof}
The endpoint range bundles in the V14 construction are both trivial rank-$r$
bundles, so a global $v$ exists.  If
$\mathfrak C_{\rm comp}(v)$ vanished, Theorem
\ref{thm:v18-complement-descent} would give a global complement isomorphism,
contradicting $F\not\simeq\underline{\mathbb C}^2$.  Stable triviality was part
of the V14 construction.
\end{proof}

Thus the secondary class is genuinely unstable.  The next theorem makes the
stable statement completely explicit and also corrects the phrase
``controlled $K_1$ obstruction'' in its stable meaning.

\subsection{The stabilized projection-pair obstruction is identically zero}
\label{sec:v18-stable-rotation}

Let
\begin{equation}
 \widehat p:=
 \begin{pmatrix}p&0\\0&q\end{pmatrix},
 \qquad
 J_v:=
 \begin{pmatrix}0&-v^*\\ v&0\end{pmatrix}
 \in M_2(\mathfrak A).
 \label{eq:v18-Jv}
\end{equation}
Then
\begin{equation}
 J_v^*=-J_v,
 \qquad
 J_v^2=-\widehat p,
 \qquad
 \widehat pJ_v=J_v\widehat p=J_v.
 \label{eq:v18-J-properties}
\end{equation}

\begin{theorem}[Universal stabilized identity-component rotation]
\label{thm:v18-stable-rotation}
For every global controlled partial isometry $v:p\to q$, define
\begin{equation}
 U_v(\theta)
 :=I-\widehat p+
   \cos\theta\,\widehat p+
   \sin\theta\,J_v,
 \qquad
 0\le\theta\le\frac\pi2.
 \label{eq:v18-Uv-theta}
\end{equation}
Then $U_v(\theta)$ is a controlled unitary path with $U_v(0)=I$ and
\begin{equation}
 U_v(\pi/2)
 \begin{pmatrix}p&0\\0&0\end{pmatrix}
 U_v(\pi/2)^*
 =
 \begin{pmatrix}0&0\\0&q\end{pmatrix}.
 \label{eq:v18-stable-move}
\end{equation}
A fixed scalar $2\times2$ rotation then moves the latter projection to
$q\oplus0$.  Hence
\begin{equation}
 \boxed{
 p\oplus0
 \text{ and }
 q\oplus0
 \text{ are conjugate by a unitary in }
 \mathcal U_0(M_2(\mathfrak A)).}
 \label{eq:v18-stable-U0}
\end{equation}
All retained source derivatives stay controlled.
\end{theorem}

\begin{proof}
Equations \eqref{eq:v18-J-properties} show that on the range of $\widehat p$
we have $U_v(\theta)=e^{\theta J_v}$, while on the orthogonal complement the
operator is the identity.  Hence the path is unitary and controlled.  For
$x\in p\mathcal H$,
\[
 U_v(\pi/2)(x,0)=(0,vx),
\]
which proves \eqref{eq:v18-stable-move}.  The scalar rotation
\[
 R(\theta)=
 \begin{pmatrix}\cos\theta&-\sin\theta\\
                 \sin\theta&\cos\theta\end{pmatrix}
\]
is onsite and connected to the identity; its inverse at $\pi/2$ moves
$0\oplus q$ to $q\oplus0$.  Concatenate the two paths.  Fixed source
derivatives are finite polynomial expressions in $v$ and $v^*$.
\end{proof}

\begin{corollary}[No stable controlled $K_1$ obstruction for the projection pair]
\label{cor:v18-no-stable-K1}
Let
\begin{equation}
 K_{1,\rm ctl}^{\rm st}(\mathfrak A)
 :=\varinjlim_n\pi_0\mathcal U(M_n(\mathfrak A))
 \label{eq:v18-stable-K1}
\end{equation}
be the stabilized controlled unitary component group, whenever this notation is
formed in the chosen source/locality category.  Once the V17 primary zero-mode
class has been trivialized, the endpoint projection pair carries no further
class in \eqref{eq:v18-stable-K1}: the implementing stabilized unitary may be
chosen in the identity component by Theorem~\ref{thm:v18-stable-rotation}.
Therefore every remaining global microscopic obstruction is \emph{unstable}.
\end{corollary}

\begin{proof}
The explicit unitary in \eqref{eq:v18-stable-U0} already represents the zero
component before taking the direct limit.
\end{proof}

\begin{firewall}[Stable $K_1$ is not the V18 obstruction]
The V17 handoff anticipated a controlled $K_1$ class.  V18 shows that the
projection-pair obstruction has trivial \emph{stable} $K_1$ shadow once the
primary zero-mode module has been globally identified.  What can remain is the
failure of cancellation and identity-component conjugacy in the original,
unamplified carrier.  This is exactly the unstable V14 obstruction, now with a
global complement-descent layer made explicit.
\end{firewall}

\subsection{The tertiary unitary-component class}
\label{sec:v18-tertiary}

Assume now that for a chosen primary trivialization $v$ one has
\begin{equation}
 \mathfrak C_{\rm comp}(v)=0.
 \label{eq:v18-comp-zero}
\end{equation}
Choose one global complement map $w_0$ and put
\begin{equation}
 u_0:=v+w_0\in\mathcal U(\mathfrak A).
 \label{eq:v18-u0-global}
\end{equation}
Every other extension of the same fixed $v$ has the form
\begin{equation}
 u_0\bigl(p+z\bigr),
 \qquad
 z\in\mathcal U(\mathfrak B_p).
 \label{eq:v18-all-extensions}
\end{equation}

\begin{definition}[Global relative component class]
\label{def:v18-kappa-global}
Define
\begin{equation}
 \boxed{
 \kappa_{\rm glob}(v)
 :=[u_0]\,
 \operatorname{im}\!\left(
 \pi_0\mathcal U(\mathfrak B_p)
 \longrightarrow
 \pi_0\mathcal U(\mathfrak A)
 \right).}
 \label{eq:v18-kappa-global}
\end{equation}
This is the V14 relative component coset, now in the global controlled algebra.
It is independent of the chosen $w_0$.
\end{definition}

\begin{theorem}[Tertiary identity-component criterion]
\label{thm:v18-kappa-criterion}
Under \eqref{eq:v18-comp-zero}, the fixed primary isomorphism $v$ is realized
by a global controlled unitary in the identity component if and only if
\begin{equation}
 \boxed{
 1\in\kappa_{\rm glob}(v).}
 \label{eq:v18-kappa-zero}
\end{equation}
Equivalently, there is $z\in\mathcal U(\mathfrak B_p)$ such that
$u_0(p+z)\in\mathcal U_0(\mathfrak A)$.
\end{theorem}

\begin{proof}
This is Proposition~\ref{prop:v14-kappa-criterion} applied after the complement
Cech class has been globally descended.  Formula
\eqref{eq:v18-all-extensions} is the global form of the V14 extension torsor.
\end{proof}

If the microscopic path is not required to realize a preselected $v$, one may
change $v$ by a global zero-mode automorphism in
$\mathcal U(p\mathfrak A p)$.  The exact endpoint-pair criterion must therefore
quantify over the torsor of primary trivializations.

\subsection{The exact three-stage obstruction tower}
\label{sec:v18-tower}

Let
\begin{equation}
 \mathcal T_Z(p,q)
 \label{eq:v18-primary-torsor}
\end{equation}
denote the set of global controlled zero-mode isomorphisms
$v:p\to q$ produced by trivializations of the V17 primary Cech class.  When
nonempty it is a torsor under the global zero-mode automorphism group
$\mathcal U(p\mathfrak A p)$.

\begin{theorem}[Exact global unstabilized path criterion]
\label{thm:v18-exact-global-criterion}
Let $A$ and $B$ be V11 protected minimal fixed-index native regulators with the
same physical principal jet, representation, positive IR/reference data and
required complex/Pfaffian line datum.  A global unstabilized protected native
path from $A$ to $B$ exists if and only if there is
$v\in\mathcal T_Z(p,q)$ such that
\begin{equation}
 \boxed{
 \mathfrak C_{\rm comp}(v)=0,
 \qquad
 1\in\kappa_{\rm glob}(v).}
 \label{eq:v18-exact-global-criterion}
\end{equation}
At the zero-mode/unitary level, the obstruction tower is therefore
\begin{equation}
 \boxed{
 \mathfrak Z_{\rm ctl}
 \quad\longrightarrow\quad
 \mathfrak C_{\rm comp}
 \quad\longrightarrow\quad
 \kappa_{\rm glob}.}
 \label{eq:v18-obstruction-tower}
\end{equation}
The first is the controlled zero-mode-module descent of V17, the second is
complement descent, and the third is the unstable identity-component class.
\end{theorem}

\begin{proof}
Necessity is immediate from any global native path.  V13 local Kato transport
gives a global controlled zero-mode isomorphism $v$, so the primary class
vanishes.  The endpoint unitary obtained from the path restricts to a global
complement partial isometry, hence
$\mathfrak C_{\rm comp}(v)=0$, and lies in the identity component, hence
$1\in\kappa_{\rm glob}(v)$.

Conversely choose $v$ satisfying \eqref{eq:v18-exact-global-criterion}.  By
Theorem~\ref{thm:v18-complement-descent} there is a global unitary extension;
by Theorem~\ref{thm:v18-kappa-criterion} it may be chosen in
$\mathcal U_0(\mathfrak A)$.  The V14 unstabilized lifting theorem then
constructs the zero-mode projection path, and the V13 affine reduced-graph
interpolation completes it to a V11-admissible native fixed-index regulator
path after the standard one-time protected-chart shrinking.
\end{proof}

\begin{corollary}[Global unstabilized coefficientwise universality under the tower criterion]
\label{cor:v18-global-universality}
Under the hypotheses of Theorem~\ref{thm:v18-exact-global-criterion}, the V11
relative-forest theorem gives an invertible finite local continuum map
\begin{equation}
 F_{\infty,\mathfrak q,k}^{B\leftarrow A}
 \label{eq:v18-global-cont-map}
\end{equation}
with
\begin{equation}
 \boxed{
 \widehat{\mathfrak C}^{B,(k),\rm sat}_{\infty,\lambda}
 =F_{\infty,\mathfrak q,k}^{B\leftarrow A}
  \widehat{\mathfrak C}^{A,(k),\rm sat}_{\infty,\lambda}.}
 \label{eq:v18-global-continuum}
\end{equation}
No spectator is introduced.  The equality is projective in perturbative order
and contains every retained connected/source/BV/contact/line coefficient.
\end{corollary}

\subsection{The stable theorem versus the unstabilized tower}
\label{sec:v18-stable-vs-tower}

The exact tower also clarifies what the V13 spectator does globally.

\begin{theorem}[One stabilization annihilates every secondary projection-pair obstruction]
\label{thm:v18-stabilization-annihilates}
Assume only that the V17 primary class vanishes, so that one global controlled
$v:p\to q$ exists.  Then, regardless of
$\mathfrak C_{\rm comp}(v)$ and $\kappa_{\rm glob}(v)$, the stabilized
projections $p\oplus0$ and $q\oplus0$ are connected by the explicit
identity-component path of Theorem~\ref{thm:v18-stable-rotation}.  Therefore:
\begin{enumerate}[label=\textup{(S\arabic*)},leftmargin=3em]
\item the V13 globally stabilized native path exists once the primary module is
descended;
\item the secondary and tertiary obstructions are purely unstable;
\item no stable controlled $K_1$ invariant of the projection pair can detect
them.
\end{enumerate}
\end{theorem}

\begin{proof}
Theorem~\ref{thm:v18-stable-rotation} constructs the required path directly
from $v$ and never uses a complement map.  The three consequences are immediate.
\end{proof}

This theorem explains why the V14 $S^6$ complement defect and the V13 spectator
are compatible: the former obstructs cancellation in the original finite
carrier, while the latter moves the problem into a matrix amplification where
that unstable defect is bypassed by an explicit rotation.

\subsection{Ordinary finite-bundle shadows and the large-cutoff stable range}
\label{sec:v18-ordinary-shadow}

Forgetting locality from $\mathfrak C_{\rm comp}(v)$ gives the ordinary
isomorphism obstruction for the complement bundles.  On a finite-dimensional
background panel this ordinary obstruction disappears in the V12 large ambient
stable range whenever the zero-mode bundles have already been matched.  V18
therefore sharpens V14's large-cutoff conclusion.

\begin{proposition}[What can remain after ordinary stable-range cancellation]
\label{prop:v18-after-stable-range}
Let $B$ be a finite CW background/source panel of dimension $d$ and suppose the
finite-cutoff ambient chiral rank is sufficiently large for the V12 stable
range.  Assume the V17 primary controlled class is zero.  Then:
\begin{enumerate}[label=\textup{(R\arabic*)},leftmargin=3em]
\item the ordinary complement bundles of $p$ and $q$ are isomorphic;
\item the ordinary topological shadow of
$\mathfrak C_{\rm comp}(v)$ vanishes;
\item any surviving nonzero complement Cech class or tertiary component class
is genuinely controlled/spatial/source-locality information, not merely an
unstable finite-dimensional Grassmannian bundle defect.
\end{enumerate}
\end{proposition}

\begin{proof}
V12 gives homotopy of the finite Grassmannian classifying maps once the ordinary
zero-mode bundles are isomorphic in the stated stable range.  An ambient unitary
homotopy then identifies the ordinary complements.  Forgetting locality from
the controlled complement cocycle therefore gives a Cech coboundary.  Any
remaining class must lie in the kernel of this forgetful map.
\end{proof}

\subsection{Relation to the V16 microscopic loop kernel}
\label{sec:v18-loop-kernel}

The V16 loop cocycle can now be decomposed by blocks.  After primary descent,
choose the local normalized extensions
\begin{equation}
 \widetilde u_\alpha=v+w_\alpha.
 \label{eq:v18-local-block-unitary}
\end{equation}
On an overlap,
\begin{equation}
 \widetilde u_\alpha^*\widetilde u_\beta
 =p+z_{\alpha\beta},
 \label{eq:v18-loop-block}
\end{equation}
because the zero-mode block is common and all cross terms vanish.

\begin{proposition}[The secondary cocycle is the complement block of the microscopic overlap loop]
\label{prop:v18-loop-complement-block}
After the V17 primary class has been trivialized, the endpoint unitary of the
V16 closed overlap loop
$\ell_{\alpha\beta}=\Gamma_\alpha*\overline{\Gamma_\beta}$ has zero-mode block
$I_p$ and complement block
\begin{equation}
 \boxed{z_{\beta\alpha}=z_{\alpha\beta}^{-1}.}
 \label{eq:v18-loop-zblock}
\end{equation}
Thus the microscopic loop filtration is block triangular in the obstruction
order: V17 kills the zero-mode block first, and V18 identifies the next block as
$\mathfrak C_{\rm comp}$.
\end{proposition}

\begin{proof}
The V16 convention for the closed loop gives continuum/microscopic endpoint
comparison $\widetilde u_\beta^*\widetilde u_\alpha$.  Using
\eqref{eq:v18-loop-block} with reversed indices yields
$p+z_{\beta\alpha}$.  The two blocks act on the orthogonal decomposition
$p\mathcal H\oplus p^\perp\mathcal H$.
\end{proof}

If the complement cocycle is killed, these block overlap loops patch to one
global unitary extension.  The only remaining microscopic information is then
whether that global unitary lies in the identity component, namely the tertiary
class $\kappa_{\rm glob}$.

\subsection{V18 anti-circularity audit}
\label{sec:v18-audit}

\begin{enumerate}[leftmargin=2.4em]
\item \textbf{The complement problem is not defined before the zero-mode module
is globally descended.}  A fixed global $v$ is assumed throughout V18.
\item \textbf{Local unitary extensions are not compared with different
zero-mode blocks.}  Lemma~\ref{lem:v18-normalize-v} first gauge-fixes every
local extension so that its $p\to q$ block is exactly the same $v$.
\item \textbf{Complement descent and unitary-component cancellation are kept
separate.}  The former is the Cech class
$\mathfrak C_{\rm comp}(v)$; the latter is the coset
$\kappa_{\rm glob}(v)$ only after a global complement map exists.
\item \textbf{Stable $K_1$ is not overclaimed.}  The explicit $2\times2$
rotation proves that the stabilized projection-pair class is already zero.
Any surviving obstruction is unstable and can disappear after one
amplification.
\item \textbf{The V14 counterexample is retained at its exact layer.}  Its
primary zero-mode class is trivial, while its complement bundle supplies a
nonzero secondary class.
\item \textbf{The finite-dimensional stable range is not confused with
controlled locality.}  Once the ordinary complement bundle obstruction has
vanished, a controlled Cech or component class may still remain.
\item \textbf{The continuum theorem is not weakened.}  V16's global physical
scheme-groupoid descent and V13's stable continuum equivalence remain valid
whether or not the secondary/tertiary microscopic classes vanish.
\end{enumerate}

\section{V18 status ledger}
\label{sec:v18-status}

\subsection*{Proved in V18}
\begin{enumerate}[leftmargin=2.2em]
\item After the V17 primary controlled zero-mode class is trivialized and one
global partial isometry $v:p\to q$ is fixed, every local V15 unitary comparison
can be normalized so that its zero-mode block is exactly $v$.
\item The remaining local complement partial isometries define the exact
nonabelian Cech cocycle
$z_{\alpha\beta}=w_\alpha^*w_\beta$ in the unitary sheaf of the global
complement corner $p^\perp\mathfrak A p^\perp$.  Its class
$\mathfrak C_{\rm comp}(v)$ is independent of local extension choices and
vanishes exactly when the complements admit one global controlled partial
isometry, equivalently when the fixed $v$ has a global controlled unitary
extension.
\item The V14 propagation-zero $S^6$ example has trivial primary zero-mode
descent but nontrivial secondary complement descent; hence the two layers are
genuinely distinct.
\item Every global controlled zero-mode partial isometry admits an explicit
$2\times2$ identity-component stabilized rotation.  Therefore the endpoint
projection pair has no stable controlled $K_1$ obstruction once the primary
module has been descended; all remaining obstructions are unstable.
\item When the complement Cech class vanishes, the set of global extensions of
a fixed $v$ is a torsor under the global complement unitary group.  The exact
tertiary obstruction to an identity-component extension is the global V14
relative component coset $\kappa_{\rm glob}(v)$.
\item A global unstabilized native path exists exactly when some primary
trivialization $v$ of the V17 class has both
$\mathfrak C_{\rm comp}(v)=0$ and
$1\in\kappa_{\rm glob}(v)$.  This gives the exact three-stage microscopic
obstruction tower
$\mathfrak Z_{\rm ctl}\to\mathfrak C_{\rm comp}\to\kappa_{\rm glob}$.
\item In the V12 large-cutoff finite-panel stable range, the ordinary complement
bundle shadow is already trivial once the zero-mode bundles are matched.  Any
remaining secondary/tertiary obstruction is therefore genuinely controlled or
spatial.
\item After primary descent, the complement Cech cocycle is exactly the
complement block of the V16 microscopic overlap loop; killing it leaves only
the global unitary-component obstruction.
\end{enumerate}

\subsection*{Inherited}
\begin{enumerate}[leftmargin=2.2em]
\item V1--V7: source-complete coefficientwise continuum, common BV restoration,
fixed-loop forests, ultraviolet Cauchy convergence, finite scheme/reference
covariance and exact loop-order projectivity.
\item V8--V12: regulator-marked two-scale comparison, fixed-index shell
calculus, global zero-mode bundle/line transport and faithful complex
saturation.
\item V13--V15: exponentially local protected zero modes, controlled module
classification, stable lifting, exact unstabilized $U_0$ criterion and local
coefficientwise internal-buffer cancellation.
\item V16: global physical continuum descent through the finite scheme groupoid
and the microscopic overlap-loop kernel.
\item V17: primary controlled zero-mode Cech descent and its relative
controlled $K_0$ shadow.
\end{enumerate}

\subsection*{Conditional}
\begin{enumerate}[leftmargin=2.2em]
\item V18 begins only after the V17 primary class has been trivialized.  If no
global controlled $v:p\to q$ exists, the secondary complement corner is not a
single globally defined comparison problem.
\item The complement Cech and component classes are defined in the chosen
positive-exponent controlled algebra/source packet.  A loss of the reduced gap
or locality assumptions changes the category.
\item Stable vanishing does not imply unstabilized vanishing.  The explicit
$2\times2$ rotation proves only that amplification removes the projection-pair
obstruction.
\item The V12 real Pfaffian mod-two orientation character remains independent.
\item All continuum statements remain coefficientwise at fixed finite
$(L,\Delta,\varrho,N)$ and positive IR/reference scale.
\end{enumerate}

\subsection*{Open beyond V18}
\begin{enumerate}[leftmargin=2.2em]
\item Determine whether the actual large-cutoff Renewal ESM complement corner
automatically has the controlled cancellation and unitary-component properties
needed to kill $\mathfrak C_{\rm comp}$ and $\kappa_{\rm glob}$ after primary
zero-mode descent.  V15 proves this locally at every finite source order; V18
isolates the global descent which remains.
\item If automatic global cancellation fails, construct a genuinely
spacetime-local witness of nonzero $\mathfrak C_{\rm comp}$ or
$\kappa_{\rm glob}$, rather than only an unstable finite-dimensional bundle
example.
\item Relate the tertiary component coset, when stable-range hypotheses make
such a reduction valid, to a computable controlled/coarse unitary invariant of
the complement corner.
\item Compute the independent real Pfaffian mod-two spectral-flow character and
analyze exits from the minimal fixed-index stratum.
\item The positive-IR removal and nonperturbative/all-orders directions remain
independent stronger programmes.
\end{enumerate}

\subsection*{Exact next load-bearing theorem}
\begin{center}
\fbox{\parbox{0.92\textwidth}{
\textbf{Global automatic complement cancellation in the Renewal ESM carrier.}
Assume the V17 primary controlled zero-mode class vanishes.  Use the actual
large-cutoff geometry of the Renewal Standard-Model hard complement, together
with the V15 finite-rank/local-source reserve construction, to descend one
complement buffer globally over a connected protected background component.
Prove that $\mathfrak C_{\rm comp}(v)$ is a coboundary for some primary
trivialization $v$, and then prove that the resulting global extension can be
chosen in $\mathcal U_0$; or exhibit an explicit native controlled complement
or unitary-component witness.  Because V18 proves that all stable projection
and stable $K_1$ shadows vanish after primary descent, any nonzero witness must
be genuinely unstable, spatial/source-local, and invisible to the already
closed coefficientwise continuum.}}
\end{center}

\section*{Append-only continuation marker: V18}
\addcontentsline{toc}{section}{Append-only continuation marker: V18}
\textbf{End of V18.}  The complete V1--V17 body above is retained verbatim.
V18 assumes the primary controlled zero-mode module descent of V17 has been
trivialized and resolves the next two microscopic layers.  After normalizing
all local unitary comparisons to one fixed global zero-mode isomorphism, their
complement blocks form a second exact nonabelian Cech cocycle; its class
vanishes exactly when the complements glue to one global controlled partial
isometry and hence the fixed zero-mode map extends to a global unitary.  If it
vanishes, the remaining obstruction is the global V14 relative
identity-component coset.  A global unstabilized native path exists exactly
when some primary trivialization kills both secondary classes.  In contrast,
an explicit $2\times2$ controlled rotation shows that the projection-pair
obstruction is identically zero after one stabilization, so no stable
controlled $K_1$ class survives primary descent.  The remaining problem is
therefore genuinely unstable global cancellation in the actual Renewal ESM
complement, not a stable continuum phase.



\clearpage
\section{V19 continuation: global finite-cover reserve, automatic complement cancellation, and collapse of the secondary tower}
\label{sec:v19-continuation}

V18 showed that, after the primary controlled zero-mode module has been
identified globally, an unstabilized microscopic path can still fail at two
secondary stages: complement Cech descent and the identity component of the
resulting ambient unitary.  V15, on the other hand, proved that on every
protected local source chart a finite-rank zero-mode packet has enough native
ultraviolet room to manufacture an internal buffer at every prescribed finite
source order.  The V18 handoff asked whether those local buffers descend
through a noncontractible background component.

The answer at the coefficientwise quantifiers of the continuum theorem is
positive.  The key observation is that a finite good cover introduces only
finitely many additional overlap constraints.  One can therefore construct the
V15 local buffers \emph{sequentially}, requiring each new buffer to be
orthogonal not only to the two endpoint zero-mode packets but also to every
previous buffer on the region where the two chart weights coexist.  The same
finite translated-copy dimension count closes after increasing its multiplicity
by a fixed amount depending on the cover and the finite source algebra.

A square partition of unity then glues the local partial isometries themselves,
not merely their range projections.  If
\(a_\alpha:p\to r_\alpha\) are the resulting pairwise-orthogonal local buffer
isometries and
\begin{equation}
 \sum_\alpha\lambda_\alpha^2=1,
 \label{eq:v19-square-partition}
\end{equation}
then
\begin{equation}
 \boxed{
 A:=\sum_\alpha\lambda_\alpha a_\alpha
 }
 \label{eq:v19-global-A}
\end{equation}
is one global controlled isometry with initial projection \(p\).  Once the
V17 primary class is zero and a global controlled
\(v:p\to q\) has been fixed, the second isometry
\begin{equation}
 \boxed{B:=Av^*}
 \label{eq:v19-global-B}
\end{equation}
has initial projection \(q\) and exactly the same final projection as \(A\).
Thus the local reserves glue to one genuine global V14 internal buffer in the
\emph{original} lattice carrier.

The resulting V14 two-rotation path extends the chosen \(v\) itself and lies in
the identity component.  Consequently the V18 secondary and tertiary classes
vanish automatically at every prescribed finite coefficient/source order once
the V17 primary controlled zero-mode class has vanished.  The microscopic
global obstruction tower therefore collapses, coefficientwise, to its primary
stage (plus the independent real Pfaffian orientation datum).

Throughout V19 fix a compact connected protected background component
\(\mathcal B\) which admits a finite good cover
\begin{equation}
 \mathfrak U=(U_1,\ldots,U_N),
 \label{eq:v19-good-cover}
\end{equation}
and fix one finite perturbative/source packet
\begin{equation}
 \mathfrak q=(L,\Delta,\varrho,N_{\rm ext},\mathcal P_{\rm ext}).
 \label{eq:v19-q}
\end{equation}
Let
\(\mathbf R_\varrho\) be the complete finite graded BV/source algebra from V15
and put
\begin{equation}
 \nu_\varrho:=\dim_{\mathbb C}\mathbf R_\varrho.
 \label{eq:v19-nu}
\end{equation}
All source identities below are identities in this truncated algebra.  The
large-cutoff threshold is permitted to depend on
\((\mathfrak q,\mathfrak U,|k|)\), as required by the coefficientwise theorem.

\subsection{A finite-constraint strengthening of the V15 reserve theorem}
\label{sec:v19-multiavoid}

The V15 dimension count never used the fact that there were exactly two
families to avoid.  We make the finite-constraint extension explicit.

\begin{theorem}[Finite-constraint local reserve]
\label{thm:v19-finite-constraint-reserve}
Let \(U\) be one protected framed source chart and let
\(p,q\) be the two rank-\(\kappa\) endpoint zero-mode packets on \(U\), with a
controlled local isomorphism \(v:p\to q\).  In addition, let
\begin{equation}
 c_1,\ldots,c_s:p\longrightarrow\mathcal H_a
 \label{eq:v19-extra-constraints}
\end{equation}
be a finite family of controlled source-jet maps of rank at most \(\kappa\);
they need not be isometries.  Then for every fixed source order \(\varrho\),
for all sufficiently large cutoffs, there is a controlled isometry
\begin{equation}
 a:p\longrightarrow r
 \label{eq:v19-local-a}
\end{equation}
with rank-\(\kappa\) final projection \(r=aa^*\) satisfying
\begin{align}
 pa&=0,
 &qa&=0,
 \label{eq:v19-pq-a-zero}\\
 c_j^*a&=0,
 &&1\le j\le s,
 \label{eq:v19-c-a-zero}
\end{align}
through the complete retained source jet.  One may take a translated-copy
multiplicity
\begin{equation}
 \boxed{
 M\ge (2+s)\nu_\varrho+1.}
 \label{eq:v19-M-multi}
\end{equation}
All locality/source bounds remain cutoff-uniform and the added translation
radius remains bounded in lattice units.
\end{theorem}

\begin{proof}
Use the V15 localized Kato frame
\(\Psi:\mathbb C^\kappa\to\mathcal H_a\) for \(p\), the corresponding
\(q\)-frame \(\mathsf X=v\Psi\), and the uniformly conditioned translated-copy
isometry
\begin{equation}
 W:\mathbb C^{M\kappa}\longrightarrow\mathcal H_a
 \label{eq:v19-W}
\end{equation}
from Corollary~\ref{cor:v15-copy-bank}, with its source-jet orthonormalization.
For each extra constraint put
\begin{equation}
 C_j:=c_j\Psi:\mathbb C^\kappa\longrightarrow\mathcal H_a.
 \label{eq:v19-Cj}
\end{equation}
The complete overlap map is
\begin{equation}
 F=
 \begin{pmatrix}
  \Psi^*W\\
  \mathsf X^*W\\
  C_1^*W\\
  \vdots\\
  C_s^*W
 \end{pmatrix}
 :\mathbb C^{M\kappa}
 \longrightarrow
 \mathbb C^{(2+s)\kappa}\otimes\mathbf R_\varrho.
 \label{eq:v19-Fmulti}
\end{equation}
Choose a homogeneous vector-space basis of
\(\mathbf R_\varrho\) and stack all coefficient maps.  This gives an ordinary
linear map
\begin{equation}
 \mathbf F:
 \mathbb C^{M\kappa}
 \longrightarrow
 \mathbb C^{(2+s)\kappa\nu_\varrho}.
 \label{eq:v19-stacked-F}
\end{equation}
Hence
\begin{equation}
 \dim\ker\mathbf F
 \ge
 M\kappa-(2+s)\kappa\nu_\varrho
 \ge\kappa.
 \label{eq:v19-dim-kernel}
\end{equation}
Choose an orthonormal \(\kappa\)-frame
\(Z:\mathbb C^\kappa\to\ker\mathbf F\), set
\begin{equation}
 \mathsf E:=WZ,
 \qquad
 a:=\mathsf E\Psi^*,
 \qquad
 r:=\mathsf E\mathsf E^*,
 \label{eq:v19-Ear}
\end{equation}
and repeat the V15 calculation.  Because \(W^*W=I\), one has
\(a^*a=p\) and \(aa^*=r\).  The stacked-kernel condition gives
\(\Psi^*\mathsf E=0\), \(\mathsf X^*\mathsf E=0\), and
\(C_j^*\mathsf E=0\), which are exactly
\eqref{eq:v19-pq-a-zero}--\eqref{eq:v19-c-a-zero} after multiplying by the
appropriate frame maps.  The locality and bounded-displacement conclusions are
unchanged because only finitely many rows were added to the finite coefficient
matrix.
\end{proof}

\begin{corollary}[Finite orthogonal family of local reserves]
\label{cor:v19-finite-orthogonal-reserves}
On one framed protected chart, for every fixed integer \(R\), one can construct
rank-\(\kappa\) controlled buffer isometries
\begin{equation}
 a_1,\ldots,a_R:p\longrightarrow\mathcal H_a
 \label{eq:v19-many-a}
\end{equation}
which are each orthogonal to \(p\) and \(q\) and satisfy
\begin{equation}
 \boxed{a_i^*a_j=\delta_{ij}p}
 \label{eq:v19-many-orthogonal}
\end{equation}
through the retained source order.
\end{corollary}

\begin{proof}
Construct them inductively.  When choosing \(a_j\), apply
Theorem~\ref{thm:v19-finite-constraint-reserve} with the previously constructed
maps \(c_i=a_i\), \(i<j\).  The number of constraints is finite and bounded by
\(R-1\), so one fixed translated-copy multiplicity suffices.
\end{proof}

\subsection{Square partitions of unity and overlap cutoffs}
\label{sec:v19-square-partition}

\begin{lemma}[Square partition subordinate to a finite good cover]
\label{lem:v19-square-partition}
There are smooth real functions
\begin{equation}
 \lambda_1,\ldots,\lambda_N\in C^\infty(\mathcal B),
 \qquad
 \operatorname{supp}\lambda_\alpha\Subset U_\alpha,
 \label{eq:v19-lambda-support}
\end{equation}
with
\begin{equation}
 \boxed{
 \sum_{\alpha=1}^N\lambda_\alpha^2=1.}
 \label{eq:v19-square-pou}
\end{equation}
For each ordered pair \(\beta<\alpha\), there is a smooth scalar cutoff
\(\chi_{\alpha\beta}\) supported in
\(U_\alpha\cap U_\beta\) and equal to one on a neighborhood of
\begin{equation}
 \operatorname{supp}(\lambda_\alpha\lambda_\beta).
 \label{eq:v19-active-overlap}
\end{equation}
All derivatives needed by the finite background/source panel are bounded.
\end{lemma}

\begin{proof}
Choose a finite smooth partition of unity \((\phi_\alpha)\) subordinate to a
slightly smaller cover with compactly contained supports.  Choose smooth
nonnegative \(\rho_\alpha\) with the same supports and with
\(\sum_\alpha\rho_\alpha^2>0\), and set
\begin{equation}
 \lambda_\alpha=
 \frac{\rho_\alpha}
      {(\sum_\gamma\rho_\gamma^2)^{1/2}}.
 \label{eq:v19-lambda-def}
\end{equation}
This gives \eqref{eq:v19-square-pou}.  The compact set
\(\operatorname{supp}(\lambda_\alpha\lambda_\beta)\) lies inside the open
overlap, so the usual smooth Urysohn construction gives
\(\chi_{\alpha\beta}\).  Compactness and finiteness of the cover bound all
prescribed derivatives.
\end{proof}

The partition functions act only in the background/source parameter space and
are scalar with respect to the spacetime kernel algebra.  Their retained jets
are therefore ordinary scalar coefficients in the V4--V7 finite source packet.

\subsection{Sequential overlap-orthogonal V15 buffers}
\label{sec:v19-sequential-buffers}

Assume from now on that the V17 primary class is zero and fix one global
controlled partial isometry
\begin{equation}
 \boxed{
 v^*v=p,
 \qquad
 vv^*=q.}
 \label{eq:v19-global-v}
\end{equation}
Here \(p\) and \(q\) denote the global kernel projections for positive index;
the negative-index cokernel case is the chiral dual and is treated identically.

\begin{theorem}[Sequential local buffers orthogonal on active overlaps]
\label{thm:v19-sequential-buffers}
For every fixed packet \(\mathfrak q\) and all sufficiently large cutoffs, one
can choose on each chart \(U_\alpha\) a controlled isometry
\begin{equation}
 a_\alpha:p|_{U_\alpha}\longrightarrow r_\alpha
 \label{eq:v19-a-alpha}
\end{equation}
with
\begin{align}
 a_\alpha^*a_\alpha&=p,
 \label{eq:v19-aa-p}\\
 p a_\alpha&=0,
 &q a_\alpha&=0,
 \label{eq:v19-local-buffer-orth}
\end{align}
and such that for every \(\alpha\ne\beta\),
\begin{equation}
 \boxed{
 a_\alpha^*a_\beta=0
 \quad\text{on a neighborhood of }
 \operatorname{supp}(\lambda_\alpha\lambda_\beta).}
 \label{eq:v19-cross-zero}
\end{equation}
All identities hold in the truncated source algebra through order \(\varrho\).
The locality constants and the required lattice translation radius are uniform
in the cutoff.
\end{theorem}

\begin{proof}
Order the finite cover.  Construct \(a_1\) by the ordinary V15 reserve theorem.
Suppose \(a_1,\ldots,a_{\alpha-1}\) have been constructed.  For each
\(\beta<\alpha\), choose the cutoff \(\chi_{\alpha\beta}\) from
Lemma~\ref{lem:v19-square-partition}.  The product
\begin{equation}
 c_{\alpha\beta}:=
 \chi_{\alpha\beta}a_\beta
 \label{eq:v19-c-alpha-beta}
\end{equation}
may be viewed on \(U_\alpha\): because the cutoff has support compactly inside
\(U_\alpha\cap U_\beta\), the product extends smoothly by zero to the rest of
\(U_\alpha\).  It is a controlled finite-source-jet map of rank at most
\(\kappa\).

Apply Theorem~\ref{thm:v19-finite-constraint-reserve} on \(U_\alpha\), using
all maps \(c_{\alpha\beta}\), \(\beta<\alpha\), as the extra constraints.  A
single multiplicity
\begin{equation}
 M\ge (N+1)\nu_\varrho+1
 \label{eq:v19-uniform-M-cover}
\end{equation}
works for every chart.  The resulting \(a_\alpha\) obeys
\begin{equation}
 a_\alpha^*c_{\alpha\beta}=0.
 \label{eq:v19-a-c-zero}
\end{equation}
Where \(\lambda_\alpha\lambda_\beta\ne0\), the cutoff equals one on a
neighborhood, so \eqref{eq:v19-a-c-zero} is exactly
\(a_\alpha^*a_\beta=0\) there, including every retained derivative.  Taking
adjoints gives the reverse order.  The cover is finite, so the maximum of the
finitely many local thresholds and locality constants is still cutoff
uniform.
\end{proof}

\subsection{Global gluing of the internal buffer}
\label{sec:v19-global-buffer}

Extend each product \(\lambda_\alpha a_\alpha\) by zero outside \(U_\alpha\)
and define the global controlled operator
\begin{equation}
 \boxed{
 A:=\sum_{\alpha=1}^N\lambda_\alpha a_\alpha.}
 \label{eq:v19-A-global}
\end{equation}

\begin{theorem}[Global coefficientwise internal buffer]
\label{thm:v19-global-buffer}
For every fixed coefficient packet and all sufficiently large cutoffs,
\eqref{eq:v19-A-global} satisfies
\begin{align}
 A^*A&=p,
 \label{eq:v19-AstarA}\\
 pA&=0,
 &qA&=0.
 \label{eq:v19-pqA-zero}
\end{align}
Hence
\begin{equation}
 r:=AA^*
 \label{eq:v19-r-global}
\end{equation}
is one globally defined rank-\(\kappa\) controlled projection orthogonal to
both endpoint zero-mode packets.  If
\begin{equation}
 B:=Av^*,
 \label{eq:v19-B-global2}
\end{equation}
then
\begin{equation}
 \boxed{
 B^*B=q,
 \qquad
 BB^*=r.}
 \label{eq:v19-B-buffer}
\end{equation}
Thus \(r\) is a global V14 internal buffer for the \emph{fixed primary
trivialization} \(v\).
\end{theorem}

\begin{proof}
Using Theorem~\ref{thm:v19-sequential-buffers},
\begin{align}
 A^*A
 &=\sum_{\alpha,\beta}
   \lambda_\alpha\lambda_\beta a_\alpha^*a_\beta
 \nonumber\\
 &=\sum_\alpha\lambda_\alpha^2p
 =p
 \label{eq:v19-AstarA-proof}
\end{align}
because every cross term vanishes wherever its scalar coefficient is nonzero
and \(\sum\lambda_\alpha^2=1\).  Likewise
\(pA=qA=0\) chartwise, so
\eqref{eq:v19-pqA-zero} holds globally.  Therefore \(A\) is a partial isometry
with initial projection \(p\) and final projection \(r=AA^*\).

For \(B=Av^*\), use
\(A^*A=p\), \(v^*v=p\), and \(vv^*=q\):
\begin{align}
 B^*B
 &=vA^*Av^*=vpv^*=vv^*=q,
 \label{eq:v19-BstarB-proof}\\
 BB^*
 &=Av^*vA^*=ApA^*=AA^*=r.
 \label{eq:v19-BBstar-proof}
\end{align}
Every operation is a finite sum/product in the controlled source algebra, so
all retained source derivatives and locality bounds are global and uniform.
\end{proof}

\begin{corollary}[The buffer remains inside the reduced massive packet]
\label{cor:v19-buffer-massive}
At both endpoint regulators, the global buffer projection \(r\) lies in the
orthogonal complement of the protected kernel/cokernel zero-mode packet.
Therefore the V11 reduced singular-value floor controls it as a nonzero/massive
internal direction.  The construction adds no new species, matrix
amplification, or spectator carrier.
\end{corollary}

\begin{proof}
This is immediate from \eqref{eq:v19-pqA-zero}: the range of \(r\) is contained
in the orthogonal complement of both zero-mode ranges.  Apply the V11 reduced
floor on those complements.
\end{proof}

\subsection{The chosen primary isomorphism extends globally in the identity component}
\label{sec:v19-extend-v}

Use the V14 internal-buffer rotations associated with \(A:p\to r\) and
\(B:q\to r\):
\begin{align}
 U_A(\theta)
 &=I-p-r+\cos\theta\,(p+r)+\sin\theta\,(A-A^*),
 \label{eq:v19-UA}\\
 U_B(\theta)
 &=I-q-r+\cos\theta\,(q+r)+\sin\theta\,(B-B^*).
 \label{eq:v19-UB}
\end{align}

\begin{theorem}[Automatic global extension of the primary map]
\label{thm:v19-global-extension}
Define
\begin{equation}
 \boxed{
 u:=U_B(\pi/2)^*U_A(\pi/2).}
 \label{eq:v19-u-global}
\end{equation}
Then
\begin{equation}
 \boxed{
 u\in\mathcal U_0(\mathfrak A),
 \qquad
 up=v,
 \qquad
 upu^*=q.}
 \label{eq:v19-u-properties}
\end{equation}
Consequently, for the fixed primary trivialization \(v\),
\begin{equation}
 \boxed{
 \mathfrak C_{\rm comp}(v)=0,
 \qquad
 1\in\kappa_{\rm glob}(v).}
 \label{eq:v19-secondary-tertiary-zero}
\end{equation}
Thus both V18 secondary stages vanish automatically, coefficientwise, on the
large-cutoff Renewal lattice once the V17 primary class has been trivialized.
\end{theorem}

\begin{proof}
Each of \(U_A(\theta)\) and \(U_B(\theta)\) is a controlled unitary path from
the identity by the V14 buffer-rotation theorem, so their endpoint product
\eqref{eq:v19-u-global} lies in the identity component.  It conjugates \(p\) to
\(q\) because the first rotation moves \(p\) to \(r\) and the inverse second
rotation moves \(r\) to \(q\).

It remains to identify the zero-mode block.  On \(p\mathcal H\), the first
endpoint rotation acts as \(A\), while on \(r\mathcal H\) the inverse second
rotation acts as \(B^*\).  Hence
\begin{equation}
 up=B^*A.
 \label{eq:v19-up-BA}
\end{equation}
Using \(B=Av^*\) and \(A^*A=p\),
\begin{equation}
 B^*A=vA^*A=vp=v.
 \label{eq:v19-BA-v}
\end{equation}
Thus \(u\) extends the chosen \(v\) itself.  The V18 complement class therefore
vanishes by Theorem~\ref{thm:v18-complement-descent}, and the tertiary class
contains the identity by Theorem~\ref{thm:v18-kappa-criterion}.
\end{proof}

\begin{corollary}[Collapse of the V18 tower after primary descent]
\label{cor:v19-tower-collapse}
At every prescribed finite coefficient/source order, for all sufficiently large
cutoffs, the V18 obstruction tower reduces to
\begin{equation}
 \boxed{
 \mathfrak Z_{\rm ctl}
 \quad\Longrightarrow\quad
 \mathfrak C_{\rm comp}=0
 \quad\Longrightarrow\quad
 \kappa_{\rm glob}=0,}
 \label{eq:v19-tower-collapse}
\end{equation}
where the last two implications hold once the primary class has been
trivialized.  Therefore the only complex microscopic descent obstruction left
at this coefficientwise level is the V17 controlled zero-mode module class.
\end{corollary}

\subsection{Global unstabilized native lifting after primary descent}
\label{sec:v19-global-native-lift}

\begin{theorem}[Global unstabilized native lifting after primary descent]
\label{thm:v19-global-native-lift}
Let the two endpoint regulators satisfy the V11 protected minimal fixed-index
hypotheses on the compact connected component \(\mathcal B\), and fix one finite
coefficient packet \(\mathfrak q\).  Assume:
\begin{enumerate}[label=\textup{(G\arabic*)},leftmargin=3em]
\item the physical Standard-Model principal Dirac--Yukawa jet and representation
packet agree;
\item the V17 primary controlled zero-mode Cech class vanishes, so a global
controlled \(v:p\to q\) exists;
\item the required V12 complex determinant-line/saturation datum agrees, and
when a real Pfaffian sector is present its declared orientation datum agrees.
\end{enumerate}
Then for all sufficiently large cutoffs there is a global unstabilized
coefficientwise V11-admissible native regulator path from the first regulator
to the second, with no auxiliary spectator copy.
\end{theorem}

\begin{proof}
Theorem~\ref{thm:v19-global-extension} gives a global controlled
\(u\in\mathcal U_0\) conjugating the zero-mode packet and extending the chosen
primary map.  Choose a controlled unitary path \(u_t\) from \(I\) to \(u\) and
put \(p_t=u_tpu_t^*\).  The V13 controlled-reduction theorem transports the
initial affine reduced graph along \(p_t\) and then linearly interpolates it to
the target graph after the zero-mode packet has reached \(q\).  All graph
coordinates remain exponentially local by the V13 reconstruction formulas.
The physical principal jet is fixed throughout.

The path is compact in the finite coefficient/source topology.  The V11
reduced singular-value floor and the full Yukawa/source pivot are open
conditions, so the standard one-time shrinking of the protected chart retains
one positive reduced floor.  The V12 line/saturation packet is transported
along the same path.  Hence every segment satisfies the fixed-index
admissibility conditions without any stabilization.
\end{proof}

\begin{theorem}[Global coefficientwise regulator universality after primary descent]
\label{thm:v19-global-universality}
Under Theorem~\ref{thm:v19-global-native-lift}, for every fixed packet
\(\mathfrak q\) there is an invertible finite local continuum matching map
\begin{equation}
 F_{\infty,\mathfrak q,k}^{B\leftarrow A}
 \label{eq:v19-continuum-map}
\end{equation}
such that
\begin{equation}
 \boxed{
 \widehat{\mathfrak C}^{B,(k),\rm sat}_{\infty,\lambda}
 =F_{\infty,\mathfrak q,k}^{B\leftarrow A}
  \widehat{\mathfrak C}^{A,(k),\rm sat}_{\infty,\lambda}.}
 \label{eq:v19-global-universality}
\end{equation}
No microscopic stabilization is used.  The equality contains every retained
connected/source coefficient, operator-mixing coefficient, stress/contact
writer, ghost/antifield descendant, BV/BRST/master relation and
Determinant/Pfaffian insertion, and is strictly projective in perturbative
order.
\end{theorem}

\begin{proof}
Apply the V11 relative-forest and two-scale comparison theorem along the global
native path of Theorem~\ref{thm:v19-global-native-lift}.  V12 supplies the
global complex saturation on the faithful branch and the declared independent
real orientation packet where applicable.  The V7 chronological construction
makes every lower-order matching immutable when a higher loop order is added,
so the continuum map is projective in terminal perturbative order.
\end{proof}

\subsection{Spatial ultraviolet reserve versus matrix stabilization}
\label{sec:v19-spatial-vs-matrix}

The V19 construction explains why the V14 finite-dimensional complement
counterexample does not contradict the positive large-cutoff theorem.

\begin{proposition}[The lattice itself supplies the unstable reserve]
\label{prop:v19-spatial-reserve}
At fixed zero-mode rank, finite source order and finite cover size, the number
of translated copies used in the V19 construction is finite and independent of
the cutoff.  They are copies of the \emph{same} native zero-mode packet moved by
bounded gauge/spin-covariant lattice shifts.  Hence the large-cutoff lattice
provides the reserve which the V14 finite-dimensional example lacked, without
adding an internal species or passing to a matrix-amplified theory.
\end{proposition}

\begin{proof}
The largest multiplicity is bounded by
\begin{equation}
 M_{\rm glob}\ge (N+1)\nu_\varrho+1,
 \label{eq:v19-Mglobal}
\end{equation}
and the V15 shift-selection lemma gives a maximal displacement depending only
on \((\kappa,M_{\rm glob})\) and the fixed exponential-locality constants.
Neither quantity depends on the cutoff.  The only role of the ultraviolet limit
is eventually to make the periodic lattice large enough to accommodate these
finitely many translated packets without wraparound conflicts.
\end{proof}

\begin{corollary}[The V14 $S^6$ defect is removed after spatial embedding at large cutoff]
\label{cor:v19-S6-spatial}
If the propagation-zero V14 $S^6$ endpoint pair is embedded as a finite-rank
packet in a sufficiently large Renewal lattice while the V19 locality/source
hypotheses hold, the spatially translated internal reserve removes its
unstable complement defect in the original enlarged lattice carrier.  This does
not contradict the finite-dimensional counterexample: the ambient carrier has
changed by adding physical lattice sites, not by asserting cancellation in the
original finite onsite fibre.
\end{corollary}

\subsection{What remains: the primary controlled zero-mode module}
\label{sec:v19-primary-remains}

V19 closes the secondary and tertiary layers only after a global controlled
zero-mode isomorphism has been supplied.  That hypothesis is substantive.
V12 already showed that determinant-line triviality does not determine a
rank-\(\kappa\) zero-mode bundle when \(\kappa>1\), and V17 strengthened the
ordinary bundle class to the controlled local module class.

\begin{theorem}[Exact coefficientwise global path criterion after V19]
\label{thm:v19-exact-criterion}
Fix a protected compact component, one finite coefficient/source packet and all
V11 reduced-gap assumptions.  On the faithful complex Standard-Model branch,
for all sufficiently large cutoffs a global unstabilized coefficientwise native
regulator path exists if and only if:
\begin{enumerate}[label=\textup{(P\arabic*)},leftmargin=3em]
\item the V17 primary controlled zero-mode module Cech class vanishes; and
\item every independently physical real Pfaffian line, if present, has the
same declared orientation component at the two endpoints.
\end{enumerate}
The V18 complement Cech class and relative unitary-component class impose no
additional obstruction at this coefficientwise level.
\end{theorem}

\begin{proof}
Necessity of the primary controlled zero-mode class is V17 Kato transport along
any global native path.  The real orientation necessity is V12.  Sufficiency is
Theorems~\ref{thm:v19-global-buffer}--\ref{thm:v19-global-native-lift}, which
construct one global internal reserve and an identity-component extension of
any chosen primary trivialization.
\end{proof}

\begin{firewall}[V19 does not trivialize the primary index bundle]
The global reserve is constructed \emph{after} a global controlled
\(v:p\to q\) has been fixed.  The square-partition gluing never identifies two
nonisomorphic zero-mode bundles and never changes the family index class.  Thus
V19 cannot be used to erase a nonzero V17 primary class, a nontrivial ordinary
kernel/cokernel bundle, or an independent real Pfaffian orientation character.
\end{firewall}

\subsection{Relation to V16 physical descent}
\label{sec:v19-v16}

Even when the primary microscopic class does not vanish, V16's physical result
remains stronger than a single-path statement: all coordinate-free normalized
continuum relations still descend through the finite scheme groupoid.  V19
therefore yields the final hierarchy
\begin{center}
\fbox{\parbox{0.90\textwidth}{\centering
\textbf{Global hierarchy.}\quad
Coordinate-free coefficientwise physical continuum equivalence holds throughout
the V16 protected component.  A single global coefficientwise microscopic
native path exists precisely when the V17 primary controlled zero-mode class
vanishes, together with the independently required real orientation datum when
such a Pfaffian sector is present.}}
\label{eq:v19-final-hierarchy}
\end{center}
The V18 secondary/tertiary classes are genuine in an arbitrary finite carrier,
but the large-cutoff spatial reserve makes them non-load-bearing for the fixed
coefficientwise Renewal ESM theorem.

\section{V19 anti-circularity audit}
\label{sec:v19-audit}

\begin{enumerate}[leftmargin=2.4em]
\item \textbf{The primary zero-mode class is assumed before complement descent.}
V19 does not use a local frame patch to manufacture a global zero-mode
isomorphism; it starts from the V17 global \(v\).
\item \textbf{The global buffer is not chosen by a discontinuous chart
selection.}  The local ranges are made pairwise orthogonal precisely where the
partition weights overlap, so the weighted sum of partial isometries is itself
an exact global isometry.
\item \textbf{The partition of unity is squared.}  The identity
\(\sum\lambda_\alpha^2=1\), rather than
\(\sum\lambda_\alpha=1\), is what makes \(A^*A=p\) exactly.
\item \textbf{Cross terms vanish before gluing.}  Sequential use of the V15
dimension count treats every previously active buffer as an additional finite
source-jet constraint; no approximate orthogonality is promoted to an exact
projection statement.
\item \textbf{The chosen primary map is extended, not replaced.}  Setting
\(B=Av^*\) forces the endpoint V14 rotation to satisfy \(up=v\) exactly.
Therefore V19 kills \(\mathfrak C_{\rm comp}(v)\) and
\(\kappa_{\rm glob}(v)\) for the same primary trivialization used in V17.
\item \textbf{Large cutoff does not mean cancellation of the whole local
operator algebra.}  Only finitely many fixed-rank translated packets are used;
the V15 Roe-algebra firewall remains in force.
\item \textbf{The theorem remains coefficientwise.}  The number of translated
copies may depend on the finite graded source dimension \(\nu_\varrho\) and the
finite cover size; no bound uniform in source order or loop order is claimed.
\item \textbf{The complex determinant and real Pfaffian sectors stay
separate.}  V12's faithful complex line trivialization does not orient an
independent real Pfaffian line.
\end{enumerate}

\section{V19 status ledger}
\label{sec:v19-status}

\subsection*{Proved in V19}
\begin{enumerate}[leftmargin=2.2em]
\item The V15 translated-copy dimension count extends to any finite number of
additional controlled rank-\(\kappa\) source-jet constraints.  With \(s\) extra
packets, \(M\ge(2+s)\nu_\varrho+1\) copies suffice.
\item On a finite good cover, local V15 buffers can be chosen sequentially so
that their range isometries are exactly orthogonal on every region where the
corresponding square-partition weights coexist.
\item A square partition of unity glues those local maps to one global
controlled isometry
\(A=\sum\lambda_\alpha a_\alpha\) with \(A^*A=p\), whose range projection is
orthogonal to both endpoint zero-mode packets.
\item For any fixed V17 primary trivialization \(v:p\to q\), the map
\(B=Av^*\) has initial projection \(q\) and the same global final projection as
\(A\).  Hence the original lattice carrier contains one global V14 internal
buffer for that exact \(v\).
\item The associated V14 two-rotation endpoint unitary lies in the controlled
identity component and restricts to \(v\).  Therefore both the V18 complement
Cech class and the tertiary global component class vanish automatically once
the V17 primary class has been trivialized.
\item The V13 affine reduced-graph construction then gives a global
unstabilized V11-admissible native path at every fixed finite source order,
with no spectator or matrix amplification.
\item Consequently, on the faithful complex Standard-Model branch, the exact
large-cutoff coefficientwise microscopic path criterion reduces to the V17
primary controlled zero-mode module class, together with the independent real
Pfaffian orientation datum when present.
\item The unstable finite-dimensional V14 complement defect is removed in the
actual large lattice by finitely many bounded spatial translations of the same
native packet; this is spatial ultraviolet reserve, not a claim of universal
Roe-algebra cancellation.
\end{enumerate}

\subsection*{Inherited}
\begin{enumerate}[leftmargin=2.2em]
\item V1--V7: source-complete coefficientwise continuum, common BV restoration,
fixed-loop forest closure, ultraviolet Cauchy convergence, scheme/reference
covariance and exact perturbative-order projectivity.
\item V8--V12: relative regulator forests, fixed-index reduced shell calculus,
global zero-mode bundle/line transport and faithful complex saturation.
\item V13--V15: protected zero-mode locality, controlled module equivalence,
stable lifting and the local finite-source internal-reserve theorem.
\item V16: global physical continuum descent in the finite scheme groupoid.
\item V17: the primary controlled zero-mode Cech class.
\item V18: the secondary complement Cech class, tertiary identity-component
class and the exact three-stage obstruction tower in an arbitrary controlled
carrier.
\end{enumerate}

\subsection*{Conditional}
\begin{enumerate}[leftmargin=2.2em]
\item V19 assumes a compact protected component admitting one finite good cover
with uniform V11/V15 locality and reduced-gap constants.  The cover may depend
on the fixed finite coefficient/source panel but not on the cutoff after the
threshold is chosen.
\item The primary controlled zero-mode class must already vanish.  V19 does not
prove ordinary or controlled equivalence of globally different index bundles.
\item The translated-copy multiplicity and lattice displacement may depend
badly on \(|k|\), the cover size and the finite graded source dimension; no
uniform source-order or loop-order estimate is claimed.
\item The V12 real Pfaffian mod-two orientation condition remains independent.
\item All continuum conclusions retain the positive IR/reference scale and the
coefficientwise perturbative interpretation.
\end{enumerate}

\subsection*{Open beyond V19}
\begin{enumerate}[leftmargin=2.2em]
\item Attack the remaining \emph{primary} controlled zero-mode module class.
Determine whether equality of the physical elliptic principal symbol and fixed
family index forces the V17 controlled kernel/cokernel modules of two native
regulators to be globally controlled-isomorphic at sufficiently large cutoff,
or exhibit a genuine controlled family-index obstruction.
\item Relate the primary controlled module class to the ordinary family index
bundle and establish a locality-preserving finite-rank representative theorem
in the stable ambient range.
\item Compute the independent real Pfaffian mod-two spectral-flow character on
any genuinely physical Majorana sector.
\item Analyze exits from the V11 minimal fixed-index stratum, including
accidental zero-mode pair creation and genuine index change.
\item The positive-IR removal, analytic finite-coupling RG trajectory,
perturbative summability and nonperturbative directions remain separate stronger
programmes.
\end{enumerate}

\subsection*{Exact next load-bearing theorem}
\begin{center}
\fbox{\parbox{0.92\textwidth}{
\textbf{Primary controlled family-index matching.}
For two protected native Renewal ESM regulators with the same Standard-Model
representation and the same physical elliptic Dirac--Yukawa principal symbol,
compare their finite-rank zero-mode families over a connected fixed-index
background component.  Prove that equality of the family Fredholm/index class
lifts, at sufficiently large cutoff and every prescribed finite source order,
to a controlled exponentially local partial isometry between the kernel
(respectively cokernel) packets; or identify a computable controlled refinement
of the family index which obstructs such a lift.  Combined with V19, a positive
theorem would eliminate the last complex microscopic regulator-path obstruction
and leave only the independently real Pfaffian orientation/zero-crossing
problems.}}
\end{center}

\section*{Append-only continuation marker: V19}
\addcontentsline{toc}{section}{Append-only continuation marker: V19}
\textbf{End of V19.}  The complete V1--V18 body above is retained verbatim.
V19 assumes the V17 primary controlled zero-mode module class has been
trivialized and proves that the two downstream V18 obstructions disappear
automatically at every fixed finite coefficient/source order on the large
Renewal lattice.  The V15 finite translated-copy argument is strengthened to
avoid any finite list of extra buffer packets; this allows the local reserves
on a finite good cover to be chosen pairwise orthogonal wherever their square
partition-of-unity weights overlap.  Their weighted sum is therefore one exact
global controlled isometry from the initial zero-mode module into a common
internal massive reserve.  Composing with the fixed primary partial isometry
gives the same reserve for the target zero-mode module.  The V14 buffer
rotations then extend the chosen primary map itself to a global unitary in the
identity component, killing both the complement Cech class and the tertiary
component class without spectators or matrix amplification.  The remaining
complex microscopic obstruction is precisely the V17 primary controlled
family-index/module class.


\clearpage
\section{V20 continuation: controlled family index, failure of ordinary index lifting, and classification of the protected complex microscopic phase}
\label{sec:v20-continuation}

V19 reduced the global unstabilized microscopic problem to one primary datum:
the V17 controlled zero-mode module.  Its proposed next theorem asked whether
equality of the ordinary family Fredholm/index class, together with the common
physical elliptic Dirac--Yukawa principal symbol, forces equality of that
controlled module at sufficiently large cutoff.

The answer is negative in the unrestricted native local-operator category.
The ordinary family index forgets 
\emph{where} a finite-dimensional zero-mode packet lives in the regulator
lattice and how much spatial transport is required to identify two such
packets.  This information survives every determinant-line and ordinary
bundle forgetful map and is exactly what the exponentially weighted native
algebra records.

This is not a setback for the programme.  Combining V17 and V19 yields a
stronger and exact classification statement: on the protected minimal
fixed-index complex branch, the correct microscopic family-index object is the
\emph{unstable controlled zero-mode module class} itself.  Equality of this
class is necessary and sufficient for one global unstabilized coefficientwise
native regulator path; its Grothendieck image is the controlled $K_0$ shadow,
and its ordinary Atiyah--Singer family index is a still coarser shadow.

V20 proves:
\begin{enumerate}[label=\textup{(V20.\arabic*)},leftmargin=3.8em]
\item a controlled family-index package for every V11 protected minimal family,
with exact semigroup and group-completed versions and a natural forgetful map
to the ordinary family index;
\item invariance of the exact controlled family index under every global native
fixed-index homotopy;
\item a minimal index-$+1$ range-zero counterexample showing that equal
ordinary family index and equal principal propagating symbol do not imply
controlled equivalence;
\item a quantitative transport functional whose boundedness is equivalent to
controlled equivalence and which evaluates exactly to the macroscopic
zero-mode displacement in the counterexample;
\item a bounded-centre matching criterion which turns the V15 localized frames
into an explicit controlled partial isometry;
\item the final protected-complex classification theorem: at every prescribed
finite coefficient/source order and sufficiently large cutoff, two regulators
are globally unstabilized-equivalent if and only if their exact controlled
family-index modules agree (together with the already-declared complex line
data; a real Pfaffian orientation remains separate).
\end{enumerate}

The external Atiyah--Singer families theorem enters only as a comparison
principle for the ordinary continuum index: the topological family index is
determined by the elliptic principal-symbol class.  V20 does \emph{not}
promote that theorem to a statement about the exponentially local native
module.

\subsection{Exact and group-completed controlled family index}
\label{sec:v20-controlled-index}

Let $D_a(b)$ be a V11 protected chiral/Yukawa family on a compact background
component $B$, with cutoff label $a$, fixed index $k$, minimal nullities and a
uniform reduced singular-value floor.  By V13, its kernel and cokernel Riesz
projectors
\begin{equation}
 p_a(b):=\Pi_{\ker D_a(b)},
 \qquad
 c_a(b):=\Pi_{\ker D_a(b)^*}
 \label{eq:v20-pc}
\end{equation}
are uniformly exponentially local, together with every fixed source derivative
required by the prescribed coefficient packet.

Let $\mathcal V_{\rm ctl}(B)$ and $K_{0,\rm ctl}(B)$ be the V17 controlled
Murray--von Neumann semigroup and its Grothendieck group, understood for the
whole regulator sequence with cutoff-uniform locality constants.

\begin{definition}[Exact controlled family index]
\label{def:v20-exact-index}
On the V11 minimal stratum define the \emph{exact controlled family-index
module}
\begin{equation}
 \boxed{
 \operatorname{Ind}^{\rm ex}_{\rm ctl}(D)
 :=
 \begin{cases}
 [p]_{\mathcal V_{\rm ctl}},&k>0,\\[0.2em]
 [0],&k=0,\\[0.2em]
 [c]_{\mathcal V_{\rm ctl}},&k<0,
 \end{cases}}
 \label{eq:v20-exact-index}
\end{equation}
where for $k<0$ the sign is remembered separately by the fixed index label.
Its group-completed shadow is
\begin{equation}
 \boxed{
 \operatorname{Ind}_{\rm ctl}(D)
 :=[p]-[c]\in K_{0,\rm ctl}(B).}
 \label{eq:v20-group-index}
\end{equation}
\end{definition}

The exact class is deliberately semigroup-valued.  On the minimal stratum one
of $p,c$ is zero, so no information is gained by subtracting a second nonzero
packet; group completion can only forget unstable cancellation data already
identified in V17--V19.

\begin{proposition}[Forgetful hierarchy]
\label{prop:v20-forgetful-index}
There are natural forgetful maps
\begin{equation}
 \boxed{
 \operatorname{Ind}^{\rm ex}_{\rm ctl}(D)
 \longmapsto
 [K_D]-[C_D]\in K^0(B)
 \longmapsto
 \det C_D\otimes(\det K_D)^*.}
 \label{eq:v20-index-hierarchy}
\end{equation}
The group-completed controlled index
\eqref{eq:v20-group-index} maps to the ordinary family Fredholm index
\begin{equation}
 \operatorname{Ind}(D):=[K_D]-[C_D]\in K^0(B).
 \label{eq:v20-ordinary-index}
\end{equation}
On the faithful complex Standard-Model branch V12 trivializes the last line
factor, but neither of the two preceding classes is thereby forced to vanish.
\end{proposition}

\begin{proof}
A controlled partial isometry is in particular a continuous fibrewise partial
isometry after forgetting its spatial kernel bounds, hence induces an ordinary
vector-bundle isomorphism.  Passing to Grothendieck groups gives the first
forgetful map.  Taking top exterior powers gives the determinant line.  The
last statement is precisely the V12/V17 distinction between the full zero-mode
bundle/module and its determinant shadow.
\end{proof}

\begin{theorem}[Native homotopy invariance of the exact controlled family index]
\label{thm:v20-index-invariance}
Let $D_{a,t}(b)$, $0\le t\le1$, be one global V11-admissible native fixed-index
homotopy with a uniform reduced floor and the declared fixed source packet.
Then
\begin{equation}
 \boxed{
 \operatorname{Ind}^{\rm ex}_{\rm ctl}(D_{\cdot,0})
 =
 \operatorname{Ind}^{\rm ex}_{\rm ctl}(D_{\cdot,1})}
 \label{eq:v20-index-invariance}
\end{equation}
in the controlled semigroup, and hence
$\operatorname{Ind}_{\rm ctl}$ and the ordinary family index are invariant as
well.
\end{theorem}

\begin{proof}
V13 proves that the Riesz zero-mode projectors and their $t$ derivatives are
uniformly exponentially local.  Its local Kato equation therefore gives a
cutoff-uniform controlled unitary transport $U_{a,t}$.  For $k>0$, the endpoint
partial isometry $v_a=U_{a,1}p_{a,0}$ satisfies
$v_a^*v_a=p_{a,0}$ and $v_av_a^*=p_{a,1}$ with all prescribed source jets in
one positive locality class.  Thus the exact semigroup classes agree.  The
$c$-projector proof for $k<0$ is identical; $k=0$ is trivial on the minimal
stratum.  Group completion and forgetting control preserve equality.
\end{proof}

\subsection{What the ordinary families index theorem can and cannot certify}
\label{sec:v20-ordinary-family-index}

For an ordinary smooth family of elliptic continuum operators, the standard
Atiyah--Singer families theorem associates to the elliptic principal-symbol
class a $K^0(B)$ family index.  Consequently homotopic elliptic principal
symbols have the same ordinary family index.  In the Renewal programme this
statement may be used only after the relevant finite-to-ordinary comparison
has identified the native principal symbol with that continuum elliptic
family.

\begin{firewall}[Principal-symbol equality is an ordinary-index statement]
\label{fire:v20-principal-symbol}
Even after such a comparison, equality of the continuum principal-symbol class
certifies at most equality in $K^0(B)$.  It does not control the exponentially
weighted spatial kernel of a zero-mode identification, and therefore does not
imply equality of
$\operatorname{Ind}^{\rm ex}_{\rm ctl}$ or even of the group-completed
$\operatorname{Ind}_{\rm ctl}$.  The next theorem shows the failure already on
a minimal index-$+1$ branch with a unit reduced gap.
\end{firewall}

\subsection{A minimal fixed-index counterexample with the same propagating principal symbol}
\label{sec:v20-minimal-counterexample}

Let $\Lambda_L=\mathbb Z/L\mathbb Z$ with $L$ even and let
\begin{equation}
 \mathcal H_{+,L}=\ell^2(\Lambda_L),
 \qquad
 \mathcal H_{-,L}=\ell^2(\Lambda_L).
 \label{eq:v20-chiral-carrier}
\end{equation}
Write $P_+$ for the projection onto $\mathcal H_{+,L}$.  For a site
$x\in\Lambda_L$ put
\begin{equation}
 K_x:=\mathbb C\delta_x\subset\mathcal H_{-,L},
 \qquad
 W_x:=\mathcal H_{+,L}\oplus K_x,
 \label{eq:v20-Wx}
\end{equation}
and let $P_x$ be the orthogonal projection onto $W_x$ in
$\mathcal H_{+,L}\oplus\mathcal H_{-,L}$.  Define the range-zero Hermitian
Wilson sign
\begin{equation}
 \varepsilon_x:=2P_x-I
 \label{eq:v20-epsilon-x}
\end{equation}
and its overlap Weyl block, with $a=1$ for this model,
\begin{equation}
 D_x:=P_+(I+\gamma_5\varepsilon_x)|_{W_x}.
 \label{eq:v20-Dx}
\end{equation}

\begin{proposition}[Minimal index-one local defect family]
\label{prop:v20-defect-family}
For every $x$:
\begin{enumerate}[label=\textup{(D\arabic*)},leftmargin=3em]
\item $\varepsilon_x$ is Hermitian, range zero and has Wilson gap one;
\item
\begin{equation}
 \boxed{D_x=2P_+|_{W_x};}
 \label{eq:v20-Dx-projection}
\end{equation}
\item
\begin{equation}
 \ker D_x=K_x,
 \qquad
 \operatorname{coker}D_x=0,
 \qquad
 \operatorname{ind}D_x=1;
 \label{eq:v20-index-one}
\end{equation}
\item the reduced nonzero singular spectrum is the single value $2$;
\item the propagating block on $\mathcal H_{+,L}$ is independent of $x$.
Hence changing $x$ moves only a finite-rank zero-mode defect and does not
change the propagating principal symbol.
\end{enumerate}
\end{proposition}

\begin{proof}
Items (D1) and the range statement follow from
$P_x=P_++|\delta_x^-\rangle\langle\delta_x^-|$.  For $v\in W_x$ one has
$\varepsilon_xv=v$, so the V11 projection identity gives
$D_xv=2P_+v$, proving (D2).  Its kernel is exactly the extra negative-chirality
line $K_x$ and its image is all of $\mathcal H_{+,L}$, proving (D3).  On the
orthogonal complement of $K_x$ inside $W_x$, the map is $2I$, proving (D4).
The same formula shows that the entire propagating plus-chirality block is
independent of $x$; only the finite-rank kernel line changes.
\end{proof}

If desired, this defect block may be direct-summed with any common physical
elliptic Dirac--Yukawa block.  The common physical principal symbol is then
unchanged while the following controlled obstruction remains in the defect
summand.

\begin{theorem}[Equal ordinary family index does not lift to a controlled module]
\label{thm:v20-index-does-not-lift}
Choose
\begin{equation}
 x_L=0,
 \qquad
 y_L=L/2.
 \label{eq:v20-far-sites}
\end{equation}
The two protected minimal families $D_{x_L}$ and $D_{y_L}$ have:
\begin{enumerate}[label=\textup{(I\arabic*)},leftmargin=3em]
\item the same propagating principal symbol;
\item the same ordinary family index $+1$ over a point (and, after taking a
constant product with any compact parameter space $B$, the same trivial
index-line bundle over $B$);
\item the same trivial determinant line and the same reduced singular floor
$2$;
\item distinct controlled zero-mode module classes for the regulator sequence
$L\to\infty$.
\end{enumerate}
Indeed any partial isometry from $K_{x_L}$ to $K_{y_L}$ is
\begin{equation}
 v_L=e^{i\theta_L}
 |\delta_{y_L}\rangle\langle\delta_{x_L}|
 \label{eq:v20-vL}
\end{equation}
and for every fixed $\beta>0$,
\begin{equation}
 \boxed{\|v_L\|_\beta=e^{\beta L/2}\longrightarrow\infty.}
 \label{eq:v20-vL-cost}
\end{equation}
Thus equality of the ordinary family index and of the physical principal
symbol does not imply
$\operatorname{Ind}^{\rm ex}_{\rm ctl}(D_{x_L})
 =\operatorname{Ind}^{\rm ex}_{\rm ctl}(D_{y_L})$.
\end{theorem}

\begin{proof}
The first three items are Proposition~\ref{prop:v20-defect-family}.  Since both
kernel lines are one dimensional, every partial isometry between them has the
form \eqref{eq:v20-vL}.  Its unique nonzero kernel entry joins sites at periodic
distance $L/2$, so the weighted row and column norm is exactly
$e^{\beta L/2}$.  Hence no cutoff-uniform positive exponential locality class
contains these partial isometries, which is precisely failure of V13
controlled equivalence.
\end{proof}

\begin{firewall}[Scope of the counterexample]
\label{fire:v20-counterexample-scope}
Theorem~\ref{thm:v20-index-does-not-lift} proves that the hypotheses
``same principal elliptic symbol + same ordinary family index'' are logically
insufficient in the unrestricted local-operator category.  It does not claim
that the actual Renewal Standard-Model gauge/background grammar realizes such
a freely movable defect.  Proving that the physical native class excludes this
transport sector is now the precise additional theorem required for a stronger
universality statement.
\end{firewall}

\subsection{A quantitative controlled transport profile}
\label{sec:v20-transport-profile}

The preceding obstruction can be measured without choosing a reference frame.

\begin{definition}[Controlled transport cost]
\label{def:v20-transport-cost}
For two finite-rank projections $p,q$ at one cutoff and exponent $\beta>0$ put
\begin{equation}
 \tau_\beta(p,q)
 :=\inf_{v^*v=p,\ vv^*=q}
 {1\over\beta}\log\max\{1,\|v\|_\beta\},
 \label{eq:v20-tau}
\end{equation}
with value $+\infty$ if no such partial isometry exists.  For a regulator
sequence define
\begin{equation}
 \boxed{
 \mathfrak T_\beta(p,q)
 :=\sup_a\tau_\beta(p_a,q_a).}
 \label{eq:v20-T}
\end{equation}
The supremum includes the compact background panel and the finite source-jet
norm when those variables are present.
\end{definition}

\begin{proposition}[Transport-cost calculus]
\label{prop:v20-transport-cost}
Whenever the products are defined,
\begin{align}
 \tau_\beta(p,q)&=\tau_\beta(q,p),
 \label{eq:v20-tau-sym}\\
 \tau_\beta(p,r)&\le
 \tau_\beta(p,q)+\tau_\beta(q,r).
 \label{eq:v20-tau-triangle}
\end{align}
Moreover
\begin{equation}
 \boxed{
 p\sim_{\rm ctl}q
 \quad\Longleftrightarrow\quad
 \mathfrak T_\beta(p,q)<\infty
 \text{ for some }\beta>0.}
 \label{eq:v20-T-equivalence}
\end{equation}
For the defect projections of Theorem~\ref{thm:v20-index-does-not-lift},
\begin{equation}
 \boxed{
 \tau_\beta(p_{x_L},p_{y_L})=L/2.}
 \label{eq:v20-tau-site}
\end{equation}
\end{proposition}

\begin{proof}
If $v:p\to q$ is a partial isometry, then $v^*:q\to p$ and
$\|v^*\|_\beta=\|v\|_\beta$, proving symmetry.  If
$v:p\to q$ and $w:q\to r$, then $wv:p\to r$ and the V13 weighted Banach
algebra estimate gives
$\|wv\|_\beta\le\|w\|_\beta\|v\|_\beta$, proving the triangle inequality
after logarithms and infima.  The equivalence
\eqref{eq:v20-T-equivalence} is exactly the definition of a cutoff-uniform
controlled partial isometry.  Equation \eqref{eq:v20-tau-site} follows from
the unique rank-one partial isometry and
\eqref{eq:v20-vL-cost}.
\end{proof}

Thus $\mathfrak T_\beta$ is a computable coarse/locality refinement invisible
to the ordinary family index.

\subsection{A positive bounded-transport criterion from localized frames}
\label{sec:v20-positive-criterion}

V15 proved that every fixed-rank protected zero-mode projector admits a
uniformly exponentially localized orthonormal frame on a protected local
chart.  When the two frames can be matched within bounded lattice distance,
controlled equivalence follows immediately.

\begin{theorem}[Bounded-bottleneck frame criterion]
\label{thm:v20-bounded-bottleneck}
Let $p_a,q_a$ be rank-$\kappa$ protected zero-mode projectors on one framed
compact source chart.  Suppose there are orthonormal localized frames
\begin{equation}
 \psi_{1,a},\ldots,\psi_{\kappa,a},
 \qquad
 \phi_{1,a},\ldots,\phi_{\kappa,a}
 \label{eq:v20-localized-frames}
\end{equation}
with the V15 exponential bounds around centres $x_{j,a}$ and $y_{j,a}$, and a
permutation $\sigma_a$ such that
\begin{equation}
 \boxed{
 \sup_{a,j}d_a(x_{j,a},y_{\sigma_a(j),a})\le R<\infty.}
 \label{eq:v20-bottleneck-bound}
\end{equation}
Assume the same matching is source-smooth on the fixed packet.  Then
\begin{equation}
 v_a:=\sum_{j=1}^{\kappa}
 |\phi_{\sigma_a(j),a}\rangle\langle\psi_{j,a}|
 \label{eq:v20-frame-v}
\end{equation}
is a controlled partial isometry from $p_a$ to $q_a$.  For every
$0<\beta'<\beta$ there is
\begin{equation}
 \boxed{
 \sup_a\|v_a\|_{\beta'}
 \le C_{\kappa,\beta'}e^{\beta'R}.}
 \label{eq:v20-frame-v-bound}
\end{equation}
Hence the exact controlled family indices agree.
\end{theorem}

\begin{proof}
Orthonormality gives directly
$v_a^*v_a=p_a$ and $v_av_a^*=q_a$.  Each rank-one kernel in
\eqref{eq:v20-frame-v} is a product of two exponentially localized vectors
whose centres are at distance at most $R$.  The same convolution estimate used
in V15's translated-frame lemma gives a weighted row/column bound
$Ce^{\beta'R}$ after reducing the exponent from $\beta$ to $\beta'$.  Summing
only $\kappa$ terms preserves a cutoff-uniform constant.  Fixed source
 derivatives obey the finite Leibniz rule with the V15 frame-jet bounds.
\end{proof}

\begin{corollary}[What an anchoring theorem would need to prove]
\label{cor:v20-anchoring-target}
For fixed-rank protected zero modes, a sufficient native theorem upgrading
ordinary family-index equality to controlled equivalence is a uniform
bounded-bottleneck matching of localized zero-mode frames.  In particular, one
does not need equality of the frames themselves; one needs only a cutoff-
uniform bound on their spatial transport.
\end{corollary}

\subsection{Complete protected-complex microscopic classification at fixed order}
\label{sec:v20-classification}

V17 proved necessity of the exact controlled zero-mode module class.  V19
proved that, once this class is matched, the complement and identity-component
obstructions disappear automatically at every fixed finite source order on the
large Renewal lattice.  Together they yield the classification theorem which
V20 was seeking.

\begin{theorem}[Controlled family-index classification theorem]
\label{thm:v20-classification}
Fix a faithful complex Standard-Model representation, a protected minimal
fixed-index background component, positive IR/reference data, and one finite
coefficient packet
\begin{equation}
 \mathfrak q=(L,\Delta,\varrho,N,\mathcal P_{\rm ext}).
 \label{eq:v20-q}
\end{equation}
Let $A$ and $B'$ be two native regulators with the same physical principal
Dirac--Yukawa jet and the same declared global complex line data.  For all
sufficiently large cutoffs, the following are equivalent:
\begin{enumerate}[label=\textup{(C\arabic*)},leftmargin=3em]
\item $A$ and $B'$ are joined coefficientwise through order $\mathfrak q$ by
one global unstabilized V11-admissible native fixed-index regulator path;
\item their V17 primary controlled zero-mode class vanishes;
\item their exact controlled family-index modules agree,
\begin{equation}
 \boxed{
 \operatorname{Ind}^{\rm ex}_{\rm ctl}(D_A)
 =
 \operatorname{Ind}^{\rm ex}_{\rm ctl}(D_{B'}).}
 \label{eq:v20-class-equality}
\end{equation}
\end{enumerate}
For $k>0$ this is exact controlled equivalence of the kernel packets; for
$k<0$ it is exact controlled equivalence of the cokernel packets; for $k=0$
the statement reduces to the V10 balanced no-doubler component.
\end{theorem}

\begin{proof}
\textup{(C1)} implies \textup{(C3)} by
Theorem~\ref{thm:v20-index-invariance}.  Exact semigroup equality in
\textup{(C3)} is, by definition, a global controlled Murray--von Neumann
partial isometry between the nonzero minimal zero-mode packets, so it is the
global trivialization required in V17; hence \textup{(C3)} implies
\textup{(C2)}.  Finally V19 proves that after the V17 primary class is zero,
the global complement Cech class and tertiary identity-component obstruction
vanish coefficientwise at sufficiently large cutoff, producing one global
unstabilized native path.  Thus \textup{(C2)} implies \textup{(C1)}.
The $k=0$ case has no zero-mode packet and is exactly the V10 theorem.
\end{proof}

\begin{corollary}[The ordinary family index is not the microscopic classifier]
\label{cor:v20-ordinary-not-classifier}
The forgetful map
\begin{equation}
 \operatorname{Ind}^{\rm ex}_{\rm ctl}
 \longmapsto
 \operatorname{Ind}\in K^0(B)
 \label{eq:v20-forgetful-map}
\end{equation}
is not injective in the unrestricted protected local-operator category.  Hence
ordinary family-index equality, even together with equality of the propagating
principal symbol and determinant line, is not sufficient for one global
microscopic regulator path.
\end{corollary}

\begin{proof}
Theorem~\ref{thm:v20-index-does-not-lift} gives two minimal index-one families
with equal ordinary index and equal principal propagating symbol but different
exact controlled classes.  Apply
Theorem~\ref{thm:v20-classification}.
\end{proof}

\begin{corollary}[Continuum universality remains coarser than microscopic classification]
\label{cor:v20-continuum-coarser}
Even when the exact controlled family indices differ, the V16 scheme-groupoid
continuum equivalence can still hold from local V15 comparisons.  Thus
\begin{equation}
 \boxed{
 \text{same coefficientwise continuum}
 \not\Rightarrow
 \text{same controlled family-index class}.}
 \label{eq:v20-cont-vs-micro}
\end{equation}
The controlled family index is a strictly post-continuum microscopic regulator
invariant.
\end{corollary}

\subsection{V20 anti-circularity audit}
\label{sec:v20-audit}

\begin{enumerate}[leftmargin=2.4em]
\item \textbf{The ordinary Atiyah--Singer family index is not promoted to a
native locality theorem.}  Its principal-symbol invariance is used only for the
ordinary $K^0$ shadow.
\item \textbf{The counterexample lies on the minimal nonzero-index stratum.}
Unlike the earlier V13 square-operator illustration, there is no accidental
kernel--cokernel pair: the cokernel is zero and the reduced singular floor is
uniform.
\item \textbf{The principal propagating symbol is held fixed.}  The two
endpoints differ only by the location of one finite-rank kernel defect; a
common physical Dirac--Yukawa block may be direct-summed without changing the
obstruction.
\item \textbf{The new controlled index is not invented downstream of the
answer.}  It is exactly the V17 semigroup class already proved necessary under
native Kato transport, now repackaged as the correct family-index object.
\item \textbf{Group completion is not mistaken for exact equivalence.}  The
classification theorem uses
$\operatorname{Ind}^{\rm ex}_{\rm ctl}$ in the Murray--von Neumann semigroup,
not merely its $K_{0,\rm ctl}$ image.
\item \textbf{A positive criterion is quantitative.}  Bounded matching distance
of V15 localized frame centres gives the explicit partial isometry
\eqref{eq:v20-frame-v} with an exponential norm bound.
\item \textbf{V19 is used at exactly its proved scope.}  Once the exact primary
controlled class is matched, only finitely many source jets are needed and the
large-cutoff internal reserve kills the downstream unstable obstructions.
\end{enumerate}

\section{V20 status ledger}
\label{sec:v20-status}

\subsection*{Proved in V20}
\begin{enumerate}[leftmargin=2.2em]
\item Every protected minimal fixed-index family has an exact controlled
family-index module in the V17 Murray--von Neumann semigroup and a
group-completed controlled index in $K_{0,\rm ctl}$.  Forgetting spatial control
recovers the ordinary family Fredholm index, and taking top exterior powers
recovers the determinant-line shadow.
\item The exact controlled family index is invariant under every global
V11-admissible native fixed-index homotopy by V13 local Kato transport.
\item Equality of the ordinary family index and equality of the propagating
principal symbol do not imply equality of the controlled family index.  A
range-zero minimal index-$+1$ defect family with zero modes at sites separated
by $L/2$ gives an explicit counterexample with reduced singular floor $2$.
\item The quantitative transport functional $\mathfrak T_\beta$ is finite
exactly when a cutoff-uniform controlled partial isometry exists.  In the
far-site defect example its value grows exactly like the zero-mode separation.
\item A bounded-bottleneck matching of the V15 localized zero-mode frames is a
sufficient condition for controlled equivalence, with an explicit locality
bound $Ce^{\beta R}$.
\item Combining V17 necessity with V19 sufficiency yields a complete
coefficientwise microscopic classification on the faithful complex protected
minimal branch: equality of the exact controlled family-index module is
necessary and sufficient for one global unstabilized native fixed-index path
at sufficiently large cutoff.
\item Consequently the ordinary family index is only a coarse shadow and the
controlled family index is a strictly post-continuum regulator invariant; the
V16 coefficientwise continuum can agree even when this microscopic class does
not.
\end{enumerate}

\subsection*{Inherited}
\begin{enumerate}[leftmargin=2.2em]
\item V1--V7: source-complete coefficientwise ESM continuum, common BV
restoration, fixed-loop forests, ultraviolet Cauchy convergence,
scheme/reference covariance and exact perturbative-order projectivity.
\item V8--V12: regulator-marked comparison, fixed-index reduced shell calculus,
global zero-mode bundle/line transport and faithful complex saturation.
\item V13--V16: exponential locality of protected zero modes, stable and local
unstabilized lifting, automatic finite-jet local reserves and global continuum
descent in the finite scheme groupoid.
\item V17: the exact primary controlled zero-mode Cech/module obstruction and
its controlled $K_0$ shadow.
\item V18--V19: the secondary complement and tertiary unitary-component tower,
and its automatic coefficientwise collapse after the primary class is
trivialized.
\end{enumerate}

\subsection*{Conditional}
\begin{enumerate}[leftmargin=2.2em]
\item The ordinary principal-symbol/family-index implication is imported only
at the continuum elliptic-family level, or after a separately proved native-to-
ordinary symbol comparison.  V20 does not infer a controlled index from the
principal symbol alone.
\item The exact microscopic classification is coefficientwise at fixed finite
$(L,\Delta,\varrho,N)$, on one protected minimal fixed-index component and for
sufficiently large cutoff.
\item The far-transport counterexample shows logical insufficiency of the
principal symbol/index data in the unrestricted local category; whether the
physical Renewal gauge/background grammar excludes such sectors is a separate
native anchoring theorem.
\item The faithful complex determinant-line assumptions of V12 remain in force;
an independently real Pfaffian sector still carries its mod-two orientation
condition.
\item All continuum statements retain the positive IR/reference scale and make
no claim of perturbative-series convergence or nonperturbative completion.
\end{enumerate}

\subsection*{Open beyond V20}
\begin{enumerate}[leftmargin=2.2em]
\item Prove or falsify a \emph{native zero-mode anchoring theorem} for the actual
Renewal ESM regulator class: show that two admissible microscopic completions
with the same physical elliptic Dirac--Yukawa symbol and the same ordinary
family index have V15 localized zero-mode frames with uniformly bounded
transport matching, or construct a genuine physical counterexample.
\item Equivalently, relate the exact controlled family index to a computable
local index-density/coarse-transport invariant of the gauge/background field,
not merely to the abstract $K^0(B)$ family index.
\item Compute the independent V12 real Pfaffian mod-two spectral-flow character
on any physical Majorana sector.
\item Analyze index-changing/reduced-gap-closing divisors and the independent
positive-IR, analytic-RG, summability and nonperturbative directions retained
by V19.
\end{enumerate}

\subsection*{Exact next load-bearing theorem}
\begin{center}
\fbox{\parbox{0.92\textwidth}{
\textbf{Native zero-mode anchoring / controlled index-density theorem.}
For two protected Renewal ESM regulator completions with the same
Standard-Model representation, the same physical elliptic Dirac--Yukawa
principal symbol and the same ordinary family index, analyze the native
zero-mode Riesz projectors against the common gauge/coframe/Higgs background.
Prove that their V15 localized frames admit a cutoff-uniform bounded-transport
matching (hence equal exact controlled family index by
Theorem~\ref{thm:v20-bounded-bottleneck}), or construct a genuine Renewal-native
coarse transport/index-density invariant which distinguishes them.  A positive
anchoring theorem, combined with
Theorem~\ref{thm:v20-classification}, would eliminate the last complex
microscopic regulator-class datum beyond the ordinary family index and leave
only the independent real Pfaffian orientation/zero-crossing problems.}}
\end{center}

\section*{Append-only continuation marker: V20}
\addcontentsline{toc}{section}{Append-only continuation marker: V20}
\textbf{End of V20.}  The complete V1--V19 body above is retained verbatim.
V20 answers the V19 family-index question negatively at the level of ordinary
$K^0$: equal principal propagating symbol and equal ordinary family index do
not force a controlled zero-mode identification.  A minimal index-$+1$
range-zero defect family places the unique kernel line at macroscopically
separated lattice sites, leaving the ordinary index and determinant data
unchanged while every identifying partial isometry has exponentially growing
locality norm.  The correct microscopic family-index object is therefore the
exact V17 controlled Murray--von Neumann module class, with
$K_{0,\rm ctl}$ and the ordinary family index as successive forgetful shadows.
Combining V17 necessity with V19 sufficiency gives a complete protected-complex
coefficientwise microscopic classification: two regulators lie in one global
unstabilized native fixed-index path component exactly when their exact
controlled family-index modules agree.  The remaining stronger gate is a
native anchoring theorem relating this controlled class to the actual common
Renewal gauge/coframe/Higgs background and ordinary principal-symbol index.


\clearpage
\section{V21 continuation: anchored zero-mode atlases, controlled index density, and the exact native anchoring gap}
\label{sec:v21-continuation}

V20 proved that the ordinary family Fredholm index and the physical elliptic
principal symbol do not determine the exact controlled zero-mode module.  The
remaining question is therefore not whether the ordinary index has been
computed correctly, but what additional \emph{spatially controlled} datum the
native regulator must supply in order to lift that index class to the V20
microscopic classifier.

The correct variable is already implicit in V15.  A fixed-rank exponentially
local zero-mode projector admits local orthonormal frames with cutoff-uniform
exponential moments about finitely many lattice anchors.  A controlled partial
isometry preserves precisely this form of localization.  Conversely, two such
local frames with uniformly bounded matched anchors produce the required
controlled partial isometry by the V20 rank-one synthesis formula.  Globally,
one must in addition match the ordinary zero-mode bundle transition cocycle.
This gives an exact geometric characterization of the V17/V20 controlled
module class.

V21 then extracts a regulator-independent spatial shadow from that exact class.
For a zero-mode projector $p$, define its diagonal index density
\begin{equation}
 \rho_p(x):=\operatorname{tr}_{\rm fib}p_{xx}\ge0,
 \qquad
 \sum_x\rho_p(x)=\operatorname{rank}p.
 \label{eq:v21-density}
\end{equation}
Controlled equivalence forces the endpoint densities to be related by a
cutoff-uniform exponentially localized transport plan.  This produces a
\emph{controlled index-density class}, strictly finer than the total ordinary
index but strictly coarser than the full controlled zero-mode module.  The
V20 far-site defect is detected exactly by this density transport cost, while
an onsite unstable bundle example shows that the density still forgets internal
bundle twisting.

Thus the V20 anchoring target can be restated sharply: the actual Renewal ESM
background grammar must either force a regulator-independent bounded anchoring
of the zero-mode packet, or it must admit a genuine native controlled
index-density/transport distinction.  Neither conclusion follows from the
present principal-symbol hypotheses alone.

Throughout V21, $\kappa=|k|$ is fixed and finite.  For $k>0$ the zero-mode
projector is the kernel projector $p=\Pi_K$; for $k<0$ it is the cokernel
projector $p=\Pi_C$.  The case $k=0$ has no primary zero-mode module and remains
covered by V10.

\subsection{Anchored local frames}
\label{sec:v21-anchored-frames}

\begin{definition}[Exponentially anchored frame]
\label{def:v21-anchored-frame}
Let $p$ be a rank-$\kappa$ projection on one protected source/background
chart.  A $(\beta,C)$-\emph{anchored frame} consists of an orthonormal frame
$\psi_1,\ldots,\psi_\kappa$ of $\operatorname{Ran}p$ and lattice sites
$x_1,\ldots,x_\kappa$ such that
\begin{equation}
 \boxed{
 \sum_z e^{\beta d(z,x_j)}\|\psi_j(z)\|\le C
 \qquad(1\le j\le\kappa).}
 \label{eq:v21-anchored-frame}
\end{equation}
For a fixed source packet, all prescribed frame derivatives are required to
satisfy the analogous estimate after one common finite loss of exponent.
\end{definition}

V15 proves existence of such frames for every protected fixed-rank Riesz
projector.  Their anchors are not unique; only their bounded coarse location is
load bearing.

\begin{definition}[Bounded matched anchored frames]
\label{def:v21-bounded-anchor-match}
Two rank-$\kappa$ projections $p,q$ on one chart admit a
\emph{$(\beta,C,R)$ bounded anchor matching} if they possess anchored frames
$(\psi_j,x_j)$ and $(\phi_j,y_j)$ and a permutation $\sigma$ such that
\begin{equation}
 d(x_j,y_{\sigma(j)})\le R
 \qquad(1\le j\le\kappa),
 \label{eq:v21-anchor-bottleneck}
\end{equation}
with constants uniform in cutoff and over the compact chart.
\end{definition}

\begin{lemma}[A controlled operator preserves an anchor]
\label{lem:v21-operator-preserves-anchor}
Let $A$ be a kernel with $\|A\|_\alpha\le M$ and let $\psi$ satisfy
\eqref{eq:v21-anchored-frame} at exponent $\beta<\alpha$.  Then
\begin{equation}
 \boxed{
 \sum_y e^{\beta d(y,x_0)}\|(A\psi)(y)\|
 \le \|A\|_\beta
      \sum_x e^{\beta d(x,x_0)}\|\psi(x)\|.}
 \label{eq:v21-A-anchor}
\end{equation}
The same estimate holds for every fixed source derivative by the finite
Leibniz rule.
\end{lemma}

\begin{proof}
Use
$d(y,x_0)\le d(y,x)+d(x,x_0)$:
\[
 \begin{aligned}
 \sum_y e^{\beta d(y,x_0)}\|(A\psi)(y)\|
 &\le
 \sum_{x,y}e^{\beta d(y,x)}\|A_{yx}\|
            e^{\beta d(x,x_0)}\|\psi(x)\|\\
 &\le
 \|A\|_\beta
 \sum_xe^{\beta d(x,x_0)}\|\psi(x)\|.
 \end{aligned}
\]
Source derivatives are finite sums of products of differentiated factors and
obey the same weighted-algebra estimate.
\end{proof}

\begin{theorem}[Local anchored-frame characterization of controlled equivalence]
\label{thm:v21-local-anchor-characterization}
Let $p,q$ be protected rank-$\kappa$ zero-mode projectors on one star-shaped
source/background chart.  The following are equivalent.
\begin{enumerate}[label=\textup{(A\arabic*)},leftmargin=3em]
\item There is a controlled partial isometry $v$ with
\begin{equation}
 v^*v=p,
 \qquad
 vv^*=q.
 \label{eq:v21-local-v}
\end{equation}
\item There are cutoff-uniform bounded matched anchored frames for $p$ and $q$.
\end{enumerate}
Moreover, under \textup{(A1)} the two frames can be chosen with the
\emph{same} anchor sites, so one may take $R=0$ after changing the target
frame.
\end{theorem}

\begin{proof}
Assume \textup{(A1)}.  By V15 choose a uniformly anchored frame
$(\psi_j,x_j)$ for $p$ and put
\begin{equation}
 \phi_j:=v\psi_j.
 \label{eq:v21-v-frame}
\end{equation}
Since $v^*v=p$ and $vv^*=q$, the vectors $\phi_j$ form an orthonormal basis of
$\operatorname{Ran}q$.  Lemma~\ref{lem:v21-operator-preserves-anchor} shows
that every $\phi_j$ is exponentially localized about the \emph{same} site
$x_j$, uniformly in the cutoff and with the prescribed source derivatives.
Thus \textup{(A2)} holds with $R=0$.

Conversely suppose \textup{(A2)} holds.  Relabel the target frame using the
permutation and define
\begin{equation}
 v:=\sum_{j=1}^{\kappa}
 |\phi_j\rangle\langle\psi_j|.
 \label{eq:v21-frame-partial-isometry}
\end{equation}
Orthonormality gives $v^*v=p$ and $vv^*=q$.  The rank-one kernel
$\phi_j(y)\psi_j(x)^*$ satisfies
\[
 d(x,y)\le d(x,x_j)+R+d(y,y_j),
\]
so the same convolution estimate as in V20 gives, for every
$0<\beta'<\beta$,
\begin{equation}
 \|v\|_{\beta'}
 \le C_{\kappa,\beta'}e^{\beta'R}.
 \label{eq:v21-v-local-bound}
\end{equation}
Fixed source derivatives are finite Leibniz sums of the anchored frame jets.
Hence $v$ is controlled.
\end{proof}

\subsection{Global anchored atlases and the exact controlled module}
\label{sec:v21-anchored-atlas}

A nontrivial zero-mode bundle need not admit one global orthonormal frame.  The
correct global version of Theorem~\ref{thm:v21-local-anchor-characterization}
is therefore an atlas statement.

\begin{definition}[Matched controlled anchor atlas]
\label{def:v21-anchor-atlas}
Let $(U_\alpha)$ be a finite good cover of the compact protected background
component $B$.  A \emph{matched controlled anchor atlas} for $p,q$ consists of
local anchored frames
\begin{equation}
 \Psi_\alpha=(\psi_{\alpha,1},\ldots,\psi_{\alpha,\kappa}),
 \qquad
 \Phi_\alpha=(\phi_{\alpha,1},\ldots,\phi_{\alpha,\kappa})
 \label{eq:v21-atlas-frames}
\end{equation}
with cutoff-uniform locality constants and bounded matched anchors, such that
after local unitary re-framing the transition functions agree exactly:
\begin{equation}
 \boxed{
 \Psi_\beta=\Psi_\alpha g_{\alpha\beta},
 \qquad
 \Phi_\beta=\Phi_\alpha g_{\alpha\beta}}
 \quad\text{on }U_{\alpha\beta}.
 \label{eq:v21-same-transition}
\end{equation}
Here $g_{\alpha\beta}:U_{\alpha\beta}\to U(\kappa)$ carries the full retained
source jet.
\end{definition}

\begin{theorem}[Global anchored-atlas characterization]
\label{thm:v21-global-anchor-characterization}
For protected rank-$\kappa$ zero-mode projector families $p,q$ on $B$, the
following are equivalent.
\begin{enumerate}[label=\textup{(G\arabic*)},leftmargin=3em]
\item Their exact V20 controlled family-index modules agree; equivalently,
there is one global controlled partial isometry
$v^*v=p$, $vv^*=q$.
\item They admit a matched controlled anchor atlas in the sense of
Definition~\ref{def:v21-anchor-atlas}.
\end{enumerate}
Thus the exact controlled family-index class is precisely the combination of
ordinary zero-mode bundle descent and bounded spatial anchoring.
\end{theorem}

\begin{proof}
Assume \textup{(G1)}.  On each chart choose the V15 anchored Kato frame
$\Psi_\alpha$ of $p$ and put
$\Phi_\alpha=v\Psi_\alpha$.  Lemma
\ref{lem:v21-operator-preserves-anchor} gives the same anchors for the target
frame.  If
$\Psi_\beta=\Psi_\alpha g_{\alpha\beta}$, then
\[
 \Phi_\beta
 =v\Psi_\beta
 =v\Psi_\alpha g_{\alpha\beta}
 =\Phi_\alpha g_{\alpha\beta},
\]
so \eqref{eq:v21-same-transition} holds.

Conversely suppose \textup{(G2)} holds and set locally
\begin{equation}
 v_\alpha:=\Phi_\alpha\Psi_\alpha^*.
 \label{eq:v21-local-atlas-v}
\end{equation}
By Theorem~\ref{thm:v21-local-anchor-characterization}, every $v_\alpha$ is
uniformly controlled.  On an overlap,
\[
 v_\beta
 =\Phi_\alpha g_{\alpha\beta}
   g_{\alpha\beta}^*\Psi_\alpha^*
 =v_\alpha.
\]
Hence the $v_\alpha$ glue to one global controlled operator $v$, and the local
partial-isometry identities glue to $v^*v=p$, $vv^*=q$.
\end{proof}

\subsection{The diagonal controlled index density}
\label{sec:v21-index-density}

The anchored-atlas theorem is exact but still contains frame and bundle data.
There is a simpler spatial shadow which can be read directly from the Riesz
projector.

\begin{definition}[Zero-mode index density]
\label{def:v21-index-density}
For a protected zero-mode projection $p$, define the positive lattice measure
\begin{equation}
 \mu_p(x):=\rho_p(x)
 :=\operatorname{tr}_{\rm fib}(p_{xx}).
 \label{eq:v21-mu}
\end{equation}
Its total mass is
\begin{equation}
 \boxed{\sum_x\mu_p(x)=\operatorname{Tr}p=\kappa.}
 \label{eq:v21-density-total}
\end{equation}
For $k>0$ this is the kernel density and for $k<0$ the cokernel density; the
sign of the Fredholm index is stored separately by $k$.
\end{definition}

\begin{definition}[Exponential transport cost]
\label{def:v21-transport-cost}
For two positive lattice measures $\mu,\nu$ of equal finite mass and
$\gamma>0$, let $\Pi(\mu,\nu)$ be the set of nonnegative couplings with the
given marginals and define
\begin{equation}
 \boxed{
 \mathcal W_\gamma(\mu,\nu)
 :=\inf_{\pi\in\Pi(\mu,\nu)}
 \sum_{x,y}e^{\gamma d(x,y)}\pi(x,y).}
 \label{eq:v21-Wgamma}
\end{equation}
For regulator sequences we call the two density families
\emph{exponentially transport equivalent} if for some $\gamma>0$ the quantity
\begin{equation}
 \sup_a\sup_{b\in B}
 \mathcal W_\gamma(\mu_{p_a(b)},\mu_{q_a(b)})
 \label{eq:v21-uniform-W}
\end{equation}
is finite.
\end{definition}

\begin{lemma}[Frame coupling]
\label{lem:v21-frame-coupling}
Let $(\psi_j,x_j)$ and $(\phi_j,y_j)$ be bounded matched anchored frames with
$d(x_j,y_j)\le R$.  Then
\begin{equation}
 \pi(x,y):=
 \sum_{j=1}^{\kappa}
 \|\psi_j(x)\|^2\|\phi_j(y)\|^2
 \label{eq:v21-frame-coupling}
\end{equation}
is a coupling of $\mu_p$ and $\mu_q$.  For every
$0<\gamma\le2\beta$,
\begin{equation}
 \boxed{
 \sum_{x,y}e^{\gamma d(x,y)}\pi(x,y)
 \le \kappa C^4 e^{\gamma R}.}
 \label{eq:v21-frame-coupling-bound}
\end{equation}
\end{lemma}

\begin{proof}
Since each frame is orthonormal,
\[
 \sum_y\pi(x,y)
 =\sum_j\|\psi_j(x)\|^2
 =\operatorname{tr}p_{xx}=\mu_p(x),
\]
and similarly the second marginal is $\mu_q$.
For each matched pair,
\[
 d(x,y)\le d(x,x_j)+R+d(y,y_j).
\]
Moreover, for $\gamma\le2\beta$,
\[
 \sum_xe^{\gamma d(x,x_j)}\|\psi_j(x)\|^2
 \le
 \left(
 \sum_xe^{\beta d(x,x_j)}\|\psi_j(x)\|
 \right)^2
 \le C^2,
\]
and the same estimate holds for $\phi_j$.  Factor the double sum and sum over
$j$ to obtain \eqref{eq:v21-frame-coupling-bound}.
\end{proof}

\begin{theorem}[Controlled module implies controlled index-density transport]
\label{thm:v21-module-implies-density}
If two protected zero-mode packets have equal exact controlled family-index
modules, then their density families are exponentially transport equivalent.
More quantitatively, any bounded anchor matching with parameters
$(\beta,C,R)$ gives
\begin{equation}
 \boxed{
 \mathcal W_\gamma(\mu_p,\mu_q)
 \le \kappa C^4e^{\gamma R}}
 \qquad(0<\gamma\le2\beta).
 \label{eq:v21-density-bound}
\end{equation}
Thus every global native fixed-index regulator homotopy preserves the
exponential transport class of the zero-mode density.
\end{theorem}

\begin{proof}
Use Theorem~\ref{thm:v21-global-anchor-characterization} on a finite good cover.
The constants are uniform on the compact cover.  At each background point pick
one chart containing it and apply Lemma~\ref{lem:v21-frame-coupling}; the
infimum in \eqref{eq:v21-Wgamma} is no larger than this explicit coupling.
Native-path invariance follows from V13 controlled Kato transport and V20
controlled-index invariance.
\end{proof}

\begin{definition}[Controlled index-density class]
\label{def:v21-density-class}
The \emph{controlled index-density class} of a minimal fixed-index family is
\begin{equation}
 \operatorname{Ind}^{\rm dens}_{\rm ctl}(D)
 :=[\mu_{\Pi_K}]_{\exp}
 \quad(k>0),
 \qquad
 [\mu_{\Pi_C}]_{\exp}
 \quad(k<0),
 \label{eq:v21-density-class}
\end{equation}
where $[\cdot]_{\exp}$ denotes equivalence under a cutoff-uniform exponential
transport bound \eqref{eq:v21-uniform-W}.  Its total mass remembers only
$|k|$.
\end{definition}

There are now canonical forgetful arrows
\begin{equation}
 \boxed{
 \operatorname{Ind}^{\rm ex}_{\rm ctl}
 \longrightarrow
 \operatorname{Ind}^{\rm dens}_{\rm ctl}
 \longrightarrow
 |\operatorname{Ind}(D)|.}
 \label{eq:v21-density-hierarchy}
\end{equation}
The first arrow forgets internal frame/bundle phases while retaining spatial
transport; the second forgets all spatial information and keeps only total
zero-mode number.

\subsection{The V20 far-site defect is exactly a density-transport defect}
\label{sec:v21-far-site-density}

For the minimal index-one family of V20,
\begin{equation}
 \mu_{p_{x_L}}=\delta_{x_L},
 \qquad
 \mu_{p_{y_L}}=\delta_{y_L}.
 \label{eq:v21-delta-densities}
\end{equation}
There is only one coupling between two unit point masses.

\begin{proposition}[Exact exponential transport cost of the far-site defect]
\label{prop:v21-far-site-density}
For the V20 endpoints,
\begin{equation}
 \boxed{
 \mathcal W_\gamma(
 \delta_{x_L},\delta_{y_L})
 =e^{\gamma d(x_L,y_L)}.}
 \label{eq:v21-far-site-W}
\end{equation}
In particular, if $d(x_L,y_L)=L/2$, the controlled index-density classes are
different for every fixed $\gamma>0$ as $L\to\infty$.
\end{proposition}

\begin{proof}
The unique coupling is
$\pi(x,y)=\mathbf1_{x=x_L}\mathbf1_{y=y_L}$.  Substitution into
\eqref{eq:v21-Wgamma} gives \eqref{eq:v21-far-site-W}.
\end{proof}

Thus V20's failure of ordinary index lifting has a completely diagonal spatial
witness.  No phase information is needed to detect that example.

\subsection{The density shadow is not complete}
\label{sec:v21-density-not-complete}

The converse to Theorem~\ref{thm:v21-module-implies-density} is false.  The
simplest obstruction is internal bundle twisting at one physical site.

\begin{proposition}[Same density and same stable index can still have different exact modules]
\label{prop:v21-onsite-unstable-bundle}
Let $B=S^6$.  Let $F\to S^6$ be the nontrivial rank-two complex bundle defined
by the nonzero clutching element
\begin{equation}
 \pi_5(U(2))\cong\mathbb Z_2.
 \label{eq:v21-pi5U2}
\end{equation}
Its stable class is trivial because the stabilization map lands in
$\pi_5(U)\cong\mathbb Z$ and therefore kills the torsion element.  Embed $F$
and the trivial rank-two bundle as onsite subbundles of one sufficiently large
trivial internal fibre at a fixed lattice site $x_0$, and let $p_F,p_0$ be the
corresponding propagation-zero projector families.  Then
\begin{equation}
 \mu_{p_F}=2\delta_{x_0}=\mu_{p_0}.
 \label{eq:v21-same-density}
\end{equation}
\begin{equation}
 [p_F]=[p_0]\quad\text{in ordinary stable }K^0(B).
 \label{eq:v21-same-stable-K}
\end{equation}
\begin{equation}
 [p_F]_{\mathcal V_{\rm ctl}}
 \ne[p_0]_{\mathcal V_{\rm ctl}}.
 \label{eq:v21-different-exact-module}
\end{equation}
Thus the controlled index-density class and the ordinary stable family index
together still need not determine the exact controlled module at fixed finite
rank.
\end{proposition}

\begin{proof}
Both projectors act at the same site, so their diagonal lattice densities are
identical and their transport cost is trivial.  Stable triviality of $F$ gives
\eqref{eq:v21-same-stable-K}.  If a global controlled partial isometry between
$p_F$ and $p_0$ existed, then because both projections are propagation zero at
$x_0$ it would restrict fibrewise to a unitary vector-bundle isomorphism
$F\simeq\underline{\mathbb C}^2$.  This contradicts the nonzero unstable
clutching class.
\end{proof}

\begin{firewall}[The density invariant is an intermediate shadow]
The diagonal density detects coarse spatial transport but forgets internal
zero-mode bundle topology.  The exact V20 classifier remains
$\operatorname{Ind}^{\rm ex}_{\rm ctl}$.  V21 introduces
$\operatorname{Ind}^{\rm dens}_{\rm ctl}$ because it is a computable native
witness and the natural object for any proposed ``index-density anchoring''
theorem, not because it replaces the full controlled module.
\end{firewall}

\subsection{A conditional background-anchoring theorem}
\label{sec:v21-background-anchoring}

We can now state exactly what additional native information would close the
V20 gap for the physical Renewal regulator class.

\begin{definition}[Common background anchor atlas]
\label{def:v21-background-anchor-atlas}
Fix a common gauge/coframe/Higgs background family on $B$.  A collection of
background-determined anchor sets
\begin{equation}
 \mathfrak A_\alpha(b)
 =\{A_{\alpha,1}(b),\ldots,A_{\alpha,\kappa}(b)\}
 \label{eq:v21-background-anchor-sets}
\end{equation}
on a finite good cover is a \emph{common background anchor atlas} if:
\begin{enumerate}[label=\textup{(B\arabic*)},leftmargin=3em]
\item every set $A_{\alpha,j}(b)$ has uniformly bounded lattice diameter;
\item the sets and their retained source jets are constructed from the common
native background data, not from a regulator-specific zero-mode frame;
\item every admissible regulator in the declared class admits V15 localized
zero-mode frames whose $j$th anchor lies within a cutoff-uniform distance $R$
of $A_{\alpha,j}(b)$ after a local permutation;
\item regulators with the same ordinary family index can choose their local
zero-mode frames so that the transition cocycles are identified on overlaps
(in the V12 stable ambient range this follows from an actual bundle
isomorphism, not merely equality of determinant lines).
\end{enumerate}
\end{definition}

\begin{theorem}[Background anchoring lifts the ordinary index to the controlled index]
\label{thm:v21-background-anchor-lift}
Suppose two protected native regulator completions have the same sign and value
of the ordinary minimal family index and satisfy Definition
\ref{def:v21-background-anchor-atlas} for one common background anchor atlas.
Then their exact controlled family-index modules agree.  Consequently, for
every fixed coefficient packet and sufficiently large cutoff, V20's
classification theorem gives one global unstabilized native fixed-index path
between them.
\end{theorem}

\begin{proof}
On every chart choose the two V15 frames supplied by (B3).  Since both are
within distance $R$ of the same background anchor set, their matched anchor
distance is at most $2R$ plus the uniform anchor-set diameter.  Condition (B4)
puts the two frame systems into the same transition cocycle.  Hence they form a
matched controlled anchor atlas.  Theorem
\ref{thm:v21-global-anchor-characterization} gives one global controlled
partial isometry, i.e. equality of the exact controlled family-index modules.
Apply Theorem~\ref{thm:v20-classification}.
\end{proof}

\begin{corollary}[A regulator-independent local index-density anchor is sufficient only with bundle descent]
\label{cor:v21-density-anchor-plus-bundle}
Suppose a regulator-independent background density or collection of bounded
background regions forces every admissible zero-mode density to stay within a
uniform exponential transport distance of the same background anchor, and
suppose independently that the ordinary zero-mode bundles are identified in
the finite-rank Grassmannian.  If this information can be upgraded to the
frame-level localization in Definition~\ref{def:v21-background-anchor-atlas},
then the exact controlled family indices agree.  The density condition alone
is insufficient by Proposition~\ref{prop:v21-onsite-unstable-bundle}.
\end{corollary}

This is the precise point at which an overlap/local-index-density theorem would
enter the Renewal programme.  One would need a regulator-independent native
background quantity which anchors not only the total index but the spatial
support of the finite zero-mode module.

\subsection{Why the present Renewal assumptions do not yet prove anchoring}
\label{sec:v21-no-current-anchor-proof}

The V2--V10 regulator hypotheses fix the physical Dirac--Yukawa principal jet,
locality, covariance and the protected Wilson/chiral gaps.  These data control
hard-shell propagation and the coefficientwise continuum, but they do not
forbid adding a regulator-specific local finite-rank defect whose position is
not fixed by the common background.  V20 exploited exactly that freedom.

\begin{proposition}[Principal-symbol data do not determine the anchor atlas]
\label{prop:v21-principal-not-anchor}
In the unrestricted protected local-operator category used by V20, no theorem
of the form
\begin{equation}
 \begin{aligned}
 &\text{same physical principal symbol}\\
 &\quad+\ \text{same ordinary family index}
 \end{aligned}
 \quad\Longrightarrow\quad
 \text{common background anchor atlas}
 \label{eq:v21-false-anchor-implication}
\end{equation}
can follow from the existing hypotheses alone.
\end{proposition}

\begin{proof}
The V20 far-site minimal index-one pair has identical propagating principal
symbol and ordinary family index, but
Proposition~\ref{prop:v21-far-site-density} shows that the exponential transport
cost between its zero-mode densities grows as $e^{\gamma L/2}$.  A common
background anchor atlas with cutoff-uniform radius would, by
Theorem~\ref{thm:v21-module-implies-density}, give a cutoff-uniform transport
bound, a contradiction.
\end{proof}

Thus a positive V21-style anchoring theorem for the \emph{actual} Renewal ESM
class must use an additional structural statement which excludes such
regulator-specific defects.  Examples of admissible forms of such a statement
would be:
\begin{enumerate}[label=\textup{(N\arabic*)},leftmargin=3em]
\item a native local index-density theorem tying the zero-mode density to a
regulator-independent topological density of the common gauge/coframe/Higgs
background, with a quantitative localization remainder;
\item a microscopic naturality/minimality axiom forbidding additional
regulator-specific relevant or defect terms beyond the declared Wilson
completion and irrelevant deformations;
\item a direct relative low-energy theorem showing that two admissible
completions have zero-mode Riesz projectors related by a uniformly controlled
partial isometry.
\end{enumerate}
None of these stronger statements is currently contained in V1--V20.

\subsection{V21 classification picture}
\label{sec:v21-classification-picture}

The hierarchy of fixed-index information is now
\begin{equation}
 \boxed{
 \begin{array}{c}
 \text{exact controlled zero-mode module}\\[0.15em]
 \Downarrow\\[-0.1em]
 \text{controlled index-density transport class}\\[0.15em]
 \Downarrow\\[-0.1em]
 \text{ordinary family index / total zero-mode number}\\[0.15em]
 \Downarrow\\[-0.1em]
 \text{determinant-line shadow}.
 \end{array}}
 \label{eq:v21-information-hierarchy}
\end{equation}
The first implication can be strict by
Proposition~\ref{prop:v21-onsite-unstable-bundle}; the second is strict by the
V20/V21 far-site example.  V12 already showed that determinant data can forget
higher bundle information.

\begin{theorem}[Anchored-atlas reformulation of the V20 microscopic classifier]
\label{thm:v21-final-characterization}
At every fixed finite coefficient/source order and sufficiently large cutoff,
two faithful-complex protected minimal fixed-index Renewal regulators lie in
one global unstabilized native path component if and only if their zero-mode
packets admit a matched controlled anchor atlas.  Equivalently, the exact
controlled family-index module is the complete microscopic classifier, and its
spatial part is precisely bounded anchor transport.
\end{theorem}

\begin{proof}
Theorem~\ref{thm:v21-global-anchor-characterization} identifies a matched
anchor atlas with equality of the exact controlled family-index modules.
Theorem~\ref{thm:v20-classification} identifies that module equality with one
global unstabilized protected native path at the stated coefficientwise scope.
\end{proof}

\section{V21 anti-circularity audit}
\label{sec:v21-audit}

\begin{enumerate}[leftmargin=2.4em]
\item \textbf{The ordinary family index is not asked to encode spatial
transport.}  V20 already disproved that implication; V21 extracts the missing
transport datum explicitly.
\item \textbf{The anchor criterion is equivalent to the already-proved
controlled module, not a new assumption disguised as a theorem.}  The forward
direction transports a V15 frame through the controlled partial isometry; the
reverse direction reconstructs that partial isometry from matched localized
frames.
\item \textbf{Global bundle twisting is retained.}  The atlas theorem requires
matching transition cocycles.  Local anchor matching alone does not trivialize
a nontrivial zero-mode bundle.
\item \textbf{The index density is not promoted to a complete invariant.}
Proposition~\ref{prop:v21-onsite-unstable-bundle} gives identical spatial
density and identical stable $K^0$ class with different unstable exact modules.
\item \textbf{The far-site obstruction is measured by a native diagonal
quantity.}  Its exponential transport cost is exactly
$e^{\gamma L/2}$, so the failure is not an artefact of a bad frame choice.
\item \textbf{A positive physical anchoring theorem is stated conditionally.}
The current Renewal assumptions do not yet tie zero-mode positions to a
regulator-independent background density or anchor set.
\item \textbf{No continuum theorem is reopened.}  All new invariants remain
strictly post-continuum microscopic data; V16 scheme-groupoid equivalence is
unchanged.
\end{enumerate}

\section{V21 status ledger}
\label{sec:v21-status}

\subsection*{Proved in V21}
\begin{enumerate}[leftmargin=2.2em]
\item On a protected source chart, two fixed-rank zero-mode projectors are
controlled Murray--von Neumann equivalent if and only if they admit uniformly
localized orthonormal frames with a bounded matching of lattice anchors.  Under
an existing controlled partial isometry, the target frame can be chosen with
exactly the same anchors as the source frame.
\item Globally, equality of the exact controlled family-index modules is
equivalent to a finite matched anchor atlas whose source and target frames have
the same zero-mode bundle transition cocycle.  This separates ordinary bundle
descent from spatial anchoring exactly.
\item The diagonal zero-mode density
$\mu_p(x)=\operatorname{tr}p_{xx}$ has total mass $\kappa$.  Matched anchored
frames construct an explicit transport plan between two endpoint densities
with exponential moment bounded by $\kappa C^4e^{\gamma R}$.
\item Hence exact controlled family-index equality implies equality of a new
controlled index-density transport class.  This class is invariant under every
native fixed-index homotopy and forgets to the ordinary index by taking total
mass.
\item The V20 far-site minimal defect has exponential density-transport cost
exactly $e^{\gamma d(x_L,y_L)}$, so its failure of controlled equivalence is
visible directly in the projector diagonal.
\item The density shadow is not complete.  An onsite nontrivial stably trivial
rank-two bundle over $S^6$ and the onsite trivial rank-two bundle have identical
spatial density and the same ordinary stable $K^0$ class but different exact
finite-rank controlled modules.
\item A common regulator-independent background anchor atlas, together with
ordinary zero-mode bundle descent, is sufficient to lift ordinary family-index
equality to equality of the exact controlled family index and hence, by V20,
to one global unstabilized native regulator path.
\item The present V1--V20 assumptions do not imply such anchoring: the V20
far-site pair is an explicit logical counterexample inside the unrestricted
protected local-operator category.
\end{enumerate}

\subsection*{Inherited}
\begin{enumerate}[leftmargin=2.2em]
\item V1--V7: source-complete coefficientwise ESM continuum, common BV
restoration, fixed-loop forests, ultraviolet Cauchy convergence,
scheme/reference covariance and exact perturbative-order projectivity.
\item V8--V12: regulator-marked comparison, fixed-index reduced shell calculus,
global zero-mode bundle/line transport and faithful complex saturation.
\item V13--V16: exponential locality of protected Riesz zero modes, stable and
unstabilized finite-jet lifting, automatic local internal reserves and global
continuum descent in the finite scheme groupoid.
\item V17--V19: the primary controlled zero-mode module obstruction, the
secondary/tertiary unstable tower and its automatic collapse once the primary
class is globally trivialized.
\item V20: exact controlled family-index classification and the minimal
far-transport counterexample showing that ordinary family-index equality is too
coarse.
\end{enumerate}

\subsection*{Conditional}
\begin{enumerate}[leftmargin=2.2em]
\item The background-anchoring theorem is conditional on a regulator-independent
anchor atlas or an equivalent native index-density localization statement.
Such a theorem is not yet derived from the present Wilson/gauge/coframe/Higgs
grammar.
\item The controlled density class is a necessary spatial shadow of the exact
module but not a complete classifier; internal bundle data must still be
matched.
\item All statements retain fixed finite zero-mode rank and fixed finite
$(L,\Delta,\varrho,N)$, together with the positive IR/reference scale.  No
uniform source order, loop order or finite-amplitude interacting path is
asserted.
\item The V12 real Pfaffian mod-two orientation condition remains independent
of the complex controlled-index hierarchy.
\end{enumerate}

\subsection*{Open beyond V21}
\begin{enumerate}[leftmargin=2.2em]
\item Prove a genuine \emph{native local index-density theorem} for the actual
Renewal overlap/Wilson completion: identify a regulator-independent local
background density or bounded family of anchor regions, and prove that the
protected zero-mode density/frame packet remains within cutoff-uniform
controlled transport of it for every admissible completion sharing the same
physical Dirac--Yukawa symbol.
\item Alternatively, construct a Renewal-native background and two admissible
microscopic completions for which the controlled index-density classes differ,
thereby showing that the V20 exact module is an unavoidable additional native
choice.
\item Compute the independent real Pfaffian mod-two spectral-flow character and
analyze reduced-gap/index-changing divisors.
\item The independent positive-IR, analytic-RG, summability and nonperturbative
directions remain outside this fixed-order programme.
\end{enumerate}

\subsection*{Exact next load-bearing theorem}
\begin{center}
\fbox{\parbox{0.92\textwidth}{
\textbf{Native local index-density / background anchoring theorem.}
For the explicit Renewal Wilson/overlap completion and its admitted local
regulator deformations, construct a regulator-independent local density or
finite family of bounded anchor regions from the common gauge/coframe/Higgs
background.  Prove that every protected fixed-index zero-mode Riesz projector
has a V15 localized frame anchored within a cutoff-uniform bounded lattice
distance of this common background datum, with all retained source derivatives,
and prove compatible zero-mode bundle descent.  By
Theorem~\ref{thm:v21-background-anchor-lift} this would force equality of the
exact controlled family index for all admissible completions with the same
ordinary family index.  If such anchoring fails, exhibit the corresponding
native controlled index-density transport class explicitly.}}
\end{center}

\section*{Append-only continuation marker: V21}
\addcontentsline{toc}{section}{Append-only continuation marker: V21}
\textbf{End of V21.}  The complete V1--V20 body above is retained verbatim.
V21 gives an exact geometric reformulation of the V20 microscopic classifier:
controlled zero-mode equivalence is equivalent to a matched atlas of
exponentially localized zero-mode frames with bounded spatial anchors and the
same bundle transition cocycle.  It then extracts the diagonal controlled
index-density class, proves that every controlled module equivalence induces a
cutoff-uniform exponentially localized transport of that density, and shows
that the V20 far-site defect is detected exactly by divergent transport cost.
The density shadow is not complete because internal unstable bundle twisting
can survive at one site.  The remaining physical gate is therefore a native
background-anchoring/local-index-density theorem for the actual Renewal
Wilson/overlap regulator grammar; it cannot be inferred from the existing
principal-symbol and ordinary family-index data alone.



\clearpage
\section{V22 continuation: soft spectral anchoring, zero-mode residual control, and the common norm-resolvent bridge}
\label{sec:v22-continuation}

V21 identified the exact controlled zero-mode module as the remaining complex
microscopic datum and reformulated its spatial part as a matched atlas of
localized zero-mode frames.  It also showed that equality of the ordinary
family index and of the hard/physical principal symbol does not, by itself,
anchor those frames.  The next question is therefore not primarily a local
index-density identity.  A more direct upstream variable is the \emph{soft
spectral subspace itself}.

The load-bearing observation of V22 is elementary but strong.  Suppose two
microscopic regulators admit one common approximate resolvent on a fixed
complex contour which separates the protected zero-mode cluster from the
nonzero spectrum.  Then their true resolvents are uniformly close on that
contour, their Riesz zero-mode projectors are close in operator norm, and for
small enough cutoff spacing the norm difference is strictly below one.  Since
V13 already proved exponential locality of each protected Riesz projector,
V14's canonical direct rotation is then a globally defined controlled local
unitary.  Thus the V17 primary Cech class vanishes without any chartwise
choice, and the V19 cancellation theorem supplies one global unstabilized
native path.

This route is strictly stronger than the V2 hard-annulus comparison and
strictly weaker than proving one finite-amplitude analytic regulator trajectory.
It is also the natural spectral condition needed to exclude the V20 far-site
minimal defect.

Throughout this continuation fix one finite coefficient/source packet
\begin{equation}
 \mathfrak q=(L,\Delta,\varrho,N,\mathcal P_{\rm ext})
 \label{eq:v22-fixed-packet}
\end{equation}
and the positive IR/reference prescription.  Let $B$ be a compact protected
background/source panel on which the relevant chiral family has constant
minimal index $k$.  For definiteness we first write the kernel case $k>0$;
the cokernel case is obtained by replacing $D^*D$ by $DD^*$, and the balanced
case $k=0$ is immediate.

For a regulator $R\in\{A,B\}$ write
\begin{equation}
 D_{a,R}(b):E_a\longrightarrow F_a,
 \qquad
 L_{a,R}(b):=D_{a,R}(b)^*D_{a,R}(b),
 \label{eq:v22-lattice-positive}
\end{equation}
and let
\begin{equation}
 p_{a,R}(b):=\mathbf 1_{\{0\}}(L_{a,R}(b))
 \label{eq:v22-zero-projector}
\end{equation}
be its rank-$\kappa=k$ kernel projector.  By the V11--V13 protected-zero-mode
hypotheses there is a common reduced singular-value floor
\begin{equation}
 L_{a,R}(b)\succeq \tau_*^2\,(I-p_{a,R}(b))
 \label{eq:v22-gap}
\end{equation}
with $\tau_*>0$ independent of $a$, $b$, and $R$ on the declared panel.

\subsection{A purely native cross-residual criterion}
\label{sec:v22-residual}

Before introducing any continuum reference, it is useful to isolate the exact
endpoint quantity which controls the angle between two protected zero-mode
spaces.

\begin{definition}[Native zero-mode residual]
\label{def:v22-residual}
For two minimal positive-index protected endpoints define
\begin{equation}
 \epsilon_{\rm zm}(A,B)
 :=\max\bigl\{
 \|D_{a,B}p_{a,A}\|,
 \|D_{a,A}p_{a,B}\|
 \bigr\}.
 \label{eq:v22-residual}
\end{equation}
For $k<0$, replace $D_{a,R}$ by $D_{a,R}^*$ and the kernel projectors by the
cokernel projectors.  On a source panel the norm means the supremum in $b$;
fixed source derivatives are retained separately as in V13.
\end{definition}

\begin{lemma}[Residual controls the principal angle]
\label{lem:v22-residual-angle}
Under the common reduced floor \eqref{eq:v22-gap},
\begin{align}
 \|(I-p_{a,B})p_{a,A}\|
 &\le \tau_*^{-1}\|D_{a,B}p_{a,A}\|,
 \label{eq:v22-angle-one}\\
 \|(I-p_{a,A})p_{a,B}\|
 &\le \tau_*^{-1}\|D_{a,A}p_{a,B}\|.
 \label{eq:v22-angle-two}
\end{align}
Consequently
\begin{equation}
 \boxed{
 \|p_{a,A}-p_{a,B}\|
 \le {2\epsilon_{\rm zm}(A,B)\over\tau_*}.}
 \label{eq:v22-projector-residual}
\end{equation}
\end{lemma}

\begin{proof}
For every $x\in\Ran p_{a,A}$, decompose
$x=p_{a,B}x+(I-p_{a,B})x$.  Since $D_{a,B}p_{a,B}=0$ and
\eqref{eq:v22-gap} holds on the orthogonal complement,
\[
 \|D_{a,B}x\|^2
 =\langle (I-p_{a,B})x,
 L_{a,B}(I-p_{a,B})x\rangle
 \ge \tau_*^2\|(I-p_{a,B})x\|^2.
\]
Taking the supremum over unit $x\in\Ran p_{a,A}$ gives
\eqref{eq:v22-angle-one}; the second estimate is symmetric.  Finally
\[
 p_{a,A}-p_{a,B}
 =(I-p_{a,B})p_{a,A}-p_{a,B}(I-p_{a,A}),
\]
and the triangle inequality gives \eqref{eq:v22-projector-residual}.
\end{proof}

\begin{theorem}[Residual criterion for exact controlled equivalence]
\label{thm:v22-residual-controlled}
Assume that both zero-mode projectors have the V13 cutoff-uniform exponential
locality/source bounds and
\begin{equation}
 \epsilon_{\rm zm}(A,B)<\frac{\tau_*}{2}.
 \label{eq:v22-residual-small}
\end{equation}
Then
\begin{equation}
 \|p_{a,A}-p_{a,B}\|<1.
 \label{eq:v22-projector-under-one}
\end{equation}
The V14 direct-rotation formula therefore gives a globally defined controlled
unitary
\begin{equation}
 u_{B\leftarrow A}
 =X(X^*X)^{-1/2},
 \qquad
 X=p_{a,B}p_{a,A}+(I-p_{a,B})(I-p_{a,A}),
 \label{eq:v22-direct-rotation}
\end{equation}
with
\begin{equation}
 \boxed{
 u_{B\leftarrow A}p_{a,A}u_{B\leftarrow A}^*=p_{a,B}.}
 \label{eq:v22-controlled-equivalence}
\end{equation}
Every prescribed fixed source derivative belongs to one common positive
exponential locality class after the usual finite exponent loss.  Hence the
exact V20 controlled family-index modules agree.
\end{theorem}

\begin{proof}
Equation \eqref{eq:v22-projector-residual} and
\eqref{eq:v22-residual-small} give
\eqref{eq:v22-projector-under-one}.  V14 then proves that $X$ is invertible,
$X^*X=I-(p_{a,A}-p_{a,B})^2$, and the polar factor
\eqref{eq:v22-direct-rotation} is a unitary in the identity component which
conjugates the two projections.  V13 makes the endpoint projectors and all
retained source derivatives exponentially local.  The V13 local resolvent
functional calculus applied to $(X^*X)^{-1/2}$ therefore gives the same fixed
source calculus for $u_{B\leftarrow A}$.  The partial isometry
$v=u_{B\leftarrow A}p_{a,A}$ is exactly the controlled Murray--von Neumann
isomorphism defining equality of the V20 exact controlled index.
\end{proof}

The residual criterion is useful because it is entirely native.  It asks only
whether each endpoint operator almost annihilates the other's protected
zero-mode packet.  The V20 far-site defect fails this test by an $O(1)$ amount.

\subsection{A common soft spectral anchor}
\label{sec:v22-soft-anchor}

We now formulate a stronger but more structural sufficient condition which can
be attacked directly from the explicit Wilson/overlap stencil.

Let $\mathcal H_0$ be a continuum comparison space for the fixed
finite-volume background panel, and let
\begin{equation}
 D_0(b):E_0\longrightarrow F_0,
 \qquad
 L_0(b):=D_0(b)^*D_0(b)
 \label{eq:v22-continuum-positive}
\end{equation}
be the common ordinary covariant Dirac--Yukawa operator determined by the
physical principal background data.  Assume its kernel rank is the same
$\kappa$ and that, on the compact panel,
\begin{equation}
 \operatorname{spec}L_0(b)
 \subset \{0\}\cup[\tau_0^2,\infty)
 \label{eq:v22-continuum-gap}
\end{equation}
for some $\tau_0>0$.  Choose once and for all a positively oriented contour
$\Gamma$ enclosing $0$ but lying inside
$|z|<\frac12\min\{\tau_*^2,\tau_0^2\}$.

Let
\begin{equation}
 J_a:E_a\longrightarrow E_0
 \label{eq:v22-identification}
\end{equation}
be one regulator-independent lattice-to-continuum identification associated
with the common Renewal carrier.  Only the compressed continuum resolvent
$J_a^*(z-L_0)^{-1}J_a$ will be used, so no assertion that $J_a$ is surjective
is required.

\begin{definition}[Common soft norm-resolvent anchor]
\label{def:v22-soft-resolvent}
The two regulators $A,B$ have a \emph{common soft norm-resolvent anchor} on
$\Gamma$ if there is $\varepsilon_a\downarrow0$ such that, for
$R\in\{A,B\}$,
\begin{equation}
 \boxed{
 \sup_{b\in B}\sup_{z\in\Gamma}
 \left\|
 (z-L_{a,R}(b))^{-1}
 -J_a^*(z-L_0(b))^{-1}J_a
 \right\|
 \le\varepsilon_a.}
 \label{eq:v22-soft-resolvent}
\end{equation}
At source order $\varrho$ we require the analogous estimate, with constants
$\varepsilon_{a,I}\to0$, for every retained fixed derivative
$\partial_I$ of the displayed resolvents.
\end{definition}

The definition is a generalized norm-resolvent statement in one fixed soft
spectral window.  It does not ask for uniform comparison of the whole UV
spectrum and it does not require a finite-amplitude interacting RG trajectory.

\begin{theorem}[Common soft anchor forces the exact controlled family index]
\label{thm:v22-soft-anchor-controlled}
Assume Definition~\ref{def:v22-soft-resolvent}.  Then
\begin{equation}
 \boxed{
 \sup_{b\in B}
 \|p_{a,A}(b)-p_{a,B}(b)\|
 \le {|\Gamma|\over\pi}\,\varepsilon_a.}
 \label{eq:v22-riesz-difference}
\end{equation}
In particular, for all sufficiently small $a$ the two projectors have norm
distance less than one and therefore are globally controlled-equivalent by the
canonical direct rotation \eqref{eq:v22-direct-rotation}.  With source-jet
version of Definition~\ref{def:v22-soft-resolvent}, every retained derivative
of the two Riesz projectors differs by $o(1)$ and the direct rotation is a
single global source-smooth controlled equivalence.  Hence
\begin{equation}
 \boxed{
 \operatorname{Ind}^{\rm ex}_{\rm ctl}(D_{a,A})
 =\operatorname{Ind}^{\rm ex}_{\rm ctl}(D_{a,B})}
 \label{eq:v22-exact-index-equality}
\end{equation}
for all sufficiently small $a$.
\end{theorem}

\begin{proof}
The Riesz formula gives
\[
 p_{a,R}(b)
 ={1\over2\pi i}\oint_\Gamma(z-L_{a,R}(b))^{-1}\,dz.
\]
Subtract the $R=A$ and $R=B$ formulas and insert the same compressed continuum
resolvent between them.  The triangle inequality and
\eqref{eq:v22-soft-resolvent} give
\[
 \|p_{a,A}-p_{a,B}\|
 \le {1\over2\pi}|\Gamma|(\varepsilon_a+\varepsilon_a)
 ={|\Gamma|\over\pi}\varepsilon_a,
\]
uniformly on $B$.  The bound tends to zero, hence is eventually less than one.
Apply Theorem~\ref{thm:v22-residual-controlled} without using the residual
estimate itself: V14's direct rotation only needs the projection-distance
bound and V13 supplies locality.  Differentiating the Riesz formula gives the
source-jet statement directly from the differentiated resolvent hypothesis.
The direct-rotation source calculus is then the V14/V13 inverse-functional
calculus.  Equality of exact controlled index modules is precisely the
existence of this global controlled partial isometry.
\end{proof}

\begin{corollary}[Soft spectral anchoring kills all complex microscopic descent classes]
\label{cor:v22-soft-kills-all}
On the faithful complex minimal fixed-index branch, the hypotheses of
Theorem~\ref{thm:v22-soft-anchor-controlled} imply, for sufficiently small
$a$:
\begin{enumerate}[label=\textup{(S\arabic*)},leftmargin=3em]
\item the V17 primary controlled zero-mode Cech class vanishes;
\item the V18 complement and identity-component classes vanish coefficientwise
by V19;
\item one global unstabilized V11-admissible native path exists at every fixed
finite coefficient/source order;
\item the endpoint coefficientwise continua are related by the global V19/V20
matching map.
\end{enumerate}
\end{corollary}

\begin{proof}
Theorem~\ref{thm:v22-soft-anchor-controlled} supplies one global controlled
partial isometry, so the V17 class is zero.  V19 proves that the secondary and
tertiary V18 obstruction layers then collapse at every fixed source order once
the cutoff is sufficiently large.  The V20 classification theorem converts
equality of the exact controlled family-index module into the global
unstabilized native path and continuum matching.
\end{proof}

\subsection{A common parametrix is enough}
\label{sec:v22-parametrix}

Definition~\ref{def:v22-soft-resolvent} is already much closer to the desired
native theorem than a statement about a local topological density, but it can
be reduced one step further to a pair of explicit error estimates for an
approximate inverse.

\begin{definition}[Soft parametrix gate]
\label{def:v22-parametrix}
For $z\in\Gamma$ define the common comparison operator
\begin{equation}
 Q_a(z,b):=J_a^*(z-L_0(b))^{-1}J_a.
 \label{eq:v22-Q}
\end{equation}
For regulator $R$, set
\begin{align}
 \mathcal E_{a,R}^{\rm L}(z,b)
 &:=(z-L_{a,R}(b))Q_a(z,b)-I,
 \label{eq:v22-left-error}\\
 \mathcal E_{a,R}^{\rm R}(z,b)
 &:=Q_a(z,b)(z-L_{a,R}(b))-I.
 \label{eq:v22-right-error}
\end{align}
The \emph{soft parametrix gate} holds if
\begin{equation}
 \boxed{
 \sup_{R,b,z}
 \bigl(
 \|\mathcal E_{a,R}^{\rm L}\|
 +\|\mathcal E_{a,R}^{\rm R}\|
 \bigr)
 \le\eta_a,
 \qquad
 \eta_a\downarrow0,}
 \label{eq:v22-parametrix-gate}
\end{equation}
and analogously for every fixed retained source derivative.
\end{definition}

\begin{theorem}[Two-sided soft parametrix implies the common norm-resolvent anchor]
\label{thm:v22-parametrix-to-resolvent}
Assume the base resolvent/compression $Q_a(z,b)$ is uniformly bounded on
$\Gamma$ and \eqref{eq:v22-parametrix-gate} holds with
$\eta_a<1/2$.  Then $z-L_{a,R}(b)$ is invertible for all
$R,b,z$ and
\begin{equation}
 \boxed{
 \left\|(z-L_{a,R})^{-1}-Q_a\right\|
 \le {2\eta_a\over 1-\eta_a}\,\|Q_a\|}
 \label{eq:v22-parametrix-resolvent}
\end{equation}
(up to replacing the harmless numerical factor by the larger of the left- and
right-parametrix bounds).  In particular Definition
\ref{def:v22-soft-resolvent} follows with $\varepsilon_a=O(\eta_a)$.
The same implication holds coefficientwise for every prescribed fixed source
jet.
\end{theorem}

\begin{proof}
Write $A=z-L_{a,R}$ and $Q=Q_a$.  The left identity is
$AQ=I+E_L$ with $\|E_L\|<1$.  Hence $AQ$ is invertible and
\[
 A\,Q(I+E_L)^{-1}=I,
\]
so $A$ has a bounded right inverse.  The right identity similarly gives a
bounded left inverse.  They coincide, proving invertibility.  Moreover
\[
 A^{-1}=Q(I+E_L)^{-1}
\]
and therefore
\[
 A^{-1}-Q
 =-Q(I+E_L)^{-1}E_L.
\]
Thus
\[
 \|A^{-1}-Q\|
 \le \|Q\|{\eta_a\over1-\eta_a}.
\]
Using both-sided estimates and one common constant gives the displayed slightly
looser bound.  Source derivatives follow by differentiating the identities and
using the ordinary inverse derivative formula.  Because only finitely many
source orders are retained, the induction is finite and consumes only the
already-declared uniform base pivots.
\end{proof}

Theorem~\ref{thm:v22-parametrix-to-resolvent} isolates the exact native work
still missing: construct $J_a$ and prove the two parametrix residuals for the
explicit overlap/Wilson/Yukawa operator on the protected nonzero-index soft
sector.

\subsection{Relation to the V21 controlled index density}
\label{sec:v22-density}

The soft spectral theorem immediately implies the spatial anchoring statement
V21 sought, but without first guessing a background density.

\begin{corollary}[Soft resolvent bridge gives a matched anchor atlas]
\label{cor:v22-soft-anchor-atlas}
Under Theorem~\ref{thm:v22-soft-anchor-controlled}, the V21 zero-mode packets
admit one global matched controlled anchor atlas for all sufficiently small
$a$.  Their controlled index-density transport class therefore agrees, and for
some positive exponent $\gamma$ the V21 exponential transport cost remains
cutoff-uniform.
\end{corollary}

\begin{proof}
The direct-rotation partial isometry is controlled.  Apply the V21 local/global
anchor characterization: transport a V15 localized frame of $p_{a,A}$ through
the controlled unitary.  The target frame may be taken with the same anchors,
and the global unitary automatically preserves the transition cocycle.  V21
then supplies the exponential transport plan for the diagonal zero-mode
densities.
\end{proof}

\begin{proposition}[Trace-norm closeness of the zero-mode density]
\label{prop:v22-density-TV}
If the common zero-mode rank is $\kappa$ and
$\|p_{a,A}-p_{a,B}\|\le\delta_a$, then
\begin{equation}
 \boxed{
 \sum_x\left|
 \mu_{p_{a,A}}(x)-\mu_{p_{a,B}}(x)
 \right|
 \le 2\kappa\,\delta_a.}
 \label{eq:v22-density-TV}
\end{equation}
Hence the soft resolvent bridge forces the two diagonal index densities to
coincide in total variation as $a\downarrow0$.
\end{proposition}

\begin{proof}
The difference of two rank-$\kappa$ projections has rank at most $2\kappa$, so
\[
 \|p_{a,A}-p_{a,B}\|_1
 \le2\kappa\|p_{a,A}-p_{a,B}\|.
\]
For any trace-class matrix kernel, the sum of absolute values of the fibre
traces of its diagonal blocks is bounded by its trace norm.  Apply this to
$p_{a,A}-p_{a,B}$.
\end{proof}

Thus a common soft spectral anchor is stronger than equality of a local index
density: it controls the entire finite-rank zero-mode projector, including the
internal bundle data which V21 showed the diagonal density can forget.

\subsection{Why the V2 hard-shell theorem does not imply the soft anchor}
\label{sec:v22-v2-gap}

The present theorem is not hidden inside E1.  V2 was designed to control hard
fermionic annuli and the complementary full Brillouin zone, not an isolated
zero eigencluster.

\begin{proposition}[Hard symbol matching does not imply soft spectral matching]
\label{prop:v22-hard-not-soft}
The V2 conclusions
\begin{equation}
 d_{\chi,a}(p)=d_{\chi,0}(p)+O(a|p|^2)
 \label{eq:v22-v2-hard-form}
\end{equation}
on physical hard annuli, together with the full-zone Wilson gap and equality of
the ordinary family index, do not imply Definition
\ref{def:v22-soft-resolvent}.
\end{proposition}

\begin{proof}
Use the V20 minimal far-site positive-index pair.  Its propagating physical
Dirac block, hard symbol and ordinary family index are identical at the two
endpoints; the only difference is the position of the protected finite-rank
kernel defect.  V21 proves that the two zero-mode projectors have divergent
controlled transport and in particular
\begin{equation}
 \|p_{a,A}-p_{a,B}\|=1
 \label{eq:v22-farsite-proj}
\end{equation}
for the orthogonal far-site kernel lines.  If a common soft norm-resolvent
anchor with $\varepsilon_a\to0$ followed from the hard conclusions,
Theorem~\ref{thm:v22-soft-anchor-controlled} would force the left-hand side of
\eqref{eq:v22-farsite-proj} to tend to zero, a contradiction.  Hence a new soft
spectral input is logically necessary.
\end{proof}

\begin{firewall}[Do not infer zero-mode convergence from hard-annulus consistency]
The V2 hard-shell theorem identifies the correct continuum kinetic branch and
controls every fixed hard/source jet required by the ultraviolet forest.  It
does not compare the spectral projector at the isolated eigenvalue zero.  V22
therefore adds no retrospective assumption to E1--E6; it identifies a genuinely
stronger post-continuum theorem needed only for microscopic regulator-phase
classification.
\end{firewall}

\subsection{External methodological comparison}
\label{sec:v22-literature}

\begin{remark}[Why norm-resolvent is the right spectral strength]
\label{rem:v22-literature}
There is independent mathematical precedent for this route.  In simplified
square-lattice Dirac models, adding an appropriate Wilson term can restore
(generalized) norm-resolvent convergence to the continuum Dirac operator,
whereas naive discretizations can fail norm-resolvent convergence because of
doubling.  Separately, classical-continuum calculations for the overlap
operator show convergence of its anomaly/index density to the continuum
topological density on suitable smooth gauge backgrounds.  These results
support the spectral-anchor strategy, but V22 does not import them as proofs:
they do not by themselves establish the required gauge/spin/coframe/Yukawa
source-complete soft parametrix for the present Renewal overlap family.
\end{remark}

The distinction is important.  A local topological density theorem certifies
the integrated index and its local Chern-character shadow.  V21 showed that the
exact controlled zero-mode module contains more information.  A common soft
Riesz-resolvent theorem, by contrast, directly controls the whole zero-mode
projection and is therefore exactly strong enough for the V20 classifier.

\subsection{A conditional native theorem from consistency and stability}
\label{sec:v22-consistency-stability}

For future use it is convenient to record the minimal analytic ingredients
which would verify the soft parametrix gate.

\begin{assumption}[Soft native consistency package]
\label{ass:v22-soft-consistency}
For one continuum identification $J_a$ and contour $\Gamma$, suppose that for
both regulator endpoints and every retained source jet:
\begin{enumerate}[label=\textup{(C\arabic*)},leftmargin=3em]
\item the continuum resolvent $(z-L_0)^{-1}$ is uniformly bounded on
$B\times\Gamma$;
\item the sampled continuum resolvent $Q_a=J_a^*(z-L_0)^{-1}J_a$ is uniformly
bounded;
\item the native consistency defect obeys
\begin{equation}
 \|(z-L_{a,R})Q_a-I\|\le C\omega_a,
 \qquad \omega_a\downarrow0;
 \label{eq:v22-consistency-left}
\end{equation}
\item the dual consistency defect obeys
\begin{equation}
 \|Q_a(z-L_{a,R})-I\|\le C\omega_a;
 \label{eq:v22-consistency-right}
\end{equation}
\item the same bounds hold after every fixed derivative consumed by
$\mathfrak q$.
\end{enumerate}
\end{assumption}

\begin{theorem}[Consistency plus soft spectral stability closes the V21 anchoring gap]
\label{thm:v22-consistency-closes}
Under Assumption~\ref{ass:v22-soft-consistency}, for all sufficiently small
$a$ the two native regulators have the same exact controlled family index, the
same controlled index-density transport class, and one global unstabilized
coefficientwise native regulator path.  Their coefficientwise continuum
matching is therefore the V19/V20 global matching with no additional complex
microscopic obstruction.
\end{theorem}

\begin{proof}
Assumption~\ref{ass:v22-soft-consistency} is precisely the soft parametrix gate
with $\eta_a=C\omega_a$.  Apply
Theorem~\ref{thm:v22-parametrix-to-resolvent}, then
Theorem~\ref{thm:v22-soft-anchor-controlled}, then
Corollary~\ref{cor:v22-soft-kills-all}.  The controlled density conclusion is
Corollary~\ref{cor:v22-soft-anchor-atlas}.
\end{proof}

\subsection{V22 anti-circularity audit}
\label{sec:v22-audit}

\begin{enumerate}[leftmargin=2.4em]
\item \textbf{Ordinary family-index equality is not promoted to zero-mode
projector convergence.}  The far-site defect remains an explicit
counterexample.
\item \textbf{Hard-annulus matching is not reused outside its proven scale
regime.}  Proposition~\ref{prop:v22-hard-not-soft} isolates the missing soft
spectral theorem explicitly.
\item \textbf{The background anchor is spectral, not guessed geometrically.}
V22 does not assume that zero modes are localized near a chosen topological
charge density.  It asks instead for one common soft resolvent/parametrix.
\item \textbf{The exact controlled module, not merely its density shadow, is
controlled.}  Riesz projector norm closeness gives the canonical direct
rotation and therefore the full finite-rank bundle/module isomorphism.
\item \textbf{No full-spectrum norm-resolvent theorem is demanded.}  Only one
fixed contour separating the protected zero cluster is used.
\item \textbf{No finite-amplitude source trajectory is introduced.}  Source
control remains coefficientwise at the finite jet order consumed by
$\mathfrak q$.
\item \textbf{External continuum-limit literature is methodological only.}
Its hypotheses have not been silently identified with the Renewal
spin/gauge/Yukawa regulator.
\end{enumerate}

\section{V22 status ledger}
\label{sec:v22-status}

\subsection*{Proved in V22}
\begin{enumerate}[leftmargin=2.2em]
\item A purely native cross-residual estimate controls the angle between two
minimal fixed-index zero-mode spaces.  If each endpoint almost annihilates the
other's zero-mode packet by less than $\tau_*/2$, the projectors are at norm
distance less than one and the V14 direct rotation gives their exact controlled
module equivalence.
\item If two regulators share one common soft norm-resolvent anchor on a fixed
contour around zero, their zero-mode Riesz projectors differ by at most
$|\Gamma|\varepsilon_a/\pi$.  Hence for small cutoff spacing they have the
same exact V20 controlled family index globally, with all retained source jets.
\item The same theorem kills the entire V17--V19 complex microscopic descent
tower: the primary controlled module isomorphism is canonical, and V19 removes
the downstream unstable complement/unitary obstructions coefficientwise.
\item A two-sided common soft parametrix with residual $o(1)$ is sufficient for
the common norm-resolvent anchor.  This reduces the native task to explicit
left/right parametrix error estimates for the Renewal overlap/Wilson/Yukawa
operator.
\item Soft projector convergence implies a matched V21 anchor atlas and a
cutoff-uniform exponential transport plan for the diagonal zero-mode densities.
For rank $\kappa$, the total-variation difference of the two diagonal densities
is bounded by $2\kappa\|p_A-p_B\|$.
\item The V2 hard-annulus/full-zone theorem and ordinary family-index equality
do not imply the new soft spectral anchor.  The V20 far-site minimal defect is
a direct obstruction.
\item Consequently the missing physical anchoring theorem can be replaced by a
more precise spectral statement: prove the V22 soft consistency/parametrix
package for the actual Renewal native regulator family.
\end{enumerate}

\subsection*{Inherited}
\begin{enumerate}[leftmargin=2.2em]
\item V1--V7: source-complete coefficientwise ESM continuum, common BV
restoration, arbitrary-fixed-loop forest theorem, ultraviolet Cauchy
convergence, scheme/reference covariance and exact perturbative-order
projectivity.
\item V8--V12: relative regulator comparison, fixed-index shell calculus,
global zero-mode bundle/line transport and faithful complex saturation.
\item V13--V19: locality of protected Riesz projectors, controlled zero-mode
modules, automatic coefficientwise cancellation after primary descent, and the
global microscopic obstruction tower.
\item V20--V21: exact controlled family-index classification, far-transport
counterexamples, matched anchor atlases and the controlled index-density
shadow.
\end{enumerate}

\subsection*{Conditional}
\begin{enumerate}[leftmargin=2.2em]
\item The positive V22 anchoring theorem is conditional on a common soft
norm-resolvent/parametrix estimate.  V22 proves its consequences but not yet the
native estimate itself for a topologically nontrivial gauge/spin/Yukawa
background family.
\item The continuum comparison operator $D_0$ must have a constant kernel rank
and a positive reduced spectral gap on the compact protected panel.  Crossing a
zero-mode divisor is outside this theorem.
\item The lattice-to-continuum identification $J_a$ must be common to the two
microscopic completions and compatible with the fixed source/background panel.
\item All statements retain a fixed finite coefficient packet and fixed
zero-mode rank, together with the positive IR/reference scale.  No uniform
loop/source order or finite-amplitude interacting trajectory is asserted.
\item The independent real Pfaffian mod-two orientation/zero-crossing problem
remains separate.
\end{enumerate}

\subsection*{Open beyond V22}
\begin{enumerate}[leftmargin=2.2em]
\item Prove the V22 soft parametrix gate for the explicit Renewal Wilson/overlap
completion on a compact protected nonzero-index background family.  The proof
must include gauge parallel transport, spin/coframe dependence, Yukawa/Higgs
terms and every retained source derivative.
\item Construct a regulator-independent lattice/continuum identification
$J_a$ and show that the Wilson high-zone gap suppresses all doubler/high-mode
feedback into the isolated soft contour while the physical finite-difference
stencil has $O(a)$ (or better) consistency there.
\item If the soft parametrix fails, isolate the exact background or microscopic
sector producing an $O(1)$ residual and relate it to the V20/V21 controlled
transport invariant.
\item Compute the independent real Pfaffian mod-two spectral-flow character and
analyze genuine reduced-gap/index-changing divisors.
\item The independent positive-IR, analytic-RG, perturbative-summability and
nonperturbative directions remain outside this fixed-order programme.
\end{enumerate}

\subsection*{Exact next load-bearing theorem}
\begin{center}
\fbox{\parbox{0.92\textwidth}{
\textbf{Native soft norm-resolvent / parametrix theorem.}
For the explicit V2 Renewal Wilson/overlap/Yukawa completion on a compact
protected fixed-index gauge/coframe/Higgs background family, construct one
regulator-independent lattice--continuum identification $J_a$ and the common
soft parametrix
$Q_a(z)=J_a^*(z-L_0)^{-1}J_a$ on a contour isolating the continuum zero-mode
cluster.  Prove, with every fixed source derivative consumed by the declared
coefficient packet,
\[
 \|(z-L_{a,R})Q_a-I\|+
 \|Q_a(z-L_{a,R})-I\|
 \le C\,\omega_a,
 \qquad \omega_a\to0,
\]
uniformly for both admissible microscopic completions.  The proof must use the
native Wilson high-zone/no-doubler gap and the actual gauge/spin/coframe/Yukawa
finite-difference consistency; it may not infer soft spectral convergence from
the V2 hard-annulus theorem alone.  Combined with
Theorem~\ref{thm:v22-consistency-closes}, this would identify the exact
controlled family index of all such completions with the common continuum soft
spectral subspace and close the last complex microscopic anchoring gap.}}
\end{center}

\section*{Append-only continuation marker: V22}
\addcontentsline{toc}{section}{Append-only continuation marker: V22}
\textbf{End of V22.}  The complete V1--V21 body above is retained verbatim.
V22 replaces the unresolved spatial-background-anchor conjecture by a sharper
soft spectral criterion.  A native cross-residual theorem shows directly that
small endpoint zero-mode residuals force controlled equivalence.  More
structurally, one common generalized norm-resolvent anchor on a contour around
the protected zero cluster forces the two native Riesz projectors to be
$o(1)$-close; V14's direct rotation then gives a canonical global controlled
partial isometry, so the V20 exact family indices agree and the full V17--V19
complex microscopic descent tower collapses.  A two-sided common soft
parametrix is sufficient for this resolvent anchor.  The V2 hard-annulus theorem
is proved insufficient by the far-site defect, so the exact next gate is the
native soft parametrix estimate for the explicit Renewal Wilson/overlap/Yukawa
regulator on protected nonzero-index backgrounds.


\clearpage
\section{V23 continuation: two-zone soft resolvent theory, Feshbach anchoring, and the exact nonconstant-background gate}
\label{sec:v23-continuation}

V22 isolated the correct spectral strength for collapsing the controlled
zero-mode obstruction, but its final ``common parametrix'' formulation was
stronger than necessary.  A pure continuum resolvent should not be expected to
invert the cutoff-scale lattice dispersion: the Wilson term is deliberately
noncontinuum there.  What the zero-mode Riesz projector needs is only that the
\emph{soft} part of the native resolvent converge to one common continuum
resolvent, while the complementary native resolvent tends to zero in operator
norm.  The high block may remain regulator dependent.

This continuation makes that correction exact.  It proves a two-zone Feshbach
theorem with three inputs:
\begin{equation}
 \boxed{
 \begin{gathered}
 \text{low-block continuum consistency}
 +\text{ high-block coercivity}\\
 +\text{ small low/high resolvent transfer}
 \end{gathered}}
 \label{eq:v23-three-inputs}
\end{equation}
These imply norm convergence of the full native resolvent to the embedded soft
continuum resolvent on the contour isolating the zero-mode cluster.  In
particular the two endpoint Riesz projectors become norm-close and the V22
controlled-family-index theorem applies.

The theorem is then verified completely for the translation-invariant V2
reference branch, with the explicit intermediate scale
\(M_a=a^{-1/3}\) and rate \(O(a^{2/3})\).  The same fixed-order source-word
power counting is recorded.  What does \emph{not} close from V1--V22 is the
corresponding high-frequency coercivity/mixing estimate on a genuinely
nonconstant fixed-index gauge--spin--coframe background.  That missing estimate
is isolated at the end as a native discrete G\aa rding/Feshbach theorem rather
than hidden inside the hard-annulus result.

\subsection{Why the full-lattice pure continuum parametrix is too strong}
\label{sec:v23-full-parametrix-too-strong}

The issue is visible already in the exact V2 free symbol.

\begin{proposition}[A pure continuum resolvent does not invert the ultraviolet Wilson branch]
\label{prop:v23-pure-parametrix-fails}
Consider the identity-coframe, identity-link, zero-Yukawa V2 branch on a periodic
lattice, and identify lattice Fourier momentum with the corresponding continuum
momentum.  Let
\begin{equation}
 L_a(p):=d_{{\rm ov},a}(p)^*d_{{\rm ov},a}(p),
 \qquad
 L_0(p):=|p|^2 I.
 \label{eq:v23-free-L}
\end{equation}
At the Brillouin point
\(p_a=(\pi/a,0,0,0)\), one has
\begin{equation}
 L_a(p_a)=\frac4{a^2}I,
 \qquad
 L_0(p_a)=\frac{\pi^2}{a^2}I.
 \label{eq:v23-UV-values}
\end{equation}
Hence for every fixed bounded \(z\),
\begin{equation}
 \boxed{
 \left\|
 (z-L_a(p_a))(z-L_0(p_a))^{-1}-I
 \right\|
 \longrightarrow
 \left|\frac4{\pi^2}-1\right|>0.}
 \label{eq:v23-UV-residual}
\end{equation}
Thus the V22 two-sided residual with a \emph{pure} continuum resolvent on the
entire lattice Hilbert space cannot hold for the canonical physical
momentum-preserving identification, even on the simplest native branch.
\end{proposition}

\begin{proof}
At \(x=ap=(\pi,0,0,0)\) one has
\(s_\mu=0\), \(W=2\), \(B=1\), and \(R=1\).  The exact V2 overlap symbol
\eqref{eq:v2-exact-overlap-symbol} is therefore
\(d_{{\rm ov},a}=2a^{-1}I\), proving the first identity in
\eqref{eq:v23-UV-values}.  The continuum squared Dirac symbol is
\(|p_a|^2=\pi^2a^{-2}\).  Substitute these values and let \(a\downarrow0\):
\[
 (z-4a^{-2})(z-\pi^2a^{-2})^{-1}-1
 \longrightarrow {4\over\pi^2}-1.
\]
Taking absolute values gives \eqref{eq:v23-UV-residual}.
\end{proof}

\begin{firewall}[Correction to the final V22 gate]
The consequences proved in V22 from a common soft resolvent remain correct.
What V23 corrects is only the sufficient \emph{full-lattice} parametrix used in
Assumption~\ref{ass:v22-soft-consistency}.  At cutoff momenta the native and
continuum dispersions need not agree and should not be forced to agree.  The
high native resolvent instead has norm \(O(M_a^{-2})\) once an intermediate
scale \(M_a\to\infty\) is chosen, which is enough for the Riesz projector near
zero.
\end{firewall}

\subsection{Abstract two-zone Feshbach theorem}
\label{sec:v23-Feshbach}

Let \(\mathcal H_a\) be the native one-particle space and let
\begin{equation}
 \mathcal H_a=\mathcal H_a^{<}\oplus\mathcal H_a^{>},
 \qquad
 P_a+Q_a=I
 \label{eq:v23-split}
\end{equation}
be one common low/high decomposition for the two microscopic completions
\(R\in\{A,B\}\).  Let \(L_{a,R}=L_{a,R}^*\ge0\) be the squared chiral/Yukawa
operator used in V22.  Write its blocks
\begin{equation}
 L_{a,R}=
 \begin{pmatrix}
 A_{a,R}&B_{a,R}\\
 B_{a,R}^*&H_{a,R}
 \end{pmatrix}
 \label{eq:v23-block-L}
\end{equation}
relative to \eqref{eq:v23-split}.

Let \(A_{0,a}=A_{0,a}^*\) be a common sampled continuum soft operator on
\(\mathcal H_a^{<}\).  Fix a contour \(\Gamma\) around the common continuum
zero-mode cluster and assume
\begin{equation}
 R_{0,a}(z):=(z-A_{0,a})^{-1},
 \qquad
 \sup_{a,z\in\Gamma}\|R_{0,a}(z)\|\le C_0.
 \label{eq:v23-R0}
\end{equation}
For each endpoint define
\begin{align}
 D_{a,R}(z)&:=z-H_{a,R},
 \label{eq:v23-Dhigh}\\
 \Sigma_{a,R}(z)&:=B_{a,R}D_{a,R}(z)^{-1}B_{a,R}^*,
 \label{eq:v23-self-energy}\\
 S_{a,R}(z)&:=z-A_{a,R}-\Sigma_{a,R}(z).
 \label{eq:v23-Schur-soft}
\end{align}

\begin{definition}[Two-zone soft-control parameters]
\label{def:v23-control-params}
Put, uniformly for \(z\in\Gamma\),
\begin{align}
 h_{a,R}&:=\|D_{a,R}(z)^{-1}\|,
 \label{eq:v23-h}\\
 \nu_{a,R}&:=
 \max\bigl\{
 \|B_{a,R}D_{a,R}(z)^{-1}\|,
 \|D_{a,R}(z)^{-1}B_{a,R}^*\|
 \bigr\},
 \label{eq:v23-nu}\\
 e_{a,R}&:=
 \left\|
 \bigl[(A_{a,R}-A_{0,a})+\Sigma_{a,R}(z)\bigr]
 R_{0,a}(z)
 \right\|.
 \label{eq:v23-e}
\end{align}
We call a family \emph{soft Feshbach controlled} if
\begin{equation}
 \boxed{
 \max_{R=A,B}\sup_{z\in\Gamma}
 \bigl(e_{a,R}+\nu_{a,R}+h_{a,R}\bigr)
 \xrightarrow[a\downarrow0]{}0.}
 \label{eq:v23-Feshbach-control}
\end{equation}
\end{definition}

\begin{theorem}[Two-zone norm-resolvent bridge]
\label{thm:v23-Feshbach-resolvent}
Assume \eqref{eq:v23-R0} and
Definition~\ref{def:v23-control-params}.  If
\(e_{a,R}\le1/2\) for small \(a\), then for each endpoint
\begin{equation}
 \boxed{
 \sup_{z\in\Gamma}
 \left\|
 (z-L_{a,R})^{-1}
 -P_aR_{0,a}(z)P_a
 \right\|
 \le
 C_\Gamma
 \bigl(e_{a,R}+\nu_{a,R}+h_{a,R}+\nu_{a,R}^2\bigr),}
 \label{eq:v23-Feshbach-resolvent-bound}
\end{equation}
where \(C_\Gamma\) depends only on the uniform soft resolvent bound.  In
particular, soft Feshbach control implies
\begin{equation}
 \boxed{
 \sup_{z\in\Gamma}
 \|(z-L_{a,A})^{-1}-(z-L_{a,B})^{-1}\|
 \longrightarrow0.}
 \label{eq:v23-endpoint-resolvent}
\end{equation}
\end{theorem}

\begin{proof}
For fixed \(R\), suppress \(a,R,z\).  The block matrix of
\(z-L\) is
\[
 \begin{pmatrix}
 z-A&-B\\
 -B^*&z-H
 \end{pmatrix}.
\]
Since \(D=z-H\) is invertible, the Schur complement on the low block is
\[
 S=z-A-BD^{-1}B^*.
\]
Write
\[
 E:=(A-A_0)+BD^{-1}B^*.
\]
Then \(S=(I-ER_0)(z-A_0)\), so
\begin{equation}
 S^{-1}=R_0(I-ER_0)^{-1}.
 \label{eq:v23-Sinverse}
\end{equation}
The hypothesis \(\|ER_0\|=e\le1/2\) gives
\begin{align}
 \|S^{-1}\|&\le2C_0,
 \label{eq:v23-Sinverse-bound}\\
 \|S^{-1}-R_0\|&\le2C_0e.
 \label{eq:v23-S-R0}
\end{align}

The exact block inverse is
\begin{equation}
 (z-L)^{-1}=
 \begin{pmatrix}
 S^{-1}&S^{-1}BD^{-1}\\
 D^{-1}B^*S^{-1}&
 D^{-1}+D^{-1}B^*S^{-1}BD^{-1}
 \end{pmatrix}.
 \label{eq:v23-block-resolvent}
\end{equation}
Subtract \(\operatorname{diag}(R_0,0)\).  Equations
\eqref{eq:v23-Sinverse-bound}--\eqref{eq:v23-S-R0} and the definitions of
\(h,\nu\) give respectively the block bounds
\[
 2C_0e,
 \qquad
 2C_0\nu,
 \qquad
 2C_0\nu,
 \qquad
 h+2C_0\nu^2.
\]
The norm of a \(2\times2\) operator block matrix is bounded by the sum of its
block norms, proving \eqref{eq:v23-Feshbach-resolvent-bound}.  Apply the same
bound to \(A\) and \(B\) and use the triangle inequality to obtain
\eqref{eq:v23-endpoint-resolvent}.
\end{proof}

\begin{corollary}[Soft Riesz projectors converge without a common ultraviolet parametrix]
\label{cor:v23-Riesz-close}
Let
\begin{equation}
 p_{a,R}:={1\over2\pi i}\oint_\Gamma(z-L_{a,R})^{-1}\,dz,
 \qquad
 p_{0,a}:={1\over2\pi i}\oint_\Gamma R_{0,a}(z)\,dz.
 \label{eq:v23-Riesz}
\end{equation}
Under Theorem~\ref{thm:v23-Feshbach-resolvent},
\begin{equation}
 \boxed{
 \|p_{a,R}-P_ap_{0,a}P_a\|
 \le {\operatorname{len}(\Gamma)\over2\pi}
 C_\Gamma
 \sup_{z\in\Gamma}
 \bigl(e_{a,R}+\nu_{a,R}+h_{a,R}+\nu_{a,R}^2\bigr).}
 \label{eq:v23-Riesz-bound}
\end{equation}
Hence
\begin{equation}
 \boxed{\|p_{a,A}-p_{a,B}\|\longrightarrow0.}
 \label{eq:v23-endpoint-projectors-close}
\end{equation}
For sufficiently small \(a\), V14's direct rotation therefore identifies the
exact controlled zero-mode modules of the two endpoints.
\end{corollary}

\begin{proof}
Integrate \eqref{eq:v23-Feshbach-resolvent-bound} around \(\Gamma\).  The
embedded continuum low resolvent has Riesz projection
\(P_ap_{0,a}P_a\).  The endpoint estimate follows by the triangle inequality.
Once the norm distance is less than one, apply
Lemma~\ref{lem:v14-direct-rotation} and V22's controlled-family-index
consequences.
\end{proof}

\subsection{A useful sufficient form of the three Feshbach inputs}
\label{sec:v23-sufficient-inputs}

The quantities in Definition~\ref{def:v23-control-params} are naturally
produced by an intermediate scale \(M_a\to\infty\) with
\(aM_a\to0\).

\begin{proposition}[Intermediate-scale sufficient conditions]
\label{prop:v23-intermediate-sufficient}
Assume there are numbers \(M_a\to\infty\), \(aM_a\to0\), and errors
\(\epsilon_a,\delta_a\to0\) such that, uniformly for
\(R\in\{A,B\}\) and \(z\in\Gamma\),
\begin{align}
 \|D_{a,R}(z)^{-1}\|&\le {C\over M_a^2},
 \label{eq:v23-high-coercive}\\
 \max\{\|B_{a,R}D_{a,R}(z)^{-1}\|,
       \|D_{a,R}(z)^{-1}B_{a,R}^*\|\}
 &\le\delta_a,
 \label{eq:v23-mixing-small}\\
 \|(A_{a,R}-A_{0,a})R_{0,a}(z)\|
 &\le\epsilon_a,
 \label{eq:v23-low-consistency}\\
 \|B_{a,R}D_{a,R}(z)^{-1}B_{a,R}^*R_{0,a}(z)\|
 &\le\epsilon_a.
 \label{eq:v23-selfenergy-small}
\end{align}
Then
\begin{equation}
 e_{a,R}\le2\epsilon_a,
 \qquad
 \nu_{a,R}\le\delta_a,
 \qquad
 h_{a,R}\le CM_a^{-2},
 \label{eq:v23-control-from-four}
\end{equation}
and Theorem~\ref{thm:v23-Feshbach-resolvent} applies whenever
\(2\epsilon_a<1/2\).
\end{proposition}

\begin{proof}
The first estimate follows from the triangle inequality and
\eqref{eq:v23-low-consistency}--\eqref{eq:v23-selfenergy-small}.  The other two
are exactly \eqref{eq:v23-high-coercive}--\eqref{eq:v23-mixing-small}.
\end{proof}

The four estimates have direct interpretations:
\begin{equation}
 \boxed{
 \begin{array}{ll}
 \eqref{eq:v23-high-coercive}:&\text{native high-frequency discrete ellipticity,}\\
 \eqref{eq:v23-mixing-small}:&\text{soft/high transfer is small after one hard resolvent,}\\
 \eqref{eq:v23-low-consistency}:&\text{the low native operator matches the common continuum operator,}\\
 \eqref{eq:v23-selfenergy-small}:&\text{the high self-energy shifts the soft contour only by }o(1).
 \end{array}}
 \label{eq:v23-meaning-four}
\end{equation}
This is the corrected native soft-parametrix programme.

\subsection{Complete verification on the translation-invariant V2 reference branch}
\label{sec:v23-flat-verification}

We now verify the mechanism using only theorems already proved in V2.  Work on
one fixed physical periodic box, so the continuum momenta form a discrete set.
Let \(Y\) be constant and bounded, and let
\begin{equation}
 D_a(p):=\widetilde d_{\chi,a}(p),
 \qquad
 D_0(p):=d_{\chi,0}(p),
 \qquad
 L_a(p):=D_a(p)^*D_a(p),
 \qquad
 L_0(p):=D_0(p)^*D_0(p).
 \label{eq:v23-flat-DL}
\end{equation}
Choose a contour \(\Gamma\) around the zero-mode cluster of \(L_0\), below its
first positive eigenvalue on the fixed box.

For an intermediate momentum \(M_a\) satisfying
\begin{equation}
 1\ll M_a\ll a^{-1},
 \label{eq:v23-M-window}
\end{equation}
let \(P_a\) be the Fourier projector onto \(|p|\le M_a\) and
\(Q_a=I-P_a\).  Translation invariance makes
\begin{equation}
 \boxed{B_a=P_aL_aQ_a=0.}
 \label{eq:v23-flat-no-mixing}
\end{equation}
Hence the self-energy and transfer terms vanish exactly.

\begin{lemma}[Low-block relative consistency on the V2 reference branch]
\label{lem:v23-flat-low-consistency}
There is a constant \(C\), independent of \(a\) and \(M_a\), such that for
small \(aM_a\),
\begin{equation}
 \boxed{
 \sup_{z\in\Gamma}
 \|(P_aL_aP_a-P_aL_0P_a)
   (z-P_aL_0P_a)^{-1}\|
 \le C aM_a.}
 \label{eq:v23-flat-low-bound}
\end{equation}
\end{lemma}

\begin{proof}
On \(|p|\le M_a\), Theorem~\ref{thm:v2-native-chiral-annular} and the
small-momentum Taylor expansion at the finitely many momenta below the first
hard annulus give
\begin{equation}
 \|D_a(p)-D_0(p)\|
 \le Ca(1+|p|)^2.
 \label{eq:v23-D-diff}
\end{equation}
Also \(\|D_0(p)\|+\|D_a(p)\|\le C(1+|p|)\).  Therefore
\begin{equation}
 \|L_a(p)-L_0(p)\|
 \le Ca(1+|p|)^3.
 \label{eq:v23-L-diff}
\end{equation}
Because \(\Gamma\) lies in a fixed spectral gap of the continuum squared
operator,
\begin{equation}
 \|(z-L_0(p))^{-1}\|
 \le {C_\Gamma\over1+|p|^2}
 \label{eq:v23-cont-res-decay}
\end{equation}
for all momenta in the fixed physical box.  Multiplying
\eqref{eq:v23-L-diff} and \eqref{eq:v23-cont-res-decay} gives
\[
 Ca{(1+|p|)^3\over1+|p|^2}
 \le C'a(1+M_a).
\]
Since \(M_a\to\infty\), absorb the harmless constant term into
\(CaM_a\).  Taking the supremum over the diagonal Fourier blocks proves
\eqref{eq:v23-flat-low-bound}.
\end{proof}

\begin{lemma}[High-block coercivity on the V2 reference branch]
\label{lem:v23-flat-high-coercive}
For all sufficiently large \(M_a\) and sufficiently small \(a\),
\begin{equation}
 \boxed{
 \sup_{z\in\Gamma}
 \|(z-Q_aL_aQ_a)^{-1}\|
 \le {C\over M_a^2}.}
 \label{eq:v23-flat-high-bound}
\end{equation}
\end{lemma}

\begin{proof}
Split \(|p|>M_a\) into the physical region \(|ap|\le\delta\) of
Theorem~\ref{thm:v2-physical-expansion} and the complementary full zone.
In the physical region, the kinetic chiral singular value is at least
\(c|p|\).  The bounded Yukawa block changes this by at most \(M_Y\), so for
\(M_a\ge2M_Y/c\),
\begin{equation}
 s_{\min}(D_a(p))\ge c'|p|\ge c'M_a.
 \label{eq:v23-high-physical-D}
\end{equation}
Hence \(L_a(p)\succeq c'^2M_a^2I\).  On the full-zone region,
Theorem~\ref{thm:v2-full-zone-pivot} gives
\(s_{\min}(D_a(p))\ge c_{\rm UV}/a-M_Y\ge c''/a\), which is still larger
than a fixed multiple of \(M_a\) because \(aM_a\to0\).  Thus the whole high
block has spectrum above \(c_*M_a^2\).  Since \(\Gamma\) is bounded, its
resolvent has norm at most \(C/M_a^2\).
\end{proof}

\begin{theorem}[Soft norm-resolvent convergence on the V2 reference branch]
\label{thm:v23-flat-soft-resolvent}
Choose
\begin{equation}
 \boxed{M_a=a^{-1/3}.}
 \label{eq:v23-optimal-M}
\end{equation}
Then
\begin{equation}
 \boxed{
 \sup_{z\in\Gamma}
 \left\|
 (z-L_a)^{-1}-P_a(z-P_aL_0P_a)^{-1}P_a
 \right\|
 \le C a^{2/3}.}
 \label{eq:v23-flat-resolvent-rate}
\end{equation}
Consequently the native soft Riesz projector converges in operator norm to the
embedded continuum soft projector at rate \(O(a^{2/3})\).
\end{theorem}

\begin{proof}
By \eqref{eq:v23-flat-no-mixing}, \(\nu_a=0\) and the self-energy vanishes.
Lemma~\ref{lem:v23-flat-low-consistency} gives
\(e_a\le CaM_a\), while Lemma~\ref{lem:v23-flat-high-coercive} gives
\(h_a\le CM_a^{-2}\).  Theorem~\ref{thm:v23-Feshbach-resolvent} therefore
bounds the resolvent error by
\begin{equation}
 C(aM_a+M_a^{-2}).
 \label{eq:v23-balance}
\end{equation}
The choice \(M_a=a^{-1/3}\) balances the two terms at \(a^{2/3}\).  Integrate
around \(\Gamma\) for the Riesz projector statement.
\end{proof}

\begin{remark}[Why the rate is not the point]
The exponent \(2/3\) is only the rate obtained from the first V2 consistency
remainder and the elementary two-zone balance.  Higher-order native consistency
could improve it.  V22 only needs convergence to zero, not an optimal rate.
\end{remark}

\subsection{Fixed source jets: the same two-zone bookkeeping}
\label{sec:v23-source-jets}

For a fixed source multi-index \(I\), differentiating a resolvent produces the
finite ordered-partition words already used in V2:
\begin{equation}
 \partial_I(z-L)^{-1}
 =\sum_{I=B_1\sqcup\cdots\sqcup B_m}^{\rm ord}
 (z-L)^{-1}L_{B_1}(z-L)^{-1}\cdots
 L_{B_m}(z-L)^{-1},
 \label{eq:v23-resolvent-source-words}
\end{equation}
with the standard signs absorbed into the ordered-word convention.  The
load-bearing fact is that every differentiated word contains one more
resolvent factor than differentiated operator vertices.

\begin{proposition}[Source-word two-zone reduction]
\label{prop:v23-source-word-reduction}
Fix the finite source/BV packet \(\mathfrak q\).  Suppose that, after the
engineering stocks of V2 are divided out, the following hold uniformly for
every source word appearing in \eqref{eq:v23-resolvent-source-words}:
\begin{enumerate}[label=\textup{(S\arabic*)},leftmargin=3em]
\item on the low block, replacing every native source vertex by the common
continuum source vertex changes the complete low resolvent word by at most
\(C_I\epsilon_a\);
\item every word containing at least one high resolvent but no low/high transfer
has norm at most \(C_I h_a\);
\item every word crossing from the low block to the high block or back carries
at least one transfer factor bounded by \(\delta_a\), while every closed high
excursion contributing to the low block is bounded by \(C_I\epsilon_a\).
\end{enumerate}
If \(\epsilon_a+\delta_a+h_a\to0\), then every fixed source derivative of the
native resolvent converges uniformly on \(\Gamma\) to the corresponding common
soft continuum derivative, and the same is true of the Riesz projector.
\end{proposition}

\begin{proof}
Differentiate the exact block inverse
\eqref{eq:v23-block-resolvent}, or equivalently expand
\eqref{eq:v23-resolvent-source-words} and insert \(P_a+Q_a=I\) between every
factor.  At fixed source order there are only finitely many resulting block
words.  A word staying entirely low is controlled by (S1).  A word staying
high contains at least one high resolvent and is controlled by (S2).  Every
mixed word contains either an open low/high transfer, giving (S3)'s
\(\delta_a\), or a closed high excursion, giving its \(\epsilon_a\).  Finite
summation preserves convergence to zero.  Contour integration gives the Riesz
projector derivatives.
\end{proof}

On the translation-invariant V2 reference branch, V2's fixed source-vertex
transfer theorem supplies (S1) with \(\epsilon_a=O(aM_a)\), while the hard
inverse/source-word estimates give (S2).  A fixed external source momentum can
move a Fourier label by only a bounded amount; choosing the low/high split with
a fixed source-order buffer around \(M_a\) turns all remaining crossings into
hard-annulus words and gives (S3) with the same vanishing scale.  Thus the
reference-branch theorem extends coefficientwise to every fixed V2 source
packet.  This statement is still a reference-branch theorem; it does not
replace the nonconstant-background estimate below.

\subsection{The exact nonconstant-background native gate}
\label{sec:v23-nonconstant-gate}

For a genuinely nonconstant protected fixed-index background family, the
correct next theorem is no longer ambiguous.  One must construct a common
intermediate-scale projector \(P_a\) (for example from a background-covariant
reference Laplacian or another regulator-independent elliptic soft operator)
and verify the four estimates of
Proposition~\ref{prop:v23-intermediate-sufficient} with all prescribed source
jets.

\begin{assumption}[Native covariant soft--hard package]
\label{ass:v23-covariant-package}
For one compact protected fixed-index gauge/spin/coframe/Higgs family, suppose
there are \(M_a\to\infty\), \(aM_a\to0\), a common low projector \(P_a\), and
a common sampled continuum soft operator \(A_{0,a}\) such that for each
admissible microscopic completion:
\begin{enumerate}[label=\textup{(G\arabic*)},leftmargin=3em]
\item \textbf{low consistency:}
\((A_{a,R}-A_{0,a})R_{0,a}=o(1)\) coefficientwise;
\item \textbf{discrete high-frequency G\aa rding:}
\(Q_aL_{a,R}Q_a\succeq cM_a^2Q_a\) modulo a bounded contour margin;
\item \textbf{low/high transfer:}
\(B_{a,R}(z-H_{a,R})^{-1}=o(1)\) and its adjoint analogue;
\item \textbf{high self-energy:}
\(B_{a,R}(z-H_{a,R})^{-1}B_{a,R}^*R_{0,a}=o(1)\);
\item every statement remains true for the finite source-resolvent words
required by \(\mathfrak q\).
\end{enumerate}
\end{assumption}

\begin{theorem}[Covariant soft--hard package closes the last complex anchoring gap]
\label{thm:v23-covariant-package-closes}
Under Assumption~\ref{ass:v23-covariant-package}, the two native completions have
one common soft Riesz projector in the V22 sense:
\begin{equation}
 \|p_{a,A}-p_{a,B}\|\longrightarrow0
 \label{eq:v23-covariant-projector-close}
\end{equation}
with every retained source derivative.  Consequently their exact V20
controlled family indices agree for sufficiently small \(a\), and V19 gives
one global unstabilized coefficientwise native regulator path on the faithful
complex branch.
\end{theorem}

\begin{proof}
Items (G1)--(G4) are exactly
Proposition~\ref{prop:v23-intermediate-sufficient}.  Apply
Theorem~\ref{thm:v23-Feshbach-resolvent} and
Corollary~\ref{cor:v23-Riesz-close}.  Item (G5) and
Proposition~\ref{prop:v23-source-word-reduction} give the source derivatives.
The V22 direct-rotation consequence identifies the controlled module, and V19
removes the downstream unstable cancellation classes.
\end{proof}

\subsection{What can already be proved natively, and what cannot}
\label{sec:v23-boundary}

The V2 stencil and V1 functional calculus already supply the following pieces
of Assumption~\ref{ass:v23-covariant-package} on any chart where their stated
hypotheses hold:
\begin{enumerate}[leftmargin=2.4em]
\item finite-range Wilson locality and exact gauge/spin covariance;
\item the cutoff-uniform Wilson sign gap needed to define the overlap operator;
\item local covariant Taylor consistency of every fixed source vertex;
\item physical-branch ellipticity on frozen local coefficients and the
full-zone no-doubler pivot;
\item quasilocality of the sign, overlap projectors and fixed source
functional-calculus words.
\end{enumerate}
What has \emph{not} yet been proved in this lineage is the global estimate
which assembles these local facts into (G2)--(G4) on a nonconstant topological
background: a covariant discrete G\aa rding inequality with an intermediate
scale, together with the associated commutator/IMS estimate showing that
freezing coefficients and patching charts produces only \(o(1)\) soft/high
mixing after one high resolvent.

\begin{firewall}[The remaining gap is not a symbol Taylor coefficient]
The missing theorem is an operator-level elliptic localization estimate on a
nonconstant fixed-index background.  Recomputing more flat Wilson Taylor
coefficients cannot prove it.  Conversely, once the discrete G\aa rding/mixing
estimate is available, no new topological or determinant-line input is needed:
V23's Feshbach theorem converts it directly into the common soft Riesz
projector required by V22.
\end{firewall}

\subsection{V23 anti-circularity audit}
\label{sec:v23-audit}

\begin{enumerate}[leftmargin=2.4em]
\item \textbf{The ultraviolet lattice dispersion is not forced to equal the
continuum dispersion.}  Proposition~\ref{prop:v23-pure-parametrix-fails}
shows explicitly why the full-lattice pure continuum parametrix is the wrong
variable.
\item \textbf{The high block is allowed to remain regulator native.}  Only its
resolvent norm and its transfer back into the soft block must vanish.
\item \textbf{Soft convergence is proved by an exact Schur/Feshbach identity.}
No spectral-projector convergence is assumed in the theorem.
\item \textbf{The V2 branch is used only at its proven scope.}  It verifies the
whole mechanism on the translation-invariant reference branch, including the
high-zone Wilson pivot; it is not silently promoted to a nonconstant
fixed-index family theorem.
\item \textbf{Source derivatives remain coefficientwise.}  The source-word
reduction uses only the finite ordered-resolvent bank already present in V2 and
V5.
\item \textbf{The remaining input is operator level.}  The next theorem is a
covariant discrete G\aa rding/IMS and mixing estimate, not another anomaly,
line, or flat-symbol calculation.
\end{enumerate}

\section{V23 status ledger}
\label{sec:v23-status}

\subsection*{Proved in V23}
\begin{enumerate}[leftmargin=2.2em]
\item The final V22 full-lattice pure continuum parametrix is unnecessarily
strong and generally false for the canonical physical identification.  At the
V2 Brillouin corner, the native squared overlap eigenvalue is \(4a^{-2}\)
while the continuum value is \(\pi^2a^{-2}\), leaving an \(O(1)\) parametrix
residual.
\item A two-zone Feshbach theorem gives an exact sufficient criterion for soft
norm-resolvent convergence: low-block relative consistency, high-block
coercivity, small resolvent-weighted low/high transfer and small high self-energy.
The high native inverse may be regulator specific.
\item Under these estimates the full native resolvent converges on the soft
contour to the embedded common low continuum resolvent, so the two endpoint
Riesz projectors converge to one another in operator norm and therefore have
the same exact controlled family index for small cutoff spacing.
\item On the translation-invariant V2 reference branch, the low/high mixing
vanishes exactly.  V2's physical expansion gives relative low error
\(O(aM_a)\), while the physical/full-zone pivots give high resolvent
\(O(M_a^{-2})\).  Choosing \(M_a=a^{-1/3}\) yields the explicit soft
norm-resolvent rate \(O(a^{2/3})\).
\item The same two-zone mechanism extends to every fixed source derivative once
the finite source-resolvent words satisfy low consistency, high suppression and
resolvent-weighted transfer bounds.  V2's source calculus verifies this
bookkeeping on the translation-invariant reference branch.
\item For a nonconstant protected fixed-index background, the exact remaining
native gate is the covariant soft--hard package of
Assumption~\ref{ass:v23-covariant-package}.  Once that package is proved, V22
and V19 close the entire faithful complex microscopic anchoring problem.
\end{enumerate}

\subsection*{Inherited}
\begin{enumerate}[leftmargin=2.2em]
\item V1--V7: the source-complete coefficientwise ESM continuum; common BV
restoration; fixed-loop forests; the ultraviolet Cauchy theorem; and exact
scheme/reference and perturbative-order compatibility.
\item V8--V19: regulator comparison, fixed-index zero-mode shell calculus,
controlled module descent, and automatic downstream cancellation after primary
zero-mode matching.
\item V20--V21: exact controlled family-index classification and its spatial
anchor/index-density interpretation.
\item V22: cross-residual and common soft-Riesz criteria showing that norm-close
soft projectors identify the exact controlled family index and collapse the
complex microscopic obstruction tower.
\end{enumerate}

\subsection*{Conditional}
\begin{enumerate}[leftmargin=2.2em]
\item The V23 nonconstant-background conclusion is conditional on the native
covariant high-frequency G\aa rding/mixing package.  That package is not yet a
theorem of the supplied notes.
\item The explicit \(a^{2/3}\) rate is proved only on the
translation-invariant V2 reference branch.  It is not claimed to be optimal or
valid on arbitrary topological backgrounds.
\item The common continuum soft operator must have constant kernel rank and a
positive reduced gap on the protected compact panel.
\item Source statements are coefficientwise at the prescribed finite packet;
no source-radius or uniform source-order theorem is asserted.
\item The independent real Pfaffian mod-two orientation/zero-crossing problem
remains separate.
\end{enumerate}

\subsection*{Open beyond V23}
\begin{enumerate}[leftmargin=2.2em]
\item Prove the native covariant discrete G\aa rding estimate on a smooth
nonconstant fixed-index gauge/spin/coframe/Higgs family, uniformly in the
cutoff, with an intermediate scale \(M_a\to\infty\), \(aM_a\to0\).
\item Prove the associated IMS/commutator and source-word estimates showing
that low/high transfer after one hard resolvent and the high-block self-energy
are \(o(1)\).
\item Combine those estimates with the local covariant Taylor consistency of
the explicit V2 Wilson/overlap/Yukawa stencil to verify
Assumption~\ref{ass:v23-covariant-package} and thereby close the last faithful
complex microscopic regulator-phase ambiguity.
\item Compute the independent real Pfaffian mod-two spectral-flow character and
analyze genuine reduced-gap/index-changing divisors.
\item The positive-IR removal, analytic finite-coupling RG, perturbative
summability and nonperturbative directions remain separate stronger problems.
\end{enumerate}

\subsection*{Exact next load-bearing theorem}
\begin{center}
\fbox{\parbox{0.92\textwidth}{
\textbf{Covariant discrete G\aa rding / Feshbach mixing theorem.}
For the explicit V2 Wilson/overlap/Yukawa regulator on a compact smooth
protected fixed-index gauge--spin--coframe--Higgs family, construct a common
background-covariant low projector \(P_a\) at an intermediate scale
\(M_a\to\infty\), \(aM_a\to0\).  Prove, with every retained source derivative,
\[
 \begin{aligned}
 Q_aL_{a,R}Q_a&\succeq cM_a^2Q_a,\\
 \|B_{a,R}(z-H_{a,R})^{-1}\|&=o(1),\\
 \|B_{a,R}(z-H_{a,R})^{-1}B_{a,R}^*R_{0,a}\|&=o(1).
 \end{aligned}
\]
and the corresponding low-block covariant consistency estimate.
The proof should use local freezing of the native Wilson symbol, an
IMS/partition argument and the existing V2 all-zone gap, while keeping the
actual gauge/spin parallel transport and Yukawa/source jets.  Combined with
Theorem~\ref{thm:v23-covariant-package-closes}, this is sufficient for the
common soft Riesz projector and therefore for equality of the exact controlled
family index of all admissible microscopic completions on the protected
faithful complex branch.}}
\end{center}

\section*{Append-only continuation marker: V23}
\addcontentsline{toc}{section}{Append-only continuation marker: V23}
\textbf{End of V23.}  The complete V1--V22 body above is retained verbatim.
V23 goes upstream from the overstrong full-lattice pure continuum parametrix
and replaces it by the correct two-zone Feshbach problem.  A common continuum
comparison is required only on a growing soft window; the regulator-specific
high resolvent is permitted to differ because its norm vanishes.  An exact
block-resolvent theorem turns low consistency, high coercivity, small transfer
and small high self-energy into soft norm-resolvent and Riesz-projector
convergence.  The entire mechanism is verified on the translation-invariant V2
reference branch with rate \(O(a^{2/3})\).  The remaining native theorem is now
operator-level and precise: prove the covariant discrete G\aa rding/IMS and
mixing estimates on a nonconstant protected fixed-index background family.



\clearpage
\section{V24 continuation: covariant discrete G\aa rding, three-zone mixing, and closure of the soft Feshbach package}
\label{sec:v24-continuation}

V23 reduced the last faithful-complex microscopic anchoring problem to one
operator-level estimate on a nonconstant protected background.  The correct
proof is not a sharper flat-symbol calculation.  It has three scales.  A common
background-covariant second-order frequency operator controls the soft and
intermediate modes; the native Wilson/overlap operator supplies full-zone
ellipticity; and repeated commutators with the common frequency operator suppress
the coupling of the soft window directly to cutoff-scale modes.

There is one additional quantifier point.  A sharp spectral projector of the
frequency operator need not depend smoothly on a source parameter when an
auxiliary high eigenvalue crosses the chosen threshold.  No such derivative is
needed.  On every finite protected source chart we freeze the auxiliary
frequency decomposition at the chart centre and differentiate only the native
operator, its source vertices and its resolvents.  The complete resolvent is
independent of this proof decomposition.  A finite cover of the compact
background panel then gives uniform constants.  This is the same
coefficientwise/frozen-reference logic already used for the V2 hard pivot and
for V15's finite source-jet reserve.

The main theorem below proves the V23 covariant soft--hard package on the smooth
sampled V2 background class.  The explicit scale choice
\begin{equation}
 \boxed{
 M_a=a^{-9/20},
 \qquad
 \Lambda_a=a^{-3/5}}
 \label{eq:v24-scales}
\end{equation}
satisfies
\begin{equation}
 1\ll M_a\ll\Lambda_a\ll a^{-1},
 \qquad
 a\Lambda_a\longrightarrow0.
 \label{eq:v24-scale-window}
\end{equation}
The worst error is $O(a^{1/10})$; no optimality is claimed.

\subsection{The common native frequency operator}
\label{sec:v24-frequency}

Fix one compact smooth protected gauge--spin--coframe--Higgs panel.  On a local
spin/gauge trivialization let $\nabla_{\mu,a}$ be the actual nearest-neighbour
covariant differences built from the same sampled parallel transports as V2.
Let $g^{\mu\nu}(x)$ be the positive coframe metric on the compact ellipticity
set.  Choose one smooth positive lattice discretization
\begin{equation}
 K_a
 :=
 \sum_{\mu,\nu}
 \nabla_{\mu,a}^*\,g^{\mu\nu}_a\,\nabla_{\nu,a}
 +I,
 \label{eq:v24-K}
\end{equation}
with the usual symmetric placement of $g^{\mu\nu}_a$ making $K_a=K_a^*$.
It is built only from the common background geometry/parallel transport and is
therefore independent of the irrelevant microscopic completion label
$R\in\{A,B\}$.  Under a change of local gauge/spin trivialization it transforms
by unitary similarity.  Hence all spectral statements below are intrinsic.

On a source chart centred at $\zeta_0$, write $K_a:=K_a(\zeta_0)$ and keep this
operator fixed while taking the finite formal source derivatives at $\zeta_0$.
For the scales \eqref{eq:v24-scales} define the mutually orthogonal spectral
projections
\begin{align}
 P_a&:=\mathbf 1_{[0,M_a^2]}(K_a),
 \label{eq:v24-P}\\
 \Pi_a&:=\mathbf 1_{(M_a^2,\Lambda_a^2]}(K_a),
 \label{eq:v24-Pi}\\
 U_a&:=\mathbf 1_{(\Lambda_a^2,\infty)}(K_a),
 \label{eq:v24-U}\\
 Q_a&:=\Pi_a+U_a=I-P_a.
 \label{eq:v24-Q}
\end{align}
The projector $P_a$ is an auxiliary frequency projector, not the physical
zero-mode Riesz projector $p_{a,R}$.

\begin{lemma}[Source-frozen decomposition is legitimate]
\label{lem:v24-frozen-decomposition}
Fix a source multi-index $I$ and a chart centre $\zeta_0$.  Every coefficient
$\partial_I(z-L_{a,R}(\zeta))^{-1}|_{\zeta_0}$ is computed correctly by
holding \eqref{eq:v24-P}--\eqref{eq:v24-Q} fixed and inserting
$P_a+Q_a=I$ between the differentiated operator vertices.  No derivative of
$P_a$ or $Q_a$ enters.  Consequently a lack of a spectral gap of $K_a$ at the
artificial thresholds $M_a^2,\Lambda_a^2$ creates no source singularity.
\end{lemma}

\begin{proof}
The source derivative of the full resolvent is the finite ordered word
\eqref{eq:v23-resolvent-source-words}.  The identity
$P_a+Q_a=I$ may be inserted algebraically between any two factors after the
derivative has been taken.  Since the auxiliary decomposition is evaluated at
$\zeta_0$ and is not part of the differentiated physical family, no derivative
of it appears.  Summing all block words recovers the exact differentiated full
resolvent.  The argument is coefficientwise and therefore applies independently
at every centre of a finite protected cover.
\end{proof}

\subsection{Native smooth-background regularity inherited from V1--V2}
\label{sec:v24-regularity}

The following package is a repackaging of the explicit V2 stencil and V1
quasilocal functional calculus in the operator orders needed below.

\begin{lemma}[V2 smooth-background operator package]
\label{lem:v24-V2-package}
On every compact smooth protected background panel, and for every fixed source
order used by $\mathfrak q$, the squared fixed-frame chiral/Yukawa operator
\begin{equation}
 L_{a,R}=D_{\chi,a,R}^*D_{\chi,a,R}
 \label{eq:v24-L}
\end{equation}
satisfies the following uniformly in cutoff and volume.
\begin{enumerate}[label=\textup{(N\arabic*)},leftmargin=3em]
\item \textbf{quasilocal order two:} $L_{a,R}$ has an exponentially summable
path/shift expansion of overall scale $a^{-2}$, with every fixed weighted
lattice moment and retained source derivative bounded after multiplication by
$a^2$;
\item \textbf{smooth coefficient variation:} translating any coefficient/path
amplitude by one lattice edge changes it by $O(a)$ after its engineering stock
is removed; the commutator of two elementary parallel shifts around one
plaquette differs from the reversed product by $O(a^2)$;
\item \textbf{frozen full-zone ellipticity:} after freezing the smooth
background coefficients at one physical point, the squared chiral symbol
$\ell_{a,R}^{\rm fr}(x,\theta)$ obeys
\begin{equation}
 \ell_{a,R}^{\rm fr}(x,\theta)
 \succeq c_* k_a^{\rm fr}(x,\theta)-C_*I
 \label{eq:v24-frozen-elliptic}
\end{equation}
against the frozen symbol of $K_a$, with $c_*,C_*>0$ uniform on the compact
panel;
\item \textbf{low/intermediate order decomposition:} on a spectral window
$K_a\le\Lambda^2$ with $1\ll\Lambda\ll a^{-1}$,
\begin{equation}
 L_{a,R}-K_a
 =T_{1,a,R}+T_{0,a,R}+E_{a,R}
 \label{eq:v24-order-decomp}
\end{equation}
in the fixed local Kato frame, where $T_1$ has native differential order one,
$T_0$ order zero, and
\begin{align}
 \|T_{1,a,R}\,\mathbf 1_{K_a\le\Lambda^2}\|
 &\le C(1+\Lambda),
 \label{eq:v24-T1}\\
 \|T_{0,a,R}\|&\le C,
 \label{eq:v24-T0}\\
 \|E_{a,R}\,\mathbf 1_{K_a\le\Lambda^2}\|
 &\le Ca(1+\Lambda)^3.
 \label{eq:v24-E3}
\end{align}
The same decomposition and estimates hold for every fixed source derivative,
with constants depending on the derivative order but not on $a$.
\item \textbf{common soft continuum jet:} there is one sampled continuum
Dirac--Yukawa square $L_{0,a}$, independent of the microscopic completion
label, such that on $K_a\le\Lambda^2$
\begin{equation}
 L_{a,R}-L_{0,a}
 =aS_{3,a,R}+aS_{2,a,R}+aS_{1,a,R}+aS_{0,a,R},
 \label{eq:v24-common-soft-jet}
\end{equation}
where $S_m$ has native differential order at most $m$ and
\begin{equation}
 \|S_{m,a,R}\mathbf 1_{K_a\le\Lambda^2}\|
 \le C_m(1+\Lambda)^m.
 \label{eq:v24-common-soft-orders}
\end{equation}
The same common continuum vertices and $a$-suppressed remainder hold for every
retained fixed source derivative.  In particular, an $O(1)$ finite-rank soft
defect such as the V20 far-site example is not in this V2-consistent class.
\end{enumerate}
\end{lemma}

\begin{proof}
(N1) is V1's weighted functional calculus applied to the gapped finite-range V2
Wilson kernel, followed by multiplication with the bounded onsite Yukawa block.
The sign, overlap projector, Kato frame and every fixed source derivative have
uniform exponential lattice moments.  The overlap operator has engineering
order one, so its square has overall scale $a^{-2}$.

For (N2), the primitive shift coefficients are sampled smooth coframe fields and
edge parallel transports.  Lemma~\ref{lem:v2-link-jets} gives the one-edge
$O(a)$ variation of the sampled coefficients.  The product of two parallel
shifts around the two orders of an elementary plaquette differs by its
curvature holonomy, $I+O(a^2)$ on a compact smooth connection panel.  The
exponential path tails from the overlap sign are summable with arbitrary fixed
polynomial path weights, so the same estimates survive the V1 functional
calculus.

For (N3), at frozen coefficients Theorem~\ref{thm:v2-physical-expansion} gives
$s_{\min}(D_{\chi,a,R})\ge c|p|-C$ on the physical branch, while
Theorem~\ref{thm:v2-full-zone-pivot} gives the $c/a$ pivot on the complementary
zone.  The frozen symbol $k_a^{\rm fr}$ is comparable to $|p|^2$ for
$|ap|\ll1$ and to $a^{-2}$ away from the physical origin.  Compactness of the
remaining annulus therefore gives \eqref{eq:v24-frozen-elliptic}; the bounded
Yukawa block contributes only the constant $C_*$.  Uniform coframe ellipticity
makes the constants independent of the frozen point.

For (N4), the continuum square is a Laplace-type operator with the same
second-order principal metric as $K_a$; the difference is first/zeroth order
(connection, spin, curvature, Yukawa and lower geometric terms).  Theorem
\ref{thm:v2-source-vertex-transfer} gives one extra factor $a|p|$ in the
native-to-continuum remainder of every fixed vertex.  Squaring the first-order
operator therefore gives an $a$-weighted differential-order-three remainder,
plus lower-order $a$ corrections.  On a $K_a$-spectral window the discrete
covariant Bernstein estimate proved in Lemma~\ref{lem:v24-Bernstein} below
turns orders $1,0,3$ into the bounds
\eqref{eq:v24-T1}--\eqref{eq:v24-E3}.  The source-derivative statement follows
from the same finite V2 source-word calculus.

Finally, (N5) is the operator form of the actual V2 finite-to-ordinary
comparison: Theorem~\ref{thm:v2-physical-expansion} and
Theorem~\ref{thm:v2-source-vertex-transfer} give one extra factor $a$ and one
extra differential order relative to the continuum Dirac--Yukawa jet.  Squaring
the first-order operator gives the order-three term in
\eqref{eq:v24-common-soft-jet}; all remaining products are of order at most two
and carry the same explicit factor $a$ (or a stronger $a^2$ factor).  Lemma
\ref{lem:v24-Bernstein} turns this into
\eqref{eq:v24-common-soft-orders}.  The continuum jet is common because the two
microscopic completions are compared at fixed physical gauge/spin/coframe/Higgs
background and fixed renormalized Dirac--Yukawa data.
\end{proof}

The proof just invoked the standard spectral regularity estimate; we record it
explicitly so no smoothness is hidden.

\begin{lemma}[Discrete covariant Bernstein estimate]
\label{lem:v24-Bernstein}
Let $\mathcal D_{m,a}$ be any fixed finite linear combination of products of at
most $m$ native covariant first differences, with smooth bounded coefficients
and fixed source derivatives.  For every fixed $m$ there are constants $C_m$
such that, for $1\le\Lambda\le c/a$,
\begin{equation}
 \boxed{
 \|\mathcal D_{m,a}\mathbf 1_{K_a\le\Lambda^2}\|
 \le C_m(1+\Lambda)^m.}
 \label{eq:v24-Bernstein}
\end{equation}
\end{lemma}

\begin{proof}
For $m=1$, the quadratic form definition \eqref{eq:v24-K} and compact coframe
ellipticity give
\(
 \sum_\mu\|\nabla_{\mu,a}\psi\|^2
 \le C\langle\psi,K_a\psi\rangle.
\)
Thus the claim follows on the spectral subspace $K_a\le\Lambda^2$.

Induct on $m$.  Commute one covariant difference through the positive operator
$K_a$.  The commutator is a finite sum of differential words of order at most
two with coefficients containing one smooth background derivative or one
curvature insertion; by (N2) those coefficients are uniformly bounded in
physical units.  Repeatedly applying the $m=1$ form bound and Cauchy--Schwarz
gives the graph-norm estimate
\begin{equation}
 \sum_{|J|\le m}\|\nabla_J\psi\|^2
 \le C_m\sum_{r=0}^m\|K_a^{r/2}\psi\|^2.
 \label{eq:v24-discrete-elliptic-regularity}
\end{equation}
The same induction is the usual elliptic-regularity proof, now finite because
$m$ is fixed.  On $\mathbf 1_{K_a\le\Lambda^2}$ the right side is at most
$C_m(1+\Lambda)^{2m}\|\psi\|^2$, proving
\eqref{eq:v24-Bernstein}.  Fixed source derivatives only change the bounded
coefficient family.
\end{proof}

\subsection{Discrete IMS and global G\aa rding}
\label{sec:v24-Garding}

\begin{lemma}[Exact discrete IMS identity]
\label{lem:v24-IMS}
Let $L=L^*$ be any lattice kernel and let real multiplication operators
$\chi_j$ obey $\sum_j\chi_j^2=I$.  Then
\begin{equation}
 \boxed{
 \sum_j\chi_jL\chi_j
 =L-\frac12\sum_j[\chi_j,[\chi_j,L]].}
 \label{eq:v24-IMS}
\end{equation}
If $L=L_{a,R}$ satisfies (N1) and the $\chi_j$ form a uniformly finite-overlap
physical partition with
\begin{equation}
 |\chi_j(x)-\chi_j(y)|
 \le C\min\left\{1,{a\,d_a(x,y)\over r}\right\},
 \label{eq:v24-chi-Lipschitz}
\end{equation}
then
\begin{equation}
 \left\|\sum_j[\chi_j,[\chi_j,L_{a,R}]]\right\|
 \le {C\over r^2}
 \label{eq:v24-IMS-error}
\end{equation}
uniformly in $a$.
\end{lemma}

\begin{proof}
Expand the double commutator:
\[
 [\chi,[\chi,L]]=\chi^2L-2\chi L\chi+L\chi^2.
\]
Summing and using $\sum\chi_j^2=I$ gives
\eqref{eq:v24-IMS}.  The kernel of one double commutator contains the factor
$(\chi_j(x)-\chi_j(y))^2L(x,y)$.  By
\eqref{eq:v24-chi-Lipschitz}, this contributes
$C a^2d_a(x,y)^2/r^2$.  The $a^{-2}$ engineering size of $L$ cancels the
$a^2$, and the second weighted lattice moment is finite by (N1).  Finite
partition overlap proves \eqref{eq:v24-IMS-error} by the Schur test.
\end{proof}

\begin{lemma}[Local freezing form estimate]
\label{lem:v24-freezing}
For every sufficiently small fixed physical patch radius $r>0$, choose local
gauge/spin trivializations and freeze the smooth background coefficients at one
point of each patch.  Then
\begin{equation}
 \left|
 \langle\chi_j\psi,
 (L_{a,R}-L_{a,R}^{(j),\rm fr})\chi_j\psi\rangle
 \right|
 \le
 Cr\,\langle\chi_j\psi,K_a\chi_j\psi\rangle
 +C_r\|\chi_j\psi\|^2,
 \label{eq:v24-freezing-form}
\end{equation}
and the same estimate holds with $K_a$ replaced by its frozen local form up to
another $Cr$ relative error.  Constants are uniform in $a$ and in the compact
background panel.
\end{lemma}

\begin{proof}
In a local trivialization, write the quasilocal shift expansion of (N1).
Changing a coefficient from $x$ to the patch centre changes its smooth part by
$O(r)$; replacing a short parallel path by the frozen path produces the same
$O(r)$ coefficient error plus curvature terms bounded by the physical patch
area.  Pair every shift with its adjoint.  The second-order part is therefore a
relative $O(r)$ perturbation of the positive covariant-difference form defining
$K_a$.  First- and zeroth-order pieces are bounded by
$\varepsilon\langle\psi,K_a\psi\rangle+C_\varepsilon\|\psi\|^2$.
Exponential path tails are summed using their fixed moments.  This gives
\eqref{eq:v24-freezing-form}.  The same argument applies to $K_a$ itself.
\end{proof}

\begin{theorem}[Covariant discrete G\aa rding inequality]
\label{thm:v24-Garding}
There exist constants $c_G>0$ and $C_G<\infty$, independent of cutoff, volume,
background point and microscopic completion in the declared smooth protected
class, such that
\begin{equation}
 \boxed{
 L_{a,R}\succeq c_GK_a-C_GI.}
 \label{eq:v24-Garding}
\end{equation}
Consequently, with $Q_a$ from \eqref{eq:v24-Q},
\begin{equation}
 Q_aL_{a,R}Q_a
 \succeq(c_GM_a^2-C_G)Q_a,
 \label{eq:v24-high-form}
\end{equation}
and uniformly on the fixed contour $\Gamma$,
\begin{equation}
 \boxed{
 \|(z-Q_aL_{a,R}Q_a)^{-1}\|
 \le CM_a^{-2}=Ca^{9/10}.}
 \label{eq:v24-high-resolvent}
\end{equation}
\end{theorem}

\begin{proof}
Choose the physical patch radius $r$ so small that the two relative freezing
errors in Lemma~\ref{lem:v24-freezing} are at most $c_*/8$, where $c_*$ is the
frozen ellipticity constant in \eqref{eq:v24-frozen-elliptic}.  On each patch,
\eqref{eq:v24-frozen-elliptic} and the freezing estimate give
\[
 \langle\chi_j\psi,L_{a,R}\chi_j\psi\rangle
 \ge {c_*\over2}
 \langle\chi_j\psi,K_a\chi_j\psi\rangle
 -C_r\|\chi_j\psi\|^2.
\]
Sum over $j$.  Apply the IMS identity both to $L_{a,R}$ and to $K_a$.
Lemma~\ref{lem:v24-IMS} bounds both commutator remainders by $C/r^2$.
Finite overlap and $\sum\chi_j^2=I$ therefore yield
\eqref{eq:v24-Garding} after renaming constants.  On $\Ran Q_a$ one has
$K_a\ge M_a^2$, proving \eqref{eq:v24-high-form}.  Since $\Gamma$ is bounded
and $M_a\to\infty$, the spectral theorem gives
\eqref{eq:v24-high-resolvent}.
\end{proof}

\subsection{Low/intermediate transfer}
\label{sec:v24-mid-transfer}

Let
\begin{equation}
 H_{a,R}:=Q_aL_{a,R}Q_a,
 \qquad
 B_{a,R}:=P_aL_{a,R}Q_a.
 \label{eq:v24-HB}
\end{equation}
Since $P_aK_aQ_a=0$, every transfer is generated by
$L_{a,R}-K_a$.

\begin{lemma}[Intermediate-zone transfer]
\label{lem:v24-mid-transfer}
With $\Pi_a$ from \eqref{eq:v24-Pi},
\begin{equation}
 \|P_aL_{a,R}\Pi_a\|
 \le C\bigl(1+\Lambda_a+a\Lambda_a^3\bigr).
 \label{eq:v24-mid-B}
\end{equation}
Hence
\begin{equation}
 \boxed{
 \|P_aL_{a,R}\Pi_a(z-H_{a,R})^{-1}\|
 \le
 C\left({\Lambda_a\over M_a^2}
       +{a\Lambda_a^3\over M_a^2}
       +M_a^{-2}\right)
 \le Ca^{1/10}.}
 \label{eq:v24-mid-transfer-rate}
\end{equation}
\end{lemma}

\begin{proof}
Because $P_aK_a\Pi_a=0$,
\[
 P_aL_{a,R}\Pi_a
 =P_a(T_{1,a,R}+T_{0,a,R}+E_{a,R})\Pi_a.
\]
Apply \eqref{eq:v24-T1}--\eqref{eq:v24-E3} with
$\Lambda=\Lambda_a$ to get \eqref{eq:v24-mid-B}.  Multiply by the high
resolvent bound \eqref{eq:v24-high-resolvent}.  With
\eqref{eq:v24-scales},
\begin{align*}
 {\Lambda_a\over M_a^2}&=a^{3/10},\\
 {a\Lambda_a^3\over M_a^2}&=a^{1/10},\\
 M_a^{-2}&=a^{9/10}.
\end{align*}
The largest term is $O(a^{1/10})$.
\end{proof}

\subsection{Ultraviolet transfer by repeated commutators}
\label{sec:v24-UV-transfer}

The cutoff-scale branch is not estimated by a continuum Taylor expansion.  It
is suppressed by spectral separation and smoothness.

\begin{lemma}[Native iterated-commutator bound]
\label{lem:v24-commutator-bound}
For every fixed integer $N\ge1$ and every retained source derivative,
\begin{equation}
 \boxed{
 \|\operatorname{ad}_{K_a}^{N}(L_{a,R})\|
 \le C_N a^{-(N+2)}.}
 \label{eq:v24-ad-bound}
\end{equation}
\end{lemma}

\begin{proof}
Use the exponentially summable shift/path expansions from (N1).  A typical
second-order term has engineering prefactor $a^{-2}$ times a smooth coefficient
and a bounded parallel shift.  Commuting once with $K_a$ introduces another
$a^{-2}$ shift factor.  If the two translations were constant commuting
translations their leading products would cancel.  The actual remainder comes
from either:
\begin{enumerate}[label=\textup{(i)},leftmargin=2.5em]
\item a one-edge difference of a smooth coefficient, which is $O(a)$ by (N2);
\item the difference between the two path orderings around a bounded plaquette,
which is $O(a^2)$ by smooth curvature.
\end{enumerate}
Thus the first commutator has size at most $Ca^{-3}$, not the crude $a^{-4}$.
Each further commutator adds one $a^{-2}$ shift but again at least one fresh
coefficient difference $O(a)$ or a curvature area factor, so it costs at most
one additional factor $a^{-1}$.  Induction gives $a^{-(N+2)}$.  The path range
grows only linearly with fixed $N$; the V1 exponential moments sum all tails.
Source differentiation changes only the finite smooth coefficient family and
therefore preserves the same power counting.
\end{proof}

\begin{lemma}[Spectral off-diagonal commutator estimate]
\label{lem:v24-Sylvester}
Let $K=K^*$, let
$E_-:=\mathbf 1_{(-\infty,\mu]}(K)$ and
$E_+:=\mathbf 1_{[\lambda,\infty)}(K)$ with $\lambda>\mu$, and let $T$ have a
bounded $N$th iterated commutator with $K$.  Then
\begin{equation}
 \boxed{
 \|E_-TE_+\|
 \le {\|\operatorname{ad}_K^N(T)\|\over(\lambda-\mu)^N}.}
 \label{eq:v24-Sylvester}
\end{equation}
\end{lemma}

\begin{proof}
On operators from $\Ran E_+$ to $\Ran E_-$ consider the Sylvester map
\(
 \mathcal S(X)=K_-X-XK_+
\), where $K_-=E_-KE_-$ and $K_+=E_+KE_+$.  Since the two spectra are separated
by $\lambda-\mu$, the standard spectral integral formula for the Sylvester
equation gives
\(
 \|\mathcal S^{-1}\|\le(\lambda-\mu)^{-1}
\).
Moreover
\[
 \mathcal S^N(E_-TE_+)=E_-\operatorname{ad}_K^N(T)E_+.
\]
Apply $\mathcal S^{-N}$ and take norms.
\end{proof}

\begin{lemma}[Cutoff-zone transfer]
\label{lem:v24-UV-transfer}
Take $N=6$.  Then
\begin{equation}
 \|P_aL_{a,R}U_a\|
 \le C a^{-8}(\Lambda_a^2-M_a^2)^{-6}
 \le Ca^{-4/5},
 \label{eq:v24-PLU}
\end{equation}
and therefore
\begin{equation}
 \boxed{
 \|P_aL_{a,R}U_a(z-H_{a,R})^{-1}\|
 \le Ca^{1/10}.}
 \label{eq:v24-UV-transfer-rate}
\end{equation}
\end{lemma}

\begin{proof}
Apply Lemma~\ref{lem:v24-Sylvester} with
$K=K_a$, $T=L_{a,R}$, $E_-=P_a$, $E_+=U_a$ and then use
Lemma~\ref{lem:v24-commutator-bound}.  Since
$M_a/\Lambda_a\to0$,
$\Lambda_a^2-M_a^2\ge\Lambda_a^2/2$ for small $a$.  Hence
\[
 \|P_aL_{a,R}U_a\|
 \le C a^{-8}\Lambda_a^{-12}
 =Ca^{-8+36/5}=Ca^{-4/5}.
\]
Multiplying by \eqref{eq:v24-high-resolvent}, which is $O(a^{9/10})$, gives
\eqref{eq:v24-UV-transfer-rate}.
\end{proof}

\begin{theorem}[Low/high transfer closes]
\label{thm:v24-transfer}
Uniformly for $z\in\Gamma$,
\begin{equation}
 \boxed{
 \|B_{a,R}(z-H_{a,R})^{-1}\|
 +\|(z-H_{a,R})^{-1}B_{a,R}^*\|
 \le Ca^{1/10}.}
 \label{eq:v24-transfer}
\end{equation}
\end{theorem}

\begin{proof}
Insert $Q_a=\Pi_a+U_a$ on the high side of $B_{a,R}$.  The two terms are
Lemmas~\ref{lem:v24-mid-transfer} and \ref{lem:v24-UV-transfer}.  The adjoint
bound follows from self-adjointness of $L_{a,R}$ and the conjugate contour
estimate; equivalently repeat the same argument with adjoints.
\end{proof}

\subsection{Low consistency and self-energy}
\label{sec:v24-low-selfenergy}

Let $A_{0,a}$ and $R_{0,a}(z)$ be the common sampled continuum soft operator and
resolvent from V23.  Standard continuum elliptic regularity on the compact
protected panel, transported through the fixed sampling map, gives
\begin{align}
 \|(K_a+I)R_{0,a}(z)\|&\le C,
 \label{eq:v24-R0-H2}\\
 \|(K_a+I)^{3/2}P_aR_{0,a}(z)\|&\le CM_a,
 \label{eq:v24-R0-H3}
\end{align}
uniformly on $\Gamma$.  The second estimate merely records that the continuum
resolvent gains two derivatives while $P_a$ limits the remaining derivative by
$M_a$.

\begin{lemma}[Covariant low-block consistency]
\label{lem:v24-low-consistency}
For each endpoint,
\begin{equation}
 \boxed{
 \sup_{z\in\Gamma}
 \|(A_{a,R}-A_{0,a})R_{0,a}(z)\|
 \le CaM_a
 =Ca^{11/20}.}
 \label{eq:v24-low-consistency}
\end{equation}
The same estimate, with a constant depending on the source multi-index, holds
for every fixed retained source derivative after the declared engineering stock
is divided out.
\end{lemma}

\begin{proof}
Use the common soft-jet decomposition
\eqref{eq:v24-common-soft-jet} from (N5).  Apply
\eqref{eq:v24-R0-H2}--\eqref{eq:v24-R0-H3}: the order-three term is bounded by
$CaM_a$, while every lower-order term is $O(a)$.  Since $M_a\to\infty$, the
latter is absorbed.  The source version of (N5) gives the identical estimate
for every retained fixed source jet, and the finite continuum resolvent
regularity is unchanged.
\end{proof}

\begin{lemma}[The return leg from the soft resolvent is uniformly bounded]
\label{lem:v24-return-leg}
Uniformly in $z\in\Gamma$,
\begin{equation}
 \boxed{
 \|B_{a,R}^*R_{0,a}(z)\|
 \le C(1+aM_a)
 \le C'.}
 \label{eq:v24-return-leg}
\end{equation}
The same statement holds coefficientwise for every retained source derivative.
\end{lemma}

\begin{proof}
Since $Q_aK_aP_a=0$,
\[
 B_{a,R}^*R_{0,a}
 =Q_a(L_{a,R}-K_a)P_aR_{0,a}.
\]
Use the order decomposition \eqref{eq:v24-order-decomp}.  The order-one term is
bounded by \eqref{eq:v24-R0-H2}; the zeroth-order term is bounded directly; the
$a$-weighted order-three term is bounded by
$a\|(K_a+I)^{3/2}P_aR_{0,a}\|\le CaM_a$.
Fixed source derivatives satisfy the same differential-order bounds from
Lemma~\ref{lem:v24-V2-package}.
\end{proof}

\begin{corollary}[High self-energy is small]
\label{cor:v24-selfenergy}
Uniformly on $\Gamma$,
\begin{equation}
 \boxed{
 \|B_{a,R}(z-H_{a,R})^{-1}B_{a,R}^*R_{0,a}(z)\|
 \le Ca^{1/10}.}
 \label{eq:v24-selfenergy}
\end{equation}
\end{corollary}

\begin{proof}
Multiply the transfer estimate \eqref{eq:v24-transfer} by the bounded return
leg \eqref{eq:v24-return-leg}.
\end{proof}

\subsection{Fixed source-resolvent words}
\label{sec:v24-source-words}

For a source derivative of the high resolvent, the useful estimate is relative,
not the crude global operator norm of a differentiated second-order vertex.

\begin{lemma}[Uniform high form bounds for source vertices]
\label{lem:v24-source-form}
For every fixed source multi-index $I$ in $\mathfrak q$ there is $C_I$ such
that
\begin{equation}
 \boxed{
 \|(K_a+I)^{-1/2}(\partial_IL_{a,R})(K_a+I)^{-1/2}\|
 \le C_I.}
 \label{eq:v24-source-form}
\end{equation}
Lower-order source vertices satisfy the corresponding stronger estimate with
one or both half-powers omitted.  The same statement holds for the common
continuum source vertices.
\end{lemma}

\begin{proof}
Coframe variations differentiate the principal second-order coefficients and
are therefore form bounded by $K_a$ on the compact ellipticity chart.
Gauge/spin connection variations are at most first order after their declared
source stock is removed; Higgs/Yukawa, antifield and line-frame insertions are
zeroth order or finite-rank/source-local.  This is precisely the engineering
degree classification in Theorem~\ref{thm:v2-source-vertex-transfer}, now
written as a quadratic-form estimate.  Fixed products and source descendants
remain bounded because the source order is finite.
\end{proof}

\begin{proposition}[Source-complete V24 package]
\label{prop:v24-source-complete}
Fix the finite source/BV packet $\mathfrak q$.  Every differentiated block word
appearing in Proposition~\ref{prop:v23-source-word-reduction} obeys the V24
estimates with a constant depending on the word but with the same vanishing
scale $O(a^{1/10})$ for every open low/high transfer and closed high excursion.
In particular the complete source-resolvent/Riesz family satisfies V23 item
\textup{(G5)}.
\end{proposition}

\begin{proof}
Differentiate the exact high resolvent.  Each derivative is a finite sum of
words
\[
 (z-H)^{-1}H_{I_1}(z-H)^{-1}\cdots
 H_{I_m}(z-H)^{-1}.
\]
Insert $(K_a+I)^{\pm1/2}$ around every $H_{I_j}$.  On $\Ran Q_a$,
Theorem~\ref{thm:v24-Garding} makes the energy-weighted high resolvent bounded:
\[
 \|(K_a+I)^{1/2}(z-H)^{-1}(K_a+I)^{1/2}\|\le C,
\]
while an unsandwiched high resolvent contributes the net factor $O(M_a^{-2})$.
Lemma~\ref{lem:v24-source-form} therefore shows that inserting any fixed number
of source vertices does not destroy the high suppression.

For low/intermediate pieces, differentiate the order decomposition
\eqref{eq:v24-order-decomp}; V2's source-vertex transfer preserves the same
orders $1,0,a\times3$.  For the ultraviolet piece, source differentiation of
Lemma~\ref{lem:v24-commutator-bound} only changes $C_N$.  Hence the proofs of
Lemmas~\ref{lem:v24-mid-transfer} and \ref{lem:v24-UV-transfer} repeat word by
word.  Every closed excursion is completed with the source-differentiated
version of Lemma~\ref{lem:v24-return-leg}.  The number of words is finite at
fixed source order, so all constants may be summed.
\end{proof}

\subsection{Closure of the V23 covariant soft--hard package}
\label{sec:v24-close-package}

\begin{theorem}[Covariant discrete G\aa rding/Feshbach mixing theorem]
\label{thm:v24-main}
Consider the explicit V2 Wilson/overlap/Yukawa completion, or any microscopic
completion obeying the same V1--V2 smooth protected native package of
Lemma~\ref{lem:v24-V2-package}, on a compact smooth fixed-index
gauge--spin--coframe--Higgs family.  On each finite protected source chart,
freeze the common frequency operator \eqref{eq:v24-K} at the chart centre and
use the scales \eqref{eq:v24-scales}.  Then, uniformly in the microscopic
completion and in $z\in\Gamma$,
\begin{align}
 Q_aL_{a,R}Q_a&\succeq cM_a^2Q_a,
 \label{eq:v24-main-Garding}\\
 \|(z-H_{a,R})^{-1}\|&\le Ca^{9/10},
 \label{eq:v24-main-high}\\
 \|B_{a,R}(z-H_{a,R})^{-1}\|
 +\|(z-H_{a,R})^{-1}B_{a,R}^*\|
 &\le Ca^{1/10},
 \label{eq:v24-main-transfer}\\
 \|(A_{a,R}-A_{0,a})R_{0,a}\|
 &\le Ca^{11/20},
 \label{eq:v24-main-low}\\
 \|B_{a,R}(z-H_{a,R})^{-1}B_{a,R}^*R_{0,a}\|
 &\le Ca^{1/10}.
 \label{eq:v24-main-self}
\end{align}
Every statement holds for all retained fixed source derivatives in the sense of
V23's source-resolvent words.  Thus Assumption
\ref{ass:v23-covariant-package} is a theorem on this smooth protected native
class.
\end{theorem}

\begin{proof}
Equations \eqref{eq:v24-main-Garding}--\eqref{eq:v24-main-high} are
Theorem~\ref{thm:v24-Garding}.  Equation
\eqref{eq:v24-main-transfer} is Theorem~\ref{thm:v24-transfer}.
Equation \eqref{eq:v24-main-low} is Lemma~\ref{lem:v24-low-consistency}, and
\eqref{eq:v24-main-self} is Corollary~\ref{cor:v24-selfenergy}.  The source
statement is Proposition~\ref{prop:v24-source-complete}.  A finite cover of the
compact background panel makes all chart constants uniform; the final full
resolvent estimates do not depend on which frozen auxiliary frequency projector
was used in a chart.
\end{proof}

\begin{corollary}[Common soft Riesz projector on the nonconstant protected family]
\label{cor:v24-common-Riesz}
Under Theorem~\ref{thm:v24-main}, the V23 Feshbach bridge gives
\begin{equation}
 \boxed{
 \sup_{z\in\Gamma}
 \|(z-L_{a,A})^{-1}-(z-L_{a,B})^{-1}\|
 \le Ca^{1/10},}
 \label{eq:v24-endpoint-resolvent}
\end{equation}
and therefore
\begin{equation}
 \boxed{
 \|p_{a,A}-p_{a,B}\|
 \le C_\Gamma a^{1/10}.}
 \label{eq:v24-endpoint-Riesz}
\end{equation}
The same convergence holds for every retained fixed source derivative.
For sufficiently small $a$, the norm distance is less than one, so V14's
canonical direct rotation supplies one controlled identification of the exact
zero-mode modules.  Hence
\begin{equation}
 \boxed{
 \operatorname{Ind}^{\rm ex}_{\rm ctl}(D_{a,A})
 =
 \operatorname{Ind}^{\rm ex}_{\rm ctl}(D_{a,B}).}
 \label{eq:v24-controlled-index-equal}
\end{equation}
\end{corollary}

\begin{proof}
The estimates in Theorem~\ref{thm:v24-main} give the V23 control parameters
$h=O(a^{9/10})$, $\nu=O(a^{1/10})$ and
$e=O(a^{1/10})$.  Apply
Theorem~\ref{thm:v23-Feshbach-resolvent} and then integrate around $\Gamma$ as
in Corollary~\ref{cor:v23-Riesz-close}.  The controlled-index statement is the
V22/V14 direct-rotation consequence.
\end{proof}

\begin{corollary}[Closure of the faithful complex microscopic anchoring problem]
\label{cor:v24-complex-anchoring-closed}
On the smooth protected faithful complex Standard-Model branch covered by
Theorem~\ref{thm:v24-main}, any two \emph{V2-consistent} microscopic
completions---meaning that they share the common soft continuum jet (N5), in
addition to the same physical Dirac--Yukawa principal branch and ordinary fixed
index---have the same exact controlled family index for sufficiently small
cutoff spacing.
Consequently V19 supplies one global unstabilized coefficientwise native
regulator path, and the corresponding normalized continuum states are related
by the finite scheme/source/line groupoid equivalence already proved in V16.
\end{corollary}

\begin{proof}
Corollary~\ref{cor:v24-common-Riesz} kills the V17 primary controlled module
class by a canonical direct rotation.  V19 then kills the secondary complement
and tertiary identity-component obstructions coefficientwise.  V16 gives the
global continuum descent statement.  The faithful complex determinant-line
holonomy was already closed in V12.
\end{proof}

\subsection{Scope and remaining stronger problems}
\label{sec:v24-scope}

\begin{firewall}[Smooth protected family, not arbitrary rough gauge fields]
Theorem~\ref{thm:v24-main} uses a compact smooth sampled background class: the
coframe and gauge/spin connection have the finite derivative bounds needed for
local freezing and for the iterated-commutator estimate, and the reduced
zero-mode cluster remains separated by the declared positive gap.  V24 does not
claim the same estimate for lattice-scale rough gauge fields, for backgrounds
approaching an index-changing divisor, or after the positive IR/reference floor
is removed.
\end{firewall}

\begin{firewall}[The auxiliary frequency projector need not be quasilocal]
The sharp spectral projections $P_a,\Pi_a,U_a$ of $K_a$ are only proof devices.
No locality of these projectors is used or claimed.  The physical Riesz
projectors $p_{a,R}$ are quasilocal by V13.  Source derivatives are taken with
the auxiliary decomposition frozen at a chart centre, so no differentiability
of the artificial spectral thresholds is required.
\end{firewall}

\begin{firewall}[The rate is nonoptimal]
The exponent $1/10$ is chosen only to make all three zones close simultaneously
with elementary fixed-order commutator bounds.  Any exponents
$0<\beta<\gamma<1$ satisfying
\begin{equation}
 \gamma>\frac12,
 \qquad
 \gamma<2\beta,
 \qquad
 1-3\gamma+2\beta>0,
 \label{eq:v24-exponent-region}
\end{equation}
and a sufficiently large fixed commutator depth $N$ yield convergence.  No
physical meaning is attached to $1/10$.
\end{firewall}

\subsection{V24 anti-circularity audit}
\label{sec:v24-audit}

\begin{enumerate}[leftmargin=2.4em]
\item \textbf{The high lattice branch is not compared to the continuum.}
It enters only through native G\aa rding coercivity and a vanishing high
resolvent.
\item \textbf{Flat Taylor expansion is used only below $\Lambda_a\ll a^{-1}$.}
The cutoff-scale low/high coupling is controlled instead by repeated
commutators and spectral separation.
\item \textbf{The sharp auxiliary projector is never source differentiated.}
The complete differentiated resolvent is expanded after freezing the proof
decomposition at the source-chart centre.
\item \textbf{The G\aa rding theorem is global.}
It is obtained by local freezing plus the exact discrete IMS identity; local
symbol ellipticity is not simply asserted to hold globally.
\item \textbf{The high self-energy is not bounded by a crude $\|B\|^2$.}
One vanishing transfer factor is multiplied by the continuum-resolvent return
leg, whose elliptic regularity is uniformly bounded.
\item \textbf{All source statements retain the native engineering stocks.}
Principal coframe derivatives are treated as high-form-bounded second-order
vertices; lower source sectors are easier.
\item \textbf{No new topology is imported.}
Once the operator estimate gives norm-close Riesz projectors, V14, V19 and V22
perform the already-proved controlled-module descent.
\end{enumerate}

\section{V24 status ledger}
\label{sec:v24-status}

\subsection*{Proved in V24}
\begin{enumerate}[leftmargin=2.2em]
\item A common gauge/spin/coframe-covariant comparison Laplacian $K_a$ gives a
three-zone decomposition with
$M_a=a^{-9/20}$ and $\Lambda_a=a^{-3/5}$; the decomposition may be frozen at
each source-chart centre without entering source derivatives.
\item The explicit V2 stencil plus V1 quasilocal functional calculus gives the
native smooth-background package: order-two exponential locality, one-edge
$O(a)$ coefficient variation, frozen full-zone ellipticity and an
order-$1/0/a\times3$ low/intermediate decomposition and one common
$a$-suppressed continuum soft jet, all with fixed source jets.
\item A discrete covariant Bernstein estimate controls every fixed covariant
difference word on a $K_a$-spectral window by the corresponding power of the
window scale.
\item An exact discrete IMS identity plus local freezing proves the global
covariant G\aa rding inequality
$L_{a,R}\succeq c_GK_a-C_G$, hence the high resolvent is
$O(M_a^{-2})=O(a^{9/10})$.
\item The intermediate-zone transfer is
$O(\Lambda_aM_a^{-2}+a\Lambda_a^3M_a^{-2})=O(a^{1/10})$.
\item Repeated commutators satisfy
$\|\operatorname{ad}_{K_a}^N L_{a,R}\|\le C_Na^{-(N+2)}$.
A Sylvester spectral-separation estimate with $N=6$ makes the direct
soft-to-cutoff transfer $O(a^{1/10})$ after one high resolvent.
\item The common low continuum consistency is $O(aM_a)=O(a^{11/20})$ and the
closed high self-energy is $O(a^{1/10})$.
\item High-form bounds for every retained source vertex propagate the same
three-zone estimates through all fixed source-resolvent words.  Thus the full
V23 covariant soft--hard package, including item (G5), is proved on the smooth
protected native class.
\item The endpoint soft resolvents and Riesz projectors differ by
$O(a^{1/10})$.  Their exact controlled family indices therefore agree for
small $a$, and the V17--V19 complex microscopic obstruction tower collapses.
\item Consequently the faithful complex microscopic anchoring problem is
closed on this smooth protected fixed-index \emph{V2-consistent} class: all
such admissible native completions have one global unstabilized coefficientwise
regulator path and the same matched continuum.  Principal-symbol/index equality
without the common soft jet remains excluded by the V20 counterexample.
\end{enumerate}

\subsection*{Inherited}
\begin{enumerate}[leftmargin=2.2em]
\item V1--V7: source-complete coefficientwise ESM continuum, common BV
restoration, fixed-loop forest induction, UV Cauchy convergence,
scheme/reference covariance and exact perturbative-order projectivity.
\item V8--V19: regulator universality machinery, fixed-index zero-mode shell
calculus, global line transport, controlled-module descent and automatic
downstream cancellation once the primary zero-mode modules match.
\item V20--V22: exact controlled family-index classification, its spatial
anchor/density shadows, and the theorem that norm-close soft Riesz projectors
identify that exact controlled class.
\item V23: the exact two-zone Feshbach resolvent bridge and complete
translation-invariant reference-branch verification.
\end{enumerate}

\subsection*{Conditional}
\begin{enumerate}[leftmargin=2.2em]
\item V24 assumes the compact background family is smooth enough for the finite
coefficient derivatives and curvature bounds used in freezing and in the six
commutators.  It does not cover lattice-scale rough fields.
\item The fixed-index reduced singular-value floor and the positive IR/reference
prescription remain in force.  Index-changing or accidental zero-mode-pair
crossings leave the protected class.
\item Source conclusions are coefficientwise at the prescribed finite packet;
no analytic source-radius theorem is asserted.
\item The independent real Pfaffian mod-two orientation/zero-crossing problem
is not solved by the complex Riesz-projector comparison.
\item The explicit $a^{1/10}$ rate is a convenient nonoptimal choice and is not
claimed outside the declared smooth native class.
\end{enumerate}

\subsection*{Open beyond V24}
\begin{enumerate}[leftmargin=2.2em]
\item Compute the real Pfaffian mod-two spectral-flow/orientation character for
any independently physical Majorana sector and determine whether the faithful
Renewal background grammar trivializes it.
\item Analyze a controlled crossing out of the minimal fixed-index stratum:
creation of an accidental kernel--cokernel pair and true index change require a
spectral-flow/divisor source packet rather than a uniformly bounded reduced
inverse.
\item Extend the V24 G\aa rding/Feshbach theorem to rougher background classes
if such a theorem is physically needed; this is stronger than the smooth
coefficientwise continuum target.
\item The independent positive-IR removal, analytic finite-coupling RG,
perturbative summability and nonperturbative completion problems remain open.
\end{enumerate}

\subsection*{Exact next load-bearing theorem}
\begin{center}
\fbox{\parbox{0.92\textwidth}{
\textbf{Real Pfaffian orientation / mod-two spectral-flow theorem.}
On any independently physical neutral Majorana sector, construct the protected
real zero-mode/Pfaffian line over the same smooth fixed-index background class
and compute its first Stiefel--Whitney character.  Prove that every closed
native regulator/background loop has even real spectral flow (hence the
Pfaffian line is orientable and the V12 mod-two character vanishes), or exhibit
a genuine loop with odd spectral flow.  The complex faithful Standard-Model
regulator ambiguity is closed by V24; this real orientation problem is now the
smallest remaining global line obstruction before one turns to index-changing
divisors or the independent infrared/nonperturbative programmes.}}
\end{center}

\section*{Append-only continuation marker: V24}
\addcontentsline{toc}{section}{Append-only continuation marker: V24}
\textbf{End of V24.}  The complete V1--V23 body above is retained verbatim.
V24 proves the covariant discrete G\aa rding/Feshbach mixing theorem isolated by
V23 on the smooth protected V2-consistent native background class.  A common covariant
frequency operator and a three-zone split separate the soft, intermediate and
cutoff branches.  Local freezing plus an exact discrete IMS identity gives the
global high-frequency lower bound; V2 Taylor consistency controls the
intermediate zone; six iterated commutators suppress direct soft-to-cutoff
mixing.  With $M_a=a^{-9/20}$ and $\Lambda_a=a^{-3/5}$ every Feshbach error is
$O(a^{1/10})$ or better, including all retained fixed source words.  The two
endpoint soft Riesz projectors are therefore norm close and have the same exact
controlled family index.  The faithful complex microscopic anchoring ambiguity
is consequently closed on this smooth fixed-index class.  The next independent
line-level gate is the real Pfaffian mod-two orientation problem.



\clearpage
\section{V25 continuation: native real Pfaffian orientation and mod-two spectral flow}
\label{sec:v25-continuation}

V24 closes the faithful \emph{complex} microscopic regulator ambiguity on the
smooth protected V2-consistent class.  The only remaining line-level issue from
V12 is genuinely real: if an independently physical neutral Majorana sector is
represented by a real skew fermionic block, its Pfaffian is a real line-valued
object and the complex determinant theorem does not choose its sign.

There is an upstream distinction which must be made before attacking that
sign.  A Majorana mass term can be handled in two different ways:
\begin{enumerate}[label=\textup{(M\arabic*)},leftmargin=3em]
\item it may be embedded in the already-controlled complex chiral/Nambu block,
      in which case there is no independent real Pfaffian orientation datum;
\item or one may supply a genuinely real skew Majorana block with its own real
      Grassmann measure, in which case a real Pfaffian line is physical and its
      orientation must be proved.
\end{enumerate}
V25 treats only the second case.  This is exactly the ``independently physical
real Pfaffian'' branch kept separate since V12.

The key finite-regulator observation is simpler than a general eta-invariant
calculation.  On a constant-rank branch, a real skew operator canonically
orients the nonzero complement by its symplectic form.  If the underlying real
Majorana carrier is itself oriented, this induced symplectic orientation then
canonically orients the kernel.  Hence the zero-mode Pfaffian line is oriented
without choosing local signs.  For the gauge-neutral native Majorana carrier,
the allowed spin/coframe and lattice identifications preserve that carrier
orientation.  The mod-two Pfaffian character therefore vanishes on the whole
protected branch.

\subsection{The oriented native real Majorana carrier}
\label{sec:v25-oriented-carrier}

Fix a cutoff $a$ and a protected smooth background component $B$.  Let
\begin{equation}
 \mathcal E_a^{\mathbb R}\longrightarrow B
 \label{eq:v25-real-carrier}
\end{equation}
be the finite real one-particle bundle of an independently supplied neutral
Majorana block.  Its real fibre dimension is even.  The microscopic Grassmann
order at one reference background fixes an orientation
$\mathfrak o_{\mathcal E}$ of this bundle.

\begin{definition}[Native oriented neutral-Majorana branch]
\label{def:v25-oriented-branch}
A real Majorana realization is \emph{natively oriented} if:
\begin{enumerate}[label=\textup{(O\arabic*)},leftmargin=3em]
\item the real carrier has even onsite fibre dimension $d_M$ and one fixed
      Berezin/site/flavour orientation at the reference background;
\item local spin transport acts through a real representation of the connected
      group $\mathrm{Spin}^+(1,3)$ or its Euclidean protected continuation, so every
      spin-frame transition has real determinant $+1$;
\item the sector is internal-gauge neutral, or more generally every internal
      real transition allowed in this block has determinant $+1$;
\item oriented lattice/background relabellings act by site permutations and
      orientation-preserving onsite maps;
\item any independent neutral flavour frame is fixed once and for all, or its
      allowed transition group is reduced to the orientation-preserving
      subgroup.
\end{enumerate}
\end{definition}

The first four items are automatic for the canonical sitewise gauge-neutral
Majorana realization once it has been physically supplied.  The fifth prevents
an arbitrary orientation-reversing flavour redefinition from being declared a
gauge redundancy after the Pfaffian sign has been fixed.

\begin{lemma}[Native transitions preserve the real carrier orientation]
\label{lem:v25-native-SO}
Every transition generated by Definition~\ref{def:v25-oriented-branch} acts on
$\mathcal E_a^{\mathbb R}$ with determinant $+1$.  Hence
$\mathfrak o_{\mathcal E}$ extends uniquely over the connected protected
background component.
\end{lemma}

\begin{proof}
The determinant of a finite-dimensional real representation of a connected
group is a continuous map to $\{\pm1\}$.  Since it equals $+1$ at the identity,
every spin transition has determinant $+1$.  The neutral internal action is the
identity, or is orientation preserving by assumption.  A permutation of two
lattice sites exchanges two blocks of real dimension $d_M$ and therefore has
determinant
\[
 (-1)^{d_M^2}=+1
\]
because $d_M$ is even.  Thus every site permutation is orientation preserving.
The same is true of the declared flavour transitions.  Products of these
transformations have determinant $+1$, proving the claim.
\end{proof}

\begin{firewall}[Orientation is physical input, not complex positivity]
The orientation in Definition~\ref{def:v25-oriented-branch} comes from the real
Majorana measure/Berezin order and the orientation-preserving native structure
group.  It is not inferred from positivity of a doubled determinant or from the
complex faithful-spin line theorem.  If the physical real branch permits a
transition of determinant $-1$, V25 does not discard that sign.
\end{firewall}

\subsection{A skew form canonically orients its nonzero complement}
\label{sec:v25-symplectic-orientation}

Let
\begin{equation}
 A_b:\mathcal E_{a,b}^{\mathbb R}\longrightarrow
     \mathcal E_{a,b}^{\mathbb R},
 \qquad
 A_b^T=-A_b,
 \label{eq:v25-skew-family}
\end{equation}
be the real skew precision operator of the Majorana Gaussian.  On a protected
constant-rank chart assume
\begin{equation}
 K_b:=\ker A_b,
 \qquad
 \dim_{\mathbb R}K_b=2r,
 \qquad
 s_{\min}^+(A_b)\ge\tau_M>0.
 \label{eq:v25-real-gap}
\end{equation}
The kernel dimension is even because the rank of a real skew matrix is even.
Let $H_b=K_b^\perp$.

\begin{definition}[Reduced symplectic form]
\label{def:v25-omega}
On $H_b$ define
\begin{equation}
 \omega_b(u,v):=\langle u,A_bv\rangle.
 \label{eq:v25-omega}
\end{equation}
It is a nondegenerate skew form.  Its \emph{symplectic orientation}
$\mathfrak o_{\omega,b}$ is the orientation for which
\begin{equation}
 \frac{1}{m!}\omega_b^{\wedge m}>0,
 \qquad \dim H_b=2m.
 \label{eq:v25-symp-orientation}
\end{equation}
\end{definition}

\begin{lemma}[Smoothness of the symplectic orientation]
\label{lem:v25-symp-smooth}
On every protected constant-rank chart, the orientation
$\mathfrak o_{\omega,b}$ varies continuously and is invariant under every
orientation-preserving real congruence of the Majorana carrier.
\end{lemma}

\begin{proof}
The reduced gap gives a smooth Riesz projector onto $K_b$ and therefore a
smooth complement bundle $H$.  Nondegeneracy of $\omega_b|_{H_b}$ follows from
invertibility of $A_b|_{H_b}$.  The top form $\omega_b^{\wedge m}$ is nowhere
zero and therefore chooses a continuous orientation.  If
$A\mapsto g^TAg$ and $g$ is orientation preserving on the ambient carrier, the
induced pullback of $\omega^{\wedge m}$ has the same symplectic orientation on
the transported complement.
\end{proof}

\subsection{Canonical orientation of the real zero-mode bundle}
\label{sec:v25-kernel-orientation}

The ambient orientation and the reduced symplectic orientation determine the
missing zero-mode orientation.

\begin{theorem}[Canonical kernel orientation on a constant-rank real skew family]
\label{thm:v25-kernel-orientation}
Under \eqref{eq:v25-real-gap}, there is a unique continuous orientation
$\mathfrak o_K$ of the kernel bundle $K\to B$ characterized fibrewise by
\begin{equation}
 \boxed{
 \mathfrak o_{\mathcal E}
 =\mathfrak o_K\wedge\mathfrak o_{\omega}.}
 \label{eq:v25-kernel-orientation}
\end{equation}
It is invariant under every native transition of
Definition~\ref{def:v25-oriented-branch}.  Consequently
\begin{equation}
 \boxed{w_1(K)=0.}
 \label{eq:v25-w1-K-zero}
\end{equation}
\end{theorem}

\begin{proof}
The orthogonal splitting
$\mathcal E^{\mathbb R}=K\oplus H$ gives
$\det\mathcal E\simeq\det K\otimes\det H$.  The ambient orientation gives a
positive ray in $\det\mathcal E$, while
Definition~\ref{def:v25-omega} gives a positive ray in $\det H$.
There is therefore a unique positive ray in $\det K$ satisfying
\eqref{eq:v25-kernel-orientation}.  Smoothness follows from smoothness of the
Riesz kernel projector and Lemma~\ref{lem:v25-symp-smooth}.  Native transition
maps preserve the ambient orientation by Lemma~\ref{lem:v25-native-SO} and the
symplectic orientation by congruence, so they preserve the induced kernel
orientation.  A real vector bundle admits an orientation exactly when its first
Stiefel--Whitney class vanishes.
\end{proof}

\begin{corollary}[Canonical orientation of the zero-mode Pfaffian line]
\label{cor:v25-pf-line-orientation}
In the inherited finite zero-mode factorization, the real zero-mode Pfaffian
factor is a top exterior line of $K$ or its dual, depending on the fixed
Berezin convention.  Either convention has the same first
Stiefel--Whitney class.  Theorem~\ref{thm:v25-kernel-orientation} therefore
orients the real zero-mode Pfaffian line globally on the protected
constant-rank branch:
\begin{equation}
 \boxed{w_1(\mathcal P)=0.}
 \label{eq:v25-w1-P-zero}
\end{equation}
\end{corollary}

\begin{proof}
The first Stiefel--Whitney class is unchanged by dualization of a real line.
Apply \eqref{eq:v25-w1-K-zero} and the inherited zero-mode/Pfaffian
factorization.
\end{proof}

\subsection{The reduced Pfaffian has a canonical positive sign}
\label{sec:v25-reduced-pf}

The preceding theorem orients the zero modes.  The nonzero modes have their own
canonical sign supplied by the skew form itself.

\begin{proposition}[Positive reduced Pfaffian in the symplectic frame]
\label{prop:v25-positive-reduced-Pf}
Let $A_b^\perp=A_b|_{H_b}$.  Relative to the symplectic orientation
$\mathfrak o_{\omega,b}$,
\begin{equation}
 \boxed{\operatorname{Pf}(A_b^\perp)>0.}
 \label{eq:v25-pf-positive}
\end{equation}
Thus the reduced nonzero Pfaffian cannot contribute a sign holonomy on the
protected component.
\end{proposition}

\begin{proof}
Choose any symplectic basis
$(e_1,f_1,\ldots,e_m,f_m)$ compatible with
$\mathfrak o_{\omega,b}$.  In such a basis the top form
$\omega_b^{\wedge m}/m!$ is positive.  By the defining identity between the
Pfaffian and the top exterior power of a skew form, the coefficient of that top
form is the Pfaffian.  Its sign is therefore positive in the symplectic
orientation.  Since the reduced gap prevents degeneracy, this positive scalar
varies continuously and never crosses zero.
\end{proof}

\begin{theorem}[Canonical scalarization of a saturated real Majorana amplitude]
\label{thm:v25-real-saturation}
Let $\omega_K$ be a zero-mode saturation writer which is positive in the
canonical kernel orientation of Theorem~\ref{thm:v25-kernel-orientation}.
Then the line-valued Majorana Gaussian factor scalarizes as
\begin{equation}
 \boxed{
 \mathcal A_A^{\rm Maj}
 =\langle\omega_K,z_K\rangle
  \operatorname{Pf}'(A),}
 \label{eq:v25-Maj-amplitude}
\end{equation}
with both factors positive in their canonical orientations.  This scalarization
is independent of local Kato-frame choices and every retained source derivative
is transported by the same finite contact calculus as in V11--V12.
\end{theorem}

\begin{proof}
The zero-mode factor is positive by the chosen saturation writer and
Theorem~\ref{thm:v25-kernel-orientation}; the reduced factor is positive by
Proposition~\ref{prop:v25-positive-reduced-Pf}.  On overlaps, a Kato-frame
change in $K$ has positive determinant because both frames represent the same
canonical orientation.  Hence the zero-mode saturation pairing is unchanged
in sign.  The reduced symplectic orientation is intrinsic, so its Pfaffian sign
is also frame independent.  Differentiated source contacts are exactly the
ones already retained by the V11 zero-mode and Pfaffian line calculus.
\end{proof}

\subsection{Mod-two spectral flow equals the endpoint orientation character}
\label{sec:v25-mod2-flow}

The constant-rank theorem already closes the protected branch.  It is useful to
record the stronger crossing statement because it identifies exactly how a
nontrivial sign could reappear outside that branch.

\begin{theorem}[Finite native mod-two crossing formula]
\label{thm:v25-mod2-crossing}
Let $A(t)$, $0\le t\le1$, be a generic $C^1$ path of real skew matrices on one
fixed oriented even-dimensional carrier, with invertible endpoints and only
regular Pfaffian crossings.  Suppose the endpoint is identified with the
initial point by a real native map $g$:
\begin{equation}
 A(1)=g^TA(0)g.
 \label{eq:v25-endpoint-congruence}
\end{equation}
Then
\begin{equation}
 \boxed{
 (-1)^{\operatorname{SF}_2(A)}
 =\operatorname{sgn}\det_{\mathbb R}g.}
 \label{eq:v25-sf2-detg}
\end{equation}
In particular, if $g$ preserves the native carrier orientation,
\begin{equation}
 \boxed{\operatorname{SF}_2(A)=0\pmod2.}
 \label{eq:v25-even-sf}
\end{equation}
\end{theorem}

\begin{proof}
Fix the ambient orientation and let
$f(t)=\operatorname{Pf}(A(t))$ in the corresponding fixed real Berezin frame.
At every regular crossing exactly one Pfaffian eigenvalue changes sign, so
\begin{equation}
 \operatorname{sgn}f(1)
 =(-1)^{\operatorname{SF}_2(A)}\operatorname{sgn}f(0).
 \label{eq:v25-sign-cross}
\end{equation}
On the other hand the congruence law for the Pfaffian gives
\begin{equation}
 f(1)=\operatorname{Pf}(g^TA(0)g)
 =\det(g)\operatorname{Pf}(A(0)).
 \label{eq:v25-pf-congruence}
\end{equation}
Comparing signs in the two identities proves
\eqref{eq:v25-sf2-detg}.  The parity is stable under generic perturbations by
creation or annihilation of crossing pairs, so the formula extends to every
path for which the mod-two spectral flow is defined.
\end{proof}

\begin{corollary}[Vanishing mod-two character on the oriented neutral branch]
\label{cor:v25-mod2-zero}
Every closed native loop in the oriented gauge-neutral Majorana branch of
Definition~\ref{def:v25-oriented-branch} has
\begin{equation}
 \boxed{
 \operatorname{Hol}_{\mathcal P}(\gamma)=+1,
 \qquad
 \operatorname{SF}_2(\gamma)=0.}
 \label{eq:v25-hol-zero}
\end{equation}
Hence the V12 real Pfaffian orientation character vanishes on this branch.
\end{corollary}

\begin{proof}
The endpoint identification of a native loop is a product of the transitions
listed in Definition~\ref{def:v25-oriented-branch}.  Lemma
\ref{lem:v25-native-SO} gives determinant $+1$.  Apply
Theorem~\ref{thm:v25-mod2-crossing} and the V12 equality between Pfaffian-line
holonomy and mod-two spectral flow.
\end{proof}

\begin{remark}[Consistency with the five-dimensional spin-bordism test]
\label{rem:v25-spin-bordism}
For a gauge-neutral Majorana sector coupled only to an ordinary spin background,
the standard spin-bordism calculation gives
$\Omega^{\mathrm{Spin}}_5(\mathrm{pt})=0$.  Thus there is no independent
five-dimensional spin-bordism character available to contradict
Corollary~\ref{cor:v25-mod2-zero}.  V25 does not use this external
classification in the proof: the native orientation and crossing arguments
already give the finite-regulator result.
\end{remark}

\subsection{Cofinal cutoff orientation and Pfaffian Schur factorization}
\label{sec:v25-cutoff-orientation}

A finite-cutoff orientation must be compatible with shell regrouping and cutoff
transport before it can enter the continuum theorem.  The exact Pfaffian Schur
cocycle inherited in V1--V8 supplies this compatibility once the hard factor is
oriented.

\begin{definition}[Positive hard Pfaffian component]
\label{def:v25-hard-positive}
For one hard shell let $A_h$ be its real skew Majorana block.  The protected
hard-pivot component is declared \emph{positive} when its Pfaffian sign, in the
fixed hard Berezin ordering, agrees with the chosen reference massive block.
Because the hard pivot never vanishes on the protected chart, this sign is
constant on the component.
\end{definition}

\begin{theorem}[Orientation-compatible Pfaffian shell cocycle]
\label{thm:v25-shell-orientation}
On the positive hard component, the inherited exact Pfaffian Schur identity
transports the canonical soft orientation of
Theorem~\ref{thm:v25-kernel-orientation} without a sign ambiguity.  For nested
hard shells, the resulting orientation maps compose exactly under finite shell
regrouping.  Hence the real Pfaffian orientation is cofinal under the same
adjacent-cutoff system used for the complex determinant line.
\end{theorem}

\begin{proof}
The inherited skew Schur factorization has the form
\begin{equation}
 \operatorname{Pf}(A_{\rm full})
 =\operatorname{Pf}(A_h)\,
  \operatorname{Pf}(A_{\rm Schur})
 \label{eq:v25-pf-Schur}
\end{equation}
with the sign fixed by the globally frozen Berezin block order.  On the
positive hard component, the first factor is positive and cannot change sign.
Therefore the orientation of the full Pfaffian factor is exactly the
orientation of the soft Schur factor.  Applying the same identity twice to
three nested scales gives the inherited Pfaffian scale cocycle; associativity
of finite block elimination proves exact composition of the orientation maps.
\end{proof}

\begin{corollary}[Cutoff-independent real orientation]
\label{cor:v25-cutoff-orientation}
For every cofinal sequence of cutoffs staying in the protected oriented neutral
Majorana component, the canonical real Pfaffian orientations form one
compatible inverse system.  A reference orientation chosen at one cutoff fixes
all finer and coarser cutoff orientations uniquely.
\end{corollary}

\subsection{Compatibility with V24 regulator comparison}
\label{sec:v25-regulator-comparison}

The V24 direct rotation which identifies two V2-consistent soft Riesz
projectors is itself orientation preserving on a real branch.

\begin{proposition}[Real direct rotation lies in the orientation-preserving component]
\label{prop:v25-real-direct-rotation}
Let $p,q$ be real orthogonal projections with $\|p-q\|<1$.  The V14 canonical
polar direct rotation
\begin{equation}
 u_{q,p}=X(X^TX)^{-1/2},
 \qquad
 X=qp+(I-q)(I-p),
 \label{eq:v25-real-rotation}
\end{equation}
is real orthogonal and is connected to the identity through real orthogonal
matrices.  Hence
\begin{equation}
 \boxed{\det_{\mathbb R}u_{q,p}=+1.}
 \label{eq:v25-real-rotation-det}
\end{equation}
It transports the canonical kernel orientation of $p$ to that of $q$.
\end{proposition}

\begin{proof}
All matrices in \eqref{eq:v25-real-rotation} are real.  The V14 path
$X_s=I+s(X-I)$ stays invertible when $\|p-q\|<1$, and its polar factors form a
continuous path in the real orthogonal group from $I$ to $u_{q,p}$.  The
determinant cannot change sign along such a path and equals $+1$ at the
identity.  The transported ambient and symplectic complement orientations are
therefore unchanged, so Theorem~\ref{thm:v25-kernel-orientation} gives the same
kernel orientation at the endpoint.
\end{proof}

\begin{theorem}[Regulator independence of the real Pfaffian orientation]
\label{thm:v25-regulator-orientation}
Let $A$ and $B$ be two smooth protected V2-consistent microscopic completions
of the same independently physical oriented neutral-Majorana sector.  For
sufficiently small $a$, V24 gives norm-close real soft Riesz projectors and the
real direct rotation of Proposition~\ref{prop:v25-real-direct-rotation}.
Therefore the two regulators determine the same real Pfaffian orientation,
the same zero-mode saturation sign and the same cofinal cutoff orientation.
\end{theorem}

\begin{proof}
V24 gives $\|p_{a,A}-p_{a,B}\|=O(a^{1/10})<1$.  Apply
Proposition~\ref{prop:v25-real-direct-rotation}.  The canonical kernel
orientation is preserved, and the positive reduced Pfaffian orientation is
intrinsic by Proposition~\ref{prop:v25-positive-reduced-Pf}.  Shell/cutoff
compatibility is Theorem~\ref{thm:v25-shell-orientation}.
\end{proof}

\subsection{Exact scope: when a real Pfaffian sign can still survive}
\label{sec:v25-scope}

The orientation theorem is strong but not universal.  Its hypotheses identify
precisely the physical mechanisms which could still generate a real sign.

\begin{theorem}[Real Pfaffian sign alternatives]
\label{thm:v25-sign-alternatives}
For an independently physical real skew fermion sector, a nontrivial V12
Pfaffian orientation character can survive only if at least one of the
following occurs:
\begin{enumerate}[label=\textup{(R\arabic*)},leftmargin=3em]
\item the real carrier is not globally orientable under the allowed native
      background identifications;
\item an allowed gauge/flavour/background transition acts with real determinant
      $-1$;
\item the path leaves the protected constant-rank component and crosses a
      Pfaffian divisor, with odd mod-two crossing number relative to the endpoint
      identification;
\item the purported ``Majorana'' term is not actually an independent real
      skew block, so the correct physical object is instead the already-treated
      complex determinant/Nambu line.
\end{enumerate}
On the natively oriented gauge-neutral branch of
Definition~\ref{def:v25-oriented-branch}, items \textup{(R1)}--\textup{(R2)}
are excluded; on the V11 protected reduced-gap component item \textup{(R3)} is
excluded as well.  Hence no real sign obstruction remains there.
\end{theorem}

\begin{proof}
If none of \textup{(R1)}--\textup{(R4)} occurs, the sector is a genuine real
skew family on an oriented carrier with orientation-preserving endpoint
identifications and constant reduced rank.  Theorem
\ref{thm:v25-kernel-orientation} orients its kernel line, Proposition
\ref{prop:v25-positive-reduced-Pf} orients the nonzero factor, and
Theorem~\ref{thm:v25-mod2-crossing} gives trivial loop parity.  This excludes a
nontrivial Pfaffian orientation character.
\end{proof}

\begin{firewall}[Charged pseudoreal sectors are different]
Theorem~\ref{thm:v25-sign-alternatives} does not trivialize the mod-two anomaly
of a charged real or pseudoreal chiral representation such as an isolated
$\mathrm{SU}(2)$ Weyl doublet.  Such a sector has additional gauge-bundle
structure and may carry a genuine five-dimensional mod-two index.  The V25
positive theorem is for the independently supplied \emph{gauge-neutral}
Majorana block with the oriented native carrier above.  The ordinary
Standard-Model charged chiral packet continues to use the anomaly-free complex
line theorem already proved in V3/V12.
\end{firewall}

\begin{firewall}[No crossing theorem for the continuum shell yet]
The parity formula \eqref{eq:v25-sf2-detg} identifies what happens to the real
sign if a path crosses a Pfaffian divisor.  The V1--V24 continuum shell theorem,
however, still assumes the reduced singular-value floor and therefore does not
transport through such a crossing.  A crossing requires a new local divisor
normal form and a source packet which retains the vanishing eigenpair rather
than a bounded pseudoinverse.
\end{firewall}

\section{V25 anti-circularity audit}
\label{sec:v25-audit}

\begin{enumerate}[leftmargin=2.4em]
\item \textbf{The complex determinant theorem was not reused as a real sign
proof.}  The Pfaffian orientation comes from the oriented real carrier and the
intrinsic symplectic orientation of the nonzero skew block.
\item \textbf{Kernel orientation is not chosen chart by chart.}  It is defined
by the global equation
$\mathfrak o_{\mathcal E}=\mathfrak o_K\wedge\mathfrak o_\omega$, so its Cech
transition signs are automatically positive.
\item \textbf{The mod-two loop statement includes the endpoint
identification.}  Equation \eqref{eq:v25-sf2-detg} shows explicitly that an
orientation-reversing native transition would produce odd parity rather than
being discarded.
\item \textbf{Shell factorization does not manufacture orientation.}  It only
propagates the already fixed real orientation because the hard Pfaffian stays
on one nonzero reference component.
\item \textbf{The optional Majorana term is not automatically declared an
independent Pfaffian sector.}  If the physical formulation uses the combined
complex/Nambu chiral block, the real-line analysis is not duplicated.
\item \textbf{Crossings remain separate.}  V25 computes their sign parity but
does not apply the bounded reduced inverse across a divisor.
\end{enumerate}

\section{V25 status ledger}
\label{sec:v25-status}

\subsection*{Proved in V25}
\begin{enumerate}[leftmargin=2.2em]
\item The canonical sitewise gauge-neutral Majorana realization carries a
cutoff-wise real orientation preserved by all connected spin/coframe
transitions and by lattice-site permutations because the real onsite Majorana
fibre is even dimensional.
\item On a constant-rank real skew family, the nonzero complement has a
canonical symplectic orientation.  Together with the ambient native
orientation this canonically orients the kernel bundle and forces
$w_1(K)=0$.
\item The inherited zero-mode Pfaffian line is therefore globally orientable on
the protected real branch; dual convention choices have the same
Stiefel--Whitney class.
\item The reduced nonzero Pfaffian is positive in its intrinsic symplectic
orientation.  A positive kernel saturation writer therefore gives a globally
well-defined saturated real Majorana amplitude on the protected branch.
\item For a generic real skew path with endpoint congruence
$A(1)=g^TA(0)g$, the mod-two Pfaffian spectral flow obeys
$(-1)^{\mathrm{SF}_2}=\operatorname{sgn}\det_{\mathbb R}g$.  Every loop whose
native endpoint identification preserves the carrier orientation has even
mod-two spectral flow.
\item Consequently the V12 real Pfaffian orientation character vanishes on the
natively oriented gauge-neutral Majorana branch.
\item The inherited exact Pfaffian Schur/cutoff cocycle transports this
orientation cofinally because every protected hard Pfaffian remains on one
positive reference component.
\item The V24 canonical real direct rotation lies in the orientation-preserving
orthogonal component.  Hence all smooth protected V2-consistent microscopic
completions of the same oriented neutral-Majorana branch carry the same
Pfaffian orientation.
\end{enumerate}

\subsection*{Inherited}
\begin{enumerate}[leftmargin=2.2em]
\item V1--V12: finite determinant/Pfaffian factorization, zero-mode
saturation, source-line calculus and the statement that an independent real
Pfaffian line carries a possible $\mathbb Z_2$ orientation character.
\item V13--V23: local Riesz/Kato control, controlled zero-mode module
classification, regulator-path descent and the soft Feshbach criterion.
\item V24: the smooth-background G\aa rding/Feshbach theorem and norm-close
soft Riesz projectors for every V2-consistent microscopic completion.
\end{enumerate}

\subsection*{Conditional}
\begin{enumerate}[leftmargin=2.2em]
\item V25 applies only when an independently physical real skew Majorana
realization is actually supplied.  A Majorana mass written inside the combined
complex/Nambu chiral block does not create a second real sign problem.
\item The positive theorem assumes the native real carrier is orientable and
all allowed gauge, spin, flavour, and background transitions preserve that
orientation.  An orientation-reversing physical symmetry must be analyzed
separately.
\item The continuum shell still stays on a constant-rank reduced-gap component.
V25 identifies but does not cross a Pfaffian divisor.
\item The V24 smooth-background and positive IR/reference hypotheses remain in
force for the regulator-comparison statement.
\end{enumerate}

\subsection*{Open beyond V25}
\begin{enumerate}[leftmargin=2.2em]
\item Construct a local normal form for a simple kernel--cokernel/Pfaffian
crossing and enlarge the source packet so the vanishing mode is retained rather
than inverted.  Prove the corresponding shell/divisor factorization across the
crossing.
\item Use that normal form to treat creation/annihilation of an accidental
zero-mode pair and then a genuine change of Fredholm index, with spectral-flow
and determinant/Pfaffian divisor data transported together.
\item If an independently physical charged real/pseudoreal fermion sector is
ever introduced, compute its own mod-two bordism character; V25 does not infer
its sign from the neutral theorem.
\item The independent positive-IR removal, analytic finite-coupling RG,
perturbative summability and nonperturbative completion problems remain open.
\end{enumerate}

\subsection*{Exact next load-bearing theorem}
\begin{center}
\fbox{\parbox{0.92\textwidth}{
\textbf{Simple divisor crossing / spectral-flow source-packet theorem.}
Let a protected real or complex chiral family approach a single regular
zero-mode crossing while all complementary singular values remain uniformly
bounded away from zero.  Construct a finite native normal form which splits the
crossing eigenpair from the reduced complement, replace the divergent inverse
by an explicit rank-one/rank-two divisor coordinate, and prove an exact shell
Schur/Pfaffian factorization with all source derivatives and BV contacts through
the crossing.  The target is a coefficientwise continuum comparison on each
side plus a controlled transition law for the determinant/Pfaffian line and
spectral-flow sign.  This is the smallest remaining gate before true
index-changing families can be added to the otherwise closed smooth protected
regulator programme.}}
\end{center}

\section*{Append-only continuation marker: V25}
\addcontentsline{toc}{section}{Append-only continuation marker: V25}
\textbf{End of V25.}  The complete V1--V24 body above is retained verbatim.
V25 closes the independent real Pfaffian orientation problem on a physically
supplied, natively oriented gauge-neutral Majorana branch.  A nondegenerate
real skew form canonically orients the reduced complement; the fixed native
carrier orientation then canonically orients the zero-mode bundle, so the
Pfaffian $w_1$ vanishes.  For a generic path, the parity of Pfaffian spectral
flow is exactly the orientation character of the endpoint identification;
all native neutral-Majorana identifications have determinant $+1$, hence every
closed loop has even parity.  Exact Pfaffian Schur factorization transports the
orientation cofinally, and V24's real direct rotation preserves it between
V2-consistent regulators.  The next gate is no longer orientation but a true
zero-mode divisor crossing, where the reduced inverse ceases to exist and the
vanishing mode must be retained explicitly.



\clearpage
\section{V26 continuation: simple divisor crossings, saturated source packets, and regular shell transport}
\label{sec:v26-continuation}

V25 closed the independent real orientation problem on a protected
constant-rank branch.  The next failure mechanism is analytically different:
a nonzero singular value can approach zero, so the V11 Moore--Penrose inverse
ceases to have a uniform bound.  The wrong response is to estimate that inverse
more aggressively.  The correct upstream variable is the vanishing finite
crossing block itself.

This continuation treats a \emph{simple accidental divisor crossing}: after
any pre-existing V11 kernel/cokernel modes have first been factored and
saturated, exactly one additional complex right/left mode becomes singular, or
one oriented real two-dimensional skew block becomes singular, while every
complementary singular value stays uniformly bounded away from zero.  The
crossing mode is retained among the soft/source variables and is never
integrated into a hard shell.

The main result is that the apparent pole is entirely an artifact of the vacuum
trivialization.  Exact Schur reduction gives the regular crossing packets
\begin{equation}
 \boxed{\mathcal P_\times^{\mathbb C}
       =\sigma+\widehat{\bar\jmath}\,\widehat\jmath,}
 \qquad
 \boxed{\mathcal P_\times^{\mathbb R}
       =\sigma+\widehat\jmath_1\widehat\jmath_2.}
 \label{eq:v26-crossing-packets}
\end{equation}
Every source/BV derivative of these packets is finite at $\sigma=0$.  Away from
the divisor, choosing the vacuum coefficient instead of the crossing-saturated
coefficient multiplies the line section by $\sigma$; its logarithmic derivative
$d\sigma/\sigma$ is therefore a line-coordinate pole, not a new ultraviolet
divergence.

The crossing treated in V26 does \emph{not} change the Fredholm index: at the
divisor one kernel and one cokernel mode are created simultaneously in the
complex case.  A true overlap-index change requires the Wilson/sign projector
itself to change rank and remains the next stronger gate.

\subsection{Regular complex divisor chart}
\label{sec:v26-complex-chart}

Let $D(\zeta):E\to F$ be a smooth finite-cutoff complex chiral/Yukawa family on
a compact parameter chart.  Factor first the pre-existing V11 protected
kernel/cokernel bundles and work on their reduced square complement.  We call
$\zeta_0$ a \emph{simple complex divisor point} if there are smooth rank-one
Riesz projectors
\begin{equation}
 P_\times(\zeta):E\to E,
 \qquad
 Q_\times(\zeta):F\to F
 \label{eq:v26-cross-projectors}
\end{equation}
for the unique small singular cluster, while all complementary singular values
obey
\begin{equation}
 s_{\min}\!\left(
 Q_\times^\perp D P_\times^\perp
 \right)\ge\tau_*>0.
 \label{eq:v26-complement-floor}
\end{equation}
At $\zeta_0$ the reduced operator has one additional kernel and one additional
cokernel dimension.  Along a declared one-parameter crossing $t\mapsto\zeta(t)$
we additionally require regularity of the scalar divisor coordinate constructed
below:
\begin{equation}
 \sigma(0)=0,
 \qquad
 \dot\sigma(0)\ne0.
 \label{eq:v26-regular-complex-crossing}
\end{equation}
For a multi-parameter complex divisor one replaces
\eqref{eq:v26-regular-complex-crossing} by the corresponding regular-zero
condition in the chosen local category.

Choose local unit Kato frames $u$ and $v$ for the crossing domain and target
lines.  With
\begin{equation}
 E=\mathbb Cu\oplus E_r,
 \qquad
 F=\mathbb Cv\oplus F_r,
 \label{eq:v26-complex-split}
\end{equation}
write
\begin{equation}
 D=
 \begin{pmatrix}
 d&b\\
 c&H
 \end{pmatrix},
 \qquad
 H:E_r\longrightarrow F_r,
 \label{eq:v26-complex-block}
\end{equation}
where $H$ is uniformly invertible by
\eqref{eq:v26-complement-floor}.

\begin{definition}[Dressed divisor coordinate]
\label{def:v26-sigma-complex}
The exact complex crossing coordinate is
\begin{equation}
 \boxed{\sigma:=d-bH^{-1}c.}
 \label{eq:v26-sigma-complex}
\end{equation}
It is naturally a local section of the crossing line
$\operatorname{Hom}(\mathbb Cu,\mathbb Cv)$; a change of Kato frames rescales
$\sigma$ by the corresponding nonzero endpoint line factor, so its zero locus
is frame independent.
\end{definition}

\begin{lemma}[Exact triangular diagonalization]
\label{lem:v26-complex-diagonal}
Let
\begin{equation}
 L=
 \begin{pmatrix}
 1&-bH^{-1}\\
 0&I
 \end{pmatrix},
 \qquad
 R=
 \begin{pmatrix}
 1&0\\
 -H^{-1}c&I
 \end{pmatrix}.
 \label{eq:v26-LR}
\end{equation}
Then
\begin{equation}
 \boxed{
 LDR=
 \begin{pmatrix}
 \sigma&0\\
 0&H
 \end{pmatrix}.}
 \label{eq:v26-complex-diagonal}
\end{equation}
Both triangular factors have determinant one.  Therefore
\begin{equation}
 \boxed{\det D=\det H\,\sigma}
 \label{eq:v26-det-divisor}
\end{equation}
in the chosen determinant-line frames.
\end{lemma}

\begin{proof}
Multiply $D$ by $R$ on the right:
\[
 DR=
 \begin{pmatrix}
 d-bH^{-1}c&b\\
 0&H
 \end{pmatrix}.
\]
Left multiplication by $L$ removes the remaining upper-right block, giving
\eqref{eq:v26-complex-diagonal}.  The determinant statement follows because
$L$ and $R$ are block triangular with identity diagonal.
\end{proof}

\begin{corollary}[Pole decomposition away from the divisor]
\label{cor:v26-pole-decomposition}
For $\sigma\ne0$,
\begin{equation}
 D^{-1}
 =R
 \begin{pmatrix}
 \sigma^{-1}&0\\
 0&H^{-1}
 \end{pmatrix}
 L.
 \label{eq:v26-inverse-factor}
\end{equation}
Equivalently, the only singular part is rank one:
\begin{equation}
 \boxed{
 D^{-1}
 =G_{\rm reg}
 +\sigma^{-1}
 \begin{pmatrix}1\\-H^{-1}c\end{pmatrix}
 \begin{pmatrix}1&-bH^{-1}\end{pmatrix},}
 \label{eq:v26-rank-one-pole}
\end{equation}
where $G_{\rm reg}$ is smooth at $\sigma=0$.
\end{corollary}

The purpose of V26 is precisely to avoid using
\eqref{eq:v26-rank-one-pole} at $\sigma=0$.

\subsection{Exact complex crossing-source packet}
\label{sec:v26-complex-source}

Use the finite Grassmann convention already fixed by the inherited determinant
identity
\begin{equation}
 Z_D[\bar J,J]
 :=\int d\bar\psi\,d\psi\,
 e^{-\bar\psi D\psi+\bar J\psi+\bar\psi J}
 =\det D\,e^{\bar J D^{-1}J}
 \label{eq:v26-complex-gaussian-convention}
\end{equation}
when $D$ is invertible.  Split the fields and sources according to
\eqref{eq:v26-complex-split}:
\begin{equation}
 \psi=(\eta,\chi),
 \quad
 \bar\psi=(\bar\eta,\bar\chi),
 \qquad
 J=(\jmath,J_r),
 \quad
 \bar J=(\bar\jmath,\bar J_r).
 \label{eq:v26-complex-field-split}
\end{equation}

\begin{definition}[Dressed crossing sources]
\label{def:v26-dressed-complex-sources}
Define
\begin{align}
 \widehat{\bar\jmath}
 &:=\bar\jmath-\bar J_rH^{-1}c,
 \label{eq:v26-barjhat}\\
 \widehat\jmath
 &:=\jmath-bH^{-1}J_r,
 \label{eq:v26-jhat}\\
 \mathcal C_r
 &:=\bar J_rH^{-1}J_r.
 \label{eq:v26-Cr}
\end{align}
\end{definition}

\begin{theorem}[Exact regular complex divisor packet]
\label{thm:v26-complex-packet}
Without assuming $\sigma\ne0$, exact integration of the uniformly invertible
complement gives
\begin{equation}
 \boxed{
 Z_D[\bar J,J]
 =\det H\,e^{\mathcal C_r}
 \bigl(
 \sigma+\widehat{\bar\jmath}\,\widehat\jmath
 \bigr).}
 \label{eq:v26-complex-packet}
\end{equation}
Thus the complete crossing contribution is polynomial and regular at the
divisor.  At $\sigma=0$ the vacuum coefficient vanishes, while the coefficient
of $\widehat{\bar\jmath}\widehat\jmath$ remains exactly one in the crossing
factor.
\end{theorem}

\begin{proof}
Integrate only the $(\bar\chi,\chi)$ variables.  Completing the finite
Grassmann Gaussian gives the reduced contact
$\bar J_rH^{-1}J_r$, the dressed sources of
Definition~\ref{def:v26-dressed-complex-sources}, and the remaining one-pair
integral
\[
 \int d\bar\eta\,d\eta\,
 e^{-\bar\eta\sigma\eta
   +\widehat{\bar\jmath}\eta
   +\bar\eta\widehat\jmath}.
\]
With the inherited Berezin orientation for which the source-free integral is
$\sigma$, direct expansion of the two Grassmann variables gives
$\sigma+\widehat{\bar\jmath}\widehat\jmath$.
No inverse of $\sigma$ is introduced.
\end{proof}

\begin{definition}[Crossing-saturation writer]
\label{def:v26-complex-writer}
Let $\omega_\times^{\mathbb C}$ be the dual writer which extracts the
coefficient of
$\widehat{\bar\jmath}\widehat\jmath$ from the crossing packet.  Then
\begin{equation}
 \boxed{
 \left\langle
 \omega_\times^{\mathbb C},
 \sigma+\widehat{\bar\jmath}\widehat\jmath
 \right\rangle=1}
 \label{eq:v26-complex-writer}
\end{equation}
for every value of $\sigma$, including the divisor.
\end{definition}

If pre-existing V11 zero modes are present, tensor
$\omega_\times^{\mathbb C}$ with the inherited zero-mode saturation writer.
The resulting scalarized amplitude is smooth across the accidental crossing.

\subsection{Vacuum versus saturated trivializations}
\label{sec:v26-trivializations}

Away from the divisor there are two useful local line trivializations.  The
vacuum writer extracts the crossing-source degree zero coefficient; the
crossing-saturated writer extracts degree two.

\begin{proposition}[Exact divisor transition]
\label{prop:v26-divisor-transition}
On the punctured chart $\{\sigma\ne0\}$, the vacuum and saturated crossing
trivializations differ by multiplication with $\sigma$.  If all other source
writers are denoted collectively by $J$, then the normalized functionals obey
\begin{equation}
 \boxed{
 W_{\rm vac}[J]
 =W_{\rm sat}[J]
 +\log\frac{\sigma(J)}{\sigma(0)} }
 \label{eq:v26-W-transition}
\end{equation}
whenever the logarithm is defined in the chosen local line frame.  Hence every
pole in a differentiated vacuum expression is generated by derivatives of
$\log\sigma$ and is a divisor-line coordinate effect.  The saturated
functional is regular at $\sigma=0$.
\end{proposition}

\begin{proof}
At vanishing crossing sources,
Theorem~\ref{thm:v26-complex-packet} gives the vacuum scalar factor
$\det H\,e^{\mathcal C_r}\sigma$, while pairing with
$\omega_\times^{\mathbb C}$ gives
$\det H\,e^{\mathcal C_r}$.  Divide each by its value at the reference source
and take the logarithm.  The difference is exactly
\eqref{eq:v26-W-transition}.
\end{proof}

\begin{firewall}[A divisor pole is not a UV divergence]
Terms such as $d\sigma/\sigma$ in the vacuum frame are singular because that
frame vanishes on the divisor.  They are not hard-shell divergences and must
not be cancelled by a regulator-dependent counterterm.  The regular object is
the line-valued generating section, or equivalently the crossing-saturated
scalarization of Definition~\ref{def:v26-complex-writer}.
\end{firewall}

\subsection{The real skew/Pfaffian normal form}
\label{sec:v26-real-crossing}

Now let $A(\zeta)$ be a smooth real skew family on an oriented real carrier.
Assume one isolated real skew eigenpair becomes soft while every complementary
singular value is bounded below by $\tau_*>0$.  The crossing Riesz bundle has
real rank two.  Choose an oriented orthonormal frame of that bundle and write
\begin{equation}
 A=
 \begin{pmatrix}
 a&b\\
 -b^T&H
 \end{pmatrix},
 \qquad
 H^T=-H,
 \label{eq:v26-real-block}
\end{equation}
where $H$ is invertible and $a$ is $2\times2$ skew.

Let
\begin{equation}
 J_2:=
 \begin{pmatrix}
 0&1\\
 -1&0
 \end{pmatrix}.
 \label{eq:v26-J2}
\end{equation}

\begin{definition}[Real Pfaffian divisor coordinate]
\label{def:v26-real-sigma}
The skew Schur complement is
\begin{equation}
 S_\times:=a+bH^{-1}b^T.
 \label{eq:v26-real-Schur}
\end{equation}
Because it is a real $2\times2$ skew matrix, there is a unique real scalar
$\sigma$ such that
\begin{equation}
 \boxed{S_\times=\sigma J_2.}
 \label{eq:v26-real-sigma}
\end{equation}
A regular real crossing satisfies
$\sigma(0)=0$ and $\dot\sigma(0)\ne0$.
\end{definition}

\begin{lemma}[Exact skew congruence reduction]
\label{lem:v26-real-diagonal}
With
\begin{equation}
 T:=
 \begin{pmatrix}
 I&0\\
 H^{-1}b^T&I
 \end{pmatrix},
 \label{eq:v26-real-T}
\end{equation}
one has
\begin{equation}
 \boxed{
 T^TAT=
 \begin{pmatrix}
 \sigma J_2&0\\
 0&H
 \end{pmatrix}.}
 \label{eq:v26-real-diagonal}
\end{equation}
In the orientation inherited from the ordered crossing/complement splitting,
\begin{equation}
 \boxed{
 \operatorname{Pf}(A)
 =\operatorname{Pf}(H)\,\sigma.}
 \label{eq:v26-Pf-divisor}
\end{equation}
\end{lemma}

\begin{proof}
First compute
\[
 AT=
 \begin{pmatrix}
 a+bH^{-1}b^T&b\\
 0&H
 \end{pmatrix}.
\]
Since $H^{-T}=-H^{-1}$,
$T^T=\begin{psmallmatrix}I&-bH^{-1}\\0&I\end{psmallmatrix}$, which removes the
upper-right block and gives
\eqref{eq:v26-real-diagonal}.  The triangular matrix $T$ has determinant one,
so the Pfaffian congruence law yields
\eqref{eq:v26-Pf-divisor}.
\end{proof}

Let the real Grassmann source split as
$j=(j_\times,j_r)$ with $j_\times=(j_1,j_2)^T$ and define
\begin{equation}
 \widehat j_\times
 :=j_\times-bH^{-1}j_r,
 \qquad
 \mathcal C_r^{\mathbb R}
 :=\frac12j_r^TH^{-1}j_r.
 \label{eq:v26-real-dressed-source}
\end{equation}

\begin{theorem}[Exact regular real Pfaffian packet]
\label{thm:v26-real-packet}
Choose the ordered Berezin measure so that
$\operatorname{Pf}(\sigma J_2)=\sigma$.  Exact integration of the invertible
real complement gives
\begin{equation}
 \boxed{
 Z_A[j]
 =\operatorname{Pf}(H)
  e^{\mathcal C_r^{\mathbb R}}
 \bigl(
 \sigma+\widehat j_1\widehat j_2
 \bigr).}
 \label{eq:v26-real-packet}
\end{equation}
The writer extracting the coefficient of
$\widehat j_1\widehat j_2$ is smooth and nonzero at the divisor.  The vacuum
Pfaffian changes sign exactly when the regular real coordinate $\sigma$ changes
sign.
\end{theorem}

\begin{proof}
Use the change of variables associated with the congruence
\eqref{eq:v26-real-diagonal}.  Its determinant is one, so the Berezin measure is
unchanged.  The complement contributes
$\operatorname{Pf}(H)e^{\mathcal C_r^{\mathbb R}}$ and the remaining oriented
two-variable Grassmann integral is
$\sigma+\widehat j_1\widehat j_2$ by direct expansion.
\end{proof}

\begin{corollary}[Local spectral-flow sign]
\label{cor:v26-local-spectral-flow}
For a regular one-parameter real crossing,
\begin{equation}
 \boxed{
 \operatorname{sgn}
 \frac{\operatorname{Pf}(A(t_+))}
      {\operatorname{Pf}(A(t_-))}
 =-1}
 \label{eq:v26-Pf-sign-flip}
\end{equation}
for $t_-<0<t_+$ sufficiently small.  Thus the local mod-two spectral-flow
contribution is one, in agreement with the V25 loop formula.  The
crossing-saturated real amplitude remains continuous.
\end{corollary}

\subsection{All source/BV descendants remain regular}
\label{sec:v26-source-bv}

The key point for the ESM continuum programme is not merely the scalar
factorization but the full fixed source packet.

\begin{theorem}[Source-complete divisor regularity]
\label{thm:v26-source-complete}
Fix a finite coefficient packet
$\mathfrak q=(L,\Delta,r,N,\mathcal P_{\rm ext})$.  On a simple crossing chart,
every retained source derivative of the following objects is uniformly bounded
at the divisor:
\begin{enumerate}[label=\textup{(S\arabic*)},leftmargin=3em]
\item the crossing Riesz projectors and their Kato frames;
\item the complement inverse $H^{-1}$;
\item the dressed divisor coordinate $\sigma$;
\item the dressed crossing sources and reduced contact term;
\item the complex or real polynomial packet in
      \eqref{eq:v26-crossing-packets};
\item the crossing-saturated determinant/Pfaffian line writer.
\end{enumerate}
No retained coefficient contains $\sigma^{-1}$.  All apparent negative powers
of $\sigma$ arise only after changing to the vacuum trivialization on
$\{\sigma\ne0\}$.
\end{theorem}

\begin{proof}
The crossing Riesz projectors are defined by contours separating the one small
singular cluster from the complementary spectrum; the contour does not collapse
when the enclosed eigenvalue reaches zero.  The V11/V13 Riesz calculus
therefore gives all fixed source derivatives.  Kato transport gives the local
frames.  The complementary floor
\eqref{eq:v26-complement-floor} gives bounded source derivatives of $H^{-1}$ by
the ordinary inverse rule.  Equations
\eqref{eq:v26-sigma-complex},
\eqref{eq:v26-barjhat}--\eqref{eq:v26-Cr}, and
\eqref{eq:v26-real-dressed-source} are finite algebraic combinations of those
bounded jets.  Finally the packets
\eqref{eq:v26-crossing-packets} are polynomial in the dressed variables.  No
step divides by $\sigma$.
\end{proof}

\begin{theorem}[BV pushforward with the crossing block retained]
\label{thm:v26-BV-crossing}
Use the inherited exact BV high/low split but place the entire crossing Riesz
block among the retained low variables.  Integrating any hard complement then
preserves the same finite quantum master/source identity as in V3--V5, now in
the enlarged state containing the crossing Grassmann writers.  The local
master vector and all stored future contraction descendants are regular at
$\sigma=0$.
\end{theorem}

\begin{proof}
The inherited BV pushforward theorem requires a canonical high/low split and
exact integration of the high Lagrangian variables.  It does not require the
retained low quadratic form to be invertible.  The only inverse used in the
present pushforward is the uniformly invertible hard/reduced complement.  The
finite crossing Grassmann block remains explicit, and
Theorems~\ref{thm:v26-complex-packet} and
\ref{thm:v26-real-packet} show that its source dependence is polynomial.
Therefore the same algebraic BV-Stokes argument preserves the master identity.
V5 contraction completeness stores every local descendant generated by later
BV/source contractions, including derivatives of the new crossing writers.
\end{proof}

\subsection{Shell composition through the divisor}
\label{sec:v26-shell-composition}

Suppose the invertible complement itself is split into retained noncrossing soft
modes and one hard shell.  There are then two possible operations: eliminate
the full complement at once, or eliminate the hard shell first and then the
remaining complement.

\begin{theorem}[Associativity of the divisor Schur coordinate]
\label{thm:v26-Schur-associativity}
For every finite nested shell decomposition which leaves the crossing block
unintegrated, the final complex scalar $\sigma$ of
\eqref{eq:v26-sigma-complex} and the final real skew scalar $\sigma$ of
\eqref{eq:v26-real-sigma} are independent of the order in which invertible
complement blocks are Schur-eliminated.  The dressed crossing sources and the
reduced source contact obey the same associativity.  Consequently the exact
line/source factors satisfy one shell cocycle:
\begin{align}
 \det D
 &=\left(\prod_j\det H_j\right)\sigma,
 \label{eq:v26-shell-det}\\
 \operatorname{Pf}(A)
 &=\left(\prod_j\operatorname{Pf}(H_j)\right)\sigma
 \label{eq:v26-shell-Pf}
\end{align}
with the declared complex or real orientations.
\end{theorem}

\begin{proof}
Finite Gaussian elimination is associative.  Algebraically, the Schur
complement of a union of two invertible blocks equals the Schur complement of
the second block in the Schur complement of the first.  Applying that identity
to the one-dimensional complex crossing block gives the same scalar
$d-bH^{-1}c$ independent of grouping.  Applying the skew congruence analogue to
the oriented two-dimensional real block gives the same scalar coefficient of
$J_2$.  The source formulas follow from the same triangular changes of
variables, and the determinant/Pfaffian products follow by multiplying the
exact shell factors.
\end{proof}

\begin{corollary}[No shell-dependent refitting of the divisor]
\label{cor:v26-no-refit}
The crossing coordinate and its saturation writer are common finite-state data.
They are not independently normalized on different shells.  Mixed-shell source
contacts generated by differentiating
\eqref{eq:v26-shell-det}--\eqref{eq:v26-shell-Pf} are retained by the same
Bell-polynomial/source-contact calculus as in V1--V7.
\end{corollary}

\subsection{Coefficientwise continuum across a simple accidental crossing}
\label{sec:v26-continuum}

The old normalized vacuum functional cannot be evaluated exactly at a divisor
because its source-free fermion factor vanishes.  The correct endpoint is the
crossing-saturated functional.

\begin{definition}[Divisor-saturated normalized functional]
\label{def:v26-Wsat}
Let $\omega_Z$ denote any pre-existing V11 zero-mode saturation writer and let
$\omega_\times$ be the complex or real crossing writer above.  Pair the
line-valued finite generating section with
$\omega_Z\otimes\omega_\times$ and denote the resulting scalar amplitude by
$\mathcal A^{\rm sat}[J]$.  On a protected reference point where this scalar is
nonzero, set
\begin{equation}
 \boxed{
 W^{\rm sat}[J]
 :=\log
 \frac{\mathcal A^{\rm sat}[J]}
      {\mathcal A^{\rm sat}[0]}.}
 \label{eq:v26-Wsat}
\end{equation}
\end{definition}

\begin{theorem}[Fixed-order continuum through a simple accidental divisor]
\label{thm:v26-continuum-crossing}
Assume the V1--V24 fixed-order native hypotheses on all hard shells and, on one
soft crossing chart, the simple complex or real divisor hypotheses above.  In
addition assume the V24 common soft chiral consistency on the finite crossing
cluster, so that the crossing projectors, the reduced complement blocks and
their prescribed source jets for any two V2-consistent microscopic completions
converge to the same continuum crossing data.

Then for every prescribed finite
$(L,\Delta,r,N,\mathcal P_{\rm ext})$:
\begin{enumerate}[label=\textup{(C\arabic*)},leftmargin=3em]
\item the shell state enlarged by the finite crossing-source packet satisfies
the same rooted/quasilocal hard-shell estimates as V5--V7;
\item all coefficients of $W^{\rm sat}$ have cutoff-independent continuum
limits uniformly across the simple divisor chart;
\item the complete retained BV/BRST/master/source vector passes to the same
limit and remains zero;
\item on each side of the divisor, where $\sigma\ne0$, the ordinary vacuum
trivialization is recovered by the exact transition
\eqref{eq:v26-W-transition};
\item in the real skew case, the two vacuum Pfaffian signs differ by the local
spectral-flow factor of Corollary~\ref{cor:v26-local-spectral-flow}, while the
saturated continuum is continuous.
\end{enumerate}
\end{theorem}

\begin{proof}
The crossing block is finite dimensional and is always retained on the soft
side.  Hence every hard internal fermion contraction still uses only uniformly
invertible hard/reduced propagators, so the V5 rooted forest estimates and V6
Cauchy sums are unchanged up to finite constants.  The enlarged local state has
only finitely many additional source coordinates, so contraction completeness
and the common-action induction remain finite at the prescribed packet.

Theorem~\ref{thm:v26-source-complete} makes every crossing coefficient regular
at $\sigma=0$.  The V24 common soft consistency identifies the finite crossing
Riesz bundles and all block/source coefficients of any two V2-consistent native
completions; bounded inversion of $H$ then identifies $\sigma$, the dressed
sources and the reduced contacts.  Thus the same multi-shell Cauchy theorem
applies coefficientwise to the saturated scalarization.  The BV statement is
Theorem~\ref{thm:v26-BV-crossing}.  Items \textup{(C4)}--\textup{(C5)} are the
exact line-transition and real-sign formulas already proved above.
\end{proof}

\begin{firewall}[The ordinary connected vacuum functional does not exist on the divisor]
At $\sigma=0$ the source-free fermion integral has an unsaturated zero mode, so
$\log(Z[J]/Z[0])$ in the vacuum trivialization is undefined.  V26 does not hide
this by dividing out the zero.  The globally regular object is the line-valued
section or the explicitly saturated functional
\eqref{eq:v26-Wsat}.  The ordinary $W$ is recovered only on the open sets where
the vacuum section is nonzero.
\end{firewall}

\subsection{Accidental pair creation is not an index change}
\label{sec:v26-index-firewall}

\begin{proposition}[Index is unchanged by a simple complex divisor]
\label{prop:v26-index-unchanged}
In the complex normal form of V26, at $\sigma=0$ one additional right kernel
mode and one additional left cokernel mode appear simultaneously.  Therefore
\begin{equation}
 \boxed{
 \Delta\bigl(\dim\ker D-\dim\operatorname{coker}D\bigr)=0.}
 \label{eq:v26-index-unchanged}
\end{equation}
The Fredholm index on the two punctured sides is the same.
\end{proposition}

\begin{proof}
At $\sigma=0$, Lemma~\ref{lem:v26-complex-diagonal} reduces the crossing block
to the zero map $\mathbb C\to\mathbb C$ while $H$ stays invertible.  Hence both
kernel and cokernel dimensions increase by one and their difference is
unchanged.
\end{proof}

\begin{firewall}[True overlap-index change is a different singularity]
A genuine change of the overlap/Fredholm index requires the dimension of the
chiral domain or target bundle to change.  In the native overlap regulator this
occurs when the Wilson/sign projector changes rank, which necessarily leaves
the protected Wilson-gap component used by V1--V26.  The simple divisor theorem
above therefore does not solve an index-changing Wilson spectral crossing.
\end{firewall}

\section{V26 anti-circularity audit}
\label{sec:v26-audit}

\begin{enumerate}[leftmargin=2.4em]
\item \textbf{The divergent pseudoinverse is never used at the divisor.}
Only the complement inverse $H^{-1}$ appears; the crossing block remains an
explicit finite Grassmann packet.
\item \textbf{The zero of the vacuum determinant/Pfaffian is retained.}
It is the divisor coordinate $\sigma$, not a quantity silently divided out or
cancelled by a counterterm.
\item \textbf{Source singularities are identified as frame singularities.}
The pole $d\sigma/\sigma$ appears only after choosing the vacuum line
trivialization; the saturated packet is polynomial.
\item \textbf{BV restoration is not recomputed graph by graph.}
The crossing variables are added to the common low/source state and the exact
inherited BV pushforward is applied to the uniformly invertible complement.
\item \textbf{Shell regrouping cannot redefine the crossing coordinate.}
Associativity of finite Schur elimination proves one common divisor/source
cocycle.
\item \textbf{Accidental pair creation is not mislabeled as index change.}
The complex crossing changes kernel and cokernel dimensions equally; a true
index jump remains outside the protected Wilson-gap component.
\end{enumerate}

\section{V26 status ledger}
\label{sec:v26-status}

\subsection*{Proved in V26}
\begin{enumerate}[leftmargin=2.2em]
\item A simple complex accidental crossing admits an exact rank-one Schur normal
form with divisor coordinate $\sigma=d-bH^{-1}c$ and determinant factorization
$\det D=\det H\,\sigma$.
\item Exact complement integration yields the regular complex source packet
$\sigma+\widehat{\bar\jmath}\widehat\jmath$.  A crossing-saturation writer
extracts its top source coefficient and remains nonzero on the divisor.
\item The vacuum and saturated line trivializations differ by $\sigma$; all
$1/\sigma$ and $d\sigma/\sigma$ terms therefore belong to the vacuum divisor
frame rather than to the hard-shell ultraviolet remainder.
\item A simple real skew crossing admits an exact oriented $2\times2$ Pfaffian
normal form with skew Schur complement $\sigma J_2$ and
$\operatorname{Pf}(A)=\operatorname{Pf}(H)\sigma$.  Its regular source packet is
$\sigma+\widehat j_1\widehat j_2$.
\item The real vacuum Pfaffian flips sign at a regular crossing, giving one unit
of mod-two local spectral flow, while the crossing-saturated real amplitude is
continuous.
\item Every retained source derivative of the crossing projectors, complement
inverse, divisor coordinate, dressed sources, contacts and saturated line
writer is bounded at the divisor.  No retained coefficient contains
$\sigma^{-1}$.
\item Keeping the crossing block among the low variables preserves the exact
BV/master/source identity under hard integration.  The enlarged finite source
packet is compatible with the V5 contraction-complete induction.
\item The divisor coordinate, dressed sources and determinant/Pfaffian factors
compose exactly under nested shell elimination by associativity of Schur
reduction.
\item Under the V24 common soft chiral-consistency hypothesis on the finite
crossing cluster, the divisor-saturated fixed-order continuum exists uniformly
through the crossing and is regulator independent at the same coefficientwise
scope.
\item A simple complex divisor creates one kernel--cokernel pair and leaves the
Fredholm index unchanged.
\end{enumerate}

\subsection*{Inherited}
\begin{enumerate}[leftmargin=2.2em]
\item V1--V7: source-complete coefficientwise continuum, common BV restoration,
forest induction, UV Cauchy convergence, scheme/reference covariance and exact
loop-order projectivity.
\item V8--V24: regulator comparison, controlled zero-mode classification,
soft spectral anchoring and the V2-consistent covariant G\aa rding/Feshbach
theorem.
\item V25: native orientation of an independently physical gauge-neutral real
Majorana/Pfaffian branch and the mod-two sign formula for regular real
crossings.
\end{enumerate}

\subsection*{Conditional}
\begin{enumerate}[leftmargin=2.2em]
\item The crossing must be simple and isolated: one complex right/left pair or
one oriented real skew two-plane, with a cutoff-uniform complementary singular
floor.
\item The fixed-order continuum comparison across the divisor uses the V24
smooth-background/V2-consistent soft chiral jet on the finite crossing cluster.
An $O(1)$ regulator-dependent soft defect is not covered.
\item The regular scalarization at the divisor requires the declared crossing
saturation writer.  The unsaturated vacuum partition function vanishes there.
\item The real theorem assumes the V25 oriented neutral-Majorana carrier; an
independent charged pseudoreal sector retains its own mod-two anomaly problem.
\item The positive IR/reference assumptions of the coefficientwise continuum
remain in force.
\end{enumerate}

\subsection*{Open beyond V26}
\begin{enumerate}[leftmargin=2.2em]
\item Treat a genuine overlap-index-changing crossing where the Wilson kernel
loses its protected gap and the sign projector changes rank.  This requires a
normal form for the Wilson spectral crossing itself, not merely a chiral
kernel--cokernel divisor.
\item Extend the divisor packet to simultaneous higher-multiplicity crossings
by a finite matrix-valued crossing block and identify the corresponding
stratified determinant/Pfaffian divisor.
\item Remove the positive IR/reference scale, and separately address analytic
finite-coupling RG trajectories, perturbative summability and nonperturbative
completion.
\end{enumerate}

\subsection*{Exact next load-bearing theorem}
\begin{center}
\fbox{\parbox{0.92\textwidth}{
\textbf{Native Wilson-sign crossing / genuine index-jump theorem.}
Let a smooth native Hermitian Wilson family develop one regular eigenvalue
crossing through zero while all other Wilson eigenvalues remain uniformly
gapped.  Construct the exact rank-one Wilson crossing projector, replace the
discontinuous sign by a two-sided oriented crossing chart, and determine the
induced change of the overlap chiral domain, determinant/Pfaffian line and
Fredholm index.  The target is a source-complete shell theorem which compares
the coefficientwise continua on the two adjacent index sectors and supplies an
explicit transition/spectral-flow law at the crossing, without pretending that
the overlap sign operator itself is continuous through the Wilson-gap zero.}}
\end{center}

\section*{Append-only continuation marker: V26}
\addcontentsline{toc}{section}{Append-only continuation marker: V26}
\textbf{End of V26.}  The complete V1--V25 body above is retained verbatim.
V26 replaces the divergent pseudoinverse at a simple accidental zero-mode
crossing by an exact finite divisor source packet.  A one-dimensional complex
crossing gives the regular factor
$\sigma+\widehat{\bar\jmath}\widehat\jmath$; an oriented real skew crossing
gives $\sigma+\widehat j_1\widehat j_2$.  The complement remains uniformly
invertible, all source/BV descendants are regular, nested shell elimination
preserves one common Schur/Pfaffian divisor cocycle, and a crossing-saturated
normalized functional has the same fixed-order continuum through the divisor.
The ordinary vacuum frame is recovered away from the zero by multiplication
with $\sigma$, so its pole is a line-coordinate effect.  The crossing is
accidental and leaves the Fredholm index unchanged.  The next gate is a true
Wilson-sign crossing where the overlap projector itself changes rank.



\clearpage
\section{V27 continuation: resolved Wilson-sign crossings, elementary domain modification, and genuine overlap-index jump}
\label{sec:v27-continuation}

V26 closed an accidental chiral divisor: the chiral operator itself acquired a
kernel--cokernel pair while the Wilson sign projector remained in one protected
gapped component, so the Fredholm index did not change.  A genuine overlap
index jump is more singular.  It occurs when the Hermitian Wilson kernel itself
has an eigenvalue cross zero.  The overlap sign is then undefined on the
crossing locus and the dimension of the moving chiral domain changes.

The correct variable is therefore not a value of $\sgn H_W$ at the crossing.
It is an \emph{oriented blow-up} of the Wilson divisor: one retains the isolated
crossing eigenline and the sign of the normal coordinate as discrete boundary
data, while the uniformly gapped Wilson complement is treated by the ordinary
V1--V25 functional calculus.  On the two boundary faces the overlap sign is
perfectly regular, but the faces differ by one rank-one projector.  This gives
an exact elementary modification of the chiral domain and an exact determinant
line transition.

V27 proves, at the native finite-regulator level:
\begin{enumerate}[label=\textup{(V27.\arabic*)},leftmargin=3.8em]
\item the resolved one-sided Wilson signs and projectors at a simple crossing;
\item the exact overlap-index jump and its equality with Wilson spectral flow;
\item the rank-one update of the overlap/Yukawa chiral map;
\item the five-term kernel/cokernel exact sequence and the induced Fredholm
 determinant-line elementary modification;
\item the source-complete tangential Kato calculus and the discrete normal jump
 calculus, with no distributional differentiation of the sign;
\item compatibility of the elementary modification with hard-shell Schur
 elimination and cutoff composition;
\item a coefficientwise continuum comparison of the two adjacent fixed-index
 sectors under the V24 soft-consistency hypotheses on the gapped complement;
\item the parity obstruction showing that an independently physical real
 Pfaffian branch cannot remain nondegenerate across an isolated odd-dimensional
 Wilson-domain jump.
\end{enumerate}

The theorem deliberately does \emph{not} define an overlap theory at
$H_W$-eigenvalue zero.  It constructs the exact transition between the two
one-sided theories and records the index-changing line/source writer carried by
the crossing eigenline.

\subsection{Simple Wilson crossing and the oriented blow-up chart}
\label{sec:v27-Wilson-crossing}

Let $s$ denote one real crossing parameter and let $y$ collect the remaining
background/source coordinates.  Consider a smooth native Hermitian Wilson
family
\begin{equation}
 H_W(s,y)=H_W(s,y)^*
 \label{eq:v27-HW-family}
\end{equation}
with the same finite-range/exponentially local gauge/spin covariance and fixed
source differentiability used in V1--V25.

\begin{definition}[Simple regular Wilson crossing]
\label{def:v27-simple-crossing}
A point $(0,y_0)$ is a \emph{simple regular Wilson crossing} if, after shrinking
one compact coordinate chart, there are:
\begin{enumerate}[label=\textup{(W\arabic*)},leftmargin=3em]
\item a smooth rank-one Riesz projector $\Pi_\times(s,y)$ and eigenvalue
 $\lambda(s,y)$ with
 \begin{equation}
  H_W\Pi_\times=\lambda\Pi_\times,
  \qquad
  \lambda(0,y_0)=0,
  \qquad
  \partial_s\lambda(0,y_0)\ne0;
  \label{eq:v27-crossing-eigenvalue}
 \end{equation}
\item the orthogonal complement $Q_\times:=I-\Pi_\times$ reduces $H_W$ and
 obeys
 \begin{equation}
  \boxed{
  |Q_\times H_WQ_\times|
  \succeq g_\times Q_\times}
  \label{eq:v27-complement-gap}
 \end{equation}
 for one cutoff- and volume-independent $g_\times>0$;
\item all fixed background/source derivatives required by the coefficient
 packet obey the V1 weighted local kernel bounds.
\end{enumerate}
The crossing line bundle is
\begin{equation}
 \mathcal L_\times:=\Ran\Pi_\times.
 \label{eq:v27-crossing-line}
\end{equation}
\end{definition}

\begin{lemma}[Locality and source calculus of the crossing line]
\label{lem:v27-crossing-projector-local}
On a simple regular Wilson crossing chart, $\Pi_\times$, $Q_\times$ and every
prescribed fixed source derivative are uniformly exponentially local after one
harmless loss of exponent.  On each star-shaped subchart, $\mathcal L_\times$
admits a unit Kato frame $\psi_\times$ with the same fixed source bounds.
\end{lemma}

\begin{proof}
Choose a small contour around $\lambda(s,y)$ which remains separated from the
complement spectrum by $g_\times/2$.  The projector is the Riesz integral
\[
 \Pi_\times={1\over2\pi i}\oint_\Gamma(z-H_W)^{-1}\,dz.
\]
The V13 gap-to-resolvent locality theorem and its differentiated version give
uniform exponential locality and all fixed source derivatives.  The V1/V13
Kato equation then supplies the local unit frame.
\end{proof}

Put
\begin{equation}
 H_W^\perp:=Q_\times H_WQ_\times,
 \qquad
 \varepsilon^\perp:=Q_\times\sgn(H_W^\perp)Q_\times .
 \label{eq:v27-complement-sign}
\end{equation}
The operator $\varepsilon^\perp$ is smooth across $\lambda=0$ because the
complement never loses its gap.

\begin{definition}[Resolved Wilson-sign blow-up]
\label{def:v27-resolved-sign}
Replace the divisor coordinate $\lambda$ by
\begin{equation}
 \lambda=\delta\rho,
 \qquad
 \rho\ge0,
 \qquad
 \delta\in\{-1,+1\}.
 \label{eq:v27-oriented-blowup}
\end{equation}
The two resolved boundary values of the Wilson sign are
\begin{equation}
 \boxed{
 \varepsilon^{[\delta]}
 :=\varepsilon^\perp+\delta\Pi_\times .}
 \label{eq:v27-resolved-epsilon}
\end{equation}
The corresponding overlap chiral projectors are
\begin{equation}
 \boxed{
 \widehat P_-^{[\delta]}
 ={1\over2}(I+\varepsilon^{[\delta]})
 =P_\perp^+
  +{1+\delta\over2}\Pi_\times,}
 \qquad
 P_\perp^+:={1\over2}(Q_\times+\varepsilon^\perp).
 \label{eq:v27-resolved-projectors}
\end{equation}
\end{definition}

\begin{theorem}[Exact one-sided sign/projector jump]
\label{thm:v27-sign-jump}
The resolved sign and chiral projector satisfy
\begin{align}
 \boxed{
 \varepsilon^{[+]}-\varepsilon^{[-]}
 =2\Pi_\times,}
 \label{eq:v27-epsilon-jump}\\
 \boxed{
 \widehat P_-^{[+]}-\widehat P_-^{[-]}
 =\Pi_\times.}
 \label{eq:v27-projector-jump}
\end{align}
In particular, with
\begin{equation}
 W_-:=\Ran\widehat P_-^{[-]},
 \qquad
 W_+:=\Ran\widehat P_-^{[+]},
 \label{eq:v27-Wpm}
\end{equation}
one has the orthogonal elementary modification
\begin{equation}
 \boxed{W_+=W_-\oplus\mathcal L_\times.}
 \label{eq:v27-domain-elementary-modification}
\end{equation}
Every tangential source derivative of the two resolved faces is bounded; the
normal transition is the finite jump
\eqref{eq:v27-epsilon-jump}--\eqref{eq:v27-projector-jump}, not a derivative of
$\sgn$ through zero.
\end{theorem}

\begin{proof}
On $Q_\times$ the two signs coincide with $\varepsilon^\perp$.  On the
crossing line they are respectively $+I$ and $-I$.  The displayed formulas are
therefore immediate.  Since $P_\perp^+\Pi_\times=0$, the range formula is an
orthogonal direct sum.  Tangential derivatives act only on the smooth Riesz
projector and complement sign.  No derivative across the discontinuity is
taken.
\end{proof}

\begin{firewall}[There is no overlap sign at the unresolved divisor]
At $\lambda=0$, $\sgn H_W$ is undefined on $\mathcal L_\times$.  V27 does not
assign $\sgn(0)=0$, does not average the two signs and does not insert an
ad hoc regulator on the crossing line.  The resolved object has two boundary
faces $\delta=\pm1$ carrying the two legitimate overlap theories.
\end{firewall}

\subsection{Exact index jump and Wilson spectral flow}
\label{sec:v27-index-jump}

Let the ambient complex fibre have dimension $2m$ so that
$\operatorname{rank} P_+=m$, as in V10--V11.  On either resolved face define the overlap
Fredholm index by the V11 convention
\begin{equation}
 k^{[\delta]}
 :=\operatorname{rank}\widehat P_-^{[\delta]}-m.
 \label{eq:v27-index-def}
\end{equation}

\begin{theorem}[Trace formula and unit index jump]
\label{thm:v27-index-jump}
On the two resolved faces,
\begin{equation}
 \boxed{
 k^{[\delta]}
 ={1\over2}\Tr\varepsilon^{[\delta]}.}
 \label{eq:v27-index-trace}
\end{equation}
Consequently
\begin{equation}
 \boxed{k^{[+]}-k^{[-]}=1.}
 \label{eq:v27-index-plusminus}
\end{equation}
If the physical parameter $s$ crosses the divisor with
\begin{equation}
 \chi_\times:=\sgn\partial_s\lambda(0,y_0),
 \label{eq:v27-crossing-orientation}
\end{equation}
then the one-sided Fredholm indices satisfy
\begin{equation}
 \boxed{
 k(s>0)-k(s<0)=\chi_\times.}
 \label{eq:v27-index-spectral-flow}
\end{equation}
Thus the overlap-index jump equals the signed spectral flow of the single
Wilson eigenvalue.
\end{theorem}

\begin{proof}
Since the ambient dimension is $2m$,
\[
 \operatorname{rank}\widehat P_-^{[\delta]}
 ={1\over2}\Tr(I+\varepsilon^{[\delta]})
 =m+{1\over2}\Tr\varepsilon^{[\delta]}.
\]
This proves \eqref{eq:v27-index-trace}.  The sign jump
\eqref{eq:v27-epsilon-jump} has trace two because $\Pi_\times$ has complex rank
one, proving \eqref{eq:v27-index-plusminus}.  Along the oriented parameter, the
sign of $\lambda$ changes from $-\chi_\times$ to $+\chi_\times$, so the index
change is $\chi_\times$.
\end{proof}

\begin{corollary}[Additivity for finitely many isolated simple crossings]
\label{cor:v27-index-additivity}
Along a path with finitely many separated simple Wilson crossings and no
simultaneous multiplicity, the total overlap-index change is
\begin{equation}
 \boxed{
 \Delta k
 =\sum_j\chi_{\times,j}.}
 \label{eq:v27-index-sum}
\end{equation}
\end{corollary}

\begin{proof}
Between crossings the Wilson gap remains open, so the rank of
$\widehat P_-$ is constant.  Sum Theorem~\ref{thm:v27-index-jump} over the
crossings.
\end{proof}

\subsection{The one-column chiral-map update}
\label{sec:v27-chiral-update}

On each face define
\begin{equation}
 D_{\rm ov}^{[\delta]}
 ={1\over a}(I+\gamma_5\varepsilon^{[\delta]}),
 \label{eq:v27-Dov-faces}
\end{equation}
and let $D_\chi^{[\delta]}$ be the complete overlap/Yukawa chiral block with the
same smooth Higgs/Yukawa map restricted to the corresponding domain
$W_\delta$.

\begin{lemma}[Rank-one overlap update]
\label{lem:v27-Dov-jump}
The ambient overlap operators obey
\begin{equation}
 \boxed{
 D_{\rm ov}^{[+]}
 -D_{\rm ov}^{[-]}
 ={2\over a}\gamma_5\Pi_\times.}
 \label{eq:v27-Dov-jump}
\end{equation}
On the common domain $W_-$ the two chiral restrictions coincide.  If
$\psi_\times$ is a unit frame of $\mathcal L_\times$, then relative to
$W_+=W_-\oplus\mathcal L_\times$ one has
\begin{equation}
 \boxed{
 D_\chi^{[+]}
 =\bigl[D_\chi^{[-]}\ \ c_\times\bigr],
 \qquad
 c_\times:=D_\chi^{[+]}\psi_\times.}
 \label{eq:v27-one-column-update}
\end{equation}
For the pure overlap kinetic block,
\begin{equation}
 c_\times^{\rm kin}={2\over a}P_+\psi_\times.
 \label{eq:v27-crossing-column-kinetic}
\end{equation}
\end{lemma}

\begin{proof}
Equation \eqref{eq:v27-Dov-jump} follows immediately from
\eqref{eq:v27-epsilon-jump}.  Every vector in $W_-$ lies in the gapped
complement of $\Pi_\times$, where the two signs and hence the two overlap
operators coincide.  The Yukawa map is the same smooth ambient native map on
both faces; only its domain restriction changes.  Therefore the plus-side map
is the minus-side map with one additional column.  The pure kinetic column
follows from the exact V11 formula
$D_P=2a^{-1}P_+|_W$.
\end{proof}

\subsection{Kernel--cokernel exact sequence and determinant-line modification}
\label{sec:v27-determinant-line}

Let
\begin{equation}
 T:=D_\chi^{[-]}:W_-\to F,
 \qquad
 T':=D_\chi^{[+]}:W_-\oplus\mathcal L_\times\to F,
 \label{eq:v27-T-Tprime}
\end{equation}
where $F$ is the fixed target chiral space.  Let
$K_\pm=\ker D_\chi^{[\pm]}$ and
$C_\pm=\operatorname{coker}D_\chi^{[\pm]}$.

\begin{theorem}[Five-term exact sequence of an index-changing crossing]
\label{thm:v27-five-term}
The one-column extension \eqref{eq:v27-one-column-update} induces a canonical
exact sequence
\begin{equation}
 \boxed{
 0\longrightarrow K_-
 \longrightarrow K_+
 \longrightarrow \mathcal L_\times
 \xrightarrow{\ \bar c_\times\ }
 C_-
 \longrightarrow C_+
 \longrightarrow0,}
 \label{eq:v27-five-term}
\end{equation}
where
\begin{equation}
 \bar c_\times(\ell)
 :=[D_\chi^{[+]}\ell]\in C_-=F/\operatorname{im}D_\chi^{[-]}.
 \label{eq:v27-cbar}
\end{equation}
In particular
\begin{equation}
 \dim K_+-\dim C_+
 =\dim K_--\dim C_-+1.
 \label{eq:v27-index-from-exact-sequence}
\end{equation}
\end{theorem}

\begin{proof}
Write an element of $W_+$ as $(w,\ell)$.  The kernel equation is
$Tw+c_\times\ell=0$.  Projection to the second component gives a map
$K_+\to\mathcal L_\times$.  Its kernel is precisely $K_-$.  Its image consists
of those $\ell$ for which $c_\times\ell\in\operatorname{im}T$, which is exactly
$\ker\bar c_\times$.  The cokernel of $\bar c_\times$ is
$F/(\operatorname{im}T+\operatorname{im}c_\times)$, which is $C_+$.
This proves exactness.  Alternating dimensions in the finite exact sequence
give \eqref{eq:v27-index-from-exact-sequence}.
\end{proof}

V12 uses the Fredholm-line convention
\begin{equation}
 \operatorname{Det}(T)
 :=\det C_T\otimes(\det K_T)^*.
 \label{eq:v27-Det-convention}
\end{equation}

\begin{theorem}[Elementary modification of the Fredholm determinant line]
\label{thm:v27-line-modification}
The exact sequence \eqref{eq:v27-five-term} gives a canonical line isomorphism
\begin{equation}
 \boxed{
 \operatorname{Det}(D_\chi^{[+]})
 \simeq
 \operatorname{Det}(D_\chi^{[-]})
 \otimes \mathcal L_\times^*.}
 \label{eq:v27-line-modification}
\end{equation}
For the reverse crossing the transition is the dual one,
\begin{equation}
 \operatorname{Det}(D_\chi^{[-]})
 \simeq
 \operatorname{Det}(D_\chi^{[+]})
 \otimes \mathcal L_\times.
 \label{eq:v27-line-modification-reverse}
\end{equation}
\end{theorem}

\begin{proof}
For an exact finite sequence
$0\to A_0\to A_1\to A_2\to A_3\to A_4\to0$, the determinant functor gives
\[
 \det A_0\otimes\det A_2\otimes\det A_4
 \simeq
 \det A_1\otimes\det A_3.
\]
Apply this to \eqref{eq:v27-five-term}:
\[
 \det K_-\otimes\mathcal L_\times\otimes\det C_+
 \simeq
 \det K_+\otimes\det C_-.
\]
Rearranging with convention \eqref{eq:v27-Det-convention} gives
\eqref{eq:v27-line-modification}.  The reverse relation is its dual.
\end{proof}

\begin{corollary}[Transition of saturation writers]
\label{cor:v27-writer-transition}
Let $\omega_-$ be a local nonzero writer in
$\operatorname{Det}(D_\chi^{[-]})^*$.  Choosing a Kato unit vector
$\psi_\times\in\mathcal L_\times$ defines the plus-side writer
\begin{equation}
 \boxed{
 \omega_+:=\omega_-\otimes\psi_\times
 \in\operatorname{Det}(D_\chi^{[+]})^*.}
 \label{eq:v27-writer-transition}
\end{equation}
A phase change $\psi_\times\mapsto e^{i\theta}\psi_\times$ is precisely the
expected determinant-line frame change and is retained with all source
contacts.
\end{corollary}

\begin{proof}
Dualize \eqref{eq:v27-line-modification}.  The dual transition line is
$\mathcal L_\times$, so tensoring the old writer with a crossing-line frame
gives the new one.  Frame covariance is the ordinary one-dimensional line
transformation law.
\end{proof}

\subsection{Source-complete tangential calculus and the discrete jump writer}
\label{sec:v27-source-jump}

The crossing coordinate is singular only in the normal direction.  Tangential
source derivatives remain ordinary differentiable data.

\begin{definition}[Wilson jump operator]
\label{def:v27-jump-operator}
For any resolved object $F^{[\delta]}$ with one-sided boundary values define
\begin{equation}
 \boxed{
 \mathfrak J_\times F
 :=F^{[+]}-F^{[-]}.}
 \label{eq:v27-jump-operator}
\end{equation}
Tangential derivatives $\partial_y^I$ are taken separately on the two resolved
faces.  The normal transition is recorded by $\mathfrak J_\times$, not by
$\partial_\lambda\sgn H_W$.
\end{definition}

\begin{lemma}[Jump algebra]
\label{lem:v27-jump-algebra}
For resolved products $F,G$,
\begin{equation}
 \boxed{
 \mathfrak J_\times(FG)
 =(\mathfrak J_\times F)G^{[-]}
 +F^{[+]}(\mathfrak J_\times G).}
 \label{eq:v27-jump-Leibniz}
\end{equation}
Moreover
\begin{align}
 \mathfrak J_\times\varepsilon&=2\Pi_\times,
 \label{eq:v27-jump-epsilon}\\
 \mathfrak J_\times\widehat P_-&=\Pi_\times,
 \label{eq:v27-jump-P}\\
 \mathfrak J_\times D_{\rm ov}&={2\over a}\gamma_5\Pi_\times.
 \label{eq:v27-jump-Dov}
\end{align}
Every tangential derivative of these jump coefficients is a finite sum of
V13-local Riesz/Kato source words.
\end{lemma}

\begin{proof}
Expand
$F^{[+]}G^{[+]}-F^{[-]}G^{[-]}$ by adding and subtracting
$F^{[+]}G^{[-]}$.  The three explicit jumps are Theorem
\ref{thm:v27-sign-jump} and Lemma~\ref{lem:v27-Dov-jump}.  Tangential
regularity follows from Lemma~\ref{lem:v27-crossing-projector-local} and the
gapped complement functional calculus.
\end{proof}

\begin{theorem}[Finite Wilson-crossing source packet]
\label{thm:v27-source-packet}
Fix a finite source/BV order $r$.  The enlarged crossing packet consisting of
\begin{equation}
 \boxed{
 \bigl(
 \Pi_\times,\ \psi_\times,\
 \varepsilon^\perp,\
 \mathfrak J_\times\varepsilon,\
 \mathfrak J_\times\widehat P_-,\
 \mathfrak J_\times D_{\rm ov},\
 \omega_-,\omega_+
 \bigr)}
 \label{eq:v27-source-packet}
\end{equation}
together with all prescribed tangential derivatives through order $r$ is a
finite, cutoff-uniform local source packet.  It is closed under the finite
Leibniz/contact operations needed by the V4--V7 coefficientwise BV calculus.
No distribution $\delta(\lambda)$ and no inverse power of $\lambda$ is
introduced.
\end{theorem}

\begin{proof}
Every listed object is either a Riesz/Kato datum of the isolated crossing line,
a gapped-complement functional calculus datum, or a finite one-sided jump of
such data.  The differentiated families are therefore finite at every fixed
order.  Lemma~\ref{lem:v27-jump-algebra} replaces the ordinary Leibniz rule in
the single normal direction.  The source algebra remains finite because the
crossing adds only one line frame and finitely many jump writers.
\end{proof}

\subsection{Compatibility with hard-shell elimination}
\label{sec:v27-shell-compatibility}

The crossing line is part of the low state.  Hard shell elimination therefore
never integrates it out.

\begin{theorem}[Elementary modification commutes with hard-shell Schur reduction]
\label{thm:v27-shell-commutes}
Let a hard shell be chosen inside the uniformly invertible complement of the
crossing line.  Denote the corresponding one-shell effective chiral maps by
$D_{\chi,j}^{[-]}$ and $D_{\chi,j}^{[+]}$.  Then:
\begin{enumerate}[label=\textup{(S\arabic*)},leftmargin=3em]
\item the same crossing line $\mathcal L_\times$ survives to the low state;
\item the effective maps still satisfy a one-column relation
\begin{equation}
 D_{\chi,j}^{[+]}
 =\bigl[D_{\chi,j}^{[-]}\ \ c_{\times,j}\bigr];
 \label{eq:v27-shell-column}
\end{equation}
\item their Fredholm lines obey
\begin{equation}
 \boxed{
 \operatorname{Det}(D_{\chi,j}^{[+]})
 \simeq
 \operatorname{Det}(D_{\chi,j}^{[-]})\otimes\mathcal L_\times^*;}
 \label{eq:v27-shell-line}
\end{equation}
\item the transition is independent of how nested hard shells are regrouped.
\end{enumerate}
\end{theorem}

\begin{proof}
Write each one-sided chiral map in low/hard block form, with
$\mathcal L_\times$ included in the low domain.  The hard block and its inverse
are identical on the common domain because the sign difference is supported on
$\mathcal L_\times$.  Schur elimination therefore transforms only the added
column:
\[
 c_\times\longmapsto
 c_{\times,j}=c_{\times,\ell}-B H^{-1}c_{\times,h},
\]
while the domain elementary modification survives unchanged.  Apply Theorem
\ref{thm:v27-line-modification} to the effective maps.  Associativity of Schur
reduction gives regrouping independence.
\end{proof}

\begin{corollary}[Cofinal cutoff compatibility of the index jump]
\label{cor:v27-cutoff-compatibility}
The crossing line, unit index jump and determinant-line elementary modification
are common soft data across all ultraviolet shell cutoffs.  In particular, the
index jump cannot be redefined by shell-dependent finite counterterms or line
normalizations.
\end{corollary}

\subsection{Adjacent-sector coefficientwise continuum}
\label{sec:v27-continuum}

V27 now compares the two one-sided continuum theories.  There is deliberately
no theory assigned at the unresolved Wilson divisor itself.

\begin{definition}[Resolved adjacent-sector continuum datum]
\label{def:v27-resolved-continuum}
A \emph{resolved adjacent-sector datum} consists of:
\begin{enumerate}[label=\textup{(R\arabic*)},leftmargin=3em]
\item the common smooth Wilson complement
$(Q_\times,H_W^\perp,\varepsilon^\perp)$;
\item the crossing eigenline $\mathcal L_\times$ with its Kato/source packet;
\item the two boundary signs $\delta=\pm1$;
\item on each face, a V11 minimal fixed-index chiral/Yukawa block with a
 positive reduced singular-value floor and its declared saturation writer;
\item the V24 V2-consistent common soft chiral jet on the gapped complement and
 on the finite crossing cluster.
\end{enumerate}
\end{definition}

\begin{theorem}[Coefficientwise continuum on the two sides of a genuine index jump]
\label{thm:v27-adjacent-continuum}
Fix a simple regular Wilson crossing and a resolved adjacent-sector datum.
For every prescribed finite
$(L,\Delta,r,N,\mathcal P_{\rm ext})$:
\begin{enumerate}[label=\textup{(C\arabic*)},leftmargin=3em]
\item each resolved face $\delta=\pm1$ has the V1--V24 source-complete
 coefficientwise continuum in its own fixed-index sector;
\item the hard-shell forest, BV/master/source and determinant-line estimates are
 uniform up to the resolved boundary because the Wilson complement and both
 reduced chiral complements retain positive spectral floors;
\item the two continuum determinant lines are related by the cutoff-independent
 elementary modification
\begin{equation}
 \boxed{
 \operatorname{Det}_\infty^{[+]}
 \simeq
 \operatorname{Det}_\infty^{[-]}
 \otimes\mathcal L_{\times,\infty}^*;}
 \label{eq:v27-continuum-line-transition}
\end{equation}
\item after choosing the crossing Kato frame, the saturated scalar coefficient
 systems are related by the writer transition
 $\omega_+=\omega_-\otimes\psi_\times$ and all retained tangential source
 descendants of that transition;
\item along an oriented physical parameter crossing, the Fredholm index jump is
 $\chi_\times$ and is independent of the ultraviolet cutoff.
\end{enumerate}
\end{theorem}

\begin{proof}
On either resolved face the Wilson sign is gapped and the chiral block lies in
one V11 protected fixed-index component.  The V24 soft consistency and hard
Feshbach estimates apply to the common complement and identify the finite
crossing Riesz data of V2-consistent microscopic completions.  Thus the V6
Cauchy theorem applies separately to each face, uniformly up to the boundary in
all quantities which remain inside that face.

The crossing line itself is a finite soft Riesz bundle and survives every hard
shell by Theorem~\ref{thm:v27-shell-commutes}.  Therefore the finite-cutoff line
isomorphisms \eqref{eq:v27-shell-line} pass directly to the projective continuum
limit, proving \eqref{eq:v27-continuum-line-transition}.  The writer transition
is Corollary~\ref{cor:v27-writer-transition}; V7 projectivity and the V27 finite
source packet pass all fixed tangential descendants to the same limit.  The
index statement is Theorem~\ref{thm:v27-index-jump}.
\end{proof}

\begin{firewall}[No scalar equality between different index sectors]
The two sides have different Fredholm determinant lines and different required
zero-mode saturation degree.  V27 therefore does not write
$W^{[+]}=W^{[-]}$ as scalar functions.  The invariant statement is the
line-valued elementary modification
\eqref{eq:v27-continuum-line-transition}, or its scalarized version after the
explicit crossing writer $\psi_\times$ is supplied.
\end{firewall}

\subsection{Real Pfaffian parity at a Wilson-domain jump}
\label{sec:v27-real-parity}

The complex rank-one theorem must not be copied blindly into an independently
physical real Majorana/Pfaffian sector.

\begin{proposition}[Odd real domain jump leaves the nondegenerate Pfaffian component]
\label{prop:v27-Pfaffian-parity}
Let an independently real Majorana quadratic form be represented on each side
by a nondegenerate real skew matrix on the occupied Majorana domain.  Then the
domain dimension is even on each side.  Consequently an isolated Wilson-sign
crossing which changes that real domain dimension by one cannot connect two
nondegenerate Pfaffian sectors.  A Pfaffian-compatible Wilson crossing must have
even total real multiplicity (or else pass through an additional real zero
mode).
\end{proposition}

\begin{proof}
A real skew-symmetric matrix in odd dimension has zero determinant because
\[
 \det A=\det A^T=\det(-A)=(-1)^n\det A=-\det A.
\]
Thus a nondegenerate real skew form exists only in even dimension.  A rank-one
real domain modification flips the parity, so one of the two sides cannot
remain in the nondegenerate Pfaffian component.
\end{proof}

\begin{corollary}[Minimal real Wilson crossing is a finite even block]
\label{cor:v27-real-even-block}
For the separately real V25 branch, the correct genuine Wilson-index transition
must be treated by an even-dimensional real crossing block compatible with the
native reality structure.  The V27 complex rank-one elementary-modification
proof applies after complexification, but a real Pfaffian transition requires
the corresponding even-block determinant/Pfaffian functor and orientation data.
\end{corollary}

\begin{firewall}[No fictitious rank-one real Pfaffian jump]
V27 proves the genuine rank-one Wilson index jump for the complex chiral
sector.  It does not force an isolated real rank-one crossing into a Pfaffian
formula.  The independently real branch has an additional parity/reality
constraint and its minimal index-changing block may have higher multiplicity.
\end{firewall}

\subsection{Multiple and higher-multiplicity Wilson crossings}
\label{sec:v27-multiple}

The rank-one result extends formally to a finite crossing block, but the
stratification is richer and is not needed for the present load-bearing gate.

\begin{proposition}[Finite crossing-block trace law]
\label{prop:v27-finite-block}
Let a finite-dimensional Wilson crossing block $E_\times$ of complex dimension
$m_\times$ be separated from a uniformly gapped complement and let its
one-sided sign change from an involution $S_-$ to $S_+$ on $E_\times$.  Then
\begin{equation}
 \boxed{
 \Delta k
 ={1\over2}\Tr_{E_\times}(S_+-S_-).}
 \label{eq:v27-finite-block-index}
\end{equation}
If the block diagonalizes into simple crossings, this is the signed sum of
their spectral flows.
\end{proposition}

\begin{proof}
The gapped complement contributes the same trace to the two one-sided signs.
Apply the trace formula \eqref{eq:v27-index-trace} to the finite crossing block.
\end{proof}

\begin{firewall}[Higher multiplicity is a stratified problem]
When several Wilson eigenvalues meet zero simultaneously, the crossing bundle
may have nontrivial internal mixing and the determinant-line modification is a
higher exterior power rather than a single line.  V27 records the trace law but
does not claim a complete normal form for that stratified crossing.
\end{firewall}

\section{V27 anti-circularity audit}
\label{sec:v27-audit}

\begin{enumerate}[leftmargin=2.4em]
\item \textbf{The sign is never evaluated at zero.}
The Wilson divisor is replaced by two oriented boundary faces carrying the two
legitimate signs.
\item \textbf{The index jump is not inferred from a continuum anomaly.}
It is the exact finite-dimensional rank/trace change of the native overlap
projector.
\item \textbf{The determinant line is not scalarized without the missing
crossing line.}
The elementary modification explicitly tensors by
$\mathcal L_\times^*$, and scalar comparison requires the Kato crossing writer.
\item \textbf{Normal differentiation of $\sgn H_W$ is not used.}
The singular normal direction is represented by the finite jump writer
$\mathfrak J_\times$; tangential source derivatives remain ordinary Riesz/Kato
derivatives.
\item \textbf{Hard shells do not integrate the index-changing mode.}
The crossing line is retained as a soft state and survives Schur elimination
unchanged.
\item \textbf{No theory is claimed on the unresolved divisor.}
V27 compares the two adjacent one-sided continuum sectors; it does not invent a
value of the overlap sign at the Wilson zero.
\item \textbf{The real Pfaffian parity constraint is kept separate.}
A complex rank-one index jump is not automatically an admissible real
rank-one Pfaffian transition.
\end{enumerate}

\section{V27 status ledger}
\label{sec:v27-status}

\subsection*{Proved in V27}
\begin{enumerate}[leftmargin=2.2em]
\item A simple isolated Wilson eigenvalue crossing admits a smooth rank-one
Riesz projector and a uniformly gapped local Wilson complement.  The crossing
projector and every fixed source derivative are exponentially local.
\item Replacing the zero eigenvalue by an oriented two-face blow-up yields exact
one-sided signs
$\varepsilon^{[\delta]}=\varepsilon^\perp+\delta\Pi_\times$ and chiral
projectors whose ranges differ by the crossing eigenline.
\item The overlap Fredholm index satisfies
$k=\frac12\Tr\varepsilon$ in the V11 convention.  A negative-to-positive
Wilson eigenvalue crossing increases the index by one; an oriented simple
crossing changes the index by its signed Wilson spectral flow.
\item The plus-side overlap/Yukawa map is the minus-side map with one additional
crossing-line column.  This gives the exact five-term kernel/cokernel sequence
\eqref{eq:v27-five-term}.
\item The Fredholm determinant line undergoes the canonical elementary
modification
$\operatorname{Det}(D_+)\simeq\operatorname{Det}(D_-)\otimes
\mathcal L_\times^*$; the dual saturation writer is multiplied by the crossing
Kato frame.
\item A finite source-complete transition calculus is obtained from ordinary
tangential Riesz/Kato derivatives plus the discrete jump operator
$\mathfrak J_\times$.  No distributional derivative of the sign and no inverse
power of the Wilson eigenvalue is required.
\item The elementary domain/line modification commutes with exact hard-shell
Schur elimination and is independent of shell regrouping.
\item Under the V24 V2-consistent smooth soft hypotheses on the gapped
complement and finite crossing cluster, both adjacent fixed-index sectors have
source-complete coefficientwise continua whose determinant lines are related by
the same elementary modification.  The index jump is cutoff independent.
\item An independently physical nondegenerate real Pfaffian branch cannot
undergo an isolated odd-dimensional real Wilson-domain jump; the minimal real
crossing block must have even multiplicity or leave the protected Pfaffian
component.
\end{enumerate}

\subsection*{Inherited}
\begin{enumerate}[leftmargin=2.2em]
\item V1--V7: source-complete coefficientwise continuum, common BV restoration,
fixed-loop forest induction, ultraviolet Cauchy convergence, scheme/reference
covariance and exact perturbative-order projectivity.
\item V8--V24: regulator comparison, controlled zero-mode classification,
common soft Riesz anchoring and the V2-consistent covariant
G\aa rding/Feshbach theorem.
\item V25: global orientation of the separately real neutral-Majorana branch and
its mod-two spectral-flow sign law.
\item V26: regular source/BV/divisor packet for accidental kernel--cokernel pair
crossings which leave the Fredholm index unchanged.
\end{enumerate}

\subsection*{Conditional}
\begin{enumerate}[leftmargin=2.2em]
\item The Wilson crossing is simple and isolated from the remaining Wilson
spectrum by a cutoff-uniform gap.  Higher multiplicity is only covered by the
trace law of Proposition~\ref{prop:v27-finite-block}.
\item Each one-sided chiral/Yukawa theory is assumed to lie in a V11 minimal
fixed-index reduced-gap component.  Additional accidental chiral divisors are
treated by the separate V26 packet.
\item The coefficientwise continuum comparison uses the V24 smooth-background
V2-consistent soft chiral jet on the gapped complement and finite crossing
cluster.
\item The real Pfaffian branch requires its native reality/orientation
structure; an odd real rank jump necessarily leaves the nondegenerate Pfaffian
component.
\item The positive IR/reference assumptions remain in force.
\end{enumerate}

\subsection*{Open beyond V27}
\begin{enumerate}[leftmargin=2.2em]
\item Construct the complete finite-multiplicity Wilson-crossing normal form,
including non-Abelian mixing inside a crossing bundle and the associated higher
exterior-power determinant/Pfaffian transition.
\item Extend the adjacent-sector comparison to a global stratified family with
many Wilson-gap walls and prove cocycle consistency around codimension-two wall
intersections.
\item Remove the positive IR/reference scale and separately address analytic
finite-coupling RG trajectories, perturbative summability and nonperturbative
completion.
\end{enumerate}

\subsection*{Exact next load-bearing theorem}
\begin{center}
\fbox{\parbox{0.92\textwidth}{
\textbf{Higher-multiplicity Wilson-wall and wall-intersection theorem.}
Let a finite-rank Wilson spectral cluster cross zero while the complementary
Wilson spectrum remains uniformly gapped.  Construct a matrix-valued resolved
sign chart on the crossing bundle, identify the elementary modification of the
chiral domain by its positive/negative crossing subbundles, and derive the
higher determinant/Pfaffian exterior-power transition with all fixed source/BV
descendants.  Then analyze codimension-two intersections of two Wilson walls
and prove that the resulting elementary-modification maps satisfy the required
cocycle/commutation law.  This would upgrade the V27 simple-wall theorem to a
global stratified index-changing regulator family.}}
\end{center}

\section*{Append-only continuation marker: V27}
\addcontentsline{toc}{section}{Append-only continuation marker: V27}
\textbf{End of V27.}  The complete V1--V26 body above is retained verbatim.
V27 resolves a genuine simple Wilson-gap crossing by replacing the undefined
sign at the divisor with two oriented boundary faces.  The one-sided overlap
projectors differ by the isolated crossing eigenline, so the native Fredholm
index jumps by the signed Wilson spectral flow.  The chiral map changes by one
column, producing a canonical five-term kernel/cokernel sequence and the
Fredholm-line elementary modification
$\operatorname{Det}(D_+)\simeq\operatorname{Det}(D_-)\otimes
\mathcal L_\times^*$.  Tangential source derivatives remain ordinary local
Riesz/Kato data while the normal singular direction is represented by a finite
jump writer.  The elementary modification commutes with hard-shell Schur
elimination and passes to the adjacent coefficientwise continuum sectors.  No
overlap theory is assigned at the unresolved Wilson zero.  The next stronger
gate is the finite-multiplicity/codimension-two stratified Wilson-wall problem.



\clearpage
\section{V28 continuation: finite-multiplicity Wilson walls, graded elementary modifications, and codimension-two coherence}
\label{sec:v28-continuation}

V27 resolved a simple complex Wilson eigenline crossing.  The next gate is not
obtained by repeating that proof eigenvalue by eigenvalue: at a simultaneous
finite-multiplicity crossing, the eigenvectors may mix non-Abelianly inside the
crossing cluster, and at an intersection of two walls the determinant-line
transitions carry the graded signs of exterior algebra.  The correct variables
are therefore the finite-dimensional Hermitian \emph{crossing form} on the
crossing bundle and the graded determinant functor.

V28 proves the finite-rank regular-crossing theorem at the native regulator
level.  It also closes the codimension-two cocycle for a clean transverse
intersection of two spectrally separated crossing bundles.  The result is more
precise than a naive commuting square: the two wall transitions commute in the
symmetric monoidal category of $\mathbb Z_2$-graded determinant lines.  If one
forgets the grading and insists on one fixed scalar ordering of crossing
writers, the two paths differ by the canonical Koszul sign.

The new statements are:
\begin{enumerate}[label=\textup{(V28.\arabic*)},leftmargin=3.8em]
\item a smooth finite-rank crossing bundle separated from the complementary
 Wilson spectrum;
\item the regular Hermitian crossing form and its matrix-valued resolved sign;
\item the gained/lost crossing subbundles and the signature formula for the
 native Fredholm-index jump;
\item the higher-rank kernel/cokernel exact sequences and the associated
 determinant-line elementary modification
 \(\det E_-\otimes(\det E_+)^*\);
\item the complete finite source/BV jump packet with no differentiation of
 $\operatorname{sgn}H_W$ through zero;
\item compatibility with hard-shell Schur elimination and coefficientwise
 continuum passage;
\item the real/Pfaffian even-block version with the V25 orientation data;
\item the graded codimension-two wall-intersection coherence law, including
 its explicit Koszul sign in scalarized writer coordinates.
\end{enumerate}

The theorem still does not cover a genuinely non-split conical intersection in
which no two smooth crossing subbundles can be separated.  That is a different
stratified blow-up problem and is isolated at the end of V28.

\subsection{Regular finite-rank Wilson crossing form}
\label{sec:v28-crossing-form}

Let $n$ be an oriented normal coordinate to a Wilson-gap wall and let $y$ denote
the tangential background/source variables.  Consider a smooth native Hermitian
Wilson family $H_W(n,y)$ with all V1--V27 locality/source bounds.

\begin{definition}[Regular finite-rank Wilson wall]
\label{def:v28-regular-wall}
A point $(0,y_0)$ lies on a \emph{regular rank-$m_\times$ Wilson wall} if,
after shrinking one compact chart, there exist:
\begin{enumerate}[label=\textup{(W\arabic*)},leftmargin=3em]
\item a smooth rank-$m_\times$ Riesz projector $\Pi_\times(n,y)$ whose
 spectral cluster is the only Wilson spectrum in $(-g_\times,g_\times)$;
\item the complementary projector $Q_\times=I-\Pi_\times$ and the uniform gap
 \begin{equation}
  \boxed{|Q_\times H_WQ_\times|\succeq g_\times Q_\times;}
  \label{eq:v28-complement-gap}
 \end{equation}
\item the wall condition
 \begin{equation}
  \Pi_\times H_W(0,y)\Pi_\times=0;
  \label{eq:v28-wall-zero}
 \end{equation}
\item an invertible Hermitian crossing form on
 $E_\times:=\operatorname{Ran}\Pi_\times|_{n=0}$,
 \begin{equation}
  \boxed{
  A_\times(y)
  :=\Pi_\times(0,y)(\partial_nH_W)(0,y)\Pi_\times(0,y),
  \qquad
  |A_\times|\succeq\mu_\times I,}
  \label{eq:v28-crossing-form}
 \end{equation}
 for one $\mu_\times>0$;
\item all fixed source/background derivatives required by the coefficient
 packet obey the native weighted local bounds.
\end{enumerate}
\end{definition}

The form $A_\times$ is allowed to be completely non-diagonal in any prechosen
crossing basis.  Its sign is the invariant replacement for choosing individual
crossing eigenvectors.

\begin{lemma}[Locality and tangential source calculus]
\label{lem:v28-crossing-bundle-local}
The crossing projector $\Pi_\times$, the complement projector $Q_\times$ and
every prescribed tangential source derivative are uniformly exponentially
local after one fixed loss of exponent.  On the wall, the Hermitian involution
\begin{equation}
 \boxed{J_\times:=\operatorname{sgn}A_\times}
 \label{eq:v28-J}
\end{equation}
and its spectral projectors
\begin{equation}
 R_\times^+:={1\over2}(I+J_\times),
 \qquad
 R_\times^-:={1\over2}(I-J_\times)
 \label{eq:v28-Rpm}
\end{equation}
form smooth finite-rank bundles with every required tangential source
 derivative bounded.  Write
\begin{equation}
 E_\times^\pm:=\operatorname{Ran}R_\times^\pm,
 \qquad
 r_\pm:=\operatorname{rank}E_\times^\pm,
 \qquad
 r_++r_-=m_\times.
 \label{eq:v28-Epm}
\end{equation}
\end{lemma}

\begin{proof}
The proof for $\Pi_\times$ and $Q_\times$ is the V27 Riesz/Combes--Thomas
argument with a contour enclosing the whole finite crossing cluster.  The
crossing form acts on this finite local bundle and has the uniform spectral gap
$\mu_\times$ from zero.  Its sign and the projectors
\eqref{eq:v28-Rpm} are therefore ordinary finite-dimensional Riesz functional
calculus, with differentiated bounds controlled by $\mu_\times^{-1}$ and the
source jets of $A_\times$.
\end{proof}

\subsection{Matrix-valued resolved sign and signature index jump}
\label{sec:v28-resolved-sign}

Put
\begin{equation}
 H_W^\perp:=Q_\times H_WQ_\times,
 \qquad
 \varepsilon^\perp:=Q_\times\operatorname{sgn}(H_W^\perp)Q_\times.
 \label{eq:v28-complement-sign}
\end{equation}
The complement sign is smooth across the wall.

\begin{theorem}[Finite-rank oriented blow-up normal form]
\label{thm:v28-resolved-sign}
For $n\to0^\pm$, the one-sided Wilson signs have the limits
\begin{equation}
 \boxed{
 \varepsilon^{[\delta]}
 =\varepsilon^\perp+\delta J_\times,
 \qquad
 \delta\in\{-1,+1\}.}
 \label{eq:v28-resolved-epsilon}
\end{equation}
Consequently the resolved overlap projectors are
\begin{equation}
 \boxed{
 \widehat P_-^{[\delta]}
 =P_\perp^+
  +{1\over2}(I_{E_\times}+\delta J_\times),}
 \qquad
 P_\perp^+:={1\over2}(Q_\times+\varepsilon^\perp),
 \label{eq:v28-resolved-P}
\end{equation}
and therefore
\begin{align}
 \widehat P_-^{[+]}|_{E_\times}&=R_\times^+,
 &
 \widehat P_-^{[-]}|_{E_\times}&=R_\times^-.
 \label{eq:v28-crossing-occupancies}
\end{align}
If
\begin{equation}
 W_0:=\operatorname{Ran}P_\perp^+,
 \label{eq:v28-W0}
\end{equation}
then the two one-sided chiral domains are
\begin{equation}
 \boxed{
 W_+=W_0\oplus E_\times^+,
 \qquad
 W_-=W_0\oplus E_\times^-.}
 \label{eq:v28-domain-replacement}
\end{equation}
The finite normal jumps are
\begin{align}
 \boxed{\mathfrak J_\times\varepsilon=2J_\times,}
 \label{eq:v28-sign-jump}\\
 \boxed{\mathfrak J_\times\widehat P_-=J_\times.}
 \label{eq:v28-projector-jump}
\end{align}
\end{theorem}

\begin{proof}
Taylor expansion in the crossing bundle gives
\[
 \Pi_\times H_W(n,y)\Pi_\times
 =nA_\times(y)+O(n^2)
\]
in operator norm.  Since $A_\times$ is invertible, the sign of this block
converges for $n\to0^\delta$ to $\delta\operatorname{sgn}A_\times$ by ordinary
gapped functional calculus.  The complementary block converges to
$\varepsilon^\perp$.  The projector and domain statements follow by taking
$(I+\varepsilon)/2$ and using the orthogonality of $E_\times$ and the
complement.
\end{proof}

\begin{theorem}[Signature formula for the genuine index jump]
\label{thm:v28-signature-index}
Let
\begin{equation}
 k^{[\delta]}
 :=\operatorname{rank}\widehat P_-^{[\delta]}-m
 \label{eq:v28-index-def}
\end{equation}
be the V11 overlap Fredholm index on each face.  Then
\begin{equation}
 \boxed{
 k^{[+]}-k^{[-]}
 =r_+-r_-
 =\operatorname{Tr}_{E_\times}J_\times
 =\operatorname{sign}(A_\times).}
 \label{eq:v28-signature-index}
\end{equation}
For an oppositely oriented normal coordinate the sign reverses.
\end{theorem}

\begin{proof}
The common domain $W_0$ contributes equally to both ranks.  Equation
\eqref{eq:v28-domain-replacement} therefore gives the first equality.  Since
$J_\times$ has $+1$ on $E_\times^+$ and $-1$ on $E_\times^-$, its trace is
$r_+-r_-$.  This is the signature of the invertible Hermitian crossing form.
\end{proof}

\begin{corollary}[Recovery of V27]
\label{cor:v28-V27}
If $m_\times=1$, then either $(r_+,r_-)=(1,0)$ or $(0,1)$ and
Theorem~\ref{thm:v28-signature-index} is exactly the V27 signed unit spectral
flow formula.
\end{corollary}

\subsection{Higher-rank chiral-map replacement and Fredholm-line functor}
\label{sec:v28-Fredholm-line}

Restrict the one-sided chiral/Yukawa maps to the common core and the crossing
subbundles.  Write
\begin{equation}
 D_0:W_0\longrightarrow\mathcal H_+,
 \qquad
 C_\pm:E_\times^\pm\longrightarrow\mathcal H_+
 \label{eq:v28-D0-Cpm}
\end{equation}
for the common-core map and the added crossing columns.  Then
\begin{equation}
 \boxed{
 D_+=[D_0\ \ C_+]:W_0\oplus E_\times^+\to\mathcal H_+,
 \qquad
 D_-=[D_0\ \ C_-]:W_0\oplus E_\times^-\to\mathcal H_+.}
 \label{eq:v28-Dpm}
\end{equation}
No nesting of $W_+$ and $W_-$ is assumed or needed.

Let $K_0=\ker D_0$ and $C_0=\operatorname{coker}D_0$.  For each sign there is
the canonical exact sequence
\begin{equation}
 0\to K_0\to K_\pm\to E_\times^\pm
 \xrightarrow{\bar C_\pm} C_0\to C_\pm^{\rm cok}\to0,
 \label{eq:v28-five-term-pm}
\end{equation}
where $K_\pm=\ker D_\pm$ and
$C_\pm^{\rm cok}=\operatorname{coker}D_\pm$.

\begin{theorem}[Higher-rank determinant-line elementary modification]
\label{thm:v28-line-transition}
Use the lineage convention
\begin{equation}
 \operatorname{Det}(D)
 :=\det(\operatorname{coker}D)\otimes(\det\ker D)^*.
 \label{eq:v28-Det-convention}
\end{equation}
Then the exact sequences \eqref{eq:v28-five-term-pm} induce canonical
isomorphisms
\begin{equation}
 \operatorname{Det}(D_\pm)
 \simeq
 \operatorname{Det}(D_0)\otimes(\det E_\times^\pm)^*.
 \label{eq:v28-line-to-core}
\end{equation}
Consequently
\begin{equation}
 \boxed{
 \operatorname{Det}(D_+)
 \simeq
 \operatorname{Det}(D_-)
 \otimes\det E_\times^-
 \otimes(\det E_\times^+)^*.}
 \label{eq:v28-line-transition}
\end{equation}
Define the graded wall line
\begin{equation}
 \boxed{
 \mathbb T_\times
 :=\det E_\times^-\otimes(\det E_\times^+)^*,
 \qquad
 |\mathbb T_\times|
 :=m_\times\pmod2.}
 \label{eq:v28-wall-line}
\end{equation}
Then the wall transition is simply
$\operatorname{Det}(D_+)\simeq\operatorname{Det}(D_-)\otimes\mathbb T_\times$
in the graded determinant-line category.
\end{theorem}

\begin{proof}
Apply the determinant functor to
\eqref{eq:v28-five-term-pm}.  For any exact finite complex, the alternating
product of determinant lines is canonically trivial.  With the convention
\eqref{eq:v28-Det-convention}, this gives
$\operatorname{Det}(D_\pm)\simeq
 \operatorname{Det}(D_0)\otimes(\det E_\times^\pm)^*$.
Eliminate $\operatorname{Det}(D_0)$ between the two formulas to obtain
\eqref{eq:v28-line-transition}.  The parity of a determinant line is the
dimension of its underlying vector space modulo two; dualization does not
change parity.  Hence
$r_-+r_+=m_\times$ gives the displayed grading.
\end{proof}

\begin{corollary}[Saturation-writer transition]
\label{cor:v28-writer-transition}
Choose local determinant frames
\begin{equation}
 \eta_\pm\in\det E_\times^\pm.
 \label{eq:v28-eta-pm}
\end{equation}
If $\omega_-$ is a local writer in
$\operatorname{Det}(D_-)^*$, then the corresponding plus-side writer is
\begin{equation}
 \boxed{
 \omega_+
 =\omega_-\otimes\eta_-^*\otimes\eta_+,}
 \label{eq:v28-writer-transition}
\end{equation}
with the order of the exterior factors fixed as displayed.  Every retained
tangential source derivative is obtained by differentiating this finite
Kato/exterior packet.
\end{corollary}

\subsection{Finite source/BV wall packet}
\label{sec:v28-source-packet}

The normal direction remains discontinuous, so V28 retains the V27 jump
operator rather than differentiating the sign through zero:
\begin{equation}
 \mathfrak J_\times F:=F^{[+]}-F^{[-]}.
 \label{eq:v28-jump-operator}
\end{equation}
The product rule is
\begin{equation}
 \mathfrak J_\times(FG)
 =(\mathfrak J_\times F)G^{[-]}
 +F^{[+]}(\mathfrak J_\times G).
 \label{eq:v28-jump-Leibniz}
\end{equation}

\begin{theorem}[Source-complete finite-rank wall packet]
\label{thm:v28-source-packet}
At every prescribed finite source/BV order, the regular rank-$m_\times$ wall is
represented by the finite packet
\begin{equation}
 \boxed{
 (\Pi_\times,A_\times,J_\times,R_\times^\pm,
  E_\times^\pm,\eta_\pm,
  \mathfrak J_\times\varepsilon,
  \mathfrak J_\times\widehat P_-,
  \mathfrak J_\times D_{\rm ov},
  \omega_-,\omega_+),}
 \label{eq:v28-source-packet}
\end{equation}
where
\begin{align}
 \mathfrak J_\times\varepsilon&=2J_\times,
 \label{eq:v28-J-epsilon}\\
 \mathfrak J_\times\widehat P_-&=J_\times,
 \label{eq:v28-J-P}\\
 \mathfrak J_\times D_{\rm ov}&={2\over a}\gamma_5J_\times.
 \label{eq:v28-J-Dov}
\end{align}
All tangential source derivatives are finite Riesz/Kato/exterior derivatives;
no retained coefficient contains an inverse crossing eigenvalue or a
distributional derivative of the Wilson sign.
\end{theorem}

\begin{proof}
The only inverse spectral operations are on the uniformly gapped Wilson
complement and on the invertible crossing form $A_\times$.  Their differentiated
Riesz calculi are bounded by $g_\times^{-1}$ and $\mu_\times^{-1}$,
respectively.  The displayed jump identities follow immediately from
Theorem~\ref{thm:v28-resolved-sign}.  Exterior powers of the finite Kato bundles
supply $\eta_\pm$ and all fixed source descendants.
\end{proof}

\subsection{Hard-shell compatibility and continuum passage}
\label{sec:v28-shell-continuum}

The entire crossing bundle is retained in the low state.  Hard shell
elimination acts only on a uniformly invertible complement.

\begin{theorem}[Finite-rank wall transition commutes with Schur elimination]
\label{thm:v28-shell-commutes}
Let one or several UV hard shells lie in the gapped complement of
$E_\times$.  After exact Schur elimination the effective one-sided maps retain
the form
\begin{equation}
 D_{j,\pm}=[D_{j,0}\ \ C_{j,\pm}]
 \label{eq:v28-effective-pm}
\end{equation}
with the same crossing subbundles $E_\times^\pm$.  Consequently
\begin{equation}
 \boxed{
 \operatorname{Det}(D_{j,+})
 \simeq
 \operatorname{Det}(D_{j,-})\otimes\mathbb T_\times,}
 \label{eq:v28-shell-line}
\end{equation}
and this isomorphism is independent of the grouping of nested UV shells.
\end{theorem}

\begin{proof}
The Schur complement modifies only the common-core map and the crossing column
maps $C_\pm$; it never changes the domain bundles $E_\times^\pm$.  Apply
Theorem~\ref{thm:v28-line-transition} to the effective maps.  Associativity of
Schur complementation gives regrouping independence exactly as in V26--V27.
\end{proof}

\begin{theorem}[Coefficientwise continuum across a regular finite-rank wall]
\label{thm:v28-continuum}
Assume the V24 V2-consistent smooth soft data on the gapped complement and the
finite crossing cluster, and assume each resolved face lies in a V11 minimal
fixed-index reduced-gap component.  Then for every prescribed finite
$(L,\Delta,r,N,\mathcal P_{\rm ext})$:
\begin{enumerate}[label=\textup{(C\arabic*)},leftmargin=3em]
\item both resolved faces have source-complete coefficientwise continua;
\item the index jump is the cutoff-independent signature
$\operatorname{sign}(A_\times)$;
\item the continuum determinant lines obey
\begin{equation}
 \boxed{
 \operatorname{Det}_\infty^{[+]}
 \simeq
 \operatorname{Det}_\infty^{[-]}\otimes\mathbb T_{\times,\infty};}
 \label{eq:v28-continuum-line}
\end{equation}
\item all fixed tangential BV/source descendants of the wall writer transition
pass to the same limit.
\end{enumerate}
\end{theorem}

\begin{proof}
The proof is V27 with a finite crossing bundle replacing the single line.  The
crossing data are soft finite-rank Riesz/Kato bundles and survive every UV shell.
The V24 Feshbach theorem applies to the uniformly gapped complement.  The
finite-cutoff line transitions \eqref{eq:v28-shell-line} are exact and
projective, so they pass to the V6/V7 continuum limit together with the finite
source packet.
\end{proof}

\subsection{Real/Pfaffian finite crossing block}
\label{sec:v28-real-block}

For an independently real Majorana branch, assume the native reality structure
preserves the finite Wilson crossing bundle.  V27 already showed that odd real
domain parity cannot connect two nondegenerate skew/Pfaffian sectors.

\begin{theorem}[Oriented even-block Pfaffian transition]
\label{thm:v28-Pfaffian-transition}
Assume the real crossing subbundles $E_{\times,\mathbb R}^\pm$ have dimensions
$r_\pm^{\mathbb R}$ such that the one-sided Majorana domains remain even, and
use the V25 native orientations.  Then the complexified determinant transition
of Theorem~\ref{thm:v28-line-transition} admits a real Pfaffian square root:
there is an oriented real wall line $\mathbb P_\times$ satisfying
\begin{equation}
 \boxed{
 \mathbb P_\times^{\otimes2}
 \simeq
 \det_{\mathbb R}E_{\times,\mathbb R}^-
 \otimes
 (\det_{\mathbb R}E_{\times,\mathbb R}^+)^*,}
 \label{eq:v28-Pfaffian-square}
\end{equation}
and
\begin{equation}
 \boxed{
 \operatorname{PfLine}(A_+)
 \simeq
 \operatorname{PfLine}(A_-)
 \otimes\mathbb P_\times.}
 \label{eq:v28-Pfaffian-transition}
\end{equation}
The sign of the square root is fixed by the V25 carrier/kernel orientation and
is compatible with hard-shell Schur elimination.
\end{theorem}

\begin{proof}
For a real skew Fredholm family the Pfaffian line squares to the real
determinant line.  Apply the determinant-line transition to the complexified
crossing block and restrict to the real structure.  The V25 orientation of the
ambient Majorana carrier and the induced orientations of the kernel/cokernel
bundles fix the real square root on each connected resolved face.  The even
parity hypothesis excludes the V27 odd-dimensional obstruction.  Exact
Pfaffian Schur factorization gives shell compatibility.
\end{proof}

\begin{firewall}[No unoriented real square root]
The existence of a formal square root of the complex determinant transition is
not enough to define a real Pfaffian transition.  V28 uses the independent V25
orientation input.  If that real orientation data are absent, the appropriate
object remains the complex determinant line and its reality involution rather
than an arbitrarily signed Pfaffian scalar.
\end{firewall}

\subsection{Clean codimension-two intersection of two Wilson walls}
\label{sec:v28-codim-two}

Let $(n_1,n_2,y)$ be local coordinates near an intersection of two Wilson-gap
walls.  A general multi-parameter Hermitian degeneracy need not split into two
independent wall bundles.  V28 treats the load-bearing \emph{clean transverse}
case.

\begin{definition}[Clean split double wall]
\label{def:v28-clean-double-wall}
A codimension-two point is a \emph{clean split double wall} if there are smooth
mutually orthogonal finite-rank Riesz bundles $E_1,E_2$ and a uniformly gapped
complement such that:
\begin{enumerate}[label=\textup{(D\arabic*)},leftmargin=3em]
\item the two crossing clusters remain spectrally separated from each other
away from the intersection and from the complement;
\item on $E_i$ the regular crossing form in the normal direction $n_i$ is an
invertible Hermitian operator $A_i$ with sign $J_i$;
\item to first normal order, the $n_1$ crossing acts only on $E_1$ and the
$n_2$ crossing acts only on $E_2$; equivalently, the off-diagonal normal
crossing forms between $E_1$ and $E_2$ vanish at the intersection;
\item every required tangential/source derivative obeys the V1 weighted local
bounds.
\end{enumerate}
Write
\begin{equation}
 E_i^\pm:=\operatorname{Ran}{1\over2}(I\pm J_i),
 \qquad
 m_i:=\dim E_i,
 \qquad
 \Delta k_i:=\dim E_i^+-\dim E_i^-.
 \label{eq:v28-double-Epm}
\end{equation}
\end{definition}

There are four resolved sectors labelled
$\delta=(\delta_1,\delta_2)\in\{\pm1\}^2$, with domain
\begin{equation}
 \boxed{
 W_{\delta_1\delta_2}
 =W_0
 \oplus E_1^{\delta_1}
 \oplus E_2^{\delta_2}.}
 \label{eq:v28-four-domains}
\end{equation}
Here $E_i^{+1}=E_i^+$ and $E_i^{-1}=E_i^-$.

For wall $i$ define the graded transition line
\begin{equation}
 \mathbb T_i
 :=\det E_i^-\otimes(\det E_i^+)^*,
 \qquad
 |\mathbb T_i|=m_i\pmod2.
 \label{eq:v28-Ti}
\end{equation}

\begin{theorem}[Graded wall-intersection coherence]
\label{thm:v28-double-wall-coherence}
At a clean split double wall, both orders of elementary modification identify
the determinant line in the $(+,+)$ sector with the $(-,-)$ sector tensored by
the two wall lines.  More precisely, the two composites are
\begin{align}
 \operatorname{Det}_{++}
 &\simeq
 \operatorname{Det}_{--}\otimes\mathbb T_2\otimes\mathbb T_1,
 \label{eq:v28-path12}\\
 \operatorname{Det}_{++}
 &\simeq
 \operatorname{Det}_{--}\otimes\mathbb T_1\otimes\mathbb T_2.
 \label{eq:v28-path21}
\end{align}
They agree after the canonical graded braiding
\begin{equation}
 \boxed{
 \beta_{12}:
 \mathbb T_2\otimes\mathbb T_1
 \longrightarrow
 \mathbb T_1\otimes\mathbb T_2,
 \qquad
 \beta_{12}(u_2\otimes u_1)
 =(-1)^{m_1m_2}u_1\otimes u_2.}
 \label{eq:v28-Koszul-braiding}
\end{equation}
Equivalently, if one scalarizes all saturation writers using one fixed ordinary
ordering of crossing frames, reversing the order of the two wall crossings
produces the unavoidable Koszul factor
\begin{equation}
 \boxed{
 (-1)^{m_1m_2}
 =(-1)^{\Delta k_1\Delta k_2}.}
 \label{eq:v28-Koszul-sign}
\end{equation}
The total index jump is order independent:
\begin{equation}
 \boxed{\Delta k=\Delta k_1+\Delta k_2.}
 \label{eq:v28-double-index}
\end{equation}
\end{theorem}

\begin{proof}
Because the crossing bundles are orthogonal and split to first normal order,
the four chiral domains are exactly \eqref{eq:v28-four-domains}.  Apply
Theorem~\ref{thm:v28-line-transition} first to one bundle and then the other.
The first order of crossing gives
$\operatorname{Det}_{++}\simeq
 \operatorname{Det}_{--}\otimes\mathbb T_2\otimes\mathbb T_1$;
the reverse order gives the second displayed tensor order.

The determinant functor is symmetric monoidal only in the category of graded
lines.  Interchanging lines of parities $m_1$ and $m_2$ therefore contributes
$(-1)^{m_1m_2}$ and makes the square commute.  Finally
$m_i=r_i^++r_i^-$ and
$\Delta k_i=r_i^+-r_i^-$ have the same parity, proving the second expression
in \eqref{eq:v28-Koszul-sign}.  Index additivity follows directly from the
rank formula for \eqref{eq:v28-four-domains}.
\end{proof}

\begin{corollary}[Cech cocycle on a clean wall arrangement]
\label{cor:v28-wall-Cech}
For a finite arrangement of pairwise clean transverse Wilson walls with no
higher non-split degeneracy, assign to each oriented wall its graded transition
line $\mathbb T_i$.  The V28 wall maps satisfy the codimension-two Cech
compatibility condition in the Picard groupoid of graded lines.  Thus the
stratified determinant-line family obtained by gluing the fixed-index sectors
is well defined independently of the order in which adjacent walls are
crossed.
\end{corollary}

\begin{proof}
The only nontrivial two-cell compatibility is the interchange of two transverse
wall modifications.  Theorem~\ref{thm:v28-double-wall-coherence} supplies the
canonical graded braiding on every such square.  Coherence of the symmetric
monoidal braiding gives consistency on iterated pairwise intersections.
\end{proof}

\begin{corollary}[Source-writer intersection law]
\label{cor:v28-writer-intersection}
Let $\eta_{i,\pm}$ be ordered determinant frames for $E_i^\pm$.  The two orders
of applying the saturation-writer transition
\eqref{eq:v28-writer-transition} differ, in an ungraded scalar convention, by
exactly the same factor $(-1)^{m_1m_2}$.  In the graded writer line they are the
same section.  All fixed tangential BV/source derivatives obey the same law.
\end{corollary}

\subsection{Why a general conical intersection is a different problem}
\label{sec:v28-nonsplit}

\begin{firewall}[No fictitious pair of walls inside a non-split degeneracy]
At a genuinely non-Abelian conical degeneracy the crossing kernel may not
admit smooth subbundles $E_1\oplus E_2$ associated with two independent wall
normals.  The matrix pencil
\[
 n_1A_1+n_2A_2
\]
may have noncommuting $A_1,A_2$, and its sign can depend on the approach angle.
In that situation there are no four canonical sectors of the form
\eqref{eq:v28-four-domains}; forcing two independent rank-one or finite-rank
wall moves would invent structure not present in the native Wilson family.
The correct object is the real oriented blow-up of the full discriminant in the
space of Hermitian crossing forms, with angular sign data on its boundary.
\end{firewall}

\begin{proposition}[Angular sign map for a non-split two-parameter crossing]
\label{prop:v28-angular-sign}
Let $E_\times$ be a finite crossing bundle and suppose
\begin{equation}
 H_W(n_1,n_2)|_{E_\times}
 =n_1A_1+n_2A_2+O(|n|^2)
 \label{eq:v28-pencil}
\end{equation}
with Hermitian $A_1,A_2$.  On every angular direction
$\theta\in S^1$ for which
\begin{equation}
 A(\theta):=\cos\theta\,A_1+\sin\theta\,A_2
 \label{eq:v28-angular-form}
\end{equation}
is invertible, the one-sided blow-up boundary carries the well-defined sign
\begin{equation}
 \boxed{J(\theta)=\operatorname{sgn}A(\theta).}
 \label{eq:v28-angular-sign}
\end{equation}
The index on that boundary arc is
$k(\theta)=k_0+\frac12\operatorname{Tr}J(\theta)$ and changes only when
$A(\theta)$ becomes singular.  Thus a non-split intersection reduces to a new
stratified problem on the angular discriminant of the Hermitian pencil.
\end{proposition}

\begin{proof}
The radial expansion is
$H_W(\rho\cos\theta,\rho\sin\theta)|_{E_\times}
 =\rho A(\theta)+O(\rho^2)$.  For invertible $A(\theta)$, gapped functional
calculus gives the radial sign limit $J(\theta)$.  Its trace is locally constant
until an eigenvalue of $A(\theta)$ crosses zero.
\end{proof}

\subsection{V28 anti-circularity audit}
\label{sec:v28-audit}

\begin{enumerate}[leftmargin=2.4em]
\item \textbf{The finite-multiplicity crossing is not diagonalized by hand.}
All internal mixing is encoded invariantly by the Hermitian crossing form
$A_\times$ and its sign $J_\times$.
\item \textbf{The two one-sided domains need not be nested.}
They are both elementary modifications of one common core by the gained and
lost bundles $E_\times^\pm$.
\item \textbf{The determinant transition uses the full exterior powers.}
The rank-one V27 line factor is replaced by
$\det E_-\otimes(\det E_+)^*$.
\item \textbf{Codimension-two commutation is graded, not naively scalar.}
The Koszul factor $(-1)^{m_1m_2}$ is retained explicitly and is the same sign
seen by ordered Grassmann saturation writers.
\item \textbf{The real Pfaffian square root is conditional on V25 orientation.}
V28 does not choose an unoriented real square root of a determinant line.
\item \textbf{A non-split conical degeneracy is not forced into the clean
 double-wall theorem.}
Its angular Hermitian pencil is recorded as a separate stratified gate.
\end{enumerate}

\section{V28 status ledger}
\label{sec:v28-status}

\subsection*{Proved in V28}
\begin{enumerate}[leftmargin=2.2em]
\item A regular finite-rank Wilson wall is controlled by the invertible
Hermitian crossing form
$A_\times=\Pi_\times(\partial_nH_W)\Pi_\times$; its sign $J_\times$ gives the
matrix-valued resolved Wilson sign without choosing individual crossing
eigenvectors.
\item The two one-sided chiral domains are
$W_0\oplus E_\times^+$ and $W_0\oplus E_\times^-$, and the genuine overlap
index jump is the signature
$\operatorname{sign}(A_\times)=\dim E_\times^+-\dim E_\times^-$.
\item The Fredholm determinant-line transition is the higher exterior-power
modification
$\operatorname{Det}(D_+)\simeq
 \operatorname{Det}(D_-)\otimes
 \det E_\times^-\otimes(\det E_\times^+)^*$.
\item The complete finite source/BV wall packet consists of the crossing Riesz
bundle, crossing form/sign projectors, gained/lost determinant frames and the
finite jump writers.  No inverse crossing eigenvalue or distributional sign
derivative occurs.
\item The higher-rank transition commutes exactly with native hard-shell Schur
elimination and passes to the coefficientwise continuum under the V24
V2-consistent soft hypotheses.
\item For an independently real branch with even one-sided Majorana domain
parity, the determinant transition admits the oriented V25 Pfaffian square root
and composes with Pfaffian Schur elimination.
\item At a clean split codimension-two intersection, the two wall-transition
orders agree in the graded determinant-line category.  In a scalarized fixed
writer ordering they differ by the explicit Koszul sign
$(-1)^{m_1m_2}=(-1)^{\Delta k_1\Delta k_2}$.
\item A finite clean wall arrangement therefore carries a coherent stratified
graded determinant-line/source cocycle independent of crossing order.
\item For a non-split two-parameter crossing, the correct boundary datum is the
angular sign map $J(\theta)=\operatorname{sgn}(\cos\theta A_1+\sin\theta A_2)$;
its further angular discriminant is isolated as a separate problem rather than
being misidentified with two independent walls.
\end{enumerate}

\subsection*{Inherited}
\begin{enumerate}[leftmargin=2.2em]
\item V1--V7: source-complete coefficientwise continuum, common BV restoration,
fixed-loop forest induction, ultraviolet Cauchy convergence, scheme/reference
covariance and exact perturbative-order projectivity.
\item V8--V24: regulator comparison, controlled family-index classification,
soft Riesz anchoring and the V2-consistent covariant G\aa rding/Feshbach theorem.
\item V25: native orientation and mod-two spectral-flow theorem for the
separately real neutral-Majorana branch.
\item V26: regular divisor packet for accidental kernel--cokernel pair
crossings at fixed Fredholm index.
\item V27: simple rank-one Wilson wall, genuine unit index jump and rank-one
Fredholm-line elementary modification.
\end{enumerate}

\subsection*{Conditional}
\begin{enumerate}[leftmargin=2.2em]
\item The finite wall is regular in the sense that the crossing form
$A_\times$ is invertible.  A tangent/higher-order Wilson contact requires a
higher-order normal form.
\item The coefficientwise continuum theorem requires the V24 smooth
V2-consistent soft chiral jet on the gapped complement and finite crossing
cluster, plus one-sided V11 reduced gaps.
\item The real Pfaffian statement requires even real parity and the V25 native
orientation data.
\item The codimension-two coherence theorem assumes a clean split intersection
with smooth spectrally separated bundles $E_1,E_2$.  A non-split Hermitian
matrix pencil is not covered by that theorem.
\item The positive IR/reference assumptions remain in force.
\end{enumerate}

\subsection*{Open beyond V28}
\begin{enumerate}[leftmargin=2.2em]
\item Resolve a genuinely non-split multi-parameter Wilson degeneracy by
blowing up the full Hermitian crossing-form discriminant and construct the
graded determinant/Pfaffian line on its angular stratification.
\item Treat higher-order/tangent Wilson contacts where the first normal crossing
form is singular and the leading nonzero crossing polynomial has degree greater
than one.
\item Globalize arbitrary finite wall arrangements with triple and higher
non-clean intersections and prove higher coherence of the resulting graded
Picard/source cocycle.
\item Remove the positive IR/reference scale and separately address analytic
finite-coupling RG trajectories, perturbative summability and nonperturbative
completion.
\end{enumerate}

\subsection*{Exact next load-bearing theorem}
\begin{center}
\fbox{\parbox{0.92\textwidth}{
\textbf{Non-split Hermitian-pencil blow-up theorem.}
Let a finite Wilson crossing bundle meet zero over two or more parameters with
leading Hermitian pencil $A(n)=\sum_i n_iA_i$ and no global splitting into
independent wall bundles.  Blow up the parameter origin, stratify the sphere of
normal directions by the singular locus of $A(\hat n)$, and construct on each
regular angular stratum the resolved sign, chiral domain, controlled index and
graded determinant/Pfaffian transition.  Prove compatibility across the angular
walls and around their intersections, with all fixed source/BV descendants.
This would extend the clean V28 wall arrangement to a general finite
multi-parameter Wilson discriminant.}}
\end{center}

\section*{Append-only continuation marker: V28}
\addcontentsline{toc}{section}{Append-only continuation marker: V28}
\textbf{End of V28.}  The complete V1--V27 body above is retained verbatim.
V28 replaces individual crossing eigenvectors by the invariant finite-rank
Hermitian crossing form.  Its sign splits the crossing bundle into gained and
lost chiral-domain subbundles, giving an index jump equal to the crossing-form
signature and a higher exterior-power determinant-line modification.  The
finite wall packet is source complete, commutes with hard-shell elimination and
passes to the coefficientwise continuum.  At a clean codimension-two
intersection the two elementary modifications satisfy the exact graded
commutation law; scalarized writer orderings differ by the explicit Koszul sign
$(-1)^{m_1m_2}$.  The next stronger gate is the non-split Hermitian-pencil
blow-up problem.



\clearpage
\section{V29 continuation: non-split Hermitian-pencil blow-up, angular Berry bundles, and constructible determinant lines}
\label{sec:v29-continuation}

V28 reduced a genuinely non-split Wilson degeneracy to the angular sign map of
one finite Hermitian pencil.  The present continuation carries out that
blow-up.  The main upstream correction is that angular-wall combinatorics are
not the whole story.  Even when the angular pencil is invertible in every
normal direction, its positive spectral bundle can be topologically nontrivial.
Thus the correct boundary object is a \emph{constructible graded determinant
line with a Kato/Berry connection}, not merely a list of index jumps.

Fix a parameter neighbourhood of a Wilson degeneracy of codimension $r\ge2$.
After removing spectator parameters tangent to the degeneracy stratum, write the
normal parameter as
\begin{equation}
 n=(n_1,\ldots,n_r)\in\mathbb R^r,
 \qquad
 \rho:=|n|,
 \qquad
 \widehat n:=n/|n|\in S^{r-1}.
 \label{eq:v29-normal-coordinates}
\end{equation}
Let $H_W(n)$ be the native Hermitian Wilson kernel.  At $n=0$ suppose that
zero is isolated from the complementary spectrum and that
\begin{equation}
 E_\times:=\ker H_W(0),
 \qquad
 \dim_{\mathbb C}E_\times=m_\times<\infty,
 \label{eq:v29-crossing-space}
\end{equation}
while
\begin{equation}
 |H_W(0)|_{E_\times^\perp}\succeq g_\times I
 \label{eq:v29-complement-gap}
\end{equation}
for some $g_\times>0$.  All tangential background/source parameters remain in
one compact V1--V28 protected chart.

\subsection{Exact finite-dimensional reduction before blowing up}
\label{sec:v29-finite-reduction}

The first step is to remove a possible ambiguity which would otherwise infect
the angular pencil: the crossing space itself moves with the background.  One
must first identify it by native Riesz/Kato transport.

\begin{lemma}[Kato-trivialized crossing block]
\label{lem:v29-Kato-block}
For $|n|$ sufficiently small there is a Riesz projector
$\Pi_\times(n)$ of rank $m_\times$ containing precisely the Wilson eigenvalues
inside $(-g_\times/3,g_\times/3)$, and a unitary Kato transport
\begin{equation}
 U_\times(n):E_\times\xrightarrow{\simeq}\Ran\Pi_\times(n).
 \label{eq:v29-Kato-transport}
\end{equation}
In the resulting fixed crossing frame,
\begin{equation}
 h(n):=
 U_\times(n)^*H_W(n)U_\times(n)\big|_{E_\times}
 \label{eq:v29-h-block}
\end{equation}
is an $m_\times\times m_\times$ Hermitian matrix family and
\begin{equation}
 H_W(n)
 \simeq
 h(n)\oplus H_W^\perp(n)
 \label{eq:v29-exact-block}
\end{equation}
by a unitary block decomposition, with
$|H_W^\perp(n)|\succeq g_\times/2$.
Moreover
\begin{equation}
 \boxed{
 h(n)=A(n)+R_2(n),
 \qquad
 A(n):=\sum_{j=1}^r n_jA_j,
 \qquad
 \|R_2(n)\|\le C|n|^2,}
 \label{eq:v29-linear-pencil}
\end{equation}
where
\begin{equation}
 A_j=\Pi_\times(0)(\partial_{n_j}H_W)(0)\Pi_\times(0)
 \big|_{E_\times}.
 \label{eq:v29-Aj}
\end{equation}
All prescribed fixed tangential source/BV derivatives satisfy the same
expansion coefficientwise.
\end{lemma}

\begin{proof}
The fixed complement gap permits a common Riesz contour around the cluster at
zero.  The V1/V13 resolvent calculus gives a smooth exponentially local Riesz
projector and Kato unitary with every fixed source derivative.  Since
$\Pi_\times(n)$ is a spectral projector of $H_W(n)$, the transported crossing
space and its orthogonal complement are exactly reducing subspaces, which gives
\eqref{eq:v29-exact-block}.  Taylor expansion of the finite matrix
\eqref{eq:v29-h-block} at $n=0$ gives \eqref{eq:v29-linear-pencil}.  The
Kato gauge can be chosen with vanishing connection at the base point, so the
linear coefficient is exactly the compression in \eqref{eq:v29-Aj}; a
different smooth Kato gauge conjugates the entire pencil and does not change any
spectral bundle or determinant statement below.
\end{proof}

\begin{definition}[Angular Hermitian pencil and discriminant]
\label{def:v29-angular-pencil}
On the oriented real blow-up of $n=0$, define
\begin{equation}
 A(\widehat n):=\sum_{j=1}^r\widehat n_jA_j,
 \qquad
 \widehat n\in S^{r-1},
 \label{eq:v29-angular-pencil}
\end{equation}
and its angular discriminant
\begin{equation}
 \boxed{
 \Sigma_\times
 :=\{\widehat n\in S^{r-1}:\det A(\widehat n)=0\}.}
 \label{eq:v29-angular-discriminant}
\end{equation}
Put
\begin{equation}
 \Omega_\times:=S^{r-1}\setminus\Sigma_\times.
 \label{eq:v29-regular-angular-set}
\end{equation}
Since the determinant is a real polynomial in the sphere coordinates,
$\Sigma_\times$ is compact semialgebraic and admits a finite Whitney
stratification.
\end{definition}

\subsection{The radial blow-up boundary and its spectral bundles}
\label{sec:v29-radial-boundary}

Write $n=\rho\widehat n$.  Lemma~\ref{lem:v29-Kato-block} gives
\begin{equation}
 \rho^{-1}h(\rho\widehat n)
 =A(\widehat n)+\rho\,R(\rho,\widehat n),
 \label{eq:v29-renormalized-block}
\end{equation}
with $R$ uniformly bounded together with every retained tangential source jet.

\begin{theorem}[Resolved angular sign]
\label{thm:v29-angular-sign}
On the regular angular set $\Omega_\times$, define
\begin{equation}
 \boxed{
 J(\widehat n):=\operatorname{sgn}A(\widehat n),
 \qquad
 P_+(\widehat n):={I+J(\widehat n)\over2},
 \qquad
 P_-(\widehat n):={I-J(\widehat n)\over2}.}
 \label{eq:v29-angular-projectors}
\end{equation}
Then:
\begin{enumerate}[label=\textup{(A\arabic*)},leftmargin=3em]
\item $J$ and $P_\pm$ are smooth on every connected component of
$\Omega_\times$ and have finite Kato/source derivatives of every prescribed
order;
\item for every compact $K\Subset\Omega_\times$ there is $\rho_K>0$ such that
$h(\rho\widehat n)$ is invertible for
$0<\rho<\rho_K$, $\widehat n\in K$, and
\begin{equation}
 \boxed{
 \operatorname{sgn}h(\rho\widehat n)
 \xrightarrow[\rho\downarrow0]{}
 J(\widehat n)}
 \label{eq:v29-sign-limit}
\end{equation}
uniformly with all retained tangential source derivatives;
\item the positive and negative angular spectral spaces
\begin{equation}
 E_\pm:=\Ran P_\pm\longrightarrow\Omega_\times
 \label{eq:v29-angular-bundles}
\end{equation}
are smooth Hermitian vector bundles, not necessarily trivial.
\end{enumerate}
\end{theorem}

\begin{proof}
On a compact subset of $\Omega_\times$ the smallest absolute eigenvalue of
$A(\widehat n)$ has a positive minimum $\mu_K$.  For sufficiently small
$\rho$, \eqref{eq:v29-renormalized-block} differs from $A(\widehat n)$ by less
than $\mu_K/2$.  Riesz functional calculus therefore gives a common positive
and negative spectral split and convergence of the corresponding Riesz
projectors.  Fixed derivatives follow from differentiated resolvent words.
\end{proof}

Let $W_0$ denote the common chiral domain after the crossing cluster has been
removed.  On a regular angular point the blow-up boundary chiral domain is
\begin{equation}
 \boxed{
 W(\widehat n)=W_0\oplus E_+(\widehat n).}
 \label{eq:v29-angular-domain}
\end{equation}
Consequently, if $k_0:=\operatorname{ind}(D_0)$ is the index of the common-core
map, then
\begin{equation}
 \boxed{
 k(\widehat n)=k_0+\operatorname{rank}E_+(\widehat n).}
 \label{eq:v29-angular-index}
\end{equation}
Equivalently, relative index differences are one-half trace differences of
$J$, in agreement with V28.

\subsection{The angular determinant superline and Berry curvature}
\label{sec:v29-angular-line}

The non-split problem already contains topology before one reaches an angular
wall.

\begin{theorem}[Sector determinant line on the angular regular set]
\label{thm:v29-sector-line}
Let $D_0:W_0\to\mathcal H_+$ be the common-core chiral map.  On
$\Omega_\times$ there is a canonical line isomorphism
\begin{equation}
 \boxed{
 \operatorname{Det}(D_{\widehat n})
 \simeq
 \operatorname{Det}(D_0)
 \otimes
 \bigl(\det E_+(\widehat n)\bigr)^*.}
 \label{eq:v29-sector-line}
\end{equation}
Thus the intrinsic angular regulator line is
\begin{equation}
 \boxed{
 \mathbb T_{\rm ang}:=(\det E_+)^*
 \longrightarrow\Omega_\times,}
 \label{eq:v29-angular-line-bundle}
\end{equation}
with parity
$|\mathbb T_{\rm ang}|=\operatorname{rank}E_+\pmod2$ on each angular chamber.
\end{theorem}

\begin{proof}
The domain is the direct sum \eqref{eq:v29-angular-domain}.  Apply the same
finite determinant-line exact-sequence argument used in V27--V28 to the
extension of $D_0$ by the finite domain bundle $E_+$.  With the lineage
convention
$\operatorname{Det}(D)=\det\operatorname{coker}D\otimes
(\det\ker D)^*$, adding the finite domain factor contributes
$(\det E_+)^*$.
\end{proof}

Choose local unitary Kato frames for $E_+$.  Let $\nabla^+$ be the induced
Berry connection and $\nabla^{\rm ang}$ the dual connection on
$\mathbb T_{\rm ang}$.

\begin{proposition}[Projector curvature formula]
\label{prop:v29-Berry-curvature}
On every smooth angular chamber,
\begin{equation}
 \boxed{
 \mathcal F_+
 =\operatorname{Tr}
 \bigl(P_+\,dP_+\wedge dP_+\bigr),
 \qquad
 \mathcal F_{\rm ang}=-\mathcal F_+.}
 \label{eq:v29-Berry-curvature}
\end{equation}
Hence
\begin{equation}
 \boxed{
 c_1(\mathbb T_{\rm ang})
 =-\left[{1\over2\pi i}
 \operatorname{Tr}
 (P_+dP_+\wedge dP_+)\right].}
 \label{eq:v29-angular-c1}
\end{equation}
In particular, a nonzero class obstructs any global nonvanishing scalar angular
saturation writer even when the index is constant and no angular wall is
present.
\end{proposition}

\begin{proof}
For an orthogonal projector $P_+$ the universal Grassmann connection on its
range has curvature $P_+dP_+\wedge dP_+P_+$.  Taking the trace gives the
curvature of its determinant line.  Dualization reverses the sign.  The Chern
class formula is the ordinary Chern--Weil normalization.
\end{proof}

\subsection{A nonsplit conical example with no angular wall and nontrivial line}
\label{sec:v29-Pauli}

The following example is the basic reason V29 cannot reduce to angular-wall
combinatorics.

\begin{proposition}[Pauli pencil: constant index, nontrivial angular line]
\label{prop:v29-Pauli}
Let $r=3$, $E_\times=\mathbb C^2$, and
\begin{equation}
 \boxed{
 A(\widehat n)
 =\widehat n_1\sigma_1
  +\widehat n_2\sigma_2
  +\widehat n_3\sigma_3,
 \qquad
 \widehat n\in S^2.}
 \label{eq:v29-Pauli-pencil}
\end{equation}
Then
\begin{equation}
 A(\widehat n)^2=I,
 \qquad
 \Sigma_\times=\varnothing,
 \qquad
 \operatorname{rank}E_+=1
 \label{eq:v29-Pauli-gap}
\end{equation}
for every angle.  Nevertheless the positive eigenline has
\begin{equation}
 \boxed{c_1(E_+)=1,\qquad c_1(\mathbb T_{\rm ang})=-1.}
 \label{eq:v29-Pauli-c1}
\end{equation}
Thus there is no global nonzero angular eigenvector and no global scalarization
of the crossing contribution, despite the complete absence of angular walls or
index jumps.
\end{proposition}

\begin{proof}
The Pauli identity gives
$(\widehat n\cdot\sigma)^2=I$, proving invertibility and rank one of
$P_+=(I+\widehat n\cdot\sigma)/2$.  Using
$\operatorname{Tr}(\sigma_a\sigma_b\sigma_c)=2i\epsilon_{abc}$,
\begin{align}
 \operatorname{Tr}(P_+dP_+\wedge dP_+)
 &= {i\over4}\epsilon_{abc}
    \widehat n_a\,d\widehat n_b\wedge d\widehat n_c
 \notag\\
 &= {i\over2}\,\Omega_{S^2},
 \label{eq:v29-Pauli-curvature}
\end{align}
where
$\Omega_{S^2}={1\over2}\epsilon_{abc}\widehat n_a
 d\widehat n_b\wedge d\widehat n_c$ has integral $4\pi$.  Therefore
\begin{equation}
 {1\over2\pi i}\int_{S^2}
 \operatorname{Tr}(P_+dP_+\wedge dP_+)=1.
 \end{equation}
Proposition~\ref{prop:v29-Berry-curvature} gives the dual-line class.
\end{proof}

\begin{firewall}[Index jumps do not classify a nonsplit conical crossing]
The Pauli pencil has constant signature and no angular discriminant, yet its
angular determinant line is topologically nontrivial.  Therefore the chamber
index function, the set of angular walls, and the V28 Koszul wall signs are not
a complete invariant of a general nonsplit Wilson blow-up.  One must retain the
full positive spectral bundle $E_+$ and its determinant/Berry data.
\end{firewall}

\subsection{Regular angular walls as a recursive V28 problem}
\label{sec:v29-angular-walls}

We now recover the V28 wall theorem intrinsically on the angular discriminant.
Let $S$ be a smooth codimension-one stratum of $\Sigma_\times$ on which
\begin{equation}
 K_S:=\ker A(\widehat n)
 \label{eq:v29-angular-kernel}
\end{equation}
has constant rank.  Choose a coorientation and a unit angular normal field
$\nu$ to $S$ inside $S^{r-1}$.

\begin{definition}[Angular crossing form]
\label{def:v29-angular-crossing-form}
The angular crossing form on $K_S$ is
\begin{equation}
 \boxed{
 B_S
 :=\Pi_{K_S}(\partial_\nu A)\Pi_{K_S}\big|_{K_S}.}
 \label{eq:v29-angular-crossing-form}
\end{equation}
The wall stratum is called \emph{regular} if $B_S$ is invertible.
Put
\begin{equation}
 J_S:=\operatorname{sgn}B_S,
 \qquad
 K_S^\pm:=\Ran{I\pm J_S\over2}.
 \label{eq:v29-angular-wall-split}
\end{equation}
\end{definition}

\begin{theorem}[Angular wall elementary modification]
\label{thm:v29-angular-wall}
Let $S$ be a regular angular wall, and orient the normal coordinate $t$ so that
$t>0$ is the plus side.  Then, after one smooth Kato identification of the
nonzero angular spectral complement,
\begin{align}
 E_+^{(+)}&=E_{+,0}\oplus K_S^+,
 \label{eq:v29-Eplus-plus}\\
 E_+^{(-)}&=E_{+,0}\oplus K_S^-.
 \label{eq:v29-Eplus-minus}
\end{align}
Consequently
\begin{equation}
 \boxed{
 \Delta_S k
 =\dim K_S^+-\dim K_S^-
 =\operatorname{sign}(B_S),}
 \label{eq:v29-angular-index-jump}
\end{equation}
and the wall transition of the angular determinant line is
\begin{equation}
 \boxed{
 \mathbb T_S
 :=\det K_S^-
   \otimes(\det K_S^+)^*,
 \qquad
 \operatorname{Det}_{+}
 \simeq
 \operatorname{Det}_{-}\otimes\mathbb T_S.}
 \label{eq:v29-angular-wall-line}
\end{equation}
The finite jump of the angular sign is
\begin{equation}
 \boxed{
 J^{(+)}-J^{(-)}=2J_S
 \quad\text{on }K_S,}
 \label{eq:v29-angular-sign-jump}
\end{equation}
with zero jump on the common gapped angular complement.
\end{theorem}

\begin{proof}
Choose a local angular coordinate $t$ normal to $S$.  Relative to
$K_S\oplus K_S^\perp$, the Hermitian matrix has the block expansion
\begin{equation}
 A(t)=
 \begin{pmatrix}
  tB_S+O(t^2)&tC+O(t^2)\\
  tC^*+O(t^2)&A_0^\perp+O(t)
 \end{pmatrix},
 \label{eq:v29-angular-block-expansion}
\end{equation}
where $A_0^\perp$ is invertible.  Its Schur complement on $K_S$ is
$tB_S+O(t^2)$.  Hence for small $t>0$ the crossing-cluster sign is $J_S$ and
for small $t<0$ it is $-J_S$.  Positive vectors on the two sides are therefore
$K_S^+$ and $K_S^-$ respectively.  The index and line formulas are exactly the
V28 finite-rank elementary-modification formulas applied on the angular sphere.
\end{proof}

\begin{corollary}[Strict transform of a regular angular wall]
\label{cor:v29-strict-transform-wall}
For the full renormalized block
$A(\widehat n)+\rho R(\rho,\widehat n)$, every compact piece of a regular
angular wall has, for sufficiently small $\rho$, a unique nearby smooth
strict-transform wall.  Its crossing form is $B_S+O(\rho)$ and hence has the
same signature and the same determinant transition line up to canonical Kato
transport.
\end{corollary}

\begin{proof}
The invertibility of $B_S$ is precisely transversality of the small eigenvalue
cluster to the zero set in the chosen normal direction.  The finite-dimensional
implicit-function theorem and Kato perturbation theory therefore persist the
wall and its spectral split for small $\rho$.
\end{proof}

\subsection{Constructible determinant superline on the angular stratification}
\label{sec:v29-constructible-line}

Choose a finite Whitney stratification of $\Sigma_\times$ refining the loci of
constant kernel rank.  Let $(C_a)_a$ be the connected components of
$\Omega_\times$.

\begin{definition}[Angular constructible determinant object]
\label{def:v29-constructible-object}
On each chamber $C_a$ retain the graded line bundle
\begin{equation}
 \mathcal L_a
 :=\operatorname{Det}(D_0)\otimes(\det E_{+,a})^*.
 \label{eq:v29-chamber-line}
\end{equation}
Across every regular codimension-one angular wall $S$ separating chambers
$a,b$, glue them by the exterior-power elementary modification
\eqref{eq:v29-angular-wall-line}, with its natural parity.  Chamberwise Kato
transport and the V28 graded braiding are part of the gluing data.
\end{definition}

The constructible object carries two kinds of information which must not be
confused:
\begin{equation}
 \boxed{
 \text{wall modification data}
 \quad+\quad
 \text{intrachamber Berry topology/holonomy}.}
 \label{eq:v29-two-angular-data}
\end{equation}

\begin{theorem}[Gallery composition and endpoint relative line]
\label{thm:v29-gallery-composition}
Let $\gamma$ be a generic angular path whose endpoints lie in regular chambers
and which crosses only regular codimension-one angular walls transversely.
Along the open chamber pieces use Kato transport of $E_+$; at each wall use
Theorem~\ref{thm:v29-angular-wall}.  The resulting determinant transition is
canonically isomorphic to
\begin{equation}
 \boxed{
 \mathbb T_\gamma
 \simeq
 \det E_+(\gamma(0))
 \otimes
 \bigl(\det E_+(\gamma(1))\bigr)^*,}
 \label{eq:v29-gallery-relative-line}
\end{equation}
where the endpoint determinant fibres are understood after the chamberwise
Kato transports.  Refining the gallery by inserting additional regular wall
crossings or subdividing chamber segments does not change the result.  At a
clean split wall intersection this theorem reduces exactly to the V28 Koszul
braiding law.
\end{theorem}

\begin{proof}
On every chamber segment the determinant line is the dual top exterior power
of the positive spectral bundle.  Crossing a regular wall replaces the lost
positive crossing space $K_S^-$ by the gained one $K_S^+$ and hence tensors by
$\det K_S^-\otimes(\det K_S^+)^*$.  Multiplying these factors along the path,
all intermediate determinant fibres cancel by evaluation, leaving only the
initial determinant and the dual of the final determinant.  In the graded
category, every interchange of finite exterior-power factors is performed by
the canonical Koszul braiding, which is exactly the V28 sign at a clean
intersection.  Associativity and symmetry coherence of graded lines make any
refinement canonical.
\end{proof}

\begin{corollary}[Closed galleries measure Berry/constructible holonomy]
\label{cor:v29-closed-gallery}
For a generic closed angular gallery, the product of all wall factors and
chamber Kato transports equals the holonomy of the constructible angular
determinant superline.  It need not be the identity.  In particular, when
$\Sigma_\times=\varnothing$, the entire obstruction is the ordinary Berry
holonomy/Chern class of $(\det E_+)^*$, as in
Proposition~\ref{prop:v29-Pauli}.
\end{corollary}

\begin{firewall}[Higher coherence does not mean trivial holonomy]
The graded determinant functor guarantees that two local decompositions of the
same sequence of elementary modifications agree after the canonical braiding.
It does \emph{not} imply that the constructible angular line is globally
trivial.  Nonzero chamber Chern classes or Berry holonomy survive all local
Koszul coherence identities.
\end{firewall}

\subsection{Angular source/BV packet}
\label{sec:v29-source-packet}

The blow-up introduces no singular source derivative.  The singularity is in
the normal direction and is handled by finite wall transitions.

\begin{definition}[Finite angular crossing packet]
\label{def:v29-angular-packet}
At fixed finite source/BV order, the V29 angular packet consists of
\begin{equation}
 \boxed{
 \left(
 \Pi_\times,
 U_\times,
 A(\widehat n),
 J(\widehat n),
 P_\pm(\widehat n),
 E_\pm,
 \nabla^+,
 \mathbb T_{\rm ang};
 \{K_S,B_S,K_S^\pm,\mathbb T_S\}_{S}
 \right),}
 \label{eq:v29-angular-packet}
\end{equation}
where $S$ ranges over the regular angular wall strata met by the chosen finite
panel.  Every prescribed tangential gauge/coframe/Higgs/antifield/line source
derivative is retained.
\end{definition}

\begin{theorem}[Source-complete regularity on the blow-up]
\label{thm:v29-source-regularity}
On each regular angular stratum, every fixed source/BV derivative of the packet
\eqref{eq:v29-angular-packet} is finite and quasilocal in the inherited native
norms.  Across a regular angular wall, normal differentiation is replaced by
the finite jump packet
\begin{equation}
 \mathfrak J_SJ=2J_S,
 \qquad
 \mathfrak J_SP_+=J_S
 \label{eq:v29-angular-jump-packet}
\end{equation}
on $K_S$, together with the determinant-line factor $\mathbb T_S$ and its
Kato/exterior-power source derivatives.  No inverse angular eigenvalue and no
distributional derivative of $\operatorname{sgn}A$ occurs.
\end{theorem}

\begin{proof}
Every regular-stratum object is defined by finite-dimensional Riesz functional
calculus separated by a positive angular spectral gap, so the V1/V13
resolvent-word argument applies.  The wall-normal singularity is exactly the
finite-rank V28 elementary modification proved in
Theorem~\ref{thm:v29-angular-wall}; retaining the jump rather than differentiating
through zero removes every would-be inverse crossing eigenvalue.
\end{proof}

\subsection{Compatibility with native shell elimination and continuum limits}
\label{sec:v29-shell-continuum}

The complete crossing cluster remains in the soft state throughout the UV
shell construction.  Hence the angular spectral bundle itself is untouched by
hard-shell integration.

\begin{theorem}[Cofinality of the angular constructible line]
\label{thm:v29-cofinal-angular-line}
Assume the V24 smooth V2-consistent hard/soft hypotheses on the gapped
complement and the finite crossing cluster.  Then:
\begin{enumerate}[label=\textup{(C\arabic*)},leftmargin=3em]
\item every exact hard-shell Schur elimination leaves the angular bundles
$E_\pm$, the chamber Berry connection and every regular-wall transition line
$\mathbb T_S$ unchanged up to the canonical Kato identification;
\item the chamberwise coefficientwise continuum exists by the inherited
V24/V27/V28 theorem;
\item the continuum determinant lines glue by the same constructible angular
superline, including its Berry holonomy and all Koszul wall signs;
\item the result is independent of regrouping the UV shells.
\end{enumerate}
\end{theorem}

\begin{proof}
The crossing block is retained among the low variables, while every hard shell
lies in its uniformly invertible complement.  Schur elimination acts only on
the common-core chiral map and therefore factors out before the finite domain
bundle $E_+$.  The determinant functor then leaves the angular factor
$(\det E_+)^*$ untouched.  V28 already proves the same statement for regular
wall exterior powers.  The V24 continuum theorem applies to the common-core
coefficients and commutes with the finite Kato/exterior-power operations.
\end{proof}

\subsection{Real/Pfaffian angular data}
\label{sec:v29-real-angular}

Suppose in addition that the crossing bundle carries a real structure
preserved by the angular pencil and that the independently physical Majorana
sector of V25 is present.  Then the positive angular spectral spaces define a
real bundle $E_{+,\mathbb R}$ on the real regular set.

\begin{proposition}[Real angular orientation obstruction]
\label{prop:v29-real-angular}
The real determinant orientation line of the angular positive bundle has
\begin{equation}
 \boxed{
 w_1\bigl(\det_{\mathbb R}E_{+,\mathbb R}\bigr)
 =w_1(E_{+,\mathbb R})\in H^1(\Omega_\times;\mathbb Z_2).}
 \label{eq:v29-real-w1}
\end{equation}
A global Pfaffian scalarization of the angular crossing contribution requires
this orientation datum, together with the V25 orientation of the full native
Majorana carrier.  Across regular real angular walls the V28 even-parity and
Pfaffian square-root conditions remain in force.
\end{proposition}

\begin{proof}
The first Stiefel--Whitney class of a real vector bundle equals that of its top
exterior power.  The Pfaffian statement is exactly the V25/V28 real-line
firewall applied to the angular positive bundle.  Orientation of the full
carrier does not, by itself, trivialize an arbitrary real spectral subbundle.
\end{proof}

\subsection{V29 anti-circularity audit}
\label{sec:v29-audit}

\begin{enumerate}[leftmargin=2.4em]
\item \textbf{The non-split pencil is not artificially diagonalized.}
All angular data are expressed through the invariant projector
$P_+=(I+\operatorname{sgn}A)/2$ and its spectral bundle.
\item \textbf{Wall combinatorics are not treated as a complete invariant.}
The Pauli pencil has no angular wall at all but carries a nontrivial Hopf
spectral line.
\item \textbf{The higher-intersection coherence statement is graded.}
V28 Koszul signs are recovered as the local braiding of exterior-power wall
factors; they are not dropped when scalar writers are ordered.
\item \textbf{Coherence is distinguished from triviality.}
A constructible line can satisfy every local cocycle identity and still carry
nonzero Berry curvature/Chern class.
\item \textbf{The blow-up uses the actual finite-regulator crossing block.}
The common Kato trivialization is obtained from the native Wilson Riesz
projector before the linear pencil is formed.
\item \textbf{Singular angular strata are not forced into the regular-wall
formula.}
Theorem~\ref{thm:v29-angular-wall} requires an invertible angular crossing
form; tangent/higher-order angular contacts remain a separate gate.
\end{enumerate}

\section{V29 status ledger}
\label{sec:v29-status}

\subsection*{Proved in V29}
\begin{enumerate}[leftmargin=2.2em]
\item The finite native Wilson crossing cluster admits an exact Kato-trivialized
block whose first normal jet is the invariant Hermitian pencil
$A(n)=\sum_jn_jA_j$; the gapped Wilson complement is removed exactly before
blowing up.
\item On the regular angular set
$\Omega_\times=S^{r-1}\setminus\{\det A=0\}$, the radial boundary sign is
$J(\widehat n)=\operatorname{sgn}A(\widehat n)$ and defines smooth positive and
negative spectral bundles $E_\pm$ with all retained source derivatives.
\item The angular contribution to the Fredholm determinant line is the graded
line bundle $(\det E_+)^*$.  Its Kato/Berry curvature is
$-\operatorname{Tr}(P_+dP_+\wedge dP_+)$, so it can be globally nontrivial even
when the index is constant.
\item The Pauli pencil $A(\widehat n)=\widehat n\cdot\sigma$ has no angular
discriminant and no index jump, but its positive eigenline has $c_1=1$ and the
angular determinant line has $c_1=-1$.  This proves that angular Berry topology
is an independent nonsplit obstruction beyond wall jumps and Koszul signs.
\item Every regular codimension-one angular wall has its own finite crossing
form $B_S=\Pi_K(\partial_\nu A)\Pi_K$.  Its signature gives the angular index
jump, and its gained/lost bundles give the V28 exterior-power transition line
$\det K_S^-\otimes(\det K_S^+)^*$.
\item The chamber line bundles, their Kato transports and the regular-wall
exterior-power maps form a constructible graded determinant object on a Whitney
stratification of the angular discriminant.  Generic gallery composition is
canonical and recovers V28 graded codimension-two coherence.
\item Closed angular galleries measure the actual Berry/constructible holonomy;
local graded coherence does not force this holonomy to vanish.
\item The complete finite angular source/BV packet is regular on every angular
stratum and uses finite wall jumps rather than singular derivatives of the sign
function.
\item The entire angular constructible line is cofinal under native hard-shell
elimination and passes to the coefficientwise continuum under the V24 smooth
V2-consistent hypotheses.
\item In a separately real branch, the angular positive bundle can carry an
independent real orientation class $w_1$, which must be combined with the V25
native Majorana orientation data.
\end{enumerate}

\subsection*{Inherited}
\begin{enumerate}[leftmargin=2.2em]
\item V1--V7: source-complete coefficientwise continuum, common BV restoration,
fixed-loop forests, ultraviolet Cauchy convergence, scheme/reference covariance
and exact perturbative-order projectivity.
\item V8--V24: regulator universality, controlled family-index classification,
soft spectral anchoring and the smooth V2-consistent covariant
G\aa rding/Feshbach theorem.
\item V25: native real orientation and mod-two spectral-flow theorem for a
separately physical gauge-neutral Majorana branch.
\item V26: regular accidental divisor crossing at fixed Fredholm index.
\item V27: simple Wilson-sign crossing and rank-one index jump.
\item V28: regular finite-rank wall, higher exterior-power line transition and
clean split codimension-two graded coherence.
\end{enumerate}

\subsection*{Conditional}
\begin{enumerate}[leftmargin=2.2em]
\item The leading nonzero normal jet of the Wilson crossing cluster is linear.
A higher-order/tangent radial contact requires replacing the pencil by its
leading homogeneous crossing polynomial.
\item Regular angular-wall formulas require an invertible angular crossing form
$B_S$.  Singular angular strata are included in the constructible
stratification but do not yet have their own native local normal form.
\item The continuum statement retains the V24 smooth V2-consistent soft data on
the common core and finite crossing block.
\item The real/Pfaffian statement requires a preserved real structure, the V25
carrier orientation and the V28 even-parity condition.
\item The positive IR/reference assumptions remain in force.
\end{enumerate}

\subsection*{Open beyond V29}
\begin{enumerate}[leftmargin=2.2em]
\item Determine whether the \emph{complete} faithful Standard-Model fermion
packet trivializes the angular Berry/constructible determinant class.  The
single-block Pauli example shows that such trivialization is not automatic from
index constancy.
\item Resolve singular/tangent strata of the angular discriminant where the
angular crossing form is degenerate, by the leading nonzero normal polynomial
on the angular kernel.
\item Extend the constructible line to arbitrary higher-codimension angular
singularities with a fully global descent theorem on the blown-up parameter
space.
\item Remove the positive IR/reference scale and separately address analytic
finite-coupling RG trajectories, perturbative summability and nonperturbative
completion.
\end{enumerate}

\subsection*{Exact next load-bearing theorem}
\begin{center}
\fbox{\parbox{0.92\textwidth}{
\textbf{Complete-packet angular Berry cancellation / regulator-loop theorem.}
For the full faithful Standard-Model chiral determinant packet, place the V29
angular constructible line inside the inherited global determinant-line family
and determine whether its Chern/holonomy class vanishes on every resolved
Wilson crossing sphere.  Match the hypotheses of the inherited faithful-spin
eta/bordism trivialization to this regulator-parameter blow-up rather than
assuming they apply automatically.  If the class vanishes, construct the global
angular saturation writer and prove that the entire nonsplit Wilson crossing is
regulator-line trivial after complete Standard-Model assembly.  If it does not,
identify the surviving microscopic Berry class and its image under the V16
continuum-holonomy map.}}
\end{center}

\section*{Append-only continuation marker: V29}
\addcontentsline{toc}{section}{Append-only continuation marker: V29}
\textbf{End of V29.}  The complete V1--V28 body above is retained verbatim.
V29 resolves a non-split finite Hermitian Wilson pencil on the oriented radial
blow-up.  The regular angular boundary carries smooth positive/negative
spectral bundles and an intrinsic determinant superline $(\det E_+)^*$ with its
Kato/Berry connection.  Regular angular walls reproduce the V28 finite-rank
crossing-form elementary modification, while generic gallery composition is
coherent in the graded line category.  The Pauli pencil proves that this wall
calculus is not complete by itself: the angular determinant line can have
nonzero Chern class even when there are no angular walls or index jumps.  The
next gate is to determine whether the complete faithful Standard-Model packet
cancels this angular Berry/constructible line class on regulator-parameter
blow-ups.


\clearpage
\section{V30 continuation: complete-packet angular Berry cancellation under a faithful geometric lift}
\label{sec:v30-continuation}

V29 showed that a nonsplit Wilson crossing can carry a genuine angular Berry
class even when there are no angular walls and no index jump.  The next question
is therefore not another local crossing calculation but whether the
\emph{complete} faithful Standard-Model determinant packet trivializes that
constructible angular line.  The inherited Grand Tensor contains a strong
answer on a smooth ordinary-spin faithful family: the complete chiral
determinant line has a unit parallel section, because the local anomaly
polynomial vanishes and the faithful $G_6$ eta character is trivial.  What is
not automatic is that the blown-up regulator family of V29 lies in the precise
geometric family category to which that theorem applies.

V30 separates those two statements.  It proves that, once a faithful geometric
lift of the resolved regulator family exists and its exact finite crossing lines
are identified with the V29 constructible transitions, the entire complete
Standard-Model angular determinant superline is globally trivial.  All chamber
Berry curvature, closed-gallery holonomy and graded wall cocycles cancel in the
complete packet.  It also proves that ordinary anomaly-trace cancellation by
itself is not enough to infer this result: the missing input is the geometric
family lift and determinant-line comparison, not another representation trace.

\subsection{The resolved angular base and its constructible determinant line}
\label{sec:v30-angular-base}

Let $\widetilde B_\times$ denote one compact resolved regulator-parameter
neighborhood obtained by radially blowing up a finite Wilson crossing locus and
then choosing a finite Whitney stratification of the angular discriminant as in
V29.  Its open regular chambers are denoted $C_a$, and its regular codimension
one angular walls by $S$.

For one chiral species $f$, V29 constructs on each chamber the line
\begin{equation}
 \mathcal L_{f,a}
 =\operatorname{Det}(D_{f,0})\otimes(\det E^+_{f,a})^*,
 \label{eq:v30-species-chamber-line}
\end{equation}
with wall transition
\begin{equation}
 \mathbb T_{f,S}
 =\det K^-_{f,S}\otimes(\det K^+_{f,S})^*.
 \label{eq:v30-species-wall-line}
\end{equation}
For the complete Standard-Model packet put
\begin{equation}
 \mathcal L^{\rm SM}_a
 :=\bigotimes_f\mathcal L_{f,a},
 \qquad
 \mathbb T^{\rm SM}_S
 :=\bigotimes_f\mathbb T_{f,S},
 \label{eq:v30-complete-angular-line}
\end{equation}
with the inherited graded tensor convention.  The resulting chamberwise lines
and wall maps form the complete constructible angular determinant superline
$\mathcal L^{\rm SM}_{\rm ang}$.

\begin{definition}[Faithful geometric lift of the resolved regulator family]
\label{def:v30-faithful-lift}
A \emph{faithful geometric lift} of the V29 resolved family consists of:
\begin{enumerate}[label=\textup{(G\arabic*)},leftmargin=3em]
\item a compact connected ordinary-spin $G_6$ family base $B_{\rm geom}$ and a
continuous map
\begin{equation}
 \iota:\widetilde B_\times\longrightarrow B_{\rm geom};
 \label{eq:v30-iota}
\end{equation}
\item a smooth full bulk Dirac--Yukawa family over
$B_{\rm geom}\times[q_0,Q]$ satisfying the hypotheses of the inherited smooth
faithful-phase theorem on every compact cutoff slab;
\item on every regular chamber $C_a$, an exact determinant-line isomorphism
\begin{equation}
 \Phi_a:\mathcal L^{\rm SM}_a
 \xrightarrow{\ \simeq\ }
 \iota^*\mathcal L^{\rm geom}|_{C_a};
 \label{eq:v30-Phi-chamber}
\end{equation}
\item across every regular angular wall, compatibility of $\Phi_a$ with the
exact finite relative determinant line: the V29 wall factor
$\mathbb T^{\rm SM}_S$ is identified with the restriction of the inherited
relative determinant-line comparison of the geometric family;
\item the same compatibility at clean higher intersections, including the V28
Koszul braiding.
\end{enumerate}
\end{definition}

This definition is deliberately stronger than equality of local anomaly
polynomials or of Fredholm indices.  It identifies the actual native resolved
line with the geometric determinant family to which the eta/bordism theorem
applies.

\subsection{Pullback of the faithful parallel section}
\label{sec:v30-parallel-pullback}

\begin{theorem}[Complete-packet angular trivialization under a faithful lift]
\label{thm:v30-angular-trivialization}
Assume Definition~\ref{def:v30-faithful-lift}.  Then the complete constructible
angular determinant superline $\mathcal L^{\rm SM}_{\rm ang}$ admits a global
unit section $s_{\rm ang}$ which is parallel on every regular chamber and is
carried across every regular wall by the V29/V28 elementary-modification map.
It is unique up to one constant phase on each connected resolved component.
Consequently:
\begin{align}
 c_1(\mathcal L^{\rm SM}_{\rm ang}|_{C_a})&=0,
 \label{eq:v30-c1-zero}\\
 \operatorname{Hol}_\gamma(\mathcal L^{\rm SM}_{\rm ang})&=1
 \label{eq:v30-hol-zero}
\end{align}
for every closed regular/constructible gallery $\gamma$.
\end{theorem}

\begin{proof}
The inherited smooth faithful-phase theorem supplies a unit parallel section
$s_{\rm geom}$ of the complete geometric chiral determinant line on the
connected ordinary-spin $G_6$ family, unique up to one constant phase.  Pull it
back by $\iota$ and then by the chamber isomorphisms
\eqref{eq:v30-Phi-chamber}:
\begin{equation}
 s_a:=\Phi_a^{-1}(\iota^*s_{\rm geom})
 \in\Gamma(\mathcal L^{\rm SM}_a).
 \label{eq:v30-local-parallel}
\end{equation}
Each $s_a$ is nonzero and parallel for the pulled-back determinant connection.
By item (G4), the two chamber sections adjacent to a regular wall are carried
into one another by the exact V29 wall transition.  Item (G5) makes this
identification coherent at higher intersections, including the graded Koszul
braiding.  Hence the $s_a$ glue to a global constructible section
$s_{\rm ang}$.

A global nonzero parallel section trivializes the complex line on every chamber,
so its first Chern class vanishes.  Parallel transport around a closed gallery
returns the same section; therefore the complete constructible holonomy is one.
Uniqueness follows from connectedness and the corresponding inherited uniqueness
of $s_{\rm geom}$.
\end{proof}

\begin{corollary}[Cancellation of complete angular Berry curvature]
\label{cor:v30-Berry-cancellation}
On every regular chamber,
\begin{equation}
 \boxed{
 \sum_f \operatorname{Tr}
 \bigl(P^+_f\,dP^+_f\wedge dP^+_f\bigr)
 +\mathcal F_{\rm core}^{\rm SM}=0,}
 \label{eq:v30-Berry-sum}
\end{equation}
where $\mathcal F_{\rm core}^{\rm SM}$ is the curvature of the complete common
core determinant factor in the same chamber trivialization.  In particular,
if the common core line is already in the inherited parallel frame, then the
sum of angular spectral-bundle curvatures vanishes identically.
\end{corollary}

\begin{proof}
The curvature of a tensor-product determinant line is the sum of the factor
curvatures.  Theorem~\ref{thm:v30-angular-trivialization} gives zero curvature
for the complete line.  V29 identifies the angular contribution of each
species with
$-\operatorname{Tr}(P_f^+dP_f^+\wedge dP_f^+)$.  Rearranging gives
\eqref{eq:v30-Berry-sum}.
\end{proof}

\subsection{Wall factors cancel as part of the same global section}
\label{sec:v30-wall-cancellation}

The phrase ``Berry cancellation'' must include the finite exterior-power
jumps.  A vanishing chamber Chern class alone would not control the wall
transitions of the constructible line.

\begin{theorem}[Triviality of the complete wall cocycle]
\label{thm:v30-wall-cocycle}
Under the faithful lift, choose local nonzero chamber frames by restricting the
global section $s_{\rm ang}$.  For every regular angular wall $S$, the complete
wall transition scalar in these frames is exactly one:
\begin{equation}
 \boxed{t_S^{\rm SM}=1.}
 \label{eq:v30-wall-scalar-one}
\end{equation}
At a clean codimension-two intersection, both wall-crossing orders give the
same scalar after the canonical graded braiding; the V28 Koszul sign is exactly
absorbed by that braiding rather than producing residual holonomy.
\end{theorem}

\begin{proof}
The chamber frames are restrictions of one global constructible section.  By
definition of gluing, the wall map carries the section from one chamber to the
section in the adjacent chamber, so its scalar coefficient in these frames is
one.  At a clean intersection, Definition~\ref{def:v30-faithful-lift}(G5)
identifies both compositions with the same geometric relative determinant map.
The graded determinant functor inserts the V28 Koszul braiding automatically,
so after this canonical identification the scalar triangle/square defect is
one.
\end{proof}

\begin{corollary}[Global angular saturation writer]
\label{cor:v30-global-writer}
Let $\omega_a$ be the dual of $s_a$ on each chamber.  The wall compatibility of
Theorem~\ref{thm:v30-wall-cocycle} makes these local writers glue to a global
constructible dual section
\begin{equation}
 \boxed{
 \omega_{\rm ang}\in
 \Gamma\bigl((\mathcal L^{\rm SM}_{\rm ang})^*\bigr),
 \qquad
 \langle\omega_{\rm ang},s_{\rm ang}\rangle=1.}
 \label{eq:v30-global-writer}
\end{equation}
Thus the complete Standard-Model nonsplit Wilson crossing admits a global
angular scalarization on the resolved family.
\end{corollary}

\subsection{Why anomaly-trace cancellation alone is insufficient}
\label{sec:v30-anomaly-firewall}

V29's Pauli example already shows that chamber Berry topology can live entirely
in regulator-parameter directions.  The ordinary four-dimensional anomaly
traces do not by themselves calculate this topology.

\begin{proposition}[Local anomaly cancellation is not the angular trivialization theorem]
\label{prop:v30-anomaly-insufficient}
The conditions
\begin{equation}
 I_6^{\rm SM}=0,
 \qquad
 \text{ordinary four-dimensional anomaly traces}=0
 \label{eq:v30-anomaly-zero}
\end{equation}
do not, by themselves, imply
$c_1(\mathcal L^{\rm SM}_{\rm ang})=0$ for an arbitrary microscopic regulator
pencil.  To infer the latter from the inherited faithful-spin result one must
also establish the geometric-family and relative-line identifications of
Definition~\ref{def:v30-faithful-lift}.
\end{proposition}

\begin{proof}
The anomaly polynomial controls the curvature of the determinant connection for
families obtained from the declared geometric Dirac--Yukawa construction.  A
finite microscopic Hermitian pencil can carry Berry curvature in regulator
parameter directions even when its physical gauge representation is held
fixed; the V29 Pauli pencil is the explicit rank-two model.  Nothing in the
ordinary anomaly traces identifies that finite pencil line with the geometric
determinant-family line.  Definition~\ref{def:v30-faithful-lift} supplies
precisely that missing identification.  Once it is supplied,
Theorem~\ref{thm:v30-angular-trivialization} applies.
\end{proof}

\begin{firewall}[No silent extension of the inherited eta theorem]
The inherited faithful-spin eta/bordism theorem is a theorem about a smooth
ordinary-spin $G_6$ determinant family with specified geometric and boundary
line data.  V30 does not declare every abstract regulator-parameter pencil to be
such a family.  The angular cancellation theorem is conditional on proving the
faithful geometric lift of Definition~\ref{def:v30-faithful-lift} for the native
Renewal regulator family.
\end{firewall}

\subsection{Relation with the V16 continuum-holonomy map}
\label{sec:v30-V16-map}

V16 separated microscopic regulator-loop topology from continuum scheme-groupoid
transport.  The V29 angular line gives a distinguished family of such loops.

\begin{theorem}[Image of angular loops under continuum holonomy]
\label{thm:v30-continuum-holonomy}
Let $\gamma$ be a closed regular/constructible loop in the resolved angular
family.  Under the faithful lift,
\begin{equation}
 \boxed{
 \operatorname{Hol}_{\rm cont}([\gamma])=1.}
 \label{eq:v30-continuum-holonomy}
\end{equation}
for the complete faithful Standard-Model determinant packet.  Thus every angular
Berry/wall loop lies in the kernel of the V16 continuum-holonomy map after
complete-packet assembly.
\end{theorem}

\begin{proof}
Theorem~\ref{thm:v30-angular-trivialization} gives unit constructible holonomy
around $\gamma$.  V16 identifies the determinant-line component of continuum
loop transport with the corresponding complete line holonomy.  Hence its image
is the identity.  Any remaining V16 loop class, if present, must therefore lie
in the purely microscopic scheme/local-operator kernel rather than in the
faithful complex determinant phase.
\end{proof}

\subsection{Finite-cutoff and continuum compatibility}
\label{sec:v30-cutoff}

The faithful lift is required uniformly on every compact cutoff slab, exactly as
in the inherited theorem.  The V29 angular line is already cofinal under hard
shell elimination, so the trivialization is cofinal as well.

\begin{theorem}[Cofinal global angular trivialization]
\label{thm:v30-cofinal-trivialization}
Assume the V24 smooth V2-consistent hypotheses and a faithful geometric lift on
every compact cutoff slab, with compatible relative determinant-line data under
cutoff inclusion.  Then the sections $s_{{\rm ang},q}$ may be chosen so that
\begin{equation}
 J^{\rm geom}_{q'q}s_{{\rm ang},q}
 =s_{{\rm ang},q'}
 \label{eq:v30-cutoff-section}
\end{equation}
exactly.  The algebraic direct-limit angular line is therefore the trivial
complex line with its unit parallel section, and the coefficientwise continuum
inherits the same global saturation writer.
\end{theorem}

\begin{proof}
The inherited smooth faithful-phase theorem gives exact unitary cutoff
comparisons with cocycle law and unit phase in the parallel frame.  The chamber
identifications and V29 cofinal wall lines are assumed compatible with those
relative determinant comparisons.  Pulling the parallel section through the
same maps yields \eqref{eq:v30-cutoff-section}.  The inherited cofinal
determinant-line theorem then gives the direct-limit section without an
infinite-product summability assumption.
\end{proof}

\subsection{Real/Pfaffian packet}
\label{sec:v30-real}

The complex faithful-spin theorem does not automatically trivialize an
independent real Pfaffian angular bundle.  V25 and V29 remain the relevant
real inputs.

\begin{proposition}[Real angular condition after complex cancellation]
\label{prop:v30-real-condition}
Suppose a separately physical real Majorana/Pfaffian angular packet is present.
Even when Theorem~\ref{thm:v30-angular-trivialization} trivializes the complete
complex determinant line, a real scalarization additionally requires
\begin{equation}
 w_1(E_{+,\mathbb R}^{\rm Maj})=0
 \label{eq:v30-real-w1}
\end{equation}
and the V25 native carrier orientation together with the V28 even-parity
Pfaffian condition across walls.  On the gauge-neutral natively oriented branch
for which these hypotheses hold, the real angular writer is globally
orientation-compatible with the complex trivialization.
\end{proposition}

\begin{proof}
The real orientation class is the first Stiefel--Whitney class of the real
positive spectral bundle, as proved in V29.  Complexification can kill this
class and therefore cannot replace the real orientation test.  When the class
vanishes and the V25/V28 orientation data are supplied, the Pfaffian square-root
line glues consistently and inherits the same wall/cutoff coherence.
\end{proof}

\subsection{V30 anti-circularity audit}
\label{sec:v30-audit}

\begin{enumerate}[leftmargin=2.4em]
\item \textbf{The faithful-spin theorem is not applied by notation alone.}
V30 states an explicit geometric-lift hypothesis identifying the native angular
line with the determinant family covered by the inherited eta theorem.
\item \textbf{Anomaly traces are not used as a substitute for regulator-line
matching.}  The Pauli pencil shows why regulator-parameter Berry curvature is a
separate datum before the geometric lift is proved.
\item \textbf{Wall cancellation and chamber Berry cancellation are both
included.}  The global section trivializes the complete constructible line,
not merely its chamber Chern class.
\item \textbf{Koszul signs are retained.}  Clean higher-wall intersections
commute in the graded determinant category; the scalar defect is one only after
that canonical braiding is included.
\item \textbf{The real Pfaffian problem is not inferred from complex
triviality.}  The V25/V29 real orientation gate remains explicit.
\item \textbf{Cutoff coherence is intrinsic.}  The inherited geometric cutoff
maps transport one parallel section exactly; no infinite phase product is
introduced.
\end{enumerate}

\section{V30 status ledger}
\label{sec:v30-status}

\subsection*{Proved in V30}
\begin{enumerate}[leftmargin=2.2em]
\item A precise faithful geometric-lift criterion is formulated for the V29
resolved regulator family.  It requires chamberwise identification with the
inherited smooth faithful determinant family and exact compatibility with every
finite wall/intersection relative line.
\item Under that criterion, the complete faithful Standard-Model angular
constructible determinant superline has a global unit parallel section, unique
up to one constant phase on each connected resolved component.
\item Consequently every chamber Chern class and every closed-gallery
constructible holonomy of the complete complex packet vanishes.
\item In the induced parallel chamber frames every complete wall transition
scalar is one; V28 Koszul signs are absorbed exactly by the graded determinant
braiding, so higher-intersection coherence is global rather than merely local.
\item The dual global angular saturation writer exists and is compatible with
all wall transitions and fixed source/BV descendants.
\item Local four-dimensional anomaly cancellation alone is not sufficient to
infer angular Berry triviality for an arbitrary microscopic pencil; the missing
input is the geometric determinant-family lift.
\item Under the faithful lift, angular regulator loops have trivial image under
the V16 determinant-line continuum-holonomy map.
\item The angular trivialization is cofinal under cutoff transport and passes to
the coefficientwise continuum without any infinite-product phase hypothesis.
\item A separately real Pfaffian packet still requires its own V25/V29 real
orientation data; complex faithful-spin trivialization does not remove that
mod-two condition.
\end{enumerate}

\subsection*{Inherited}
\begin{enumerate}[leftmargin=2.2em]
\item V1--V24: source-complete coefficientwise continuum, native regulator
universality, controlled family-index classification and the smooth
V2-consistent soft spectral theorem.
\item V25: native real Majorana/Pfaffian orientation and mod-two transport on
the gauge-neutral oriented branch.
\item V26--V28: accidental divisor packets, genuine Wilson index jumps,
finite-rank wall exterior-power transitions and graded clean-intersection
coherence.
\item V29: nonsplit Hermitian-pencil blow-up, chamber Berry bundle,
constructible angular determinant line and the Pauli/Hopf obstruction.
\item The Grand Tensor faithful ordinary-spin $G_6$ theorem: vanishing complete
local anomaly polynomial, trivial faithful eta character, and one global unit
parallel determinant section for its stated smooth geometric family.
\end{enumerate}

\subsection*{Conditional}
\begin{enumerate}[leftmargin=2.2em]
\item The main V30 cancellation theorem is conditional on proving
Definition~\ref{def:v30-faithful-lift} for the actual V29 Renewal regulator
blow-up.  V30 does not infer this lift from anomaly traces alone.
\item The resolved family must remain in the ordinary-spin faithful $G_6$
category with the inherited bulk/boundary determinant-family and filling data.
\item Singular angular strata beyond the regular V29/V28 local normal forms
still require their own geometric relative-line identification.
\item The real/Pfaffian angular packet retains its independent orientation
condition.
\item The V24 smooth V2-consistent and positive IR/reference assumptions remain
in force.
\end{enumerate}

\subsection*{Open beyond V30}
\begin{enumerate}[leftmargin=2.2em]
\item Prove the \emph{native faithful geometric-lift theorem}: construct, from
the explicit Renewal Wilson/overlap regulator blow-up, the smooth bulk
Dirac--Yukawa determinant family and exact relative-line comparison required by
Definition~\ref{def:v30-faithful-lift}.
\item Resolve singular/tangent angular strata and verify that their native
higher-order transition lines agree with the corresponding geometric relative
determinant lines.
\item If the faithful lift fails for some microscopic regulator loop, compute
the surviving angular Berry/Chern class and its V16 microscopic-kernel class
explicitly.
\item Remove the positive IR/reference scale and separately address analytic
finite-coupling RG trajectories, perturbative summability and nonperturbative
completion.
\end{enumerate}

\subsection*{Exact next load-bearing theorem}
\begin{center}
\fbox{\parbox{0.92\textwidth}{
\textbf{Native faithful geometric lift of the resolved Wilson blow-up.}
Starting from the explicit V2 Renewal Wilson/overlap/Yukawa regulator and a V29
finite Hermitian crossing pencil, construct a smooth ordinary-spin $G_6$
bulk/boundary Dirac--Yukawa family over the resolved regulator base and prove an
exact determinant-line comparison between the native constructible angular line
and the inherited geometric determinant family, including every V28/V29 wall
and higher-intersection transition.  The proof must identify the finite native
crossing complex with the geometric relative determinant line rather than use
local anomaly cancellation as a surrogate.  Once this lift is proved,
Theorem~\ref{thm:v30-angular-trivialization} closes the complete faithful
Standard-Model angular Berry obstruction globally.}}
\end{center}

\section*{Append-only continuation marker: V30}
\addcontentsline{toc}{section}{Append-only continuation marker: V30}
\textbf{End of V30.}  The complete V1--V29 body above is retained verbatim.
V30 proves a complete-packet angular cancellation theorem conditional on one
precise missing bridge: a faithful geometric lift identifying the native V29
constructible regulator line with the smooth ordinary-spin $G_6$ determinant
family already known to have trivial eta holonomy.  Under that bridge the full
Standard-Model angular Berry curvature, wall cocycle and closed-gallery holonomy
all vanish, and a global angular saturation writer exists cofinally in the
cutoff.  The Pauli example shows why anomaly-trace cancellation alone cannot
supply this conclusion.  The next gate is therefore the native construction of
that geometric lift and exact relative determinant-line comparison.



\clearpage
\section{V31 continuation: index-constancy obstruction, stratified geometric lift, and the relative phase line}
\label{sec:v31-continuation}

V30 isolated one missing bridge between the native V29 angular determinant
superline and the inherited smooth faithful geometric determinant family.  A
direct attempt to construct the V30 lift exposes two distinct upstream
corrections.

First, a smooth family of Fredholm operators on a connected parameter space
cannot realize a nonzero jump of the Fredholm index.  Therefore a resolved
native Wilson family which contains a V27/V28 wall with
$\Delta k\ne0$ cannot be identified with one smooth connected geometric
Fredholm family across that wall.  The geometric comparison must be
\emph{stratified}: one smooth geometric family on each fixed-index chamber,
with a relative determinant-line/cobordism comparison across the wall.

Second, an isomorphism of complex line bundles does not identify their natural
connections.  Pulling back a geometric parallel section through an arbitrary
line isomorphism produces a global nonzero native section, but it need not be
parallel for the native Kato/Quillen connection and need not have the same loop
holonomy.  V30's chamber item (G3) must therefore be strengthened from an exact
line-bundle isomorphism to an exact \emph{phase-faithful} comparison, or else
its topological and holonomy conclusions must be separated.

V31 proves both corrections precisely.  It also proves that the V24 soft
norm-resolvent theorem supplies the complete \emph{topological} chamber bridge:
on every V2-consistent fixed-index chamber the native and geometric soft
zero-mode bundles are canonically isomorphic by the direct rotation.  The
remaining load-bearing object is one relative unitary line with connection,
whose curvature and loop character measure exactly the missing native-to-
geometric phase transgression.

\subsection{Index constancy obstructs a single smooth geometric lift across a Wilson wall}
\label{sec:v31-index-obstruction}

Let $B$ be connected and let
\begin{equation}
 D_b:\mathcal H_b^+\longrightarrow\mathcal H_b^-,
 \qquad b\in B,
 \label{eq:v31-Fredholm-family}
\end{equation}
be a norm-continuous family of Fredholm operators between Hilbert bundles (or a
smooth elliptic family in the geometric setting).

\begin{lemma}[Local constancy of the numerical Fredholm index]
\label{lem:v31-index-constant}
The integer
\begin{equation}
 \operatorname{ind}D_b
 =\dim\ker D_b-\dim\operatorname{coker}D_b
 \label{eq:v31-index}
\end{equation}
is locally constant in $b$, and hence constant on connected $B$.
\end{lemma}

\begin{proof}
Fix $b_0$.  Choose finite-dimensional subspaces which split the kernel and
cokernel of $D_{b_0}$ and invert the restriction of $D_{b_0}$ on their
orthogonal complements.  Invertibility is open in operator norm, so the same
finite-dimensional stabilization makes $D_b$ invertible for all $b$ near
$b_0$.  The stabilized finite-dimensional defect therefore has constant
rank-difference in that neighborhood.  This equals the Fredholm index.
Connectedness gives the global statement.
\end{proof}

\begin{theorem}[Index-jump obstruction to V30's connected smooth lift]
\label{thm:v31-index-obstruction}
Let two V29/V28 chambers $C_-$ and $C_+$ meet along a regular Wilson wall with
\begin{equation}
 \Delta k:=k_+-k_-\ne0.
 \label{eq:v31-nonzero-jump}
\end{equation}
There does not exist one connected smooth Fredholm family
\eqref{eq:v31-Fredholm-family} together with chamberwise determinant-line
identifications which realizes both native chiral families and preserves their
Fredholm indices on the two sides.

In particular, Definition~\ref{def:v30-faithful-lift} cannot be interpreted as
one smooth connected geometric Fredholm family across a genuine V27/V28 index
wall.  Its correct replacement is a stratified family with one fixed-index
geometric component on each chamber and a separate relative determinant-line
comparison across the wall.
\end{theorem}

\begin{proof}
If such a connected smooth family existed, Lemma~\ref{lem:v31-index-constant}
would give one constant integer index throughout the connected geometric base.
The chamberwise identifications would therefore force $k_-=k_+$, contradicting
\eqref{eq:v31-nonzero-jump}.
\end{proof}

\begin{corollary}[Fixed-index nonsplit chambers are the maximal smooth pieces]
\label{cor:v31-fixed-index-chambers}
A V29 resolved angular chamber (or any union of chambers joined only across
index-neutral accidental-divisor strata) may admit one smooth geometric
Fredholm lift.  Crossing a wall with nonzero V28 signature requires changing
geometric Fredholm component, operator domain/topological sector, or passing
through a non-Fredholm/geometric singularity.  These alternatives must be
recorded explicitly in the relative-line data.
\end{corollary}

\subsection{A line isomorphism does not identify determinant connections}
\label{sec:v31-connection-obstruction}

The next correction concerns the word ``parallel'' in V30.

\begin{proposition}[Topological line matching does not imply phase matching]
\label{prop:v31-line-not-phase}
Let $L=S^1\times\mathbb C$ be the trivial Hermitian line.  Give it two unitary
connections
\begin{equation}
 \nabla_0=d,
 \qquad
 \nabla_\theta=d+\frac{i\theta}{2\pi}\,d\varphi,
 \label{eq:v31-two-connections}
\end{equation}
where $\varphi$ is the angular coordinate.  The identity map is an exact
unitary line-bundle isomorphism, and both connections are flat, but
\begin{equation}
 \operatorname{Hol}_{\nabla_0}(S^1)=1,
 \qquad
 \operatorname{Hol}_{\nabla_\theta}(S^1)=e^{-i\theta}
 \label{eq:v31-different-holonomy}
\end{equation}
with the convention $\nabla s=0\Rightarrow
s(2\pi)=\exp(-\int A)s(0)$.
Hence an exact line-bundle isomorphism alone does not carry a parallel section
for one connection to a parallel section for the other.
\end{proposition}

\begin{proof}
The underlying line is trivial in both cases.  Parallel transport for
$\nabla_\theta$ solves
$\partial_\varphi s+(i\theta/2\pi)s=0$, so one turn multiplies $s$ by
$e^{-i\theta}$.  The result is not removable by a single-valued unitary
rephasing unless $e^{-i\theta}=1$.
\end{proof}

\begin{firewall}[Correction to the V30 lift criterion]
The chamberwise isomorphism in Definition~\ref{def:v30-faithful-lift}(G3)
proves equality of the underlying determinant-line bundle classes.  By itself
it does \emph{not} prove equality of the native and geometric determinant
connections or their loop holonomies.  Therefore every V30 conclusion involving
``parallel'', Berry curvature, or holonomy is to be read after the stronger
phase-faithful condition introduced below.  The purely topological consequence
(cancellation of the first Chern class when the geometric line is trivial)
requires only the line-bundle isomorphism.
\end{firewall}

\subsection{V24 supplies the topological chamber bridge}
\label{sec:v31-soft-bridge}

Fix one regular native chamber $C$ of constant Fredholm index.  Let
$D_{a}^{\rm nat}$ be the V2-consistent native chiral/Yukawa family and let
$D^{\rm geom}$ be the continuum geometric boundary family whose squared soft
operator is the common $L_0$ used in V22--V24.  Denote the protected kernel and
cokernel projectors by
\begin{equation}
 p_{a}^{K},\ p_a^C,
 \qquad
 p_0^K,\ p_0^C.
 \label{eq:v31-soft-projectors}
\end{equation}
Assume the V24 soft norm-resolvent estimate holds uniformly on compact subsets
of $C$, with the same fixed source derivatives.

\begin{lemma}[Canonical direct-rotation bundle isomorphisms]
\label{lem:v31-direct-rotation}
For sufficiently small $a$,
\begin{equation}
 \|p_a^K-p_0^K\|<1,
 \qquad
 \|p_a^C-p_0^C\|<1.
 \label{eq:v31-projector-close}
\end{equation}
The V14 direct-rotation formula therefore gives canonical source-smooth
unitaries
\begin{equation}
 U_a^K:\operatorname{Ran}p_a^K\xrightarrow{\simeq}
          \operatorname{Ran}p_0^K,
 \qquad
 U_a^C:\operatorname{Ran}p_a^C\xrightarrow{\simeq}
          \operatorname{Ran}p_0^C,
 \label{eq:v31-KC-rotations}
\end{equation}
with all retained source derivatives bounded in the same local calculus.
\end{lemma}

\begin{proof}
V24 and V22 give norm convergence of the native Riesz projectors to the common
continuum soft projectors.  Once the norm is below one, apply the explicit V14
polar direct rotation.  Its functional-calculus denominators have a uniform
positive floor, so the fixed source derivatives follow from the V13 weighted
resolvent calculus.
\end{proof}

Use the lineage convention
\begin{equation}
 \operatorname{Det}(D)
 =\det C\otimes(\det K)^*.
 \label{eq:v31-det-convention}
\end{equation}
Taking top exterior powers in \eqref{eq:v31-KC-rotations} gives
\begin{equation}
 \boxed{
 \Phi_a^{\rm top}:
 \operatorname{Det}(D_a^{\rm nat})
 \xrightarrow{\ \simeq\ }
 \operatorname{Det}(D^{\rm geom}).}
 \label{eq:v31-top-line-iso}
\end{equation}

\begin{theorem}[Topological faithful bridge on a fixed-index chamber]
\label{thm:v31-topological-bridge}
Under the V24 smooth V2-consistent hypotheses, the canonical map
\eqref{eq:v31-top-line-iso} is a smooth unitary line-bundle isomorphism on every
compact fixed-index chamber, compatible with all retained source derivatives.
Consequently
\begin{equation}
 \boxed{
 c_1\bigl(\operatorname{Det}D_a^{\rm nat}\bigr)
 =c_1\bigl(\operatorname{Det}D^{\rm geom}\bigr).}
 \label{eq:v31-c1-match}
\end{equation}
In particular, if the complete faithful geometric determinant line is
topologically trivial on that chamber, then so is the complete native line.
\end{theorem}

\begin{proof}
The direct rotations identify the kernel and cokernel bundles.  The determinant
functor applied to these finite-rank bundle isomorphisms gives
\eqref{eq:v31-top-line-iso}.  First Chern classes are natural under line-bundle
isomorphism.  Source compatibility follows from the source-smooth direct
rotations of Lemma~\ref{lem:v31-direct-rotation}.
\end{proof}

This proves the topological part of V30's desired chamber bridge without any
bulk interpolation.  It deliberately says nothing yet about the natural
connections.

\subsection{The relative phase line}
\label{sec:v31-relative-phase}

Transport the geometric determinant connection through
$\Phi_a^{\rm top}$ and compare it with the native Kato/determinant connection.
Define the unitary relative line
\begin{equation}
 \mathcal R_a
 :=\operatorname{Det}(D_a^{\rm nat})
   \otimes
   \bigl(\operatorname{Det}(D^{\rm geom})\bigr)^*.
 \label{eq:v31-relative-line}
\end{equation}
The topological bridge trivializes $\mathcal R_a$ as a complex line bundle, but
it carries the relative connection
\begin{equation}
 \nabla^{\rm rel}_a
 :=\nabla^{\rm nat}_a
   \otimes1
   +1\otimes(\nabla^{\rm geom})^*.
 \label{eq:v31-relative-connection}
\end{equation}
which need not be trivial.

Choose the canonical unit section $r_a$ of the topologically trivialized
relative line induced by $\Phi_a^{\rm top}$ and write
\begin{equation}
 \nabla_a^{\rm rel}r_a
 =\alpha_a\,r_a,
 \qquad
 \alpha_a\in\Omega^1(C;i\mathbb R).
 \label{eq:v31-alpha}
\end{equation}

\begin{proposition}[Relative curvature and holonomy]
\label{prop:v31-relative-phase}
The relative curvature is
\begin{equation}
 \boxed{
 \mathcal F_a^{\rm rel}
 =d\alpha_a
 =\mathcal F_a^{\rm nat}-\mathcal F^{\rm geom}.}
 \label{eq:v31-relative-curvature}
\end{equation}
For every closed loop $\gamma\subset C$,
\begin{equation}
 \boxed{
 \operatorname{Hol}_{\rm rel}(\gamma)
 :=\operatorname{Hol}_{\rm nat}(\gamma)
   \operatorname{Hol}_{\rm geom}(\gamma)^{-1}}
 \label{eq:v31-relative-hol-def}
\end{equation}
obeys
\begin{equation}
 \operatorname{Hol}_{\rm rel}(\gamma)
 =\exp\!\left(-\oint_\gamma\alpha_a\right)
 \label{eq:v31-relative-hol}
\end{equation}
with the convention of Proposition~\ref{prop:v31-line-not-phase}.
Changing the topological bridge by a unitary rephasing
$r_a\mapsto e^{i\chi}r_a$ changes
$\alpha_a\mapsto\alpha_a+i\,d\chi$ but leaves
\eqref{eq:v31-relative-curvature} and
\eqref{eq:v31-relative-hol-def} invariant.
\end{proposition}

\begin{proof}
Connections on a fixed complex line differ by an imaginary-valued one-form.
Taking the exterior derivative gives the curvature difference.  Parallel
transport in the relative line gives \eqref{eq:v31-relative-hol}.  The gauge
transformation formula for a unitary frame is standard and proves the final
statement.
\end{proof}

\begin{theorem}[Exact phase-faithfulness criterion on a fixed-index chamber]
\label{thm:v31-phase-faithful-criterion}
The topological bridge $\Phi_a^{\rm top}$ can be rephased to a
connection-preserving unitary isomorphism
\begin{equation}
 \Phi_a^{\rm ph}:
 (\operatorname{Det}D_a^{\rm nat},\nabla_a^{\rm nat})
 \xrightarrow{\ \simeq\ }
 (\operatorname{Det}D^{\rm geom},\nabla^{\rm geom})
 \label{eq:v31-phase-iso}
\end{equation}
if and only if
\begin{equation}
 \boxed{
 \mathcal F_a^{\rm rel}=0
 \quad\text{and}\quad
 \operatorname{Hol}_{\rm rel}(\gamma)=1
 \ \text{for every closed }\gamma\subset C.}
 \label{eq:v31-phase-criterion}
\end{equation}
Equivalently, the relative line $\mathcal R_a$ has a global nonzero parallel
section.
\end{theorem}

\begin{proof}
A connection-preserving isomorphism is exactly a nonzero parallel section of the
relative line.  Such a section forces zero curvature and unit loop holonomy.
Conversely, if the relative connection is flat and has trivial holonomy,
parallel transport of one nonzero vector from a base point is path independent
and gives a global parallel unit section.  Multiplying the original bridge by
the corresponding unitary function yields \eqref{eq:v31-phase-iso}.
\end{proof}

Thus V24 has already solved the \emph{bundle} part of the native-to-geometric
bridge.  The only chamberwise phase information still missing is the relative
connection/holonomy of \eqref{eq:v31-relative-line}.

\subsection{Stratified faithful geometric lift}
\label{sec:v31-stratified-lift}

The index obstruction now dictates the correct global definition.

\begin{definition}[Stratified phase-faithful geometric lift]
\label{def:v31-stratified-lift}
Let the V29 resolved native base be decomposed into regular fixed-index chambers
$C_a$ and their V28/V29 wall strata.  A \emph{stratified phase-faithful
geometric lift} consists of:
\begin{enumerate}[label=\textup{(S\arabic*)},leftmargin=3em]
\item for each chamber $C_a$ of index $k_a$, a smooth ordinary-spin faithful
$G_6$ geometric Dirac--Yukawa family on a connected geometric component
$B_a^{\rm geom}$ of the same index and a map
$\iota_a:C_a\to B_a^{\rm geom}$;
\item a connection-preserving unitary chamber isomorphism
\begin{equation}
 \Phi_a^{\rm ph}:
 \mathcal L_a^{\rm SM}
 \xrightarrow{\ \simeq\ }
 \iota_a^*\mathcal L_a^{\rm geom};
 \label{eq:v31-stratified-Phi}
\end{equation}
\item for every regular native wall $S$ between $C_a$ and $C_b$, a geometric
relative determinant line $\mathbb T_S^{\rm geom}$ and a unitary
connection-compatible identification
\begin{equation}
 \mathbb T_S^{\rm nat}
 \xrightarrow{\ \simeq\ }
 \mathbb T_S^{\rm geom}
 \label{eq:v31-wall-bridge}
\end{equation}
which intertwines the two determinant-line elementary modifications;
\item the corresponding graded compatibility at every clean higher
intersection, including the exact V28 Koszul braiding;
\item compatibility with the inherited cutoff maps and with all retained
source/BV derivatives.
\end{enumerate}
\end{definition}

The geometric wall comparison in (S3) may come from a cobordism, a change of
boundary condition/domain, or another declared relative determinant
construction.  The point is that it is \emph{not} one smooth Fredholm family
through a nonzero index jump.

\begin{theorem}[Corrected complete-packet angular trivialization]
\label{thm:v31-stratified-trivialization}
Assume a stratified phase-faithful geometric lift in the sense of
Definition~\ref{def:v31-stratified-lift}.  Assume further that the inherited
faithful-spin theorem supplies a unit parallel section on every geometric
fixed-index chamber component and that the declared geometric wall comparisons
carry these sections into one another, with the inherited graded coherence at
higher intersections.  Then the complete native constructible angular
determinant superline admits one global constructible unit section which is
parallel on every chamber and is transported exactly across every wall.
Consequently:
\begin{align}
 \mathcal F^{\rm SM}_{{\rm ang},a}&=0
 \quad\text{on every chamber }C_a,
 \label{eq:v31-angular-curvature-zero}\\
 \operatorname{Hol}^{\rm SM}_{\gamma}&=1
 \quad\text{for every closed constructible gallery }\gamma.
 \label{eq:v31-angular-hol-zero}
\end{align}
\end{theorem}

\begin{proof}
Pull the unit parallel geometric section back to each chamber using
$\Phi_a^{\rm ph}$.  Connection preservation makes the resulting native chamber
section parallel, not merely nonzero.  Item (S3) identifies the adjacent
sections under the native V29 wall transition.  Item (S4) gives equality of all
higher compositions in the graded determinant category.  Hence the local
sections glue to one constructible section.  Curvature and closed-gallery
holonomy vanish because this section is parallel chamberwise and globally
coherent across the walls.
\end{proof}

\begin{corollary}[V30 recovered with the strengthened hypotheses]
\label{cor:v31-V30-recovered}
On a resolved component with no nonzero-index wall, Definition
\ref{def:v31-stratified-lift} reduces to V30 after strengthening the V30
chamber isomorphism from a bare line-bundle map to a connection-preserving map.
Under that strengthened reading, Theorem~\ref{thm:v30-angular-trivialization}
and all of its Berry/holonomy consequences are valid exactly as stated.
\end{corollary}

\subsection{What V24 proves and what it does not}
\label{sec:v31-V24-scope}

The distinction is now exact.

\begin{theorem}[V24 closes the topological chamber gate, not the phase gate]
\label{thm:v31-V24-scope}
On every smooth V2-consistent fixed-index chamber for which the continuum
geometric boundary operator is the common V24 soft reference, the V24
norm-resolvent theorem implies the canonical line-bundle isomorphism
\eqref{eq:v31-top-line-iso} and therefore equality of determinant first Chern
classes.  It does not, by itself, imply
\begin{equation}
 \mathcal F_a^{\rm rel}=0
 \quad\text{or}\quad
 \operatorname{Hol}_{\rm rel}=1.
 \label{eq:v31-V24-not-phase}
\end{equation}
Those are precisely the additional phase-faithfulness conditions of
Theorem~\ref{thm:v31-phase-faithful-criterion}.
\end{theorem}

\begin{proof}
The first statement is Theorem~\ref{thm:v31-topological-bridge}.  Norm-closeness
of Riesz projectors supplies a canonical bundle isomorphism but does not impose
an equality of the two natural connections; Proposition
\ref{prop:v31-line-not-phase} shows that even flat isomorphic lines can have
different holonomy.  Therefore \eqref{eq:v31-V24-not-phase} is an independent
condition.
\end{proof}

\subsection{The minimal remaining bridge is a relative phase theorem}
\label{sec:v31-minimal-gate}

The preceding reduction is substantially smaller than constructing an entirely
new geometric bulk family from scratch.  The common continuum geometric family
already exists in the Grand Tensor programme.  V24 supplies its topological
soft identification with the native family.  What remains is to identify the
\emph{relative phase connection}.

\begin{definition}[Native-to-geometric relative phase problem]
\label{def:v31-relative-phase-problem}
For a fixed-index chamber, the native-to-geometric relative phase problem is to
prove that the relative determinant line \eqref{eq:v31-relative-line} has unit
parallel section, equivalently
\begin{equation}
 \boxed{
 \mathcal F_a^{\rm rel}=0,
 \qquad
 \operatorname{Hol}_{\rm rel}(\gamma)=1
 \ \forall\gamma.}
 \label{eq:v31-relative-phase-goal}
\end{equation}
Across an index wall, one must in addition identify the native V28/V29 wall line
with the declared geometric relative determinant line and verify the graded
higher-intersection coherence.
\end{definition}

A natural route is a five-dimensional native-to-geometric interpolation: the
overlap sign already suggests a domain-wall/APS cylinder, and the inherited
geometric theorem already expresses the desired phase through an eta invariant.
However, V31 does not import a domain-wall identity whose locality, source
calculus and boundary-line conventions have not been matched to the explicit
V2 regulator.

\begin{firewall}[No phase theorem from topological closeness alone]
The direct rotation of V24 proves that the native and continuum zero-mode
bundles are isomorphic on a fixed-index chamber.  It cannot by itself prove the
relative eta/holonomy statement \eqref{eq:v31-relative-phase-goal}.  Any future
proof must calculate the relative determinant connection or construct a
geometric interpolation whose determinant line is known to carry that
connection.
\end{firewall}

\subsection{A useful transgression identity for the next attack}
\label{sec:v31-transgression}

Although V31 does not yet prove \eqref{eq:v31-relative-phase-goal}, it records
the exact finite-dimensional transgression which the native calculation must
match.

Let $P_s$, $0\le s\le1$, be a smooth path of finite-rank orthogonal projectors
on one fixed Hilbert space, with $P_0=p_a$ and $P_1=p_0$, and let
\begin{equation}
 \Omega(P):=\operatorname{Tr}(P\,dP\wedge dP)
 \label{eq:v31-projector-curvature}
\end{equation}
be the determinant-bundle Berry curvature form on the external parameter base.

\begin{proposition}[Projector-curvature transgression]
\label{prop:v31-projector-transgression}
There exists the globally defined one-form
\begin{equation}
 \mathcal T(P_s)
 :=\int_0^1
 \operatorname{Tr}\bigl(
 P_s[\partial_sP_s,dP_s]
 \bigr)\,ds
 \label{eq:v31-transgression-form}
\end{equation}
(up to the conventional overall factor of $i$) such that
\begin{equation}
 \boxed{
 \Omega(P_1)-\Omega(P_0)
 =d\mathcal T(P_s).}
 \label{eq:v31-curvature-transgression}
\end{equation}
The same identity holds for the kernel-minus-cokernel determinant curvature by
subtracting the two transgressions.
\end{proposition}

\begin{proof}
Differentiate $P_s^2=P_s$ to obtain
$P_s(\partial_sP_s)P_s=0$ and similarly for $dP_s$.  Differentiating
$\operatorname{Tr}(P_s dP_s\wedge dP_s)$ in $s$, using cyclicity of the trace
and the projector identities, rearranges the result into the exterior
derivative of
$\operatorname{Tr}(P_s[\partial_sP_s,dP_s])$.  Integrate from $s=0$ to $1$.
The determinant convention is the cokernel contribution minus the kernel
contribution, so the final assertion follows by linearity.
\end{proof}

This proposition proves that the native/geometric curvature difference is
always exact once the projectors are joined by the V24 direct-rotation path.
It does \emph{not} make it zero.  The next theorem must show that the complete
Standard-Model relative transgression itself is the differential of the
appropriate native/geometric determinant phase with trivial loop periods.

\subsection{V31 anti-circularity audit}
\label{sec:v31-audit}

\begin{enumerate}[leftmargin=2.4em]
\item \textbf{A nonzero index wall is not forced into one smooth Fredholm
family.}  Lemma~\ref{lem:v31-index-constant} makes that impossible; V31 replaces
V30's connected lift by a stratified one.
\item \textbf{A line-bundle isomorphism is not called a connection
isomorphism.}  Proposition~\ref{prop:v31-line-not-phase} is the explicit
counterexample.
\item \textbf{V24 is used at its proved strength.}  Norm-resolvent control gives
canonical soft projector and determinant-line bundle isomorphisms, but V31 does
not infer eta/holonomy from them.
\item \textbf{The remaining phase datum is isolated as one relative line.}
Its curvature and loop character are gauge invariant and are exactly the
missing phase-faithfulness test.
\item \textbf{Wall data remain explicit.}  Across index-changing strata the
comparison is a relative determinant line/cobordism, not a fictitious smooth
continuation of the same Fredholm index.
\item \textbf{The next domain-wall/APS route is only a proposed proof strategy.}
V31 does not import it without matching the explicit V2 regulator, source jets
and determinant-line conventions.
\end{enumerate}

\section{V31 status ledger}
\label{sec:v31-status}

\subsection*{Proved in V31}
\begin{enumerate}[leftmargin=2.2em]
\item The numerical Fredholm index is locally constant in every smooth
connected Fredholm family.  Therefore a genuine V27/V28 Wilson wall with
$\Delta k\ne0$ cannot be realized inside one smooth connected geometric
Fredholm family preserving the native chamber indices.
\item The correct global replacement of V30's connected lift is a stratified
geometric lift: one fixed-index geometric family per chamber plus exact relative
determinant-line data across the walls.
\item A bare determinant-line bundle isomorphism does not identify determinant
connections or loop holonomies.  V30's phase conclusions require a
connection-preserving (phase-faithful) chamber comparison.
\item V24 norm-resolvent control gives canonical direct-rotation isomorphisms of
native and geometric kernel/cokernel bundles on every smooth V2-consistent
fixed-index chamber.  Taking determinant exterior powers gives a canonical
unitary line-bundle isomorphism and equality of first Chern classes.
\item The remaining chamber obstruction is the relative determinant line
$\mathcal R_a$ with connection.  A phase-faithful bridge exists if and only if
its curvature vanishes and every relative loop holonomy is one.
\item Under a stratified phase-faithful lift, the complete native angular
constructible determinant line has one global chamberwise-parallel section and
unit constructible holonomy.  This is the corrected form of V30's global
trivialization theorem.
\item The difference of native and geometric projector Berry curvatures is an
exact transgression along any V24 direct-rotation path.  Exactness alone does
not imply vanishing or trivial periods.
\end{enumerate}

\subsection*{Inherited}
\begin{enumerate}[leftmargin=2.2em]
\item V1--V24: the coefficientwise ESM continuum, regulator universality,
controlled fixed-index classification and the smooth V2-consistent soft
norm-resolvent theorem.
\item V25--V29: real Pfaffian orientation, accidental and genuine Wilson
crossings, finite-rank wall lines, clean graded coherence and nonsplit angular
Berry/constructible determinant data.
\item V30: the complete-packet angular cancellation mechanism once the native
constructible line is phase-faithfully identified with the inherited faithful
geometric determinant family.
\item The Grand Tensor smooth faithful-spin theorem: on each declared connected
ordinary-spin $G_6$ geometric family component the complete determinant line has
one unit parallel section and exact cofinal cutoff transport.
\end{enumerate}

\subsection*{Conditional}
\begin{enumerate}[leftmargin=2.2em]
\item V31 proves the topological native-to-geometric chamber bridge only on the
V24 smooth V2-consistent fixed-index class where both sides use the same
continuum soft operator.
\item The corrected global angular cancellation theorem remains conditional on
the relative phase condition \eqref{eq:v31-relative-phase-goal} in every
chamber and on exact geometric relative-line matching across index-changing
walls.
\item Different geometric index sectors require declared relative
cobordism/domain data; they cannot be merged into one smooth Fredholm family by
notation.
\item The independent real/Pfaffian orientation conditions of V25/V28/V30 remain
in force.
\item The positive IR/reference and fixed finite coefficient/source-order
quantifiers remain unchanged.
\end{enumerate}

\subsection*{Open beyond V31}
\begin{enumerate}[leftmargin=2.2em]
\item Prove the native-to-geometric \emph{relative phase theorem}: construct a
V2-compatible local five-dimensional/domain-wall or APS interpolation and show
that the relative line $\mathcal R_a$ has zero curvature and trivial loop
holonomy for the complete faithful Standard-Model packet.
\item Match each V27/V28/V29 native wall complex to the corresponding geometric
relative determinant/cobordism line across changes of Fredholm sector, including
all graded higher-intersection signs.
\item If the relative phase line is nontrivial for some admissible microscopic
completion, compute its flat character/curvature and its image in the V16
microscopic-loop kernel explicitly.
\item Retain the separate real Pfaffian orientation/zero-crossing analysis and
the stronger positive-IR/nonperturbative problems.
\end{enumerate}

\subsection*{Exact next load-bearing theorem}
\begin{center}
\fbox{\parbox{0.92\textwidth}{
\textbf{Native-to-geometric relative eta/domain-wall phase theorem.}
On one V24 smooth V2-consistent fixed-index chamber, construct an explicit
local interpolation from the native Wilson/overlap/Yukawa determinant family to
the common geometric Dirac--Yukawa boundary family, with the full retained
source/BV packet.  Prove that the induced relative determinant line has
\(\mathcal F^{\rm rel}=0\) and trivial paired loop holonomy, equivalently a unit
parallel section.  Across a V27/V28 index wall, identify the native exterior-
power wall line with the geometric relative determinant/cobordism line and prove
graded coherence at intersections.  This theorem, rather than a bare line
isomorphism, is the precise bridge needed to invoke the inherited faithful-spin
eta trivialization on the native resolved regulator family.}}
\end{center}

\section*{Append-only continuation marker: V31}
\addcontentsline{toc}{section}{Append-only continuation marker: V31}
\textbf{End of V31.}  The complete V1--V30 body above is retained verbatim.
V31 goes upstream on the missing V30 geometric bridge.  It proves that genuine
index-changing Wilson walls cannot live inside one smooth connected Fredholm
family, so the faithful lift must be stratified by fixed-index chambers with
relative determinant data across the walls.  It also corrects the phase
criterion: V24 already gives a canonical topological determinant-line bridge by
norm-close soft Riesz projectors, but a line isomorphism alone does not identify
Berry/Quillen connections or holonomy.  The remaining chamber datum is the
native-to-geometric relative phase line; phase-faithfulness is exactly the
vanishing of its curvature and loop character.  The next gate is therefore a
native domain-wall/APS relative eta theorem which trivializes that relative
phase line and matches the wall complexes geometrically.



\clearpage
\section{V32 continuation: coefficientwise phase bridge and exact-eta firewall}
\label{sec:v32-continuation}

V31 isolated the exact native-to-geometric phase problem: after V24 has
identified the native and geometric determinant \emph{bundles}, one still has
to compare their natural determinant connections and loop holonomies.  The
strongest formulation requested an exact five-dimensional domain-wall/APS
identity at finite cutoff.  Going upstream shows that this is stronger than the
coefficientwise continuum theorem which governs the entire V1--V24 programme.
At every prescribed finite source/BV order, the already-proved native hard-shell
and forest calculus controls the determinant connection through a finite set of
trace/Kato words.  Their native-to-geometric difference has exactly the same
structure as the V8 line transgression:
\begin{equation}
 \boxed{
 \text{finite local source/reference jet}
 \; + \;
 \text{quasilocal remainder tending to zero}.}
 \label{eq:v32-local-plus-rem}
\end{equation}
The local jet is an admissible V7 line/reference coordinate and is fixed once by
the same chronological physical normalization as the rest of the common action.
After that one finite rephasing, every retained determinant-connection and
curvature coefficient converges to the geometric coefficient.  This closes the
phase bridge at the exact quantifier required by the source-complete
coefficientwise continuum.

V32 deliberately does \emph{not} promote this statement to an exact finite-cutoff
eta theorem.  A nontrivial flat loop character can be invisible to every finite
local source jet.  Determining that character still requires the stronger
domain-wall/APS comparison proposed in V31.  The distinction is load bearing:
local coefficient matching and global exact holonomy are not equivalent.

Fix throughout a finite packet
\begin{equation}
 \mathfrak q=(L,\Delta,r,N,\mathcal P_{\rm ext})
 \label{eq:v32-q}
\end{equation}
and one smooth V24 V2-consistent fixed-index chamber $C$.  Choose one protected
source chart $U\subset C$ with centre $\zeta_0$.  Let
$\mathcal L_a^{\rm nat}$ and $\mathcal L^{\rm geom}$ be the native and geometric
determinant lines from V31, and use the V24 direct rotations to identify their
kernel and cokernel bundles topologically.

\subsection{Finite determinant-connection packets}
\label{sec:v32-connection-packet}

In a local Kato frame, the determinant connection through source order $r$ is
not an infinite-dimensional object.  It is a finite family of coefficients.

\begin{definition}[Retained determinant-connection jet]
\label{def:v32-connection-jet}
Let $\mathcal I_r$ be the finite set of graded source/BV multi-indices retained
by $\mathfrak q$.  In a local unit frame $s$ of a determinant line write
\begin{equation}
 \nabla s=\mathcal A\,s.
 \label{eq:v32-local-connection}
\end{equation}
The \emph{retained connection jet} at $\zeta_0$ is
\begin{equation}
 \mathbf J_r\mathcal A
 :=\{\partial_I\mathcal A(\zeta_0):I\in\mathcal I_r\},
 \label{eq:v32-connection-jet}
\end{equation}
with the usual Koszul convention for odd writers.  Curvature jets are obtained
by graded exterior differentiation of the same finite packet.
\end{definition}

For an invertible finite block with propagator $G=D^{-1}$, repeated
differentiation of Jacobi's identity expresses every coefficient in
\eqref{eq:v32-connection-jet} as a finite sum of trace words of the form
\begin{equation}
 \Tr\bigl(
 G D_{I_1}G D_{I_2}\cdots G D_{I_m}
 \bigr),
 \label{eq:v32-trace-word}
\end{equation}
together with the Kato-frame terms from kernel/cokernel transport.  In a
fixed-index zero-mode chart, $G$ is replaced by the V11 reduced inverse
$D^\dagger$ and the zero-mode Kato projectors are retained explicitly.  These
are precisely the line/source words already controlled in V1--V8 and V11.

\begin{lemma}[Finite word reduction]
\label{lem:v32-finite-word}
At fixed $\mathfrak q$, every coefficient of the native determinant connection,
its curvature, and all retained determinant-line source descendants is a finite
linear combination of:
\begin{enumerate}[label=\textup{(W\arabic*)},leftmargin=3em]
\item reduced trace words \eqref{eq:v32-trace-word};
\item source derivatives of the kernel/cokernel Riesz projectors and Kato
frames;
\item finite local reference/contact coefficients already present in the V5
contraction-complete state.
\end{enumerate}
The same statement holds for the geometric reference family on the common soft
panel.
\end{lemma}

\begin{proof}
Use the differentiated determinant identity of
Lemma~\ref{lem:v8-line-transgression} and the V11 zero-mode factorization.
Repeated source differentiation gives a finite Bell/Leibniz sum because the
source order is fixed.  Every derivative of an inverse is a finite resolvent
word, while every derivative of a zero-mode frame is a finite Riesz/Kato word.
The V5 packet was defined precisely so that any later BV/source contraction
which becomes local is already represented by a retained local coordinate.
The geometric family has the same finite source list, so the same formal
reduction applies there.
\end{proof}

\subsection{Native-to-geometric two-zone reduction of the line words}
\label{sec:v32-two-zone}

The V24 frequency projectors are proof devices and are frozen at $\zeta_0$.
Insert $P_a+Q_a=I$ between every factor of each trace word, exactly as in
Lemma~\ref{lem:v24-frozen-decomposition}.  This separates each retained line
coefficient into soft words, open low/high transfers, closed hard excursions,
and purely ultraviolet local pieces.

\begin{theorem}[Coefficientwise native--geometric line decomposition]
\label{thm:v32-line-decomposition}
On the smooth V24 V2-consistent class, for every retained source/BV index
$I\in\mathcal I_r$ there is a coefficient
\begin{equation}
 \ell^{\rm loc}_{a,I}
 \label{eq:v32-local-coeff}
\end{equation}
belonging to the finite V7 scalar line/reference/contact bank such that
\begin{equation}
 \boxed{
 \partial_I\bigl(
 \mathcal A_a^{\rm nat}-\mathcal A^{\rm geom}
 \bigr)(\zeta_0)
 =\ell^{\rm loc}_{a,I}+\mathcal E_{a,I},
 \qquad
 |\mathcal E_{a,I}|\le C_I\,\omega_a,}
 \label{eq:v32-line-decomp}
\end{equation}
where one may take
\begin{equation}
 \omega_a=O(a^{1/10})
 \label{eq:v32-omega}
\end{equation}
on the V24 scale choice.  The same decomposition holds after every retained
BV/source contraction and for the determinant-line curvature coefficients.

The local coefficients $\ell^{\rm loc}_{a,I}$ are common-action coordinates:
they are chosen once for the complete Standard-Model packet and not separately
for species, graphs, Ward rows, or wall charts.
\end{theorem}

\begin{proof}
Apply Lemma~\ref{lem:v32-finite-word} and insert the frozen V24 decomposition in
every native trace/Riesz word.

For a word entirely in the soft block, V24 common-soft consistency and the
source-complete Feshbach theorem identify every native reduced resolvent/source
vertex with the corresponding geometric soft word up to
$O(a^{1/10})$.  The kernel/cokernel Kato words satisfy the same estimate because
V24 gives norm convergence of the Riesz projectors with all retained source
derivatives.

Every word containing an open low/high transfer carries one of the factors
bounded in \eqref{eq:v24-main-transfer}; every closed hard excursion carries the
self-energy factor \eqref{eq:v24-main-self}.  Proposition
\ref{prop:v24-source-complete} shows that fixed source insertions do not remove
that suppression.  These pieces therefore contribute $O(a^{1/10})$.

The remaining terms are purely ultraviolet words with all external/source
momenta Taylor expanded before the first omitted soft order.  By the V2
full-zone calculus and Theorem~\ref{thm:v8-line-local-plus-rem}, their
nondecaying part is a finite local line/reference/source polynomial; assign its
coefficient to $\ell^{\rm loc}_{a,I}$.  The Taylor complement has the usual
one-excess soft-to-hard factor and is absorbed into the same vanishing error.
Because the complete packet and forest ownership are fixed before the
comparison, these local coefficients belong to the one common V7 normalization
bank.  The V5 contraction-complete packet gives the identical statement after
all retained BV/source descendants.
\end{proof}

\subsection{One finite phase normalization removes the local packet}
\label{sec:v32-phase-normalization}

The coefficients in \eqref{eq:v32-local-coeff} are not independent.  They are
the Taylor/source coefficients of an actual local determinant transgression and
therefore satisfy the graded integrability relations of one scalar local
functional.  This prevents a graphwise or derivative-by-derivative phase fit.

\begin{lemma}[Integrable local phase jet]
\label{lem:v32-integrable-jet}
There exists one finite graded source polynomial
\begin{equation}
 \chi_{a,\mathfrak q}^{\rm loc}(J)
 \label{eq:v32-chi}
\end{equation}
in the V7 scalar line/reference bank, unique after its value at $J=0$ is fixed,
such that through retained source degree $r$
\begin{equation}
 \partial_I d\chi_{a,\mathfrak q}^{\rm loc}(\zeta_0)
 =\ell^{\rm loc}_{a,I}
 \qquad(I\in\mathcal I_r).
 \label{eq:v32-chi-jet}
\end{equation}
Here $d$ denotes the exterior differential on the declared background/source
parameter panel and the equality is coefficientwise in the finite graded
source algebra.
\end{lemma}

\begin{proof}
Take the local Taylor polynomial, through the retained source order, of the
actual native-to-geometric logarithmic determinant transgression after all hard
local Taylor subtractions have been made.  Its exterior derivative is exactly
the local one-form packet in Theorem~\ref{thm:v32-line-decomposition}.  Equality
of mixed graded derivatives follows from the fact that this packet arose by
differentiating one scalar determinant/line functional, not by fitting the
coefficients independently.  Fixing the source-free constant gives uniqueness.
\end{proof}

Rephase the native local line frame by
\begin{equation}
 s_a^{\rm ren}
 :=e^{-\chi_{a,\mathfrak q}^{\rm loc}}s_a^{\rm nat}.
 \label{eq:v32-rephase}
\end{equation}
This is exactly a V7 finite line/reference scheme transformation.  It must be
retained when source dependent; it is not discarded as a field-independent
constant.

\begin{theorem}[Coefficientwise phase-faithful chamber bridge]
\label{thm:v32-coefficientwise-phase}
In the rephased frame \eqref{eq:v32-rephase}, for every retained source/BV
multi-index $I$,
\begin{equation}
 \boxed{
 \partial_I\bigl(
 \mathcal A_a^{\rm ren}-\mathcal A^{\rm geom}
 \bigr)(\zeta_0)
 =O_I(\omega_a),
 \qquad \omega_a\to0.}
 \label{eq:v32-connection-convergence}
\end{equation}
Consequently every retained curvature/source coefficient satisfies
\begin{equation}
 \boxed{
 \partial_I\bigl(
 \mathcal F_a^{\rm ren}-\mathcal F^{\rm geom}
 \bigr)(\zeta_0)
 =O_I(\omega_a).}
 \label{eq:v32-curvature-convergence}
\end{equation}
The limiting native connection jet is therefore the geometric connection jet in
one coherent admissible V7 scheme.  The construction is projective in terminal
loop order.
\end{theorem}

\begin{proof}
Subtract \eqref{eq:v32-chi-jet} from
\eqref{eq:v32-line-decomp}.  This leaves exactly the vanishing remainder in
every retained coefficient, proving \eqref{eq:v32-connection-convergence}.
Exterior differentiation and the finite source algebra give
\eqref{eq:v32-curvature-convergence}.  The phase polynomial is one finite local
coordinate in the chronological V7 scheme family; once its order-$q$ component
is fixed, higher terminal loop orders do not recompute it.  Hence the same
argument used for the V7 scheme maps proves exact perturbative-order
projectivity.
\end{proof}

\begin{corollary}[Faithful Standard-Model local phase cancellation in the continuum]
\label{cor:v32-SM-local-flat}
On a complete faithful Standard-Model chamber for which the inherited geometric
determinant connection is parallel,
\begin{equation}
 \mathcal F^{\rm geom}=0,
 \label{eq:v32-geom-flat}
\end{equation}
Theorem~\ref{thm:v32-coefficientwise-phase} gives, in the matched continuum
scheme,
\begin{equation}
 \boxed{
 \partial_I\mathcal F_{\infty}^{\rm nat}(\zeta_0)=0
 \qquad(I\in\mathcal I_r).}
 \label{eq:v32-native-curvature-zero}
\end{equation}
Thus every local Berry/connection/source coefficient carried by the complete
packet agrees with the faithful geometric value through the prescribed finite
order.
\end{corollary}

\subsection{Wall and higher-intersection phase jets}
\label{sec:v32-wall-jets}

The same distinction between local coefficient matching and exact global phase
appears at an index-changing wall.  V27--V29 already identify the native wall
line topologically.  V24 consistency on the finite crossing cluster identifies
the native crossing subbundles with their common geometric soft counterparts by
a direct rotation whose norm tends to the identity.

\begin{theorem}[Coefficientwise wall-line phase matching]
\label{thm:v32-wall-phase}
Let $S$ be a regular V27/V28/V29 wall on the smooth V24-consistent class.  After
the chamber phase normalization of Theorem~\ref{thm:v32-coefficientwise-phase},
there is one finite local wall/reference polynomial such that the retained
source/BV jets of the native exterior-power wall line and the geometric relative
wall line differ by $O(\omega_a)$.  After the corresponding V7 wall-frame
rephasing,
\begin{equation}
 \boxed{
 \mathbf J_r\mathcal A_{S,a}^{\rm nat}
 -\mathbf J_r\mathcal A_{S}^{\rm geom}
 =O(\omega_a).}
 \label{eq:v32-wall-phase-match}
\end{equation}
At a clean higher intersection the V28 Koszul braiding is identical on the two
sides, so no additional scalar sign is available to obstruct the matched
coefficientwise limit.
\end{theorem}

\begin{proof}
The crossing Riesz projectors and crossing forms are finite-dimensional soft
data.  The V24 source-complete norm-resolvent estimate, applied to the finite
crossing cluster, makes the native and geometric positive/negative crossing
subbundles norm close with all retained source derivatives.  Their V14 direct
rotations therefore identify the exterior-power wall bundles.  The logarithmic
phase of the resulting finite line comparison has a well-defined finite Taylor
jet; assign that jet to the V7 wall/reference coordinate.  The remaining
source coefficients are $O(\omega_a)$.  Exterior powers are graded functorially,
so the V28 braiding sign depends only on the ranks and is exactly the same for
native and geometric wall packets.
\end{proof}

\subsection{What is still not proved: exact finite-cutoff eta transport}
\label{sec:v32-exact-firewall}

The preceding theorems close the native--geometric phase bridge at the same
coefficientwise quantifier as the continuum theorem.  They do not prove the
stronger V31 condition
\begin{equation}
 \mathcal F_a^{\rm rel}=0,
 \qquad
 \operatorname{Hol}_{\rm rel}(\gamma)=1
 \quad\text{at each finite cutoff}. 
 \label{eq:v32-not-exact}
\end{equation}
A flat connection can have nontrivial holonomy while all of its local curvature
coefficients vanish.  More generally, finite Taylor/source jets on a chart do
not determine a character of $\pi_1(C)$.

\begin{definition}[Exact eta residue]
\label{def:v32-eta-residue}
After the coefficientwise local phase normalization, define the \emph{exact eta
residue} at cutoff $a$ to be the gauge-invariant relative loop character
\begin{equation}
 \boxed{
 \chi_{a}^{\eta}(\gamma)
 :=\operatorname{Hol}_{\rm nat}^{\rm ren}(\gamma)
   \operatorname{Hol}_{\rm geom}(\gamma)^{-1},
 \qquad \gamma\in\pi_1(C),}
 \label{eq:v32-eta-character}
\end{equation}
whenever the exact natural connections are defined on the whole chamber.  It is
unchanged by any single-valued V7 local rephasing.  V32 makes no claim that
$\chi_a^{\eta}$ is cutoff independent or trivial.
\end{definition}

\begin{proposition}[The exact eta residue is invisible to the fixed local coefficient packet]
\label{prop:v32-eta-invisible}
Suppose two exact relative connections have the same retained local connection
and curvature jets on every chart of the finite source atlas but differ by a
flat unitary line with character $\chi:\pi_1(C)\to U(1)$.  Then every
coefficientwise source/BV statement proved in V1--V32 is identical for the two
connections, while their global loop holonomies differ by $\chi$.
Consequently no finite local source-jet argument can, by itself, prove
$\chi_a^{\eta}=1$ on a noncontractible chamber.
\end{proposition}

\begin{proof}
A flat line has zero curvature and, in a local parallel frame, zero connection
one-form.  Tensoring by it therefore changes none of the local source/curvature
jets or their BV/contact descendants.  Around a noncontractible loop its
parallel transport multiplies by the character $\chi$.  Hence all local
coefficient packets agree while the global holonomy can differ.
\end{proof}

This proposition is the exact reason an APS/domain-wall theorem remains a
stronger post-continuum question even after the coefficientwise phase bridge is
closed.

\subsection{Coefficientwise angular cancellation after V32}
\label{sec:v32-angular-consequence}

The corrected V30/V31 angular picture now has two levels.

\begin{theorem}[Complete-packet angular cancellation at coefficientwise continuum order]
\label{thm:v32-angular-coefficientwise}
Assume the V24 smooth V2-consistent native class, the V29 resolved angular
stratification, and the complete faithful Standard-Model packet.  For every
fixed coefficient packet $\mathfrak q$, choose the one common V32 phase/reference
normalization on each chamber and the matched wall normalization of
Theorem~\ref{thm:v32-wall-phase}.  Then in the cutoff-independent continuum:
\begin{enumerate}[label=\textup{(A\arabic*)},leftmargin=3em]
\item every retained chamber Berry-curvature/source coefficient equals its
faithful geometric value and hence vanishes in the inherited parallel frame;
\item every retained wall/source coefficient agrees with the geometric relative
wall coefficient;
\item all V28 graded higher-intersection signs agree exactly;
\item the resulting local constructible determinant/source system is independent
of the microscopic V2-consistent completion up to the already-proved V7 finite
scheme groupoid.
\end{enumerate}
No exact statement about the finite-cutoff global loop character
$\chi_a^{\eta}$ is included.
\end{theorem}

\begin{proof}
Apply Corollary~\ref{cor:v32-SM-local-flat} chamberwise,
Theorem~\ref{thm:v32-wall-phase} on every wall, and V28 graded coherence at
higher intersections.  V7 scheme covariance identifies two different local
phase normalizations by a finite continuum scheme map.  The exact-loop firewall
is Definition~\ref{def:v32-eta-residue} and Proposition
\ref{prop:v32-eta-invisible}.
\end{proof}

\subsection{V32 anti-circularity audit}
\label{sec:v32-audit}

\begin{enumerate}[leftmargin=2.4em]
\item \textbf{The exact eta theorem was not inferred from a topological line
isomorphism.}  V31's counterexample remains decisive.
\item \textbf{The native-to-geometric comparison is made only at the finite
coefficient/source quantifier actually consumed by the continuum proof.}
Every retained connection coefficient is reduced to finite trace/Kato words.
\item \textbf{Ultraviolet phase terms are not discarded.}  Their nondecaying
Taylor/source part is retained as one finite V7 line/reference polynomial; only
the quasilocal Taylor complement is estimated away.
\item \textbf{The local phase polynomial is not fitted coefficient by
coefficient.}  It is the finite Taylor jet of one determinant transgression, so
graded integrability is automatic.
\item \textbf{Wall phases and Koszul signs remain explicit.}  V32 matches the
finite wall connection jets and preserves the exact V28 graded braiding.
\item \textbf{Global flat holonomy is not claimed from local source jets.}
The possible exact eta residue \eqref{eq:v32-eta-character} is retained as a
separate stronger datum.
\end{enumerate}

\section{V32 status ledger}
\label{sec:v32-status}

\subsection*{Proved in V32}
\begin{enumerate}[leftmargin=2.2em]
\item At every prescribed finite loop/EFT/source order, the native determinant
connection, curvature and line-source descendants reduce to a finite family of
reduced trace words, Riesz/Kato words and local reference/contact coefficients.
\item On the V24 smooth V2-consistent class, each retained native-to-geometric
connection coefficient decomposes into one finite local V7 line/reference jet
plus a quasilocal remainder $O(a^{1/10})$; the same holds for every retained
BV/source descendant.
\item The complete local coefficient packet is integrable: it is the jet of one
scalar local determinant transgression.  One chronological V7 phase/reference
polynomial removes it simultaneously, rather than by graphwise or derivative-
wise rephasing.
\item After that finite normalization, all retained native determinant-
connection and curvature coefficients converge to their geometric values.  On
the faithful Standard-Model branch the continuum local curvature/source packet
therefore vanishes in the inherited geometric parallel frame.
\item V27--V29 wall lines admit the same coefficientwise native-to-geometric
phase matching; higher-intersection Koszul signs coincide exactly.
\item The complete V29 angular constructible determinant/source system is
therefore microscopic-regulator independent at every fixed coefficient order,
up to the V7 finite scheme groupoid.
\item An exact finite-cutoff global relative holonomy is a stronger datum.  A
flat character can be invisible to every finite local coefficient/source jet,
so the exact eta residue cannot be killed by the coefficientwise argument
alone.
\end{enumerate}

\subsection*{Inherited}
\begin{enumerate}[leftmargin=2.2em]
\item V1--V8: source-complete hard-shell/forest calculus, exact line
transgression, local-plus-quasilocal line decomposition, and finite scheme
covariance.
\item V9--V24: regulator-class phase lifting, controlled zero-mode modules,
soft norm-resolvent comparison and the source-complete covariant Feshbach
package on the V2-consistent smooth native class.
\item V25--V29: real Pfaffian orientation, divisor and Wilson-wall transitions,
finite-rank graded wall coherence and nonsplit angular Berry/constructible line
data.
\item V30--V31: the faithful geometric angular cancellation mechanism, the
stratified fixed-index correction, the topological chamber bridge and the exact
relative phase/holonomy criterion.
\item The Grand Tensor faithful-spin theorem: the complete geometric determinant
line has a unit parallel section and trivial paired loop holonomy on every
declared smooth ordinary-spin $G_6$ component.
\end{enumerate}

\subsection*{Conditional}
\begin{enumerate}[leftmargin=2.2em]
\item V32's quantitative $O(a^{1/10})$ remainder uses the V24 smooth
V2-consistent class and its common soft continuum jet; rough backgrounds and
$O(1)$ soft defects remain excluded.
\item The phase matching is coefficientwise in one finite source/BV packet.
V32 does not construct one analytic finite-amplitude phase function uniform in
source order.
\item Global exact natural-connection holonomy at finite cutoff is not closed.
The eta residue \eqref{eq:v32-eta-character} may in principle survive on a
noncontractible chamber.
\item The independent real/Pfaffian orientation hypotheses of V25/V28 remain
in force.
\item Positive IR/reference and fixed finite perturbative-order hypotheses are
unchanged.
\end{enumerate}

\subsection*{Open beyond V32}
\begin{enumerate}[leftmargin=2.2em]
\item Prove the exact native domain-wall/APS identity which evaluates
$\chi_a^{\eta}$ and shows whether it is identically one on the faithful native
branch, rather than merely invisible to every fixed local coefficient packet.
\item Construct the corresponding exact geometric relative determinant line
across a V27/V28 nonzero-index wall and match its finite-cutoff phase, not only
its coefficientwise wall/source jets.
\item If a nontrivial exact eta residue survives for an admissible completion,
compute its image in the V16 microscopic-loop kernel and determine whether it
has any observable effect beyond the coefficientwise scheme groupoid.
\item Retain the separate real Pfaffian zero-crossing problem.  The stronger\
positive-IR and nonperturbative directions also remain open.
\end{enumerate}

\subsection*{Exact next load-bearing theorem}
\begin{center}
\fbox{\parbox{0.92\textwidth}{
\textbf{Exact native domain-wall / APS eta-character theorem.}
For the V2 Wilson/overlap/Yukawa family on one smooth fixed-index chamber,
construct a genuine five-dimensional native interpolation whose boundary
comparison is the V31 relative determinant line.  Prove an exact finite-cutoff
formula identifying its exponentiated eta/spectral-flow invariant with the
relative loop character $\chi_a^{\eta}$, with all retained source/BV writers and
with the V27/V28 relative wall line at an index-changing stratum.  Then use the
faithful Standard-Model anomaly/bordism input to decide whether this exact
character is one.  This is now a strictly stronger global phase theorem: V32
has already closed every local coefficientwise connection/source consequence
needed by the perturbative continuum programme.}}
\end{center}

\section*{Append-only continuation marker: V32}
\addcontentsline{toc}{section}{Append-only continuation marker: V32}
\textbf{End of V32.}  The complete V1--V31 body above is retained verbatim.
V32 goes upstream on the exact V31 eta target and separates it from the
coefficientwise continuum quantifier.  Using the V8 local-plus-quasilocal line
transgression, the V24 source-complete Feshbach package and V7 chronological
scheme covariance, it proves that every retained native-to-geometric
determinant-connection/source coefficient equals one finite local
line/reference jet plus a vanishing quasilocal remainder.  One common finite
phase normalization removes the local jet, so the cutoff-independent native
connection and curvature coefficients agree with the faithful geometric values
through every prescribed finite source/BV order; the same holds for V27--V29
wall packets and their graded intersection signs.  What remains outside the
coefficientwise theorem is an exact finite-cutoff flat eta/holonomy character
on noncontractible chambers, which requires a genuine native domain-wall/APS
identity rather than another local source estimate.


\clearpage
\section{V33 continuation: exact determinant-line suspension, overlap integrability, and the native five-dimensional bordism gate}
\label{sec:v33-continuation}

V32 deliberately left one strictly stronger datum outside the coefficientwise
continuum theorem: the exact finite-cutoff relative loop character
$\chi_a^\eta$.  Its proposed next step was a five-dimensional
Wilson/domain-wall APS identity.  Going upstream separates that request into
two mathematically different statements.

First, the exact eta interpretation of the determinant-line holonomy is already
available intrinsically at finite cutoff.  Restrict the relative determinant
line to a loop.  It is a unitary line over $S^1$, and the one-dimensional
covariant Dirac operator on that line has an eta invariant whose exponent is
exactly its holonomy.  In a local overlap frame the same holonomy is the ratio
between the measure Wilson line and the Kato/overlap clutching twist.  Thus
$\chi_a^\eta$ is precisely the overlap-measure integrability defect.

Second, identifying this intrinsic suspension eta with the eta invariant of a
\emph{local five-dimensional Wilson/domain-wall operator} is an additional
realization theorem.  Bordism triviality of the faithful Standard-Model branch
can be applied to the native finite-cutoff phase only after this five-dimensional
character has been shown to be local, gluing compatible, and a genuine
faithful-spin bordism character.  The V1--V32 source-local anomaly theorem does
not by itself supply those global finite-cutoff properties.

V33 closes the first statement completely, including the V27--V29 wall
transitions, and reduces the second statement to an exact native lattice
integrability/bordism gate.

\subsection{The determinant-line clutching scalar of a closed native loop}
\label{sec:v33-clutching}

Let $C$ be a smooth fixed-index chamber and let
\begin{equation}
 \gamma:[0,1]\longrightarrow C,
 \qquad
 \gamma(0)=\gamma(1)=b_0
 \label{eq:v33-loop}
\end{equation}
be a smooth closed loop.  At finite cutoff $a$, let
\begin{equation}
 (\mathcal L_a,\nabla_a)
 :=\bigl(\operatorname{Det}(D_{a}^{\rm nat}),
          \nabla_{a}^{\rm nat,ren}\bigr)
 \label{eq:v33-native-line}
\end{equation}
be the V32-renormalized native determinant line with its exact natural
connection.  Pull it back to $[0,1]$ along $\gamma$.

Choose a smooth unit local frame $s(t)$ on the interval.  Since the endpoints
lie over the same base point, there is a unique clutching phase
\begin{equation}
 \boxed{s(1)=T_a(\gamma)\,s(0),
 \qquad |T_a(\gamma)|=1.}
 \label{eq:v33-clutching-phase}
\end{equation}
Write the connection in this frame as
\begin{equation}
 \nabla_a s=i\,\ell_a(t)\,dt\otimes s,
 \qquad \ell_a(t)\in\mathbb R,
 \label{eq:v33-connection-frame}
\end{equation}
and define its interval Wilson factor by
\begin{equation}
 \boxed{
 W_a(\gamma)
 :=\exp\left(-i\int_0^1\ell_a(t)\,dt\right).}
 \label{eq:v33-Wilson-factor}
\end{equation}
The sign convention is chosen so that parallel transport of the coefficient of
$s$ multiplies by $W_a$.

\begin{lemma}[Frame-independent integrability defect]
\label{lem:v33-frame-independent}
The quotient
\begin{equation}
 \boxed{
 \mathfrak I_a(\gamma)
 :=W_a(\gamma)\,T_a(\gamma)^{-1}}
 \label{eq:v33-integrability-defect}
\end{equation}
is independent of the interval frame $s(t)$ and equals the exact holonomy of
$\nabla_a$ around the closed loop $\gamma$.
\end{lemma}

\begin{proof}
If $s'(t)=e^{i\varphi(t)}s(t)$, then
$\ell'_a=\ell_a+\dot\varphi$ and
\[
 W'_a=e^{-i(\varphi(1)-\varphi(0))}W_a.
\]
The endpoint clutching changes by
\[
 T'_a=e^{i(\varphi(1)-\varphi(0))}T_a.
\]
With the convention in \eqref{eq:v33-Wilson-factor}, the coefficient obtained
by parallel transport from $t=0$ to $t=1$ is $W_a$; returning from the endpoint
frame to the original fibre frame multiplies by $T_a^{-1}$.  Hence their
product is the closed-loop holonomy and is frame independent.
\end{proof}

For the overlap chiral bundle one may choose the endpoint clutching from Kato
transport.  Let $P_t$ denote the relevant moving finite-dimensional projector
and let $Q_t$ solve
\begin{equation}
 \dot Q_t=[\dot P_t,P_t]Q_t,
 \qquad Q_0=I.
 \label{eq:v33-Kato-Q}
\end{equation}
Then $P_tQ_t=Q_tP_0$.  On a closed loop $P_1=P_0$, and the familiar finite
Kato twist is
\begin{equation}
 \boxed{
 T_P(\gamma)
 =\det\bigl(I-P_0+P_0Q_1\bigr).}
 \label{eq:v33-overlap-twist}
\end{equation}
For a general fixed-index determinant line the analogous clutching factor is
the product of the cokernel Kato determinant and the inverse kernel Kato
determinant, together with the chosen reduced determinant frame.  Equation
\eqref{eq:v33-clutching-phase} is the coordinate-free version of this same
construction.

\begin{theorem}[Exact overlap integrability identity]
\label{thm:v33-integrability-identity}
A single-valued global unit frame with connection Wilson factor compatible with
the Kato/overlap clutching exists around $\gamma$ if and only if
\begin{equation}
 \boxed{W_a(\gamma)=T_a(\gamma).}
 \label{eq:v33-W-equals-T}
\end{equation}
Equivalently,
\begin{equation}
 \boxed{\mathfrak I_a(\gamma)=1.}
 \label{eq:v33-integrability-one}
\end{equation}
Thus the exact finite-cutoff global measure-integrability condition is exactly
trivial determinant-line holonomy.
\end{theorem}

\begin{proof}
By Lemma~\ref{lem:v33-frame-independent}, the obstruction to returning a
parallel frame to itself is exactly $W_aT_a^{-1}$.  It vanishes precisely when
\eqref{eq:v33-W-equals-T} holds.  Rephasing the interval frame changes $W_a$ and
$T_a$ by the same endpoint factor and cannot change their quotient.
\end{proof}

This theorem is the finite determinant-line form of the overlap reconstruction
integrability condition.  It requires no continuum approximation and no
five-dimensional bulk.

\subsection{Exact one-dimensional eta representative of a unitary line holonomy}
\label{sec:v33-eta-suspension}

We next construct an eta invariant which represents
\eqref{eq:v33-integrability-defect} exactly.

\begin{definition}[Twisted circle Dirac operator]
\label{def:v33-twisted-circle}
For $u\in U(1)$ define
\begin{equation}
 \mathcal D_u:=-i\frac{d}{dt}
 \label{eq:v33-Du}
\end{equation}
on $L^2([0,1])$ with domain
\begin{equation}
 \operatorname{Dom}(\mathcal D_u)
 =\{f\in H^1([0,1]):f(1)=u f(0)\}.
 \label{eq:v33-Du-domain}
\end{equation}
Let
\begin{equation}
 \bar\eta(\mathcal D_u)
 :=\frac{\eta(\mathcal D_u)+\dim\ker\mathcal D_u}{2}
 \label{eq:v33-reduced-eta}
\end{equation}
be the reduced eta invariant and put
\begin{equation}
 \xi(u):=\bar\eta(\mathcal D_u)-\frac12
 \quad\in\mathbb R/\mathbb Z.
 \label{eq:v33-xi}
\end{equation}
\end{definition}

\begin{lemma}[Exact eta--holonomy identity]
\label{lem:v33-eta-holonomy}
For every $u\in U(1)$,
\begin{equation}
 \boxed{u=\exp\bigl(-2\pi i\,\xi(u)\bigr).}
 \label{eq:v33-eta-u}
\end{equation}
\end{lemma}

\begin{proof}
Write $u=e^{2\pi i\alpha}$ with $0\le\alpha<1$.  The spectrum of
$\mathcal D_u$ is $2\pi(n+\alpha)$, $n\in\mathbb Z$.  For
$0<\alpha<1$, Hurwitz-zeta continuation gives
\[
 \eta(\mathcal D_u;0)=1-2\alpha,
 \qquad \ker\mathcal D_u=0.
\]
Thus $\xi(u)=-\alpha$ modulo integers.  For $\alpha=0$, the spectrum is
symmetric with a one-dimensional kernel, so
$\bar\eta=1/2$ and again $\xi=0$.  Equation
\eqref{eq:v33-eta-u} follows in both cases.
\end{proof}

\begin{definition}[Intrinsic native eta character]
\label{def:v33-native-eta-character}
For a native loop $\gamma$ define
\begin{equation}
 \boxed{
 \Xi_a^{\rm nat}(\gamma)
 :=\xi\bigl(\mathfrak I_a(\gamma)\bigr)
 \in\mathbb R/\mathbb Z.}
 \label{eq:v33-native-Xi}
\end{equation}
Equivalently, it is the reduced eta invariant of the unitary line obtained by
clutching the pulled-back determinant line around the loop, normalized as in
\eqref{eq:v33-xi}.
\end{definition}

\begin{theorem}[Exact finite-cutoff eta representation of determinant holonomy]
\label{thm:v33-exact-native-eta}
For every smooth closed loop in a fixed-index native chamber,
\begin{equation}
 \boxed{
 \operatorname{Hol}_{\rm nat}^{\rm ren}(\gamma)
 =\mathfrak I_a(\gamma)
 =\exp\bigl(-2\pi i\,\Xi_a^{\rm nat}(\gamma)\bigr).}
 \label{eq:v33-native-hol-eta}
\end{equation}
This is an exact statement at fixed spatial cutoff $a$; no continuum limit is
used.
\end{theorem}

\begin{proof}
The first equality is Lemma~\ref{lem:v33-frame-independent}; the second is
Lemma~\ref{lem:v33-eta-holonomy} with
$u=\mathfrak I_a(\gamma)$.
\end{proof}

The construction is intrinsic to the determinant line.  It should not yet be
confused with the eta invariant of a five-dimensional local Wilson/domain-wall
operator.  V33 will isolate that extra realization below.

\subsection{Exact eta formula for the V32 relative residue}
\label{sec:v33-relative-eta}

Let $(\mathcal L_a^{\rm nat},\nabla_a^{\rm ren})$ and
$(\mathcal L^{\rm geom},\nabla^{\rm geom})$ be the native and geometric lines
on one V31 fixed-index chamber.  Their relative line is
\begin{equation}
 \mathcal R_a
 =\mathcal L_a^{\rm nat}\otimes(\mathcal L^{\rm geom})^*.
 \label{eq:v33-relative-line}
\end{equation}
Restricting to $\gamma$ gives a unitary line over $S^1$ with holonomy exactly
$\chi_a^\eta(\gamma)$ from V32.

\begin{definition}[Relative suspension eta]
\label{def:v33-relative-suspension}
Define
\begin{equation}
 \boxed{
 \Xi_a^{\rm rel}(\gamma)
 :=\xi\bigl(\chi_a^\eta(\gamma)\bigr)
 \in\mathbb R/\mathbb Z.}
 \label{eq:v33-relative-Xi}
\end{equation}
Equivalently it is the normalized reduced eta invariant of the twisted circle
Dirac operator on $\gamma^*\mathcal R_a$.
\end{definition}

\begin{corollary}[Exact eta formula for the V32 eta residue]
\label{cor:v33-relative-eta}
For every loop for which the exact native and geometric connections are
defined,
\begin{equation}
 \boxed{
 \chi_a^\eta(\gamma)
 =\exp\bigl(-2\pi i\,\Xi_a^{\rm rel}(\gamma)\bigr).}
 \label{eq:v33-chi-eta}
\end{equation}
On the inherited faithful geometric branch,
$\operatorname{Hol}_{\rm geom}(\gamma)=1$, so
\begin{equation}
 \Xi_a^{\rm rel}(\gamma)=\Xi_a^{\rm nat}(\gamma)
 \quad\text{in }\mathbb R/\mathbb Z.
 \label{eq:v33-rel-native-faithful}
\end{equation}
\end{corollary}

\begin{proof}
Apply Lemma~\ref{lem:v33-eta-holonomy} to the relative line.  The faithful
geometric equality is the inherited unit parallel-section theorem.
\end{proof}

Thus the exact global datum left open by V32 already has an exact eta
interpretation.  The unresolved question is not the existence of an eta
character but whether this intrinsic character is computed by a local
five-dimensional native domain-wall operator and therefore inherits the
faithful-spin bordism trivialization.

\subsection{Constructible loops crossing Wilson walls}
\label{sec:v33-wall-loop}

The preceding discussion assumed a fixed-index chamber.  We now include the
V27--V29 wall data.

Let a closed constructible gallery $\Gamma$ consist of chamber paths
$\gamma_j$ separated by regular Wilson-wall crossings $S_j$.  On each chamber
use the exact connection Wilson factor $W_j$ and determinant-line clutching
factor $T_j$.  At a wall insert the exact V28/V29 graded wall isomorphism
\begin{equation}
 \Theta_{S_j}:\mathcal L_{j,-}
 \longrightarrow
 \mathcal L_{j,+}\otimes\mathbb T_{S_j}.
 \label{eq:v33-wall-theta}
\end{equation}
At higher clean intersections use the canonical graded braiding already proved
in V28.

\begin{definition}[Constructible integrability defect]
\label{def:v33-constructible-defect}
The scalar $\mathfrak I_a^{\rm con}(\Gamma)\in U(1)$ is the closed product of
all chamber parallel transports and all wall identifications, expressed back in
the starting graded determinant line.  Equivalently it is the holonomy of the
V29 constructible determinant line along $\Gamma$.
\end{definition}

\begin{theorem}[Exact wall-glued eta character]
\label{thm:v33-wall-glued-eta}
There is a canonically defined class
\begin{equation}
 \Xi_a^{\rm con}(\Gamma)
 :=\xi\bigl(\mathfrak I_a^{\rm con}(\Gamma)\bigr)
 \in\mathbb R/\mathbb Z
 \label{eq:v33-wall-Xi}
\end{equation}
such that
\begin{equation}
 \boxed{
 \mathfrak I_a^{\rm con}(\Gamma)
 =\exp\bigl(-2\pi i\,\Xi_a^{\rm con}(\Gamma)\bigr).}
 \label{eq:v33-wall-eta}
\end{equation}
The value is independent of a refinement of the gallery.  At a clean
codimension-two intersection the two crossing orders agree exactly in the
graded determinant category; scalarizing in the opposite order produces the
V28 Koszul sign and no additional eta defect.
\end{theorem}

\begin{proof}
V27--V29 prove that the wall maps and higher graded braidings form a coherent
constructible line.  Hence their closed product is one unitary scalar in the
starting line and is unchanged by inserting or deleting artificial subdivision
points.  Apply Lemma~\ref{lem:v33-eta-holonomy} to that scalar.  The V28 Koszul
sign is already part of the canonical graded braiding and therefore is not a
new failure of coherence.
\end{proof}

This theorem supplies the exact eta character of the full stratified native
line without ever differentiating the Wilson sign through a wall.

\subsection{What a genuine five-dimensional native realization must prove}
\label{sec:v33-five-dimensional-gate}

The intrinsic eta character above is a one-dimensional suspension of the
determinant line.  A genuine lattice APS/domain-wall theorem should instead
construct a local five-dimensional fermion operator whose phase computes the
same character.

\begin{definition}[Native five-dimensional eta realization]
\label{def:v33-5d-realization}
A \emph{native five-dimensional eta realization} for a loop or constructible
gallery $\Gamma$ consists of:
\begin{enumerate}[label=\textup{(F\arabic*)},leftmargin=3em]
\item a local gauge/spin-covariant five-dimensional Wilson/domain-wall operator
$\mathbb D_{5,a,\Gamma}$ on the four-dimensional Renewal lattice crossed with
an auxiliary circle/cylinder;
\item a Pauli--Villars/reference operator $\mathbb D_{5,a,\Gamma}^{\rm PV}$
with the same ultraviolet principal data;
\item a unit-modulus exponentiated eta/domain-wall phase
\begin{equation}
 \tau_{5,a}(\Gamma)
 :=\exp\bigl(-2\pi i\,\bar\eta_{\rm PV}
 (\mathbb D_{5,a,\Gamma})\bigr),
 \label{eq:v33-tau5}
\end{equation}
with the convention absorbed into $\bar\eta_{\rm PV}$;
\item exact gluing under concatenation and exact agreement with the
V27--V29 relative wall line at every index-changing stratum;
\item all fixed source/BV derivatives demanded by the coefficient packet.
\end{enumerate}
It \emph{realizes the native determinant phase} when
\begin{equation}
 \boxed{
 \tau_{5,a}(\Gamma)=\mathfrak I_a^{\rm con}(\Gamma)}
 \label{eq:v33-5d-realizes}
\end{equation}
for every declared gallery.
\end{definition}

\begin{theorem}[Bordism reduction of the exact eta problem]
\label{thm:v33-bordism-reduction}
Assume the complete faithful Standard-Model packet admits a native
five-dimensional eta realization in the sense of
Definition~\ref{def:v33-5d-realization}.  Assume further that, after inclusion
of the exact finite local anomaly/Chern--Simons counterterm, the phase
$\tau_{5,a}$ depends only on the ordinary-spin $G_6$ bordism class of the closed
five-dimensional mapping torus.  Then
\begin{equation}
 \boxed{
 \mathfrak I_{a,\rm SM}^{\rm con}(\Gamma)=1,
 \qquad
 \chi_{a,\rm SM}^{\eta}(\gamma)=1}
 \label{eq:v33-faithful-trivial}
\end{equation}
for every declared closed gallery and fixed-index loop.
\end{theorem}

\begin{proof}
By the realization hypothesis,
$\mathfrak I_{a,\rm SM}^{\rm con}=\tau_{5,a}^{\rm SM}$.  By the second
hypothesis this is a character of
$\Omega_5^{\mathrm{Spin}}(BG_6)$.  The inherited faithful-spin theorem proves
that this bordism group is zero for the declared ordinary-spin quotient.
Therefore the character is identically one.  The fixed-index statement is the
special case with no wall crossings.
\end{proof}

This theorem is exact but conditional for a precise reason: V1--V32 prove local
anomaly cancellation and coefficientwise phase matching, whereas a
finite-cutoff bordism character requires exact global five-dimensional gluing
and integrability.

\subsection{The native lattice integrability defect}
\label{sec:v33-lattice-integrability}

The five-dimensional gate can be stated without any continuum terminology.
Let $T_a(\gamma)$ be the Kato/overlap clutching twist and let
$W_a^{\rm ren}(\gamma)$ be the Wilson factor of the exact V32-renormalized
measure connection.  Define
\begin{equation}
 \boxed{
 \Delta_a^{\rm int}(\gamma)
 :=W_a^{\rm ren}(\gamma)T_a(\gamma)^{-1}.}
 \label{eq:v33-lattice-integrability-defect}
\end{equation}

\begin{proposition}[Eta residue equals the exact lattice integrability defect]
\label{prop:v33-eta-integrability}
On the faithful geometric branch, after choosing the inherited geometric
parallel frame,
\begin{equation}
 \boxed{
 \chi_a^\eta(\gamma)
 =\Delta_a^{\rm int}(\gamma)
 =\exp\bigl(-2\pi i\Xi_a^{\rm rel}(\gamma)\bigr).}
 \label{eq:v33-chi-integrability}
\end{equation}
Thus the exact finite-cutoff eta problem of V32 is equivalent to the global
integrability problem for the native overlap measure.
\end{proposition}

\begin{proof}
In the geometric parallel frame the geometric holonomy is one.  The native
closed-loop holonomy is Lemma~\ref{lem:v33-frame-independent}, namely the
Wilson factor divided by the clutching factor.  Corollary
\ref{cor:v33-relative-eta} gives the eta expression.
\end{proof}

This proposition is the exact point of contact with the standard overlap
measure reconstruction problem: a local measure current can define a global
fermion measure only when its Wilson line equals the geometric twist on every
closed loop.

\subsection{Why V3/V32 do not yet imply exact integrability}
\label{sec:v33-not-yet}

V3 proved the complete Standard-Model local anomaly projection is zero and
constructed one common local restoring primitive.  V32 proved that the native
and geometric determinant connections agree through every prescribed finite
local source/BV order after one finite phase normalization.  Neither theorem
computes the global quantity \eqref{eq:v33-lattice-integrability-defect} on a
noncontractible loop.

\begin{proposition}[Local curvature control does not determine the exact lattice twist]
\label{prop:v33-local-not-global}
There exist unitary line connections with identical local curvature and all
finite local connection jets but different values of
$\Delta^{\rm int}$ on a noncontractible loop.  Hence local anomaly cancellation,
local measure-current reconstruction and coefficientwise phase matching do not
by themselves prove
\begin{equation}
 W_a^{\rm ren}(\gamma)=T_a(\gamma)
 \label{eq:v33-not-implied}
\end{equation}
at finite cutoff.
\end{proposition}

\begin{proof}
Tensor one of the lines by a nontrivial flat unitary line.  This changes no
local curvature or local parallel-frame jet, but multiplies the closed-loop
holonomy, and therefore $W/T$, by the flat character.  This is the same
finite-cutoff obstruction isolated abstractly in V32.
\end{proof}

Thus an exact five-dimensional lattice integrability theorem is not optional if
one wants the stronger finite-cutoff character, even though it is unnecessary
for the already-closed coefficientwise continuum.

\subsection{External methodology check and exact scope}
\label{sec:v33-literature-check}

The structure isolated above agrees with two established external frameworks.
The Dai--Freed holonomy theorem identifies determinant-line holonomy of a
geometric Dirac family with an exponentiated eta invariant of its mapping
torus, including gluing through relative determinant lines.  In overlap chiral
lattice gauge theory, Kato transport of the Ginsparg--Wilson chiral projector
defines the twist
$\det(I-P_0+P_0Q_1)$, and global measure reconstruction requires equality of
that twist with the Wilson line of the local measure current.  Five-dimensional
domain-wall formulations recast this integrability condition in terms of
higher-dimensional lattice phases.

V33 uses none of those external results to infer a conclusion whose hypotheses
have not been proved for the Renewal regulator.  Instead, it proves the exact
finite determinant-line identities directly and uses the literature only to
identify the appropriate next native theorem.

\subsection{V33 anti-circularity audit}
\label{sec:v33-audit}

\begin{enumerate}[leftmargin=2.4em]
\item \textbf{An eta character is actually constructed.}  The V32 loop
holonomy is represented exactly by the reduced eta invariant of a twisted
circle Dirac operator; it is not merely renamed ``eta.''
\item \textbf{The one-dimensional suspension is not called the physical
five-dimensional domain-wall operator.}  The latter remains the explicit gate
of Definition~\ref{def:v33-5d-realization}.
\item \textbf{The overlap integrability condition is exact.}  The quotient
$W/T$ is frame independent and equals determinant-line holonomy at finite
cutoff.
\item \textbf{Wall crossings are retained.}  V27--V29 exterior-power wall
lines and V28 Koszul braiding enter the constructible eta character exactly;
no sign derivative through a Wilson zero is used.
\item \textbf{Bordism triviality requires a prior bordism-character theorem.}
The inherited vanishing of
$\Omega_5^{\mathrm{Spin}}(BG_6)$ is not used to manufacture native lattice
bordism invariance.
\item \textbf{Local anomaly cancellation is not promoted to global
integrability.}  V3/V32 control local curvature/source data; V33 leaves the
flat global character visible.
\end{enumerate}

\section{V33 status ledger}
\label{sec:v33-status}

\subsection*{Proved in V33}
\begin{enumerate}[leftmargin=2.2em]
\item For every closed native loop, the exact determinant-line holonomy is the
frame-independent quotient of the connection Wilson factor and the endpoint
Kato/overlap clutching factor.  The equality $W=T$ is exactly the global
measure-integrability condition.
\item Every unitary line holonomy has an exact reduced-eta representation via
the twisted circle operator $-i\,d/dt$; applied to the native determinant line,
this gives an exact finite-cutoff intrinsic eta character.
\item The V32 relative eta residue satisfies
$\chi_a^\eta=\exp(-2\pi i\Xi_a^{\rm rel})$ exactly, and on the faithful
geometric branch it equals the native intrinsic eta character.
\item V27--V29 chamber paths, wall exterior-power maps and higher Koszul
braidings define one exact constructible eta character for a stratified closed
gallery.
\item On the faithful geometric branch the exact eta residue is identical to
the native overlap-measure integrability defect
$W_a^{\rm ren}T_a^{-1}$.
\item If a local five-dimensional native domain-wall realization computes this
same character and is a faithful-spin $G_6$ bordism character, then the
inherited vanishing of $\Omega_5^{\mathrm{Spin}}(BG_6)$ forces the exact
finite-cutoff character to be one.
\end{enumerate}

\subsection*{Inherited}
\begin{enumerate}[leftmargin=2.2em]
\item V1--V8: native hard-shell, determinant/Pfaffian line calculus, exact
Schur cocycles, common BV restoration and finite scheme/reference covariance.
\item V9--V24: regulator-class lifting, controlled zero-mode modules, soft
norm-resolvent comparison and the source-complete V2-consistent Feshbach
package.
\item V25--V29: native real/Pfaffian orientation, divisor and Wilson-wall
transitions, finite-rank crossing forms, graded wall coherence and nonsplit
angular constructible lines.
\item V30--V32: faithful geometric cancellation mechanism, corrected
stratified/phase-faithful bridge, and coefficientwise native-to-geometric phase
matching with an explicitly retained exact eta residue.
\item Grand Tensor: the geometric complete Standard-Model determinant family
has trivial faithful-spin eta character and a unit parallel section on its
declared smooth ordinary-spin $G_6$ components.
\end{enumerate}

\subsection*{Conditional}
\begin{enumerate}[leftmargin=2.2em]
\item The intrinsic eta character of V33 is exact.  Its realization by a local
five-dimensional Wilson/domain-wall operator is not yet proved for the full
Renewal gauge, spin/coframe, Higgs, and Yukawa packet.
\item Exact triviality of the native character requires a proof that the
five-dimensional native phase is a bordism character after the complete finite
local anomaly/Chern--Simons counterterm is included.
\item V32 coefficientwise phase matching remains sufficient for the
perturbative continuum regardless of whether this stronger exact character is
trivial.
\item The real/Pfaffian branches retain the V25--V29 orientation and parity
conditions; V33's complex line suspension does not replace them.
\item Positive IR/reference and fixed finite perturbative/source-order
hypotheses remain unchanged.
\end{enumerate}

\subsection*{Open beyond V33}
\begin{enumerate}[leftmargin=2.2em]
\item Construct the actual five-dimensional Renewal Wilson/domain-wall operator
for a closed native loop and prove that its Pauli--Villars-regulated phase equals
$W_a^{\rm ren}T_a^{-1}$ exactly at finite cutoff.
\item Prove the exact lattice integrability/bordism theorem for the complete
faithful Standard-Model packet: after the finite local anomaly/Chern--Simons
counterterm, show that the native five-dimensional phase depends only on the
ordinary-spin $G_6$ bordism class.
\item Extend that theorem through the V27--V29 index-changing wall strata and
show that its relative determinant line is exactly the already-constructed
exterior-power wall line.
\item If exact native bordism invariance fails for an admissible completion,
compute the residual flat character and place it in the V16 microscopic-loop
kernel.  This would remain invisible to the already-closed coefficientwise
continuum.
\end{enumerate}

\subsection*{Exact next load-bearing theorem}
\begin{center}
\fbox{\parbox{0.92\textwidth}{
\textbf{Exact five-dimensional lattice integrability / bordism theorem.}
For the V2 Wilson/overlap/Yukawa completion and every declared smooth closed
native loop, construct a local five-dimensional Wilson/domain-wall operator
with Pauli--Villars reference and prove that its exponentiated eta/domain-wall
phase is exactly the V33 determinant-line integrability defect
$W_a^{\rm ren}T_a^{-1}$, including all retained source/BV writers and the
V27--V29 relative wall lines.  Then prove that after the complete native local
anomaly/Chern--Simons counterterm this phase is an ordinary-spin $G_6$ bordism
character.  The inherited faithful-spin bordism theorem would then force the
exact finite-cutoff native eta residue to be one.  This is now a strictly
global lattice-integrability theorem; every local coefficientwise phase
consequence required by the perturbative continuum was already closed in V32.}}
\end{center}

\section*{Append-only continuation marker: V33}
\addcontentsline{toc}{section}{Append-only continuation marker: V33}
\textbf{End of V33.}  The complete V1--V32 body above is retained verbatim.
V33 separates the exact global eta problem into an intrinsic finite-cutoff
statement and a genuinely five-dimensional realization problem.  It proves
that the exact determinant-line holonomy is the overlap integrability defect
$W/T$, gives that phase an exact reduced-eta representation by a twisted circle
Dirac operator, and extends the construction through the V27--V29 stratified
wall cocycle.  On the faithful geometric branch this intrinsic eta character
is exactly the V32 eta residue.  Triviality at finite cutoff is reduced to one
precise remaining theorem: a local five-dimensional Renewal domain-wall phase
must compute the same character and, after the native anomaly counterterm, must
define a faithful-spin $G_6$ bordism character.  Only after that theorem may the
inherited bordism vanishing be used to set the native exact eta residue to one.



\clearpage
\section{V34 continuation: finite-wall obstruction, local Cayley domain wall, and exact infinite-wall realization of the overlap eta character}
\label{sec:v34-continuation}

V33 isolated one stronger finite-cutoff question beyond the coefficientwise
continuum theorem: can the exact native integrability defect
$W_a^{\rm ren}T_a^{-1}$ be realized as the phase of a local five-dimensional
Wilson/domain-wall system?  The most literal formulation asked for an exact
identity at fixed finite fifth extent.  Going upstream shows that this is too
strong for a generic open protected family.  A finite wall separation produces
an analytic rational approximation to the sign function, whereas the exact
overlap projector contains the discontinuous spectral sign.  The correct native
statement therefore has two layers:
\begin{equation}
 \boxed{
 \begin{array}{c}
 \text{finite fifth extent: local truncated-overlap/domain-wall phase,}\\[0.2em]
 \text{infinite wall separation: exact overlap/Kato phase.}
 \end{array}}
 \label{eq:v34-two-layers}
\end{equation}
The limit is taken at fixed four-dimensional cutoff $a$ and is uniform on every
compact protected Wilson-gap/source chart.  Thus it is still an exact
finite-$a$ realization of the V33 character, but not by one fixed finite number
of fifth-direction slices.

V34 constructs this realization directly from the V2 Hermitian Wilson kernel.
It uses only a nearest-neighbour fifth-direction recurrence and the native
four-dimensional kernel on each slice.  The phase after eliminating the bulk is
the determinant of a boundary transfer map.  Its reference-normalized polar
phase converges to the Kato clutching twist.  Multiplication by the already
fixed V32 local measure Wilson factor then converges to the V33 integrability
defect.  The entire construction is source differentiable at every prescribed
finite order and extends chamberwise through the V27--V29 constructible wall
packet.

What remains after V34 is no longer the five-dimensional \emph{realization}
of the eta residue.  The remaining theorem is the genuinely global one:
after the complete local lattice Chern--Simons/anomaly counterterm, prove that
the resulting infinite-wall five-dimensional phase is independent of the
interpolation and factors through the faithful-spin $G_6$ bordism class.

\subsection{The Cayley transfer function and the finite-wall no-go theorem}
\label{sec:v34-finite-wall}

Let $H=H^*$ be one finite-cutoff V2 Hermitian Wilson kernel on a protected
chart, with
\begin{equation}
 \operatorname{spec}H\subset[-M,-g]\cup[g,M],
 \qquad 0<g<M<\infty.
 \label{eq:v34-gap-window}
\end{equation}
Choose a fifth-direction spacing $a_5>0$ such that
\begin{equation}
 a_5M<1.
 \label{eq:v34-a5}
\end{equation}
Define the positive Cayley transfer operator
\begin{equation}
 \mathsf T_{a_5}(H)
 :=(I-a_5H)(I+a_5H)^{-1}
 \label{eq:v34-Cayley-T}
\end{equation}
and, for an integer $N\ge1$, the truncated sign
\begin{equation}
 \boxed{
 \varepsilon_{N,a_5}(H)
 :=\frac{I-\mathsf T_{a_5}(H)^N}
         {I+\mathsf T_{a_5}(H)^N}.}
 \label{eq:v34-truncated-sign}
\end{equation}
All quotients are functional-calculus expressions and are well defined by
\eqref{eq:v34-gap-window}--\eqref{eq:v34-a5}.

\begin{lemma}[Scalar transfer formula]
\label{lem:v34-scalar-transfer}
For a scalar $x$ with $|a_5x|<1$,
\begin{equation}
 \varepsilon_{N,a_5}(x)
 =\tanh\!\left(
 N\,\operatorname{arctanh}(a_5x)
 \right).
 \label{eq:v34-tanh}
\end{equation}
In particular $\varepsilon_{N,a_5}$ is real analytic on
$(-a_5^{-1},a_5^{-1})$ and
\begin{equation}
 \varepsilon_{N,a_5}(0)=0.
 \label{eq:v34-epsN-zero}
\end{equation}
\end{lemma}

\begin{proof}
Put $t=(1-a_5x)/(1+a_5x)$.  Then
$t=\exp[-2\operatorname{arctanh}(a_5x)]$ and
$(1-t^N)/(1+t^N)=\tanh[N\operatorname{arctanh}(a_5x)]$.
\end{proof}

\begin{theorem}[No fixed finite fifth extent gives the exact overlap sign on an open family]
\label{thm:v34-finite-N-no-go}
Fix $a_5>0$ and a finite $N$.  There is no open interval
$U\subset(-a_5^{-1},a_5^{-1})$ meeting both signs on which
\begin{equation}
 \varepsilon_{N,a_5}(x)=\operatorname{sgn}x
 \label{eq:v34-no-go-equality}
\end{equation}
for every nonzero $x\in U$.  More strongly, the equality cannot hold on any
nonempty open positive interval and any nonempty open negative interval within
one connected analytic domain.

Consequently, no fixed finite fifth extent can reproduce the V2 overlap sign
exactly on an open protected regulator/source family whose Wilson spectrum is
allowed to vary continuously inside the positive and negative gapped spectral
bands.  Exact overlap realization at fixed four-dimensional cutoff requires
$N\to\infty$ (or a different fifth-direction construction which already
contains the nonanalytic sign by hand).
\end{theorem}

\begin{proof}
By Lemma~\ref{lem:v34-scalar-transfer},
$\varepsilon_{N,a_5}(x)$ is real analytic across $x=0$ and vanishes there.  If
it equalled $+1$ on a nonempty positive open interval, the analytic identity
theorem would force it to equal $+1$ on the connected analytic domain, in
contradiction with \eqref{eq:v34-epsN-zero}.  The negative side is identical.
A matrix identity on an open operator family would imply the corresponding
scalar identity on every continuously varying spectral branch, so the operator
statement follows.
\end{proof}

\begin{firewall}[Finite wall separation is an approximation, not the exact overlap theory]
At a fixed finite fifth extent one may tune a rational approximation to be
extremely accurate on the protected Wilson bands, or even interpolate a finite
set of eigenvalues.  V34 does not call such an approximation the exact overlap
operator on an open native family.  The exact V2 sign is recovered only after
the wall-separation limit used below.
\end{firewall}

\subsection{Exponential recovery of the native overlap sign}
\label{sec:v34-sign-limit}

Let
\begin{equation}
 T_5:=Na_5
 \label{eq:v34-T5}
\end{equation}
be the physical fifth-direction length in the transfer normalization.

\begin{theorem}[Uniform sign convergence at fixed four-dimensional cutoff]
\label{thm:v34-sign-convergence}
Under \eqref{eq:v34-gap-window}--\eqref{eq:v34-a5}, let
\begin{equation}
 q_g:=\frac{1-a_5g}{1+a_5g}\in(0,1).
 \label{eq:v34-qg}
\end{equation}
Then
\begin{equation}
 \boxed{
 \|\varepsilon_{N,a_5}(H)-\operatorname{sgn}H\|
 \le
 \frac{2q_g^N}{1-q_g^N}
 \le
 \frac{2e^{-2gT_5}}{1-e^{-2gT_5}}.}
 \label{eq:v34-sign-rate}
\end{equation}
The bound is uniform in the four-dimensional volume and cutoff whenever the
protected gap $g$ and spectral ceiling $M$ are uniform.
\end{theorem}

\begin{proof}
For $x\ge g$, the transfer eigenvalue satisfies
$0<t(x)\le q_g<1$ and
\[
 1-\varepsilon_{N,a_5}(x)
 =\frac{2t(x)^N}{1+t(x)^N}
 \le2q_g^N.
\]
For $x\le-g$, apply the same estimate to $-x$ and use oddness.  The displayed
slightly weaker symmetric fraction follows uniformly.  Finally
\[
 \log q_g
 =\log(1-a_5g)-\log(1+a_5g)
 \le-2a_5g,
\]
which gives $q_g^N\le e^{-2gT_5}$.  Functional calculus lifts the scalar bound
to $H$.
\end{proof}

\begin{theorem}[Fixed source-jet convergence]
\label{thm:v34-source-sign-convergence}
Let $H(\zeta)$ be a protected source/background family with a common Wilson
gap $g$, spectral ceiling $M$, and all source derivatives through a prescribed
finite order $r$ uniformly bounded in the V1/V13 local algebra.  For every
source multi-index $I$ with $|I|\le r$ there exist constants $C_I,c_I>0$ such
that
\begin{equation}
 \boxed{
 \left\|
 \partial_I\bigl(
 \varepsilon_{N,a_5}(H(\zeta))-\operatorname{sgn}H(\zeta)
 \bigr)
 \right\|_{\beta_I}
 \le C_I e^{-c_IT_5}.}
 \label{eq:v34-source-sign-rate}
\end{equation}
Here $\beta_I>0$ is one fixed exponentially local weight obtained after the
finite V13 functional-calculus exponent loss.  The same estimate holds for the
corresponding overlap projectors and overlap Dirac/Yukawa source words.
\end{theorem}

\begin{proof}
Choose disjoint compact complex neighborhoods of the positive and negative
spectral bands, separated from zero and from the poles at $\pm a_5^{-1}$.  On
these neighborhoods the scalar functions
$f_N(z)=\varepsilon_{N,a_5}(z)$ converge exponentially, together with every
fixed complex derivative, to the locally constant sign function.  The Cauchy
functional-calculus formula writes each source derivative as a finite sum of
resolvent words weighted by derivatives of $f_N$; equivalently one can
differentiate the Riesz/holomorphic representation.  The V13 resolvent-locality
theorem controls each word in a fixed positive weighted norm.  The finite
source order absorbs the polynomial factors generated by differentiation into
a slightly smaller exponential rate $c_I$.  The overlap and Yukawa statements
follow by the fixed algebraic formulas of V2.
\end{proof}

This theorem already gives the standard infinite-wall recovery of the overlap
operator, but V33 requires a statement about the \emph{phase} of a closed
native loop.  We now construct the required local five-dimensional transfer
system.

\subsection{A local five-dimensional Cayley domain-wall operator}
\label{sec:v34-5d-operator}

Let
\begin{equation}
 \gamma:[0,1]\to C,
 \qquad
 s_j=\frac{j}{N},
 \qquad
 H_j:=H_W(\gamma(s_j))
 \label{eq:v34-discrete-loop}
\end{equation}
be a smooth loop inside one protected fixed-index chamber.  Put
\begin{equation}
 A_j:=I+a_5H_j,
 \qquad
 B_j:=I-a_5H_j.
 \label{eq:v34-AjBj}
\end{equation}
On the five-dimensional strip with slices $j=1,\ldots,N$, define the
nearest-neighbour transfer-domain-wall operator
\begin{equation}
 \boxed{
 (\mathbb K_{5,N}\Psi)_j
 :=A_j\Psi_j-B_j\Psi_{j-1}.}
 \label{eq:v34-K5}
\end{equation}
The boundary variables $\Psi_0$ and $\Psi_N$ are retained rather than
integrated out.  Because each $H_j$ is a native gauge/spin-covariant local
Wilson kernel, $\mathbb K_{5,N}$ is local in four dimensions and strictly
nearest-neighbour in the fifth direction.  Its fixed source derivatives have
the same property.

\begin{lemma}[Exact transfer obtained by Schur elimination]
\label{lem:v34-transfer-Schur}
The homogeneous equation $\mathbb K_{5,N}\Psi=0$ obeys
\begin{equation}
 \Psi_j=\mathsf T_j\Psi_{j-1},
 \qquad
 \mathsf T_j:=A_j^{-1}B_j
 =\mathsf T_{a_5}(H_j),
 \label{eq:v34-recursion}
\end{equation}
and hence
\begin{equation}
 \boxed{
 \Psi_N=\mathcal T_{N}\Psi_0,
 \qquad
 \mathcal T_N:=\mathsf T_N\cdots\mathsf T_1.}
 \label{eq:v34-transfer-product}
\end{equation}
Gaussian/Schur elimination of the interior five-dimensional variables produces
this same boundary transfer map.  The eliminated bulk determinant factor is
\begin{equation}
 \prod_{j=1}^N\det A_j>0,
 \label{eq:v34-positive-bulk}
\end{equation}
so it carries no phase.
\end{lemma}

\begin{proof}
Since $a_5\|H_j\|<1$, every $A_j$ is positive and invertible.  The recursion is
immediate.  Iteration gives \eqref{eq:v34-transfer-product}.  Eliminating the
interior block variables in ascending $j$ performs exactly the same recursion;
at each step the Schur pivot is $A_j$, giving the product
\eqref{eq:v34-positive-bulk}.  Positivity follows from
$I+a_5H_j\succ0$.
\end{proof}

Thus the only nontrivial phase of the boundary amplitude lies in the
finite-dimensional transfer between the moving low-energy subspaces.

\subsection{The boundary determinant line and its reference-normalized phase}
\label{sec:v34-boundary-phase}

Let $P(s)$ denote the spectral projector of $H_W(\gamma(s))$ defining the
chosen overlap many-body/zero-mode side of the determinant-line clutching, and
assume its rank is constant on the chamber.  Let
\begin{equation}
 P_j:=P(s_j).
 \label{eq:v34-Pj}
\end{equation}
The transfer product defines the compressed boundary map
\begin{equation}
 M_N(\gamma)
 :=P_N\mathcal T_NP_0:
 \operatorname{Ran}P_0\longrightarrow\operatorname{Ran}P_N.
 \label{eq:v34-MN}
\end{equation}
For a closed loop $P_N=P_0$.  For sufficiently large $T_5$, the adiabatic
theorem below makes $M_N$ invertible on the finite-rank occupied space.

\begin{definition}[Reference-normalized five-dimensional boundary phase]
\label{def:v34-boundary-phase}
Let
\begin{equation}
 \operatorname{pol}(M_N)
 :=M_N(M_N^*M_N)^{-1/2}
 \label{eq:v34-polar-M}
\end{equation}
be the polar unitary.  Define
\begin{equation}
 \boxed{
 \Theta_{5,N}(\gamma)
 :=\det_{\operatorname{Ran}P_0}\operatorname{pol}(M_N(\gamma))
 \in U(1).}
 \label{eq:v34-Theta5}
\end{equation}
Equivalently,
\begin{equation}
 \Theta_{5,N}(\gamma)
 =\frac{\det M_N(\gamma)}{|\det M_N(\gamma)|}.
 \label{eq:v34-phase-det}
\end{equation}
The denominator is the positive reference/PV normalization.  It may be
implemented by pairing the path with its reflected reverse, for which the
boundary transfer is $M_N^*$.
\end{definition}

\begin{lemma}[Locality of the five-dimensional phase data]
\label{lem:v34-phase-locality}
At every finite $N$, the full five-dimensional operator
$\mathbb K_{5,N}$ is native local.  The boundary phase
\eqref{eq:v34-Theta5} depends on the bulk only through Schur elimination of
that local operator.  Every prescribed fixed source derivative of the boundary
map and its polar factor is finite and uniformly controlled for sufficiently
large $T_5$ on a compact protected loop family.
\end{lemma}

\begin{proof}
Locality of $\mathbb K_{5,N}$ follows directly from
\eqref{eq:v34-K5}.  Source differentiation of the transfer product is a finite
Leibniz sum of source derivatives of $A_j^{-1}B_j$; $A_j^{-1}$ is uniformly
local because $A_j\succeq(1-a_5M)I$.  On the finite-rank occupied boundary
space, invertibility of $M_N$ for large $T_5$ follows from the adiabatic theorem
below.  Functional calculus of the positive matrix $M_N^*M_N$ then gives all
fixed derivatives of the polar factor.
\end{proof}

\subsection{Discrete Kato adiabatic theorem for the native transfer product}
\label{sec:v34-adiabatic}

Let $Q(s)$ denote Kato parallel transport for $P(s)$:
\begin{equation}
 \partial_sQ(s)=[\partial_sP(s),P(s)]Q(s),
 \qquad
 Q(0)=I.
 \label{eq:v34-Kato}
\end{equation}
For a closed loop the V33 clutching twist is
\begin{equation}
 T_P(\gamma)
 =\det_{\operatorname{Ran}P_0}
 \bigl(P_0Q(1)P_0\bigr).
 \label{eq:v34-Kato-twist}
\end{equation}

\begin{theorem}[Adiabatic polar limit of the local five-dimensional transfer]
\label{thm:v34-adiabatic-transfer}
Assume the loop $H_W(\gamma(s))$ is $C^{r+2}$ in $s$ and in the retained
source variables, with uniform spectral gap $g$ and ceiling $M$.  Fix
$a_5M<1$ and let $N\to\infty$, so $T_5=Na_5\to\infty$.  Then there are
constants $C_I$ such that for every retained source multi-index $I$,
\begin{equation}
 \boxed{
 \left\|
 \partial_I\left[
 \operatorname{pol}(M_N)-P_NQ(1)P_0
 \right]
 \right\|
 \le \frac{C_I}{T_5}.}
 \label{eq:v34-adiabatic-rate}
\end{equation}
After restricting to the occupied boundary spaces this gives
\begin{equation}
 \boxed{
 \partial_I\left[
 \Theta_{5,N}(\gamma)-T_P(\gamma)
 \right]
 =O_I(T_5^{-1}).}
 \label{eq:v34-phase-Kato-rate}
\end{equation}
The same conclusion holds in the exponentially local/source-differentiable
norm on the finite boundary packet.
\end{theorem}

\begin{proof}
Write
\begin{equation}
 K_{a_5}(H)
 :=a_5^{-1}\operatorname{arctanh}(a_5H),
 \qquad
 \mathsf T_{a_5}(H)=e^{-2a_5K_{a_5}(H)}.
 \label{eq:v34-effective-generator}
\end{equation}
The Hermitian generator $K_{a_5}(H)$ has exactly the same spectral projectors
as $H$ and a uniform gap bounded below by
$a_5^{-1}\operatorname{arctanh}(a_5g)\ge g$.

Pass to the Kato frame $R(s)=Q(s)$, in which the projector is constant:
$R(s)^*P(s)R(s)=P_0$.  Between consecutive slices,
\[
 R(s_j)^*R(s_{j-1})
 =I-N^{-1}R^*\partial_sR(s_j)+O(N^{-2}).
\]
The Kato generator is off diagonal with respect to $P_0\oplus(I-P_0)$, so
its diagonal $P_0$ block vanishes.  In this frame each transfer step is block
diagonal to leading order, with an exponential separation between the two
spectral sectors controlled by $e^{-2a_5g}$.  The off-diagonal leakage at one
step is $O(N^{-1})$.  Propagating such a leakage away from the occupied block
produces a geometric series with denominator
$1-e^{-2a_5g}$, while the slowly varying path contributes the factor
$N^{-1}=(a_5/T_5)$.  Thus the total relative leakage is $O(T_5^{-1})$.

The occupied-to-occupied block equals a positive dynamical factor times
$I+O(T_5^{-1})$ in the Kato frame.  Polar normalization removes that positive
factor, leaving $I+O(T_5^{-1})$.  Returning to the original frame gives
$P_NQ(1)P_0+O(T_5^{-1})$.  Differentiating the discrete recurrence a fixed
number of times gives the same geometric-series estimate with constants
$C_I$ depending on the finite path/source jet.  This proves
\eqref{eq:v34-adiabatic-rate}; the determinant estimate follows because the
occupied rank is finite and fixed.
\end{proof}

\begin{corollary}[Exact overlap clutching as an infinite-wall local 5D limit]
\label{cor:v34-exact-twist-limit}
For every closed protected loop,
\begin{equation}
 \boxed{
 T_P(\gamma)
 =\lim_{T_5\to\infty}\Theta_{5,N}(\gamma),
 \qquad T_5=Na_5,}
 \label{eq:v34-twist-limit}
\end{equation}
with convergence through every prescribed fixed source/BV derivative.
\end{corollary}

\subsection{Exact realization of the V33 integrability defect as a local 5D limit}
\label{sec:v34-integrability-limit}

Let $W_a^{\rm ren}(\gamma)$ be the exact V32-renormalized measure Wilson factor
of V33.  Its local counterterm/source content is already fixed once by the V7
chronological normalization.  Define the finite-wall five-dimensional phase
\begin{equation}
 \boxed{
 \tau_{5,a}^{(N)}(\gamma)
 :=W_a^{\rm ren}(\gamma)\,
   \Theta_{5,N}(\gamma)^{-1}.}
 \label{eq:v34-tauN}
\end{equation}
Every factor in \eqref{eq:v34-tauN} is defined at the same four-dimensional
cutoff $a$; the first is the native local measure phase and the second is the
reference-normalized boundary phase of the local five-dimensional operator
\eqref{eq:v34-K5}.

\begin{theorem}[Exact infinite-wall realization of the native eta character]
\label{thm:v34-exact-5d-limit}
For every smooth closed loop in a protected fixed-index chamber,
\begin{equation}
 \boxed{
 \chi_a^\eta(\gamma)
 =\Delta_a^{\rm int}(\gamma)
 =\lim_{T_5\to\infty}\tau_{5,a}^{(N)}(\gamma).}
 \label{eq:v34-exact-realization-limit}
\end{equation}
The equality is exact at the fixed four-dimensional cutoff $a$.  Moreover, for
every retained source/BV multi-index $I$,
\begin{equation}
 \boxed{
 \partial_I\left[
 \tau_{5,a}^{(N)}(\gamma)-\chi_a^\eta(\gamma)
 \right]
 =O_I(T_5^{-1}).}
 \label{eq:v34-source-phase-limit}
\end{equation}
Thus the V33 eta residue is realized exactly as the infinite-wall-separation
limit of phases generated by genuinely local five-dimensional native operators.
\end{theorem}

\begin{proof}
V33 proves
$\chi_a^\eta=W_a^{\rm ren}T_P^{-1}$.  Corollary
\ref{cor:v34-exact-twist-limit} gives
$\Theta_{5,N}\to T_P$ with all fixed derivatives.  Substitution into
\eqref{eq:v34-tauN} proves both displayed formulas.
\end{proof}

This theorem closes the five-dimensional \emph{realization} problem after
correcting the finite-thickness quantifier.  It does not yet prove that the
resulting closed-loop phase depends only on a five-dimensional bordism class.

\subsection{General Fredholm determinant line rather than one projector}
\label{sec:v34-general-line}

The preceding formulas were written for one occupied projector.  The V11/V33
determinant line is the quotient of cokernel and kernel exterior powers.  Apply
the same transfer construction separately to the two Kato bundles.  Let
$\Theta_{5,N}^{C}$ and $\Theta_{5,N}^{K}$ denote their boundary polar phases
and define
\begin{equation}
 \Theta_{5,N}^{\operatorname{Det}}
 :=\Theta_{5,N}^{C}
   \left(\Theta_{5,N}^{K}\right)^{-1}.
 \label{eq:v34-det-transfer-phase}
\end{equation}
Then
\begin{equation}
 \boxed{
 \Theta_{5,N}^{\operatorname{Det}}
 \longrightarrow
 T_a(\gamma),}
 \label{eq:v34-general-twist}
\end{equation}
where $T_a$ is the exact determinant-line clutching scalar used in V33.  Hence
Theorem~\ref{thm:v34-exact-5d-limit} holds for the complete fixed-index
Fredholm determinant line, not merely a single spectral subspace.

\begin{proof}
Both kernel and cokernel bundles are finite rank on the protected chamber and
have the V11 Riesz/Kato source calculus.  Apply
Theorem~\ref{thm:v34-adiabatic-transfer} to each and take the determinant
quotient dictated by the V11 line convention.
\end{proof}

\subsection{Constructible galleries crossing Wilson walls}
\label{sec:v34-wall-gallery}

A genuine V27--V29 wall closes the Wilson gap, so one cannot run the gapped
adiabatic theorem through that point.  V34 therefore uses the stratified
construction already forced by V27--V31.

Let a closed gallery $\Gamma$ consist of fixed-index chamber pieces
$\gamma_1,\ldots,\gamma_m$ separated by regular Wilson walls.  On each chamber
piece define the local five-dimensional phase
$\tau_{5,a}^{(N)}(\gamma_j)$ as above.  At each wall insert the exact V27/V28
exterior-power line map, and at higher clean intersections insert the V28
Koszul braiding.

\begin{definition}[Constructible local-five-dimensional phase]
\label{def:v34-constructible-5d}
Let $\tau_{5,a}^{(N),\rm con}(\Gamma)$ be the closed graded product of the
finite-wall chamber phases and the exact V27--V29 wall maps, expressed in the
starting determinant superline.
\end{definition}

\begin{theorem}[Constructible infinite-wall realization]
\label{thm:v34-constructible-limit}
For every declared regular constructible gallery,
\begin{equation}
 \boxed{
 \mathfrak I_a^{\rm con}(\Gamma)
 =\lim_{T_5\to\infty}
  \tau_{5,a}^{(N),\rm con}(\Gamma).}
 \label{eq:v34-constructible-limit}
\end{equation}
The limit is independent of a refinement of the gallery and retains the V28
graded Koszul coherence.  Every fixed tangential source/BV derivative converges
with the same chamberwise $O(T_5^{-1})$ rate.
\end{theorem}

\begin{proof}
On each chamber piece use
Theorem~\ref{thm:v34-exact-5d-limit}.  The wall maps are independent of
$T_5$ and are already exact finite-regulator data.  V28--V29 prove their graded
composition and refinement coherence.  A finite product of the chamber limits
therefore gives the exact V33 constructible character.
\end{proof}

Thus the wall problem does not obstruct the five-dimensional realization.  It
only requires the already-proved stratified determinant-line category rather
than one globally gapped transfer cylinder.

\subsection{Relation to standard domain-wall transfer matrices}
\label{sec:v34-standard-DW}

The Cayley transfer function used above is the same finite-dimensional
rational transfer mechanism that appears in truncated-overlap/domain-wall
constructions.  In the standard notation one often writes the effective
four-dimensional sign approximation as
\begin{equation}
 \varepsilon_N(H)
 =\frac{1-T(H)^N}{1+T(H)^N},
 \label{eq:v34-standard-transfer}
\end{equation}
with $T$ the one-step fifth-direction transfer matrix.  Finite $N$ therefore
gives a rational approximation, while the exact overlap sign is obtained only
in the infinite-wall limit.  V34 proves the needed statements directly for the
native Cayley realization \eqref{eq:v34-K5}; no external domain-wall theorem is
used as a substitute for the Renewal proof.

The external lattice literature provides an important scope check: five-dimensional
Wilson/domain-wall determinants reproduce overlap/Ginsparg--Wilson chiral
phases after the appropriate wall-separation and Pauli--Villars limits, but
path integrability of an anomaly-free chiral theory is a further condition and
is not automatic at finite lattice spacing.  This is exactly the distinction
left open below.

\subsection{The remaining global problem is the lattice Chern--Simons/bordism descent}
\label{sec:v34-bordism-gate}

Theorem~\ref{thm:v34-exact-5d-limit} gives a local five-dimensional
realization of the exact V33 phase as a wall-separation limit.  What it does
not prove is that this phase depends only on the five-dimensional
ordinary-spin $G_6$ bordism class.

For two five-dimensional interpolations $c_1,c_2$ between the same boundary
families, define the closed-loop phase
\begin{equation}
 \mathcal Q_{5,a}(c_2-c_1)
 :=\lim_{T_5\to\infty}
 \frac{\Theta_{5,N}(c_2)}{\Theta_{5,N}(c_1)}.
 \label{eq:v34-Q5}
\end{equation}
After including the V32/V3 local phase normalization, exact global
integrability requires this loop phase to be a coboundary of local
interpolation functionals:
\begin{equation}
 \boxed{
 \mathcal Q_{5,a}(c_2-c_1)
 =\exp\bigl(iC_{5,a}(c_2)-iC_{5,a}(c_1)\bigr),}
 \label{eq:v34-CS-coboundary}
\end{equation}
with $C_{5,a}$ local, gauge/spin covariant and source complete.  Equivalently,
the renormalized five-dimensional loop phase must factor through the declared
bordism class.

\begin{proposition}[V34 reduces exact finite-cutoff integrability to a 5+1-dimensional local cohomology problem]
\label{prop:v34-cohomology-reduction}
On the V2 protected family, all nonlocal boundary/overlap dependence of the
exact V33 character is reproduced by the local five-dimensional transfer limit
of Theorem~\ref{thm:v34-exact-5d-limit}.  Therefore the only remaining
obstruction to
\begin{equation}
 \chi_{a,\rm SM}^{\eta}(\gamma)=1
 \label{eq:v34-final-character}
\end{equation}
for the complete faithful Standard-Model packet is whether the induced closed
five-dimensional phase admits the local decomposition
\eqref{eq:v34-CS-coboundary} and hence defines the faithful-spin bordism
character assumed in V33.
\end{proposition}

\begin{proof}
The boundary transfer phase converges to the exact Kato clutching factor, while
the V32 local measure phase supplies the complete locally normalized Wilson
factor.  Their quotient is exactly the V33 character.  Any residual dependence
on the choice of five-dimensional interpolation therefore appears solely as a
closed five-dimensional phase.  Removing it by a local interpolation
functional is precisely the coboundary/bordism condition
\eqref{eq:v34-CS-coboundary}.
\end{proof}

\begin{firewall}[Bordism triviality is still not imported]
The inherited vanishing of the faithful-spin bordism group can be used only
after \eqref{eq:v34-CS-coboundary}, or an equivalent lattice bordism theorem,
has been proved for the native five-dimensional phase.  V34 closes the local
five-dimensional realization but does not infer interpolation independence
from continuum anomaly cancellation.
\end{firewall}

\subsection{V34 anti-circularity audit}
\label{sec:v34-audit}

\begin{enumerate}[leftmargin=2.4em]
\item \textbf{Finite fifth extent is not called exact overlap.}
Theorem~\ref{thm:v34-finite-N-no-go} proves that a fixed finite transfer
thickness cannot equal the nonanalytic sign on an open protected family.
\item \textbf{The five-dimensional realization is genuinely local.}
Every finite-$N$ operator \eqref{eq:v34-K5} is nearest-neighbour in the fifth
direction and uses only the native local four-dimensional Wilson kernel.
\item \textbf{The exact phase is obtained as an infinite-wall limit at fixed
four-dimensional cutoff.}  No $a\downarrow0$ limit is used in
Theorem~\ref{thm:v34-exact-5d-limit}.
\item \textbf{The positive bulk determinant is removed explicitly.}
The phase comes from the boundary transfer minor after positive
reference/PV normalization; it is not contaminated by a hidden dynamical
magnitude.
\item \textbf{Kato transport is derived from the transfer product.}
The adiabatic theorem shows that polar-normalized boundary transfer converges
to the exact overlap clutching map, including fixed source derivatives.
\item \textbf{Index-changing walls are not crossed by a false gapped
adiabatic theorem.}  Chamber transfer phases are glued with the exact V27--V29
wall lines.
\item \textbf{Bordism invariance remains a separate global theorem.}
The local five-dimensional realization does not itself prove the lattice
Chern--Simons coboundary condition.
\end{enumerate}

\section{V34 status ledger}
\label{sec:v34-status}

\subsection*{Proved in V34}
\begin{enumerate}[leftmargin=2.2em]
\item A fixed finite fifth extent produces an analytic rational transfer
approximation to the Wilson sign and cannot equal the exact overlap sign on an
open protected spectral family.  Exact realization generically requires the
infinite-wall limit.
\item The Cayley transfer approximation converges exponentially to
$\operatorname{sgn}H_W$ at fixed four-dimensional cutoff, uniformly on every
protected Wilson-gap chart and through every prescribed fixed source derivative.
\item The native nearest-neighbour five-dimensional operator
$\mathbb K_{5,N}$ has exact one-particle transfer
$\mathcal T_N=\prod_j(I+a_5H_j)^{-1}(I-a_5H_j)$; its eliminated bulk determinant
is positive, so the nontrivial phase is carried by the boundary transfer map.
\item The polar-normalized occupied boundary transfer converges to Kato
parallel transport with $O(T_5^{-1})$ error, including every retained fixed
source/BV derivative.  Therefore its determinant converges to the exact V33
overlap clutching factor.
\item Multiplying the inverse boundary transfer phase by the exact V32
renormalized measure Wilson factor gives a local five-dimensional phase whose
infinite-wall limit is exactly the V33 integrability defect
$W_a^{\rm ren}T_a^{-1}=\chi_a^\eta$ at fixed four-dimensional cutoff.
\item The construction extends to the full Fredholm determinant line by taking
the cokernel/kernel transfer quotient and extends to V27--V29 constructible
wall galleries by inserting their exact graded exterior-power wall maps.
\item Thus the five-dimensional \emph{realization} part of the V33 gate is
closed after correcting the finite-wall quantifier.  The remaining exact global
problem is interpolation independence/bordism descent of the induced
five-dimensional phase.
\end{enumerate}

\subsection*{Inherited}
\begin{enumerate}[leftmargin=2.2em]
\item V1--V24: native Wilson/overlap hard and soft functional calculus,
source-complete Feshbach comparison, controlled zero-mode modules and
coefficientwise regulator universality on V2-consistent protected charts.
\item V25--V29: real/Pfaffian orientation, divisor and Wilson-wall transitions,
higher-rank crossing forms, graded wall coherence and nonsplit angular lines.
\item V30--V32: faithful geometric cancellation mechanism and coefficientwise
native-to-geometric phase matching with an explicitly retained exact flat
character.
\item V33: the exact native determinant holonomy is the overlap measure
integrability defect $W/T$ and has an intrinsic finite-cutoff reduced-eta
representation.
\end{enumerate}

\subsection*{Conditional}
\begin{enumerate}[leftmargin=2.2em]
\item The V34 five-dimensional operator is the native Cayley-transfer
realization built directly from $H_W$.  Identifying it by a local triangular
field redefinition with a preferred conventional Shamir/Moebius 5D stencil is
a representation theorem, not needed for the boundary phase result.
\item Exact equality with the overlap phase holds after the infinite-wall
limit $T_5\to\infty$, not at fixed finite fifth extent.
\item The local transfer adiabatic theorem assumes a smooth protected fixed-index
chamber with a uniform Wilson gap.  Genuine Wilson walls are handled only by
the separate V27--V29 constructible line maps.
\item The real/Pfaffian sector retains the V25--V29 orientation and parity
conditions; V34's complex boundary-transfer determinant does not replace them.
\item Positive IR/reference and fixed finite perturbative/source-order
hypotheses remain unchanged.
\end{enumerate}

\subsection*{Open beyond V34}
\begin{enumerate}[leftmargin=2.2em]
\item Prove the native five-dimensional lattice Chern--Simons coboundary
condition \eqref{eq:v34-CS-coboundary} for the complete faithful
Standard-Model packet, including all declared geometric, Yukawa, and source/BV
writers.
\item Show that the resulting renormalized infinite-wall phase is invariant
under five-dimensional bordism and therefore defines the faithful-spin $G_6$
bordism character required by V33.
\item Extend that interpolation-independence theorem through the V27--V29 wall
strata, where the boundary determinant line changes by the exact exterior-power
crossing factors.
\item If the coboundary theorem fails for an admissible native completion,
compute the residual closed five-dimensional flat character and place it in the
V16 microscopic-loop kernel; it remains invisible to the already-closed
coefficientwise continuum.
\end{enumerate}

\subsection*{Exact next load-bearing theorem}
\begin{center}
\fbox{\parbox{0.92\textwidth}{
\textbf{Native 5+1-dimensional lattice Chern--Simons / bordism descent.}
For the complete V2 faithful Standard-Model packet, take the V34 local
five-dimensional Cayley/domain-wall realization of a closed interpolation loop.
Prove that its renormalized infinite-wall phase has a local, gauge/spin-covariant
source-complete Chern--Simons density whose closed-loop value is a coboundary of
local interpolation functionals, equivalently that the phase depends only on
the ordinary-spin $G_6$ five-dimensional bordism class.  The proof must include
Higgs/Yukawa dependence, V27--V29 wall exterior-power factors and every retained
BV/source descendant.  Once this exact lattice descent theorem is proved, the
inherited faithful-spin bordism vanishing immediately forces
$\chi_{a,\rm SM}^{\eta}=1$ for every native loop.}}
\end{center}

\section*{Append-only continuation marker: V34}
\addcontentsline{toc}{section}{Append-only continuation marker: V34}
\textbf{End of V34.}  The complete V1--V33 body above is retained verbatim.
V34 goes upstream from the impossible demand of an exact fixed finite fifth
extent.  It proves that finite wall separation yields only a rational analytic
approximation to the overlap sign, constructs instead a nearest-neighbour local
five-dimensional Cayley domain-wall operator directly from the V2 Wilson
kernel, and proves that its reference-normalized boundary transfer phase
converges with all fixed source derivatives to the exact Kato/overlap clutching
factor as the wall separation tends to infinity.  Multiplying by the already
fixed V32 local measure phase therefore realizes the exact V33 eta/integrability
character as a limit of genuinely local five-dimensional native phases at each
fixed four-dimensional cutoff.  The sole remaining exact global gate is now the
5+1-dimensional lattice Chern--Simons/bordism descent of that phase.



\clearpage
\section{V35 continuation: coefficientwise $5{+}1$ lattice Chern--Simons descent and the exact-global-character firewall}
\label{sec:v35-continuation}

V34 realized the exact native determinant-line integrability character as the
infinite-wall limit of a genuinely local five-dimensional Cayley/domain-wall
transfer system.  Its handoff asked for an exact finite-spacing
Chern--Simons/bordism descent theorem.  The first upstream question is whether
that exact global statement is actually required by the coefficientwise
continuum programme.

It is not.  At every prescribed finite perturbative/EFT/source order, the V34
five-dimensional phase and all of its source/BV derivatives lie in one finite
local coefficient bank.  On a star-shaped interpolation chart the positive
interpolation-form degrees are explicitly contractible by the radial Poincare
homotopy, and this contraction acts only on the finite interpolation
coefficients, not on spacetime.  Therefore it preserves native locality and the
complete source/BV packet.  Once the V3--V5 Standard-Model anomaly restoration
has made the complete local phase curvature a covariant writer, that radial
homotopy produces one common local five-dimensional Chern--Simons primitive and
then one local scalar endpoint functional.

The result is exact at the quantifier actually used by the Renewal continuum
theorem: for every fixed coefficient packet, the renormalized closed
five-dimensional phase is a coboundary of one local interpolation functional.
What remains open is stronger and genuinely global: a flat or torsion character
on a noncontractible exact interpolation component can have zero local Taylor
jet to every finite order and still be nontrivial.  V35 keeps that possible
exact character separate rather than inferring its absence from perturbative
local cohomology.

Throughout, fix the explicit V2 Wilson/overlap branch and finite data
\begin{equation}
 \mathfrak q=(L,\Delta,r,N,\mathcal P_{\rm ext})
 \label{eq:v35-q}
\end{equation}
on one compact protected chart.  The positive IR/reference prescription of
V1--V34 remains in force.  Write \(\mathsf P_{\mathfrak q}\) for the complete
coefficient projector onto this finite packet.

\subsection{Interpolation coefficient space and the V34 phase connection}
\label{sec:v35-interpolation-space}

Let \(\mathcal C_{5,\mathfrak q}\) be a star-shaped finite-dimensional affine
chart of five-dimensional interpolation coefficients with fixed boundary data.
Its coordinates are the finitely many background/source jets retained by
\(\mathfrak q\).  A point \(c\in\mathcal C_{5,\mathfrak q}\) denotes one
coefficientwise five-dimensional interpolation, not an unrestricted
finite-amplitude field configuration.

For a sufficiently short oriented segment
\(c\to c+\epsilon\dot c\), let
\(\tau_{5,a}[c,c+\epsilon\dot c]\in U(1)\) be the V34 infinite-wall phase of
that segment after the already-fixed V32 local phase normalization.  Define the
local phase-connection one-form by
\begin{equation}
 \boxed{
 \mathbb A_{5,a;c}(\dot c)
 :=-i\left.\frac{d}{d\epsilon}\right|_{\epsilon=0}
 \log\tau_{5,a}[c,c+\epsilon\dot c].}
 \label{eq:v35-A5}
\end{equation}
Its curvature is
\begin{equation}
 \mathbb F_{6,a}:=d_{\mathcal C}\mathbb A_{5,a}.
 \label{eq:v35-F6}
\end{equation}
The notation emphasizes that one interpolation variation of the local
five-dimensional V34 phase produces the usual $4{+}2$ descent curvature.

\begin{theorem}[Locality and source completeness of the interpolation phase]
\label{thm:v35-local-phase-form}
For every fixed packet \(\mathfrak q\),
\begin{equation}
 \mathsf P_{\mathfrak q}\mathbb A_{5,a},
 \qquad
 \mathsf P_{\mathfrak q}\mathbb F_{6,a}
 \label{eq:v35-local-forms}
\end{equation}
are finite sums of native exponentially local five-dimensional kernels and
finite crossing-line terms.  They retain every prescribed geometric, gauge,
Higgs/Yukawa, BV, line/reference and contact-source derivative.  Their locality
constants may depend on \(\mathfrak q\), but are uniform in the ultraviolet
cutoff on the declared protected chart.
\end{theorem}

\begin{proof}
At finite fifth extent the V34 Cayley/domain-wall operator is finite range in
the fifth direction and belongs to the same four-dimensional weighted local
algebra as the native Wilson kernel.  The derivative of its logarithmic
reference-normalized determinant phase is a finite trace word containing
five-dimensional resolvents, native source vertices and the already-retained
line/reference writers.  V34 proves convergence of every fixed such derivative
as the wall separation tends to infinity.  The V1/V13 weighted Schur algebra
is closed under the finite products and gapped inverses entering these words.
Thus every retained coefficient of \(\mathbb A_{5,a}\) is local.

Differentiating once more in interpolation coefficient space gives
\(\mathbb F_{6,a}\).  V5's contraction-complete packet already contains every
later BV/source contraction of such a local writer.  At V27--V29 walls one does
not differentiate through the Wilson sign; instead one inserts the finite
crossing exterior-power line map, whose source derivatives are the Kato and
exterior-power derivatives already controlled there.
\end{proof}

\begin{corollary}[Bianchi identity]
\label{cor:v35-Bianchi}
On every fixed-index interpolation chamber,
\begin{equation}
 \boxed{d_{\mathcal C}\mathbb F_{6,a}=0.}
 \label{eq:v35-Bianchi}
\end{equation}
The constructible version remains true after inserting the exact V27--V29 wall
maps.
\end{corollary}

\begin{proof}
Inside a chamber this is \(d_{\mathcal C}^2\mathbb A_{5,a}=0\).  Across a
wall, V27--V29 identify the two determinant lines by an exact functorial
elementary modification; gallery refinement changes neither the line map nor
its graded composition.  Hence the constructible description adds no
spurious distributional curvature.
\end{proof}

\subsection{Relative interpolation forms are contractible without losing locality}
\label{sec:v35-relative-complex}

Let \(E\) be the Euler vector field on the star-shaped affine chart
\(\mathcal C_{5,\mathfrak q}\).  On a polynomial \(p\)-form with \(p>0\), set
\begin{equation}
 (h_{\mathcal C}\omega)(c)
 :=\int_0^1
 t^{p-1}\,\iota_E\omega(tc)\,dt.
 \label{eq:v35-radial-h}
\end{equation}
Let \(\pi_0\) evaluate interpolation coefficients at the chart origin and kill
positive interpolation form degree.

\begin{lemma}[Source-complete radial contraction]
\label{lem:v35-radial-contraction}
On the finite polynomial-jet form algebra retained by \(\mathfrak q\),
\begin{equation}
 \boxed{
 d_{\mathcal C}h_{\mathcal C}
 +h_{\mathcal C}d_{\mathcal C}
 =I-\pi_0.}
 \label{eq:v35-radial-identity}
\end{equation}
The operator \(h_{\mathcal C}\):
\begin{enumerate}[label=\textup{(R\arabic*)},leftmargin=3em]
\item preserves native spacetime locality and exponential kernel bounds;
\item commutes with every retained source derivative;
\item commutes with the fixed V5 classical BV differential in the
chronological V7 coefficient coordinates;
\item preserves exact gauge/spin covariance when the interpolation coefficients
are written in the same covariant local tensor basis used by V2--V7.
\end{enumerate}
\end{lemma}

\begin{proof}
The homotopy identity is the ordinary radial Poincare formula.  The integral
acts only on the finitely many interpolation coefficients and never integrates
over the physical lattice, so it cannot enlarge a spacetime kernel.  Source
and BV differentiations act on the coefficient values, while the radial
integration acts on their interpolation arguments; hence the operations
commute.  In the V2--V7 local tensor basis, gauge/spin transformations act by
finite representation matrices on the coefficient tensors, so the same radial
integral is equivariant.
\end{proof}

\begin{remark}[Finite $4{+}2$ descent complex]
\label{rem:v35-4plus2}
Equivalently, tensor the V5 finite local ESM bank with the finite polynomial
form algebra of \(\mathcal C_{5,\mathfrak q}\).  Lemma
\ref{lem:v35-radial-contraction} is an explicit contraction of every positive
interpolation-form degree.  Thus the interpolation directions expose the
$4{+}2$ descent but create no new local perturbative cohomology by themselves.
The only possible gauge/BV obstruction is inherited from the V5 local ESM
complex.
\end{remark}

\subsection{Why the complete Standard-Model curvature is a covariant local writer}
\label{sec:v35-covariance}

The radial contraction is useful only after the complete determinant/measure
packet has been restored into the common gauge/BV-covariant local bank.  This
is exactly what V3--V5 already provide.

\begin{theorem}[Native anomaly restoration supplies the covariant $4{+}2$ curvature packet]
\label{thm:v35-covariant-curvature}
After the chronological V5/V7 restoring primitives have been installed, the
complete three-generation Standard-Model coefficient
\begin{equation}
 \mathsf P_{\mathfrak q}\mathbb F_{6,a}
 \label{eq:v35-F6-projected}
\end{equation}
and all of its retained gauge, diffeomorphism, Lorentz/spin, Higgs/Yukawa,
BV, line/reference and contact descendants belong to the same finite covariant
local bank.  No nontrivial Standard-Model anomaly descendant remains.
\end{theorem}

\begin{proof}
V5 constructs one common restored finite local action/source packet and proves
that the only relevant ghost-number-one harmonic quotient is the standard ESM
anomaly space of Theorem~\ref{thm:v5-complex}.  The complete native packet has
zero projection onto that quotient by
Theorem~\ref{thm:v5-anomaly-zero}.  The V34 interpolation phase is built from
that same restored determinant/measure family, not from an independently
chosen fermion phase.  Consequently its local curvature and source/BV
variations transform in the restored covariant bank.  Any noncovariant harmonic
remainder would descend to a nonzero V5 anomaly coordinate, which is excluded.
\end{proof}

This is the precise point at which Standard-Model anomaly cancellation enters
V35.  The Poincare contraction below is purely local algebra; anomaly
cancellation guarantees that its output is an admissible common covariant
counterterm rather than a gauge-breaking primitive.

\subsection{One common coefficientwise lattice Chern--Simons primitive}
\label{sec:v35-CS-primitive}

Define
\begin{equation}
 \boxed{
 \mathbb B_{5,a,\mathfrak q}^{\rm CS}
 :=h_{\mathcal C}
   \bigl(\mathsf P_{\mathfrak q}\mathbb F_{6,a}\bigr).}
 \label{eq:v35-B5}
\end{equation}

\begin{theorem}[Source-complete coefficientwise $5{+}1$ Chern--Simons descent]
\label{thm:v35-common-CS}
The one-form \(\mathbb B_{5,a,\mathfrak q}^{\rm CS}\) is a native local,
gauge/spin-covariant five-dimensional Chern--Simons/measure counterterm with the
complete retained source/BV packet, and
\begin{equation}
 \boxed{
 d_{\mathcal C}\mathbb B_{5,a,\mathfrak q}^{\rm CS}
 =\mathsf P_{\mathfrak q}\mathbb F_{6,a}.}
 \label{eq:v35-curvature-exact}
\end{equation}
All descendants are determined by this same primitive and are not independently
fitted.
\end{theorem}

\begin{proof}
Apply the radial identity
\eqref{eq:v35-radial-identity} to the closed interpolation two-form
\(\mathsf P_{\mathfrak q}\mathbb F_{6,a}\).  Since \(\pi_0\) kills positive
form degree and \(d_{\mathcal C}\mathbb F_{6,a}=0\),
\[
 d_{\mathcal C}h_{\mathcal C}
 \bigl(\mathsf P_{\mathfrak q}\mathbb F_{6,a}\bigr)
 =\mathsf P_{\mathfrak q}\mathbb F_{6,a}.
\]
Locality, source completeness and covariance follow from
Lemma~\ref{lem:v35-radial-contraction} and
Theorem~\ref{thm:v35-covariant-curvature}.  Because every later BV/source
contraction is already present in the V5 packet, the one primitive carries all
descendants simultaneously.
\end{proof}

\subsection{From curvature descent to an interpolation coboundary}
\label{sec:v35-coboundary}

Subtract the Chern--Simons one-form from the V34 phase connection:
\begin{equation}
 \mathbb A_{5,a}^{\rm flat}
 :=\mathsf P_{\mathfrak q}\mathbb A_{5,a}
  -\mathbb B_{5,a,\mathfrak q}^{\rm CS}.
 \label{eq:v35-Aflat}
\end{equation}
Theorem~\ref{thm:v35-common-CS} gives
\begin{equation}
 d_{\mathcal C}\mathbb A_{5,a}^{\rm flat}=0.
 \label{eq:v35-Aflat-closed}
\end{equation}
Define the scalar local interpolation functional
\begin{equation}
 \boxed{
 C_{5,a,\mathfrak q}^{\rm loc}
 :=h_{\mathcal C}\mathbb A_{5,a}^{\rm flat}.}
 \label{eq:v35-C5loc}
\end{equation}
Since \(\pi_0\) kills one-forms,
\begin{equation}
 \boxed{
 d_{\mathcal C}C_{5,a,\mathfrak q}^{\rm loc}
 =\mathbb A_{5,a}^{\rm flat}.}
 \label{eq:v35-C5loc-d}
\end{equation}

\begin{theorem}[Coefficientwise lattice Chern--Simons coboundary]
\label{thm:v35-coboundary}
Let \(c_1,c_2\in\mathcal C_{5,\mathfrak q}\) be two declared five-dimensional
interpolations with the same boundary data.  After adding the single common
local Chern--Simons primitive of
Theorem~\ref{thm:v35-common-CS}, the retained coefficient of their V34 relative
phase obeys
\begin{equation}
 \boxed{
 \mathsf P_{\mathfrak q}\log
 \mathcal Q_{5,a}^{\rm ren}(c_2-c_1)
 =i\Bigl[
 C_{5,a,\mathfrak q}^{\rm loc}(c_2)
 -C_{5,a,\mathfrak q}^{\rm loc}(c_1)
 \Bigr].}
 \label{eq:v35-coboundary}
\end{equation}
Equivalently, after the endpoint interpolation functional is included in the
common V7 reference prescription,
\begin{equation}
 \boxed{
 \mathsf P_{\mathfrak q}\log
 \mathcal Q_{5,a}^{\rm ren}=0}
 \label{eq:v35-loop-zero}
\end{equation}
for every closed interpolation loop contractible inside the protected
coefficient chart.
\end{theorem}

\begin{proof}
The logarithmic derivative of the renormalized phase is
\(\mathbb A_{5,a}^{\rm flat}=d_{\mathcal C}C_{5,a,\mathfrak q}^{\rm loc}\).
Integrate this exact one-form between the two interpolation points.  The local
Chern--Simons term and scalar endpoint functional were each constructed once on
the complete packet, so no pathwise or graphwise finite refit is introduced.
The second statement is the same identity after absorbing the endpoint
functional into the common V7 reference coordinate.
\end{proof}

\begin{corollary}[All retained coefficients of the exact eta residue vanish locally]
\label{cor:v35-eta-coeff-zero}
For every prescribed finite packet \(\mathfrak q\), the V33/V34 native eta
character on a contractible protected interpolation loop satisfies
\begin{equation}
 \boxed{
 \mathsf P_{\mathfrak q}
 \log\chi_{a,\rm SM}^{\eta}=0}
 \label{eq:v35-eta-coeff-zero}
\end{equation}
after the common V35 local Chern--Simons/reference normalization.  The statement
is exact coefficientwise at fixed lattice spacing; no \(a\downarrow0\) limit
is required.
\end{corollary}

\begin{proof}
V34 identifies the infinite-wall five-dimensional phase with the exact V33
integrability character at fixed four-dimensional cutoff.  Apply
Theorem~\ref{thm:v35-coboundary}.
\end{proof}

\subsection{Constructible Wilson-wall galleries}
\label{sec:v35-walls}

A V27--V29 gallery is a finite sequence of fixed-index interpolation chambers
separated by finite crossing-line maps.  V35 does not smooth those walls into
an artificial gapped family.

\begin{theorem}[Coefficientwise descent is compatible with Wilson-wall gluing]
\label{thm:v35-wall-descent}
For a regular constructible gallery, apply
Theorem~\ref{thm:v35-common-CS} on every fixed-index chamber and retain the
exact V27--V29 exterior-power wall maps.  Then:
\begin{enumerate}[label=\textup{(W\arabic*)},leftmargin=3em]
\item chamber Chern--Simons endpoint terms telescope under refinement;
\item wall exterior-power factors are unchanged by the local Chern--Simons
normalization;
\item the V28 Koszul braiding at higher wall intersections remains exact;
\item every retained source/BV descendant satisfies the same constructible
cocycle.
\end{enumerate}
Hence \eqref{eq:v35-coboundary} holds for every constructible interpolation
loop contractible in the resolved coefficient chart.
\end{theorem}

\begin{proof}
The V35 primitive is defined separately in each fixed-index local bank.  On a
common chamber overlap two choices differ by a closed local one-form; the radial
homotopy makes this difference an admissible V7 scalar reference coboundary.
When a wall is crossed, its finite determinant-line elementary modification is
an independent soft factor already proved in V27--V29 and is not altered by a
local phase counterterm inside the chamber.  Concatenating the chamber
identities therefore telescopes all artificial subdivision terms.  V28 already
proves the graded coherence of the remaining wall factors.
\end{proof}

\subsection{Relation to the perturbative $4{+}2$ lattice cohomology method}
\label{sec:v35-external-comparison}

The structure above is the native Renewal realization of the standard lattice
chiral-gauge strategy in which anomalous conservation and measure integrability
are combined into a local $4{+}2$-dimensional cohomology problem.  For
Ginsparg--Wilson fermions the corresponding non-Abelian local cohomology problem
is known perturbatively to all orders in the gauge potential, whereas exact
global integrability is a stronger question.  V35 does not import that theorem
as a substitute for the Renewal proof: the positive interpolation-form complex
is contracted explicitly by \eqref{eq:v35-radial-h}, while V3--V5 supply the
native Standard-Model anomaly restoration needed for covariance.

The comparison is nevertheless a useful scope check.  It confirms that
all-orders perturbative local integrability and exact nonperturbative global
integrability are distinct mathematical targets.

\subsection{Why V35 does not prove the exact finite-cutoff global character}
\label{sec:v35-global-firewall}

Theorem~\ref{thm:v35-coboundary} is local on a star-shaped coefficient chart.
It annihilates every retained Taylor/source coefficient of the closed-loop
phase.  It does not exclude a globally flat character living on a
noncontractible component of the exact interpolation space.

\begin{proposition}[Perturbative blindness to a flat global character]
\label{prop:v35-flat-blindness}
Let \(\rho:\pi_1(\mathcal M)\to U(1)\) be a nontrivial flat unitary character
on a noncontractible exact interpolation space \(\mathcal M\).  Every local
curvature coefficient and every finite Taylor jet of a local parallel frame can
vanish while \(\rho\) remains nontrivial on a large loop.  Therefore
\eqref{eq:v35-eta-coeff-zero} does not imply
\begin{equation}
 \chi_{a,\rm SM}^{\eta}(\gamma)=1
 \label{eq:v35-exact-global-not-implied}
\end{equation}
for an arbitrary exact noncontractible loop.
\end{proposition}

\begin{proof}
A flat line bundle has zero curvature.  On every simply connected coordinate
chart it admits a parallel frame, so all local connection coefficients and all
of their Taylor derivatives vanish in that frame.  Its transition constants
around a noncontractible loop can nevertheless multiply to a nontrivial unitary
number.  This is precisely the information stored in \(\rho\) and invisible to
local perturbative jets.
\end{proof}

Lattice global anomalies provide concrete examples of this distinction in
chiral gauge theory: local anomaly cancellation and local measure equations do
not by themselves eliminate every possible global obstruction.  V35 therefore
keeps the exact flat/torsion character as a separate stronger problem.

\begin{firewall}[Coefficientwise Chern--Simons descent is not a nonperturbative bordism theorem]
V35 proves the V34 coboundary condition at every prescribed finite
perturbative/EFT/source order, with all descendants, on protected coefficient
charts.  It does not prove that the exact finite-cutoff infinite-wall phase is
a globally defined bordism character on the full nonperturbative lattice
configuration space.  The inherited faithful-spin bordism vanishing may not be
invoked to remove a flat character until that exact global bridge is proved.
\end{firewall}

\subsection{Coefficientwise continuum consequence}
\label{sec:v35-continuum}

For the original Renewal continuum programme this distinction is already
sufficient.

\begin{theorem}[The five-dimensional phase introduces no new coefficientwise continuum datum]
\label{thm:v35-no-new-continuum-datum}
Fix \(\mathfrak q\).  On the V2-consistent protected Standard-Model branch,
the V34 five-dimensional phase, after the single common V35 local
Chern--Simons/reference normalization, contributes no independent retained
coefficient to the ultraviolet continuum state.  In particular, the common
BV/BRST/master vector, determinant-line source packet, stress/composite/contact
coefficients and all normalized connected coefficients obtained in V7 remain
unchanged.
\end{theorem}

\begin{proof}
Corollary~\ref{cor:v35-eta-coeff-zero} removes every retained coefficient of
the closed five-dimensional interpolation phase.  The counterterm and endpoint
functional belong to the same finite local reference/source bank whose scheme
covariance was proved in V7.  They are therefore a finite scheme/reference
change rather than a new observable coordinate.  V7 projectivity freezes the
lower-order choices when higher terminal loop orders are appended.
\end{proof}

\begin{corollary}[Order projectivity of the V35 descent]
\label{cor:v35-projective}
The V35 Chern--Simons primitive can be chosen chronologically so that increasing
the terminal loop/source order appends only new higher-order coefficients and
does not change the previously fixed local Chern--Simons/reference terms.
\end{corollary}

\begin{proof}
Use the same chronological coherent-family convention as V7.  At each new
order, the finite local curvature/source packet is assembled after all lower
coefficients are fixed, and the radial homotopy acts only on that new finite
coefficient layer.  V5 contraction-completeness prevents a later source/BV
contraction from forcing a refit of an earlier primitive.
\end{proof}

\subsection{V35 anti-circularity audit}
\label{sec:v35-audit}

\begin{enumerate}[leftmargin=2.4em]
\item \textbf{The exact global theorem was not forced from local anomaly
cancellation.}  V35 restricts its positive theorem to finite coefficient
packets on protected interpolation charts.
\item \textbf{The $4{+}2$ interpolation directions are contracted explicitly.}
Equation~\eqref{eq:v35-radial-identity} is the load-bearing local homotopy; no
new higher-dimensional cohomology classification is assumed.
\item \textbf{Standard-Model anomaly cancellation enters only where needed.}
V3--V5 guarantee that the complete local phase curvature is an admissible
covariant writer before the radial primitive is taken.
\item \textbf{One common primitive is used.}  Gauge, gravity, spin,
Higgs/Yukawa, ghost/antifield, line/reference and contact descendants are
carried by the same local interpolation functional rather than fitted
separately.
\item \textbf{Wilson walls are not smoothed away.}  V27--V29 exterior-power
crossing factors remain exact soft data and are inserted between fixed-index
chamber descents.
\item \textbf{Coefficientwise exactness is distinguished from analytic
finite-amplitude exactness.}  No convergence of the perturbative/source series
is asserted.
\item \textbf{A flat global character remains in the ledger.}
Proposition~\ref{prop:v35-flat-blindness} prevents the inherited bordism theorem
from being applied prematurely.
\end{enumerate}

\section{V35 status ledger}
\label{sec:v35-status}

\subsection*{Proved in V35}
\begin{enumerate}[leftmargin=2.2em]
\item The V34 infinite-wall interpolation phase defines, at every prescribed
finite coefficient/source order, a local source-complete phase connection and
local $4{+}2$ curvature packet.
\item The positive interpolation-form complex on a star-shaped coefficient
chart has an explicit radial contraction which preserves spacetime locality,
source derivatives, the fixed BV differential and exact gauge/spin covariance.
\item After the inherited V3--V5 Standard-Model anomaly restoration, the
complete curvature packet is a covariant local writer with no surviving
nontrivial anomaly descendant.
\item One common radial primitive therefore gives a source-complete local
$5{+}1$ lattice Chern--Simons/measure counterterm satisfying
$d_{\mathcal C}\mathbb B_5^{\rm CS}=\mathsf P_{\mathfrak q}\mathbb F_6$.
\item After subtracting that primitive, the retained interpolation phase
connection is closed; a second radial contraction gives one local scalar
endpoint functional.  Hence the V34 closed-loop phase is a local coboundary
coefficient by coefficient at fixed lattice spacing.
\item Equivalently, after one common local Chern--Simons/reference
normalization, every retained perturbative/source coefficient of the exact V33
eta/integrability character vanishes on contractible protected interpolation
loops.
\item The construction is compatible with V27--V29 constructible Wilson-wall
galleries and preserves their exact exterior-power/Koszul transition data.
\item The V35 descent is projective in perturbative order and introduces no new
coefficientwise continuum observable beyond the V7 finite scheme/reference
groupoid.
\end{enumerate}

\subsection*{Inherited}
\begin{enumerate}[leftmargin=2.2em]
\item V1--V7: native source-complete fixed-order ESM continuum, common BV
restoration, Standard-Model anomaly cancellation, forest induction, UV Cauchy
convergence and scheme/order projectivity.
\item V8--V24: regulator comparison, fixed-index zero-mode geometry,
controlled module analysis, soft norm-resolvent anchoring and the V2-consistent
covariant Feshbach theorem.
\item V25--V29: real/Pfaffian orientation, accidental divisors, genuine Wilson
index jumps, higher-rank wall lines and nonsplit angular constructible
superlines.
\item V30--V33: faithful geometric cancellation mechanism, corrected
phase-faithful bridge, coefficientwise phase matching and the exact intrinsic
eta/integrability character.
\item V34: local five-dimensional Cayley/domain-wall realization whose
infinite-wall phase equals the exact V33 character at fixed four-dimensional
cutoff.
\end{enumerate}

\subsection*{Conditional}
\begin{enumerate}[leftmargin=2.2em]
\item The positive V35 theorem is coefficientwise on a protected star-shaped
interpolation/source chart.  It is not one finite-amplitude theorem on the full
lattice configuration space.
\item The interpolation coefficients must be expressed in the same finite
covariant tensor coordinates used by the V2--V7 local bank so that the radial
homotopy preserves covariance.
\item The V27--V29 wall strata must satisfy their declared regularity,
orientation and reduced-gap hypotheses.
\item The V25 real/Pfaffian orientation conditions remain separate from the
complex Chern--Simons descent.
\item Positive IR/reference and fixed finite loop/EFT/source-order hypotheses
remain unchanged.
\end{enumerate}

\subsection*{Open beyond V35}
\begin{enumerate}[leftmargin=2.2em]
\item Prove or falsify the \emph{exact nonperturbative finite-cutoff} lattice
Chern--Simons integrability theorem on the full admissible faithful
Standard-Model configuration space.  V35 kills every perturbative coefficient
but does not eliminate a flat/torsion holonomy character.
\item Construct an exact native bordism map for the V34 infinite-wall phase and
show that its residual flat character factors through
$\Omega_5^{\mathrm{Spin}}(BG_6)$, including the V27--V29 wall strata.
\item If a residual exact character survives, compute it explicitly and place
it in the V16 microscopic-loop kernel; it remains invisible to the
coefficientwise continuum theorem.
\item The independent stronger problems remain: removal of the positive IR
scale, analytic finite-coupling RG trajectories, perturbative summability and a
nonperturbative ESM completion.
\end{enumerate}

\subsection*{Exact next load-bearing theorem}
\begin{center}
\fbox{\parbox{0.92\textwidth}{
\textbf{Exact global lattice bordism / torsion-character theorem.}
The coefficientwise $5{+}1$ local descent problem is closed in V35.  The
remaining stronger question is purely global: on the full admissible
ordinary-spin faithful $G_6$ lattice configuration/interpolation space, prove
that the V34 infinite-wall phase after the V35 local Chern--Simons
normalization defines an exact bordism character, including Wilson-wall
strata, and compute its possible flat/torsion component.  If this exact native
character factors through the inherited group
$\Omega_5^{\mathrm{Spin}}(BG_6)=0$, then the finite-cutoff equality
$\chi_{a,\rm SM}^{\eta}=1$ follows globally.  Otherwise the surviving flat
character is a genuinely nonperturbative microscopic invariant which does not
alter the already-proved coefficientwise ESM continuum.}}
\end{center}

\section*{Append-only continuation marker: V35}
\addcontentsline{toc}{section}{Append-only continuation marker: V35}
\textbf{End of V35.}  The complete V1--V34 body above is retained verbatim.
V35 goes upstream from the stronger demand of exact nonperturbative
finite-spacing interpolation independence and closes the load-bearing problem
at the coefficientwise quantifiers of the Renewal continuum theorem.  It
constructs the local interpolation-phase connection of the V34 five-dimensional
system, proves that the positive $4{+}2$ interpolation-form complex is
contractible by a source-complete radial homotopy, uses the inherited V3--V5
anomaly restoration to keep the resulting primitive in the common covariant
local bank, and builds one $5{+}1$ lattice Chern--Simons counterterm followed by
one local scalar endpoint functional.  The V34 closed-loop phase is therefore a
local coboundary at every prescribed fixed perturbative/EFT/source order and
introduces no new coefficientwise continuum datum.  A possible exact
flat/torsion character on a noncontractible finite-cutoff interpolation
component is retained as a separate stronger global problem rather than being
inferred away from local anomaly cancellation.



\clearpage
\section{V36 continuation: exact local flattening, fibered bordism, and the native global residue}
\label{sec:v36-continuation}

V35 closed the $5{+}1$ interpolation-phase problem at every prescribed finite
perturbative/EFT/source order and correctly retained a possible exact flat or
torsion residue.  The next question is not answered by the bare statement
$\Omega_5^{\mathrm{Spin}}(BG_6)=0$.  A loop in the native configuration space
has two different notions of ``bounding'': its mapping torus may bound an
arbitrary six-dimensional spin $G_6$ manifold, while the loop itself may fail
to bound a disk of admissible four-dimensional native backgrounds.  The latter
is a strictly stronger, fibered statement.

Before separating these two bordism notions, V36 first removes one ambiguity
left in V35.  At fixed lattice spacing and on one protected star-shaped chart,
the exact relative determinant curvature is itself an exponentially local
smooth two-form.  The ordinary radial homotopy therefore gives an \emph{exact}
quasilocal primitive, without expanding in loop or source order.  This exact
primitive is stronger than the finite local EFT counterterms of the original
continuum programme and is used only to expose the remaining global topology.
After this local flattening, the entire exact residue is a constant Cech
$U(1)$ cocycle, equivalently a flat character of the native protected
configuration space.

The new results are:
\begin{enumerate}[label=\textup{(V36.\arabic*)},leftmargin=3.8em]
\item exact quasilocal flattening of the native--geometric relative determinant
connection on every protected star-shaped chart;
\item construction of the residual flat Cech character and proof that V35
kills all of its perturbative local jets but not necessarily the character;
\item distinction between ordinary spin bordism of the mapping torus and a
fibered/native null-homotopy through admissible four-dimensional backgrounds;
\item definition of the native fibered-bordism defect group on which the exact
flat residue lives;
\item a discretization/lifting theorem: every smooth, spectrally admissible
fibered geometric filling has a native V24--V29 resolved filling for all
sufficiently small lattice spacing;
\item exact triviality of the V34/V35 residue on every loop admitting such a
fibered filling.
\end{enumerate}

Thus the remaining exact obstruction is no longer ``some torsion phase'' in an
unstructured sense.  It is a character of a precise native/fibered bordism
defect group.  Ordinary faithful-spin bordism vanishing kills it only after
one proves that this defect group is trivial, or that the native phase factors
through the ordinary bordism forgetful map.

\subsection{Exact relative curvature on a protected chart}
\label{sec:v36-exact-local}

Fix a compact star-shaped protected chart $U$ on a V24 V2-consistent
fixed-index branch.  Use the V31 canonical topological determinant-line
identification to regard the native and geometric determinant connections as
connections on one complex line.  Their exact difference is an imaginary
one-form
\begin{equation}
 \mathcal A_a^{\rm rel}
 :=\mathcal A_a^{\rm nat}-\mathcal A^{\rm geom},
 \label{eq:v36-Arel}
\end{equation}
with exact curvature
\begin{equation}
 \mathcal F_a^{\rm rel}
 :=d\mathcal A_a^{\rm rel}.
 \label{eq:v36-Frel}
\end{equation}

\begin{lemma}[Exact quasilocality of the relative curvature]
\label{lem:v36-relative-curvature-local}
On the declared protected chart, $\mathcal F_a^{\rm rel}$ is a smooth
imaginary-valued two-form whose coefficient kernels and every prescribed fixed
source/BV derivative are cutoff-uniform exponentially local after one fixed
loss of exponent.  The same statement holds on each regular V27--V29 wall
stratum after the crossing block is retained among the soft variables.
\end{lemma}

\begin{proof}
At finite cutoff the determinant connection and curvature are finite-dimensional
Riesz/Kato trace words in the chiral block, its reduced inverse, and the
kernel/cokernel projectors.  V13 proves exponential locality of the protected
Riesz projectors and Moore--Penrose inverse; V24 proves the corresponding
source-complete soft/hard bounds for the V2-consistent native block.  The
geometric comparison uses the same fixed low-energy carrier under the V31
topological bridge.  Products and traces of the resulting weighted kernels are
again exponentially local by the V13 weighted Schur algebra.  On a wall the
finite crossing projector is retained and only the uniformly invertible
complement is differentiated, exactly as in V26--V29.
\end{proof}

Let $E$ be the radial vector field on the finite chart coordinates and define,
for a positive-degree form $\omega$,
\begin{equation}
 (h_U\omega)_x
 :=\int_0^1 t^{p-1}\,\iota_E\omega(tx)\,dt,
 \qquad \omega\in\Omega^p(U),\ p>0.
 \label{eq:v36-radial-h}
\end{equation}
Because $U$ is star-shaped,
\begin{equation}
 \boxed{d h_U+h_Ud=I\quad\text{on positive-degree forms}.}
 \label{eq:v36-radial-identity}
\end{equation}

\begin{theorem}[Exact local quasilocal flattening]
\label{thm:v36-exact-local-flattening}
Define
\begin{equation}
 \boxed{
 \mathcal B_{a,U}^{\rm ex}
 :=h_U\mathcal F_a^{\rm rel}.}
 \label{eq:v36-Bex}
\end{equation}
Then $\mathcal B_{a,U}^{\rm ex}$ is an exact exponentially local
source-complete one-form and
\begin{equation}
 \boxed{d\mathcal B_{a,U}^{\rm ex}=\mathcal F_a^{\rm rel}.}
 \label{eq:v36-dB-F}
\end{equation}
Consequently the locally renormalized relative connection
\begin{equation}
 \boxed{
 \mathcal A_{a,U}^{\rm flat}
 :=\mathcal A_a^{\rm rel}-\mathcal B_{a,U}^{\rm ex}}
 \label{eq:v36-Aflat}
\end{equation}
is exactly flat on $U$:
\begin{equation}
 d\mathcal A_{a,U}^{\rm flat}=0.
 \label{eq:v36-flat}
\end{equation}
Every fixed source/BV derivative of the construction is obtained by
differentiating the same radial integral and remains exponentially local.
\end{theorem}

\begin{proof}
The curvature is closed by the Bianchi identity for a line connection, so
\eqref{eq:v36-radial-identity} gives
$d h_U\mathcal F_a^{\rm rel}=\mathcal F_a^{\rm rel}$.  Locality and source
bounds follow from Lemma~\ref{lem:v36-relative-curvature-local}: the radial
integral acts only on the finite background/interpolation coefficients and
does not enlarge spacetime support.  Subtracting the primitive therefore gives
\eqref{eq:v36-flat} exactly.
\end{proof}

\begin{firewall}[Exact quasilocal flattening is stronger than the original renormalization scheme]
The one-form $\mathcal B_{a,U}^{\rm ex}$ is generally an exact finite-amplitude
quasilocal functional, not a finite EFT polynomial belonging to the original
V7 local counterterm bank.  It is introduced only to isolate the exact global
phase of the finite-cutoff determinant line.  V35 remains the correct theorem
for the coefficientwise perturbative continuum, where one keeps only the
finite local jet of this construction.
\end{firewall}

\subsection{The residual exact flat Cech character}
\label{sec:v36-flat-cech}

Choose a finite good cover $(U_\alpha)_\alpha$ of a connected protected
component and perform Theorem~\ref{thm:v36-exact-local-flattening} on each
chart.  Since each $U_\alpha$ is contractible, choose a unit parallel frame
$s_\alpha$ for the flattened relative connection.

\begin{proposition}[Constant overlap cocycle]
\label{prop:v36-constant-cocycle}
On every nonempty connected overlap $U_{\alpha\beta}$ there is a constant
$c_{\alpha\beta}\in U(1)$ such that
\begin{equation}
 s_\beta=c_{\alpha\beta}s_\alpha,
 \label{eq:v36-c-alpha-beta}
\end{equation}
and on triple overlaps
\begin{equation}
 \boxed{c_{\alpha\beta}c_{\beta\gamma}=c_{\alpha\gamma}.}
 \label{eq:v36-cech-cocycle}
\end{equation}
The resulting class
\begin{equation}
 \boxed{
 [c_a]\in
 \check H^1(\mathcal M_a^{\rm prot};U(1)_{\rm const})}
 \label{eq:v36-flat-class}
\end{equation}
is independent of the chosen local flattening primitives and frames.
Its loop evaluation equals the exact relative determinant holonomy after the
local quasilocal flattening.
\end{proposition}

\begin{proof}
Both frames are parallel for the same flat relative connection on the overlap,
so their ratio has zero differential and is constant.  The triple-overlap
identity is immediate.  Changing a local parallel frame multiplies
$c_{\alpha\beta}$ by a Cech coboundary.  Changing the local primitive by an
exact one-form is likewise absorbed into a local scalar rephasing.  The product
of the overlap constants encountered around a loop is precisely the holonomy
of the flat connection.
\end{proof}

On a path-connected component this class can equivalently be viewed as a flat
character
\begin{equation}
 \rho_a:\pi_1(\mathcal M_a^{\rm prot})\longrightarrow U(1).
 \label{eq:v36-rho}
\end{equation}
V35 proves that every prescribed finite perturbative/source jet of this exact
residue is invisible to the coefficientwise continuum.  V36 now identifies the
residue itself without claiming it is zero.

\subsection{Ordinary bordism versus fibered/native null-homotopy}
\label{sec:v36-fibered-vs-ordinary}

A loop of four-dimensional backgrounds determines a five-dimensional mapping
torus.  However, nullbordism of that mapping torus is weaker than contraction
of the loop through four-dimensional backgrounds.

\begin{definition}[Fibered/native filling]
\label{def:v36-fibered-filling}
Let $\gamma:S^1\to\widetilde{\mathcal M}_a^{\rm prot}$ be a loop in the
resolved native configuration space, where V27--V29 wall strata are allowed.
A \emph{fibered/native filling} is a continuous map
\begin{equation}
 \Gamma:D^2\longrightarrow\widetilde{\mathcal M}_a^{\rm prot},
 \qquad
 \Gamma|_{\partial D^2}=\gamma,
 \label{eq:v36-native-filling}
\end{equation}
which is smooth on each fixed-index stratum and meets the declared wall strata
with the V26--V29 regularity conditions.  Equivalently, its total space is a
family of four-dimensional native backgrounds over the disk, not merely an
arbitrary six-dimensional manifold bounding the mapping torus.
\end{definition}

\begin{proposition}[Fibered filling kills the exact flat residue]
\label{prop:v36-fibered-kills}
If $\gamma$ admits a fibered/native filling, then
\begin{equation}
 \boxed{\rho_a(\gamma)=1.}
 \label{eq:v36-rho-filling-zero}
\end{equation}
The same conclusion holds for a constructible loop crossing V27--V29 walls,
provided the exact wall line maps are included in the boundary transport.
\end{proposition}

\begin{proof}
Pull the flattened relative line and its constructible wall identifications
back to the disk.  Since $D^2$ is simply connected and the connection is flat
on every chamber, parallel transport from one interior base point defines a
single frame on each chamber.  The V27--V29 wall maps identify these frames
coherently across strata.  The boundary holonomy of the resulting global
constructible line is therefore one.
\end{proof}

By contrast, ordinary spin bordism supplies a six-manifold $X^6$ with
\begin{equation}
 \partial X^6=Y_\gamma,
 \label{eq:v36-ordinary-bordism}
\end{equation}
where $Y_\gamma$ is the mapping torus.  It need not supply a projection
$X^6\to D^2$ with four-dimensional fibres realizing
\eqref{eq:v36-native-filling}.

\begin{proposition}[Ordinary bordism vanishing does not imply native loop contraction]
\label{prop:v36-bordism-not-fibered}
The implication
\begin{equation}
 [Y_\gamma]=0\in\Omega_5^{\mathrm{Spin}}(BG_6)
 \quad\Longrightarrow\quad
 \gamma\text{ admits a fibered/native filling}
 \label{eq:v36-false-implication}
\end{equation}
is not a formal consequence of bordism theory.  Therefore the inherited
identity
$\Omega_5^{\mathrm{Spin}}(BG_6)=0$ cannot by itself trivialize the character
\eqref{eq:v36-rho}.
\end{proposition}

\begin{proof}
Ordinary bordism forgets the presentation of a five-manifold as a mapping
torus of a fixed four-dimensional carrier and allows an arbitrary bounding
six-manifold.  A native filling retains the much stronger datum of a disk of
four-dimensional configurations.  Abstractly, a map from a loop space to the
zero bordism group can annihilate every bordism class while the source still
has nontrivial fundamental group.  Hence no converse lifting follows from the
vanishing of the target group.
\end{proof}

This is the precise reason the inherited faithful-spin theorem could not be
invoked directly in V35.

\subsection{The native fibered-bordism defect group}
\label{sec:v36-defect-group}

Let $\mathcal N_a^{\rm fill}$ be the normal subgroup of
$\pi_1(\widetilde{\mathcal M}_a^{\rm prot})$ generated by loops admitting
fibered/native fillings.  Define
\begin{equation}
 \boxed{
 \mathfrak B_{5,a}^{\rm nat}
 :=
 \pi_1(\widetilde{\mathcal M}_a^{\rm prot})/
 \mathcal N_a^{\rm fill}.}
 \label{eq:v36-native-bordism-defect}
\end{equation}

\begin{theorem}[Exact residue factors through the native defect group]
\label{thm:v36-residue-defect}
The exact flat relative character descends uniquely to
\begin{equation}
 \boxed{
 \bar\rho_a:\mathfrak B_{5,a}^{\rm nat}\longrightarrow U(1).}
 \label{eq:v36-rho-bar}
\end{equation}
The exact finite-cutoff native phase is globally trivial if and only if
$\bar\rho_a$ is the trivial character.  In particular,
\begin{equation}
 \mathfrak B_{5,a}^{\rm nat}=0
 \quad\Longrightarrow\quad
 \rho_a\equiv1.
 \label{eq:v36-defect-zero}
\end{equation}
\end{theorem}

\begin{proof}
Proposition~\ref{prop:v36-fibered-kills} shows that $\rho_a$ is one on every
generator of $\mathcal N_a^{\rm fill}$ and hence on the normal subgroup they
generate.  The universal property of the quotient gives
\eqref{eq:v36-rho-bar}.  The remaining statements are immediate.
\end{proof}

Thus the exact global problem has been converted into a concrete topological
question about the native resolved configuration space.  The ordinary bordism
group is no longer confused with this stronger quotient.

\subsection{Spectrally admissible geometric fillings lift to native fillings}
\label{sec:v36-lifting}

The preceding distinction is useful only if one can recognize a large class of
fibered fillings.  V24 supplies exactly such a lifting theorem.

\begin{definition}[Spectrally admissible geometric filling]
\label{def:v36-spectral-filling}
A smooth disk family of geometric gauge/spin/coframe/Higgs backgrounds
\begin{equation}
 \Gamma_0:D^2\to\widetilde{\mathcal M}_0
 \label{eq:v36-geometric-disk}
\end{equation}
is \emph{spectrally admissible} if:
\begin{enumerate}[label=\textup{(S\arabic*)},leftmargin=3em]
\item on every fixed-index chamber the common geometric Dirac--Yukawa family
has a uniform reduced singular-value floor;
\item every Wilson-sign or accidental chiral crossing is contained in a finite
stratification satisfying the V26--V29 crossing-form hypotheses;
\item the background tensors and all derivatives required by the fixed source
packet are uniformly bounded on the compact disk;
\item the V2 soft continuum jet and V24 covariant frequency/Garding hypotheses
hold uniformly after discretization.
\end{enumerate}
\end{definition}

\begin{theorem}[Discretization of a spectrally admissible fibered filling]
\label{thm:v36-discretize-filling}
Let $\Gamma_0$ be a spectrally admissible geometric filling of a boundary loop
which is the continuum realization of a V2-consistent native loop.  Then there
exists $a_0>0$ such that for every $0<a<a_0$ the canonical V2 discretization of
$\Gamma_0$ defines a resolved native filling
\begin{equation}
 \Gamma_a:D^2\to\widetilde{\mathcal M}_a^{\rm prot}
 \label{eq:v36-discrete-filling}
\end{equation}
with the same boundary loop up to the V8--V24 controlled endpoint matching.
All declared wall signatures, determinant/Pfaffian exterior-power factors and
fixed source/BV derivatives agree with the geometric filling for sufficiently
small $a$.
\end{theorem}

\begin{proof}
Compactness of $D^2$ turns all geometric lower bounds and tensor derivative
bounds into uniform constants.  On each fixed-index chamber, V24 gives uniform
soft norm-resolvent convergence and high-block coercivity; hence the reduced
singular-value floor persists for small $a$.  V22 then identifies the soft
Riesz projectors.  Near a regular wall, V27--V29 express the crossing data by a
finite Riesz projector and an invertible crossing form.  Norm convergence of
the finite crossing block preserves its rank, crossing-form signature and
angular spectral split.  The finitely many strata are covered by finitely many
compact neighborhoods, so one common $a_0$ works.  Source/BV derivatives obey
the same conclusion by the source-complete V24 estimates.  Boundary endpoint
matching is the controlled direct rotation already used in V22--V24.
\end{proof}

\begin{corollary}[Exact triviality on spectrally fillable loops]
\label{cor:v36-spectrally-fillable-trivial}
For every sufficiently small lattice spacing, if a native loop admits a
spectrally admissible geometric disk filling, then its exact locally-flattened
relative phase satisfies
\begin{equation}
 \boxed{\rho_a(\gamma)=1.}
 \label{eq:v36-rho-spectral-fill}
\end{equation}
\end{corollary}

\begin{proof}
Theorem~\ref{thm:v36-discretize-filling} gives a fibered/native filling and
Proposition~\ref{prop:v36-fibered-kills} applies.
\end{proof}

\subsection{Relation to ordinary faithful-spin bordism}
\label{sec:v36-faithful-bordism}

The inherited Grand Tensor theorem proves that every closed ordinary-spin
$G_6$ five-manifold is zero in
$\Omega_5^{\mathrm{Spin}}(BG_6)$ and that the complete geometric
Standard-Model eta character is trivial.  V36 now states exactly what extra
input is required to transfer this statement to the native finite-cutoff
character.

\begin{definition}[Native bordism-lifting property]
\label{def:v36-bordism-lifting}
The resolved V2-consistent native configuration component has the
\emph{bordism-lifting property} if every closed native loop admits, after the
usual endpoint controlled identifications, a spectrally admissible geometric
fibered filling in the sense of
Definition~\ref{def:v36-spectral-filling}.
\end{definition}

\begin{theorem}[Faithful-spin triviality under native bordism lifting]
\label{thm:v36-faithful-under-lift}
If the ordinary-spin faithful $G_6$ native component has the bordism-lifting
property, then for all sufficiently small cutoffs
\begin{equation}
 \boxed{
 \rho_{a,\rm SM}\equiv1,
 \qquad
 \chi_{a,\rm SM}^{\eta}\equiv1}
 \label{eq:v36-exact-trivial-under-lift}
\end{equation}
after the exact local quasilocal flattening of
Theorem~\ref{thm:v36-exact-local-flattening}.  The same conclusion holds for
constructible loops containing regular V27--V29 Wilson walls.
\end{theorem}

\begin{proof}
Every loop has a spectrally admissible geometric filling by hypothesis.
Corollary~\ref{cor:v36-spectrally-fillable-trivial} makes the exact residual
character one.  V33 identifies that residual with the exact eta/integrability
character.  Wall loops are included because the filling theorem retains the
V27--V29 constructible line data.
\end{proof}

The role of the inherited bordism theorem is conceptual but precise: it says
there is no \emph{ordinary} faithful-spin obstruction to choosing a filling.
V36 shows that the remaining native question is whether such a filling can be
chosen in the stronger fibered, spectrally admissible class required by the
microscopic regulator.

\subsection{A no-go theorem for deriving exact triviality from existing inputs alone}
\label{sec:v36-no-go}

\begin{theorem}[Current Renewal hypotheses do not force exact global triviality]
\label{thm:v36-no-go}
The results inherited through V35 do not imply
$\rho_a\equiv1$ on an arbitrary noncontractible protected component.  More
precisely, they determine the exact relative connection up to the flat Cech
class \eqref{eq:v36-flat-class} and prove triviality on the normal subgroup of
fibered-fillable loops, but they contain no theorem forcing
\begin{equation}
 \mathfrak B_{5,a}^{\rm nat}=0
 \label{eq:v36-not-proved-defect-zero}
\end{equation}
or forcing $\bar\rho_a$ to vanish on a nonzero native defect group.
\end{theorem}

\begin{proof}
V35 controls every finite perturbative/source jet on star-shaped coefficient
charts, while V36 upgrades the exact local connection to a flat one by allowing
an exact quasilocal primitive.  Neither statement constructs a disk of native
backgrounds for a noncontractible loop.  Proposition
\ref{prop:v36-bordism-not-fibered} shows that ordinary bordism vanishing cannot
supply such a disk formally.  Hence the quotient
\eqref{eq:v36-native-bordism-defect} and its character remain logically
undetermined by the present hypotheses.
\end{proof}

\begin{firewall}[Exact local flattening does not re-enter the perturbative counterterm bank]
The V36 quasilocal primitive is a diagnostic device for the exact global line.
It is not inserted into the V1--V7 finite EFT renormalized action and therefore
does not change the already-proved coefficientwise continuum or its scheme
classification.  The perturbative theorem continues to use the finite V35
Chern--Simons/reference jets.
\end{firewall}

\subsection{V36 anti-circularity audit}
\label{sec:v36-audit}

\begin{enumerate}[leftmargin=2.4em]
\item \textbf{Ordinary bordism was not confused with configuration-space
contractibility.}  The difference between an arbitrary six-dimensional
filling and a disk of four-dimensional native backgrounds is explicit.
\item \textbf{The exact local curvature was not left mixed with the global
character.}  A radial quasilocal primitive removes it exactly on every
protected chart, leaving one constant Cech class.
\item \textbf{The exact quasilocal primitive was not promoted to an allowed
finite EFT counterterm.}  It is used only to expose the global phase; V35
remains the coefficientwise renormalization theorem.
\item \textbf{Wilson walls are retained.}  Fibered fillings may cross only the
resolved V26--V29 strata, with their exact determinant/Pfaffian line maps.
\item \textbf{The inherited $\Omega_5^{\mathrm{Spin}}(BG_6)=0$ theorem is used
only after the stronger native bordism-lifting property is stated.}
\item \textbf{A positive lifting theorem is proved where its hypotheses are
actually available.}  Compact spectrally admissible geometric disk families
discretize to native fillings by V24--V29 stability.
\end{enumerate}

\section{V36 status ledger}
\label{sec:v36-status}

\subsection*{Proved in V36}
\begin{enumerate}[leftmargin=2.2em]
\item On every protected star-shaped fixed-index chart, the exact
native--geometric relative determinant curvature is an exponentially local,
source-complete two-form and has an exact quasilocal radial primitive.
\item Subtracting that primitive makes the exact relative connection flat on
the chart.  On overlaps, local flat frames differ by constant $U(1)$
transitions defining a canonical exact Cech/flat-holonomy class.
\item Every loop admitting a resolved fibered/native disk filling has trivial
exact residual holonomy, including loops crossing regular V27--V29 wall
strata.
\item The exact global residue factors through the native fibered-bordism
defect group
$\mathfrak B_{5,a}^{\rm nat}$.
\item Ordinary vanishing of $\Omega_5^{\mathrm{Spin}}(BG_6)$ does not formally
imply vanishing of this native defect group; arbitrary six-dimensional bordism
is weaker than a fibered null-homotopy through four-dimensional backgrounds.
\item Every compact spectrally admissible geometric disk filling lifts, for all
sufficiently small lattice spacings, to a resolved native filling with the same
fixed-index chambers, crossing signatures, wall determinant/Pfaffian factors
and retained source/BV derivatives.
\item Consequently the exact residual eta/integrability character is one on
every spectrally fillable loop.
\item If the resolved faithful Standard-Model native component satisfies the
native bordism-lifting property, then the exact finite-cutoff faithful
character is globally trivial.
\end{enumerate}

\subsection*{Inherited}
\begin{enumerate}[leftmargin=2.2em]
\item V1--V7: source-complete coefficientwise ESM continuum, anomaly-free
common BV restoration, fixed-loop forest theorem, UV Cauchy convergence and
scheme/order projectivity.
\item V8--V24: regulator universality, fixed-index controlled modules, global
line transport, soft norm-resolvent anchoring and the V2-consistent covariant
Feshbach theorem.
\item V25--V29: native real/Pfaffian orientation, accidental divisors, Wilson
index jumps, higher-rank wall lines and nonsplit angular constructible lines.
\item V30--V34: faithful geometric cancellation mechanism, corrected
stratified/phase-faithful bridge, exact intrinsic eta character and local
five-dimensional infinite-wall realization.
\item V35: coefficientwise source-complete $5{+}1$ Chern--Simons descent and
proof that no new perturbative continuum datum is carried by the exact
five-dimensional phase.
\item Grand Tensor: the complete geometric ordinary-spin $G_6$ determinant
family has trivial eta character because
$\Omega_5^{\mathrm{Spin}}(BG_6)=0$.
\end{enumerate}

\subsection*{Conditional}
\begin{enumerate}[leftmargin=2.2em]
\item The exact local V36 primitive is quasilocal and finite-amplitude; it is
not asserted to belong to the original finite local EFT counterterm bank.
\item The geometric-to-native filling theorem requires a compact smooth
fibered disk with uniform V24 soft/Garding bounds and only regular V26--V29
strata.
\item Exact global triviality on the full native component remains conditional
on the native bordism-lifting property.
\item The independent V25 real/Pfaffian orientation hypotheses remain in force
on real wall sectors.
\item Positive IR/reference and V2-consistent smooth-background assumptions
remain unchanged.
\end{enumerate}

\subsection*{Open beyond V36}
\begin{enumerate}[leftmargin=2.2em]
\item Prove or falsify the native bordism-lifting property on the full
ordinary-spin faithful $G_6$ admissible component.  This is now the only exact
global complex-phase obstruction left by V36.
\item Compute the topology of
$\mathfrak B_{5,a}^{\rm nat}$ for the actual Renewal admissibility/gap
conditions.  If nonzero, determine the exact character $\bar\rho_a$.
\item Determine whether every ordinary faithful-spin nullbordism can be
modified to a spectrally admissible fibered filling after allowing the resolved
V26--V29 wall strata.  Ordinary bordism theory alone does not decide this.
\item The independent stronger problems remain: positive-IR removal, analytic
finite-coupling RG trajectories, perturbative summability and nonperturbative
ESM completion.
\end{enumerate}

\subsection*{Exact next load-bearing theorem}
\begin{center}
\fbox{\parbox{0.92\textwidth}{
\textbf{Native fibered-bordism lifting / defect-group theorem.}
The exact local phase has been flattened in V36 and the residual character now
lives on the native fibered-bordism defect group
$\mathfrak B_{5,a}^{\rm nat}$.  Determine this group for the actual V2
ordinary-spin faithful $G_6$ admissible configuration space, with resolved
V26--V29 wall strata.  Prove that every loop admits a spectrally admissible
fibered filling (hence $\mathfrak B_{5,a}^{\rm nat}=0$), or construct a
nontrivial defect class and compute its exact flat character.  A positive
lifting theorem, combined with
Theorem~\ref{thm:v36-faithful-under-lift}, would finally upgrade the inherited
ordinary-spin bordism vanishing to the exact finite-cutoff native equality
$\chi_{a,\rm SM}^{\eta}=1$ on the full protected component.}}
\end{center}

\section*{Append-only continuation marker: V36}
\addcontentsline{toc}{section}{Append-only continuation marker: V36}
\textbf{End of V36.}  The complete V1--V35 body above is retained verbatim.
V36 goes upstream from the bare use of ordinary spin bordism and separates
exact local phase curvature from the genuinely global native obstruction.  On
every protected star-shaped chart it constructs an exact exponentially local
radial primitive for the native--geometric relative curvature, reducing the
finite-cutoff residue to one flat Cech character.  It then distinguishes an
arbitrary six-dimensional spin $G_6$ nullbordism from the stronger fibered disk
of admissible four-dimensional native backgrounds required to kill that
character.  The residual phase factors through a precise native
fibered-bordism defect group.  Compact spectrally admissible geometric disk
families lift to native resolved fillings by the V24--V29 stability theory, so
all such loops have trivial exact phase.  The remaining stronger gate is to
compute or annihilate the native defect group on the full faithful admissible
component.



\clearpage
\section{V37 continuation: nontrivial native defect, electromagnetic winding, and the exact phase kernel}
\label{sec:v37-continuation}

V36 reduced the exact finite-cutoff global phase to a character of the native
fibered-bordism defect group
$\mathfrak B_{5,a}^{\rm nat}$ and left two logically possible outcomes: either
that group vanishes, or it is nonzero but the exact determinant character
vanishes on it.  The first possibility is already too strong on the periodic
native carrier.  The compact link variables retain genuine one-cycle topology
which an arbitrary six-dimensional bordism forgets.

The present continuation proves this directly and, at the same time, evaluates
a large explicit subgroup of the residual phase character.  The new results
are:
\begin{enumerate}[label=\textup{(V37.\arabic*)},leftmargin=3.8em]
\item the faithful Standard-Model group
$G_6=(SU(3)\times SU(2)\times U(1))/\mathbb Z_6$ has
$\pi_1(G_6)\cong\mathbb Z$, and the Higgs-vacuum stabilizer
$U(3)\hookrightarrow G_6$ induces an isomorphism on $\pi_1$;
\item every loop of native link variables carries an integer edge-winding
cochain; on a V2-consistent smooth plaquette chart this is a cellular
$1$-cocycle, well defined modulo vertex-gauge coboundaries;
\item the resulting cohomology class is annihilated by every fibered/native
filling and therefore descends to a quotient of the V36 defect group;
\item on a periodic four-torus, explicit flat Higgs-vacuum loops realize all
four harmonic winding generators, proving that the native defect group is
nonzero and that the V36 universal bordism-lifting property fails on the full
periodic component;
\item exact finite gauge invariance kills the vertex-gauge coboundary subgroup;
\item the complete Standard-Model determinant/Pfaffian phase is exactly one on
the harmonic electromagnetic winding subgroup because the charged packet is
vector-like under $U(1)_{\rm em}$ and the neutral real packet has the V25
orientation;
\item consequently the only possible exact faithful phase left after V37 lies
in the zero-winding defect subgroup.  The next load-bearing question is the
non-Abelian/global-anomaly content of that subgroup, not the vanishing of
$\mathfrak B_{5,a}^{\rm nat}$ itself.
\end{enumerate}

Thus V37 changes the exact global target: one should compute the character of a
nonzero native defect group, rather than try to prove that the group itself is
zero.

\subsection{The fundamental group of the faithful Standard-Model gauge group}
\label{sec:v37-pi1-G6}

Let
\begin{equation}
 \widetilde G
 :=SU(3)\times SU(2)\times\mathbb R
 \label{eq:v37-universal-cover}
\end{equation}
with covering map
\begin{equation}
 \pi(g_3,g_2,t)
 =\bigl[g_3,g_2,e^{it}\bigr]\in G_6.
 \label{eq:v37-cover-map}
\end{equation}
Choose the conventional generator of the diagonal $\mathbb Z_6$ quotient,
\begin{equation}
 z_6=
 \left(e^{2\pi i/3}I_3,-I_2,e^{i\pi/3}\right).
 \label{eq:v37-z6}
\end{equation}

\begin{theorem}[Fundamental group of the faithful quotient]
\label{thm:v37-pi1-G6}
The kernel of \eqref{eq:v37-cover-map} is the infinite cyclic group generated
by
\begin{equation}
 \widetilde z
 =\left(e^{2\pi i/3}I_3,-I_2,\pi/3\right).
 \label{eq:v37-ztilde}
\end{equation}
Consequently
\begin{equation}
 \boxed{\pi_1(G_6)\cong\mathbb Z.}
 \label{eq:v37-pi1-G6}
\end{equation}
\end{theorem}

\begin{proof}
Both $SU(3)$ and $SU(2)$ are simply connected, so
$SU(3)\times SU(2)\times\mathbb R$ is simply connected.  An element lies in the
kernel of \eqref{eq:v37-cover-map} exactly when its projection to
$SU(3)\times SU(2)\times U(1)$ is one of the six elements generated by
$z_6$.  Thus every kernel element has the form $\widetilde z^{\,n}$ for one
$n\in\mathbb Z$.  Since
$\widetilde z^{\,6}=(I_3,I_2,2\pi)$, this also contains the ordinary kernel of
$\mathbb R\to U(1)$.  Hence the displayed cyclic group is the full deck group
and therefore the fundamental group of $G_6$.
\end{proof}

The Grand Tensor Higgs vacuum has stabilizer $U(3)\subset G_6$.  The following
fact identifies which loop in that stabilizer represents the generator above.

\begin{theorem}[The Higgs stabilizer carries the faithful fundamental loop]
\label{thm:v37-U3-pi1}
Let $h_0$ be a unit Higgs vacuum and let
$G_{h_0}\cong U(3)$ be its stabilizer.  The orbit of $h_0$ is the unit sphere
$S^3\subset\mathbb C^2$, so there is a fibration
\begin{equation}
 U(3)\longrightarrow G_6\longrightarrow S^3.
 \label{eq:v37-Higgs-fibration}
\end{equation}
The induced map
\begin{equation}
 \boxed{
 \pi_1(U(3))\xrightarrow{\ \simeq\ }\pi_1(G_6)
 \cong\mathbb Z}
 \label{eq:v37-U3-pi1-iso}
\end{equation}
is an isomorphism.  In particular one may choose a loop
$u_{\rm em}:S^1\to G_{h_0}$ generating $\pi_1(G_6)$ and fixing $h_0$
pointwise.
\end{theorem}

\begin{proof}
The Higgs representation is transitive on the unit sphere and its stabilizer
at $h_0$ is the inherited $U(3)$ subgroup, giving
\eqref{eq:v37-Higgs-fibration}.  The long exact homotopy sequence contains
\[
 \pi_2(S^3)\longrightarrow\pi_1(U(3))
 \longrightarrow\pi_1(G_6)\longrightarrow\pi_1(S^3).
\]
Both outer groups vanish, so the middle map is an isomorphism.  The generator
may therefore be represented inside the unbroken electromagnetic/color
stabilizer and, by definition of the stabilizer, fixes $h_0$.
\end{proof}

We use the notation $U(1)_{\rm em}$ for one such one-parameter subgroup
representing this fundamental loop.  Only its homotopy class and the fact that
it fixes the Higgs vacuum are load bearing below.

\subsection{Edge winding and the smooth V2 cocycle}
\label{sec:v37-edge-winding}

Let $X_a$ be one periodic finite spacetime cell complex with oriented edge set
$E_a$, vertex set $V_a$ and plaquette set $F_a$.  Let
\begin{equation}
 U_e(t)\in G_6,
 \qquad 0\le t\le1,
 \qquad U_e(0)=U_e(1)
 \label{eq:v37-link-loop}
\end{equation}
be the link component of a closed native background loop $\gamma$.  Left or
right multiplication by the fixed base value identifies this loop with a
based loop in $G_6$, so Theorem~\ref{thm:v37-pi1-G6} assigns an integer
\begin{equation}
 w_e(\gamma)\in\mathbb Z.
 \label{eq:v37-edge-winding}
\end{equation}
Collect these integers into the cellular $1$-cochain
\begin{equation}
 w(\gamma)\in C^1(X_a;\mathbb Z).
 \label{eq:v37-w-cochain}
\end{equation}

\begin{lemma}[Plaquette-small loops have closed winding]
\label{lem:v37-w-closed}
Suppose the loop is V2-consistent in the sense that every plaquette holonomy
loop stays in one contractible principal-logarithm neighborhood of the
identity.  Then
\begin{equation}
 \boxed{d_{\rm cell}w(\gamma)=0.}
 \label{eq:v37-w-closed}
\end{equation}
Hence $w(\gamma)$ defines a class
\begin{equation}
 [w(\gamma)]\in H^1(X_a;\mathbb Z).
 \label{eq:v37-w-H1}
\end{equation}
\end{lemma}

\begin{proof}
The fundamental group of a topological group is abelian and multiplication of
loops adds their $\pi_1$ classes.  For an oriented plaquette $p$, the winding
of the loop $U_{\partial p}(t)$ is therefore the oriented sum of the edge
windings around $p$, namely $(d_{\rm cell}w)(p)$.  By hypothesis
$U_{\partial p}(t)$ stays in a contractible logarithm chart, so its winding is
zero.  This holds for every plaquette.
\end{proof}

\begin{lemma}[Vertex gauge loops are cellular coboundaries]
\label{lem:v37-gauge-coboundary}
Let $g_v(t)$ be closed loops in $G_6$ at the vertices with winding integers
$n_v\in\mathbb Z$, and act on links by
\begin{equation}
 U_e(t)\longmapsto
 g_{t(e)}(t)U_e(t)g_{s(e)}(t)^{-1}.
 \label{eq:v37-gauge-loop-action}
\end{equation}
Then
\begin{equation}
 \boxed{
 w_e\longmapsto w_e+n_{t(e)}-n_{s(e)}.}
 \label{eq:v37-w-gauge-shift}
\end{equation}
Thus the cohomology class \eqref{eq:v37-w-H1} is invariant under closed
vertex-gauge loops.
\end{lemma}

\begin{proof}
Apply additivity of $\pi_1(G_6)$ to the three loop factors in
\eqref{eq:v37-gauge-loop-action}.  Inversion changes the sign of the winding.
The shift is exactly the cellular coboundary of the vertex $0$-cochain
$n=(n_v)$.
\end{proof}

\begin{definition}[Smooth native winding shadow]
\label{def:v37-winding-shadow}
On the V2-consistent resolved loop space define
\begin{equation}
 \boxed{
 \mathfrak w_a([\gamma])
 :=[w(\gamma)]\in H^1(X_a;\mathbb Z).}
 \label{eq:v37-winding-shadow}
\end{equation}
\end{definition}

\begin{proposition}[Fibered fillings have zero winding shadow]
\label{prop:v37-fill-zero-winding}
If a V2-consistent resolved loop admits a fibered/native filling in the sense
of V36, then
\begin{equation}
 \boxed{\mathfrak w_a([\gamma])=0.}
 \label{eq:v37-fill-w-zero}
\end{equation}
Consequently the winding shadow descends to a homomorphism on the corresponding
V2-consistent part of the native defect group,
\begin{equation}
 \boxed{
 \overline{\mathfrak w}_a:
 \mathfrak B_{5,a}^{\rm nat,V2}
 \longrightarrow H^1(X_a;\mathbb Z).}
 \label{eq:v37-wbar}
\end{equation}
\end{proposition}

\begin{proof}
Compose a fibered filling with the projection onto any fixed edge-link factor.
This extends the boundary loop $U_e:S^1\to G_6$ to a disk in $G_6$, so its
winding integer is zero.  Thus every edge winding vanishes for a genuinely
closed raw native filling.  If the filling is described in overlapping gauge
charts, the changes are vertex-gauge coboundaries by
Lemma~\ref{lem:v37-gauge-coboundary}; therefore the gauge-invariant class
$[w]$ still vanishes.  Hence the normal subgroup generated by fibered-fillable
loops lies in the kernel of $\mathfrak w_a$, and the quotient map follows.
\end{proof}

\subsection{Explicit nonfillable harmonic loops on the periodic carrier}
\label{sec:v37-harmonic-loops}

Assume now that $X_a$ is the standard periodic cellulation of a four-torus.
For each coordinate direction $\mu$, choose a dual three-torus cut
$\mathcal C_\mu$ consisting of the positively oriented $\mu$-edges which cross
one fixed periodic boundary.  Let
$u(t)=u_{\rm em}(t)$ be the generator of
Theorem~\ref{thm:v37-U3-pi1}, parametrized so that $u(0)=u(1)=I$.
For $n\in\mathbb Z$ define a link loop
\begin{equation}
 U_e^{(\mu,n)}(t)
 =
 \begin{cases}
 u(t)^n,&e\in\mathcal C_\mu,\\
 I,&e\notin\mathcal C_\mu.
 \end{cases}
 \label{eq:v37-harmonic-link-loop}
\end{equation}
Keep the coframe fixed, put $H_v=v_Hh_0$, and set every other background writer
at its reference value.

\begin{lemma}[The harmonic loop is classically flat and Higgs compatible]
\label{lem:v37-harmonic-flat}
For every $t$, every plaquette holonomy of
\eqref{eq:v37-harmonic-link-loop} is the identity and every covariant Higgs
edge difference vanishes.  Hence the gauge and Higgs-gradient pieces of the
native classical action stay at their vacuum values along the whole loop.
\end{lemma}

\begin{proof}
Every plaquette crossing the chosen dual cut contains two cut edges with
opposite induced orientations, so the two factors $u(t)^n$ and $u(t)^{-n}$
cancel; all other plaquettes contain no cut edge.  Since $u(t)$ lies in the
Higgs stabilizer, $u(t)h_0=h_0$, and therefore every covariant Higgs difference
also vanishes.
\end{proof}

The fermion spectrum may encounter Wilson or chiral zero crossings as the flat
holonomy varies.  Those crossings are precisely the soft strata already
resolved in V26--V29.  A finite-dimensional generic perturbation, arbitrarily
small within the same winding class, makes the crossings Whitney-regular; the
resolved homotopy class and the integer winding are unchanged.  We henceforth
understand \eqref{eq:v37-harmonic-link-loop} in this resolved sense.

\begin{theorem}[The native fibered-bordism defect group is nonzero]
\label{thm:v37-defect-nonzero}
For the periodic four-torus carrier,
\begin{equation}
 H^1(X_a;\mathbb Z)\cong\mathbb Z^4.
 \label{eq:v37-H1-T4}
\end{equation}
The four resolved loops
$\gamma_\mu:=\gamma_{\mu,1}$ defined by
\eqref{eq:v37-harmonic-link-loop} satisfy
\begin{equation}
 \boxed{
 \mathfrak w_a(\gamma_\mu)=e_\mu,}
 \label{eq:v37-w-basis}
\end{equation}
where $(e_\mu)$ is the standard basis of $H^1(T^4;\mathbb Z)$.  Therefore
\begin{equation}
 \boxed{
 \overline{\mathfrak w}_a:
 \mathfrak B_{5,a}^{\rm nat,V2}	woheadrightarrow\mathbb Z^4}
 \label{eq:v37-w-surjective}
\end{equation}
is surjective.  In particular
\begin{equation}
 \boxed{\mathfrak B_{5,a}^{\rm nat,V2}\ne0,}
 \label{eq:v37-defect-nonzero}
\end{equation}
and the V36 bordism-lifting property is false on the full periodic native
component.
\end{theorem}

\begin{proof}
The edge winding of \eqref{eq:v37-harmonic-link-loop} is the integer cocycle
Poincare dual to the chosen dual cut.  Its cohomology class is the corresponding
basis vector of $H^1(T^4;\mathbb Z)$.  Products of the four commuting
$U(1)_{\rm em}$ loops realize arbitrary integer combinations, so the winding
shadow is onto.  Proposition~\ref{prop:v37-fill-zero-winding} shows that a loop
with nonzero image cannot be fibered-fillable.  Hence the defect group is
nonzero and no property asserting that every loop has a native filling can hold
on this component.
\end{proof}

\begin{firewall}[V36 is corrected, not contradicted]
V36 proved that fibered-fillable loops have trivial exact phase and gave a
positive discretization theorem for spectrally admissible geometric fillings.
Those results remain valid.  V37 shows only that the additional universal
hypothesis ``every loop has such a filling'' is false on the periodic native
carrier.  Exact global triviality, if true, must therefore come from vanishing
of the character on a nonzero defect group rather than from vanishing of the
defect group itself.
\end{firewall}

\subsection{Exact gauge-orbit loops lie in the phase kernel}
\label{sec:v37-gauge-kernel}

The nonzero defect group already contains many loops which are physically pure
gauge.  Their exact phase is fixed by the finite gauge covariance inherited in
V3 and the Grand Tensor determinant section.

\begin{theorem}[Closed finite gauge loops have unit exact phase]
\label{thm:v37-gauge-loop-phase}
Let $g_v(t)$ be arbitrary closed vertex-gauge loops and let
$\gamma_g(t)=g(t)\cdot b_0$ be the induced loop through one protected native
background, resolved through any regular wall crossings.  For the complete
anomaly-free Standard-Model determinant packet,
\begin{equation}
 \boxed{\rho_{a,\rm SM}(\gamma_g)=1.}
 \label{eq:v37-gauge-phase-one}
\end{equation}
Its edge-winding cochain is the cellular coboundary
$w=dn$ of Lemma~\ref{lem:v37-gauge-coboundary}.
\end{theorem}

\begin{proof}
Finite gauge covariance gives
$D_{\chi,t}=R(g_t)D_{\chi,0}R(g_t)^{-1}$ on the complete packet, with the
corresponding action on the kernel/cokernel and determinant lines.  The
inherited complete Standard-Model determinant section is gauge invariant after
the common V3 restoring choice, including all line/source descendants.  Its
parallel transport along the gauge orbit is therefore the gauge action on that
same section.  Since the gauge transformation is closed,
$g(1)=g(0)=I$, the endpoint action is the identity and the holonomy is one.
Resolved Wilson or accidental crossings do not change the conclusion because
V26--V29 wall maps are the exact determinant functor of the same gauge-covariant
finite complexes.
\end{proof}

Thus the exact phase character already vanishes on a nontrivial subgroup of the
native defect group which ordinary bordism theory does not identify.

\subsection{Exact electromagnetic vector-pair triviality}
\label{sec:v37-em-phase}

The harmonic winding loops are not pure vertex-gauge loops.  Nevertheless, on
the Higgs vacuum they lie entirely in the unbroken electromagnetic subgroup,
where the Standard-Model charged fermion packet is vector-like.

We first record the finite overlap positivity statement used below.

\begin{lemma}[Positive phase of a massive vector-like overlap determinant]
\label{lem:v37-vectorlike-positive}
Let $D_0$ be a finite overlap Dirac operator satisfying
\begin{equation}
 D_0^\dagger=\gamma_5D_0\gamma_5
 \label{eq:v37-gamma5-Herm}
\end{equation}
and the usual Ginsparg--Wilson spectral circle, whose real spectral points are
$0$ and $2/a$.  For
\begin{equation}
 D_m=m+\left(1-\frac{am}{2}\right)D_0,
 \qquad 0<m<2/a,
 \label{eq:v37-massive-overlap}
\end{equation}
all nonreal eigenvalues occur in complex-conjugate pairs and the two possible
real eigenvalues of $D_m$ are $m$ and $2/a$, both positive.  Hence
\begin{equation}
 \boxed{\det D_m>0}
 \label{eq:v37-vectorlike-det-positive}
\end{equation}
whenever $D_m$ is finite-dimensional.  Its determinant line therefore has a
canonical positive real frame and trivial phase holonomy along every closed
background loop.
\end{lemma}

\begin{proof}
Equation~\eqref{eq:v37-gamma5-Herm} implies that the characteristic polynomial
has conjugate symmetry.  The Ginsparg--Wilson circle has only the real points
$0$ and $2/a$.  The affine map
$\lambda\mapsto m+(1-am/2)\lambda$ sends these to $m$ and $2/a$ respectively.
Every nonreal pair contributes the positive number $|m+(1-am/2)\lambda|^2$ to
the determinant, while each real mode contributes a positive factor.
\end{proof}

\begin{theorem}[The complete electromagnetic winding subgroup has trivial exact phase]
\label{thm:v37-em-trivial}
Assume the Higgs-vacuum branch with nonzero charged-fermion Yukawa singular
values in the reference packet.  Restrict the complete Standard-Model native
fermion family to backgrounds with links in the unbroken
$U(1)_{\rm em}\subset G_{h_0}$, fixed Higgs vacuum and fixed coframe.  Then the
charged chiral packet reorganizes exactly into vector-like Dirac sectors, one
for every electric charge and colour/flavour multiplicity; the neutral sector
is independent of the electromagnetic links.  Consequently
\begin{equation}
 \boxed{
 \rho_{a,\rm SM}(\gamma_{\mu,n})=1
 \qquad
 (\mu=1,\ldots,4,\ n\in\mathbb Z).}
 \label{eq:v37-em-rho-one}
\end{equation}
The same statement holds with an independently physical neutral Majorana block
when it carries the V25 native orientation.
\end{theorem}

\begin{proof}
On the Higgs vacuum, electric charge is unbroken and every charged left/right
Standard-Model pair has the same electromagnetic charge.  Equivalently, in a
left-handed-only convention the right-handed field is represented by its
charge-conjugate Weyl field, so the two chiral determinant factors are dual.
The Yukawa matrices preserve electric charge; constant flavour biunitary
rotations therefore reduce each charged mass block to positive singular values
without changing the electromagnetic links.  Each charged pair is then a
finite massive vector-like overlap determinant of the form in
Lemma~\ref{lem:v37-vectorlike-positive}, possibly with colour/flavour
multiplicity.  Its phase is identically one.  Multiplying all charged sectors
preserves this property.

The neutrino sector is electrically neutral and hence independent of
$\gamma_{\mu,n}$.  If a separate real Majorana Pfaffian block is present, V25
supplies its global positive orientation on the neutral branch, so it also
contributes no electromagnetic phase.  The complete determinant/Pfaffian
holonomy is therefore one.
\end{proof}

\begin{corollary}[A split nonzero defect subgroup lies in the exact phase kernel]
\label{cor:v37-Z4-kernel}
Let $\mathfrak H_a\subset\mathfrak B_{5,a}^{\rm nat,V2}$ be the subgroup
generated by the four commuting harmonic loops $\gamma_\mu$.  Then
\begin{equation}
 \boxed{
 \mathfrak H_a\cong\mathbb Z^4,
 \qquad
 \bar\rho_{a,\rm SM}|_{\mathfrak H_a}=1.}
 \label{eq:v37-Z4-kernel}
\end{equation}
Moreover $\overline{\mathfrak w}_a|_{\mathfrak H_a}$ is the identity
isomorphism onto $H^1(T^4;\mathbb Z)$.
\end{corollary}

\begin{proof}
The winding-shadow theorem shows that the four generators are independent and
that their subgroup maps isomorphically onto $H^1(T^4;\mathbb Z)$.  The phase
statement is Theorem~\ref{thm:v37-em-trivial} and multiplicativity of the flat
character.
\end{proof}

Thus the exact faithful character can be nontrivial only on defect classes
whose smooth winding shadow vanishes after removal of the electromagnetic
harmonic subgroup and the exact gauge-orbit subgroup.

\subsection{The zero-winding defect is the new exact global target}
\label{sec:v37-zero-winding}

Define
\begin{equation}
 \boxed{
 \mathfrak B_{5,a}^{0}
 :=\ker\overline{\mathfrak w}_a
 \subset\mathfrak B_{5,a}^{\rm nat,V2}.}
 \label{eq:v37-zero-winding-defect}
\end{equation}
Because the harmonic subgroup gives a section of
$\overline{\mathfrak w}_a$, the V2-consistent defect group contains a split
extension
\begin{equation}
 1\longrightarrow\mathfrak B_{5,a}^{0}
 \longrightarrow\mathfrak B_{5,a}^{\rm nat,V2}
 \longrightarrow\mathbb Z^4
 \longrightarrow1.
 \label{eq:v37-defect-extension}
\end{equation}
The exact faithful character is trivial on the displayed $\mathbb Z^4$ section
and on the subgroup generated by finite gauge-orbit loops.  Therefore all
remaining exact phase information is carried by the zero-winding subgroup.

\begin{theorem}[Corrected exact-global reduction]
\label{thm:v37-corrected-global-reduction}
On the periodic V2-consistent faithful Standard-Model component, exact finite
cutoff triviality of the determinant/Pfaffian character is equivalent to
triviality of its restriction to the zero-winding defect subgroup:
\begin{equation}
 \boxed{
 \bar\rho_{a,\rm SM}\equiv1
 \quad\Longleftrightarrow\quad
 \bar\rho_{a,\rm SM}|_{\mathfrak B_{5,a}^{0}}\equiv1.}
 \label{eq:v37-rho-reduction}
\end{equation}
Here the equality is understood after quotienting the already-trivial finite
gauge-orbit subgroup.
\end{theorem}

\begin{proof}
The forward implication is immediate.  Conversely, every defect class can be
written as a zero-winding class times an element of the harmonic
$\mathbb Z^4$ section, modulo the finite gauge-orbit subgroup, using the split
extension \eqref{eq:v37-defect-extension}.  The exact character is one on the
latter two factors by Theorems~\ref{thm:v37-gauge-loop-phase} and
\ref{thm:v37-em-trivial}; by hypothesis it is one on the zero-winding factor.
Multiplicativity gives the result.
\end{proof}

This theorem replaces the false goal
$\mathfrak B_{5,a}^{\rm nat}=0$ by the smaller and physically relevant problem
of calculating one exact character on the zero-winding sector.

\subsection{What can live in the zero-winding sector}
\label{sec:v37-zero-winding-content}

The edge-winding shadow detects only the $\pi_1(G_6)$ part of the native link
variables.  It cannot see:
\begin{enumerate}[label=\textup{(Z\arabic*)},leftmargin=3em]
\item non-Abelian large-gauge topology with zero $\pi_1(G_6)$ winding,
including the four-dimensional $SU(2)$ mod-two class;
\item V29 angular spectral-bundle loops whose underlying link winding is zero;
\item loops generated by resolved higher-codimension Wilson/chiral strata;
\item any purely microscopic local-operator loop in the V16 continuum-holonomy
kernel.
\end{enumerate}
These are precisely the classes on which the inherited even-$SU(2)$-doublet
parity and the faithful-spin eta theorem can carry information not visible to
V37's $H^1$ invariant.

\begin{firewall}[Nonzero defect group does not mean a global anomaly]
Theorem~\ref{thm:v37-defect-nonzero} proves that the native fibered-bordism
defect group is nonzero.  Corollary~\ref{cor:v37-Z4-kernel} simultaneously
proves that a nonzero subgroup has exact phase one.  Therefore non-fillability
through native four-dimensional backgrounds is not itself an anomaly.  The
physical question is the character $\bar\rho_a$, not whether the defect group
vanishes.
\end{firewall}

\subsection{V37 anti-circularity audit}
\label{sec:v37-audit}

\begin{enumerate}[leftmargin=2.4em]
\item \textbf{The V36 bordism-lifting property was tested rather than assumed.}
An explicit $H^1$ winding obstruction shows that it fails on the periodic
native component.
\item \textbf{The link winding is native and finite-cutoff.}
It is defined directly from $\pi_1(G_6)$ on each compact link variable and does
not use a continuum replacement of the regulator.
\item \textbf{Gauge dependence is removed explicitly.}
Vertex-gauge loops shift the edge winding by cellular coboundaries, so the
smooth V2 shadow lies in $H^1(X_a;\mathbb Z)$.
\item \textbf{The nonfillable harmonic loops remain in the physical Higgs
branch.}  Their links lie in the Higgs stabilizer, every plaquette is trivial
and every covariant Higgs difference vanishes.
\item \textbf{Phase triviality is not inferred from non-fillability.}
It is proved separately from exact gauge invariance and the vector-like
$U(1)_{\rm em}$ decomposition of the complete Standard-Model fermion packet.
\item \textbf{The zero-winding sector is retained.}
V37 does not use the $H^1$ calculation to dismiss non-Abelian or angular global
anomalies which the winding shadow cannot detect.
\end{enumerate}

\section{V37 status ledger}
\label{sec:v37-status}

\subsection*{Proved in V37}
\begin{enumerate}[leftmargin=2.2em]
\item The faithful Standard-Model gauge group has
$\pi_1(G_6)\cong\mathbb Z$, generated by the diagonal quotient lift.  The Higgs
stabilizer $U(3)$ induces an isomorphism on $\pi_1$, so the generator may be
represented in the unbroken electromagnetic stabilizer.
\item Every closed native link loop has an integer edge-winding cochain.  On a
V2-consistent principal-plaquette chart the cochain is closed, and closed
vertex-gauge loops change it only by cellular coboundaries.  Hence there is a
native smooth winding shadow in $H^1(X_a;\mathbb Z)$.
\item Every fibered/native filling has zero winding shadow, so the winding map
descends to the V36 native defect group.
\item On the periodic four-torus, four explicit flat Higgs-vacuum
$U(1)_{\rm em}$ loops realize a basis of $H^1(T^4;\mathbb Z)\cong\mathbb Z^4$.
Therefore the V2-consistent native defect group is nonzero and the V36 universal
bordism-lifting property is false on the full periodic component.
\item Exact complete-packet finite gauge invariance makes the residual flat
phase one on all closed vertex-gauge loops, including those with nonzero edge
coboundary winding.
\item On the Higgs vacuum the complete charged Standard-Model packet is
vector-like under $U(1)_{\rm em}$.  The finite massive overlap determinants
have positive phase, while the neutral sector is electromagnetic-blind and an
independent real Majorana factor has the V25 orientation.  Hence the exact
phase is one on the full harmonic $\mathbb Z^4$ winding subgroup.
\item Exact global phase triviality is therefore reduced to the zero-winding
defect subgroup $\mathfrak B_{5,a}^{0}$ rather than to vanishing of the full
defect group.
\end{enumerate}

\subsection*{Inherited}
\begin{enumerate}[leftmargin=2.2em]
\item V1--V7: source-complete coefficientwise ESM continuum, common BV
restoration, anomaly-free Standard-Model packet, fixed-loop forests, ultraviolet
Cauchy convergence and scheme/order projectivity.
\item V8--V24: regulator universality, controlled zero-mode modules, global
line transport, soft norm-resolvent anchoring and the V2-consistent covariant
Feshbach theorem.
\item V25--V29: real/Pfaffian orientation, accidental divisors, Wilson index
jumps, higher-rank wall lines and nonsplit angular constructible superlines.
\item V30--V35: faithful geometric cancellation mechanism, phase-faithful
bridge corrections, exact intrinsic eta character, local 5D infinite-wall
realization and coefficientwise $5{+}1$ Chern--Simons descent.
\item V36: exact quasilocal local flattening, the residual flat character and
the native fibered-bordism defect group, together with positive lifting for
spectrally admissible geometric disk families.
\end{enumerate}

\subsection*{Conditional}
\begin{enumerate}[leftmargin=2.2em]
\item The harmonic-loop construction is stated on the periodic Higgs-vacuum
branch.  Other spacetime topology replaces $\mathbb Z^4$ by the corresponding
$H^1(X_a;\mathbb Z)$ image.
\item The smooth winding class requires the V2 principal-plaquette condition;
rough link loops outside that class can carry additional edge winding not
captured by cellular cohomology.
\item The exact electromagnetic positivity statement assumes the charged
Yukawa singular values in the chosen Higgs-vacuum packet are nonzero; a massless
charged sector should instead be treated with the V26 saturated divisor packet.
\item The generic resolved representative of a harmonic holonomy loop uses the
V26--V29 finite crossing stratification.  No claim is made across an
unresolved nonregular infinite-multiplicity degeneration.
\item Positive IR/reference and all V2-consistent smooth-background assumptions
remain unchanged.
\end{enumerate}

\subsection*{Open beyond V37}
\begin{enumerate}[leftmargin=2.2em]
\item Compute the zero-winding defect subgroup
$\mathfrak B_{5,a}^{0}$ and the restriction of the exact flat character to it.
This is now the only global complex-phase problem on the periodic V2-consistent
branch.
\item Construct and evaluate the native representative of the four-dimensional
$SU(2)$ mod-two/Witten class inside $\mathfrak B_{5,a}^{0}$.  The inherited even
number of Standard-Model doublets strongly suggests a trivial complete-packet
phase, but the exact native loop calculation is not yet in this lineage.
\item Determine the contribution of V29 angular spectral-bundle loops with zero
edge winding and their relation to the faithful complete-packet eta character.
\item The independent stronger problems remain: positive-IR removal, analytic
finite-coupling RG trajectories, perturbative summability and nonperturbative
ESM completion.
\end{enumerate}

\subsection*{Exact next load-bearing theorem}
\begin{center}
\fbox{\parbox{0.92\textwidth}{
\textbf{Zero-winding non-Abelian defect / native Witten-generator theorem.}
V37 proves that the native fibered-bordism defect group is nonzero but that its
entire electromagnetic $H^1$ winding subgroup and finite gauge-coboundary
subgroup lie in the kernel of the exact faithful phase.  The remaining exact
character lives on
$\mathfrak B_{5,a}^{0}=\ker\overline{\mathfrak w}_a$.  Construct the native
resolved loop corresponding to the nontrivial four-dimensional $SU(2)$
large-gauge class, compute its finite overlap/Pfaffian determinant-line
transport, and prove that the complete Standard-Model phase is
$(-1)^{N_{\rm dbl}}=1$ from the inherited even weak-doublet parity.  Then show
that all remaining smooth gauge-transform classes with zero edge winding are
generated by this mod-two class together with fibered-fillable loops, or
identify the next non-Abelian/angular generator.}}
\end{center}

\section*{Append-only continuation marker: V37}
\addcontentsline{toc}{section}{Append-only continuation marker: V37}
\textbf{End of V37.}  The complete V1--V36 body above is retained verbatim.
V37 goes upstream from the attempted proof that the native fibered-bordism
defect group vanishes.  The faithful link group has nontrivial fundamental
group, and on the periodic Higgs-vacuum branch this produces four explicit
non-fillable smooth electromagnetic winding loops.  Their cellular winding
classes give a surjection from the V2-consistent native defect group onto
$H^1(T^4;\mathbb Z)\cong\mathbb Z^4$, so the V36 universal bordism-lifting
property is false.  Nevertheless exact finite gauge invariance kills the gauge
coboundary subgroup and the vector-like electromagnetic decomposition of the
complete Standard-Model packet makes the exact determinant/Pfaffian phase one
on the whole harmonic $\mathbb Z^4$ subgroup.  The exact global problem is
therefore reduced to the zero-winding defect subgroup, whose first expected
non-Abelian generator is the native $SU(2)$ mod-two/Witten loop.



\clearpage
\section{V38 continuation: the native $SU(2)$ Witten generator and the zero-winding gauge-transform 4-type}
\label{sec:v38-continuation}

V37 reduced the exact finite-cutoff faithful phase problem on the periodic
V2-consistent component to the zero-edge-winding defect subgroup
$\mathfrak B_{5,a}^{0}$.  The first expected non-Abelian generator is the
four-dimensional $SU(2)$ mod-two/Witten class.  This continuation constructs
that class directly in the V2 overlap regulator, matches the native Kato twist
to the established finite-overlap global-anomaly theorem, and proves exact
cancellation in the complete Standard-Model packet from the inherited even
weak-doublet parity.

There is one further upstream point.  Once the $H^1$ winding of V37 is removed,
the smooth large-gauge transformation space still contains free $\pi_3$
sectors in addition to the Witten $\mathbb Z_2$ class.  V38 identifies the
complete four-dimensional Postnikov shadow explicitly.  Thus cancellation of
the Witten generator is strong but does not yet classify the whole
zero-winding defect group.

Throughout this continuation, the spacetime carrier is the periodic four-torus
$X_a$ at sufficiently small lattice spacing, the spin/coframe background is
held fixed, and the gauge links are restricted to the smooth V2-consistent
admissible component.  The only varying gauge factor in the Witten construction
is the weak $SU(2)$ subgroup of the faithful group
\begin{equation}
 G_6=\frac{SU(3)\times SU(2)\times U(1)}{\mathbb Z_6}.
 \label{eq:v38-G6}
\end{equation}
The discrete quotient does not change homotopy groups in degree at least two.

\subsection{A native representative of the nontrivial $\pi_4(SU(2))$ class}
\label{sec:v38-native-loop}

Choose an embedded closed four-ball $B^4\Subset T^4$.  Collapsing
$T^4\setminus\operatorname{int}B^4$ to one point gives a degree-one quotient
map
\begin{equation}
 q:T^4\longrightarrow S^4.
 \label{eq:v38-collapse}
\end{equation}
Let
\begin{equation}
 g_*:S^4\longrightarrow SU(2)
 \label{eq:v38-generator-map}
\end{equation}
represent the nonzero element of
$\pi_4(SU(2))\cong\mathbb Z_2$, and choose a smooth based representative.  Put
\begin{equation}
 g:=g_*\circ q:T^4\longrightarrow SU(2)\subset G_6.
 \label{eq:v38-g}
\end{equation}
Thus $g=I$ outside a compact subset of $B^4$.

For each sufficiently fine lattice, sample $g$ on the lattice sites and denote
the resulting vertex gauge transformation by $g_a(x)$.  The endpoint pure-gauge
link field is
\begin{equation}
 U_a^g(x,\mu)
 :=g_a(x)g_a(x+a\hat\mu)^{-1}.
 \label{eq:v38-pure-gauge-endpoint}
\end{equation}
Choose any smooth native path of weak links
\begin{equation}
 \Gamma_a:t\longmapsto U_{a,t},
 \qquad
 U_{a,0}=I,
 \quad
 U_{a,1}=U_a^g,
 \label{eq:v38-Gamma}
\end{equation}
which is a discretization of a fixed smooth continuum interpolation and remains
inside the V2 Wilson-gap/admissibility chart for sufficiently small $a$.
Other gauge, coframe and source backgrounds are transported by any smooth
endpoint-compatible path inside their protected chart.  The endpoint is closed
in the physical background quotient by the gauge identification $g_a$.

\begin{definition}[Native Witten loop]
\label{def:v38-Witten-loop}
The resulting loop in the resolved gauge-background quotient is denoted
\begin{equation}
 \boxed{\Gamma_{W,a}(g).}
 \label{eq:v38-Witten-loop}
\end{equation}
Its orientation is fixed by the path from $I$ to $U^g$ followed by the endpoint
gauge identification $g_a^{-1}$.
\end{definition}

\begin{proposition}[The Witten loop lies in the V37 zero-winding sector]
\label{prop:v38-zero-winding}
For the loop $\Gamma_{W,a}(g)$,
\begin{equation}
 \boxed{
 \mathfrak w_a(\Gamma_{W,a}(g))=0.}
 \label{eq:v38-zero-winding}
\end{equation}
Hence its defect class lies in $\mathfrak B_{5,a}^{0}$.
\end{proposition}

\begin{proof}
The loop varies only the embedded $SU(2)$ link factor.  The V37 edge-winding
cochain records the $\pi_1(G_6)\cong\mathbb Z$ class of the corresponding link
loop.  The inclusion
$SU(2)\hookrightarrow G_6$ induces the zero map on $\pi_1$ because $SU(2)$ is
simply connected.  Therefore every edge winding is zero and so is its cellular
cohomology class.
\end{proof}

\begin{proposition}[The native class is the nontrivial weak $\pi_4$ class]
\label{prop:v38-pi4-class}
The continuum gauge transformation underlying
$\Gamma_{W,a}(g)$ represents the nonzero element of
\begin{equation}
 \pi_4(G_6)\cong\pi_4(SU(2))\cong\mathbb Z_2.
 \label{eq:v38-pi4-G6}
\end{equation}
For sufficiently small $a$, any two smooth admissible discretizations of the
same based continuum class are homotopic through V2-admissible lattice loops.
\end{proposition}

\begin{proof}
A quotient by a discrete central subgroup is a covering and therefore preserves
all homotopy groups in degree at least two.  Since
$\pi_4(SU(3))=0$, $\pi_4(U(1))=0$ and
$\pi_4(SU(2))\cong\mathbb Z_2$, equation
\eqref{eq:v38-pi4-G6} follows.  The continuum map was chosen to be its
nontrivial element.  A compact homotopy between two smooth representatives has
uniformly bounded curvature and finitely many derivatives.  V24's smooth
background stability and V2 plaquette/Wilson-gap openness then imply that all
sufficiently fine lattice discretizations of that homotopy remain in the same
admissible component.
\end{proof}

\subsection{The finite overlap twist}
\label{sec:v38-twist}

Consider first one left-handed weak doublet and suppress spectator colour and
hypercharge factors.  Let
\begin{equation}
 P_t:=\widehat P_-[U_{a,t}]
 \label{eq:v38-Pt}
\end{equation}
be the V2 overlap chiral projector, and let $Q_t$ be its Kato transport,
\begin{equation}
 \dot Q_t=[\dot P_t,P_t]Q_t,
 \qquad
 Q_0=I.
 \label{eq:v38-Kato}
\end{equation}
Gauge covariance gives
\begin{equation}
 P_1=R(g_a)P_0R(g_a)^{-1}.
 \label{eq:v38-endpoint-covariance}
\end{equation}
Therefore
$R(g_a)^{-1}Q_1$ preserves $\Ran P_0$.

\begin{definition}[Gauge-closed overlap twist]
\label{def:v38-twist}
The finite native twist of the Witten loop in a representation $R$ is
\begin{equation}
 \boxed{
 \mathcal T_{a,R}(g)
 :=\det_{\Ran P_0}
 \left(
 P_0R(g_a)^{-1}Q_1P_0
 \right).}
 \label{eq:v38-twist}
\end{equation}
Equivalently one may take the determinant of
$I-P_0+P_0R(g_a)^{-1}Q_1P_0$ on the full finite fermion space.
\end{definition}

\begin{lemma}[Reality and homotopy invariance for a fundamental doublet]
\label{lem:v38-real-twist}
For the fundamental representation of $SU(2)$,
\begin{equation}
 \boxed{
 \mathcal T_{a,\mathbf 2}(g)\in\{+1,-1\}.}
 \label{eq:v38-real-twist}
\end{equation}
On the admissible Wilson-gap component it is constant under every continuous
homotopy of the closed loop.
\end{lemma}

\begin{proof}
The fundamental representation is pseudoreal.  Combining its invariant
antisymmetric tensor with the Euclidean spin charge-conjugation operation gives
the standard antiunitary reality structure of the $SU(2)$ overlap projector.
The Kato equation is equivariant under the same structure.  Consequently the
finite determinant in \eqref{eq:v38-twist} is real.  Kato transport and the
endpoint gauge action are unitary, so its modulus is one; hence it is $\pm1$.
The twist depends continuously on an admissible loop.  A continuous map into
the discrete set $\{\pm1\}$ is constant on each loop-homotopy class.
\end{proof}

\begin{theorem}[Finite-overlap Witten twist]
\label{thm:v38-fundamental-minus}
For the V2 overlap regulator in the fundamental weak representation and the
native Witten loop of Definition~\ref{def:v38-Witten-loop}, there exists
$a_*>0$ such that
\begin{equation}
 \boxed{
 \mathcal T_{a,\mathbf 2}(g)=-1
 \qquad(0<a<a_*).}
 \label{eq:v38-fundamental-minus}
\end{equation}
\end{theorem}

\begin{proof}
This is the established finite-overlap realization of Witten's global anomaly,
with hypotheses matched to the present native operator.  Neuberger proved that
the overlap reproduces the $SU(2)$ global anomaly through the Berry phase of
the Hermitian Wilson operator, and B\"ar--Campos formulated the same statement
in the Ginsparg--Wilson/Kato language used here: for a smooth admissible loop
built from a gauge transformation in the nontrivial class of
$\pi_4(SU(2))$, the fundamental-representation twist is nontrivial.  Their
construction uses precisely a gauge-covariant Wilson kernel, the overlap
projector, and Kato transport; these are the V2 objects in
\eqref{eq:v38-Pt}--\eqref{eq:v38-twist}.

For completeness, the part of the argument needed to turn their continuum
calculation into the exact finite-lattice sign is short.  Embed the weak
$SU(2)$ loop into $SU(3)$, where the corresponding $\pi_4$ class is trivial,
and span it by the standard smooth two-parameter surface.  The Kato twist is
then the exponential of the projector-curvature surface integral.  Its
small-$a$ expansion is the continuum Wess--Zumino integral and tends to $-1$
for the nontrivial $SU(2)$ class.  Lemma~\ref{lem:v38-real-twist} says that the
finite-lattice twist is already restricted exactly to $\{\pm1\}$.  Therefore,
once the continuum expansion lies in a sufficiently small neighbourhood of
$-1$, the exact lattice value must equal $-1$.  V24 supplies the smooth
background/admissibility control required to stay on the same overlap branch
for all sufficiently small $a$.
\end{proof}

\begin{remark}[Matched external input]
\label{rem:v38-external}
The finite-overlap theorem used in
Theorem~\ref{thm:v38-fundamental-minus} is the lattice result of
H.~Neuberger, \emph{Witten's $SU(2)$ anomaly on the lattice} (1998), and the
Ginsparg--Wilson/Kato formulation of O.~B\"ar and I.~Campos,
\emph{Global anomalies in chiral gauge theories on the lattice},
Nucl. Phys. B581 (2000) 499--519.  V38 does not import a continuum Pfaffian sign
and rename it a lattice theorem; it uses the published finite-overlap twist
after matching its projector, Wilson-kernel, smoothness and admissibility
hypotheses to V2/V24.
\end{remark}

\subsection{Direct sums and the exact Standard-Model sign}
\label{sec:v38-SM-sign}

\begin{lemma}[Multiplicativity of the native twist]
\label{lem:v38-multiplicativity}
For a finite direct sum of weak representations,
\begin{equation}
 \boxed{
 \mathcal T_{a,R_1\oplus R_2}
 =\mathcal T_{a,R_1}\mathcal T_{a,R_2}.}
 \label{eq:v38-multiplicativity}
\end{equation}
The weak singlet has twist one.  Consequently $N$ identical fundamental
left-handed doublets have
\begin{equation}
 \boxed{
 \mathcal T_{a,N\mathbf2}=(-1)^N}
 \label{eq:v38-N-doublets}
\end{equation}
for sufficiently small $a$.
\end{lemma}

\begin{proof}
For a direct sum, the Wilson kernel, overlap sign, chiral projector and Kato
transport are block diagonal.  The determinant in
\eqref{eq:v38-twist} therefore factors.  A weak singlet is constant along the
pure $SU(2)$ loop, so its Kato/end-point operator is the identity.  Apply
Theorem~\ref{thm:v38-fundamental-minus} to every fundamental block.
\end{proof}

\begin{theorem}[Exact cancellation of the Witten generator in the complete Standard Model]
\label{thm:v38-SM-cancellation}
For one Standard-Model generation, the weak left-handed doublets contribute
three colour copies of the quark doublet and one lepton doublet.  Hence
\begin{equation}
 N_{\rm dbl}^{(1)}=3+1=4.
 \label{eq:v38-one-generation}
\end{equation}
For three generations,
\begin{equation}
 \boxed{N_{\rm dbl}=12.}
 \label{eq:v38-twelve}
\end{equation}
Therefore the complete native Standard-Model determinant-line phase on the
Witten generator is
\begin{equation}
 \boxed{
 \rho_{a,\rm SM}(\Gamma_{W,a})
 =(-1)^{12}=+1}
 \label{eq:v38-SM-Witten-one}
\end{equation}
for every sufficiently fine V2-consistent lattice representative.
\end{theorem}

\begin{proof}
The left-handed quark representation $(\mathbf3,\mathbf2)$ is three copies of
the weak fundamental when the colour background is held fixed.  The
left-handed lepton representation $(\mathbf1,\mathbf2)$ gives one further
copy.  All right-handed Standard-Model fermions are weak singlets.  Hypercharge
is a spectator along the pure weak loop, and colour supplies only the
multiplicity three.  Thus Lemma~\ref{lem:v38-multiplicativity} gives four
minus signs per generation and twelve for three generations.  Their product is
one.

The same conclusion is the inherited ``even $SU(2)$ doublet parity'' stated in
V1--V7, now realized as an explicit finite-overlap transport theorem.  Gauge
and Yukawa interactions do not alter this measure-line sign: the moving chiral
measure is determined by the overlap projector, while the finite
Higgs/Yukawa/action block is gauge covariant and the endpoints are identified by
$g_a$.  Therefore the full action/source section acquires exactly the same
measure-line product sign, namely $+1$.
\end{proof}

\begin{corollary}[All retained source/BV descendants have trivial Witten sign]
\label{cor:v38-source-sign}
Let $I$ be any retained gauge, coframe, Higgs/Yukawa, ghost/antifield,
reference or determinant-line source multi-index.  The complete
Standard-Model Witten-loop transport acts trivially on the corresponding
source coefficient:
\begin{equation}
 \boxed{
 \rho_{a,\rm SM}^{(I)}(\Gamma_{W,a})=1.}
 \label{eq:v38-source-sign}
\end{equation}
Equivalently, the Witten generator contributes no nontrivial transition to the
V7/V16 coefficientwise continuum packet.
\end{corollary}

\begin{proof}
The sign in Theorem~\ref{thm:v38-SM-cancellation} is an endpoint transition of
the finite determinant line, not a source-dependent scalar function.  Source
differentiation acts on the Kato frames, determinant/Pfaffian factors and
covariant action coefficients before the endpoint transition is taken.  The
transition of a direct sum is the same product of representation twists for
every differentiated coefficient.  Since that product is exactly one for the
complete Standard-Model packet, every retained descendant is unchanged.
\end{proof}

\subsection{Pseudoreal Pfaffian interpretation}
\label{sec:v38-pseudoreal}

The weak fundamental is pseudoreal, so one may equivalently describe the
single-doublet obstruction as the orientation of a real/pseudoreal Pfaffian
line.  This is the charged analogue deliberately excluded from V25's
positive neutral-Majorana theorem.

\begin{proposition}[Single-doublet orientation reversal]
\label{prop:v38-pseudoreal-sign}
For one fundamental left-handed weak doublet, the real/pseudoreal orientation
line transported around $\Gamma_{W,a}$ reverses orientation:
\begin{equation}
 \boxed{
 \operatorname{Hol}_{\rm Pf}^{\mathbf2}(\Gamma_{W,a})=-1.}
 \label{eq:v38-pf-minus}
\end{equation}
For $N$ such doublets the orientation sign is $(-1)^N$.
\end{proposition}

\begin{proof}
The pseudoreal structure identifies the complex overlap determinant transition
with the square/orientation datum of the corresponding real skew formulation.
The nontrivial complex Kato twist of
Theorem~\ref{thm:v38-fundamental-minus} is therefore precisely the orientation
reversal of the pseudoreal Pfaffian line.  Direct sums tensor the real
orientation lines, so the signs multiply.  This is the charged case anticipated
in the V25 firewall and is distinct from the natively oriented neutral
Majorana branch treated there.
\end{proof}

\subsection{The complete smooth zero-winding gauge-transform 4-type}
\label{sec:v38-Postnikov}

Cancellation of the Witten generator does not yet exhaust the smooth
zero-winding sector.  We now identify the remaining large-gauge components
before asking for their exact phase.

Let $\mathcal G_0^{\rm sm}$ be the based smooth gauge transformations
$T^4\to G_6$ with vanishing V37 $H^1$ winding.  Such a transformation lifts to
the universal cover
\begin{equation}
 SU(3)\times SU(2)\times\mathbb R,
 \label{eq:v38-universal-cover}
\end{equation}
and the $\mathbb R$ factor is contractible.

\begin{theorem}[Four-dimensional Postnikov classification of the zero-winding gauge group]
\label{thm:v38-Postnikov}
The based homotopy classes of smooth zero-winding gauge transformations have
four-dimensional Postnikov shadow
\begin{equation}
 \boxed{
 \pi_0(\mathcal G_0^{\rm sm})
 \cong
 H^3(T^4;\mathbb Z\oplus\mathbb Z)
 \oplus
 H^4(T^4;\mathbb Z_2)
 \cong
 \mathbb Z^8\oplus\mathbb Z_2.}
 \label{eq:v38-Postnikov}
\end{equation}
The two $\mathbb Z$ coefficients in degree three come from
$\pi_3(SU(3))\cong\mathbb Z$ and
$\pi_3(SU(2))\cong\mathbb Z$.  The $\mathbb Z_2$ factor is
$\pi_4(SU(2))$ and is generated by the Witten transformation
\eqref{eq:v38-g}.
\end{theorem}

\begin{proof}
Vanishing $H^1$ winding is exactly the obstruction-free condition for lifting
the map through the universal cover.  After discarding its contractible real
factor, the target is the simply connected group
$SU(3)\times SU(2)$.  Its homotopy groups through degree four are
\begin{equation}
 \pi_2=0,
 \qquad
 \pi_3\cong\mathbb Z\oplus\mathbb Z,
 \qquad
 \pi_4\cong\mathbb Z_2.
 \label{eq:v38-low-homotopy}
\end{equation}
Since the target is $2$-connected, its Postnikov $4$-stage begins with
$K(\mathbb Z^2,3)$ and a $K(\mathbb Z_2,4)$ fibre.  The first possible
Postnikov $k$-invariant affecting this extension lies in degree five and
therefore cannot be detected by the four-dimensional source $T^4$.  Obstruction
theory consequently gives
\begin{equation}
 [T^4,SU(3)\times SU(2)]_*
 \cong
 H^3(T^4;\mathbb Z^2)
 \oplus H^4(T^4;\mathbb Z_2).
 \label{eq:v38-obstruction-theory}
\end{equation}
Finally
$H^3(T^4;\mathbb Z)\cong\mathbb Z^4$ and
$H^4(T^4;\mathbb Z_2)\cong\mathbb Z_2$, yielding
\eqref{eq:v38-Postnikov}.  The $\pi_4(SU(2))$ generator is exactly the map
used in Definition~\ref{def:v38-Witten-loop}.
\end{proof}

\begin{corollary}[What remains after Witten cancellation]
\label{cor:v38-free-sector}
After the V37 $H^1$ electromagnetic winding sector and the V38 Witten
$\mathbb Z_2$ generator have been shown to lie in the exact complete-packet
phase kernel, the remaining smooth large-gauge transformation shadow of the
zero-winding sector is the free group
\begin{equation}
 \boxed{
 H^3(T^4;\mathbb Z^2)\cong\mathbb Z^8.}
 \label{eq:v38-Z8}
\end{equation}
It consists of four $SU(3)$ and four $SU(2)$ three-cycle winding generators.
\end{corollary}

\begin{proof}
Combine Theorem~\ref{thm:v38-Postnikov} with
Theorem~\ref{thm:v38-SM-cancellation}.  The latter kills the only torsion
$\mathbb Z_2$ factor in the smooth zero-winding large-gauge shadow, leaving the
degree-three free factors.
\end{proof}

\begin{firewall}[The $\mathbb Z^8$ phase is not inferred from local anomaly cancellation]
The free $H^3$ generators are natural candidates for phases controlled by
four-dimensional perturbative anomaly/Chern--Simons data, but V38 does not
promote the V3 local anomaly cancellation or the V35 coefficientwise descent to
an exact finite-cutoff statement about their global flat character.  A flat
$U(1)$ character on $\mathbb Z^8$ can be invisible to every local curvature
coefficient.  Its exact evaluation is the next global theorem.
\end{firewall}

\subsection{V38 anti-circularity audit}
\label{sec:v38-audit}

\begin{enumerate}[leftmargin=2.4em]
\item \textbf{The continuum Witten sign was not silently substituted for the
native regulator.}  Theorem~\ref{thm:v38-fundamental-minus} uses the established
finite-overlap twist theorem after matching its Wilson, overlap, Kato and
admissibility hypotheses to V2/V24.
\item \textbf{The sign is exact at finite lattice spacing for sufficiently
small $a$.}  Pseudoreality restricts the twist to $\{\pm1\}$ before the
continuum expansion is taken; the expansion selects which discrete value
occurs.
\item \textbf{The loop is genuinely in V37's zero-winding sector.}
The $SU(2)$ inclusion contributes no $\pi_1(G_6)$ edge winding.
\item \textbf{Yukawa/Higgs terms are not used to define the anomaly sign.}
The sign belongs to the moving overlap measure line; the finite interaction
action is gauge covariant and does not alter the representation twist.
\item \textbf{Even doublet parity is realized rather than merely quoted.}
The complete packet has twelve fundamental weak copies, so the finite overlap
sign is explicitly $(-1)^{12}=1$.
\item \textbf{The remaining zero-winding gauge group is not declared
trivial.}  The Postnikov calculation exposes the free $\mathbb Z^8$ sector
which survives after the Witten $\mathbb Z_2$ sign is cancelled.
\end{enumerate}

\section{V38 status ledger}
\label{sec:v38-status}

\subsection*{Proved in V38}
\begin{enumerate}[leftmargin=2.2em]
\item A based nontrivial map $S^4\to SU(2)$ produces a smooth V2-admissible
native Witten loop whose V37 edge-winding class vanishes and whose underlying
large-gauge class is the nonzero element of $\pi_4(G_6)\cong\mathbb Z_2$.
\item For one fundamental weak doublet, the finite overlap/Kato twist is real,
unit modulus and homotopy invariant.  Matching the V2 operator to the
established finite-overlap Witten theorem gives the exact sign $-1$ for all
sufficiently fine smooth admissible lattice representatives.
\item The twist is multiplicative under direct sums and trivial on weak
singlets.  Therefore $N$ fundamental doublets contribute $(-1)^N$.
\item The complete three-generation Standard Model contains twelve weak
fundamental copies: three colour copies of the quark doublet plus one lepton
doublet per generation.  Hence the exact native Witten phase is
$(-1)^{12}=+1$, including every retained source/BV descendant.
\item In the pseudoreal real-line formulation, the same result is the
orientation reversal of one charged fundamental Pfaffian line and its exact
cancellation for an even number of doublets.  This closes the charged mod-two
firewall left open by V25 for the Standard-Model packet.
\item The based smooth zero-$H^1$ gauge-transform components have four-type
$\mathbb Z^8\oplus\mathbb Z_2$.  The torsion $\mathbb Z_2$ is the Witten class;
a free $\mathbb Z^8=H^3(T^4;\mathbb Z^2)$ sector remains after its complete
Standard-Model phase is cancelled.
\end{enumerate}

\subsection*{Inherited}
\begin{enumerate}[leftmargin=2.2em]
\item V1--V7: source-complete coefficientwise ESM continuum, the complete local
anomaly cancellation, inherited even weak-doublet parity, common BV restoration
and exact loop-order projectivity.
\item V8--V24: regulator universality, global controlled zero-mode/module
classification and the V2-consistent soft norm-resolvent/Feshbach theorem.
\item V25--V29: neutral real-Pfaffian orientation, divisor crossings, Wilson
index jumps, higher-rank wall lines and nonsplit angular spectral bundles.
\item V30--V36: faithful geometric cancellation mechanism, relative phase
correction, exact determinant-line eta character, local five-dimensional
infinite-wall realization, coefficientwise Chern--Simons descent and the native
fibered-bordism defect group.
\item V37: nonzero electromagnetic $H^1$ defect subgroup, exact phase
triviality on that subgroup, and reduction of the remaining global character to
$\mathfrak B_{5,a}^{0}$.
\end{enumerate}

\subsection*{Matched external input}
\begin{enumerate}[leftmargin=2.2em]
\item The finite-overlap $SU(2)$ global-anomaly theorem of
Neuberger and of B\"ar--Campos: the overlap twist on the nontrivial
$\pi_4(SU(2))$ loop equals $-1$ in the fundamental representation for
sufficiently fine smooth admissible lattices.  V38 matches this theorem to the
explicit V2 Wilson/overlap/Kato objects and uses its discreteness to obtain an
exact finite-cutoff sign.
\end{enumerate}

\subsection*{Conditional}
\begin{enumerate}[leftmargin=2.2em]
\item The exact $-1$ theorem is stated for sufficiently fine discretizations of
a fixed smooth admissible continuum Witten loop; rough lattice representatives
outside the V2 Wilson-gap component are not covered.
\item Other background fields must be transported in an endpoint-compatible
protected path.  They do not affect the weak overlap twist but can matter for
whether a chosen full ESM path stays inside the declared source chart.
\item The Postnikov calculation classifies the smooth based large-gauge shadow.
The full microscopic zero-winding defect group can still contain V29 angular or
other purely local-operator classes not represented by a smooth gauge
transformation.
\item Positive IR/reference and all V2-consistent smooth-background assumptions
remain in force.
\end{enumerate}

\subsection*{Open beyond V38}
\begin{enumerate}[leftmargin=2.2em]
\item Compute the exact native phase on the free
$H^3(T^4;\mathbb Z^2)\cong\mathbb Z^8$ large-gauge sector.  Local anomaly
cancellation strongly constrains it but does not by itself determine an exact
flat character.
\item Determine whether the $\mathbb Z^8$ generators become fibered-fillable
after adjoining the V35 local Chern--Simons primitive, or whether their phase
must be evaluated through a separate exact five-dimensional transgression.
\item After the full smooth gauge-transform shadow is exhausted, compare the
remaining zero-winding defect with the V29 angular/microscopic loop sector.
\item The independent stronger problems remain: positive-IR removal, analytic
finite-coupling RG trajectories, perturbative summability and nonperturbative
ESM completion.
\end{enumerate}

\subsection*{Exact next load-bearing theorem}
\begin{center}
\fbox{\parbox{0.92\textwidth}{
\textbf{Free $H^3$ large-gauge / exact Chern--Simons phase theorem.}
After removal of the V37 electromagnetic $H^1$ subgroup and the V38 Witten
$\mathbb Z_2$ class, the smooth zero-winding large-gauge shadow is
$H^3(T^4;\mathbb Z^2)\cong\mathbb Z^8$, with four generators from $SU(3)$ and
four from $SU(2)$.  Construct native resolved loops for these generators and
compute the exact V33/V34 determinant-line integrability character.  Prove that
the complete Standard-Model phase is one using the full anomaly-free
representation packet and an exact native five-dimensional transgression, or
exhibit the residual flat character.  Only after this free sector is settled
should one return to purely angular/microscopic classes in
$\mathfrak B_{5,a}^{0}$.}}
\end{center}

\section*{Append-only continuation marker: V38}
\addcontentsline{toc}{section}{Append-only continuation marker: V38}
\textbf{End of V38.}  The complete V1--V37 body above is retained verbatim.
V38 constructs the native zero-winding loop associated with the nontrivial
$\pi_4(SU(2))$ gauge transformation, matches its V2 Kato twist to the
established finite-overlap Witten theorem, and obtains the exact fundamental
sign $-1$ for sufficiently fine admissible lattices.  Direct-sum
multiplicativity then realizes the inherited even-doublet criterion explicitly:
the three-generation Standard Model has twelve weak fundamental copies and
therefore phase $+1$ on the Witten generator, including all retained source/BV
descendants.  Going upstream, V38 also computes the four-dimensional Postnikov
shadow of smooth zero-$H^1$ gauge transformations as
$\mathbb Z^8\oplus\mathbb Z_2$; after the Witten factor is cancelled, the next
exact global gate is the free $H^3(T^4;\mathbb Z^2)$ large-gauge sector.



\clearpage
\section{V39 continuation: free $H^3$ large-gauge generators, exact conjugate/pseudoreal cancellation, and exhaustion of the smooth gauge-transform shadow}
\label{sec:v39-continuation}

V38 isolated the free
\(H^3(T^4;\mathbb Z^2)\cong\mathbb Z^8\) sector of smooth zero-$H^1$
large gauge transformations.  The tempting route would be to compute an exact
five-dimensional Chern--Simons phase for every generator.  That is unnecessary
for the complete Standard-Model packet.  The native overlap line has two exact
finite-dimensional reality properties which are stronger than the local
anomaly calculation:
\begin{enumerate}[label=\textup{(V39.\arabic*)},leftmargin=3.8em]
\item a representation and its complex conjugate have conjugate Kato twists;
\item a pseudoreal fundamental $SU(2)$ representation has a real unit twist,
so every closed pure-$SU(2)$ loop contributes only a sign.
\end{enumerate}
The Standard Model pairs the colour representations as two fundamentals with
two antifundamentals per generation, while it contains four weak fundamental
copies per generation.  These two native facts kill the complete $SU(3)$ and
$SU(2)$ phases on \emph{every} smooth pure-factor large-gauge loop, not merely
on one chosen $H^3$ representative.

V39 first constructs explicit native representatives of the eight free
$H^3$ generators, then proves their exact finite-cutoff phase is one.  Combining
this with V37 and V38 shows that the entire smooth large-gauge transformation
shadow of the faithful Standard-Model defect group lies in the exact phase
kernel.  The next global problem is therefore purely microscopic: it is the
quotient by smooth gauge-transform classes, whose first known representatives
are the V29 angular/Wilson-operator loops.

\subsection{Smooth representatives of the eight $H^3$ generators}
\label{sec:v39-generators}

Let $G$ be either $SU(3)$ or $SU(2)$.  For each
$\mu\in\{1,2,3,4\}$ let
$T^3_{\widehat\mu}\subset T^4$ denote the complementary coordinate three-torus.
Choose a smooth based map
\begin{equation}
 g_{G,\mu}:T^4\longrightarrow G
 \label{eq:v39-g-generator}
\end{equation}
which is independent of $x_\mu$ and whose restriction to
$T^3_{\widehat\mu}$ has degree one.  With the conventional normalized trace,
\begin{equation}
 \nu_{G,\mu}(g)
 :=\frac{1}{24\pi^2}
 \int_{T^3_{\widehat\mu}}
 \operatorname{tr}\bigl((g^{-1}dg)^3\bigr)\in\mathbb Z
 \label{eq:v39-H3-winding}
\end{equation}
is the corresponding $H^3$ coordinate, and
$\nu_{G,\mu}(g_{G,\mu})=1$.

\begin{lemma}[Smooth native path to the pure-gauge endpoint]
\label{lem:v39-native-H3-path}
Put $\alpha=g^{-1}dg$ and define the continuum interpolation
\begin{equation}
 A_t=t\alpha,
 \qquad 0\le t\le1.
 \label{eq:v39-At}
\end{equation}
Let $U_{a,t}$ be its native link parallel transports.  Then
\begin{equation}
 U_{a,0}=I,
 \qquad
 U_{a,1}(x,x+a\hat\nu)
 =g(x)^{-1}g(x+a\hat\nu),
 \label{eq:v39-pure-gauge-endpoint}
\end{equation}
so the endpoint is gauge-equivalent to the vacuum.  For a fixed smooth $g$,
all links obey $U_{a,t}=I+O(a)$ and all plaquettes obey
\begin{equation}
 U_{\square,a,t}=I+O(a^2)
 \label{eq:v39-plaquette-small}
\end{equation}
uniformly in $t$.  Hence for sufficiently small $a$ the path lies in the V2
smooth admissible Wilson-gap component.  Its V37 edge-winding class vanishes.
\end{lemma}

\begin{proof}
The endpoint formula is the standard parallel-transport identity for the flat
connection $g^{-1}dg$.  Smoothness of $g$ on the compact torus gives uniformly
bounded connection and curvature coefficients.  The link and plaquette
estimates are therefore the usual short-edge expansions.  V2 gap stability
then applies for all sufficiently small $a$.  Since the path and endpoint lie
inside the simply connected $SU(3)$ or $SU(2)$ factor, no edge loop represents
the V37 $\pi_1(G_6)$ generator, so the cellular $H^1$ winding vanishes.
\end{proof}

After endpoint gauge identification, the path of
Lemma~\ref{lem:v39-native-H3-path} is a closed native determinant-line loop,
denoted $\Gamma_{G,\mu}$.  Its smooth large-gauge component is the corresponding
basis element of
$H^3(T^4;\pi_3(G))$.

\begin{proposition}[The eight loops realize the free Postnikov subgroup]
\label{prop:v39-Z8-generators}
The four loops $\Gamma_{SU(3),\mu}$ and the four loops
$\Gamma_{SU(2),\mu}$ generate the free group
\begin{equation}
 \boxed{
 H^3(T^4;\mathbb Z^2)
 \cong
 \mathbb Z^4_{SU(3)}\oplus\mathbb Z^4_{SU(2)}.}
 \label{eq:v39-Z8-generators}
\end{equation}
Their products and inverses realize arbitrary integral winding vectors.
\end{proposition}

\begin{proof}
The classes in \eqref{eq:v39-H3-winding} are the four standard basis elements
of $H^3(T^4;\mathbb Z)$ for each simply connected factor.  The V38 Postnikov
theorem identifies these eight integers with the entire free part of the
smooth zero-$H^1$ component group.  Concatenation adds the winding numbers and
path reversal changes their signs.
\end{proof}

\subsection{Exact conjugation of the finite overlap twist}
\label{sec:v39-conjugate}

The colour cancellation is an exact finite-regulator duality, not an anomaly
expansion.

\begin{lemma}[Conjugate representations have conjugate overlap twists]
\label{lem:v39-conjugate-twist}
Let $R$ be a finite-dimensional unitary representation of a compact gauge
factor and let $\overline R$ be its complex-conjugate representation.  For any
smooth admissible gauge-closed native loop $\Gamma$, the V2 Wilson kernels,
overlap projectors and Kato transports obey a fixed antiunitary conjugation
relation.  Consequently their determinant-line twists satisfy
\begin{equation}
 \boxed{
 \mathcal T_{a,\overline R}(\Gamma)
 =\overline{\mathcal T_{a,R}(\Gamma)}.}
 \label{eq:v39-conjugate-twist}
\end{equation}
Since every Kato twist is unitary,
\begin{equation}
 \boxed{
 \mathcal T_{a,R}(\Gamma)
 \mathcal T_{a,\overline R}(\Gamma)=1.}
 \label{eq:v39-conjugate-product}
\end{equation}
The same identities hold for every retained source/BV descendant.
\end{lemma}

\begin{proof}
Complex conjugation of the internal gauge matrices changes $R(U)$ to
$\overline R(U)$.  Euclidean spin charge conjugation supplies one fixed
unitary matrix $C_{\rm sp}$ which restores the chosen Clifford basis after
complex conjugation.  Thus, with
$\mathcal C=C_{\rm sp}K$ and $K$ ordinary coefficient conjugation,
\begin{equation}
 H_{W,\overline R}(U)
 =\mathcal C H_{W,R}(U)\mathcal C^{-1},
 \label{eq:v39-H-conjugate}
\end{equation}
where the anti-linearity of $\mathcal C$ is understood.  Riesz functional
calculus then gives the same relation for the overlap sign and chiral
projector.  The Kato ODE is polynomial in the projector and its derivative, so
its solution is conjugated in the same way.  Endpoint gauge actions are also
complex conjugates.  Taking the determinant of the gauge-closed Kato operator
therefore gives \eqref{eq:v39-conjugate-twist}.  Its modulus is one because the
Kato transport and endpoint gauge identification are unitary.  Source
coefficients are obtained by differentiating these exact finite identities, so
the conjugation/product laws remain valid coefficientwise.
\end{proof}

\subsection{Exact colour cancellation on the full $H^3$ sector}
\label{sec:v39-SU3}

\begin{theorem}[The complete Standard-Model colour phase is one on every pure $SU(3)$ loop]
\label{thm:v39-SU3-one}
Let $\Gamma$ be any smooth admissible closed native loop obtained from a pure
$SU(3)$ gauge-field path, with the electroweak links and spectator backgrounds
held fixed in a compatible protected packet.  Then for one Standard-Model
generation
\begin{equation}
 \boxed{\rho_{a,{\rm gen}}^{SU(3)}(\Gamma)=1,}
 \label{eq:v39-SU3-one-gen}
\end{equation}
and hence for all three generations
\begin{equation}
 \boxed{\rho_{a,{\rm SM}}^{SU(3)}(\Gamma)=1.}
 \label{eq:v39-SU3-one}
\end{equation}
The statement is exact at finite cutoff on the admissible loop and includes all
retained source/BV descendants.
\end{theorem}

\begin{proof}
In a left-handed convention, one generation contains two colour fundamentals
from the two weak components of $Q_L$ and two colour antifundamentals from
$u_R^c$ and $d_R^c$.  Along a pure colour loop the electroweak representation
labels are spectators, so if
\begin{equation}
 \tau_3(\Gamma):=\mathcal T_{a,\mathbf3}(\Gamma),
 \label{eq:v39-tau3}
\end{equation}
then Lemma~\ref{lem:v39-conjugate-twist} gives
\begin{equation}
 \rho_{a,{\rm gen}}^{SU(3)}(\Gamma)
 =\tau_3^2\overline{\tau_3}^{\,2}=1.
 \label{eq:v39-SU3-cancel-calc}
\end{equation}
No mass or Yukawa assumption is needed for this measure-line identity; the
finite Higgs/Yukawa action is gauge covariant and does not change the chiral
measure twist.  Three generations merely cube the unit result.  Differentiated
source coefficients inherit the same exact representation-product transition.
\end{proof}

\begin{corollary}[All four colour $H^3$ generators lie in the phase kernel]
\label{cor:v39-SU3-H3}
For every $\mu$ and $n\in\mathbb Z$,
\begin{equation}
 \boxed{
 \rho_{a,{\rm SM}}
 \bigl(\Gamma_{SU(3),\mu}^{\,n}\bigr)=1.}
 \label{eq:v39-SU3-H3}
\end{equation}
\end{corollary}

\begin{proof}
Apply Theorem~\ref{thm:v39-SU3-one} to the generators of
Proposition~\ref{prop:v39-Z8-generators} and use multiplicativity of the flat
character under concatenation.
\end{proof}

\subsection{Exact weak cancellation on every pure $SU(2)$ loop}
\label{sec:v39-SU2}

V38 proved more than the sign of its specific Witten representative: its
pseudoreality lemma says that the fundamental weak twist of any admissible
closed pure-$SU(2)$ overlap loop is real and unitary.  We record the resulting
complete-packet consequence.

\begin{theorem}[Even-doublet parity kills every pure $SU(2)$ overlap phase]
\label{thm:v39-SU2-all}
Let $\Gamma$ be any smooth admissible gauge-closed loop whose moving gauge links
lie only in the weak $SU(2)$ factor.  For one fundamental left-handed doublet,
\begin{equation}
 \mathcal T_{a,\mathbf2}(\Gamma)\in\{+1,-1\}.
 \label{eq:v39-SU2-sign-any}
\end{equation}
The complete three-generation Standard-Model packet therefore obeys
\begin{equation}
 \boxed{
 \rho_{a,{\rm SM}}^{SU(2)}(\Gamma)
 =\mathcal T_{a,\mathbf2}(\Gamma)^{12}
 =1.}
 \label{eq:v39-SU2-all}
\end{equation}
This includes all free $H^3$ generators, the V38 Witten generator, and every
retained source/BV descendant.
\end{theorem}

\begin{proof}
The first statement is exactly the antiunitary pseudoreality argument of
Lemma~\ref{lem:v38-real-twist}; that proof uses only the fundamental $SU(2)$
reality structure, unitarity of Kato transport and gauge closure, not the
special homotopy class of the Witten loop.  For a pure weak loop, each quark
colour copy has the same fundamental overlap twist and the lepton doublet gives
one more identical copy.  Thus one generation contributes
$\mathcal T_{a,\mathbf2}^4=1$, and three generations contribute the twelfth
power.  Weak singlets contribute one.  Hypercharge and colour are spectators
on the loop.  As in V38, Yukawa/Higgs terms are finite gauge-covariant action
data and do not alter the measure-line twist.  The endpoint transition is the
same exact sign for every differentiated source coefficient, so all retained
descendants are trivial as well.
\end{proof}

\begin{corollary}[All four weak $H^3$ generators lie in the phase kernel]
\label{cor:v39-SU2-H3}
For every $\mu$ and $n\in\mathbb Z$,
\begin{equation}
 \boxed{
 \rho_{a,{\rm SM}}
 \bigl(\Gamma_{SU(2),\mu}^{\,n}\bigr)=1.}
 \label{eq:v39-SU2-H3}
\end{equation}
\end{corollary}

\begin{proof}
The loops of Proposition~\ref{prop:v39-Z8-generators} are pure weak loops.
Apply Theorem~\ref{thm:v39-SU2-all} and multiplicativity.
\end{proof}

\subsection{Exact triviality of the full free $H^3$ character}
\label{sec:v39-Z8}

\begin{theorem}[The complete Standard-Model character is trivial on $H^3(T^4;\mathbb Z^2)$]
\label{thm:v39-Z8-trivial}
For every smooth V2-admissible zero-$H^1$ large gauge transformation whose V38
Postnikov class lies in
\begin{equation}
 H^3(T^4;\mathbb Z^2)
 \cong
 \mathbb Z^8,
 \label{eq:v39-H3-group}
\end{equation}
the exact complete-packet native determinant/Pfaffian character is one:
\begin{equation}
 \boxed{
 \rho_{a,{\rm SM}}(\Gamma_g)=1.}
 \label{eq:v39-Z8-trivial}
\end{equation}
The equality holds at finite cutoff for every sufficiently fine smooth
admissible representative and for every retained source/BV coefficient.
\end{theorem}

\begin{proof}
By V38's Postnikov classification, the free class decomposes uniquely into an
$SU(3)$ $H^3$ winding vector and an $SU(2)$ $H^3$ winding vector.  The exact
phase character is multiplicative under concatenation.  Corollaries
\ref{cor:v39-SU3-H3} and \ref{cor:v39-SU2-H3} kill each of the eight basis
generators and their inverses.  Hence the character is one on the whole free
group.  The source/BV statement follows from the corresponding exact
coefficientwise transition theorems in the two factors.
\end{proof}

\begin{remark}[Why no exact Chern--Simons calculation was needed]
\label{rem:v39-no-CS}
The V38 firewall was correct: local anomaly cancellation alone would not have
determined an arbitrary flat character on $\mathbb Z^8$.  V39 instead uses two
stronger finite-regulator representation identities.  Colour is paired with
its complex conjugate, so the phase cancels as $\tau\bar\tau$.  Weak
pseudoreality quantizes the phase to $\pm1$, and even doublet multiplicity kills
that sign.  These are exact statements about the native overlap line, not a
perturbative Chern--Simons expansion.
\end{remark}

\subsection{The entire smooth large-gauge shadow lies in the exact phase kernel}
\label{sec:v39-all-smooth-gauge}

We can now combine all smooth gauge-transform components which have appeared in
V37--V39.

\begin{theorem}[Exact triviality on the complete smooth large-gauge shadow]
\label{thm:v39-all-smooth}
Let $g:T^4\to G_6$ be any based smooth gauge transformation whose native
representative is V2-admissible at sufficiently fine cutoff.  Decompose its
four-dimensional Postnikov data into:
\begin{enumerate}[label=\textup{(G\arabic*)},leftmargin=3em]
\item the $H^1(T^4;\mathbb Z)$ winding shadow of V37;
\item the free $H^3(T^4;\mathbb Z^2)$ class of V38--V39;
\item the $H^4(T^4;\mathbb Z_2)$ Witten class of V38.
\end{enumerate}
Then the complete Standard-Model exact native character on the corresponding
closed regulator loop is
\begin{equation}
 \boxed{\rho_{a,{\rm SM}}(\Gamma_g)=1.}
 \label{eq:v39-all-smooth}
\end{equation}
\end{theorem}

\begin{proof}
The $H^1$ generator may be represented in the unbroken electromagnetic
stabilizer and has phase one by Theorem~\ref{thm:v37-em-trivial}.  The free
$H^3$ class has phase one by Theorem~\ref{thm:v39-Z8-trivial}.  The sole
$H^4$ torsion class has phase one in the complete packet by
Theorem~\ref{thm:v38-SM-cancellation}.  The exact determinant-line character is
a homomorphism under concatenation, so the product class also has phase one.
\end{proof}

\begin{corollary}[No smooth large-gauge class can carry the remaining exact phase]
\label{cor:v39-smooth-kernel}
Let $\mathfrak G_a^{\rm sm}$ denote the subgroup of the V37 native defect group
generated by smooth large-gauge transformation loops.  Then
\begin{equation}
 \boxed{
 \bar\rho_{a,{\rm SM}}|_{\mathfrak G_a^{\rm sm}}=1.}
 \label{eq:v39-smooth-kernel}
\end{equation}
Consequently the exact residual character factors through the quotient
\begin{equation}
 \boxed{
 \mathfrak B_{5,a}^{\rm mic}
 :=
 \mathfrak B_{5,a}^{\rm nat,V2}/\mathfrak G_a^{\rm sm}.}
 \label{eq:v39-microscopic-quotient}
\end{equation}
Every nontrivial class detected by the resulting character is genuinely
microscopic/local-operator data rather than a smooth large gauge
transformation.
\end{corollary}

\begin{proof}
The phase statement is Theorem~\ref{thm:v39-all-smooth}.  A homomorphism which
is trivial on a subgroup factors uniquely through the group quotient by the
normal subgroup it generates.  Smooth gauge-transform loops form such a
normal phase-kernel subgroup after taking their normal closure in the defect
group.  The quotient therefore carries the residual exact character.
\end{proof}

\subsection{Relation to the V29 angular sector}
\label{sec:v39-angular-target}

The quotient \eqref{eq:v39-microscopic-quotient} removes every smooth
large-gauge component classified by V37--V39.  The first explicit classes which
remain in the lineage are the nonsplit Wilson/angular loops of V29: their
spectral bundles can carry Berry/Chern data even when there is no gauge
transformation and no index jump.

\begin{proposition}[V29 angular loops are invisible to the smooth gauge-transform classification]
\label{prop:v39-angular-not-gauge}
A V29 Pauli-type angular loop with nontrivial positive spectral bundle can have
\begin{equation}
 c_1(E_+)\ne0
 \label{eq:v39-angular-c1}
\end{equation}
while its physical gauge background is fixed.  Hence such a loop maps to zero
under all V37--V39 smooth gauge-transform invariants but can remain nontrivial
in $\mathfrak B_{5,a}^{\rm mic}$.
\end{proposition}

\begin{proof}
The Pauli pencil of V29 varies only the finite Wilson crossing matrix and has
$E_+$ equal to the Hopf line bundle on the resolved $S^2$, with
$c_1(E_+)=1$.  No physical large gauge transformation is involved, so its
$H^1$, $H^3$ and $H^4$ gauge-transform shadows vanish.  Thus V37--V39 do not
remove the class.  Whether its complete faithful Standard-Model determinant
phase is trivial is a separate microscopic angular problem, exactly as V29--V32
anticipated.
\end{proof}

\begin{firewall}[Smooth gauge-transform exhaustion is not defect-group exhaustion]
Theorem~\ref{thm:v39-all-smooth} proves that every smooth large-gauge class has
exact complete-packet phase one.  It does not prove
$\mathfrak B_{5,a}^{\rm mic}=0$.  V29 supplies explicit regulator-parameter
spectral topology which is not represented by a smooth gauge transformation.
The next theorem must therefore analyze this microscopic quotient directly,
not repeat the large-gauge anomaly calculation.
\end{firewall}

\section{V39 anti-circularity audit}
\label{sec:v39-audit}

\begin{enumerate}[leftmargin=2.4em]
\item \textbf{The free $H^3$ character was not inferred from V3 local anomaly
cancellation.}  The exact $SU(3)$ cancellation uses conjugate-representation
Kato twists; the exact $SU(2)$ cancellation uses pseudoreality plus even
multiplicity.
\item \textbf{The $SU(2)$ argument does not require computing the sign of each
$H^3$ generator.}  Pseudoreality quantizes every fundamental twist to
$\pm1$, and the complete packet contains twelve identical weak fundamentals.
\item \textbf{The colour argument is genuinely finite-regulator.}  The V2
Wilson kernels and Kato transports in conjugate representations are related by
one antiunitary conjugation, so their twists are complex conjugates before any
continuum limit is taken.
\item \textbf{The native $H^3$ loops are actual smooth admissible lattice
paths.}  The interpolation $A_t=tg^{-1}dg$ has links $I+O(a)$ and plaquettes
$I+O(a^2)$, so the V2 Wilson gap survives at sufficiently small $a$.
\item \textbf{Higgs/Yukawa data are not used to manufacture the measure
phase.}  They are gauge-covariant finite action data; the large-gauge phase is
the overlap measure-line transition and therefore depends only on the chiral
representation packet along the pure gauge-factor loop.
\item \textbf{The remaining defect is not declared zero.}  V39 explicitly
passes to the microscopic quotient, where V29 angular/spectral classes can
survive after every smooth gauge-transform class has been removed.
\end{enumerate}

\section{V39 status ledger}
\label{sec:v39-status}

\subsection*{Proved in V39}
\begin{enumerate}[leftmargin=2.2em]
\item Four smooth native $SU(3)$ and four smooth native $SU(2)$ loops realize
the basis of the V38 free
$H^3(T^4;\mathbb Z^2)\cong\mathbb Z^8$ component group.  The paths
$A_t=tg^{-1}dg$ are V2-admissible at sufficiently fine cutoff and have zero
V37 $H^1$ winding.
\item For every admissible closed loop, the overlap/Kato twist in the
complex-conjugate representation is the complex conjugate of the original
twist.  Hence a representation paired with its conjugate has exact unit phase
at finite cutoff, with all retained source/BV descendants.
\item The Standard-Model colour packet is exactly such a paired system: per
generation the two colour fundamentals in $Q_L$ pair with the two
antifundamentals $u_R^c,d_R^c$.  Therefore every pure-$SU(3)$ large-gauge loop
has exact complete-packet phase one.
\item The fundamental weak overlap twist of every admissible closed pure-$SU(2)$
loop is $\pm1$ by pseudoreality.  The complete Standard Model has twelve weak
fundamental copies, so every pure-$SU(2)$ large-gauge loop has exact phase one,
not only the V38 Witten representative.
\item Consequently the complete exact character is trivial on the entire free
$H^3(T^4;\mathbb Z^2)\cong\mathbb Z^8$ sector.
\item Combining V37 electromagnetic $H^1$ triviality, V39 free-$H^3$
triviality and V38 Witten-$H^4$ triviality shows that the exact complete
Standard-Model phase is one on the entire smooth large-gauge transformation
shadow of the periodic faithful carrier.
\item The residual exact character therefore factors through a genuinely
microscopic quotient
$\mathfrak B_{5,a}^{\rm mic}$ in which smooth large-gauge loops are removed.
V29 angular/Wilson-operator loops are the first explicit candidate classes in
that quotient.
\end{enumerate}

\subsection*{Inherited}
\begin{enumerate}[leftmargin=2.2em]
\item V1--V7: source-complete coefficientwise ESM continuum, complete local
anomaly cancellation, common BV restoration and exact loop-order projectivity.
\item V8--V24: regulator universality, controlled zero-mode/module
classification and the V2-consistent covariant soft norm-resolvent theorem.
\item V25--V29: real/Pfaffian orientation, accidental and Wilson index
crossings, higher-rank wall lines and nonsplit angular spectral topology.
\item V30--V36: faithful geometric cancellation mechanism, exact determinant
eta/integrability character, native five-dimensional realization,
coefficientwise Chern--Simons descent and the native fibered-bordism defect
group.
\item V37--V38: nontrivial electromagnetic defect winding with exact phase one,
finite-overlap Witten sign, exact cancellation for twelve weak doublets, and the
$\mathbb Z^8\oplus\mathbb Z_2$ smooth zero-$H^1$ Postnikov shadow.
\end{enumerate}

\subsection*{Conditional}
\begin{enumerate}[leftmargin=2.2em]
\item The explicit $H^3$ representatives are sufficiently fine
discretizations of fixed smooth continuum maps; rough lattice large-gauge
representatives outside the V2 admissible Wilson-gap component are not covered.
\item A complete full-ESM path must transport Higgs/coframe/spectator data in an
endpoint-compatible protected manner.  The exact measure-line cancellations
proved here do not depend on the details of that finite action interpolation.
\item The V38 Postnikov classification concerns the smooth large-gauge shadow.
Purely microscopic regulator loops and V29 angular classes are not classified
by $H^1\oplus H^3\oplus H^4$.
\item Positive IR/reference and the V2-consistent smooth-background hypotheses
remain in force.
\end{enumerate}

\subsection*{Open beyond V39}
\begin{enumerate}[leftmargin=2.2em]
\item Compute the exact faithful Standard-Model character on the microscopic
quotient $\mathfrak B_{5,a}^{\rm mic}$ after every smooth large-gauge class has
been removed.  The V29 angular/Hopf spectral bundle is the first explicit
candidate.
\item Determine whether V30's complete-packet angular Berry cancellation can be
proved natively for these microscopic loops without assuming the stronger
phase-faithful geometric lift corrected in V31.
\item If an angular class survives, compute its image under the V16 continuum
holonomy map and decide whether it is an exact finite-cutoff microscopic phase
or lies entirely in the continuum-invisible kernel.
\item The independent stronger problems remain: positive-IR removal, analytic
finite-coupling RG trajectories, perturbative summability and nonperturbative
ESM completion.
\end{enumerate}

\subsection*{Exact next load-bearing theorem}
\begin{center}
\fbox{\parbox{0.92\textwidth}{
\textbf{Purely microscopic/angular exact-character theorem.}
V37--V39 show that every smooth large-gauge component of the faithful periodic
Standard-Model carrier lies in the exact determinant/Pfaffian phase kernel.
Pass to
$\mathfrak B_{5,a}^{\rm mic}$, the quotient by those smooth gauge-transform
classes.  Construct a native representative of the V29 Pauli/Hopf angular
spectral class inside the complete Standard-Model regulator packet and compute
its exact V33/V34 character.  Prove directly that complete-packet
representation pairing trivializes the angular Berry/wall line, or exhibit a
nontrivial microscopic phase which is invisible to all smooth large-gauge
invariants.  This is now the first remaining exact global phase gate; repeating
$H^1$, $H^3$ or Witten-anomaly calculations cannot affect it.}}
\end{center}

\section*{Append-only continuation marker: V39}
\addcontentsline{toc}{section}{Append-only continuation marker: V39}
\textbf{End of V39.}  The complete V1--V38 body above is retained verbatim.
V39 constructs the eight smooth free-$H^3$ native large-gauge generators and
computes their exact complete-packet phases without relying on perturbative
anomaly descent.  Colour representations cancel by exact conjugate-representation
Kato duality, while every weak fundamental loop has a pseudoreal sign whose
twelfth power is one.  Hence the complete Standard-Model character is trivial
on all of $H^3(T^4;\mathbb Z^2)\cong\mathbb Z^8$.  Together with V37's
$H^1$ electromagnetic result and V38's Witten $H^4$ result, this exhausts the
smooth large-gauge transformation shadow: every such class lies in the exact
phase kernel.  The remaining exact problem is genuinely microscopic and is
captured by the quotient $\mathfrak B_{5,a}^{\rm mic}$, whose first explicit
candidate classes are the V29 angular/Wilson-operator spectral loops.


\clearpage
\section{V40 continuation: native Pauli embedding, microscopic Picard class, and the limit of the loop-character picture}
\label{sec:v40-continuation}

V39 removed every smooth large-gauge transformation component from the exact
complete Standard-Model phase problem and proposed the V29 Pauli/Hopf angular
family as the first candidate in the residual microscopic quotient.  Going
upstream reveals that two different kinds of global microscopic data were being
mixed.  The V37--V39 defect group is a one-dimensional loop invariant after the
relevant line has been flattened.  By contrast, the Pauli/Hopf phenomenon is
already a two-dimensional Picard/Chern obstruction: its positive spectral
bundle can have nonzero first Chern class even though the angular parameter
sphere is simply connected.

The present continuation makes that distinction exact and, at the same time,
constructs an explicit native finite-cutoff realization of the Pauli block
inside the common three-generation V2 Wilson carrier.  The construction uses
only an irrelevant finite-range generation-space deformation, preserves exact
gauge/spin covariance and the V2 physical principal symbol, and acts identically
on every Standard-Model representation block.  Its complete angular line is
\emph{not} cancelled by Standard-Model representation pairing: for the minimal
complex chiral packet it has first Chern number $-30$.

There is, however, an equally important limit to the construction.  On a
refining translation-invariant torus, the angular gap protecting this isolated
ultraviolet cone shrinks with the momentum spacing.  Thus V40 proves an exact
finite-cutoff microscopic Picard obstruction and a correction to the V39 loop
interpretation; it does not yet construct a cutoff-uniform V24 protected Hopf
sphere.

\subsection{A native irrelevant profile which selects one Wilson corner}
\label{sec:v40-profile}

Retain the V2 flat/reference notation
\begin{equation}
 c_\mu(x):=1-\cos x_\mu,
 \qquad
 L_\mu^{\mathcal U}
 =I-\frac12\bigl(T_\mu^{\mathcal U}+(T_\mu^{\mathcal U})^*\bigr).
 \label{eq:v40-cmu}
\end{equation}
Choose the Wilson corner
\begin{equation}
 x_*=(\pi,0,0,0).
 \label{eq:v40-corner}
\end{equation}
At the flat identity-link background define the scalar trigonometric profile
\begin{equation}
 \boxed{
 r(x)
 :=\frac{c_1(x)}2
   \prod_{\nu=2}^4
   \left(1-\frac{c_\nu(x)}2\right)
 =\frac{1-\cos x_1}{2}
  \prod_{\nu=2}^4\frac{1+\cos x_\nu}{2}.}
 \label{eq:v40-r-symbol}
\end{equation}
Then
\begin{equation}
 0\le r(x)\le1,
 \qquad
 r(x_*)=1,
 \qquad
 r(x)=\frac{x_1^2}{4}+O(|x|^4)
 \quad(x\to0).
 \label{eq:v40-r-properties}
\end{equation}

\begin{definition}[Native corner profile]
\label{def:v40-native-profile}
Put
\begin{equation}
 \mathcal P_a^{\mathcal U}
 :=\frac{L_1^{\mathcal U}}2
   \prod_{\nu=2}^4
   \left(I-\frac{L_\nu^{\mathcal U}}2\right)
 \label{eq:v40-P-profile}
\end{equation}
in the displayed order, and define its Hermitian symmetrization
\begin{equation}
 \boxed{
 R_a^{\mathcal U}
 :=\frac12
 \left(
  \mathcal P_a^{\mathcal U}
  +(\mathcal P_a^{\mathcal U})^*
 \right).}
 \label{eq:v40-R-profile}
\end{equation}
\end{definition}

\begin{lemma}[Locality, covariance and soft irrelevance]
\label{lem:v40-profile-properties}
The profile $R_a^{\mathcal U}$ is Hermitian, exactly gauge/spin covariant and
has fixed finite range at most four lattice steps.  Every prescribed fixed
source derivative has the same finite-range property.  At the flat reference
background its Fourier symbol is exactly \eqref{eq:v40-r-symbol}.  Hence on the
physical branch
\begin{equation}
 R_a^{\mathcal U}(p)=O(a^2|p|^2)
 \label{eq:v40-profile-soft}
\end{equation}
with every fixed physical momentum/source derivative.
\end{lemma}

\begin{proof}
Each $L_\mu^{\mathcal U}$ is Hermitian, range one and transforms by unitary
similarity.  A product of four such factors has range at most four and the same
covariance; averaging with its adjoint gives Hermiticity.  At identity links the
$L_\mu$ commute and have symbol $c_\mu(x)$, giving
\eqref{eq:v40-r-symbol}.  The Taylor expansion in
\eqref{eq:v40-r-properties} proves \eqref{eq:v40-profile-soft}.  Fixed source
derivatives are finite Leibniz sums of the same one-edge factors.
\end{proof}

\subsection{A common two-generation Pauli deformation of the V2 kernel}
\label{sec:v40-generation-Pauli}

Let $\mathcal G\cong\mathbb C^3$ be generation space.  Fix the first-two-
generation summand $\mathcal G_{12}\cong\mathbb C^2$ and Pauli matrices
$\tau_1,\tau_2,\tau_3$ on it.  Extend them by zero on the third generation.
For $\mathbf y\in\mathbb R^3$ put
\begin{equation}
 M(\mathbf y)
 :=\left(-I_2+\mu\,\mathbf y\cdot\boldsymbol\tau\right)
   \oplus 0,
 \qquad 0<\mu<\frac12.
 \label{eq:v40-generation-matrix}
\end{equation}
This matrix acts only on generation space and therefore commutes with every
internal gauge representation and local spin action.

\begin{definition}[Common Pauli--Wilson completion]
\label{def:v40-Pauli-Wilson}
For every Standard-Model chiral multiplet use the \emph{same} microscopic
completion
\begin{equation}
 \boxed{
 K_{W,a}(\mathbf y)
 :=K_{W,a}^{(0)}
   +R_a^{\mathcal U}\gamma_5\otimes M(\mathbf y),
 \qquad
 H_{W,a}(\mathbf y):=\gamma_5K_{W,a}(\mathbf y).}
 \label{eq:v40-Pauli-Wilson}
\end{equation}
The generation factor is understood to be tensored with the corresponding
fixed internal representation space.
\end{definition}

\begin{proposition}[Native admissibility and unchanged physical principal jet]
\label{prop:v40-Pauli-admissible}
The family \eqref{eq:v40-Pauli-Wilson} is finite range, Hermitian and exactly
gauge/spin covariant.  At every fixed value of $\mathbf y$ in a compact set,
the physical overlap/Dirac principal symbol is the V2 one; the new term enters
only at the irrelevant order
\begin{equation}
 \delta D_{\rm ov}(p)=O(a|p|^2)
 \label{eq:v40-irrelevant-overlap}
\end{equation}
on the physical branch, with every retained fixed source derivative.
\end{proposition}

\begin{proof}
The structural statements follow from Lemma~\ref{lem:v40-profile-properties},
$[R_a^{\mathcal U},\gamma_5]=0$ and Hermiticity of $M(\mathbf y)$.  Since
\[
 \gamma_5
 R_a^{\mathcal U}\gamma_5
 =R_a^{\mathcal U},
\]
the Hermitian Wilson kernel changes by
\begin{equation}
 H_{W,a}(\mathbf y)-H_{W,a}^{(0)}
 =R_a^{\mathcal U}\otimes M(\mathbf y).
 \label{eq:v40-H-shift}
\end{equation}
Lemma~\ref{lem:v40-profile-properties} makes this $O(a^2|p|^2)$ on the
physical branch.  The protected V2 sign functional calculus then changes the
overlap sign by the same dimensionless order, and the prefactor $a^{-1}$ in
$D_{\rm ov}$ gives \eqref{eq:v40-irrelevant-overlap}.  The differentiated
statement follows from the V2 resolvent/source calculus.
\end{proof}

\subsection{Exact finite-cutoff Pauli cone}
\label{sec:v40-exact-cone}

At the flat identity-link/reference-coframe background, V2 gives
\begin{equation}
 H_{W,a}^{(0)}(x)^2=R_W(x)^2I,
 \qquad
 R_W(x)^2
 =1+2\sum_{\mu<\nu}c_\mu(x)c_\nu(x)
 \ge1.
 \label{eq:v40-RW}
\end{equation}
At $x_*$ one has
\begin{equation}
 H_{W,a}^{(0)}(x_*)=\gamma_5,
 \qquad
 r(x_*)=1.
 \label{eq:v40-corner-H}
\end{equation}

\begin{lemma}[The undeformed divisor is isolated at the selected corner]
\label{lem:v40-unique-zero}
At $\mathbf y=0$, on the first-two-generation block the flat symbol is
\begin{equation}
 H_W(x,0)=H_W^{(0)}(x)-r(x)I.
 \label{eq:v40-y0-symbol}
\end{equation}
It has a zero eigenvalue if and only if $x=x_*$ and the spinor lies in the
positive-chirality subspace $S_+:=\Ran\frac{1+\gamma_5}{2}$.
\end{lemma}

\begin{proof}
The eigenvalues of $H_W^{(0)}(x)$ are $\pm R_W(x)$.  The negative branch of
\eqref{eq:v40-y0-symbol} is $-R_W-r<0$.  A zero on the positive branch requires
$R_W(x)=r(x)$.  Since $R_W\ge1$ and $0\le r\le1$, equality forces
$R_W=r=1$.  From \eqref{eq:v40-RW}, $R_W=1$ means
$c_\mu c_\nu=0$ for every $\mu<\nu$, so at most one $c_\mu$ is nonzero.
Equation \eqref{eq:v40-r-symbol} has value one only when
$c_1=2$ and $c_2=c_3=c_4=0$, namely at $x_*$.  There
$H_W^{(0)}=\gamma_5$ and subtracting the identity kills precisely $S_+$.
\end{proof}

\begin{theorem}[Exact native Pauli crossing block at finite cutoff]
\label{thm:v40-native-Pauli-block}
Fix one finite periodic lattice for which $x_*$ belongs to the momentum set.
There exists $\rho_a>0$ such that on the sphere
$|\mathbf y|=\rho_a$ the flat family
\eqref{eq:v40-Pauli-Wilson} is Wilson-gapped and its only resolved crossing
cluster is the one issuing from $(x_*,\mathbf y=0)$.  On a chiral multiplet of
internal complex dimension $d_R$, the crossing form is
\begin{equation}
 \boxed{
 A_R(\mathbf y)
 =\mu\,I_{S_+}\otimes
   (\mathbf y\cdot\boldsymbol\tau)
   \otimes I_R.}
 \label{eq:v40-crossing-form}
\end{equation}
Consequently its positive angular spectral bundle is
\begin{equation}
 \boxed{
 E_{+,R}
 \cong
 S_+\otimes L_{\rm H}\otimes\mathbb C^{d_R},}
 \label{eq:v40-Eplus-R}
\end{equation}
where $L_{\rm H}\to S^2$ is the Hopf line of
Proposition~\ref{prop:v29-Pauli}.
\end{theorem}

\begin{proof}
At $x_*$, equations \eqref{eq:v40-corner-H} and
\eqref{eq:v40-generation-matrix} give, on the first two generations,
\[
 H_W(x_*,\mathbf y)
 =\gamma_5-I_2+\mu\,\mathbf y\cdot\boldsymbol\tau.
\]
On $S_+$ this is exactly
$\mu\,\mathbf y\cdot\boldsymbol\tau$, while on the negative-chirality
subspace it is
$-2I+\mu\,\mathbf y\cdot\boldsymbol\tau$, which stays invertible for
$\mu|\mathbf y|<1$.  The third generation is undeformed at the selected
corner.

By Lemma~\ref{lem:v40-unique-zero}, all other finite-lattice momentum/spin
blocks are nonzero at $\mathbf y=0$.  There are finitely many of them, so their
minimum singular value is a positive number $\delta_a$.  Continuity in
$\mathbf y$ gives a radius $\rho_a>0$ on which none of those complementary
blocks closes its gap.  The displayed crossing form and
\eqref{eq:v40-Eplus-R} then follow from the V29 Pauli calculation and
$\dim_{\mathbb C}S_+=2$.
\end{proof}

\subsection{The complete Standard-Model angular line does not cancel}
\label{sec:v40-SM-Chern}

For one minimal Standard-Model generation, the dimensions of the complex
internal representation blocks are
\begin{equation}
 6,
 \quad 3,
 \quad 3,
 \quad 2,
 \quad 1,
 \label{eq:v40-SM-dims}
\end{equation}
for $Q_L,u^c,d^c,L,e^c$, respectively.  Their sum is
\begin{equation}
 \boxed{d_{\rm SM}^{(1)}=15.}
 \label{eq:v40-fifteen}
\end{equation}
The deformation \eqref{eq:v40-Pauli-Wilson} acts on the first two generation
slots, so the third generation is a spectator for the angular line.

\begin{theorem}[Noncancellation of the native Hopf line in the complete packet]
\label{thm:v40-SM-noncancellation}
For the native family of Theorem~\ref{thm:v40-native-Pauli-block}, the complete
minimal complex Standard-Model positive crossing bundle satisfies
\begin{equation}
 E_+^{\rm SM}
 \cong
 \bigoplus_R
 \left(
  S_+\otimes L_{\rm H}\otimes\mathbb C^{d_R}
 \right),
 \label{eq:v40-Eplus-SM}
\end{equation}
and hence
\begin{equation}
 \boxed{
 c_1(E_+^{\rm SM})
 =2\sum_Rd_R
 =30.}
 \label{eq:v40-c1-E-SM}
\end{equation}
The complete angular determinant line therefore has
\begin{equation}
 \boxed{
 c_1(\mathcal L_{\rm ang}^{\rm SM})=-30\ne0.}
 \label{eq:v40-c1-line-SM}
\end{equation}
In particular, Standard-Model representation pairing does not cancel this
microscopic regulator-parameter Hopf class.
\end{theorem}

\begin{proof}
For each representation block, Theorem~\ref{thm:v40-native-Pauli-block} gives a
constant $2d_R$-dimensional multiplicity tensoring the same Hopf line.  Therefore
\[
 \det E_{+,R}
 \cong
 L_{\rm H}^{\otimes2d_R}
\]
up to a constant one-dimensional factor.  First Chern classes add under direct
sum and tensor product, and $c_1(L_{\rm H})=1$ by
Proposition~\ref{prop:v29-Pauli}.  Summing the dimensions
\eqref{eq:v40-SM-dims} gives \eqref{eq:v40-c1-E-SM}.  V29 identifies the
angular contribution to the chiral determinant line with
$(\det E_+)^*$, which reverses the sign and proves
\eqref{eq:v40-c1-line-SM}.
\end{proof}

\begin{remark}[Optional neutral sector]
\label{rem:v40-neutral}
An optional independently real neutral Majorana/Pfaffian sector cannot cancel a
nonzero complex first Chern class.  If instead one adds a complex right-handed
neutrino representation and applies the same generation deformation to it, the
integer $30$ is replaced by $32$; it remains nonzero.  Thus the noncancellation
statement is not an artefact of omitting the optional neutral sector.
\end{remark}

\begin{proposition}[Why the V39 conjugate-representation cancellation does not apply]
\label{prop:v40-no-conjugate-cancel}
The exact V39 identity
$\mathcal T_{\bar R}=\overline{\mathcal T_R}$ for a pure gauge loop does not
force cancellation of \eqref{eq:v40-c1-line-SM}.  Along the family
\eqref{eq:v40-Pauli-Wilson}, the regulator parameter is inserted as the
\emph{same} Hermitian generation matrix $M(\mathbf y)$ in both $R$ and
$\bar R$ blocks.  The crossing bundle therefore tensors the same Hopf line with
both representation spaces.  Complex conjugation would instead replace the
Pauli coefficient $\tau_2$ by $-\tau_2$ and pull back the Hopf line by an
orientation-reversing angular map.  Imposing that additional conjugation rule
is an extra regulator-reality constraint, not a consequence of the declared
V1--V9 locality, covariance and Hermiticity gates.
\end{proposition}

\begin{proof}
At the flat gauge background the gauge representation enters the native Wilson
kernel only as an identity multiplicity.  Equation
\eqref{eq:v40-crossing-form} is therefore literally the same matrix pencil in
$R$ and $\bar R$.  Their Chern contributions have the same sign.  The V39
antiunitary proof applies when the entire one-particle family in $\bar R$ is the
complex conjugate of that in $R$; the common regulator deformation used here
does not conjugate the generation parameter.  Hence the hypotheses of the V39
duality statement are not met.
\end{proof}

\begin{firewall}[A stronger regulator-reality subcategory could remove the example]
V40 proves noncancellation in the microscopic regulator class actually declared
by the lineage: finite-range/exponentially local, Hermitian, gauge/spin
covariant and V2-soft-consistent.  If one additionally requires every
microscopic flavour/regulator deformation to respect a fixed antiunitary
reality relation between conjugate representation blocks, the family
\eqref{eq:v40-Pauli-Wilson} with its complex $\tau_2$ parameter is excluded.
Such a reality axiom would be a new restriction on the regulator class and must
be stated and justified separately; it cannot be imported retroactively from
pure-gauge conjugation identities.
\end{firewall}

\subsection{An exact finite-cutoff Berry phase}
\label{sec:v40-Berry-phase}

V29 gives the Hopf curvature
\begin{equation}
 \operatorname{Tr}(P_+dP_+\wedge dP_+)
 =\frac{i}{2}\,\Omega_{S^2}.
 \label{eq:v40-Hopf-curvature}
\end{equation}
Therefore the complete determinant line of
Theorem~\ref{thm:v40-SM-noncancellation} has curvature
\begin{equation}
 \boxed{
 \mathcal F_{\rm ang}^{\rm SM}
 =-15i\,\Omega_{S^2}.}
 \label{eq:v40-SM-curvature}
\end{equation}

\begin{corollary}[Explicit nontrivial raw angular holonomy]
\label{cor:v40-raw-holonomy}
Let $\Sigma\subset S^2$ be any spherical cap with solid angle
\begin{equation}
 \int_\Sigma\Omega_{S^2}=\frac{\pi}{15},
 \label{eq:v40-cap-area}
\end{equation}
and let $\gamma=\partial\Sigma$.  In the native Kato/Berry connection of the
complete angular determinant line,
\begin{equation}
 \boxed{
 \operatorname{Hol}_{\rm raw}^{\rm SM}(\gamma)=-1.}
 \label{eq:v40-raw-minus}
\end{equation}
\end{corollary}

\begin{proof}
The line is trivial over the disk $\Sigma$, so Stokes' theorem applies in any
local frame.  Equation \eqref{eq:v40-SM-curvature} gives
\[
 \operatorname{Hol}_{\rm raw}^{\rm SM}(\partial\Sigma)
 =\exp\left(\int_\Sigma\mathcal F_{\rm ang}^{\rm SM}\right)
 =\exp(-i\pi)=-1.
\]
Changing the overall sign convention for the unitary Berry connection replaces
the exponent by its complex conjugate and leaves the value $-1$ unchanged.
\end{proof}

\subsection{Going upstream: this is an $H^2$ obstruction, not a defect-group character}
\label{sec:v40-Picard-vs-loop}

The Pauli sphere is simply connected:
\begin{equation}
 \pi_1(S^2)=0.
 \label{eq:v40-pi1-S2}
\end{equation}
Thus the loop in Corollary~\ref{cor:v40-raw-holonomy} is null-homotopic on the
resolved angular sphere.  Its nontrivial holonomy is generated by curvature and
changes under a homotopy which sweeps nonzero Berry flux.  It is therefore not
a homomorphism of the fundamental group.

\begin{theorem}[Microscopic Picard class precedes the flat loop character]
\label{thm:v40-Picard-first}
Let $B_{\rm reg}$ be a compact resolved microscopic regulator-parameter family
and let $\mathcal L_{\rm reg}^{\rm SM}$ be its complete determinant line.  The
first global obstruction to a single scalar regulator phase is
\begin{equation}
 \boxed{
 \mathfrak P_a^{\rm mic}
 :=c_1(\mathcal L_{\rm reg}^{\rm SM})
 \in H^2(B_{\rm reg};\mathbb Z).}
 \label{eq:v40-mic-Picard}
\end{equation}
Only after $\mathfrak P_a^{\rm mic}=0$ and a global line trivialization has
been chosen does the residual flat phase define a class in
\begin{equation}
 H^1(B_{\rm reg};U(1)_{\rm const})
 \cong
 \operatorname{Hom}(\pi_1(B_{\rm reg}),U(1)).
 \label{eq:v40-flat-H1}
\end{equation}
For the native Pauli family of V40,
\begin{equation}
 \boxed{
 \mathfrak P_a^{\rm mic}=-30\,[S^2]^*\ne0,}
 \label{eq:v40-Picard-Pauli}
\end{equation}
so there is no globally flattened line and no associated flat defect-group
character on the whole angular sphere.
\end{theorem}

\begin{proof}
A global nonzero scalar phase frame is precisely a trivialization of the
complex determinant line, whose primary topological obstruction is its first
Chern class.  If the class vanishes, the line is topologically trivial and one
may compare two connections on that fixed trivial line; subtracting the local
curvature primitive leaves a flat connection whose gauge-equivalence class is
$H^1(B;U(1)_{\rm const})$.  Conversely a nonzero first Chern class forbids any
global flat unitary connection on the line.  Equation
\eqref{eq:v40-Picard-Pauli} is Theorem~\ref{thm:v40-SM-noncancellation}.
\end{proof}

\begin{corollary}[Correction to the V39 angular-target interpretation]
\label{cor:v40-correct-V39}
The nonzero Chern class of the V29 Pauli/Hopf sphere does not by itself define a
nontrivial element of the V37--V39 loop defect group
$\mathfrak B_{5,a}^{\rm mic}$.  The latter is a quotient of loop homotopy data
after the relevant determinant line has been flattened.  The Pauli/Hopf class
belongs one level earlier, to the microscopic Picard group of regulator
families.  Any claim that it survives as a flat loop character requires an
additional topologically trivial family or a separate loop construction.
\end{corollary}

\begin{proof}
Equation \eqref{eq:v40-pi1-S2} makes every angular loop homotopically trivial on
the Pauli sphere, whereas
\eqref{eq:v40-Picard-Pauli} is nonzero.  The raw holonomy of
Corollary~\ref{cor:v40-raw-holonomy} is therefore curvature holonomy, not a
fundamental-group character.  This is exactly the distinction asserted.
\end{proof}

\subsection{Why the explicit translation-invariant sphere is not V24-uniform}
\label{sec:v40-gap-scaling}

The preceding finite-cutoff construction is exact, but the protected angular
radius obtained from Theorem~\ref{thm:v40-native-Pauli-block} is not uniform
under a standard continuum refinement at fixed physical volume.

Let the periodic lattice have $N$ cells in the first direction.  The nearest
momentum to $x_*$ along the $x_1$ axis is
\begin{equation}
 x_N=\left(\pi-\frac{2\pi}{N},0,0,0\right).
 \label{eq:v40-near-corner}
\end{equation}
Along this axis, the V2 identity gives
\begin{equation}
 R_W(x_N)=1,
 \label{eq:v40-R-axis}
\end{equation}
while
\begin{equation}
 r(x_N)
 =\frac{1+\cos(2\pi/N)}2
 =1-\frac{\pi^2}{N^2}+O(N^{-4}).
 \label{eq:v40-r-axis}
\end{equation}
Hence the complementary gap at $\mathbf y=0$ obeys
\begin{equation}
 \boxed{
 \delta_a
 \le
 R_W(x_N)-r(x_N)
 =\frac{\pi^2}{N^2}+O(N^{-4}).}
 \label{eq:v40-gap-N2}
\end{equation}
At fixed physical period, $N\asymp a^{-1}$, so
\begin{equation}
 \delta_a=O(a^2)
 \label{eq:v40-gap-a2}
\end{equation}
as an upper bound.  Any angular radius which keeps all neighboring momentum
blocks gapped must therefore satisfy
\begin{equation}
 \rho_a\lesssim \delta_a/\mu=O(a^2).
 \label{eq:v40-rho-a2}
\end{equation}

\begin{proposition}[Finite-cutoff Picard class without a uniform V24 wall gap]
\label{prop:v40-not-V24-uniform}
The native family of Definition~\ref{def:v40-Pauli-Wilson} gives an exact
finite-cutoff angular Picard class with
$c_1=-30$, but the simple translation-invariant realization above does not
supply a cutoff-uniform Wilson/complement gap on a fixed-radius regulator
sphere.  Hence it is not, by itself, a counterexample to the V24 smooth
protected soft theorem or to coefficientwise regulator universality.
\end{proposition}

\begin{proof}
The exact finite-cutoff Chern calculation is
Theorem~\ref{thm:v40-SM-noncancellation}.  Estimate
\eqref{eq:v40-rho-a2} shows that the radius on which the isolated cone is the
only Wilson crossing shrinks to zero under refinement.  V24 requires one
cutoff-uniform protected soft/high decomposition with positive spectral
margins; that hypothesis is not supplied by this particular family.
\end{proof}

\begin{firewall}[No contradiction with the coefficientwise continuum]
V40 exhibits a genuine finite-cutoff microscopic regulator Picard class in the
local/covariant V2-soft-consistent grammar.  Its protected angular gap shrinks
under the elementary translation-invariant refinement, so it need not survive
as a cofinal protected continuum family.  V16/V32 coefficientwise regulator
universality is therefore unaffected.  Conversely, V40 proves that Standard-
Model representation pairing alone cannot be invoked to kill arbitrary
microscopic regulator-parameter Berry topology.
\end{firewall}

\subsection{Revised global hierarchy after V40}
\label{sec:v40-hierarchy}

The exact global microscopic data exposed by V29--V40 now occur in the order
\begin{equation}
 \boxed{
 \begin{array}{c}
 \text{microscopic Picard/Chern class }H^2\\
 \Downarrow\\
 \text{flat determinant character }H^1(U(1))\\
 \Downarrow\\
 \text{controlled native path / defect-group data}.
 \end{array}}
 \label{eq:v40-global-hierarchy}
\end{equation}
The lower levels are meaningful globally only after the higher obstruction has
been removed or restricted away.

\begin{theorem}[V40 microscopic angular classification at the proved scope]
\label{thm:v40-summary}
Within the microscopic regulator class defined only by the lineage's declared
locality, Hermiticity, gauge/spin covariance and V2 soft-principal-symbol
conditions:
\begin{enumerate}[label=\textup{(A\arabic*)},leftmargin=3em]
\item a native finite-cutoff nonsplit Pauli/Hopf crossing can be embedded in the
common three-generation Wilson carrier;
\item the complete minimal Standard-Model angular determinant line need not
cancel and in the explicit family has $c_1=-30$;
\item this nonzero class is an $H^2$ Picard obstruction, not a flat loop
character;
\item the explicit translation-invariant realization does not have a
cutoff-uniform angular gap and hence does not yet define a cofinal V24 protected
regulator family.
\end{enumerate}
\end{theorem}

\begin{proof}
Items (A1)--(A4) are respectively
Theorem~\ref{thm:v40-native-Pauli-block},
Theorem~\ref{thm:v40-SM-noncancellation},
Theorem~\ref{thm:v40-Picard-first}, and
Proposition~\ref{prop:v40-not-V24-uniform}.
\end{proof}

\section{V40 anti-circularity audit}
\label{sec:v40-audit}

\begin{enumerate}[leftmargin=2.4em]
\item \textbf{The Pauli block was embedded natively rather than appended as an
abstract finite matrix.}  The profile $R_a^{\mathcal U}$ is finite range,
gauge/spin covariant and physically irrelevant to first order, and the
generation deformation is one common kinetic regulator term across all
Standard-Model representation blocks.
\item \textbf{Representation pairing was tested at the actual regulator
parameter.}  The V39 conjugate-representation cancellation concerns conjugate
gauge families.  V40's common generation parameter tensors the same Hopf line
with $R$ and $\bar R$, so their Chern contributions add unless an extra
antiunitary regulator-reality axiom is imposed.
\item \textbf{A nonzero Berry holonomy was not promoted to a flat global
character.}  The Pauli sphere is simply connected; its nontrivial holonomy is
curvature holonomy and the true invariant is $c_1\in H^2$.
\item \textbf{The finite-cutoff example was not promoted to a V24-uniform
continuum family.}  The protected angular radius shrinks at least as fast as
the neighboring UV momentum separation squared in the explicit
translation-invariant realization.
\item \textbf{No contradiction with V30 was asserted.}  Theorem V30 is
conditional on a phase-faithful geometric lift; the nonzero V40 Chern class
proves precisely that such a lift cannot exist for this microscopic family
without additional compensating data or an extra regulator restriction.
\end{enumerate}

\section{V40 status ledger}
\label{sec:v40-status}

\subsection*{Proved in V40}
\begin{enumerate}[leftmargin=2.2em]
\item A fixed-range gauge/spin-covariant native profile with flat symbol
$r(x)=\frac{1-\cos x_1}{2}\prod_{\nu>1}\frac{1+\cos x_\nu}{2}$ selects the
Wilson corner $(\pi,0,0,0)$ while vanishing as $O((ap)^2)$ on the physical
branch.
\item Adding this profile times $\gamma_5$ and a common two-generation Pauli
matrix gives a finite-range Hermitian common Wilson completion which preserves
the V2 physical principal Dirac symbol but has an exact finite-cutoff nonsplit
Pauli crossing at the selected UV corner.
\item On a representation block of internal dimension $d_R$, the positive
crossing bundle is $S_+\otimes L_{\rm H}\otimes\mathbb C^{d_R}$ and contributes
$2d_R$ times the Hopf Chern class.
\item Summing the minimal Standard-Model representation dimensions
$6+3+3+2+1=15$ gives
$c_1(E_+^{\rm SM})=30$ and
$c_1(\mathcal L_{\rm ang}^{\rm SM})=-30$.  Hence complete Standard-Model
representation pairing does not generically cancel microscopic angular Berry
topology.
\item The complete angular Berry curvature is
$-15i\,\Omega_{S^2}$, so a cap of solid angle $\pi/15$ has exact raw determinant
holonomy $-1$.
\item The Pauli/Hopf obstruction is a microscopic Picard class in $H^2$, not a
flat $\pi_1$ character.  Thus the V39 suggestion that its nonzero Chern class
alone supplies a class in $\mathfrak B_{5,a}^{\rm mic}$ is corrected.
\item In the explicit translation-invariant realization the isolated-cone gap
above neighboring UV momenta is at most $O(a^2)$ at fixed physical volume, so
the angular radius required to isolate the cone shrinks and the construction is
not a cutoff-uniform V24 protected family.
\end{enumerate}

\subsection*{Inherited}
\begin{enumerate}[leftmargin=2.2em]
\item V1--V7: the source-complete coefficientwise ESM continuum, common native
Wilson profile, anomaly cancellation, BV restoration and exact order
projectivity.
\item V8--V24: microscopic regulator comparison, controlled zero-mode/module
classification and the V2-consistent soft/Feshbach theorem.
\item V25--V29: Pfaffian orientation, Wilson index walls, higher-rank graded
wall lines and the nonsplit Pauli/Hopf angular determinant bundle.
\item V30--V36: conditional complete-packet geometric cancellation, exact
integrability/eta character, local five-dimensional realization and the
fibered-bordism/flat-character separation.
\item V37--V39: electromagnetic, Witten and free-$H^3$ smooth large-gauge
classes all have exact complete Standard-Model phase one, leaving only genuinely
microscopic regulator data after the smooth gauge-transform shadow is removed.
\end{enumerate}

\subsection*{Conditional}
\begin{enumerate}[leftmargin=2.2em]
\item The explicit V40 Pauli sphere is an exact finite-cutoff construction.
Its protected angular radius is cutoff dependent and does not satisfy the
cutoff-uniform V24 gap hypothesis.
\item The noncancellation theorem assumes only the regulator axioms actually
declared in the lineage.  Imposing an additional antiunitary
flavour/regulator-reality symmetry can exclude the complex generation-space
Pauli family and must be analyzed as a separate restricted regulator class.
\item The minimal complex Standard-Model count is $15$ internal Weyl
components per generation type.  Optional independently real neutral Majorana
sectors remain governed by V25 and do not cancel the nonzero complex Chern
class.
\item Positive IR/reference and all fixed-order/source quantifiers of the
original continuum theorem remain in force.
\end{enumerate}

\subsection*{Open beyond V40}
\begin{enumerate}[leftmargin=2.2em]
\item Determine whether a nonzero microscopic regulator Picard class can be
realized with a \emph{cutoff-uniform} Wilson/complement gap inside the actual
V2-consistent finite-range or exponentially local regulator class.  This is a
stronger native embedding problem than the finite-cutoff construction proved
here.
\item Alternatively, prove that every uniformly protected V2-consistent
regulator family has trivial complete microscopic Picard class, perhaps under a
natural antiunitary regulator-reality condition.
\item Once the Picard class is zero, return to the V36--V39 flat-character
problem on the topologically trivial regulator component; the smooth gauge
shadow is already in its exact kernel.
\item The independent stronger problems remain: positive-IR removal, analytic
finite-coupling RG trajectories, perturbative summability and nonperturbative
ESM completion.
\end{enumerate}

\subsection*{Exact next load-bearing theorem}
\begin{center}
\fbox{\parbox{0.92\textwidth}{
\textbf{Uniform microscopic Picard embedding/no-go theorem.}
V40 constructs an exact native Pauli/Hopf regulator sphere and proves that the
complete minimal Standard-Model angular line has Chern number $-30$, but the
simple translation-invariant realization loses its isolated-cone gap under
continuum refinement.  Determine whether there exists a finite-range or
uniformly exponentially local, gauge/spin-covariant, V2-soft-consistent
microscopic regulator family with the same nonzero angular Picard class and a
cutoff-uniform Wilson/complement gap.  Either construct such a uniformly
protected local defect family, proving that a genuine microscopic $H^2$
regulator phase survives all complete-packet assembly, or prove a no-go theorem
showing that uniform V2 protection (possibly together with a justified native
antiunitary reality axiom) forces the complete microscopic Picard class to
vanish.  Only after this degree-two question is settled is the residual flat
loop character of V36--V39 the correct next global invariant.}}
\end{center}

\section*{Append-only continuation marker: V40}
\addcontentsline{toc}{section}{Append-only continuation marker: V40}
\textbf{End of V40.}  The complete V1--V39 body above is retained verbatim.
V40 constructs an explicit finite-range, gauge/spin-covariant and V2-soft-
consistent two-generation Pauli deformation of the common Wilson regulator.
At a selected ultraviolet corner it realizes the V29 Hopf crossing natively.
Because the same crossing line tensors every Standard-Model representation
block rather than its complex conjugate, the complete minimal packet has
$c_1(\mathcal L_{\rm ang}^{\rm SM})=-30$ and an explicit raw Berry loop with
phase $-1$; Standard-Model representation pairing does not cancel this generic
microscopic regulator-parameter topology.  Going upstream, V40 corrects the
V39 interpretation: the Hopf sphere is simply connected, so its primary
obstruction is a microscopic Picard class in $H^2$, not a flat defect-group
character.  The explicit translation-invariant cone is not cutoff-uniformly
protected, leaving as the next gate the existence or no-go theorem for a
uniformly gapped native realization of this nonzero Picard class.


\clearpage
\section{V41 continuation: uniform-protection no-go for microscopic regulator Picard classes}
\label{sec:v41-continuation}

V40 exhibited an exact finite-cutoff Pauli/Hopf regulator sphere with a
nonzero complete Standard-Model angular Chern class, but its isolated-cone
radius collapsed under continuum refinement.  The natural next question was
whether that loss of protection was merely an artifact of the elementary
translation-invariant construction, or whether a nonzero regulator-parameter
Picard class is fundamentally incompatible with the uniform V24 protected
class.

The answer is the latter at the scope actually used by the continuum theorem.
No additional antiunitary regulator-reality axiom is needed.  The load-bearing
input is instead the full V24 package: one common soft continuum
Dirac--Yukawa jet, fixed minimal Fredholm index, a uniform reduced singular
floor, and the covariant soft--hard estimates.  Uniformly on any compact
regulator-parameter base, V24 makes the native kernel/cokernel Riesz projectors
norm-close to those of one reference completion.  Once the norm distance is
strictly less than one, the V14 direct rotation trivializes the entire
regulator-parameter dependence globally.  Consequently the native Fredholm
determinant line is pulled back from the physical-background base and carries
no Chern class in a pure regulator direction.

This continuation proves that statement abstractly, gives a quantitative
converse for any attempted nonzero microscopic Picard class, and records an
independent translation-invariant symbol argument which explains the V40
$O(a^2)$ gap collapse directly.  The conclusion is deliberately narrower than
a no-go theorem for every local finite-cutoff regulator: V40 remains a valid
counterexample outside the uniform V24 protected class.

\subsection{Uniformly protected regulator families}
\label{sec:v41-family}

Let $B$ be a compact physical background/source panel and let $Y$ be a compact
connected regulator-parameter manifold.  Write
\begin{equation}
 \pi_B:B\times Y\longrightarrow B
 \label{eq:v41-piB}
\end{equation}
for the projection.  For every $y\in Y$ let
\begin{equation}
 D_{a,y}(b):E_a\longrightarrow F_a,
 \qquad
 L_{a,y}(b)=D_{a,y}(b)^*D_{a,y}(b)
 \label{eq:v41-family-D}
\end{equation}
be the native chiral/Yukawa family.  The cokernel formulation for negative
index uses $D_{a,y}D_{a,y}^*$ exactly as in V11--V24.

\begin{definition}[Uniform V24-protected regulator family]
\label{def:v41-uniform-family}
The family $D_{a,y}$ is \emph{uniformly V24-protected over $B\times Y$} if:
\begin{enumerate}[label=\textup{(U\arabic*)},leftmargin=3em]
\item the V1--V2 locality, covariance and fixed source-jet bounds are uniform
in $(b,y)$;
\item the Fredholm index is fixed and the V11 minimal kernel/cokernel ranks are
constant on $B\times Y$;
\item the nonzero singular spectrum has one common floor $\tau_*>0$;
\item every microscopic completion shares the same V24 soft continuum jet
(N5), and that continuum operator is independent of $y$;
\item the constants in Theorem~\ref{thm:v24-main} may be chosen uniformly on
$B\times Y$.
\end{enumerate}
The regulator parameter $y$ is an auxiliary microscopic coordinate; it is not
promoted to an additional physical source writer.
\end{definition}

Choose once and for all $y_*\in Y$.  Let
$p^K_{a,y}(b)$ and $p^C_{a,y}(b)$ denote the kernel and cokernel Riesz
projectors, with the zero projector used when the corresponding minimal rank
vanishes.

\begin{lemma}[Uniform Riesz closeness over the regulator base]
\label{lem:v41-uniform-Riesz}
For a uniformly V24-protected family there is $C$ independent of $a,b,y$ such
that
\begin{align}
 \sup_{(b,y)\in B\times Y}
 \|p^K_{a,y}(b)-p^K_{a,y_*}(b)\|
 &\le Ca^{1/10},
 \label{eq:v41-K-close}\\
 \sup_{(b,y)\in B\times Y}
 \|p^C_{a,y}(b)-p^C_{a,y_*}(b)\|
 &\le Ca^{1/10}.
 \label{eq:v41-C-close}
\end{align}
The same conclusion holds for every prescribed fixed source derivative.
\end{lemma}

\begin{proof}
Regard the completion label in Theorem~\ref{thm:v24-main} as the compact
parameter $y$.  Definition~\ref{def:v41-uniform-family} makes every constant
in that theorem uniform in $y$.  Apply
Corollary~\ref{cor:v24-common-Riesz} to the pair of completions $(y,y_*)$.
The kernel statement is exactly \eqref{eq:v41-K-close}; the cokernel statement
is its $DD^*$ version.  The source-derivative assertion is inherited from the
source-complete clause of Theorem~\ref{thm:v24-main}.
\end{proof}

\subsection{Global direct rotation and pullback of the zero-mode bundles}
\label{sec:v41-pullback}

For sufficiently small $a$, both bounds in
Lemma~\ref{lem:v41-uniform-Riesz} are strictly below one.  Therefore the V14
polar direct rotation is available simultaneously for every $(b,y)$.

\begin{theorem}[Regulator-parameter pullback theorem]
\label{thm:v41-pullback}
For all sufficiently small $a$, there exist source-smooth exponentially local
unitaries
\begin{equation}
 U_a^K(b,y),\qquad U_a^C(b,y)
 \label{eq:v41-direct-unitaries}
\end{equation}
with $U_a^{K,C}(b,y_*)=I$ and
\begin{align}
 p^K_{a,y}(b)
 &=U_a^K(b,y)p^K_{a,y_*}(b)U_a^K(b,y)^*,
 \label{eq:v41-K-conj}\\
 p^C_{a,y}(b)
 &=U_a^C(b,y)p^C_{a,y_*}(b)U_a^C(b,y)^*.
 \label{eq:v41-C-conj}
\end{align}
Consequently the kernel and cokernel bundles over $B\times Y$ satisfy
\begin{equation}
 \boxed{
 \mathcal K_a\simeq \pi_B^*\mathcal K_{a,*},
 \qquad
 \mathcal C_a\simeq \pi_B^*\mathcal C_{a,*}.}
 \label{eq:v41-bundle-pullback}
\end{equation}
The isomorphisms are controlled in the V13 sense and preserve every retained
source jet.
\end{theorem}

\begin{proof}
For two projections $p,q$ with $\|p-q\|<1$, V14 defines
\[
 X=qp+(I-q)(I-p),
 \qquad
 U=X(X^*X)^{-1/2}.
\]
Apply this formula pointwise to
$(p^K_{a,y_*},p^K_{a,y})$ and separately to the cokernel pair.  The uniform
strict norm bound keeps $X^*X$ uniformly positive.  V13 functional calculus
therefore makes the rotations exponentially local with all fixed source
derivatives.  Their smooth dependence on $(b,y)$ follows from the same
holomorphic functional calculus.  Restriction to $y=y_*$ gives the identity.
The displayed unitaries are global bundle isomorphisms from the pullbacks of
the reference bundles to the full families, proving
\eqref{eq:v41-bundle-pullback}.
\end{proof}

\begin{corollary}[All pure regulator characteristic classes vanish]
\label{cor:v41-characteristic}
For every characteristic class functor $c$ of the complex kernel/cokernel
bundles,
\begin{equation}
 \boxed{
 c(\mathcal K_a)=\pi_B^*c(\mathcal K_{a,*}),
 \qquad
 c(\mathcal C_a)=\pi_B^*c(\mathcal C_{a,*}).}
 \label{eq:v41-characteristic-pullback}
\end{equation}
In particular every component containing a positive cohomological degree in the
regulator factor $Y$ vanishes.  At fixed $b\in B$, both restricted bundles over
$\{b\}\times Y$ are trivial.
\end{corollary}

\begin{proof}
Characteristic classes are natural under bundle isomorphism and pullback.
Apply Theorem~\ref{thm:v41-pullback}.  On the slice $\{b\}\times Y$ the
pullback of a bundle on the point $b$ is a constant vector bundle.
\end{proof}

\subsection{No-go theorem for a uniformly protected microscopic Picard class}
\label{sec:v41-Picard-no-go}

Let $\mathcal L_a^{\rm Fred}$ be the complex Fredholm determinant line of the
full chiral packet in the V11 convention,
\begin{equation}
 \mathcal L_a^{\rm Fred}
 =\det\mathcal C_a\otimes(\det\mathcal K_a)^*.
 \label{eq:v41-Fredholm-line}
\end{equation}
The invertible reduced determinant factor is already a global nonzero section,
as proved in Proposition~\ref{prop:v11-residual-line}; it contributes no
additional topological line class.

\begin{theorem}[Uniform microscopic Picard no-go]
\label{thm:v41-Picard-no-go}
Let $D_{a,y}$ be a uniformly V24-protected regulator family with one
$y$-independent continuum soft operator.  Then, for all sufficiently small
$a$,
\begin{equation}
 \boxed{
 \mathcal L_a^{\rm Fred}
 \simeq
 \pi_B^*\mathcal L_{a,*}^{\rm Fred}.}
 \label{eq:v41-line-pullback}
\end{equation}
Hence
\begin{equation}
 \boxed{
 c_1(\mathcal L_a^{\rm Fred})
 =\pi_B^*c_1(\mathcal L_{a,*}^{\rm Fred}).}
 \label{eq:v41-c1-pullback}
\end{equation}
For every fixed physical background $b$,
\begin{equation}
 \boxed{
 c_1\!\left(
 \mathcal L_a^{\rm Fred}|_{\{b\}\times Y}
 \right)=0.}
 \label{eq:v41-pure-reg-c1-zero}
\end{equation}
In particular, if $Y=S^2$,
\begin{equation}
 \boxed{
 \int_{S^2}c_1(\mathcal L_a^{\rm Fred})=0.}
 \label{eq:v41-Picard-number-zero}
\end{equation}
Thus no nonzero complete microscopic regulator Picard class can survive inside
a cutoff-uniform V24-protected V2-consistent family with parameter-independent
soft continuum data.
\end{theorem}

\begin{proof}
Take determinant exterior powers in
\eqref{eq:v41-bundle-pullback}.  This gives
\[
 \det\mathcal C_a\otimes(\det\mathcal K_a)^*
 \simeq
 \pi_B^*
 \left[
 \det\mathcal C_{a,*}\otimes(\det\mathcal K_{a,*})^*
 \right].
\]
By Proposition~\ref{prop:v11-residual-line}, the reduced invertible block does
not change the line topology, so the displayed isomorphism is precisely
\eqref{eq:v41-line-pullback}.  Naturality of $c_1$ gives
\eqref{eq:v41-c1-pullback}.  Restricting a pullback from $B$ to a regulator
slice $\{b\}\times Y$ gives a constant line, proving
\eqref{eq:v41-pure-reg-c1-zero} and its $S^2$ specialization.
\end{proof}

\begin{corollary}[V40 cannot be made uniformly V24-protected without changing its soft sector]
\label{cor:v41-V40-excluded}
The V40 family with
\begin{equation}
 c_1(\mathcal L_{\rm ang}^{\rm SM})=-30
 \label{eq:v41-V40-c1}
\end{equation}
cannot satisfy Definition~\ref{def:v41-uniform-family} on a cutoff-independent
angular sphere.  More generally, any attempted nonzero microscopic angular
Picard class must violate at least one of:
\begin{enumerate}[label=\textup{(F\arabic*)},leftmargin=3em]
\item the common V2 soft continuum jet;
\item the fixed minimal kernel/cokernel ranks and reduced singular floor;
\item the uniform V24 soft/high estimates;
\item compact uniformity of the regulator parameter family.
\end{enumerate}
No antiunitary regulator-reality axiom is required for this exclusion.
\end{corollary}

\begin{proof}
If all four items held, Theorem~\ref{thm:v41-Picard-no-go} would force the
left-hand side of \eqref{eq:v41-V40-c1} to vanish.  This contradicts the exact
V40 computation.  The final sentence follows because the proof of
Theorem~\ref{thm:v41-Picard-no-go} used only the V24 protected package and the
V14 direct rotation.
\end{proof}

\subsection{A quantitative obstruction to any nonzero Picard family}
\label{sec:v41-quantitative}

The no-go can be stated without characteristic-class language.

\begin{proposition}[Nonzero regulator topology forces a norm-one excursion]
\label{prop:v41-norm-one}
Let $Y$ be connected and let $p_y$ be a continuous finite-rank projection
family over $Y$.  Fix $y_*\in Y$.  If the bundle $\Ran p\to Y$ is not
isomorphic to the constant bundle $Y\times\Ran p_{y_*}$, then
\begin{equation}
 \boxed{
 \sup_{y\in Y}\|p_y-p_{y_*}\|\ge1.}
 \label{eq:v41-norm-one}
\end{equation}
Consequently, if the Fredholm determinant line has nonzero pure-$Y$ first
Chern class, at least one of the kernel/cokernel projector families obeys the
same norm-one lower bound relative to every fixed reference completion.
\end{proposition}

\begin{proof}
If the supremum were strictly less than one, V14's direct rotation would give a
global unitary $U_y$ with
$p_y=U_yp_{y_*}U_y^*$, trivializing the bundle.  This proves the contrapositive.
If both kernel and cokernel bundles were trivialized in this way, their
determinant ratio would be trivial by
\eqref{eq:v41-Fredholm-line}; hence a nonzero Fredholm $c_1$ forces the lower
bound for at least one packet.
\end{proof}

\begin{corollary}[Failure of a common soft anchor is quantitatively necessary]
\label{cor:v41-anchor-failure}
If a regulator family has nonzero pure-$Y$ Fredholm Picard class, then it
cannot obey a common soft norm-resolvent anchor with error
$\varepsilon_a\to0$ in the sense of V22.  Indeed, for sufficiently small $a$,
Theorem~\ref{thm:v22-soft-anchor-controlled} would imply
\begin{equation}
 \sup_y\|p_{a,y}-p_{a,y_*}\|<1,
 \label{eq:v41-anchor-under-one}
\end{equation}
contradicting Proposition~\ref{prop:v41-norm-one}.
\end{corollary}

\subsection{Flat translation-invariant symbol cross-check}
\label{sec:v41-symbol-check}

There is a simple independent explanation of the V40 gap collapse in the
translation-invariant setting.

\begin{theorem}[Momentum-homotopy no-go]
\label{thm:v41-momentum-no-go}
Let
\begin{equation}
 P:T^d\times Y\longrightarrow\operatorname{Gr}(r,\mathbb C^N)
 \label{eq:v41-P-symbol}
\end{equation}
be the positive spectral projector of a continuous uniformly gapped Hermitian
symbol family.  Suppose the regulator parameter is soft-irrelevant at one
physical momentum $k_0$, in the strong sense that
\begin{equation}
 P(k_0,y)=P_0
 \quad\text{is independent of }y.
 \label{eq:v41-soft-independent}
\end{equation}
Then for every $k\in T^d$ the bundle
\begin{equation}
 E_k:=\Ran P(k,\cdot)\longrightarrow Y
 \label{eq:v41-Ek}
\end{equation}
is isomorphic to the constant bundle $Y\times\Ran P_0$.  Hence all its Chern
classes vanish.  In particular a Hopf bundle at a UV momentum cannot coexist
with a trivial physical-momentum bundle while the full symbol remains uniformly
gapped and continuous in momentum.
\end{theorem}

\begin{proof}
The Brillouin torus is path connected.  Choose a continuous path
$k(t)$ from $k_0$ to $k$.  Then
\[
 (t,y)\longmapsto P(k(t),y)
\]
is a homotopy of rank-$r$ projectors, hence defines one vector bundle over
$[0,1]\times Y$ whose endpoint restrictions are $E_{k_0}$ and $E_k$.
Because the interval is contractible, the two endpoint bundles are isomorphic.
The first is constant by \eqref{eq:v41-soft-independent}, so the second is
constant as well.
\end{proof}

\begin{corollary}[Why the V40 discrete corner had to lose its gap]
\label{cor:v41-V40-gap-explained}
At each finite lattice the V40 momentum set is discrete, so one selected corner
can carry a Hopf line while the physical momentum does not.  Under continuum
refinement the momentum mesh becomes dense in the connected Brillouin torus.
If the complementary gap remained cutoff-uniform, the projectors would converge
to a continuous uniformly gapped symbol family contradicting
Theorem~\ref{thm:v41-momentum-no-go}.  Hence a gap collapse such as
\eqref{eq:v40-gap-a2} is structurally necessary for that translation-invariant
mechanism.
\end{corollary}

\subsection{Scope: local defects and the role of the full V24 hypothesis}
\label{sec:v41-defect-scope}

Theorem~\ref{thm:v41-Picard-no-go} is stronger than the momentum argument but
also has a sharper hypothesis.  Locality and a uniform Wilson-kernel gap alone
do not force the regulator Picard class to vanish.  A localized microscopic
defect can evade momentum homotopy because it is not described by one
translation-invariant symbol.  What V41 proves is that if such a defect carries
a nonzero complete chiral determinant Picard class, it cannot simultaneously
remain in the same V24 soft Fredholm family as the reference continuum theory.
It must inject additional soft zero-mode/bundle data or otherwise violate the
common-soft spectral estimate.

\begin{proposition}[Uniform Wilson gap is not the load-bearing condition]
\label{prop:v41-Wilson-not-enough}
A hypothetical microscopic defect family may keep the Hermitian Wilson kernel
uniformly gapped on its regulator sphere and still evade
Theorem~\ref{thm:v41-Picard-no-go} only by leaving the V24 common-soft
fixed-index class.  Thus the true no-go hypothesis is
\begin{equation}
 \boxed{
 \text{uniform Wilson/locality control}
 +
 \text{common soft Fredholm anchoring},}
 \label{eq:v41-true-hyp}
\end{equation}
not the Wilson gap by itself.
\end{proposition}

\begin{proof}
The proof of Theorem~\ref{thm:v41-Picard-no-go} uses only the V24 consequence
\eqref{eq:v41-K-close}--\eqref{eq:v41-C-close}.  A Wilson-gapped family which
still has nonzero Picard class must therefore fail the assumptions from which
those estimates are derived, namely the common soft jet/minimal reduced-gap
package.  No stronger statement about all local Wilson kernels is required.
\end{proof}

\begin{firewall}[No universal finite-cutoff no-go]
V41 does not retract V40.  The V40 Pauli sphere remains an exact finite-cutoff
local/covariant regulator Picard class.  V41 proves that such a class cannot be
made into a cofinal uniformly V24-protected family with the same
parameter-independent continuum soft operator.  A broader microscopic
regulator category which allows additional soft defect sectors can still carry
nonzero Picard topology.
\end{firewall}

\subsection{Return to the flat-character layer}
\label{sec:v41-return-H1}

On the V24-protected V2-consistent class, the degree-two obstruction has now
been removed.  The determinant line may therefore be globally trivialized in
the regulator parameter directions.  This does not by itself fix the exact
finite-cutoff connection: a topologically trivial line can still carry a flat
$U(1)$ character around a noncontractible regulator loop.  That is precisely
the stronger residue isolated in V32--V36.

\begin{theorem}[V41 hierarchy reduction]
\label{thm:v41-hierarchy-reduction}
For every compact uniformly V24-protected regulator family with one common
soft continuum jet:
\begin{enumerate}[label=\textup{(H\arabic*)},leftmargin=3em]
\item the complete complex Fredholm determinant line has no pure regulator
Picard/Chern class;
\item all exact controlled kernel/cokernel bundle classes are pulled back from
the physical-background base;
\item the first possible genuinely global regulator-only datum is therefore a
flat line character in the already-defined V32/V36 sense.
\end{enumerate}
Thus the V40 $H^2$ layer disappears on the actual cofinal protected class, and
the exact-global problem returns to the flat $H^1$ layer only after this no-go
has been imposed.
\end{theorem}

\begin{proof}
Items (H1)--(H2) are Theorem~\ref{thm:v41-Picard-no-go} and
Theorem~\ref{thm:v41-pullback}.  Once a complex line is topologically trivial,
its remaining isomorphism class as a flat unitary line is determined by its
holonomy character.  V32--V36 already isolated exactly that residue after local
phase flattening.
\end{proof}

\section{V41 anti-circularity audit}
\label{sec:v41-audit}

\begin{enumerate}[leftmargin=2.4em]
\item \textbf{The V40 finite-cutoff example was not ignored.}  V41 keeps it as
a valid local/covariant microscopic Picard class and identifies precisely which
additional V24 hypothesis prevents it from surviving cofinally.
\item \textbf{The no-go does not come from Standard-Model anomaly
cancellation.}  It comes from norm-close Riesz projectors and the canonical
V14 direct rotation.  No representation trace is used.
\item \textbf{A Wilson-kernel gap was not confused with a chiral soft gap.}
The theorem requires the full common-soft minimal-Fredholm package.  A local
Wilson-gapped defect may still carry nontrivial topology by creating extra soft
zero-mode data.
\item \textbf{No antiunitary regulator-reality axiom was silently imposed.}
Such an axiom would define a smaller regulator class but is unnecessary for the
uniform-protection no-go proved here.
\item \textbf{Topological triviality was not promoted to exact phase
triviality.}  V41 kills the $H^2$ Picard layer; a residual flat $H^1$ character
remains the next stronger exact-global question.
\end{enumerate}

\section{V41 status ledger}
\label{sec:v41-status}

\subsection*{Proved in V41}
\begin{enumerate}[leftmargin=2.2em]
\item A compact regulator-parameter family satisfying the V24 smooth protected
package uniformly has kernel and cokernel Riesz projectors within
$O(a^{1/10})$ of those of one fixed reference completion, uniformly over the
entire regulator base and with every retained source derivative.
\item The V14 polar direct rotation therefore gives global controlled bundle
isomorphisms
$\mathcal K_a\simeq\pi_B^*\mathcal K_{a,*}$ and
$\mathcal C_a\simeq\pi_B^*\mathcal C_{a,*}$ for sufficiently small cutoff.
All pure regulator-parameter characteristic classes vanish.
\item The full complex Fredholm determinant line is pulled back from the
physical-background base.  In particular its restriction to a pure regulator
sphere has $c_1=0$; no nonzero microscopic regulator Picard class can survive
inside a cutoff-uniform V24-protected V2-consistent family with
parameter-independent continuum soft data.
\item Any nonzero regulator Picard class forces a norm-one excursion of at
least one zero-mode Riesz projector from every fixed reference bundle and hence
precludes a common V22 soft norm-resolvent anchor.
\item Independently, a continuous uniformly gapped translation-invariant
projector symbol on the connected Brillouin torus cannot carry a trivial
regulator bundle at the physical momentum and a Hopf bundle at a UV momentum.
This explains structurally why the V40 isolated-corner gap collapses under
momentum refinement.
\item The V40 $c_1=-30$ construction is therefore an exact finite-cutoff
microscopic Picard class but cannot be upgraded to the cofinal uniformly V24
protected class without changing its soft Fredholm sector.  No additional
antiunitary reality assumption is needed for this exclusion.
\item After imposing the V24 protected hypotheses, the degree-two regulator
Picard layer disappears and the first remaining exact regulator-only datum is
the flat line character already isolated by V32--V36.
\end{enumerate}

\subsection*{Inherited}
\begin{enumerate}[leftmargin=2.2em]
\item V1--V24: the coefficientwise continuum and controlled regulator\
comparison, together with the soft norm-resolvent/Feshbach theorem and equality\
of exact controlled zero-mode modules for V2-consistent completions.
\item V25--V32: Pfaffian orientation, Wilson-wall/constructible determinant
lines, nonsplit angular Berry data, the corrected geometric bridge and
coefficientwise phase matching.
\item V33--V39: exact integrability/eta character, local 5D transfer
realization, coefficientwise Chern--Simons descent, native fibered-bordism
separation and exact cancellation of all smooth large-gauge components in the
complete Standard Model.
\item V40: an explicit finite-cutoff common-generation Pauli/Hopf deformation
with complete Standard-Model angular Chern class $-30$, together with the
observed collapse of its elementary translation-invariant protected radius.
\end{enumerate}

\subsection*{Conditional}
\begin{enumerate}[leftmargin=2.2em]
\item The V41 no-go applies to regulator families satisfying the full V24
uniform protected package, including the common soft continuum jet and fixed
minimal reduced-gap hypothesis.  It is not a theorem about every local finite-
cutoff Wilson kernel.
\item The regulator parameter base is compact and its constants are required to
be uniform before the cutoff is removed.  Cutoff-dependent shrinking spheres,
such as the explicit V40 realization, lie outside this theorem.
\item A localized microscopic defect can evade the momentum-symbol argument,
but if it carries nonzero determinant Picard topology it must also evade the
V24 common-soft Fredholm anchoring theorem.
\item The V12 real Pfaffian orientation and all positive IR/reference
assumptions remain separate and unchanged.
\end{enumerate}

\subsection*{Open beyond V41}
\begin{enumerate}[leftmargin=2.2em]
\item Determine the exact flat $U(1)$ character of a topologically trivial,
uniformly V24-protected microscopic regulator loop.  V32 removes every local
coefficient of the relative connection, but a noncontractible regulator base
can still support a flat finite-cutoff eta residue.
\item Decide whether the V34 local five-dimensional transfer construction and
the V35--V36 descent machinery force that purely regulator flat character to be
trivial, or whether the native completion space admits a genuine heavy-sector
flat loop invisible to the soft continuum and to all smooth large-gauge
classes.
\item The independent stronger problems remain: exact global treatment of any
non-V24 defect sectors, positive-IR removal, analytic finite-coupling RG
trajectories, perturbative summability and nonperturbative ESM completion.
\end{enumerate}

\subsection*{Exact next load-bearing theorem}
\begin{center}
\fbox{\parbox{0.92\textwidth}{
\textbf{Exact flat microscopic regulator-loop theorem on the V24-protected
class.}
V41 removes the degree-two microscopic Picard obstruction whenever the native
completion family is cutoff-uniformly V24 protected and shares one continuum
soft jet.  Let $Y$ now be a noncontractible regulator-parameter space whose
complex determinant line is topologically trivial by V41.  Compute the exact
finite-cutoff flat holonomy character left after the V32 local phase
normalization.  Either prove, using the V34 five-dimensional local transfer and
V35--V36 descent, that every such purely regulator loop has unit character, or
construct a genuinely heavy-sector flat loop with nontrivial eta residue while
all V24 soft data and all V37--V39 smooth gauge-transform invariants remain
fixed.  This is now the first unresolved exact complex microscopic invariant on
the cofinal protected regulator class.}}
\end{center}

\section*{Append-only continuation marker: V41}
\addcontentsline{toc}{section}{Append-only continuation marker: V41}
\textbf{End of V41.}  The complete V1--V40 body above is retained verbatim.
V41 proves a no-go theorem for nonzero complete microscopic regulator Picard
classes on the cofinal uniformly V24-protected V2-consistent regulator class.
Uniform soft/Feshbach control makes every kernel and cokernel Riesz bundle over
the regulator parameter base globally direct-rotation equivalent to one fixed
reference bundle; the full Fredholm determinant line is therefore pulled back
from the physical-background base and has no pure regulator Chern class.  This
explains structurally why the explicit V40 Hopf sphere had to lose its protected
gap under refinement and shows that no extra antiunitary regulator-reality axiom
is needed at the V24 scope.  A translation-invariant momentum-homotopy theorem
gives an independent flat-symbol proof of the same obstruction.  The degree-two
microscopic hierarchy is therefore closed on the cofinal protected class; the
next stronger question is the exact flat finite-cutoff regulator-loop character
on the resulting topologically trivial determinant line.


\clearpage
\section{V42 continuation: exact contraction of the V24-protected regulator fibre and triviality of the residual flat character}
\label{sec:v42-continuation}

V41 removed the degree-two microscopic regulator obstruction on every compact
cutoff-uniform V24-protected family with one common soft continuum operator.
It left open a logically stronger possibility: even when the complete complex
Fredholm line is topologically trivial in the regulator directions, a
noncontractible regulator parameter space could carry an exact flat character
\(
 \rho_a^{\rm reg}:\pi_1(Y)\to U(1)
\)
after the V36 exact local flattening.  Searching directly for a new heavy-sector
eta invariant is, however, one categorical level too late.  The V11 geometry of
the minimal overlap stratum already says that every nonzero-mode degree of
freedom is an affine graph once the forced zero-mode plane is fixed.  V41 gives
precisely the missing global input: the forced kernel/cokernel plane is
uniformly within norm distance strictly less than one from one reference
completion.  The canonical V14 direct rotation therefore aligns the zero-mode
plane globally, after which the affine graph contracts.

The conclusion is stronger than flat-holonomy triviality.  On the cofinal
V24-protected class, the entire pure-regulator overlap family is natively
homotopic, in the exponentially local protected category used from V13 onward,
to the pullback of one reference regulator.  Consequently every regulator
loop admits an explicit fibered/native disk filling and the V36 residual flat
character is exactly one.  No independent complex ``heavy \(K_1\)'' phase
survives inside the declared minimal overlap geometry.

The new results are:
\begin{enumerate}[label=\textup{(V42.\arabic*)},leftmargin=3.8em]
\item a global controlled alignment of the V11 kernel/cokernel plane over an
arbitrary compact regulator base \(Y\), relative to one reference slice;
\item an explicit source-complete contraction of the full V11 chiral
projector family over \(B\times Y\) to the pullback of that reference slice;
\item an exponentially local sign-flattened native lift of this contraction,
with exact Wilson-sign gap and fixed physical overlap principal jet;
\item an explicit fibered/native disk filling for every pure-regulator loop;
\item exact triviality of the V36 regulator flat character and of the V33 eta
residue on the complete faithful complex Standard-Model packet;
\item a no-go statement for a separate heavy-sector flat/K1 invariant within
the V24 minimal protected class.
\end{enumerate}

\subsection{Setup: a compact pure-regulator family over one physical panel}
\label{sec:v42-setup}

Let \(B\) be the fixed compact physical gauge/spin/coframe/Higgs/source panel
and let \(Y\) be a compact regulator-parameter space.  For each sufficiently
small lattice spacing let
\begin{equation}
 (b,y)\longmapsto P_{a}(b,y)
 \label{eq:v42-P-family}
\end{equation}
be the positive Wilson-sign projector defining the native overlap chiral
plane
\(
 W_a(b,y)=\Ran P_a(b,y)
\).
Assume uniformly on \(B\times Y\):
\begin{enumerate}[label=\textup{(R\arabic*)},leftmargin=3em]
\item the V24 V2-consistent common-soft package with one regulator-independent
continuum Dirac--Yukawa operator;
\item one fixed minimal index \(k\) and one positive reduced singular-value
floor;
\item the V13 exponentially local/source-complete functional calculus;
\item the same physical overlap principal Dirac--Yukawa jet and the same
faithful complex Standard-Model representation grading.
\end{enumerate}
Choose a base point \(y_*\in Y\) and write
\begin{equation}
 P_{a,*}(b):=P_a(b,y_*).
 \label{eq:v42-P-star}
\end{equation}
For \(k>0\) let \(p^K_a(b,y)\) be the forced kernel projector; for \(k<0\)
let \(p^C_a(b,y)\) be the forced cokernel projector.  The index-zero case has
no forced zero-mode plane.

V41 gives, with every retained source derivative,
\begin{equation}
 \sup_{B\times Y}
 \|p^{Z}_a(b,y)-p^{Z}_{a,*}(b)\|
 \le C a^{1/10}<1,
 \qquad
 Z=K\ \text{or}\ C,
 \label{eq:v42-zero-close}
\end{equation}
for sufficiently small \(a\).

\subsection{Global controlled alignment of the forced zero-mode plane}
\label{sec:v42-zero-align}

\begin{lemma}[Canonical global zero-mode alignment relative to the reference slice]
\label{lem:v42-global-zero-align}
For sufficiently small \(a\), there is a unique canonical V14 polar direct
rotation
\begin{equation}
 u_a(b,y)\in\mathcal A_{\exp,\mathfrak q}
 \label{eq:v42-u-align}
\end{equation}
with
\begin{equation}
 u_a p_a^Z u_a^*=p_{a,*}^Z,
 \qquad
 u_a(b,y_*)=I,
 \label{eq:v42-u-align-prop}
\end{equation}
where \(Z=K\) for \(k>0\) and \(Z=C\) for \(k<0\).  The family \(u_a\) is
smooth in \((b,y)\), exactly gauge/spin covariant, exponentially local with all
retained source/BV derivatives, and lies in the identity component through a
canonical controlled path \(u_{a,t}\), \(0\le t\le1\), with
\(u_{a,0}=I\), \(u_{a,1}=u_a\).
\end{lemma}

\begin{proof}
Equation~\eqref{eq:v42-zero-close} is exactly the norm-\(<1\) hypothesis of the
V14 direct-rotation theorem.  With \(p=p_a^Z\), \(q=p_{a,*}^Z\), put
\[
 X=qp+(1-q)(1-p),
 \qquad
 u=X(X^*X)^{-1/2}.
\]
The inverse square root is uniformly defined because
\(X^*X=I-(p-q)^2\) has a positive floor.  V13 weighted functional calculus
therefore makes \(u\) exponentially local, and differentiation of the same
formula gives all fixed source/BV jets.  Gauge/spin covariance is preserved by
products and functional calculus.  At \(y=y_*\) one has \(p=q\), hence
\(u=I\).  The V14 identity-component path is obtained from the polar factors of
\(X_t=I+t(X-I)\), which remain invertible under the same norm-\(<1\) bound.
\end{proof}

Extend \(u_a\) block-diagonally with respect to the fixed chirality splitting:
for \(k>0\), it acts on \(\mathcal H_-\) and is the identity on
\(\mathcal H_+\); for \(k<0\), it acts on \(\mathcal H_+\) and is the identity
on \(\mathcal H_-\).  Denote the resulting unitary again by \(u_a\).

\begin{corollary}[Alignment stays in the minimal stratum and fixes the physical quotient]
\label{cor:v42-alignment-minimal}
The conjugated family
\begin{equation}
 \widetilde P_a(b,y):=u_a(b,y)P_a(b,y)u_a(b,y)^*
 \label{eq:v42-P-tilde}
\end{equation}
has the same fixed minimal index as \(P_a\), its forced kernel/cokernel plane
is exactly the reference plane \(p^Z_{a,*}(b)\), and the annular homotopy
\begin{equation}
 P^{\rm ann}_{a,t}:=u_{a,t}P_a u_{a,t}^*
 \label{eq:v42-annular-path}
\end{equation}
is an exact exponentially local protected native homotopy.  Its change from
the original family is confined to the controlled soft/compact module and does
not alter the frozen physical principal quotient.
\end{corollary}

\begin{proof}
A block-diagonal chirality-preserving unitary maps the V11 kernel or cokernel
plane isomorphically and cannot create an accidental opposite-chirality zero
pair.  Minimality and the numerical index are therefore preserved.  The
locality and source statements are Lemma~\ref{lem:v42-global-zero-align}.
The last statement is the same compact/soft-ideal observation used in the V13
controlled-reduction theorem: controlled finite-rank zero-mode transport does
not modify the principal Dirac--Yukawa quotient.
\end{proof}

\subsection{The affine graph kills every remaining pure-regulator homotopy class}
\label{sec:v42-affine-contraction}

After Corollary~\ref{cor:v42-alignment-minimal}, the zero-mode plane no longer
depends on \(y\).  The V11 minimal-stratum theorem can therefore be applied
globally over \(B\times Y\).

For \(k>0\), write
\begin{equation}
 \Ran\widetilde P_a(b,y)
 =K_{a,*}(b)\oplus
 \{v+\widetilde A_a(b,y)v:v\in\mathcal H_+\},
 \label{eq:v42-positive-graph}
\end{equation}
with
\begin{equation}
 \widetilde A_a(b,y):\mathcal H_+\to K_{a,*}(b)^\perp.
 \label{eq:v42-positive-A}
\end{equation}
For \(k<0\), write
\begin{equation}
 \Ran\widetilde P_a(b,y)
 =\{v+\widetilde A_a(b,y)v:
 v\in C_{a,*}(b)^\perp\},
 \label{eq:v42-negative-graph}
\end{equation}
with
\begin{equation}
 \widetilde A_a(b,y):C_{a,*}(b)^\perp\to\mathcal H_-.
 \label{eq:v42-negative-A}
\end{equation}
For \(k=0\) this is the V10 balanced graph coordinate.  Let
\(A_{a,*}(b)\) be the corresponding reference graph at \(y_*\).

\begin{lemma}[Global source-complete graph contraction]
\label{lem:v42-graph-contraction}
For \(0\le s\le1\), define
\begin{equation}
 A_{a,s}(b,y)
 :=(1-s)\widetilde A_a(b,y)+sA_{a,*}(b).
 \label{eq:v42-A-contract}
\end{equation}
Then the V11 graph formula defines a smooth family of projectors
\begin{equation}
 \widetilde P_{a,s}(b,y)
 \label{eq:v42-P-contract}
\end{equation}
which remains in the same minimal index-\(k\) stratum for every
\((b,y,s)\), satisfies
\begin{equation}
 \widetilde P_{a,0}=\widetilde P_a,
 \qquad
 \widetilde P_{a,1}(b,y)=P_{a,*}(b),
 \label{eq:v42-P-contract-endpoints}
\end{equation}
and is exponentially local with every retained source/BV derivative.  The
reduced kinetic singular values have a positive uniform floor on the compact
\(B\times Y\times[0,1]\) family.  After the same protected interacting-chart
shrinking already allowed in V11--V13, the full Dirac--Yukawa reduced floor is
also positive.
\end{lemma}

\begin{proof}
Once the kernel/cokernel plane is fixed, the V11 minimal stratum is the total
space of a vector bundle whose fibre is precisely the graph space in
\eqref{eq:v42-positive-A} or \eqref{eq:v42-negative-A}.  Hence the convex
interpolation \eqref{eq:v42-A-contract} never leaves the minimal stratum.
V13 proves that the graph coordinate is exponentially local and source smooth;
finite convex combinations preserve the same weighted algebra.  The graph
projector is obtained from \(A\) by products and
\((I+A^*A)^{-1}\), so V13 functional calculus gives the asserted locality.
The V11 reduced kinetic block is explicitly
\[
 {2\over a}(I+A_{a,s}^*A_{a,s})^{-1/2}.
\]
Compactness bounds \(\|A_{a,s}\|\), yielding a positive reduced floor.  The
last statement is exactly the interacting-chart step in the V13 controlled
reduction theorem: one common protected shrinking keeps the bounded
Dirac--Yukawa perturbation away from an accidental paired-zero divisor along
the compact graph homotopy.
\end{proof}

\begin{theorem}[Exact contraction of the full V24-protected regulator fibre]
\label{thm:v42-full-contraction}
Under the V42 setup, for all sufficiently small \(a\) there is an exact
source-complete exponentially local gauge/spin-covariant homotopy
\begin{equation}
 \boxed{
 P_{a,s}:B\times Y\longrightarrow\mathcal M_{k,a}^{\rm prot},
 \qquad 0\le s\le1,}
 \label{eq:v42-full-homotopy}
\end{equation}
through the same V11 minimal fixed-index native component such that
\begin{equation}
 \boxed{
 P_{a,0}(b,y)=P_a(b,y),
 \qquad
 P_{a,1}(b,y)=P_{a,*}(b).}
 \label{eq:v42-full-homotopy-endpoints}
\end{equation}
The homotopy is relative to the reference slice \(B\times\{y_*\}\).  Thus the
entire chiral-projector family is natively homotopic to the pullback
\begin{equation}
 \boxed{
 P_a\simeq_{\rm nat}\pi_B^*P_{a,*}.}
 \label{eq:v42-pullback-full}
\end{equation}
In particular every pure regulator homotopy invariant of the full minimal
projector family, not merely its determinant first Chern class, is pulled back
from \(B\).
\end{theorem}

\begin{proof}
Concatenate the annular alignment path
\eqref{eq:v42-annular-path} with the graph contraction
\eqref{eq:v42-P-contract}.  At \(y=y_*\), the direct rotation is the identity
and the graph is already \(A_{a,*}\), so the concatenation fixes the reference
slice.  A small reparameterization smooths the concatenation seam without
leaving the protected open set.  Corollary~\ref{cor:v42-alignment-minimal} and
Lemma~\ref{lem:v42-graph-contraction} give exact locality, covariance, source
smoothness and the reduced-gap statement throughout.
\end{proof}

\begin{corollary}[No higher pure-regulator projector obstruction in the minimal V24 class]
\label{cor:v42-no-higher-projector}
For any compact CW regulator base \(Y\), the classifying map of the full native
minimal overlap projector over \(B\times Y\) factors through \(B\) up to a
native protected homotopy.  Hence all pure-\(Y\) characteristic classes,
clutching classes and higher homotopy invariants of that projector vanish.
The V41 Picard no-go is the determinant-line shadow of this stronger result.
\end{corollary}

\begin{proof}
Equation~\eqref{eq:v42-pullback-full} is a homotopy relative to the reference
slice.  Every homotopy invariant functor applied to \(P_a\) therefore equals
the pullback of its value on \(P_{a,*}\).  Restriction to a pure regulator
fibre \(\{b\}\times Y\) is constant up to homotopy.
\end{proof}

\subsection{Native lift by sign flattening}
\label{sec:v42-sign-lift}

The preceding theorem is stated at the overlap-projector level.  For the exact
flat-character problem this is already the relevant datum, but it is useful to
record an explicit Hermitian-kernel lift.

\begin{proposition}[Sign-flattened exponentially local native lift]
\label{prop:v42-sign-lift}
Let \(P_{a,s}\) be the homotopy of
Theorem~\ref{thm:v42-full-contraction}.  Define
\begin{equation}
 \boxed{
 \widehat H_{W,a,s}:=2P_{a,s}-I,
 \qquad
 \widehat K_{W,a,s}:=\gamma_5\widehat H_{W,a,s}.}
 \label{eq:v42-sign-flattened-H}
\end{equation}
Then:
\begin{enumerate}[label=\textup{(L\arabic*)},leftmargin=3em]
\item \(\widehat H_{W,a,s}\) is Hermitian, gauge/spin covariant and
exponentially local with all retained source/BV derivatives;
\item its Wilson-sign gap is exactly one and
\begin{equation}
 \operatorname{sgn}\widehat H_{W,a,s}=2P_{a,s}-I;
 \label{eq:v42-sign-exact}
\end{equation}
\item the induced overlap operator is exactly
\begin{equation}
 \widehat D_{{\rm ov},a,s}
 ={1\over a}\bigl(I+\gamma_5(2P_{a,s}-I)\bigr),
 \label{eq:v42-Dov-lift}
\end{equation}
and remains in the same minimal fixed-index protected component with the same
physical overlap principal jet.
\end{enumerate}
Thus Theorem~\ref{thm:v42-full-contraction} is a genuine native filling in the
exponentially local regulator category used by V13--V24, not merely an
abstract Grassmannian homotopy.
\end{proposition}

\begin{proof}
Items (L1)--(L2) are immediate from the projector identities and the weighted
locality of \(P_{a,s}\).  Equation~\eqref{eq:v42-Dov-lift} is the defining
overlap formula with the flattened sign.  The physical quotient and minimal
reduced-gap properties were preserved in
Theorem~\ref{thm:v42-full-contraction}; sign flattening changes no overlap
projector and therefore changes no overlap Dirac operator associated with that
projector.
\end{proof}

\begin{firewall}[Finite range is not required for the exact global phase classification]
The lift \eqref{eq:v42-sign-flattened-H} is generally exponentially local rather
than ultralocal.  This is the native local-operator category used from V13
onward for controlled modules, Kato transport and exact determinant-line
classification.  V42 therefore does not claim a contraction through one fixed
finite stencil or one fixed hopping range.  Imposing such an extra microscopic
symmetry can create additional pump-like invariants, but they are not
invariants of the declared exponentially local Renewal regulator category and
cannot obstruct the determinant-line character studied here.
\end{firewall}

\subsection{Every pure-regulator loop has a fibered/native disk filling}
\label{sec:v42-loop-fill}

Let \(\gamma:S^1\to Y\) be any smooth regulator loop.  Pull
Theorem~\ref{thm:v42-full-contraction} back along \(\gamma\).  The resulting
homotopy contracts the loop of native overlap completions to the constant
reference completion.  To match the precise V36 definition one may convert
this cylinder contraction into a smooth disk as follows.

\begin{proposition}[Explicit fibered filling of a regulator loop]
\label{prop:v42-regulator-filling}
For every smooth \(\gamma:S^1\to Y\) and all sufficiently small \(a\), the
native loop
\begin{equation}
 \Gamma_{a,\gamma}:S^1\to\widetilde{\mathcal M}_a^{\rm prot},
 \qquad
 \theta\mapsto P_a(\cdot,\gamma(\theta))
 \label{eq:v42-native-loop}
\end{equation}
admits a V36 fibered/native filling.  The filling is source-complete over the
whole physical panel \(B\) and remains in one fixed-index chamber; no V27--V29
wall insertion is required.
\end{proposition}

\begin{proof}
Use the first half of the disk as an annulus implementing the canonical
identity-component direct rotation from the original loop to the aligned loop.
On the inner disk keep the zero-mode plane fixed and fill the graph loop inside
its vector space.  For smoothness at the disk centre choose a radial cutoff
\(\chi(r)\) that is flat at \(r=0\) and equals one at the inner-annulus
boundary, and put
\[
 A(r,\theta)
 =A_{a,*}+\chi(r)
 \bigl(\widetilde A_a(\gamma(\theta))-A_{a,*}\bigr).
\]
The V11 graph formula then gives a smooth centre because every angular
derivative is multiplied by a function vanishing to infinite order there.
The annulus/disk seam can be smoothed inside a small collar since the minimal
reduced-gap inequalities have a compact positive margin.  Lemma
\ref{lem:v42-global-zero-align}, Lemma~\ref{lem:v42-graph-contraction} and
Proposition~\ref{prop:v42-sign-lift} give locality, covariance and the fixed
source/BV packet.  This is exactly Definition~\ref{def:v36-fibered-filling}.
\end{proof}

\subsection{Exact triviality of the residual flat regulator character}
\label{sec:v42-flat-trivial}

\begin{theorem}[Exact flat microscopic regulator-loop theorem]
\label{thm:v42-flat-trivial}
Let \(Y\) be any compact regulator-parameter space satisfying the V42/V41
uniform V24 protected hypotheses.  After the V36 exact local quasilocal
flattening of the native--geometric relative determinant connection, the pure
regulator holonomy character is exactly trivial:
\begin{equation}
 \boxed{
 \rho_a^{\rm reg}:\pi_1(Y)\longrightarrow U(1),
 \qquad
 \rho_a^{\rm reg}\equiv1}
 \label{eq:v42-rho-trivial}
\end{equation}
for all sufficiently small \(a\).  On the inherited faithful complex
Standard-Model geometric branch this is equivalently
\begin{equation}
 \boxed{
 \chi_{a,{\rm SM}}^{\eta,{\rm reg}}(\gamma)=1
 \quad\text{for every }\gamma\in\pi_1(Y).}
 \label{eq:v42-eta-trivial}
\end{equation}
The conclusion holds simultaneously with every retained source/BV descendant.
\end{theorem}

\begin{proof}
Proposition~\ref{prop:v42-regulator-filling} gives a V36 fibered/native disk
filling for every regulator loop.  Proposition~\ref{prop:v36-fibered-kills}
therefore gives \(\rho_a^{\rm reg}(\gamma)=1\).  V33 identifies the relative
flat holonomy with the exact eta/integrability residue, and the inherited
faithful geometric Standard-Model line has unit parallel holonomy.  The whole
filling was constructed with the source-complete V13/V41 direct-rotation and
graph calculus, so the conclusion holds for the retained descendant packet as
well.
\end{proof}

\begin{corollary}[The V34 five-dimensional regulator phase has no residual torsion character]
\label{cor:v42-v34-no-torsion}
For every pure regulator loop in the V42 class, the infinite-wall V34 local
five-dimensional phase, after the exact V36 local curvature subtraction, has
unit closed-loop value.  Equivalently, all of its nontrivial regulator
variation is carried by the local/quasilocal curvature primitive already
separated in V32, V35 and V36; no additional flat/torsion factor remains.
\end{corollary}

\begin{proof}
V34 realizes the V33 eta character as the infinite-wall limit of local
five-dimensional transfer phases.  Theorem~\ref{thm:v42-flat-trivial} makes
the exact residual character one, leaving only the locally exact connection
part removed by the V36 primitive.
\end{proof}

\begin{firewall}[Raw Berry holonomy is not claimed to vanish]
Theorem~\ref{thm:v42-flat-trivial} concerns the \emph{residual flat character}
after the exact V36 local curvature has been separated.  A raw natural
connection can have nonzero Berry holonomy around a contractible loop because
its curvature has nonzero flux through a filling disk.  V42 does not set that
curvature to zero.  It proves that once the local/quasilocal curvature piece is
removed, no additional regulator-only flat holonomy survives.  At the original
finite-EFT quantifiers, V32 and V35 provide the corresponding coefficientwise
local normalization.
\end{firewall}

\subsection{Why there is no separate heavy-sector \texorpdfstring{$K_1$}{K1} phase here}
\label{sec:v42-no-heavy-k1}

The V41 handoff allowed the possibility of a purely heavy loop invisible to all
soft data.  The preceding contraction shows why such a loop cannot produce a
new complex determinant character while the family remains in the V11 minimal
protected overlap geometry.

\begin{theorem}[No independent heavy-sector flat invariant in the minimal overlap component]
\label{thm:v42-no-heavy-k1}
Suppose a closed microscopic Hermitian Wilson-kernel loop satisfies the V42
hypotheses and may vary arbitrarily in its ultraviolet coefficients while
keeping the common physical overlap jet.  Then its overlap determinant-line
flat character depends only on the corresponding sign-projector family
\(P_a\), and that family is null-homotopic by
Theorem~\ref{thm:v42-full-contraction}.  Hence no additional stable or unstable
``heavy \(K_1\)'' invariant can be detected by the complex overlap determinant
line on this component.

A nontrivial microscopic loop character must therefore violate at least one of:
\begin{enumerate}[label=\textup{(K\arabic*)},leftmargin=3em]
\item the V11 minimal fixed-index condition, by crossing an accidental
paired-zero or Wilson wall;
\item the V41 norm-\(<1\) controlled zero-mode alignment;
\item the common physical overlap principal jet/V24 soft anchoring;
\item the exponentially local gauge/spin-covariant native category.
\end{enumerate}
\end{theorem}

\begin{proof}
The overlap operator and its chiral measure line depend on the microscopic
Hermitian kernel through its spectral sign, equivalently through the positive
projector \(P_a\).  Spectral flattening
\(H_W\mapsto\operatorname{sgn}H_W=2P_a-I\) preserves that data exactly and has
a unit Hermitian gap.  Theorem~\ref{thm:v42-full-contraction} then contracts the
projector family after the forced zero-mode module is aligned.  The V36
fibered-filling theorem kills the residual flat phase.  Thus any proposed
counterexample must leave one of the hypotheses used by that contraction.
\end{proof}

\begin{remark}[Extra microscopic symmetries can define a different problem]
A translation-invariant or fixed-range subcategory can support pump-like
indices which disappear once one allows the exponentially local covariant
homotopies of the Renewal native category.  V42 makes no claim that those
stronger symmetry-protected classifications are trivial.  It shows only that
they are not obstructions to the declared overlap determinant character unless
that extra symmetry is promoted to an additional physical axiom.
\end{remark}

\subsection{Real/Pfaffian sector}
\label{sec:v42-real}

The V42 contraction is complex-linear and does not replace the V25 real
orientation theorem.  On the natively oriented neutral Majorana branch, the
V14 direct rotations lie in the orientation-preserving component by V25, so the
same contraction introduces no new regulator-only \(w_1\) class.  Charged
pseudoreal weak-doublet transport remains governed by the exact V38 parity
theorem.  Thus V42 closes the new \emph{complex} flat regulator layer without
silently identifying it with the independent real/Pfaffian orientation data.

\section{V42 anti-circularity audit}
\label{sec:v42-audit}

\begin{enumerate}[leftmargin=2.4em]
\item \textbf{Topological triviality of the line was not mistaken for flat
triviality.}  V42 constructs an actual native disk filling of every regulator
loop and invokes the V36 filling theorem; it does not infer unit holonomy from
\(c_1=0\) alone.
\item \textbf{The heavy sector was not assumed small in operator norm.}  Only
the forced zero-mode projector is norm-close to the reference by V41.  The
remaining full overlap family is contracted using the exact V11 affine graph
geometry, which allows order-one ultraviolet variation.
\item \textbf{No global graph chart was assumed before aligning zero modes.}
The V14 direct rotation first makes the kernel/cokernel plane identical to the
reference plane; only then is the V11 graph fibre used globally.
\item \textbf{Locality is preserved.}  Zero-mode rotation uses V13/V14 weighted
functional calculus, and the graph contraction uses the V13 local graph
formula.  The sign-flattened lift has exact Hermitian gap one.
\item \textbf{The contraction need not stay inside the quantitative V24
Feshbach subclass.}  V24 is used to obtain the uniform zero-mode alignment at
the endpoints.  The filling itself lies in the larger exact V11/V13 protected
native component, which is exactly what V36 requires.
\item \textbf{Raw Berry curvature was not erased.}  Only the residual flat
character after V36 exact local curvature subtraction is proved trivial.
\item \textbf{No fixed-range or translation-invariance constraint was silently
added.}  The theorem works in the exponentially local native category already
used for the controlled module classification.
\end{enumerate}

\section{V42 status ledger}
\label{sec:v42-status}

\subsection*{Proved in V42}
\begin{enumerate}[leftmargin=2.2em]
\item On every compact uniformly V24-protected regulator family, the V41
norm-\(<1\) estimate gives a canonical source-complete controlled direct
rotation aligning the forced kernel/cokernel plane globally to one reference
completion.
\item After this alignment, the full V11 minimal overlap projector is an affine
graph over a fixed zero-mode plane.  Linear graph contraction therefore gives
an exact exponentially local gauge/spin-covariant homotopy of the whole
projector family to the pullback of the reference completion.
\item The contraction is relative to the reference regulator slice and
preserves the fixed minimal index, the physical overlap principal jet and a
positive reduced floor.  A sign-flattened Hermitian lift has exact Wilson-sign
gap one.
\item Consequently the entire pure-regulator projector family over
\(B\times Y\) is natively homotopic to \(\pi_B^*P_{a,*}\).  All pure regulator
projector characteristic and clutching data vanish, not only the V41 Picard
class.
\item Every regulator loop admits an explicit source-complete V36
fibered/native disk filling.  Hence the exact locally-flattened regulator
character satisfies \(\rho_a^{\rm reg}\equiv1\) for all sufficiently small
cutoffs.
\item On the faithful complex Standard-Model branch the exact V33 eta residue
of every pure regulator loop is therefore one.  The V34 infinite-wall phase has
no additional regulator-only flat/torsion factor after the V36 local curvature
subtraction.
\item No separate complex heavy-sector K1-like invariant is visible to the
overlap determinant line inside the minimal V24-protected component.  Any such
candidate must leave the minimal stratum, the V41 controlled alignment, the
common soft/principal class, or the declared exponentially local native
category.
\end{enumerate}

\subsection*{Inherited}
\begin{enumerate}[leftmargin=2.2em]
\item V1--V24: the source-complete coefficientwise ESM continuum, the V11
minimal-stratum affine geometry, V13 locality of zero-mode/graph data, V14
direct rotation and the V24 common soft anchoring theorem.
\item V25--V32: real Pfaffian orientation, resolved Wilson-wall determinant
lines, angular Berry data, corrected geometric bridge and coefficientwise phase
matching.
\item V33--V39: exact eta/integrability character, local five-dimensional
realization, exact local/global separation, native defect group and exact
cancellation of all smooth large-gauge Standard-Model classes.
\item V40--V41: an exact finite-cutoff regulator Picard counterexample outside
the cofinal protected class, followed by the uniform V24 Picard no-go and
global zero-mode pullback theorem.
\end{enumerate}

\subsection*{Conditional}
\begin{enumerate}[leftmargin=2.2em]
\item V42 assumes the full V41/V24 compact uniform protected package and one
minimal fixed-index branch.  A finite-cutoff defect sector with a nonzero V40
Picard class is outside this theorem.
\item The exact native contraction is constructed in the exponentially local
operator category used from V13 onward.  A smaller category imposing one fixed
ultralocal stencil, exact translation invariance or another microscopic
symmetry may have additional symmetry-protected invariants.
\item The interacting Dirac--Yukawa/source panel is restricted to the same
protected compact chart on which the V11--V13 reduced-floor theorem applies.
\item The V25 real/Pfaffian orientation hypotheses remain separate and are not
replaced by the complex contraction theorem.
\end{enumerate}

\subsection*{Open beyond V42}
\begin{enumerate}[leftmargin=2.2em]
\item Exact global phases in non-V24 defect sectors remain possible, but V40--
V42 show that such sectors cannot contaminate the cofinal common-soft regulator
class without leaving its protected hypotheses.
\item The independent stronger problems now move beyond microscopic complex
regulator topology: removal of the positive IR/reference scale at an
appropriately infrared-safe set of observables/kinematics; analytic
finite-coupling RG trajectories; perturbative summability; and a genuinely
nonperturbative ESM completion.
\item A real/Pfaffian extension beyond the already oriented V25 neutral branch
would require a separately declared real regulator category if new unoriented
Majorana sectors are introduced.
\end{enumerate}

\subsection*{Exact next load-bearing theorem}
\begin{center}
\fbox{\parbox{0.92\textwidth}{
\textbf{Infrared-safe zero-reference-scale extension.}
V42 closes the first unresolved exact complex microscopic invariant on the
cofinal V24-protected regulator class: after exact local curvature subtraction,
every pure-regulator flat loop has unit character.  The next genuinely
load-bearing strengthening is therefore no longer microscopic topology.
Choose a precisely infrared-safe observable/source class (for example
nonexceptional Euclidean momentum configurations, or an inclusive/gauge-
invariant class with the necessary soft subtractions) and determine whether
the positive IR/reference scale retained throughout V1--V42 can be removed
coefficientwise at every fixed loop/EFT/source order.  The theorem must identify
the exact soft factorization/subtraction operator, prove uniform compatibility
with the V5 forest/BV recursion and V7 projective normalization, and state
explicitly which massless Green functions or on-shell amplitudes remain
excluded.  No claim of an unrestricted massless S-matrix limit should be made
without the corresponding KLN/inclusive hypotheses.}}
\end{center}

\section*{Append-only continuation marker: V42}
\addcontentsline{toc}{section}{Append-only continuation marker: V42}
\textbf{End of V42.}  The complete V1--V41 body above is retained verbatim.
V42 proves that the residual exact flat regulator character left after V41 is
actually trivial on the cofinal uniformly V24-protected minimal overlap class.
The V41 norm-\(<1\) zero-mode alignment is upgraded, using the V11 affine graph
geometry and V13 locality, to an exact source-complete native contraction of
the entire regulator projector family to one reference completion.  Every
pure-regulator loop therefore has a V36 fibered/native disk filling, so its
locally-flattened V33 eta/integrability residue is exactly one.  This removes a
separate complex heavy-sector K1-like phase from the declared exponentially
local native category while retaining the finite-cutoff V40 defect sectors
outside the V24 class and the independent real/Pfaffian hypotheses.  The next
stronger programme can therefore move upstream from microscopic regulator
topology to the positive-IR/reference-scale assumption.




% ============================================================================
% V43 APPEND-ONLY CONTINUATION
% ============================================================================
\clearpage
\part{V43 append-only continuation:\\
infrared-safe zero-floor completion at nonexceptional Euclidean momenta}
\label{part:v43}

\section{Purpose and the upstream two-scale correction}
\label{sec:v43-purpose}

V42 closed the pure microscopic complex regulator topology on the cofinal
V24-protected class.  The strongest remaining assumption already visible in
V6--V7 is the positive ``IR/reference scale.''  Before attempting to send that
quantity to zero, one must separate two roles which were deliberately bundled
in the earlier fixed-scale theorem:
\begin{enumerate}[label=\textup{(R\arabic*)},leftmargin=3em]
\item a \emph{renormalization/reference point} $\mu_R>0$, at which the finite
physical/EFT/source normalization coordinates are defined;
\item an \emph{infrared floor} $\rho>0$, below which massless loop modes have
not yet been included in the low-shell completion.
\end{enumerate}
These are not the same datum.  A massless renormalized theory generally retains
an arbitrary nonzero subtraction point $\mu_R$; sending the subtraction point
to zero is neither required for an infrared limit nor desirable at exceptional
kinematics.  The actual stronger problem is to remove $\rho$ while keeping
$\mu_R$ fixed and proving covariance under the already-established V7 changes
of positive reference point.

Accordingly V43 does \emph{not} attempt a $\mu_R\downarrow0$ theorem.  It
constructs the zero-floor limit
\[
 \rho\downarrow0
\]
for a precise infrared-safe class of Euclidean connected/source coefficients.
The output is a massless coefficientwise continuum with no auxiliary positive
infrared floor and with one ordinary nonzero renormalization point, exactly as
in standard massless perturbative renormalization.

\begin{firewall}[Reference scale versus infrared regulator]
\label{fw:v43-reference-not-ir}
The symbol $\mu$ in V1--V42 controlled several fixed positive scales at once:
the external/reference panel, the starting scale of the hard-shell comparison,
and the weighted long-distance bounds.  V43 refines that notation rather than
changing any previous theorem.  We write $\mu_R>0$ for the normalization point,
$Q<\infty$ for an upper external/source momentum bound, $\kappa>0$ for
nonexceptionality, and $\rho>0$ for an auxiliary lower soft-shell floor.  Only
$\rho$ is removed.
\end{firewall}

\section{Strictly nonexceptional Euclidean source panels}
\label{sec:v43-panel}

Let a connected coefficient have external and source momenta
$\mathbf q=(q_1,\ldots,q_n)$, including momenta carried by composite, stress,
ghost/antifield and line-source insertions which are treated as genuine
external insertions.  Momentum conservation gives $\sum_iq_i=0$.

\begin{definition}[Nonexceptionality modulus]
\label{def:v43-exceptionality}
For $n\ge2$ define
\begin{equation}
 \eta(\mathbf q)
 :=\min_{\varnothing\ne S\subsetneq\{1,\ldots,n\}}
 \left|\sum_{i\in S}q_i\right|.
 \label{eq:v43-eta}
\end{equation}
For $0<\kappa\le Q<\infty$, the compact panel
$\mathcal K_{\kappa,Q}^{(n)}$ consists of Euclidean configurations satisfying
\begin{equation}
 \sum_iq_i=0,
 \qquad
 |q_i|\le Q,
 \qquad
 \eta(\mathbf q)\ge\kappa.
 \label{eq:v43-panel}
\end{equation}
A finite source/observable bank is \emph{strictly nonexceptional} if every
nonlocal connected coefficient in the bank is evaluated on one such compact
panel, with one common $\kappa>0$.  Purely local contact coefficients are kept
in the V7 local bank and do not require a soft momentum integral.
\end{definition}

Thus every external leg is Euclidean off shell and no proper channel carries
zero total momentum.  Zero-momentum integrated operator insertions, threshold
limits and on-shell configurations are not part of this panel.

\begin{remark}[Why the panel is the right fixed quantifier]
At fixed $\kappa>0$ a massless soft cluster can never be confused with a hard
external transfer: every nonzero partial external momentum is separated from
zero by $\kappa$.  This is the infrared analogue of the fixed finite external
panel used in V2 and V6.  The limit $\kappa\downarrow0$ is a distinct
exceptional-momentum problem and is not taken in V43.
\end{remark}

\section{Downward soft shells and the proof cutoff}
\label{sec:v43-soft-shells}

Fix $b>1$ and a normalization scale $\mu_R$.  Below $\mu_R$ introduce downward
scales
\begin{equation}
 \lambda_s=\mu_R b^{-s},
 \qquad s=0,1,2,\ldots.
 \label{eq:v43-soft-scales}
\end{equation}
For a floor $\rho$, stop after the largest $s$ with $\lambda_s\gtrsim\rho$.
This does not define a new renormalization scheme.  It is a coefficientwise
organization of the already-renormalized V5 forest expansion by the smallest
massless momentum scale appearing in a term.

Equivalently one may insert, only for the purpose of proving convergence, a
smooth radial cutoff $\chi_\rho(p)$ on each massless internal line with
$\chi_\rho(p)=0$ for $|p|\le\rho/2$ and $\chi_\rho(p)=1$ for $|p|\ge\rho$.
The cutoff is never used as a gauge-breaking definition of the theory; it is a
majorant device on the common renormalized graph integrand.  The physical
coefficient is the limit as the device is removed.

At prescribed $(L,\Delta,r,N)$, V5--V6 give only finitely many decorated graph,
forest, source and line topologies.  Denote this finite set by
$\mathfrak G_{L,\Delta,r,N}$.  For a graph $\Gamma$ let
$R\Gamma(\mathbf q)$ denote its already UV-forest-prepared Euclidean integrand,
including all V5 local counterterm insertions and source/contact descendants.

\section{The finite-order Euclidean soft-sector lemma}
\label{sec:v43-soft-sector}

The load-bearing point is not a new ultraviolet estimate.  It is the standard
Euclidean fact, now stated in the exact finite packet needed here, that
nonexceptional external momentum creates a positive scaling margin in every
massless soft sector after UV forests and source-free vacuum pieces have been
removed.

\begin{lemma}[Strict soft-sector margin on a nonexceptional panel]
\label{lem:v43-soft-margin}
Fix $(L,\Delta,r,N)$ and a compact strictly nonexceptional panel
$\mathcal K_{\kappa,Q}$.  For every
$\Gamma\in\mathfrak G_{L,\Delta,r,N}$ contributing to a normalized connected
coefficient, and every fixed retained external/source derivative, the
forest-renormalized Euclidean integral admits a finite infrared sector
decomposition for which every genuine soft radial variable $t$ occurs with
power
\begin{equation}
 t^{\sigma_{\Gamma,\alpha}-1}
 (1+|\log t|)^{m_{\Gamma,\alpha}},
 \qquad
 \sigma_{\Gamma,\alpha}>0.
 \label{eq:v43-sector-power}
\end{equation}
The constants and sector charts are uniform for
$\mathbf q\in\mathcal K_{\kappa,Q}$.  Hence, because the graph/source packet is
finite, there exist
\begin{equation}
 \boxed{
 \sigma_*:=\min_{\Gamma,\alpha}\sigma_{\Gamma,\alpha}>0,
 \qquad
 M_*:=\max_{\Gamma,\alpha}m_{\Gamma,\alpha}<\infty.}
 \label{eq:v43-common-margin}
\end{equation}
The same statement holds after every retained momentum derivative through
$\Delta$ and source/BV derivative through $r$.
\end{lemma}

\begin{proof}
We give the finite-packet argument because this is exactly where a hidden
massless assumption could otherwise enter.  After the V5 UV forest operation,
a graph integrand is a finite sum of terms whose denominators are ordinary
Euclidean massless or massive propagator denominators and whose numerators are
local polynomials in loop/external momenta and fixed source tensors.  No EFT,
BV or source insertion introduces a new infrared denominator.

Choose a Hepp-type sector for the loop variables and successively scale every
nested set of momenta which tends to zero.  A denominator in a given scaling
cluster is of one of two types.  Either its momentum is homogeneous in the
soft variables, in which case its canonical propagator power contributes to
the soft degree, or it carries a nonzero partial external transfer.  In the
second case nonexceptionality gives, once the soft variables are below a fixed
fraction of $\kappa$,
\[
 |p_{\rm soft}+Q_S|\ge c\kappa,
 \qquad |Q_S|\ge\kappa,
\]
so that line is uniformly hard and cannot contribute an infrared pole.

A marginal soft sector with no such hard denominator would separate a proper
set of external/source legs with zero total transfer.  This is excluded by
\eqref{eq:v43-panel}.  A component carrying no external/source leg is a vacuum
soft component; source-free vacuum components cancel in the normalized
functional $W[J]-W[0]$, while a local vacuum descendant already belongs to the
finite V7 normalization bank.  Thus no logarithmically marginal zero-transfer
soft component remains in a nonlocal normalized coefficient.

The renormalizable Standard-Model vertices are local and have the usual
nonnegative soft numerator degree once a hard transfer is held fixed.  At a
fixed EFT weight, higher-dimensional and gravitational vertices are again local
polynomials; they can increase a numerator power or add fields, but cannot add
a massless denominator.  Hence they cannot convert a strictly convergent soft
sector into a divergent one.  Massive internal lines are uniformly harmless.
External momentum derivatives can raise powers only on denominators which
carry the same nonzero transfer and therefore remain bounded below by
$c\kappa$; the fixed source/BV alphabet likewise only inserts finitely many
local numerators and already-stored contact terms.

The radial Jacobians therefore leave a strictly positive exponent in every
actual soft sector.  There are finitely many sectors for each graph and
finitely many graphs/derivatives at the prescribed packet, so their minimum is
the positive number $\sigma_*$ in \eqref{eq:v43-common-margin}.  Coincident
sector exponents can generate only a finite power of logarithms, absorbed in
$M_*$.  Compactness of $\mathcal K_{\kappa,Q}$ makes all hard-denominator
constants uniform.
\end{proof}

\begin{firewall}[What is doing the infrared work]
\label{fw:v43-pq}
The strict inequality $\sigma_*>0$ is a property of \emph{nonexceptional
Euclidean kinematics plus the normalized/UV-renormalized forest integrand}.
It is not supplied by the positive normalization point $\mu_R$, and it is not
a consequence of ultraviolet renormalizability alone.  If
$\eta(\mathbf q)=0$, the hard-denominator step in the proof can fail and V43's
conclusion need not hold.
\end{firewall}

\section{One-soft-shell estimate and summability}
\label{sec:v43-one-shell}

Let $\Delta_s W_\Gamma$ denote the contribution of a fixed graph/descendant
whose smallest active massless scale lies in
$[\lambda_{s+1},\lambda_s]$.  Sector rescaling by $\kappa$ gives the direct
soft analogue of the V2/V6 ultraviolet one-excess estimate.

\begin{theorem}[Infrared one-excess estimate]
\label{thm:v43-ir-one-excess}
Under Lemma~\ref{lem:v43-soft-margin}, for every retained derivative label
$I$ and momentum multi-index $\alpha$ there is a constant $C_{I,\alpha}$ such
that for all sufficiently soft shells $\lambda_s\le c\kappa$,
\begin{equation}
 \boxed{
 \sup_{\mathbf q\in\mathcal K_{\kappa,Q}}
 \left|
 \partial_I\partial_{\mathbf q}^{\alpha}
 \Delta_s W_\Gamma(\mathbf q)
 \right|
 \le
 C_{I,\alpha}
 \left({\lambda_s\over\kappa}\right)^{\sigma_*}
 (1+s)^{M_*}.}
 \label{eq:v43-ir-one-excess}
\end{equation}
The same bound holds after summing the finite graph/forest/source packet at
fixed $(L,\Delta,r,N)$.
\end{theorem}

\begin{proof}
Use the sector coordinates of Lemma~\ref{lem:v43-soft-margin} and scale every
soft radial variable by $\kappa$.  Restricting the smallest variable to one
geometric shell replaces its radial integral by
\[
 \int_{\lambda_{s+1}/\kappa}^{\lambda_s/\kappa}
 t^{\sigma-1}(1+|\log t|)^M\,dt
 \le
 C\left({\lambda_s\over\kappa}\right)^\sigma(1+s)^M.
\]
Replace $\sigma$ by the finite common minimum $\sigma_*$ and $M$ by $M_*$.  All
remaining sector variables have integrable majorants from the forest-prepared
integral.  The finite number of terms is absorbed into $C_{I,\alpha}$.
\end{proof}

\begin{corollary}[Summability of all low shells]
\label{cor:v43-ir-summable}
For every fixed packet,
\begin{equation}
 \sum_{s\ge s_0}
 \sup_{\mathbf q\in\mathcal K_{\kappa,Q}}
 \left|
 \partial_I\partial_{\mathbf q}^{\alpha}
 \Delta_s W(\mathbf q)
 \right|<\infty.
 \label{eq:v43-ir-sum}
\end{equation}
\end{corollary}

\begin{proof}
Because $\lambda_s=\mu_Rb^{-s}$, the right side of
\eqref{eq:v43-ir-one-excess} is a geometric sequence times a fixed polynomial
in $s$.
\end{proof}

\section{Zero-floor Cauchy theorem}
\label{sec:v43-zero-floor}

Let $W_{\infty,\rho}$ denote the V6 UV-continuum connected coefficient at fixed
normalization point $\mu_R$, followed by the coefficientwise soft-shell
completion down to floor $\rho$.  The subscript $\infty$ refers only to the
ultraviolet cutoff removal.  No all-orders perturbative sum is implied.

\begin{theorem}[Infrared-safe zero-floor limit]
\label{thm:v43-zero-floor}
Fix finite $(L,\Delta,r,N)$, a protected V24/V42 native branch, one coherent V7
normalization scheme at $\mu_R>0$, and a finite strictly nonexceptional
Euclidean source/observable panel $\mathcal K_{\kappa,Q}$.  Then every retained
normalized connected coefficient has a unique limit
\begin{equation}
 \boxed{
 W_{\infty,0}(\mathbf q)
 :=\lim_{\rho\downarrow0}W_{\infty,\rho}(\mathbf q)}
 \label{eq:v43-zero-floor-limit}
\end{equation}
uniformly on the panel and in every retained momentum/source derivative.
Moreover, for $0<\rho\le c\kappa$,
\begin{equation}
 \boxed{
 \|W_{\infty,\rho}-W_{\infty,0}\|_{L,\Delta,r,N;\mathcal K}
 \le
 C
 \left({\rho\over\kappa}\right)^{\sigma_*}
 \left[1+\left|\log{\rho\over\kappa}\right|^{M_*}\right].}
 \label{eq:v43-zero-floor-rate}
\end{equation}
No new infrared counterterm or lower-order refit is required.
\end{theorem}

\begin{proof}
Choose $s(\rho)$ with
$\lambda_{s(\rho)+1}<\rho\le\lambda_{s(\rho)}$.  The difference between two
floors is the sum of the intervening low-shell increments.  Apply
Theorem~\ref{thm:v43-ir-one-excess} and sum the geometric-polynomial tail.
This proves the Cauchy property and the stated rate.  The V7 normalization
coordinates are fixed at $\mu_R$ before the soft completion; since the low-shell
tail is absolutely convergent, no local retuning is needed to define the
limit.
\end{proof}

\begin{corollary}[No auxiliary massless infrared regulator on the declared panel]
\label{cor:v43-no-ir-reg}
On $\mathcal K_{\kappa,Q}$ the coefficientwise continuum of V7 may be completed
to zero loop momentum with no positive auxiliary infrared floor.  The only
remaining positive scale is the ordinary renormalization/reference point
$\mu_R$, which labels coordinates rather than regulating soft momentum.
\end{corollary}

\section{Compatibility with the ultraviolet continuum}
\label{sec:v43-uv-ir}

The V43 estimate is independent of the ultraviolet tail once the V5 forests
have been prepared.  Conversely the V6 ultraviolet one-excess estimate only
uses a finite hard scale associated with the fixed external/reference panel and
does not require the infrared floor to stay positive.  Let
\begin{equation}
 Q_R:=1+Q+\mu_R+m_{\rm ref}
 \label{eq:v43-hardscale}
\end{equation}
stand schematically for the finite maximum of the declared external,
normalization and massive reference scales entering the packet.

\begin{theorem}[Commuting UV removal and nonexceptional zero-floor limit]
\label{thm:v43-uv-ir-commute}
For a finite UV cutoff with top shell $\Lambda_J$ and a soft floor $\rho$, the
same normalized coefficient satisfies, after enlarging the packet constant,
\begin{equation}
 \boxed{
 \|W_{J,\rho}-W_{\infty,0}\|
 \le
 C_{\rm UV}{Q_R\over\Lambda_J}
 +C_{\rm IR}
 \left({\rho\over\kappa}\right)^{\sigma_*}
 \left[1+\left|\log{\rho\over\kappa}\right|^{M_*}\right].}
 \label{eq:v43-two-tail}
\end{equation}
Hence the ultraviolet and infrared-floor limits commute coefficientwise on the
strictly nonexceptional panel.
\end{theorem}

\begin{proof}
Insert the common intermediate coefficient $W_{\infty,\rho}$.  The first
difference is the V6 ultraviolet Cauchy tail, with the finite reference/external
scale denoted here by $Q_R$ rather than the overloaded V1--V42 symbol $\mu$.
The second difference is Theorem~\ref{thm:v43-zero-floor}.  The two bounds are
independent, so either limit may be taken first.
\end{proof}

\section{BV/master/source identities and perturbative projectivity}
\label{sec:v43-identities}

The low-shell completion is performed with the same graph, forest, source and
canonical-pair alphabet as V5.  At every positive stopping floor the finite
perturbative shell sum therefore satisfies the same projected common
master/source relation.  V43 adds only a summable tail.

\begin{theorem}[Common identities survive the zero-floor limit]
\label{thm:v43-master}
On a strictly nonexceptional panel,
\begin{equation}
 \boxed{
 \mathbf M_{L,\Delta,r,N}(W_{\infty,0})=0,}
 \label{eq:v43-master}
\end{equation}
where $\mathbf M_{L,\Delta,r,N}$ is the same finite common
BV/BRST/master/source vector used in V6.  The statement includes every retained
stress, composite, contact, ghost/antifield and determinant/Pfaffian line
coefficient whose nonlocal momentum arguments belong to the declared panel.
\end{theorem}

\begin{proof}
At every positive soft stopping scale, forest preparation and the common V5
restorer are unchanged; the coefficientwise identity is therefore the same
finite algebraic relation.  The vector $\mathbf M$ is a finite polynomial in
the retained coefficients and canonical contractions.  Theorem
\ref{thm:v43-zero-floor} gives uniform convergence for every coordinate entering
it, so continuity passes the identity to $\rho=0$.
\end{proof}

\begin{theorem}[Exact loop-order projectivity after IR completion]
\label{thm:v43-projective}
Let $\pi_L^{L+1}$ be the V7 perturbative truncation.  The zero-floor completed
states satisfy
\begin{equation}
 \boxed{
 \pi_L^{L+1}W_{\infty,0}^{[\le L+1]}
 =W_{\infty,0}^{[\le L]}.}
 \label{eq:v43-projective}
\end{equation}
No lower-order IR counterterm is rediscovered at higher terminal order.
\end{theorem}

\begin{proof}
At every positive floor the graph/forest expansion commutes exactly with loop
truncation by V7.  The low-shell tails at orders $\le L$ are uniformly
summable by Theorem~\ref{thm:v43-ir-one-excess}.  Take the zero-floor limit on
both sides of the finite-floor truncation identity.
\end{proof}

\section{Reference-scale covariance after removing the IR floor}
\label{sec:v43-reference}

A central consequence of the upstream split is that a massless limit does not
require a zero subtraction point.

\begin{theorem}[Positive reference covariance commutes with the zero-floor limit]
\label{thm:v43-reference-cov}
Let $\mu_R$ and $\mu_R'$ lie in one compact positive V7 reference interval and
let
$F_{\infty,L}^{\mu_R'\leftarrow\mu_R}$ be the V7 continuum scheme/reference
map.  On every common strictly nonexceptional panel,
\begin{equation}
 \boxed{
 W_{\infty,0}^{\mu_R'}
 =F_{\infty,L}^{\mu_R'\leftarrow\mu_R}
 \bigl(W_{\infty,0}^{\mu_R}\bigr).}
 \label{eq:v43-reference-cov}
\end{equation}
Thus the zero-floor limit is independent of the chosen positive subtraction
point up to the already-proved finite coordinate change.  No
$\mu_R\downarrow0$ operation is required or asserted.
\end{theorem}

\begin{proof}
The V7 finite reference map acts on the finite normalized local/source/line
coordinates and is smooth on the compact positive interval.  At every positive
soft floor it intertwines the two coherent normalized descriptions.  Theorem
\ref{thm:v43-zero-floor} gives uniform convergence of both sides and all
retained coordinates, so the finite continuous map passes through the limit.
\end{proof}

\begin{firewall}[A renormalization scale is not an infrared divergence]
The appearance of $\mu_R$ in logarithms or running couplings after
Theorem~\ref{thm:v43-reference-cov} is ordinary renormalization-scale
dependence.  It is related between schemes by the V7 finite maps and by the
usual coefficientwise RG equations.  V43 removes the auxiliary soft floor; it
does not attempt to erase physical running by setting the subtraction point to
zero.
\end{firewall}

\section{Thermodynamic order and long-range massless kernels}
\label{sec:v43-volume}

V6 proved a particularly strong position-space periodization estimate using
arbitrary polynomial moments weighted by a positive scale.  Such a bound is
not uniform as the massless infrared floor is removed: genuine massless
correlators have long-range power tails.  V43 therefore keeps the logically
correct staged order for its strongest theorem.

\begin{proposition}[Safe staged order for the V43 theorem]
\label{prop:v43-staged-volume}
Take the V6 ultraviolet and thermodynamic limits at one fixed positive
reference scale, obtaining the infinite-volume renormalized Euclidean graph
coefficients.  Then perform the V43 zero-floor completion on a compact
nonexceptional momentum panel.  The resulting coefficient is unique and obeys
Theorems~\ref{thm:v43-zero-floor}--\ref{thm:v43-reference-cov}.

V43 does not claim a uniform position-space periodization estimate as
$\rho\downarrow0$ and $R\to\infty$ simultaneously.
\end{proposition}

\begin{proof}
The V6 volume theorem supplies the infinite-volume forest-renormalized
coefficient at fixed positive reference data.  Lemma~\ref{lem:v43-soft-margin}
is an infinite-volume Euclidean momentum estimate and therefore applies to
that coefficient directly.  Long-distance power tails affect uniform
periodization, not convergence of the nonexceptional momentum integral.
\end{proof}

\section{What remains excluded}
\label{sec:v43-exclusions}

\begin{firewall}[Exceptional and on-shell kinematics]
\label{fw:v43-exceptional}
V43 does not cover configurations for which
\[
 \eta(\mathbf q)\to0.
\]
This includes zero-momentum channels, soft external gauge/gravity legs,
collinear/on-shell limits after analytic continuation, and generic massless
S-matrix elements.  In those regimes a soft sector can become marginal and
$\sigma_*$ may vanish.  Physical on-shell probabilities require their own
soft-factorization/dressing or inclusive KLN-type theorem; that problem is not
hidden inside Theorem~\ref{thm:v43-zero-floor}.
\end{firewall}

\begin{firewall}[Zero-momentum local insertions]
A local normalization/contact coefficient evaluated at zero momentum remains a
coordinate of the finite V7 bank.  V43 does not infer its massless infrared
finiteness from a nonexceptional theorem by continuity.  If a zero-momentum
composite observable is desired, it must either be protected by an independent
Ward identity or be treated in the later exceptional-momentum programme.
\end{firewall}

\section{V43 closure theorem}
\label{sec:v43-close}

\begin{theorem}[Source-complete nonexceptional massless IR completion]
\label{thm:v43-close}
Fix:
\begin{enumerate}[label=\textup{(I\arabic*)},leftmargin=3em]
\item the V42 cofinal V24-protected native overlap branch;
\item finite $(L,\Delta,r,N)$ and one coherent V7 normalization/source/reference
scheme at a nonzero subtraction point $\mu_R$;
\item the V6 infinite-volume coefficientwise UV continuum;
\item a finite strictly nonexceptional Euclidean external/source panel with
$0<\kappa\le\eta(\mathbf q)$ and $|q_i|\le Q$.
\end{enumerate}
Then the auxiliary massless infrared floor can be removed coefficientwise:
\begin{equation}
 \boxed{
 \rho\downarrow0:
 \quad
 W_{\infty,\rho}^{[\le L]}
 \longrightarrow
 W_{\infty,0}^{[\le L]}}
 \label{eq:v43-final-limit}
\end{equation}
uniformly with every retained momentum/source derivative.  The low-shell tail
obeys the power-log estimate \eqref{eq:v43-zero-floor-rate}; the UV and IR
limits commute as in \eqref{eq:v43-two-tail}; the common
BV/BRST/master/source identities survive; loop-order projectivity remains
exact; and positive renormalization points remain related by the V7 finite
reference maps.

Accordingly the V1--V42 phrase ``positive IR/reference scale'' can be sharpened
on this observable class to:
\begin{equation}
 \boxed{
 \text{no auxiliary IR floor, but one ordinary nonzero renormalization point}.}
 \label{eq:v43-refined-scale}
\end{equation}
The theorem does not extend to exceptional or on-shell massless kinematics.
\end{theorem}

\begin{proof}
Lemma~\ref{lem:v43-soft-margin} gives a common strictly positive infrared
sector margin.  Theorems~\ref{thm:v43-ir-one-excess} and
\ref{thm:v43-zero-floor} sum all downward soft shells.  Theorem
\ref{thm:v43-uv-ir-commute} combines this with the V6 ultraviolet tail.
Theorems~\ref{thm:v43-master} and \ref{thm:v43-projective} pass the common
identities and loop-order chronology through the zero-floor limit.  Theorem
\ref{thm:v43-reference-cov} keeps the renormalization point as a finite
coordinate choice rather than misidentifying it with an infrared regulator.
\end{proof}

\section{V43 anti-circularity audit}
\label{sec:v43-audit}

\begin{enumerate}[leftmargin=2.4em]
\item \textbf{The subtraction point was not sent to zero.}  V43 removes the
auxiliary soft floor $\rho$, not the positive renormalization point $\mu_R$.
\item \textbf{Nonexceptionality is quantitative.}  The proof uses the fixed
margin $\eta(\mathbf q)\ge\kappa>0$; it does not infer exceptional-momentum
control by continuity.
\item \textbf{No new IR counterterms are postulated.}  The V5 UV forests and V7
normalization data are frozen.  The soft tail is absolutely summable at the
fixed packet.
\item \textbf{Gravity as an EFT is retained.}  Only finitely many local EFT
vertices occur at fixed $(L,\Delta,N)$; they add polynomial numerators and no
new massless denominators, so the finite soft-sector proof remains valid.
\item \textbf{The proof cutoff is not a gauge regulator.}  A smooth $\rho$
cutoff may be used to dominate the integral, but the common identities are
those of the unmodified forest/shell algebra and are passed by convergence.
\item \textbf{The V6 position-space volume bound is not extrapolated to zero
mass.}  V43 takes the infinite-volume coefficient first and proves the
zero-floor limit in Euclidean momentum space.
\item \textbf{On-shell physics is not claimed.}  No unrestricted massless
S-matrix or KLN cancellation is inferred from off-shell Euclidean finiteness.
\end{enumerate}

\section{V43 status ledger}
\label{sec:v43-status}

\subsection*{Proved in V43}
\begin{enumerate}[leftmargin=2.2em]
\item The former positive ``IR/reference'' datum is separated into an ordinary
nonzero renormalization point $\mu_R$ and an auxiliary soft floor $\rho$; only
the latter is a candidate for removal.
\item A compact strictly nonexceptional Euclidean source/observable panel is
defined by a uniform lower bound $\kappa$ on every nontrivial partial external
momentum sum.
\item For every fixed finite graph/source/BV/line packet, the V5
forest-renormalized Euclidean integrands have a common strictly positive soft
sector exponent $\sigma_*>0$ on such a panel, with at most a finite logarithmic
power $M_*$.  Fixed EFT/gravity insertions do not introduce additional
infrared denominators.
\item Each downward massless shell obeys the source-complete bound
$C(\lambda_s/\kappa)^{\sigma_*}(1+s)^{M_*}$; the complete low-shell tail is
summable.
\item Every retained normalized connected coefficient therefore has a unique
zero-floor limit with the uniform power-log rate
\eqref{eq:v43-zero-floor-rate}.
\item UV removal and zero-floor completion commute coefficientwise, with the
two independent tails displayed in \eqref{eq:v43-two-tail}.
\item The common projected BV/BRST/master/source identities and V7 exact
perturbative-order projectivity survive the zero-floor limit.
\item V7 positive reference-scale covariance commutes with zero-floor
completion.  No $\mu_R\downarrow0$ limit is needed to define the massless
nonexceptional coefficients.
\end{enumerate}

\subsection*{Inherited}
\begin{enumerate}[leftmargin=2.2em]
\item V1--V7: the native source-complete arbitrary-fixed-order UV continuum,
common BV restoration, multi-shell Cauchy theorem, finite reference covariance
and exact perturbative projectivity.
\item V8--V24: regulator universality/classification, controlled zero-mode
modules and the V24 common-soft Feshbach theorem.
\item V25--V42: real orientation, Wilson-wall/eta/global-phase analysis,
large-gauge cancellation, microscopic Picard no-go and exact contraction of
the cofinal complex regulator family.
\end{enumerate}

\subsection*{Conditional}
\begin{enumerate}[leftmargin=2.2em]
\item V43 is an off-shell Euclidean theorem on a compact strictly
nonexceptional panel.  Its constants may diverge as $\kappa\downarrow0$.
\item The graph/source packet is finite at prescribed $(L,\Delta,r,N)$; no
uniformity as perturbative or EFT order tends to infinity is asserted.
\item The volume limit is staged as in Proposition~\ref{prop:v43-staged-volume};
a joint $R\to\infty$, $\rho\downarrow0$ position-space bound for massless
long-range kernels is not proved.
\item Purely local zero-momentum operator coordinates remain in the V7
normalization bank unless an independent infrared theorem applies to them.
\end{enumerate}

\subsection*{Open beyond V43}
\begin{enumerate}[leftmargin=2.2em]
\item Quantify the approach to exceptional Euclidean momentum strata
$\kappa\downarrow0$ and identify the exact soft/collinear factorization
operators responsible for any logarithmic or power singularities.
\item For physical on-shell amplitudes, formulate and prove an inclusive or
dressed soft theorem compatible with the common BV/source packet before any
massless S-matrix claim is made.
\item Analytic finite-coupling RG trajectories, perturbative summability and a
nonperturbative Einstein--Standard Model completion remain independent stronger
problems.
\end{enumerate}

\subsection*{Exact next load-bearing theorem}
\begin{center}
\fbox{\parbox{0.92\textwidth}{
\textbf{Exceptional-momentum stratification and soft-factorization theorem.}
V43 removes the auxiliary infrared floor at every fixed perturbative/EFT/source
order for strictly nonexceptional Euclidean coefficients while retaining only
an ordinary nonzero renormalization point.  The next load-bearing infrared gate
is to let the exceptionality modulus $\kappa$ approach zero in a controlled
stratum.  For each soft or collinear forest, identify the universal factor
multiplying the hard V43 coefficient, prove a sharp power/log bound in
$\kappa$, and show compatibility with the V5 BV/source forest recursion and V7
projective normalization.  Only after that factorization is established should
one formulate an inclusive/dressed on-shell cancellation theorem; no raw
massless S-matrix limit is to be inferred from V43.}}
\end{center}

\section*{Append-only continuation marker: V43}
\addcontentsline{toc}{section}{Append-only continuation marker: V43}
\textbf{End of V43.}  The complete V1--V42 body above is retained verbatim.
V43 goes upstream on the long-standing positive ``IR/reference'' assumption by
separating an ordinary nonzero renormalization point from the auxiliary
infrared floor.  On a compact strictly nonexceptional Euclidean source panel,
a finite soft-sector decomposition of the already UV-forest-renormalized graph
packet has one common positive infrared exponent.  This gives a summable
downward-shell estimate, a unique zero-floor limit with an explicit power-log
tail, commuting UV and IR removals, preservation of the common BV/master/source
identities and exact loop-order projectivity, and compatibility with the V7
positive-reference maps.  The theorem removes the auxiliary massless IR floor
but deliberately retains the renormalization point and excludes exceptional or
on-shell massless kinematics.  The next stronger gate is the stratified
$\kappa\downarrow0$ factorization problem.


% ============================================================================
% V44 APPEND-ONLY CONTINUATION
% ============================================================================
\clearpage
\part{V44 append-only continuation:\\
resolved exceptional Euclidean strata and source-complete soft factorization}
\label{part:v44}

\section{Purpose and the Euclidean upstream correction}
\label{sec:v44-purpose}

V43 removed the auxiliary infrared floor on every compact strictly
nonexceptional Euclidean panel.  Its exact next gate was phrased as an
``exceptional-momentum stratification and soft/collinear factorization
problem.''  Before attacking that gate, one further separation is necessary.
In real Euclidean signature there is no independent collinear pinch surface of
the massless on-shell type: a Euclidean vector satisfying $p^2=0$ is the zero
vector.  Hence, while external momenta remain in a compact real Euclidean
panel, every new infrared degeneration is generated by one or more channel
sums tending to zero.  Genuine soft--collinear overlap is a later
Minkowski-boundary-value problem and is not needed to resolve the Euclidean
exceptional set.

V44 therefore studies the already-defined V43 coefficient
$W_{\infty,0}(\mathbf q)$ as one approaches the arrangement of exceptional
Euclidean channel hyperplanes.  The result is not a claim that all exceptional
limits are finite.  Instead V44 resolves the exceptional arrangement, proves a
polyhomogeneous power--log expansion on every boundary face, and identifies
each coefficient as a universal soft forest multiplying a hard V43 quotient
coefficient.  Thus divergences, logarithms and finite exceptional limits are
classified rather than hidden behind a vanishing nonexceptionality constant.

\begin{firewall}[Euclidean exceptional momentum is not yet a collinear theorem]
\label{fw:v44-euclidean-not-collinear}
All momenta in V44 are real Euclidean momenta.  The theorem does not perform
analytic continuation to a physical Minkowski sheet, does not impose
$p_i^2=0$ with nonzero null $p_i$, and does not invoke KLN, jet or dressed-state
cancellation.  Those questions require a separate boundary-value/pinch-surface
theorem after V44.
\end{firewall}

\section{Exceptional channel arrangement and its resolved strata}
\label{sec:v44-arrangement}

Let $I_n=\{1,\ldots,n\}$ label the nonlocal external and source insertions of a
connected coefficient, with $\sum_iq_i=0$.  For a proper nonempty subset
$A\subset I_n$ define the channel momentum
\begin{equation}
 Q_A(\mathbf q):=\sum_{i\in A}q_i.
 \label{eq:v44-channel}
\end{equation}
The complementary subsets $A$ and $A^c$ define the same channel up to sign.
Write $[A]$ for this unoriented channel and
\begin{equation}
 \Sigma_{[A]}:=\{\mathbf q:Q_A(\mathbf q)=0\}
 \label{eq:v44-channel-face}
\end{equation}
for its exceptional locus.

\begin{definition}[Compatible exceptional family]
\label{def:v44-compatible}
Choose one rooted representative for every channel class.  A finite family
$\mathcal N$ of channel classes is \emph{compatible} if its rooted
representatives are pairwise nested or disjoint.  Equivalently, $\mathcal N$
can be represented by the vertices of a forest whose leaves are external
insertions.  Such a family records channel sums which may become small at
independent nested scales without forcing an unrecorded overlapping channel to
vanish.
\end{definition}

If several overlapping channels vanish simultaneously, enlarge them to the
finite nested building set generated by their intersections/unions.  Thus no
actual exceptional configuration is lost by working with compatible families;
one merely resolves a singular intersection into finitely many nested charts.

\begin{definition}[Resolved exceptional chart]
\label{def:v44-resolved-chart}
Fix a compatible family $\mathcal N$.  A compact resolved chart
$\widetilde{\mathcal K}_{\mathcal N}(Q,\kappa_0)$ consists of coordinates
\begin{equation}
 \left(
  \{\rho_A,\widehat Q_A\}_{[A]\in\mathcal N},
  \mathbf q_{\rm hard}
 \right),
 \qquad 0\le\rho_A\le\rho_0,
 \label{eq:v44-resolved-coordinates}
\end{equation}
with
\begin{equation}
 Q_A=\rho_A\widehat Q_A,
 \qquad |\widehat Q_A|=1,
 \label{eq:v44-polar-channel}
\end{equation}
and with every channel not generated by $\mathcal N$ bounded below by one
fixed number $\kappa_0>0$.  The angular variables and hard coordinates are
restricted to one compact set away from all further exceptionalities.
\end{definition}

These are the standard local coordinates of the iterated real blow-up of the
finite channel arrangement.  V44 only needs the finite collection of such
charts associated with the fixed source/observable packet; no global
algebro-geometric machinery is required.

\begin{lemma}[No independent Euclidean collinear boundary face]
\label{lem:v44-no-euclidean-collinear}
On a compact real Euclidean momentum panel, after the ultraviolet regions have
already been removed by V5--V6, every infrared degeneration of a massless
propagator sector is represented in a resolved chart by one or more
$\rho_A\downarrow0$.  There is no additional boundary hypersurface describing
nonzero massless collinearity.
\end{lemma}

\begin{proof}
A Euclidean massless denominator can become singular only when its Euclidean
momentum tends to zero.  If a collection of loop momenta remains bounded while
a denominator momentum tends to zero, momentum conservation expresses that
limit as a vanishing linear combination of soft loop variables and a partial
external/source sum $Q_A$.  After the soft loop variables are included in the
sector scaling, the remaining external condition is $Q_A\to0$.  By contrast,
a nonzero collinear on-shell pinch would require a nonzero vector with
$p^2=0$, impossible over real Euclidean momenta.  The finite resolved channel
arrangement therefore contains every new Euclidean infrared boundary face.
\end{proof}

\section{Soft subgraphs, quotient graphs and the infrared forest packet}
\label{sec:v44-soft-forest}

Fix one V5 forest-prepared decorated graph/source descendant $\Gamma$ in the
finite packet $\mathfrak G_{L,\Delta,r,N}$.  A \emph{soft subgraph} $\gamma$
of $\Gamma$ is a connected set of massless internal lines and vertices whose
internal loop variables scale to zero together and whose net external transfer
is one of the channel sums $Q_A$.  Massive lines may end on $\gamma$ but are
never part of its homogeneous soft denominator set.

A \emph{$\mathcal N$-soft forest} $F$ is a nested/disjoint family of such
subgraphs compatible with the channel forest $\mathcal N$.  Contracting every
maximal soft component to a local marked vertex gives a hard quotient graph
$\Gamma/F$.  The mark remembers the finite Taylor tensor with which the soft
component couples to the quotient graph.  Because the V5 packet was explicitly
made contraction-complete, every such marked quotient is already an allowed
local/source/contact descendant of the same finite packet after enlarging only
the finite Taylor order required for the chosen asymptotic accuracy.

\begin{definition}[Infrared radial degree]
\label{def:v44-ir-degree}
For a soft subgraph $\gamma$ and a retained Taylor/source label $\beta$, let
$\omega_{\rm IR}(\gamma,\beta)$ be the radial homogeneity exponent of its
already UV-forest-prepared Euclidean soft integral, including loop measures,
propagator homogeneities and the soft degree of its local numerator, so that a
single soft scale $t$ begins as
\begin{equation}
 t^{\omega_{\rm IR}(\gamma,\beta)-1}
 (\log t)^m\,dt.
 \label{eq:v44-ir-degree}
\end{equation}
The exponent may now be positive, zero or negative.  V43 was precisely the
case in which the relevant packet minimum was strictly positive.
\end{definition}

For a nested forest, each node has its own radial variable.  The local hard
Taylor expansion shifts $\omega_{\rm IR}$ by nonnegative integers.  Coincident
or resonant exponents generate only finite logarithmic powers at fixed graph
and fixed Taylor order.

\section{Single-channel soft factorization}
\label{sec:v44-single-channel}

The basic factorization mechanism is clearest when exactly one primitive
channel tends to zero.  Let
\begin{equation}
 Q_A=\rho\widehat Q,
 \qquad \rho\downarrow0,
 \label{eq:v44-single-scaling}
\end{equation}
while every other primitive channel remains at least $\kappa_0>0$.

For a soft forest $F$, denote by
$\mathcal S_{F,\beta}(\rho,\widehat Q)$ the renormalized integral of its soft
components after all hard dependence has been Taylor-expanded into a finite
local tensor labelled by $\beta$.  Denote by
$\mathcal H_{\Gamma/F,\beta}(\mathbf q_{\rm hard})$ the V43 zero-floor hard
quotient coefficient carrying that local tensor insertion.

\begin{theorem}[Single-channel source-complete soft factorization]
\label{thm:v44-single-factorization}
Fix finite $(L,\Delta,r,N)$, one compact resolved one-channel chart and every
retained source/BV/momentum derivative.  For every real accuracy threshold
$A$ there is a finite set $\mathfrak E_{\Gamma,A}$ of soft-forest/Taylor labels
such that
\begin{align}
 \partial_I W_{\Gamma}(\rho,\widehat Q,\mathbf q_{\rm hard})
 &={}
 \sum_{(F,\beta,j,m)\in\mathfrak E_{\Gamma,A}}
 \rho^{\lambda_{F,\beta}+j}
 (\log\rho)^m
 \mathcal S_{F,\beta,j,m}(\widehat Q)
 \mathcal H_{\Gamma/F,\beta,j,I}(\mathbf q_{\rm hard})
 \notag\\
 &\qquad
 +\mathcal R_{\Gamma,A,I}(\rho),
 \label{eq:v44-single-factorization}
\end{align}
where
\begin{equation}
 \lambda_{F,\beta}
 =\sum_{\gamma\in F}\omega_{\rm IR}(\gamma,\beta_\gamma)
 \label{eq:v44-single-exponent}
\end{equation}
with the appropriate nested relative scaling convention, $j\in\mathbb N_0$
comes from hard Taylor order, and
\begin{equation}
 \boxed{
 |\mathcal R_{\Gamma,A,I}(\rho)|
 \le C_{A,I}\rho^A
 \left(1+|\log\rho|^{M_{A,I}}\right).}
 \label{eq:v44-single-remainder}
\end{equation}
All constants are uniform in the compact angular/hard chart.

The factor $\mathcal S_{F,\beta,j,m}$ depends only on the corresponding soft
subgraph packet, its representation/source labels and the angular soft data.
The global hard graph enters only through the contracted V43 coefficient
$\mathcal H_{\Gamma/F,\beta,j,I}$.
\end{theorem}

\begin{proof}
Choose the finite Hepp/sector decomposition already used in V43, but do not use
nonexceptionality to declare the $Q_A$-carrying lines hard.  In each sector,
collect the loop variables which scale together with $Q_A$.  Every denominator
outside the resulting soft forest carries either a fixed hard loop scale or a
channel bounded below by $\kappa_0$; it is therefore analytic in the soft loop
variables and $Q_A$ on the compact chart.  Taylor-expand those hard
denominators and the local numerator to the finite order required to make the
remainder have radial degree $>A$.

Each Taylor coefficient separates into: (i) a homogeneous soft integral over
the soft forest and (ii) the forest-prepared quotient graph with a local tensor
insertion.  The first is $\mathcal S$; the second is exactly the declared
$\mathcal H$.  The V5 contraction-complete packet is what prevents the local
Taylor tensor from creating an unowned counterterm/source coordinate.

Radial integration of a homogeneous soft term gives powers determined by
Definition~\ref{def:v44-ir-degree}; equality of radial exponents or integration
of $t^{-1}$ generates a finite logarithmic power.  Taylor's integral remainder
has its soft degree raised by the chosen finite number of derivatives, hence
obeys \eqref{eq:v44-single-remainder}.  The fixed graph/source packet and the
compact angular/hard chart make every constant uniform.  Source/BV and retained
momentum derivatives may be taken before the finite Taylor expansion; they
only enlarge the finite numerator/Taylor labels and therefore satisfy the same
argument.
\end{proof}

\begin{corollary}[Exceptional trichotomy]
\label{cor:v44-trichotomy}
Let $\lambda_*$ be the smallest exponent appearing in the complete fixed-packet
single-channel expansion after summing graphs.
\begin{enumerate}[label=\textup{(\roman*)},leftmargin=2.7em]
\item If $\lambda_*>0$, the nonlocal coefficient extends continuously to the
exceptional face and the soft correction vanishes there.
\item If $\lambda_*=0$, the only possible leading singularity is a finite
polynomial in $\log\rho$.
\item If $\lambda_*<0$, the coefficient has at most the corresponding power--log
soft singularity $\rho^{\lambda_*}(1+|\log\rho|)^M$.
\end{enumerate}
Cancellations in the complete Standard-Model packet may improve the actual
leading exponent; they can never make the sector bound worse.
\end{corollary}

\section{Nested exceptional faces and polyhomogeneous resolution}
\label{sec:v44-polyhomogeneous}

The single-channel result iterates because a compatible channel family is a
forest.  Start with the smallest soft cluster, Taylor-expand its hard
complement, contract it, and proceed outward through the nesting tree.  Disjoint
clusters factor independently before their parent cluster is treated.

\begin{theorem}[Polyhomogeneous extension on the resolved exceptional space]
\label{thm:v44-polyhomogeneous}
Let $\beta:\widetilde{\mathcal K}\to\mathcal K$ denote the finite iterated
real blow-up of all exceptional channel strata meeting the declared compact
Euclidean panel.  For every fixed finite $(L,\Delta,r,N)$ and every retained
source/BV/momentum derivative, the pullback of the V43 zero-floor coefficient
has, on each resolved chart, a polyhomogeneous expansion
\begin{equation}
 \boxed{
 \beta^*\partial_I W_{\infty,0}
 \sim
 \sum_{\boldsymbol\lambda,\mathbf m}
 a_{\boldsymbol\lambda,\mathbf m,I}
 (\widehat{\mathbf Q},\mathbf q_{\rm hard})
 \prod_{A\in\mathcal N}
 \rho_A^{\lambda_A}(\log\rho_A)^{m_A}.}
 \label{eq:v44-polyhomogeneous}
\end{equation}
For every finite multiweight cutoff, only finitely many terms lie below that
cutoff and the remainder gains the prescribed positive power in every radial
variable.  The coefficient
$a_{\boldsymbol\lambda,\mathbf m,I}$ is a finite sum of products of universal
soft-forest factors and V43 hard quotient coefficients with local V5 packet
insertions.
\end{theorem}

\begin{proof}
Induct on the depth of the compatible channel forest.  Depth one is
Theorem~\ref{thm:v44-single-factorization}, applied independently to disjoint
soft clusters.  After contracting all depth-one clusters, the quotient graph
again belongs to the contraction-complete finite V5 packet and the remaining
channel variables form a compatible forest of one smaller depth.  Apply the
same Taylor/sector decomposition at the next radial scale.

At each stage the new exponents are the previous infrared degrees shifted by
nonnegative Taylor integers.  Hence below any prescribed multiweight there are
only finitely many terms.  Repeated integration of equal exponents produces
finite logarithmic powers but no new functional class.  The induction ends
because the external/source packet is finite.  Compatibility on overlaps of
resolved charts follows from uniqueness of asymptotic expansions: two sector
representations of the same coefficient have identical coefficients after
re-expansion on their common subchart.
\end{proof}

\begin{remark}[What ``universal soft factor'' means here]
Universality in Theorems~\ref{thm:v44-single-factorization} and
\ref{thm:v44-polyhomogeneous} is deliberately local to the declared
perturbative theory.  A soft factor depends on the soft subgraph topology,
field representations, couplings and source labels, but not on the detailed
hard graph in which that subgraph is embedded.  All hard dependence is carried
by the contracted coefficient and the finite local Taylor tensor.  No claim of
an all-orders eikonal exponentiation theorem is made.
\end{remark}

\section{Sharp one-parameter approach bounds}
\label{sec:v44-one-parameter}

A resolved chart permits arbitrary relative rates among exceptional channels.
To obtain a scalar asymptotic bound, choose positive weights $\nu_A$ and a
one-parameter approach
\begin{equation}
 \rho_A=\varepsilon^{\nu_A},
 \qquad \varepsilon\downarrow0.
 \label{eq:v44-weighted-approach}
\end{equation}
For an expansion monomial define its weighted exponent
\begin{equation}
 \Lambda(\boldsymbol\lambda;\boldsymbol\nu)
 :=\sum_{A\in\mathcal N}\nu_A\lambda_A.
 \label{eq:v44-weighted-exponent}
\end{equation}
Let $\Lambda_*$ be the minimum over the finite set of leading soft-forest
monomials contributing to the retained packet and let $M_*$ be the largest
associated total logarithmic degree.

\begin{theorem}[Stratified exceptional power--log bound]
\label{thm:v44-power-log}
Along \eqref{eq:v44-weighted-approach}, every retained coefficient satisfies
\begin{equation}
 \boxed{
 |\partial_IW_{\infty,0}(\varepsilon)|
 \le
 C_I\,\varepsilon^{\Lambda_*}
 \left(1+|\log\varepsilon|^{M_*}\right)}
 \label{eq:v44-power-log}
\end{equation}
for sufficiently small $\varepsilon$.  If the complete coefficient of at least
one monomial with exponent $\Lambda_*$ is nonzero, this exponent is the actual
leading power.  If complete-packet Ward or representation cancellations remove
all such monomials, recomputing the finite minimum gives the improved exponent.
\end{theorem}

\begin{proof}
Insert $\rho_A=\varepsilon^{\nu_A}$ into the finite leading portion of
\eqref{eq:v44-polyhomogeneous}; every monomial becomes
$\varepsilon^{\Lambda}(\log\varepsilon)^m$ times a bounded angular/hard
coefficient.  Choose the asymptotic remainder cutoff strictly above the
smallest weighted exponent.  Compactness bounds all coefficient functions and
gives \eqref{eq:v44-power-log}.  The final assertion follows from uniqueness
of the one-variable power--log asymptotic expansion.
\end{proof}

\section{No new counterterm and the exceptional hard coefficients}
\label{sec:v44-no-new-ct}

The exceptional factors in V44 are nonlocal infrared asymptotics.  They are
not new UV counterterms and are not to be subtracted into the V7 normalization
bank.

\begin{proposition}[Hard coefficients are already V43 coefficients]
\label{prop:v44-hard-v43}
Every $\mathcal H_{\Gamma/F,\beta}$ occurring in V44 is one of the following:
\begin{enumerate}[label=\textup{(\roman*)},leftmargin=2.7em]
\item a nonlocal V43 zero-floor quotient coefficient evaluated on a compact
strictly nonexceptional hard panel;
\item a local/contact descendant already owned by the finite V7 normalization
bank;
\item a normalized vacuum descendant which cancels from the connected
functional or reduces to case \textup{(ii)}.
\end{enumerate}
Consequently V44 introduces no new renormalization coordinate.
\end{proposition}

\begin{proof}
Contracting a maximal soft forest collapses every exceptional transfer carried
by that forest.  By construction of the resolved chart, every remaining
nonlocal channel of the quotient graph is bounded below by $\kappa_0$, so V43
applies to it.  A quotient component with no remaining nonlocal external
transfer is local and belongs to the V5/V7 contraction-complete bank.  Pure
vacuum pieces cancel in the normalized connected functional exactly as in V43.
\end{proof}

\begin{firewall}[Soft factorization is not infrared counterterm subtraction]
\label{fw:v44-not-ir-renormalization}
A negative or zero exponent in V44 is a physical exceptional-momentum
singularity of the renormalized coefficient.  V44 identifies and factors it;
it does not remove it by adding a local counterterm.  Local UV normalization
remains exactly the V5--V7 one.
\end{firewall}

\section{Source, BV and master-equation compatibility}
\label{sec:v44-bv}

Let $\operatorname{Asymp}_{\mathcal N}$ denote the coefficientwise
polyhomogeneous expansion map of Theorem~\ref{thm:v44-polyhomogeneous}.  It is
defined graph by graph before the finite graph/source sum.

\begin{theorem}[Source-complete factorization intertwines the common identities]
\label{thm:v44-master-factorization}
For every retained source/BV derivative $\partial_I$,
\begin{equation}
 \boxed{
 \operatorname{Asymp}_{\mathcal N}(\partial_IW)
 =\partial_I\operatorname{Asymp}_{\mathcal N}(W).}
 \label{eq:v44-source-commute}
\end{equation}
Moreover, if $\mathbf M_{L,\Delta,r,N}$ is the common V5/V43 projected
BV/BRST/master/source vector, then
\begin{equation}
 \boxed{
 \operatorname{Asymp}_{\mathcal N}
 \bigl(\mathbf M_{L,\Delta,r,N}(W_{\infty,0})\bigr)
 =0,}
 \label{eq:v44-master-asymp}
\end{equation}
and therefore every individual power--log coefficient of the factorized
expansion satisfies the induced soft/hard master relation, including the
contact terms generated when a canonical contraction collapses into a soft
forest vertex.
\end{theorem}

\begin{proof}
Source/BV differentiation acts on the same finite local numerator, propagator
and marked-vertex alphabet used in the sector/Taylor construction, so it may be
commuted through every finite asymptotic Taylor operation and through the
absolutely convergent hard quotient integrals.  The contraction-complete V5
packet stores precisely the local descendants created when differentiated or
contracted endpoints enter the same soft component.

At every nonexceptional point, Theorem~\ref{thm:v43-master} gives
$\mathbf M(W_{\infty,0})=0$.  Apply the finite sector/Taylor expansion to this
identically vanishing function.  The master vector is a finite polynomial in
the retained coefficient coordinates and canonical contractions, so products
of polyhomogeneous series are again polyhomogeneous.  Uniqueness of the
resolved power--log expansion forces every coefficient to vanish.  This is
exactly the induced factorized master relation.
\end{proof}

\section{Perturbative projectivity and reference covariance}
\label{sec:v44-projective-reference}

\begin{theorem}[V7 chronology survives exceptional factorization]
\label{thm:v44-projective}
Let $\pi_L^{L+1}$ be the V7 loop-order truncation.  Then
\begin{equation}
 \boxed{
 \pi_L^{L+1}
 \operatorname{Asymp}_{\mathcal N}
 W_{\infty,0}^{[\le L+1]}
 =
 \operatorname{Asymp}_{\mathcal N}
 W_{\infty,0}^{[\le L]}.}
 \label{eq:v44-projective}
\end{equation}
No exceptional soft factor at order $L+1$ refits a lower-order hard or local
coefficient.
\end{theorem}

\begin{proof}
At every strictly nonexceptional point V43 already satisfies exact V7
projectivity.  The exceptional expansion is obtained by finite graphwise
Taylor/sector operations which preserve loop order.  Expanding the exact
truncation identity on a resolved chart and using uniqueness of the
polyhomogeneous coefficients proves \eqref{eq:v44-projective}.
\end{proof}

\begin{theorem}[Positive reference covariance of the factorized coefficients]
\label{thm:v44-reference}
For two positive V7 normalization points $\mu_R,\mu_R'$ in one compact
reference interval, the complete V44 expansion is intertwined by the same
finite V7 reference map:
\begin{equation}
 \boxed{
 \operatorname{Asymp}_{\mathcal N}
 W_{\infty,0}^{\mu_R'}
 =
 F_{\infty,L}^{\mu_R'\leftarrow\mu_R}
 \left(
  \operatorname{Asymp}_{\mathcal N}
  W_{\infty,0}^{\mu_R}
 \right).}
 \label{eq:v44-reference}
\end{equation}
The map acts on the hard/local coefficient coordinates; the geometric
exceptional variables $\rho_A,\widehat Q_A$ are untouched.
\end{theorem}

\begin{proof}
Use Theorem~\ref{thm:v43-reference-cov} away from the exceptional arrangement
and expand both sides in the same resolved chart.  The finite smooth V7 map is
polynomial/analytic in the retained finite coordinate bank at fixed order, so
it acts coefficientwise on the power--log series.  Uniqueness of the expansion
identifies the two sides term by term.
\end{proof}

\section{Subtracted exceptional remainder and finite boundary data}
\label{sec:v44-subtracted}

Although an exceptional coefficient need not itself be finite, V44 provides a
canonical finite remainder to any prescribed asymptotic order without changing
the theory.  Let $\mathcal P_{<A}^{\mathcal N}W$ be the finite sum of all
resolved monomials whose chosen weighted exponent is $<A$.

\begin{corollary}[Finite asymptotic subtraction]
\label{cor:v44-asymp-subtraction}
For every fixed weighted approach and every real $A$,
\begin{equation}
 \boxed{
 W_{\infty,0}
 -\mathcal P_{<A}^{\mathcal N}W_{\infty,0}
 =O\!\left(
 \varepsilon^A(1+|\log\varepsilon|^{M_A})
 \right).}
 \label{eq:v44-asymp-subtraction}
\end{equation}
This subtraction is an asymptotic decomposition into universal soft factors and
hard coefficients, not a new local renormalization prescription.
\end{corollary}

In particular, after all nonpositive soft monomials have been displayed
explicitly, the remainder extends continuously to the chosen resolved boundary
face.  Thus V44 supplies well-defined finite \emph{hard boundary data} even
when the full exceptional coefficient diverges.

\section{V44 closure theorem}
\label{sec:v44-close}

\begin{theorem}[Resolved exceptional Euclidean factorization]
\label{thm:v44-close}
Fix:
\begin{enumerate}[label=\textup{(E\arabic*)},leftmargin=3em]
\item the V42 cofinal V24-protected native branch and the V43 zero-floor
coefficientwise Euclidean continuum;
\item finite $(L,\Delta,r,N)$ and one coherent positive V7 normalization point;
\item a compact Euclidean external/source panel with bounded momenta, allowing
specified channel sums to vanish;
\item one resolved compatible exceptional chart $\mathcal N$ and any finite
set of retained source/BV/momentum derivatives.
\end{enumerate}
Then the pullback of every retained normalized coefficient to the resolved
exceptional chart has the source-complete polyhomogeneous expansion
\eqref{eq:v44-polyhomogeneous}.  Its coefficients factor into universal soft
forest tensors times V43 hard quotient coefficients; its approach to any
weighted exceptional stratum obeys the power--log bound
\eqref{eq:v44-power-log}; no new UV or IR counterterm is introduced; the common
BV/BRST/master/source identities hold coefficientwise in the expansion; exact
V7 perturbative projectivity and positive-reference covariance are preserved.

Consequently V43's loss of a positive nonexceptionality margin as
$\kappa\downarrow0$ is replaced by an explicit finite soft-forest asymptotic
calculus rather than by an uncontrolled divergence estimate.
\end{theorem}

\begin{proof}
Resolve the finite exceptional channel arrangement by compatible nested charts.
Theorem~\ref{thm:v44-single-factorization} supplies the one-channel
sector/Taylor factorization.  Iterating through the finite nested forest proves
Theorem~\ref{thm:v44-polyhomogeneous}.  The weighted bound is Theorem
\ref{thm:v44-power-log}; Proposition~\ref{prop:v44-hard-v43} identifies every
hard factor with the already-renormalized V43/V7 bank.  Theorems
\ref{thm:v44-master-factorization}, \ref{thm:v44-projective} and
\ref{thm:v44-reference} preserve the common algebraic identities, chronology
and reference covariance.  Lemma~\ref{lem:v44-no-euclidean-collinear} explains
why no separate collinear face is missing in real Euclidean signature.
\end{proof}

\section{V44 anti-circularity audit}
\label{sec:v44-audit}

\begin{enumerate}[leftmargin=2.4em]
\item \textbf{The theorem does not assume a positive nonexceptionality margin.}
Channels in $\mathcal N$ are allowed to reach zero; their singular behavior is
resolved explicitly by radial boundary variables.
\item \textbf{Soft divergences are not renamed counterterms.}  Negative and
zero infrared degrees appear as explicit power/log factors.  The local V5--V7
normalization bank is not enlarged to cancel them.
\item \textbf{Hard factors are genuinely hard.}  After contraction of the
chosen soft forest, every remaining nonlocal quotient channel is separated
from zero by the chart constant $\kappa_0$ and therefore falls under V43.
\item \textbf{Nested exceptional limits are not collapsed to one scalar
$\kappa$.}  The resolved variables $\rho_A$ retain independent relative rates;
a one-parameter $\varepsilon$ is introduced only after the multi-scale theorem
has been proved.
\item \textbf{No all-orders soft exponentiation is claimed.}  Universality is
proved graphwise and coefficientwise at the fixed finite perturbative/EFT/source
packet.
\item \textbf{No Euclidean collinear fiction is introduced.}  Collinear jet
pinches require the later Minkowski boundary-value problem.
\item \textbf{BV/source compatibility is not inferred from a scalar amplitude.}
The factorization is performed on the full contraction-complete V5 packet, and
uniqueness of the polyhomogeneous master expansion gives the induced identities
coefficient by coefficient.
\end{enumerate}

\section{V44 status ledger}
\label{sec:v44-status}

\subsection*{Proved in V44}
\begin{enumerate}[leftmargin=2.2em]
\item The real Euclidean exceptional set is the finite arrangement of vanishing
partial external/source channel sums; there is no independent nonzero
massless-collinear boundary face before analytic continuation.
\item The channel arrangement admits a finite resolved atlas whose boundary
faces are indexed by compatible nested exceptional families.
\item Every UV-forest-prepared fixed-packet coefficient has a finite soft-forest
factorization at one exceptional channel, with explicit power/log exponents and
a remainder of arbitrarily prescribed positive asymptotic order.
\item Iterating over nested/disjoint channels gives a polyhomogeneous expansion
on the resolved exceptional space.  Every coefficient is a finite sum of
universal soft subgraph factors times contracted V43 hard coefficients.
\item Along every weighted approach to a resolved stratum there is a sharp
sector power--log upper bound; complete-packet cancellations can only improve
the leading exponent.
\item No new renormalization coordinate is required.  All hard quotient terms
are V43 nonexceptional coefficients or already-owned V7 local/contact data.
\item The full retained source/BV derivative packet factorizes, and every
power--log coefficient satisfies the induced common master/source identity.
\item Exact V7 loop-order projectivity and positive-reference covariance
survive exceptional factorization.
\end{enumerate}

\subsection*{Inherited}
\begin{enumerate}[leftmargin=2.2em]
\item V1--V7: native arbitrary-fixed-order UV continuum, contraction-complete
forest recursion, common BV restoration, shell convergence, finite reference
maps and exact perturbative projectivity.
\item V8--V42: regulator universality/topology, controlled zero modes,
Wilson-wall/global-phase analysis, microscopic Picard no-go and contraction of
the cofinal regulator family.
\item V43: removal of the auxiliary infrared floor with a quantitative
source-complete bound on every compact strictly nonexceptional Euclidean panel.
\end{enumerate}

\subsection*{Conditional}
\begin{enumerate}[leftmargin=2.2em]
\item V44 is coefficientwise at finite $(L,\Delta,r,N)$; it does not assert a
uniform exceptional expansion as perturbative/EFT order tends to infinity.
\item The hard angular coordinates of each resolved chart are kept in a compact
set away from further unrecorded exceptionalities.  Additional degenerations
must be included as extra boundary faces of the nested set.
\item The asymptotic exponents are those of the fixed UV-renormalized soft
forest packet.  Ward or representation cancellations may improve the complete
leading exponent, but V44 does not assume such cancellations in advance.
\end{enumerate}

\subsection*{Open beyond V44}
\begin{enumerate}[leftmargin=2.2em]
\item Construct the analytic continuation from the V44 resolved Euclidean hard
coefficients to physical Minkowski boundary values and classify the resulting
Landau pinch surfaces.
\item On those physical boundary values, separate soft from genuinely
collinear jet regions and prove a source/BV-complete soft--jet factorization
formula with overlap subtraction.
\item Only after that pinch/factorization theorem, formulate an inclusive or
dressed cancellation theorem for physical massless transition probabilities.
\item Analytic finite-coupling RG trajectories, perturbative summability and a
nonperturbative Einstein--Standard Model completion remain independent stronger
problems.
\end{enumerate}

\subsection*{Exact next load-bearing theorem}
\begin{center}
\fbox{\parbox{0.92\textwidth}{
\textbf{Minkowski boundary-value and pinch-surface factorization theorem.}
V44 completely resolves the real Euclidean exceptional arrangement at every
fixed perturbative/EFT/source order: all vanishing-channel limits are
polyhomogeneous soft-forest expansions with V43 hard factors.  The next gate is
not another Euclidean $\kappa$ estimate.  Analytically continue the hard and
soft factors to a specified physical sheet, prove existence of the relevant
boundary values away from thresholds, classify the Landau pinch surfaces that
appear there, and construct a non-overcounting factorization into hard, soft
and genuinely collinear jet regions compatible with the V5 BV/source packet
and V7 projectivity.  An inclusive/dressed KLN-type theorem should be attempted
only after this physical pinch-surface decomposition is closed.}}
\end{center}

\section*{Append-only continuation marker: V44}
\addcontentsline{toc}{section}{Append-only continuation marker: V44}
\textbf{End of V44.}  The complete V1--V43 body above is retained verbatim.
V44 goes upstream on the phrase ``soft/collinear exceptional factorization'' by
separating the real Euclidean exceptional problem from the later Minkowski
collinear problem.  The finite arrangement of vanishing Euclidean channel sums
is resolved into nested radial boundary faces.  At every fixed
perturbative/EFT/source packet, the V43 zero-floor coefficients pull back to
polyhomogeneous power--log expansions whose coefficients are universal soft
forests times contraction-complete V43 hard quotient coefficients.  The common
BV/master/source identities, exact loop-order projectivity and positive
reference covariance survive coefficientwise, while negative/zero infrared
degrees remain explicit physical soft singularities rather than being absorbed
into new counterterms.  The next stronger gate is analytic continuation to
physical Minkowski boundary values and the resulting Landau soft--collinear
pinch-surface factorization.


\end{document}
